\documentclass[11pt,oneside,openany]{book}
\usepackage[letterpaper,top=.82in,bottom=.85in,left=.88in,right=.88in,headheight=15pt]{geometry}
\usepackage[T1]{fontenc}
\usepackage[utf8]{inputenc}
\usepackage{amsmath,amssymb,amsthm,mathtools,bm,mathrsfs,stmaryrd,mathpazo,microtype}
\usepackage{graphicx,float,booktabs,longtable,array,tabularx,multirow,enumitem}
\usepackage[table]{xcolor}
\usepackage{tikz,tikz-cd}
\usetikzlibrary{arrows.meta,positioning,calc,fit,backgrounds,decorations.pathreplacing,matrix,patterns}
\tikzset{box/.style={draw=KomorebiPine!70,rounded corners=2pt,fill=KomorebiSoft!50,align=center,inner sep=8pt,font=\small},arr/.style={-{Latex[length=2.2mm]},draw=KomorebiPine,thick},thinarr/.style={-{Latex[length=1.8mm]},draw=KomorebiMuted},lab/.style={font=\small,align=center},dot/.style={circle,fill=KomorebiPine,inner sep=2.3pt}}
\usepackage[most]{tcolorbox}
\usepackage{listings,fvextra,caption,subcaption,placeins,multicol,setspace,centernot,xurl}
\usepackage{fancyhdr,titlesec,needspace,etoolbox,hyperref,bookmark}
\usepackage[nameinlink,noabbrev]{cleveref}
\definecolor{Ink}{HTML}{142D4E}\definecolor{Accent}{HTML}{B08D57}
\definecolor{KomorebiPine}{HTML}{285A50}\definecolor{KomorebiTeal}{HTML}{285A50}
\definecolor{KomorebiInk}{HTML}{142D4E}\definecolor{KomorebiSoft}{HTML}{EAF0ED}
\definecolor{KomorebiWarm}{HTML}{F5F2E9}\definecolor{KomorebiOchre}{HTML}{B57F2A}
\definecolor{KomorebiMuted}{HTML}{65706D}\definecolor{KomorebiPale}{HTML}{EAF0ED}
\definecolor{KomorebiGold}{HTML}{B57F2A}\definecolor{KomorebiGoldPale}{HTML}{EFE2BE}
\definecolor{KomorebiBlue}{HTML}{4C6A92}\definecolor{KomorebiBluePale}{HTML}{E4EBF4}
\definecolor{KomorebiRose}{HTML}{D98673}\definecolor{KomorebiRosePale}{HTML}{F3E2DD}
\definecolor{KomorebiLine}{HTML}{7A867F}\definecolor{Slate}{HTML}{596965}
\definecolor{Pine}{HTML}{244D43}\definecolor{PalePine}{HTML}{F0F5F2}
\definecolor{RuleGray}{HTML}{D8E0DB}\definecolor{PartColor}{HTML}{223A5E}
\definecolor{AccentColor}{HTML}{6B3F8C}\definecolor{RuleColor}{HTML}{B8C1CC}
\definecolor{teal}{HTML}{146B72}\definecolor{gold}{HTML}{AA7832}
\definecolor{ink}{HTML}{172A44}\definecolor{muted}{HTML}{5B6778}
\newcolumntype{P}[1]{>{\raggedright\arraybackslash}p{#1}}
\newcolumntype{L}[1]{>{\raggedright\arraybackslash}p{#1}}
\newcolumntype{Y}{>{\raggedright\arraybackslash}X}
\newtcolorbox{scopebox}[1][]{breakable,colback=PalePine,colframe=RuleGray,boxrule=.3pt,arc=0pt,#1}
\newenvironment{scopepanel}{\begin{quote}\small}{\end{quote}}
\newenvironment{insight}{\begin{quote}}{\end{quote}}
\newenvironment{statusbox}{\begin{quote}\small}{\end{quote}}
\newenvironment{interpretation}{\par\smallskip\textit{Interpretation.}}{\par\smallskip}
\newcommand{\tightlist}{\setlength{\itemsep}{0pt}\setlength{\parskip}{0pt}}
\newcommand{\passthrough}[1]{#1}
\newcommand{\topic}[2]{\section{#2}}
\newcommand{\subtopic}[1]{\subsection*{#1}}
\newcommand{\centerbox}[1]{\begin{center}\fbox{\parbox{.9\linewidth}{\centering #1}}\end{center}}
\newcommand{\chap}[1]{\section{#1}}
\newcommand{\written}[1]{}
\counterwithout{section}{chapter}
\numberwithin{equation}{section}
\setcounter{tocdepth}{0}\setcounter{secnumdepth}{2}
\setlength{\parindent}{0pt}\setlength{\parskip}{5pt}\emergencystretch=3em
\raggedbottom\allowdisplaybreaks[2]
\setlist{leftmargin=1.6em,itemsep=2pt,topsep=4pt}
\titleformat{\chapter}[display]{\normalfont\Large\bfseries\color{Ink}}{\chaptername\ \thechapter}{8pt}{\huge}
\titlespacing*{\chapter}{0pt}{0pt}{18pt}
\pagestyle{fancy}\fancyhf{}\fancyhead[L]{\small\scshape Machine-Native Reasoning}
\fancyhead[R]{\small\scshape Realization Calculus}\fancyfoot[C]{\thepage}
\fancypagestyle{plain}{\fancyhf{}\fancyfoot[C]{\thepage}\renewcommand{\headrulewidth}{0pt}}
\hypersetup{colorlinks=true,linkcolor=Ink,citecolor=Ink,urlcolor=Ink,
pdftitle={Machine-Native Reasoning: Reconstructed Editorial Edition v1.0.0},pdfauthor={Komorebi P.},bookmarksdepth=2}
\lstset{basicstyle=\ttfamily\footnotesize,breaklines=true,columns=fullflexible,keepspaces=true}
\lstdefinestyle{code}{basicstyle=\ttfamily\small,breaklines=true,columns=fullflexible,keepspaces=true,showstringspaces=false}

% Semantic definitions retained from research/machine-native-reasoning-wider/book.tex
\newcommand{\CoreTV}{d_{\mathrm{TV}}}
\newcommand{\Coreim}{\operatorname{im}}
\newcommand{\CoreSpan}{\operatorname{span}}
\newcommand{\CorePow}{\mathcal P}
\newcommand{\Coresource}[1]{\par{\small\textit{Source relationship.} #1}\par}
\newcommand{\Coreboundary}[1]{\par{\small\textit{Boundary.} #1}\par}
\newcommand{\Corequestion}[1]{\par\medskip{\large\bfseries #1}\par\smallskip}
\newtheorem{Coretheorem}{Theorem}[chapter]
\newtheorem{Corelemma}[Coretheorem]{Lemma}
\newtheorem{Coreproposition}[Coretheorem]{Proposition}
\newtheorem{Corecorollary}[Coretheorem]{Corollary}
\newtheorem{Coredefinition}[Coretheorem]{Definition}
\newtheorem{Corecounterexample}[Coretheorem]{Counterexample}

% Semantic definitions retained from monographs/komorebi-space/source/Komorebi_Space_Monograph.tex
\newcommand{\KSR}{\mathbb R}
\newcommand{\KSZ}{\mathbb Z}
\newcommand{\KSQ}{\mathbb Q}
\newcommand{\KSN}{\mathbb N}
\newcommand{\KSK}{\mathbb K}
\newcommand{\KSV}{\mathcal V}
\newcommand{\KSW}{\mathcal W}
\newcommand{\KSid}{\operatorname{id}}
\newcommand{\KSim}{\operatorname{im}}
\newcommand{\KSran}{\operatorname{ran}}
\newcommand{\KSHom}{\operatorname{Hom}}
\newcommand{\KSOb}{\operatorname{Ob}}
\newcommand{\KSExt}{\operatorname{Ext}}
\newcommand{\KSMatch}{\operatorname{Match}}
\newcommand{\KSExact}{\operatorname{Exact}}
\newcommand{\KSdist}{\operatorname{dist}}
\newcommand{\KSAut}{\operatorname{Aut}}
\newcommand{\KSStab}{\operatorname{Stab}}
\newcommand{\KSSet}{\mathbf{Set}}
\newcommand{\KSCat}{\mathbf{Cat}}
\newcommand{\KSnorm}[1]{\left\lVert#1\right\rVert}
\newcommand{\KSabs}[1]{\left\lvert#1\right\rvert}
\newcommand{\KSsource}[1]{\phantomsection\label{KS:src:#1}}
\newcommand{\KScode}[1]{\texttt{\detokenize{#1}}}
\newcommand{\KSfixedplate}[3]{\begin{figure}[H]\centering\resizebox{.96\linewidth}{!}{\input{sources/monographs/komorebi-space/source/figures/#1}}\caption{#2}\label{KS:#3}\end{figure}}
\newcommand{\KSplate}[3]{\begin{figure}[!htbp]\centering\resizebox{.96\linewidth}{!}{\input{sources/monographs/komorebi-space/source/figures/#1}}\caption{#2}\label{KS:#3}\end{figure}}
\newcommand{\KSEditionProductionModules}{162}
\newcommand{\KSEditionRoots}{4}
\newcommand{\KSEditionTargets}{166}
\newcommand{\KSEditionTheorems}{1,859}
\newcommand{\KSEditionDeclarations}{5,339}
\newcommand{\KSEditionControls}{56}
\newcommand{\KSEditionPositiveControls}{18}
\newcommand{\KSEditionNegativeControls}{35}
\newcommand{\KSEditionPolicyControls}{3}
\newcommand{\KSEditionVerifiedGroups}{16}
\newcommand{\KSEditionWrittenGroups}{0}
\newcommand{\KSEditionPrincipalGroups}{16}
\newcommand{\KSEditionSelectedMapEntries}{114}
\newcommand{\KSEditionEntrypoints}{\KSEditionRoots}
\newcommand{\KSEditionExplicitTheorems}{\KSEditionTheorems}
\newcommand{\KSEditionExportedDeclarations}{\KSEditionDeclarations}
\newcommand{\KSwritten}[1]{}
\newtheorem{KStheorem}{Theorem}[section]
\newtheorem{KSproposition}[KStheorem]{Proposition}
\newtheorem{KSlemma}[KStheorem]{Lemma}
\newtheorem{KScorollary}[KStheorem]{Corollary}
\newtheorem{KSdefinition}[KStheorem]{Definition}
\newtheorem{KSexample}[KStheorem]{Example}
\newtheorem{KSremark}[KStheorem]{Remark}

% Semantic definitions retained from monographs/grounded-abstraction/Grounded_Abstraction_v1.0.0.tex
\newcommand{\GAX}{\mathcal X}
\newcommand{\GAY}{\mathcal Y}
\newcommand{\GAZ}{\mathcal Z}
\newcommand{\GAA}{\mathcal A}
\newcommand{\GAC}{\mathcal C}
\newcommand{\GAE}{\mathcal E}
\newcommand{\GAHh}{\mathcal H}
\newcommand{\GAPp}{\mathcal P}
\newcommand{\GATV}{d_{\mathrm{TV}}}
\newcommand{\GAim}{\operatorname{im}}
\newcommand{\GAsupp}{\operatorname{supp}}
\newcommand{\GASpan}{\operatorname{span}}
\newcommand{\GAPow}{\mathcal P}
\newtheorem{GAtheorem}{Theorem}[section]
\newtheorem{GAproposition}[GAtheorem]{Proposition}
\newtheorem{GAcorollary}[GAtheorem]{Corollary}
\newtheorem{GAdefinition}[GAtheorem]{Definition}
\newtheorem{GAcounterexample}[GAtheorem]{Counterexample}

% Semantic definitions retained from research/machine-understanding/Machine_Understanding_v1.0.0.tex
\newcommand{\MUX}{\mathcal X}
\newcommand{\MUZ}{\mathcal Z}
\newcommand{\MUU}{\mathcal U}
\newcommand{\MUQ}{\mathcal Q}
\newcommand{\MUY}{\mathcal Y}
\newcommand{\MUC}{\mathcal C}
\newcommand{\MUM}{\mathcal M}
\newcommand{\MUE}{\mathbb E}
\newcommand{\MUFfive}{\mathbb F_5}
\newcommand{\MUTV}{d_{\mathrm{TV}}}
\newcommand{\MUimg}[1]{\operatorname{im}(#1)}
\newcommand{\MUEq}{\equiv_{\MUC}}
\newcommand{\MUsupp}{\operatorname{supp}}
\newcommand{\MUcode}[1]{\texttt{\small #1}}
\newtheorem{MUtheorem}{Theorem}[section]
\newtheorem{MUproposition}[MUtheorem]{Proposition}
\newtheorem{MUdefinition}[MUtheorem]{Definition}
\newtheorem{MUcounterexample}[MUtheorem]{Counterexample}
\newtheorem{MUexample}[MUtheorem]{Example}

% Semantic definitions retained from monographs/self-improving-intelligence/Self_Improving_Intelligence_v1.0.0.tex
\newcommand{\SIIX}{\mathcal X}
\newcommand{\SIIZ}{\mathcal Z}
\newcommand{\SIIU}{\mathcal U}
\newcommand{\SIIQ}{\mathcal Q}
\newcommand{\SIIY}{\mathcal Y}
\newcommand{\SIIC}{\mathcal C}
\newcommand{\SIIM}{\mathcal M}
\newcommand{\SIIE}{\mathbb E}
\newcommand{\SIIFfive}{\mathbb F_5}
\newcommand{\SIITV}{d_{\mathrm{TV}}}
\newcommand{\SIIimg}[1]{\operatorname{im}(#1)}
\newcommand{\SIIEq}{\equiv_{\SIIC}}
\newcommand{\SIIsupp}{\operatorname{supp}}
\newcommand{\SIIcode}[1]{\texttt{\small #1}}
\newcommand{\SIIB}{\mathcal B}
\newcommand{\SIIRcal}{\mathcal R}
\newcommand{\SIIGcal}{\mathcal G}
\newcommand{\SIISafe}{\mathsf{Safe}}
\newcommand{\SIIPost}{\operatorname{Post}}
\newcommand{\SIIPareto}{\operatorname{Pareto}}
\newtheorem{SIItheorem}{Theorem}[section]
\newtheorem{SIIproposition}[SIItheorem]{Proposition}
\newtheorem{SIIdefinition}[SIItheorem]{Definition}
\newtheorem{SIIcounterexample}[SIItheorem]{Counterexample}
\newtheorem{SIIexample}[SIItheorem]{Example}

% Semantic definitions retained from monographs/machine-societies/publication/Machine_Societies_v1.0.0.tex
\newcommand{\MSTV}{d_{\mathrm{TV}}}
\newcommand{\MSLaw}{\mathcal L}
\newcommand{\MSX}{\mathcal X}
\newcommand{\MSZ}{\mathcal Z}
\newcommand{\MSU}{\mathcal U}
\newcommand{\MSY}{\mathcal Y}
\newcommand{\MSHist}{\mathcal H}
\newcommand{\MSE}{\mathbb E}
\newcommand{\MSProb}{\mathbb P}
\newcommand{\MSFtwo}{\mathbb F_2}
\newcommand{\MSpush}{\pi_{\#}}
\newcommand{\MSnorm}[1]{\left\lVert#1\right\rVert}
\newcommand{\MSabs}[1]{\left\lvert#1\right\rvert}
\newcommand{\MScode}[1]{\texttt{\small #1}}
\newcommand{\MSStatus}[1]{\par\noindent\textbf{Status:} #1\par}
\newtheorem{MSproposition}{Proposition}[section]
\newtheorem{MStheorem}[MSproposition]{Theorem}
\newtheorem{MScorollary}[MSproposition]{Corollary}
\newtheorem{MSdefinition}[MSproposition]{Definition}
\newtheorem{MSexample}[MSproposition]{Example}
\newtheorem{MScounterexample}[MSproposition]{Counterexample}
\newtheorem{MShypothesis}[MSproposition]{Research hypothesis}
\newtheorem*{MSinterpretation}{Interpretation}

% Semantic definitions retained from monographs/large-signal-model/Reference_Laboratory/theory/FOUNDATIONS.tex
\newcommand{\LSMtightlist}{\setlength{\itemsep}{0pt}\setlength{\parskip}{0pt}}
\newcommand{\LSMpassthrough}[1]{#1}

% Semantic definitions retained from research/hypergraph-bayesian-refinement/Hypergraph_Lumpability_Bayesian_Refinement.tex
\newcommand{\PRIQR}{\mathbb R}
\newcommand{\PRIQC}{\mathbb C}
\newcommand{\PRIQN}{\mathbb N}
\newcommand{\PRIQcA}{\mathcal A}
\newcommand{\PRIQcB}{\mathcal B}
\newcommand{\PRIQcC}{\mathcal C}
\newcommand{\PRIQcD}{\mathcal D}
\newcommand{\PRIQcE}{\mathcal E}
\newcommand{\PRIQcF}{\mathcal F}
\newcommand{\PRIQcG}{\mathcal G}
\newcommand{\PRIQcH}{\mathcal H}
\newcommand{\PRIQcL}{\mathcal L}
\newcommand{\PRIQcP}{\mathcal P}
\newcommand{\PRIQcR}{\mathcal R}
\newcommand{\PRIQcS}{\mathcal S}
\newcommand{\PRIQcT}{\mathcal T}
\newcommand{\PRIQcV}{\mathcal V}
\newcommand{\PRIQcY}{\mathcal Y}
\newcommand{\PRIQid}{\operatorname{id}}
\newcommand{\PRIQRan}{\operatorname{Ran}}
\newcommand{\PRIQKer}{\operatorname{Ker}}
\newcommand{\PRIQSpan}{\operatorname{span}}
\newcommand{\PRIQdiag}{\operatorname{diag}}
\newcommand{\PRIQosc}{\operatorname{osc}}
\newcommand{\PRIQsupp}{\operatorname{supp}}
\newcommand{\PRIQTV}{\operatorname{TV}}
\newcommand{\PRIQLip}{\operatorname{Lip}}
\newcommand{\PRIQtr}{\operatorname{tr}}
\newcommand{\PRIQrank}{\operatorname{rank}}
\newcommand{\PRIQdd}{\,\mathrm d}
\newcommand{\PRIQnorm}[1]{\left\lVert #1\right\rVert}
\newcommand{\PRIQabs}[1]{\left\lvert #1\right\rvert}
\newcommand{\PRIQip}[2]{\left\langle #1,#2\right\rangle}
\newcommand{\PRIQfine}{\mathrm f}
\newcommand{\PRIQcoarse}{\mathrm c}
\newcommand{\PRIQBup}{\mathsf B}
\newcommand{\PRIQNrm}{\mathsf N}
\newcommand{\PRIQDef}{\mathfrak D}
\newcommand{\PRIQErr}{\mathfrak E}
\newcommand{\PRIQone}{\mathbf 1}
\newtheorem{PRIQtheorem}{Theorem}[section]
\newtheorem{PRIQproposition}[PRIQtheorem]{Proposition}
\newtheorem{PRIQlemma}[PRIQtheorem]{Lemma}
\newtheorem{PRIQcorollary}[PRIQtheorem]{Corollary}
\newtheorem{PRIQconjecture}[PRIQtheorem]{Conjecture}
\newtheorem{PRIQdefinition}[PRIQtheorem]{Definition}
\newtheorem{PRIQexample}[PRIQtheorem]{Example}
\newtheorem{PRIQalgorithm}[PRIQtheorem]{Algorithm}
\newtheorem{PRIQcertificate}[PRIQtheorem]{Numerical Certificate}
\newtheorem{PRIQboundary}[PRIQtheorem]{Claim Boundary}
\newtheorem{PRIQremark}[PRIQtheorem]{Remark}

% Semantic definitions retained from monographs/proof-path-theory/docs/Proof_Path_Theory.tex
\newcommand{\PPTcalloutbox}[2]{\par\medskip
 {\sffamily\small\bfseries\color{KomorebiTeal}#1\par}
 \smallskip #2\par\medskip}
\newcommand{\PPTPcal}{\mathcal P}
\newcommand{\PPTJcal}{\mathcal J}
\newcommand{\PPTQcal}{\mathcal Q}
\newcommand{\PPTRcal}{\mathcal R}
\newcommand{\PPTEcal}{\mathcal E}
\newcommand{\PPTVcal}{\mathcal V}
\newcommand{\PPTCcal}{\mathcal C}
\newcommand{\PPTBcal}{\mathcal B}
\newcommand{\PPTTcal}{\mathcal T}
\newcommand{\PPTFcal}{\mathcal F}
\newcommand{\PPTAnc}{\operatorname{Anc}}
\newcommand{\PPTAncL}{\operatorname{Anc}_L}
\newcommand{\PPTDesc}{\operatorname{Desc}}
\newcommand{\PPTReach}{\operatorname{Reach}}
\newcommand{\PPTCost}{\operatorname{Cost}}
\newcommand{\PPTLoss}{\operatorname{Loss}}
\newcommand{\PPTDebt}{\operatorname{Debt}}
\newcommand{\PPTargmin}{\operatorname*{arg\,min}}
\newcommand{\PPTargmax}{\operatorname*{arg\,max}}
\newcommand{\PPTopeningfigure}[3]{%
\begin{figure}[H]\centering
\includegraphics[width=\linewidth]{sources/monographs/proof-path-theory/docs/figures/#1}
\caption{#2}\label{PPT:#3}
\end{figure}}
\newcommand{\PPTstudyhead}[2]{%
 \FloatBarrier\clearpage
 {\sffamily\small\bfseries\color{KomorebiTeal}VISUAL STUDY\quad #1\par}
 \vspace{7pt}{\sffamily\fontsize{23}{27}\selectfont\bfseries\color{KomorebiInk}#2\par}\vspace{8pt}}
\newtheorem{PPTtheorem}{Theorem}[section]
\newtheorem{PPTproposition}[PPTtheorem]{Proposition}
\newtheorem{PPTlemma}[PPTtheorem]{Lemma}
\newtheorem{PPTcorollary}[PPTtheorem]{Corollary}
\newtheorem{PPTdefinition}[PPTtheorem]{Definition}
\newtheorem{PPTexample}[PPTtheorem]{Example}
\newtheorem{PPTproblem}[PPTtheorem]{Problem}
\newtheorem{PPTremark}[PPTtheorem]{Remark}
\newtheorem{PPTwarning}[PPTtheorem]{Warning}

% Semantic definitions retained from packages/constructive-repair/monograph/Constructive_Repair.tex
\newcommand{\CRkk}{\Bbbk}
\newcommand{\CRQ}{\mathbb{Q}}
\newcommand{\CRR}{\mathbb{R}}
\newcommand{\CRZ}{\mathbb{Z}}
\newcommand{\CRN}{\mathbb{N}}
\newcommand{\CRF}{\mathbb{F}}
\DeclareMathOperator{\CRim}{im}
\DeclareMathOperator{\CRcoker}{coker}
\DeclareMathOperator{\CRSpan}{span}
\DeclareMathOperator{\CRrank}{rank}
\DeclareMathOperator{\CRcost}{cost}
\DeclareMathOperator{\CRid}{id}
\DeclareMathOperator{\CRargmin}{argmin}
\newcommand{\CROb}{\mathcal{O}}
\newcommand{\CRcode}[1]{\texttt{\detokenize{#1}}}
\newcommand{\CRdecl}[1]{{\small\nolinkurl{#1}}}
\newcommand{\CRformal}[1]{\par\nobreak\begingroup\footnotesize\color{Slate}\sffamily Formal counterpart: \normalfont\nolinkurl{#1}.\par\endgroup}
\newcommand{\CRsourcescope}[1]{\par\begingroup\small\color{Slate}#1\par\endgroup}
\newcommand{\CRrecord}[3]{\noindent\begin{minipage}{\linewidth}\small\textbf{\nolinkurl{#1}}\hfill\textsf{line #2}\par\smallskip #3\end{minipage}\par\vspace{7pt}}
\newcommand{\CRmoduleheading}[2]{\section{#1}\label{CR:#2}\noindent\small Source: \nolinkurl{lean/KomorebiRepair/#1.lean}.\normalsize\par}
\newtheorem{CRtheorem}{Theorem}[section]
\newtheorem{CRproposition}[CRtheorem]{Proposition}
\newtheorem{CRlemma}[CRtheorem]{Lemma}
\newtheorem{CRcorollary}[CRtheorem]{Corollary}
\newtheorem{CRdefinition}[CRtheorem]{Definition}
\newtheorem{CRexample}[CRtheorem]{Example}
\newtheorem{CRremark}[CRtheorem]{Remark}

% Semantic definitions retained from research/stratified-relational-workspace/monograph/Stratified_Relational_Workspace_Theory.tex
\newcommand{\SRWTWorkspace}{\mathbf W}
\newcommand{\SRWTSRW}{\operatorname{SRW}}
\newcommand{\SRWTCE}{\operatorname{CE}}
\newcommand{\SRWTRisk}{\mathcal R}
\newcommand{\SRWTHol}{\operatorname{Hol}}
\newcommand{\SRWTPath}{\operatorname{Path}}
\newcommand{\SRWTRet}{\operatorname{Ret}}
\newcommand{\SRWTDep}{\operatorname{Dep}}
\newcommand{\SRWTOut}{\operatorname{Out}}
\newcommand{\SRWTCOut}{\operatorname{COut}}
\newcommand{\SRWTCIn}{\operatorname{CIn}}
\newcommand{\SRWTInstability}{\operatorname{Instability}}
\newcommand{\SRWTRec}{\operatorname{Rec}}
\newcommand{\SRWTPass}{\operatorname{Pass}}
\newcommand{\SRWTindep}{\perp\!\!\!\perp}
\newcommand{\SRWTdd}{\,\mathrm d}
\DeclareMathOperator{\SRWTran}{ran}
\DeclareMathOperator{\SRWTspec}{spec}
\DeclareMathOperator{\SRWTEnd}{End}
\DeclareMathOperator{\SRWTFix}{Fix}
\DeclareMathOperator{\SRWTComm}{Comm}
\DeclareMathOperator{\SRWTev}{ev}
\DeclareMathOperator{\SRWTim}{im}
\DeclareMathOperator{\SRWTNat}{Nat}
\DeclareMathOperator{\SRWTKL}{KL}
\DeclareMathOperator{\SRWTTV}{TV}
\DeclareMathOperator{\SRWTpr}{pr}
\DeclareMathOperator*{\SRWTargmax}{arg\,max}
\DeclareMathOperator*{\SRWTargmin}{arg\,min}
\newcommand{\SRWTcenterbox}[1]{%
\begin{center}
\fbox{\parbox{0.90\textwidth}{\centering #1}}
\end{center}
}
\newtheorem{SRWTtheorem}{Theorem}
\newtheorem{SRWTproposition}[SRWTtheorem]{Proposition}
\newtheorem{SRWTlemma}[SRWTtheorem]{Lemma}
\newtheorem{SRWTcorollary}[SRWTtheorem]{Corollary}
\newtheorem{SRWTdefinition}[SRWTtheorem]{Definition}
\newtheorem{SRWTexample}[SRWTtheorem]{Example}
\newtheorem{SRWTremark}[SRWTtheorem]{Remark}
\renewcommand{\theSRWTtheorem}{SRWT.\arabic{SRWTtheorem}}

% Semantic definitions retained from monographs/large-signal-model/Reference_Laboratory/theory/PROOF_REGISTER.tex
\newcommand{\LSMRegistertightlist}{\setlength{\itemsep}{0pt}\setlength{\parskip}{0pt}}
\newcommand{\LSMRegisterpassthrough}[1]{#1}

% Semantic definitions retained from monographs/large-signal-model/Visual_Companion/atlas.tex
\newcommand{\LSMAtlasR}{\mathbb R}
\newcommand{\LSMAtlasim}{\operatorname{im}}
\newcommand{\LSMAtlasTV}{\operatorname{TV}}
\newcommand{\LSMAtlasdiag}{\operatorname{diag}}
\newcommand{\LSMAtlasplate}[4]{\clearpage\pdfbookmark[1]{#2}{plate#1}\textcolor{teal}{\small\textsc{#1\quad #4}}\par\vspace{3pt}{\fontsize{23}{26}\selectfont #2}\par\vspace{3pt}{\color{muted}#3}\par\vspace{8pt}}
\newcommand{\LSMAtlasdiagram}[2][0.98]{\begin{center}\includegraphics[width=#1\linewidth,height=2.45in,keepaspectratio]{sources/monographs/large-signal-model/Visual_Companion/figures/#2.pdf}\end{center}\vspace{-3pt}}
\newcommand{\LSMAtlasclaim}[1]{\par\vspace{3pt}\textcolor{teal}{\textbf{The mathematical point.}} #1\par}
\newcommand{\LSMAtlasboundary}[1]{\par\vspace{3pt}{\small\textcolor{muted}{\textbf{Scope.} #1}}\par}
\newcommand{\LSMAtlasnav}[1]{\par{\small\textcolor{muted}{\textbf{Formal route.} #1}}\par}

\begin{document}
\frontmatter
\hypersetup{pageanchor=false}
\begin{titlepage}\thispagestyle{empty}
{\small\scshape Realization Calculus\hfill Mathematical Reasoning}\par
\vspace*{1.1in}{\color{Ink}\fontsize{36}{42}\selectfont\bfseries Machine-Native\\Reasoning\par}
\vspace{.3in}{\Large Representation, Inference, Discovery, and Self-Correction\par}
\vspace{.4in}{\color{Accent}\rule{1.9in}{2pt}}\par
\vspace{.4in}{\Large Komorebi P.\par}\vfill
{\Large Reconstructed Editorial Edition v1.0.0\par}\medskip
{4 November 2026\par}\medskip
{\small An editable reconstruction from retained component manuscripts and the recovered prose-refined reading edition.}\par
\end{titlepage}
\hypersetup{pageanchor=true}
\chapter{Preface}
A reasoner can predict a device's display without knowing how it relates to the physical state. It can find a successful repair without retaining the evidence needed to justify it. It can store a correct proof whose premises no longer describe the system on which it is about to act. This book asks what a representation must retain for knowledge to remain usable.

Realization Calculus studies representations together with their operations, transformations and residuals. Komorebi Space supplies a setting in which several views may describe one coherent object, or fail to do so. Applied to reasoning, these ideas lead to concrete questions. Can the evidence initialize a predictive state? Does a compressed record support later correction? Can another program execute the method a message describes? Which guarantees survive a change of model or rule?

The opening chapters follow these questions through a finite system. A sensor, a calibration procedure and a repair task provide a running example. The specialist treatments then develop coherent realizations, grounding, signal models, nonlinear quotients, implemented understanding, proof fields, constructive repair, self-improvement, relational workspaces and machine societies. Combined conclusions must concern the same system, records, operations and resource account.

Written proofs, computational checks and evidence about an external system remain distinct. The edition and verification notes describe what the source reconstruction establishes and where result-level correspondence remains necessary.
\chapter{Edition, provenance, and verification}
This edition reconstructs an editable manuscript from the retained component papers and monographs. It is a new assembly, not the recovered historical TeX archive. The source register pins the contributing files and the reconstruction map identifies each chapter. The original 23 October reading edition remains the comparison witness.

The separate Reviewed Mathematical Core v1.0.0 accepted all twenty-seven selected original source results on 25 October 2026. In particular, the four selected source results formerly pending in the older book checkpoint are accepted. This count concerns those original source statements. Expanded formulations, integration arguments and specialist treatments retain their own correspondence obligations.

No new Lean run or independent whole-book proof audit is claimed by this reconstruction. The evidence distinguishes source preservation, transcription, compilation, editorial review, and result-level mathematical verification.

The admission record is available at \url{https://github.com/KomorebiP/Operation-KOMOREBI/blob/main/monographs/machine-native-reasoning/release/ADMISSION_RECORD.json}.
\tableofcontents
\renewcommand{\theCoretheorem}{R.\arabic{Coretheorem}}
\chapter*{The Realization Setting}


The mathematical question is not whether one machine can store an account of its environment. It is whether the account supports the questions, operations, corrections, and changes of representation to which the machine is committed. A description can be internally consistent and still fail to determine the information its consumer requires. It can encode an exact local state while being incompatible with another view of the same purported object. Conversely, a coarse representation can discard most of the ambient variables and preserve every continuation relevant to a declared task.

Realization Calculus treats these distinctions as mathematical obligations. Komorebi Space supplies an ambient setting in which several descriptions can be compared without assuming that their compatibility or global realization is automatic. This prologue introduces the portion of that setting needed for the finite reasoning systems developed in the chapters that follow. It does not import the entire category-theoretic or analytic theory of the separate Komorebi Space monograph.

\section*{Views, restrictions, and a coherent realization}
Let $\mathcal V$ be a small category of views. Associate a set $F_v$ to each view $v$ and a restriction $f^*:F_w\to F_v$ to each arrow $f:v\to w$. Identity and composition satisfy $(1_v)^*=\mathrm{id}$ and $(g\circ f)^*=f^*\circ g^*$. Thus $F$ is a contravariant set-valued functor. A coherent realization is a family of elements $x_v\in F_v$ satisfying all the restriction equations:
\[
\Gamma(F)=\left\{(x_v)_v\in\prod_v F_v:\ f^*(x_w)=x_v\text{ for every }f:v\to w\right\}.
\]
The set $\Gamma(F)$ is the strict realization set. It may be empty. A nonempty fibre at every individual view does not force a coherent family to exist.

For example, take the one-object category whose automorphisms form the two-element group, and let its nontrivial arrow exchange the two elements of $F_v=\{0,1\}$. A coherent family would be a fixed point of that exchange. There is none. No amount of relabeling can turn either local element into a global fixed point without changing the specified action. This illustrates a basic obstruction: local admissibility is weaker than coherent realization.

The evaluation $\operatorname{ev}_v:\Gamma(F)\to F_v$ records the information visible in one view. When $\Gamma(F)$ is nonempty and a target $q:\Gamma(F)\to Y$ is specified, that view is sufficient exactly when the target is constant on its fibres.
\begin{Coreproposition}[Target factorization through a realization view]\label{Core:r:viewfactor}
Suppose $\Gamma(F)$ and $Y$ are nonempty. There is a function $d:\operatorname{ev}_v(\Gamma(F))\to Y$ such that $q=d\circ\operatorname{ev}_v$ if and only if
\[
\operatorname{ev}_v(x)=\operatorname{ev}_v(x')\quad\Longrightarrow\quad q(x)=q(x')
\]
for all coherent realizations $x,x'\in\Gamma(F)$.
\end{Coreproposition}
\begin{proof}
A function of one attained record must assign the same value to equal records, proving necessity. Conversely, each attained record $a$ has a coherent preimage $x$. Define $d(a)=q(x)$. The fibre hypothesis makes this independent of the chosen preimage and supplies the required factorization. There is no need to define $d$ on unattained, possibly impossible view values.
\end{proof}

This proposition is the view-theoretic instance of the first fibre criterion in the reviewed core. It is not a new claim of priority for elementary factorization. Its significance here is that the domain consists of \emph{coherent realizations}, not all arbitrary families of local coordinates. The criterion cannot identify a physical system unless the accepted model class, the actual observation procedure, and the correspondence between observed data and the chosen view are supplied.

\section*{Protected operations and what a quotient must retain}
For an observed deterministic system, let $X$ be a set of realizations, $r:X\to Z$ a retained record, $F_a:X\to X$ an action, and $q:X\to Y$ a protected output. The coarse description is closed under action $a$ only if there is $\bar F_a:Z\to Z$ satisfying $rF_a=\bar F_a r$ on the attained record domain. Preservation of the current output $q=\bar q r$ is a separate condition. The first protects subsequent computation. The second protects present interpretation.

When these conditions hold for each action, induction transports every finite operation word. When they fail, the mismatch can be recorded as a residual rather than silently suppressed. Linear resolutions make the defect explicit. Let $C$ compress a state, let $R$ reconstruct a retained coordinate with $CR=I$, and let $T,U$ be two linear operations. The excursion $CT(I-RC)UR$ measures how discarded information can return to an observed coordinate after composition. Chapter~\ref{Core:ch:operations} derives this identity.

Stochastic systems require more structure than an arbitrary packet decomposition. A signed detail component of a probability vector may be needed for exact reconstruction but need not itself be a probability distribution. Bayesian updating is defined only where the appropriate evidence mass is positive. Legal operation systems require enabled-action sets as well as successor maps. A compression of outputs that loses action permissions is not a legal-operation quotient.

The remainder of the book follows these obligations from evidence through prediction, correction, decision, proof, and revision. The common point is preservation under a \emph{declared} change. The corresponding hypotheses differ across set-valued evidence, linear dynamics, probability laws, executable methods, and institutions. A common word such as ``packet'' or ``certificate'' does not identify those distinct types.

\section*{What the setting does not provide automatically}
Existence of a coherent realization does not identify it from finite observations. A commuting representation map does not produce a learned translator. An abstract action catalogue does not establish physical executability or external authority. A formal certificate validates a statement in a specified model and proof environment, not the model's adequacy to an unobserved world. Resource units cannot be added before a common cost interpretation is selected. These distinctions will remain explicit, particularly when the reasoning chain crosses from a mathematical result to a proposed implementation.

\renewcommand{\theCoretheorem}{\thechapter.\arabic{Coretheorem}}

\mainmatter

\setcounter{section}{0}\renewcommand{\thesection}{Core.\arabic{section}}\renewcommand{\theHsection}{Core.\thechapter.\arabic{section}}
\part{Evidence, methods, and correction}
\chapter{From evidence to a reusable state}\label{Core:book:chapter-1}
\setcounter{section}{0}\renewcommand{\thesection}{\thechapter.\arabic{section}}\label{Core:ch:evidence}
\Corequestion{The first question is not how much state to store. It is which distinctions are accessible.}
Let $\Omega$ be a finite nonempty set of possible worlds. A world may include a physical state, a fixed model index, and an origin tag. An observation record is a function $r:\Omega\to H$. A target is a function $t:\Omega\to Y$. All codomains used below are nonempty. We work on attained records $r(\Omega)$, so no arbitrary interpretation of an impossible record is needed.

\begin{Coreproposition}[A record determines a target exactly on its fibres]\label{Core:r:factor}
There exists $d:r(\Omega)\to Y$ with $d(r(\omega))=t(\omega)$ for every world if and only if
\[
r(\omega)=r(\omega')\ \Longrightarrow\ t(\omega)=t(\omega').
\]
\end{Coreproposition}
\begin{proof}
Necessity follows because a function assigns one value to one input. For sufficiency, define $d(h)$ to be $t(\omega)$ for any world with $r(\omega)=h$. Such a world exists for attained $h$, and the hypothesis makes the choice immaterial. The resulting decoder satisfies the required equation.
\end{proof}
\Coresource{Grounded Abstraction [GA], Proposition 3.2 and Theorem 3.4. Machine Understanding [MU], the representation and sender-access criteria. This is the shared elementary factorization argument, stated once.}

The proposition is an information statement, not a runtime bound. The table defining $d$ could be large. Nor does the mere existence of $d$ teach a learner how to find it. A proof that a suitable decoder exists and a justified account of its acquisition remain separate.

\section{The sensor used throughout the book}
Write a world as
\[
\omega=(b,s,n)\in\{0,1\}^3.
\]
The physical target bit is $b$. The sensor convention is $s$. The nuisance bit $n$ never affects the operations or outputs of the initial contract. The display is
\[
y=b\oplus s,
\]
where $\oplus$ is addition modulo two. A permitted control $f$ flips the physical bit and leaves $s,n$ unchanged. Consequently the next display is $1-y$.

An agent given only $y$ can predict every future display under repeated flips: after $k$ flips it is $y\oplus(k\bmod2)$. It cannot determine the current physical bit. Worlds $(0,0,0)$ and $(1,1,0)$ share display zero and have different targets. The obstruction is the fibre of the available record, not a shortage of computational effort.

Calibration is a separate, nondestructive observation $c(\omega)=s$. When calibration is permitted and trustworthy, the joint record $(y,s)$ determines $b=y\oplus s$. Adding a blank field called ``sensor convention'' does not provide this information. The observation operation, its interpretation, and its cost must be part of the contract.

\begin{center}\begin{tabular}{llll}\toprule
Available record & Number of values & Determines display? & Determines physical bit?\\\midrule
$y$ & 2 & Yes & No\\
$(y,s)$ & 4 & Yes & Yes\\
$(b,s,n)$ & 8 & Yes & Yes, with unused detail\\\bottomrule\end{tabular}\end{center}

\begin{figure}[ht]
\centering\includegraphics[width=\linewidth]{sources/research/machine-native-reasoning-wider/figures/sensor_chain.pdf}
\caption{One exact observation sequence. The last two worlds differ only in the nuisance bit. Knowing every hidden coordinate is unnecessary for the protected physical target.}
\end{figure}

\section{One answer versus every permitted continuation}
Let $X$ be a finite nonempty state space, and fix an available record $r:X\to H$. Let $\mathcal A$ be a finite alphabet of total deterministic operations $F_a:X\to X$, and let $q:X\to Y$ be the protected output. A word $w=a_1\cdots a_k$ executes $a_1$ first. Thus $F_w=F_{a_k}\cdots F_{a_1}$, and the empty word denotes the identity.

Define behavioral equivalence by
\[
x\sim x'\quad\Longleftrightarrow\quad q(F_wx)=q(F_wx')\text{ for every finite word }w.
\]
The empty word protects the present output. The other words protect its permitted future uses.

\begin{Coretheorem}[Evidence-initialized behavioral state]\label{Core:r:behavior}
A finite predictor initialized only from $r(x)$, updated only by the operation names, and exact for all protected continuation outputs exists if and only if
\[
r(x)=r(x')\ \Longrightarrow\ x\sim x'.
\]
When it exists, $X/{\sim}$ supplies a predictor with the smallest possible number of initialized states. Its updates are $\bar F_a([x])=[F_ax]$ and its readout is $\bar q([x])=q(x)$.
\end{Coretheorem}
\begin{proof}
If two initial records are equal, a deterministic predictor has the same initialized state. Applying the same subsequent operations produces the same predicted outputs. Exactness therefore requires behavioral equivalence.

For the converse, equivalent states have equivalent successors: an output after $F_a$ followed by $w$ is one of the outputs in the defining family. Hence $\bar F_a$ is well defined. The empty word makes $\bar q$ well defined. The fibre premise gives a well-defined initializer $d(r(x))=[x]$. Induction on word length proves exactness. Finally, behaviorally distinct states cannot initialize the same predictor state, since some continuation separates their outputs. Every exact predictor therefore needs at least one initialized value for each class, and the quotient attains that number.
\end{proof}
\Coresource{[GA], Theorem 3.4, and [MU], Theorem 4.1. The two source versions are consolidated, not counted as independent discoveries.}

In the sensor, the display task has two behavioral states. Changing the protected output to $(y,b)$ yields four behavioral states. The old record still has only two values and cannot initialize the enlarged quotient. The new task has exposed a real information requirement.

\section{Why a statistical law is not an observed record}
A parameter can determine a distinct probability law while every finite sample remains possible under several parameters. For instance, Bernoulli parameters $1/3$ and $2/3$ assign positive probability to every finite binary string. No estimator based on one such string can be certainly correct for both parameters. The laws differ, but their finite supports overlap.

A correct uncertainty record can be the appropriate output in this setting. Calling the most probable world ``identified'' would silently replace an error-free claim with a probabilistic one. Chapter~\ref{Core:ch:memory} gives the corresponding probability update explicitly.

\section*{What is now available}
We have a criterion for an initialized predictive state, but no automatic guarantee that an arbitrary coordinate manipulation preserves it. The next chapter studies operations on a representation. Its input is a declared state map and operation family, not an assumed learned encoder.


\chapter{Operations must survive their representations}\label{Core:book:chapter-2}
\setcounter{section}{0}\renewcommand{\thesection}{\thechapter.\arabic{section}}\label{Core:ch:operations}
\Corequestion{A representation is useful only to the extent that its intended operations remain meaningful.}
There are two different ways to change representation. An exact split representation retains the whole mathematical state in another space. A compression retains only part of it. The first can preserve arbitrary composition. The second needs additional conditions.

\section{Exact packets and the hidden-return defect}
Let $V,Z$ be vector spaces over a field. Let $A:V\to Z$ be analysis and $S:Z\to V$ synthesis, with $SA=I_V$. Then $AS$ is an idempotent on $Z$. The map $\mathcal E(T)=ATS$ represents a full-state operator $T:V\to V$ on resolved coordinates. Direct multiplication gives
\[
\mathcal E(T)\mathcal E(U)=ATSAUS=ATUS=\mathcal E(TU).
\]
The resolved identity is $AS$ on the compatible image $A(V)$, not necessarily $I_Z$ on arbitrary coordinates.

For a compression, reverse the splitting. Let $C:V\to Z$ and $R:Z\to V$ satisfy $CR=I_Z$, and put $P=RC$. The compressed operator is $\Phi(T)=CTR$.

\begin{Coreproposition}[Exact accounting of the discarded excursion]\label{Core:r:packet}
For any endomorphisms $T,U$ of $V$,
\[
\Phi(TU)-\Phi(T)\Phi(U)=CT(I-P)UR.
\]
Thus compression preserves this composition precisely when the discarded excursion on the right vanishes. Full split encoding $T\mapsto ATS$ preserves composition without that extra condition.
\end{Coreproposition}
\begin{proof}
Expand the product $\Phi(T)\Phi(U)=CTRCUR=CTPUR$ and subtract it from $CTUR$. The identity for full split encoding follows from $SA=I_V$ as computed above. These computations make no approximation and require no orthogonality.
\end{proof}
\Coresource{Sufficient Resolution [SR], Theorems 2.1 and 2.3. The present proof reviews only these algebraic identities, not the complete forty-seven-page source or its full formal implementation.}

\begin{Corecounterexample}[A discarded coordinate returns]\label{Core:ce:return}
Take $V=\mathbb R^2$, let $C(x_1,x_2)=x_1$, let $R(z)=(z,0)$, and let $T(x_1,x_2)=(x_2,x_1)$. The compressed one-step map is zero. Yet $T^2=I_V$, so $\Phi(T^2)=I_Z$, while $\Phi(T)^2=0$. The discarded second coordinate was not irrelevant to the two-step question.
\end{Corecounterexample}

This is the mathematical reason that a short successful readout is not enough. A representation may preserve what is visible now and lose a route by which hidden information becomes visible later. Packets identify the retained and residual components, but the decomposition alone does not prove the residual harmless.

\section{A method can be transported without reconstructing everything}
Suppose source states $z\in Z$ have operations $F_a$, target states $u\in U$ have operations $G_{\tau(a)}$, and $T:Z\to U$ is a translation. The operation-name map $\tau$ is fixed. Let $q$ and $\bar q$ be matched readouts.

\begin{Coretheorem}[Forward transport of a reusable method]\label{Core:r:transport}
If the source domain is invariant, and
\[
TF_a=G_{\tau(a)}T,\qquad q=\bar qT,
\]
then for every finite word $w$,
\[
TF_w=G_{\tau(w)}T,\qquad qF_w=\bar qG_{\tau(w)}T.
\]
\end{Coretheorem}
\begin{proof}
The empty word gives $TI_Z=I_UT$. Assume the first equality for a prefix. Invariance allows the next local equality to be applied on its successor, yielding
$TF_aF_w=G_{\tau(a)}TF_w=G_{\tau(a)}G_{\tau(w)}T$.
Induction proves the first claim. Composition with the matched readout proves the second.
\end{proof}
\Coresource{[MU], Theorem 6.1. The same local-to-continuation argument appears in OSL and Neuralese. No claim about a learned latent translator is imported.}
\Coreboundary{The theorem is forward. It does not imply an inverse, matching of new target operations, cheap computation of $T$, or that the source solver causally uses the represented state.}

The sensor representation $(y,s)$ translates to $(b,s)$ through $T(y,s)=(y\oplus s,s)$. Flipping the physical bit complements the first coordinate in either representation. The square commutes, and so does every repeated flip. The nuisance bit is never reconstructed because it has no role in this contract.

\section{One quotient for several possible dynamics}
A stronger situation arises when the operator itself is uncertain. Let $V$ be a finite-dimensional vector space over a fixed field, and let $Y$ be a vector space over the same field. Let $T_a^\theta:V\to V$ be a finite family of candidate linear operators indexed by an action $a$ and a model $\theta$. Choose one model $\theta_0$ as the reference and put $A_a=T_a^{\theta_0}$. A model index $\theta$ is fixed during the initial comparison. Let $Q:V\to Y$ be the same linear readout for every candidate.

Define $D_a^\theta=T_a^\theta-A_a$ and the propagated discrepancy space
\[
E=\CoreSpan\{A_wD_a^\theta v: v\in V,\ a,\theta,\ w\text{ finite}\}.
\]
It can be constructed by starting with the sum of the discrepancy ranges and repeatedly adding their images under the reference actions. In finite dimension this increasing sequence stops. Immediate invisibility $QD_a^\theta=0$ alone does not guarantee $QE=0$.

\begin{Coretheorem}[A common observable quotient]\label{Core:r:common}
The following conditions are equivalent:
\begin{enumerate}
\item $QT_w^\theta=QA_w$ on all of $V$, for every fixed candidate $\theta$ and every word $w$.
\item $E\subseteq\ker Q$.
\item There is a vector space $Z$, a surjective linear map $r:V\to Z$, common linear updates $G_a:Z\to Z$, and a linear readout $\bar Q:Z\to Y$ with $rT_a^\theta=G_ar$ and $Q=\bar Qr$ for every $a,\theta$.
\end{enumerate}
Under these conditions, $V/E$ is a common quotient and also predicts the output when the candidate index changes between operations.
\end{Coretheorem}
\begin{proof}
Assume the first condition. For every $a,\theta,w$, equality as operators on all of $V$ gives
\[
QA_wD_a^\theta=QT_w^\theta T_a^\theta-QA_wA_a=0.
\]
The last equality applies the premise to the word executing $a$ and then $w$. Spanning proves $E\subseteq\ker Q$.

Conversely, $E$ is invariant under every $A_a$ by construction, and contains the entire range of each difference $D_a^\theta$. Hence $T_a^\theta E\subseteq E$, and every candidate induces the same map $G_a([v])=[A_av]$ on $V/E$. Since $QE=0$, the readout descends. This proves the third condition. Composing its equalities yields the first condition. The same induction permits a different $\theta$ at each step because the quotient map and update remain shared.
\end{proof}
\Coresource{[GA], Theorem 13.1, ``Exact common quotient for structured model uncertainty.'' The all-initial-vectors hypothesis is retained.}

For example, let $Q(x_1,x_2,x_3)=x_1$ and let every candidate act as
\[
T^\theta(x_1,x_2,x_3)=(x_1+x_2,x_2,\theta x_3).
\]
Then $E$ is contained in the third coordinate and remains invisible. The two-coordinate method $(x_1,x_2)\mapsto(x_1+x_2,x_2)$ is exact, although the hidden dynamics are not identified.

The conclusion would be different if a later action moved the third coordinate into the first. That change enlarges $E$ beyond $\ker Q$. It invalidates the use of this quotient, not the theorem that justified the former one.


\section{An algebraic packet is not automatically a probability}
The previous sections concern linear representations. Bayesian correction will require more: a nonnegative state with a well-defined total mass. Those requirements belong to a declared domain, not to the spelling of its coordinates.

Let $\mathcal K$ be a convex cone in a real vector space, and let $\lambda$ be linear and strictly positive on every nonzero member of $\mathcal K$. Its normalized states satisfy $v\in\mathcal K$ and $\lambda v=1$. A coordinate change can send $\mathcal K$ to a cone containing negative coordinates without changing those states' probability meaning. Conversely, an arbitrary signed component of a valid state need not itself be valid.

\begin{Coreproposition}[When complete packets are genuine probability components]\label{Core:r:positive-packets}
Let $P_1,\ldots,P_m$ be linear maps with $\sum_iP_i=I$. If every $P_i$ preserves $\mathcal K$, each normalized $v\in\mathcal K$ has the mixture decomposition
\[
 v=\sum_{i:w_i>0}w_i v_i,\qquad
 w_i=\lambda P_iv\ge0,\quad \sum_iw_i=1,\quad v_i=P_iv/w_i\quad(w_i>0).
\]
Every zero-weight component is zero. If $\mathcal K=\mathbb R^d_{\ge0}$ and the $P_i$ are also idempotent, they must be disjoint diagonal coordinate selections. Orthogonal completeness alone does not imply a probability mixture.
\end{Coreproposition}
\begin{proof}
Cone preservation places $P_iv$ in $\mathcal K$, so its mass is nonnegative. Strict positivity says that a member of this cone with zero mass is zero. Applying $\lambda$ to the complete decomposition gives $\sum_iw_i=\lambda v=1$. Dividing every nonzero component by its mass proves the mixture identity.

On the nonnegative coordinate cone, preservation is equivalent to nonnegative matrix entries, as is seen by applying a matrix to every standard basis vector. Because the matrices sum to $I$, their nonnegative off-diagonal entries must all vanish. An idempotent diagonal entry satisfies $x^2=x$, hence is zero or one. The diagonal sums are one, so each coordinate is selected by exactly one map. Finally, completeness and orthogonality without positivity do not give the first premise, as the following example shows.
\end{proof}
\Coresource{Large Signal Model Foundations [LSM], Proposition 30, physical page 40. Only this finite packet argument and the joint-closure result below are admitted from that source at this stage.}

\begin{Corecounterexample}[An exact detail component has negative probability mass]\label{Core:ce:signed-packet}
For $p=(3/4,1/4)^\top$, the orthogonal average and detail projectors are
\[
 P_{\rm av}=\tfrac12\begin{pmatrix}1&1\\1&1\end{pmatrix},\qquad
 P_{\rm det}=\tfrac12\begin{pmatrix}1&-1\\-1&1\end{pmatrix}.
\]
They sum to the identity and produce $(1/2,1/2)^\top$ and $(1/4,-1/4)^\top$. The components reconstruct $p$ exactly, but the detail has zero total mass without being zero. Normalizing it as a belief is undefined. Clipping its negative entry and normalizing the resulting sum would instead produce $(3/5,2/5)^\top$, a different belief.
\end{Corecounterexample}

\begin{figure}[ht]
\centering\includegraphics[width=.92\linewidth]{sources/research/machine-native-reasoning-wider/figures/signed_packets.pdf}
\caption{Exact reconstruction and probability interpretation are different requirements. The signed detail is useful algebraically but is not a separate probability distribution. Values are from Counterexample~\ref{Core:ce:signed-packet}.}
\end{figure}

This distinction supplies the interface to the next chapter. A signal packet may describe a discrepancy, a coordinate, or a component of an operator. A belief represents nonnegative mass over alternatives. Moving between them requires a domain-preserving map. No learned signal architecture is inferred from this algebraic admission.

\section*{What is now available}
We have a declared method that preserves the selected outputs. But a predictive quotient need not retain the relationships needed for new evidence. The next chapter changes the retained object from a current physical coordinate to a set or distribution of candidate worlds.


\chapter{Memory that remains open to correction}\label{Core:book:chapter-3}
\setcounter{section}{0}\renewcommand{\thesection}{\thechapter.\arabic{section}}\label{Core:ch:memory}
\Corequestion{A useful memory preserves the distinctions by which its own conclusions can be revised.}
Let $t:\Omega\to Y$ be the current target. A candidate set $V\subseteq\Omega$ records worlds still possible. The empty set is allowed and means that the represented candidates contradict the admitted evidence. It is not a successful identification.

Let $\mathcal E\subseteq\CorePow(\Omega)$ contain $\Omega$ and be closed under intersection. In this chapter the world does not change when an evidence event is received. The update is $V\mapsto V\cap E$. Physical actions and their effects will return in Chapter~\ref{Core:ch:inquiry}.

\section{A state sufficient for every later filter}
Define the correction signature
\[
\Sigma(V)=\bigl(t(V\cap E)\bigr)_{E\in\mathcal E}.
\]
It records what each admitted future filter would leave of the possible target values.

\begin{Coretheorem}[Minimal exact memory for a correction contract]\label{Core:r:memory}
Equality of $\Sigma$ is a congruence for every update $V\mapsto V\cap E_0$. The signature supports the current target set and its exact update after every finite evidence sequence. Any deterministic retained state supporting the same readouts and updates on all candidate subsets must distinguish every two distinct signatures. Thus the image of $\Sigma$ has the least possible cardinality for this contract.
\end{Coretheorem}
\begin{proof}
If $\Sigma(V)=\Sigma(W)$, then for each $E$,
\[
t((V\cap E_0)\cap E)=t(V\cap(E_0\cap E))=t(W\cap(E_0\cap E)).
\]
Intersection closure supplies the coordinate $E_0\cap E$, so the updated signatures agree. The $\Omega$ coordinate is the present target set. Induction gives every finite sequence.

If signatures differ, one admitted event produces distinct target sets. A memory that merged those candidates would update identically and return the same answer on that event, contradicting exactness. Hence each signature needs its own state. The signature construction attains the lower bound.
\end{proof}
\Coresource{[GA], Theorem 5.2. This is a behavioral quotient of the correction process, not a claim that all candidate sets must be stored literally.}

\section{An exact count for partition evidence}
Suppose $C_1,\ldots,C_m$ partition $\Omega$ into nonempty cells and the admitted events are \emph{exactly} their unions. Let
\[
D_t=\{(i,y): C_i\cap t^{-1}(y)\ne\varnothing\},\qquad r=|D_t|.
\]
The occupancy memory retains whether a candidate exists in each attainable cell--target pair.

\begin{Corecorollary}[The size of partition correction memory]\label{Core:r:occupancy}
On all subsets of $\Omega$, correction signatures are exactly distinguished by these occupancy records. There are $2^r$ signatures, and $r$ fixed binary bits are necessary and sufficient to encode them.
\end{Corecorollary}
\begin{proof}
An individual-cell filter reveals the target values occupied in that cell, so the signature determines occupancy. Conversely, occupancy determines the target values left by any union of cells. Every subset of $D_t$ occurs: select one world from each included nonempty cell--target intersection and none from excluded intersections. Those intersections are disjoint. Thus there are precisely $2^r$ possible records. The preceding theorem and binary counting give necessity, while the occupancy vector supplies sufficiency.
\end{proof}
\Coresource{[GA], Corollary 5.3. The word ``exactly'' in the evidence premise is essential.}

For this sensor contract, the four cells are fixed pairs $(b,s)$ with either nuisance value. We admit the entire Boolean algebra generated by calibration and display events, including every union of its cells. This is a choice of evidence language, not a claim that a single calibration physically implements every union event. The target is $b$, so there are four attainable cell--target pairs. Candidate membership uses eight bits for eight worlds. Correction memory uses four bits for sixteen possible occupancies. The current target set alone has four possibilities and needs only two bits, but is not sufficient for later calibration.

For example,
\[
V=\{(0,0,0),(1,1,0)\},\qquad W=\{(0,1,0),(1,0,0)\}
\]
have the same possible target set $\{0,1\}$. After the event $s=0$, one leaves target zero and the other target one. Their different associations are operationally relevant.

\begin{Corecounterexample}[A new evidence event changes the required memory]\label{Core:ce:split}
Take one cell $\Omega=\{u,v\}$ and a constant target. The singleton candidate sets have the same old occupancy. If the event $\{u\}$ is added, filtering one singleton gives a target value while filtering the other gives contradiction. A record sufficient for the old event algebra is insufficient for the larger one.
\end{Corecounterexample}

\section{Noisy correction is not elimination}
A likelihood can be small without being zero. Deleting a candidate because an observation was surprising can invalidate a correct noise model. We therefore specify a separate probability contract.

Let $\mathcal M$ be a finite nonempty model set with prior $p(m)$. The sampled model remains fixed during the experiment. For each $m$, let $L_m$ be a probability distribution on a finite nonempty observation alphabet $\mathcal O$. Let $\mathcal Z$ be a finite nonempty report alphabet. A stochastic reporting channel $B(o,z)$, conditionally independent of the model given $o$, creates a retained report $Z\in\mathcal Z$. The joint law is
\[
\Pr(m,o,z)=p(m)L_m(o)B(o,z),\qquad Q_m(z)=\sum_oL_m(o)B(o,z).
\]
Posterior distributions are defined only on positive-probability events.

\begin{Coretheorem}[When a report preserves every Bayesian correction]\label{Core:r:bayes}
The following are equivalent:
\begin{enumerate}
\item For every prior, the posterior after $o$ equals that after $z$ on every positive-probability pair $(o,z)$.
\item The equality holds for one prior positive on every model.
\item There is a stochastic kernel $\rho(o\mid z)$ such that
\[
L_m(o)B(o,z)=Q_m(z)\rho(o\mid z)\quad\text{for all }m,o,z.
\]
\end{enumerate}
The reconstruction in the last condition refers to the actual joint experiment, not an unrelated simulation with matching marginals.
\end{Coretheorem}
\begin{proof}
The first condition implies the second. Under a positive prior, define $\rho(o\mid z)=\Pr(o\mid z)$ whenever $z$ has positive probability. For a positive pair, posterior equality yields
\[
p(m)L_m(o)B(o,z)=\Pr(o,z)p(m\mid o)
=\rho(o\mid z)\Pr(z)p(m\mid z).
\]
Substituting Bayes' formula and cancelling positive $p(m)$ gives the factorization. A zero-probability pair under a positive prior has $L_m(o)B(o,z)=0$ for all $m$. Its conditional weight is zero. If $z$ is impossible for every model, assign any probability row to $\rho(\cdot\mid z)$ because both sides vanish.

Conversely, the factorization makes the likelihood of $(o,z)$ a model-independent multiple of the likelihood of $z$. Normalization under any prior therefore gives $p(m\mid o,z)=p(m\mid z)$. Since the common channel only reads $o$, the same positive-pair posterior equals $p(m\mid o)$. This proves the first condition.
\end{proof}
\Coresource{[MU], Theorem 12.1. The full-support and zero-probability cases are included in the proof.}

The theorem permits removal of an independent nuisance flag. It does not permit removal of every association. Consider two priors on $(b,s)$, with rows indexed by $b$ and columns by $s$:
\[
P=\begin{pmatrix}1/2&0\\0&1/2\end{pmatrix},\qquad
P'=\begin{pmatrix}0&1/2\\1/2&0\end{pmatrix}.
\]
Both individual marginals are uniform. A noisy calibration reports $s$ correctly with probability $3/4$, independently conditional on the fixed world. The probability of report one is $1/2$ under both priors. After that report, however, the target posteriors are $3/4$ and $1/4$. Separate target and sensor marginals cannot support this correction. Their association matters even though neither marginal changed.

These examples do not identify a universally optimal memory. They identify the least information required by a specified correction contract. A new likelihood, a newly admitted model, or a different target can change that requirement.


\section{A record must close under prediction as well as correction}
The preceding Bayesian theorem preserves the posterior over a fixed model family after one experiment. A continuing reasoner also predicts what happens between observations. The same compressed record must support both operations. Otherwise a valid posterior can become an inadequate input to the very next prediction.

First consider a known finite state set $X$ with row-stochastic transition matrices $K_a(x,x')$. We retain column-vector probability distributions. A partition $\pi$ has aggregation $C_\pi(z,x)=1_{\{\pi(x)=z\}}$, so prediction is $p\mapsto K_a^\top p$. For each permitted observation symbol $o$, let $\ell_o(x)>0$ be its specified likelihood. When these are a complete report alphabet, assume also $\sum_o\ell_o(x)=1$.

A mandatory partition $\pi_0$ with at least two blocks protects the target distinctions. We seek a refinement that supports every prediction and each separate Bayesian update for every prior. The observation is an operation available in its own right. This is stronger than inspecting only a chosen combined transition-and-report schedule.

\begin{Coretheorem}[Least refinement closed under prediction and correction]\label{Core:r:joint-refinement}
First split $\pi_0$ by the complete vectors $(\ell_o(x))_o$, obtaining $\pi_{\rm init}$. Then repeatedly split current blocks by
\[
 \left(\sum_{x'\in B}K_a(x,x')\right)_{a,B},
\]
where $B$ runs over the current blocks. The process stops after at most $|X|-|\pi_{\rm init}|$ strict refinement rounds. Its final partition is the unique coarsest refinement of $\pi_0$ for which all specified predictions admit stochastic coarse transitions and all separate Bayesian updates admit fixed strictly positive coarse likelihoods for every prior.
\end{Coretheorem}
\begin{proof}
At a stable partition the probability of reaching each coarse block depends only on the current block. These sums define a stochastic coarse transition and yield $C_\pi K_a^\top=\bar K_a^\top C_\pi$, first on point masses and then by linearity. Likelihoods are constant within final blocks. Grouping the numerator and normalizer by those blocks therefore commutes with every Bayesian update.

For minimality, take any admissible refinement $\eta$ of $\pi_0$. It has at least two blocks. Fix a strictly positive fine likelihood $\ell$ and a strictly positive coarse likelihood $\bar\ell$ inducing its assumed update. For states $x,y$ in one block and $z$ in a different block, compare priors assigning half their mass to $x,z$ and to $y,z$. Their coarse priors agree. Equality of the posterior odds gives
\[
 \ell(x)/\ell(z)=\bar\ell(\eta(x))/\bar\ell(\eta(z))=\ell(y)/\ell(z).
\]
Thus $\ell(x)=\ell(y)$, and $\eta$ refines the first likelihood split. Coarse likelihoods may differ by a positive scale for each observation without changing this conclusion.

If $\eta$ refines a current partition, each current block is a union of $\eta$ blocks. Exact prediction under $\eta$ makes the probability of reaching each such union constant inside every $\eta$ block. Hence $\eta$ refines the next split too. Induction proves that every admissible refinement refines the terminal one. Each strict round increases the block count, proving the bound. Coarseness determines the partition uniquely, up to naming its blocks.
\end{proof}
\Coresource{[LSM], Proposition 9 and Proposition 12, physical pages 11--14. Bayesian necessity also appears in PRIQ [PRIQ], Theorem 5.6. The proof includes the two-block and strict-positivity premises rather than importing them by a cross-reference.}

For the sensor, the mandatory display partition has two blocks. A noisy calibration that reports $s$ correctly with probability $3/4$ splits it into the four $(b,s)$ blocks. The nuisance bit remains merged. The physical flip maps these four blocks among themselves, so the transition stage introduces no further split. The representation is now closed for both prediction and correction, not merely for today's display.

Refinement can also propagate backward through dynamics. In a separate four-state example, the mandatory partition is $\{0,1,2\}\mid\{3\}$ and the deterministic successors are $(0,2,3,3)$. A constant likelihood adds no information. Transition signatures first isolate state 2, then separate 0 from 1. The final partition is discrete. The reason for the second split is a future observable difference, not a newly sampled measurement.

The construction uses complete tables. It does not acquire those tables, certify their physical accuracy, or guarantee a finite encoding of every posterior coordinate. Its least-state statement concerns the specified separate operations and all priors. Different admissible schedules or reachable-prior restrictions can justify another quotient.

\section{Retain the joint law of the report and the next state}
A general experiment can change the state and report correlated information about the change. Describing its transition and report marginals separately is then insufficient. Let
\[
 K_{a,o}(x',x)=\Pr(X_{t+1}=x',O_{t+1}=o\mid X_t=x,A_t=a)
\]
be a nonnegative finite instrument with $\sum_{o,x'}K_{a,o}(x',x)=1$ for every $a,x$. This notation is column-oriented: $K_{a,o}p$ is an unnormalized post-experiment distribution. It specifies an experiment's physical change and its observation in one object.

For a surjection $\pi:X\to Z$, let $C$ sum probabilities over its fibres. Write
\[
 z_{a,o}(p)=\mathbf1^\top K_{a,o}p,\qquad
 \Psi_{a,o}(p)=K_{a,o}p/z_{a,o}(p)
\]
when the denominator is positive. Zero mass means an impossible report under that prior, not an arbitrarily assigned posterior.

\begin{Coretheorem}[One quotient for reports, posteriors, and feedback]\label{Core:r:instrument}
A nonnegative normalized coarse instrument $\bar K_{a,o}$ satisfies
\[
 CK_{a,o}=\bar K_{a,o}C\quad\text{for every }a,o
\]
if and only if, for all $a,o,z'$,
\[
 \sum_{x':\pi(x')=z'}K_{a,o}(x',x)
\]
is constant as $x$ varies within any one input fibre. Such a coarse instrument is unique. Equivalently, it preserves for every prior both the report probabilities and the aggregated positive-probability posteriors:
\[
 z_{a,o}(p)=\bar z_{a,o}(Cp),\qquad C\Psi_{a,o}(p)=\bar\Psi_{a,o}(Cp).
\]
Starting from $Cp_0$, these identities preserve every finite joint action/report law and every corresponding coarse posterior under one common policy that uses only the previously observed action/report history. This includes randomized policies with the same history-conditional action probabilities.
\end{Coretheorem}
\begin{proof}
Applying the matrix identity to a point mass proves necessity of the fibre condition. Conversely, use the common sums to define each column of $\bar K_{a,o}$. They are nonnegative. Summing across $z',o$ partitions the fine successor set and gives column mass one. The definition proves the matrix equality on point masses, hence on every vector. Surjectivity gives uniqueness.

Because $\mathbf1_Z^\top C=\mathbf1_X^\top$, the matrix equality gives equal report masses. Dividing it by a common positive mass gives equal aggregated posteriors. Conversely, equality of masses and positive-mass posteriors gives the unnormalized matrix equality on every prior. If the common mass is zero, both unnormalized vectors are zero by nonnegativity. Point masses finish that direction.

For a history $h$, keep unnormalized states $u_h,\bar u_h$. Initially $\bar u_\varnothing=Cu_\varnothing$. After choosing action $a$ with common probability $\gamma(a\mid h)$ and receiving $o$, the next states are $\gamma(a\mid h)K_{a,o}u_h$ and $\gamma(a\mid h)\bar K_{a,o}\bar u_h$. The intertwining identity preserves their aggregation. Induction gives equal mass for every history, including zero-mass histories. Normalization on positive histories gives the posterior statement. Stopping decisions based on the same history can be included as matching terminal choices.
\end{proof}
\Coresource{Selected synthesis of the positive-transfer and filtering arguments in [PRIQ], Theorems 5.2, 6.1 and 6.3, and the projected-law construction in AVOC [AVOC], Theorem 4.5. The joint instrument and feedback formulation is proved here. No entire nonlinear source theorem or Lean implementation is imported.}

\begin{Corecounterexample}[Correct marginals with an incorrect correction]\label{Core:ce:joint-marginals}
Two fine current states have the same retained label. From the first, the next retained state and report are $(0,0)$ or $(1,1)$, each with probability $1/2$. From the second they are $(0,1)$ or $(1,0)$, again equally likely. In both cases the next-state marginal is uniform and the report marginal is uniform. After report one, however, the next retained state is certainly one in the first case and certainly zero in the second. No common filter based only on the old label can realize both. The joint fibre condition detects exactly the lost association.
\end{Corecounterexample}

Preserving only normalized posteriors is weaker still. A report coin independent of the hidden state leaves every prior unchanged whether its success probability is $1/4$ or $3/4$. Those two instruments have identical posterior maps but different report laws. A planner that prices later branches needs those probabilities. In the exact instrument theorem the evidence scale is fixed, rather than silently discarded during normalization.

\section{Small law error and a rare conditional history}
Approximate models must be compared before declaring their posteriors close. Let $\widetilde K$ be a normalized instrument on the same coarse state and report alphabets. Suppose, at step $t$, every relevant action and fine state satisfy
\[
 \CoreTV\left(\left(\sum_{x':\pi(x')=z'}K_{a,o}(x',x)\right)_{o,z'},
             (\widetilde K_{a,o}(z',\pi(x)))_{o,z'}\right)\le\epsilon_t.
\]
This is a bound on a complete report/next-state distribution, not on an isolated matrix entry. Assume $\CoreTV(Cp_0,\widetilde p_0)\le\delta_0$ with nonnegative $\delta_0,\epsilon_t$. Both models run the same history-based policy. Coarse states recorded in the comparison below are analytical coordinates, not extra observations supplied to that policy.

\begin{Coreproposition}[A whole-history bound with an explicit conditioning cost]\label{Core:r:history-bound}
At horizon $H$, the distributions $P_H,\widetilde P_H$ of the full coarse-state and action/report histories satisfy
\[
 \CoreTV(P_H,\widetilde P_H)\le d_H:=\min\left\{1,\delta_0+\sum_{t=0}^{H-1}\epsilon_t\right\}.
\]
The same bound holds after forgetting any of those history coordinates. For an observed history event $E$ with probabilities $p_E,\widetilde p_E>0$, the conditional final coarse-state distributions obey
\[
 \CoreTV\bigl(P_H(Z_H\in\cdot\mid E),\widetilde P_H(Z_H\in\cdot\mid E)\bigr)
 \le\min\left\{1,\frac{d_H}{\max(p_E,\widetilde p_E)}\right\}.
\]
If either event probability exceeds $d_H$, the other is necessarily positive.
\end{Coreproposition}
\begin{proof}
Condition on a complete coarse history with positive fine probability. Its hidden current-state distribution is a mixture of states in the last coarse fibre. Convexity of total variation makes the true next report/coarse-state row differ from the declared row by at most $\epsilon_t$. Mixing over a common action kernel preserves this bound. Extend the two old history laws first through the same declared row. Total variation cannot increase under a common stochastic kernel, by the triangle inequality and row sums. Replacing those rows with the true conditional rows adds at most $\epsilon_t$, weighted by a total history mass of one. Hence the history distance at the next step is at most the old distance plus $\epsilon_t$. Induction and the cap one prove the first claim. Summing over forgotten coordinates is also a stochastic map, proving the second.

For conditioning, let $u,v$ be the restrictions of the two history laws to $E$, with masses $p_E,\widetilde p_E$. The difference on the complement has $\ell^1$ norm at least $|p_E-\widetilde p_E|$. Thus
\[
 \|u-v\|_1+|p_E-\widetilde p_E|\le2d_H.
\]
Adding and subtracting $v/p_E$ bounds the conditional distance by $d_H/p_E$. Interchanging the laws gives $d_H/\widetilde p_E$, so the smaller of these bounds has the displayed maximum denominator. Projection to the final state cannot increase the distance. Finally, the probability of any event changes by at most total variation, which proves positivity when one mass exceeds $d_H$.
\end{proof}
\Coresource{Finite probability-instrument specialization of [LSM], Propositions 43--44, with the joint retained-state record made explicit. The symmetric conditioning form follows by applying the same elementary bound in both directions. [PRIQ], Section 9 supplies the complementary stepwise posterior-error perspective.}

\begin{Corecounterexample}[A rare report carries the entire discrepancy]\label{Core:ce:rare-history}
Two joint laws place mass $1-\delta$ on the same ordinary outcome. On a report event $E$ of mass $\delta>0$, one puts all mass on retained state zero and the other on retained state one. Their whole-law distance is $\delta$, while their conditional distance after $E$ is one. The factor involving event probability is therefore essential. This is a finite-law statement, not an empirical estimate of a sensor error.
\end{Corecounterexample}

\begin{figure}[ht]
\centering\includegraphics[width=.92\linewidth]{sources/research/machine-native-reasoning-wider/figures/rare_history.pdf}
\caption{In Counterexample~\ref{Core:ce:rare-history}, shrinking the rare-event mass reduces whole-law error but leaves conditional error maximal. The points are exact specified probability laws, not sampled observations.}
\end{figure}

The bound compares one policy across two models, not independently optimized controllers. Chapter~\ref{Core:ch:inquiry} supplies the uniform comparison needed for re-optimization. First, the next chapter asks what decisions and explanations the corrected record can justify.


\chapter{From corrected records to decisions and explanations}\label{Core:book:chapter-4}
\setcounter{section}{0}\renewcommand{\thesection}{\thechapter.\arabic{section}}\label{Core:ch:decision}
\Corequestion{The same record, model, and continuation must support the decision being certified.}
Correctness of a probability calculation is not a guarantee that a real process follows its likelihood model. Within the declared model, however, a small error in the retained distribution can support a precise decision bound. We state that bridge before introducing communication, so that an explanation has a definite object to preserve.

\section{One quantitative decision contract}
Let $p$ and $\widehat p$ be distributions on the same finite world set. For each world let $v_\omega$ be a response and $\widehat v_\omega$ its approximation in one metric space. Assume
\[
\CoreTV(p,\widehat p)\le e,\qquad
\sup_\omega d(v_\omega,\widehat v_\omega)\le d_0.
\]
For each action $a$ in a finite nonempty set, let the loss $\ell(a,v)$ lie in $[0,1]$ and be $K$-Lipschitz in $v$. Define the model risks
\[
R(a)=\sum_\omega p(\omega)\ell(a,v_\omega),\qquad
\widehat R(a)=\sum_\omega\widehat p(\omega)\ell(a,\widehat v_\omega).
\]

\begin{Coretheorem}[A decision bound for one actual approximation chain]\label{Core:r:risk}
Put $\xi=\min\{1,e+Kd_0\}$. For every action, $|R(a)-\widehat R(a)|\le\xi$. If $\widehat a$ minimizes the approximate risk, then
\[
R(\widehat a)-\min_a R(a)\le\min\{1,2\xi\}.
\]
If the approximate gap between its unique minimizer and every alternative exceeds $2\xi$, that action is also the unique true model-risk minimizer.
\end{Coretheorem}
\begin{proof}
Add and subtract $\sum_\omega p(\omega)\ell(a,\widehat v_\omega)$. The response term has absolute value at most $Kd_0$. The expectation of a function in $[0,1]$ changes by at most total variation, giving the other term at most $e$. Risks lie in $[0,1]$, so their difference is also at most one.

For a true minimizer $a_*$,
\[
R(\widehat a)\le\widehat R(\widehat a)+\xi
\le\widehat R(a_*)+\xi\le R(a_*)+2\xi.
\]
The loss range supplies the cap one. Finally, for any other $a$, subtract the two error bounds to obtain
$R(a)-R(\widehat a)\ge\widehat R(a)-\widehat R(\widehat a)-2\xi>0$.
\end{proof}
\Coresource{[GA], ``A common quantitative grounding certificate'' and [MU], ``A joint posterior and recipient decision certificate.'' This is a specialization to a common finite candidate set and bounded loss.}

The phrase ``model risk'' matters. The theorem cannot repair an omitted candidate, a false sensor convention, or an action outside the stated loss function. It also cannot combine a posterior bound from one history with a response bound from another. The record must bind the actual calculation and its premises.

For the noisy calibration example, use zero-one loss for selecting the target bit. Under $P$ and report one, choosing one has risk $1/4$, while choosing zero has risk $3/4$. An approximation with total variation at most $1/10$ and no response error preserves the unique decision because the gap $1/2$ exceeds $2/10$. Under $P'$ the correct decision reverses. A valid error bound around the wrong prior is not a substitute for retaining the correct association.

\section{An explanation belongs to a recipient}
An artifact is not an explanation merely because its producer can interpret it. Fix a finite case set $X$, a sender-access map $\alpha:X\to H$, and a finite artifact catalogue $\mathcal C$. Recipient $i$ has a background record $\beta_i(x)$, a demanded output $t_i(x)$, and a fixed bounded procedure $L_i$. A timeout, malformed artifact, wrong output, or budget overrun counts as failure.

Let $\mathcal C(x)$ be the artifacts that succeed for every recipient in case $x$. For an attained sender value $h$, put
\[
\mathcal C_h=\bigcap_{x:\alpha(x)=h}\mathcal C(x).
\]

\begin{Coretheorem}[A common artifact for fixed recipients]\label{Core:r:recipient}
A sender restricted to $\alpha(x)$ can provide one exact artifact for all recipients if and only if every $\mathcal C_h$ is nonempty. When feasible, the smallest possible number of distinct transmitted artifacts is
\[
\min\{|C|:C\subseteq\mathcal C,\ C\cap\mathcal C_h\ne\varnothing
\text{ for every attained }h\}.
\]
\end{Coretheorem}
\begin{proof}
A sender seeing $h$ must choose the same artifact for all cases in that fibre. It must therefore select a member of $\mathcal C_h$, proving necessity. Choosing one member of each nonempty finite intersection gives a valid encoder. The image of any encoder intersects every $\mathcal C_h$. Conversely, any catalogue meeting all those sets supplies such an encoder by choosing from the intersections. Finite minimization gives the formula.
\end{proof}
\Coresource{[MU], ``The fixed-recipient covering criterion.'' The earlier designed-decoder coloring criterion is not substituted for the installed-recipient problem.}

\begin{Corecounterexample}[Individual acceptability is not a common message]\label{Core:ce:message}
Two cases have the same sender record. In one, the sole acceptable artifact is $c_0$. In the other it is $c_1\ne c_0$. Each case has a valid explanation if its hidden identity is known. Their intersection is empty, so the sender cannot select a valid common explanation from its actual information.
\end{Corecounterexample}

In the sensor, a sender observing $y$ can transmit one bit to a recipient holding $s$. The recipient's fixed XOR procedure returns $b=y\oplus s$. The sender need not know $b$: its message combines with the recipient's background.

A decoder that ignores this message fails where background alone cannot distinguish the target. More strings do not fix that procedure. Installed information, interpreter work, and message size are separate resources.


\chapter{One investigation, one resource budget}\label{Core:book:chapter-5}
\setcounter{section}{0}\renewcommand{\thesection}{\thechapter.\arabic{section}}\label{Core:ch:inquiry}
\Corequestion{A possible answer, an affordable observation, and a small program must belong to the same workflow.}
Let $X$ be a finite nonempty state space. It includes every origin tag, permission, background record, and history variable that affects the stated behavior. An experiment $e$ is allowed on a domain $D_e\subseteq X$, has a deterministic successor $f_e:D_e\to X$, a report $o_e:D_e\to O_e$, and a fixed nonnegative integer resource vector $c_e$.

A current support $V\ne\varnothing$ records possible states. The experiment is permitted only if $V\subseteq D_e$. After report $o$, the possible \emph{post-experiment} states are
\[
V_{e,o}=\{f_e(x):x\in V,\ o_e(x)=o\}.
\]
Only nonempty report branches are retained. If the target concerns a past value, that value must survive as an immutable tag. Resetting the current state is not recovering its history.

The experiment catalogue and the terminal-artifact catalogue are finite, and the resource vector has a fixed finite dimension. A terminal artifact $a$ is acceptable only if every fixed recipient succeeds in every state of $V$. Its worst-case terminal cost is $g_a(V)$, computed componentwise. A controller is a literal finite tree with one record per experiment or terminal node. Stored memory counts all branches, including those not visited in one execution.

\section{The complete finite frontier}
Let $\mathcal P_h(V)$ be the nondominated pairs $(v,m)$ of worst-case online resource vectors and stored-record counts attainable by successful trees with at most $h$ experiment steps. At depth zero the available pairs are
\[
\operatorname{Pareto}\{(g_a(V),1):a\text{ succeeds throughout }V\}.
\]
At greater depth retain those stops and, for each uniformly permitted experiment and one feasible child pair $(v_o,m_o)$ for every nonempty branch, form
\[
\left(c_e+\max_o v_o,\quad1+\sum_o m_o\right).
\]
A profile dominates another when it is no larger in every resource coordinate and strictly smaller in at least one. The maximum is componentwise. An impossible child contributes no tree. It is not a zero-cost branch.

\begin{Coretheorem}[Exact joint feasibility for explicit controllers]\label{Core:r:inquiry}
This recurrence gives exactly the nondominated resource profiles of successful depth-bounded controllers. Every retained pair has a concrete tree witness. A common online and storage budget is feasible if and only if it bounds one retained pair. Removing dominated child profiles does not change feasibility under any componentwise budget.
\end{Coretheorem}
\begin{proof}
At zero depth the only possibilities are acceptable terminals, whose profiles are exactly the stated stops. At positive depth a successful tree either stops or begins with one uniformly allowed experiment. Every possible report needs a successful child on precisely $V_{e,o}$. By induction those children have profiles in the lower-depth attainable set. Worst-case online cost is the root cost plus their componentwise maximum. Complete tree storage is one plus the sum of all child records.

Conversely, combining any allowed root with a successful child for every report constructs a successful tree with that profile. This proves exact enumeration before pruning. Replacing a dominated child by a dominating one cannot increase either a coordinatewise maximum or a sum of record counts, and does not change its supported initial set. Thus Pareto pruning preserves every feasible budget and every retained profile has a witness.
\end{proof}
\Coresource{[MU], ``A joint resource certificate for fixed recipients,'' and [SII], Theorem 6.2. The scalar-cost repair case is a specialization, not a different composition theorem.}

The recurrence can be expensive. It is a finite mathematical characterization, not an efficient general learner. It assumes its transition and acceptance tables. Their discovery and checking are preparation costs, not free consequences of the existence theorem.

\begin{Corecounterexample}[Two separate optima do not form one controller]\label{Core:ce:pareto}
There are four immutable cases, each requiring its own terminal answer. A direct observation of the case costs four units and gives a five-record tree including four terminals. Unit-cost equality tests can identify the case in at most three tests with seven records. Terminal answers have zero online cost in this particular example. The frontier is $(4,5)$ and $(3,7)$. The mixed budget $(3,5)$ is infeasible: an all-equality successful tree needs four leaves and at least seven records. Any tree using the direct observation on a reachable path costs at least four.
\end{Corecounterexample}

\begin{figure}[ht]
\centering\includegraphics[width=.88\linewidth]{sources/research/machine-native-reasoning-wider/figures/joint_frontier.pdf}
\caption{The four-origin example has two nondominated profiles. Combining the lowest coordinate from each does not construct a feasible controller. Terminal answers have zero online cost in this example.}
\end{figure}

\section{When a checked candidate can end the search}
Consider $n\ge1$ immutable cases with total nondestructive binary unit-cost queries. Every unit-cost repair succeeds in exactly one case, and every case has a repair. A controller must succeed certainly in every case.

\begin{Coretheorem}[A candidate attaining both lower bounds is optimal]\label{Core:r:lowerbound}
Every successful controller has worst-case work at least $1+\lceil\log_2n\rceil$ and literal storage at least $2n-1$. A checked controller attaining both bounds and the requested horizon realizes the entire nondominated work/storage frontier.
\end{Coretheorem}
\begin{proof}
Distinct cases cannot end at the same successful repair leaf, so at least $n$ leaves are reachable. A binary tree of query depth $d$ has at most $2^d$ leaves, giving $d\ge\lceil\log_2n\rceil$. The terminal repair adds one unit.

Let $l$ be the number of leaves and $i_1,i_2$ the counts of unary and binary internal nodes. Counting tree edges gives $i_1+2i_2=i_1+i_2+l-1$, hence $i_2=l-1$. Total records are $2l-1+i_1\ge2n-1$. A successful candidate at both coordinatewise lower bounds dominates every other attainable cost profile. Its validity and horizon check make it an actual member of the frontier, leaving that pair as the whole nondominated frontier.
\end{proof}
\Coresource{[SII], Theorem 11.1. The assumptions exclude destructive resets and repairs that succeed on several cases.}

For eight possible defective cells, equality-first diagnosis requires seven queries in the worst case, while three coordinate-bit queries identify the cell. Including repair gives costs eight and four, respectively. Both literal trees have fifteen records. Faster execution has not eliminated the cost of storing the procedure.

This eight-cell example is separate from the eight-world sensor. The sensor has only four action-relevant $(b,s)$ pairs and an irrelevant nuisance coordinate. Confusing hidden-world cardinality with the number of required repair distinctions would import a false lower bound.


\section{Expected-cost inquiry uses the report probabilities}
The preceding frontier certifies worst-case execution and complete tree storage for deterministic observations. A noisy report requires a different optimization. Expected cost weights each future branch by its probability. A low expectation is not a hard execution budget, and neither replaces the cost of storing the whole controller.

Use the finite instrument of Theorem~\ref{Core:r:instrument}, with fine belief $p$ and coarse belief $q=Cp$. At $q$, a finite set of experiments $A(q)$ is permitted, and a nonempty finite set of terminal decisions $D(q)$ is available. These catalogues must be identical for every fine belief with the same $q$. Stopping with decision $d$ has loss $\ell_d(\pi(x))$ in fine state $x$, so its conditional risk is $\langle q,\ell_d\rangle$. An experiment has a finite nonnegative charge $c_a(q)$, identical on the fine and coarse side. All history variables affecting permissions, charges, or targets must be included in the declared state.

Let $J_h(p)$ be the least expected charge plus terminal loss attainable using at most $h$ further experiments, with terminal decisions permitted at any step. It is defined by finite observation trees, not by a hidden-state controller. Its coarse recurrence is
\begin{align*}
 \bar J_0(q)&=\min_{d\in D(q)}\langle q,\ell_d\rangle,\\
 \mathcal Q_h(q,a)&=c_a(q)+\sum_{o:\bar z_{a,o}(q)>0}
       \bar z_{a,o}(q)\bar J_{h-1}\bigl(\bar\Psi_{a,o}(q)\bigr),\\
 \bar J_h(q)&=\min\left\{\bar J_0(q),\ \min_{a\in A(q)}\mathcal Q_h(q,a)\right\}.
\end{align*}
An empty experiment catalogue leaves only stopping. Impossible reports do not receive invented posteriors.

\begin{Coretheorem}[The same compressed belief supports the same inquiry]\label{Core:r:belief-planning}
Under the exact instrument quotient and the common loss, charge, and catalogue assumptions above,
\[
 J_h(p)=\bar J_h(Cp)\quad\text{for every prior and every finite }h.
\]
Every first-action value and optimal first-action set agrees. A minimizing coarse-belief policy produces a fine-model optimal history-based policy, including its stopping decision. Finite minimizing trees exist. The result supplies neither an efficient synthesis algorithm nor a worst-case storage or resource guarantee.
\end{Coretheorem}
\begin{proof}
At zero horizon the terminal risks agree by fibrewise summation. Suppose the equality holds at $h-1$. For any permitted first experiment, Theorem~\ref{Core:r:instrument} gives the same probability of each report and the same coarse posterior on every positive branch. The induction hypothesis gives equal optimal continuation costs on each such branch. The first charge agrees, so the entire action value agrees. Stopping has the same value as well. Taking the minimum over the identical catalogues proves both value equality and equality of minimizing first choices.

Choose one minimizing first choice and attach minimizing lower-horizon trees to its finitely many positive branches. If the choice is terminal, no child is required. Induction constructs a finite tree attaining the value. The same induction bounds every other admissible tree below by that minimum. An observed history initializes the same coarse belief in the fine execution, so selecting from that belief realizes the coarse policy without observing the fine hidden state.
\end{proof}
\Coresource{The finite backup and projected-law proof mechanisms of [AVOC], Theorems 4.5 and 9.2, applied here to the belief transition supplied by Theorem~\ref{Core:r:instrument}. The source-facing Machine Theory correspondence covers this selected result. That correspondence is separate from any claim that the complete AVOC implementation is verified. The established partially observed planning context is [KLC].}

\begin{Corecounterexample}[A correct probability quotient can hide a forbidden action]\label{Core:ce:permission-loss}
Two fine states share all report and transition probabilities and are given the same coarse label. A terminal repair is permitted in one state but forbidden in the other. The instrument quotient is valid, but a coarse interface that offers the repair in both states is not. Exact belief and value formulas do not repair a missing permission distinction. The catalogues must factor through the retained state, or the state must be refined.
\end{Corecounterexample}

\section{The sensor chooses when another calibration is worth its cost}
Return to the sensor after display $y=0$. Conditional on that display, assign equal probability to $(b,s)=(0,0)$ and $(1,1)$, with any nuisance distribution inside each pair. Only noisy calibration and terminal bit selection are offered in this example. A calibration reports $s$ correctly with probability $3/4$, with fresh conditionally independent noise and no change to the hidden world. It costs $1/20$ in the same stated loss units as a wrong terminal choice, whose loss is one. A correct terminal choice has loss zero.

The retained belief on the four $(b,s)$ pairs is an exact quotient of the eight-world belief. Its two diagonal masses initially equal $1/2$. After one calibration the favored target has probability $3/4$. A second matching report raises that probability to $9/10$, while conflicting reports restore $1/2$.

\begin{center}\begin{tabular}{@{}>{\raggedright\arraybackslash}p{.16\linewidth}>{\raggedright\arraybackslash}p{.24\linewidth}>{\raggedright\arraybackslash}p{.54\linewidth}@{}}\toprule
Report limit & Optimal expected cost & An optimal initial choice\\\midrule
0 & $1/2$ & Stop and select either bit.\\
1 & $3/10$ & Calibrate once, then select.\\
2 & $3/10$ & Calibrate once, then stop.\\
3 & $11/40$ & Calibrate, then use an adaptive continuation.\\\bottomrule\end{tabular}\end{center}

\begin{figure}[H]
\centering\includegraphics[width=.94\linewidth]{sources/research/machine-native-reasoning-wider/figures/calibration_policy.pdf}
\caption{An optimal three-report policy in the unchanged eight-world sensor with the explicitly added noisy calibration channel. All masses and costs are exact. The diagram groups sign-symmetric report histories, not distinct hidden states supplied to the agent.}
\end{figure}

With three reports available, the optimal policy takes a second report after the first. If the first two agree, it stops. If they disagree, it takes the third. The disagreement probability is $3/8$. Thus the expected report count is $2+3/8=19/8$, the expected charge is $19/160$, and the error probability is $5/32$. The total is $19/160+5/32=11/40$.

A second report alone does not reduce expected classification error enough to pay for itself. A third-report allowance makes it worthwhile to investigate only the unresolved branch. This gain is achieved by one contingent policy, without access to the hidden bit.



The two-valued display state is too coarse for this inquiry. The four-pair probability record supports calibration and decisions while still discarding the nuisance bit.

\section{An approximate model supports a different, quantitative claim}
Exact policy transport and approximate planning must not share one unqualified guarantee. Fix a finite horizon $H$ and one nonempty finite class $\mathfrak P$ of admitted action/report trees, shared by the true and approximate models. The complete loss of a tree is the same function of its coarse-state and action/report path in both models, taking values in $[0,B]$. For example, a terminal loss in $[0,1]$ and at most $H$ charges bounded by $c_{\max}$ give $B=1+Hc_{\max}$.

Assume Theorem~\ref{Core:r:history-bound}'s bound $d_H$ holds uniformly over every tree in that class. An early stop is padded by a common absorbing terminal step for the history comparison. State-dependent charges are permitted when included in the same coarse-path loss. A model-dependent posterior score used as an uncharged surrogate loss would be a different problem.

\begin{Coreproposition}[Uniform model error controls a changed policy]\label{Core:r:policy-error}
Let $J(\pi)$ and $\widetilde J(\pi)$ be a tree's expected losses in the true and approximate models. Then
\[
 |J(\pi)-\widetilde J(\pi)|\le Bd_H\quad(\pi\in\mathfrak P).
\]
If $\widetilde\pi$ minimizes the approximate loss and $\pi_*$ minimizes the true loss over this same class, then
\[
 J(\widetilde\pi)-J(\pi_*)\le\min\{B,2Bd_H\}.
\]
These conclusions do not certify a zero-risk constraint or preservation of admissibility outside the common class.
\end{Coreproposition}
\begin{proof}
The expectation of a function in $[0,B]$ changes by at most $B$ times total variation. Apply the common-policy history bound separately to each fixed tree, proving the first assertion uniformly. Now
\[
 J(\widetilde\pi)\le\widetilde J(\widetilde\pi)+Bd_H
 \le\widetilde J(\pi_*)+Bd_H\le J(\pi_*)+2Bd_H.
\]
The loss range gives the cap $B$. Each tree is first compared with itself across models, before applying optimality to different controllers. A zero-probability bad event in one model can have positive probability in the other, so the average bound does not certify sure feasibility.
\end{proof}
\Coresource{Direct combination of the whole-history bound with the selection argument already used in Theorem~\ref{Core:r:risk}. [AVOC], Sections 8--10, provide the source context for guarded planning and model-error value bounds. This finite fixed-policy-class formulation is stated and proved in full here.}

With true calibration accuracy $37/50$ and planning accuracy $3/4$, the one-step defect is $1/100$. Taking $H=3$ and $B=23/20$ gives the uniform regret bound $69/1000$. A computed performance value does not replace that guarantee. Expected-risk and hard-budget claims remain distinct.


\chapter{Improving discovery without changing its meaning}\label{Core:book:chapter-6}
\setcounter{section}{0}\renewcommand{\thesection}{\thechapter.\arabic{section}}\label{Core:ch:discovery}
\Corequestion{Removing redundant work is an improvement only when the removed work could not change the justified result.}
Let $\mathcal G=(p_0,\ldots,p_{N-1})$ be a finite ordered grammar of pure, total candidate programs. Let $D$ be a finite context set and $t:D\to Y$ a fixed labeling. A teacher either accepts a candidate that agrees everywhere or returns a context where it disagrees. The teacher's work is part of the development account.

\section{Finite counterexample-guided acquisition}
The learner retains every returned counterexample. It proposes the earliest candidate consistent with all retained examples, or reports exhaustion of the grammar when none remains.

\begin{Coreproposition}[What finite counterexamples establish]\label{Core:r:cegis}
The learner never submits the same rejected candidate twice. Within at most $N$ teacher calls, it returns the earliest globally correct candidate in the declared grammar or exhausts that grammar. A timeout before this process finishes is neither acceptance nor exhaustion.
\end{Coreproposition}
\begin{proof}
Each counterexample refutes the submitted program and remains a constraint forever. That program cannot be submitted again. A globally correct candidate never violates a truthful counterexample. With only $N$ candidates, either an accepted one is found or all are ruled out. If the accepted candidate were preceded by a globally correct candidate, that earlier candidate would still be consistent and would have been selected first. This proves the ordering claim. Interruption changes none of these logical implications and cannot supply a completed search result.
\end{proof}
\Coresource{[SII], Theorem 15.1. This is finite exact synthesis in a supplied language, not a statistical generalization guarantee.}

Agreement on $D$ is not validity outside $D$. That distinction is especially important when a candidate is used to prune future searches. A wrong proposal can be rejected by a final checker. A wrong pruning rule may remove the only candidate the checker ought to see.

\section{A learned search representation can preserve the exact investigation}
Suppose a deterministic normalizing map $n$ is sound in the following direction:
\[
n(p)=n(q)\ \Longrightarrow\ p(z)=q(z)\quad\text{for every admitted context }z.
\]
The map need not detect every semantic equivalence or produce a mathematically canonical normal form. Each successful contextual rewrite must be justified under the actual pure, total semantics. Keep the earliest original AST for every attained normal form, not a replacement AST with an uncharged new interpreter.

\begin{Coretheorem}[Ordered semantic compression preserves the teacher trace]\label{Core:r:ordered}
Assume all counterexamples are retained, candidate and example orders are fixed, and the teacher is deterministic and extensional, returning the first disagreement in a fixed context order. Searching only the earliest representatives yields exactly the same submitted ASTs, counterexamples, and final acceptance or grammar exhaustion as the original learner. With the same AST interpreter it uses no more retained-example evaluations. Preparation, dispatch, and stored-index costs are separate.
\end{Coretheorem}
\begin{proof}
A removed expression has an earlier representative with the same values on every admitted context. If that representative was accepted, the original search ended before reaching the removed expression. Otherwise an already retained example refuted the representative, either during a consistency test or from the teacher. Persistence of examples means that the same example refutes the later equivalent expression. The original learner therefore never submits that removed expression to the teacher.

Deleting only those unsuccessful consistency tests preserves the next submitted AST. A deterministic teacher then returns the same counterexample or acceptance. Induction over submissions proves equality of the complete trace, including exhaustion. Every deleted evaluation has nonnegative cost in the named evaluation count. The argument makes no claim that constructing or consulting the representative index is free.
\end{proof}
\Coresource{[SII], Theorems 21.1 and 21.2. Sound equational replacement is the premise. Empirical agreement on a few probe points does not supply it.}

\begin{Corecounterexample}[A probe-equivalent equation is not a general identity]\label{Core:ce:probe}
The expressions $\max(x,1)$ and $x$ agree for probes $x\in\{1,2,3\}$. At zero they differ. Grouping them by those probe values alone can remove a needed candidate. The identity $\max(x,0)=x$ is valid over nonnegative integers, but not over all integers. The semantic domain is part of the equation's certificate.
\end{Corecounterexample}

The retained object may therefore improve later search without containing its answers. This is the precise process-level role of a checked equation library. The grammar, primitive semantics, teacher, and proof rules can still be supplied. Their presence must remain visible instead of being redescribed as discoveries.

\section{A reusable rule for stopping an unpromising read}
A diagnostic compiler inspects a binary query on $m$ immutable cases. In the prefix already read, $a$ cases give zero and $b$ give one. Each child is allowed at most $C$ cases because the remaining binary depth has that leaf capacity. Assume integers
\[
C\ge1,\quad 2\le m\le2C,\quad a,b\ge0,\quad a+b\le m.
\]

\begin{Coreproposition}[A sound prefix exclusion]\label{Core:r:prefix}
The unread reports can be completed with both final branch counts at most $C$ if and only if $\max(a,b)\le C$. A scan may therefore stop when one count first exceeds $C$. In a fixed ordered depth-first compiler with identical scope, terminal, and fallback rules, replacing full root scans by this early rejection preserves the first returned tree or the completed depth-bounded failure result.
\end{Coreproposition}
\begin{proof}
Counts only increase, so an overfull child cannot be rescued. Otherwise the two remaining capacities are $C-a$ and $C-b$, both nonnegative. Their sum $2C-a-b$ is at least the number $m-a-b$ of unread rows. Distributing the unread cases between those capacities gives a feasible completion.

Every root rejected early would also be rejected after its full scan. A root not rejected early is read completely and produces the same ordered child supports. Induction on the finite remaining depth therefore yields identical recursive searches, backtracking, and first completion. The same failure condition invokes the same fallback. The skipped rows are only suffixes whose contents could not change rejection. A resource timeout is excluded from the completed-failure claim.
\end{proof}
\Coresource{[SII], Proposition 16.1 and Theorem 16.3. The conclusion concerns capacity feasibility, not the existence of a completing query in the installed language.}

For $C=4$, a fifth zero already closes a root. Reading more rows cannot reduce that branch. By contrast, a four-versus-four split merely survives this necessary test. If no remaining query distinguishes either child, the root still fails.

\begin{Corecounterexample}[A changed repair defeats the old premise]\label{Core:ce:universal}
Four cases previously required four different repairs. Admit one new repair successful in all four. No identifying query is now needed. The old singleton-leaf lower bound and its pruning interpretation are not applicable. The correct response is to recheck the repair scope, not to preserve an obsolete optimization because it once passed.
\end{Corecounterexample}

\section*{What is now available}
The resulting method is semantically justified only under its premises. We must still determine whether it saves total resources, and which proof records remain usable when a later revision changes those premises. Both questions belong to the next chapter rather than being inferred from a successful synthesis.


\chapter{What remains valid after a revision}\label{Core:book:chapter-7}
\setcounter{section}{0}\renewcommand{\thesection}{\thechapter.\arabic{section}}\label{Core:ch:revision}
\Corequestion{A mathematical proof can remain true while no longer applying to the installed system.}
A certificate for ordinary integer addition is not a certificate for an operation with the same printed name and a different implementation. A grammar quotient depends on the grammar order. A comparison depends on the actual controller and the prices used to value its events. These are distinct dependencies.

\section{Reuse follows semantic dependencies}
Let an artifact state be $a=(a_1,\ldots,a_d)$. A finite directed acyclic graph represents claimed guarantees. For vertex $v$, let $D_v$ be its direct artifact dependencies and let $\operatorname{pa}(v)$ be its parent guarantees. The declaration is semantically complete if the meaning of its claim and local verification rule is determined by those components and parent claims. External assumptions are included explicitly as premises.

After changing $a$ to $a'$, mark a vertex stale when a direct input changes and propagate staleness to all descendants. Call the resulting set $\Delta$.

\begin{Coretheorem}[Selective reuse with complete dependencies]\label{Core:r:dependencies}
Assume the graph is finite and acyclic, direct dependencies are semantically complete, local rules are sound, and the applicable external premises hold. Assume also that every old witness established its old claim before the revision. Old witnesses still establish every claim outside $\Delta$. Re-establishing affected claims in topological order, using revised values and established parents, restores the corresponding revised claims. Without an affected parent's re-establishment, this reuse rule supplies no justification for its dependent claim.
\end{Coretheorem}
\begin{proof}
A vertex outside $\Delta$ has unchanged direct values and no stale parent. In topological order, induction shows that every parent retains its meaning and proof. Completeness of the dependency declaration leaves the local argument applicable. For a stale vertex, once its revised parents are established, soundness of the newly checked local rule gives the revised conclusion. A second topological induction proves restoration. An unchecked parent leaves a missing premise, so the rule alone cannot justify a dependent conclusion. This is not a proof that the conclusion is false.
\end{proof}
\Coresource{[SII], Theorem 24.1. It is a theorem about a declared complete dependency graph, not a method that infers complete dependencies from arbitrary programs.}

\begin{Corecounterexample}[A missing dependency defeats a perfect replay]\label{Core:ce:dependency}
A stored claim is $u=v$, initially with both values zero. Its registry records only $u$. Change $v$ to one. A flawless comparison with the recorded $u$ still passes, while the claim is false. A hash or replay cannot recover the premise the registry omitted.
\end{Corecounterexample}

Adding the universal repair from the preceding chapter changes a successful-repair set. It invalidates a cardinality certificate and any dependent performance statement. It need not invalidate an arithmetic equation used elsewhere. A new valid repair also does not inherit the old controller's speed comparison. Rechecking can legitimately conclude that the new route is slower, cheaper, or incomparable under the stated resources.

\begin{figure}[ht]
\centering\includegraphics[width=\linewidth]{sources/research/machine-native-reasoning-wider/figures/revision_scope.pdf}
\caption{A change to the repair language affects the diagnostic justification but need not change arithmetic semantics. The depicted graph is a declared example, not an automatically discovered dependency registry.}
\end{figure}

\section{Semantic improvement and resource improvement are separate}
Fix named additive event types. Let $K\ge0$ be an extra preparation vector, let $s_i$ be the old-minus-new per-use vector for workload type $i$, and let $n_i$ count its uses. The net event difference is
\[
\Delta=\sum_i n_i s_i-K.
\]
A nonnegative price vector $w$ converts these events into one declared additive measure. Storage capacity, peak memory, and elapsed time need their own models unless a justified conversion is specified.

\begin{Coreproposition}[Which gains survive every admitted price?]\label{Core:r:prices}
The saving $w\cdot\Delta$ is nonnegative for every nonnegative price vector if and only if every coordinate of $\Delta$ is nonnegative. If prices lie in the convex hull of a finite nonempty set $W$, their minimum saving equals $\min_{w\in W}w\cdot\Delta$.
\end{Coreproposition}
\begin{proof}
A componentwise nonnegative vector has nonnegative inner product with every nonnegative vector. If coordinate $j$ is negative, assigning price one only to that coordinate gives a negative result. This proves the first assertion. A point in the stated convex hull is a convex combination of its listed prices, so its saving is the same convex combination of their savings. It cannot be smaller than the minimum listed value, and a listed minimizer attains that value.
\end{proof}
\Coresource{[SII], Proposition 25.1. No hardware performance law follows from this additive accounting identity.}

For the source's checked weighted-search comparison, the event difference is $(1552,1552,-1868)$ in reads, rational additions, and comparisons. Equal prices give $1236$ saved units. Prices $(1,0,1)$ give a loss of $316$. Both conclusions concern the same valid controller. The main text uses this as a transparently rechecked source calculation, not as a new measured timing result.

\section{Evidence across an adaptive sequence}
A procedure that searches many candidates can select a favorable error in its evaluator. The right statement must therefore condition on the information that produced the candidate, not merely repeat a marginal error bound for a candidate fixed in advance.

\begin{Coreproposition}[A conditional error budget for a lineage]\label{Core:r:lineage}
Let $\mathcal H_{j-1}$ contain the information available before attempt $j$. Suppose the event $A_j$ of accepting a false claim satisfies
\[
\Pr(A_j\mid\mathcal H_{j-1})\le\delta_j\quad\text{almost surely},
\qquad\sum_{j\ge1}\delta_j\le\delta,
\]
for deterministic nonnegative $\delta_j$. Then the probability of ever accepting any false claim is at most $\delta$, even when candidates and the stopping time depend on earlier history.
\end{Coreproposition}
\begin{proof}
Taking expectations gives $\Pr(A_j)\le\delta_j$. A union bound gives the result for every finite set of attempts. The finite unions increase to the event of any false acceptance, so continuity of probability gives the infinite statement. Stopping removes attempted events and cannot enlarge their union.
\end{proof}
\Coresource{[SII], Proposition 26.1. The mathematical input is conditional validity of each actual attempt. A numerical error schedule alone does not establish that premise.}

This proposition is not an empirical result about a population of agents. It explains what evidence would need to remain valid when an evaluator is used adaptively. Deterministic exact checks in this book are not treated as statistical samples and are not assigned confidence intervals through this formula.

\section*{What is now available}
A revision can preserve its meaning, identify stale obligations, and state its resource advantage under explicit prices. None of this yet determines what several participants can know together or which authority they should exercise. The next chapter keeps information sharing and institutional permissions separate while applying the same preservation discipline.


\chapter{Shared work and maintained obligations}\label{Core:book:chapter-8}
\setcounter{section}{0}\renewcommand{\thesection}{\thechapter.\arabic{section}}\label{Core:ch:shared}
\Corequestion{Joint information can support a task without creating a single decision maker.}
One agent may have a display record and another a calibration record. Neither need know the physical target bit alone. Together their records can determine it. This is an informational relation before it is any claim about collective agency.

\section{What several accessible records can reconstruct}
Let $V$ be a finite-dimensional vector space over a field. Participant $i$ has a linear record $A_i v$. Stack the records into $Av=(A_1v,\ldots,A_mv)$, and let $Qv$ be the protected linear target.

\begin{Coreproposition}[Target-relative joint reconstruction]\label{Core:r:reconstruction}
There is a linear decoder $D$ with $Q=DA$ if and only if $\ker A\subseteq\ker Q$. Full-state reconstruction is possible exactly when $A$ is injective. The intersection identity $\ker A=\bigcap_i\ker A_i$ determines which ambiguities sharing removes.
\end{Coreproposition}
\begin{proof}
If $Q=DA$, a vector annihilated by $A$ is also annihilated by $Q$. Conversely, define $D(Av)=Qv$ on $\Coreim A$. If $Av=Aw$, then $v-w\in\ker A\subseteq\ker Q$, so this definition is independent of the representative. It is linear and extends to the full record space by extending a basis. For $Q=I_V$ the kernel condition is injectivity. A stacked record vanishes exactly when each component record vanishes, giving the intersection identity.
\end{proof}
\Coresource{Target-relative specialization of the reconstruction and factorization machinery in Sufficient Resolution [SR] and Stratified Relational Workspace Theory [SRWT]. The proof here does not import their complete noise, dynamics, or identity theories.}

Over $\mathbb F_2$, Agent A observes $A_A(b,s)=b+s$ and Agent B observes $A_B(b,s)=s$. The joint matrix is invertible. Either agent alone lacks a direction needed to recover $b$. The nuisance bit may remain unknown because the target map ignores it.

The theorem concerns accessible records. It does not send them. A communication protocol, truthful interpretation, permitted disclosure, and the recipient's executable decoder remain separate requirements. Chapter~\ref{Core:ch:decision} states the corresponding fixed-recipient criterion.

\section{Permissions are part of an operational state}
Now let a finite institutional state $x$ have a protected label $\ell(x)$, an enabled-action set $D(x)$ within a fixed finite alphabet, and a deterministic successor $F(x,a)$ for enabled actions. Labels may include resource balances, outstanding obligations, and offices, but need not retain every historical byte.

\begin{Coretheorem}[A quotient that preserves legal operations]\label{Core:r:governance}
For a surjection $\pi:X\to Z$, quotient labels, enabled sets, and successors satisfying
\[
\ell=\bar\ell\pi,\qquad D=\bar D\pi,\qquad
\pi F(x,a)=\bar F(\pi x,a)
\]
exist exactly when equal fibres have equal labels, equal enabled sets, and equal projected successors for every enabled action. The quotient labels, enabled sets, and successors on enabled actions are unique. They preserve every legal finite action and label sequence. Values assigned by a totalized implementation to disabled actions are outside this uniqueness statement.
\end{Coretheorem}
\begin{proof}
Factoring each quantity proves necessity. For sufficiency choose a representative of each fibre to define its label, enabled set, and successor. The three equalities make those definitions independent of the representative. Surjectivity gives uniqueness. Equal enabled sets admit the same first action, and successor equality preserves membership in one fibre. Induction on the sequence length preserves legality and all labels.
\end{proof}
\Coresource{[MS], ``Exact governance quotient.'' This extends the deterministic behavioral construction by explicitly protecting action availability. It does not derive moral or external legal authority from the state description.}

A summary that preserves only the next sensor response can merge states in which a repair is authorized with states in which it is forbidden. It may remain a good predictor while becoming an invalid action interface. The necessary correction is to preserve the enabled-set distinction, not merely to attach a reassuring name to the old code.

\section{Continuing should not erase the ability to finish}
Use a finite \emph{observed} configuration space $Z$. Under partial observation, the configurations would instead have to be justified public possibility states. At $z$ there are finite catalogues of running actions and ending actions, each with a nonempty successor set. A declared safe set is $S$, and an ending target is $G\subseteq S$. A safe hold action is available at $G$.

A stop command is assumed authentic and observable before the next action is selected. Its delivery and authentication are premises, not consequences of the following fixed-point calculation. Define
\begin{align*}
T_0&=G,\\
T_{k+1}&=T_k\cup\{z\in S:\exists a\in A_{\rm end}(z),\ \operatorname{Post}(z,a)\subseteq T_k\}.
\end{align*}
Let $T_*$ be the stable set. Then define
\[
L_0=T_*,\qquad
L_{k+1}=\{z\in L_k:\exists a\in A_{\rm run}(z),\ \operatorname{Post}(z,a)\subseteq L_k\},
\]
and let $L_*$ be its stable value. Both sequences stabilize because $Z$ is finite.

\begin{Coretheorem}[Safe continuation that retains an ending strategy]\label{Core:r:ending}
A starting configuration belongs to $L_*$ exactly when one policy can run safely indefinitely without a stop command and, after a stop command at any decision epoch, reach $G$ on every possible successor path using only the admitted ending actions. Before the command a policy preserving $L_*$ suffices. After the command it can reach $G$ within the first-entry rank of the current state in the $T_k$ sequence.
\end{Coretheorem}
\begin{proof}
Induction shows that membership in $T_k$ is equivalent to an ending-only policy reaching $G$ within $k$ transitions while staying safe. The chosen action moves every successor to a smaller rank. Outside $T_*$, every ending action has a successor outside $T_*$ or outside safety. Choosing such a successor at every step defeats any purported surely safe ending. Thus $T_*$ is exactly the sure-ending region.

At $L_*$ a running action keeps every successor inside $L_*$. Repeatedly using such actions is safe and retains membership in $T_*$. A command can therefore switch to the decreasing-rank ending policy at any epoch.

Conversely, if a policy has the stated guarantee, every pre-command reachable configuration must admit an ending immediately and hence lie in $T_*$. For any finite number of further running steps, it must keep all paths in that set. Induction in the decreasing construction places the initial state in every $L_k$, hence in $L_*$. No action is chosen after secretly observing its successor.
\end{proof}
\Coresource{Fully observed finite specialization of [MS], ``Continued operation with an available authorized ending.'' The more general public-possibility construction is not silently assumed in this version.}

\begin{Corecounterexample}[Safe inactivity without an ending]\label{Core:ce:ending}
A state $z$ has a safe idle self-loop and no ending action reaching $G$. It is safe forever but outside $T_*$. In a second state $u$, an ending is available, while the only running action spends its reserved resource and moves to $z$. Then $u$ belongs to $T_*$ but not $L_*$. Being able to stop now does not establish that ordinary operation preserves that ability.
\end{Corecounterexample}

These results do not establish society formation, strategic compliance, or a just constitution. They establish information, action, and continuation properties of a specified finite institution. Their place in this finite reasoning chain follows from the preceding account of accessible evidence, common artifacts, feasible controllers, and dependency-preserving revision.


\chapter{One finite reasoning chain}\label{Core:book:chapter-9}
\setcounter{section}{0}\renewcommand{\thesection}{\thechapter.\arabic{section}}\label{Core:ch:close}
The sensor story can now be followed without changing the meaning of a record midway through the argument.

Begin with eight candidate worlds $(b,s,n)$. The agent's initial display $y$ determines its two-state predictive class under flips, but not the physical target bit. The first chapter proves both claims. The second gives the operator on the retained class and the coordinate map that becomes usable after calibration. Neither requires the nuisance bit.

When display and calibration observations arrive, candidate sets are filtered. The third chapter proves why the joint association must be retained, and why sixteen occupancy states suffice for the exact partition-evidence contract. Under a separately declared noisy calibration, four probability masses replace the occupancy bits. The same numerical posterior cannot be inferred from separate marginals. The correction argument requires a joint report/next-state law, not only correct separate marginals. Its finite refinement makes the required distinctions explicit, and its history bound states the separate price of conditioning on a rare report.

The fourth chapter places that posterior and any response approximation under one loss function. It also specifies what a recipient can decode from a message and its prior record. The fifth asks whether the sender can acquire the missing observation and transmit a successful artifact within a single online and stored-controller budget. The complete tree is counted even if an execution visits only one path. A separate noisy-calibration calculation then uses the exact same retained belief to select a finite expected-cost inquiry. Uniform whole-history error supports a quantitative policy comparison, but not an automatic hard-budget or zero-risk guarantee.

The sixth chapter then changes the discovery procedure rather than the hidden world. Checked equations can remove equivalent candidates while preserving every teacher proposal and counterexample. A separately justified capacity rule can stop reading an impossible diagnostic root. Neither step invents the primitive grammar, the teacher, or the theorem used to authorize pruning.

The seventh chapter binds those uses to their premises. Adding a universal repair changes the required diagnostic distinctions. The old lower bound becomes inapplicable. Arithmetic identities can remain applicable. A fresh cost comparison may find a saving or a loss. The eighth chapter finally combines different participants' records and asks what actions they are entitled to perform under their declared institutional state. Safe indefinite operation and preservation of an authorized ending are checked separately.

\section{A worked end-to-end record}
Take the actual illustrative world $(b,s,n)=(0,1,1)$. The complete protocol begins with all eight worlds possible:
\begin{center}\begin{tabular}{p{0.25\linewidth}p{0.27\linewidth}p{0.35\linewidth}}\toprule
Event & Remaining worlds & What is justified\\\midrule
Before observation & All eight & No physical target identified.\\
Display $y=1$ & Four worlds & Predict every later display flip.\\
Calibration $s=1$ & Two worlds differing only in $n$ & Physical target $b=0$.\\
Send $y$ to a recipient holding $s$ & Same two worlds & Recipient computes $b=0$ through its fixed XOR decoder.\\
Commit the authorized target action & As declared by its effect & One action justified across the current support.\\
Record the fulfilled obligation & Includes the action's receipt & Continuation or ending can refer to the actual completed act.\\\bottomrule\end{tabular}\end{center}

The witness does not claim to have discovered its sensor semantics. Its exact checking confirms the stated scenario. Another physical process requires its own evidence. A failed receipt, an untrusted calibration, or an altered action effect changes the guarantee rather than being averaged away as a small numerical error.

\section{Why the core is narrower than the archive}
This book does not append the old Neuralese experimental history, proposed PRISM architecture, full SRIC lineage, or every formal guide. Their presence would make the reader repeatedly switch between a proved conditional statement, a proposed mechanism, a historical execution, and an implementation plan. Those are useful source objects, but they are not one continuous proof.

Selected source results are restated and proved where they carry the argument. Selected stochastic-value arguments from AVOC enter through complete proofs and explicitly matched finite models. Several resulting statements have source-facing counterparts in the archived Machine Theory development. The full AVOC implementation and the specialist Certified Bayesian Inference construction remain outside this selection. The correspondence appendix records the difference between those archived counterparts and the present editorial work.

Nor is a source excluded merely because it contains open questions. Its proved results can be considered individually. A claim enters when its exact form, dependencies, and evidence have been reviewed at the level promised for this edition. This is a stricter and more useful rule than treating every page of a final-named document as equally checked.

\section{The result of the compilation}
The core establishes a connected conditional account of accessible state, meaningful operations, corrective memory, usable explanations, bounded investigations, justified search changes, and maintained obligations. The arguments share several elementary constructions, but each use preserves its distinct type and hypotheses.

No theorem here turns a semantic interface into a trained neural model, a finite investigation into general intelligence, a shared record into a collective subject, or an invariant into legitimate external authority. The mathematical contribution is the explicit chain and the failures that break it. Those failures can be investigated without changing the standard of proof or requiring a large implementation campaign.


\part{Composition, control, and revision}
\chapter{A common contract for stochastic and legal operations}\label{Core:book:chapter-10}
\setcounter{section}{0}\renewcommand{\thesection}{\thechapter.\arabic{section}}\label{Core:ch:jointcontract}

A probability model describes what an action may do. An action catalogue describes whether that action is available. A useful abstraction may need to preserve both. Keeping a correct next-state marginal while losing the report correlation can destroy a filter. Keeping the correct filter while losing a permission distinction can admit an action that the underlying model forbids.

The two preservation problems use one state map in this chapter. Their composition is exact for a supplied finite model. The proof also specifies what the policy can observe, because a comparison between two models is meaningless if their controllers secretly receive different information.

\section{The joint data of a decision process}
Let $X$ and $Z$ be finite nonempty state sets and let $\pi:X\to Z$ be surjective. Let $A$ and $O$ be finite nonempty action and report alphabets. At a fine state $x$, the set $D(x)\subseteq A$ contains the admitted actions, and $\ell(x)\in Y$ is a protected label. For $a\in D(x)$, the nonnegative numbers
\[
 K_{a,o}(x'\mid x),\qquad
 \sum_{o\in O}\sum_{x'\in X}K_{a,o}(x'\mid x)=1
\]
specify the joint distribution of the next state and report. No transition is specified for a disabled action. An implementation may fill unused entries, but those values have no role in the contract.

For a nonnegative vector $u$ on $X$, write
\[
 (\pi_*u)(z)=\sum_{\pi x=z}u(x).
\]
This notation applies to probability distributions and to unnormalized history masses. A coarse model has labels $\bar\ell$, enabled sets $\bar D$, and joint instruments $\bar K$ on its enabled actions.

\begin{Coretheorem}[Joint legal and stochastic quotient criterion]\label{Core:r:joint-quotient}
There are coarse data satisfying
\[
 D(x)=\bar D(\pi x),\quad \ell(x)=\bar\ell(\pi x),\quad
 \sum_{\pi x'=z'}K_{a,o}(x'\mid x)=\bar K_{a,o}(z'\mid\pi x)
\]
for every $x$, every $a\in D(x)$, and every $o,z'$, if and only if states in the same fibre have identical labels and enabled sets and identical projected joint successor/report rows. Explicitly, $\pi x=\pi y$ must imply
\begin{align}
 D(x)&=D(y),\qquad \ell(x)=\ell(y),\label{Core:eq:fibre-legal}\\
 \sum_{\pi x'=z'}K_{a,o}(x'\mid x)
 &=\sum_{\pi x'=z'}K_{a,o}(x'\mid y)
 \quad(a\in D(x),\ o\in O,\ z'\in Z).\label{Core:eq:fibre-joint}
\end{align}
When these conditions hold, the coarse labels, enabled sets, and instruments on enabled actions are unique.
\end{Coretheorem}
\begin{proof}
The factorization identities give necessity immediately. For sufficiency, choose a representative $x$ of each fibre. Define its coarse label and enabled set by their values at $x$, and define the coarse instrument by the displayed fibre sum. The hypotheses make all three choices independent of the representative. The new entries are nonnegative. Summing them over $o,z'$ sums every fine successor exactly once and therefore gives one. Surjectivity gives uniqueness at every coarse state. Uniqueness says nothing about padded values for disabled actions.
\end{proof}

\section{Observable histories and common admissibility}
A history available to the policy is
\[
 h=(a_1,o_1,\ldots,a_t,o_t).
\]
The fine hidden state and the coarse state are not additional policy inputs. Protected labels are recorded for comparison, but are not assumed observed merely because the theorem preserves them. An application that exposes a label must include it in the observation interface of both models.

Start from a fine prior $p_0$ and the coarse prior $\bar p_0=\pi_*p_0$. Let $u_h(x)$ be the unnormalized probability of history $h$ and current state $x$. A common policy assigns the same probability $\alpha(a\mid h)$ in both models. At every positive-mass history it may select only actions in
\[
 D_h=\bigcap_{x:u_h(x)>0}D(x).
\]
The intersection is over a nonempty support. If it is empty, no policy can make another admitted action at that history. A stopping convention can be included as a shared terminal action, but is not silently supplied by a stochastic row.

\begin{Corelemma}[Support and admissible actions agree]\label{Core:r:joint-support}
If $\bar u_h=\pi_*u_h$, then
\[
 \operatorname{supp}\bar u_h=\pi(\operatorname{supp}u_h),\qquad
 \bigcap_{x:u_h(x)>0}D(x)=\bigcap_{z:\bar u_h(z)>0}\bar D(z)
\]
whenever the history has positive mass.
\end{Corelemma}
\begin{proof}
A finite sum of nonnegative terms is positive exactly when at least one summand is positive. This gives the support identity. The enabled sets are constant on each fibre, so repeating a fibre's action set in an intersection changes nothing. The two intersections are therefore equal.
\end{proof}

\begin{Coretheorem}[Preservation of complete controlled histories]\label{Core:r:joint-history}
Under Theorem~\ref{Core:r:joint-quotient}, a common admissible history policy preserves the law of every finite coarse-state/action/report path. Consequently it preserves every action/report/label history. At every positive-mass history it also preserves the aggregated posterior and the set of actions admitted throughout its support. These claims hold at every horizon for which the policy is defined on all positive-mass prefixes.
\end{Coretheorem}
\begin{proof}
First prove $\bar u_h=\pi_*u_h$ by induction. It holds initially by the prior assumption. At a positive history, Lemma~\ref{Core:r:joint-support} gives the same available actions on both sides. For any action with positive policy mass, every current fine state of positive mass enables that action. Thus
\[
 \pi_*u_{hao}
 =\pi_*\bigl(\alpha(a\mid h)K_{a,o}u_h\bigr)
 =\alpha(a\mid h)\bar K_{a,o}\pi_*u_h
 =\bar u_{hao}.
\]
Here multiplication uses only enabled columns. Values in all other columns are irrelevant because their input mass is zero. Zero-mass prefixes stay zero. Summation proves equality of history probabilities, and division by the common positive mass proves the posterior assertion.

This induction alone records only the current state and the observed history. To retain the entire state path, fix a coarse prefix
\[
 (z_0,a_1,o_1,z_1,\ldots,a_t,o_t,z_t).
\]
Sum fine path weights over all fine prefixes mapping to this prefix. For one more step, each fine current state in the last fibre contributes the same projected successor/report row by Equation~\eqref{Core:eq:fibre-joint}. The policy factor depends only on the already fixed action/report prefix. It can therefore be taken outside the fine-state sum. Induction gives the exact path mass
\[
 \bar p_0(z_0)\prod_{i=1}^{t}
 \alpha(a_i\mid a_1,o_1,\ldots,a_{i-1},o_{i-1})
 \bar K_{a_i,o_i}(z_i\mid z_{i-1}).
\]
An inadmissible or zero-mass prefix contributes zero. This is precisely the coarse path law. Applying $\bar\ell$ to each $z_i$ gives the full label history, including its correlations across time. Labels need not be disclosed to the policy for this pushforward comparison to be valid.
\end{proof}

\Coresource{This is a joint finite composition of Theorems~\ref{Core:r:instrument} and~\ref{Core:r:governance}. The path argument and the common-support lemma are explicit parts of the present statement. The result is a written proof, not a newly compiled Lean theorem.}

\section{Why preserving two marginals is insufficient}
Consider two fine states in one proposed coarse input fibre. From the first, the pairs $(z',o)$ are $(0,0)$ and $(1,1)$ with probability $1/2$ each. From the second, they are $(0,1)$ and $(1,0)$ with probability $1/2$ each. The successor marginal and the report marginal are both uniform in either case. Their association is different, however, so the same report leads to opposite posterior states. Equation~\eqref{Core:eq:fibre-joint} fails.
\begin{figure}[htbp]\centering\includegraphics[width=.87\linewidth]{sources/research/machine-native-reasoning-wider/figures/joint_instrument_obstruction.pdf}\caption{Equal marginals do not determine the joint instrument. Each table is an exact probability law. The report needed for correction is correlated with opposite successor states.}\label{Core:fig:joint-obstruction}\end{figure}

Final-state checks have a different limitation. Two deterministic paths can visit different intermediate labels and return to the same final state. A test that stores only the terminal state and the report word cannot distinguish them when the reports omit those labels. Complete path preservation requires the prefix argument above, not just a terminal distribution comparison.

\section{An exact sensor realization}
Return to $x=(b,s,n)\in\{0,1\}^3$ and retain $z=(b,s)$. The actions \emph{hold} and \emph{flip} are available everywhere. A \emph{repair} action is available only where $b=0$ and leaves that target value fixed. The nuisance bit is retained by every fine transition but ignored by the coarse state.

After an action, report the new display $b\oplus s$ correctly with probability $3/4$ and incorrectly with probability $1/4$. Both the label and the joint instrument ignore $n$. The ten coarse state/action rows satisfy the quotient conditions exactly.

From a uniform prior, a report leaves both target values possible, so repair is not uniformly admitted. A policy that reacts to those reports must continue with hold or flip. From a prior supported on $b=0$, repair is available until an action changes that support. The distinction exercises the support condition rather than assuming all actions are always available.

The reference checker compares complete projected state paths, action words, report words, and label words under deterministic and randomized report policies. It also compares every positive-history posterior and available-action intersection. The exact finite examples supplement the general proof. Their probabilities and action semantics are supplied model data.

\section{Three different failures}
A quotient may lose a protected label or an enabled-action distinction. It may preserve separate probability marginals and lose their joint law. Or it may be mathematically exact but unavailable to a learner whose evidence cannot initialize it. The first two are failures of preservation. The third is a failure of acquisition. A correction must address the failed condition rather than merely increase the size of the retained vector.


\chapter{Preserving an ending strategy under safe continuation}\label{Core:book:chapter-11}
\setcounter{section}{0}\renewcommand{\thesection}{\thechapter.\arabic{section}}\label{Core:ch:endinganalysis}

A system may remain within a declared safe region indefinitely and still lose every way to end its operation. Conversely, an authorized ending procedure may exist at the present state but become unavailable after the next ordinary action. These are distinct requirements. The relevant mathematical object is therefore not a safety invariant alone, but a safe-running region entirely contained in a sure-ending region.

The following development expands the finite result stated earlier in Chapter~\ref{Core:ch:shared}. It supplies an explicit rank proof, treats strategies that depend on observed history, and separates deadlock from an infinite unsuccessful path. Formal correspondence is recorded in the appendix.

\section{Finite game data and the command boundary}
Fix a finite, fully observed configuration set $Z$. Each $z\in Z$ has finite collections $A_{\rm run}(z)$ of running actions and $A_{\rm end}(z)$ of ending actions. For each admitted action $a$ there is a nonempty set $\operatorname{Post}(z,a)\subseteq Z$ of possible successors. A controller chooses an action \emph{before} the successor is determined. The quantifier over all successors is adversarial and is stronger than an expected-value objective.

Let $S\subseteq Z$ be a safe set and $G\subseteq S$ an ending target. At every state of $G$, a hold operation with successors contained in $G$ is available after ending. Running forever means that the controller can choose an admitted action at every pre-command epoch. Deadlock is not counted as successful indefinite operation. A stop command is authenticated and observed before the next action is selected. Whether such an authentication protocol exists in an external institution is not proved by the game calculation. It is part of the model supplied to this theorem.

A strategy may inspect the finite history of observed states and chosen actions. Before the stop command it selects running actions. After the command it selects ending actions. The desired guarantee consists of two clauses: every running path without a command remains safe forever, and every path following a command at any decision epoch reaches $G$ in finitely many ending steps while remaining safe. The strategy must satisfy both clauses under the same command-conditioned policy.

\section{The sure-ending attractor}
Define
\begin{align*}
T_0&=G,\\
T_{k+1}&=T_k\cup\left\{z\in S:\text{some }a\in A_{\rm end}(z)
\text{ satisfies }\operatorname{Post}(z,a)\subseteq T_k\right\}.
\end{align*}
The sequence is increasing and every term is contained in $S$. Because the state set is finite, it becomes constant after finitely many strict additions. Let $T_*$ denote its stable value. For $z\in T_*$, define the first-entry rank $\rho(z)$ as the least $k$ with $z\in T_k$.

\begin{Corelemma}[Finite-horizon ending criterion]\label{Core:r:end-horizon}
A configuration $z$ belongs to $T_k$ exactly when there exists an ending-only strategy that, whatever the successors, reaches $G$ within at most $k$ transitions while remaining in $S$. In particular, if $\rho(z)>0$, there is an ending action whose every successor has strictly smaller first-entry rank.
\end{Corelemma}
\begin{proof}
The assertion at $k=0$ says precisely that the initial state is already in $G$. Assume it holds for $k$. If $z\in T_{k+1}\setminus T_k$, the construction supplies an ending action with all successors in $T_k$. After the action the controller uses the strategy given by induction for the observed successor. Every path reaches $G$ within another $k$ transitions. States already in $T_k$ retain their strategy. Conversely, a policy that reaches $G$ in at most $k+1$ steps must either start in $G$, or select some first ending action such that the continuation from every possible successor reaches $G$ within $k$ steps. By induction every such successor belongs to $T_k$, hence the initial state lies in $T_{k+1}$. At the first-entry stage $\rho(z)$, the selected action sends every successor into $T_{\rho(z)-1}$, establishing strict rank decrease.
\end{proof}

\begin{Corelemma}[Adversarial completeness of the ending region]\label{Core:r:ending-complete}
From a configuration outside $T_*$, no strategy, including one depending on arbitrary finite observed history, can guarantee safe eventual arrival at $G$ along every successor path.
\end{Corelemma}
\begin{proof}
If the initial configuration is outside $S$, safety has already failed. Otherwise, because $T_*$ is a fixed point, no ending action at a state $z\in S\setminus T_*$ can have all its successors in $T_*$. An adversary can therefore select a successor outside $T_*$ after every proposed ending action. If the selected successor is outside $S$, safety fails. If it remains within $S\setminus T_*$, the same argument can be repeated. The adversary does not need to know the controller's future choices, only its chosen action at each current state. If the controller reaches a non-goal state with no ending action, it deadlocks before reaching the target and loses. Otherwise, either safety fails at a finite step or an infinite path never reaches $G$, which is contained in $T_*$. This defeats any sure-ending guarantee. The argument permits history-dependent policies and requires nonempty successor sets so that the proposed actions represent actual transitions.
\end{proof}

These lemmas show that $T_*$ is not merely the region reached by a convenient algorithm. It is \emph{exactly} the set from which an ending strategy can guarantee success under all modeled successor choices. The rank bounds the required ending transitions by a finite number no greater than the number of strict expansion rounds.

\section{A safety invariant that retains the ending region}
The ending region is not necessarily closed under ordinary running. Define the decreasing sequence
\[
L_0=T_*,\qquad
L_{k+1}=\left\{z\in L_k:\exists a\in A_{\rm run}(z),\ \operatorname{Post}(z,a)\subseteq L_k\right\}.
\]
Let $L_*$ be its stable value. Each stage removes states for which ordinary operation cannot be guaranteed to remain in the preceding stage.

\begin{Corelemma}[Greatest controlled invariant inside the ending region]\label{Core:r:running-invariant}
The set $L_*$ is the greatest subset of $T_*$ admitting a running action at every member with all successors again in that subset. A state belongs to $L_k$ exactly when a strategy can safely take $k$ running actions while leaving a sure-ending strategy available at every reached decision epoch, including the initial one.
\end{Corelemma}
\begin{proof}
The bounded-horizon claim follows by induction. At horizon zero only membership in $T_*$ is needed. For the next step, a first running action must send every successor to a state that permits the remaining $k$ running steps and an ending option at every intermediate epoch. Truncating a successful $(k+1)$-step strategy after $k$ steps also shows that its initial state lies in $L_k$. Thus the initial state and all successors of its first action lie in $L_k$, exactly as required by $L_{k+1}$. Conversely, membership supplies the first action and the inductive strategies from its successors.

At stabilization, the defining equation gives one running action for each state in $L_*$ whose successors stay in $L_*$. Selecting these actions proves controlled invariance. If $B\subseteq T_*$ is any other subset with the same property, induction shows $B\subseteq L_k$ for every $k$: inclusion at zero is given. The witnessing running action for any $z\in B$ has all successors in $B\subseteq L_k$, so $z\in L_{k+1}$. Thus $B\subseteq L_*$. No larger controlled invariant within $T_*$ exists.
\end{proof}

\begin{Coretheorem}[Safe running with an available ending at every epoch]\label{Core:r:joint-ending}
A state $z$ belongs to $L_*$ if and only if a single command-conditioned policy can run safely forever in the absence of a command and, upon receiving an authentic stop command at any decision epoch, can reach $G$ on every possible successor path using only ending actions. The stopping phase from a state $w$ reaches $G$ within $\rho(w)$ transitions. The combined policy can be chosen positional in the pair consisting of the current state and the command mode.
\end{Coretheorem}
\begin{proof}
For sufficiency, at each running state in $L_*$ choose an action with all successors in $L_*$. Such actions exist by the preceding lemma and keep every running path within $L_*\subseteq T_*\subseteq S$ for any duration. If a stop command arrives at a decision epoch, the current state still lies in $T_*$. Switch modes and choose a first-entry-rank-decreasing ending action. The rank is a nonnegative integer, decreases strictly after every ending transition outside $G$, and must reach zero in no more than $\rho(w)$ steps. At rank zero the system is in $G$ and can hold safely. This is one policy with the command as its observable mode switch, not two incompatible policies chosen after seeing a future path.

For necessity, consider any policy with the stated guarantee, including one that depends on the entire available history. The stop command may arrive at the present epoch, so by the adversarial completeness lemma the starting state must belong to $T_*$. It may instead arrive after any number $k$ of running steps. Every state reachable immediately before such a command must therefore remain in $T_*$. The running prefix of length $k$ consequently satisfies the bounded-horizon condition of the preceding lemma, so the initial state is in $L_k$. This holds for every $k$. The decreasing sequence stabilizes on a finite set, hence the state belongs to its eventual value $L_*$. No operation is selected with advance knowledge of which successor an adversary will return.
\end{proof}

\section{A three-state obstruction to naive safety}
Consider the safe configurations $S=\{u,z,g\}$ with ending target $G=\{g\}$. The state $g$ has a hold action in both modes, $z$ has a safe running self-loop but cannot end, and $u$ can end immediately at $g$. Yet the only running action at $u$ moves to $z$. This gives
\[
T_*=\{u,g\},\qquad L_*=\{g\}.
\]
It is possible to end from $u$ now, but not to run from $u$ while retaining the option to end later. The obstruction is about \emph{future control access}, not immediate safety.
\begin{figure}[htbp]\centering\includegraphics[width=.95\linewidth]{sources/research/machine-native-reasoning-wider/figures/safe_ending_kernel.pdf}\caption{Exact finite safe-running and sure-ending example. The initial running action from $u$ loses the ending option. The sets were checked against explicit strategy enumeration.}\label{Core:fig:ending-kernel}\end{figure}

The companion verifier compares the fixed-point calculation with explicit positional-policy enumeration and checks the rank-decreasing witnesses. The exact input domains and results appear in the verification guide. The calculation treats arrival at $G$ as termination. The stipulated hold concerns subsequent behavior and does not change the reachability calculation.

\section{What changes with partial observation}
The proof assumes the actual configuration is observed when an action is chosen. If only a public report is available, strategies generally act on a belief or possibility state whose update must be justified by the observation protocol. Replacing $Z$ with arbitrary hidden-state names does not produce a valid partially observed ending theorem. One would need a common public-state construction, exact reporting and successor semantics, and a proof that the stop command is observed before action selection. The finite theorem is a solid component of that larger problem without silently solving it.


\chapter{Transporting safe continuation through a quotient}\label{Core:book:chapter-12}
\setcounter{section}{0}\renewcommand{\thesection}{\thechapter.\arabic{section}}\label{Core:ch:controltransport}

The previous chapters establish a stochastic quotient and a safe-ending game separately. To use both in one system, the retained state must also preserve the successor choices relevant to a sure guarantee. Probabilities and supports have different roles here. A small positive transition probability still creates a possible path. It cannot be discarded in a claim that quantifies over every path.

This chapter gives an exact interface between representation change and the two fixed-point constructions. It then weakens equality to a one-sided enclosure that can certify an entire family of models.

\section{The controlled predecessor}
Let $X,Z$ be finite nonempty sets and $\pi:X\to Z$ a surjection. Use a common finite action alphabet with separate running and ending catalogues $A^X_m(x)$ and $A^Z_m(z)$, where $m\in\{\mathrm{run},\mathrm{end}\}$. Every admitted action has a nonempty successor set. The mode can be included in the action name, so the same printed action in two modes need not denote the same transition.

Assume the fine and coarse safe and goal sets agree by inverse image:
\[
 S_X=\pi^{-1}(S_Z),\qquad G_X=\pi^{-1}(G_Z),\qquad G_Z\subseteq S_Z.
\]
This is a substantive restriction. A label that distinguishes a safe state from an unsafe state cannot be omitted from a safety-preserving quotient.

For $B\subseteq X$, define the mode-specific predecessor
\[
 \operatorname{CPre}^X_m(B)
 =\{x:\exists a\in A^X_m(x),\ \operatorname{Post}_X(x,a)\subseteq B\}.
\]
Define $\operatorname{CPre}^Z_m$ in the same way. These operators do not themselves include a safety test. Safety enters the ending recurrence through its intersection with $S$, and the running recurrence through its current set.

\begin{Corelemma}[Predecessors commute with an exact controlled quotient]\label{Core:r:control-predecessor}
Suppose, for each mode and fine state,
\begin{align*}
 A^X_m(x)&=A^Z_m(\pi x),\\
 \pi(\operatorname{Post}_X(x,a))&=\operatorname{Post}_Z(\pi x,a)
 \qquad(a\in A^X_m(x)).
\end{align*}
Then for every $B\subseteq Z$,
\[
 \operatorname{CPre}^X_m(\pi^{-1}B)
 =\pi^{-1}(\operatorname{CPre}^Z_m(B)).
\]
\end{Corelemma}
\begin{proof}
For an admitted action,
\[
 \operatorname{Post}_X(x,a)\subseteq\pi^{-1}B
 \quad\Longleftrightarrow\quad
 \pi(\operatorname{Post}_X(x,a))\subseteq B.
\]
Replace the projected successor set by its assumed coarse value. The action ranges in the two existential quantifiers are identical, so existence of a witnessing action is equivalent on both sides.
\end{proof}

\section{Transport of both fixed-point sequences}
Construct $T^X_k,L^X_k$ on the fine game and $T^Z_k,L^Z_k$ on the coarse game as in Chapter~\ref{Core:ch:endinganalysis}. In particular, the running sequence starts from the stable ending region in its own model.

\begin{Coretheorem}[Ending regions, running kernels, and ranks are preserved]\label{Core:r:control-transport}
Under the exact controlled-quotient assumptions,
\[
 T^X_k=\pi^{-1}(T^Z_k),\qquad
 T^X_*=\pi^{-1}(T^Z_*),\qquad
 L^X_k=\pi^{-1}(L^Z_k),\qquad
 L^X_*=\pi^{-1}(L^Z_*)
\]
for every finite $k$. For every $x\in T^X_*$, the first-entry ranks agree:
\[
 \rho_X(x)=\rho_Z(\pi x).
\]
A positional coarse policy for safe continuation and ending lifts to a fine policy with the same rank bound. To execute that lift, the controller needs the current coarse state and command mode, but not the remaining fine-state coordinates.
\end{Coretheorem}
\begin{proof}
The goal-set assumption gives equality at $T_0$. Inverse images commute with union and intersection. Applying Lemma~\ref{Core:r:control-predecessor} to the ending predecessor gives the induction step
\begin{align*}
 T^X_{k+1}
 &=T^X_k\cup\bigl(S_X\cap\operatorname{CPre}^X_{\rm end}(T^X_k)\bigr)\\
 &=\pi^{-1}\!\left(T^Z_k\cup
       (S_Z\cap\operatorname{CPre}^Z_{\rm end}(T^Z_k))\right).
\end{align*}
The finite increasing sequences stabilize, so their stable sets also agree by inverse image. The first stage at which $x$ appears is exactly the first stage at which $\pi x$ appears, proving rank equality.

The equality of stable ending sets gives the base of the running induction. The recurrence is intersection with the running predecessor. The same inverse-image calculation proves equality at every running stage and hence at stabilization.

At a coarse state in $L^Z_*$, choose a running action whose successors remain in that set. It is admitted at every fine representative, and all its fine successors lie in $\pi^{-1}(L^Z_*)=L^X_*$. After the command, choose a coarse ending action that decreases rank. Its projected successors have smaller coarse rank, so every fine successor has smaller fine rank. These choices define the lifted policy.
\end{proof}

The equality is stronger than soundness of one selected policy. A fine state is winning exactly when its coarse image is winning, even if a proposed fine policy uses additional hidden coordinates. Chapter~\ref{Core:ch:endinganalysis} characterizes all winning states through the fixed points, including those achievable by history-dependent policies. An exact quotient therefore does not conceal a better sure-ending strategy that depends on discarded state.

This does not give the controller an observation it did not have. A report-only policy from Chapter~\ref{Core:ch:jointcontract} need not know its current coarse state. Applying the policy-lifting conclusion requires that $\pi x$ be observed or reconstructed from an available record. Equality of the abstract winning regions alone does not solve a partially observed control problem.

\section{From stochastic rows to successor supports}
Suppose a mode's game successors are exactly the support of the stochastic instrument:
\[
 \operatorname{Post}_X(x,a)
 =\{x':\exists o,\ K_{a,o}(x'\mid x)>0\}.
\]
Define the coarse successor set from $\bar K$ in the same way.

\begin{Corecorollary}[An exact stochastic quotient preserves the support game]\label{Core:r:stochastic-support-game}
The joint quotient of Theorem~\ref{Core:r:joint-quotient} gives
\[
 \pi(\operatorname{Post}_X(x,a))=\operatorname{Post}_Z(\pi x,a)
\]
for every enabled action. If the mode catalogues and the safe and goal sets satisfy the preceding compatibility conditions, Theorem~\ref{Core:r:control-transport} applies to the resulting support games.
\end{Corecorollary}
\begin{proof}
For each $o,z'$, the coarse entry is a finite sum of nonnegative fine entries. It is positive exactly when at least one fine successor in that fibre has a positive entry for that report. Taking the union over reports gives the support identity. Normalization makes the support of every enabled action nonempty. The remaining hypotheses are precisely those of Theorem~\ref{Core:r:control-transport}.
\end{proof}

A sure guarantee in this support game covers all supported successor paths. It is stronger than almost-sure success under a probability law. An infinite path can have probability zero while every one of its finite prefixes has positive probability. The support-game interpretation deliberately retains that path.

\section{Conservative enclosures give one-sided certificates}
Exact equality is not needed to certify a policy conservatively. Suppose instead that a coarse envelope admits only actions available at every fine representative and includes all projected actual successors:
\[
 A^Z_m(\pi x)\subseteq A^X_m(x),\qquad
 \pi(\operatorname{Post}_X(x,a))
 \subseteq\operatorname{Post}_Z(\pi x,a)
 \quad(a\in A^Z_m(\pi x)).
\]
The envelope may contain additional adverse successors and may omit some fine actions. Retain the exact inverse-image assumptions for safety and goals.

\begin{Coreproposition}[A winning envelope policy remains winning]\label{Core:r:control-envelope}
Under these one-sided assumptions,
\[
 \pi^{-1}(T^Z_*)\subseteq T^X_*,\qquad
 \pi^{-1}(L^Z_*)\subseteq L^X_*.
\]
A coarse ending rank bounds the ending time of the lifted policy. The same coarse policy works for every fine model satisfying the enclosure assumptions. The reverse inclusions need not hold.
\end{Coreproposition}
\begin{proof}
Induction on the coarse ending stages proves
$\pi^{-1}(T^Z_k)\subseteq T^X_k$. A witnessing coarse action is available at each fine representative. Its actual successors project into the envelope's successors and hence into the previous coarse stage. The fine action therefore witnesses the next fine stage. The same argument proves rank decrease along every lifted ending path.

The inverse image of $L^Z_*$ lies inside $T^X_*$ by the first inclusion. Choose the coarse invariant running action at each coarse state. Every fine successor lies in the inverse image of $L^Z_*$. Thus that inverse image is a controlled invariant subset of the fine ending region. The greatest-invariant characterization places it inside $L^X_*$. The argument uses only the enclosure relations and therefore applies to every model satisfying them.
\end{proof}

A conservative rejection is not a proof that the fine problem is impossible. For example, a coarse envelope may include an unsafe successor that the actual fine model never reaches. The envelope can then reject a state from which the actual model has an immediate safe ending. The conclusion is intentionally one-sided.

\section{Why a small probability error is not a sure guarantee}
Take three states $s,g,b$, with safe set $\{s,g\}$ and goal $\{g\}$. Running holds $s$ or $g$ in place. The nominal ending action at $s$ moves to $g$ with probability one. The perturbed action instead moves to $g$ with probability $1-\epsilon$ and to unsafe state $b$ with probability $\epsilon$, where $0<\epsilon<1$. A hold action is available at $g$.

The row total variation distance is exactly $\epsilon$. In the nominal game, $s$ belongs to both the ending region and the running kernel. In the perturbed support game, the only ending action from $s$ has an unsafe successor, so $s$ belongs to neither. This happens for every positive $\epsilon$, however small.

\begin{figure}[htbp]\centering
\includegraphics[width=.91\linewidth]{sources/research/machine-native-reasoning-wider/figures/support_margin.pdf}
\caption{An arbitrarily small unsafe branch changes a sure-ending certificate. The nominal and perturbed rows are exact declared distributions. Total variation controls probability discrepancy, not the disappearance of a supported adverse path.}\label{Core:fig:support-margin}
\end{figure}

The appropriate robust certificate is an enclosure containing every possible successor, together with an action that succeeds throughout that enclosure. A small numerical residual is useful evidence for other purposes, but it cannot replace that quantified condition.

\section{The resulting integration}
We can now state the connection without changing the meaning of any interface. The stochastic quotient preserves controlled history laws. Its nonnegative rows induce an exact quotient of support games. Compatible safety and target labels then preserve the ending attractor, running kernel, and first-entry rank. Implementing the coarse policy still requires access to its current coarse state. Keeping these conditions separate is what makes the combined statement usable.


\chapter{Changing a representation without losing correctability}\label{Core:book:chapter-13}
\setcounter{section}{0}\renewcommand{\thesection}{\thechapter.\arabic{section}}\label{Core:ch:migration}

A representation can be sufficient for yesterday's questions and insufficient for today's. The old theorem need not be false. Its target may have changed. A revision procedure must therefore distinguish two tasks: translating information already retained and acquiring information that the old record discarded.

The simplest exact model makes this distinction precise. It also gives a constructive lower bound on the extra information required when a direct migration is impossible.

\section{A migration map exists exactly on constant fibres}
Let $H$ be a finite nonempty set of possible source histories. A history may include a model tag or a past observation whose value must remain available for later correction. Let
\[
 M:H\to U,\qquad N:H\to V
\]
be the old record and the required new record. Work on their attained images $U_0=M(H)$ and $V_0=N(H)$. The new record must be specified before sufficiency is claimed. It can be an answer, a correction signature, or the complete profile of future outputs in a declared finite contract.

\begin{Coretheorem}[Exact migration criterion]\label{Core:r:migration}
There is a map $J:U_0\to V_0$ satisfying $JM=N$ if and only if
\[
 M(h)=M(h')\quad\Longrightarrow\quad N(h)=N(h')
 \qquad(h,h'\in H).
\]
When it exists, $J$ is unique on attained old records.
\end{Coretheorem}
\begin{proof}
Necessity follows by applying $J$ to the equal old records. Conversely, define $J(M(h))=N(h)$. The fibre condition makes the definition independent of the representative. Every $u\in U_0$ has such a representative, so the same equation determines $J$ uniquely.
\end{proof}

\Coresource{This is the attained-record factorization of Proposition~\ref{Core:r:factor}, applied to a representation revision. Grounded Abstraction [GA] supplies the correction-memory and evidence-migration setting. The argument does not assert a separate novelty claim for the fibre criterion.}

A collision with unequal new records is a complete obstruction to an exact deterministic migration from the old record alone. Adding more computation cannot distinguish the two histories. A randomized translator cannot be certainly correct on both either: the same old input gives the same output law, which cannot be concentrated at two different required values.

\section{What additional evidence must provide}
Suppose an allowed observation produces $E:H\to W$. A translator can use the pair $(M(h),E(h))$ rather than $M(h)$ alone.

\begin{Corecorollary}[Evidence-assisted migration]\label{Core:r:migration-evidence}
An exact translator from the available pair $(M,E)$ to $N$ exists if and only if
\[
 M(h)=M(h'),\ E(h)=E(h')\quad\Longrightarrow\quad N(h)=N(h').
\]
If a collision remains after the observation, that particular evidence interface cannot support the requested migration.
\end{Corecorollary}
\begin{proof}
Apply Theorem~\ref{Core:r:migration} to the record $h\mapsto(M(h),E(h))$.
\end{proof}

This result does not choose a physical observation or show that the observation is affordable. It tests the information supplied by a specified observation. A protocol with several actions can be treated similarly once its complete attained transcript is defined. Its admissibility, state changes, and costs then require the controller analysis of Chapter~\ref{Core:ch:inquiry}.

\section{The least joint refinement}
When both the old and new uses must remain available, retain
\[
 P(h)=(M(h),N(h))\in P(H).
\]
The projections recover each record. No ordering between the old and new information is required.

\begin{Coreproposition}[Universal property of the joint record]\label{Core:r:joint-refinement-record}
If a record $R:H\to W$ determines both $M$ and $N$, then it determines $P$. In particular, $|R(H)|\ge |P(H)|$. The pair record attains this lower bound among records required to preserve both uses.
\end{Coreproposition}
\begin{proof}
Let $fR=M$ and $gR=N$ on $H$. Then $w\mapsto(f(w),g(w))$ maps every attained value of $R$ to the corresponding attained pair. This map is onto $P(H)$, since every pair comes from a source history. Finiteness gives the cardinality inequality. The record $P$ itself has the required two projection maps.
\end{proof}

This is a statement about distinct attained values, not the byte size of a particular data structure. It allows the new state to discard duplicated encodings while preserving both information contracts.

\begin{Coretheorem}[Minimum fixed-length supplementary information]\label{Core:r:migration-bits}
For $u\in U_0$, put
\[
 k(u)=|N(M^{-1}\{u\})|,\qquad k_* =\max_{u\in U_0}k(u).
\]
Suppose an auxiliary encoder may inspect the complete source history and sends one fixed-length binary word in addition to the old record. Exact recovery of $N$ is possible with $b$ supplementary bits if and only if $2^b\ge k_*$. Thus the minimum is $\lceil\log_2 k_*\rceil$.
\end{Coretheorem}
\begin{proof}
For a fixed old record $u$, distinct required new values must receive distinct supplementary words. Otherwise the decoder would receive identical inputs with different required outputs. Hence $2^b\ge k(u)$ for every $u$.

Conversely, within each old fibre assign distinct binary words to the distinct attained new values, using at most $k_*\le 2^b$ words. The same word can be reused in different old fibres because the decoder receives $u$ as well. Define the encoder by this fibre-specific table and decode by the inverse table. Every source history is recovered correctly.
\end{proof}

The encoder's access assumption is essential. The bound does not allow a controller holding only $M(h)$ to manufacture the missing bits. It describes the amount of supplementary information when an encoder actually has the source information. With restricted observations, Corollary~\ref{Core:r:migration-evidence} must be checked first. Code tables, acquisition, transmission, and validation also carry costs not counted by the fixed binary payload.

\section{Migration must commute with subsequent operations}
Recovering one new record is not yet preserving a method. Let a common finite action alphabet act on $H$ through total maps $F_a:H\to H$. Suppose both representations close under those same source operations:
\[
 M F_a=P_a M,\qquad N F_a=Q_a N.
\]
The maps $P_a$ and $Q_a$ preserve the attained record sets. These equations refer to the same source action, not merely to operations sharing a name.

\begin{Coretheorem}[A valid migration intertwines the common operations]\label{Core:r:migration-operations}
If $JM=N$, then on $U_0$,
\[
 JP_a=Q_aJ.
\]
Consequently $JP_w=Q_wJ$ for every finite action word $w$, with the execution-order convention of Chapter~\ref{Core:ch:evidence}. Any protected readout that factors through $N$ agrees along the two migrated executions.
\end{Coretheorem}
\begin{proof}
Take $u=M(h)$. Then
\[
 JP_a(u)=JM(F_a h)=N(F_a h)=Q_aN(h)=Q_aJ(u).
\]
Every attained $u$ has a representative, so the equality holds throughout $U_0$. Closure keeps later values within that domain. Induction on the word length gives the continuation identity. Composing with the protected readout gives equality of outputs.
\end{proof}

If a revision changes the underlying operation, the second closure equation must be proved for the changed semantics. A successful initial decoding does not imply it. For example, the identity translator on a bit preserves its current value, but cannot intertwine an old hold operation with a new flip operation. The same stored bit and the same printed operation name do not supply the missing commuting equation.

\section{A worked change of task}
Use the eight histories $(b,s,n)\in\{0,1\}^3$ and the old display record $M=b\oplus s$. For display prediction under physical-bit flips, this record is sufficient. Suppose the new task requires both the display and the physical target, so $N=(b\oplus s,b)$.

There are two possible new records inside each old display fibre. Direct migration fails, but $k_*=2$, so one supplementary bit is information-theoretically sufficient. Calibration supplies $E=s$, and
\[
 b=M\oplus E
\]
gives the actual decoder. The joint old/new record has four attained states. The nuisance bit remains discarded because the new task still does not use it.

Now require the nuisance bit as well. The new record $N'=(b\oplus s,b,n)$ has four possible values inside each old fibre, giving a two-bit supplementary lower bound. Calibration alone no longer suffices. Two histories with equal $b,s$ and different $n$ collide under $(M,E)$ but have different $N'$. The lower bound identifies a missing distinction rather than a computational bottleneck.

\begin{figure}[htbp]\centering
\includegraphics[width=.93\linewidth]{sources/research/machine-native-reasoning-wider/figures/migration_fibres.pdf}
\caption{The required information depends on the revised task. Display alone has two attained values. Adding the target bit requires four joint values. Adding the nuisance bit requires eight. Calibration supplies the first missing distinction, not the second.}\label{Core:fig:migration-fibres}
\end{figure}

\section{Certificates after a change of contract}
A migration map transports data. A certificate authorizes a conclusion only relative to its premises. To reuse a certificate across a revision, specify an old predicate $\varphi$ on $U_0$, a new predicate $\psi$ on $V_0$, and a proved implication
\[
 \varphi(u)\ \Longrightarrow\ \psi(Ju).
\]
An old witness for $\varphi(u)$ and this implication then yield the new conclusion. The implication may express semantic preservation, a stronger guard, or a deliberately weaker postcondition. It cannot be replaced by equality of file names or by an unchanged serialization format.

Chapter~\ref{Core:ch:revision} handles the dependency graph needed to know which such implications remain applicable. The present chapter handles information loss and operation transport. Together they separate three questions that a revision must answer: can the new record be recovered, does it support the required operations, and do the old proof premises still justify the intended conclusions?

\part{Proof and interpretation}
\chapter{Proof fields, closure, and executable certificates}\label{Core:book:chapter-14}
\setcounter{section}{0}\renewcommand{\thesection}{\thechapter.\arabic{section}}\label{Core:ch:prooffields}

A reasoning archive contains more than its final proofs. It may retain alternative derivations, conjectures, unsuccessful proposals, and observations that later become obsolete. A checker must know which of those records it is permitted to use. Otherwise an archive that is useful for discovery can become an unsound source of premises.

Proof-Path Theory supplies a finite distinction between the surrounding field and the certified derivation. This chapter admits its closure and trace-checking results in a self-contained form. It does not identify acceptance of a data structure with truth of an arbitrary sentence attached to that structure.

\section{A finite field with explicitly admitted rules}
Let $V=\{0,\ldots,n-1\}$ be a finite set of claim addresses. Let $A\subseteq V$ be admitted leaves, $S\subseteq V$ positive observations, and $N\subseteq V$ forbidden consequences. These sets need not be disjoint. A finite rule list $R$ has rows
\[
 r_j=(P_j,c_j),\qquad P_j\subseteq V,\quad c_j\in V.
\]
Only admitted rules are included. Identical rows may occupy different addresses. A conjectural rule elsewhere in the discovery archive is not admitted merely because it has a convenient conclusion.

A model of the field is a set $B\subseteq V$ containing $A$ and satisfying every implication $P_j\subseteq B\Rightarrow c_j\in B$. A protected extension also contains $S$ and avoids $N$. Define
\[
 C_0=A\cup S,\qquad
 C_{k+1}=C_k\cup\{c_j:P_j\subseteq C_k\}.
\]
The terms increase and lie in $V$, so the sequence stabilizes. Write $C$ for its stable value.

\begin{Coreproposition}[The least seeded model]\label{Core:r:field-closure}
The set $C$ is the least model containing $S$. There are at most $n-|A\cup S|$ strict expansion rounds. A protected extension exists if and only if $C\cap N=\varnothing$.
\end{Coreproposition}
\begin{proof}
At a fixed point every enabled rule has its conclusion in $C$, so $C$ is a model containing the seeds. If $B$ is any other model containing $S$, induction gives $C_k\subseteq B$: the seeds lie in $B$, and every rule enabled on $C_k$ has its conclusion in $B$. Hence $C\subseteq B$. Each strict expansion adds at least one of the finitely many nonseed nodes, giving the bound. If $C$ avoids $N$, it is itself a protected extension. If it meets $N$, every seeded model contains that forbidden node, so no protected extension exists.
\end{proof}

A cycle does not create a seed. With rules $u\to v$ and $v\to u$ and no admitted leaves or observations, the least model is empty. The set $\{u,v\}$ is closed, but its members have no derivation from the empty seed. This is why closedness alone is not an adequate proof certificate.

\section{Replay checks a finite causal sequence}
A trace $\sigma=(j_1,\ldots,j_m)$ lists rule addresses. Start from $K_0=A\cup S$. At each step, reject an unknown rule address or a rule whose premises are not all in the current known set. Otherwise adjoin its conclusion. \emph{Strict replay} also rejects a conclusion already known. \emph{Causal replay} allows it.

The distinction concerns the certificate's contract, not the truth of a repeated conclusion. Strict replay records a sequence of new insertions. Causal replay permits reuse of an existing derivation after stronger observations have already supplied one of its intermediate nodes.

\begin{Coretheorem}[Replay is sound and compositional]\label{Core:r:field-replay}
If replay accepts a trace from a set $K\subseteq C$ and returns $K'$, then
\[
 K\subseteq K'\subseteq C.
\]
For either fixed replay mode, replay of a concatenated trace is replay of its first part followed by replay of its second part from the returned set. Failure in either part is failure of the concatenated replay.
\end{Coretheorem}
\begin{proof}
The empty trace changes nothing. At a successful next step, all premises of an admitted rule already lie in $C$. Closedness of $C$ places its conclusion in $C$. Adjoining it preserves the old known set. Induction proves the inclusions. A strict novelty test can reject a step, but does not alter this reasoning for an accepted step.

For concatenation, induct on the first trace. Both computations inspect the same first address, the same current set, and the same enabling and novelty conditions. A rejection ends both computations. A successful step produces identical new sets, after which the induction hypothesis applies.
\end{proof}

\begin{Coretheorem}[Replay and saturation identify exact closure]\label{Core:r:field-exact}
A complete-closure certificate consists of a strict trace and a reported set $B$. Accept it only when replay from $A\cup S$ succeeds with output $B$ and $B$ satisfies every admitted rule. Acceptance implies $B=C$.
\end{Coretheorem}
\begin{proof}
Replay soundness gives $B\subseteq C$. Its extensivity ensures that $B$ contains both seed sets. The saturation check makes $B$ a seeded model, so leastness gives $C\subseteq B$. The two inclusions prove equality.
\end{proof}

The tests serve different purposes. Replay prevents unsupported extra nodes. Saturation prevents omission of a consequence whose rule has not yet been used. Neither test subsumes the other.

\section{A complete negative answer has a witness}
An obstruction certificate contains a forbidden target $v\in N$ and a causal trace deriving $v$ from $A\cup S$.

\begin{Coretheorem}[Finite obstruction witnesses are sound and complete]\label{Core:r:field-obstruction}
An accepted obstruction certificate exists exactly when no protected extension exists. No acyclicity assumption on the rule graph is required.
\end{Coretheorem}
\begin{proof}
An accepted certificate derives a node of $C\cap N$, so Proposition~\ref{Core:r:field-closure} excludes a protected extension.

Conversely, if there is no extension, choose $v\in C\cap N$. Every node in $C$ is reached at a finite stage. Starting with the seeds, process the rules that add the nodes of each successive stage, appending a rule only when its conclusion is not yet known. All its premises belong to an earlier stage and are already known. This yields a finite strict trace reaching the whole closure, and in particular $v$. Its prefix ending when $v$ is first present is also accepted in causal mode. If $v$ is a seed, the empty trace suffices. Cycles cause no difficulty because only rules with already known premises are applied.
\end{proof}

The argument constructs some witness, not a shortest one. Nor does it prove that every implementation of a witness producer terminates or faithfully indexes its source document. A complete result format can require both exact closure and a feasibility flag. A bare negative witness needs only its own replay and forbidden-target check.

\section{A worked four-node certificate}
Let $V=\{0,1,2,3\}$, $A=\{0\}$, $S=\varnothing$, and let the two rules be $0\to1$ and $1\to2$. The exact closure is $\{0,1,2\}$.

The trace containing only the first rule is a valid partial derivation of $\{0,1\}$, but it fails saturation. The set $\{0,1,2,3\}$ is saturated, but node $3$ cannot be produced by replay. The trace of both rules and the exact three-node set pass both tests.

\begin{figure}[htbp]\centering
\includegraphics[width=.94\linewidth]{sources/research/machine-native-reasoning-wider/figures/closure_contract.pdf}
\caption{Replay and saturation certify opposite inclusions. A partial derivation omits a consequence. A closed superset can contain an unsupported node. Their conjunction identifies the least seeded model.}\label{Core:fig:closure-contract}
\end{figure}

With forbidden set $N=\{3\}$, the exact closure is a positive feasibility certificate. With $N=\{2\}$, the same closure and the trace to node $2$ certify infeasibility. The request changed, even though the derivation did not. A receiving checker must therefore be given the intended request separately from the producer's proposed answer.

\section{Logical interpretation is an additional map}
Let $\mu(v)$ be the proposition represented by claim address $v$. Suppose every admitted leaf and positive observation is true under $\mu$, and each admitted rule satisfies
\[
 \left(\bigwedge_{p\in P_j}\mu(p)\right)\Longrightarrow\mu(c_j).
\]
Then every node in $C$ is true under $\mu$. Indeed, induction over the closure stages propagates truth from the seeds through the admitted implications.

The finite checker proves derivability relative to those inputs. It does not establish the truth of an empirical observation merely because that observation is placed in $S$. A node excluded by a protected request is likewise not necessarily a logically false proposition. Its appearance shows incompatibility with that request, whose meaning must be stated separately.

\section{The serialization and implementation boundary}
An exchange format must distinguish a problem from its proposed certificate, preserve rule addresses, and reject malformed node identifiers. In Python, for example, a Boolean passes a broad integer-instance test. An address decoder that intends exact integer tokens must exclude that case explicitly. Parsing, range checking, replay, saturation, and interpretation are separate operations.

The Proof-Path Theory source supplies typed problem and certificate records, proved encoder/decoder round trips, an executable Lean replay function, and an accepted-answer theorem at the Lean JSON syntax-tree boundary. Its text parser, file operations, Python implementation, and adapter from named source claims remain separate assurance layers. This chapter uses the finite closure and replay statements, not a blanket verification claim for that entire path.

\Coresource{Proof-Path Theory [PPT], ``Executable certificates for the declared closure,'' ``Finite witnesses for forbidden consequences,'' and ``Typed exchange and the boundary of serialization.'' The selected source statements are \texttt{res:trace:replay}, \texttt{res:trace:least}, \texttt{res:trace:complete}, and \texttt{res:trace:verdict}. Their original formal concordance is preserved with the source excerpt. The present notation and proofs were reviewed as a restatement, not recompiled into a new Lean development.}

\section{Where proof checking enters the reasoning calculus}
The machine-reasoning chain now has a concrete certificate interface. A discovery process can propose a finite derivation without becoming the authority on its validity. The checker verifies the derivation against a fixed problem. A revision can change the problem or its admitted rules, in which case the dependency and migration conditions must be checked again. Reusing the same bytes does not establish that the old certificate answers the new question.

% Reconstructed from physical pages 74--87 of the recovered 23 October PDF.
% Mathematical statements/proofs transcribed; explanatory prose adapted for this
% new assembly. This is not an authenticated recovery of the original TeX bytes.

\chapter{Transporting derivations between languages}\label{Core:book:chapter-15}
\setcounter{section}{0}\renewcommand{\thesection}{\thechapter.\arabic{section}}
Changing a proof language can introduce auxiliary claims, replace a rule by a derivation, or merge claims that the source language distinguishes. A dictionary of names does not justify these changes. A translation must preserve the inference being used. We distinguish forward proof transport, reflection of old consequences, and preservation of a protected request.

\section{A rule may translate to a derivation}
Let $V,W$ be finite claim-address sets. A source field has rules $R$ and seeds $K\subseteq V$; a target field has rules $R'$ and seeds $L\subseteq W$. Write
\[
C=\operatorname{Cl}_{R}(K),\qquad C'=\operatorname{Cl}_{R'}(L)
\]
for their least seeded closures. The seeds combine admitted leaves and positive observations. No logical interpretation is built into an address.
For an address map $f:V\to W$, require, for every source rule $P\Rightarrow c$,
\begin{equation}\label{Core:integration:rule-transport}
f(c)\in\operatorname{Cl}_{R'}(f(P)).
\end{equation}
The translated conclusion need not follow in one target step. It must be derivable from the translated premises; an empty-premise rule requires a derivation with no target assumptions. A finite causal trace witnessing this condition is a \emph{rule witness}. Its rule addresses refer to the fixed target table, including each row's actual premises and conclusion.

\begin{Corelemma}[Causal replay tolerates stronger premises]\label{Core:integration:stronger-premises}
Suppose a causal trace is accepted from $B\subseteq W$ and returns $D$. If $B\subseteq\widetilde B$, the same trace is accepted from $\widetilde B$ and returns a set $\widetilde D$ containing $D$.
\end{Corelemma}
\begin{proof}
At the first step, every premise present in the smaller known set is present in the larger one. The same rule address is valid in both computations. Adjoining its conclusion preserves inclusion of known sets. Induction over the finite trace proves acceptance and the final inclusion. A conclusion already present in the larger set causes no rejection in causal mode.
\end{proof}
Strict replay has a different contract. For $u\Rightarrow v$, the one-step trace is strict from $\{u\}$ but not from $\{u,v\}$: the promised new insertion has already occurred.

\section{Forward transport preserves derivability}
\begin{Coretheorem}[Derivation-valued rule transport]\label{Core:integration:forward-transport}
Suppose every source rule satisfies \eqref{Core:integration:rule-transport} and $f(K)\subseteq\operatorname{Cl}_{R'}(L)$. Then
\[
f\bigl(\operatorname{Cl}_{R}(K)\bigr)\subseteq\operatorname{Cl}_{R'}(L).
\]
Supplied target traces for the seed claims and rule witnesses constructively translate every accepted source causal trace into an accepted target causal trace.
\end{Coretheorem}
\begin{proof}
Induct over the increasing source closure sequence. The initial inclusion is the seed hypothesis. If a source rule adds $c$ after all members of $P$ have been reached, the induction hypothesis puts $f(P)$ in $C'$. A closed set containing $f(P)$ contains its least closure, so
\[
f(c)\in\operatorname{Cl}_{R'}(f(P))\subseteq C'.
\]
This proves the next-stage inclusion and hence the inclusion at stabilization.

For the constructive assertion, concatenate target traces deriving the members of $f(K)$ from $L$ in a fixed order. Lemma~\ref{Core:integration:stronger-premises} permits each trace after previous traces enlarge the known set. Traverse the source trace in its given order. Before $P\Rightarrow c$ is used, the target known set contains $f(P)$. Append its rule witness. The lemma ensures acceptance and adds $f(c)$. Induction on the source trace gives a target trace containing every translated source conclusion.
\end{proof}
The target may contain auxiliary claims with no source name. Forward transport gives inclusion, and does not require $f$ to be injective.

\begin{Corecorollary}[An explicit certificate-length bound]
Let a combined target seed-initialization trace have length $s$, and let the witness for rule $j$ have length $\lambda_j$. A source trace using $j_1,\ldots,j_m$ translates by the preceding construction to length
\[
s+\sum_{t=1}^{m}\lambda_{j_t}.
\]
If every used witness has length at most $\lambda_*$, this is at most $s+m\lambda_*$.
\end{Corecorollary}
\begin{proof}
The construction concatenates the seed trace and exactly one stored witness per source step. List-concatenation lengths add.
\end{proof}
This measures listed rule applications, including repetitions. It does not measure processor time, bytes, address-map construction, or witness discovery. Removing redundant steps requires checking the shorter trace in its chosen replay mode.

\section{Reflection requires a second interface}
Suppose $g:W\to V$ satisfies $gf=\operatorname{id}_V$. Require
\[
g(d)\in\operatorname{Cl}_{R}(g(Q))\quad\text{for every target rule }Q\Rightarrow d,
\qquad g(L)\subseteq C.
\]
These are local conditions on rules and seeds.
\begin{Coretheorem}[A reflective translation preserves old consequences]
Under the forward and backward conditions, for every $v\in V$,
\[
v\in C\quad\Longleftrightarrow\quad f(v)\in C'.
\]
For a forbidden set $B\subseteq V$ and target request forbidding $g^{-1}(B)$,
\[
C\cap B=\varnothing\quad\Longleftrightarrow\quad C'\cap g^{-1}(B)=\varnothing.
\]
Consequently the corresponding protected-extension feasibility is preserved.
\end{Coretheorem}
\begin{proof}
Apply Theorem~\ref{Core:integration:forward-transport} to $f$ and to $g$. It gives $f(C)\subseteq C'$ and $g(C')\subseteq C$. If $f(v)\in C'$, then $v=gf(v)\in C$; the converse is the first inclusion. If $w\in C'\cap g^{-1}(B)$, then $g(w)\in C\cap B$. Conversely, if $v\in C\cap B$, then $f(v)\in C'$ and $g(f(v))=v\in B$. Thus either forbidden intersection is nonempty exactly when the other is. The protected-extension criterion of Chapter~14 gives feasibility.
\end{proof}
A prohibition involving a genuinely new target distinction is a different request.
\begin{Corecounterexample}[A forward translation need not preserve feasibility]
Use source claims $\{u,v\}$, seed $\{u\}$, no rules, and forbid $v$. Its closure is $\{u\}$ and the request is feasible. Keep the claims and seed in the target and add $u\Rightarrow v$. The identity satisfies the vacuous forward rule condition and the seed condition, but the target closure contains $v$. The backward rule condition fails.
\end{Corecounterexample}

\section{A worked expansion of one inference step}
Take source seed $a$ and rules $a\Rightarrow b$, $b\Rightarrow c$. The target has addresses $a',z,b',c'$ and rules
\[
a'\Rightarrow z,\qquad z\Rightarrow b',\qquad b'\Rightarrow c'.
\]
The map $f$ sends each source address to its primed address. Its first rule witness passes through $z$; its second is the last target rule. Two source steps become three target steps. Define $g(a')=a$, $g(z)=g(b')=b$, and $g(c')=c$. The first and third target rules map to source rules. The middle rule maps to the already available $b\Rightarrow b$. All backward rule conditions hold. Forbidding $c$ corresponds to forbidding $c'$; forbidding $b$ corresponds to excluding both $z$ and $b'$.
\begin{figure}[htbp]\centering
\begin{tikzpicture}[node distance=17mm,>=Latex]
\node (a) {$a$};\node[right=of a] (b) {$b$};\node[right=of b] (c) {$c$};
\draw[->] (a)--(b);\draw[->] (b)--(c);
\node[below=15mm of a] (aa) {$a'$};\node[right=of aa] (z) {$z$};\node[right=of z] (bb) {$b'$};\node[right=of bb] (cc) {$c'$};
\draw[->] (aa)--(z);\draw[->] (z)--(bb);\draw[->] (bb)--(cc);
\end{tikzpicture}
\caption{A source inference expands to a checked target derivation. The backward interpretation places both $z$ and $b'$ over $b$.}
\end{figure}

\section{What a transport package must carry}
A usable package supplies both rule tables, both seed sets, the address map, and actual finite witnesses. Reflection additionally needs the backward map and its witnesses. A recipient can replay these local obligations before using a translated proof. Correct prose parsing, intended interpretation, and discovery of suitable witnesses remain separate interfaces.

\chapter{A finite language of executable methods}\label{Core:book:chapter-16}
\setcounter{section}{0}
A program can branch, stop after failure, invoke operations, and consume resources along its actual path. Primitive commuting equations must extend to these control decisions. We use finite syntax trees with partial primitives, sequencing and Boolean branching, without loops or recursive calls.

\section{Partial operations and observable tests}
Let $X$ be a finite nonempty state space and $A,B$ finite alphabets of primitives and tests. A primitive has partial semantics $F_a:X\to X_\bot$, where $X_\bot=X\sqcup\{\bot\}$. A test $b:X\to\{0,1\}$ is total and leaves the state unchanged. Failure does not imply rollback: the execution records states reached before failure.
Charges lie in $\mathbb N^d$ with fixed coordinate meanings. An attempted primitive, including a failed attempt, costs $c_a(x)$; evaluating a test costs $t_b(x)$. Programs have grammar
\[
p::=\mathsf{skip}\mid a\mid p;q\mid\mathsf{if}\ b\ \mathsf{then}\ p\ \mathsf{else}\ q.
\]
The alphabets, tables and resource meanings belong to the declared language.

\section{An execution records its path}
Evaluation returns a result in $X_\bot$, a resource vector, an event list and a state path. The path starts at the input and appends each successful primitive's successor. A test adds its Boolean outcome to the event list without extending the state path.

The program $\mathsf{skip}$ returns its input with zero charge and no events. A primitive records and charges its attempt, recording success or failure and appending a successor exactly on success. A sequence evaluates its left term, stops if that fails, and otherwise evaluates its right term from the returned state. Charges add and paths join at the common state. A conditional records and charges its test and evaluates only the selected branch. Stored syntax in an untaken branch contributes no events to that execution.

\begin{Coreproposition}[A structural bound on event count]
Define
\begin{align*}
w(\mathsf{skip})&=0,&w(a)&=1,\\
w(p;q)&=w(p)+w(q),&w(\mathsf{if}\ b\ \mathsf{then}\ p\ \mathsf{else}\ q)&=1+\max\{w(p),w(q)\}.
\end{align*}
Every evaluation of $p$, successful or failed, has at most $w(p)$ events.
\end{Coreproposition}
\begin{proof}
For a primitive and for $\mathsf{skip}$ the count is exact. A sequence evaluates its first term and either stops or adds an evaluation of its second; the induction bounds add. A conditional contributes one test event and one branch whose induction bound is at most the larger branch bound. These cases cover the grammar.
\end{proof}
Parsing, construction, certificate production and external waiting require separate charges in a larger workflow contract.

\section{Representation change must preserve the failure boundary}
Let $\pi:X\to Z$ be a surjection onto a finite coarse state space, and extend it by $\pi_\bot(\bot)=\bot$. Give each symbol a coarse primitive $\overline F_a:Z\to Z_\bot$ and test $\overline b:Z\to\{0,1\}$. Require
\begin{equation}\label{Core:integration:partial-semantics}
\pi_\bot F_a=\overline F_a\pi,\qquad b=\overline b\pi.
\end{equation}
This preserves and reflects success and compares projected successors. Replacing disabled primitives by successful idle operations violates the first equation. For resource transport additionally require
\[
c_a=\overline c_a\pi,\qquad t_b=\overline t_b\pi.
\]

\begin{Coretheorem}[Compositional execution through a quotient]\label{Core:integration:execution-quotient}
Under these hypotheses, evaluations of any program $p$ from $x$ and $\pi(x)$ have the same success or failure tag, event list, and resource vector. On success the coarse final state is the projection of the fine final state. In every case the complete coarse state path is the pointwise projection of the fine path.
\end{Coretheorem}
\begin{proof}
Proceed by structural induction on $p$. For $\mathsf{skip}$, paths contain only initial states and other records are empty or zero. For a primitive, \eqref{Core:integration:partial-semantics} gives the same tag, primitive event, and, on success, projected successor; the resource premise gives equal charges.

For $p;q$, apply the induction hypothesis to $p$. On failure both interpreters stop with corresponding records. On success its returned states correspond under $\pi$, so apply the hypothesis to $q$ from those states. Concatenation gives equal event lists and projected paths, counting the joining state once; addition gives equal resource vectors. The second evaluation gives the result tags and final states.

For a conditional, corresponding states give the same Boolean test outcome and the same test charge. Both interpreters select the same branch. Apply the induction hypothesis to that branch and append its record to the common test event.
\end{proof}
The coarse primitive and actual test are fixed before evaluation; neither a favorable branch nor a coarse successor is selected afterward.

\begin{Corecorollary}[Path predicates and coordinate budgets]
Suppose $S=\pi^{-1}(\overline S)$ and $Q=\pi^{-1}(\overline Q)$. Successful execution stays in $S$ and finishes in $Q$ exactly when coarse execution stays in $\overline S$ and finishes in $\overline Q$. Under cost transport, a coordinatewise bound $u\in\mathbb N^d$ holds in the fine execution exactly when it holds in the coarse execution.
\end{Corecorollary}
\begin{proof}
Each fine path state's membership in $S$ is equivalent to its projection's membership in $\overline S$. Theorem~\ref{Core:integration:execution-quotient} supplies this projected path and the successful endpoint. Apply the target pullback to the endpoint. Equal resource vectors give both directions of the budget assertion.
\end{proof}

\section{Three ways the contract can fail}
\begin{Corecounterexample}[A hidden test changes the branch]
Let $X=\{0,1\}$ and let $\pi$ merge both states. One primitive always succeeds without changing state and another always fails. Both partial primitives descend. A test reading the fine bit does not descend. The program choosing success at zero and failure at one therefore has no common coarse execution from the merged record.
\end{Corecounterexample}
A primitive enabled at one fine state and disabled at another in the same fibre also cannot descend. Costs varying within a fibre can invalidate equal-cost transport even when final results transport. These are distinct defects.

\section{A complete calibrated repair program}
Extend the sensor state to $(b,s,n,d)\in\{0,1\}^4$, where $d$ records repair completion. Retain
\[
\pi(b,s,n,d)=(b\mathbin\oplus s,s,d).
\]
This determines $b$ without retaining $n$. A test $\mathsf{zero}$ asks whether $(b\oplus s)\oplus s=0$. Primitive $R_i$ succeeds exactly when $b=i$ and sets $d=1$. Then
\[
\mathsf{if\ zero\ then}\ R_0\ \mathsf{else}\ R_1
\]
succeeds everywhere. Charge $(1,0)$ per primitive and $(0,1)$ per test: every run costs $(1,1)$. All primitive, test and cost equations hold in the eight-state coarse model. The theorem applies to every expression in the language. Without calibration, retaining only the display and $d$ loses the test factorization; a correct program has not thereby acquired missing input.

\section{Mathematical semantics and a host implementation}
For supplied finite tables, the local equations can be checked exhaustively; the theorem extends them to arbitrary finite syntax depth. The original reference interpreter is a concrete implementation rather than a formally verified compiler. Parsing, object representation, exception behavior and external action refinement retain their own obligations.

\chapter{From checked claims to guarded execution}\label{Core:book:chapter-17}
\setcounter{section}{0}
A certificate establishes a derivation; an action changes a state. Their connection requires a claim interpretation, sound starting observations, and a command bound to the same program and context used by the proof.

\section{Claims interpreted as state predicates}
Let $X$ be a finite state space, $r:X\to Y$ an observation map, and $(V,R)$ a finite claim field with interpretation $I:V\to\mathcal P(X)$. For attained $y$, let $K_y\subseteq V$ be its admitted seeds. Require
\begin{equation}\label{Core:integration:sound-observation}
r^{-1}\{y\}\subseteq I(v)\quad(v\in K_y),
\end{equation}
and for every rule $P\Rightarrow c$,
\begin{equation}\label{Core:integration:sound-rule}
\bigcap_{v\in P}I(v)\subseteq I(c).
\end{equation}
The empty intersection is $X$. Seeds must hold throughout the observation fibre, rather than at a selected representative.

\begin{Coretheorem}[Interpreted replay from accessible evidence]\label{Core:integration:interpreted-replay}
Under \eqref{Core:integration:sound-observation} and \eqref{Core:integration:sound-rule}, every claim in $\operatorname{Cl}_R(K_y)$ holds at every state in $r^{-1}\{y\}$. In particular, an accepted trace deriving $q$ certifies $r^{-1}\{y\}\subseteq I(q)$.
\end{Coretheorem}
\begin{proof}
Fix $x\in r^{-1}\{y\}$. Every seed holds at $x$. Whenever all premises of a rule hold at $x$, the rule inclusion puts $x$ in the conclusion's interpretation. Induction on closure stages proves the assertion throughout the closure. Replay soundness puts every accepted trace conclusion in that closure. The choice of $x$ was arbitrary.
\end{proof}

\section{A guard names an actual execution contract}
Choose a finite program catalogue from Chapter~16. For each $p$, specify a safety set $S_p$, target $Q_p$, and coordinatewise budget $u_p$. Define $G_p$ as precisely the starting states from which evaluation succeeds, all recorded states lie in $S_p$, the final state lies in $Q_p$, and cost is at most $u_p$. For a designated guard claim $q_p$, require
\begin{equation}\label{Core:integration:guard}
I(q_p)\subseteq G_p.
\end{equation}
The bridge concerns the actual syntax, primitive tables, tests and cost account.

\begin{Coretheorem}[One guarded dispatch from a common witness]
Suppose the observation, rule and guard conditions hold. Let the current state be $x$, the current record $y=r(x)$, and suppose replay from $K_y$ derives $q_p$. If dispatch evaluates the same $p$ in the same modeled state $x$, execution succeeds, remains in $S_p$, finishes in $Q_p$, and costs at most $u_p$.
\end{Coretheorem}
\begin{proof}
Theorem~\ref{Core:integration:interpreted-replay} gives $x\in I(q_p)$. The guard inclusion gives $x\in G_p$. Its definition gives all four conclusions for this evaluation.
\end{proof}
The result uses one observation context, interpretation, program and resource account. Separate existence results with incompatible witnesses do not supply these premises.

\section{What a stale record invalidates}
For a bit $b$, let $R_i$ be enabled exactly when $b=i$. An observation of zero can justify $R_0$, but an unobserved flip before dispatch makes the recorded context stale. The derivation remains valid in its original context. Fresh observation or a guard robust to intervening changes is needed.
Let $\Delta\subseteq X\times X$ describe those changes and define
\[
\Delta(B)=\{x':\exists x\in B,\ (x,x')\in\Delta\}.
\]
\begin{Coreproposition}[A guard robust to declared intervening changes]
If a checked claim establishes $r^{-1}\{y\}\subseteq I(q_p)$ and $\Delta(I(q_p))\subseteq G_p$, dispatch is valid from every current state reached by one declared change from the observed fibre. For a bounded sequence, replace $\Delta$ by the corresponding relational composition.
\end{Coreproposition}
\begin{proof}
Monotonicity of relational image gives
\[
\Delta(r^{-1}\{y\})\subseteq\Delta(I(q_p))\subseteq G_p.
\]
Membership in $G_p$ supplies the contract. A composed relation records the permitted finite sequences under its composition convention, so the same inclusion applies.
\end{proof}
If both holding and flipping the bit are possible, neither state-specific repair has the robust guard. A changed command binding is another failure: executing $R_1$ for a certificate naming $R_0$ violates program identity. A checksum helps identify bytes; it does not prove their external operational meaning.

\section{Repeated dispatch and the boundary of a safety claim}
Let $E\subseteq X$ be a decision-epoch invariant and $S\supseteq E$ a safety set. Every admitted dispatch from the actual current state in $E$ has a checked guard, a path in $S$, and endpoint in $E$. Rejected proposals execute no program in this model.
\begin{Coretheorem}[Safety under repeated guarded dispatch]
Starting in $E$, every finite sequence of admitted dispatches remains in $S$ at every recorded state and returns to $E$ after every completed dispatch. The same safety statement holds for every finite prefix of an infinite sequence.
\end{Coretheorem}
\begin{proof}
The first guard gives a safe execution and an endpoint in $E$, satisfying the next dispatch's invariant. Induct on the number of completed dispatches. Each recorded state of an infinite execution belongs to a finite prefix, giving the prefix assertion.
\end{proof}
This does not guarantee another admitted proposal, progress, or interruption at an internal instruction. Macro-action interruption needs its own semantics. To preserve an ending strategy, $E$ must be the relevant safe-running region and the specified ending controller must remain available.

\section{The full ledger must charge rejected work too}
Let $b_j\in\mathbb N^d$ be remaining resources before workflow step $j$ and $c_j$ its complete actual charge. Evidence acquisition, proposals, checking, execution and rejection all fall within the chosen accounting contract. Require coordinatewise
\[
c_j\le b_j,\qquad b_{j+1}=b_j-c_j.
\]
\begin{Coreproposition}[A finite-prefix resource invariant]
For each completed prefix of length $m$,
\[
b_m+\sum_{j=0}^{m-1}c_j=b_0.
\]
Thus accumulated charges do not exceed the initial reserve in any coordinate.
\end{Coreproposition}
\begin{proof}
For one step this is the accounting rule, with subtraction defined because $c_j\le b_j$. Sum $b_j=b_{j+1}+c_j$ over the prefix. Intermediate reserves cancel. Nonnegativity of $b_m$ gives the bound.
\end{proof}
Reserving a proved upper bound before a step and returning unused resources is compatible with this invariant when every performed operation is included. Zero-cost modeled steps remain possible, so this is not a termination theorem.

\section{A worked certificate-to-program chain}
Use the calibrated sensor $(b,s,n,d)$ and record $(b\oplus s,s,d)$. Interpret $z_0,z_1$ by $b=0,b=1$; the attained record determines the sound seed. Add guards $g_i$ with the same interpretations and sound rules $z_i\Rightarrow g_i$. The catalogue contains $R_0,R_1$. Each matching guard ensures completion with $d=1$ and one primitive attempt.

A producer proposes a program and trace. The checker derives that program's own guard from the current record's seeds; dispatch executes that exact program. The wrong repair has no enabled derivation from the correct seed. A forged seed fails the observation interface.
\begin{figure}[htbp]\centering
\begin{tikzpicture}[>=Latex,node distance=9mm,
box/.style={draw=Ink,rounded corners,align=center,text width=34mm,minimum height=13mm,font=\small}]
\node[box] (r) {Current record\\$y=r(x)$};
\node[box,right=of r] (t) {Accepted trace\\$K_y\vdash q_p$};
\node[box,right=of t] (p) {Same program\\execute $p$};
\draw[->] (r)--(t);\draw[->] (t)--(p);
\end{tikzpicture}
\caption{A valid derivation and dispatch share an observation, interpretation, program and guard. State changes require revalidation; changed bindings require a new semantic bridge.}
\end{figure}
The finite chain closes a mathematical interface. Named-source adapters, learned model acquisition, host-language verification and external dispatcher refinement retain separate obligations.


\chapter{From the finite chain to the wider theory}\label{Core:book:chapter-18}
\label{Core:ch:frontier}

The selected reasoning chain is coherent where its maps share an explicit domain, operational semantics, and evidence language. That does not imply that every proposed component of a machine-native reasoning system has already been joined to every other one. It is useful to state what remains separate, because those boundaries suggest the most precise next theorems.

\section{Existence and acquisition are different problems}
The representation theorems identify when a smaller state \emph{can} preserve selected future observations. They do not generally tell an agent which smaller state to construct from sparse experience. Learning such a state requires a declared hypothesis family, a policy for obtaining observations, and conditions under which incompatible hypotheses can be separated. Statistical identifiability is not certain finite identification. Every finite binary sequence can occur with positive probability under two distinct Bernoulli parameters in $(0,1)$.

A meaningful acquisition result would specify the loss or property of a proposed representation, the permitted query or intervention family, the distribution of observations and the tolerable failure probability. A failure certificate may show the available evidence cannot separate two task-relevant worlds. The resulting impossibility is relative to that evidence interface, not an unconditional impossibility of machine understanding.

\section{A learned semantic map is not a written commuting equation}
Conditional method transport has the form $TF_a=G_{\tau(a)}T$. This is meaningful when $T$, the operator translation $\tau$, and the applicable domains are already specified. A learned system must acquire them, ensure they cover its actual state trajectories, and show the diagrams continue to commute under the relevant interventions. A correct table for a finite declared task is a useful positive control, but it cannot certify that a large recurrent network learned a stable operator language across changed domains.

The Neuralese research documents both finite positive controls and failed stability or scientific gates. Those negative results delimit the next semantic acquisition contracts. An honest mathematical volume keeps the theorem about \emph{given} translators separate from an experiment asking whether one was discovered.

\section{A certificate is not a validated world action}
Chapter~\ref{Core:ch:prooffields} gives the finite closure and replay interface from Proof-Path Theory. Its typed justification fields also distinguish mathematical memory from a completed certificate. A checked proof records a consequence of its premises. Passing a certificate through an executable planning or action system requires additional semantics: which world facts correspond to the propositions, who may perform the selected action, and what actual state change the dispatcher implements.

The Typed Proof Learner and Resolution Completion branches retain relevant constructions, but the recovered handoff does not claim a complete connection from checked propositional derivation to world-action dispatch. Even a fully proved checker theorem would not automatically verify its external parser, message transport, or control interface. Those are additional refinement obligations.

\section{Revision requires dependency completeness}
A revision may change the observation vocabulary, the task, the permitted actions, the model family or the representation used to express an old result. The statement of a theorem can remain true while its application becomes invalid. A retained dependency graph makes selective invalidation possible if the graph is complete for the applicable premises. Matching a recorded digest checks an integrity claim about bytes under its hashing assumptions. It does not establish the adequacy of the recorded dependency set.

Chapter~\ref{Core:ch:migration} gives an exact criterion for migrating finite records and transporting common operations. A broader revision theorem must additionally carry a translation of certificate premises and prove that every retained conclusion remains justified when the model, evidence language, and allowed operations change together. In the absence of those data, an old result should remain accessible as historical mathematics but should not automatically authorize a new action. Formalizing that rule is more valuable than assigning another name to a ``self-correcting agent.''

\section{Resources, workspaces, and authority}
Improvement claims also require a common cost ledger. Reduced operator calls may be outweighed by the preparation, validation or storage needed to obtain a reusable method. Different price weights can change the sign of a proposed saving even when semantic correctness remains unchanged. The reviewed core already formalizes a finite version of this principle. Broad claims of scalable or indefinite improvement require additional assumptions about task distributions, resource limits and repeated acquisition.

Likewise, jointly reconstructible information is not sufficient for a functioning workspace or a machine society. Causal relevance, persistence, identity, access rules, permissions, institutional formation and accountable endings require their own theorems. The safe-ending calculation of Chapter~\ref{Core:ch:endinganalysis} uses a stipulated action system. It does not show that authority was legitimately established or that real institutions will comply with a command.

\section{A bounded agenda}
The strongest next steps are mathematical and testable. Complete the four selected Lean correspondence targets. Extend the exact migration and support-enclosure contracts to revisions of observation models and certificate semantics. Add explicit no-go examples for proposed learned semantic translators when an observation interface is insufficient. Then establish restricted refinement theorems linking a certificate checker to a producer and to a declared runtime, rather than promising end-to-end verification of every historical component. These are substantial research tasks, but each admits finite exact examples, sharp hypotheses and ordinary mathematical review before any large-scale training campaign.

\part{Komorebi Space: realization, obstruction, and repair}
\setcounter{section}{0}\renewcommand{\thesection}{KS.\arabic{section}}\renewcommand{\theHsection}{KS.\arabic{section}}

% BEGIN retained input sources/monographs/komorebi-space/source/chapters/01_views.tex

\chapter{Realizations and their obstructions}\label{KS:book:chapter-19}
A global object is not added to local data by terminology. It is a family satisfying every prescribed transition equation. The first task is to construct the space of such families and to keep its possible emptiness visible.

\section{Views and coherent realizations}\label{KS:sec:views}
\subsection{The data of a view system}
Fix a small category \(\KSV\), or work in fixed universes so that the products below exist in the chosen ambient universe. A \emph{view} is an object of \(\KSV\). It can encode a resolution level, a coordinate presentation, or a context in which part of an object is visible. These interpretations are mathematical choices, not additional axioms.

\begin{KSdefinition}[View system]
A view system consists of sets \(F_v\) and maps
\[
 f^*:F_w\longrightarrow F_v\qquad(f:v\longrightarrow w)
\]
satisfying \((1_v)^*=\KSid\) and \((g\circ f)^*=f^*g^*\). Equivalently, it is a functor \(F:\KSV^{\mathrm{op}}\to\KSSet\).
\end{KSdefinition}
The direction matters. An arrow \(v\to w\) induces a restriction from the fiber at \(w\) to the fiber at \(v\). A restriction need not be invertible. A category with several arrows between two views can retain different routes between their descriptions. Replacing it by a preorder can lose those routes and the constraints they impose.

\begin{KSdefinition}[Strict Komorebi realization]
The ambient realization set is
\begin{equation}\label{KS:eq:ambient}
 \mathfrak K(F)=\Gamma(F)=
 \left\{x\in\prod_{v\in\KSOb\KSV}F_v:
 f^*(x_w)=x_v\text{ for every }f:v\to w\right\}.
\end{equation}
A point of \(\mathfrak K(F)\) is a whole coherent family. Its value in view \(v\) is the evaluation \(\operatorname{ev}_v(x)=x_v\).
\end{KSdefinition}
The raw product contains all families, including inconsistent ones. Equation~\eqref{KS:eq:ambient} cuts out the exact families. No nonemptiness hypothesis is built into the definition.

\begin{KStheorem}[Joint determination and the universal property]\label{KS:KS01}
The evaluations jointly determine a realization. If a set \(T\) carries maps \(c_v:T\to F_v\) satisfying \(f^*c_w=c_v\), there is a unique map \(c:T\to\Gamma(F)\) satisfying \(\operatorname{ev}_v c=c_v\) for every view.
\end{KStheorem}
\begin{proof}
We first verify that the evaluation maps themselves form a compatible family. If $x\in\Gamma(F)$ and $f:v\to w$, the defining equation for $x$ gives
\[
 f^*(\operatorname{ev}_w x)=f^*(x_w)=x_v=\operatorname{ev}_v x.
\]
If $x,y\in\Gamma(F)$ have equal evaluations, then $x_v=y_v$ for every $v$. Equality of functions into a dependent product gives equality of their underlying families. Membership in $\Gamma(F)$ is a property of a family, not a further choice of coordinate. Thus $x=y$.

Now fix the compatible maps $c_v:T\to F_v$. For each $t\in T$ define a family by $x(t)_v=c_v(t)$. For every arrow,
\[
 f^*x(t)_w=f^*c_w(t)=c_v(t)=x(t)_v.
\]
This calculation verifies membership in $\Gamma(F)$, so the rule $c(t)=x(t)$ defines an actual function $c:T\to\Gamma(F)$. Evaluating it gives $\operatorname{ev}_vc(t)=c_v(t)$, as required.

To prove uniqueness, let $b:T\to\Gamma(F)$ be another such function. At every $t$ and $v$,
\[
 (b(t))_v=\operatorname{ev}_v b(t)=c_v(t)=\operatorname{ev}_v c(t).
\]
Joint determination first yields $b(t)=c(t)$, and function extensionality then yields $b=c$. Nothing in this argument assumes that $T$ or $\Gamma(F)$ is inhabited. For example, when $T$ is empty there is the unique empty function. When $T$ is inhabited, existence of the compatible family $(c_v)$ is already substantive data.
\end{proof}
\KSsource{KS01}
Thus the strict construction is precisely the ordinary limit of the view diagram, with the usual limit interpretation in the sequential case~\cite{KS:stacks-limits}. This identifies its mathematical content and its universal property. It does not supply a universal view diagram for every possible mathematical question.

\begin{figure}[htbp]\centering\resizebox{.96\linewidth}{!}{
% BEGIN retained input sources/monographs/komorebi-space/source/figures/01_views.tex
\begin{tikzpicture}[x=1cm,y=1cm]
\node[box,minimum width=4.2cm] (g) at (0,2.7) {$\Gamma(F)$\\a whole coherent family};
\node[box,minimum width=2.6cm] (v) at (-4.7,0) {$F_v$\\appearance $x_v$};
\node[box,minimum width=2.6cm] (w) at (0,0) {$F_w$\\appearance $x_w$};
\node[box,minimum width=2.6cm] (z) at (4.7,0) {$F_z$\\appearance $x_z$};
\draw[arr] (g.south west) -- node[lab,above left] {$\operatorname{ev}_v$} (v.north);
\draw[arr] (g) -- node[lab,right] {$\operatorname{ev}_w$} (w);
\draw[arr] (g.south east) -- node[lab,above right] {$\operatorname{ev}_z$} (z.north);
\draw[arr] (w) -- node[lab,above] {$f^*$} (v);
\draw[arr] (z) -- node[lab,above] {$g^*$} (w);
\draw[arr] (z.south) to[out=-90,in=-90] node[lab,below] {$(g\circ f)^*=f^*g^*$} (v.south);
\end{tikzpicture}

}\caption{The same realization is evaluated in several views. Compatibility is imposed along the arrows, not inferred from the fact that all displayed fibers exist. The ambient family belongs to the simultaneous solution set of the transition equations.}\label{KS:fig:views}\end{figure}

\subsection{Maps of realizations}
Let \(\alpha:F\to G\) consist of maps \(\alpha_v:F_v\to G_v\) with
\[
 \alpha_v f_F^*=f_G^*\alpha_w.
\]
The induced map on realizations is \(\Gamma(\alpha)(x)_v=\alpha_v(x_v)\). The naturality square is exactly what is needed to show that the image family is coherent.

\begin{KSproposition}[Functorial transport]\label{KS:KS02}
These induced maps preserve identities and composition:
\[
 \Gamma(\KSid_F)=\KSid_{\Gamma(F)},\qquad
 \Gamma(\beta\alpha)=\Gamma(\beta)\Gamma(\alpha).
\]
\end{KSproposition}
\begin{proof}
Evaluate each side at an arbitrary view and apply the corresponding identity for the coordinate maps. Theorem~\ref{KS:KS01} then identifies the maps on the whole realization set.
\end{proof}
\KSsource{KS02}
An arbitrary collection of fiber maps need not define a map on realizations. Its compatibility equations are part of the theorem, not a convention that can be omitted after a notation has been chosen.

\subsection{A global view and a useful warning}
\begin{KSproposition}[Terminal evaluation]\label{KS:KS03}
If \(t\) is terminal in \(\KSV\), evaluation induces an equivalence \(\Gamma(F)\simeq F_t\).
\end{KSproposition}
\begin{proof}
For \(a\in F_t\), take its restriction along the unique arrow \(v\to t\) at each view. Uniqueness of the arrow to \(t\), together with the functor laws, proves coherence. Evaluation at \(t\) recovers \(a\). Conversely, any coherent family is the restriction of its coordinate at \(t\).
\end{proof}
\KSsource{KS03}
This proposition is both useful and diagnostic. If the supposedly reconstructed global object has already been placed in a terminal fiber, the limit recovers that fiber. The proposition does not prove that arbitrary partial data have a preimage there.

The category of elements \(\int_\KSV F\) and \(\Gamma(F)\) also have different roles. The former organizes individual appearances and their transitions. The latter selects coherent families across all views. The category of elements serves here as explanatory notation, while the Lean formalization directly represents families and their coherence.

\chap{Partial data, gluing, and correspondence fibers}\label{KS:sec:glue}
\subsection{Prescribed views}
For \(U\subseteq\KSOb\KSV\), use the full subcategory on \(U\). A partial section \(s\in\Gamma_U(F)\) satisfies all transition equations whose endpoints lie in \(U\). Define
\begin{equation}\label{KS:eq:extension}
 \KSExt_U(s)=\{x\in\Gamma(F):x|_U=s\}.
\end{equation}
This fiber can be empty, a singleton, or much larger. A partial section is already coherent on its prescribed domain. A raw assignment that fails an internal transition equation is not a partial section of this kind.

Let \((U_i)_{i\in I}\) be full patches. A matching family consists of \(s_i\in\Gamma_{U_i}(F)\) with equal values on every overlap. The relevant covering condition is
\begin{equation}\label{KS:eq:transitioncover}
 \text{for every }f:v\to w,\quad
 \text{some }U_i\text{ contains both }v\text{ and }w.
\end{equation}
Identities imply that this condition also covers every object.

\begin{KStheorem}[Strict gluing]\label{KS:KS04}
Under \eqref{KS:eq:transitioncover}, restriction induces an equivalence
\[
 \Gamma(F)\simeq\KSMatch(F,(U_i)).
\]
Every fixed matching family has a unique global glue.
\end{KStheorem}
\begin{proof}
\emph{Step 1: every coordinate can be selected.} Apply the cover condition to the identity arrow of $v$. It gives an index containing $v$. Choose one such index $i(v)$ and put
\[
 x_v=(s_{i(v)})_v.
\]
The choice is an existence construction. It is not an algorithm selecting patches in a preferred order. If $v\in U_j$, matching on the overlap $U_{i(v)}\cap U_j$ gives
\[
 x_v=(s_j)_v.                                                   \tag{*}
\]
In particular the value is independent of the selected index.

\emph{Step 2: the selected coordinates satisfy every arrow equation.} For $f:v\to w$, choose a single index $j$ with $v,w\in U_j$. Formula (*) applies to both endpoints. Since the patch is full, it contains this actual arrow, even when other arrows have the same endpoints. Coherence of $s_j$ gives
\[
 f^*x_w=f^*(s_j)_w=(s_j)_v=x_v.
\]
Thus $x\in\Gamma(F)$. Object coverage alone would not supply the common patch used in this step.

\emph{Step 3: assembly and restriction are inverse.} Formula (*) proves $x|_{U_j}=s_j$ for every $j$. Conversely, start with $y\in\Gamma(F)$ and assemble its restrictions. The resulting coordinate at $v$ is $y_v$, regardless of which containing patch is chosen. Joint determination gives the original $y$. These two identities are exactly the inverse identities for the displayed equivalence.

For uniqueness with a fixed matching family, any other glue has coordinate $(s_{i(v)})_v=x_v$ at every $v$, hence equals $x$. If the view category is empty, its product has one empty family and the same conclusion follows without making any choices.
\end{proof}
\KSsource{KS04}
The proof explains why object coverage alone is insufficient. Coherence is an equation attached to an arrow, so existence of a glue must see that equation somewhere in the local data.

\begin{KSexample}[An object cover with no glue]\label{KS:KS05}
Take two views \(0\to1\), Boolean fibers, and the identity restriction. The singleton patches cover all objects. A local value false at \(0\) and true at \(1\) agrees on the empty overlap, yet no global section realizes both values. A global family would have to assign the same Boolean to both views.
\end{KSexample}
\KSsource{KS05}
Object coverage does prove uniqueness of a possible glue. It does not prove its existence. This distinction persists when sets are replaced by objects with symmetries.

\subsection{Correspondences must also be coherent}
Suppose natural maps form a span \(F\xleftarrow{p}M\xrightarrow{q}G\), and fix \(a\in\Gamma(F)\), \(b\in\Gamma(G)\). At each view define the ordinary fiber
\[
 B_{a,b}(v)=\{m\in M_v:p_v(m)=a_v,\ q_v(m)=b_v\}.
\]
Naturality and coherence provide its restrictions. A global bridge is an element of
\[
 \mathfrak B(a,b)=\{z\in\Gamma(M):\Gamma(p)z=a,\ \Gamma(q)z=b\}.
\]

\begin{KStheorem}[Fiber and section compatibility]\label{KS:KS06}
There is a coordinate preserving equivalence
\[
 \Gamma(B_{a,b})\simeq\mathfrak B(a,b).
\]
Under a transition cover, matching local bridge sections glue to a unique global bridge with both prescribed projections unchanged.
\end{KStheorem}
\begin{proof}
We must first check that the displayed fibers really form a view system. Let $m\in B_{a,b}(w)$ and $f:v\to w$. Naturality of $p$ and coherence of $a$ give
\[
 p_v(f_M^*m)=f_F^*p_w(m)=f_F^*a_w=a_v.
\]
The same calculation with $q,b$ gives $q_v(f_M^*m)=b_v$. Therefore $f_M^*m$ belongs to $B_{a,b}(v)$. Identity and composition of these restricted maps are inherited from $M$, because equality in a subset is equality of its underlying elements.

Given $u\in\Gamma(B_{a,b})$, forget each fiber-membership condition and call the underlying family $z_v\in M_v$. Coherence of $u$ is precisely coherence of $z$, so $z\in\Gamma(M)$. The equations $p_vz_v=a_v$ hold at every coordinate. Joint determination in $\Gamma(F)$ yields $\Gamma(p)z=a$, and the equations for $q$ similarly yield $\Gamma(q)z=b$. Thus $z$ is a global bridge.

In the opposite direction, let $z\in\mathfrak B(a,b)$. Evaluate its two global projection equations at $v$. They show that $z_v$ belongs to $B_{a,b}(v)$, and its original coherence shows that these elements form a section of the fiber system. Forgetting and reattaching membership conditions preserves every coordinate, so the two constructions are inverse.

Finally apply strict gluing to this specific fiber system. The glued section retains its membership in every $B_{a,b}(v)$, which is why both prescribed projections survive gluing. Uniqueness is for the fixed matching bridge family. Different matching families for the same $(a,b)$ can still yield different global bridges.
\end{proof}
\KSsource{KS06}
These are ordinary fibers. They do not retain comparison isomorphisms or automorphisms. Section~\ref{KS:sec:relative} constructs the appropriate groupoid version when this information matters.

\subsection{Locally inhabited does not mean globally inhabited}
\begin{KSexample}[The loop with no coherent choice]\label{KS:KS07}
Let the view category have one object and automorphism group \(C_2\). Let its nonidentity arrow act on a set with two elements by swapping the elements. The fiber is nonempty and every transport is bijective. Nevertheless a global section would be a fixed point of the swap, and there is none.
\end{KSexample}
\KSsource{KS07}
The problem is not missing local elements or noninvertible transition maps. It is the absence of an element satisfying the loop equation. Inserting a point in each fiber does not certify a coherent selection of those points.

\begin{insight}
The strict theory already distinguishes three levels: possible local appearances, compatible partial descriptions, and an actual global realization. A correspondence must pass through the same levels. Local nonemptiness of its fibers is not a globalization theorem.
\end{insight}



% BEGIN retained input sources/monographs/komorebi-space/source/chapters/02_obstructions.tex
\chap{The additive obstruction and the space of extensions}\label{KS:sec:obstruction}
The strict definition gives a system of equations. In an additive setting, that system has an actual differential, and the failure to extend prescribed data has a quotient valued obstruction.

\subsection{The compatibility differential}
Assume each fiber \(A_v\) is an abelian group and each restriction is additive. Set
\[
 E=\prod_v A_v,\qquad R=\prod_{f:v\to w}A_v,
 \qquad (\delta x)_f=f^*(x_w)-x_v.
\]
The group \(R\) retains an entry for every arrow. It is not silently identified with a cohomology group.

\begin{KSproposition}[Differential identities]\label{KS:KS08}
The map \(\delta:E\to R\) is additive, \(\Gamma(A)=\ker\delta\), and
\[
 (\delta x)_{1_v}=0,\qquad
 (\delta x)_{g\circ f}=f^*((\delta x)_g)+(\delta x)_f.
\]
\end{KSproposition}
\begin{proof}
Additivity follows coordinate by coordinate. Vanishing is exactly the defining coherence equation. For composition, add and subtract \(f^*x_w\):
\[
 f^*g^*x_z-x_v
 =f^*(g^*x_z-x_w)+(f^*x_w-x_v).
\]
The identity equation uses \((1_v)^*=\KSid\).
\end{proof}
\KSsource{KS08}
This exact formula is more informative than a scalar error score. It specifies how defects transform and combine. Norms can be added later, but they do not replace the typed residual.

\subsection{The admissible corrections}
Fix \(U\subseteq\KSOb\KSV\). Write \(p_U:E\to\prod_{v\in U}A_v\) for raw restriction. A correction preserving the prescribed values belongs to \(\ker p_U\). Its possible defects form the subgroup
\[
 B_U=\delta(\ker p_U)\le R.
\]
For a partial section \(s\), choose any raw extension \(\widetilde s\), for example zero off \(U\), and define
\begin{equation}\label{KS:eq:obstruction}
 Q_U=R/B_U,\qquad
 \mathfrak o_U(s)=[\delta\widetilde s]\in Q_U.
\end{equation}
If another raw extension is chosen, their difference vanishes on \(U\), so its differential belongs to \(B_U\). Thus the class is independent of the lift.

\begin{KStheorem}[Extension criterion]\label{KS:KS09}
For every coherent partial section,
\[
 \mathfrak o_U(s)=0\quad\Longleftrightarrow\quad\KSExt_U(s)\ne\varnothing.
\]
\end{KStheorem}
\begin{proof}
Let $\pi_U:R\to R/B_U$ be the quotient homomorphism. By the definition of a quotient by a subgroup,
\[
 \pi_U(r)=0\quad\Longleftrightarrow\quad r\in B_U.
\]
Here $B_U$ is the \emph{image} of $\ker p_U$ under $\delta$. Consequently membership of $\delta\widetilde s$ means the existence of a specific $h\in E$ satisfying both $p_Uh=0$ and $\delta h=\delta\widetilde s$. The subgroup need not be closed in any topology. No topology has been introduced.

If such $h$ exists, set $x=\widetilde s-h$. Additivity gives the two equations separately:
\[
 \delta x=\delta\widetilde s-\delta h=0,
 \qquad p_Ux=p_U\widetilde s-p_Uh=s.
\]
The first equation puts $x$ in $\Gamma(A)$, and the second puts that coherent family in $\KSExt_U(s)$.

Conversely, take $x\in\KSExt_U(s)$ and set $h=\widetilde s-x$. Since both raw families restrict to $s$, one has $p_Uh=0$. Since $\delta x=0$, one has $\delta h=\delta\widetilde s$. This witnesses $\delta\widetilde s\in B_U$, hence vanishing of its quotient class.

The proof is independent of the selected raw lift. Indeed, if $\widetilde s'$ is another lift, $k=\widetilde s'-\widetilde s$ lies in $\ker p_U$, and
\[
 \delta\widetilde s'-\delta\widetilde s=\delta k\in B_U.
\]
Thus the two quotient classes agree before either feasibility direction is invoked. A raw lift always exists in this additive model by using zero outside $U$. Constructing a raw lift does not solve the coherence equations.
\end{proof}
\KSsource{KS09}
The quotient is constructed before the solvability test. The theorem therefore does not define obstruction circularly by saying that it vanishes when a solution exists. It also does not give a decision procedure for membership in an arbitrary subgroup of an infinite product.

\begin{figure}[htbp]\centering\resizebox{.96\linewidth}{!}{
% BEGIN retained input sources/monographs/komorebi-space/source/figures/02_extension.tex
\begin{tikzpicture}[x=1cm,y=1cm]
\node[box,text width=3.15cm,minimum height=1.3cm] (raw) at (-5,1.6) {raw extension\\$p_U\widetilde s=s$};
\node[box,text width=2.8cm,minimum height=1.3cm] (res) at (0,1.6) {measured defect\\$\delta\widetilde s$};
\node[box,text width=3.55cm,minimum height=1.3cm] (obs) at (5,1.6) {obstruction class\\$[\delta\widetilde s]\in R/\delta(\ker p_U)$};
\node[box,text width=3.15cm,minimum height=1.3cm] (allow) at (-5,-1.5) {permitted change\\$h\in\ker p_U$};
\node[box,text width=2.8cm,minimum height=1.3cm] (image) at (0,-1.5) {removable defect\\$\delta h$};
\node[box,text width=3.55cm,minimum height=1.3cm] (zero) at (5,-1.5) {zero in the quotient\\$[\delta h]=0$};
\draw[arr] (raw) -- node[lab,above] {$\delta$} (res);
\draw[arr] (res) -- node[lab,above,font=\footnotesize] {quotient} (obs);
\draw[arr] (allow) -- node[lab,above] {$\delta$} (image);
\draw[arr] (image) -- node[lab,above,font=\footnotesize] {quotient} (zero);
\draw[<->,thick,draw=KomorebiOchre] (res.south)--(image.north);
\node[lab,fill=white,inner sep=4pt] at (0,.05) {$\delta\widetilde s=\delta h$};
\node[lab,text width=4.0cm] at (-4.8,.05) {This equality is the\\repair equation.};
\node[lab,text width=4.1cm] at (5,.05) {Solvable exactly when\\the top class is zero.};
\node[lab,text width=13.7cm] at (0,-2.8) {Then $x=\widetilde s-h$ satisfies both $\delta x=0$ and $p_Ux=s$.};
\end{tikzpicture}

}\caption{The raw family is repaired only along directions invisible on the prescribed views. The obstruction records whether its defect lies in the image of those admissible directions. Once an extension exists, a different kernel measures the remaining freedom.}\label{KS:fig:extension}\end{figure}

\subsection{Existence is not uniqueness}
Define the ambiguity group
\[
 H_U=\ker\delta\cap\ker p_U.
\]
It consists of coherent changes that are invisible on the prescribed domain.

\begin{KStheorem}[Classification of a nonempty extension fiber]\label{KS:KS10}
If \(x_0\in\KSExt_U(s)\), translation gives an equivalence
\[
 H_U\simeq\KSExt_U(s),\qquad h\longmapsto x_0+h.
\]
Consequently the extension is unique exactly when \(\mathfrak o_U(s)=0\) and \(H_U=0\).
\end{KStheorem}
\begin{proof}
Define $\Theta:H_U\to\KSExt_U(s)$ by $\Theta(h)=x_0+h$. It lands in the extension set because
\[
 \delta(x_0+h)=\delta x_0+\delta h=0,
 \qquad p_U(x_0+h)=s+0=s.
\]
Define $\Psi:\KSExt_U(s)\to H_U$ by $\Psi(x)=x-x_0$. Its differential is zero, and its restriction is $s-s=0$, so the claimed target is correct. The formulas give
\[
 \Psi\Theta(h)=x_0+h-x_0=h,\qquad
 \Theta\Psi(x)=x_0+(x-x_0)=x.
\]
This proves the equivalence rather than only a comparison of cardinalities.

If $H_U=\{0\}$, the parameter of every extension is zero, so every extension equals $x_0$. Conversely, if $x_0$ is the only extension, then for any $h\in H_U$ the equality $x_0+h=x_0$ forces $h=0$ by cancellation. Combining this observation with the existence criterion proves the stated necessary and sufficient condition for a unique extension.

The choice of $x_0$ only supplies an origin for the parameters. If $x_1=x_0+h_1$ is used instead, the parameter of $x$ changes from $x-x_0$ to $x-x_1=(x-x_0)-h_1$. The affine set of answers itself has not changed.
\end{proof}
\KSsource{KS10}
An extension fiber is therefore an affine copy, or torsor, of the invisible coherent group. Selecting one solution gives coordinates on the fiber. Without a selected solution, the affine structure does not designate an origin.

\begin{KSexample}[Parity and an independent ambiguity]
For two views with \(f^*(m)=2m\) on \(\KSZ\), a prescribed value \(k\) at the coarse view extends exactly when \(k\) is even. The extension is \((k,k/2)\) and is unique. In particular \(1\) is obstructed and \(2\) extends nontrivially. The parity statement is an instance of Theorem~\ref{KS:KS09}. If a separate, disconnected integer view is added and left unprescribed, every extendable seed has infinitely many completions. This last variant is a direct application of Theorem~\ref{KS:KS10}, not a claim that the parity obstruction has changed.
\end{KSexample}
\KSsource{KS11}

\subsection{The exact sequence behind the criterion}
The additive structures make restriction a homomorphism
\(r_U:\Gamma(A)\to\Gamma_U(A)\). Using the zero extension, or lift independence, \(\mathfrak o_U\) is also a homomorphism.

\begin{KSproposition}[Exactness at partial sections]\label{KS:KS12}
The sequence
\[
 \Gamma(A)\xrightarrow{r_U}\Gamma_U(A)
 \xrightarrow{\mathfrak o_U}Q_U
\]
is exact at \(\Gamma_U(A)\). Thus \(\ker\mathfrak o_U=\KSim r_U\), and restriction is onto precisely when \(\mathfrak o_U\) is the zero map.
\end{KSproposition}
\begin{proof}
First verify that $\mathfrak o_U$ is a homomorphism, since exactness is being asserted in an additive setting. Given $s,t\in\Gamma_U(A)$, choose raw lifts $\widetilde s,\widetilde t$. Their sum lifts $s+t$. Lift independence and additivity imply
\[
 \mathfrak o_U(s+t)
 =[\delta(\widetilde s+\widetilde t)]
 =[\delta\widetilde s]+[\delta\widetilde t]
 =\mathfrak o_U(s)+\mathfrak o_U(t).
\]
The zero lift gives $\mathfrak o_U(0)=0$. Restriction $r_U$ is a homomorphism because addition of global sections is coordinatewise.

If $s=r_Ux$ for a global section $x$, use $x$ itself as a raw lift. Its differential vanishes, so $\mathfrak o_U(s)=0$. This proves $\KSim r_U\subseteq\ker\mathfrak o_U$. If $s\in\ker\mathfrak o_U$, the extension criterion constructs a global $x$ restricting to $s$, proving the reverse inclusion.

Finally, $r_U$ is surjective exactly when its image is all of $\Gamma_U(A)$. By the equality just proved, this is equivalent to every partial section having zero obstruction, which says that $\mathfrak o_U$ is the zero homomorphism. Nothing here claims that its image is all of $Q_U$.
\end{proof}
\KSsource{KS12}
No surjectivity onto the entire residual quotient is asserted. The residual group contains many entries that need not arise as the differential of a partial lift.

\subsection{Extension in stages}
For \(U\subseteq W\subseteq\KSOb\KSV\), an extension first to \(W\) and then globally must retain the intermediate choice. There is an explicit equivalence
\begin{equation}\label{KS:eq:stages}
 \KSExt_U(s)\simeq
 \coprod_{\substack{t\in\Gamma_W(A)\\t|_U=s}}\KSExt_W(t).
\end{equation}
It sends a global extension to its restriction to \(W\), together with itself as an extension of that restriction. The inverse forgets the intermediate coordinate. These operations are inverse because the required intermediate coordinate of a global section is its actual restriction. A coherent intermediate choice can still be obstructed, even when a different intermediate choice completes. Equation~\eqref{KS:eq:stages} is why extension by stages is a dependent sum rather than a statement that every branch succeeds.
\KSsource{KS13}

\chap{Protected correction as a general principle}\label{KS:sec:protected}
The additive extension theorem is a particular instance of a useful general construction. Let \(E,R,S\) be abelian groups, with additive maps \(d:E\to R\) and \(p:E\to S\). The residual \(d\) measures compatibility, and \(p\) specifies the information that repair must preserve.

\begin{KStheorem}[Semantics preserving correction]\label{KS:KS14}
For \(x\in E\), the following conditions are equivalent:
\[
\begin{aligned}
 &[dx]=0\text{ in }R/d(\ker p),\\
 &\exists h\in E:\ p(h)=0,\quad d(h)=d(x),\\
 &\exists y\in E:\ d(y)=0,\quad p(y)=p(x).
\end{aligned}
\]
The class of \(dx\) is unchanged when an admissible correction is subtracted from \(x\).
\end{KStheorem}
\begin{proof}
The image $d(\ker p)$ is a subgroup of $R$. The defining property of the quotient says
\[
 [dx]=0\quad\Longleftrightarrow\quad
 dx\in d(\ker p)
 \quad\Longleftrightarrow\quad
 \exists h\in E:\ ph=0,\ dh=dx.
\]
This is the first equivalence with its witness made explicit.

For the second, a permitted $h$ gives $y=x-h$. Then
\[
 dy=dx-dh=0,\qquad py=px-ph=px.
\]
Conversely an exact $y$ with $py=px$ gives $h=x-y$, which satisfies $ph=0$ and $dh=dx$. The two witness transformations are inverse: subtracting $h$ from $x$ and then subtracting that answer from $x$ recovers $h$, and conversely.

To verify invariance, fix any $h$ with $ph=0$. Then $d(x-h)=dx-dh$, and $dh$ lies in the subgroup divided out. Therefore $[d(x-h)]=[dx]$. This does not assert that the residual itself is unchanged. It asserts that permitted changes can alter the residual only within its already correctable coset.
\end{proof}
\KSsource{KS14}
Here \emph{semantics} means the declared mathematical observable \(p\). It is not an assertion that one homomorphism captures every meaning of every object. In a tower, \(p\) might preserve a seed operator and every new output component. In a graph, it might preserve designated vertex values. Different choices give different problems.

\begin{KSproposition}[A certificate of impossibility]\label{KS:KS15}
Let \(L:R\to D\) be additive. If \(L\) vanishes on \(d(\ker p)\) but \(L(dx)\ne0\), no exact \(y\) can preserve \(p(x)\).
\end{KSproposition}
\begin{proof}
Suppose, for contradiction, that $dy=0$ and $py=px$. Then $h=x-y$ satisfies $ph=0$ and $dh=dx$. By the annihilation hypothesis,
\[
 L(dx)=L(dh)=0,
\]
contrary to the assumed nonzero value. This already proves impossibility directly.

Equivalently, if $r-r'\in d(\ker p)$, then $L(r-r')=0$, so $L(r)=L(r')$. Hence $\overline L([r])=L(r)$ is a well-defined homomorphism on $R/d(\ker p)$. Its value on $[dx]$ is nonzero. This supplies the quotient interpretation of the same contradiction, without assuming that every nonzero quotient class can be detected by a map into one fixed preferred coefficient group.
\end{proof}
\KSsource{KS15}
This is often the most economical negative proof. One need not classify every possible correction. It is enough to find an invariant annihilating all admissible corrections but not the observed defect.

\subsection{Why unrestricted cancellation is empty information}
If every correction is permitted, taking \(h=x\) leaves the exact candidate zero. Thus \(\exists h:d(x-h)=0\) is automatic. It says nothing about recovering an object of interest. A repair theorem becomes informative only when it preserves something, is quantitatively controlled, or satisfies a comparison property that excludes arbitrary replacement.

\begin{KSexample}[A rigid obstruction]
Take \(E=R=S=\KSZ\), \(d=p=\KSid\), and \(x=1\). The only admissible correction is zero, while the residual is nonzero. Hence there is no exact answer preserving the observed value. This is not a failure of the formalism. It is the correct certificate for an incompatible demand.
\end{KSexample}

\subsection{What an obstruction does not yet provide}
An algebraic quotient does not automatically carry a useful norm. A zero infimum of distances to its boundary subgroup does not imply membership unless closure is controlled. Nor does vanishing supply a preferred correction when \(\ker d\cap\ker p\) is nontrivial. These two issues lead to different parts of the book: ambiguity is treated algebraically and through groupoids, while closed range and condition numbers enter the analytic theory.

\begin{insight}
A repair problem is not determined by its defect alone. It is determined by the defect \emph{and} the admissible changes. Keeping the pair \((d,p)\) explicit prevents a successful cancellation from being mistaken for a successful reconstruction.
\end{insight}



% BEGIN retained input sources/monographs/komorebi-space/source/chapters/02a_prescription.tex
\chap{Changing what is prescribed}\label{KS:sec:prescription}
The observable is part of a realization problem, not an annotation added after it has been solved. Prescribing more views can remove ambiguity, but it can also exclude every completion. This section follows both effects through the actual obstruction groups. It then works through a three-view model in which an initially harmless choice creates a genuine dead end.

\subsection{Three maps that must not be confused}
Keep the additive view system of Section~\ref{KS:sec:obstruction}. For a prescribed domain \(U\), there are three distinct maps:
\[
 p_U:E\longrightarrow\prod_{u\in U}A_u,\qquad
 r_U:\Gamma(A)\longrightarrow\Gamma_U(A),\qquad
 \mathfrak o_U:\Gamma_U(A)\longrightarrow Q_U.
\]
The first restricts a raw family, whether or not it is coherent. The second restricts a global coherent family. The third detects whether an already coherent partial family has a global completion. Their domains are different. In particular, surjectivity of the raw coordinate projection does not imply surjectivity of \(r_U\).

The obstruction group is \(Q_U=R/B_U\), where \(B_U=\delta(\ker p_U)\). The ambiguity group is \(H_U=\ker\delta\cap\ker p_U\). Both depend on what has been fixed. Enlarging the prescription reduces the permitted correction space and the invisible exact space at the same time.

\subsection{Forgetting a prescription gives a quotient map}
Let \(U\subseteq W\). Then
\[
 \ker p_W\subseteq\ker p_U,\qquad B_W\subseteq B_U,
 \qquad H_W\subseteq H_U.
\]
A correction invisible on all of \(W\) is certainly invisible on \(U\). The direction of the quotient map is consequently
\[
 \pi_{W,U}:Q_W\longrightarrow Q_U,\qquad
 [r]_{B_W}\longmapsto[r]_{B_U}.
\]
The larger prescribed domain appears in the source. More protection means a smaller subgroup is divided out, so more obstruction information can remain.

\begin{KStheorem}[Functorial forgetting of protected views]\label{KS:KS59}
The homomorphism \(\pi_{W,U}\) is well defined and surjective. For \(U\subseteq W\subseteq Z\), it satisfies
\[
 \pi_{U,U}=I,\qquad \pi_{W,U}\pi_{Z,W}=\pi_{Z,U}.
\]
For every \(s\in\Gamma_W(A)\),
\begin{equation}\label{KS:eq:prescribed-naturality}
 \pi_{W,U}\mathfrak o_W(s)=\mathfrak o_U(s|_U).
\end{equation}
In particular an obstructed restriction remains obstructed when more compatible prescribed data are retained. A successful extension on the larger prescribed domain gives a successful extension of its restriction.
\end{KStheorem}
\begin{proof}
If \(r-r'\in B_W\), then \(r-r'\in B_U\), proving well-definedness. Every class modulo \(B_U\) has a representative \(r\), and the class of that same representative modulo \(B_W\) maps to it. This proves surjectivity. Identity and composition follow by applying the formulas to a representative.

Choose a raw extension \(\widetilde s\) of the partial section on \(W\). It is also a raw extension of \(s|_U\). Lift independence of the two extension obstructions therefore identifies both sides of \eqref{KS:eq:prescribed-naturality} with the class of \(\delta\widetilde s\) modulo \(B_U\). A homomorphism sends zero to zero. Thus nonvanishing on the smaller domain rules out vanishing on the larger one, while vanishing on the larger domain implies vanishing after restriction.
\end{proof}
\KSsource{KS59}
The converse implication does not hold. The quotient map can annihilate the obstruction created by a newly fixed value. Surjectivity of \(\pi_{W,U}\) is also not a theorem that arbitrary larger partial sections extend. It is a statement about representatives in two obstruction groups, not about the image of restriction from global objects.

\subsection{An initially free parameter can become an obstruction}
Consider the chain of views \(0\to1\to2\), with fiber \(\KSZ^2\) at each view. Write the maps in the direction in which data are restricted:
\[
 \rho_{0,1}(a,b)=(a,0),\qquad
 \rho_{1,2}(a,b)=(a,2b),\qquad
 \rho_{0,2}(a,b)=(a,0).
\]
The composite identity \(\rho_{0,1}\rho_{1,2}=\rho_{0,2}\) is literal. A global family is exactly
\begin{equation}\label{KS:eq:branch-family}
 x_0=(k,0),\qquad x_1=(k,2n),\qquad x_2=(k,n),
 \qquad k,n\in\KSZ.
\end{equation}
Indeed the first coordinate must agree at every view. The second coordinate at view zero is forced to vanish, and the second coordinate at view one must be twice the second coordinate at view two. These conditions also suffice.

First prescribe only \(U=\{0\}\), with seed \((k,0)\). Every integer \(n\) in \eqref{KS:eq:branch-family} gives a completion. Now prescribe \(W=\{0,1\}\), with values \((k,0)\) and \((k,m)\). This is a coherent partial section for every integer \(m\), because \(\rho_{0,1}\) forgets its second coordinate. Global extendability is a stronger requirement: \(m\) must be even.

\begin{KSproposition}[A coherent branch can be a dead end]\label{KS:KS60}
Every initial seed \((k,0)\) in this model has a completion set in bijection with \(\KSZ\). A fixed intermediate value \((k,m)\) extends precisely when \(m\in2\KSZ\), and then has a unique completion. In particular the nonzero initial seed \((7,0)\) extends, while the coherent intermediate choice \((7,1)\) does not.
\end{KSproposition}
\begin{proof}
Equation~\eqref{KS:eq:branch-family} supplies a completion for each \(n\), and the second coordinate of \(x_2\) recovers the parameter. Fixing the intermediate value imposes \(2n=m\). This equation has no integer solution for odd \(m\), and exactly one for even \(m\). Taking \(k=7\) and \(m=1\) proves the final assertion.
\end{proof}
\KSsource{KS60}
The invisible coherent changes for the initial prescription are
\[
 ((0,0),(0,2n),(0,n)),\qquad n\in\KSZ.
\]
Thus \(H_U\simeq\KSZ\). Requiring invisibility also at view one forces \(2n=0\), hence \(n=0\), so \(H_W=0\). More observation removes all ambiguity in the successful branches. It does not turn the unsuccessful branches into solutions.

\begin{figure}[htbp]\centering\resizebox{.96\linewidth}{!}{
% BEGIN retained input sources/monographs/komorebi-space/source/figures/19_branches.tex
\includegraphics[width=15cm]{sources/monographs/komorebi-space/source/figures/generated/19_branches.pdf}

}\caption{Every displayed middle value is internally coherent with the initial seed seven. Only even middle values lift to the last view. The initial problem has many global answers, but choosing an odd middle value destroys extendability. The arrows show actual integer preimages, not a probabilistic search policy.}\label{KS:fig:branches}\end{figure}

\subsection{Why the extension space is a dependent sum}
Let \(s\) be the initial partial section. The stage theorem gives
\[
 \KSExt_U(s)\simeq
 \coprod_{\substack{t\in\Gamma_W(A)\\t|_U=s}}\KSExt_W(t).
\]
In the present example the index \(t\) is the integer \(m\). The summand is empty for odd \(m\) and a singleton for even \(m\). Taking their disjoint union gives one integer's worth of completions again. No contradiction appears between many initial completions and unique completion after a successful intermediate choice.

This formula also explains a logical error in greedy reconstruction. The fact that some branch extends does not justify replacing the existential choice by an arbitrary coherent branch. A method that fixes an intermediate value must either retain the remaining extension obligation or supply a rule proved to preserve extendability. Coherence on the current patch alone provides no such guarantee.

For an increasing sequence of prescribed domains, the same warning applies repeatedly. The finite stage formula keeps each intermediate choice explicit. It does not by itself manufacture one compatible infinite sequence of choices. Proposition~\ref{KS:KS58} supplies an existence criterion for towers of nonempty finite sets.

\subsection{Changing the coefficients changes both questions}
The equations behind the example can be read over a general abelian coefficient group \(T\). With the initial first coordinate fixed, the possible completions are parametrized by \(n\in T\). Prescribing the middle value \(m\) gives the equation \(2n=m\). Direct substitution therefore gives
\[
 \text{feasible middle values}=2T,\qquad
 \text{ambiguity after prescribing the middle}=\ker(2:T\to T).
\]
This calculation specializes as follows.
\begin{center}
\begin{tabular}{@{}lll@{}}
\toprule
Coefficient group & Which middle values extend? & Completions of an extendable value\\
\midrule
\(\KSZ\) & Even integers & One\\
\(\KSQ\) & Every rational & One\\
\(\mathbb F_2\) & Only zero & Two\\
\(\mathbb F_3\) & Every element & One\\
\bottomrule
\end{tabular}
\end{center}
The integer instance is the actual formal example cited above. The other rows are worked coefficient calculations from the same two equations. In characteristic two, stronger prescription does not eliminate the hidden parameter because doubling annihilates it. Over the rationals it creates no parity obstruction because doubling is invertible.

A solver that changes the coefficient domain while leaving its user-facing description untouched changes the mathematical problem. The rational answer \(1/2\) is not an integral completion of the odd middle value one. Conversely, a failure over \(\KSZ\) is not evidence of a failure over \(\KSQ\). The coefficient structure belongs beside the views, restrictions, and preserved values in the problem specification.

\begin{insight}
More information can reduce the set of answers without making the problem easier. The obstruction quotient records which new demands become impossible, while the ambiguity kernel records which freedom survives the demands that can be met.
\end{insight}



% BEGIN retained input sources/monographs/komorebi-space/source/chapters/03_towers.tex

\chapter{Resolution, repair, and retained detail}\label{KS:book:chapter-20}
Resolution gives a concrete family of views. Its compatibility differential connects exact maps, protected information, and residual coordinates. Recovering every finite coordinate is then separated from convergence in a chosen norm.

\section{Resolution towers and exact transport}\label{KS:sec:towers}
\subsection{The sequential model}
A tower in a category \(\mathcal C\) is a sequence
\[
 \cdots\xrightarrow{\rho_2}X_2\xrightarrow{\rho_1}X_1
 \xrightarrow{\rho_0}X_0.
\]
For sets, a point is a sequence \(x_n\in X_n\) with \(\rho_n(x_{n+1})=x_n\). Composite restrictions turn the tower into a view system on \(\KSN\), with an arrow \(i\to j\) when \(i\le j\). Its section space is the same coherent point type, with the same coordinates.

A candidate transport from \(X\) to \(Y\) consists of maps \(f_n:X_n\to Y_n\). It is exact when
\begin{equation}\label{KS:eq:exacttower}
 \rho^Y_n f_{n+1}=f_n\rho^X_n\quad\text{for all }n.
\end{equation}

\begin{KSproposition}[The exact tower category]\label{KS:KS16}
Identity maps and composites of exact transports are exact. An exact transport induces a map of coherent point spaces by \(x\mapsto(f_nx_n)_n\), preserving identities and composition.
\end{KSproposition}
\begin{proof}
For two exact transports, substitute \(\rho^Z_ng_{n+1}=g_n\rho^Y_n\) and \(\rho^Y_nf_{n+1}=f_n\rho^X_n\) into the composite square. For a coherent point,
\[
 \rho^Y_n f_{n+1}x_{n+1}
 =f_n\rho^X_nx_{n+1}=f_nx_n.
\]
Identity, associativity, and composition on points are checked coordinatewise.
\end{proof}
\KSsource{KS16}
This is a category on a fixed schedule. A full procategory allows additional reindexing and equivalence relations on maps. Those constructions must be supplied rather than identified with the fixed schedule by terminology.

\subsection{The actual residual of a transport}
Assume \(\mathcal C\) is preadditive, so differences of maps make sense. The tower differential is
\[
 (\partial f)_n=\rho^Y_nf_{n+1}-f_n\rho^X_n.
\]
Its target has components \(\KSHom(X_{n+1},Y_n)\), not \(\KSHom(X_n,Y_n)\).

\begin{KSproposition}[The transport Leibniz identity]\label{KS:KS17}
For composable candidates,
\[
 \partial(gf)_n=(\partial g)_nf_{n+1}+g_n(\partial f)_n.
\]
In particular exactness is equivalent to vanishing of the differential.
\end{KSproposition}
\begin{proof}
Expand the right side. The terms \(-g_n\rho^Y_nf_{n+1}\) and \(+g_n\rho^Y_nf_{n+1}\) cancel, leaving \(\rho^Z_ng_{n+1}f_{n+1}-g_nf_n\rho^X_n\). This is the differential of the composite. The exactness statement is \eqref{KS:eq:exacttower}.
\end{proof}
\KSsource{KS17}
There is no additional mixed defect term in this identity. Such a term would require a different definition of the defect or a different expansion. Typed cancellation is a useful safeguard against importing an attractive but incorrect error formula.

\subsection{Cofinal views preserve points, not every measurement}
Let \(\phi:\KSN\to\KSN\) be monotone and unbounded. The sampled tower uses levels \(X_{\phi(n)}\) and their composite restrictions.

\begin{KStheorem}[Cofinal reconstruction of points]\label{KS:KS18}
Sampling induces an equivalence between the coherent points of \(X\) and of its sampled tower.
\end{KStheorem}
\begin{proof}
Write $\rho_{i,j}:X_j\to X_i$ for the composite restriction when $i\le j$, with $\rho_{i,i}=I$. These maps satisfy $\rho_{i,j}\rho_{j,k}=\rho_{i,k}$. A coherent sampled point obeys
\[
 \rho_{\phi(n),\phi(m)}y_m=y_n\qquad(n\le m),
\]
by iterating its adjacent coherence equations.

For each $k$, unboundedness gives an index $n$ with $k\le\phi(n)$. Set $x_k=\rho_{k,\phi(n)}y_n$. If $m$ is another permissible index, choose $\ell\ge n,m$. Monotonicity permits both composite restrictions, and sampled coherence gives
\[
 \rho_{k,\phi(n)}y_n
 =\rho_{k,\phi(n)}\rho_{\phi(n),\phi(\ell)}y_\ell
 =\rho_{k,\phi(\ell)}y_\ell
 =\rho_{k,\phi(m)}y_m.
\]
Thus $x_k$ is independent of the choice, including when the schedule repeats a level.

For the adjacent equation at $k$, choose one $n$ with $k+1\le\phi(n)$. Then
\[
 \rho_{k,k+1}x_{k+1}
 =\rho_{k,k+1}\rho_{k+1,\phi(n)}y_n
 =\rho_{k,\phi(n)}y_n=x_k.
\]
The recovered family is therefore a point of the original tower.

Sampling a full point and recovering its $k$th coordinate gives $\rho_{k,\phi(n)}x_{\phi(n)}=x_k$ by its original coherence. Recovering a sampled point and evaluating at $k=\phi(n)$ allows the index $n$ itself, giving $\rho_{\phi(n),\phi(n)}y_n=y_n$. These are both inverse identities. No surjectivity of the bonding maps or of the point evaluations was used.
\end{proof}
\KSsource{KS18}
This theorem is about exact coherent points. Raw residual size is a different quantity. On identity scalar towers, \(f_n=(-1)^n\) has nonzero fine differences, but its even subsequence is constant. The coarse residual vanishes because adjacent fine defects cancel. A cofinal equivalence of limits does not imply equality of residual norms or an absence of aliasing in uncorrected data.

\begin{figure}[htbp]\centering\resizebox{.96\linewidth}{!}{
% BEGIN retained input sources/monographs/komorebi-space/source/figures/03_tower.tex
\begin{tikzpicture}[x=1cm,y=1cm]
\foreach \i/\x in {0/-4.8,1/-1.6,2/1.6,3/4.8}{
 \node[box,minimum width=1.5cm] (x\i) at (\x,2) {$X_{\i}$};
 \node[box,minimum width=1.5cm] (y\i) at (\x,0) {$Y_{\i}$};
 \draw[arr] (x\i) -- node[lab,right] {$f_{\i}$} (y\i);
}
\foreach \j/\i in {1/0,2/1,3/2}{
 \draw[arr] (x\j) -- node[lab,above] {$\rho^X_{\i}$} (x\i);
 \draw[arr] (y\j) -- node[lab,below] {$\rho^Y_{\i}$} (y\i);
}
\node[lab,text width=12.2cm] at (0,-1.5) {$\partial f_n=\rho^Y_nf_{n+1}-f_n\rho^X_n$\\[5pt] Exactness means every displayed square commutes.};
\end{tikzpicture}

}\caption{Exact transport commutes with every adjacent restriction. A cofinal subsequence retains coherent points, but a raw defect can cancel during coarsening. The retained coarse value alone need not record the omitted fine discrepancy.}\label{KS:fig:tower}\end{figure}


% BEGIN retained input sources/monographs/komorebi-space/source/chapters/03a_observation.tex
\chap{Finite observations and realized transformations}\label{KS:sec:observation}
The point-space construction has a converse question. Given a function between coherent realization spaces, can it be represented by compatible maps between finite views? The answer is not automatically yes. This section constructs the relevant category of reindexed maps and identifies its realized functions under an explicit observability hypothesis. Here a finite view means a view at a finite resolution level. Its set of possible values need not be finite.

\subsection{An output level may require a finer input level}
Let \(X\) and \(Y\) be sequential set-valued resolutions. Write \(\rho^X_{i,j}:X_j\to X_i\) for composite restriction. A reindexed transport consists of a monotone schedule \(\phi:\KSN\to\KSN\) and maps
\[
 f_n:X_{\phi(n)}\longrightarrow Y_n
\]
satisfying, for \(i\le j\),
\begin{equation}\label{KS:eq:scheduled-naturality}
 \rho^Y_{i,j}f_j=f_i\rho^X_{\phi(i),\phi(j)}.
\end{equation}
The schedule records how much input is used for each output level. Monotonicity means that asking for more output never forces the representation to discard input already available. It does not require strict growth, a bounded look-ahead, or an efficiency estimate.

The realized map is
\[
 |f|:\operatorname{Point}(X)\to\operatorname{Point}(Y),\qquad
 |f|(x)_n=f_n(x_{\phi(n)}).
\]
Equation~\eqref{KS:eq:scheduled-naturality} proves that this is a coherent output. A second transport \(g_n:Y_{\psi(n)}\to Z_n\) composes with \(f\) using schedule \(\phi\circ\psi\) and coordinate map \(g_n f_{\psi(n)}\). The order of the schedules is fixed by the types, not by the order in which their letters are written in prose.

\subsection{Equality is witnessed at finite stages}
Two different schedules can represent the same transformation. To compare \((\phi,f)\) and \((\psi,g)\), require that for every output level \(n\) there be a finite level \(k\ge\phi(n),\psi(n)\) with
\begin{equation}\label{KS:eq:eventual-equality}
 f_n\rho^X_{\phi(n),k}=g_n\rho^X_{\psi(n),k}
 \quad\text{as functions on all of }X_k.
\end{equation}
Call this eventual equality. The witnessing \(k\) may depend on \(n\). Equality is required on the whole finite-level fiber, not only on states that happen to extend globally.

\begin{KStheorem}[The eventual transport category]\label{KS:KS61}
Eventual equality is an equivalence relation and is preserved by composition. Its quotient therefore gives a category of resolutions and eventual transports. Realization descends to a functor into sets, preserving identities and composition.
\end{KStheorem}
\begin{proof}
A finite-stage equality remains true after moving to a later stage. Indeed, if it holds on $X_k$ and $\ell\ge k$, precompose both sides with $\rho^X_{k,\ell}$ and use the composite-restriction laws. Reflexivity uses $k=\phi(n)$, symmetry reverses the equality, and transitivity moves two witnesses to their maximum. The shared middle presentation then cancels as a literal function on that common fiber. This proves an equivalence relation without referring to global points.

We check compatibility with composition with all indices displayed. Write $f,f'$ with schedules $\phi,\phi'$ and $g,g'$ with schedules $\psi,\psi'$. Fix an output level $n$. Choose $j\ge\psi(n),\psi'(n)$ with
\[
 g_n\rho^Y_{\psi(n),j}=g'_n\rho^Y_{\psi'(n),j}.
\]
Next choose $k\ge\phi(j),\phi'(j)$ with
\[
 f_j\rho^X_{\phi(j),k}=f'_j\rho^X_{\phi'(j),k}.
\]
Monotonicity implies that $k$ is also at least both composite input levels. Applying the naturality equations and then the two witness equations gives
\[
\begin{aligned}
 g_nf_{\psi(n)}\rho^X_{\phi(\psi(n)),k}
 &=g_n\rho^Y_{\psi(n),j}f_j\rho^X_{\phi(j),k}\\
 &=g'_n\rho^Y_{\psi'(n),j}f_j\rho^X_{\phi(j),k}\\
 &=g'_n\rho^Y_{\psi'(n),j}f'_j\rho^X_{\phi'(j),k}\\
 &=g'_nf'_{\psi'(n)}\rho^X_{\phi'(\psi'(n)),k}.
\end{aligned}
\]
This is the required equality on every element of $X_k$. Thus composition descends to classes. The representative composition is associative because function composition and schedule composition are associative. The schedule $n\mapsto n$ with identity coordinate maps gives both identity laws. Hence the quotient has the claimed category structure.

Finally, evaluate a witness equation at the $k$th coordinate of a coherent point $x$. Its two restricted arguments become $x_{\phi(n)}$ and $x_{\psi(n)}$, proving equality of the realized outputs. Thus realization is well defined on the quotient. For a composite, direct evaluation gives
\[
 |g\circ f|(x)_n=g_n\bigl(f_{\psi(n)}(x_{\phi(\psi(n))})\bigr)
 =(|g|(|f|(x)))_n.
\]
The identity law is equally coordinatewise. Joint determination of output points proves the functor laws.
\end{proof}
\KSsource{KS61}
This is an actual quotient category in the formal source. No identification with the full procategory of an arbitrary base category is needed or claimed. The objects here are sequential set-valued resolutions, and the equivalence relation is the finite-stage relation just defined.

\subsection{When global behavior determines the finite presentation}
Say that \(X\) has surjective point projections when
\[
 \forall k\quad \operatorname{ev}_k:\operatorname{Point}(X)\to X_k
 \quad\text{is onto}.
\]
This means that every finite input state really occurs in some coherent global input. It is stronger than nonemptiness of the global space.

\begin{KStheorem}[Faithfulness on an observable source]\label{KS:KS62}
If \(X\) has surjective point projections, two reindexed maps out of \(X\) are eventually equal exactly when they induce the same realized function. Surjective adjacent bonding maps are a sufficient condition for surjective point projections.
\end{KStheorem}
\begin{proof}
Realized equality follows from eventual equality by the preceding theorem, so only its converse remains. Suppose \(|f|=|g|\). Fix the output level \(n\), and let \(k=\max\{\phi(n),\psi(n)\}\). To prove equality of two maps defined on \(X_k\), fix an arbitrary element \(a\in X_k\). The assumed surjectivity of \(\operatorname{ev}_k\) supplies a coherent point \(x\) with \(x_k=a\). Equality of the two realized maps gives equality of their output coordinates:
\[
 f_n(x_{\phi(n)})=g_n(x_{\psi(n)}).
\]
Coherence of \(x\) rewrites this as
\[
 f_n\bigl(\rho^X_{\phi(n),k}(a)\bigr)
 =g_n\bigl(\rho^X_{\psi(n),k}(a)\bigr).
\]
The index \(k\) was selected before \(a\), so this is one common finite-stage equality for \emph{all} elements of \(X_k\). It is exactly the witness required by eventual equality at output \(n\). Repeating for each \(n\) proves faithfulness.

For the sufficient condition on bonding maps, fix \(k\) and \(a\in X_k\). Define coordinates below \(k\) by restriction of \(a\). At each later stage use surjectivity of the adjacent bonding map to choose a preimage of the coordinate already chosen. Recursion with classical choice constructs the upper tail. Adjacent equations hold by construction, and induction on the interval length gives every composite equation. This produces a point evaluating to \(a\). If a fiber is empty, the statement that each of its elements lifts is vacuous, as required by surjectivity onto that fiber. Without the projection-surjectivity assumption, equality on realizable inputs alone need not identify finite maps on values that never occur in a global point.
\end{proof}
\KSsource{KS62}
The hypothesis is sufficient, not presented as a characterization of every source on which faithfulness happens to hold. Its role in the proof is precise: it converts a statement about globally realized inputs into a statement about all finite inputs.

Without it, global behavior can lose information. For example, use \(X_n=\{m\in\KSN:m\ge n\}\) with inclusion bonds and a constant Boolean output tower. The source point space is empty. The two transports that always output false and always output true induce the same function from the empty set, but disagree at every finite stage. Every \(X_k\) is nonempty, so no eventual equality witness exists. This calculation uses the finite-prefix boundary of Example~\ref{KS:KS54} and does not identify empty global behavior with equality of finite presentations.

\subsection{The exact image of realization}
For a function \(F:\operatorname{Point}(X)\to\operatorname{Point}(Y)\), define finite observation by the quantifiers
\begin{equation}\label{KS:eq:finite-observation}
 \forall n\ \exists k\ \exists a:X_k\to Y_n\ \forall x:\quad
 F(x)_n=a(x_k).
\end{equation}
The input level is uniform over all global inputs for each fixed output level. It may vary with the output level. This is not the weaker condition that each individual input can be inspected to some input-dependent finite depth.

\begin{KStheorem}[Characterization of realized transformations]\label{KS:KS63}
If \(X\) has surjective point projections, realization induces a bijection
\[
 \operatorname{EventualHom}(X,Y)
 \simeq
 \{F:\operatorname{Point}(X)\to\operatorname{Point}(Y):F
 \text{ is finitely observed}\}.
\]
\end{KStheorem}
\begin{proof}
For a reindexed map $(\phi,f)$, the formula $|f|(x)_n=f_n(x_{\phi(n)})$ immediately supplies a finite-observation witness for each $n$. Eventual equality does not affect the realized function by Theorem~\ref{KS:KS61}, so this defines a map from the quotient to the stated set of functions.

For surjectivity, let $F$ be finitely observed. Choose one pair $(k_n,a_n)$ at each output level, satisfying $F(x)_n=a_n(x_{k_n})$ for every coherent input. Define
\[
 \phi(n)=\max_{0\le i\le n}k_i,
 \qquad f_n=a_n\rho^X_{k_n,\phi(n)}.
\]
The maximum is of a nonempty finite set. It gives a monotone schedule with $k_n\le\phi(n)$. For every point $x$, coherence in $X$ gives
\[
 f_n(x_{\phi(n)})
 =a_n\bigl(\rho^X_{k_n,\phi(n)}x_{\phi(n)}\bigr)
 =a_n(x_{k_n})=F(x)_n.                                      \tag{*}
\]

The remaining obligation is naturality on \emph{all} finite states. Fix $i\le j$ and $u\in X_{\phi(j)}$. By surjectivity of that point projection, choose $x$ with $x_{\phi(j)}=u$. Equation (*) and output coherence yield
\[
\begin{aligned}
 \rho^Y_{i,j}f_j(u)&=\rho^Y_{i,j}F(x)_j=F(x)_i,\\
 f_i\rho^X_{\phi(i),\phi(j)}(u)&=f_i(x_{\phi(i)})=F(x)_i.
\end{aligned}
\]
The state $u$ was arbitrary. Hence the two finite-level functions are equal, which proves scheduled naturality. The resulting transport realizes $F$ by (*).

Different witness choices can give different schedules and coordinate maps, but the realized function is still $F$. By Theorem~\ref{KS:KS62} these transports are eventually equal, so their quotient class is independent of all choices. The same faithfulness result says that two classes with equal realized functions are identical. We have proved both surjectivity and injectivity, including the well-definedness of the inverse choice construction. No bound on the growth of $\phi$ follows from these quantifiers.
\end{proof}
\KSsource{KS63}
The finite-level compatibility is derived in this proof. It is not inserted as an additional field in the hypothesis that \(F\) is finitely observed. This is one of the useful bridges between a global description of a transformation and an actual presentation in the resolution calculus.

\begin{figure}[htbp]\centering\resizebox{.96\linewidth}{!}{
% BEGIN retained input sources/monographs/komorebi-space/source/figures/20_schedule.tex
\includegraphics[width=15cm]{sources/monographs/komorebi-space/source/figures/generated/20_schedule.pdf}

}\caption{Independent finite-observation witnesses need not use a monotone input schedule. Taking the running maximum produces one without changing the realized function. The map at the enlarged level first restricts to the originally sufficient level. This construction gives no efficiency bound on how large those levels are.}\label{KS:fig:schedule}\end{figure}

\subsection{Binary prefixes: a complete interpretation}
Let \(X_n\) consist of binary words of length \(n\), with prefix restriction. A coherent point is exactly an infinite binary sequence: its length-\(n\) view is its first \(n\) bits. Every prefix extends, for example by appending zeros. Hence the source satisfies the hypothesis of the characterization theorem.

Bitwise complementation is represented at level \(n\) by complementing the first \(n\) bits. The left shift \(F(x)_i=x_{i+1}\) uses the schedule \(\phi(n)=n+1\). A nearest-neighbor difference over \(\mathbb F_2\),
\[
 G(x)_i=x_i+x_{i+1},
\]
also uses \(n+1\) input bits for \(n\) output bits. Applying \(G\) twice uses \(n+2\) bits, and the middle terms cancel:
\[
 G(G(x))_i=x_i+x_{i+2}.
\]
The compatible finite maps make each statement meaningful on the entire coherent sequence, not merely on one chosen truncation. These are worked instances of the scheduled formulas. The finite regression suite checks their prefix compatibility independently.

Not every set function on these points is finitely observed. Consider the Boolean output that is one exactly when every input bit is zero, regarded as a point of the constant Boolean tower. Given any proposed finite level \(k\), compare the all-zero sequence with a sequence whose first one occurs at position \(k\). Their first \(k\) bits agree, but the outputs differ. Therefore no finite witness \(a:X_k\to\{0,1\}\) can compute even that one output coordinate. This rules out a reindexed presentation by the forward direction of the theorem.

The obstruction in this example is one of information, not a runtime benchmark. No amount of computation on the specified finite prefix distinguishes those two inputs. The theorem does not prohibit a different mathematical representation that supplies more information as part of each view. Changing the available view is changing the input model.

\subsection{Cofinal changes also preserve the morphism calculus}
A monotone unbounded schedule \(\theta\) gives the sampled resolution \(X_\theta\). There is a transport from \(X\) to \(X_\theta\) that reads the selected level. For the reverse direction choose, for each \(n\), a level \(j(n)\) with \(n\le\theta(j(n))\), and restrict the available value down to \(X_n\). Taking a least such index gives a monotone schedule.

\begin{KStheorem}[Cofinal invariance of eventual transports]\label{KS:KS64}
Sampling and recovery are inverse in the eventual transport category. Consequently, for monotone unbounded schedules \(\theta\) and \(\eta\), there is a bijection
\[
 \operatorname{EventualHom}(X,Y)
 \simeq
 \operatorname{EventualHom}(X_\theta,Y_\eta).
\]
On point spaces it agrees with the previously constructed cofinal realization equivalences.
\end{KStheorem}
\begin{proof}
Let \(s:X\to X_\theta\) be sampling, with schedule \(\theta\) and identity level maps. For recovery choose the least \(j(n)\) with \(n\le\theta(j(n))\). Unboundedness makes this set nonempty. Increasing \(n\) only strengthens the inequality, so \(j\) is monotone. The level map of \(r:X_\theta\to X\) is \(\rho^X_{n,\theta(j(n))}\). Composite restriction proves its naturality.

The composite \(rs:X\to X\) reads level \(\theta(j(n))\) and restricts to \(n\). At that very stage its map agrees with the identity transport after restriction from \(\theta(j(n))\) to \(n\). Thus \(rs\) is eventually equal to the identity.

For \(sr:X_\theta\to X_\theta\), the schedule at output \(n\) is \(j(\theta(n))\), and the level map is \(\rho^X_{\theta(n),\theta(j(\theta(n)))}\). Compare it with the identity at a sampled index
\[
 k\ge\max\{n,j(\theta(n))\}.
\]
Restricting either route from level \(k\) of the sampled system gives the same original restriction \(\rho^X_{\theta(n),\theta(k)}\). This proves eventual equality even when \(\theta\) has repeated values. Hence the two classes are inverse isomorphisms, without requiring them to be literal inverse schedules.

Write \(s_X,r_X\) and \(s_Y,r_Y\) for the sampling and recovery isomorphisms of source and target. The asserted Hom map is
\[
 [f]\longmapsto s_Y[f]r_X,
\]
with inverse \([g]\mapsto r_Y[g]s_X\). Associativity and the two inverse equations cancel the inserted pairs, proving both round trips. Applying the realization functor transports these equations to point spaces, where sampling and recovery are the cofinal point equivalences already constructed. No norm invariance or surjectivity of source point projections is needed for this categorical conjugation argument.
\end{proof}
\KSsource{KS64}
This strengthens cofinal invariance from points to a specified category of transformations. It still does not make raw residual norms invariant, identify every global set function with a finite presentation, or preserve the numerical conditioning of a repair. Those are separate questions, addressed by the defect maps in the next sections.

\Needspace{7\baselineskip}
\begin{insight}
For a projection-surjective source, eventual finite presentations correspond exactly to uniformly finite observations at each output level. This is a precise machine-oriented representation theorem. It is not a claim that every represented map is computationally efficient, or that every arbitrary function of an infinite input has such a presentation.
\end{insight}



\chap{Banach repair and constructed split coordinates}\label{KS:sec:split}
\subsection{An actual complete candidate space}
Let \(X_n,Y_n\) be Banach spaces over a complete nontrivially normed field, with continuous linear bonding maps and uniform bounds \(\KSnorm{\rho^X_n}\le B_X\), \(\KSnorm{\rho^Y_n}\le B_Y\). Use the spaces of bounded operator families
\[
 E=\ell^\infty_n\mathcal L(X_n,Y_n),\qquad
 R=\ell^\infty_n\mathcal L(X_{n+1},Y_n).
\]
The source formalization derives the required completeness from the operator spaces and their bounded product. The same differential satisfies
\[
 \KSnorm{\partial f}_\infty\le(B_X+B_Y)\KSnorm f_\infty.
\]
Thus it is a continuous linear map, and evaluation at a fixed level is continuous. These facts connect algebraic compatibility to an analytic limit argument.

\subsection{From a correction operator to a projection}
The next statement is general. Let \(E\) be Banach, let \(d:E\to R\), \(p:E\to S\), and \(A:R\to E\) be continuous linear maps. Suppose \(pA=0\) and, for \(0\le q<1\),
\begin{equation}\label{KS:eq:contraction}
 \KSnorm{df-dA(df)}\le q\KSnorm{df}\qquad(f\in E).
\end{equation}
The hypothesis concerns residuals of actual candidates. It need not hold for every independent vector in \(R\).

\begin{KStheorem}[Constructed geometric repair]\label{KS:KS19}
Starting from \(f^{(0)}=f\), define
\[
 f^{(m+1)}=f^{(m)}-A(df^{(m)}).
\]
There is a bounded linear projection \(P_A:E\to\ker d\) with
\[
 f^{(m)}\longrightarrow P_Af,\qquad p(P_Af)=p(f),
\]
and
\[
 \KSnorm{f^{(m)}-P_Af}
 \le\frac{\KSnorm A}{1-q}\KSnorm{df^{(m)}}.
\]
It fixes every exact candidate and satisfies
\[
 \KSnorm{P_A}\le1+\frac{\KSnorm A\KSnorm d}{1-q}.
\]
\end{KStheorem}
\begin{proof}
Put $r_m=df^{(m)}$. Additivity of $d$ gives
\[
 r_{m+1}=r_m-dAr_m.
\]
Each $r_m$ is the residual of an actual candidate, so the stated hypothesis applies at every iteration. Induction yields, for $j\ge0$,
\[
 \KSnorm{r_{m+j}}\le q^j\KSnorm{r_m},\qquad
 \KSnorm{f^{(m+j+1)}-f^{(m+j)}}
 \le\KSnorm A\,q^j\KSnorm{r_m}.
\]
For $n>m$, the finite triangle inequality therefore gives
\[
 \KSnorm{f^{(n)}-f^{(m)}}
 \le\KSnorm A\,\KSnorm{r_m}\sum_{j=0}^{n-m-1}q^j
 \le\frac{\KSnorm A}{1-q}\KSnorm{r_m}.
\]
Since $\KSnorm{r_m}\le q^m\KSnorm{df}$ tends to zero, this is a Cauchy estimate. Completeness of $E$ supplies a limit, denoted $P_Af$. Passing $n\to\infty$ in the same norm inequality proves the required tail estimate.

Continuity of $d$ implies $d(P_Af)=\lim_m r_m=0$. Protection holds before taking a limit: $pA=0$ gives $pf^{(m+1)}=pf^{(m)}$, hence $pf^{(m)}=pf$ for every $m$. Continuity of $p$ gives $p(P_Af)=pf$. This explains why continuity of the observable belongs to this analytic statement.

For linearity, write $L=I-Ad$, so $f^{(m)}=L^mf$. For fixed $f,g$ and scalars $\lambda,\mu$,
\[
 L^m(\lambda f+\mu g)=\lambda L^mf+\mu L^mg.
\]
Continuity of addition and scalar multiplication allows passage to the norm limit and proves linearity of $P_A$. The estimate with $m=0$ gives
\[
 \KSnorm{P_Af}\le\KSnorm f+\frac{\KSnorm A}{1-q}\KSnorm{df}
 \le\left(1+\frac{\KSnorm A\KSnorm d}{1-q}\right)\KSnorm f,
\]
so the linear map is bounded. If $df=0$, every iterate equals $f$, hence $P_Af=f$. Its image is contained in $\ker d$ and it fixes that kernel, proving surjectivity onto its stated image and $P_A^2=P_A$. For $q=0$ the residual vanishes after the first update and the same inequalities remain valid.
\end{proof}
\KSsource{KS19}
This theorem constructs the dynamics from \(A\). It does not construct a suitable \(A\) for every residual problem. That missing input can, however, be derived for a useful class of actual towers.

\subsection{A right inverse from target splittings}
Suppose each target bonding map has a continuous linear section
\[
 s_n:Y_n\to Y_{n+1},\qquad \rho^Y_ns_n=I,\qquad \KSnorm{s_n}\le S.
\]
Assume the propagation gap
\begin{equation}\label{KS:eq:splitgap}
 SB_X<1,\qquad K=\frac{S}{1-SB_X}.
\end{equation}
For \(r\in R\), recursively set
\begin{equation}\label{KS:eq:greenrecurrence}
 (Ar)_0=0,\qquad
 (Ar)_{n+1}=s_n\bigl(r_n+(Ar)_n\rho^X_n\bigr).
\end{equation}

\begin{KStheorem}[Split tower correction]\label{KS:KS20}
The recurrence defines a continuous linear map \(A:R\to E\) with
\[
 \partial A=I_R,\qquad \KSnorm A\le K.
\]
The operator \(P=I-A\partial\) is an exact projection and
\[
 (Pf)_0=f_0,\qquad
 \KSnorm{f-Pf}_\infty\le K\KSnorm{\partial f}_\infty.
\]
\end{KStheorem}
\begin{proof}
Let $\theta=SB_X<1$. The number $K=S/(1-\theta)$ is nonnegative and satisfies
\[
 S+SB_XK=K.
\]
Fix $r\in R$ and define its coordinates recursively as in \eqref{KS:eq:greenrecurrence}. The initial coordinate has norm zero. If $\KSnorm{(Ar)_n}\le K\KSnorm r_\infty$, submultiplicativity of operator norms gives
\[
 \begin{aligned}
 \KSnorm{(Ar)_{n+1}}
 &\le\KSnorm{s_n}\bigl(\KSnorm{r_n}+\KSnorm{(Ar)_n}\KSnorm{\rho^X_n}\bigr)\\
 &\le S(1+KB_X)\KSnorm r_\infty=K\KSnorm r_\infty.
 \end{aligned}
\]
Induction proves the estimate at every coordinate, so the recursively defined family actually belongs to the bounded product $E$.

The coordinate recurrence preserves addition and scalar multiplication, starting from the zero initial coordinate. Induction therefore proves these laws for the family map $A$. Its uniform estimate proves $\KSnorm A\le K$, giving a continuous linear map with the claimed domain and codomain.

For every $n$, apply the target restriction to the recurrence. The equation $\rho^Y_ns_n=I$ gives
\[
 \rho^Y_n(Ar)_{n+1}=r_n+(Ar)_n\rho^X_n,
 \qquad (\partial Ar)_n=r_n.
\]
Coordinate equality yields the operator identity $\partial A=I_R$. In particular this identity holds for every bounded residual, not just residuals already known to lie in the image.

Set $P=I-A\partial$. Then
\[
 \partial P=\partial-\partial A\partial=0,\qquad
 P^2=I-2A\partial+A(\partial A)\partial=I-A\partial=P.
\]
The expression $2A\partial$ abbreviates addition of the same operator twice and uses no division by two. Every exact $f$ is fixed because $\partial f=0$. Since $(Ar)_0=0$, the seed of $Pf$ is $f_0$. Finally $f-Pf=A\partial f$, and the bound on $A$ gives the asserted residual estimate. The denominators above are justified exactly by the strict propagation gap.
\end{proof}
\KSsource{KS20}
The surjectivity is onto all bounded residual families. It is stronger than the assertion that some exact output can be assigned to each candidate. The gap and the chosen splittings are genuine inputs.

\subsection{Protected information beyond the seed}
When the target grows, its new kernel directions are independent information. Define
\[
 N_n(f)=(I-s_n\rho^Y_n)f_{n+1},\qquad
 (p_Lf)_0=f_0,\quad(p_Lf)_{n+1}=N_n(f).
\]
Each \(N_n(f)\) lands in \(\ker\rho^Y_n\). The observable \(p_L\) retains both the seed and every new output component, relative to the chosen sections.

\begin{KStheorem}[Split coordinate theorem]\label{KS:KS21}
The observable is a bounded projection with \(p_LA=0\) and \(p_LP=p_L\). It separates exact candidates. If
\[
 \mathcal S_L=\{s\in E:p_Ls=s\},
\]
then
\[
 \Phi:E\simeq_L\mathcal S_L\times R,
 \qquad \Phi(f)=(p_Lf,\partial f),
 \qquad \Phi^{-1}(s,r)=Ps+Ar.
\]
Repair becomes \((s,r)\mapsto(s,0)\). It is the unique exact repair preserving \(p_Lf\).
\end{KStheorem}
\begin{proof}
Put \(B_n=I-s_n\rho^Y_n\). Since \(\rho^Y_ns_n=I\), direct multiplication gives
\[
 B_n^2=B_n,\qquad \rho^Y_nB_n=0,\qquad B_ns_n=0.
\]
Moreover, \(B_ny=y\) whenever \(\rho^Y_ny=0\). Thus \(B_n\) really projects onto the target kernel. Its norm is at most \(1+SB_Y\). Coordinate zero of \(p_L\) is the identity, and every later coordinate applies \(B_n\), so \(p_L^2=p_L\) and \(\KSnorm{p_L}\le1+SB_Y\). In particular the formula defines a map on the bounded product, not just on arbitrary coordinate families.

The recurrence for \(Ar\) gives zero seed and places \((Ar)_{n+1}\) in the image of \(s_n\). The identities just proved make all its innovations zero. Therefore \(p_LA=0\), and
\[
 p_LP=p_L(I-A\partial)=p_L.
\]
We next prove separation on the exact space. If \(f\) is exact, its restriction equation rewrites as
\[
 f_{n+1}=s_n\rho^Y_nf_{n+1}+B_nf_{n+1}
          =s_nf_n\rho^X_n+N_n(f).
\]
Two exact families with the same seed and innovations agree first at level zero and then at every later level by this recurrence. More generally, if two arbitrary families have the same observable and differential, their difference is exact with zero observable, so separation on the exact space makes their difference zero.

The proposed chart has the claimed target because \(p_L(p_Lf)=p_Lf\). Given \(s\in\mathcal S_L\) and \(r\in R\), set \(f=Ps+Ar\). The preceding identities and \(\partial A=I\) give
\[
 p_Lf=p_Ls=s,\qquad \partial f=r.
\]
This proves one round trip. For the other, \(f\) and \(P(p_Lf)+A\partial f\) have the same observable and residual, so the joint separation just proved identifies them. All maps are linear and bounded on their indicated subspaces and finite products, establishing a continuous linear equivalence.

Finally, \(Pf\) has coordinates \((p_Lf,0)\). Any exact answer with the same observable has those same two coordinates, hence equals \(Pf\). This uniqueness is relative to the complete observable of seed and innovations. It is not obtained by assuming that the seed alone determines an exact growing-target family.
\end{proof}
\KSsource{KS21}
This is an actual analytic coordinate equivalence, not generally an isometry. The protected coordinate depends on the declared splitting. It should not be called an intrinsic semantic coordinate without a further comparison theorem.

\begin{figure}[htbp]\centering\resizebox{.96\linewidth}{!}{
% BEGIN retained input sources/monographs/komorebi-space/source/figures/04_split.tex
\begin{tikzpicture}[x=1cm,y=1cm]
\node[box,minimum width=3.3cm,minimum height=1.8cm] (raw) at (-4.5,1.2) {candidate $f\in E$\\[6pt] actual tower operators};
\draw[arr] (raw.east) -- node[lab,above] {$\Phi$} (-1.1,1.2);
\begin{scope}[shift={(.1,-.4)}]
 \draw[thinarr] (-.5,0)--(5,0) node[right,lab] {$s$};
 \draw[thinarr] (0,-.5)--(0,3.8) node[above,lab] {$r$};
 \draw[KomorebiPine,line width=2pt] (0,0)--(4.7,0);
 \node[lab,below] at (2.7,-.15) {exact slice $r=0$};
 \node[dot] (c) at (3.2,2.7) {};
 \node[lab,right] at (c.east) {$(p_Lf,\partial f)$};
 \node[dot] (p) at (3.2,0) {};
 \draw[arr,dashed] (c)--node[lab,left] {repair} (p);
 \node[lab,above left] at (p) {$(p_Lf,0)$};
\end{scope}
\node[lab,text width=12cm] at (0,-1.8) {$\Phi^{-1}(s,r)=Ps+Ar$\qquad The protected coordinate does not move.};
\end{tikzpicture}

}\caption{In the split coordinate chart, a candidate separates into retained information and compatibility residual. The exact subspace is the zero residual slice. The chart and its protected coordinate depend on the chosen target sections.}\label{KS:fig:split}\end{figure}

\subsection{A growing example and the gap boundary}
Use \(X_n=Y_n=\KSR^{n+1}\) with maximum norms. Let \(\pi_n\) delete the last coordinate, take \(\rho^X_n=\tfrac12\pi_n\), \(\rho^Y_n=\pi_n\), and let \(s_n\) append zero. Then \(S=B_Y=1\), \(B_X=1/2\), and \(K=2\).

\begin{KSexample}[Repair of the identity family]\label{KS:KS22}
The bounded candidate \(f_n=I\) is not exact. Its repair is
\[
 (Pf)_n=\operatorname{diag}(2^{-n},2^{-(n-1)},\ldots,2^{-1},1).
\]
The new last output row is retained at every stage while inherited rows are propagated by the factor \(1/2\).
\end{KSexample}
\begin{proof}
At level zero the value is one. Exactness forces every inherited row at the next level to be one half of its previous value. Preservation of the innovation forces the new row to remain the corresponding identity row. Induction gives the displayed diagonal.
\end{proof}
\KSsource{KS22}
For constant real scalar identity towers, by contrast, \((\partial f)_n=f_{n+1}-f_n\). The constant residual \(r_n=1\) has no bounded primitive, since any primitive satisfies \(f_n=f_0+n\). Here \(SB_X=1\). This counterexample excludes discarding the gap indiscriminately. It does not prove that the gap is necessary for every individual tower. Exact repair preserving the initial coordinate still exists in this constant example, by replacing a candidate with its initial constant family. Surjectivity onto all bounded residuals does not follow from that repair.
\KSsource{KS23}

\subsection{What is canonical}
For fixed splitting and protected observable, the repair above is unique. Changing the splitting can change both the observable and the repaired answer. In the growing example, replace \(s_n(y)=(y,0)\) by \(s_n'(y)=(y,y_0)\). At the first positive level, repair of the identity becomes
\[
 \begin{pmatrix}1/2&0\\-1/2&1\end{pmatrix}
 \quad\text{instead of}\quad
 \begin{pmatrix}1/2&0\\0&1\end{pmatrix}.
\]
Both are exact continuations for their respective innovation profiles and have the same seed. Thus a unique answer for a declared preservation problem is not an answer independent of all preservation choices.
\KSsource{KS24}


% BEGIN retained input sources/monographs/komorebi-space/source/chapters/03b_repair_transport.tex
\chap{Repair under a change of representation}\label{KS:sec:repair-transport}
An exact transformation should carry coherent inputs to coherent outputs. It need not carry the selected repair of an incoherent input to the selected repair of its transformed image. The distinction has an exact formula. In split coordinates, a compatible change of representation is generally block triangular rather than block diagonal. Its off-diagonal term records precisely the discrepancy between the two repair routes.

\subsection{A section is a choice of residual coordinates}
Let \(d:E\to R\) be a continuous linear map between normed spaces over a nontrivially normed field. Suppose it has a continuous linear right inverse \(A:R\to E\), so \(dA=I_R\). Set
\[
 K=\ker d,\qquad P_A=I_E-Ad.
\]
Then \(dP_A=0\), \(P_A\) fixes \(K\), and \(P_A^2=P_A\). The right inverse chooses where each residual is represented inside the candidate space. Different right inverses can choose different representatives of the same residual.

\begin{KStheorem}[The split defect chart]\label{KS:KS65}
The maps
\[
 \chi_A:E\simeq_L K\times R,\qquad
 \chi_A(x)=(P_Ax,dx),\qquad
 \chi_A^{-1}(u,r)=u+Ar
\]
are inverse continuous linear equivalences. In these coordinates, \(P_A\) sends \((u,r)\) to \((u,0)\). If \(B\) is another right inverse, then \(A-B\) takes values in \(K\), and
\[
 P_Bx-P_Ax=(A-B)dx.
\]
The projections commute if and only if they are equal, and they are equal if and only if \(A=B\).
\end{KStheorem}
\begin{proof}
The identity \(dA=I\) gives \(dP_A=d-dAd=0\), so the proposed first chart coordinate actually belongs to \(K\). If \(u\in K\), then \(du=0\) and \(P_Au=u\). Therefore, for \((u,r)\in K\times R\),
\[
 d(u+Ar)=r,\qquad P_A(u+Ar)=u+Ar-A(r)=u.
\]
In the other direction, \(P_Ax+A(dx)=x-Adx+Adx=x\). These are both inverse equations, including the membership condition for the kernel coordinate. All component maps are bounded linear, so they define a continuous linear equivalence on the standard finite product. Applied to \(P_Ax\), the chart gives \((P_Ax,0)\), proving the repair formula.

For a second right inverse \(B\), the equation \(d(A-B)=I-I=0\) places its difference from \(A\) in \(K\). Direct subtraction gives
\[
 P_B-P_A=(I-Bd)-(I-Ad)=(A-B)d.
\]
Both projections have range \(K\) and fix every point of \(K\). Thus applying \(P_A\) after \(P_B\) gives \(P_B\), while applying \(P_B\) after \(P_A\) gives \(P_A\). These two compositions are equal if and only if the projections are equal.

If \(P_A=P_B\), evaluate their difference at \(Ar\). Since \(dAr=r\), the difference formula yields \((A-B)r=0\) for every \(r\), hence \(A=B\). Equality of sections conversely makes their projection formulas identical. The argument uses the right inverse to test every residual. It does not infer equality of sections merely from agreement on some restricted set of observed residuals.
\end{proof}
\KSsource{KS65}
This is a linear splitting, not necessarily an orthogonal one. The exact coordinate \(u\) depends on the chosen section. The residual coordinate \(r=dx\) does not. A specified Hilbert geometry can make one section distinguished through orthogonality, but the algebraic split alone does not do so.

The tower construction is stronger in a different direction. Its recurrence builds a right inverse from the target splittings, and its observable identifies what that repair preserves. The abstract chart here isolates the mechanism shared by those concrete constructions.

\subsection{A change of representation has two components}
Let \(D:F\to T\) be a second continuous linear defect map. A compatible change consists of bounded linear maps
\[
 S:E\to F,\qquad C:R\to T,\qquad DS=Cd.
\]
The equation is the commuting defect square. It says that transforming a candidate and then measuring its defect agrees with transforming its actual defect. No sections have yet been used.

In particular, \(S\) carries \(\ker d\) into \(\ker D\). It may still have a kernel, lose information, or fail to be onto. Calling it a change of representation does not assert an equivalence unless inverse maps have also been supplied.

Choose sections \(A\) of \(d\) and \(B\) of \(D\). Define
\begin{equation}\label{KS:eq:repair-mismatch}
 M=SA-BC:R\longrightarrow F.
\end{equation}
The commuting square and section identities give \(DM=0\). Thus this mismatch takes residual information into an exact direction of the target.

\begin{KStheorem}[Triangular transport and repair mismatch]\label{KS:KS66}
In split defect coordinates, the change of representation is
\begin{equation}\label{KS:eq:triangular-chart}
 \chi_BS\chi_A^{-1}(u,r)=(Su+Mr,Cr).
\end{equation}
The two repair routes differ by
\begin{equation}\label{KS:eq:repair-change}
 P_BSx-SP_Ax=M(dx).
\end{equation}
Consequently they agree on exact inputs, and
\[
 \KSnorm{P_BSx-SP_Ax}\le\KSnorm M\KSnorm{dx},
 \qquad \KSnorm M\le\KSnorm S\KSnorm A+\KSnorm B\KSnorm C.
\]
They agree on all inputs if and only if \(SA=BC\).

For two consecutive compatible changes \((S_1,C_1)\), \((S_2,C_2)\), with compatible choices of sections at the three ends, their mismatches satisfy
\[
 M_{21}=S_2M_1+M_2C_1.
\]
If \(dx_j\to0\) in norm for a fixed change and fixed sections, the repair discrepancy tends to zero.
\end{KStheorem}
\begin{proof}
First \(DS=Cd\) sends \(\ker d\) into \(\ker D\), and
\[
 DM=DSA-DBC=CdA-C=C-C=0.
\]
Thus both terms in the proposed exact coordinate have the correct codomain. Given \(u\in\ker d\), compute the residual and exact coordinates of \(S(u+Ar)\) separately:
\[
 DS(u+Ar)=Cr,\qquad
 S(u+Ar)-BCr=Su+(SA-BC)r.
\]
This proves the triangular formula as an equality in \(\ker D\times T\).

For a general input, use the same commuting square to obtain
\[
 \begin{aligned}
 P_BSx-SP_Ax
 &=Sx-BDSx-Sx+SAdx\\
 &=(SA-BC)dx=M(dx).
 \end{aligned}
\]
Its norm is at most \(\KSnorm M\KSnorm{dx}\). The triangle inequality and submultiplicativity give \(\KSnorm M\le\KSnorm S\KSnorm A+\KSnorm B\KSnorm C\). Exact inputs have \(dx=0\), so the two routes agree there regardless of the section choices. If they agree on \emph{every} input, substitute \(x=Ar\). Then \(dAr=r\) makes \(Mr=0\) for each residual, hence \(SA=BC\). That equality conversely makes the discrepancy vanish everywhere.

For composition, denote the three sections by \(A,B,A_2\). The composite residual map is \(C_2C_1\), so its mismatch is \(S_2S_1A-A_2C_2C_1\). Insert and subtract \(S_2BC_1\):
\[
 \begin{aligned}
 S_2S_1A-A_2C_2C_1
 &=S_2(S_1A-BC_1)+(S_2B-A_2C_2)C_1\\
 &=S_2M_1+M_2C_1.
 \end{aligned}
\]
All terms map the first residual space into the last candidate space, in the written order. Finally, for fixed bounded \(M\), the estimate \(\KSnorm{M(dx_j)}\le\KSnorm M\KSnorm{dx_j}\) proves convergence of the discrepancy when \(dx_j\to0\). Allowing the section or representation norms to grow with \(j\) would require a separate uniform bound.
\end{proof}
\KSsource{KS66}
The off-diagonal term in \eqref{KS:eq:triangular-chart} has a specific meaning. A raw residual can become part of the target's chosen exact coordinate after a change of presentation. Deleting the residual before changing presentation removes that contribution. Deleting it afterwards does not. This is not a failure of exactness preservation. Both routes end in the target exact space, but they can end at different points in it.

\subsection{A two-coordinate model in which the routes differ}
Use \(E=F=\KSR^2\), write candidates as \((u,r)\), and take
\[
 d(u,r)=D(u,r)=r,\qquad A(r)=B(r)=(0,r).
\]
The exact space is the horizontal axis and both repairs are \(P(u,r)=(u,0)\). For a real parameter \(c\), set
\[
 S_c(u,r)=(u+cr,r),\qquad C=I.
\]
The defect square commutes because the second coordinate is unchanged. The mismatch is \(M_c(r)=(cr,0)\). Hence
\[
 PS_c(u,r)=(u+cr,0),\qquad S_cP(u,r)=(u,0).
\]
With \(c=2\) and input \((0,1)\), one route gives \((2,0)\) and the other gives \((0,0)\). Both have zero residual. A report containing only their final residuals would miss the entire difference.

\begin{figure}[htbp]\centering\resizebox{.96\linewidth}{!}{
% BEGIN retained input sources/monographs/komorebi-space/source/figures/21_mismatch.tex
\includegraphics[width=15cm]{sources/monographs/komorebi-space/source/figures/generated/21_mismatch.pdf}

}\caption{The shear preserves the residual coordinate but mixes it into the chosen exact coordinate. Repair after the shear reaches \((2,0)\). Repair before the shear reaches \((0,0)\). Their difference is the exact-valued mismatch applied to the original residual. The example uses the ordinary Euclidean drawing, while the algebraic identity does not depend on that drawing.}\label{KS:fig:mismatch}\end{figure}

Two shears compose by adding their parameters. In the mismatch law this appears as \(M_{c_2+c_1}=S_{c_2}M_{c_1}+M_{c_2}\), and \(S_{c_2}\) fixes the exact horizontal axis. The identity is therefore a simple instance of the same composition rule, not an unrelated coordinate trick.

\subsection{Actual tower resampling gives a defect map}
Return to Banach towers with contractive source and target bonds. Let \(\phi\) be a monotone schedule. Candidate sampling is
\[
 (S_\phi f)_n=f_{\phi(n)}.
\]
The sampled bonding map is the actual composite restriction between the selected levels. The residual map \(C_\phi\) is not pointwise sampling of a scalar error score. It accumulates transported fine defects across each omitted block.

For a block from \(a\) to \(b\), its value on a residual family \(r\) is
\begin{equation}\label{KS:eq:block-defect-sum}
 (C_{a,b}r)=\sum_{k=a}^{b-1}
 \rho^Y_{a,k}\,r_k\,\rho^X_{k+1,b}.
\end{equation}
For the schedule \(\phi\), set \((C_\phi r)_n=C_{\phi(n),\phi(n+1)}r\). Every summand maps \(X_b\) into \(Y_a\). These surrounding restrictions are necessary when the spaces change with the level. For an actual differential, the sum telescopes to
\[
 C_{a,b}(\partial f)=\rho^Y_{a,b}f_b-f_a\rho^X_{a,b}.
\]
The dyadic coarsener in Section~\ref{KS:sec:cascade} is the case \(b=a+2\).

\begin{KStheorem}[Bounded-gap resampling and selected repair]\label{KS:KS67}
Suppose \(\phi(n+1)-\phi(n)\le m\) for a fixed natural number \(m\). Then the actual sampling maps are bounded and satisfy
\[
 \KSnorm{S_\phi}\le1,\qquad \KSnorm{C_\phi}\le m,
 \qquad \partial_\phi S_\phi=C_\phi\partial.
\]
Exact candidates remain exact under every monotone common resampling, even without a gap bound. If \(A\) and \(B\) are bounded right inverses for the fine and sampled differentials, then
\[
 \KSnorm{P_BS_\phi f-S_\phi P_Af}
 \le(\KSnorm A+m\KSnorm B)\KSnorm{\partial f}.
\]
\end{KStheorem}
\begin{proof}
For a bounded candidate family, \(\sup_n\KSnorm{f_{\phi(n)}}\le\sup_k\KSnorm{f_k}\), proving the sampling bound. For a bounded residual family, each summand in \eqref{KS:eq:block-defect-sum} is a map from \(X_b\) to \(Y_a\) and has norm at most \(\KSnorm r_\infty\), because every composite bond is contractive. A selected block has length at most \(m\), so the triangle inequality bounds its sum by \(m\KSnorm r_\infty\). Taking the supremum over blocks proves that \(C_\phi\) maps into the bounded residual product and has the asserted operator norm. If a block has length zero its sum is zero, in agreement with the identity restriction.

To verify the commuting square, insert
\(r_k=\rho^Y_kf_{k+1}-f_k\rho^X_k\). The \(k\)-th term becomes
\[
 \rho^Y_{a,k+1}f_{k+1}\rho^X_{k+1,b}
 -\rho^Y_{a,k}f_k\rho^X_{k,b}.
\]
Successive internal terms cancel in the finite sum, leaving \(\rho^Y_{a,b}f_b-f_a\rho^X_{a,b}\). This is exactly the sampled differential. No infinite series interchange is involved.

If the fine differential vanishes, the finite telescope is zero on every selected interval. That conclusion remains true for unbounded gap lengths, because each individual interval is finite. What may fail without a gap bound is boundedness of the residual transformation on \emph{arbitrary} bounded residual families, not preservation of exact inputs.

When the two bounded sections exist, Theorem~\ref{KS:KS66} applies to the commuting square just established. Substitute \(\KSnorm{S_\phi}\le1\) and \(\KSnorm{C_\phi}\le m\) in its estimate to get \((\KSnorm A+m\KSnorm B)\KSnorm{\partial f}\). Existence of those sections is an additional hypothesis, not a consequence of the sampling estimate.
\end{proof}
\KSsource{KS67}
The existence of the two right inverses is not inferred merely from contractivity of the bonds. The split-tower theorem provides them under its additional section and propagation-gap conditions. Keeping these hypotheses separate prevents a resampling theorem from being mistaken for an unconditional reconstruction algorithm.

\subsection{Two distinct losses under coarsening}
On scalar identity towers, a fine residual block \((1,-1)\) has zero two-step sum. Thus \(C_\phi r=0\) need not imply \(r=0\), even for a bounded-gap schedule. This is cancellation. Retaining an independent detail coordinate remedies that information loss in a split coarsener, as the next section shows.

Unbounded gaps cause a different issue. With \(\phi(n)=n^2\), the bounded residual family \(r_k=1\) has block sums \(2n+1\). These are unbounded. Therefore block accumulation cannot be a map from every bounded residual family to a bounded sampled residual family in this example. The constant residual is not claimed to be the differential of a bounded candidate on the identity tower. It is a witness about the proposed operator on the entire residual space.

Neither example contradicts cofinal invariance of coherent points or eventual transports. Those theorems concern exact compatibility. A norm on all raw defects, the retained fine detail, and the choice of a repair section are additional structures.

\begin{insight}
A compatible change of representation transports equations correctly. It transports the selected repair correctly only when the chosen residual sections also commute. The mismatch is an exact-valued operator, so it can remain invisible to a residual-only check.
\end{insight}



\chap{Multiscale residuals and the limits of reconstruction}\label{KS:sec:cascade}
\subsection{A split coarsener retains its missing detail}
Let \(C:R\to\overline R\) and \(J:\overline R\to R\) be linear maps with \(CJ=I\). Put \(H=I-JC\). Then
\[
 r=JCr+Hr,\qquad CH=0,\qquad H^2=H.
\]
The two coordinates \((Cr,Hr)\) retain all information. The coarse coordinate alone generally does not.

The actual dyadic tower coarsener has the typed formula
\begin{equation}\label{KS:eq:dyadic}
 (Cr)_n=\rho^Y_{2n}r_{2n+1}+r_{2n}\rho^X_{2n+1}.
\end{equation}
It is obtained by telescoping two adjacent compatibility squares. With target sections, a lift is
\[
 (Jc)_{2n}=0,\qquad (Jc)_{2n+1}=s_{2n}c_n.
\]
The section identity gives \(CJ=I\), so the missing detail is explicitly \(Hr\), not an unspecified reconstruction error. The normed construction imposes the bounds needed to make these maps bounded on its operator family spaces.
\KSsource{KS25}

\subsection{A whole cascade}
At scale \(j\), let \(C_j:R_j\to R_{j+1}\) and \(J_j:R_{j+1}\to R_j\) be bounded linear maps with \(C_jJ_j=I\). For \(x\in R_0\), define
\[
 x_0=x,\quad x_{j+1}=C_jx_j,\quad
 L_0=I,\quad L_N=J_0\cdots J_{N-1},
\]
\[
 h_j=x_j-J_jx_{j+1},\qquad Q_Nx=L_Nx_N,
 \qquad S_N=I-Q_N.
\]
The map \(Q_N\) restores the surviving coarse remainder to the original space. It is this restored remainder, not just \(x_N\), that controls reconstruction.

\begin{KStheorem}[Finite reconstruction and nested remainders]\label{KS:KS26}
For every \(N\),
\[
 x=Q_Nx+\sum_{j<N}L_jh_j,
 \qquad S_Nx=\sum_{j<N}L_jh_j.
\]
Each \(Q_N\) is a projection, and for \(N\le M\),
\[
 Q_NQ_M=Q_MQ_N=Q_M.
\]
The vectors with all details zero are exactly the persistent subspace
\[
 \mathcal P_\infty=\{x:Q_Nx=x\text{ for every }N\}.
\]
\end{KStheorem}
\begin{proof}
Let $D_0=I$ and $D_N=C_{N-1}\cdots C_0:R_0\to R_N$. Thus $x_N=D_Nx$, $Q_N=L_ND_N$, and repeated use of $C_jJ_j=I$ gives
\[
 D_NL_N=I_{R_N}.
\]
This cancellation takes place between adjacent maps with matching domains. No identification of different residual spaces is needed.

Since $L_{j+1}=L_jJ_j$ and $x_{j+1}=C_jx_j$,
\[
 L_jh_j=L_j(x_j-J_jx_{j+1})=Q_jx-Q_{j+1}x.
\]
Summing from $j=0$ to $N-1$ leaves $x-Q_Nx$ because $Q_0=I$. This proves both finite reconstruction identities, including $N=0$, when the detail sum is empty.

Idempotence follows from $D_NL_N=I$:
\[
 Q_N^2=L_N(D_NL_N)D_N=Q_N.
\]
For $M\ge N$, put $C_{N,M}=C_{M-1}\cdots C_N$ and $J_{N,M}=J_N\cdots J_{M-1}$, using identities for $M=N$. Then
\[
 D_M=C_{N,M}D_N,\quad L_M=L_NJ_{N,M},\quad C_{N,M}J_{N,M}=I.
\]
The identities $D_NL_M=J_{N,M}$ and $D_ML_N=C_{N,M}$ now give
\[
 Q_NQ_M=L_NJ_{N,M}D_M=Q_M,\qquad
 Q_MQ_N=L_MC_{N,M}D_N=Q_M.
\]
Both orders have been checked, rather than inferred from one range inclusion.

If all $h_j$ vanish, the finite telescope gives $Q_Nx=x$ for every $N$. Conversely, if every $Q_Nx=x$, then $L_jh_j=Q_jx-Q_{j+1}x=0$. Apply the left inverse $D_j$ of $L_j$ to obtain $h_j=0$. The injectivity of the restoration composite is the step that recovers the original detail from its transported value.
\end{proof}
\KSsource{KS26}

\begin{figure}[htbp]\centering\resizebox{.96\linewidth}{!}{
% BEGIN retained input sources/monographs/komorebi-space/source/figures/05_cascade.tex
\begin{tikzpicture}[x=1cm,y=1cm]
\foreach \i/\x in {0/-4.8,1/-1.6,2/1.6,3/4.8}{\node[box,minimum width=1.5cm] (r\i) at (\x,1.7) {$x_{\i}\in R_{\i}$};}
\foreach \i/\j in {0/1,1/2,2/3}{
 \draw[arr] (r\i.east) -- node[lab,above] {$C_{\i}$} (r\j.west);
 \draw[thinarr] (r\j.north) to[out=130,in=50] node[lab,above] {$J_{\i}$} (r\i.north);
}
\foreach \i/\x in {0/-4.8,1/-1.6,2/1.6}{
 \node[box,minimum width=1.6cm] (h\i) at (\x,-.3) {$h_{\i}$};
 \draw[arr] (r\i) -- node[lab,right] {detail} (h\i);
}
\node[box,draw=KomorebiOchre,fill=KomorebiWarm,minimum width=2.1cm] (q) at (4.8,-.3) {$Q_3x=L_3x_3$};
\draw[arr] (r3)--node[lab,right] {restore} (q);
\node[lab,text width=13cm] at (0,-1.6) {$x=h_0+J_0h_1+J_0J_1h_2+Q_3x$\\[6pt] Finite exactness includes the boundary remainder.};
\end{tikzpicture}

}\caption{A split cascade retains every discarded component as a detail. The finite sum is always exact after the restored coarse remainder is included. An infinite detail expansion requires that remainder to vanish in the topology being used.}\label{KS:fig:cascade}\end{figure}

\subsection{Convergence has its own theorem}
Call \(x\) norm recoverable when \(Q_Nx\to0\). This is equivalent to \(S_Nx\to x\), but it does not follow merely from vanishing persistent subspace. Let
\[
 \mathcal F=\bigcup_N\ker Q_N
\]
be the finite core.

\begin{KStheorem}[Density criterion for norm recovery]\label{KS:KS27}
Suppose \(\sup_N\KSnorm{Q_N}\le K<\infty\). Then the norm recoverable vectors are exactly \(\overline{\mathcal F}\). In particular all vectors are norm recoverable if and only if the finite core is dense.
\end{KStheorem}
\begin{proof}
The nested identities first imply that the kernels increase: if $Q_My=0$ and $N\ge M$, then $Q_Ny=Q_NQ_My=0$. Thus the finite core really is an increasing union of linear subspaces.

Suppose $Q_Nx\to0$. Put $y_N=S_Nx=x-Q_Nx$. Idempotence gives $Q_Ny_N=Q_Nx-Q_N^2x=0$, so $y_N\in\mathcal F$. Since $x-y_N=Q_Nx\to0$, these are finite-core approximants converging to $x$. Hence $x\in\overline{\mathcal F}$.

Conversely, fix $x\in\overline{\mathcal F}$ and $\epsilon>0$. If $K=0$, the common operator bound gives $\KSnorm{Q_Nx}=0$ for every $N$, so convergence is immediate. Suppose $K>0$. Choose $y\in\mathcal F$ satisfying $\KSnorm{x-y}<\epsilon/(K+1)$. Some $M$ has $Q_My=0$. For every $N\ge M$, nesting and the common operator bound give
\[
 \KSnorm{Q_Nx}=\KSnorm{Q_N(x-y)}\le K\KSnorm{x-y}
 <\epsilon\frac{K}{K+1}\le\epsilon.
\]
The approximation $y$ is chosen before the depth $N$, giving precisely the quantifier order in norm convergence. This proves equality of the two sets. Applying it to every $x\in R_0$ gives the density criterion.
\end{proof}
\KSsource{KS27}
There are other sufficient conditions. If the transported details are absolutely summable in the Banach space \(R_0\), their sum exists and the remaining boundary lies in \(\mathcal P_\infty\). If that persistent subspace is zero, the sum equals \(x\). Alternatively,
\[
 \KSnorm{x-S_Nx}\le\left(\prod_{j<N}\KSnorm{J_j}\right)\KSnorm{x_N}
\]
gives a direct sufficient decay test. A decreasing coarse norm by itself is not enough when restoration gains grow.

\subsection{The actual dyadic result}
For the dyadic residual hierarchy defined above, the persistent residual is zero. Zero detail at a scale reveals every even coordinate, while an odd coordinate is passed to a smaller index at the next scale. Strong induction on the coordinate index therefore proves that a residual with all details zero is zero.

\begin{KStheorem}[Dyadic determination and finite coordinate recovery]\label{KS:KS28}
In the actual dyadic system, all detail coordinates jointly determine the residual. Moreover,
\[
 k+1<2^N\quad\Longrightarrow\quad
 (Q_Nr)_k=0,\qquad(S_Nr)_k=r_k.
\]
Thus every fixed residual coordinate is recovered exactly at a sufficiently large finite depth. The statement does not assert convergence in the supremum norm over all coordinates.
\end{KStheorem}
\begin{proof}
Write \(R_j\) for the residual space at scale \(j\). Its lift has the actual coordinate formula
\[
 (J_jz)_{2m}=0,\qquad (J_jz)_{2m+1}=s^{(j)}_{2m}z_m,
\]
where \(s^{(j)}_{2m}\) is the corresponding target section at that scale. The expression on the right is composition of the section with the residual operator. Its source and target are those of the odd coordinate at scale \(j\).

Set \(L_{j,0}=I_{R_j}\) and \(L_{j,N+1}=J_jL_{j+1,N}\). We prove, simultaneously for all starting scales \(j\), that
\[
 k+1<2^N\quad\Longrightarrow\quad (L_{j,N}z)_k=0
 \qquad(z\in R_{j+N}).
\]
For \(N=0\), the inequality \(k+1<1\) has no solution. Suppose the assertion holds at depth \(N\). At depth \(N+1\), an even coordinate vanishes by the displayed lift formula. For an odd coordinate \(k=2m+1\), the hypothesis gives
\[
 2(m+1)<2^{N+1},\qquad m+1<2^N.
\]
The induction hypothesis at starting scale \(j+1\) makes \((L_{j+1,N}z)_m\) zero. Composing that zero operator with \(s^{(j)}_{2m}\) gives zero at coordinate \(k\). This completes the induction without identifying residual spaces at different scales.

Now \(Q_Nr=L_{0,N}x_N\), so the support assertion applies to the restored remainder. Since \(S_Nr=r-Q_Nr\), it gives the reconstruction equation. Finally, if two residuals have identical details at every scale, the finite telescope makes their \(S_N\) reconstructions equal for every \(N\). At any fixed coordinate choose \(N\) with \(k+1<2^N\). The reconstruction equation then makes their original coordinates equal. Hence the residuals are equal.
\end{proof}
\KSsource{KS28}
The strict inequality matters: the next coordinate, \(k=2^N-1\), can still have a nonzero remainder. There is also a sharper, nonconsecutive description of the support.

\begin{KSproposition}[The dyadic support pattern]\label{KS:KS55}
For the same lifts,
\[
 2^N\nmid(k+1)\quad\Longrightarrow\quad(Q_Nr)_k=0.
\]
Thus coordinate \(k\) is certainly recovered once \(N>\nu_2(k+1)\), where \(\nu_2(m)\) is the largest power-of-two exponent dividing the positive integer \(m\). The only possible unrecovered coordinates at depth \(N\) are
\[
 2^N-1,\quad 2\cdot2^N-1,\quad3\cdot2^N-1,\quad\ldots.
\]
\end{KSproposition}
\begin{proof}
Prove the assertion for \(L_{j,N}\) by induction, uniformly in \(j\). At depth zero every positive integer is divisible by one, so there is nothing to prove. At depth \(N+1\), even coordinates vanish. If \(k=2m+1\), then
\[
 2^{N+1}\mid(k+1)\quad\Longleftrightarrow\quad2^N\mid(m+1).
\]
When divisibility fails, the inner coordinate in \(J_jL_{j+1,N}\) vanishes by induction. The formula for the odd lift finishes the proof. Apply the result to \(Q_Nr=L_{0,N}x_N\).
\end{proof}
\KSsource{KS55}
This refinement is formalized by a scale-uniform induction on the actual varying-fiber residual system in \KScode{DyadicSupport.lean}. It asserts where a remainder \emph{may} survive, not that every allowed coordinate is nonzero. In the scalar identity-bonding model with \(C(a,b)=a+b\), the residual constantly equal to one has value \(2^N\) at every surviving position. That example proves sharpness of the uniform prefix cutoff.

\begin{figure}[htbp]\centering\resizebox{.96\linewidth}{!}{
% BEGIN retained input sources/monographs/komorebi-space/source/figures/16_dyadic_support.tex
\includegraphics[width=15cm]{sources/monographs/komorebi-space/source/figures/generated/16_dyadic_support.pdf}

}\caption{The possible support of the restored dyadic remainder is sparse and periodic. At depth \(N\), only positions whose index plus one is divisible by \(2^N\) can remain. Open cells are guaranteed recovered. Marked cells are allowed to survive, not asserted nonzero for every residual.}\label{KS:fig:dyadicsupport}\end{figure}

\subsection{Two boundary examples}
On \(\ell^\infty(\KSN,\KSR)\), let \(C\) remove the first coordinate and let \(J\) prepend zero. Then \(Q_Nx\) is the tail of \(x\), placed at its original coordinates with zeros before \(N\). Every fixed coordinate is eventually zero in the remainder, and \(\KSnorm{Q_N}\le1\). Nevertheless the constant sequence \(\mathbf1\) satisfies \(\KSnorm{Q_N\mathbf1}=1\) for every \(N\).

\begin{KSproposition}[The recovery domain is \(c_0\)]\label{KS:KS29}
For this shift cascade,
\[
 S_Nx\to x\text{ in }\ell^\infty
 \quad\Longleftrightarrow\quad x_n\to0.
\]
For \(x_n=q^n\), \(0\le q<1\), the error is exactly \(q^N\).
\end{KSproposition}
\begin{proof}
The error norm is \(\sup_{n\ge N}|x_n|\). It tends to zero precisely when the sequence tends to zero. For a geometric sequence its largest tail value is the first one, \(q^N\).
\end{proof}
\KSsource{KS29}
As a second example, take scalar spaces with \(C=\tfrac12I\), \(J=2I\). The coarse value \(x_N=2^{-N}x\) tends to zero, but restoration gives \(Q_Nx=x\). Every detail is zero and the entire input is persistent. The example belongs to this scalar cascade, not to the actual dyadic hierarchy, whose persistent subspace was proved to vanish.

\begin{insight}
Retained packets make omission explicit. They do not guarantee convergence. Coordinate recovery, bounded reconstruction, summability, and attainment remain different mathematical properties.
\end{insight}



% BEGIN retained input sources/monographs/komorebi-space/source/chapters/04_symmetry.tex

\chapter{Symmetry and effective descent}\label{KS:book:chapter-21}
A set of solutions forgets the witnesses relating equivalent solutions. The next layer retains those witnesses, their composition, and their automorphisms. Gluing then has to respect the overlap identifications themselves.

\section{Nonabelian coherence and relative trivialization}\label{KS:sec:gauge}
\subsection{Coherent transition labels}
Fix groups \(G_v\) and restriction homomorphisms \(f^*:G_w\to G_v\), with the same identity and composition laws as before. The groups need not be abelian and the homomorphisms need not be bijective.

\begin{KSdefinition}[Cocycle]
A cocycle is a family \(c_f\in G_v\) for every \(f:v\to w\), satisfying
\[
 c_{1_v}=1,\qquad c_{g\circ f}=c_f f^*(c_g).
\]
It determines a twisted view system with maps
\[
 T_c(f)(x)=c_f f^*(x).
\]
A section of the twisted system satisfies \(x_v=c_f f^*(x_w)\).
\end{KSdefinition}
The cocycle law says that the transition rules compose coherently. It does not assert that the twisted system has a section. The coherent transition system and a global trivialization of that system are different objects.

For a family \(a=(a_v)\in\prod_vG_v\), define
\begin{equation}\label{KS:eq:gauge}
 (a\cdot c)_f=a_vc_f[f^*(a_w)]^{-1}.
\end{equation}

\begin{KSproposition}[Gauge action and its groupoid]\label{KS:KS30}
Equation~\eqref{KS:eq:gauge} defines an action on cocycles. A gauge carrying \(c\) to \(d\) induces an equivalence of twisted section sets by \(x_v\mapsto a_vx_v\). The cocycles and actual gauges form a groupoid: an arrow \(a:c\to d\) followed by \(b:d\to e\) has label \(ba\).
\end{KSproposition}
\begin{proof}
First, $T_c$ really is a view system. For $f:v\to w$, $g:w\to z$, and $x\in G_z$, the homomorphism and cocycle laws give
\[
 T_c(f)T_c(g)(x)=c_ff^*(c_gg^*x)
 =c_ff^*(c_g)f^*g^*x=c_{g\circ f}(g\circ f)^*x.
\]
The identity equation follows from $c_{1_v}=1$. No rearrangement of factors is involved.

Put $d=a\cdot c$. Its identity label is $a_v1a_v^{-1}=1$. For a composite, expand all factors in their original order:
\[
\begin{aligned}
 d_f f^*(d_g)
 &=a_vc_f[f^*(a_w)]^{-1}
   f^*(a_wc_g[g^*(a_z)]^{-1})\\
 &=a_vc_f[f^*(a_w)]^{-1}f^*(a_w)f^*(c_g)
       [f^*g^*(a_z)]^{-1}\\
 &=a_vc_{g\circ f}[(g\circ f)^*(a_z)]^{-1}
 =d_{g\circ f}.
\end{aligned}
\]
Thus the gauge formula preserves the cocycle condition. For two gauges $a,b$, substitution gives
\[
 (a\cdot(b\cdot c))_f
 =a_vb_vc_f[f^*(b_w)]^{-1}[f^*(a_w)]^{-1}
 =((ab)\cdot c)_f.
\]
The reversed order in the two inverse factors is exactly the formula for the inverse of a product. The identity gauge acts trivially.

If $d=a\cdot c$ and $x$ is a section of $T_c$, put $y_v=a_vx_v$. Then
\[
 d_ff^*(y_w)
 =a_vc_f[f^*(a_w)]^{-1}f^*(a_w)f^*(x_w)
 =a_vx_v=y_v.
\]
The inverse map is multiplication by $a_v^{-1}$ at each view, proving an equivalence of the actual section sets.

For the groupoid, an arrow $a:c\to d$ means the equation $a\cdot c=d$. If $b:d\to e$, then $(ba)\cdot c=b\cdot(a\cdot c)=e$, giving composition. The identity label is $1$ and the inverse label is $a^{-1}$. Associativity and the two inverse identities now follow from the corresponding multiplication laws in the product group. The label $ba$ is therefore forced by the order of application, not an arbitrary notational convention.
\end{proof}
\KSsource{KS30}
The action groupoid \(Z^1(\KSV;G)//\prod_vG_v\) retains comparison witnesses. Its orbit set \(Z^1(\KSV;G)/\prod_vG_v\) retains only isomorphism classes. There is no reason for the orbit set to be a group in this nonabelian setting.

\subsection{Automorphisms and ambiguity are not the same question}
The stabilizer of a cocycle is the group of families with
\[
 a_v=c_f f^*(a_w)c_f^{-1}.
\]
If one gauge \(a:c\to d\) is known, all gauges from \(c\) to \(d\) are uniquely of the form \(ah\), with \(h\in\KSStab(c)\). Indeed \(a^{-1}b\) stabilizes \(c\) for every other gauge \(b\). This is the ordinary Hom torsor principle in a groupoid.

That symmetry group must not be confused with the set of global trivializations of one fixed cocycle. The latter will be parametrized by ordinary coherent group elements, after one trivialization has been selected.

\subsection{The relative obstruction without commutativity}
Prescribe a coherent partial section \(s\) of \(T_c\) on \(U\). The based gauge group is
\[
 \mathcal G_U=\{a\in\prod_vG_v:a_u=1\text{ for }u\in U\}.
\]
Choose a raw family \(q\) equal to \(s\) on \(U\). Identity elements can be used elsewhere. Normalize the labels by
\[
 c^q_f=q_v^{-1}c_f f^*(q_w).
\]
Its labels on arrows internal to \(U\) are one. Define the pointed orbit obstruction
\[
 \mathfrak o_U(c,s)=[c^q]\in Z^1(\KSV;G)/\mathcal G_U,
\]
whose neutral point is the class of the trivial cocycle.

\begin{KStheorem}[Nonabelian relative extension]\label{KS:KS31}
The obstruction is independent of the raw lift. It is neutral exactly when \(s\) extends to a global section of \(T_c\).
\end{KStheorem}
\begin{proof}
The normalized label $c^q$ is the gauge transform of $c$ by $q^{-1}$, so Proposition~\ref{KS:KS30} makes it a cocycle. If $f:v\to w$ is internal to $U$, the partial-section equation says $q_v=c_ff^*(q_w)$, and therefore $c_f^q=1$ on that arrow.

If $t$ is another raw lift, define $a_v=t_v^{-1}q_v$. On $U$ both lifts equal $s$, so $a_v=1$ and $a\in\mathcal G_U$. Its action is
\[
\begin{aligned}
 (a\cdot c^q)_f
 &=t_v^{-1}q_v\,q_v^{-1}c_ff^*(q_w)
       [f^*(t_w^{-1}q_w)]^{-1}\\
 &=t_v^{-1}c_ff^*(q_w)
       [f^*(q_w)]^{-1}f^*(t_w)
 =t_v^{-1}c_ff^*(t_w)=c_f^t.
\end{aligned}
\]
Thus changing the raw lift leaves the based orbit unchanged.

Suppose the orbit is neutral. Equality of orbit classes means there exists an actual based gauge $b$ with $b\cdot c^q=1$. At an arrow $f:v\to w$, this equation is
\[
 b_vq_v^{-1}c_ff^*(q_w)[f^*(b_w)]^{-1}=1.
\]
Multiply on the left by $q_vb_v^{-1}$. Homomorphic preservation of products then gives
\[
 c_ff^*(q_wb_w^{-1})=q_vb_v^{-1}.
\]
Therefore $x_v=q_vb_v^{-1}$ is a twisted coherent family. On $U$, $b_v=1$ and $q_v=s_v$, so it extends the prescribed section.

Conversely, suppose $x$ is a global twisted section extending $s$. Put $b_v=x_v^{-1}q_v$. This is based, and direct substitution gives
\[
 (b\cdot c^q)_f
 =x_v^{-1}c_ff^*(x_w)=1,
\]
where the last equality is precisely the section equation for $x$. Hence the obstruction is neutral. These arguments establish the result using group cancellation only. In particular they do not assume an additive quotient, a norm on labels, or a decision procedure for finding $b$.
\end{proof}
\KSsource{KS31}
The orbit is defined from actual cocycles and an actual group action. It is not a Boolean test defined to be neutral precisely when an extension exists. Deciding orbit membership in arbitrary groups remains a separate problem.

Let
\[
 H_U=\{h:f^*(h_w)=h_v\text{ for every }f,\ h|_U=1\}.
\]

\begin{KStheorem}[The nonabelian extension torsor]\label{KS:KS32}
Given one extension \(x_0\), the map \(h\mapsto x_0h\) is an equivalence from \(H_U\) to the extension set. The extension is unique if and only if the obstruction is neutral and \(H_U\) is trivial.
\end{KStheorem}
\begin{proof}
Take \(h\in H_U\). Its pointwise product with the chosen extension \(x_0\) agrees with the prescribed data on \(U\), because \(h|_U=1\). For an arrow \(f:v\to w\), the coefficient homomorphism preserves multiplication and the two coherence equations give
\[
 c_f f^*((x_0)_wh_w)
 =c_f f^*((x_0)_w)f^*(h_w)
 =(x_0)_vh_v.
\]
Thus \(x_0h\) is another twisted extension.

Conversely, if \(x\) is a twisted extension of the same prescribed data, define \(h_v=(x_0)_v^{-1}x_v\). On \(U\) its value is one. For every arrow,
\[
 \begin{aligned}
 h_v
 &=[c_f f^*((x_0)_w)]^{-1}[c_f f^*(x_w)]\\
 &=f^*((x_0)_w)^{-1}c_f^{-1}c_f f^*(x_w)\\
 &=f^*((x_0)_w^{-1}x_w)=f^*(h_w).
 \end{aligned}
\]
The cancellation uses inverse multiplication in its actual order. No commuting factors have been moved past one another. Hence \(h\in H_U\). Pointwise group cancellation shows that \(h\mapsto x_0h\) and \(x\mapsto x_0^{-1}x\) are inverse maps.

The extension set is nonempty precisely when the earlier orbit obstruction is neutral. Once it is nonempty, the displayed bijection makes it a singleton precisely when \(H_U\) has only its identity element. Both conditions are needed for unique existence. Changing the chosen extension changes the parametrization, not the extension set or the criterion.
\end{proof}
\KSsource{KS32}
The additive theory is recovered by applying a multiplicative type tag to an additive view system and choosing the trivial cocycle. The normalized label then becomes \(f^*q_w-q_v\). The source proves agreement of the actual extension sets and vanishing criteria. It does not identify the whole nonabelian orbit set with every element of the larger additive arrow residual quotient.

\subsection{A small noncommutative model}
\begin{KSexample}[A \(C_2\) loop with \(S_3\) coefficients]\label{KS:KS33}
Take one view with a loop \(\sigma\) satisfying \(\sigma^2=1\), and constant coefficients in a group \(H\). A cocycle is determined by an element \(p\in H\) whose square is the identity. Gauge equivalence is conjugacy, and its automorphisms form the centralizer of \(p\).

For \(H=S_3\), these elements are the identity and the three transpositions. There are exactly two orbit classes. The identity cocycle has automorphism group \(S_3\), which is noncommutative. A nonidentity involution gives a coherent cocycle whose regular twisted section equation \(x=px\) has no solution.
\end{KSexample}
\begin{proof}
The only nontrivial cocycle relation is \(p^2=1\). At a single view the gauge action is \(p\mapsto apa^{-1}\). All three transpositions are conjugate and none is conjugate to the identity. The stabilizer is therefore the centralizer. Finally \(x=px\) implies \(p=1\) after multiplication by \(x^{-1}\).
\end{proof}
\KSsource{KS33}
The example makes an important boundary visible. A cocycle can exist and glue as an object even when it admits no global trivializing section. Stackhood will concern the former gluing problem, not a claim that every cocycle becomes trivial.

\begin{figure}[htbp]\centering\resizebox{.96\linewidth}{!}{
% BEGIN retained input sources/monographs/komorebi-space/source/figures/06_groupoid.tex
\begin{tikzpicture}[x=1cm,y=1cm]
\fill[KomorebiSoft!65,rounded corners=7pt] (-6.7,-.8) rectangle (-1.3,3.2);
\fill[KomorebiSoft!65,rounded corners=7pt] (.1,-.8) rectangle (6.7,3.2);
\node[box,circle,minimum size=1cm] (id) at (-4,1.15) {$1$};
\draw[arr] (id) to[loop above,min distance=12mm] (id);
\node[lab] at (-4,2.7) {$\KSAut(1)=S_3$};
\node[lab] at (-4,-.25) {identity cocycle};
\node[box,circle,inner sep=5pt] (a) at (3.4,2.15) {$(12)$};
\node[box,circle,inner sep=5pt] (b) at (1.4,.35) {$(13)$};
\node[box,circle,inner sep=5pt] (c) at (5.4,.35) {$(23)$};
\draw[<->,thick,KomorebiPine] (a)--(b);
\draw[<->,thick,KomorebiPine] (b)--(c);
\draw[<->,thick,KomorebiPine] (c)--(a);
\node[lab,fill=KomorebiSoft!65,inner sep=3pt] at (3.4,1) {conjugation};
\node[lab] at (3.4,-.35) {each stabilizer is $C_2$};
\node[dot] (oi) at (-4,-2.05) {};
\node[dot] (ot) at (3.4,-2.05) {};
\draw[thinarr,dashed] (-4,-.85)--(oi);
\draw[thinarr,dashed] (3.4,-.85)--(ot);
\node[lab,below=4pt] at (oi) {one orbit};
\node[lab,below=4pt] at (ot) {one orbit};
\node[lab,text width=13.5cm] at (0,-3.4) {Two orbit classes, but nontrivial groups of comparison witnesses.\\The arrows and automorphisms disappear in the coarse orbit set.};
\end{tikzpicture}

}\caption{The orbit set sees two classes in the \(C_2\)-to-\(S_3\) example. The groupoid also retains the arrows between conjugate transpositions and all automorphisms. Collapsing the arrows loses information that matters in descent.}\label{KS:fig:groupoid}\end{figure}

\chap{Effective gluing of objects and their identifications}\label{KS:sec:descent}
\subsection{Why pairwise agreement is not the right datum}
For objects with symmetry, local descriptions need not be literally equal on overlaps. They may instead come with specified isomorphisms. Those isomorphisms must satisfy their own composition law. It is not enough to know that some isomorphism exists for each pair.

Let \(U_i\) be full patches in \(\KSV\). Assume the composition cover condition
\begin{equation}\label{KS:eq:compositioncover}
 \forall(v\xrightarrow f w\xrightarrow g z),\quad
 \exists i:\ v,w,z\in U_i.
\end{equation}
By inserting identity arrows, this implies coverage of objects and arrows.

A descent datum consists of local cocycles \(c_i\) and overlap gauges \(a_{ij}(v)\) with
\begin{align}
 c_j(f)&=a_{ij}(v)c_i(f)[f^*a_{ij}(w)]^{-1},\label{KS:eq:localgauge}\\
 a_{ik}(v)&=a_{jk}(v)a_{ij}(v).\label{KS:eq:triple}
\end{align}
The triple equation implies \(a_{ii}=1\) and \(a_{ji}=a_{ij}^{-1}\). A morphism of descent data is a family of local gauges \(r_i\) satisfying
\[
 r_j(v)a_{ij}(v)=a'_{ij}(v)r_i(v).
\]
These objects and arrows form a groupoid.

\begin{KStheorem}[Effective cocycle descent]\label{KS:KS34}
Under \eqref{KS:eq:compositioncover}, restriction is an equivalence between the global cocycle gauge groupoid and the groupoid of descent data. The overlap gauges supplied in the datum are realized, not discarded.
\end{KStheorem}
\begin{proof}
\emph{Step 1: normalize the supplied identifications.} The triple-overlap equation with all indices equal gives $a_{ii}=a_{ii}a_{ii}$, hence $a_{ii}=1$. With indices $i,j,i$ it gives $a_{ji}a_{ij}=1$, hence $a_{ji}=a_{ij}^{-1}$. Composition coverage implies object coverage by inserting identity arrows. Choose a containing patch $\kappa(v)$ for each view and put
\[
 b_i(v)=a_{\kappa(v),i}(v)\quad(v\in U_i).
\]
At a point of $U_i\cap U_j$, the triple law, now on patches $\kappa(v),i,j$, gives
\[
 b_j(v)=a_{ij}(v)b_i(v).                                    \tag{*}
\]
The two endpoints of an arrow may have different reference patches. No common reference patch is being assumed.

\emph{Step 2: obtain literal matching labels.} Define
\[
 \widetilde c_i(f)=b_i(v)^{-1}c_i(f)f^*(b_i(w))
 \quad(f:v\to w\text{ in }U_i).
\]
This is a gauge transform and hence a cocycle on $U_i$. If the arrow lies also in $U_j$, substitute the original overlap equation and (*) at its two endpoints:
\[
\begin{aligned}
 \widetilde c_j(f)
 &=b_i(v)^{-1}a_{ij}(v)^{-1}
   a_{ij}(v)c_i(f)[f^*a_{ij}(w)]^{-1}
   f^*(a_{ij}(w)b_i(w))\\
 &=b_i(v)^{-1}c_i(f)f^*(b_i(w))
 =\widetilde c_i(f).
\end{aligned}
\]
Every cancellation is between neighboring inverse factors. Commutativity is not required.

\emph{Step 3: construct the global cocycle.} Every arrow is contained in a patch, again by inserting an identity in a composable pair. Define $c_f$ to be $\widetilde c_i(f)$ in any such patch. Step 2 proves independence. For an identity arrow, choose a patch through its view to obtain $c_{1_v}=1$. For $f:v\to w$ and $g:w\to z$, composition coverage supplies a patch containing all three vertices. Fullness places $f$, $g$, and $g\circ f$ in that same patch. Its cocycle law therefore proves
\[
 c_{g\circ f}=\widetilde c_i(g\circ f)
 =\widetilde c_i(f)f^*(\widetilde c_i(g))=c_ff^*(c_g).
\]
This is the exact point where composition coverage, rather than arrow coverage alone, is used.

Undoing the normalization proves
\[
 b_i(v)c_f[f^*(b_i(w))]^{-1}=c_i(f).
\]
Together with (*), these gauges constitute an isomorphism from the restriction of $c$ to the supplied descent datum. They implement its original overlaps, proving essential surjectivity of restriction.

\emph{Step 4: retain every morphism.} Let $c,d$ be global cocycles. Their canonical restrictions have identity overlap gauges. A morphism between those restrictions is consequently a family $r_i$ satisfying $r_j(v)=r_i(v)$ on each overlap. Assemble $r_v$ from any containing patch. For an arrow choose a patch containing both endpoints. Its gauge equation proves
\[
 d_f=r_vc_f[f^*(r_w)]^{-1}.
\]
Thus $r$ is a global gauge. Its restriction is the original family. Conversely, assembling the restrictions of a global gauge preserves every coordinate and returns that gauge. Restriction is therefore bijective on each Hom set, not merely injective on isomorphism classes. This is full faithfulness.

\emph{Step 5: obtain an equivalence.} For each descent object choose the global cocycle and isomorphism constructed above. For an arrow between descent objects, conjugate it by these two chosen isomorphisms and use the unique preimage under the Hom bijection of Step 4. This defines the arrow map of a quasi-inverse. Uniqueness of the preimage verifies its identity and composition laws, and the chosen isomorphisms provide the comparison to the identity functor. Hence the restriction functor is an equivalence of groupoids. Different patch choices can alter the representative but not this conclusion.
\end{proof}
\KSsource{KS34}

\subsection{The correct uniqueness statement}
Suppose \((c,b_i)\) and \((c',b'_i)\) realize the same descent datum. The only possible compatible comparison is
\[
 g_v=b'_i(v)^{-1}b_i(v),
\]
using any patch through \(v\). The overlap identities make it independent of the patch, and \(b'_ig=b_i\). It is therefore the unique global comparison respecting the fixed local identifications.

This is not a statement that the underlying cocycle has no automorphisms. Those automorphisms can reappear when the local identifications are forgotten. Nor does the construction choose one representative independently of every reference patch. It proves that the choices are related by the appropriate unique comparison.

\subsection{A failed triple}
Let \(t\ne1\) be an involution and ask for frames \(b_0,b_1,b_2\) satisfying
\[
 b_1=tb_0,\qquad b_2=tb_1,\qquad b_2=tb_0.
\]
The first two give \(b_2=b_0\), while the third forces \(t=1\). Each displayed transition is individually invertible, but they cannot be the overlap identifications of a global frame. The defect lies in the triple relation.

\begin{figure}[htbp]\centering\resizebox{.96\linewidth}{!}{
% BEGIN retained input sources/monographs/komorebi-space/source/figures/07_cover.tex
\begin{tikzpicture}[x=1cm,y=1cm]
\node[lab,font=\bfseries\small] at (-3.8,3.65) {Three views};
\node[lab,font=\bfseries\small] at (3.5,3.65) {Four views};
\foreach \i/\x in {0/-5.4,1/-3.8,2/-2.2}{\node[box,circle,inner sep=4pt] (a\i) at (\x,2.3) {$\i$};}
\draw[arr] (a0)--(a1);\draw[arr] (a1)--(a2);
\draw[thinarr] (a0.north) to[out=60,in=120] (a2.north);
\foreach \i/\x in {0/1.4,1/2.8,2/4.2,3/5.6}{\node[box,circle,inner sep=4pt] (b\i) at (\x,2.3) {$\i$};}
\draw[arr](b0)--(b1);\draw[arr](b1)--(b2);\draw[arr](b2)--(b3);
\foreach \i/\x in {0/-5.4,1/-3.8,2/-2.2}{\node[lab] at (\x,1.25) {$\i$};}
\foreach \j/\y/\name in {0/.55/{01},1/-.15/{02},2/-.85/{12}}{\node[lab] at (-6.25,\y) {$U_{\name}$};}
\foreach \x/\y in {-5.4/.55,-3.8/.55,-5.4/-.15,-2.2/-.15,-3.8/-.85,-2.2/-.85}{\fill[KomorebiPine] (\x,\y) circle (2.5pt);}
\foreach \i/\x in {0/1.4,1/2.8,2/4.2,3/5.6}{\node[lab] at (\x,1.25) {$\i$};}
\foreach \i/\y in {0/.55,1/-.15,2/-.85,3/-1.55}{
 \node[lab] at (.45,\y) {$U_{\i}$};
 \foreach \j/\x in {0/1.4,1/2.8,2/4.2,3/5.6}{\ifnum\i=\j\draw[KomorebiMuted] (\x,\y) circle (2.5pt);\else\fill[KomorebiPine] (\x,\y) circle (2.5pt);\fi}
}
\node[lab,text width=5.6cm,text=KomorebiOchre] at (-3.8,-2.35) {Every arrow is covered.\\The three-vertex composition is not.};
\node[lab,text width=5.6cm,text=KomorebiPine] at (3.5,-2.75) {Every composition support is covered.\\No patch is the whole category.};
\end{tikzpicture}

}\caption{Object coverage, arrow coverage, and composition coverage are distinct. The proper two-view patches on a three-view chain see every arrow but no patch sees the relation \(0\to1\to2\). The four proper three-view patches on a four-view chain do cover every composable pair.}\label{KS:fig:cover}\end{figure}

\chap{A genuine stack on the composition site}\label{KS:sec:stack}
The preceding theorem can be placed in a standard site, rather than called stack-like. We use the usual effective-descent meaning of a stack~\cite{KS:stacks-stack,KS:mathlib-stack}. The base of this site is the inclusion poset of full patches
\[
 \mathcal P(\KSV)=(\operatorname{Set}(\KSOb\KSV),\subseteq).
\]
This patch category is different from the original category of views.

\subsection{The topology is defined without assuming effectivity}
For a composable pair \(q=(v\xrightarrow f w\xrightarrow g z)\), let its support be \(|q|=\{v,w,z\}\). A sieve \(\mathcal R\) on \(W\) is covering when
\begin{equation}\label{KS:eq:site}
 |q|\subseteq W\quad\Longrightarrow\quad
 (|q|\hookrightarrow W)\in\mathcal R
 \quad\text{for every composable }q.
\end{equation}
This condition depends only on the view category, not on the coefficient groups or on a promised global realization.

\begin{KSproposition}[The composition topology]\label{KS:KS35}
Condition~\eqref{KS:eq:site} defines a Grothendieck topology \(J_{\mathrm{comp}}\). The sieve generated by a family of inclusions \(U_i\subseteq W\) is covering exactly when the family covers every composition support in \(W\).
\end{KSproposition}
\begin{proof}
In the inclusion poset, identify a sieve on $W$ with a downward-closed collection of subsets of $W$. Pullback along $Z\subseteq W$ consists of those subsets of $Z$ that were already members of the original sieve. We verify the three topology axioms using this description.

The maximal sieve contains every subset of $W$, hence every composition support. For stability under pullback, let $\mathcal R$ cover $W$ and let $q$ be a composable pair whose support lies in $Z\subseteq W$. The same pair lies in the full patch $W$, so $|q|\in\mathcal R$. Since $|q|\subseteq Z$, it belongs to the pulled-back sieve on $Z$. Thus that sieve covers $Z$.

For transitivity, suppose $\mathcal R$ covers $W$, and that for every $U\in\mathcal R$ the pullback of a sieve $\mathcal T$ to $U$ covers $U$. Fix a composable pair $q$ in $W$. Its support $S=|q|$ belongs to $\mathcal R$. Since patches are full, $q$ is also a composable pair in $S$. The covering condition for the pullback of $\mathcal T$ to $S$ therefore puts the identity inclusion $S\subseteq S$ in that pullback. By the definition of pullback, this says $S\in\mathcal T$. It holds for every $q$, so $\mathcal T$ covers $W$.

Finally, the sieve generated by $U_i\subseteq W$ consists exactly of the subpatches $Z$ contained in at least one $U_i$. One inclusion follows because a sieve contains all restrictions of its generators. The reverse inclusion follows because this collection is itself downward closed and contains all generators, hence is the smallest such sieve. Consequently a support belongs to the generated sieve exactly when some $U_i$ contains all its endpoints. This proves the family characterization. When $W$ is empty, there are no composable pairs to test, so even the empty sieve is covering.
\end{proof}
\KSsource{KS35}
Intersections of a covering family with a subpatch remain covering. Replacing each member by a covering family again gives a cover. Identities imply coverage of single views and arrows. The empty patch admits the empty cover.

\subsection{The actual patch pseudofunctor}
On a patch \(W\), use
\[
 \mathfrak K_G(W)=Z^1(\KSV|_W;G|_W)//\prod_{v\in W}G_v.
\]
Restriction keeps the existing cocycle values and gauge values on the smaller patch. It first gives a strict category valued functor on the opposite patch poset. Its standard pseudofunctor form has identity comparison isomorphisms, with the required coherence laws.

\begin{KStheorem}[Stackhood on the composition site]\label{KS:KS36}
The patch pseudofunctor \(\mathfrak K_G\) is a stack for \(J_{\mathrm{comp}}\). It satisfies the standard condition of descent for morphisms and effective descent for objects on every covering sieve.
\end{KStheorem}
\begin{proof}
We use the equivalence-of-descent-categories formulation of stackhood for the already constructed patch pseudofunctor. Its restriction maps act on actual cocycle labels and gauge values. Their identity and composition laws have been proved, so the descent condition is applied to a specified pseudofunctor, not assumed as part of its definition.

\emph{From standard descent data to overlap labels.} For a covering family $(U_i\subseteq W)$, standard descent data provide local objects and comparison gauges on common subpatches, natural under further restriction. At a view $v\in U_i\cap U_j$, restrict the comparison to the singleton patch $\{v\}$ and call its value $a_{ij}(v)$. To verify the gauge equation along $f:v\to w$ in an overlap, use the comparison on $U_i\cap U_j$, which is a full patch containing the arrow. Naturality identifies its two endpoint values with the singleton values. Its own gauge equation then becomes \eqref{KS:eq:localgauge}. Evaluating the standard triple-comparison identity at a singleton gives \eqref{KS:eq:triple}. Thus these labels form the concrete descent datum of Theorem~\ref{KS:KS34}.

\emph{From the concrete realization back to standard data.} By Theorem~\ref{KS:KS34} there is a global cocycle $c$ and local gauges $b_i$ implementing these overlaps. On any common subpatch $Z\subseteq U_i\cap U_j$, the required equation between gauge arrows is an equality of their values at every $v\in Z$. At such a view it is $b_j(v)=a_{ij}(v)b_i(v)$, already proved by the construction. Equality of coordinate functions proves equality of the arrows on $Z$. Thus every standard common-subpatch equation holds, not only the equations on singleton tests. This gives an isomorphism from the standard restriction of $c$ to the original standard descent object.

\emph{Morphisms.} A standard morphism between restrictions of global objects has local gauge components. Its comparison equations at singleton patches show that these components agree on overlaps. The Hom gluing map from Theorem~\ref{KS:KS34} yields a unique global gauge. Conversely, a global gauge restricts to those components and satisfies all common-subpatch equations. Restriction and gluing are inverse on every Hom set. Standard restriction is therefore fully faithful as well as essentially surjective, hence an equivalence.

\emph{Every covering sieve.} Let $\mathcal R$ be an arbitrary covering sieve on $W$. Index a family by all its member inclusions $U\subseteq W$. Each composition support belongs to this family by the covering definition, so it is a composition cover. Its generated sieve is exactly $\mathcal R$: one inclusion uses that $\mathcal R$ is downward closed, and the other uses each member as its own generator. The family argument therefore proves the descent equivalence on the original sieve. Applying it to every covering sieve proves both descent of morphisms and effectiveness of objects. These are the required standard stack conditions \cite{KS:stacks-stack,KS:mathlib-stack}.
\end{proof}
\KSsource{KS36}
The principal source declaration is \KScode{GroupViewSystem.patch_isStack}. Its target is mathlib's \KScode{Pseudofunctor.IsStack}, not a project predicate that repeats the conclusion as an assumption.

\subsection{The cover is nontrivial, but specific}
On a four-view chain, take the four patches obtained by deleting one vertex. All patches are proper, yet every composable pair has at most three endpoints and lies in one of them. The theorem genuinely glues pieces smaller than the whole object.

On the three-view chain, the analogous proper patches cover each object and each arrow but miss the support of \(0\to1\to2\). They do not form a cover in this topology. Consequently the theorem is neither vacuous nor an unrestricted gluing principle.

\begin{scopebox}
\textbf{Exact scope.} This is an ordinary stack of strict group valued cocycle realizations on a specified combinatorial site. It is not an infinity-stack, a universal heterogeneous atlas, or an arithmetic existence theorem. The stack theorem glues coherent cocycles with their overlap isomorphisms. It does not assert that every resulting cocycle has a global trivialization.
\end{scopebox}

\chap{Equivalent presentations and identified boundaries}\label{KS:sec:relative}
\subsection{Transporting the same question}
Let \(H:\KSV\to\KSW\) be part of a category equivalence, with inverse \(L\) and counit \(\varepsilon:HL\Rightarrow\KSid_\KSW\). Pull a view system \(F\) on \(\KSW\) back along \(H\). Sampling a global section gives \((H^*x)_v=x_{Hv}\).

\begin{KSproposition}[Equivalence of section spaces]\label{KS:KS37}
Pullback induces \(\Gamma(F)\simeq\Gamma(H^*F)\). If \(U\subseteq\KSOb\KSW\) is closed under isomorphism and \(U'=H^{-1}U\), it also induces
\[
 \KSExt_U^F(s)\simeq\KSExt_{U'}^{H^*F}(H^*s).
\]
Existence and uniqueness of the prescribed extension are preserved and reflected.
\end{KSproposition}
\begin{proof}
Write the unit of the equivalence as $\eta:I_\KSV\Rightarrow LH$, and take the unit and counit with their triangle identities. Set $t_w=\varepsilon_w^{-1}:w\to HLw$. Given $y\in\Gamma(H^*F)$, define
\[
 (Ry)_w=F(t_w)y_{Lw}.
\]
For an arrow $f:w\to w'$, naturality of the inverse counit states
\[
 HL(f)\,t_w=t_{w'}f.
\]
Contravariance of $F$ and coherence of $y$ therefore give
\[
\begin{aligned}
 F(f)(Ry)_{w'}
 &=F(f)F(t_{w'})y_{Lw'}\\
 &=F(t_w)F(HL(f))y_{Lw'}
 =F(t_w)y_{Lw}=(Ry)_w.
\end{aligned}
\]
Thus recovery gives an actual coherent section.

If $x\in\Gamma(F)$, coherence along $t_w$ gives
\[
 (RH^*x)_w=F(t_w)x_{HLw}=x_w.
\]
For the other composite, the triangle identity says $t_{Hv}=H(\eta_v)$. Coherence of $y$ along $\eta_v:v\to LHv$ yields
\[
 (H^*Ry)_v=F(t_{Hv})y_{LHv}=F(H\eta_v)y_{LHv}=y_v.
\]
These coordinate equations prove both inverse identities.

For the relative claim, pullback clearly sends an extension of $s$ to an extension of $H^*s$ on $U'=H^{-1}U$. Conversely suppose $y$ extends $H^*s$. If $w\in U$, the isomorphism $w\cong HLw$ and isomorphism closure imply $HLw\in U$. Hence $Lw\in U'$ and $y_{Lw}=s_{HLw}$. Both endpoints of $t_w$ lie in $U$, so partial coherence gives
\[
 (Ry)_w=F(t_w)s_{HLw}=s_w.
\]
Recovery therefore preserves the required boundary. The two inverse identities remain valid on the subsets of extensions. Bijections preserve nonemptiness and identify any two elements exactly when their images are equal, proving preservation and reflection of existence and uniqueness.
\end{proof}
\KSsource{KS37}
The coefficients really are pulled back. This does not compare unrelated coefficient systems simply because their diagrams look alike. Isomorphism closure is a sufficient general hypothesis for the relative statement. The theorem does not assert that every unsaturated domain fails.

For group coefficients, the invisible coherent groups \(H_U\) are transported by a multiplicative equivalence, and the parameterizations \(x_0h\) commute with pullback. Thus feasibility and ambiguity are both retained. This is more than a comparison of the number of extensions. It is not an assertion that every possible based gauge orbit set is presentation independent.

\subsection{Why a boundary must retain its identification}
Let \(R:\mathcal C\to\mathcal D\) be a functor between groupoids, and fix \(b\in\mathcal D\). A realization over \(b\) should include a chosen isomorphism to \(b\), rather than require literal equality of representatives.

\begin{KSdefinition}[Marked realization groupoid]
The objects of \(\mathcal E_b(R)\) are pairs \((x,\phi)\), with \(x\in\mathcal C\) and \(\phi:Rx\xrightarrow{\sim}b\). An arrow \((x,\phi)\to(y,\psi)\) is an arrow \(a:x\to y\) with \(\psi R(a)=\phi\).
\end{KSdefinition}
Composition and inverses preserve this equation, so this is a groupoid. It is the explicit groupoid homotopy fiber construction, without invoking a higher categorical interface.

\begin{KSproposition}[Existence and invisible automorphisms]\label{KS:KS38}
The groupoid \(\mathcal E_b(R)\) is nonempty exactly when \(b\) is in the essential image of \(R\). An automorphism of \((x,\phi)\) is an automorphism \(a\) of \(x\) satisfying \(R(a)=1_{Rx}\).
\end{KSproposition}
\begin{proof}
The first assertion is the object definition. For the second, the boundary equation is \(\phi R(a)=\phi\). Cancel the isomorphism \(\phi\). Conversely \(R(a)=1\) makes the equation hold.
\end{proof}
\KSsource{KS38}
This retains precisely the symmetries invisible at the identified boundary. It avoids freezing arbitrary representative coordinates and then mistaking the result for an invariant realization problem.

\begin{figure}[htbp]\centering\resizebox{.96\linewidth}{!}{
% BEGIN retained input sources/monographs/komorebi-space/source/figures/08_marked.tex
\begin{tikzpicture}[x=1cm,y=1cm]
\node[box] (x) at (-5,2) {$x$};\node[box] (y) at (-1.9,2) {$y$};
\node[box] (rx) at (-5,0) {$Rx$};\node[box] (ry) at (-1.9,0) {$Ry$};
\node[box] (b) at (-3.45,-2) {$b$};
\draw[arr](x)--node[lab,above] {$a$}(y);
\draw[thinarr](x)--node[lab,left] {$R$}(rx);\draw[thinarr](y)--node[lab,right] {$R$}(ry);
\draw[arr](rx)--node[lab,above] {$R(a)$}(ry);
\draw[arr](rx)--node[lab,left] {$\phi$}(b);\draw[arr](ry)--node[lab,right] {$\psi$}(b);
\node[lab] at (-3.45,-3) {$\psi R(a)=\phi$};
\node[box] (cc) at (1.1,2) {$\mathcal C$};\node[box] (cp) at (5,2) {$\mathcal C'$};
\node[box] (dd) at (1.1,-.2) {$\mathcal D$};\node[box] (dp) at (5,-.2) {$\mathcal D'$};
\draw[arr](cc)--node[lab,above] {$E$}(cp);\draw[arr](dd)--node[lab,below] {$K$}(dp);
\draw[arr](cc)--node[lab,left] {$R$}(dd);\draw[arr](cp)--node[lab,right] {$R'$}(dp);
\node[lab] at (3.05,.9) {$KR\cong R'E$};
\node[lab,text width=5.8cm] at (3.05,-2) {An isomorphism fills the square.\\The marked fiber is transported.};
\end{tikzpicture}

}\caption{The chosen boundary identification is part of a realization. A comparison must commute with it. Under an equivalent presentation, the comparison square may commute through an isomorphism rather than a literal equality, and the marked realization groupoid is still transported correctly.}\label{KS:fig:marked}\end{figure}

\subsection{Weakly commuting squares are enough}
Suppose \(E:\mathcal C\to\mathcal C'\) and \(K:\mathcal D\to\mathcal D'\) are equivalences, with functors \(R,R'\) and a natural isomorphism
\[
 \sigma:KR\xRightarrow{\sim}R'E.
\]
The square is not required to commute literally.

\begin{KStheorem}[Invariance of the marked realization problem]\label{KS:KS39}
The map
\[
 (x,\phi)\longmapsto(Ex,K(\phi)\sigma_x^{-1}),\qquad a\longmapsto E(a)
\]
defines an equivalence \(\mathcal E_b(R)\simeq\mathcal E_{Kb}(R')\).
\end{KStheorem}
\begin{proof}
For an object $(x,\phi)$ put $\beta_x=K(\phi)\sigma_x^{-1}:R'Ex\to Kb$. This is an isomorphism of the required type. If $a:(x,\phi)\to(y,\psi)$, naturality of $\sigma$ says
\[
 R'E(a)\sigma_x=\sigma_yKR(a),
 \qquad \sigma_y^{-1}R'E(a)=KR(a)\sigma_x^{-1}.
\]
Consequently
\[
 \beta_yR'E(a)
 =K(\psi)KR(a)\sigma_x^{-1}
 =K(\psi R(a))\sigma_x^{-1}
 =K(\phi)\sigma_x^{-1}=\beta_x.
\]
So $E(a)$ is a marked arrow. Identity and composition are preserved because the underlying arrow map is that of the functor $E$.

To prove full faithfulness, start with a marked arrow $u:Ex\to Ey$ between the transported objects. Full faithfulness of $E$ gives a unique $a:x\to y$ with $E(a)=u$. Its marked equation, rewritten by the naturality identity above, is
\[
 K(\psi R(a))\sigma_x^{-1}=K(\phi)\sigma_x^{-1}.
\]
Cancel the isomorphism on the right and use faithfulness of $K$ to obtain $\psi R(a)=\phi$. Thus the lifted arrow is marked. Uniqueness of the underlying lift also proves uniqueness as a marked arrow, giving a bijection on every Hom set.

For essential surjectivity, let $(y,\psi)$ be a target realization. Choose $x$ and an isomorphism $i:Ex\to y$, using essential surjectivity of $E$. Fullness of $K$ lifts the well-typed composite
\[
 KRx\xrightarrow{\sigma_x}R'Ex
 \xrightarrow{R'(i)}R'y\xrightarrow{\psi}Kb
\]
to an arrow $\phi:Rx\to b$. Since $\mathcal D$ is a groupoid, $\phi$ is automatically an isomorphism. Its image satisfies $K(\phi)=\psi R'(i)\sigma_x$, hence $\beta_x=\psi R'(i)$. This is precisely the condition for $i$ to be an arrow from the transported realization to $(y,\psi)$. It is an isomorphism because all categories involved are groupoids.

We have verified full faithfulness and essential surjectivity of the specified transport functor. Choosing these preimages and lifting arrows through the Hom bijections constructs a quasi-inverse, as in Theorem~\ref{KS:KS34}, proving the claimed equivalence. No literal equality $KR=R'E$ was used.
\end{proof}
\KSsource{KS39}
Apply this to the actual restriction functors between cocycle patch groupoids and the equivalences induced by an equivalent view presentation. The full realization problem, including comparison witnesses and automorphisms, is transported.

Replacing \(b\) by an isomorphic boundary \(b'\) also gives an equivalence, by composing the identification with \(b\to b'\). For \(R=\KSid\), the marked realization groupoid has center \((b,1_b)\), and between \((x,\phi)\) and \((y,\psi)\) there is exactly one compatible arrow \(\psi^{-1}\phi\). This is a useful example of contractible realization data, not a claim that the original groupoid has no symmetries.




% BEGIN retained input sources/monographs/komorebi-space/source/chapters/04a_boundary_lab.tex
\chap{A boundary identification is data}\label{KS:sec:boundary-lab}
The groupoid formulation can seem more elaborate than a set of equivalence classes until a change of presentation is examined closely. A two-view example is enough to show what is lost by freezing representatives instead of retaining their comparison maps. It also explains why the marked realization construction uses an identification at the boundary.

\subsection{Duplicate a view without adding information}
Let the original view category have one object and only its identity arrow. Replace it by a category with two objects, \(0\) and \(1\), and exactly one arrow between each ordered pair. These arrows are inverse isomorphisms. The two categories are equivalent: selecting object zero and forgetting the duplicate give inverse equivalences up to the evident isomorphism.

Use a constant group \(H\) as coefficients, with identity pullback homomorphisms. A cocycle on the duplicated category is specified by its label \(r\in H\) on the arrow \(0\to1\). Its reverse label is \(r^{-1}\), and identity labels are one. A gauge \((a_0,a_1)\) carries \(r\) to
\begin{equation}\label{KS:eq:two-copy-gauge}
 r'=a_0ra_1^{-1}.
\end{equation}
Every \(r\) is isomorphic to the trivial cocycle, for example with gauge \((1,r)\). The automorphisms of \(r\) are the pairs
\[
 (a_0,r^{-1}a_0r),\qquad a_0\in H.
\]
Thus the unframed groupoid has one isomorphism class but retains a copy of \(H\) as the automorphism group of each object. This is compatible with the original one-view coefficient groupoid, which has one object and automorphism group \(H\).

\subsection{Freezing both representative gauges breaks that comparison}
Now require a based gauge to be the identity at every view. On the duplicated category the only permitted gauge is \((1,1)\). Equation~\eqref{KS:eq:two-copy-gauge} then carries \(r\) only to itself. Distinct labels cannot be compared.

Restrict to the selected view zero. Every cocycle becomes the same one-view identity cocycle. The distinction between \(r\) and \(1\) has disappeared in the target, even though the fully framed source has no comparison arrow between them.

\begin{KSproposition}[Full representative framing is not invariant]\label{KS:KS68}
For a nontrivial coefficient group, the restriction from the fully framed duplicated model to the fully framed single-view model need not be full, and hence need not be an equivalence. This failure occurs even though the underlying view categories are equivalent and the full prescribed domains are closed under isomorphism.
\end{KSproposition}
\begin{proof}
Choose \(r\ne1\). In the fully framed duplicated model there is no arrow from \(r\) to \(1\), because the only allowed gauge leaves \(r\) unchanged. After restriction to view zero, both become the same cocycle. There is an identity arrow between their images. That target arrow has no preimage, so the functor is not full. Fullness is necessary for an equivalence.
\end{proof}
\KSsource{KS68}
The failure is not repaired by declaring the two objects equivalent as sets of labels. The issue is exactly the missing arrow. Nor does it contradict the theorem for marked relative realizations. That theorem does not freeze every arbitrary representative coordinate. It transports an actual boundary object together with an identification to it.

\subsection{The marked fiber keeps the comparison that matters}
In the same example, let \(R\) restrict an unframed duplicated cocycle to view zero. Fix the unique boundary object. A marked realization has two pieces of data:
\[
 (r,\phi),\qquad r\in H,\quad\phi\in H.
\]
The element \(\phi\) is the chosen identification of the restricted object with the fixed boundary. It is not discarded just because the boundary has only one object: that object still has many automorphisms.

A comparison from \((r,\phi)\) to \((r',\psi)\) is a gauge \((a_0,a_1)\) satisfying both the cocycle equation and the boundary equation
\[
 r'=a_0ra_1^{-1},\qquad \psi a_0=\phi.
\]
These equations determine it uniquely:
\begin{equation}\label{KS:eq:marked-two-copy}
 a_0=\psi^{-1}\phi,\qquad a_1=(r')^{-1}a_0r.
\end{equation}
Substituting the formulas proves both equations without commuting any factors. Thus in this example there is exactly one identification-preserving arrow between any two marked realizations. The marked groupoid is nonempty and equivalent to the terminal groupoid, although its raw object set has many elements.

This is a worked two-view model of the marked transport mechanism. The generic groupoid construction and its equivalence theorem are Proposition~\ref{KS:KS38} and Theorem~\ref{KS:KS39}. Formula~\eqref{KS:eq:marked-two-copy} is checked in the independent finite suite for \(H=S_3\), retaining the order of multiplication. This finite check is separate from the general formal theorem.

\subsection{Four meanings of uniqueness}
The example separates four statements that are often compressed into the word unique.
\begin{center}
\begin{tabularx}{\linewidth}{@{}p{.28\linewidth}X@{}}
\toprule
Statement & What it retains\\
\midrule
One literal object & No different representative is present in the specified object set.\\
One isomorphism class & Different representatives may remain, as may their automorphisms.\\
A chosen realization & One object is selected, but no uniqueness follows from the choice.\\
One compatible comparison & The comparison respecting fixed identifications is unique, even when many representatives exist.\\
\bottomrule
\end{tabularx}
\end{center}
The unframed duplicated groupoid has one isomorphism class and nontrivial automorphisms. Its marked relative fiber has unique compatible comparisons. Its fully framed representative groupoid can instead have many mutually incomparable objects. These are three different constructions, not three descriptions of the same object with different terminology.

A descent theorem typically uses the last notion: two global realizations of the same specified local datum are related by a unique comparison respecting that datum. It does not force their underlying labels to be literally equal, or erase automorphisms after the local identifications have been forgotten.

\begin{insight}
Boundary data are not only values. In a symmetry-sensitive problem they include the identification with those values. Retaining that identification is what allows a relative realization problem to survive an equivalence of presentations.
\end{insight}



% BEGIN retained input sources/monographs/komorebi-space/source/chapters/04b_phase.tex
\chap{Coherence need not mean removable phase}\label{KS:sec:phase}
The preceding descent theory starts with transition data that already satisfy their composition laws. It is useful to see, in a completely finite calculation, how those laws can hold while a further simplification remains impossible. The following model separates \emph{coherent composition} from \emph{a choice of coordinates that removes every phase}.

The coefficient group and the degree of the data matter. The cocycles in Section~\ref{KS:sec:stack} label individual arrows in a view category. Here a phase is attached to a \emph{pair} of group elements. These are different constructions. The finite example is not obtained by renaming the earlier section obstruction as a general second cohomology group.

\subsection{A two-bit phase law}
Write \(\mathbf F_2=\KSZ/2\KSZ\) additively, and let
\[
 G=\mathbf F_2^2,\qquad x=(x_1,x_2),\qquad
 c(x,y)=x_2y_1.
\]
All entries and operations in this subsection lie in \(\mathbf F_2\). In particular minus and plus agree, although retaining a minus sign in the coboundary formula helps explain its origin.

\begin{KSproposition}[A normalized associative phase]\label{KS:KS80}
The phase vanishes when either argument is zero and satisfies
\begin{equation}\label{KS:eq:phase-cocycle}
 c(x,y)+c(x+y,z)=c(y,z)+c(x,y+z).
\end{equation}
There is no function \(a:G\to\mathbf F_2\) such that
\begin{equation}\label{KS:eq:phase-flatten}
 c(x,y)=a(x)+a(y)-a(x+y)\quad\text{for every }x,y\in G.
\end{equation}
\end{KSproposition}
\begin{proof}
Normalization follows from multiplication by zero. Expanding the two sides of \eqref{KS:eq:phase-cocycle} gives, respectively,
\[
 x_2y_1+(x_2+y_2)z_1,
 \qquad y_2z_1+x_2(y_1+z_1),
\]
which are equal by distributivity.

Every expression \(a(x)+a(y)-a(x+y)\) is symmetric in \(x,y\), because \(G\) is abelian. With \(e_1=(1,0)\) and \(e_2=(0,1)\), however,
\[
 c(e_1,e_2)=0,\qquad c(e_2,e_1)=1.
\]
This contradicts the symmetry required by \eqref{KS:eq:phase-flatten}.
\end{proof}
\KSsource{KS80}
The proof rules out all sixteen possible scalar functions \(a\), not just linear functions or selected trial gauges. The source \KScode{ProjectiveExample.lean} proves the normalization, cocycle, asymmetry, and no-scalar-flattening statements. The matrix interpretation and group construction below are worked derivations of the finite model, not separate formal declarations.

\begin{figure}[htbp]\centering\resizebox{.96\linewidth}{!}{
% BEGIN retained input sources/monographs/komorebi-space/source/figures/30_phase.tex
\includegraphics[width=15cm]{sources/monographs/komorebi-space/source/figures/generated/30_phase.pdf}

}\caption{The entry in row $x$ and column $y$ is $c(x,y)=x_2y_1$. The two entries singled out in the text, $c(e_1,e_2)=0$ and $c(e_2,e_1)=1$, already obstruct every scalar flattening. Associativity still holds for every triple.}\label{KS:fig:phase}\end{figure}

\subsection{The phase is visible in actual matrices}
Over \(\KSR\), define
\[
 X=\begin{pmatrix}0&1\\1&0\end{pmatrix},\qquad
 Z=\begin{pmatrix}1&0\\0&-1\end{pmatrix},\qquad
 U_{(p,q)}=X^pZ^q\quad(p,q\in\{0,1\}).
\]
Direct multiplication gives \(X^2=Z^2=I\) and \(ZX=-XZ\). Moving \(Z^{x_2}\) past \(X^{y_1}\) therefore yields
\begin{equation}\label{KS:eq:phase-matrices}
 U_xU_y=(-1)^{c(x,y)}U_{x+y}.
\end{equation}
For example \(U_{e_1}U_{e_2}=XZ\), while \(U_{e_2}U_{e_1}=-XZ\). The two-bit labels add commutatively, but their chosen matrix lifts do not commute.

Associativity of three lifted matrices requires the same total sign on both parenthesizations. Equation~\eqref{KS:eq:phase-cocycle} is exactly that requirement:
\[
 (U_xU_y)U_z
 =(-1)^{c(x,y)+c(x+y,z)}U_{x+y+z},
\]
\[
 U_x(U_yU_z)
 =(-1)^{c(y,z)+c(x,y+z)}U_{x+y+z}.
\]
Thus a nonzero phase is not automatically a failed associativity law. It can be the data that make the composition law explicit.

Changing lifts to \(U'_x=(-1)^{a(x)}U_x\) changes the phase to
\[
 c'(x,y)=c(x,y)+a(x)+a(y)-a(x+y).
\]
No such change makes \(c'\) vanish identically. The antisymmetric part is unchanged by every scalar change of lift. In this model it is
\[
 \kappa(x,y)=c(x,y)-c(y,x)
            =x_2y_1-y_2x_1.
\]
This is a directly observable obstruction certificate for the model. Its nonzero value at \((e_1,e_2)\) is enough, without a search over gauges.

The same symmetry argument rules out removing this sign multiplier by arbitrary nonzero complex scalar lifts as well. A multiplicative scalar coboundary \(b(x)b(y)/b(x+y)\) on an abelian group is symmetric, whereas the multiplier in \eqref{KS:eq:phase-matrices} is not. This last sentence is an algebraic consequence written here. The retained Lean example uses the stated \(\mathbf F_2\)-valued formulation.

\subsection{Retaining the phase restores an ordinary multiplication law}
Instead of trying to discard the phase, include it in the state. On the eight-element set \(\widetilde G=\mathbf F_2\times G\), put
\begin{equation}\label{KS:eq:central-extension}
 (a,x)\cdot(b,y)=(a+b+c(x,y),x+y).
\end{equation}
The first coordinate of the left-associated product of three elements is
\(a+b+d+c(x,y)+c(x+y,z)\). The right-associated coordinate is
\(a+b+d+c(y,z)+c(x,y+z)\). Equation~\eqref{KS:eq:phase-cocycle} proves equality. Normalization gives identity \((0,0)\), and
\[
 (a,x)^{-1}=(a+c(x,x),x).
\]
Both inverse products are the identity because \(x+x=0\) and arithmetic is modulo two. Hence \eqref{KS:eq:central-extension} defines a group.

Projection \((a,x)\mapsto x\) is a homomorphism onto \(G\). Its kernel is \(\{(0,0),(1,0)\}\), which commutes with every element. The map
\[
 (a,x)\longmapsto(-1)^aU_x
\]
is a homomorphism by \eqref{KS:eq:phase-matrices}. Its eight matrices are distinct: they are \(\pm I,\pm X,\pm Z,\pm XZ\). It is therefore faithful. The phase extension is not an assertion that \(G\) itself has acquired an ordinary representation with these same lifts. The represented group has changed.

\begin{insight}
\textbf{A precise form of retaining the obstruction.} The original four labels do not support a phase-free multiplication of the chosen matrices. Keeping one additional bit gives an eight-element multiplication that does. This repairs the description by retaining information, not by proving that the original phase was a coboundary.
\end{insight}

\subsection{How this example fits the realization calculus}
There are three different questions in this small model. First, the phase satisfies the associativity relation. Second, scalar changes cannot trivialize it. Third, an enlarged state that retains the phase admits an ordinary multiplication law. Each question has a complete answer, but their answers are not interchangeable.

The example illustrates why an obstruction should be attached to the allowed changes and to a specified presentation problem. Enlarging the presentation can be mathematically productive, but it changes what is being realized and what information is carried. No general theorem that every obstruction can be removed by a finite enlargement is claimed. The finite checker enumerates all group triples, all scalar gauges, and all multiplication triples in \(\widetilde G\), independently of the symmetry proof above.




% BEGIN retained input sources/monographs/komorebi-space/source/chapters/05_geometry.tex

\chapter{The geometry of exact realization}\label{KS:book:chapter-22}
A categorical space of realizations does not come with a numerical distance by default. In a declared linear model, the compatibility equations become a bounded operator. Closed range, holonomy, and conditioning then determine what small residuals actually guarantee.

\section{Graph views, holonomy, and reconstruction}\label{KS:sec:graph}
\subsection{Actual transports on a finite graph}
Let a finite connected simple graph have vertex set \(V\), root \(o\), and a chosen route \(p_i:o\to i\) for every vertex, with the route at \(o\) empty. Use a common real normed fiber \(H\). An inverse consistent connection assigns continuous linear equivalences \(\tau_{ij}\), with \(\tau_{ji}=\tau_{ij}^{-1}\). Our convention is that \(\tau_{ij}\) carries the value at \(j\) into the coordinates at \(i\).

The candidate and residual spaces are finite products with maximum norms:
\[
 E=\prod_{i\in V}H,\qquad
 R=\prod_{(i,j)\in\vec E}H,\qquad
 (dx)_{ij}=\tau_{ij}x_j-x_i.
\]
Both orientations of an edge are retained. If \(M\) bounds the norms of all measured edge maps, then \(\KSnorm d\le M+1\). The empty edge case has the usual zero residual.

Distinct walks are retained as arrows of the view category. The product along a walk \(q=(i_0,\ldots,i_m)\) is
\[
 T_q=\tau_{i_0i_1}\cdots\tau_{i_{m-1}i_m}.
\]
Concatenation gives the restriction composition law. An exact vertex family defines a coherent walk-view section, and a coherent walk-view section satisfies every edge equation. Thus this graph construction returns to the original ambient realization type rather than creating a disconnected numerical model.

\subsection{Transported residuals and path gains}
\begin{KSproposition}[The path telescope]\label{KS:KS40}
For a walk \(q=(i_0,\ldots,i_m)\),
\[
 T_qx_{i_m}-x_{i_0}
 =\sum_{k=0}^{m-1}
 (\tau_{i_0i_1}\cdots\tau_{i_{k-1}i_k})(dx)_{i_ki_{k+1}}.
\]
The empty prefix is the identity. If \(\gamma(q)\) is the sum of the products of preceding edge norms, then
\[
 \KSnorm{T_qx_{i_m}-x_{i_0}}\le\gamma(q)\KSnorm{dx}_\infty.
\]
\end{KSproposition}
\begin{proof}
Write \(U_0=I\) and \(U_k=\tau_{i_0i_1}\cdots\tau_{i_{k-1}i_k}\) for \(1\le k\le m\). The transport convention makes every \(U_k\) a map from the fiber at \(i_k\) to that at \(i_0\). Since \(U_{k+1}=U_k\tau_{i_ki_{k+1}}\), the \(k\)-th transported defect is
\[
 U_k(dx)_{i_ki_{k+1}}
 =U_{k+1}x_{i_{k+1}}-U_kx_{i_k}.
\]
Summing over \(k=0,\ldots,m-1\) cancels each interior value once with each sign. The remaining terms are \(U_mx_{i_m}-U_0x_{i_0}=T_qx_{i_m}-x_{i_0}\), proving the identity with its actual ordered products. For an empty walk both sides are zero.

The triangle inequality and the operator norm bound now give
\[
 \KSnorm{T_qx_{i_m}-x_{i_0}}
 \le\sum_{k=0}^{m-1}\KSnorm{U_k}\KSnorm{(dx)_{i_ki_{k+1}}}
 \le\left(\sum_{k=0}^{m-1}\prod_{j=0}^{k-1}\KSnorm{\tau_{i_ji_{j+1}}}\right)\KSnorm{dx}_\infty.
\]
The bracket is \(\gamma(q)\), with the empty product equal to one. Repeated vertices or edges do not affect the cancellation. Reordering the edge maps would, in general, invalidate it.
\end{proof}
\KSsource{KS40}
For isometric maps, \(\gamma(q)\le|q|\). For arbitrary bounded equivalences, route length cannot replace the amplification gain. Reversing a route can have a different gain, even though its transport is the inverse map.

\subsection{Holonomy is the remaining exactness condition}
Write \(T_i=T_{p_i}\). Reconstruction from a proposed root value is
\[
 (R_pz)_i=T_i^{-1}z.
\]
For each directed edge define the root loop operator
\[
 K_{ij}=T_i\tau_{ij}T_j^{-1},\qquad
 K=\bigcap_{(i,j)\in\vec E}\ker(K_{ij}-I).
\]
Each \(K_{ij}\) is the transport around an actual rooted loop. No flatness assumption has been made.

\begin{KStheorem}[Holonomy classification]\label{KS:KS41}
Evaluation at the root and reconstruction give an equivalence
\[
 K\simeq_L\ker d.
\]
A root value extends to an exact vertex family precisely when it lies in \(K\), and that extension is unique. The space \(K\) is independent of the chosen routes. Reconstruction is route independent on \(K\).
\end{KStheorem}
\begin{proof}
Every \(T_i\) is a continuous linear equivalence, being a finite ordered composition of such maps. Reconstruction is therefore well defined. For an edge \((i,j)\), apply \(T_i\) to its defect:
\[
 T_i\bigl(\tau_{ij}T_j^{-1}z-T_i^{-1}z\bigr)
 =T_i\tau_{ij}T_j^{-1}z-z=K_{ij}z-z.
\]
Because \(T_i\) is injective, that edge defect is zero if and only if the corresponding root loop fixes \(z\). Thus \(dR_pz=0\) holds exactly for \(z\in K\). The root route is empty, so \((R_pz)_o=z\).

Conversely, let \(x\in\ker d\). The path telescope along \(p_i\) has all zero summands, hence \(T_ix_i=x_o\) for every vertex. Applying \(T_i^{-1}\) gives \(x_i=T_i^{-1}x_o\), so \(x=R_p(x_o)\). The first part then places \(x_o\) in \(K\). We have proved both inverse equations and uniqueness of an exact extension of a root value.

The maps are linear. On the finite vertex product,
\[
 \KSnorm{R_pz}_\infty\le\left(\max_i\KSnorm{T_i^{-1}}\right)\KSnorm z,
 \qquad \KSnorm{x_o}\le\KSnorm x_\infty.
\]
These bounds also hold on the subspaces \(K\) and \(\ker d\), proving continuity of both directions. No completeness is needed for these finite bounds.

For route independence, the proved equivalence identifies \(K\) with the image of root evaluation from the intrinsically defined space \(\ker d\). That image does not refer to any chosen routes, so every route system yields the same \(K\). For \(z\in K\), reconstruction using either system is an exact family with root \(z\), and uniqueness makes the families equal. For \(z\notin K\), this uniqueness argument does not assert equality of the two generally inconsistent reconstructed families.
\end{proof}
\KSsource{KS41}
With isometric transports the equivalence is an isometry for the vertex maximum norm, because every reconstructed coordinate has the root norm. With general bounded equivalences, it need not be an isometry.

\begin{figure}[htbp]\centering\resizebox{.96\linewidth}{!}{
% BEGIN retained input sources/monographs/komorebi-space/source/figures/09_holonomy.tex
\begin{tikzpicture}[x=1cm,y=1cm]
\node[box,circle] (o) at (-3.6,2.3) {$o$};
\node[box,circle] (i) at (-5.5,-.7) {$i$};\node[box,circle] (j) at (-1.7,-.7) {$j$};
\draw[arr] (i)--node[lab,left] {$T_i$}(o);\draw[arr] (j)--node[lab,right] {$T_j$}(o);
\draw[arr,draw=KomorebiOchre] (j)--node[lab,below] {$\tau_{ij}$}(i);
\draw[arr,dashed] (o) to[loop above,min distance=12mm] node[lab,above] {$K_{ij}=T_i\tau_{ij}T_j^{-1}$}(o);
\node[lab,text width=5.5cm] at (-3.6,-2) {Exact families correspond to\\fixed anchor vectors.};
\begin{scope}[shift={(3.2,.4)}]
 \draw[thinarr](-2.1,0)--(2.6,0) node[right,lab] {$z_1$};
 \draw[thinarr](0,-2)--(0,2.3) node[above,lab] {$z_2$};
 \draw[KomorebiPine,line width=2pt](-2,0)--(2.3,0);
 \node[dot] (p) at (1.4,1.4) {};\node[dot] (r) at (1.4,-1.4) {};
 \draw[arr,dashed](p)--(r);
 \node[lab,anchor=east] at (-.35,1.2) {$H=\operatorname{diag}(1,-1)$};
 \node[lab,left] at (p) {$(u,v)$};\node[lab,left] at (r) {$(u,-v)$};
 \node[lab] at (0,-2.6) {$K=\{(u,0):u\in\KSR\}$};
\end{scope}
\end{tikzpicture}

}\caption{Tree routes reconstruct from an anchor. The remaining edges close loops and impose fixed-vector equations at that anchor. A nontrivial loop can obstruct an arbitrary seed without preventing every nonzero seed from extending. The fixed space, rather than the identity of the whole loop operator, is the relevant realization condition.}\label{KS:fig:holonomy}\end{figure}

\subsection{Prescribed observations remain visible}
If a prescribed vertex set \(U\) contains \(o\), a candidate \(x\) has an exact repair preserving \(x|_U\) precisely when
\[
 x_o\in K\quad\text{and}\quad (R_px_o)|_U=x|_U.
\]
Necessity follows from the classification. Sufficiency is the reconstructed family itself. The source identifies this criterion with the earlier obstruction class in \(R/d(\ker p_U)\). A graph seed test is therefore an instance of the same protected correction problem, not a replacement for it.

\chap{Closed range is the stability boundary}\label{KS:sec:closed}
\subsection{Projection exists before a residual estimate}
Let \(E\) be a real Hilbert space and \(R\) a real Banach space. Let \(D:E\to R\) be continuous linear. The kernel \(K=\ker D\) is closed, so its orthogonal projection \(q=P_K\) exists~\cite{KS:mathlib-projection}.

\begin{KSproposition}[Exact orthogonal repair]\label{KS:KS42}
For every \(x\in E\),
\[
 D(qx)=0,\quad q^2=q,\quad\KSnorm{qx}\le\KSnorm x,
\]
and \(qx\) is the unique closest point to \(x\) in \(K\).
\end{KSproposition}
\begin{proof}
Boundedness of \(D\) makes \(K=\ker D\) closed. Lemma~\ref{KS:lem:projection-tool} constructs \(q=P_K\) and gives the orthogonal decomposition \(x=qx+(x-qx)\). The first term belongs to \(K\), so \(Dqx=0\). Projection fixes every point of \(K\), giving \(q(qx)=qx\).

Pythagoras applied to that decomposition gives \(\KSnorm x^2=\KSnorm{qx}^2+\KSnorm{x-qx}^2\), hence \(\KSnorm{qx}\le\KSnorm x\). For any \(y\in K\), the vector \(qx-y\) lies in \(K\), whereas \(x-qx\) lies in its orthogonal complement. Consequently
\[
 \KSnorm{x-y}^2=\KSnorm{x-qx}^2+\KSnorm{qx-y}^2\ge\KSnorm{x-qx}^2.
\]
Equality holds only when \(\KSnorm{qx-y}=0\), or \(y=qx\). This proves nearest-point uniqueness without any closedness assumption on \(\KSran D\). Closedness of the kernel and of the range serve different purposes.
\end{proof}
\KSsource{KS42}
No closed range assumption on \(D\) is needed for this projection. Such an assumption enters when the size of the residual is asked to control the distance to the exact space.

\begin{KStheorem}[Closed range and quantitative exactness]\label{KS:KS43}
The following are equivalent:
\begin{enumerate}
\item \(\KSran D\) is closed.
\item There is \(C\ge0\) with \(\KSnorm{x-qx}\le C\KSnorm{Dx}\) for all \(x\).
\item There is \(C\ge0\) such that every \(x\) has some exact \(y\) with \(\KSnorm{x-y}\le C\KSnorm{Dx}\).
\end{enumerate}
When the range is closed, the least admissible constant is the operator norm of the inverse of \(D:K^\perp\to\KSran D\). In the zero operator case that inverse has zero dimensional domain and its norm is zero.
\end{KStheorem}
\begin{proof}
Let $K=\ker D$ and $K^\perp$ carry the induced Hilbert norm. Define $D_\perp u=Du$, but regard its target as $\KSran D$ with the norm inherited from $R$. It is injective: $Du=0$ puts $u$ in both $K$ and $K^\perp$, so $u=0$. It is surjective onto that target: if $r=Dx$, then $u=x-qx\in K^\perp$ and $Du=Dx=r$. Thus it is a bounded linear bijection before any continuity of its inverse is claimed.

Suppose first that the range is closed. Because $R$ is Banach, $\KSran D$ is then Banach. The domain $K^\perp$ is also complete. The bounded inverse theorem, proved in Appendix~\ref{KS:app:proof-tools}, gives a bounded inverse $A=D_\perp^{-1}$. Applying it to $Dx=D(x-qx)$ yields
\[
 x-qx=A(Dx),\qquad \KSnorm{x-qx}\le\KSnorm A\KSnorm{Dx}.
\]
Hence condition 2 holds. If any $C\ge0$ satisfies condition 2, then for every $r\in\KSran D$ apply it to $u=Ar\in K^\perp$. Since $qu=0$, it gives $\KSnorm{Ar}\le C\KSnorm r$. Taking the operator norm proves $\KSnorm A\le C$. This establishes the least constant, not just a sufficient one. If $D=0$, both $K^\perp$ and the range are zero spaces and $\KSnorm A=0$.

Now assume condition 2. To prove closedness, take an arbitrary sequence $r_n\in\KSran D$ with $r_n\to r$ in $R$. Choose $x_n$ with $Dx_n=r_n$ and put $u_n=x_n-qx_n$. Then $u_n\in K^\perp$ and $Du_n=r_n$. Applying the estimate to $u_n-u_m$ gives
\[
 \KSnorm{u_n-u_m}\le C\KSnorm{r_n-r_m}.
\]
Thus $(u_n)$ is Cauchy in the complete space $K^\perp$. Let its limit be $u$. Continuity gives $Du=\lim_n r_n=r$, so $r\in\KSran D$. In a normed space, a set containing the limits of all its convergent sequences is closed. This proves condition 1 without presupposing that the initially unbounded inverse is continuous.

Condition 2 implies condition 3 by choosing $y=qx$. Conversely, if condition 3 supplies an exact $y$ for a given $x$, the nearest-point property of $qx$ gives
\[
 \KSnorm{x-qx}\le\KSnorm{x-y}\le C\KSnorm{Dx}.
\]
The same $C$ works for all $x$, proving condition 2. This last implication does not require the selected answers $y$ to depend linearly or continuously on $x$.
\end{proof}
\KSsource{KS43}
The last equivalence is important. A nonlinear choice of exact answers cannot bypass the failure of a uniform residual bound. Failure of closed range is not merely a defect of one selected projection algorithm.

\subsection{A small residual far from exactness}
On \(\ell^2(\KSN)\), index the standard basis from one and define \(De_n=n^{-1}e_n\). The map is bounded and injective, so its exact space is zero. Yet
\[
 \KSnorm{e_n}=1,\qquad\KSnorm{De_n}=1/n\longrightarrow0.
\]
No finite \(C\) can bound \(\KSdist(e_n,\ker D)\) by \(C\KSnorm{De_n}\). The range is not closed. For example finite truncations of \((1/n)_{n\ge1}\) belong to the range and converge in \(\ell^2\), but a preimage of the limit would be the constant sequence, which does not belong to \(\ell^2\).

This diagonal model is formalized on the infinite Hilbert space used in Section~\ref{KS:sec:diagonal-regularization}. In the namespace \path{KomorebiSpace.DiagonalRegularization}, \path{T_not_closed_range} applies the general obstruction theorem to actual unit coordinate vectors. The declaration \path{w_mem_closure_not_range} proves that the displayed limiting datum belongs to the closure of the range but not to the range itself. Formal coordinate $j$ represents the printed coordinate $n=j+1$.

\begin{figure}[htbp]\centering\resizebox{.96\linewidth}{!}{
% BEGIN retained input sources/monographs/komorebi-space/source/figures/10_closed_range.tex
\includegraphics[width=15cm]{sources/monographs/komorebi-space/source/figures/generated/10_closed_range_plot.pdf}

}\caption{Residual size controls distance to the exact space only when transverse directions have a positive lower bound. The diagonal example sends unit vectors to residuals tending to zero, so no uniform inverse gain exists. Exact projection still exists.}\label{KS:fig:closed}\end{figure}

\subsection{Norm zero versus algebraic zero}
If an additive correction image is not closed, distance to that image can vanish without actual membership. Thus an infimum of residual error can be zero while the exact protected correction is unattained. The algebraic obstruction quotient and its possible seminorm answer different questions. One must identify a closed image, compactness condition, or other attainment principle before replacing one question by the other.

\subsection{Stable repair with protected data}
The preceding orthogonal projection does not generally preserve a separately declared observable. To join the analytic theory to Section~\ref{KS:sec:protected}, let \(p:E\to S\) be continuous linear into a normed space. Define
\[
 M=\ker p,\qquad B=D|_M:M\to R,\qquad H=M\cap\ker D.
\]
The closed subspace \(M\) inherits its Hilbert structure from \(E\). A candidate \(x\) is \emph{feasible} for protected repair when an exact \(y\) with \(p(y)=p(x)\) exists. The algebraic criterion is still \(Dx\in\KSran B\).

\begin{KStheorem}[Nearest feasible protected repair]\label{KS:KS56}
Every feasible \(x\) has a unique closest exact answer \(y_*\) preserving \(p(x)\). Its correction \(h_*=x-y_*\) is the unique vector in the orthogonal complement of \(\ker B\) inside \(M\) satisfying \(Bh_*=Dx\). For every other feasible answer \(y\),
\[
 \KSnorm{x-y}^2=\KSnorm{x-y_*}^2+\KSnorm{y_*-y}^2.
\]
Moreover, the following conditions are equivalent:
\[
 \KSran B\text{ is closed}
 \quad\Longleftrightarrow\quad
 \exists C\ge0\ \forall\text{ feasible }x:
 \KSnorm{x-y_*}\le C\KSnorm{Dx}.
\]
When these conditions hold, the least nonnegative constant is the transverse inverse norm of \(B\), including the zero-range case.
\end{KStheorem}
\begin{proof}
Feasibility gives some $h_0=x-y_0\in M$ with $Bh_0=Dx$. Conversely, every $h\in M$ satisfying $Bh=Dx$ produces an exact protected answer $y=x-h$. Thus the answer problem and the correction problem are equivalent, with correction set
\[
 \{h\in M:Bh=Dx\}=h_0+\ker B.
\]
This is a nonempty closed affine subset of the Hilbert space $M$. Closedness of $\ker B$ follows from continuity of $B$, independently of whether its image is closed.

Let $k_0=P^M_{\ker B}h_0$ and $h_*=h_0-k_0$. Then $Bh_*=Bh_0=Dx$ and $h_*\perp\ker B$ in $M$. If a different initial correction is chosen, it differs from $h_0$ by an element of $\ker B$. Its orthogonal-complement component is therefore again $h_*$. This proves that the construction is independent of the feasibility witness.

Every other correction can be written uniquely as $h_*+k$ with $k\in\ker B$. Orthogonality gives
\[
 \KSnorm{h_*+k}^2=\KSnorm{h_*}^2+\KSnorm k^2.
\]
Writing $y_*=x-h_*$ and $y=x-h_*-k$, we have $x-y=h_*+k$ and $y_*-y=k$. Substitution proves the stated Pythagorean identity. Its last term vanishes only when $y=y_*$, establishing the unique closest answer.

If $\KSran B$ is closed, Theorem~\ref{KS:KS43} applied on $M$ gives a bounded inverse of $B$ on $(\ker B)^{\perp_M}$. The vector $h_*$ is exactly that inverse applied to $Dx$, so its norm obeys the required bound.

For the converse, suppose the bound holds for every feasible raw candidate. Every $x\in M$ is feasible, because $y=0$ has $Dy=0$ and $py=px=0$. Its closest protected exact answer is $P^M_{\ker B}x$. Thus the assumed estimate reduces on $M$ to
\[
 \KSnorm{x-P^M_{\ker B}x}\le C\KSnorm{Bx}\qquad(x\in M).
\]
Theorem~\ref{KS:KS43} forces $\KSran B$ to be closed. The same restriction also proves that no smaller nonnegative constant than the transverse inverse norm can work. If the range is zero, feasible candidates are already exact under permitted correction and the least correction is zero, so the optimal gain is zero.
\end{proof}
\KSsource{KS56}
The relevant stability boundary is therefore \(D(\ker p)\), not automatically \(D(E)\). Restricting permitted changes can turn a well-conditioned unconstrained problem into an obstructed or poorly conditioned protected one. The nearest protected answer is generally not \(P_{\ker D}x\).

For \(D(x_1,x_2,x_3)=x_1-x_2\), \(p(x)=(x_1,x_3)\), and \(x=(2,0,5)\), the ordinary orthogonal answer is \((1,1,5)\). The protected answer is \((2,2,5)\). The respective squared errors are two and four. Both have zero residual, but only the latter preserves the declared data. The theorem is formalized in \KScode{ProtectedHilbertRepair.lean}, using the actual operator on \(\ker p\). The displayed three-coordinate values are an illustrative calculation, not a separate universal algorithm.

\begin{figure}[htbp]\centering\resizebox{.96\linewidth}{!}{
% BEGIN retained input sources/monographs/komorebi-space/source/figures/15_protected_geometry.tex
\includegraphics[width=15cm]{sources/monographs/komorebi-space/source/figures/generated/15_protected_geometry.pdf}

}\caption{The same compatibility equation permits different repair problems. Orthogonal repair chooses the nearest point on the diagonal. Preserving the first coordinate forces a vertical correction to a different point. The third coordinate is fixed at five throughout.}\label{KS:fig:protectedgeometry}\end{figure}

\chap{Holonomy repair and transport amplification}\label{KS:sec:holonomyrepair}
\subsection{An orthogonal anchor, not a global least squares answer}
Return to the graph model with a real Hilbert fiber. The root residual is
\[
 \mathcal H z=(K_{ij}z-z)_{ij},\qquad K=\ker\mathcal H.
\]
Let \(q=P_K\), and define
\[
 Px=R_p(qx_o).
\]
This is always exact, fixes exact families, and is idempotent. It is the unique exact family preserving the feasible component \(qx_o\) of the anchor. Its orthogonality is on the anchor fiber. The vertex and edge norms are still maximum norms.

For example on a two-vertex identity connection, \((0,2)\) with root value zero is repaired to \((0,0)\). The constant family \((1,1)\) is closer in the maximum norm, but it does not preserve the anchor. The operator \(P\) is not being claimed to be a global nearest-point projection in that norm.

\subsection{Every factor in the bound has a meaning}
Define the route gain
\[
 A_p=\max_i\gamma(\overline{p_i}),
\]
using reverse routes so that the error is measured in each destination view. The path telescope gives
\[
 \KSnorm{x-R_p(x_o)}_\infty\le A_p\KSnorm{dx}_\infty.
\]
Let \(N_p\) be the residual normalization \((N_pr)_{ij}=T_ir_{ij}\). The exact identity \(\mathcal H=N_pdR_p\) implies
\[
 \KSnorm{\mathcal H(x_o)}_\infty
 \le\KSnorm{N_p}(\KSnorm d A_p+1)\KSnorm{dx}_\infty.
\]

\begin{KStheorem}[Amplified holonomy repair]\label{KS:KS44}
If \(\KSran\mathcal H\) is closed and \(\kappa\) is its transverse inverse gain, then
\begin{equation}\label{KS:eq:graphgain}
 \KSnorm{x-Px}_\infty\le
 \left[A_p+\KSnorm{R_p}\kappa\KSnorm{N_p}(\KSnorm d A_p+1)\right]
 \KSnorm{dx}_\infty.
\end{equation}
The range condition is necessary for the existence of any uniform bound of distance to an exact family by \(\KSnorm{dx}_\infty\). It is automatic for a fixed finite dimensional fiber, but not uniformly quantitative across all connections.
\end{KStheorem}
\begin{proof}
First quantify the residual seen at the root. The identity \(\mathcal H=N_pdR_p\), applied to \(x_o\), gives
\[
 \begin{aligned}
 \KSnorm{\mathcal H(x_o)}_\infty
 &=\KSnorm{N_p\bigl(d(R_px_o-x)+dx\bigr)}_\infty\\
 &\le\KSnorm{N_p}\bigl(\KSnorm d\KSnorm{R_px_o-x}_\infty+\KSnorm{dx}_\infty\bigr)\\
 &\le\KSnorm{N_p}(\KSnorm d A_p+1)\KSnorm{dx}_\infty.
 \end{aligned}
\]
Closed range of \(\mathcal H\) lets Theorem~\ref{KS:KS43} control the anchor correction by
\(\KSnorm{x_o-qx_o}\le\kappa\KSnorm{\mathcal H(x_o)}_\infty\). Insert the uncorrected reconstruction between \(x\) and \(Px=R_pqx_o\):
\[
 \KSnorm{x-Px}_\infty
 \le\KSnorm{x-R_px_o}_\infty+\KSnorm{R_p}\KSnorm{x_o-qx_o}.
\]
Substituting the two preceding bounds yields \eqref{KS:eq:graphgain}. The finite maximum-norm residual product is a Banach space because its fiber is Hilbert, so the target completeness needed in Theorem~\ref{KS:KS43} is available.

For necessity, suppose a constant \(C\) satisfies
\(\KSdist(w,\ker d)\le C\KSnorm{dw}_\infty\) for every vertex family \(w\). Fix \(z\in K^\perp\) and take \(w=R_pz\). Every exact family \(e\) has \(e_o\in K\), and root evaluation has norm at most one. Therefore
\[
 \KSnorm z=\KSdist(z,K)\le\KSnorm{z-e_o}\le\KSnorm{R_pz-e}_\infty.
\]
Taking the infimum over \(e\in\ker d\) is legitimate without assuming that this global nearest point is attained. It gives
\[
 \KSnorm z\le C\KSnorm{dR_pz}_\infty
 \le C\KSnorm{N_p^{-1}}\KSnorm{\mathcal Hz}_\infty.
\]
Here \(N_p\) is a continuous linear equivalence on the finite edge product because every component \(T_i\) is one. For a general root vector, remove its projection onto \(K\); this leaves the same \(\mathcal H\)-image and puts it in \(K^\perp\). Thus the displayed estimate is a uniform residual bound of the kind in Theorem~\ref{KS:KS43}, which forces \(\KSran\mathcal H\) to be closed. A finite-dimensional fiber makes that image finite dimensional and hence closed. It does not give a uniform numerical value of the inverse gain across different connections.
\end{proof}
\KSsource{KS44}
For isometric transports, one can use \(\KSnorm{R_p}\le1\), \(\KSnorm{N_p}\le1\), \(\KSnorm d\le2\), and \(A_p\le r\), where \(r\) is the route radius. The bound reduces to
\[
 \KSnorm{x-Px}_\infty\le[r+\kappa(2r+1)]\KSnorm{dx}_\infty.
\]
The exact repair is route independent for a fixed root and inner product. A chosen upper bound involving route gains need not have the same numerical value for different routes.

\subsection{Flatness does not remove amplification}
Consider a scalar chain with \(n+1\) vertices, where \(n\ge1\), and coordinates \(e_i(z)=2^iz\). Set \(\tau_{ij}=e_ie_j^{-1}\). This connection is flat, but reverse transport along one edge multiplies norms by two. Let
\[
 x_i=8\cdot2^i-1,\qquad y_i=7\cdot2^i.
\]
The family \(y\) is exact and preserves \(x_o=7\). On a forward edge the residual of \(x\) is \(1/2\), and on its reverse it is \(-1\). Thus
\[
 \KSnorm{dx}_\infty=1,\qquad
 \KSnorm{x-Px}_\infty=2^n-1.
\]
For \(n\ge2\), a route length bound of \(n\) fails. This is amplification with trivial holonomy, distinct from poor conditioning of the fixed space of a nearly trivial loop.
\KSsource{KS45}
\begin{figure}[htbp]\centering\resizebox{.96\linewidth}{!}{
% BEGIN retained input sources/monographs/komorebi-space/source/figures/17_amplification.tex
\includegraphics[width=15cm]{sources/monographs/komorebi-space/source/figures/generated/17_amplification.pdf}

}\caption{A flat connection can still amplify a fixed residual. Both orientations of every edge are measured, giving residual norm one. Root-preserving repair has error \(2^n-1\), which exceeds the route-length estimate once \(n\ge2\). The vertical scale is logarithmic.}\label{KS:fig:amplification}\end{figure}

\subsection{Small holonomy can still be poorly conditioned}
In \(\KSR^2\), the rational rotation
\[
 U_t=\frac1{1+t^2}
 \begin{pmatrix}1-t^2&-2t\\2t&1-t^2\end{pmatrix}
\]
satisfies
\[
 \KSnorm{(U_t-I)z}^2=\frac{4t^2}{1+t^2}\KSnorm z^2.
\]
For \(t\ne0\) its fixed space is zero and its residual gain is \(\sqrt{1+t^2}/(2|t|)\), which diverges as \(t\to0\). At \(t=0\), the entire fiber is fixed. Therefore small loop discrepancy does not supply a uniform stable realization estimate. The fixed directions and their conditioning must be analyzed together.

\chap{Quantitative descent in a common Hilbert model}\label{KS:sec:quantitative}
\subsection{Local exactness and matching are separate operators}
Let \(E\) be a real Hilbert space of raw local candidates. Let \(D:E\to R\) and \(C:E\to Y\) be continuous linear maps into Banach spaces, measuring local equation defects and overlap mismatch, respectively. Put
\[
 L=\ker D,\qquad B=C|_L,
 \qquad M=\ker D\cap\ker C.
\]
The exact matched families form \(M\). First project to \(L\), then project inside \(L\) to \(\ker B\). Denote the resulting map by \(P\).

\begin{KStheorem}[The two-stage repair is the joint projection]\label{KS:KS46}
The operator \(P\) is the orthogonal projection onto \(M\), equivalently onto the kernel of \(\Delta=(D,C)\). If \(u=x-P_Lx\) and \(v=P_Lx-Px\), then
\[
 u\perp v,\qquad\KSnorm{x-Px}^2=\KSnorm u^2+\KSnorm v^2.
\]
\end{KStheorem}
\begin{proof}
The subspace \(L=\ker D\) is closed, hence Hilbert with the inherited inner product. The restriction \(B=C|_L\) is continuous, so its kernel is a closed subspace of \(L\). Under the inclusion into \(E\), this kernel is exactly \(M=L\cap\ker C\). Thus both projections used to define \(P\) exist.

For an input \(x\), write \(\ell=P_Lx\) and let \(Px\) be the projection of \(\ell\) onto \(\ker B\) inside \(L\). Then \(Px\in M\), \(u=x-\ell\in L^\perp\), and \(v=\ell-Px\in L\) is orthogonal to \(M\). Since \(M\subseteq L\), the vector \(u\) is also orthogonal to \(M\). Their sum \(x-Px=u+v\) is therefore orthogonal to \(M\). A vector in \(M\) with that error is the unique orthogonal projection of \(x\), so \(P=P_M\).

The two errors themselves are orthogonal because \(u\in L^\perp\) and \(v\in L\). Expanding the squared norm of their sum gives \(\KSnorm{x-Px}^2=\KSnorm u^2+\KSnorm v^2\). Finally \(\ker(D,C)=\ker D\cap\ker C=M\), so this is also the projection onto the kernel of the joint defect. The proof uses a second projection \emph{inside} \(L\). Replacing it by the ambient projection onto \(\ker C\) would invalidate the membership step and need not preserve local exactness.
\end{proof}
\KSsource{KS46}
This identity does not require the local and overlap operators to commute. The second projection is taken in the already exact subspace, not in the whole raw space without restriction.

For a concrete failure of the unrestricted procedure, take \(L=\{(a,0)\}\) and \(\ker C=\{(a,a)\}\) in \(\KSR^2\). Their intersection is zero. Starting from \((1,0)\), projection to \(L\) does nothing, while subsequent projection to \(\ker C\) gives \((1/2,1/2)\), which is no longer in \(L\). The restricted second stage instead gives the correct joint answer zero.

\begin{figure}[htbp]\centering\resizebox{.96\linewidth}{!}{
% BEGIN retained input sources/monographs/komorebi-space/source/figures/11_descent.tex
\includegraphics[width=15cm]{sources/monographs/komorebi-space/source/figures/generated/11_nested_projection_3d.pdf}

}\caption{Local correction and overlap correction are orthogonal errors when the second problem is restricted to the locally exact subspace. Their composition is the projection onto the joint exact space. The overlap range inside that subspace is the additional stability obligation.}\label{KS:fig:quantdescent}\end{figure}

\subsection{A stable repair requires a stable matching problem}
Suppose \(\KSran D\) and \(\KSran B\) are closed, with residual gains \(a,b\). Then
\[
 \KSnorm u\le a\KSnorm{Dx},\qquad
 \KSnorm v\le b\KSnorm{CP_Lx}
 \le b(\KSnorm{Cx}+\KSnorm C\,a\KSnorm{Dx}).
\]
Therefore
\begin{equation}\label{KS:eq:jointgain}
 \KSnorm{x-Px}\le
 a(1+b\KSnorm C)\KSnorm{Dx}+b\KSnorm{Cx}.
\end{equation}
The Pythagorean identity gives the sharper estimate
\[
 \KSnorm{x-Px}^2\le
 a^2\KSnorm{Dx}^2+b^2\bigl(\KSnorm{Cx}+a\KSnorm C\KSnorm{Dx}\bigr)^2.
\]
For the maximum norm on \(R\times Y\), two sufficient gains are therefore
\[
 G_2=\sqrt{a^2+b^2(1+a\KSnorm C)^2}
 \le G_1=a(1+b\KSnorm C)+b.
\]
The inequality follows from \(\sqrt{u^2+v^2}\le u+v\) for nonnegative \(u,v\). The formal interface supplies the sum-type gain \(G_1\), while its squared-error theorem also yields \(G_2\). Neither is claimed to be the least joint gain. In either bound, \(\KSnorm C\) may be replaced by any certified bound on \(C\) restricted to \((\ker D)^\perp\).

\begin{KStheorem}[Closed range for the joint defect]\label{KS:KS47}
Assume \(\KSran D\) is closed. Then
\[
 \KSran(D,C)\text{ is closed}\quad\Longleftrightarrow\quad
 \KSran(C|_{\ker D})\text{ is closed}.
\]
Under either condition, a finite joint residual-to-repair bound exists.
\end{KStheorem}
\begin{proof}
Suppose the range of the joint map is closed. To test closedness of $\KSran B$, take a sequence $Bk_n\to y$ with $k_n\in\ker D$. The joint images are
\[
 (Dk_n,Ck_n)=(0,Bk_n)\longrightarrow(0,y).
\]
Closedness supplies $x\in E$ with $(Dx,Cx)=(0,y)$. The first coordinate gives $x\in\ker D$, and the second gives $Bx=y$. Hence the matching range is closed. This direction did not require closedness of $\KSran D$.

For the reverse direction, assume both $\KSran D$ and $\KSran B$ are closed, and take joint images converging to $(r,y)$. Thus $Dx_n\to r$ and $Cx_n\to y$. Decompose each input orthogonally as
\[
 x_n=u_n+k_n,\quad
 u_n\in(\ker D)^\perp,\quad k_n\in\ker D.
\]
Because $r\in\KSran D$ and the transverse inverse $A$ of $D$ is bounded by Theorem~\ref{KS:KS43}, the first components determine
\[
 u_n=A(Dx_n)\longrightarrow A(r)=u,
 \qquad Du=r.
\]
Continuity of $C$ now gives
\[
 Bk_n=Ck_n=Cx_n-Cu_n\longrightarrow y-Cu.
\]
Closedness of $\KSran B$ supplies $k\in\ker D$ with $Bk=y-Cu$. There is no need for the sequence $k_n$ itself to converge. The single vector $u+k$ satisfies
\[
 D(u+k)=r,
 \qquad C(u+k)=Cu+Bk=y.
\]
Thus the joint limit lies in the joint range, proving closedness. The target product is Banach under its maximum norm, so Theorem~\ref{KS:KS43} applies to the joint map and gives a finite repair gain. The earlier two-stage estimate supplies an explicit, though not necessarily least, value.
\end{proof}
\KSsource{KS47}


This Hilbert model is a quantitative counterpart of gluing, but it is not a norm on an arbitrary categorical stack. It starts from a fixed raw candidate space, fixed norms, and actual local and overlap operators. Changing those choices requires its own comparison.

\chap{Controlled motion of exact spaces}\label{KS:sec:motion}
\subsection{A two-sided condition number theorem}
Let \(D,D':E\to R\) be continuous linear maps from a real Hilbert space to the same normed target. Write \(q=P_{\ker D}\), \(q'=P_{\ker D'}\). Suppose
\[
 \KSnorm{(I-q)x}\le C\KSnorm{Dx},\qquad
 \KSnorm{(I-q')x}\le C'\KSnorm{D'x},\qquad C,C'\ge0,
\]
and \(\KSnorm{D-D'}\le\varepsilon\).

\begin{KStheorem}[Motion of the exact projector]\label{KS:KS48}
Without either kernel inclusion as a hypothesis,
\begin{equation}\label{KS:eq:motion}
 \KSnorm{q-q'}\le\max(C,C')\varepsilon.
\end{equation}
In particular a uniform residual gain \(M\) makes \(D\mapsto P_{\ker D}\) \(M\)-Lipschitz on the class of operators carrying that gain.
\end{KStheorem}
\begin{proof}
For $z\in\ker D'$, the old residual bound gives
\[
 \KSnorm{(I-q)z}\le C\KSnorm{Dz}
 =C\KSnorm{(D-D')z}\le C\varepsilon\KSnorm z.
\]
Interchanging the roles gives
\[
 \KSnorm{(I-q')w}\le C'\varepsilon\KSnorm w
 \qquad(w\in\ker D).                                      \tag{*}
\]
We need a dual form of (*). For $v\in(\ker D')^\perp$ and any $y\in E$, self-adjointness of $q$ and $v=(I-q')v$ give
\[
 |\langle qv,y\rangle|
 =|\langle v,qy\rangle|
 =|\langle v,(I-q')qy\rangle|
 \le C'\varepsilon\KSnorm v\KSnorm{qy}
 \le C'\varepsilon\KSnorm v\KSnorm y.
\]
The Hilbert norm equals the supremum of these pairings on the unit ball, so
$\KSnorm{qv}\le C'\varepsilon\KSnorm v$. This proves the precise restriction estimate needed on $(I-q')x$, not just a coarser bound in terms of $\KSnorm x$.

Now write
\[
 (q-q')x=q(I-q')x-(I-q)q'x.
\]
The two displayed terms lie in $\ker D$ and $(\ker D)^\perp$, respectively, so they are orthogonal. The input components $(I-q')x$ and $q'x$ are also orthogonal. Applying the dual estimate to the first and the original cross estimate to the second gives
\[
\begin{aligned}
 \KSnorm{(q-q')x}^2
 &=\KSnorm{q(I-q')x}^2+\KSnorm{(I-q)q'x}^2\\
 &\le (C'\varepsilon)^2\KSnorm{(I-q')x}^2
      +(C\varepsilon)^2\KSnorm{q'x}^2\\
 &\le\varepsilon^2\max(C,C')^2
       \bigl(\KSnorm{(I-q')x}^2+\KSnorm{q'x}^2\bigr)\\
 &=\varepsilon^2\max(C,C')^2\KSnorm x^2.
\end{aligned}
\]
The bound $\KSnorm{D-D'}\le\varepsilon$ implies $\varepsilon\ge0$, so taking square roots yields the operator estimate. If the common gain is $M$, substitute $C=C'=M$ and $\varepsilon=\KSnorm{D-D'}$. This is exactly the Lipschitz inequality on the class carrying that gain. The use of two orthogonal decompositions explains the maximum coefficient instead of a sum.
\end{proof}
\KSsource{KS48}
There is no smallness threshold when both certificates are already supplied. The estimate can be uninformative when their gains are large. Knowing only the old certificate is a different situation.

For changing input as well as model, the same theorem gives
\[
 \KSnorm{qx-q'y}\le\KSnorm{x-y}+\max(C,C')\varepsilon\KSnorm y.
\]
It also bounds the change in distance to the exact space by \(\max(C,C')\varepsilon\KSnorm x\).

\subsection{Constructing the new certificate}
Assume the target \(R\) is Banach and \(A:R\to E\) is a continuous linear right inverse of \(D\), with \(DA=I\), \(\KSnorm A\le a\), and \(a\varepsilon<1\). Set
\[
 T=(D-D')A,\qquad Z=(I-T)^{-1},\qquad A'=AZ.
\]
The inverse is supplied by the convergent Neumann series because \(\KSnorm T\le a\varepsilon<1\).

\begin{KStheorem}[Stable surjectivity and derived repair]\label{KS:KS49}
The perturbed map is surjective, with
\[
 D'A'=I,\qquad \KSnorm{A'}\le\frac{a}{1-a\varepsilon}.
\]
Its exact projection satisfies
\[
 \KSnorm{x-q'x}\le\frac{a}{1-a\varepsilon}\KSnorm{D'x},\qquad
 \KSnorm{q-q'}\le\frac{a\varepsilon}{1-a\varepsilon}.
\]
\end{KStheorem}
\begin{proof}
The map $T=(D-D')A$ acts on the Banach target $R$ and satisfies $\KSnorm T\le a\varepsilon<1$. The operator series
\[
 Z=\sum_{n=0}^\infty T^n
\]
converges in operator norm and obeys $\KSnorm Z\le(1-a\varepsilon)^{-1}$. Multiplying its finite partial sums on either side by $I-T$ gives $I-T^{N+1}$. Passing to the operator-norm limit proves $(I-T)Z=Z(I-T)=I$. Appendix~\ref{KS:app:proof-tools} supplies the full series argument, including completeness of the operator space.

Since $DA=I$, one has $D'A=DA-(D-D')A=I-T$. Therefore $A'=AZ$ satisfies
\[
 D'A'=(I-T)Z=I,\qquad
 \KSnorm{A'}\le\KSnorm A\KSnorm Z\le\frac a{1-a\varepsilon}.
\]
The first equation explicitly supplies a preimage of every target vector, proving surjectivity.

For any $x$, put $y=x-A'D'x$. Then $D'y=D'x-D'A'D'x=0$, so $y$ is a valid exact competitor for orthogonal projection. Its error gives
\[
 \KSnorm{x-q'x}\le\KSnorm{x-y}=\KSnorm{A'D'x}
 \le\frac a{1-a\varepsilon}\KSnorm{D'x}.
\]
The same argument using $A$ gives $\KSnorm{x-qx}\le a\KSnorm{Dx}$. Both coefficients are nonnegative and $a/(1-a\varepsilon)\ge a$. Substitution into Theorem~\ref{KS:KS48} yields
\[
 \KSnorm{q-q'}\le\frac{a\varepsilon}{1-a\varepsilon}.
\]
No assertion that $I-A'D'$ is orthogonal was needed. It supplies an exact comparison vector, while best approximation supplies the orthogonal bound.
\end{proof}
\KSsource{KS49}
If a measured residual \(y\) satisfies \(\KSnorm{D'x-y}\le\eta\), the distance to the new exact space is at most \(a(\KSnorm y+\eta)/(1-a\varepsilon)\). The uncertainty radius remains a quantity to justify in the application. The theorem does not estimate it statistically.

For a surjective Hilbert-source operator, the right inverse need not be a separate unexplained input. The inverse of \(D:(\ker D)^\perp\to R\), followed by inclusion, constructs it. A known residual certificate bounds its norm. Surjectivity is still substantive, and cannot be assumed for a redundant residual map whose target contains unattainable directions.

\subsection{An equivalence between the moving kernels}
\begin{KStheorem}[Explicit kernel transport]\label{KS:KS50}
Under the hypotheses of Theorem~\ref{KS:KS49},
\[
 J:\ker D\simeq_L\ker D',\qquad Jx=x-A'D'x,
\]
with inverse \(y\mapsto y-ADy\). Moreover
\[
 \KSnorm{Jx-x}\le\frac{a\varepsilon}{1-a\varepsilon}\KSnorm x.
\]
\end{KStheorem}
\begin{proof}
Let $R_0=I-AD$ and $R_1=I-A'D'$. From $DA=D'A'=I$ we obtain $DR_0=0$ and $D'R_1=0$. Each $R_i$ fixes its corresponding kernel. Thus $R_1$ restricts to a map from $\ker D$ to $\ker D'$, while $R_0$ restricts in the other direction.

The inverse identities require two additional cancellations. First,
\[
 R_0A'=(I-AD)AZ=(A-ADA)Z=0.
\]
Second, using $D'A=I-T$ and $Z(I-T)=I$,
\[
 R_1A=A-AZD'A=A-AZ(I-T)=0.
\]
For $x\in\ker D$, these identities give
\[
 R_0R_1x=R_0(x-A'D'x)=R_0x=x.
\]
For $y\in\ker D'$, they give
\[
 R_1R_0y=R_1(y-ADy)=R_1y=y.
\]
Both restricted maps are bounded and linear as restrictions of bounded linear operators. They are therefore inverse continuous linear maps, establishing the asserted equivalence.

Finally, if $Dx=0$, then $D'x=(D'-D)x$. Hence
\[
 \KSnorm{Jx-x}=\KSnorm{A'D'x}
 \le\KSnorm{A'}\varepsilon\KSnorm x
 \le\frac{a\varepsilon}{1-a\varepsilon}\KSnorm x.
\]
This derivation uses the selected right inverse $A$ throughout. It proves a controlled transport of exact vectors, not a claim that the selected map is the nearest-point projection onto the new kernel.
\end{proof}
\KSsource{KS50}
The map \(J\) depends on the selected splitting. It is not generally the restriction of the orthogonal projection \(q'\). Equivalence of exact spaces and nearest-point repair answer different questions.

\subsection{An all-real worked model}
Let \(D_t(x,y)=x+ty\). Its kernel is the line \(\{(-tb,b):b\in\KSR\}\), and its orthogonal projector is
\begin{equation}\label{KS:eq:lineprojector}
 q_t=\frac1{1+t^2}
 \begin{pmatrix}t^2&-t\\-t&1\end{pmatrix}.
\end{equation}
The common right inverse \(r\mapsto(r,0)\) supplies gain at most one for every \(t\), and \(\KSnorm{D_s-D_t}=|s-t|\). The general theorem gives \(\KSnorm{q_s-q_t}\le|s-t|\).

\begin{KSproposition}[Exact motion of the line]\label{KS:KS51}
For every real \(t\),
\[
 \KSnorm{(q_t-q_0)x}^2=\frac{t^2}{1+t^2}\KSnorm x^2,
 \qquad \KSnorm{q_t-q_0}=\frac{|t|}{\sqrt{1+t^2}}.
\]
In the small perturbation interface, kernel transport sends \((0,b)\) to \((-tb,b)\). Orthogonal projection instead sends it to \((-tb/(1+t^2),b/(1+t^2))\).
\end{KSproposition}
\begin{proof}
The vector \(v_t=(-t,1)\) spans \(\ker D_t\) and has squared norm \(1+t^2>0\). Projection onto this line is
\[
 x\longmapsto\frac{\langle x,v_t\rangle}{1+t^2}v_t,
\]
which gives the matrix in \eqref{KS:eq:lineprojector}. In particular it is the orthogonal projector, not merely a map whose image is contained in the kernel.

Subtract \(q_0=\operatorname{diag}(0,1)\), obtaining the symmetric matrix
\[
 B_t=q_t-q_0=\frac1{1+t^2}
 \begin{pmatrix}t^2&-t\\-t&-t^2\end{pmatrix}.
\]
The two off-diagonal entries of its square are zero, since \(-t^3+t^3=0\). Both diagonal entries are \((t^4+t^2)/(1+t^2)^2=t^2/(1+t^2)\). Thus
\[
 B_t^{\mathsf T}B_t=B_t^2=\frac{t^2}{1+t^2}I.
\]
Pairing this identity with an arbitrary vector proves the all-vector norm formula. It bounds the operator norm above by \(|t|/\sqrt{1+t^2}\), and evaluating any unit vector gives equality. The argument includes \(t=0\), where both sides vanish.

For the common section \(A(r)=(r,0)\), the perturbation inverse leaves \(A\) unchanged because \(D_tA=I\). The constructed transport is therefore \(Jx=x-AD_tx\) on the old kernel. At \(x=(0,b)\), it gives \((-tb,b)\). Substitution into the projection formula gives instead \((-tb/(1+t^2),b/(1+t^2))\). The strict-radius theorem applies when \(|t|<1\), but these explicit line and projection calculations are valid for every real \(t\).
\end{proof}
\KSsource{KS51}

\begin{figure}[htbp]\centering\resizebox{.96\linewidth}{!}{
% BEGIN retained input sources/monographs/komorebi-space/source/figures/12_moving_line.tex
\begin{tabular}{@{}c@{\hspace{.4cm}}c@{}}
\includegraphics[width=7.3cm]{sources/monographs/komorebi-space/source/figures/generated/12a_moving_line_geometry.pdf}&
\includegraphics[width=7.3cm]{sources/monographs/komorebi-space/source/figures/generated/12b_moving_line_norm.pdf}
\end{tabular}

}\caption{The exact space can move continuously without remaining equal to the original kernel. Transport along a chosen splitting preserves a coordinate, while orthogonal projection minimizes distance. In the line example these two maps have different formulas even though both land in the same exact space.}\label{KS:fig:movingline}\end{figure}

\subsection{What fails without the hypotheses}
Take \(D_0=\operatorname{diag}(1,0)\) and \(D_t=\operatorname{diag}(1,t)\). The operator difference has norm \(|t|\). For every nonzero \(t\), however, the new kernel is zero, while the old kernel is the vertical axis. Their projections differ in norm by one. Operator closeness alone does not imply kernel continuity.

More generally, suppose a unit vector \(x\) is exact for the original operator but orthogonal to the new kernel. Every new residual gain \(C'\) obeys
\[
 1=\KSnorm{x-q'x}\le C'\KSnorm{D'x}\le C'\varepsilon.
\]
Losing an exact direction under a small operator change forces the new gain to grow. At the strict radius boundary, \(D=I\), \(A=I\), and \(D'=0\) have \(a=\varepsilon=1\), yet \(D'\) is not surjective. Thus the strict condition in Theorem~\ref{KS:KS49} cannot be replaced universally by a non-strict one.
\KSsource{KS52}
\begin{figure}[htbp]\centering\resizebox{.96\linewidth}{!}{
% BEGIN retained input sources/monographs/komorebi-space/source/figures/18_rank_loss.tex
\includegraphics[width=15cm]{sources/monographs/komorebi-space/source/figures/generated/18_rank_loss.pdf}

}\caption{Two families can have the same small operator-change norm and very different exact-space motion. A moving line retains a uniform residual gain. The rank-changing diagonal family loses an exact direction for every nonzero parameter, so its projector gap remains one.}\label{KS:fig:rankloss}\end{figure}

\subsection{Motion of the actual descent repair}
Compare pairs \((D,C)\) and \((D',C')\) on the same raw Hilbert space. Suppose each pair has closed local range and closed overlap range on its own locally exact subspace. Let \(G,G'\) be the resulting joint residual gains. If
\[
 \KSnorm{D-D'}\le\varepsilon_D,\qquad
 \KSnorm{C-C'}\le\varepsilon_C,
\]
then the maximum norm on the target product gives
\[
 \KSnorm{(D,C)-(D',C')}\le\max(\varepsilon_D,\varepsilon_C).
\]
Applying Theorem~\ref{KS:KS48} to the joint operators, and using the actual two-stage identity of Theorem~\ref{KS:KS46}, gives
\begin{equation}\label{KS:eq:descentmotion}
 \KSnorm{P_{D,C}-P_{D',C'}}
 \le\max(G,G')\max(\varepsilon_D,\varepsilon_C).
\end{equation}
This is a comparison of the original repair maps, not a replacement definition that happens to have a similar kernel. Both the local equations and overlap conditions can change. The raw space and its inner product remain fixed.
\KSsource{KS53}



% BEGIN retained input sources/monographs/komorebi-space/source/chapters/05a_certificates.tex
\chap{Reading a stability certificate}\label{KS:sec:certificate-reading}
A numerical residual becomes useful only after it is tied to an operator, a norm, and a conclusion. This section separates uncertainty in a measured residual from uncertainty in the operator itself. It also distinguishes closeness to a reference exact space from closeness to a changed one. All estimates below refer to the actual canonical Hilbert projection used earlier.

\subsection{The statement carried by a certificate}
Let \(D:E\to R\) be bounded linear from a real Hilbert space into a real normed space. A residual certificate with gain \(C\ge0\) is the universally quantified estimate
\[
 \KSnorm{x-qx}\le C\KSnorm{Dx}\quad\text{for every }x\in E,
 \qquad q=P_{\ker D}.
\]
A reported number \(C\) is not a certificate unless this estimate has been justified in the stated model. A successful test at one input is likewise not the universally quantified estimate.

Suppose the measured vector is \(y\), with a certified error radius \(\eta\) such that \(\KSnorm{Dx-y}\le\eta\). The triangle inequality converts the measured ball into a bound for the true residual.

\begin{KSproposition}[Measurement and operator uncertainty]\label{KS:KS69}
With the reference certificate above,
\[
 \KSnorm{x-qx}\le C(\KSnorm y+\eta)
 \quad\text{if }\KSnorm{Dx-y}\le\eta.
\]
More generally, let \(\widetilde D:E\to R\) be the bounded linear operator actually used in the measurement. If
\[
 \KSnorm{D-\widetilde D}\le\varepsilon,
 \qquad\KSnorm{\widetilde D x-y}\le\eta,
\]
then
\begin{equation}\label{KS:eq:measured-reference}
 \KSnorm{x-qx}\le C(\KSnorm y+\eta+\varepsilon\KSnorm x).
\end{equation}
The projection in this conclusion is still the projection for the reference operator \(D\).
\end{KSproposition}
\begin{proof}
The first assertion follows from \(\KSnorm{Dx}\le\KSnorm y+\eta\). For the second, decompose
\[
 Dx=(D-\widetilde D)x+(\widetilde D x-y)+y.
\]
Its norm is at most \(\varepsilon\KSnorm x+\eta+\KSnorm y\). Apply the reference residual certificate.
\end{proof}
\KSsource{KS69}
The input-size term has a mathematical purpose. A small error in an operator can produce a large error on a large input. It cannot generally be replaced by a number depending only on the measured residual. Even a measurement of zero may reflect a difference between the measured kernel and the reference kernel.

\subsection{A vanishing measured residual can still miss the reference space}
For a concrete family, take \(E=\KSR^2\) with the Euclidean norm and \(R=\KSR\), and define
\[
 D(a,b)=a,\qquad \widetilde D(a,b)=a+tb.
\]
The reference projection is \(q(a,b)=(0,b)\), its least residual gain is one, and \(\KSnorm{D-\widetilde D}\le|t|\). On the input \(x=(-tB,B)\), the measured operator returns exactly zero. Its distance to the reference kernel is nevertheless \(|tB|\).

Setting \(\eta=0\) and \(y=0\) in \eqref{KS:eq:measured-reference} gives the valid upper bound \(|t|\sqrt{1+t^2}|B|\). Deleting the input-size term would give zero and would be false whenever \(tB\ne0\). The input is exact for the changed model and not exact for the reference model. This is precisely the distinction that the estimate is designed to retain.

There is no statistical assertion here. The error radii are mathematical inputs. Estimating them from noisy observations, unknown dynamics, or a fitted model is an additional task that requires its own assumptions and evidence.

\subsection{A preserved kernel gives a stronger conclusion}
A useful regime occurs when the changed operator annihilates every old exact vector:
\[
 \widetilde D q=0.
\]
This says that the old kernel is contained in the new kernel. Equality is not assumed. Let \(\KSnorm{D-\widetilde D}\le\varepsilon\) and suppose the strict budget \(C\varepsilon<1\) holds.

\begin{KStheorem}[Kernel preservation and a strict error budget]\label{KS:KS70}
Under these hypotheses,
\[
 \ker\widetilde D=\ker D,
 \qquad
 \KSnorm{x-qx}\le\frac{C}{1-C\varepsilon}\KSnorm{\widetilde D x}.
\]
For a measurement with \(\KSnorm{\widetilde D x-y}\le\eta\), the bound becomes
\[
 \KSnorm{x-qx}\le\frac{C}{1-C\varepsilon}(\KSnorm y+\eta).
\]
Thus the exact space is unchanged and the residual gain degrades by an explicit factor.
\end{KStheorem}
\begin{proof}
Let \(q=P_{\ker D}\), and write \(u=x-qx\). The stated kernel inclusion makes both \(Dqx\) and \(\widetilde Dqx\) zero. Hence \(Du=Dx\) and \(\widetilde Du=\widetilde Dx\). Applying the reference residual bound to \(u\), whose projection onto \(\ker D\) is zero, gives
\[
 \begin{aligned}
 \KSnorm u&\le C\KSnorm{Du}
 \le C\KSnorm{\widetilde Du}+C\KSnorm{(D-\widetilde D)u}\\
 &\le C\KSnorm{\widetilde Dx}+C\varepsilon\KSnorm u.
 \end{aligned}
\]
Subtract the final term. The strict inequality \(C\varepsilon<1\) makes \(1-C\varepsilon\) positive, so division preserves the inequality and gives the claimed residual estimate. This also covers \(C=0\) without division by \(C\).

If \(\widetilde Dx=0\), the same estimate forces \(u=0\), hence \(x=qx\in\ker D\). This is the inclusion not initially assumed. Combining it with the supplied inclusion proves \(\ker\widetilde D=\ker D\), and their orthogonal projections coincide in the fixed Hilbert geometry. Finally \(\KSnorm{\widetilde Dx-y}\le\eta\) gives \(\KSnorm{\widetilde Dx}\le\KSnorm y+\eta\). Substitution proves the measured certificate. The kernel-preservation hypothesis is what removed \(qx\) before applying the perturbation estimate. Without it, one obtains a term proportional to \(\KSnorm x\), not the absorbable term proportional to \(\KSnorm u\).
\end{proof}
\KSsource{KS70}
This result does not require the reference operator to be surjective. It is different from the split-surjective perturbation theorem, which allows the kernel to move and constructs an actual transport between the old and new kernels. Both theorems have a strict quantitative radius, but they retain different structural data.

For \(D(a,b)=a\) and \(\widetilde D(a,b)=(1-t)a\), with \(0\le t<1\), the kernel is preserved and the exact gain is \(1/(1-t)\). The bound is attained on transverse inputs. At \(t=1\), the changed operator is zero and its kernel becomes all of \(\KSR^2\). This illustrates why the strict endpoint cannot simply be included.

\begin{figure}[htbp]\centering\resizebox{.96\linewidth}{!}{
% BEGIN retained input sources/monographs/komorebi-space/source/figures/22_gain_budget.tex
\includegraphics[width=15cm]{sources/monographs/komorebi-space/source/figures/generated/22_gain_budget.pdf}

}\caption{For the kernel-preserving family \(\widetilde D_t(a,b)=(1-t)a\), the exact residual gain is \(1/(1-t)\). It deteriorates as the strict budget is exhausted. At the excluded endpoint the kernel changes, so the finite-radius estimate does not extend by substituting \(t=1\).}\label{KS:fig:gain-budget}\end{figure}

\subsection{Select the theorem by its conclusion}
The following distinctions organize the estimates without turning them into competing formulas.
\begin{center}
\begin{tabularx}{\linewidth}{@{}p{.29\linewidth}p{.32\linewidth}X@{}}
\toprule
Available information & Structural hypothesis & Certified conclusion\\
\midrule
Reference gain and measured error & None on a changed kernel & Distance to the reference kernel, with an input-size term\\
Reference gain and small operator change & Old kernel preserved & Same kernel and a degraded gain\\
Two residual gains & Same source Hilbert geometry & Distance between the two orthogonal projectors\\
A right inverse and a strict radius & Split surjectivity & A new right inverse and an equivalence of moving kernels\\
\bottomrule
\end{tabularx}
\end{center}
A claim about the new exact space cannot be obtained by silently replacing the reference projector in the first row. Conversely, a theorem that establishes equality of kernels under preservation need not be weakened to a moving-kernel estimate when the stronger hypothesis is actually available.

\subsection{Norm conventions are part of the certificate}
On a finite family of real Hilbert fibers, the maximum norm and the Euclidean product norm satisfy
\[
 \KSnorm x_\infty\le\KSnorm x_2\le\sqrt{n}\KSnorm x_\infty,
\]
where \(n\) is the number of components. The first inequality follows because each squared component is one term in the sum. The second follows by bounding every component by the maximum. These inequalities transfer bounds with explicit dimension factors. They do not assert that the same numerical best constant remains optimal.

The graph chapters use the maximum norm on vertex and directed-edge products unless a different geometry is declared. The nearest-point statements use a Hilbert source norm. In the worked graph calculations that follow, selecting one orientation and using a sum of squared residuals defines a different numerical operator model. It has the same edge-exactness equations for an isometric connection, but a different residual norm.

A repair certificate should therefore identify its candidate space, residual map, preserved data, norms, gain, and error budgets. The exact equation alone does not specify these quantities. Explicitly retaining them is what makes a numerical conclusion transportable between examples.



% BEGIN retained input sources/monographs/komorebi-space/source/chapters/05b_transport_cost.tex

\chap{A geometry of transformations needs its own cost}\label{KS:sec:transport-cost}
The preceding analytic chapters measure the size of a residual or the distance from a candidate to an exact subspace. Neither measurement automatically defines a distance between mathematical objects. In an additive category, the zero map is often perfectly compatible and perfectly uninformative. A geometry that assigns zero cost to every compatible map can therefore identify all objects at distance zero.

A cost of transformation must be declared separately. It may measure an admissible computational expense, a loss of retained information, or another specified quantity. The interpretation alone proves none of its mathematical laws. The construction below begins only after the required laws have been established.

\subsection{Extended directed distance}
Let \(\mathcal C\) be a category and let \(W\) assign each arrow a value in the extended nonnegative reals \([0,\infty]\). Suppose
\[
 W(\KSid_A)=0,\qquad W(gf)\le W(f)+W(g).
\]
The order of ordinary composition is first \(f\), then \(g\). Define
\[
 d_W(A,B)=\inf_{f:A\to B}W(f).
\]
The infimum of an empty family is \(\infty\). No assumption of connectedness, finite cost, symmetry, or attainment is made.

\begin{KSproposition}[Directed cost geometry]\label{KS:KS78}
For every arrow \(f:A\to B\),
\[
 d_W(A,B)\le W(f),\qquad d_W(A,A)=0,
\]
and for all objects,
\[
 d_W(A,C)\le d_W(A,B)+d_W(B,C).
\]
The symmetric function
\[
 d_W^{\mathrm{sym}}(A,B)=\max\{d_W(A,B),d_W(B,A)\}
\]
has zero diagonal, is symmetric, and satisfies the extended triangle inequality. The relation \(d_W^{\mathrm{sym}}(A,B)=0\) is an equivalence relation.
\end{KSproposition}
\begin{proof}
The infimum is bounded above by each of its entries. Apply this to the identity to obtain the zero diagonal. If either term on the right of the triangle inequality is infinite, there is nothing to prove. Otherwise, for every \(\varepsilon>0\), choose arrows with costs less than the corresponding infima plus \(\varepsilon\). Their composite has cost less than
\(d_W(A,B)+d_W(B,C)+2\varepsilon\). Letting \(\varepsilon\) tend to zero gives the result. This argument does not require an optimizing arrow.

Each of the two directed distances from \(A\) to \(C\) and from \(C\) to \(A\) is bounded by \(d_W^{\mathrm{sym}}(A,B)+d_W^{\mathrm{sym}}(B,C)\). Taking their maximum proves the symmetric triangle law. Zero diagonal and symmetry give reflexivity and symmetry of the zero-distance relation. The triangle inequality and nonnegativity give transitivity.
\end{proof}
\KSsource{KS78}
The formal source proves these laws directly in \(\KSR_{\ge0}^{\infty}\) and constructs the zero-distance setoid. It does not identify zero distance with categorical isomorphism. Passing to the quotient can separate zero-distance classes, but would still not prove that those classes are the intended semantic objects.

\subsection{A concrete path category}
Consider four vertices \(A,B,C,D\), and the free path category generated by
\[
 A\xrightarrow{0}B,\qquad B\xrightarrow{0}A,\qquad
 A\xrightarrow{2}C,\qquad C\xrightarrow{5}A.
\]
There are no arrows incident to \(D\), apart from its identity. A path costs the sum of its edge labels. Concatenation is additive and empty paths cost zero, so all cost axioms hold. The directed distances are
\[
 \begin{array}{c|rrrr}
 d_W&A&B&C&D\\\hline
 A&0&0&2&\infty\\
 B&0&0&2&\infty\\
 C&5&5&0&\infty\\
 D&\infty&\infty&\infty&0
 \end{array}
\]
For example, the path from \(B\) through \(A\) to \(C\) costs two. Every path from \(C\) to \(B\) must first pay the edge cost five and can then move from \(A\) to \(B\) for free. Nonnegative loops cannot improve either cost.

\begin{figure}[htbp]\centering\resizebox{.96\linewidth}{!}{
% BEGIN retained input sources/monographs/komorebi-space/source/figures/29_transport_cost.tex
\includegraphics[width=15cm]{sources/monographs/komorebi-space/source/figures/generated/29_transport_cost.pdf}

}\caption{Directed path costs need not be symmetric or finite. Vertices $A$ and $B$ have zero cost in both directions, but they are not isomorphic in the free path category. The distinct zero-cost arrows have no imposed inverse relations.}\label{KS:fig:transport-cost}\end{figure}

The symmetric zero-distance classes are \(\{A,B\}\), \(\{C\}\), and \(\{D\}\). The distance between the first two classes is five. The third class is infinitely far from each of them. Yet \(A\) and \(B\) are not isomorphic in this category. Any path from one to the other has positive length, so concatenating it with a return path cannot give the length-zero identity. This is a categorical distinction even though the numerical cost is zero.

One can instead impose relations declaring the two arrows inverses. That changes the category. The distance table alone does not record whether those relations were imposed.

\subsection{Why compatibility alone can collapse the geometry}
\begin{KSproposition}[The zero-arrow collapse]\label{KS:KS79}
Suppose \(\mathcal C\) has zero morphisms and \(W(0:A\to B)=0\) for every pair of objects. Then
\[
 d_W(A,B)=0\quad\text{for all }A,B.
\]
\end{KSproposition}
\begin{proof}
The zero arrow is one member of the family over which the infimum is taken. Thus \(0\le d_W(A,B)\le W(0)=0\).
\end{proof}
\KSsource{KS79}
For the preadditive tower differential, the zero family of maps is exact. Consequently any proposed nonnegative arrow cost that vanishes on every exact map also vanishes on every zero map. If it satisfies the cost axioms at all, its object distance collapses by the proposition.

The qualification is important. A raw operator defect norm need not satisfy the unweighted composition inequality required of \(W\). The transport Leibniz identity gives coefficients involving the surrounding maps. It is not legitimate to first call an arbitrary defect norm a subadditive cost and then use the metric construction.

\subsection{Preservation changes the admissible arrows}
One response to collapse is to restrict arrows to those satisfying a declared preservation contract. For objects equipped with maps \(p_A:A\to S\), a morphism can be required to satisfy
\[
 p_Bf=p_A.
\]
Identities and composition preserve this equation. Unless \(p_A=0\), the zero arrow is not admissible. In other applications the arrows might be required to be isomorphisms or to retain a particular observable. These are different choices of category, not automatic consequences of small residuals.

Even after this restriction, separation remains a theorem to prove rather than a consequence of terminology. Distinct objects can still be connected by arbitrarily cheap preserving maps. Symmetry can fail, and the infimum can fail to be attained. A quantitative comparison therefore needs four pieces of information: the objects, the admissible arrows, the actual cost, and the meaning of zero distance.

The role of this construction is modest but useful. It prevents a compatibility measure from being silently promoted to a universal geometry. It also gives a valid directed geometry when the separate categorical and cost assumptions really hold.



% BEGIN retained input sources/monographs/komorebi-space/source/chapters/05c_iteration.tex

\chapter{From approximation to exact realization}\label{KS:book:chapter-23}
A vanishing defect is a statement about measurements. A limiting realization also requires control of the candidate itself. This part follows the movement of approximate answers, the closure of admissible correction images, and the meaning of the ambient completion.

\chap{Exactification without a fixed contraction rate}\label{KS:sec:variable-exactification}
The constructed linear repair in Section~\ref{KS:sec:split} is a useful special case of a more general principle. One can prove convergence without one fixed contraction factor, without a linear update rule, and without a globally Lipschitz iteration. What must not be lost is control of the total movement.

\subsection{The two limits that must be established}
Let \((M,\KSdist)\) be a complete metric space. Let \(\omega:M\to[0,\infty)\) be a nonnegative residual function, and call a point exact when \(\omega(x)=0\). We regard \(\omega\) as a real-valued function when discussing continuity. Assume it is lower semicontinuous, so that every sublevel set \(\{x:\omega(x)\le a\}\) is closed.

For an approximation sequence \(f_n\), the statements
\[
 \omega(f_n)\longrightarrow0
 \quad\text{and}\quad
 f_n\longrightarrow x
\]
are different. The first concerns the observed discrepancy. The second concerns the mathematical candidate in the chosen metric. Neither implies the other without additional hypotheses.

\begin{KStheorem}[Summable movement gives exactification]\label{KS:KS71}
Suppose \(b_n\ge0\), \(\sum_{n=0}^{\infty}b_n<\infty\), and
\[
 \KSdist(f_n,f_{n+1})\le b_n,\qquad \omega(f_n)\longrightarrow0.
\]
Then there is a limit \(x\in M\) satisfying
\[
 f_n\longrightarrow x,\qquad\omega(x)=0,\qquad
 \KSdist(f_n,x)\le\sum_{k=n}^{\infty}b_k.
\]
\end{KStheorem}
\begin{proof}
For integers \(m>n\), apply the triangle inequality successively along the finite chain \(f_n,f_{n+1},\ldots,f_m\). It gives
\[
 \KSdist(f_n,f_m)\le\sum_{k=n}^{m-1}\KSdist(f_k,f_{k+1})
 \le\sum_{k=n}^{m-1}b_k\le\sum_{k=n}^{\infty}b_k.
\]
Given \(\varepsilon>0\), convergence of the nonnegative series supplies \(N\) for which the final tail is below \(\varepsilon\) whenever \(n\ge N\). Interchanging \(m,n\) when necessary proves the Cauchy condition for all pairs of indices at least \(N\). Completeness, not residual decay, is the hypothesis that now supplies a limit \(x\in M\).

Keep \(n\) fixed. The reverse triangle inequality says
\[
 \bigl|\KSdist(f_n,f_m)-\KSdist(f_n,x)\bigr|\le\KSdist(f_m,x)\longrightarrow0.
\]
Passing to the limit in the uniform tail bound proves the claimed distance certificate at that \(n\).

It remains to prove exactness of the limit. Fix \(\varepsilon>0\). Since \(\omega(f_m)\to0\), an entire tail of the sequence belongs to the sublevel set \(C_\varepsilon=\{z:\omega(z)\le\varepsilon\}\). Lower semicontinuity makes this set closed, so the limit of that tail belongs to it. Hence \(\omega(x)\le\varepsilon\) for every positive \(\varepsilon\). If \(\omega(x)>0\), choosing \(\varepsilon=\omega(x)/2\) contradicts this inequality. Since \(\omega\) is nonnegative, \(\omega(x)=0\). We have not used, or asserted, continuity of \(\omega\) or the stronger equality \(\omega(f_m)\to\omega(x)\) for arbitrary convergent inputs.
\end{proof}
\KSsource{KS71}
The formal theorem is slightly more general: it is stated in a complete pseudometric space. In a metric space the limit of this fixed sequence is unique. That does not mean different admissible approximation procedures select the same exact point.

\subsection{Variable rates and a movement budget}
Let \(q_n,C_n\ge0\), and define the prefix product
\[
 Q_0=1,\qquad Q_{n+1}=q_nQ_n.
\]
Assume
\begin{equation}\label{KS:eq:variable-rate}
 \omega(f_{n+1})\le q_n\omega(f_n),\qquad
 \KSdist(f_n,f_{n+1})\le C_n\omega(f_n).
\end{equation}

\begin{KStheorem}[Variable-rate exactification]\label{KS:KS72}
If \(Q_n\to0\) and \(\sum_n C_nQ_n<\infty\), then \((f_n)\) has an exact limit and
\begin{equation}\label{KS:eq:variable-tail}
 \KSdist(f_n,x)\le\omega(f_0)\sum_{k=n}^{\infty}C_kQ_k.
\end{equation}
For constant \(q_n=q\in[0,1)\) and constant \(C_n=C\), the sharper restarted estimate is
\[
 \KSdist(f_n,x)\le\frac{C}{1-q}\,\omega(f_n).
\]
\end{KStheorem}
\begin{proof}
The estimate \(\omega(f_n)\le Q_n\omega(f_0)\) follows by induction. At \(n=0\) it is equality. If it holds at \(n\), nonnegativity of \(q_n\) permits multiplication and gives
\[
 \omega(f_{n+1})\le q_n\omega(f_n)
 \le q_nQ_n\omega(f_0)=Q_{n+1}\omega(f_0).
\]
Because \(Q_n\to0\) and \(\omega(f_0)\) is a fixed finite number, the residual tends to zero. The movement assumption then gives
\[
 \KSdist(f_n,f_{n+1})\le C_nQ_n\omega(f_0).
\]
This is a summable nonnegative budget by the hypothesis on \(\sum C_nQ_n\). All assumptions of Theorem~\ref{KS:KS71} now hold, and its tail certificate is exactly \eqref{KS:eq:variable-tail}. If \(\omega(f_0)=0\), the same estimates make every step zero and every residual zero. No cancellation of that value was used.

For constant rates, the global estimate can be sharpened by restarting at the current index. Fix \(n\). A new induction on \(j\) gives
\[
 \omega(f_{n+j})\le q^j\omega(f_n),\qquad
 \KSdist(f_{n+j},f_{n+j+1})\le Cq^j\omega(f_n).
\]
Sum the finite chain from \(n\) to \(n+m\), and then let \(m\to\infty\) using the already established limit. The geometric series equals \(1/(1-q)\), including \(q=0\), where its first term is one and all later terms vanish. This proves the restarted bound. The variable-rate proof itself requires only nonnegative \(q_n\), decay of their cumulative products, and summability of the movement weights. It does not insert an unstated upper bound \(q_n<1\) at each step.
\end{proof}
\KSsource{KS72}
The individual factors \(q_n\) need not all be below one. Temporary residual increases are permitted if the prefix products decay and the weighted movements remain summable. Conversely, a decaying product is not enough when its movement budget diverges.

\subsection{A convergent path with no uniform geometric factor}
In \(\KSR^2\) with the maximum norm, fix \(s\in\KSR\) and take
\[
 f_n=(s,(n+1)^{-1}),\qquad \omega(u,r)=|r|.
\]
The protected coordinate is the first one. Direct calculation gives
\[
 q_n=\frac{n+1}{n+2},\quad Q_n=\frac1{n+1},\quad
 C_n=\frac1{n+2},\quad
 \KSdist(f_n,f_{n+1})=\frac1{(n+1)(n+2)}.
\]
No constant \(q<1\) bounds all the ratios \(q_n\), since they tend to one. Nevertheless
\[
 \sum_{k=n}^{\infty}C_kQ_k
 =\sum_{k=n}^{\infty}\left(\frac1{k+1}-\frac1{k+2}\right)
 =\frac1{n+1}.
\]
The limit is \((s,0)\), and the error bound is an equality. The model shows why the variable-rate theorem is not merely a rephrasing of a fixed geometric contraction.

\begin{figure}[htbp]\centering\resizebox{.96\linewidth}{!}{
% BEGIN retained input sources/monographs/komorebi-space/source/figures/25_summable_path.tex
\includegraphics[width=15cm]{sources/monographs/komorebi-space/source/figures/generated/25_summable_path.pdf}

}\caption{Polynomial residual decay can still give an exact limit with a sharp error certificate. The residual $1/(n+1)$ is not summable, but the actual movement $1/((n+1)(n+2))$ is. The theorem uses the movement series.}\label{KS:fig:summable-path}\end{figure}

\subsection{Residual bursts can be compatible with convergence}
Set \(q_{2k}=2\), \(q_{2k+1}=1/8\). Then
\[
 Q_{2k}=4^{-k},\qquad Q_{2k+1}=2\cdot4^{-k}.
\]
For \(f_n=(s,Q_n)\), the residual doubles at the start of every pair of steps. Choose \(C_{2k}=1\) and \(C_{2k+1}=7/8\). These give equality in the movement bound, and
\[
 \sum_{n=0}^{\infty}C_nQ_n
 =\sum_{k=0}^{\infty}\left(1+\frac74\right)4^{-k}
 =\frac{11}{3}.
\]
Thus \(f_n\to(s,0)\) despite infinitely many increases. Starting at an even index \(2k\), the remaining movement budget is \(\frac{11}{3}4^{-k}\). At an odd index \(2k+1\), it is \(\frac83 4^{-k}\). These bounds are larger than the actual distance because the path repeatedly moves away and returns.

\begin{figure}[htbp]\centering\resizebox{.96\linewidth}{!}{
% BEGIN retained input sources/monographs/komorebi-space/source/figures/26_residual_bursts.tex
\includegraphics[width=15cm]{sources/monographs/komorebi-space/source/figures/generated/26_residual_bursts.pdf}

}\caption{The residual increases at every even-to-odd step, but each two-step product is $1/4$. A convergent product and a summable movement budget control the whole path. Monotonicity of individual residual values is not assumed by the variable-rate theorem.}\label{KS:fig:residual-bursts}\end{figure}

\subsection{Protected meaning survives a controlled limit}
Return to Banach and normed linear spaces, with continuous linear maps \(d:E\to R\) and \(p:E\to S\). Put \(\omega(f)=\KSnorm{df}\). The residual is continuous, hence lower semicontinuous.

\begin{KSproposition}[Controlled paths kill the actual obstruction]\label{KS:KS73}
Suppose \(E\) is complete, \(\KSnorm{df_n}\to0\), \(p(f_n)=p(f_0)\), and \(\KSdist(f_n,f_{n+1})\le b_n\) for a summable nonnegative budget. Then there is \(x\) with
\[
 f_n\to x,\qquad dx=0,\qquad p(x)=p(f_0),
\]
and the constrained obstruction \([df_0]\in R/d(\ker p)\) vanishes.
\end{KSproposition}
\begin{proof}
The function \(f\mapsto\KSnorm{df}\) is continuous because \(d\) is bounded linear and the norm is continuous. It is therefore an admissible lower semicontinuous residual for Theorem~\ref{KS:KS71}. That theorem supplies a limit \(x\in E\), together with \(\KSnorm{dx}=0\), which means \(dx=0\).

Continuity of \(p\) gives \(p(f_n)\to p(x)\). The left side is the constant sequence with value \(p(f_0)\). Uniqueness of limits in the normed target therefore gives \(p(x)=p(f_0)\). Define the actual vector \(a=f_0-x\). Linearity, rather than a closure argument, now yields
\[
 p(a)=p(f_0)-p(x)=0,\qquad da=df_0-dx=df_0.
\]
Thus \(a\in\ker p\) and \(df_0\in d(\ker p)\). Membership in this algebraic image is precisely what makes the class \([df_0]\) zero. Merely knowing membership in the closure of that image would not justify the same conclusion.
\end{proof}
\KSsource{KS73}
This proposition converts a controlled approximation argument into a statement about the algebraic quotient defined earlier. Merely showing that \(df_n\to0\) would not produce the admissible vector \(f_0-x\), because there might be no candidate limit \(x\).

\subsection{When two repair paths select the same point}
\begin{KSproposition}[Path independence requires a comparison]\label{KS:KS74}
Let \(f_n\to x\) and \(g_n\to y\) in a metric space. If
\[
 \KSdist(f_n,g_n)\le\beta_n,\qquad\beta_n\to0,
\]
then \(x=y\). In contrast, exactness and preservation of a nonzero observable alone do not force a unique answer.
\end{KSproposition}
\begin{proof}
For each \(n\), the triangle inequality gives
\[
 \KSdist(x,y)\le\KSdist(x,f_n)+\KSdist(f_n,g_n)+\KSdist(g_n,y)
 \le\KSdist(x,f_n)+\beta_n+\KSdist(g_n,y).
\]
Given any \(\varepsilon>0\), choose \(n\) large enough that the three terms on the right are each below \(\varepsilon/3\). Thus \(\KSdist(x,y)\le\varepsilon\) for every positive \(\varepsilon\), so the distance is zero. The metric separation axiom gives \(x=y\). In a pseudometric space only zero distance would follow without an additional separation hypothesis.

For the negative assertion let \(E=\KSR\times(\KSR\times\KSR)\), \(d(s,u,r)=r\), and \(p(s,u,r)=s\). Both \(x=(1,0,0)\) and \(y=(1,1,0)\) have zero defect and protected value one, but their middle coordinates differ. They are distinct points of the same feasible fiber. Constant sequences at those points meet exactness and preservation but not a vanishing comparison bound, so these conditions alone cannot imply path independence.
\end{proof}
\KSsource{KS74}
A stronger worked example uses actual convergent paths from the \emph{same} initial input:
\[
 f_n^{(a)}=(s,a(1-2^{-n}),2^{-n})\qquad(a\in\KSR).
\]
All paths begin at \((s,0,1)\), keep \(p=s\), and have residual \(2^{-n}\). Their maximum-norm movements are \(\frac12\max\{|a|,1\}2^{-n}\), so every path satisfies geometric exactification. Their limits are \((s,a,0)\). For different \(a,b\),
\[
 \KSdist(f_n^{(a)},f_n^{(b)})=|a-b|(1-2^{-n})\longrightarrow|a-b|,
\]
not zero. The missing uniqueness hypothesis is visible in the exact direction along which the paths separate.

\begin{figure}[htbp]\centering\resizebox{.96\linewidth}{!}{
% BEGIN retained input sources/monographs/komorebi-space/source/figures/27_path_limits.tex
\includegraphics[width=15cm]{sources/monographs/komorebi-space/source/figures/generated/27_path_limits.pdf}

}\caption{All paths start at the same candidate, have the same residual at each step, and preserve the same coordinate $s$. They converge to different exact points because the coordinate $u$ is unprotected. A vanishing residual is not a path-comparison budget.}\label{KS:fig:path-limits}\end{figure}

\subsection{Changing the sampling schedule does not preserve a rate}
A residual profile \(\eta_n\) indexed by resolution is not an iteration trajectory. Nevertheless a useful convergence property is invariant under a cofinal schedule when the profile is monotone.

\begin{KSproposition}[Cofinal sampling of a decreasing profile]\label{KS:KS75}
Suppose \(\eta_n\ge0\) is nonincreasing and \(\phi:\KSN\to\KSN\) is monotone and unbounded. Then
\[
 \eta_n\to0\quad\Longleftrightarrow\quad\eta_{\phi(n)}\to0.
\]
Without the monotonicity of \(\eta\), the reverse implication can fail even for a cofinal schedule.
\end{KSproposition}
\begin{proof}
The schedule tends to infinity, so convergence of the full profile implies convergence of the sampled one. Conversely, choose a sampled value below any prescribed \(\varepsilon>0\). Every later index of the original profile has no larger value, by monotonicity. This proves convergence of the full profile.

For failure without monotonicity, take \(\eta_{2n}=0\), \(\eta_{2n+1}=1\), and \(\phi(n)=2n\). All sampled values vanish, but infinitely many original values equal one.
\end{proof}
\KSsource{KS75}
Even in the decreasing case, the numerical rate changes. The profile \(\eta_n=1/(n+1)\), sampled at \(\phi(n)=2^n-1\), becomes \(2^{-n}\). This is exponential decay in the \emph{number of observations}, not a proof of a faster correction process at the original resolution scale.

\subsection{Three hypotheses that cannot simply be dropped}
\textbf{Summable movement.} In maximum-norm \(\KSR^3\), let
\[
 f_n=(s,H_n,(n+1)^{-1}),\qquad
 H_n=\sum_{j=1}^{n}\frac1j.
\]
The first coordinate is protected and the residual is the last coordinate. It tends to zero with the same ratios as the polynomial example. The movement equals \(1/(n+1)\), however, and the unprotected coordinate diverges. Here \(C_n=1\), so \(\sum C_nQ_n\) is the divergent harmonic series. Residual decay and semantic preservation hold, but exactification of this path fails.

\textbf{Completeness.} In the metric space \((0,\infty)\), take \(f_n=1/(n+1)\) and \(\omega(x)=x\). The steps are summable and the residual tends to zero, but there is no limit in this space. Completing the space adds the missing zero. It is essential to say whether that completion belongs to the intended model.

\textbf{Lower semicontinuity.} In \(\KSR\), define \(\omega(0)=1\) and \(\omega(x)=0\) for \(x\ne0\). The same positive sequence converges with summable movement and residual zero at every step, yet its limit is not exact. The zero sublevel set is not closed. This example isolates why convergence of measurements along a sequence must interact correctly with the ambient topology.



% BEGIN retained input sources/monographs/komorebi-space/source/chapters/05d_floors.tex

\chap{The residual floor and the closure of admissible changes}\label{KS:sec:floor}
An algebraic obstruction asks whether a residual is exactly the defect of an allowed correction. A numerical floor asks how closely allowed defects can approximate it. These questions agree under a closedness hypothesis. Without that hypothesis they can disagree even in a complete Hilbert setting.

\subsection{Distance to the actual correction image}
Let \(R\) be a normed additive group and \(B\le R\) an additive subgroup. In the protected linear model, \(B=d(\ker p)\). For a residual \(\Omega\in R\), define
\[
 \nu_B(\Omega)=\inf_{b\in B}\KSnorm{\Omega-b}.
\]
The subgroup contains zero, so this is a finite nonnegative real number. It is an infimum, not an assertion that a nearest admissible residual exists.

\begin{KSproposition}[The floor detects closure]\label{KS:KS76}
For every \(b\in B\),
\[
 0\le\nu_B(\Omega)\le\KSnorm{\Omega-b}.
\]
Moreover,
\[
 \nu_B(\Omega)=0\quad\Longleftrightarrow\quad\Omega\in\overline B.
\]
If \(B\) is closed, this is equivalent to \(\Omega\in B\).
\end{KSproposition}
\begin{proof}
The set of candidate distances is nonempty because \(0\in B\), bounded below by zero, and contains the finite value \(\KSnorm\Omega\). Its infimum is therefore a finite nonnegative real, and the defining lower-bound property gives \(\nu_B(\Omega)\le\KSnorm{\Omega-b}\) for every \(b\in B\).

Suppose \(\nu_B(\Omega)=0\) and fix \(\varepsilon>0\). There must be \(b\in B\) with \(\KSnorm{\Omega-b}<\varepsilon\). Otherwise every candidate distance would be at least \(\varepsilon\), making \(\varepsilon\) a lower bound and forcing the infimum to be at least \(\varepsilon\), a contradiction. Thus every open ball about \(\Omega\) meets \(B\), which is the closure criterion.

Conversely, if \(\Omega\in\overline B\), every such ball supplies \(b\) with \(\KSnorm{\Omega-b}<\varepsilon\). The infimum then lies between zero and \(\varepsilon\) for every positive \(\varepsilon\), so it is zero. If \(B\) is closed, membership in its closure is exactly membership in \(B\). No step asserts that the infimum has a minimizing element when \(B\) is not closed, and closedness alone is being used here only for the zero-floor membership conclusion.
\end{proof}
\KSsource{KS76}
Translation by a member \(b_0\in B\) does not change the floor: replacing \(b\) by \(b-b_0\) gives the same set of distances. Thus the floor is well defined on the algebraic quotient \(R/B\). When \(B\) is a linear subspace, this is the usual quotient seminorm. A nonzero algebraic class can have seminorm zero when \(B\) is not closed. The retained source implements the floor on residual representatives and proves the closure criterion. This quotient interpretation is its elementary translation-invariance calculation, not a new claim that every obstruction quotient carries a separated norm.

\Needspace{12\baselineskip}
\subsection{A precise three-way distinction in this fixed model}
For this specified subgroup and residual there are exactly three possibilities:
\begin{center}
\begin{tabularx}{\linewidth}{@{}p{.29\linewidth}p{.20\linewidth}X@{}}
\toprule Residual location & Algebraic class & Meaning of the floor\\\midrule
\(\Omega\in B\) & Zero & Exact cancellation exists, and the floor is zero.\\
\(\Omega\in\overline B\setminus B\) & Nonzero & Arbitrarily small corrected residuals exist, but no exact cancellation exists.\\
\(\Omega\notin\overline B\) & Nonzero & The floor is strictly positive and blocks every allowed correction.\\\bottomrule
\end{tabularx}
\end{center}
This is a classification for the stated normed-group correction problem, not a universal trichotomy for every possible Komorebi model. Testing which case holds can itself be difficult. The definitions do not supply a decision algorithm.

In finite-dimensional real or complex linear models, every linear subspace is closed. The middle row is therefore absent. It becomes possible in infinite dimensions, or for additive subgroups that are not closed real-linear subspaces. For example the rational subgroup of \(\KSR\) is dense, so the distance from any irrational real to that subgroup is zero even though its coset is nonzero.

Closedness also should not be confused with a numerical error bound back in the candidate space. The statement \(\Omega\in B\) supplies an admissible correction, but not a bound on its size. The closed-range estimates of Sections~\ref{KS:sec:closed} and~\ref{KS:sec:motion}, or their protected counterparts, are the additional analytic tools that control that size in their declared settings.

\subsection{A persistent floor blocks every correction sequence}
The correction image can vary with the resolution or the admissibility rule. Let \(B_n\le R\) and let \(\Omega_n\in R\) be given.

\begin{KSproposition}[Uniformly positive floors obstruct asymptotic repair]\label{KS:KS77}
Suppose \(c>0\) and eventually
\[
 c\le\nu_{B_n}(\Omega_n).
\]
For any choices \(b_n\in B_n\), it is impossible that
\[
 \KSnorm{\Omega_n-b_n}\longrightarrow0.
\]
\end{KSproposition}
\begin{proof}
Every allowed correction satisfies
\(\KSnorm{\Omega_n-b_n}\ge\nu_{B_n}(\Omega_n)\). The corrected norms are therefore eventually at least \(c\). They cannot converge to zero.
\end{proof}
\KSsource{KS77}
The quantifier covers all admissible sequences, not just the output of one repair routine. On the other hand, the hypothesis is an eventually \emph{uniform} positive lower bound. Having a nonzero obstruction class at each stage is not the same statement. The classes might have zero floors, or their positive floors might tend to zero.

A finite-dimensional illustration uses \(R=\KSR^2\), \(B=\KSR\times\{0\}\), and \(\Omega_n=(n,1)\). Every correction leaves the second coordinate equal to one, so the floor is one independently of \(n\). Replacing the residuals by \((n,1/(n+1))\) gives nonzero algebraic classes at every stage, but their floors tend to zero. The persistent-floor theorem applies to the first sequence and not the second.

\subsection{The bridge back to exactification}
Proposition~\ref{KS:KS73} and Proposition~\ref{KS:KS76} fit together closely. If corrections produce a norm-convergent, protected path with vanishing residual, the actual quotient class vanishes. If only the corrected residuals tend to zero, all that follows is membership in \(\overline{d(\ker p)}\).

In particular, when the class is nonzero but its floor is zero, \emph{no} protected path with vanishing residual can satisfy all the exactification hypotheses. Something must fail: the path may escape, its movements may not be summable, the candidate space may not be complete, or its residual and semantic maps may not have the required continuity. The next model isolates escape while keeping the spaces Hilbert and both maps continuous.



% BEGIN retained input sources/monographs/komorebi-space/source/chapters/05e_escape.tex

\chap{Exact finite repairs can escape the intended ambient space}\label{KS:sec:escape}
This model links the algebraic obstruction, the numerical floor, the movement of approximations, and the semantic image of an inverse limit. It uses a diagonal operator whose entries are completely explicit. Every finite problem is solvable. The protected infinite problem is not.

\subsection{The complete spaces and the preservation rule}
Index sequences by the positive integers. Set
\[
 E=\ell^2(\KSN_{>0};\KSR)\oplus\KSR,\qquad
 R=\ell^2(\KSN_{>0};\KSR),\qquad
 \KSnorm{(u,t)}_E^2=\sum_{n\ge1}|u_n|^2+|t|^2.
\]
Let \(w=(1/n)_{n\ge1}\). It belongs to \(\ell^2\), because
\[
 \sum_{n\ge1}\frac1{n^2}
 \le1+\sum_{n\ge2}\frac1{n(n-1)}=2.
\]
Define the defect and protected observable by
\[
 d(u,t)_n=\frac{u_n+t}{n},\qquad p(u,t)=t.
\]
The diagonal map \(T:u\mapsto(u_n/n)_n\) has operator norm one. Since \(d(u,t)=Tu+tw\),
\[
 \KSnorm{d(u,t)}\le\KSnorm u+|t|\KSnorm w
 \le\sqrt{1+\KSnorm w^2}\,\KSnorm{(u,t)}_E.
\]
Thus both \(d\) and \(p\) are continuous linear maps between the declared complete spaces. No failure of completeness or continuity is built into the model.

Take the raw candidate \(f_0=(0,1)\). Its residual is \(w\), and the coordinate that must be preserved is \(t=1\). An exact protected answer would require
\[
 u_n=-1\quad\text{for every }n.
\]
That sequence does not lie in \(\ell^2\). Consequently no answer in \(E\) is both exact and protected.

In fact \(\ker d=\{0\}\). If \(d(u,t)=0\), then \(u_n=-t\) for every \(n\). A constant square-summable sequence must be zero, so \(t=0\) and \(u=0\). There is an exact answer after abandoning the preservation rule, namely zero, but it changes \(p(f_0)=1\) to zero. Even a trivial ambiguity space does not imply feasibility.

\subsection{The obstruction class is nonzero while its floor is zero}
The admissible correction space is \(\ker p=\ell^2\oplus\{0\}\), with correction image
\[
 B=d(\ker p)=\KSran T
  =\{r\in\ell^2:(nr_n)_{n\ge1}\in\ell^2\}.
\]
It contains every finitely supported sequence, so it is dense in \(R\). But \(w\notin B\), because \((nw_n)_n=(1,1,\ldots)\) is not square summable. Therefore
\[
 [w]\ne0\text{ in }R/B,\qquad \nu_B(w)=0.
\]
This is the middle row of the preceding classification. The example is a worked Hilbert-space derivation without a separate Lean instance. Its finite-coordinate calculations are independently tested. The general floor and exactification results in Sections~\ref{KS:sec:variable-exactification} and~\ref{KS:sec:floor} have formal counterparts.

\subsection{A complete calculation of the finite repairs}
Let \(\mathbf 1_N\) be one in coordinates \(1,\ldots,N\) and zero elsewhere. Put
\[
 a_N=(\mathbf 1_N,0),\qquad f_N=f_0-a_N=(-\mathbf 1_N,1).
\]
Then \(p(f_N)=1\), and the first \(N\) residual coordinates vanish exactly. The remaining residual and movement satisfy
\[
 d(f_N)_n=\begin{cases}0&n\le N,\\1/n&n>N,\end{cases}
 \qquad \KSnorm{a_N}^2=N,\qquad \KSnorm{f_N}^2=N+1.
\]
For every \(N\ge1\), telescoping comparisons give
\begin{equation}\label{KS:eq:escape-tail}
 \frac1{N+1}\le\KSnorm{d(f_N)}^2
 =\sum_{n>N}\frac1{n^2}\le\frac1N.
\end{equation}
Indeed \(1/[n(n+1)]\le1/n^2\le1/[n(n-1)]\) for \(n\ge2\). Summing the two telescoping bounds proves \eqref{KS:eq:escape-tail} without numerically evaluating an infinite series.

The residual tends to zero, but the corrections grow like \(\sqrt N\). Successive candidates differ by a unit coordinate, so
\[
 \KSnorm{f_{N+1}-f_N}=1.
\]
The movement series diverges and the sequence is not Cauchy. This is exactly the hypothesis missing from the summable exactification theorem.

\begin{center}
\begin{tabular}{@{}rrrr@{}}
\toprule Repaired coordinates \(N\) & \(\KSnorm{a_N}\) & Residual square interval & Finite inverse gain\\\midrule
1 & 1 & \([1/2,1]\) & 1\\
4 & 2 & \([1/5,1/4]\) & 4\\
16 & 4 & \([1/17,1/16]\) & 16\\
64 & 8 & \([1/65,1/64]\) & 64\\\bottomrule
\end{tabular}
\end{center}
The intervals bound the \emph{infinite} residual remaining after the displayed repair. The inverse-gain column concerns a different finite operator, defined below. Keeping those meanings separate prevents a finite success from being mistaken for a uniform infinite certificate.

\begin{figure}[htbp]\centering\resizebox{.96\linewidth}{!}{
% BEGIN retained input sources/monographs/komorebi-space/source/figures/28_escaping_repair.tex
\includegraphics[width=15cm]{sources/monographs/komorebi-space/source/figures/generated/28_escaping_repair.pdf}

}\caption{The squared correction size is $N$, while the infinite residual square lies between $1/(N+1)$ and $1/N$. The plotted band consists of proved bounds, not a numerical truncation of the tail. Smaller residuals require larger changes and do not produce a convergent candidate.}\label{KS:fig:escaping-repair}\end{figure}

\subsection{Escape is not an accident of this particular algorithm}
Suppose an arbitrary protected candidate \((u,1)\) satisfies \(\KSnorm{d(u,1)}\le\varepsilon\). On its first \(N\) coordinates,
\[
 \sum_{n=1}^{N}|u_n+1|^2
 \le N^2\sum_{n=1}^{N}\frac{|u_n+1|^2}{n^2}
 \le N^2\varepsilon^2.
\]
The reverse triangle inequality in \(\KSR^N\) gives
\begin{equation}\label{KS:eq:escape-necessity}
 \KSnorm u\ge\sqrt N-N\varepsilon\qquad(N\ge1).
\end{equation}
If a sequence of protected candidates has residual tending to zero, then for each fixed \(N\), its eventual correction norms are at least, for example, \(\sqrt N/2\). Taking \(N\) arbitrarily large shows that these norms must tend to infinity. There is no uniformly bounded family of protected corrections with arbitrarily small residual.

Thus changing the iteration rule cannot fix the obstruction while preserving all the declared requirements. One must change the input, the preservation contract, or the intended ambient class. The inequality is a quantitative obstruction to a bounded repair procedure, rather than a failure report from one chosen recurrence.

\subsection{Every finite operator is invertible, but the gains diverge}
On \(\KSR^N\), let
\[
 T_N=\operatorname{diag}(1,1/2,\ldots,1/N).
\]
It is invertible and \(\KSnorm{T_N^{-1}}=N\) in the Euclidean norm. The upper bound follows by multiplying each coordinate by an integer at most \(N\), and equality is attained at the last coordinate vector. Full finite rank therefore supplies no uniform gain as \(N\to\infty\).

In the infinite space, \(\KSnorm{e_n}=1\) but \(\KSnorm{Te_n}=1/n\), while \(\ker T=\{0\}\). A uniform estimate
\(\KSdist(u,\ker T)\le C\KSnorm{Tu}\) would force \(1\le C/n\) for every \(n\), which is impossible. This agrees with the nonclosed range already established directly. The low residual directions do not reflect new exact vectors. They reflect increasingly ill-conditioned approximate directions.

\subsection{A coherent inverse-limit point still exists in a larger space}
At each finite level, the protected exact equations have the unique solution
\[
 (-1,\ldots,-1;1)\in\KSR^N\oplus\KSR.
\]
Truncation maps the solution at level \(N+1\) to the one at level \(N\). These finite solution sets are singletons, so their inverse limit is also a singleton. Its coordinate description is
\[
 (-1,-1,\ldots;1)\in\KSR^{\KSN_{>0}}\times\KSR.
\]
The coherent point is not missing. It lies outside the intended \(\ell^2\)-based semantic space.

This example does not contradict finite-fiber existence. It distinguishes the ordinary inverse limit of finite views from the image of a separate comparison map that sends square-summable semantic objects to their finite prefixes. The prescribed coherent point is not in that image. Replacing \(\ell^2\) by all sequences makes the equation solvable, but changes the topology and the class of allowed objects.

\begin{insight}
\textbf{Four statements, four answers.} Every finite prefix is exactly solvable. The ordinary coherent inverse-limit point exists. Its protected Hilbert realization does not exist. Approximate protected Hilbert realizations have residuals tending to zero only at the cost of unbounded correction size. All four facts hold simultaneously in one explicit model.
\end{insight}



% BEGIN retained input sources/monographs/komorebi-space/source/chapters/05f_bounded_realization.tex

\chap{A bounded repair budget can recover existence}\label{KS:sec:bounded-realization}
The preceding example does not merely lack a convergent iteration. It has no bounded family of increasingly accurate protected corrections. This suggests a converse question. If approximate corrections stay in one bounded region, does an exact answer exist?

For bounded linear equations in a Hilbert space, the answer is yes. This is an existence statement, not a claim that the given sequence converges in norm. A stronger conclusion becomes available when the candidates are the least norm solutions of nested finite constraint systems.

\begin{scopebox}
The two bounded-realization theorems in this chapter have formal counterparts in the accompanying Lean project. The printed proof of bounded approximate feasibility uses the weak-subsequence lemma below, whose argument remains written background material. A formal counterpart does not by itself verify that printed proof route. The later budget, regularization-limit, noise, and solver-certificate results also have formal counterparts under their stated hypotheses, as do the explicit infinite diagonal estimates, asymptotic budget constant, and rounded noisy rate. Exact finite models have separate executable checks. The Hilbert and regularization background is classical \cite{KS:clason-inverse}. No novelty is claimed for weak compactness or quadratic regularization.
\end{scopebox}

\subsection{The correction space is the protected kernel}
Let $E,Y,Z$ be real Hilbert spaces, let $d:E\to Y$ and $p:E\to Z$ be bounded linear maps, and fix a raw candidate $f\in E$. Write
\[
 H=\ker p,\qquad T=d|_H:H\to Y,\qquad r=df.
\]
The space $H$ is a closed Hilbert subspace with the norm inherited from $E$. An allowed correction is $a\in H$, and its repaired candidate is $f-a$. Thus
\[
 p(f-a)=pf,\qquad d(f-a)=r-Ta.
\]
The protected exact problem is precisely $Ta=r$. Restricting the domain before optimizing is important. A penalty on $\KSnorm{pa}$ would only discourage a change in the protected data, not prohibit it.

The adjoint below always means $T^*:Y\to H$, taken with the inherited Hilbert metric on $H$. In ambient coordinates this is the projection of $d^*$ onto $H$. A coordinate basis of $\ker p$ need not be orthonormal. Its Gram matrix must then be included in any calculation of correction energy.

\subsection{The compactness needed for the existence argument}
\begin{KSlemma}[Bounded sequences have weakly convergent subsequences]\label{KS:lem:weak-subsequence}
Every bounded sequence $(a_n)$ in a real Hilbert space has a subsequence converging weakly to some $a$. Moreover,
\[
 \KSnorm a\le\liminf_k\KSnorm{a_{n_k}}.
\]
\end{KSlemma}
\begin{proof}
Choose $M\ge0$ with $\KSnorm{a_n}\le M$. If $M=0$, every term is zero and the conclusion is immediate. Otherwise let $H_0$ be the closed linear span of the countably many vectors $a_n$. Finite rational linear combinations of those vectors form a countable dense set, so $H_0$ is separable. Apply Gram--Schmidt to a countable dense sequence, omitting zero remainders. This gives either a finite orthonormal basis or a countable orthonormal sequence $(e_j)$ whose closed span is $H_0$. The ambient Hilbert space itself need not be separable.

Each scalar sequence $\langle a_n,e_j\rangle$ is bounded by $M$. For the first coordinate choose a convergent subsequence, for the second choose a subsequence of that subsequence, and continue. Taking the $k$th term from the $k$th nested subsequence, with increasing original indices, yields $a_{n_k}$ for which every fixed coordinate converges. Call the coordinate limits $c_j$. In the finite-dimensional case only finitely many successive selections are required.

For each finite $J$, Bessel's inequality and passage to the limit give
\[
 \sum_{j=1}^Jc_j^2
 =\lim_k\sum_{j=1}^J\langle a_{n_k},e_j\rangle^2\le M^2.
\]
The nondecreasing partial sums are therefore bounded. Orthogonality makes the partial vector sums $\sum_{j\le J}c_je_j$ Cauchy, so completeness defines $a\in H_0$ with $\KSnorm a\le M$.

To verify weak convergence against $h\in H_0$, choose a finite basis combination $h_J$ close to $h$. Then
\[
 |\langle a_{n_k}-a,h\rangle|
 \le |\langle a_{n_k}-a,h_J\rangle|
      +2M\KSnorm{h-h_J}.
\]
For a prescribed tolerance, choose $h_J$ first to control the second term and then $k$ to control the first by finitely many coordinate limits. For a general $h$ in the ambient space, write $h=P_{H_0}h+h_\perp$. The orthogonal component pairs to zero with every $a_{n_k}$ and with $a$, reducing to the preceding case. This proves weak convergence against every test vector.

If $a\ne0$, testing against $a/\KSnorm a$ gives
\[
 \KSnorm a=\lim_k\langle a_{n_k},a/\KSnorm a\rangle
 \le\liminf_k\KSnorm{a_{n_k}}.
\]
The zero limit case follows from nonnegativity. Thus the lower-semicontinuity conclusion concerns the actual norms on the selected subsequence, not only the original bound $M$.
\end{proof}

\begin{KStheorem}[Bounded approximate feasibility]\label{KS:WR01}
For a bounded linear map $T:H\to Y$ between real Hilbert spaces and $r\in Y$, the following are equivalent:
\[
 \begin{split}
 &r\in\KSran T,\\
 &\text{there is a bounded sequence }(a_n)\subset H
       \text{ with }\KSnorm{Ta_n-r}\longrightarrow0.
 \end{split}
\]
If all $\KSnorm{a_n}\le M$, an exact solution exists with norm at most $M$.
\end{KStheorem}
\begin{proof}
If $Ta=r$, the constant sequence $a_n=a$ is bounded and has zero residual. For the converse, let $\KSnorm{a_n}\le M$ and $Ta_n\to r$ in norm. By the preceding lemma there is a subsequence $a_{n_k}$ converging weakly to $a\in H$ with $\KSnorm a\le M$.

For every fixed $y\in Y$, the adjoint identity and weak convergence give
\[
 \langle Ta,y\rangle
 =\langle a,T^*y\rangle
 =\lim_k\langle a_{n_k},T^*y\rangle
 =\lim_k\langle Ta_{n_k},y\rangle.
\]
The last limit equals $\langle r,y\rangle$ because norm convergence implies convergence of each continuous pairing. Thus $\langle Ta-r,y\rangle=0$ for every $y$. Taking $y=Ta-r$ gives $\KSnorm{Ta-r}^2=0$, hence $Ta=r$.

This proof does not claim that $Ta_{n_k}$ converges in norm merely because $a_{n_k}$ converges weakly. Its norm convergence to $r$ was a separate hypothesis. The boundedness of the candidate corrections is used to obtain a weakly convergent subsequence, and their accurate residuals are used to identify its image exactly. The bound $\KSnorm a\le M$ comes from weak lower semicontinuity, so the exact answer stays inside the same closed budget ball.
\end{proof}
\KSwritten{WR01}

No closed range assumption appears. The theorem concerns a single datum and a bounded family of its corrections. Closed range is the stronger, uniform statement that every attainable datum admits a correction controlled by one residual gain.

An immediate contrapositive is useful. If $r\in\overline{\KSran T}\setminus\KSran T$, then every sequence with $Ta_n\to r$ satisfies $\KSnorm{a_n}\to\infty$. Otherwise infinitely many terms would lie in one bounded ball, contradicting the theorem along that subsequence. This recovers the escape mechanism of Section~\ref{KS:sec:escape} without using the diagonal formula.

\subsection{Existence does not make an arbitrary sequence converge strongly}
Take $H=Y=\ell^2$ and $Te_n=e_n/n$. The vectors $a_n=e_n$ have norm one and satisfy $Ta_n\to0$. An exact answer exists, namely zero. Yet $\KSnorm{e_n-e_m}=\sqrt2$ for $n\ne m$, so no subsequence converges in norm.

The conclusion of Theorem~\ref{KS:WR01} is the existence of a weak limit along a subsequence that solves the equation. It is not summable movement, a strong convergence rate, or a selection rule for the original candidates. The minimum correction for the zero datum is zero, which this particular sequence never approaches in norm.

\subsection{Minimum energy finite repairs do converge strongly}
\begin{figure}[htbp]\centering\resizebox{.96\linewidth}{!}{
% BEGIN retained input sources/monographs/komorebi-space/source/figures/31_nested_energy.tex
\includegraphics[width=\linewidth]{sources/monographs/komorebi-space/source/figures/31_nested_energy.pdf}

}\caption{Minimum energy finite corrections behave differently for two data under the same diagonal operator. Bounded partial energy produces a strong Hilbert realization. Escaping partial energy accompanies a coherent product-space point outside the Hilbert space. The attainable energy lies below the explicit bound two.}\label{KS:fig:nested-energy}\end{figure}

Suppose now that $A_1\supseteq A_2\supseteq\cdots$ are nonempty closed affine subsets of a real Hilbert correction space $H$. A typical $A_n$ consists of corrections satisfying the first $n$ bounded linear observation equations. The sets are affine because the equations are linear, and closed because the observations are continuous.

Let $a_n$ be the unique point of $A_n$ nearest to zero, and put $m_n=\KSnorm{a_n}$. The Hilbert projection theorem supplies these points. It can also be proved directly: a sequence approaching the minimum norm is Cauchy by the parallelogram identity, its limit belongs to the closed convex set, and the same identity gives uniqueness.

\begin{KStheorem}[A uniform finite budget yields a strong realization]\label{KS:WR02}
For the nested affine sets above, $m_n$ is nondecreasing, and for $m\ge n$,
\begin{equation}\label{KS:eq:nested-energy}
 \KSnorm{a_m-a_n}^2=m_m^2-m_n^2.
\end{equation}
Their intersection is nonempty if and only if $\sup_n m_n<\infty$. In that case $a_n$ converges in norm to the unique least norm point $a_*\in\bigcap_n A_n$, with
\[
 \KSnorm{a_*-a_n}^2=\KSnorm{a_*}^2-\KSnorm{a_n}^2.
\]
If the intersection is empty, then $m_n\to\infty$.
\end{KStheorem}
\begin{proof}
The nearest-point theorem for nonempty closed convex subsets of a Hilbert space, proved in Appendix~\ref{KS:app:proof-tools}, supplies $a_n$. Write $A_n=a_n+L_n$, where $L_n$ is its closed direction subspace. For $\ell\in L_n$ and every real $t$, the vector $a_n+t\ell$ belongs to $A_n$. Minimality implies
\[
 0\le\KSnorm{a_n+t\ell}^2-\KSnorm{a_n}^2
 =2t\langle a_n,\ell\rangle+t^2\KSnorm\ell^2.
\]
Allowing both signs of arbitrarily small $t$ forces $\langle a_n,\ell\rangle=0$.

If $m\ge n$, nesting gives $a_m\in A_n$, so $a_m-a_n\in L_n$. Expanding its norm therefore gives
\[
 \KSnorm{a_m}^2
 =\KSnorm{a_n+(a_m-a_n)}^2
 =\KSnorm{a_n}^2+\KSnorm{a_m-a_n}^2.
\]
This is the exact energy identity. It also makes $m_n^2$ nondecreasing.

Suppose these energies are bounded, and let $L=\lim_n m_n^2$. Given $\epsilon>0$, choose $N$ with $L-m_N^2<\epsilon^2$. For $m\ge n\ge N$,
\[
 \KSnorm{a_m-a_n}^2=m_m^2-m_n^2\le L-m_N^2<\epsilon^2.
\]
Thus $(a_n)$ is Cauchy and has a norm limit $a_*$. For each fixed $n$, all terms from that index onward belong to the closed set $A_n$, so $a_*\in A_n$. This places $a_*$ in their entire intersection.

If $a$ lies in the intersection, then $\KSnorm a\ge\KSnorm{a_n}$ for every $n$. Norm continuity gives $\KSnorm a\ge\KSnorm{a_*}$. Uniqueness follows from the nearest-point theorem applied to the intersection, which is closed and convex and now known to be nonempty. Alternatively, the midpoint of two different minimizers would have strictly smaller norm by the parallelogram identity. Taking $m\to\infty$ in the finite energy identity yields
\[
 \KSnorm{a_*-a_n}^2=\KSnorm{a_*}^2-\KSnorm{a_n}^2.
\]
Conversely, any element of the intersection bounds every minimum norm from above. Therefore an empty intersection forces the nondecreasing sequence $m_n$ to be unbounded. A nondecreasing unbounded nonnegative sequence tends to infinity, giving the last conclusion.
\end{proof}
\KSwritten{WR02}

The finite answers need not extend one another coordinate by coordinate. Their shared feature is minimum energy in one common ambient Hilbert space. This is the extra geometry that forces their differences to be orthogonal to the earlier minimizers.

For the attainable diagonal datum $r_n=1/n^2$ under $Tu=(u_n/n)$, with $n\ge1$, the first $N$ equations force correction coordinate $n$ to equal $1/n$ for $n\le N$. Their minimum energy completion is zero beyond $N$, and
\[
 m_N^2=\sum_{n=1}^{N}\frac1{n^2}\le2.
\]
These corrections converge strongly to $(1/n)$. For the escaping datum $r_n=1/n$, the corresponding energy is $m_N^2=N$. Every finite equation is again solved, but the budget test now rules out an ambient solution.


\begin{insight}
Finite consistency is not enough by itself. For nested closed affine problems in one Hilbert space, a uniform bound on the minimum correction norms is exactly what upgrades finite feasibility to ambient feasibility. This is a datum-specific test, not a uniform inverse bound for the operator.
\end{insight}



% BEGIN retained input sources/monographs/komorebi-space/source/chapters/05g_budget_certificates.tex
\chap{The attainable residual at a fixed correction budget}\label{KS:sec:budget-frontier}
The unrestricted floor $\KSdist(r,\KSran T)$ forgets how large a correction must be. To retain that information, fix a nonnegative budget $M$ and define
\[
 \beta_{T,r}(M)=\inf_{\KSnorm a\le M}\KSnorm{r-Ta}.
\]
The norm is the actual correction norm in $H=\ker p$. It is not the norm of arbitrary coordinates used to represent that kernel. The function $\beta$ describes the smallest residual compatible with a declared amount of allowed change.

\subsection{At every finite budget the residual minimum exists}
\begin{KStheorem}[Attainment and shape of the budget profile]\label{KS:WR03}
For a bounded linear map between real Hilbert spaces, the infimum defining $\beta(M)$ is attained for every $M\ge0$. The profile is nonincreasing and convex. For $0\le M_1\le M_2$,
\[
 0\le\beta(M_1)-\beta(M_2)\le\KSnorm T(M_2-M_1).
\]
Writing $C=\overline{\KSran T}$, $b=P_Cr$, and $q=\KSnorm{r-b}$, one has
\[
 \lim_{M\to\infty}\beta(M)=q.
\]
The equality $\beta(M)=q$ at a finite budget is equivalent to the existence of $a$ with $Ta=b$ and $\KSnorm a\le M$. In particular,
\[
 \beta(M)=0\quad\Longleftrightarrow\quad
 \exists a\in H:\ Ta=r,\ \KSnorm a\le M.
\]
\end{KStheorem}
\begin{proof}
The feasible ball is nonempty because it contains zero, and the infimum lies between zero and $\KSnorm r$. Choose $a_n$ in that ball with $\KSnorm{r-Ta_n}\to\beta(M)$. A weakly convergent subsequence has limit $a$ with $\KSnorm a\le M$. Bounded linearity makes $r-Ta_{n_k}$ converge weakly to $r-Ta$. Applying weak lower semicontinuity in the target gives
\[
 \KSnorm{r-Ta}\le\liminf_k\KSnorm{r-Ta_{n_k}}=\beta(M).
\]
The reverse inequality holds by the definition of the infimum, so $a$ attains it. When $M=0$ the only feasible correction is zero, and the same conclusion reads $\beta(0)=\KSnorm r$.

For convexity, take minimizers $a_1,a_2$ at budgets $M_1,M_2$ and $t\in[0,1]$. The candidate $a_t=ta_1+(1-t)a_2$ satisfies
\[
 \KSnorm{a_t}\le tM_1+(1-t)M_2,
\]
and its residual is $t(r-Ta_1)+(1-t)(r-Ta_2)$. The triangle inequality therefore gives
\[
 \beta(tM_1+(1-t)M_2)\le t\beta(M_1)+(1-t)\beta(M_2).
\]
Monotonicity follows from inclusion of balls. If $0\le M_1\le M_2$ and $M_2>0$, choose a minimizer $a_2$ at $M_2$ and put $a_1'=(M_1/M_2)a_2$. Then $\KSnorm{a_1'}\le M_1$ and $\KSnorm{a_2-a_1'}\le M_2-M_1$, so
\[
 \beta(M_1)\le\KSnorm{r-Ta_1'}
 \le\beta(M_2)+\KSnorm T(M_2-M_1).
\]
Together with monotonicity this proves the Lipschitz estimate. The omitted case $M_2=0$ has $M_1=0$ and is equality.

Since $r-b$ is orthogonal to $C$ and $b-Ta\in C$, every correction satisfies
\begin{equation}\label{KS:eq:range-decomposition}
 \KSnorm{r-Ta}^2=q^2+\KSnorm{b-Ta}^2.
\end{equation}
Thus $\beta(M)\ge q$. Given $\epsilon>0$, range density in $C$ supplies one correction $a$ with $\KSnorm{b-Ta}<\epsilon$. For every $M\ge\KSnorm a$ it follows that
\[
 q\le\beta(M)\le\sqrt{q^2+\epsilon^2}.
\]
First letting the budget increase and then letting $\epsilon\downarrow0$ proves the limit.

Finally, if $\beta(M)=q$, choose an attaining correction in that ball. The orthogonal identity forces $\KSnorm{b-Ta}^2=0$, hence $Ta=b$. Conversely such a correction attains residual $q$. Setting $q=0$ proves the exact-feasibility criterion. Attainment at fixed budget is essential to the forward implication: a zero unrestricted infimum would not itself supply a preimage.
\end{proof}
\KSwritten{WR03}

Although $\KSran T$ may not be closed, the image $T(\{a:\KSnorm a\le M\})$ of each closed budget ball is closed and convex. To see closedness, take a convergent sequence in that image, lift it to the bounded ball, and apply the weak limit argument. The full range is the increasing union of these closed images. Passing to an unlimited budget can therefore lose closedness even though every fixed-budget optimization is well posed.

The nearest reachable image $y_M=P_{T(\{a:\KSnorm a\le M\})}r$ is unique, so the best residual vector $r-y_M$ is unique as well. Corrections mapping to $y_M$ can differ by vectors in $\ker T$. Among those corrections there is a unique least norm correction. These two uniqueness statements should not be confused with injectivity of $T$.

\subsection{A dual witness can rule out every correction in the budget}
\begin{KStheorem}[A sharp budget certificate]\label{KS:WR04}
For every $M\ge0$,
\begin{equation}\label{KS:eq:budget-dual}
 \beta_{T,r}(M)=
 \max_{\KSnorm z\le1}\bigl(\langle r,z\rangle-M\KSnorm{T^*z}\bigr).
\end{equation}
Thus any single $z$ with $\KSnorm z\le1$ and
\[
 \langle r,z\rangle-M\KSnorm{T^*z}>\varepsilon
\]
certifies that no correction of norm at most $M$ has residual at most $\varepsilon$.
\end{KStheorem}
\begin{proof}
For every $\KSnorm a\le M$ and $\KSnorm z\le1$, the adjoint identity and Cauchy--Schwarz give
\[
\begin{aligned}
 \langle r,z\rangle-M\KSnorm{T^*z}
 &=\langle r-Ta,z\rangle+\langle a,T^*z\rangle-M\KSnorm{T^*z}\\
 &\le\langle r-Ta,z\rangle
 \le\KSnorm{r-Ta}.
\end{aligned}
\]
Taking the infimum over corrections proves that every proposed dual value is at most $\beta(M)$. This part alone already proves the impossibility certificate.

To prove sharpness, let $a_M$ be an attaining correction from Theorem~\ref{KS:WR03} and put $v=r-Ta_M$. If $v=0$, then $\beta(M)=0$ and $z=0$ attains the dual bound. Assume henceforth $v\ne0$ and set $z=v/\KSnorm v$.

Fix any feasible $a$ and put $h=a-a_M$. The segment $a_M+th$ is feasible for $0\le t\le1$. Minimality gives, for $t>0$,
\[
 0\le\KSnorm{v-tTh}^2-\KSnorm v^2
   =-2t\langle v,Th\rangle+t^2\KSnorm{Th}^2.
\]
Divide by $t$ and let $t\downarrow0$. It follows that $\langle v,T(a-a_M)\rangle\le0$, hence
\[
 \langle T^*z,a\rangle\le\langle T^*z,a_M\rangle
 \quad(\KSnorm a\le M).
\]
The support value of a Hilbert ball can be computed directly. If $T^*z\ne0$ and $M>0$, it is attained at $a=M T^*z/\KSnorm{T^*z}$ and equals $M\KSnorm{T^*z}$. If either factor is zero, every feasible pairing is zero and the same formula holds. Thus
\[
 \langle T^*z,a_M\rangle=M\KSnorm{T^*z}.
\]
Substituting this equality gives
\[
 \langle r,z\rangle-M\KSnorm{T^*z}
 =\langle r-Ta_M,z\rangle=\KSnorm v=\beta(M).
\]
We have constructed an attaining witness and thereby justified the word ``maximum.'' The proof uses one-sided variations in a convex ball, not an unproved exchange of an infimum and a supremum.
\end{proof}
\KSwritten{WR04}

In finite dimensions this is a checkable inequality, not merely an optimization prescription. A program may propose the correction or the witness. The displayed inequality can be verified independently of the proposal procedure. In a protected ambient calculation, use the restricted adjoint, not $d^*$ without projection onto $\ker p$.

\subsection{A two-dimensional calculation of the supporting witness}
Take
\[
 T=\begin{pmatrix}1&0\\0&1/2\end{pmatrix},\qquad r=(1,1),\qquad
 a=(1/2,2/5),\qquad v=r-Ta=(1/2,4/5).
\]
Then $T^*v=a$. At the budget $M=\KSnorm a=\sqrt{41}/10$, the normalized vector $z=v/\KSnorm v$ gives
\[
 \langle r,v\rangle-\KSnorm a\KSnorm{T^*v}
 =\frac{13}{10}-\frac{41}{100}=\frac{89}{100}=\KSnorm v^2.
\]
Consequently $a$ is a globally optimal budgeted correction and $\beta(M)=\sqrt{89}/10$. It is not necessary to compare it with a dense numerical grid of competing vectors. The supporting witness proves optimality over the entire ball.

\begin{figure}[htbp]\centering\resizebox{.96\linewidth}{!}{
% BEGIN retained input sources/monographs/komorebi-space/source/figures/32_budget_certificate.tex
\includegraphics[width=\linewidth]{sources/monographs/komorebi-space/source/figures/32_budget_certificate.pdf}

}\caption{The image of a Euclidean correction ball is an ellipse. The nearest reachable point has a residual normal to a supporting line. That same residual, normalized, is the dual witness in the exact budget formula. The plotted point and budget are the rational example in the text.}\label{KS:fig:budget-certificate}\end{figure}

\subsection{How measurement and model uncertainty affect the budget}
\begin{KStheorem}[A robust budget profile]\label{KS:WR05}
Suppose the bounded linear map $\widetilde T:H\to Y$ and the datum $\widetilde r\in Y$ satisfy
\[
 \KSnorm{T-\widetilde T}\le\eta,\qquad
 \KSnorm{r-\widetilde r}\le\delta,\qquad\eta,\delta\ge0.
\]
For every $M\ge0$,
\[
 \KSabs{\beta_{T,r}(M)-\beta_{\widetilde T,\widetilde r}(M)}
 \le\delta+\eta M.
\]
In particular, for $\KSnorm z\le1$,
\[
 \beta_{T,r}(M)\ge
 \langle\widetilde r,z\rangle-M\KSnorm{\widetilde T^*z}
       -\delta-\eta M.
\]
\end{KStheorem}
\begin{proof}
Fix a feasible correction $a$. The two residual vectors have difference
\[
 (r-Ta)-(\widetilde r-\widetilde Ta)
 =(r-\widetilde r)-(T-\widetilde T)a,
\]
whose norm is at most $c=\delta+\eta M$. The reverse triangle inequality therefore gives both inequalities
\[
 \KSnorm{r-Ta}\le\KSnorm{\widetilde r-\widetilde Ta}+c,
 \qquad
 \KSnorm{\widetilde r-\widetilde Ta}\le\KSnorm{r-Ta}+c.
\]
The domain of feasible corrections is identical in the two problems. Taking the infimum of each inequality over that common domain yields
\[
 \beta_{T,r}(M)\le\beta_{\widetilde T,\widetilde r}(M)+c,
 \qquad
 \beta_{\widetilde T,\widetilde r}(M)\le\beta_{T,r}(M)+c.
\]
These are precisely the absolute-value bound. One could equivalently evaluate at the two attained minimizers, but no relation between those minimizers is needed.

For any admissible $z$, Theorem~\ref{KS:WR04} for the nominal problem gives
\[
 \beta_{\widetilde T,\widetilde r}(M)
 \ge\langle\widetilde r,z\rangle-M\KSnorm{\widetilde T^*z}.
\]
Combine this with the lower half of the preceding comparison. The result is a certificate against every correction in the original fixed ball, reduced by the explicit uncertainty allowance. If the correction space itself changed, the common-domain step would fail and this proof would no longer apply.
\end{proof}
\KSwritten{WR05}

This estimate identifies a practical boundary of increasingly large corrections. A model error $\eta$ is amplified by the correction budget even if the measured residual is tiny. The candidate-level inequality is
\[
 \KSnorm{r-Ta}\le\KSnorm{\widetilde r-\widetilde Ta}+\delta+\eta\KSnorm a.
\]
The theorem compares maps on one fixed correction space. It does not cover an uncertain preservation map whose kernel itself changes.

\begin{figure}[htbp]\centering\resizebox{.96\linewidth}{!}{
% BEGIN retained input sources/monographs/komorebi-space/source/figures/33_budget_uncertainty.tex
\includegraphics[width=\linewidth]{sources/monographs/komorebi-space/source/figures/33_budget_uncertainty.pdf}

}\caption{The nominal budget profile is surrounded by the rigorous uncertainty band $\delta+\eta M$. A large budget can improve the nominal fit while enlarging the uncertainty in that conclusion. The example uses $T=1$, $r=1$, $\delta=1/20$, and $\eta=1/10$.}\label{KS:fig:budget-uncertainty}\end{figure}



% BEGIN retained input sources/monographs/komorebi-space/source/chapters/05h_regularized_repair.tex

\chap{Regularization selects a protected correction}\label{KS:sec:regularized-repair}
A fixed budget asks how small the residual can become without exceeding an allowed correction size. Quadratic regularization reverses the emphasis. It asks for the best compromise between the residual and the size of the change. The resulting problem has a unique answer even when the exact protected equation has no solution.

That last sentence is both its usefulness and its limitation. The existence of a regularized answer is not a proof of exact realizability. The limiting behavior, rather than the mere existence of the minimizer, supplies the diagnosis.

\subsection{The normal equation is derived in the protected space}
For $\alpha>0$, define on $H=\ker p$
\[
 J_\alpha(a)=\KSnorm{Ta-r}^2+\alpha\KSnorm a^2.
\]
Both norms are the declared Hilbert norms. All minimizers below are corrections in $H$, so preservation is exact for every parameter, not just in the limit. This is ordinary quadratic Tikhonov regularization applied to the actual protected correction operator. The standard regularization background is discussed in \cite{KS:clason-inverse}. The proof here does not require $T$ to be compact.

\begin{KStheorem}[A unique regularized correction]\label{KS:WR06}
For every $\alpha>0$, there is a unique minimizer
\[
 a_\alpha=(T^*T+\alpha I)^{-1}T^*r.
\]
It satisfies $T^*(r-Ta_\alpha)=\alpha a_\alpha$. For every $a\in H$,
\begin{equation}\label{KS:eq:regularization-energy}
 J_\alpha(a)-J_\alpha(a_\alpha)
 =\KSnorm{T(a-a_\alpha)}^2+\alpha\KSnorm{a-a_\alpha}^2.
\end{equation}
\end{KStheorem}
\begin{proof}
\emph{Step 1: the normal operator has an actual bounded inverse.} Let $B=T^*T+\alpha I:H\to H$. It is bounded and self-adjoint. For every $h\in H$,
\[
 \langle Bh,h\rangle=\KSnorm{Th}^2+\alpha\KSnorm h^2
 \ge\alpha\KSnorm h^2.
\]
If $h\ne0$, Cauchy--Schwarz and division by $\KSnorm h$ give $\KSnorm{Bh}\ge\alpha\KSnorm h$. The same inequality holds at zero. It implies injectivity of $B$.

To prove closedness of its range, suppose $Bh_n\to y$. Applying the lower bound to $h_n-h_m$ gives
\[
 \KSnorm{h_n-h_m}\le\alpha^{-1}\KSnorm{Bh_n-Bh_m}.
\]
Thus $h_n$ is Cauchy. Completeness of $H$ supplies a limit $h$, and continuity gives $Bh=y$. If $y$ is orthogonal to the range of $B$, then for every $h$,
\[
 0=\langle Bh,y\rangle=\langle h,By\rangle.
\]
Taking $h=By$ gives $By=0$, and injectivity gives $y=0$. A closed subspace with zero orthogonal complement is the whole Hilbert space, so $\KSran B=H$.

The inverse therefore exists. Its linearity follows from uniqueness of preimages, and its bound is $\KSnorm{B^{-1}y}\le\alpha^{-1}\KSnorm y$. This construction does not assume closedness of $\KSran T$. It uses the positive term $\alpha I$ in the normal operator.

\emph{Step 2: define the optimizer and prove its normal equation.} Set $a_\alpha=B^{-1}T^*r$. Then $Ba_\alpha=T^*r$, or equivalently
\[
 T^*(Ta_\alpha-r)+\alpha a_\alpha=0.
\]
Because the construction takes place in $H=\ker p$, it already satisfies the exact protection requirement.

\emph{Step 3: prove global minimality, rather than just stationarity.} Write any proposed correction as $a=a_\alpha+h$ and set $s=Ta_\alpha-r$. Expanding both squared norms gives
\[
\begin{aligned}
 J_\alpha(a_\alpha+h)-J_\alpha(a_\alpha)
 &=\KSnorm{Th}^2+\alpha\KSnorm h^2
     +2\langle s,Th\rangle+2\alpha\langle a_\alpha,h\rangle\\
 &=\KSnorm{Th}^2+\alpha\KSnorm h^2
     +2\langle T^*s+\alpha a_\alpha,h\rangle\\
 &=\KSnorm{Th}^2+\alpha\KSnorm h^2.
\end{aligned}
\]
This is the asserted energy identity. It is nonnegative and, since $\alpha>0$, vanishes only for $h=0$. Therefore $a_\alpha$ is the unique global minimizer. In the zero correction space, the same construction and identities reduce to its single element, zero.
\end{proof}
\KSwritten{WR06}

A declared regularization parameter is part of the problem. A calculation that changes $\alpha$ is solving a different optimization problem, even if the protected data and forward operator stay fixed. On the other hand, the preservation condition itself does not change with $\alpha$.

\subsection{Coordinate bases must retain the correction metric}
In finite dimensions, choose orthonormal coordinates in the ambient candidate and residual spaces. Let $B$ have linearly independent columns spanning $\ker p$. The columns of $B$ need not be orthonormal. Write $a=Bc$. The coefficient vector $c$ is a representation of the correction, not necessarily an isometric one. Set
\[
 S=dB,\qquad G=B^{\mathsf T}B.
\]
The correct coordinate equation is
\begin{equation}\label{KS:eq:gram-regularization}
 (S^{\mathsf T}S+\alpha G)c=S^{\mathsf T}r.
\end{equation}
The matrix $G$ is positive definite because the columns are independent. Replacing $G$ by the identity changes the meaning of correction size unless the chosen basis is orthonormal.

For example, let $d(a_1,a_2)=a_1$, $p(a_1,a_2)=a_1-a_2$, $r=1$, and $\alpha=1$. A basis of $\ker p$ is the column $B=(1,1)^{\mathsf T}$. Here $S=1$ and $G=2$, so the correction is $(1/3,1/3)$. Penalizing the coordinate with $G=1$ instead gives $(1/2,1/2)$. The latter is not the optimizer for the ambient Euclidean norm. The difference is a metric change, not a discrepancy between equivalent calculations.

The executable example in this edition constructs the nullspace basis over the rationals, retains its Gram matrix, solves \eqref{KS:eq:gram-regularization}, and returns both the correction and an exact normal-equation residual. It also handles a zero-dimensional protected kernel without pretending that a nonzero correction exists.

\subsection{Parameter monotonicity and the fixed-budget problem}
\begin{KSproposition}[Regularized answers lie on the budget frontier]\label{KS:WR07}
If $0<\alpha<\gamma$, then
\[
 \KSnorm{a_\alpha}\ge\KSnorm{a_\gamma},\qquad
 \KSnorm{r-Ta_\alpha}\le\KSnorm{r-Ta_\gamma}.
\]
At $M_\alpha=\KSnorm{a_\alpha}$,
\[
 \beta(M_\alpha)=\KSnorm{r-Ta_\alpha}.
\]
When $v_\alpha=r-Ta_\alpha\ne0$, its normalization is an optimizing witness in Theorem~\ref{KS:WR04} at that budget.
\end{KSproposition}
\begin{proof}
Let $A_\alpha=\KSnorm{a_\alpha}^2$ and $R_\alpha=\KSnorm{r-Ta_\alpha}^2$. Comparing the two admissible minimizers in each other's objectives gives
\[
 R_\alpha+\alpha A_\alpha\le R_\gamma+\alpha A_\gamma,
 \qquad R_\gamma+\gamma A_\gamma\le R_\alpha+\gamma A_\alpha.
\]
Add these inequalities and cancel the residual terms. The result is
\[
 (\gamma-\alpha)(A_\alpha-A_\gamma)\ge0.
\]
Since $\gamma>\alpha$, it follows that $A_\alpha\ge A_\gamma$. Return to the first comparison to obtain
\[
 R_\alpha-R_\gamma\le\alpha(A_\gamma-A_\alpha)\le0.
\]
Taking square roots, which preserves order for nonnegative numbers, proves both monotonicity statements.

At $M_\alpha=\KSnorm{a_\alpha}$, the vector $a_\alpha$ is budget-feasible. If another correction $a$ in that ball had strictly smaller residual norm, then
\[
 J_\alpha(a)
 =\KSnorm{r-Ta}^2+\alpha\KSnorm a^2
 <\KSnorm{r-Ta_\alpha}^2+\alpha\KSnorm{a_\alpha}^2,
\]
contradicting its defining minimality. Thus it attains the budget frontier.

For the witness statement, let $v=r-Ta_\alpha\ne0$. The normal equation gives $T^*v=\alpha a_\alpha$. Also $r=v+Ta_\alpha$, hence
\[
 \langle r,v\rangle
 =\KSnorm v^2+\langle a_\alpha,T^*v\rangle
 =\KSnorm v^2+\alpha\KSnorm{a_\alpha}^2.
\]
For $z=v/\KSnorm v$, the support penalty is
\[
 M_\alpha\KSnorm{T^*z}
 =\frac{\alpha\KSnorm{a_\alpha}^2}{\KSnorm v}.
\]
Subtracting it from $\langle r,z\rangle$ yields $\KSnorm v=\beta(M_\alpha)$, so this is an optimizing witness. This calculation also covers a zero correction $a_\alpha$ with nonzero residual. If $v=0$, normalization is unavailable, but the zero witness attains the zero budget residual.
\end{proof}
\KSwritten{WR07}

This links a penalized calculation to a hard-budget certificate. It does not say that every larger budget must produce a different answer. Once the least-squares floor is attained, excess budget can be irrelevant. If $T^*r=0$, every regularized correction is zero and the entire profile is constant.

\subsection{The limit separates projection from actual realizability}
Keep $C=\overline{\KSran T}$, $b=P_Cr$, and $q=\KSnorm{r-b}$. Since $r-b\in\ker T^*$, the regularized corrections depend on $b$, not on the component that is orthogonal to the closure of all attainable residuals.

\begin{KStheorem}[The regularization limit alternative]\label{KS:WR08}
As $\alpha\downarrow0$,
\[
 Ta_\alpha\longrightarrow b,\qquad
 \KSnorm{r-Ta_\alpha}\longrightarrow q,\qquad
 \alpha\KSnorm{a_\alpha}^2\longrightarrow0.
\]
If $b\in\KSran T$, let $a^\dagger$ be its unique minimum norm preimage. Then
\[
 a_\alpha\longrightarrow a^\dagger\quad\text{in norm}.
\]
If $b\notin\KSran T$, then $\KSnorm{a_\alpha}\to\infty$.
\end{KStheorem}
\begin{proof}
\emph{Step 1: identify the unavoidable component.} Since $r-b\perp\overline{\KSran T}$, the adjoint identity gives $T^*(r-b)=0$. Thus $a_\alpha(r)=a_\alpha(b)$. For every correction the orthogonal decomposition is
\[
 \KSnorm{r-Ta}^2=q^2+\KSnorm{b-Ta}^2.
\]
In particular both the residual square and the minimum objective are at least $q^2$.

\emph{Step 2: take the small-parameter limit in the correct order.} Let $\epsilon>0$. Since $b$ lies in the closure of the range, choose a single $a\in H$ with $\KSnorm{r-Ta}^2<q^2+\epsilon$. This $a$ is chosen before $\alpha$. Minimality gives
\[
 q^2\le\KSnorm{r-Ta_\alpha}^2
 \le J_\alpha(a_\alpha)
 \le q^2+\epsilon+\alpha\KSnorm a^2.
\]
For sufficiently small $\alpha$, the last term is below $q^2+2\epsilon$ (and if $a=0$ it is already below $q^2+\epsilon$). Hence both intervening quantities tend to $q^2$. Their difference is $\alpha\KSnorm{a_\alpha}^2$, which tends to zero. The orthogonal decomposition now implies $\KSnorm{Ta_\alpha-b}^2\to0$, and taking square roots gives the first two limits.

\emph{Step 3: when the projected datum is attained, obtain a uniform bound.} Suppose $b\in\KSran T$. The preimage set is a nonempty closed affine subset, so projection of zero supplies its unique minimum-norm element $a^\dagger$. Comparing with it and subtracting $q^2$ gives
\[
 \KSnorm{Ta_\alpha-b}^2+\alpha\KSnorm{a_\alpha}^2
 \le\alpha\KSnorm{a^\dagger}^2.
\]
Thus $\KSnorm{a_\alpha}\le\KSnorm{a^\dagger}$ for every $\alpha>0$. This bound does not follow merely from the previously proved vanishing weighted penalty.

\emph{Step 4: strengthen weak convergence to norm convergence.} Take an arbitrary sequence $\alpha_n\downarrow0$, or more generally positive $\alpha_n\to0$. The bound supplies a weakly convergent subsequence $a_{\alpha_{n_k}}\rightharpoonup a$. Since its images converge in norm to $b$, the adjoint pairing argument of Theorem~\ref{KS:WR01} yields $Ta=b$. Weak lower semicontinuity gives
\[
 \KSnorm a\le\liminf_k\KSnorm{a_{\alpha_{n_k}}}
 \le\KSnorm{a^\dagger}.
\]
Minimality and uniqueness of $a^\dagger$ force $a=a^\dagger$. Moreover the same lower bound, together with the uniform upper bound, gives convergence of the norms to $\KSnorm{a^\dagger}$. In a real Hilbert space,
\[
 \KSnorm{a_{\alpha_{n_k}}-a^\dagger}^2
 =\KSnorm{a_{\alpha_{n_k}}}^2+\KSnorm{a^\dagger}^2
       -2\langle a_{\alpha_{n_k}},a^\dagger\rangle\longrightarrow0.
\]
Thus this subsequence converges strongly. If the full positive-parameter family failed to converge strongly, there would be an $\epsilon>0$ and parameters $\alpha_n\to0$ whose corrections all stayed at distance at least $\epsilon$ from $a^\dagger$. The argument just given would produce a strongly convergent subsequence of them, a contradiction. This proves the limit for the entire parameter family.

\emph{Step 5: in the unattained case the norms tend to infinity.} If $b\notin\KSran T$ and the norms failed to tend to infinity, there would exist $M<\infty$ and positive $\alpha_n\to0$ with $\KSnorm{a_{\alpha_n}}\le M$. Their images still tend to $b$ by Step 2. Theorem~\ref{KS:WR01} would then put $b$ in the actual range, a contradiction. This proves divergence of the norms, not just unboundedness somewhere along the path.
\end{proof}
\KSwritten{WR08}

The vanishing penalty $\alpha\KSnorm{a_\alpha}^2$ is not a boundedness certificate for $a_\alpha$. Its coefficient is tending to zero. In the escaping case, the penalty and residual can both tend to zero while the actual corrections diverge.

\begin{center}
\begin{tabularx}{\linewidth}{@{}p{.17\linewidth}p{.26\linewidth}X@{}}
\toprule Residual floor & Is $b$ attained? & Limit of the protected corrections\\\midrule
$q=0$ & Yes & Strong convergence to the nearest exact protected correction.\\[3pt]
$q=0$ & No & Residual vanishes, correction norms diverge, no exact protected answer.\\[3pt]
$q>0$ & Yes & Strong convergence to the least-squares correction of minimum norm, not an exact answer.\\[3pt]
$q>0$ & No & A positive residual floor and diverging correction norms occur together.\\\bottomrule
\end{tabularx}
\end{center}

The last case is not hypothetical. Let $Tu=(u_n/n)$ on $\ell^2$, define $Su=(Tu,0)$ into $\ell^2\oplus\KSR$, and take $r=((1/n),1)$. Then $b=((1/n),0)$ is in the closure but not the range, while $q=1$. The same escaping corrections as before coexist with an irreducible orthogonal residual. Positive floor and unbounded correction cost are independent features.

\subsection{Noise is controlled at a fixed parameter}
\begin{figure}[htbp]\centering\resizebox{.96\linewidth}{!}{
% BEGIN retained input sources/monographs/komorebi-space/source/figures/34_noise_filter.tex
\includegraphics[width=\linewidth]{sources/monographs/komorebi-space/source/figures/34_noise_filter.pdf}

}\caption{For one singular value $\sigma$, the response to residual-data noise is $\sigma/(\sigma^2+\alpha)$. Its maximum occurs at $\sigma=\sqrt\alpha$ and equals $1/(2\sqrt\alpha)$. The plot evaluates this formula and the universal horizontal bounds at three explicit positive parameters.}\label{KS:fig:noise-filter}\end{figure}

\begin{KStheorem}[Data sensitivity and a convergent parameter choice]\label{KS:WR09}
For two residual data vectors $r,\widetilde r$ with $\KSnorm{r-\widetilde r}\le\delta$ and the same $T,\alpha$,
\[
 \KSnorm{a_\alpha(r)-a_\alpha(\widetilde r)}
 \le\frac{\delta}{2\sqrt\alpha},\qquad
 \KSnorm{T(a_\alpha(r)-a_\alpha(\widetilde r))}\le\delta.
\]
The coefficient $1/(2\sqrt\alpha)$ is sharp over bounded linear Hilbert operators. If $b=P_Cr\in\KSran T$, $\alpha(\delta)\to0$, and $\delta/\sqrt{\alpha(\delta)}\to0$, then the noisy corrections converge strongly to $a^\dagger$ as $\delta\to0$.
\end{KStheorem}
\begin{proof}
Both corrections are constructed at the same positive parameter and in the same Hilbert correction space. Let $h=a_\alpha(\widetilde r)-a_\alpha(r)$ and $e=\widetilde r-r$. Subtract the two normal equations to obtain
\[
 (T^*T+\alpha I)h=T^*e.
\]
Pair with $h$ and use the adjoint identity:
\[
 \KSnorm{Th}^2+\alpha\KSnorm h^2
 =\langle e,Th\rangle\le\delta\KSnorm{Th}.
\]
Put $t=\KSnorm{Th}\ge0$. If $t>0$, dropping the nonnegative penalty and dividing by $t$ gives $t\le\delta$. If $t=0$, that inequality is already true. Completing the scalar square gives
\[
 \alpha\KSnorm h^2\le\delta t-t^2
 =\frac{\delta^2}{4}-(t-\delta/2)^2
 \le\frac{\delta^2}{4}.
\]
Because $\alpha>0$, taking square roots proves the correction bound. This also explains why the coefficient is $1/(2\sqrt\alpha)$ rather than the coarser $\KSnorm T/\alpha$ obtained by bounding the inverse alone.

For sharpness at a fixed parameter, take $H=Y=\KSR$ and $T=\sqrt\alpha\,I$. The solution map is multiplication by
\[
 \frac{\sqrt\alpha}{\alpha+\alpha}=\frac1{2\sqrt\alpha}.
\]
Any nonzero data perturbation attains the first bound. This is worst-case sharpness across operators at that parameter, not a claim of equality for every fixed operator and every parameter.

For the limiting statement, insert the exact-data correction at the selected parameter:
\[
 \KSnorm{a_{\alpha(\delta)}(\widetilde r)-a^\dagger}
 \le\frac{\delta}{2\sqrt{\alpha(\delta)}}
       +\KSnorm{a_{\alpha(\delta)}(r)-a^\dagger}.
\]
The first term tends to zero by the noise/parameter hypothesis. The second tends to zero by the attained-projection case of Theorem~\ref{KS:WR08}, because the selected parameter tends to zero. Both statements require the same fixed operator and Hilbert metrics. The conclusion is a least-squares correction when the unavoidable residual component is nonzero.
\end{proof}
\KSwritten{WR09}

The theorem does not manufacture an exact solution when $r$ is outside the range. For the positive-floor attainable case it converges to a least-squares correction. For the unattainable projected datum it cannot produce a bounded exact-data limit. The range hypothesis must therefore precede the parameter-choice conclusion.


\begin{insight}
Regularization gives a unique, exactly protected answer at every positive parameter. Whether those answers converge to an exact realization, converge only to a best approximation, or escape is determined by the projected datum's actual range membership. A smooth regularized path does not erase that distinction.
\end{insight}



% BEGIN retained input sources/monographs/komorebi-space/source/chapters/05i_diagonal_regularization.tex
\chap{A complete regularization diagnostic for the escaping model}\label{KS:sec:diagonal-regularization}
The abstract alternatives become explicit for the diagonal model. They also answer a natural objection to the hard prefix corrections: perhaps a smoother correction could fit increasingly well without escaping. It can fit more efficiently, but it cannot avoid the obstruction.

\subsection{The entire regularized family is explicit}
Use $Tu=(u_n/n)$ and $r=w=(1/n)$ as in Section~\ref{KS:sec:escape}. Choose $\alpha=m^{-2}$ for an integer $m\ge1$. The normal equation in each coordinate gives
\[
 a^{(m)}_n=\frac{m^2}{m^2+n^2},\qquad
 v^{(m)}_n=r_n-(Ta^{(m)})_n=\frac{n}{m^2+n^2}.
\]
The repaired ambient candidate is $(-a^{(m)},1)$. Its protected coordinate is one for every $m$.

\begin{KSproposition}[Energy, residual, and certified tails]\label{KS:WR10}
The displayed $a^{(m)}$ belongs to $\ell^2$ and is the unique regularized correction at $\alpha=m^{-2}$. Set
\[
 A_m=\KSnorm{a^{(m)}}^2=\sum_{n\ge1}\frac{m^4}{(m^2+n^2)^2},\qquad
 B_m=\KSnorm{v^{(m)}}^2=\sum_{n\ge1}\frac{n^2}{(m^2+n^2)^2}.
\]
Then
\begin{equation}\label{KS:eq:soft-bounds}
 \frac m4\le A_m\le\frac{4m}{3},\qquad
 \frac1{4(m+1)}\le B_m\le\frac2m.
\end{equation}
For the sums $A_{m,K},B_{m,K}$ through any integer $K\ge m$,
\begin{equation}\label{KS:eq:soft-tails}
 0\le A_m-A_{m,K}\le\frac{m^4}{3K^3},\qquad
 0\le B_m-B_{m,K}\le\frac1K.
\end{equation}
\end{KSproposition}
\begin{proof}
For $n>m$, one has $(a_n^{(m)})^2\le m^4/n^4$, so the correction sequence is square-summable. The diagonal operator is bounded with norm one, and $w=(1/n)\in\ell^2$, hence the residual sequence is also in $\ell^2$. The operator is self-adjoint, and the coordinate equality
\[
 \frac1n\frac{n}{m^2+n^2}
 =m^{-2}\frac{m^2}{m^2+n^2}
\]
proves $T^*v^{(m)}=\alpha a^{(m)}$ as an equality of Hilbert vectors. Theorem~\ref{KS:WR06} therefore identifies this actual infinite vector as the unique minimizer, rather than just a coordinatewise candidate.

For $1\le n\le m$, $m^2/(m^2+n^2)\ge1/2$. There are $m$ such coordinates, so $A_m\ge m/4$. For the upper bound, the first $m$ squares are at most one, and the decreasing function $x^{-4}$ gives
\[
 \sum_{n>m}\frac1{n^4}\le\int_m^\infty x^{-4}\,dx=\frac1{3m^3}.
\]
Thus $A_m\le m+m^4/(3m^3)=4m/3$.

For $n\le m$, the residual square is at most $n^2/m^4$. The sum-of-squares estimate
\[
 \sum_{n=1}^m n^2=\frac{m(m+1)(2m+1)}6\le m^3
\]
holds for $m\ge1$, so this part contributes at most $1/m$. For $n>m$, each residual square is at most $1/n^2$, whose tail is at most $1/m$. This proves $B_m\le2/m$.

For the lower bound, $n>m$ implies $m^2+n^2\le2n^2$, hence
\[
 \frac{n^2}{(m^2+n^2)^2}\ge\frac1{4n^2}.
\]
The other direction of the decreasing-function integral comparison gives
\[
 \sum_{n>m}\frac1{n^2}\ge\int_{m+1}^\infty x^{-2}\,dx=\frac1{m+1}.
\]
This proves $B_m\ge1/(4(m+1))$.

Finally, for $K\ge m$, both omitted tails consist of nonnegative terms. Bound the first by $m^4\sum_{n>K}n^{-4}$ and the second by $\sum_{n>K}n^{-2}$. The same upper integral comparisons give the two certified remainders in \eqref{KS:eq:soft-tails}. No finite truncation is substituted for an equality of infinite energies.
\end{proof}
\KSwritten{WR10}

Thus the correction norm grows on the order of $\sqrt m$, the residual norm decays on the order of $m^{-1/2}$, and the penalty $\alpha A_m=A_m/m^2$ tends to zero. The smooth correction avoids the hard discontinuity at a prefix boundary, but it still approaches the forbidden constant sequence coordinate by coordinate.

\begin{figure}[htbp]\centering\resizebox{.96\linewidth}{!}{
% BEGIN retained input sources/monographs/komorebi-space/source/figures/35_regularized_escape.tex
\includegraphics[width=\linewidth]{sources/monographs/komorebi-space/source/figures/35_regularized_escape.pdf}

}\caption{At each fixed coordinate, the soft correction approaches one as $m$ tends to infinity. Its transition moves outward as $m$ increases. The chart compares three exact coordinate formulas $m^2/(m^2+n^2)$, not finite approximations to an unspecified optimization procedure. Every displayed sequence has finite energy, but the family's energies are unbounded.}\label{KS:fig:regularized-escape}\end{figure}

\subsection{The sharp asymptotic correction cost}
The preceding bounds establish divergence and decay. The leading constant can also be computed, which permits a meaningful comparison with hard prefix repair.

\begin{KStheorem}[A sharp budget asymptotic for the diagonal datum]\label{KS:WR11}
As $m\to\infty$,
\[
 \frac{A_m}{m}\longrightarrow\frac\pi4,\qquad
 mB_m\longrightarrow\frac\pi4.
\]
For the exact budget profile of this infinite problem,
\begin{equation}\label{KS:eq:sharp-budget}
 \lim_{M\to\infty}M\beta_{T,w}(M)=\frac\pi4.
\end{equation}
The hard prefix correction at budget $M=\sqrt N$ instead has
\[
 \sqrt N\,\KSnorm{w-T\mathbf1_N}\longrightarrow1.
\]
\end{KStheorem}
\begin{proof}
\emph{Step 1: control the infinite Riemann sums uniformly.} Define
\[
 f(x)=\frac1{(1+x^2)^2},\qquad
 g(x)=\frac{x^2}{(1+x^2)^2}.
\]
Then $A_m/m=m^{-1}\sum_{n\ge1}f(n/m)$ and $mB_m=m^{-1}\sum_{n\ge1}g(n/m)$. Fix $L\ge2$ and put $N=\lfloor mL\rfloor$. On $[0,L]$, both functions are uniformly continuous. If $\omega_f(h)$ is a modulus of continuity there, comparison on the $N$ intervals of length $1/m$ gives
\[
 \left|\frac1m\sum_{n=1}^N f(n/m)-\int_0^{N/m}f(x)\,dx\right|
 \le L\omega_f(1/m).
\]
The missing interval $[N/m,L]$ has length below $1/m$ and a bounded integrand, so it contributes a term tending to zero. The same argument applies to $g$. This proves convergence of each truncated sum to its integral on $[0,L]$.

The tails require a bound independent of sufficiently fine discretization. Since $N\ge mL/2$, the inequalities $f(x)\le x^{-4}$ and $g(x)\le x^{-2}$ give
\[
 \frac1m\sum_{n>N}f(n/m)
 \le m^3\sum_{n>N}n^{-4}\le\frac{m^3}{3N^3}
 \le\frac8{3L^3},
\]
\[
 \frac1m\sum_{n>N}g(n/m)
 \le m\sum_{n>N}n^{-2}\le\frac mN\le\frac2L.
\]
The integral tails are at most $1/(3L^3)$ and $1/L$. Given an error tolerance, choose $L$ first to control both discrete and integral tails, then choose $m$ to control the compact-interval sums. This justifies passing from finite-interval Riemann convergence to the full half-line sums. No uncontrolled interchange of two infinite limits is involved.

\emph{Step 2: evaluate the integrals.} The substitution $x=\tan\theta$ gives $dx=\sec^2\theta\,d\theta$ and $1+x^2=\sec^2\theta$, so
\[
 \int_0^\infty f(x)\,dx=\int_0^{\pi/2}\cos^2\theta\,d\theta=\frac\pi4,
 \qquad
 \int_0^\infty g(x)\,dx=\int_0^{\pi/2}\sin^2\theta\,d\theta=\frac\pi4.
\]
Step 1 now proves the two stated asymptotics.

\emph{Step 3: move from regularization parameters to arbitrary budgets.} Set $M_m=\sqrt{A_m}$. Every positive coordinate $m^2/(m^2+n^2)$ increases strictly with $m$, so $A_m$, and hence $M_m$, increases strictly. Since $A_m/m\to c=\pi/4>0$, these budgets tend to infinity and
\[
 \left(\frac{M_{m+1}}{M_m}\right)^2
 =\frac{A_{m+1}/(m+1)}{A_m/m}\frac{m+1}{m}\longrightarrow1.
\]
Proposition~\ref{KS:WR07} identifies $\beta(M_m)=\sqrt{B_m}$. Thus
\[
 M_m\beta(M_m)=\sqrt{(A_m/m)(mB_m)}\longrightarrow c.
\]
For every sufficiently large $M$, choose $m$ with $M_m\le M\le M_{m+1}$. Monotonicity of $\beta$ gives
\[
 M_m\beta(M_{m+1})\le M\beta(M)\le M_{m+1}\beta(M_m).
\]
The lower bound is $(M_m/M_{m+1})[M_{m+1}\beta(M_{m+1})]$ and the upper is $(M_{m+1}/M_m)[M_m\beta(M_m)]$. Both tend to $c$, proving the full budget limit, not merely convergence on a selected subsequence of budgets.

\emph{Step 4: compare the hard-prefix method.} That correction has norm $\sqrt N$ and residual square $\sum_{n>N}n^{-2}$. Integral comparison yields
\[
 \frac{N}{N+1}\le N\sum_{n>N}\frac1{n^2}\le1.
\]
Taking the limit and then a square root gives its leading product one. The two constants are consequently compared at large actual correction budgets, rather than at arbitrarily equated parameter labels.
\end{proof}
\KSwritten{WR11}

At large matched correction budgets, the optimal soft repair therefore has about $\pi/4$ times the hard-prefix residual. This improves the efficiency of approximation, not the existence result. Both residuals tend to zero only as the budget tends to infinity. The limiting constant is a calculation for this declared model, not a universal regularization efficiency factor.

\begin{figure}[htbp]\centering\resizebox{.96\linewidth}{!}{
% BEGIN retained input sources/monographs/komorebi-space/source/figures/36_budget_product.tex
\includegraphics[width=\linewidth]{sources/monographs/komorebi-space/source/figures/36_budget_product.pdf}

}\caption{The plotted intervals enclose the actual infinite quantities $\sqrt{A_mB_m}$, using finite sums plus the proved tails. The corresponding hard-prefix intervals enclose $\sqrt{N\sum_{n>N}n^{-2}}$ with $N=m$. These have different natural budget sequences. Their limits give the matched-budget asymptotic constants $\pi/4$ and one.}\label{KS:fig:budget-product}\end{figure}

\subsection{The same operator has an attainable datum}
Change only the residual to $b_n=1/n^2$. Its minimum norm preimage is $a^\dagger_n=1/n\in\ell^2$. At $\alpha=m^{-2}$, the regularized correction is
\[
 c^{(m)}_n=\frac{m^2}{n(m^2+n^2)},\qquad
 a^\dagger_n-c^{(m)}_n=\frac{n}{m^2+n^2}.
\]
The bias square is exactly the previously computed $B_m$. This links the two cases without conflating their existence hypotheses: the escaping residual of one problem is the regularization bias of an attainable problem.

\begin{KScorollary}[An explicit noisy attainable repair]\label{KS:WR12}
For the attainable datum $b=(1/n^2)$, suppose $\KSnorm{\widetilde b-b}\le\delta$ with $0<\delta\le1$. Choose
\[
 m=\lceil\delta^{-2/3}\rceil,\qquad\alpha=m^{-2}.
\]
Then the noisy regularized correction satisfies
\[
 \KSnorm{a_\alpha(\widetilde b)-a^\dagger}
 \le\sqrt{\frac2m}+\frac{\delta m}{2}
 \le(\sqrt2+1)\delta^{1/3}.
\]
\end{KScorollary}
\begin{proof}
The exact-data bias has coordinates $n/(m^2+n^2)$, so its squared norm is precisely $B_m$. By Proposition~\ref{KS:WR10},
\[
 \KSnorm{a_{m^{-2}}(b)-a^\dagger}\le\sqrt{2/m}.
\]
Theorem~\ref{KS:WR09} at the same parameter gives
\[
 \KSnorm{a_{m^{-2}}(\widetilde b)-a_{m^{-2}}(b)}
 \le\frac{\delta}{2\sqrt{m^{-2}}}=\frac{\delta m}{2}.
\]
The triangle inequality proves the first bound.

Because $0<\delta\le1$, the number $u=\delta^{-2/3}$ is at least one. The ceiling satisfies $u\le m<u+1\le2u$. Therefore
\[
 \sqrt{2/m}\le\sqrt2\,\delta^{1/3},
 \qquad \frac{\delta m}{2}\le\delta u=\delta^{1/3}.
\]
Adding these estimates proves the displayed rate with every ceiling factor accounted for. The positive-noise condition is needed to define this particular finite integer choice. The statement does not prescribe an infinite $m$ at $\delta=0$, nor does it use a probabilistic assumption on the perturbation.
\end{proof}
\KSwritten{WR12}

This is a deterministic guarantee for the specified datum and Hilbert noise bound. It is not a statistical confidence statement. The identical parameter rule in the escaping problem cannot produce convergence to an exact protected Hilbert solution, because no such solution exists.

\subsection{What the calculation permits a machine to certify}
The finite program can check a normal equation, exact preservation, and a quadratic optimality identity. It can return a dual witness ruling out a smaller residual at a given budget. In the infinite diagonal example it can return partial sums together with rational tail enclosures. Those are distinct certificates with distinct conclusions.

None of them alone establishes unrestricted semantic realizability. The bounded-budget theorem needs a genuinely uniform bound. The exactification theorem needs control of candidate movement. The regularization limit needs actual attainment of the projected datum. The repeated appearance of a small residual does not replace any of those hypotheses.

The useful next diagnostic is therefore not just ``Did the residual decrease?'' It is ``What happened to the correction norm, what quantity was preserved, and what existence theorem applies to this datum?''



% BEGIN retained input sources/monographs/komorebi-space/source/chapters/05j_inexact_certificates.tex
\chap{Certifying an inexact protected correction}\label{KS:sec:inexact-repair}
An existence theorem for a regularized optimizer does not guarantee that a numerical procedure has reached it. The procedure may stop early, use a poorly scaled basis, or report the residual of the wrong equation. These are mathematical distinctions, not implementation details that can be deferred until after a certificate has been interpreted.

The relevant question is now more specific. Given a proposed correction, what quantities can be checked without first computing the exact optimizer? The answer is a normal-equation residual and, independently, a primal and dual objective gap. Both give quantitative statements about the declared regularized problem. Neither promises that the original unregularized realization problem is solvable.

\begin{scopebox}
The accompanying formalization constructs the positive-parameter optimizer and proves the energy identity, global uniqueness, and the estimates in this chapter, with a separate bridge to the actual protected kernel. Regularization limits and the combined inexact stopping theorem also have formal counterparts. The explicit infinite diagonal results also have formal counterparts. The printed weak-subsequence argument remains written background material. The exact Python certificate verifier is tested, not universally verified as a refinement of Lean.
\end{scopebox}

\subsection{There are three different residuals}
Keep $H=\ker p$, $T=d|_H$, and a residual datum $r$. Fix $\alpha>0$. Write $a_\alpha$ for the unique minimizer of Theorem~\ref{KS:WR06} and let $a\in H$ be any proposed correction. Three quantities now have different meanings:
\[
 v=r-Ta,\qquad
 g=T^*(Ta-r)+\alpha a,\qquad
 e=a-a_\alpha.
\]
The vector $v$ is the fit residual. The vector $g$ is the normal-equation residual, equal to one half of the derivative of the quadratic objective after identifying the dual with $H$. The vector $e$ is the unknown correction error. The exact optimizer generally has $v\ne0$ but always has $g=0$.

The adjoint in $g$ is the restricted adjoint on $H$, not the unrestricted ambient adjoint. Equivalently,
\[
 T^*v=P_Hd^*v.
\]
Thus the correction certificate must measure the part of the normal equation that belongs to the permitted correction space. A numerical program that forgets this projection can incorrectly reject an exact protected optimizer.

\subsection{A posteriori correction and fit bounds}
\begin{KStheorem}[Normal residual as a solver certificate]\label{KS:WR13}
For every proposed $a\in H$,
\begin{align}
 \KSnorm{a-a_\alpha}&\le\frac{\KSnorm g}{\alpha},\label{KS:eq:solver-error}\\
 \KSnorm{T(a-a_\alpha)}&\le\frac{\KSnorm g}{2\sqrt\alpha},\label{KS:eq:solver-fit}\\
 0\le J_\alpha(a)-J_\alpha(a_\alpha)&\le\frac{\KSnorm g^2}{\alpha}.
 \label{KS:eq:solver-objective}
\end{align}
\end{KStheorem}
\begin{proof}
Let $e=a-a_\alpha$ and $B=T^*T+\alpha I$. The normal equation of the constructed optimizer gives
\[
 Be=(T^*T+\alpha I)a-T^*r=g.
\]
Pairing with $e$ yields the exact identity
\begin{equation}\label{KS:eq:solver-energy}
 \KSnorm{Te}^2+\alpha\KSnorm e^2=\langle g,e\rangle
 \le\KSnorm g\KSnorm e.
\end{equation}
If $e=0$, all three conclusions hold with zero error. Otherwise drop the first nonnegative term and divide by $\alpha\KSnorm e>0$ to obtain $\KSnorm e\le\KSnorm g/\alpha$.

To obtain the sharper fitted-data bound, retain the quadratic dependence on $\KSnorm e$ instead of substituting the first estimate. Multiplying \eqref{KS:eq:solver-energy} by $\alpha$ and rearranging gives
\[
\begin{aligned}
 \alpha\KSnorm{Te}^2
 &\le\alpha\KSnorm g\KSnorm e-\alpha^2\KSnorm e^2\\
 &=\frac{\KSnorm g^2}{4}
     -\left(\alpha\KSnorm e-\frac{\KSnorm g}{2}\right)^2
 \le\frac{\KSnorm g^2}{4}.
\end{aligned}
\]
Division by $\alpha$ and square roots give the factor $1/(2\sqrt\alpha)$.

By Theorem~\ref{KS:WR06} the objective gap is exactly $\KSnorm{Te}^2+\alpha\KSnorm e^2$, hence by \eqref{KS:eq:solver-energy} it is also $\langle g,e\rangle$. It is nonnegative and satisfies
\[
 J_\alpha(a)-J_\alpha(a_\alpha)
 \le\KSnorm g\KSnorm e\le\frac{\KSnorm g^2}{\alpha}.
\]
Every certificate quantity on the right is computed from the proposal $a$ and the declared problem. The unknown optimizer was used to prove the guarantee, but it is not required to evaluate the certificate.
\end{proof}
\KSwritten{WR13}

\paragraph{Sharp scalar examples.} For sharpness of the correction-error bound, take $T=0$ on $\KSR$ and a nonzero error. Then $g=\alpha e$. For sharpness of the fitted-data bound, take $T=\sqrt\alpha\,I$ on $\KSR$, so $g=2\alpha e$ and $\KSnorm{Te}=\KSnorm g/(2\sqrt\alpha)$. These examples establish worst-case sharpness, not simultaneous equality for every operator.

For a fixed raw vector $f$, the corresponding repaired candidates are $f-a$ and $f-a_\alpha$. Their distance is exactly $\KSnorm e$. Both preserve $pf$ because both corrections belong to $H$. The theorem therefore controls a genuine protected candidate error, not just an intermediate coordinate vector.

\begin{figure}[htbp]\centering\resizebox{.96\linewidth}{!}{
% BEGIN retained input sources/monographs/komorebi-space/source/figures/37_solver_amplification.tex
\includegraphics[width=\linewidth]{sources/monographs/komorebi-space/source/figures/37_solver_amplification.pdf}

}\caption{A fixed normal-equation tolerance has different consequences at different regularization parameters. The curves are universal certified upper bounds for a declared normal-residual norm, not observed errors of a particular experiment. The correction estimate grows as $\alpha^{-1}$ while the fitted-data estimate grows as $\alpha^{-1/2}$.}\label{KS:fig:solver-amplification}\end{figure}

\subsection{Why a vanishing normal residual is not enough}
Let $T=0$ on $\KSR$, let $r=0$, and take $\alpha_n=n^{-2}$ for $n\ge1$. The exact optimizer is zero for every such $n$. The fixed proposal $a_n=1$ instead has
\[
 v_n=0,\qquad g_n=\alpha_n\longrightarrow0,\qquad
 J_{\alpha_n}(a_n)=\alpha_n\longrightarrow0,
 \qquad\KSabs{a_n-a_{\alpha_n}}=1.
\]
Thus the fit residual, the normal residual, and even the absolute objective value can all tend to zero without convergence to the selected optimizer. The certificate correctly reports $\KSabs{g_n}/\alpha_n=1$.

This is not an infinite-dimensional obstruction. It occurs in one dimension because the regularizing penalty itself disappears. The question is not whether the reported residual is small in isolation, but whether it is small relative to the parameter that converts it into an error bound.

\subsection{A dual certificate does not require solving the normal equation}
For an arbitrary $z\in Y$, define
\[
 \mathcal D_\alpha(z)=2\langle r,z\rangle-\KSnorm z^2
                    -\frac{\KSnorm{T^*z}^2}{\alpha}.
\]
Unlike the fixed-budget witness in Theorem~\ref{KS:WR04}, this dual witness has no unit-norm constraint. Its quadratic penalties supply the control.

\begin{KStheorem}[Exact regularized primal and dual gap]\label{KS:WR14}
For every $a\in H$ and $z\in Y$,
\begin{equation}\label{KS:eq:regularized-duality-gap}
 J_\alpha(a)-\mathcal D_\alpha(z)
 =\KSnorm{Ta-r+z}^2+
   \frac{\KSnorm{T^*z-\alpha a}^2}{\alpha}.
\end{equation}
Consequently $\mathcal D_\alpha(z)\le J_\alpha(a)$ for every pair. At
$a=a_\alpha$ and $z=z_\alpha=r-Ta_\alpha$, equality holds, so the dual maximum is attained and equals the primal minimum. Moreover,
\[
 \alpha\KSnorm{a-a_\alpha}^2
 \le J_\alpha(a)-\mathcal D_\alpha(z).
\]
For the easily computed choice $z=r-Ta$,
\begin{equation}\label{KS:eq:residual-dual-gap}
 J_\alpha(a)-\mathcal D_\alpha(r-Ta)=\frac{\KSnorm g^2}{\alpha}.
\end{equation}
\end{KStheorem}
\begin{proof}
For arbitrary $a,z$, expand the first square as
\[
 \KSnorm{Ta-r+z}^2
 =\KSnorm{Ta-r}^2+\KSnorm z^2+2\langle Ta,z\rangle-2\langle r,z\rangle,
\]
and the scaled second square as
\[
 \frac{\KSnorm{T^*z-\alpha a}^2}{\alpha}
 =\frac{\KSnorm{T^*z}^2}{\alpha}+\alpha\KSnorm a^2
       -2\langle T^*z,a\rangle.
\]
The adjoint identity cancels $2\langle Ta,z\rangle$ against the last term. The remaining expression is
\[
 \KSnorm{Ta-r}^2+\alpha\KSnorm a^2
 -\left(2\langle r,z\rangle-\KSnorm z^2
                    -\alpha^{-1}\KSnorm{T^*z}^2\right),
\]
which is the asserted gap. Its nonnegativity gives weak duality for every pair, without any constraint on the witness norm.

Set $a=a_\alpha$ and $z_\alpha=r-Ta_\alpha$. The first square is zero by this definition. The normal equation gives $T^*z_\alpha=\alpha a_\alpha$, so the second is zero as well. Therefore
\[
 \mathcal D_\alpha(z_\alpha)=J_\alpha(a_\alpha).
\]
For every other $z$, weak duality with $a_\alpha$ gives $\mathcal D_\alpha(z)\le J_\alpha(a_\alpha)$. This proves attainment of the dual maximum and equality with the primal minimum directly.

For a general proposed pair $(a,z)$, decompose its gap as
\[
 J_\alpha(a)-\mathcal D_\alpha(z)
 =[J_\alpha(a)-J_\alpha(a_\alpha)]
  +[J_\alpha(a_\alpha)-\mathcal D_\alpha(z)].
\]
The second bracket is nonnegative by weak duality, while Theorem~\ref{KS:WR06} bounds the first below by $\alpha\KSnorm{a-a_\alpha}^2$. This proves the correction-error certificate.

Finally take $z=r-Ta$. The first square in the gap is zero, while $T^*z-\alpha a=-g$. The second becomes $\KSnorm g^2/\alpha$, proving the residual-witness formula. A small duality gap controls optimizer error only after division by the specified positive parameter.
\end{proof}
\KSwritten{WR14}

A proposed correction and a proposed dual witness can come from different procedures. The verifier only needs the mathematical problem, the two vectors, and the displayed identities. A poor witness remains valid as a lower bound, but can give a loose gap. The residual witness avoids a second optimization routine and recovers the bound in Theorem~\ref{KS:WR13} exactly.

\subsection{A protected example with every certificate calculated}
Take
\[
 d(x,y)=x,\qquad p(x,y)=x-y,\qquad f=(1,0),\qquad r=df=1,\qquad\alpha=1.
\]
Every allowed correction has the form $(t,t)$. The objective is
\[
 J(t,t)=(t-1)^2+2t^2,
\]
so $a_\alpha=(1/3,1/3)$ and $J(a_\alpha)=2/3$. Suppose a procedure returns only $a=(1/4,1/4)$. Then
\[
 v=\frac34,\qquad T^*v=(3/8,3/8),\qquad
 g=(-1/8,-1/8),\qquad\KSnorm g^2=\frac1{32}.
\]
The actual correction-error square is $1/72$, while the certified upper bound is $1/32$. The actual fitted-error square is $1/144$, while the sharper bound in Theorem~\ref{KS:WR13} is $1/128$.

The proposed objective is $11/16$. Its actual excess over the optimum is $1/48$. The residual witness gives
\[
 \mathcal D_1(3/4)=\frac{21}{32},\qquad
 J(a)-\mathcal D_1(3/4)=\frac1{32}.
\]
The duality gap is an upper bound for suboptimality, not necessarily the exact suboptimality. Both repaired candidates preserve $p(f)=1$.

\begin{figure}[htbp]\centering\resizebox{.96\linewidth}{!}{
% BEGIN retained input sources/monographs/komorebi-space/source/figures/38_objective_certificate.tex
\includegraphics[width=\linewidth]{sources/monographs/komorebi-space/source/figures/38_objective_certificate.pdf}

}\caption{The objective on the actual protected line is $3t^2-2t+1$. The proposed correction at $t=1/4$, exact minimizer at $t=1/3$, and residual-witness lower bound are distinct. The vertical gap to the dual level is the checkable certificate. The smaller gap to the minimum is the true suboptimality.}\label{KS:fig:objective-certificate}\end{figure}

The supplied program \path{examples/inexact_repair_certificate.py} verifies this calculation without calling the optimizer. It also accepts a user-supplied dual witness. Every reported scalar is rational, and all norm error bounds are stored as squared quantities so that no rounding of a square root is required.



% BEGIN retained input sources/monographs/komorebi-space/source/chapters/05k_stopping_rules.tex
\chap{Coordinates, noise, and a justified stopping rule}\label{KS:sec:stopping-rules}
A certificate can be exact and still be interpreted incorrectly. Two common errors are to measure a coordinate covector with the wrong norm, or to keep a solver tolerance fixed while sending the regularization parameter to zero. Both change what a small reported number means.

\subsection{The residual of the coordinate equation is a covector}
In finite dimensions, use orthonormal coordinates in the ambient candidate and residual spaces. Let the columns of $B$ form an arbitrary basis of $H=\ker p$, not necessarily an orthonormal one. Write $a=Bc$, $S=dB$, and $G=B^{\mathsf T}B$. As before, $G$ is positive definite. For an inexact coefficient vector $c$, the residual of the coordinate normal equation is
\[
 q=(S^{\mathsf T}S+\alpha G)c-S^{\mathsf T}r.
\]
The vector $q$ stores pairings against basis vectors. It is not the coordinate vector of the Hilbert normal residual unless the basis is orthonormal.

\begin{KSproposition}[The invariant norm of a coordinate certificate]\label{KS:WR15}
The actual normal residual is
\begin{equation}\label{KS:eq:covector-norm}
 g=BG^{-1}q,\qquad \KSnorm g^2=q^{\mathsf T}G^{-1}q.
\end{equation}
For any invertible change of basis $B'=BC$, the same proposed correction has
\[
 c'=C^{-1}c,\quad G'=C^{\mathsf T}GC,\quad q'=C^{\mathsf T}q,
\]
and
\[
 (q')^{\mathsf T}(G')^{-1}q'=q^{\mathsf T}G^{-1}q.
\]
\end{KSproposition}
\begin{proof}
Because $B$ has independent columns, for every nonzero coefficient vector $c$,
\[
 c^{\mathsf T}Gc=c^{\mathsf T}B^{\mathsf T}Bc=\KSnorm{Bc}^2>0.
\]
Thus $G$ is positive definite and invertible. The matrix $P=BG^{-1}B^{\mathsf T}$ is symmetric and obeys
\[
 P^2=BG^{-1}(B^{\mathsf T}B)G^{-1}B^{\mathsf T}=P,
 \qquad PB=B.
\]
Its range is precisely the column space $H$, so it is the orthogonal projection onto the entire permitted correction space.

The restricted Hilbert normal residual is
\[
 g=P d^{\mathsf T}(dBc-r)+\alpha Bc.
\]
Substitute $S=dB$, $G=B^{\mathsf T}B$ and the expression for $P$:
\[
 g=BG^{-1}\bigl(S^{\mathsf T}(Sc-r)+\alpha Gc\bigr)=BG^{-1}q.
\]
Equivalently $q=B^{\mathsf T}g$ records the pairings of $g$ with the basis columns. Its squared Hilbert norm is
\[
 g^{\mathsf T}g
 =q^{\mathsf T}G^{-1}B^{\mathsf T}BG^{-1}q
 =q^{\mathsf T}G^{-1}q,
\]
using symmetry of $G^{-1}$. This proves both coordinate formulas.

For the basis change $B'=BC$, the same ambient vector is represented by $c'=C^{-1}c$. Also $S'=SC$ and $G'=C^{\mathsf T}GC$. Substitute these formulas into the new coordinate normal equation:
\[
\begin{aligned}
 q'&=((S')^{\mathsf T}S'+\alpha G')c'-(S')^{\mathsf T}r\\
 &=C^{\mathsf T}(S^{\mathsf T}S+\alpha G)CC^{-1}c
      -C^{\mathsf T}S^{\mathsf T}r=C^{\mathsf T}q.
\end{aligned}
\]
The inverse metric is $(G')^{-1}=C^{-1}G^{-1}(C^{\mathsf T})^{-1}$. Hence
\[
 (q')^{\mathsf T}(G')^{-1}q'
 =q^{\mathsf T}C C^{-1}G^{-1}(C^{\mathsf T})^{-1}C^{\mathsf T}q
 =q^{\mathsf T}G^{-1}q.
\]
For a zero-dimensional kernel these are empty matrices and zero vectors, and the resulting scalar norm is zero. The spanning assumption is indispensable: a basis of a proper subspace would compute a projection to a different correction problem.
\end{proof}
\KSwritten{WR15}

In the previous example, $B=(1,1)^{\mathsf T}$, $G=2$, and the proposal $c=1/4$ gives $q=-1/4$. Thus $q^{\mathsf T}G^{-1}q=1/32$, as required. Replacing $B$ by $sB$ changes $q$ to $sq$ and $G$ to $s^2G$. The Euclidean norm of the coordinate covector can be made arbitrarily small or large by changing the nonzero scalar $s$. The actual correction, normal residual, and certified error do not change.

\begin{figure}[htbp]\centering\resizebox{.96\linewidth}{!}{
% BEGIN retained input sources/monographs/komorebi-space/source/figures/39_covector_metric.tex
\includegraphics[width=\linewidth]{sources/monographs/komorebi-space/source/figures/39_covector_metric.pdf}

}\caption{Rescaling the basis of the same protected line rescales the raw coordinate covector. Its Euclidean norm is therefore not a representation-invariant stopping criterion. The Gram-corrected norm stays equal to $1/\sqrt{32}$ for every displayed scale. All systems represent the same proposed correction.}\label{KS:fig:covector-metric}\end{figure}

The finite verifier checks that a supplied basis spans the \emph{whole} protected kernel. A basis of a smaller subspace might give a correct certificate for that smaller optimization problem while missing an allowed direction that decreases the objective. The rank, kernel-membership, and independence checks therefore belong to the certificate contract.

\subsection{Separate data error, numerical error, and regularization bias}
Let $r$ be the intended datum and $\widetilde r$ the measured datum. Suppose
\[
 \KSnorm{\widetilde r-r}\le\delta,
 \qquad \KSnorm{T^*(Ta-\widetilde r)+\alpha a}\le\varepsilon.
\]
The maps $T,p$ and their Hilbert metrics remain fixed. By Theorem~\ref{KS:WR13} and the fixed-parameter part of Theorem~\ref{KS:WR09},
\begin{equation}\label{KS:eq:combined-solver-data}
 \KSnorm{a-a_\alpha(r)}
 \le \frac{\varepsilon}{\alpha}+\frac{\delta}{2\sqrt\alpha}.
\end{equation}
The first term is the solver's contribution. The second is the worst-case propagation of data error. They should not be collapsed into one unqualified residual tolerance.

\begin{KStheorem}[A consistent inexact stopping rule]\label{KS:WR16}
Assume the projected intended datum belongs to the actual range of $T$, with minimum-norm preimage $a^\dagger$. Along a sequence of positive parameters and measured problems, suppose
\[
 \alpha_n\to0,\qquad
 \frac{\delta_n}{\sqrt{\alpha_n}}\to0,\qquad
 \frac{\varepsilon_n}{\alpha_n}\to0.
\]
If the proposed corrections $a_n\in H$ satisfy the displayed data and normal-residual bounds, then $a_n\to a^\dagger$ in norm. At every index,
\[
 \KSnorm{a_n-a^\dagger}
 \le\frac{\varepsilon_n}{\alpha_n}
      +\frac{\delta_n}{2\sqrt{\alpha_n}}
      +\KSnorm{a_{\alpha_n}(r)-a^\dagger}.
\]
\end{KStheorem}
\begin{proof}
At index $n$, distinguish two exact minimizers: let $\widetilde a_n=a_{\alpha_n}(\widetilde r_n)$ solve the measured problem and let $a_n^0=a_{\alpha_n}(r)$ solve the intended-data problem. All three corrections, including the proposal $a_n$, belong to the same space $H$.

The supplied normal-residual bound is for the measured datum. Therefore Theorem~\ref{KS:WR13} applies first to that problem and gives
\[
 \KSnorm{a_n-\widetilde a_n}\le\frac{\varepsilon_n}{\alpha_n}.
\]
The pointwise data estimate in Theorem~\ref{KS:WR09}, at the same parameter, then gives
\[
 \KSnorm{\widetilde a_n-a_n^0}\le\frac{\delta_n}{2\sqrt{\alpha_n}}.
\]
Insert these two intermediate vectors between $a_n$ and $a^\dagger$ in the triangle inequality to obtain the displayed three-term estimate.

The first term tends to zero by the solver-tolerance hypothesis and the second by the data/parameter hypothesis. Since the projected intended datum is attained and $\alpha_n\to0$, the last term tends to zero by Theorem~\ref{KS:WR08}. For a direct epsilon argument, make each term smaller than one third of a prescribed tolerance beyond its own threshold and take the maximum of the three thresholds. This gives norm convergence of the proposals.

Protection is not merely asymptotic: $a_n\in\ker p$ implies $p(f-a_n)=pf$ at every index. The proof compares no changing operators or changing kernels. Both the pointwise certificate and the convergence conclusion have formal counterparts under these hypotheses.
\end{proof}
\KSwritten{WR16}

The accompanying Lean formalization includes the pointwise estimate \eqref{KS:eq:combined-solver-data}, the regularization limit, and the resulting convergence theorem. Their statements retain the same fixed operator, Hilbert geometry, and attainment requirements.

\subsection{Retaining the explicit diagonal rate with an inexact solver}
For the attainable datum of Corollary~\ref{KS:WR12} and $0<\delta\le1$, take $m=\lceil\delta^{-2/3}\rceil$ and $\alpha=m^{-2}$ as before. If an inexact solver satisfies
\[
 \varepsilon\le\kappa\alpha\delta^{1/3}
 \quad(\kappa\ge0),
\]
then the same argument gives
\[
 \KSnorm{a-a^\dagger}\le(\sqrt2+1+\kappa)\delta^{1/3}.
\]
Thus the inner solve can retain the proven outer accuracy if its normal residual decreases on the scale $\alpha\delta^{1/3}$. Since $\alpha$ is on the scale $\delta^{4/3}$, the corresponding normal-residual scale is $\delta^{5/3}$, with the stated ceiling factors retained in the exact rule. This is a sufficient rule for this problem, not a universal optimal stopping law.

The declaration \path{DiagonalRegularization.diagonal_inexact_rounded_rate} formalizes this specialization for the actual infinite diagonal operator. Its hypotheses bound both the data error and the normal residual evaluated against the measured datum. It retains the literal ceiling in $m$ and derives the exact-data bias estimate within the model.

\begin{figure}[htbp]\centering\resizebox{.96\linewidth}{!}{
% BEGIN retained input sources/monographs/komorebi-space/source/figures/40_stopping_schedule.tex
\includegraphics[width=\linewidth]{sources/monographs/komorebi-space/source/figures/40_stopping_schedule.pdf}

}\caption{For the attainable diagonal example, the three displayed terms are the proved bias bound, propagated data-noise bound, and a solver-error budget chosen to preserve the same $\delta^{1/3}$ rate. The ceiling in $m=\lceil\delta^{-2/3}\rceil$ is evaluated explicitly. The plot shows bounds, not measurements from a trained system.}\label{KS:fig:stopping-schedule}\end{figure}

For the escaping datum, the same rule does not create convergence to a nonexistent protected Hilbert answer. It only makes the numerical proposals track the regularized path accurately. Tracking an escaping path and converging to a realization are different conclusions.

\subsection{The executable verification contract}
The program takes $d,p,f,\alpha$ and a proposed ambient correction, with optional measured residuals, a full kernel basis, and an independent dual witness. It rejects a correction that changes the protected data and checks the basis's dimensions, rank, and kernel membership. It does not silently project a rejected proposal into another problem.

The output records $f-a$, protection before and after, measured fit, actual model residual, the Gram-corrected normal certificate, objective values, and squared error bounds. Verification recomputes the record without calling an optimizer. Independent determinant-based tests compare the bounds with actual errors in small ambient constrained systems. A certificate establishes consistency with its included problem, not authentication that this is the externally intended task.

\begin{insight}
A certificate needs a declared target. Fit measures agreement with data. The normal residual bounds optimizer error at a specified parameter. The duality gap bounds objective error. A realization theorem is still needed to conclude convergence to an intended exact object.
\end{insight}



% BEGIN retained input sources/monographs/komorebi-space/source/chapters/06_cases.tex

\chapter{Worked realization problems}\label{KS:book:chapter-24}
The constructions become easier to distinguish when a complete example is carried through from its view equations to its obstruction, its ambiguity, and its numerical repair. The following calculations instantiate the preceding principles. They do not introduce a second, silently enlarged formal theory.

\section{A closed signal loop}\label{KS:sec:signal-loop}
\subsection{Prescribed differences as affine view transitions}
Consider four scalar views arranged around a cycle. Their values are \(x_0,x_1,x_2,x_3\), and the prescribed differences are \(r_0,r_1,r_2,r_3\). Indices are taken modulo four. The compatibility equations are
\begin{equation}\label{KS:eq:cycle-difference}
 x_{i+1}-x_i=r_i.
\end{equation}
Work over \(\KSQ\) or \(\KSR\). These equations can be presented as a strict set-valued view system: an oriented edge has pullback \(z\mapsto z-r_i\), and a path has the composite affine pullback. Use the category of paths, retaining the full loop as a path rather than imposing that it is an identity arrow in advance.

Every edge map is bijective. Around the cycle their composite is translation by \(-\sum_i r_i\). A global realization must be fixed by that translation. Thus the example makes local invertibility and global compatibility visibly different without needing a noncommutative group.

\subsection{The complete solvability calculation}
Summing \eqref{KS:eq:cycle-difference} gives the necessary condition
\begin{equation}\label{KS:eq:cycle-obstruction}
 r_0+r_1+r_2+r_3=0.
\end{equation}
It is sufficient as well. Given an anchor \(a\), set
\[
 x_0=a,\quad x_1=a+r_0,\quad
 x_2=a+r_0+r_1,\quad x_3=a+r_0+r_1+r_2.
\]
The first three equations hold by construction, and the last is exactly \eqref{KS:eq:cycle-obstruction}. Every solution has this form because the first three equations recursively determine it from \(x_0\). Therefore a compatible residual has a one-dimensional family of primitives, with constant shifts as its ambiguity. Prescribing the anchor removes that ambiguity.

In matrix form, the difference operator is
\[
 D=\begin{pmatrix}
 -1&1&0&0\\
 0&-1&1&0\\
 0&0&-1&1\\
 1&0&0&-1
 \end{pmatrix},\qquad
 \ker D=\operatorname{span}\{(1,1,1,1)\}.
\]
The explicit construction proves
\[
 \KSran D=\ker\ell,\qquad \ell(r)=\sum_i r_i.
\]
Thus \(\ell\) is a complete feasibility test for this prescribed-difference problem. It annihilates every actual difference family, but takes a nonzero value on an incompatible one. This is the finite obstruction-certificate pattern of Proposition~\ref{KS:KS15} and Proposition~\ref{KS:KS57} in a fully visible model.

The same proof works on a cycle of any positive length: partial sums give a primitive, and the final wraparound equation is the zero-sum condition. No Fourier transform or numerical approximation is needed for this solvability argument.

\subsection{Repairing an incompatible residual is a new problem}
Take
\[
 r=(1,0,0,0).
\]
Each of the first three edge equations can be satisfied, but the full sum is one. There is no coherent closed-loop signal with exactly these four differences. Computing a long list of partial values does not remove the missing wraparound equation.

To repair the residual, now regard the vector of edge data as a candidate in its own right. Its defect is \(\ell(r)\). The exact subspace is the zero-sum hyperplane. With the ordinary Euclidean norm and no protected edge coordinates, the nearest compatible residual is
\[
 \widehat r=r-\frac{\ell(r)}4\mathbf1
 =\left(\frac34,-\frac14,-\frac14,-\frac14\right).
\]
The correction is perpendicular to the hyperplane, so the projection theorem proves it is the unique nearest vector there. Its squared size is \(1/4\). This is a repair of the edge data, not a solution of the original incompatible equations with those data unchanged.

Anchor the repaired signal at zero. Its primitive is
\begin{equation}\label{KS:eq:cycle-anchor-primitive}
 \widehat x=\left(0,\frac34,\frac12,\frac14\right).
\end{equation}
The next partial sum is zero, so it closes. If instead the desired representative has mean zero, subtract the mean \(3/8\):
\[
 \widehat x_{\rm mean}
 =\left(-\frac38,\frac38,\frac18,-\frac18\right).
\]
Both primitives realize exactly \(\widehat r\). They differ by an element of \(\ker D\). The anchor condition and the mean condition are two different selections from the same affine solution fiber.

\begin{figure}[htbp]\centering\resizebox{.96\linewidth}{!}{
% BEGIN retained input sources/monographs/komorebi-space/source/figures/23_cycle_primitive.tex
\includegraphics[width=15cm]{sources/monographs/komorebi-space/source/figures/generated/23_cycle_primitive.pdf}

}\caption{The repaired edge data integrate to a closed signal. The value at position four repeats the anchor at position zero. The mean-zero and anchor-zero primitives differ by a constant and therefore have the same four differences. The horizontal axis labels vertices, not successive optimization iterations.}\label{KS:fig:cycle-primitive}\end{figure}

\subsection{Protecting three edges forces a different repair}
Suppose the first three values \((r_0,r_1,r_2)=(1,0,0)\) are required to remain fixed. The admissible corrections then affect only the final edge. Compatibility forces
\[
 \widetilde r=(1,0,0,-1),\qquad
 \widetilde x=(0,1,1,1)
\]
for anchor zero. This is the unique repair preserving those three edges. Its squared correction size is one rather than \(1/4\). The larger error is not a failure of nearest repair: the feasible set has changed.

Here the general protected construction is completely explicit. With candidate space \(E=\KSR^4\), defect \(d=\ell\), and protected observable \(p(r)=(r_0,r_1,r_2)\), the allowed direction is the last coordinate axis. The restricted operator \(d|_{\ker p}\) has gain one. Without protection, the least gain for distance to \(\ker\ell\) is \(1/2\), since \(\KSnorm{\ell}=2\). The different constants measure different correction spaces.

\subsection{Where the packet and transfer viewpoints enter}
The raw edge family contains information about four transitions. Its zero-sum constraint links all four. Summing a block of differences transfers information to a coarser transition, while retaining the differences within the block preserves the detail that the sum omits. For example,
\[
 r_0+r_1=x_2-x_0
\]
contains the two-step transition but does not determine \(r_0\) and \(r_1\) separately. Replacing \((r_0,r_1)\) by \((r_0+t,r_1-t)\) leaves that coarse transition unchanged.

A packet of coarse value and independent detail retains both directions. A residual-only summary can detect failure of the wraparound equation but cannot tell which protected edge may be changed. An anchor or a mean fixes the remaining signal ambiguity after feasible edge data have been chosen. These are three different roles: transition compression, feasibility repair, and representative selection.

This example is a complete finite calculation built from the strict realization, additive obstruction, and protected Hilbert principles. The companion program checks the displayed rational values, the zero-sum criterion, the primitive reconstruction, and the orthogonality equations. The calculation does not claim formal verification of a Fourier or signal-processing library.

\chap{A loop with a nontrivial fixed space}\label{KS:sec:twisted-loop}
The previous cycle has scalar translations and a scalar total obstruction. A linear connection can behave differently: its loop may be nontrivial while retaining a whole subspace of feasible anchors. The following three-vertex model distinguishes full anchor preservation, feasible-anchor preservation, and global nearest-point repair in one calculation.

\subsection{The connection and its exact families}
Let the fiber at each of three vertices be \(\KSR^2\). Use identity transport along the route \(0\to1\to2\), and let the closing edge carry the reflection
\[
 J=\begin{pmatrix}1&0\\0&-1\end{pmatrix}.
\]
For one chosen orientation of each edge, the residuals are
\[
 r_{01}=x_1-x_0,\qquad r_{12}=x_2-x_1,\qquad
 r_{20}=Jx_0-x_2.
\]
The reverse transports are inverses, so this defines an inverse-consistent isometric connection. Exactness first forces \(x_1=x_0\) and \(x_2=x_0\), and the remaining edge then forces \(Jx_0=x_0\). Hence every exact family is
\begin{equation}\label{KS:eq:reflection-exact}
 ((a,0),(a,0),(a,0)),\qquad a\in\KSR.
\end{equation}
The loop is not the identity on the whole fiber. Its fixed space is the first coordinate axis. Every anchor on that axis extends uniquely, and every anchor with nonzero second coordinate is obstructed if its full value is to be preserved.

The three-edge residual also gives a direct telescope:
\[
 r_{01}+r_{12}+r_{20}=(J-I)x_0.
\]
Its second component is \(-2(x_0)_2\). In this chosen route normalization, the loop residual determines the transverse anchor component. It does not constrain the first component, which is the exact freedom seen in \eqref{KS:eq:reflection-exact}.

\begin{figure}[htbp]\centering\resizebox{.96\linewidth}{!}{
% BEGIN retained input sources/monographs/komorebi-space/source/figures/24_reflection_cycle.tex
\includegraphics[width=15cm]{sources/monographs/komorebi-space/source/figures/generated/24_reflection_cycle.pdf}

}\caption{The tree edges identify the three fiber values. The reflection on the closing edge then removes only the transverse anchor component. The loop is nontrivial, yet all nonzero vectors on its fixed line define nonzero coherent families. The displayed connection uses a two-dimensional fiber at every vertex.}\label{KS:fig:reflection-cycle}\end{figure}

\subsection{Three different questions for the same raw data}
Take
\[
 x_0=(2,1),\qquad x_1=(4,-1),\qquad x_2=(3,2).
\]
The measured oriented residuals are
\[
 r_{01}=(2,-2),\qquad r_{12}=(-1,3),\qquad r_{20}=(-1,-3).
\]
Preserving the entire original anchor is impossible, because \((2,1)\) is not fixed by \(J\). The earlier obstruction criterion therefore gives a negative answer to that particular preservation problem.

Preserving only the feasible projection of the anchor gives a different answer. Orthogonally project \((2,1)\) to \((2,0)\), then reconstruct along the tree. The result is
\[
 P_{\rm anchor}x=((2,0),(2,0),(2,0)).
\]
It is the unique exact family preserving the first component of the original anchor. It does not preserve its second component, and makes no claim to do so.

For global nearest-point repair, declare the Euclidean product norm on \((\KSR^2)^3\). The closest member of \eqref{KS:eq:reflection-exact} uses the mean of the three first coordinates, namely three:
\[
 P_{\rm nearest}x=((3,0),(3,0),(3,0)).
\]
To verify this without appealing only to a general theorem, write the squared error at parameter \(a\):
\[
 (2-a)^2+(4-a)^2+(3-a)^2+1^2+(-1)^2+2^2
 =3(a-3)^2+8.
\]
It is minimized uniquely at \(a=3\). The anchor-based answer has squared error eleven. The globally nearest answer has squared error eight. Their difference has squared norm three, giving the expected Pythagorean identity \(11=8+3\).

No two of these conclusions conflict. Full anchor preservation is infeasible. Feasible-anchor preservation has a unique answer. Unprotected global nearest repair selects a different answer. A zero-residual output without its preservation contract would not distinguish them.

\subsection{An exact conditioning calculation}
For this paragraph use one orientation per edge and the Euclidean norm on the six residual coordinates. Group candidate coordinates by fiber direction. The two scalar blocks of the defect operator are
\[
 A=\begin{pmatrix}-1&1&0\\0&-1&1\\1&0&-1\end{pmatrix},\qquad
 B=\begin{pmatrix}-1&1&0\\0&-1&1\\-1&0&-1\end{pmatrix}.
\]
The ordinary cycle block satisfies
\[
 A^{\mathsf T}A=3I-\mathbf1\mathbf1^{\mathsf T}.
\]
It is zero on the constant line and equals \(3I\) on its orthogonal complement. The reflected block satisfies
\[
 B^{\mathsf T}B=I+vv^{\mathsf T},\qquad v=(1,-1,1)^{\mathsf T}.
\]
Since \(\KSnorm v^2=3\), this matrix has eigenvalue four in the direction of \(v\) and eigenvalue one on its orthogonal complement. In particular it is positive definite.

These two identities show that the least nonnegative gain for the full defect-to-distance estimate in this Euclidean model is one. The lower bound is attained, for example, on a candidate with zero first components and second-component vector \((1,1,0)\), which lies in \(v^\perp\). The squared input distance and squared residual are both two.

For the displayed raw candidate, the squared residual is
\[
 (2^2+(-2)^2)+((-1)^2+3^2)+((-1)^2+(-3)^2)=28.
\]
Its squared nearest-point error is eight, consistent with the estimate. The certificate is universal for this fixed matrix because the Gram identities hold for all vectors, not because that particular example passed.

If both orientations are retained with a Euclidean sum-of-squares residual norm, each isometric edge contributes its squared residual norm twice. The residual norm is multiplied by \(\sqrt2\), and the least gain becomes \(1/\sqrt2\). With the maximum residual norm used in the general graph theorem, that duplication does not change the maximum. This is a concrete illustration of why the norm convention must travel with a certificate.

\subsection{What this example connects}
The strict section equations identify the exact families. The holonomy fixed space determines which anchors are feasible. The protected observable selects which exact answer is admissible. The Hilbert geometry selects a nearest representative when that is the question being asked. The Gram matrices determine the quantitative gain in one declared coordinate model.

The example does not require the loop to be trivial on the whole fiber. Nor does it require the anchor repair to be globally orthogonal. Both of those stronger conclusions would be false. The realization calculus is useful here because it keeps the correct weaker structures available rather than forcing the model into a stronger but incorrect description.



% BEGIN retained input sources/monographs/komorebi-space/source/chapters/06_boundaries.tex
\chap{Mathematical meaning and the remaining boundaries}\label{KS:sec:meaning}
\subsection{A relative theory can be complete without being absolute}
The strict space \(\Gamma(F)\) is canonical relative to its diagram: its universal property determines it up to the appropriate unique isomorphism. The additive obstruction is canonical relative to its admissible corrections. A marked realization problem is invariant under the specified comparison equivalences. A repair can be canonical relative to a protected observable, or orthogonal relative to a fixed inner product.

These are meaningful and different kinds of canonicity. They need not culminate in one object independent of every choice of views, coefficients, topology, semantics, and geometry. Some changes of view intentionally forget information. Others strengthen the question. Treating every such change as an equivalence would remove the very distinctions the calculus is designed to preserve.

The appropriate central statement is therefore not that every possible mathematical presentation lives in one fixed universal space. It is that a declared system of representations can be organized into coherent realizations, with explicit extension, comparison, and repair problems. A proposed bridge between two such systems has to show what structure it preserves and what information it loses.

\subsection{Internal realization and intended objects}
Suppose a separate category \(\mathcal S\) contains the intended semantic objects of an application, and a realization construction gives a comparison functor
\[
 \mathcal R:\mathcal S\longrightarrow\mathcal K,
\]
where \(\mathcal K\) is the relevant category or groupoid of coherent data. Defining \(\mathcal K\) does not establish that \(\mathcal R\) is faithful, full, or essentially surjective. These are new mathematical questions about the actual semantic comparison.

Faithfulness asks whether distinct semantic arrows with the same source and target can become indistinguishable. Fullness asks whether every arrow between the images of two semantic objects comes from an arrow between those objects. Essential surjectivity asks whether every object of \(\mathcal K\) is isomorphic to the image of an actual semantic object. Depending on the application, one may seek only a restricted image or a classified obstruction to belonging to it.

This distinction prevents a formal globalization of encoded data from being mistaken for an unrestricted arithmetic or physical existence theorem. The abstract construction does not manufacture missing semantic objects.

\subsection{Every finite test can pass while the infinite problem fails}
\begin{KSexample}[Finite prefixes without a global point]\label{KS:KS54}
Let \(X_n=\{m\in\KSN:m\ge n\}\), with bonding maps given by inclusion. Every level is nonempty, and every finite prefix has a coherent point. The inverse limit is nevertheless empty.
\end{KSexample}
\begin{proof}
For a fixed terminal level \(N\), the integer \(N\) belongs to every fiber \(X_n\) with \(n\le N\). Assign that same integer at each such level. Every restriction is an inclusion, so it preserves the assigned integer and all finite coherence equations hold. This proves finite solvability with the quantifiers in the order ``for every \(N\), there exists a prefix.''

Suppose instead that \(x=(x_n)\) is a global point. Equality along each inclusion says that the underlying natural numbers of \(x_{n+1}\) and \(x_n\) are equal. Induction therefore makes every underlying coordinate equal to the number \(m\) underlying \(x_0\). Membership of \(x_{m+1}\) in \(X_{m+1}\) now gives \(m\ge m+1\), a contradiction. Thus no global point exists.

There is no conflict with a gluing theorem for prefixes. The prefix chosen above for length \(N+1\) restricts to the value \(N+1\), not to our chosen value \(N\) on the shorter prefix. Nor could another choice of one solution for every length give a matching family: its value at level zero would force the same forbidden global integer. Separate existence at every length is strictly weaker than existence of a compatible family of prefix solutions.
\end{proof}
\KSsource{KS54}
This counterexample is formalized in \KScode{FinitePrefixBoundary.lean} using the original resolution and partial-section types. It explains why inverse-limit existence needs more than separate finite solvability. Finiteness, compactness, surjectivity, or another coherence principle must be proved in the relevant application. The example does not contradict the strict gluing theorem: a compatible matching family of all prefixes was never supplied.

\subsection{A positive infinite-existence criterion}
The failed example isolates a missing compactness or lifting hypothesis. Finite fibers supply a concrete positive criterion without requiring every bonding map to be surjective.

\begin{KSproposition}[Finite fibers and persistent images]\label{KS:KS58}
Let \((X_n,\rho_n)\) be a sequential inverse system of nonempty finite sets. Its inverse limit is nonempty. More precisely, if \(\rho_{n,m}:X_m\to X_n\) is composite restriction, then
\[
 \KSim(\operatorname{ev}_n)=Y_n,
 \qquad Y_n=\bigcap_{m\ge n}\KSim\rho_{n,m}.
\]
Thus a prescribed value at level \(n\) extends precisely when it survives every later image.
\end{KSproposition}
\begin{proof}
Write \(I_{n,m}=\rho_{n,m}(X_m)\) for \(m\ge n\). Each \(I_{n,m}\) is nonempty. The composition law gives
\[
 I_{n,m+1}=\rho_{n,m}\bigl(\rho_m(X_{m+1})\bigr)\subseteq I_{n,m}.
\]
These are subsets of the finite set \(X_n\). Every strict decrease removes an element permanently, so only finitely many strict decreases are possible. There is an index \(s_n\ge n\) such that \(I_{n,m}=I_{n,s_n}\) for all \(m\ge s_n\). Consequently \(Y_n=I_{n,s_n}\) is nonempty.

We next prove that restriction induces a \emph{surjection} \(Y_{n+1}\to Y_n\), even though \(X_{n+1}\to X_n\) need not be surjective. For membership, let \(z\in Y_{n+1}\). Given \(m\ge n+1\), write \(z=\rho_{n+1,m}(w)\) for some \(w\in X_m\). Then \(\rho_n z=\rho_{n,m}(w)\). Together with the automatic condition at \(m=n\), this shows \(\rho_nz\in Y_n\).

For surjectivity, fix \(y\in Y_n\) and choose
\[
 m\ge\max\{s_n,s_{n+1},n+1\}.
\]
Since \(y\in I_{n,m}\), there is \(w\in X_m\) with \(\rho_{n,m}(w)=y\). Put \(z=\rho_{n+1,m}(w)\). Stabilization at level \(n+1\) gives \(z\in Y_{n+1}\), and composition gives \(\rho_nz=y\).

Choose \(x_0\in Y_0\) and recursively choose \(x_{n+1}\in Y_{n+1}\) above \(x_n\). The preceding surjectivity justifies each step. Adjacent coherence implies every composite coherence equation, so this gives a point. This is an existence construction using countably many choices, not an asserted effective procedure for locating the stabilization indices.

Finally, fix a level \(n\). Every global point satisfies \(x_n=\rho_{n,m}(x_m)\) for each \(m\ge n\), hence \(x_n\in Y_n\). Conversely, start with any \(y\in Y_n\), lift recursively above level \(n\) through the surjections just proved, and define the lower coordinates by \(x_i=\rho_{i,n}(y)\). The composite restriction laws verify coherence across and below level \(n\). The resulting point evaluates to \(y\). These two inclusions prove the asserted equality of evaluation images, not merely nonemptiness of the inverse limit.
\end{proof}
\KSsource{KS58}
This is the standard finite-set inverse-limit principle, here with its sequential persistent-image proof~\cite{KS:stacks-limits}. Its formalization in \KScode{FiniteRealizationExistence.lean} constructs an explicit equivalence with the pinned library's section type, then applies the finite cofiltered-system and eventual-range theorems. The printed sequential proof and this formal route establish the same stated criterion. Both use ordinary classical choice. Surjective bonding maps on arbitrary nonempty sets give a second sufficient condition by direct recursive lifting. Neither criterion says that arbitrary data in a non-surjective finite tower extend.

\subsection{A finite computation that exposes the preservation contract}
Consider a raw vector \(x=(2,0,5)\in\KSR^3\) with defect \(d(x)=x_1-x_2\). Unconstrained orthogonal repair gives \((1,1,5)\), with squared error two. If the first and third coordinates must be preserved, the unique exact repair is \((2,2,5)\), with squared error four.

Both outputs are exact. Only one satisfies the second preservation problem. Neither theorem contradicts the other because the admissible changes differ. An implementation that reports only zero residual would not distinguish the two tasks.

For a finite linear system over a field, the same question admits independently checkable positive and negative certificates. Let \(A\in\KSK^{m\times n}\) encode compatibility, \(P\in\KSK^{p\times n}\) the protected values, and \(x\in\KSK^n\) the raw candidate. Let the columns of \(N\) be a basis of \(\ker P\). The correction equation is
\[
 ANz=Ax,\qquad y=x-Nz.
\]
It is equivalent to the stacked system \(Ay=0\), \(Py=Px\).

\begin{KSproposition}[A finite certificate alternative]\label{KS:KS57}
Exactly one of the following occurs:
\begin{enumerate}
\item There is \(y\) with \(Ay=0\) and \(Py=Px\). The full answer set is \(y+H\), where \(H=\ker A\cap\ker P\).
\item There are row vectors \(\lambda\in\KSK^{1\times m}\) and \(\mu\in\KSK^{1\times p}\) with
\[
 \lambda A=\mu P,\qquad\lambda Ax=1.
\]
\end{enumerate}
The negative witness annihilates every admissible correction defect, since \(\lambda AN=\mu PN=0\). The ambiguity dimension in the positive case is
\[
 \dim H=n-\operatorname{rank}\begin{pmatrix}A\\P\end{pmatrix}.
\]
\end{KSproposition}
\begin{proof}
Set \(T=\begin{pmatrix}A\\P\end{pmatrix}\) and \(b=\begin{pmatrix}0\\Px\end{pmatrix}\). A vector \(y\) is an exact protected answer if and only if \(Ty=b\). Perform Gaussian row operations on the augmented matrix, keeping their product \(U\). Each elementary row operation is invertible over the given field. Thus \(Ty=b\) and \(UTy=Ub\) have exactly the same solutions.

If no row of the reduced augmented system has zero left side and nonzero right side, assign arbitrary values to its free variables and solve for its pivot variables. This supplies a solution \(y_0\). For any vector \(y\),
\[
 Ty=b\quad\Longleftrightarrow\quad T(y-y_0)=0.
\]
Translation by \(y_0\) is therefore a bijection from \(\ker T\) to the answer set, with inverse \(y\mapsto y-y_0\). Rank--nullity applied to the actual linear map \(T:\KSK^n\to\KSK^{m+p}\) gives the ambiguity dimension. Dependent prescription rows do not alter this argument.

If the reduced system is inconsistent, choose a row with zero left side and right side \(\alpha\ne0\). Let \(w\) be the corresponding row of \(U\), and split it according to the two blocks as \(w=(u,v)\). The retained row operations give the literal identities
\[
 uA+vP=0,\qquad vPx=\alpha.
\]
Division is permitted because \(\alpha\) is a nonzero field element. Define \(\lambda=-\alpha^{-1}u\) and \(\mu=\alpha^{-1}v\). Multiplying the displayed equations by \(\alpha^{-1}\) gives
\[
 \lambda A=\mu P,\qquad \lambda Ax=\mu Px=1.
\]
For every admissible correction \(a\in\ker P\), these same row equations imply \(\lambda Aa=\mu Pa=0\). The witness therefore annihilates the entire correction image while taking value one on the required defect.

The alternatives cannot occur together. If \(y\) were an exact protected answer in the presence of such a witness, then
\[
 1=\lambda Ax=\mu Px=\mu Py=\lambda Ay=0.
\]
A field has \(1\ne0\), so this is impossible. Consistency or inconsistency of the reduced system proves exhaustiveness, and the contradiction proves exclusivity. Empty row or column blocks cause no extra division: normalization is used only in the inconsistent case, where the nonzero scalar has already been obtained. The proof does not transfer to integral coefficients by silently allowing the same division.
\end{proof}
\KSsource{KS57}
The formal proof gives these literal row-vector equations over an arbitrary field. It uses algebraic dual separation of the stacked image rather than verifying the row-reduction program. The affine answer equivalence and rank-nullity formula are proved separately. Thus a normalized obstruction is an actual checkable witness, while universal correctness of the supplied Python elimination implementation remains a different question.

\subsection{One equation, three preservation problems}
For the following three cases work over \(\KSQ\), with \(A=(1,-1,0)\) and \(x=(2,0,5)\). With \(P(x)=(x_1,x_3)\), the unique answer is \((2,2,5)\). With \(P(x)=x_3\), the entire answer set is \(\{(t,t,5):t\in\KSQ\}\). With \(P(x)=(x_1,x_2)\), no answer exists. In this last case
\[
 \lambda=(1/2),\qquad\mu=(1/2,-1/2)
\]
is a normalized negative certificate over \(\KSQ\) or \(\KSR\): \(\lambda A=\mu P\) and \(\lambda Ax=1\).

The package implements this finite interface over \(\KSQ\) in \KScode{examples/protected_repair.py}. It returns a particular exact answer and an ambiguity basis, or a normalized dual witness. Its verifier checks the certificate equations directly. A Gaussian particular answer is \emph{not} advertised as the nearest answer of Theorem~\ref{KS:KS56}. That is a separate geometric selection.

Exact field elimination does not solve the integer parity problem by changing its coefficient domain. Over \(\KSQ\), \(2m=1\) has the solution \(1/2\). Over \(\KSZ\), it has none. The rational interface rejects floating-point inputs rather than silently treating an approximate decimal as exact arithmetic. Its finite tests establish the tested implementation behavior, not an arbitrary-input Lean refinement theorem.

\subsection{Why this is relevant to machine reasoning}
The framework gives a machine explicit mathematical objects to manipulate: views, restrictions, partial realizations, typed defects, admissible corrections, and comparison witnesses. A successful step can carry both an output and a certificate of what remained fixed. A failed step can return an obstruction in a specified quotient rather than an unexplained failure flag.

The split coordinate theorem is especially direct. In its stated class, the same candidate is represented by protected information and residual data, and repair zeros only the residual coordinate. The descent theorem explains which local certificates can be assembled. The analytic results explain when a small measured defect has a quantitatively reliable meaning.

These observations motivate a machine oriented mathematical language. They do not establish an autonomous reasoning system, an empirical intelligence gain, or a general efficient algorithm. Classical choice, infinite products, quotient membership, and bounded inverse constructions can be mathematically valid without yielding executable finite procedures. Implementation claims need their own specifications and tests.

\subsection{What this monograph closes, and what it leaves open}
The book supplies a coherent account of strict realization spaces, additive and nonabelian extension problems, split tower repair, retained multiscale detail, a genuine composition-site stack, presentation transport, graph holonomy, closed-range stability, and moving exact spaces. The convergence and obstruction theory separates small measurements from small movements, coherent coordinate families from their semantic image, and zero residual from zero transformation cost. The bounded-budget theory supplies a positive existence criterion and keeps range membership separate from regularized solvability. Its fixed-budget dual witness can certify a limit on every admissible correction, while its diagonal example computes the actual leading cost of approximation. Each layer is attached to the data that make it meaningful.

The following broader targets are not conclusions of this edition: arbitrary higher categorical coefficients and infinity-stack descent, unrestricted changes of view doctrine, norm reconstruction for every bounded infinite residual hierarchy, nonlinear or unknown-connection repair, perturbation across changing Hilbert geometries without new comparison data, and semantic essential surjectivity for unrestricted arithmetic or physical applications.

The accompanying Lean project contains \KSEditionProductionModules\ production modules and \KSEditionRoots\ entry points, with compiler records, declaration axiom reports, and rejection controls. The bounded-budget, regularization-limit, noise, inexact-certificate, and coordinate-invariance results have formal counterparts under their stated hypotheses. The explicit infinite diagonal estimates, leading budget constant over all real budgets, and rounded noisy rate also have formal counterparts. The budget results do not require finite-dimensionality, separability, injectivity, or closed range. Their robustness comparison uses one fixed correction space. Further extensions require their own hypotheses and evidence.

\begin{insight}
\textbf{The organizing principle.} Coherent realization, obstruction to extension, freedom of completion, and stability of repair belong to one mathematical conversation. Their unification is useful precisely when it does not identify them with one another.
\end{insight}



\part{Grounded abstraction and the limits of learning}
\setcounter{section}{0}\renewcommand{\thesection}{GA.\arabic{section}}\renewcommand{\theHsection}{GA.\arabic{section}}

\chapter{Evidence before representation}\label{GA:book:chapter-25}

\section{What does experience justify?}\label{GA:sec:introduction}
A device displays a bit. A reasoner learns to predict that display perfectly, including what happens when a permitted switch is pressed. It still may not know which physical bit the display describes. The code linking the display to the physical state could be reversed. A second, equally accurate reasoner might use the opposite internal convention. Prediction is real, yet the interpretation remains underdetermined.

That observation does not force a choice between calling every representation meaningless and treating every successful prediction as a complete understanding. It suggests a more precise question. What claims are common to all systems compatible with the evidence? Which of those claims can be realized by one representation that the observer can actually initialize and update? What additional information becomes necessary when the questions change?

\subsection{Three separations}
Identification of an underlying system, identification of a task-relevant property, and acquisition of a usable representation are distinct. A nuisance parameter can remain unknown while every required prediction is fixed. Conversely, knowing that some accurate representation exists in each possible world does not supply one representation that can be selected from the observer's record.

There is a second separation between a probability law and a finite observation drawn from that law. Two candidate models can have different laws while sharing every finite record in their supports. In that setting the models are identifiable at the level of their laws, but an individual finite record cannot yield error-free discrimination.

The third separation concerns time. A summary may preserve everything required for a present decision while losing the relationship needed to interpret a later measurement. More data arriving tomorrow need not repair what the memory policy erased today. This paper treats that relationship as part of the retained mathematical state.

\subsection{Position among existing theories}
Classical statistical identification separates parameters that induce different observational laws~\cite{GA:rothenberg}. Predictive state representations use action-conditional predictions as a state description~\cite{GA:psr}. Query learning makes the teacher and the information it returns explicit~\cite{GA:angluin}. These are neighboring answers to different parts of the present problem, not interchangeable definitions of grounding.

Data informativity is especially close to our organizing principle. It asks whether one property or one controller is justified for every system consistent with the data, without necessarily identifying a unique system~\cite{GA:informativity}. Our later linear construction uses this compatible-extension viewpoint for the exact domain of repeated prediction. It is not offered as a replacement for data-driven control theory.

Causal abstraction concerns faithful relationships between mechanisms and their higher-level descriptions~\cite{GA:causal}. Here we ask what evidence makes such a description selectable and which continuations remain determined when parts of the underlying mechanism are unknown. We do not derive a learned encoder from the existence of a commuting equation. State-identification theory also distinguishes individually separable states from a realizable adaptive test~\cite{GA:stateid}. That distinction will matter when measurements alter the system.

\subsection{The intended contribution}
The main results form a connected argument. Exact grounding is an information factorization through a behavioral quotient. New task distinctions are classified relative to an actual diagnostic interface. Future correction requires a richer memory signature than present target uncertainty. Partial operator calibration can nevertheless justify every future composition on a computable invariant domain. Approximate claims require both a local discrepancy and a bound on the dynamics that may amplify it.

The positive message is not that the observer must recover everything. It is that the observer can often recover exactly the structure its intended operations need, provided the evidence contract is stated honestly. The negative message is that an attractive coordinate system cannot manufacture missing distinctions, missing experiments, or missing bounds.

\section{A calibrated bit with an irrelevant hidden register}\label{GA:sec:sensor}
Consider eight possible worlds
\[
\omega=(s,\theta,k)\in\{0,1\}^3.
\]
The physical bit is $s$. The sensor convention is $\theta$, and $k$ is a nuisance bit that none of the admitted operations or questions uses. The display is
\[
y=s\oplus\theta.
\]
The ordinary action $F$ flips $s$ and leaves the other bits fixed. Its effect on the display is exactly $y\mapsto1-y$, whatever the unknown convention. Reading the current display therefore initializes a two-state representation that predicts every finite sequence of flips.

Neither $s$ nor $\theta$ is identified separately. The two worlds $(0,0,k)$ and $(1,1,k)$ both display zero, and every flip sequence produces the same display in both. That ambiguity is harmless for the original display-prediction contract.

\subsection{A new question}
Now ask for the physical bit $s$ as well as the display. The old representation is insufficient. The two previously equivalent worlds require different answers. A permitted nondestructive reference measurement returning $\theta$ makes the physical bit recoverable by $s=y\oplus\theta$. The enriched representation can retain $(y,\theta)$, or equivalently $(s,\theta)$. It has four states, not eight. The nuisance bit remains unidentified because no required continuation can detect it.

A new field named ``sensor convention'' does not provide the reference measurement. Nor does repeatedly flipping the physical bit resolve the ambiguity: the complete display history remains the same in the two worlds. Acquisition depends on an additional information channel, not merely a larger memory allocation.

\subsection{A state-changing measurement}
A different operation $Z$ sets $s$ to zero and then reports the new display. The report is $\theta$. If the old display $y_{\mathrm{old}}$ was retained, then
\[
s_{\mathrm{old}}=y_{\mathrm{old}}\oplus y_{\mathrm{after}\ Z}.
\]
The current physical state is zero, not $s_{\mathrm{old}}$. A correct account distinguishes the recovered historical value from the new state it will act on.

If the old display was discarded, worlds with opposite old physical bits but the same $\theta$ reach the same current state and give the same future reports. The old physical bit is then irrecoverable through this unchanged interface. A known present state has replaced, rather than identified, the unknown past.

\begin{figure}[tbp]\centering
\includegraphics[width=.95\linewidth]{sources/monographs/grounded-abstraction/figures/sensor_grounding.pdf}
\caption{One finite sensor model supports different exact abstractions. The current display and flip action need two behavioral states. Adding the physical-bit question requires four, but not the irrelevant nuisance register. Calibration is an information source. Resetting is a physical change.}\label{GA:fig:sensor}
\end{figure}

\section{Laws, records, and reusable predictive state}\label{GA:sec:grounding}
\begin{GAdefinition}[Finite evidence contract]\label{GA:def:evidence}
A finite evidence contract has a nonempty candidate set $\Omega$, a family $\Pi$ of permitted protocols, a finite record space $H_\pi$ for each protocol, and kernels $K^\pi(h\mid\omega)$. A target $t:\Omega\to Y$ states what is to be identified. Any hidden model index that is fixed during a protocol is included in $\omega$. The protocol's action choices are functions of its permitted records and independent randomization, not of the unobserved world.
\end{GAdefinition}

Define $\omega\sim_\Pi\omega'$ when their record laws agree for every $\pi\in\Pi$. This relation concerns the complete records, including action choices when those are observable. Equality of individual observation marginals is insufficient.

\begin{GAproposition}[Law identification and finite-record identification]\label{GA:prop:law}
The target is a function of the complete family of protocol laws if and only if it is constant on the $\sim_\Pi$ classes. For a fixed protocol $\pi$, an error-free rule $d:H_\pi\to Y$ exists if and only if $t$ is constant on every nonempty support set
\[
B_\pi(h)=\{\omega:K^\pi(h\mid\omega)>0\}.
\]
Allowing a randomized output rule does not weaken this condition for probability-one correctness in every world.
\end{GAproposition}
\begin{proof}
For the first claim, factor $t$ through the map sending a world to its tuple of protocol laws. A factor exists precisely when worlds with the same tuple have the same target. For the second claim, a record with positive probability in two worlds must be given the same output. Distinct targets on its support therefore prevent certainty. Conversely, assign each attainable record its support's common target. Every record generated in a world receives that world's target. A randomized error-free rule must place probability one on the required target at every positive-probability record. It cannot do so for two distinct targets at the same record.
\end{proof}

This elementary factorization does not supply the laws to the learner or bound the sample cost of estimating them. The two clauses answer different questions. A family of distinct laws can be identifiable in the first sense while every available short record remains ambiguous in the second.

\subsection{A canonical meaning for an operational state}
Let $X$ be a finite nonempty set, let $\GAA$ be a finite operation alphabet, and let each $F_a:X\to X$ be total and deterministic. Let $q:X\to Y$ be a finite vector of required observations. Different candidate models can be included by taking a disjoint union of their state sets and preserving each model tag under every action.

For a word $w=a_1\cdots a_m$, $F_w$ applies $a_1$ first and $a_m$ last. The empty word is the identity. Define
\[
x\sim_\GAC x'\quad\Longleftrightarrow\quad
q(F_wx)=q(F_wx')\quad\text{for every finite }w.
\]
This is a comparison in a common observation language. It does not require that different models share arbitrary latent coordinates. The behavioral class is a canonical response profile even when its implementation is not identifiable.

\begin{GAdefinition}[Exact initialization from available evidence]\label{GA:def:initialization}
Given an available initial code $r:X\to H$, an exact reusable predictor consists of a finite state set $Z$, an initialization $d:H\to Z$, updates $G_a:Z\to Z$, and an output map $\bar q:Z\to Y$ such that
\[
\bar q\,G_w\,d\,r(x)=q(F_wx)
\]
for all $x$ and all finite words. After initialization the predictor receives the operation names, but no additional world observations or hidden state.
\end{GAdefinition}

\begin{GAtheorem}[Evidence sufficient for one exact predictive realization]\label{GA:thm:grounding}
An exact reusable predictor exists if and only if
\[
r(x)=r(x')\quad\Longrightarrow\quad x\sim_\GAC x'.\tag{1}
\]
When it exists, the behavioral quotient $X/{\sim_\GAC}$ supplies a predictor with the fewest initialized states. Every predictor initialized for all $x\in X$ requires at least $|X/{\sim_\GAC}|$ distinct values in $d(r(X))$.
\end{GAtheorem}
\begin{proof}
Equal initial codes yield equal initialized states. The same operation sequence must then produce equal predicted outputs. Exactness forces equality of all true continuation outputs, proving necessity.

For sufficiency, note that $x\sim_\GAC x'$ implies $F_ax\sim_\GAC F_ax'$ because any subsequent word can be prefixed by $a$. The quotient consequently supports $G_a([x])=[F_ax]$ and $\bar q([x])=q(x)$. Condition (1) makes $d(r(x))=[x]$ well defined on the attained initial codes. Arbitrary unused codes can be assigned any quotient state. Induction on word length gives the required equality. Finally, two behaviorally different states cannot have the same initialized predictor state, or a distinguishing continuation would receive the same prediction. Thus the number of initialized values is at least the quotient cardinality.
\end{proof}

The underlying quotient argument is classical and is developed elsewhere in the collection~\cite{GA:understanding,GA:resolution}. Here its role is to attach a separate initialization condition to a common predictor across the surviving candidate worlds. A quotient can exist while the available record fails to initialize it.

\subsection{Finite checking does not require an infinite list of words}
Begin with the partition by $q$. Split any block whose members have successors in different current blocks under some action. A strict split increases the number of nonempty blocks, so at most $|X|-|q(X)|$ strict rounds occur. At a fixed point the quotient updates are well defined and preserve every later observation by induction. Any operation-stable refinement of the original output partition refines each round, so the terminal partition is the behavioral equivalence above. This supplies a finite computation when the complete candidate tables are given. It does not make their empirical acquisition free.

\subsection{State identity is relative to the chosen questions}
In the sensor example, $r=y$ satisfies (1) for display prediction. The quotient has two states even though eight worlds remain in the candidate universe. Adding physical-bit output refines the quotient to four classes and invalidates (1) for that same initial code. Adding the nuisance bit to the target would refine again to eight, but no admitted sensor action reveals it.

Only adding requirements yields a refinement relation automatically. Replacing one task with a different task can produce crossing partitions rather than a finer version of the old one. For example, the partition by display and the partition by physical bit are neither finer nor coarser than one another. A claim of cumulative learning must specify which old questions remain protected.

\subsection{An available predictor is stronger than a worldwise predictor}
Consider two immutable worlds, one with a permanently zero output and one with a permanently one output. The available initial record is the same symbol in both. Within each world separately, a one-state predictor exists: it simply prints that world's output. There is nevertheless no predictor that can be initialized from the common record and be correct in both. The obstruction is the order of quantifiers,
\[
\forall\omega\ \exists P_\omega
\qquad\text{versus}\qquad
\exists P\ \forall\omega.
\]
The first does not provide a selector for the second. The behavioral quotient of the combined family has two states, but the evidence cannot choose between them. Adding an output observation repairs the initialization problem without changing the mathematical existence of either predictor.

Likewise, a coordinate convention is not an observation. Relabeling a latent state and conjugating its updates can leave every operational prediction unchanged. Without a separately meaningful reference, evidence about those predictions need not favor one coordinate system. A canonical response profile avoids claiming that an arbitrary coordinate name was discovered as a physical entity. It does not prove that the reasoner's actual implementation uses the inferred profile causally.

Thus two distinct questions remain after fitting a predictor. Is it a common valid representation of the possible worlds, and does the observer possess the evidence needed to initialize it? The theorem answers their exact conjunction in a finite model. Efficient acquisition and causal attribution require further premises.


\FloatBarrier
\chapter{New questions and the memory needed to correct them}\label{GA:book:chapter-26}

\section{Reinterpretation, acquisition, and irrecoverable distinctions}\label{GA:sec:trichotomy}
A revised task often invites the vague instruction to learn a better representation. That instruction hides three different possibilities. The old code may already contain everything needed. A new observation may be required. Or the remaining interface may be unable to recover the needed distinction at all.

We first separate these cases in a setting where observations do not interfere with one another. Let $\Omega$ be finite, let $r:\Omega\to Z$ be retained evidence, and let $t:\Omega\to Y$ be the revised target. Let $\GAE=\{e_1,\ldots,e_m\}$ be a finite family of total, deterministic, nondestructive tests. Each test reads the same fixed world and leaves it unchanged. All tests remain available throughout the investigation. Equivalently, a reset may be allowed if it genuinely restores the same unknown world, including the relevant history. Resetting into a new known state does not satisfy this premise.

For $J\subseteq\{1,\ldots,m\}$, write
\[
R_J(\omega)=\bigl(r(\omega),(e_j(\omega))_{j\in J}\bigr).
\]
Let the unresolved target pairs be
\[
D=\{\{\omega,\omega'\}:r(\omega)=r(\omega'),\ t(\omega)\ne t(\omega')\}.
\]
A test covers such a pair when it gives different reports in the two worlds.

\begin{GAtheorem}[Exact revision trichotomy for a reusable test interface]\label{GA:thm:trichotomy}
Exactly one of the following applies.
\begin{enumerate}
\item $D$ is empty. A new decoder of the old code recovers $t$ without new observations.
\item $D$ is nonempty, but every pair in $D$ is separated by at least one admitted test. Executing all tests recovers $t$, and a selected nonadaptive set $J$ suffices exactly when it covers $D$.
\item Some pair in $D$ has identical reports under every admitted test. No finite adaptive protocol using this unchanged test interface and $r$ guarantees $t$ in every world.
\end{enumerate}
If test $j$ has nonnegative fixed cost $c_j$, the minimum cost of one fixed nonadaptive test set, used for all initial codes in case 2, is precisely
\[
\min_{J:\,J\text{ covers }D}\ \sum_{j\in J}c_j.\tag{2}
\]
This is not a claim that an optimal adaptive protocol has the same cost.
\end{GAtheorem}
\begin{proof}
If $D$ is empty, the target is constant on each old-code fibre, giving a decoder. Otherwise a set $J$ permits a decoder exactly when no pair with the same old code and different target remains equal on all its test reports. That is exactly pair coverage. In particular, the full finite test family succeeds when every unresolved pair has a separating test.

If a pair survives every test, induction on an adaptive history shows that the two worlds give the same report after each chosen action. The protocol chooses the same next test because its available history is the same. Its stopping decision and output are therefore the same, so it cannot be correct on both targets. Independent randomization does not restore a probability-one guarantee, by conditioning on its random tape. The three cases are exclusive and exhaustive. Formula (2) follows from the equivalence between successful nonadaptive test sets and covers, with exactly the stated additive charges.
\end{proof}

The cover formulation exposes a concrete resource question without making a large learner part of the theorem. It also identifies what a negative certificate says. An uncovered pair is impossible to resolve through the declared interface. It is not impossible to resolve after a new sensor, retained record, or physical intervention changes that interface. Finding the cheapest cover can itself be difficult. Its formulation is not an efficient universal discovery algorithm.

\subsection{Why experimentwise evidence need not compose}
\begin{GAcounterexample}[Pairwise separation without a successful investigation]\label{GA:ce:destructive}
There are three possible initial cases $0,1,2$. Test $a$ reports whether the case is $0$, and test $b$ reports whether it is $1$. Either first test then destroys the accessible case state. Every later test returns a fixed zero report. The target remains the initial case, stored only as an inaccessible historical tag.

Every pair is distinguishable by some permitted first test. Yet beginning with $a$ leaves cases $1$ and $2$ indistinguishable whenever its answer is no. Beginning with $b$ leaves cases $0$ and $2$ indistinguishable. No finite adaptive protocol identifies every initial case. Randomizing the first test cannot guarantee success in every case either.
\end{GAcounterexample}

The missing premise is not intelligence. It is the ability to combine evidence from mutually destructive inquiries. State-identification work studies related restrictions on adaptive tests~\cite{GA:stateid}. The example makes clear why the target map, physical transition, and record cannot be separated in a general acquisition problem. The full state-changing repair constructions in \emph{Self-Improving Intelligence} address that different model~\cite{GA:selfimproving}.

\subsection{The sensor revisited}
For physical-bit recovery, an old display and a nondestructive reference measurement satisfy case 2. Without the reference, flip sequences leave the convention ambiguity unresolved. For recovery of the pre-reset physical bit after losing the old display, the surviving states can have the same current physical state, convention, and every future sensor response while having opposite historical targets. That is case 3 relative to the post-reset interface.

A learner should record which kind of failure occurred. A failed search over decoders cannot establish that no new experiment would help. Conversely, repeatedly searching for a decoder cannot reconstruct a target that is not a function of its inputs. Formula length and information access are separate resources.

\section{A minimal memory for future correction}\label{GA:sec:memory}
The old record can matter even when it does not change today's set of possible answers. The essential information may lie in a relation between a possible answer and an observation that will arrive later.

Let $\Omega$ be a finite model set and $t:\Omega\to Y$ a fixed target. For a candidate set $V\subseteq\Omega$, let
\[
Q(V)=t(V)=\{t(\omega):\omega\in V\}.
\]
The empty set is permitted. It represents a contradiction between the admitted model family and the evidence, not a successful identification. Let $\GAE\subseteq\GAPow(\Omega)$ be a finite family of evidence events containing $\Omega$ and closed under intersections. Filtering by $E$ sends $V$ to $V\cap E$. The events are known, truthfully interpreted, and do not change the underlying model in this section.

\begin{GAdefinition}[Correction signature]\label{GA:def:signature}
The correction signature of $V$ is
\[
\Sigma_{\GAE,t}(V)=\bigl(Q(V\cap E)\bigr)_{E\in\GAE}.\tag{3}
\]
Two candidate sets have the same signature when every admitted evidence filter leaves the same possible target values, including whether no candidate remains.
\end{GAdefinition}

\begin{GAtheorem}[The coarsest exact memory for admitted corrections]\label{GA:thm:correction}
Equality of correction signatures is a congruence for all updates $V\mapsto V\cap E_0$. The signature therefore supports well-defined updates and a readout of $Q(V)$. It preserves that readout after every finite sequence of admitted filters.

Conversely, any retained representation with deterministic filter updates and a readout reproducing $Q$ after every such sequence must distinguish every two candidate sets with different signatures. Thus the image of $\Sigma_{\GAE,t}$ is the smallest exact state space, in cardinality, for this correction contract on all candidate subsets.
\end{GAtheorem}
\begin{proof}
If $\Sigma(V)=\Sigma(W)$, then for every $E\in\GAE$,
\[
Q((V\cap E_0)\cap E)=Q(V\cap(E_0\cap E))
=Q(W\cap(E_0\cap E)).
\]
Intersection closure puts $E_0\cap E$ in $\GAE$, so the equality is available from the original signatures. It proves $\Sigma(V\cap E_0)=\Sigma(W\cap E_0)$. The $E=\Omega$ coordinate gives $Q(V)$, and induction gives all finite filter sequences.

If two signatures differ, one event $E$ produces different target sets. A retained representation that merged the two candidate sets would update identically on that event and return the same readout, contradicting exactness. Every distinct signature therefore requires a distinct retained value. The signature itself attains that count with the updates just constructed.
\end{proof}

This is a behavioral quotient of a powerset system. It preserves later task-valued correction, not merely the present prediction. Candidate identities can be discarded when no admitted correction distinguishes them. Contradiction with the entire candidate set remains observable.

\subsection{An exact count when evidence identifies a partition cell}
Suppose $A_1,\ldots,A_m$ are a partition of $\Omega$ into nonempty cells, and $\GAE$ consists exactly of all unions of those cells. Define the attainable evidence--target pairs
\[
D_t=\{(i,y): y\in t(A_i)\},\qquad r=|D_t|=\sum_i|t(A_i)|.
\]
A candidate set induces an occupancy record
\[
O(V)=\{(i,t(\omega)):\omega\in V\cap A_i\}.
\]

\begin{GAcorollary}[Exact correction-memory size for partition evidence]\label{GA:cor:occupancy}
Two candidate sets have equal correction signatures if and only if their occupancy records agree. On the domain of all subsets $V\subseteq\Omega$, there are exactly $2^r$ distinct signatures. A fixed-length exact binary memory therefore requires at least $r$ bits, and an $r$-bit occupancy vector attains this bound.
\end{GAcorollary}
\begin{proof}
Filtering by cell $A_i$ returns exactly the target values occurring in the $i$th part of $O(V)$. Thus equal signatures imply equal occupancy. Conversely, the target set after filtering by any union of cells is the union of their occupied target values, so occupancy determines every signature coordinate.

Every subset of $D_t$ is realized: for each included pair choose one world in its nonempty intersection $A_i\cap t^{-1}(y)$, and choose no worlds for excluded pairs. These intersections are disjoint for different pairs. Hence the occupancy map has $2^r$ values. Distinct values need distinct memory states by the preceding theorem. Binary counting gives the lower bound, and direct occupancy bits attain it.
\end{proof}

A particular inquiry may reach fewer candidate sets, reducing the number of attained occupancy records. The bound counts logical states, not interpreter bytes, update work, or evidence acquisition.

The exact evidence family matters. Two worlds $a,b$ in one cell with the same target have identical singleton occupancies. Adding the event $\{a\}$ distinguishes them: one filtered singleton retains the target and the other is empty. A new event that splits a cell can therefore require finer memory.

\subsection{Association can be the missing state}
\begin{GAcounterexample}[Equal marginals with different corrected answers]\label{GA:ce:marginals}
In the sensor model, let
\[
V_0=\{(s,\theta,k):s\oplus\theta=0\},\qquad
V_1=\{(s,\theta,k):s\oplus\theta=1\}.
\]
Take the target to be $s$ and the evidence cells to be the two values of $\theta$. Both candidate sets currently allow physical bits $\{0,1\}$, and both allow conventions $\{0,1\}$. Yet after the reference reports $\theta=0$, $V_0$ requires $s=0$ and $V_1$ requires $s=1$.

Retaining both marginal possibility sets loses the association needed to use the next measurement. Their equality is not a sufficient criterion for merging the records.
\end{GAcounterexample}

Here the nuisance bit duplicates every attainable $(\theta,s)$ pair. There are $2^8=256$ raw candidate subsets, but only $2^4=16$ correction signatures. Four occupancy bits preserve every admitted future correction. The present physical-target set alone has only four possibilities, requiring two bits, but that smaller code is inadequate for the correction contract.

Four association bits suffice where eight raw membership bits are redundant and two marginal bits are inadequate.

\begin{figure}[tbp]\centering
\includegraphics[width=.88\linewidth]{sources/monographs/grounded-abstraction/figures/correction_association.pdf}
\caption{Two evidence histories preserve identical physical-bit and convention marginals. Their joint possibility patterns differ. A reference report selects a column, producing opposite physical-bit conclusions. The nuisance bit is omitted because it changes neither pattern.}\label{GA:fig:association}
\end{figure}

\subsection{Correction memory has a scope}
The signature is minimal for its specified target and evidence family. New targets, events, candidate models, state-changing observations, or revised event meanings can split its classes. Old minimality remains a statement about the old contract.

If a future question asks for $k$, the occupancy is insufficient. When every available measurement ignores $k$, expanding the code cannot recover it. A finer signature can exist mathematically without being obtainable from the current interface.

\FloatBarrier
\chapter{Noisy evidence and quantitative justification}\label{GA:book:chapter-27}

\section{Identifiable does not mean certainly identified}\label{GA:sec:noise}
The earlier exact filters remove a world only when it cannot produce the observation. With noisy channels that criterion can retain every candidate indefinitely. A statistical procedure may nevertheless become increasingly reliable. These are different guarantees, not contradictory descriptions of learning.

For a fixed protocol, let two worlds with different targets generate finite transcript laws $P_0$ and $P_1$. An estimator sees only that transcript and independent randomness. Its average success under equal prior probabilities is bounded by the discrimination in the entire record.

\begin{GAproposition}[A finite-record discrimination bound]\label{GA:prop:testing}
For two distinct target values and equal world probabilities, every transcript-based estimator has average success at most
\[
\frac{1+\GATV(P_0,P_1)}{2},\qquad
\GATV(P_0,P_1)=\frac12\sum_h|P_0(h)-P_1(h)|.\tag{4}
\]
A deterministic function or a common stochastic processing channel applied to the transcript cannot increase this upper bound.
\end{GAproposition}
\begin{proof}
Convert the estimator into a binary test by reporting world zero when it outputs the first target and world one otherwise. A correct target output always gives a correct binary decision. If $f(h)\in[0,1]$ is the test's chance of declaring world zero, its success is
\[
\frac12\left(1+\sum_h f(h)(P_0(h)-P_1(h))\right).
\]
The sum is at most the sum of positive differences, which equals total variation. This proves (4). If $B(h,z)$ is a common processing channel and $g(z)\in[0,1]$ a processed-record test, then $f(h)=\sum_zB(h,z)g(z)$ is a raw-record test with the same two expectations. Thus every processed decision was already available from the original record.
\end{proof}

This is an elementary testing and data-processing argument, related to the classical comparison of experiments~\cite{GA:blackwell}. It does not identify the best protocol, establish the correctness of a sensor model, or make independent randomness into additional evidence. If an adaptive protocol is compared, its complete history must be the transcript in (4).

\begin{GAcounterexample}[Distinct laws with no finite certain identification]\label{GA:ce:bernoulli}
Let the unknown world be $\omega\in\{0,1\}$. A repeated observation is conditionally independent Bernoulli with success probability $1/4$ in world zero and $3/4$ in world one. The two one-observation laws differ, so the world is identified by its law. Every finite binary string nevertheless has positive probability under both worlds. No fixed finite sample supplies an error-free world classifier.

For odd $n$, majority voting is the likelihood-ratio rule. Its error in either world is
\[
e_n=\sum_{j=0}^{(n-1)/2}\binom{n}{j}(3/4)^j(1/4)^{n-j}.
\]
The first values are $e_1=1/4$, $e_3=5/32$, and $e_5=53/512$. These are exact analytical probabilities, not observed error rates from an agent experiment.
\end{GAcounterexample}

The example illustrates why logical candidate elimination and statistical confidence should not share a Boolean field named ``known.'' Both candidate worlds remain in the support of every fixed finite record. Their relative likelihoods can become very different. Treating the most likely world as certain discards the distinction that the exact support test preserves.

\begin{figure}[tbp]\centering
\includegraphics[width=.88\linewidth]{sources/monographs/grounded-abstraction/figures/noisy_identification.pdf}
\caption{Exact majority-test error for two stipulated Bernoulli worlds. The laws differ, and repeated independent evidence reduces error, but every finite record remains compatible with both worlds. The curve is evaluated from the displayed binomial formula.}\label{GA:fig:noise}
\end{figure}

\subsection{Identifiability is not a uniform rate}
Distinct models may induce arbitrarily similar laws. Unless a model class imposes a separation condition, a bounded number of observations need not distinguish every pair reliably. The total-variation term exposes that obstruction directly. A data-dependent summary cannot increase the discrimination contained in the actual record.

Conversely, a small class with known, well-separated distributions can admit a useful finite-sample procedure even when exact certainty is impossible. The statistical question is then how much error is acceptable and which assumptions justify its probability bound. We do not add a sample-complexity claim merely by labeling a model family finite. Separation, admissible protocols, and the sampling law have to enter the argument.

\section{An explicit radius for what remains unknown}\label{GA:sec:radius}
A set of candidate worlds may disagree on some future outputs while agreeing closely on others. It is useful to quantify that disagreement without selecting one world as if it had been identified.

Let $v:\Omega\to\mathbb R^d$ record a fixed finite collection of target quantities, where $d\ge1$. These can be future predictions on specified operation words. For a nonempty $B\subseteq\Omega$, define
\[
\ell_j(B)=\min_{\omega\in B}v_j(\omega),\qquad
u_j(B)=\max_{\omega\in B}v_j(\omega),
\]
\[
m_j(B)=\frac{\ell_j(B)+u_j(B)}2,\qquad
\rho(B)=\frac12\max_j\bigl(u_j(B)-\ell_j(B)\bigr).\tag{5}
\]
The radius concerns coordinatewise prediction in a common frame. It is not an uncertainty measure invariant under arbitrary nonlinear reparameterization.

\begin{GAtheorem}[Ambiguity radius and an explicit prediction certificate]\label{GA:thm:radius}
For every nonempty finite $B$,
\[
\inf_{z\in\mathbb R^d}\ \max_{\omega\in B}\|z-v(\omega)\|_\infty=\rho(B),
\]
and $m(B)$ attains the infimum. If $B'\subseteq B$ is nonempty then $\rho(B')\le\rho(B)$.

Fix nonnegative tolerances $\eta,\varepsilon$ and $\delta_{\mathrm{id}},\delta_{\mathrm{model}}\in[0,1]$. Suppose a data record $H$ produces a candidate set $B(H)$ with
\[
\Pr\{\omega_*\in B(H)\}\ge1-\delta_{\mathrm{id}}.
\]
Suppose a second event, of probability at least $1-\delta_{\mathrm{model}}$, guarantees
\[
\sup_{\omega\in B(H)}\|v(\omega)-\widehat v_H(\omega)\|_\infty\le\eta.
\]
Read the last supremum as zero when $B(H)$ is empty. When $B(H)$ is nonempty, form the midpoint $\widehat m$ and radius $\widehat\rho$ from $\widehat v_H$. If the returned numerical vector $z$ satisfies $\|z-\widehat m\|_\infty\le\varepsilon$, then
\[
\|z-v(\omega_*)\|_\infty\le\widehat\rho+\eta+\varepsilon\tag{6}
\]
on an event of probability at least $1-\delta_{\mathrm{id}}-\delta_{\mathrm{model}}$. The rule abstains on an empty candidate set. No independence between the two stated events is required.
\end{GAtheorem}
\begin{proof}
For each coordinate $j$, any point $z_j$ is at distance at least $(u_j-\ell_j)/2$ from one of the two extremes. Taking the largest coordinate width gives the lower bound. The midpoint is within half of every coordinate's width of every candidate, so it attains that bound. Inclusion can only narrow the coordinate ranges.

On the intersection of the coverage and model-accuracy events, the true world lies in the nonempty set and
\[
\|z-v(\omega_*)\|_\infty
\le\|z-\widehat m\|_\infty
+\|\widehat m-\widehat v_H(\omega_*)\|_\infty
+\|\widehat v_H(\omega_*)-v(\omega_*)\|_\infty.
\]
The three terms are bounded by $\varepsilon$, $\widehat\rho$, and $\eta$. The union bound gives the probability statement.
\end{proof}

The center is allowed to be any vector in $\mathbb R^d$. Requiring it to be the prediction of one admissible model can increase the optimum. Nor does a coordinatewise midpoint automatically define a coherent stochastic process or preserve logical invariants. Those are additional constraints on the permitted approximation, not consequences of (5).

\subsection{Three errors and two failure probabilities}
The radius measures unresolved disagreement among retained candidates. The model discrepancy controls how well their represented predictions match the intended quantities. Numerical error concerns the computation of the returned answer. The coverage probabilities express the chance that the premises fail. They are not added to the numerical tolerance as if they had the same units.

The probability assertion is unconditional over the stated evidence procedure. It is not automatically the same bound conditional on a rare event such as issuance of a certificate. Likewise, repeatedly changing the candidate family or stopping rule requires a validity argument for the actual procedure. A confidence set cannot be assumed valid merely because its contents look plausible.

\subsection{What compression does to the candidate set}
For a fixed raw-record channel, let $B(h)$ be its exact support set. If only $c(h)$ is retained, the correct support is
\[
B_c(z)=\bigcup_{h:c(h)=z}B(h).
\]
Indeed, a world can produce the coarse record exactly when it can produce at least one raw record in that fibre. Thus $B(h)\subseteq B_c(c(h))$, and the ambiguity radius cannot decrease for the same target vector. This is an information statement, not a command to retain every byte. The correction signature in Section~\ref{GA:sec:memory} identifies one principled way to compress while retaining the specified future uses.

For general confidence sets rather than exact supports, taking the same union still preserves the original coverage guarantee. It may be conservative. Selecting a smaller coarse-record set requires a new coverage argument, not the assertion that compression should make the model simpler.

\begin{GAcounterexample}[Agreement inside a misspecified family]\label{GA:ce:misspec}
Suppose every candidate in a supplied family predicts zero for the target, but the actual unmodeled process produces one. The family has radius zero, and its internal arithmetic can be exact. The conclusion is nevertheless wrong because the true process is outside the represented family. Neither a narrow interval nor a successful proof about the family establishes that membership premise.
\end{GAcounterexample}

This is why model adequacy is not treated as a reward for confidence. In the present paper it is a stated premise, or a separate statistical event, rather than an automatic consequence of a concentrated representation.

\section{A probability memory for repeated correction}\label{GA:sec:likelihoodmemory}
The exact signature in Section~\ref{GA:sec:memory} records which candidates remain possible. With noisy observations, all candidates can remain possible while their relative weights change substantially. Keeping only their support would then preserve logical possibility but discard the result of Bayesian correction. We now specify a different memory contract: retain both the probabilities of future reports and the posterior distribution of the selected target.

Let $\Omega$ be finite and nonempty, let $t:\Omega\to Y$ have finite image, and let $\mathcal B$ be a finite collection of experiments. Experiment $e$ has a finite outcome set and specified likelihoods
\[
\ell_{e,o}(\omega)\ge0,\qquad \sum_o\ell_{e,o}(\omega)=1.
\]
The world is fixed. Every experiment remains available, and each new report has this conditional law given the world, the chosen experiment, and all earlier reports. This is a fresh conditional-observation premise, not a claim that copied readings become new samples. No experiment changes the physical state in this section.

Partition $\Omega$ into the nonempty cells $A_1,\ldots,A_m$ on which the complete likelihood vector $(\ell_{e,o}(\omega))_{e,o}$ is constant. Write its common value as $\lambda_{e,o,i}$. The profile vectors of different cells are distinct by construction. For a prior $\mu$, retain
\[
p_{i,y}(\mu)=\sum_{\omega\in A_i,\ t(\omega)=y}\mu(\omega),
\qquad (i,y)\in D_t=\{(i,y):A_i\cap t^{-1}(y)\ne\varnothing\}.
\]
These are joint masses, not the two separate marginals. There are $r=|D_t|$ coordinates with nonnegative entries summing to one. Exact real probabilities are not a fixed $r$-bit code. A bit-complexity claim would also have to account for their representation and growing denominators.

\begin{GAtheorem}[Exact target correction from likelihood-profile memory]\label{GA:thm:likelihoodmemory}
Under the fixed-world likelihood model, the record $p=(p_{i,y})$ supports every finite sequence of report predictions and target-posterior corrections. Its next-report probability and update are
\[
z_{e,o}=\sum_{i,y}\lambda_{e,o,i}p_{i,y},\qquad
p'_{i,y}=\frac{\lambda_{e,o,i}p_{i,y}}{z_{e,o}}\quad(z_{e,o}>0),
\]
and its target readout is $\Pr(t=y)=\sum_i p_{i,y}$. A zero-probability report is marked impossible rather than assigned a normalized posterior.

Conversely, any deterministic memory, initialized for every prior on $\Omega$, that preserves all next-report probabilities and all target posteriors after every finite history must distinguish priors with different joint-mass records. In fact, the joint probabilities of the target and all experiment-report words of length at most $m-1$ already determine the joint-mass record. The memory is therefore minimal in the sense of its induced equivalence on priors, not in machine words or physical storage.
\end{GAtheorem}
\begin{proof}
Likelihoods are constant on each profile cell, so summing the ordinary Bayesian update over $A_i\cap t^{-1}(y)$ gives the displayed update. Its normalization and target readout are exact. Repeating this argument proves sufficiency for any fixed experiment sequence, and also for choices made from the retained report history without extra world-dependent information.

For necessity, a valid predictor and updater determine, for every word $w=((e_1,o_1),\ldots,(e_k,o_k))$, the masses
\[
M_{\mu,y}(w)=\sum_{t(\omega)=y}\mu(\omega)
                  \prod_{j=1}^k\ell_{e_j,o_j}(\omega).
\]
Indeed, multiply the successive predicted report probabilities by the final target posterior. If an earlier factor is zero, the joint mass is zero and no impossible conditional posterior is needed. The empty word gives the current target distribution.

For each cell $i$ and each $j\ne i$, choose one likelihood coordinate $b(i,j)=(e,o)$ for which $\lambda_{b(i,j),i}\ne\lambda_{b(i,j),j}$. Such a coordinate exists because profile vectors are distinct. On $\Omega$ the polynomial
\[
P_i(\omega)=\prod_{j\ne i}
\frac{\ell_{b(i,j)}(\omega)-\lambda_{b(i,j),j}}
     {\lambda_{b(i,j),i}-\lambda_{b(i,j),j}}
\]
is exactly $\mathbf1_{A_i}(\omega)$. It is a linear combination of likelihood monomials of degree at most $m-1$. Each monomial is the likelihood product of an admitted fixed experiment-report word. Consequently $p_{i,y}(\mu)=\sum_{t(\omega)=y}\mu(\omega)P_i(\omega)$ is a linear combination of the corresponding $M_{\mu,y}(w)$. For $m=1$ the empty product is one and the same argument uses only the current target masses. Two priors merged by an exact memory have equal predicted joint masses, hence equal joint-mass records. This proves the claimed separation and finite word-length bound.
\end{proof}

The proof uses a finite interpolation argument. Its polynomial coefficients are an identification device, not negative-probability experiments that a learner must execute. All likelihood values are supplied exactly. Very similar profile vectors can make the interpolation numerically ill-conditioned, so the theorem does not provide a stable estimator of $p$ from noisy frequency estimates. Nor does it infer the true likelihood model from a narrow posterior. The distinction between sufficient experiments and compressed statistical evidence is classical~\cite{GA:blackwell}; the present contract additionally keeps the selected target and all future update operations explicit.

\subsection{Two identical marginals, two opposite corrections}
Return to the sensor worlds $\omega=(s,\theta,n)\in\{0,1\}^3$, where $s$ is the target bit, $\theta$ is a convention, and $n$ is irrelevant. Suppose the only future experiment is a nondestructive noisy reference with
\[
\Pr(o=1\mid\theta=0)=\tfrac14,\qquad
\Pr(o=1\mid\theta=1)=\tfrac34.
\]
There are two likelihood-profile cells, one for each convention, and four nonempty profile--target pairs. The eight-world prior can be reduced to four joint masses without losing any specified prediction or correction. The nuisance coordinate disappears.

Consider a prior $\mu_+$ assigning mass $1/2$ to $(s,\theta)=(0,0),(1,1)$ and a prior $\mu_-$ assigning mass $1/2$ to $(1,0),(0,1)$, with any fixed nuisance convention. Both have uniform $s$ and uniform $\theta$. Both predict the next report $1$ with probability $1/2$. After that report, however,
\[
\Pr_{\mu_+}(s=1\mid o=1)=\tfrac34,\qquad
\Pr_{\mu_-}(s=1\mid o=1)=\tfrac14.
\]
After $k$ fresh reports all equal to one, these become $3^k/(1+3^k)$ and $1/(1+3^k)$. The difference is not a numerical error. It is the association that the separate marginals failed to retain. At every finite $k$, both target values remain possible under either prior.

\begin{figure}[tbp]\centering
\includegraphics[width=.95\linewidth]{sources/monographs/grounded-abstraction/figures/noisy_association.pdf}
\caption{Exact posterior target probabilities for two priors with identical target and sensor marginals but opposite association. Each horizontal step is a fresh noisy-reference report equal to one. The curves are conditional analytical values, not measured learning performance. Four joint masses support the update; two separate marginals do not.}\label{GA:fig:noisyassociation}
\end{figure}

\subsection{What a revised evidence language requires}
Let $\mathcal P_0$ be the partition of worlds into the nonempty old profile--target cells. Let $\mathcal P_1$ be the corresponding partition after likelihoods or the target have changed, still on the same finite $\Omega$. The old record stores the probability of each $\mathcal P_0$ cell. The new record would store the probability of each $\mathcal P_1$ cell.

\begin{GAcorollary}[Exact migration of a probability memory]\label{GA:cor:memorymigration}
A deterministic map reconstructing the new record from the old record for every prior exists if and only if every $\mathcal P_0$ cell is contained in a single $\mathcal P_1$ cell. When this condition holds, summing the relevant old masses gives the new record. If the new evidence or target strictly splits an old cell, no such universal reconstruction exists.
\end{GAcorollary}
\begin{proof}
When every old cell lies inside a new one, each new mass is the sum of the old-cell masses it contains. Conversely, suppose two worlds lie in the same old cell but in different new cells. Point-mass priors on those worlds yield the same old record and distinct new records. No deterministic map can send that one old record to both required new records.
\end{proof}

\begin{GAcounterexample}[An old exact memory does not upgrade itself]\label{GA:ce:memorymigration}
Two worlds have the same target and the same old report law, so point masses on them give an identical old record. A newly admitted experiment reports one with probabilities $1/4$ and $3/4$ in the respective worlds. The old record cannot predict its probability, although it was exact for every old question and update. Adding a field for the new likelihood profile does not determine its contents.
\end{GAcounterexample}

A restricted family of reachable priors can admit a reconstruction that fails for arbitrary priors. That restriction must be stated. Likewise, an archived raw record, a justified prior relation, or fresh evidence may help initialize the new memory, but its adequacy must be checked separately. Enlarging the candidate-world set is a different operation again: the old probabilities do not decide how much mass newly admitted worlds should receive.

\subsection{Exact sufficiency is not a finite-precision guarantee}
Even when two profiles differ only slightly, enough repeated evidence can make their posterior effects very different. Near-equal likelihoods are not interchangeable for an unbounded continuation contract. A finite-precision theory needs a horizon, a metric, a likelihood or evidence-mass condition, and a policy for outcomes of very small probability. Section~\ref{GA:sec:quantized} supplies a finite-horizon approximation contract with an explicit likelihood floor and retained-record denominator. It does not make rounding an exact sufficient statistic or provide an unbounded-time guarantee.


\FloatBarrier
\chapter{Grounded continuation without complete identification}\label{GA:book:chapter-28}

\section{What a partial operator calibration determines}\label{GA:sec:linear}
The finite-world examples show that complete identification can be unnecessary. A linear setting makes the positive result more structural. We will determine the exact domain on which every future operation sequence is already fixed by partial calibration data.

Let $V$ be a finite-dimensional vector space over a field, and let $\GAA$ be a finite nonempty set of operation names. The unknown operations are linear maps $T_a:V\to V$. Calibration supplies input vectors $x_1,\ldots,x_m$ and exact outputs $y_{a,1},\ldots,y_{a,m}$ for every $a$. The operator family is fixed, rather than being reselected after each action. Let
\[
U=\GASpan\{x_1,\ldots,x_m\}.
\]
The calibration is consistent precisely when every linear relation among the inputs is respected by the outputs, for each action. Under that condition it defines maps $S_a:U\to V$ with $S_ax_i=y_{a,i}$. Every $S_a$ has linear extensions to $V$ by extending a basis. Unless stated otherwise, \emph{all} such extensions are admissible. There is no hidden norm bound, symmetry, invertibility condition, or preferred matrix completion.

The compatible-model viewpoint is shared with the data-informativity literature~\cite{GA:informativity}. Our goal here is exact full-state prediction from a known initial vector, not stabilization, optimal control, or prediction from an unmeasured initial condition.

\begin{GAproposition}[One-step identification from exact calibration]\label{GA:prop:onestep}
For a fixed action $a$, the value $T_av$ is the same for every linear extension of $S_a$ if and only if $v\in U$. The common value is then $S_av$.
\end{GAproposition}
\begin{proof}
If $v\in U$, all extensions agree by definition. If $v\notin U$, choose a linear functional $\ell$ that vanishes on $U$ and satisfies $\ell(v)=1$, and choose a nonzero vector $b\in V$. Given any extension $T_a$, the map $T_a'(x)=T_a(x)+\ell(x)b$ is another extension, but $T_a'v=T_av+b\ne T_av$. Such a functional exists by extending a basis of $U$ with $v$.
\end{proof}

The proposition uses the complete unrestricted extension class. A justified structural prior can determine more. For example, knowing an operator belongs to a smaller parametric family can constrain its action outside $U$. Those constraints must be included as prior information, not inferred from the calibration span alone.

\subsection{A known first step can leave the known domain}
Even if $v\in U$, the vector $S_av$ can lie outside $U$. A second action then acts at an input whose response was not calibrated. One-step accuracy is therefore insufficient for reusable operator meaning.

Define a descending sequence of subspaces
\[
U_0=U,\qquad
U_{j+1}=\{v\in U:S_av\in U_j\text{ for every }a\in\GAA\}.\tag{7}
\]
The sequence is descending by induction. Every strict inclusion decreases dimension, so a fixed point is reached after at most $\dim U$ strict drops. Denote it by $U_\infty$.

\begin{GAtheorem}[The maximal calibrated continuation domain]\label{GA:thm:core}
For every integer $n\ge1$ and every $v\in V$, the values $T_wv$ for all operation words of length at most $n$ are independent of the chosen compatible extension family if and only if
\[
v\in U_{n-1}.
\]
Consequently, every finite full-state continuation is identifiable exactly on $U_\infty$. This is the largest subspace $L\subseteq U$ such that $S_aL\subseteq L$ for every $a$. On this subspace the known restrictions themselves realize the complete operation family.
\end{GAtheorem}
\begin{proof}
Membership in $U_{n-1}$ means that every path of at most $n-1$ calibrated operations stays inside $U$. This follows by induction from (7), beginning with $U_0=U$. All inputs on which the first $n$ operations must act are consequently calibrated. Their outputs are determined by the $S_a$, independently of every completion outside $U$.

For necessity, suppose $v\notin U_{n-1}$. If $v\notin U$, the one-step proposition supplies two extensions that disagree within one action. Otherwise the recursive definition gives a word of length $j\le n-1$ with a first exit from $U$. Its intermediate vectors before exit lie in $U$, so its exit vector $z\notin U$ is the same under every compatible extension family. Choose any next action $b$. Two extensions of $S_b$ can be made to disagree on $z$ by adding a linear map that vanishes on $U$. Keep the other action maps the same. Even if action $b$ occurred earlier, that modification has no effect before the exit because the earlier inputs lie in $U$. The two fixed extension families therefore produce different vectors after the exit word followed by $b$, whose total length is at most $n$. This contradicts uniform continuation identification.

At a fixed point, (7) gives $S_aU_\infty\subseteq U_\infty$. If $L\subseteq U$ is invariant under every restriction, induction gives $L\subseteq U_j$ for every $j$, so $L\subseteq U_\infty$. Applying the finite-horizon equivalence at all $n$ proves the all-word assertion.
\end{proof}

This theorem supplies an identifiable \emph{domain}, not necessarily an identified whole operator. It is exact for the stated completion class and full-state output. A smaller observation does not automatically enlarge the predictable domain: Proposition~\ref{GA:prop:noenlargement} proves that it does not under the present unrestricted completion class when the observation is nonzero. Sections~\ref{GA:sec:pairtest}--\ref{GA:sec:linearpredictor} characterize the gains that become possible under explicitly structured uncertainty. Similarly, a learned nonlinear latent representation may compress the identifiable core further for its own tasks, but that is not part of this linear result.

\subsection{A two-dimensional calibration with a one-dimensional reusable core}
\begin{GAcounterexample}[Calibration does not close under continuation]\label{GA:ce:linear}
Let $V$ have basis $e_1,e_2,e_3$, let $U=\GASpan(e_1,e_2)$, and let the single calibrated operation satisfy
\[
Se_1=e_1,\qquad Se_2=e_3.
\]
The first continuation is known on all of $U$. The next calibrated domain is $U_1=\GASpan(e_1)$, which is already invariant. One completion has $T_0e_3=0$ and another has $T_1e_3=e_1$. Both agree on every supplied calibration. Starting from $e_2$, their first outputs agree at $e_3$, but their second outputs are respectively zero and $e_1$.
\end{GAcounterexample}

Every continuation from $e_1$ is nevertheless determined. The system can ground a genuinely reusable operation on that line while remaining unable to predict repeated operations from $e_2$. A new exact calibration of $e_3$ expands the known input span to $V$ and identifies the whole operator. Obtaining that calibration is a new experiment, not an algebraic consequence of the old data.

\begin{figure}[tbp]\centering
\includegraphics[width=.95\linewidth]{sources/monographs/grounded-abstraction/figures/calibrated_core.pdf}
\caption{A partial calibration can support unlimited exact continuation on a proper invariant core. The known image of $e_2$ leaves the calibrated input subspace, so its next image is unrestricted. Agreement with every supplied one-step calibration is not agreement on two-step behavior.}\label{GA:fig:core}
\end{figure}

\subsection{Several operators and the scope of the initial vector}
With several actions, (7) retains only vectors that remain calibrated under \emph{every} permitted word. An added action can therefore shrink the core even if the old operators and data are unchanged. Restricting the permitted grammar to a specific sequence can enlarge the domain of predictable use. The full free-word contract is deliberately stronger than one selected trajectory.

The initial vector $v$ is known to the predictor in this section. If the actual input is hidden behind an evidence code, its initialization must satisfy an additional factorization condition like (1). A calibrated operation and an unobserved initial state are different missing inputs. They cannot be exchanged merely because both appear as uncertainty about the next output.

\section{Why small local error needs an ambient growth bound}\label{GA:sec:approx}
Exact calibration can justify a sharp domain. Approximate calibration needs a quantitative statement about errors that leave it. Let $V$ be a normed vector space and $L\subseteq V$ a linear subspace invariant under known bounded linear maps $S_a:L\to L$. Suppose the actual bounded linear maps $T_a:V\to V$ satisfy
\[
\|T_a\|\le A,\qquad \|S_a\|\le B,\qquad
\|(T_a-S_a)v\|\le\eta\|v\|\quad(v\in L),\tag{8}
\]
for common nonnegative bounds $A,B,\eta$. The norm on $T_a$ is an ambient bound. Its purpose is to control propagation of errors outside $L$, not only behavior on calibrated inputs.

\begin{GAproposition}[Approximate continuation with explicit leakage control]\label{GA:prop:approx}
For every $v\in L$ and every word $w$ of length $n\ge1$,
\[
\|T_wv-S_wv\|
\le\eta\|v\|\sum_{j=0}^{n-1}A^{n-1-j}B^j.\tag{9}
\]
For the empty word the discrepancy is zero. Integer zero powers in the displayed finite sum have their usual value one. The conclusion does not require the actual trajectory to remain in $L$.
\end{GAproposition}
\begin{proof}
Write $z_j=S_{a_j}\cdots S_{a_1}v\in L$ and $x_j=T_{a_j}\cdots T_{a_1}v$. For the next operation,
\[
\|x_{j+1}-z_{j+1}\|
\le A\|x_j-z_j\|+\eta\|z_j\|
\le A\|x_j-z_j\|+\eta B^j\|v\|.
\]
The initial discrepancy is zero. Iterating this scalar recurrence gives (9). Ambient boundedness applies to $x_j-z_j$ even when that error has left $L$.
\end{proof}

This is a telescoping error estimate, not a new perturbation principle. Similar continuation bounds already appear in the collection~\cite{GA:understanding,GA:resolution}. Its role here is to identify a premise that calibration on a proper subspace cannot certify by itself.

\begin{GAcounterexample}[An uncalibrated direction amplifies a small residual]\label{GA:ce:growth}
Let $L=\GASpan(e_1)$, let $Se_1=e_1$, and in the $\ell^1$ norm define
\[
T=\begin{pmatrix}1&0\\ \eta&M\end{pmatrix},\qquad \eta>0,\ M\ge0.
\]
The one-step error on the unit calibrated vector is $\eta$, independent of $M$. After two steps the error is $\eta(1+M)$. It can be arbitrarily large while the local calibration error is fixed. For $n$ steps the second coordinate is $\eta\sum_{j=0}^{n-1}M^j$.
\end{GAcounterexample}

A learned model can therefore match calibrated transitions closely and still fail rapidly under repetition. Increasing one-step accuracy helps only in relation to the length of use and the growth of the unseen directions. A measured small residual does not imply an ambient norm bound. The bound can come from a declared operator class, a separate global calibration argument, or a stronger structural certificate, but not from renaming the local fit ``stable meaning.''

\section{A smaller output is not automatically a weaker calibration requirement}\label{GA:sec:noenlargement}
The preceding full-state theorem suggests asking whether a small observable can remain predictable outside the calibrated core. The answer depends on the uncertainty class. Under the unrestricted extension class of Section~\ref{GA:sec:linear}, simply requesting fewer output coordinates does not enlarge the all-word identifiable domain unless the requested observation is identically zero.

Let $Q:V\to Y$ be a fixed linear observation. The initial vector $v$ remains known. All linear extensions of every calibrated restriction $S_a$ remain permitted, independently across action names but fixed throughout an execution.

\begin{GAproposition}[No enlargement under unrestricted linear completion]\label{GA:prop:noenlargement}
If $Q\ne0$, then for every $h\ge1$, the vectors $QT_wv$ are independent of the compatible extension family for every word of length at most $h$ if and only if $v\in U_{h-1}$. In particular, the all-word output-identifiable domain is exactly $U_\infty$. If $Q=0$, every initial vector is trivially output-identifiable.
\end{GAproposition}
\begin{proof}
On $U_{h-1}$ the full vectors are determined by Theorem~\ref{GA:thm:core}, so their $Q$ outputs are determined. Outside that subspace, follow the same first-exit construction as in that proof. At an uncalibrated input $z$, choose a linear functional $\ell$ vanishing on $U$ with $\ell(z)=1$. Because $Q\ne0$, choose $b$ with $Qb\ne0$. Replacing the next action by $T_a+\ell(\cdot)b$ preserves all calibration data and every earlier step of the first-exit path, but changes the next observed output by $Qb$. This produces disagreement within $h$ actions. Each modified action map is fixed, not chosen anew during execution. The zero-output case follows immediately.
\end{proof}

This corrects an overly broad suggestion in the first draft's surrounding prose. Output-relative prediction can be less demanding under \emph{structured} uncertainty. It is not automatically less demanding when every uncalibrated direction is allowed to return arbitrarily into the output. The previous full-state theorem and its proof are unchanged.

The next sections therefore change the model class explicitly. We first consider a finite family of complete candidate systems, retaining the fact that one unknown candidate is fixed. We then characterize the stronger situation in which an entire uncertainty family admits one common quotient. These are different questions and have different answers.

\section{A finite test for agreement under a fixed unknown model}\label{GA:sec:pairtest}
Let $V$ be $n$-dimensional over a field, let $\Theta$ be a finite nonempty model set, and let $T_a^\theta:V\to V$ be specified linear maps for each model $\theta$ and action $a$ in a finite nonempty alphabet. All words are allowed. Choose a reference model $0\in\Theta$ and write $A_a=T_a^0$. The shared linear observation is $Q:V\to Y$. Define
\[
I_Q=\{v\in V:QT_w^\theta v=QA_wv\text{ for every }\theta\in\Theta\text{ and every word }w\}.
\]
This is an initial-condition subspace. It is not assumed invariant under the physical maps. The common coordinates are part of this linear comparison model. With unrelated latent carriers, a justified identification of the carriers would be an additional premise.

For each $\theta$, form the paired maps and difference observation
\[
B_a^\theta=\begin{pmatrix}T_a^\theta&0\\0&A_a\end{pmatrix},\qquad
C=\begin{pmatrix}Q&-Q\end{pmatrix},\qquad
\Delta v=(v,v).
\]
The pair is evolved by the \emph{same} action word in both components. Define
\[
K_0^\theta=\ker C,\qquad
K_{j+1}^\theta=K_j^\theta\cap\bigcap_a(B_a^\theta)^{-1}K_j^\theta.
\]

\begin{GAtheorem}[Finite output-agreement test with persistent model tags]\label{GA:thm:pairtest}
The descending paired-kernel sequence stabilizes in at most $2n$ strict dimension drops. Its terminal subspace is
\[
K_\infty^\theta=\bigcap_w\ker(CB_w^\theta),\qquad
I_Q=\bigcap_{\theta\in\Theta}\Delta^{-1}K_\infty^\theta.
\]
Thus exact all-word output agreement for the finite candidate family is decidable by finite linear algebra. If an initial vector is excluded, a disagreeing candidate and a distinguishing word can be extracted. Words of length at most $2n$ suffice for this test. The conclusion compares fixed candidate models; it does not admit arbitrary switching between candidates during a word.
\end{GAtheorem}
\begin{proof}
By induction, $K_j^\theta$ consists precisely of paired states with zero difference output for all words of length at most $j$. The initial case is the empty word. Membership in the next intersection says that the present difference is zero and that applying any first action leaves a pair whose next $j$ actions have zero difference. This is exactly the length-$j+1$ condition.

Every strict inclusion lowers dimension in $V\oplus V$, which has dimension $2n$. Once two successive subspaces agree, the subspace is invariant under every $B_a^\theta$ and lies in $\ker C$, so all later intersections agree. This proves the all-word formula. At the diagonal initial pair, $CB_w^\theta\Delta v=QT_w^\theta v-QA_wv$, giving the paired-kernel formula. Row-space construction retains the operation word producing each tested row. If a vector fails the terminal kernel test, one such row is nonzero on it and supplies a witness. The dimension argument gives the stated conservative length bound.
\end{proof}

Finite linear-algebra equivalence testing has established precedents, including algorithms for probabilistic automata~\cite{GA:tzeng}. Reachability, observability, and minimal realization for multiple linear operations are also established subjects~\cite{GA:petreczky}. The contribution here is the fixed-compatible-model interpretation and its connection to evidence initialization and correction. We do not claim a new general automata-equivalence algorithm.

\subsection{Why the unknown model must not change between steps}
\begin{GAcounterexample}[A switching adversary creates a nonexistent failure]\label{GA:ce:switching}
Over $\mathbb Q$, let there be one action and two possible fixed matrices
\[
A=\begin{pmatrix}0&1\\0&0\end{pmatrix},\qquad
B=\begin{pmatrix}0&0\\1&0\end{pmatrix},\qquad
Q=\begin{pmatrix}1&1\end{pmatrix},\qquad v=\binom11.
\]
Both fixed systems have output sequence $2,1,0,0,\ldots$ from $v$, since $A^2=B^2=0$. But applying $A$ first and $B$ second gives $QBAv=1$, not zero. That hybrid trajectory belongs to neither candidate model. Treating fixed uncertainty as a freely changing matrix would incorrectly reject the valid output prediction.
\end{GAcounterexample}

Here $I_Q=\GASpan\{(1,1)\}$, obtained already from the one-step equality. It is not invariant under $A$ or $B$: $Av=(1,0)$ and $Bv=(0,1)$ leave that line. Those new physical vectors cannot be treated as fresh initial conditions with a forgotten model history. For example, starting both models at $(1,0)$ gives different next outputs. The original fixed-model executions nevertheless remain in output agreement.

This is why the identifiable initial set is not automatically a common physical invariant core. A reference predictor may still work: it follows $A$ from the initial vector and reports its outputs. It represents the common response sequence, not a claim that the unknown physical state always equals the reference state.

\begin{figure}[tbp]\centering
\includegraphics[width=.95\linewidth]{sources/monographs/grounded-abstraction/figures/fixed_model_outputs.pdf}
\caption{Two fixed candidate systems produce the same complete output history from $(1,1)$ even though their hidden states differ. Switching from one candidate to the other after the first step creates an inadmissible history with a different second output. The fixed-model tag is information about the allowed trajectories.}\label{GA:fig:fixedmodels}
\end{figure}

\subsection{More data restricts the model family, not the logical quantifier}
If new exact evidence removes candidates while preserving the reference model, the intersection defining $I_Q$ can only enlarge. Its practical value depends on whether the learner can justify that removal and initialize the resulting predictor. Adding a new candidate has the opposite effect. Neither operation warrants replacing the model label at each step. A candidate family that omits the true system can still yield an impressive but inapplicable agreement subspace.

\section{When uncertain dynamics admit one common observable quotient}\label{GA:sec:commonquotient}
Agreement on selected initial conditions is weaker than agreement on the whole space. The stronger problem has a useful exact solution. Retain the finite family and reference matrices of Section~\ref{GA:sec:pairtest}, and define the difference maps
\[
D_a^\theta=T_a^\theta-A_a.
\]
Their ranges contain all possible one-step discrepancies. Let
\[
E_0=\sum_{a,\theta}\GAim D_a^\theta,\qquad
E_{j+1}=E_j+\sum_a A_aE_j,\qquad E_\infty=\bigcup_jE_j.
\]
This increasing sequence stabilizes after at most $n$ strict dimension increases. Its limit is the smallest subspace invariant under all reference actions and containing every discrepancy range. It includes discrepancies that are invisible now but return to the output later.

\begin{GAtheorem}[Exact common quotient for structured model uncertainty]\label{GA:thm:commonquotient}
The following conditions are equivalent.
\begin{enumerate}
\item $QT_w^\theta=QA_w$ on all of $V$, for every fixed model $\theta$ and every word $w$.
\item $E_\infty\subseteq\ker Q$.
\item There exists a subspace $N\subseteq\ker Q$, invariant under every $A_a$, containing every $\GAim D_a^\theta$.
\item There is a surjective linear representation $r:V\to Z$, shared linear updates $G_a$, and a readout $\bar Q$ such that
\[
rT_a^\theta=G_ar\quad\text{for all }a,\theta,\qquad Q=\bar Qr.
\]
\end{enumerate}
When the conditions hold, $V/E_\infty$ gives one such quotient. The smallest-dimensional common linear quotient is $V/N_0$, where $N_0=\bigcap_w\ker(QA_w)$ is the reference system's unobservable subspace. Its dimension is $n-\dim N_0$. The shared quotient also preserves outputs if the candidate model is changed between actions, a stronger conclusion than fixed-model agreement on a restricted initial set.
\end{GAtheorem}
\begin{proof}
Assume condition 1. For any $a,\theta,w$, agreement as linear maps on all of $V$ gives
\[
QA_wD_a^\theta
=QT_w^\theta T_a^\theta-QA_wA_a=0.
\]
The last equality is condition 1 for the word that executes $a$ and then $w$. Thus $Q$ annihilates the range of every $A_wD_a^\theta$, whose span is $E_\infty$, proving condition 2. Condition 2 supplies condition 3 with $N=E_\infty$.

Under condition 3, $T_a^\theta N\subseteq N$ because $T_a^\theta=A_a+D_a^\theta$, $A_aN\subseteq N$, and the entire range of $D_a^\theta$ lies in $N$. Every candidate therefore induces the same map $G_a([v])=[A_av]$ on $V/N$. Since $QN=0$, the output factors through that quotient. This proves condition 4. Composing its commuting equations proves condition 1. The same composition works with a different model index at each action because $r$ and $G_a$ are shared.

The subspace $N_0$ is invariant under every $A_a$ and lies in $\ker Q$. Every invariant subspace in $\ker Q$, including $E_\infty$ when condition 2 holds, lies in $N_0$. Hence $N_0$ gives a valid common quotient. Conversely, the kernel of any representation in condition 4 is invariant under the reference actions and contained in $\ker Q$, so it lies in $N_0$. Rank-nullity gives the lower bound $\dim Z\ge n-\dim N_0$, attained by $V/N_0$.
\end{proof}

\subsection{Large hidden variation with exact visible behavior}
For a structured example over $\mathbb Q$, take either of the fixed maps
\[
T_0=\begin{pmatrix}1&1&0\\0&1&0\\0&0&0\end{pmatrix},\qquad
T_1=\begin{pmatrix}1&1&0\\0&1&0\\2&-1&3\end{pmatrix},\qquad
Q=\begin{pmatrix}1&0&0\end{pmatrix}.
\]
The hidden third-coordinate dynamics are different, and $T_1$ can make that coordinate grow. Nevertheless the discrepancy range is $\GASpan\{e_3\}$, which the reference leaves inside itself and $Q$ annihilates. The common quotient retains $(x_1,x_2)$, with update
\[
G(x_1,x_2)=(x_1+x_2,x_2).
\]
The output at step $k$ is $x_1+kx_2$ under both models, and even under switching between them. The output alone is not an initialized reusable state: equal $x_1$ with different $x_2$ gives different next outputs. Ignoring $x_3$ is justified by the no-return structure, not by its small size or its being temporarily unobserved.

\begin{GAcounterexample}[An invisible difference returns one step later]\label{GA:ce:hiddenreturn}
Let
\[
A=\begin{pmatrix}0&1\\0&0\end{pmatrix},\qquad
D=\begin{pmatrix}0&0\\1&0\end{pmatrix},\qquad Q=\begin{pmatrix}1&0\end{pmatrix}.
\]
The candidates are $A$ and $A+D$. Since $QD=0$, every one-step output agrees for every input. But $QAD=(1,0)\ne0$. Starting at $e_1$, the second outputs are zero and one. The initial discrepancy range is invisible, while its propagated span $E_\infty=V$ is not. Checking only $QD=0$ misses the failure.
\end{GAcounterexample}

\begin{figure}[tbp]\centering
\includegraphics[width=.95\linewidth]{sources/monographs/grounded-abstraction/figures/hidden_uncertainty.pdf}
\caption{Structured uncertainty may be harmless when every difference remains in an invariant, output-invisible subspace. The opposite example has zero first-step output discrepancy but a visible second-step return. The figure represents exact algebraic alternatives, not a claim that small learned residuals certify the invariant subspace.}\label{GA:fig:hiddenuncertainty}
\end{figure}

\subsection{What the criterion does not learn}
The family $\{T_a^\theta\}$, its shared carrier, and $Q$ are given. The closure in the discrepancy-span recursion uses exact linear algebra and decides whether those models share the stated quotient. It neither discovers the model class from finite observations nor proves that the true mechanism belongs to it. A structural prior such as an invariant hidden block can justify a smaller uncertainty family, but that prior must be carried as evidence or an explicit assumption.

\section{Initializing the smallest common linear predictor}\label{GA:sec:linearpredictor}
The paired-state criterion supplies a domain of identifiable output behavior. A learner may still lack the information needed to select the correct behavior within that domain. We now combine the two obligations without requiring the identifiable initial subspace to be physically invariant.

Let $L\subseteq I_Q$ be a specified linear initial-condition subspace and let the available initial record be a linear map $R:L\to H$. Define the reference reachable span and its unobservable part by
\[
W=\GASpan\{A_wv:v\in L,\ w\text{ any word}\},\qquad
N=\{z\in W:QA_wz=0\text{ for every word }w\}.
\]
Both $W$ and $N$ are invariant under the reference actions, and $N\subseteq\ker Q$. These spaces can be computed by finite reachable-span and row-kernel iterations. They concern a reference predictor, not a claim that every physical candidate remains inside $W$.

\begin{GAtheorem}[Evidence-initialized minimal common linear prediction]\label{GA:thm:linearpredictor}
There is a finite-dimensional linear predictor with initialization $d:R(L)\to Z$, updates $G_a:Z\to Z$, and readout $\bar Q:Z\to Y$ satisfying
\[
\bar QG_wdR(v)=QT_w^\theta v
\quad\text{for every }v\in L,\theta\in\Theta,w
\]
if and only if $\ker R\subseteq N$. When it exists, the smallest possible state-space dimension among such \emph{linear} predictors is
\[
\dim W-\dim N.
\]
The quotient $W/N$, initialized by $d(Rv)=[v]$, attains that minimum.
\end{GAtheorem}
\begin{proof}
If $v\in\ker R$, linear initialization gives the same state as input zero. Exactness therefore requires $QA_wv=0$ for every word, so $v\in N$. Conversely, if $\ker R\subseteq N$, the rule $d(Rv)=[v]$ is well defined and linear. Invariance supplies $G_a([z])=[A_az]$, and $QN=0$ supplies $\bar Q([z])=Qz$. Their iteration gives $QA_wv$, which agrees with every candidate by $L\subseteq I_Q$.

For minimality, let any competing linear predictor be given. Its reachable span is generated by vectors $G_wdR(v)$. Associate such a vector with $[A_wv]\in W/N$ and extend by linearity. To see this association is well defined, suppose a finite linear combination of the predictor vectors is zero. Applying any later word and the predictor readout shows that the matching combination of reference vectors has zero $QA_u$ output for every $u$. That combination therefore belongs to $N$. The induced map from the competitor's reachable span onto $W/N$ is a well-defined linear surjection, since the classes of $A_wv$ span $W/N$. Its dimension cannot be smaller than the displayed quotient dimension. The constructed quotient attains equality.
\end{proof}

The dimension statement is a linear-realization theorem in a declared algebraic model, not a lower bound for arbitrary nonlinear encodings, unbounded-precision tricks, or shortest source programs. Reachability and observability are classical tools~\cite{GA:petreczky}. The distinctive conjunction here is compatibility across the retained models, initialization from the actual record, and the smallest common linear predictor.

\subsection{Initial dimension, dynamic memory, and hidden state}
In Counterexample~\ref{GA:ce:switching}, set $L=\GASpan\{(1,1)\}$ and $R(x)=x_1$. The reference reaches $(1,1)$ and $(1,0)$, so $W=\mathbb Q^2$. The rows $Q=(1,1)$ and $QA=(0,1)$ have trivial common kernel, giving $N=0$. A two-dimensional linear predictor is therefore necessary even for this one-dimensional initial family: it must reproduce $2t,t,0,0,\ldots$ for every multiplier $t$.

For the three-coordinate structured example, $L=V$, $W=V$, and $N=\GASpan\{e_3\}$. Two predictor coordinates suffice. The record $(x_1,x_2)$ initializes them; $x_1$ alone does not, since its kernel contains $e_2$, which becomes visible after one action. Physical dimension, initial-family dimension, and minimum predictive dimension are different quantities.

\FloatBarrier
\chapter{Quantitative grounding with finite retained memory}\label{GA:book:chapter-29}

\section{From exact probability coordinates to a finite record}\label{GA:sec:quantized}
The likelihood-profile theorem identifies what an exact statistical memory must retain. It does not make its probability coordinates finitely representable. A rational posterior can acquire larger denominators after each update. Rounding those coordinates changes the object, so a finite-record theory needs a stated error, horizon, and observation model.

We now give one such theory. The result deliberately uses a conservative positive-likelihood condition rather than claiming uniform stability for every observation process. Stability and numerical approximation of nonlinear filters have an established literature~\cite{GA:legland}. Our construction concerns a static hidden world, a finite cell record, and an explicit finite-horizon error account. It is not a particle-filter theorem and does not inherit a contraction hypothesis from that literature.

\subsection{A concrete probability record}
For an integer $k\ge1$, write $\Delta_k$ for the probability simplex on $k$ cells. Fix a denominator $N\ge1$. A retained record has the form
\[
 \widehat p=(n_1/N,\ldots,n_k/N),\qquad n_i\in\{0,\ldots,N\},\quad \sum_i n_i=N.
\]
The cell identities, denominator, and decoding rule are fixed background. The $k$ integer fields can be stored literally in $k\lceil\log_2(N+1)\rceil$ bits. This is an upper bound for this specified encoding, not a minimum description length. It excludes the likelihood table, program, temporary arithmetic, and any retained raw observation history.

Define $\mathcal Q_N(p)$ by first setting $n_i=\lfloor Np_i\rfloor$. Distribute the remaining $N-\sum_i n_i$ units to the largest fractional remainders, with ties broken by a fixed cell order. This makes the normalization exact after every rounding. It is not independent coordinate rounding followed by an unspecified repair.

For a likelihood vector $\ell=(\ell_i)$ define the Bayes map
\[
 \mathcal B_\ell(p)_i=\frac{p_i\ell_i}{\sum_jp_j\ell_j}
\]
whenever its denominator is positive.

\begin{GAproposition}[Rounding and positive-likelihood update bounds]\label{GA:prop:roundbayes}
For every $p\in\Delta_k$,
\[
 \GATV(p,\mathcal Q_N(p))\le\frac{k}{2N}.
\]
Let $p,q\in\Delta_k$, let $0<\beta\le1$, and let every entry of $\ell$ and $\widetilde\ell$ lie in $[\beta,1]$. If $\max_i|\ell_i-\widetilde\ell_i|\le\varepsilon$, then
\[
 \GATV\bigl(\mathcal B_\ell(p),\mathcal B_{\widetilde\ell}(q)\bigr)
 \le \frac{\GATV(p,q)+\varepsilon}{\beta}.
\]
Each right-hand side may also be replaced by its minimum with one. The probability-field storage bound above applies to the retained rounded record, not to all working memory.
\end{GAproposition}
\begin{proof}
The remainder rule produces nonnegative integers summing to $N$. Each rounded coordinate differs from the original by at most $1/N$, so half the sum of absolute differences is at most $k/(2N)$.

For the update, first hold $\ell$ fixed. Given $0\le f_i\le1$, put $a=\sum_i\mathcal B_\ell(q)_if_i$ and $Z_p=\sum_ip_i\ell_i$. Then
\[
 \sum_i\bigl(\mathcal B_\ell(p)_i-\mathcal B_\ell(q)_i\bigr)f_i
 =\frac{\sum_i(p_i-q_i)\ell_i(f_i-a)}{Z_p}.
\]
Every value $\ell_i(f_i-a)$ lies in $[-a,1-a]$, an interval of length one. Since $p-q$ has total mass zero, its pairing with a function of oscillation at most one is bounded in absolute value by $\GATV(p,q)$. Also $Z_p\ge\beta$. Taking the supremum over $f$ bounds the fixed-likelihood difference by $\GATV(p,q)/\beta$.

Next hold $q$ fixed and set $a=\sum_i\mathcal B_{\widetilde\ell}(q)_if_i$. The corresponding numerator is $\sum_iq_i(\ell_i-\widetilde\ell_i)(f_i-a)$, whose absolute value is at most $\varepsilon$. Its denominator is at least $\beta$. A triangle inequality proves the displayed update bound. The cap by one is the universal range of total variation. The storage assertion counts the stated integer fields only.
\end{proof}

The estimate is sufficient and can be loose. In particular, it need not credit favorable observations that reduce an existing error. Its purpose is a guarantee whose entire cost and probability assumptions are visible, not an assertion that rounding usually produces its worst possible error.

\subsection{Approximate likelihood profiles}
Let $\Omega$ be a finite static world set with prior $\mu_0$. A fixed partition $C_1,\ldots,C_k$ refines the target: $t(\omega)=t_i$ throughout $C_i$. Each admitted experiment $e$ is nondestructive and has a finite report alphabet. Its actual conditional probabilities are $\ell_{e,o}(\omega)$. Reports are freshly generated from those likelihoods conditional on the fixed world and the chosen experiments.

Choose one valid representative report law $\bar\ell_{e,o,i}$ per cell. Assume, for every admitted experiment, report, cell, and world in that cell,
\[
 \beta\le \ell_{e,o}(\omega),\bar\ell_{e,o,i}\le1,
 \qquad |\ell_{e,o}(\omega)-\bar\ell_{e,o,i}|\le\varepsilon. \tag{10}
\]
Both the actual and representative report laws sum to one over reports. Exact profile cells satisfy this condition with $\varepsilon=0$ when their likelihoods have the positive floor. A coarser partition can be used with a nonzero $\varepsilon$, but target-constant cells are still required for the target readout.

After a history, let $p_t$ be the cell masses of the exact world posterior. Initialize the retained record by $\widehat p_0=\mathcal Q_N(p_0)$. For the next selected experiment and observed report use
\[
 \widehat p_{t+1}=\mathcal Q_N\!\left(\mathcal B_{\bar\ell_{e,o}}(\widehat p_t)\right).
\]
The reference posterior and rounded record are compared along the same actual action/report history. Experiment selection may depend on earlier retained records and independent randomness. It may not introduce a hidden world-dependent selector whose information is omitted from that history. The likelihood condition applies to every experiment the policy can select.

\begin{GAtheorem}[Finite retained memory for a declared correction horizon]\label{GA:thm:finiteprecision}
Under (10), set $\eta_N=k/(2N)$, $\gamma=\beta^{-1}$, and let $e_0$ be any valid bound on $\GATV(p_0,\widehat p_0)$. For every possible action/report history of length $t$,
\[
 \GATV(p_t,\widehat p_t)\le e_t
 :=\min\!\left\{1,\ \gamma^t e_0+(\gamma\varepsilon+\eta_N)
                         \sum_{j=0}^{t-1}\gamma^j\right\}. \tag{11}
\]
The same bound holds after mapping cells to target values. With the generic initialization bound $e_0\le\eta_N$, define
\[
 A_H=\sum_{j=0}^{H}\gamma^j,\qquad
 B_H=\gamma\sum_{j=0}^{H-1}\gamma^j.
\]
For any horizon $H\ge0$ and tolerance $\tau>\varepsilon B_H$, every integer
\[
 N\ge \max\!\left\{1,\left\lceil
       \frac{kA_H}{2(\tau-\varepsilon B_H)}\right\rceil\right\} \tag{12}
\]
is sufficient for target and cell error at most $\tau$ at every time through $H$. Empty sums are zero. This is a finite-horizon retained-record guarantee, not an unbounded-time or total-computation guarantee.
\end{GAtheorem}
\begin{proof}
Condition on a particular attainable history and its next experiment. In each positive-mass cell, the exact next likelihood is the conditional average of $\ell_{e,o}(\omega)$ within that cell. It remains in $[\beta,1]$ and within $\varepsilon$ of $\bar\ell_{e,o,i}$. For a zero-mass cell any value obeying these bounds gives the same zero contribution. Thus the exact next cell mass is a Bayes update with such an effective likelihood vector, even when the worlds within a cell do not have identical laws.

Proposition~\ref{GA:prop:roundbayes} and the rounding bound give
\[
 \GATV(p_{t+1},\widehat p_{t+1})
 \le\gamma\GATV(p_t,\widehat p_t)+\gamma\varepsilon+\eta_N.
\]
Iterating the uncapped scalar bound proves (11); total variation is also at most one. Aggregating cell coordinates into target coordinates cannot increase total variation, by the triangle inequality. Substituting $e_0\le\eta_N$ gives $\eta_N A_t+\varepsilon B_t$. These quantities are nondecreasing in $t$ because $\gamma\ge1$ and all terms are nonnegative. Inequality (12) makes the horizon-$H$ value at most $\tau$. It therefore suffices at every earlier step as well. Every report has positive probability under both updates by the floor assumption, so no zero denominator is silently normalized.
\end{proof}

Equation (12) is a sufficient precision requirement, not an exact minimum. An infeasible or large bound does not prove that the particular rounded filter is inaccurate. It means that this uniform certificate is insufficient. Conversely, increasing precision cannot eliminate the model term $\varepsilon B_H$. A false likelihood model does not become correct because its probabilities are stored with more bits.

\begin{GAcounterexample}[A tiny rounded-away mass can control a later posterior]\label{GA:ce:rareprecision}
Take two worlds with prior $(1-d,d)$, where $0<d<1/(2N)$. The remainder rule gives the retained record $(1,0)$. A binary report has likelihoods $(d^2,1)$ for outcome one. After that outcome, the actual posterior mass of the second world is
\[
 \frac{d}{d+(1-d)d^2}=\frac{1}{1+d-d^2},
\]
which tends to one as $d$ tends to zero. The rounded posterior remains zero on that world. The initial total-variation error is only $d$. There is no common positive likelihood floor across this family: the first likelihood tends to zero, and the other report is impossible in the second world. No nontrivial uniform continuation conclusion follows from the small initial rounding error alone.
\end{GAcounterexample}

The counterexample concerns evidence amplification, not a failure of the remainder arithmetic. Replacing the likelihood one by $1-d^2$ makes both outcomes possible in both worlds while preserving the limiting phenomenon. The minimum likelihood still tends to zero. Exact support, approximate probability, and justified abstention should therefore remain distinct data.

\begin{figure}[tbp]\centering
\includegraphics[width=.93\linewidth]{sources/monographs/grounded-abstraction/figures/finite_memory_horizon.pdf}
\caption{The noisy sensor with four retained cell masses, three future reports, and several denominator choices. The actual maximum error is enumerated over all eight three-report histories and their prefixes. The upper curve is the uniform positive-likelihood bound using generic rounded initialization. The gap is conservatism, not evidence of an uncounted learning gain. The installed likelihoods are exact; retained-record size does not count the program or temporary arithmetic.}\label{GA:fig:finiteprecision}
\end{figure}

\section{Approximate invisibility must survive feedback}\label{GA:sec:blockleakage}
The exact common-quotient theorem allows arbitrarily large hidden dynamics when those directions never return to the selected output. Approximation changes that fact. A small hidden-to-visible coupling can make hidden growth consequential. This section states the relevant dependence explicitly rather than replacing exact invariance by a small residual with no continuation bound.

Approximate simulation and bisimulation offer established ways to compare observable behavior with a declared precision~\cite{GA:girard}. The following is a narrower block comparison argument for discrete linear operations. It does not claim that its matrix bound is a minimal simulation metric or that stability of every candidate separately supplies a common switching bound.

Write the physical state as $(v_t,h_t)$ in a direct sum of two finite-dimensional normed spaces. The required output is the visible component $v_t$. Along an admitted execution suppose
\[
 \begin{pmatrix}v_{t+1}\\h_{t+1}\end{pmatrix}
 =\begin{pmatrix}G_t&B_t\\GAC_t&H_t\end{pmatrix}
  \begin{pmatrix}v_t\\h_t\end{pmatrix},
 \qquad \bar v_{t+1}=G_t\bar v_t. \tag{13}
\]
The reduced predictor knows the chosen $G_t$ and receives the same action sequence. It does not observe $h_t$. The blocks may come from a fixed unknown candidate or from an explicitly admitted varying family. Assume common operator-norm bounds
\[
 \|G_t\|\le g,\quad \|B_t\|\le b,\quad
 \|C_t\|\le c,\quad \|H_t\|\le d,
 \qquad g,b,c,d\ge0.
\]
Let $\|v_0\|\le V_0$, $\|h_0\|\le H_0$, and $\|v_0-\bar v_0\|\le e_0$. These initial bounds are supplied evidence or declared model assumptions. Put
\[
 M=\begin{pmatrix}g&b\\c&d\end{pmatrix},\qquad
 u_0=\begin{pmatrix}V_0\\H_0\end{pmatrix},\qquad u_j=M^ju_0.
\]
Powers of this nonnegative comparison matrix are bounds on norms. They are not a replacement physical dynamics inferred from measurements.

\begin{GAtheorem}[Visible-error bounds from a common block comparison]\label{GA:thm:blockleakage}
Every admitted execution satisfying (13) obeys
\[
 \|v_t-\bar v_t\|\le d_t
 :=g^te_0+b\sum_{j=0}^{t-1}g^{t-1-j}(u_j)_2. \tag{14}
\]
In particular, when $b=0$ and $e_0=0$, the visible prediction is exact for every finite time, without a stability requirement on $d$ or $c$.

If $g<1$, $d<1$, and
\[
 D=(1-g)(1-d)-bc>0,
\]
then the total visible discrepancy is bounded by
\[
 \sum_{t=0}^{\infty}\|v_t-\bar v_t\|
 \le \frac{e_0}{1-g}
       +\frac{b\bigl(cV_0+(1-g)H_0\bigr)}{(1-g)D}. \tag{15}
\]
These are sufficient common bounds. Failure of the inequality for $D$ is not a proof that a particular signed system is unstable. Conversely, the two diagonal conditions alone are insufficient for a uniform conclusion over the stated block class.
\end{GAtheorem}
\begin{proof}
The norm triangle inequality and the common block bounds give
\[
 \begin{pmatrix}\|v_{t+1}\|\\\|h_{t+1}\|\end{pmatrix}
 \le M\begin{pmatrix}\|v_t\|\\\|h_t\|\end{pmatrix}
 \le M^{t+1}u_0
\]
componentwise. The error satisfies
\[
 v_{t+1}-\bar v_{t+1}=G_t(v_t-\bar v_t)+B_th_t.
\]
Induction gives (14), including its empty sum at $t=0$. If $b=0$ and the initial error is zero, the recurrence has no forcing term and remains zero, regardless of hidden growth.

Under the three strict inequalities, $(I-M)^{-1}=D^{-1}\begin{pmatrix}1-d&b\\c&1-g\end{pmatrix}$ has nonnegative entries. Let $a=(I-M)^{-1}(1,1)^T$. Both coordinates are positive and $Ma=a-(1,1)^T<a$. Hence $M$ is a strict contraction in the weighted maximum norm with weights $a$, and its geometric series converges to $(I-M)^{-1}$. Summing the nonnegative right-hand side of (14) gives
\[
 \frac{e_0}{1-g}+\frac{b}{1-g}\sum_{j\ge0}(M^ju_0)_2.
\]
The second coordinate of $(I-M)^{-1}u_0$ is $(cV_0+(1-g)H_0)/D$, proving (15). The common comparison applies at each step and therefore also covers switching when that switching is explicitly allowed by the contract. It need not cover a family supplied only with unrelated per-model stability certificates.
\end{proof}

The result connects exact and approximate abstraction without collapsing them. Exact no-return can protect an output even when the hidden coordinate grows rapidly. With nonzero return, hidden growth and two-way feedback enter the error account. A small $b$ is useful only in relation to the horizon, the initial hidden bound, and the other blocks.

\begin{GAcounterexample}[An arbitrarily small return channel exposes an unstable hidden mode]\label{GA:ce:observableleak}
Let $v_0=\bar v_0=0$, $h_0=1$, and use scalar blocks $G=1/2$, $B=\delta>0$, $C=0$, $H=2$. The reduced output remains zero, while
\[
 v_t=\delta\sum_{j=0}^{t-1}(1/2)^{t-1-j}2^j
     =\frac{\delta(2^t-2^{-t})}{3/2}.
\]
The first error is $\delta$, yet the later error is unbounded for every fixed positive $\delta$. At $\delta=0$, every output is exactly zero. Thus a tiny positive coupling is not uniformly equivalent to exact no-return on an unbounded horizon.
\end{GAcounterexample}

\begin{GAcounterexample}[Stable diagonal blocks can have unstable feedback]\label{GA:ce:feedbackgain}
Take scalar blocks $G=H=1/2$ and $B=C=3/4$. Both diagonal subsystems are contractions. The full matrix sends $(1,1)$ to $(5/4)(1,1)$ and consequently grows on that direction. Here $D=1/4-9/16<0$. A certificate that checks only the two diagonal bounds misses the feedback loop.
\end{GAcounterexample}

\begin{figure}[tbp]\centering
\includegraphics[width=.93\linewidth]{sources/monographs/grounded-abstraction/figures/observable_leakage.pdf}
\caption{Exact scalar output discrepancies with the same small return coupling $1/64$ and visible multiplier $1/2$. The stable hidden multiplier $1/2$ gives $t2^{-(t-1)}/64$; the hidden multiplier $2$ gives increasing error. The vertical axis is logarithmic. At exactly zero return coupling, the discrepancy is zero at every time in both cases and is not plotted on the log scale. These are explicit linear models, not neural activation traces.}\label{GA:fig:visibleleakage}
\end{figure}

\section{One justified decision from memory and dynamics}\label{GA:sec:jointquant}
A probability-memory error and an output error concern different objects. The former is a distance between distributions on retained cells. The latter is a distance between a candidate response and its surrogate. They cannot be added without a loss function connecting both to the same decision.

Fix an actual evidence history. Let $\mu$ be the exact posterior on the stipulated finite world set and let $p_i=\mu(C_i)$. Suppose the retained distribution is $\widehat p$ with $\GATV(p,\widehat p)\le e$. For a declared later action word or response horizon, let $y(\omega)$ be the actual model response and let $\bar y_i$ be the response used for cell $i$. Assume
\[
 \sup_{\omega\in C_i}\|y(\omega)-\bar y_i\|\le d
 \quad\text{for every cell }i. \tag{16}
\]
The cell need not identify the physical state or candidate operator. It must support this common response allowance. A block bound such as (14) can supply it when the initial records, common visible updates, and hidden-state bounds match that cell.

Let the permitted decision set $\mathcal D$ be finite and nonempty. Its loss $L(y,a)$ takes values in $[0,1]$ and is $K$-Lipschitz in $y$, uniformly over $a$. Define
\[
 R(a)=\sum_{\omega}\mu(\omega)L(y(\omega),a),\qquad
 \widehat R(a)=\sum_i\widehat p_iL(\bar y_i,a).
\]
These are risks under the stated model posterior. They are not frequentist coverage probabilities or guarantees on the loss of a single realized world.

\begin{GAtheorem}[A common quantitative grounding certificate]\label{GA:thm:jointrisk}
Under (16), put $\xi=\min\{1,e+Kd\}$. Then
\[
 |R(a)-\widehat R(a)|\le\xi\quad\text{for every }a\in\mathcal D.
\]
If $\widehat a$ minimizes $\widehat R$, its true posterior regret is at most $2\xi$. If it is the unique approximate minimizer and
\[
 \widehat R(a)-\widehat R(\widehat a)>2\xi\quad\text{for every }a\ne\widehat a,
\]
then it is the unique true minimizer. A decision issued only when $\widehat R(a)+\xi\le\tau$ has true posterior risk at most the declared threshold $\tau$.

The conclusion can combine (11) and (14) only when they refer to the same history, cell meanings, future action word, initialization, and admitted candidate family. No guarantee that the true external process belongs to that family is supplied by this composition.
\end{GAtheorem}
\begin{proof}
Insert the intermediate risk $\sum_ip_iL(\bar y_i,a)$. Its difference from $R(a)$ is at most $Kd$ by (16), the Lipschitz condition, and summing posterior masses. Its difference from $\widehat R(a)$ is at most $e$ because the function $i\mapsto L(\bar y_i,a)$ lies in $[0,1]$. A triangle inequality gives $e+Kd$; both risks also lie in $[0,1]$.

Let $a_*$ minimize $R$. Then
\[
 R(\widehat a)\le\widehat R(\widehat a)+\xi
 \le\widehat R(a_*)+\xi\le R(a_*)+2\xi.
\]
For any other decision, the approximate gap loses at most $2\xi$ when the two risks are replaced by their true values, proving the strict-gap assertion. The threshold assertion is the upper risk bound itself. Each step uses the same posterior and surrogate responses, which is why separately valid certificates about different histories or targets cannot be substituted.
\end{proof}

This is a bounded-loss composition of two error accounts, not a newly derived general decision principle. Related recipient-relative bounds appear in \emph{Machine Understanding}~\cite{GA:understanding}. The additional content here is the instantiated chain from a finite retained probability record and a separately controlled hidden-return error to one declared decision.

\subsection{A complete sensor-to-predictor example}
Use the eight worlds $(s,\theta,n)\in\{0,1\}^3$. The four joint $(s,\theta)$ masses in the order $(0,0),(0,1),(1,0),(1,1)$ are initially
\[
 p_0=(3,1,1,3)/8,
\]
and $n$ is conditionally uniform. A noisy nondestructive reference has probability $1/4$ of report one under $\theta=0$ and $3/4$ under $\theta=1$. It does not depend on $n$. After three fresh reports all equal to one, the exact target probability is $\Pr(s=1\mid111)=41/56$.

The learner retains the four masses using $N=65\,536$ and the stated remainder rule after each update. There are 68 literal probability-field bits because each of the four counts includes values through $65\,536$. The initial masses are exactly representable, so $e_0=0$. With $\beta=1/4$, $\varepsilon=0$, and three observations, (11) gives the sharper bound
\[
 e_3\le\frac{1+4+16}{32\,768}=\frac{21}{32\,768}.
\]
This uses the actual initialization, not the larger generic bound in (12).

The later response is produced by three applications of one fixed candidate matrix,
\[
 T_n=\begin{pmatrix}1/2&(-1)^n/64\\1/8&1/2\end{pmatrix},
 \qquad (v_0,h_0)=(s,1).
\tag{17}
\]
The predictor for cell $(s,\theta)$ uses $\bar v_t=2^{-t}s$ and therefore reports $\bar y_i=s/8$. The nuisance bit $n$ affects the physical response but remains unidentified by the reference observations. With $g=d=1/2$, $b=1/64$, $c=1/8$, and $V_0=H_0=1$, the comparison determinant is $127/512>0$. Directly evaluating (14) at three prediction steps yields
\[
 d_3=\frac{481}{32\,768}.
\]
Choose between numerical predictions $a=0$ and $a=1/8$, with loss $L(y,a)=\min\{1,|y-a|\}$. It is bounded and 1-Lipschitz. The common radius is consequently
\[
 \xi=\frac{21+481}{32\,768}=\frac{251}{16\,384}.
\]
The exact companion computes the actual retained masses, approximate risks, both true risks, and their comparison. It verifies that the retained predictor selects $a=1/8$, that its approximate advantage exceeds $2\xi$, and that $\widehat R(1/8)+\xi<3/50$. Thus the same finite construction supports a unique posterior-optimal decision and a posterior-risk threshold of $0.06$ under the supplied model. The complete prior, sensor laws, candidate matrices, bit record, and loss function all belong to this one example.

\begin{center}\small
\begin{tabularx}{\textwidth}{@{}lXX@{}}\toprule
Decision & Retained-record risk & Exact model-posterior risk\\\midrule
$a=0$ & $23991/262144\ \approx0.09152$ & $173737/1835008\ \approx0.09468$\\
$a=1/8$ & $8777/262144\ \approx0.03348$ & $81105/1835008\ \approx0.04420$\\\bottomrule
\end{tabularx}
\end{center}
The certified upper risk for the selected decision is $12793/262144\approx0.04880$, below $3/50$. The actual cell-mass error is $1/114688$, smaller than its uniform allowance. The risk certificate uses the proved allowance, not a substituted post hoc exact error. Exact enumeration is available here only because the example is fully supplied and finite.

The probability coordinates do not encode $n$, and the dynamical allowance does not pretend they do. The result remains approximate precisely where the evidence leaves the physical response unresolved. Reducing the denominator, admitting an observation that depends on $n$, or increasing the hidden return coupling changes a different premise. These revisions require separate responses rather than an undifferentiated instruction to trust the old certificate.


\FloatBarrier
\chapter{A consolidated grounding contract}\label{GA:book:chapter-30}

\section{A common theoretical contract for grounded revision}\label{GA:sec:integration}
The results do not combine into one universal scalar score. They address different parts of a continuing claim.

An evidence contract says what the observer receives and which worlds remain possible. A predictive contract says what outputs must remain correct after which operations. A correction contract says which relationships between future evidence and target values must remain available. A calibration contract says where an operation is determined, and a quantitative contract controls any approximation and its growth.

A combined statement needs one representation and consistent premises. An unavailable input prevents initialization. An uncalibrated continuation prevents exact prediction. A lost association prevents correction. A narrow confidence set does not establish its own coverage.

\subsection{The complete sensor argument}
The eight-world sensor illustrates the chain without a large implementation. Initially the display and flip operation admit a two-state exact predictor. That predictor is grounded by the current display, even though the physical bit and sensor convention are not separately identified. Adding the physical-bit question refines the needed behavioral state, and the old code fails the initialization condition. A reference measurement repairs the information deficit while leaving the nuisance bit irrelevant.

Before calibration, the records corresponding to the two possible display values have the same marginal physical uncertainty. They must not be merged if later reference evidence is expected to identify the physical bit. The correction-memory theorem identifies the missing association. Resetting supplies convention information only after changing the physical state. Recovering the original state then depends on the retained old record. If that record is gone, the unchanged future interface cannot reconstruct it.

Noise changes the strength of the conclusion. A highly probable physical bit is not an exact state certificate. A quantitative response set can carry the remaining ambiguity into a tolerance, but only under coverage and model-accuracy premises. Every clause refers to the same task and same record history.

\subsection{What this adds to the machine theory}
The framework does not demand that an agent identify every detail of a world before it can reason. It gives restricted positive conditions for doing less. An unidentifiable physical realization can support an identified predictive quotient. An underdetermined operator can support a fully determined invariant domain. A large candidate set can be compressed exactly for future correction without preserving every candidate identity.

The additional discipline is that each reduction is indexed by what it preserves. One cannot transfer the word ``grounded'' from a successful prediction to an unrelated physical interpretation, from an exact model to a learned approximation, or from a present decision to every later correction. Such transfers require new evidence or a proved bridge.

\subsection{Three different reductions should not share one name}
The exact set-valued memory merges worlds only when the selected target and every admitted filter agree. The likelihood-profile memory preserves joint target probabilities and the future statistical correction law. The observable linear quotient suppresses directions that cannot return to the chosen output under its declared model class. They are related examples of operation-compatible reduction, but their carriers, updates, and information requirements are different.

The fixed-model agreement subspace is a fourth object. It describes allowed initial vectors, not necessarily a physical invariant state space. A reference predictor can realize its common outputs even when the hidden candidate trajectories diverge. Forgetting that model history and restarting from a current physical vector can create a different inference problem.

\subsection{What changes when the evidence contract changes}
A new likelihood can split an old probability cell. A new candidate operator can shrink $I_Q$. A new action can propagate a previously invisible discrepancy into the output. A new initial-code restriction can fail to initialize a quotient that still exists mathematically. Each revision therefore needs a named compatibility check. The smallest representation for an old contract is not an entitlement to make new claims without acquiring the missing distinctions.

None of these facts obliges a learner to retain every detail. In the noisy sensor, four joint masses replace an eight-world prior. In the structured linear example, two predictor coordinates replace three physical coordinates. Both savings are exact under their respective contracts. What must be preserved is the relation between the discarded information and the future operations, not discarded information for its own sake.

\subsection{A revision changes specific claims, not every claim at once}
The sensor-to-predictor construction provides one ledger for revision. Its probability cells are indexed by $(s,\theta)$, while the hidden candidate index $n$ is represented only through the response allowance. Each row below describes a change to that same example, not a newly selected favorable model.

\begin{center}\small
\begin{tabularx}{\textwidth}{@{}P{.25\textwidth}YY@{}}\toprule
Revision & What remains justified & What must be re-established\\\midrule
Change only $N$ & The likelihood-profile partition and physical model. & Rounded history, memory error, probability-field size, and the resulting decision margin.\\
Let the reference depend on $n$ & Old predictions under the old reference. & New profile cells and their masses. The old four masses cannot universally supply a split.\\
Admit a larger return coupling & The original observation update when its likelihoods are unchanged. & The future-response allowance and posterior risk. A former threshold may no longer be certified.\\
Retain only the two marginals & The present target and convention distributions. & The association required by later evidence. Equal marginals need not determine an update.\\
Treat a copied report as fresh & The recorded original observation. & The claimed conditional sampling law. Repetition in storage is not another independent measurement.\\
Replace a fixed model by switching & Results explicitly proved for the larger switching class, when their common bounds hold. & A restricted-initial-state fixed-model prediction need not survive this change.\\\bottomrule
\end{tabularx}
\end{center}

A failed old certificate can mean that its premises no longer apply, not that the earlier theorem was false. A valid new risk assessment may conclude that no action meets the requested threshold. Abstention then preserves the stated claim instead of relabeling an uncertain answer as certain. No automatic dependency extraction from arbitrary programs is asserted here.

\section{Open theoretical problems beyond the first edition}\label{GA:sec:open}
The first edition closes a restricted account of evidence, corrective memory, uncertain dynamics, and finite-horizon quantitative use. A complete theorem does not require an implemented general agent. It does require that the assumptions delimit the conclusion rather than merely decorate it.

\paragraph{Sharper approximate observable agreement.} The block bound gives a sufficient norm-based certificate, not an optimal abstraction metric. Cancellation, structured candidate correlations, and action-language restrictions can permit smaller errors than the common nonnegative comparison detects. Characterizing the least attainable error or the smallest approximate state remains a distinct problem.

\paragraph{Economical long-horizon probability memory.} The positive-likelihood estimate can grow rapidly with the horizon. It does not exploit statistical forgetting, contraction, or favorable evidence. Adaptive precision, support-aware safeguards, and stronger model conditions may permit better retained-memory guarantees. Such refinements must preserve the distinction between posterior risk, realized error, and frequentist coverage. The current record-size bound does not bound temporary arithmetic or total runtime.

\paragraph{Restricted operation languages and infinite model families.} The exact linear results admit every action word. An automaton restricting those words changes the observable-equivalence problem and the relevant minimum. Infinite parametric families require a certified finite spanning argument, polynomial identity, or another uniform device. Sampling a few parameter values does not supply that device.

\paragraph{Model admission and state-changing evidence.} Adding candidate worlds requires an explicit allocation of prior probability. A previous posterior does not determine that allocation. Noisy observations that change the hidden state require transition filtering, not only multiplication by a static likelihood. Those extensions need their own correction-memory and approximation theorems.

\paragraph{Learning the premises.} Candidate systems, likelihoods, calibration, and structural bounds remain supplied mathematical data in this edition. Determining them from interaction, proving useful sample and search complexity bounds, and establishing a learned implementation's causal use are further research questions. They are not prerequisites that erase the completed conditional results.

\paragraph{Independent assurance.} Prose proofs, exact finite illustrations, source comparison, and PDF rebuilding provide different forms of internal evidence. They do not establish independent peer review, a proof-kernel certificate, or an externally authenticated run. Those assurance layers can be developed without changing what the present statements assert.


\section{Conclusion}\label{GA:sec:conclusion}
Grounded abstraction relates a target, its possible realizations, accessible evidence, and the future operations a retained state must support. It is neither complete reconstruction by default nor confidence in a convenient coordinate system.

Incomplete knowledge can support exact continuing use. A likelihood-profile record preserves correction while dropping nuisance identity. Different linear systems can share one observable quotient. A reference predictor can reproduce common outputs without matching the unknown physical trajectory, provided the available record can initialize it.

Approximate use needs additional structure. Finite probability fields introduce rounding error at each update. A likelihood allowance and positive floor control its finite-horizon amplification. A hidden coordinate can be arbitrarily unstable when return is exactly absent, yet even a small positive return can become consequential. The common block bound and the rounded-memory bound meet only through a specified loss for one response and one posterior.

The integrated sensor illustrates that conjunction. Its probability record does not recover the hidden candidate, and its predictor does not silently identify the physical state. It nevertheless supports an explicit posterior decision within a certified allowance. Changing the observation law, precision, or return coupling changes a named obligation rather than granting universal reuse of an old certificate.

The first edition therefore supplies a theory of which claims survive what experience leaves unresolved, how a useful memory can preserve later correction, and when a reduced predictor remains justified across continued use. Its unresolved learning and implementation questions are visible applications of that theory, not claims supplied by the word grounded.

\FloatBarrier

\part{Large Signal Model: realization and predictive inference}
\setcounter{section}{0}\renewcommand{\thesection}{LSM.\arabic{section}}\renewcommand{\theHsection}{LSM.\arabic{section}}

\chapter{Realization and predictive closure}\label{LSM:book:chapter-31}
\section{Three objects that must remain distinct}\label{LSM:three-objects-that-must-remain-distinct}

An \textbf{observation} is an element \(x\) of a declared space \(X\).
Examples include a finite real signal, a complex field on a finite grid,
or an array of categorical labels. Its type includes shape, scalar field
or palette, and boundary convention. A zero-extended finite sequence and
a periodic sequence are different objects even when their displayed
samples agree.

An \textbf{archive} is a lossless representation of the accepted
observation. In a mathematical model this means an injective analysis
map with an explicit left inverse. In an implementation, conversion and
floating-point error must be accounted for separately. An archive of a
measured signal does not recover physical variables that were never
observed.

A \textbf{predictive state} is a summary \(z=\gamma(x)\), or more
generally a summary of an observation history. It is sufficient only
relative to a specified family of future quantities, actions and
dynamics. There is no implication from exact storage of \(x\) to
sufficiency of \(\gamma(x)\).

The historical model's macro, meso and micro vectors are nested
projections of a feature vector. They are not, merely by being nested,
state variables on which the signal dynamics close. The distinction is
retained in the present theory.

The term \textbf{large} denotes the structured hypothesis family:
representations, resolutions, typed operators, compositions, memory
mechanisms and their parameters. It is not a theorem about capacity,
sample complexity or scale. Such a theorem needs a declared grammar, a
complexity measure and a learning protocol.

\section{Exact realization and atlas consistency}\label{LSM:exact-realization-and-atlas-consistency}

\subsection{2.1 A left inverse is the reconstruction
contract}\label{LSM:a-left-inverse-is-the-reconstruction-contract}

Let \(X\) and \(C_a\) be finite-dimensional Hilbert spaces. A linear
atlas consists of analysis and synthesis maps

\[
A_a:X\longrightarrow C_a,\qquad S_a:C_a\longrightarrow X,
\qquad S_aA_a=I_X.
\]

The coefficients representing actual signals form
\(\mathcal C_a=\operatorname{im}A_a\). This distinction matters for
redundant representations, where \(\mathcal C_a\) is a proper subspace
of the ambient coefficient space.

\Needspace{5\baselineskip}
\textbf{Proposition 1. Exact realization and consistency.} The analysis
map is injective, \(\Pi_a=A_aS_a\) is an idempotent with image
\(\mathcal C_a\), and \(\Pi_ac=c\) exactly when \(c\) is a realizable
coefficient vector. Moreover,

\[
\|A_ax\|\geq \frac{\|x\|}{\|S_a\|},\qquad
\|S_a(c+e)-S_ac\|\leq\|S_a\|\,\|e\|.
\]

The first inequality is understood for nonzero \(X\). No assertion of
orthogonality follows from a left inverse alone.

\textbf{Proof.} If \(A_ax=0\), then \(x=S_aA_ax=0\). Also
\(\Pi_a^2=A_a(S_aA_a)S_a=\Pi_a\). Its image is contained in
\(\mathcal C_a\) and it fixes every \(A_ax\), so its image is exactly
\(\mathcal C_a\). Both norm inequalities follow from the operator-norm
inequality applied to \(S_a\). \(\square\)

For a unitary atlas \(U_a\), the synthesis is \(U_a^*\) and the
condition number is one. For a learned invertible atlas, exact
invertibility is not enough: the synthesis norm controls amplification
of coefficient errors. An arbitrarily ill-conditioned exact transform is
not a stable representation.

\subsection{2.2 Several exact atlases do not create new
observations}\label{LSM:several-exact-atlases-do-not-create-new-observations}

For unitary \(U_a\) and positive weights \(w_a\), put

\[
\mathcal Ax=(w_aU_ax)_a,\qquad
\mathcal S(c_a)_a=\frac{1}{c}\sum_aw_aU_a^*c_a,
\qquad c=\sum_aw_a^2.
\]

Then \(\mathcal S\mathcal A=I\) and \(\|\mathcal Ax\|^2=c\|x\|^2\), as
proved in {[}K1{]}. Here
\(\mathcal A\mathcal S=c^{-1}\mathcal A\mathcal A^*\) is the orthogonal
projection onto the consistent coefficient subspace.

Each complete atlas contains the same observed information because it
has a left inverse. Multiple atlases offer different sparsity, locality,
conditioning and operator descriptions. They do not provide independent
observations of the world. In particular, multiplying likelihoods from
deterministic recodings as though they were independent measurements
generally counts the same evidence more than once.

Categorical one-hot encoding is a separately typed injective map of
sets. It is not a unitary map on the scalar-valued signal space. Its
exact decoder is valid on one-hot arrays, not on arbitrary coefficient
vectors. The linear tight-frame theorem does not apply to it without a
separately specified linear embedding and inner product. This separation
is already explicit in the capstone's categorical section.

\section{Coordinate transport and one underlying operator}\label{LSM:coordinate-transport-and-one-underlying-operator}

On realizable coefficients, define transport from atlas \(a\) to atlas
\(b\) by

\[
T_{ba}=A_bS_a:\mathcal C_a\longrightarrow\mathcal C_b.
\]

For an endomorphism \(F:X\to X\), define its coordinate realization by
\(F_a=A_aFS_a\) on \(\mathcal C_a\). In this paragraph \(F\) may be
nonlinear. For a partial map, every formula is restricted to the
transported admissible domain. Maps between different signal types use
separate input and output atlases.

\Needspace{5\baselineskip}
\textbf{Proposition 2. Coherent transport and operator realization.} On
the consistent coefficient spaces,

\[
T_{aa}=I,\qquad T_{cb}T_{ba}=T_{ca},\qquad
F_bT_{ba}=T_{ba}F_a.
\]

Conversely, suppose maps \(G_a:\mathcal C_a\to\mathcal C_a\) satisfy
\(G_bT_{ba}=T_{ba}G_a\) for every pair. Then there is a unique map
\(F:X\to X\) with \(G_a=A_aFS_a\) for every \(a\).

\textbf{Proof.} The transport identities follow by inserting
\(S_bA_b=I\). The operator identity follows because either path is
\(A_bFS_a\). For the converse, choose a reference atlas \(r\) and put
\(F=S_rG_rA_r\). Compatibility gives \(G_aA_a=A_aF\). Every consistent
coefficient is \(A_ax\), which proves the asserted realization. Applying
\(S_a\) proves uniqueness. \(\square\)

An independently fitted family of atlas predictors need not satisfy this
proposition. Its compatibility residual is a well-defined quantity,

\[
d_{ba}(x)=G_b(A_bx)-T_{ba}G_a(A_ax).
\]

After synthesis it becomes an object-space discrepancy:

\[
S_bd_{ba}(x)=S_bG_bA_bx-S_aG_aA_ax.
\]

Thus \(\|S_bd_{ba}(x)\|\leq\|S_b\|\,\|d_{ba}(x)\|\). Zero discrepancy
establishes that the chart descriptions agree, not that their common
prediction is correct. All charts can agree on the same wrong operator.

Redundant coefficient outputs must also be tested for consistency.
Synthesizing an arbitrary coefficient vector always produces some
signal, but reanalysis returns \(\Pi_ac\), not necessarily the original
\(c\). Reconstruction and coefficient consistency are different checks.

\section{Packets, refinement and transfer blocks}\label{LSM:packets-refinement-and-transfer-blocks}

\subsection{4.1 Packets belong to a declared
resolution}\label{LSM:packets-belong-to-a-declared-resolution}

Within one orthogonal atlas, let \(E_i\) be pairwise orthogonal
coordinate selectors with \(\sum_iE_i=I\). Pull them back by a unitary
map:

\[
P_i=U^*E_iU,\qquad P_iP_j=\delta_{ij}P_i,
\qquad P_i^*=P_i,\qquad\sum_iP_i=I_X.
\]

A packet is the component \(P_ix\) together with its projector, domain,
atlas and resolution. A bare coefficient block without this information
is not a fully specified packet.

Projectors from different atlases need not commute and must not be
treated as one larger mutually orthogonal family. Weighted frame packets
can sum to the identity while failing to be idempotent. These are
different structures.

\Needspace{5\baselineskip}
\textbf{Proposition 3. Complete refinement preserves reconstruction.}
Suppose a finite rooted tree has root projector \(I\), and every
internal projector is the sum of pairwise orthogonal child projectors.
Every complete antichain \(\Lambda\), meeting each root-to-leaf path
exactly once, satisfies

\[
\sum_{v\in\Lambda}P_v=I,\qquad
x=\sum_{v\in\Lambda}P_vx,\qquad
\|x\|^2=\sum_{v\in\Lambda}\|P_vx\|^2.
\]

\textbf{Proof.} Replace an internal node by the sum of its children.
Repeated replacement reaches any complete antichain and preserves the
total projector. Different selected nodes have disjoint descendant
branches, hence orthogonal ranges. The norm identity is Pythagoras.
\(\square\)

This is a conservation law for information under complete refinement.
Deleting a child, storing only its energy, or imposing a coefficient
threshold changes the contract. For a retained subset \(J\) of an
orthogonal decomposition, the lost signal is exactly
\(\sum_{i\notin J}P_ix\), and its squared norm is the sum of the omitted
packet energies.

The finite best-basis recursion

\[
V(v)=\min\left(c(v),\sum_{u\in\operatorname{ch}(v)}V(u)\right)
\]

minimizes an additive cost over the declared finite tree {[}K1, R1{]}.
The recurrence itself does not make a nonadditive predictive objective
additive. If a prediction cost couples two leaves, a scalar cost at each
node is generally insufficient for global optimization. Packet
construction and cost evaluation must also be included in an
implementation's complexity account.

\subsection{4.2 Operations, not just signals,
decompose}\label{LSM:operations-not-just-signals-decompose}

Let \(L:X\to Y\) be linear, with orthogonal input projectors \(P_i\) and
output projectors \(Q_j\). Its transfer blocks are \(L_{ji}=Q_jLP_i\).

\Needspace{5\baselineskip}
\textbf{Proposition 4. Transfer reconstruction and refinement.} One has

\[
L=\sum_{j,i}Q_jLP_i.
\]

If \(P_i=\sum_rP_{i,r}\) and \(Q_j=\sum_sQ_{j,s}\) are complete
refinements, then

\[
Q_jLP_i=\sum_{s,r}Q_{j,s}LP_{i,r}.
\]

For another linear map \(M:Y\to Z\) and output projectors \(R_k\),

\[
R_kMLP_i=\sum_j(R_kMQ_j)(Q_jLP_i).
\]

\textbf{Proof.} Insert the appropriate identity decomposition on the
input, output or intermediate space. Use \(Q_j^2=Q_j\) in the
composition formula. \(\square\)

This is the algebraic reason to make transfer blocks first-class
objects. The same operation can be expressed at several resolutions,
while the coarse and fine expressions remain linked by exact identities.
Approximations acquire explicit omitted-block terms rather than an
unnamed representation error.

\subsection{4.3 Nonlinear operations need
interactions}\label{LSM:nonlinear-operations-need-interactions}

For a nonlinear map, \(F(\sum_iP_ix)\) generally differs from
\(\sum_iF(P_ix)\). Applying a linear transfer formula to such a map
silently loses interactions.

\Needspace{5\baselineskip}
\textbf{Proposition 5. Polynomial packet interactions.} Let

\[
F(x)=b+\sum_{m=1}^{p}B_m(x,\ldots,x),
\]

where each \(B_m:X^m\to Y\) is multilinear. Then

\[
Q_jF(x)=Q_jb+
\sum_{m=1}^{p}\ \sum_{i_1,\ldots,i_m}
Q_jB_m(P_{i_1}x,\ldots,P_{i_m}x).
\]

\textbf{Proof.} Substitute \(x=\sum_iP_ix\) in every argument and
distribute by multilinearity. \(\square\)

For example, \((u+v)^2=u^2+2uv+v^2\). A grammar containing only
independent packet squares misses \(2uv\). For nonpolynomial maps, one
needs direct nonlinear semantics or a local expansion with a stated
remainder bound. The word algebraic does not remove that obligation.

\section{Predictive closure is the central contract}\label{LSM:predictive-closure-is-the-central-contract}

Let \(\gamma:X\to Z=\gamma(X)\) be a summary and let \(F_u:X\to X\)
describe a deterministic transition for action \(u\). The reduced
description closes when there is \(\bar F_u:Z\to Z\) such that

\[
\gamma F_u=\bar F_u\gamma.
\]

This is an equality of functions on the declared domain, not a fit
statistic.

\Needspace{5\baselineskip}
\textbf{Proposition 6. Exact quotient criterion.} Such a reduced
transition exists if and only if

\[
\gamma(x)=\gamma(x')\ \Longrightarrow\
\gamma(F_u(x))=\gamma(F_u(x'))
\]

for every admissible \(x,x'\) and every action. When it exists it is
unique on \(\gamma(X)\).

\textbf{Proof.} Necessity follows by applying \(\bar F_u\) to equal
summaries. For sufficiency, define \(\bar F_u(z)=\gamma(F_u(x))\) for
any representative with \(\gamma(x)=z\). The stated condition makes this
well-defined and also proves uniqueness. \(\square\)

This criterion also applies when \(X\) is a latent physical state and
\(\gamma\) is an observation map. Two physically different states can
produce identical images and different futures. Exact image
reconstruction does not resolve that hidden-state ambiguity. A history,
a belief state or additional observations may be needed.

For a stochastic kernel, the corresponding condition is equality of the
distributions of the next summary for states with the same current
summary. Section 7 makes this finite criterion explicit.

\subsection{5.1 The obstruction in the linear
case}\label{LSM:the-obstruction-in-the-linear-case}

Suppose \(X\) is a Hilbert space, \(P\) is an orthogonal projection,
\(Q=I-P\) and \(x_{t+1}=Lx_t\). The retained state is \(z_t=Px_t\). Then

\[
z_{t+1}=PLPz_t+PLQx_t.
\]

\Needspace{5\baselineskip}
\textbf{Proposition 7. Closure defect and compositional defect.} The
retained state closes for every initial signal if and only if \(PLQ=0\).
If it does not close, the error in using only \(PLP\) is

\[
e_t=PLQx_t,\qquad
\|e_t\|\leq\|PLQ\|\,\|Qx_t\|.
\]

For two linear operators \(L\) and \(M\),

\[
\boxed{PLMP-(PLP)(PMP)=PLQMP.}
\]

Consequently, compressing two operations separately can discard a path
that leaves the retained subspace and later returns.

\textbf{Proof.} If \(PLQ=0\), the displayed recurrence depends only on
\(z_t\). Conversely, apply Proposition 6 to \(x\) and \(Px\), which have
the same retained state, to obtain \(PLQx=0\) for every \(x\). The norm
bound is submultiplicativity. The compositional identity follows by
inserting \(I=P+Q\) between \(L\) and \(M\). \(\square\)

The relevant quantity is not omitted energy alone. A very small omitted
packet can have a large effect if its transfer into retained packets is
strongly amplified. Conversely, an energetic packet can be irrelevant to
a specified prediction if every corresponding return block is zero.

A residual is therefore typed. Reconstruction residual,
coefficient-consistency residual, predictive residual, closure residual
and statistical uncertainty should not be collapsed into one scalar
without a specified objective.

\section{Constructing the smallest closed linear refinement}\label{LSM:constructing-the-smallest-closed-linear-refinement}

There is a constructive finite solution when the dynamics are a known
finite family of linear maps \(L_u\) on a \(d\)-dimensional space. It
identifies which additional observables are required, rather than merely
requesting a larger state vector.

Let the initial summary be \(C_0x\). Regard its rows as linear
observables. Starting from \(W_0=\operatorname{row}(C_0)\), define

\[
W_{n+1}=W_n+\sum_uW_nL_u,
\qquad W_nL_u=\{wL_u:w\in W_n\}.
\]

\Needspace{5\baselineskip}
\textbf{Proposition 8. Minimal linear predictive completion.} The
sequence stabilizes after at most \(d-\dim W_0\) strict enlargements.
Its limit \(W_*\) is the smallest row space containing \(W_0\) and
invariant under right multiplication by every \(L_u\). If \(R\) is any
full-row-rank matrix whose rows form a basis of \(W_*\), then

\[
C_0=HR,\qquad RL_u=B_uR
\]

for suitable matrices \(H\) and \(B_u\). Therefore \(z=Rx\) retains the
original summary and has exact reduced dynamics \(z_{t+1}=B_uz_t\).
Every linear summary with those two properties has row space containing
\(W_*\).

\textbf{Proof.} The spaces are nested and have dimension at most \(d\).
Equality at one step means invariance under every \(L_u\), so it is a
fixed point. Containment in any invariant row space containing \(W_0\)
follows by induction. The matrix identities express inclusion of the
relevant rows in the basis \(R\). Conversely, if \(C_0=HR\) and
\(RL_u=B_uR\), the row space of \(R\) is invariant and contains \(W_0\),
so it contains every \(W_n\). \(\square\)

For one transition matrix this reduces to the span of the rows of

\[
C_0,\ C_0L,\ \ldots,\ C_0L^{d-1}.
\]

Its annihilator is the largest \(L\)-invariant subspace inside
\(\ker C_0\), the standard unobservable-subspace construction {[}R2{]}.
With several actions, observables pulled back through action words are
required, but the linear span still stabilizes in finite dimension.

This gives a precise role for obstruction-driven refinement: add
observable directions created by the dynamics until the chosen
prediction variables close. It does not guarantee that the completion is
small. In some systems \(W_*=X^*\) and exact memoryless closure needs
the full state.

The theorem assumes the operators and relevant state are available.
Applying it to estimated operators constructs closure for those
estimates. It is not a proof that the real system has the same closure.
If only \(C_0x\) is measured and the full \(x\) is unavailable, the
augmented state \(Rx\) is not automatically accessible.

\section{Joint closure of prediction and Bayesian inference}\label{LSM:joint-closure-of-prediction-and-bayesian-inference}

The finite categorical case admits an equally explicit construction that
joins dynamics with Paper VI's inference refinement. This is useful
because a state can be closed for prediction but too coarse for a new
observation, or sufficient for an update but not for the next
transition.

Let \(\Omega\) be a finite state set. Let \(K_u(i,j)\) be known
row-stochastic transition matrices. Let \(\ell_y(i)>0\) be a finite
family of observation likelihoods. A partition \(\pi\) has aggregation
matrix

\[
C_\pi(a,i)=\mathbf 1_{i\in B_a}.
\]

Distributions are column vectors. Prediction is \(p\mapsto K_u^\top p\),
and a Bayesian update is

\[
\mathcal B_{\ell_y}(p)
=\frac{p\odot\ell_y}{\ell_y^\top p}.
\]

The coarse state must commute with both operations:

\[
C_\pi K_u^\top=\bar K_u^\top C_\pi,
\qquad
C_\pi\mathcal B_{\ell_y}(p)
=\mathcal B_{\bar\ell_y}(C_\pi p).
\]

Fix a mandatory initial partition \(\pi_0\) with at least two blocks.
This preserves distinctions that the intended application requires and
excludes the trivial one-state quotient.

\Needspace{5\baselineskip}
\textbf{Proposition 9. Coarsest joint finite closure.} First split
\(\pi_0\) by the likelihood vectors \((\ell_y(i))_y\). Then repeatedly
split each current block by the signatures

\[
\left(\sum_{j\in B}K_u(i,j)\right)_{u,B},
\]

where \(B\) ranges over the current partition. The process terminates
after at most \(|\Omega|-|\pi_{\mathrm{init}}|\) strict refinement
rounds. The final partition is the unique coarsest refinement of
\(\pi_0\) on which all declared predictions and Bayesian updates close
for every prior.

\textbf{Proof.} At a stable partition, states in a common block have the
same transition probability into every coarse block. Those probabilities
define \(\bar K_u\) and give the first commutation identity, by checking
it on point-mass priors. Likelihood constancy on each block defines
\(\bar\ell_y\) and gives the second identity by direct summation.

For minimality, every admissible partition refining \(\pi_0\) has at
least two blocks. By Proposition 12 below, exact updating for all priors
requires each fine likelihood to be constant within each block. Thus it
refines the first likelihood split. Suppose it refines a current
partition. Each current block is a union of its blocks. Exact prediction
makes the probability of reaching each such union constant within any
admissible block, so it also refines the next split. Induction proves it
refines the final partition. Each strict refinement increases the number
of blocks, which proves termination. \(\square\)

The transition condition is the finite lumpability criterion {[}R4{]},
expressed directly as a commutation identity. The joint construction
adds the explicitly declared likelihood distinctions. Its output is a
quotient of a known finite model, not an automatically learned model of
an unknown world.

The accompanying rational example starts with six states. Likelihoods
create three blocks, transition distinctions force two further rounds of
splitting, and the final model has five states. Two original states
remain genuinely equivalent for every action and observation in the
declared family. The example verifies both commuting identities exactly.

\section{Memory is an exact alternative to adding state}\label{LSM:memory-is-an-exact-alternative-to-adding-state}

Refinement and memory are two ways to represent the influence of
unresolved components. This can be derived rather than merely proposed
as an architectural option.

For fixed linear dynamics and an orthogonal split, write

\[
z_t=Px_t,\qquad q_t=Qx_t,
\]

and define the restricted blocks

\[
A=PLP,\quad B=PLQ,\quad C=QLP,\quad D=QLQ.
\]

Thus \(z_{t+1}=Az_t+Bq_t\) and \(q_{t+1}=Cz_t+Dq_t\).

\Needspace{5\baselineskip}
\textbf{Proposition 10. Exact projected memory and a conditional tail
bound.} For every \(t\geq0\),

\[
\boxed{
z_{t+1}=Az_t+
\sum_{k=0}^{t-1}BD^kCz_{t-1-k}+BD^tq_0.
}
\]

If \(\|D\|\leq\rho<1\) and \(\|z_s\|\leq Z\) over the relevant history,
omitting terms with \(k\geq H\) changes the one-step right-hand side by
at most

\[
\|B\|\,\|C\|Z\frac{\rho^H}{1-\rho}.
\]

Dropping the initial unresolved-state term adds at most
\(\|B\|\rho^t\|q_0\|\).

\textbf{Proof.} Iteration of the unresolved recurrence gives

\[
q_t=D^tq_0+\sum_{s=0}^{t-1}D^{t-1-s}Cz_s.
\]

Substitute this into the retained recurrence and change indices. Bound
the omitted terms by the geometric series. \(\square\)

The memory kernel is \(BD^kC\). It records a path from retained
variables into unresolved variables, evolution there, and return. The
first kernel \(BC\) is precisely the return path appearing in the
two-step composition defect. The term \(BD^tq_0\) is deterministic when
\(q_0\) is specified. It becomes uncertain only relative to incomplete
information or a probability model for \(q_0\).

This is the elementary linear instance of projection-memory reduction,
associated with Mori-Zwanzig theory {[}R3{]}. A finite hand-chosen
power-law or exponential memory feature is an approximation class for
such effects, not a general derivation of the exact kernel.

The tail bound compares right-hand sides using the same resolved
history. Replacing that history by recursively generated predictions
introduces additional rollout error. Stability of that recursion needs
its own bound. Also, a spectral-radius condition alone does not imply
\(\|D\|\leq\rho\) in the norm already being used, especially for
nonnormal operators.

\section{Statistical refinement needs its own theorem}\label{LSM:statistical-refinement-needs-its-own-theorem}

Nested feature vectors are algebraically nested when the smaller vector
can be recovered from the larger by a fixed map. Predictive risk is a
separate matter.

\Needspace{5\baselineskip}
\textbf{Proposition 11. Ideal statistical benefit of refinement.} Let a
target \(Y\) take values in a finite-dimensional Hilbert space with
\(\mathbb E\|Y\|^2<\infty\). Let \(\mathcal F\subseteq\mathcal G\) be
information sets and define
\(m_{\mathcal F}=\mathbb E[Y\mid\mathcal F]\) and
\(m_{\mathcal G}=\mathbb E[Y\mid\mathcal G]\). Then

\[
\mathbb E\|Y-m_{\mathcal F}\|^2
=
\mathbb E\|Y-m_{\mathcal G}\|^2
+\mathbb E\|m_{\mathcal G}-m_{\mathcal F}\|^2.
\]

\textbf{Proof.} Expand
\(Y-m_{\mathcal F}=(Y-m_{\mathcal G})+(m_{\mathcal G}-m_{\mathcal F})\).
The cross term has zero conditional expectation given \(\mathcal G\).
\(\square\)

Thus finer information cannot worsen the optimal squared-error predictor
under a fixed distribution. It can worsen an estimated predictor with
finite data, restricted function classes, changed regularization or
shifted distributions. The theorem supplies the ideal target for
refinement, not a guarantee for every fitting procedure.

An exact archive makes later reanalysis possible without reconstructing
discarded observations from a lossy state. It does not eliminate
estimation error or distribution shift.

\section{Bayesian coarsening and the information-loss packet}\label{LSM:bayesian-coarsening-and-the-information-loss-packet}

This section imports a specific existing Komorebi foundation rather than
inventing another inference heuristic. Let \(p_i\) be a prior on a
finite fine state set and let \(B_a\) be coarse blocks. Define

\[
P_a=\sum_{i\in B_a}p_i,\qquad
r_a(p)=\frac{\sum_{i\in B_a}p_i\ell_i}{P_a}
\quad(P_a>0).
\]

On zero-mass blocks the value of \(r_a\) is immaterial. Updating finely
and then aggregating is exactly updating \(P\) with the effective
likelihood \(r(p)\). The effective likelihood generally depends on the
fine prior inside each coarse block, information that \(P\) alone does
not retain.

\Needspace{5\baselineskip}
\textbf{Proposition 12. Bayesian closure and its exact defect.} Suppose
there are at least two coarse blocks and the fine and declared coarse
likelihoods are strictly positive. Updating and coarsening commute for
every prior if and only if

\[
\ell_i=c\bar\ell_a\quad\text{for every }i\in B_a
\]

with one common \(c>0\). For a fixed prior, put
\(\delta=r(p)-\bar\ell\), \(\bar b=\mathcal B_{\bar\ell}(P)\) and
\(Z_r=r(p)^\top P\). Then

\[
\boxed{
C_\pi\mathcal B_{\ell}(p)-\bar b
=
\frac{1}{Z_r}\left(\operatorname{diag}P-\bar bP^\top\right)\delta.
}
\]

The fixed-prior identity also holds for nonnegative likelihoods when
both normalizing evidences are positive.

\textbf{Proof.} Constancy up to one global factor gives commutation by
summing the fine posterior. Conversely, compare priors supported on one
state in each of two distinct coarse blocks. Equality of coarse
posterior odds forces the ratio of those states' likelihoods to equal
the ratio of the coarse likelihoods. Varying the representatives proves
constancy within each block and the common factor across blocks. For the
defect, use \(Z_r=\bar\ell^\top P+P^\top\delta\) and subtract the two
normalized vectors. \(\square\)

With only one coarse block, both posteriors aggregate to one regardless
of the likelihood, so the necessity criterion is false there. This
boundary is explicitly retained from the corrected Paper VI statement
{[}K2{]}.

The matrix in the defect formula annihilates \(\bar\ell\), recording the
irrelevant common-scale direction. It does not make every likelihood
discrepancy irrelevant. The construction in Proposition 9 refines until
both the likelihood and transition defects vanish for a specified finite
family.

\subsection{10.1 Parameter inference is not program
truth}\label{LSM:parameter-inference-is-not-program-truth}

Within one declared regression model, a normal-inverse-gamma prior gives
a standard conditional Bayesian parameter posterior. The software's
program weights instead use penalized held-out predictive scores and
then may refit coefficients on all available training records. These are
finite-candidate predictive weights. Without a corresponding generative
derivation, they should not be conflated with the marginal-likelihood
posterior over programs under one fixed prior.

Every candidate must be scored on the same target event and base
measure. Densities on differently sized coordinate spaces are not
directly comparable. An invertible nonlinear coordinate change requires
the appropriate Jacobian. A redundant deterministic embedding induces a
measure on its consistent subspace, not an independent full-dimensional
density on all coefficient entries.

The law of total variance separates within-candidate uncertainty from
disagreement of candidate means. It does not supply a probability that
the truth lies outside the grammar. An explicit outside-grammar
component needs its own prior and predictive distribution. Assigning an
arbitrary leftover mass does not solve open-world inference.

\section{What a typed operator language must guarantee}\label{LSM:what-a-typed-operator-language-must-guarantee}

A program has a source type, target type, admissible domain, evaluation
rule and declared approximation status. A nonlinear partial operation is
not a total matrix merely because its inputs are stored in a vector. The
domain of a composite is

\[
\operatorname{dom}(g\circ f)
=\{x\in\operatorname{dom}f:f(x)\in\operatorname{dom}g\}.
\]

Composition must preserve those types and domains. Coordinate changes
must preserve the represented operation as in Proposition 2. Packet
refinement must preserve exact transfer semantics as in Proposition 4,
or carry an explicit approximation residual. Nonlinear interactions must
retain their cross-packet terms as in Proposition 5.

A grammar can be closed under a specified set of constructors without
containing every possible physical law. It can also be expressive while
computationally intractable. Universal discovery, practical search and
semantic closure are distinct claims.

For a known operator, Proposition 8 can build its minimal linear
summary. For an unknown operator, learning must address identifiability:
which different candidates make identical predictions under the
available observations and actions? More coordinate systems alone do not
distinguish candidates that agree on the entire observed signal.
Additional interventions or data may be needed.

The foundation should therefore treat a learned hypothesis as an
operator with an evidence record, not just a matrix with a high
normalized weight. Its record includes fitting assumptions, admissible
inputs, compared alternatives, unresolved equivalences and the
representation in which its residual is measured.

\section{A bounded infinite-resolution bridge}\label{LSM:a-bounded-infinite-resolution-bridge}

Finite reconstruction identities do not alone prove convergence of an
unbounded or infinite-dimensional model. A useful bounded extension can
nevertheless be established directly.

\Needspace{5\baselineskip}
\textbf{Proposition 13. Strong fixed-horizon refinement limit.} Let
\(X\) be a Hilbert space, \(L:X\to X\) bounded, and \(P_n\) orthogonal
projections converging strongly to \(I\). Put \(L_n=P_nLP_n\). For every
fixed \(x\in X\) and every fixed nonnegative integer \(k\),

\[
L_nx\longrightarrow Lx,\qquad
L_n^kx\longrightarrow L^kx.
\]

\textbf{Proof.} Since \(\|P_n\|\leq1\),

\[
\|(L_n-L)x\|
\leq\|L\|\,\|(P_n-I)x\|+\|(P_n-I)Lx\|
\longrightarrow0.
\]

Also \(\|L_n\|\leq\|L\|\). Induct on \(k\), using

\[
(L_n^{k+1}-L^{k+1})x
=L_n(L_n^k-L^k)x+(L_n-L)L^kx.
\]

The first term tends to zero by the uniform norm bound and induction,
and the second by strong convergence on the fixed vector \(L^kx\).
\(\square\)

If estimated operators \(\widehat L_n\) are uniformly bounded and
\((\widehat L_n-L_n)x\to0\) for every fixed \(x\), the same argument
gives fixed-horizon convergence for \(\widehat L_n\). These are explicit
assumptions, not conclusions of the finite learning procedure.

This result is strong convergence, not operator-norm convergence and not
a bound uniform in time. On \(\ell^2\), finite coordinate projections
converge strongly to \(I\) while \(\|I-P_n\|=1\) for every \(n\).
Unbounded operators require domain control. Adaptive data-dependent
partitions and fitted coefficients require additional consistency
arguments. These remain separate mathematical tasks.

\section{A complete small signal example}\label{LSM:a-complete-small-signal-example}

Consider the rational linear signal dynamics

\[
x_{t+1}=Lx_t,\qquad
L=\begin{pmatrix}\tfrac12&1\\[2pt]\tfrac14&\tfrac13\end{pmatrix},
\qquad x_0=\begin{pmatrix}1\\0\end{pmatrix}.
\]

The spatial archive stores \(x_t\). A second atlas uses the orthogonal
Haar matrix

\[
U=\frac{1}{\sqrt2}\begin{pmatrix}1&1\\1&-1\end{pmatrix}.
\]

Its coefficients \(c_t=Ux_t\) reconstruct as \(x_t=U^*c_t\). The same
dynamics in that atlas are \(ULU^*\), not a separate physical law. Thus
coordinate transport preserves the full dynamics.

Now retain only the first spatial coordinate. With
\(P=\operatorname{diag}(1,0)\), the retained scalar \(z_t\) and
discarded scalar \(q_t\) obey

\[
z_{t+1}=\tfrac12z_t+q_t,\qquad
q_{t+1}=\tfrac14z_t+\tfrac13q_t.
\]

The exact first-coordinate sequence begins

\[
z_0=1,\qquad z_1=\tfrac12,\qquad z_2=\tfrac12,
\qquad z_3=\tfrac{11}{24}.
\]

The compressed one-step matrix predicts \(z_2=1/4\) instead of \(1/2\).
The lost return path contributes exactly

\[
BC=1\cdot\tfrac14=\tfrac14.
\]

There are two exact repairs. One is state refinement: start with
\(C_0=(1,0)\) and add \(C_0L=(1/2,1)\). The resulting row space is the
full two-dimensional space, so it closes. The other is memory:

\[
z_{t+1}=\tfrac12z_t+
\tfrac14\sum_{k=0}^{t-1}\left(\tfrac13\right)^kz_{t-1-k},
\]

since \(q_0=0\). This is the exact same system expressed without
retaining \(q_t\) as a separate current-state variable. At \(t=2\), it
yields

\[
z_3=\tfrac12\cdot\tfrac12
+\tfrac14\cdot\tfrac12
+\tfrac1{12}\cdot1
=\tfrac{11}{24}.
\]

To see a nontrivial reduced closure, adjoin a third independent
coordinate with multiplier \(5/6\). The smallest closed refinement of
the first coordinate then contains the first two coordinates, not the
third. Exact prediction of the chosen summary does not require
reconstructing every component of the enlarged signal.

Finally, low energy is not the same as low predictive importance. For

\[
L'=\begin{pmatrix}0&1000\\0&0\end{pmatrix},
\qquad x=(0,1/1000)^\top,
\]

omitting the second coordinate loses norm \(1/1000\), but changes the
next first coordinate by exactly one. The correct certificate is the
product of omitted magnitude and transfer amplification.

\section{Connection to observable learning}\label{LSM:connection-to-observable-learning}

Exact realization, complete decomposition and operator transport provide
the representation laws. Propositions 6--10 separate predictive closure,
constructive refinement and memory. The remaining requirement for
learning is access to a sufficient predictive state. Chapter~\ref{LSM:book:chapter-32} addresses a
finite linear observable sector; it does not infer its rank or coverage
assumptions from a successful fit.


\label{LSM:part-ii.-learning-stability-and-coherent-expansion}

The first stage constructed predictive closures for declared operators.
The second stage asks how those constructions relate to what a learner
can actually observe. A closure that depends on an unavailable hidden
coordinate is a mathematical description, not yet an accessible state. A
fitted recurrence can be accessible but wrong outside the span
represented by its data. A one-step approximation can be accurate and
still fail under repeated composition.

The results below connect these distinctions. They give an exact
learning result for a finite linear observable sector, deterministic
perturbation bounds, a propagation law for typed defects, a sufficient
stability condition for truncated memory, controlled nonlinear
interaction bounds, coherent atlas reconciliation and posterior
accounting for an expanding grammar. These are finite or explicitly
bounded results. They do not assume that every unknown process has a
short linear state.

The observable-state perspective has classical antecedents in predictive
state representations {[}R5{]} and observable-operator learning
{[}R7{]}. Finite invariant observable spaces also underlie Koopman
reductions {[}R6{]}. The present development uses elementary statements
proved below and identifies their specific place in Komorebi's
realization and obstruction calculus. It does not import those papers'
sample-complexity guarantees into this model.


\chapter{Learning, stability, and coherent expansion}\label{LSM:book:chapter-32}
\section{From hidden state to an accessible predictive state}\label{LSM:from-hidden-state-to-an-accessible-predictive-state}

Let a finite-dimensional autonomous linear process have hidden state
\(x_t\in\mathbb R^d\) and observed output \(y_t\in\mathbb R^p\):

\[
x_{t+1}=Lx_t,\qquad y_t=Cx_t.
\]

The observation map need not be injective. Define the length-\(h\)
observation map

\[
O_h=\begin{pmatrix}C\\CL\\\vdots\\CL^{h-1}\end{pmatrix},
\qquad V_h=\operatorname{im}O_h.
\]

The coordinates of \(O_hx_s\) are future observations relative to time
\(s\). They are not available at time \(s\) merely because this matrix
is written down. At time \(t\geq h-1\), however, the same vector is the
observed history

\[
z_t=\begin{pmatrix}y_{t-h+1}\\\vdots\\y_t\end{pmatrix}
=O_hx_{t-h+1}.
\]

Thus a forward observation map becomes a causal state after a declared
warm-up interval. The time index is part of the contract.

\Needspace{5\baselineskip}
\textbf{Proposition 14. Exact causal delay closure.} The following
statements are equivalent:

\begin{enumerate}
\def\labelenumi{\arabic{enumi}.}
\LSMtightlist
\item
  There is a linear map \(D_h:V_h\to V_h\) satisfying \(O_hL=D_hO_h\).
\item
  \(\ker O_h\subseteq\ker(O_hL)\).
\item
  \(\operatorname{rank}O_{h+1}=\operatorname{rank}O_h\).
\end{enumerate}

When these statements hold, \(D_h\) is unique on \(V_h\). The history
state obeys \(z_{t+1}=D_hz_t\), and its last output block is \(y_t\).
Its iterates therefore determine every subsequent observed output. If
\(d\) is an upper bound on the hidden dimension, \(h=d\) is sufficient
for closure, though a smaller \(h\) may suffice.

\textbf{Proof.} The equivalence of the first two statements is the
linear factorization criterion: define \(D_h(O_hx)=O_hLx\) and require
independence from the representative \(x\). The rows of \(O_hL\) are
\(CL,\ldots,CL^h\). All except the final block already occur in \(O_h\).
The kernel condition is therefore equivalent to \(CL^h\) being in the
row space of \(O_h\), which is exactly the rank equality. It also
implies \(O_hLx\in V_h\) because this vector is \(O_h\) applied to
\(Lx\). Uniqueness follows because every element of \(V_h\) has a
representative. The history recurrence follows directly. Cayley-Hamilton
expresses \(L^d\) as a linear combination of lower powers, so \(CL^d\)
is in the row space of \(O_d\). \(\square\)

This does not recover hidden variables in \(\ker O_h\). It identifies
the quotient needed for the declared outputs. An independent hidden
component can remain unobservable forever without preventing exact
output prediction.

The rank equality here is a property of the \emph{true observation
maps}. A plateau in a finite empirical data matrix is not the same
statement. It may reflect unexcited directions rather than a closed
observation space. Section 17 makes the difference exact.

For several actions \(L_u\), a fixed delay vector does not automatically
have one autonomous transition matrix. Products of the actual action
matrices determine its observation map, which can change with the action
history. One needs an action-aware state, a common closed observable
family, or a separately justified belief-state construction. Proposition
14 is not silently generalized to arbitrary controlled or stochastic
dynamics.

\section{What finite data identify}\label{LSM:what-finite-data-identify}

Suppose candidate states and their successors have been recorded as
columns of matrices \(Z_-,Z_+\in\mathbb R^{m\times N}\). Their
construction may use output histories. No hidden-state coordinates need
be supplied to the fitting procedure.

Write \(W=\operatorname{im}Z_-\). Exact interpolation of these columns
is a question about a linear map on \(W\), not yet about every
physically realizable state.

\Needspace{5\baselineskip}
\textbf{Proposition 15. Data-span identification and its boundary.} A
linear map \(B\) on the ambient space with \(BZ_-=Z_+\) exists if and
only if

\[
\ker Z_-\subseteq\ker Z_+.
\]

Its restriction to \(W\) is unique. It is an endomorphism of \(W\) if
and only if \(\operatorname{im}Z_+\subseteq W\). A particular ambient
extension is

\[
B_0=Z_+Z_-^\dagger,
\]

where \(\dagger\) denotes the Moore-Penrose inverse. Every other fitted
extension agrees with \(B_0\) on \(W\). If \(W\) is a proper subspace of
the true realizable state space \(V\), the data alone do not identify
the operator on \(V\setminus W\).

\textbf{Proof.} Necessity follows by applying \(B\) to every linear
dependence among the input columns. Under the kernel condition, define
\(B(Z_-a)=Z_+a\). Two choices of \(a\) differ by a kernel vector and
give the same answer. This defines and uniquely determines the map on
\(W\). It maps \(W\) into \(W\) exactly under the image condition. The
pseudoinverse formula gives that map because
\(a-Z_-^\dagger Z_-a\in\ker Z_-\). If \(W\subsetneq V\), choose a
nonzero linear functional \(\varphi\) on \(V\) vanishing on \(W\) and a
nonzero \(v\in V\). The operators \(B\) and \(B+v\varphi\) agree on
every observed column and differ elsewhere. \(\square\)

The three checks are different. The kernel test checks consistency of a
deterministic linear fit. The image test checks whether the observed
span is closed under the fitted transitions. Equality \(W=V\) is a
\emph{coverage condition}. Neither of the first two checks establishes
it.

For a known fixed feature space of dimension \(r\), coverage is
witnessed by \(r\) independent input states. For an unknown larger
feature space, the same rank only gives a lower bound on dimension. A
model card must say which interpretation is being used.

If several action-specific operators act on a common declared
\(r\)-dimensional state, each action's paired data must determine that
operator on the required span. Pooling all actions into one high-rank
matrix does not identify an action that has never been tried on some
direction.

\subsection{16.1 Perturbations have an explicit amplification
factor}\label{LSM:perturbations-have-an-explicit-amplification-factor}

A noisy rank test is not an exact certificate. Even after the state
space is fixed, fitting can be ill-conditioned. The following
deterministic bound isolates that issue without introducing a
probability model for noise.

Work in a fixed \(r\)-dimensional coordinate system, where
\(Z_-\in\mathbb R^{r\times N}\) has full row rank. Let

\[
Z_+=BZ_-+R,\qquad
\widetilde Z_-=Z_-+E,\qquad
\widetilde Z_+=Z_++F.
\]

The matrix \(R\) records deviations from the declared linear recurrence
on these samples. The matrices \(E\) and \(F\) record input and
successor perturbations. All norms in the next proposition are spectral
norms.

\Needspace{5\baselineskip}
\textbf{Proposition 16. Fixed-space identification bound.} Put
\(s=\sigma_{\min}(Z_-)>0\). If \(\|E\|<s\), then \(\widetilde Z_-\) has
full row rank, and

\[
\widehat B=\widetilde Z_+\widetilde Z_-^\dagger
\]

satisfies the exact error identity and bound

\[
\begin{aligned}
\widehat B-B
&=(R+F-BE)\widetilde Z_-^\dagger,\\
\|\widehat B-B\|
&\leq\frac{\|R\|+\|F\|+\|B\|\,\|E\|}
{s-\|E\|}.
\end{aligned}
\]

\textbf{Proof.} For every unit \(v\in\mathbb R^r\),
\(\|\widetilde Z_-^\top v\|\geq\|Z_-^\top v\|-\|E^\top v\|\geq s-\|E\|\).
Thus the perturbed matrix has full row rank and its pseudoinverse has
norm at most \((s-\|E\|)^{-1}\). Since
\(\widetilde Z_-\widetilde Z_-^\dagger=I_r\), subtract \(B\) from the
fitted expression and collect terms. The bound follows by
submultiplicativity and the triangle inequality. \(\square\)

This is a bound for a fixed state space and a fixed coordinate basis. If
the basis or its dimension is also estimated from noisy data,
subspace-estimation error must be added. No such error is hidden inside
\(R\) without a derivation. Correlated and adversarial perturbations are
allowed by the displayed deterministic inequality, but there is no
sample-complexity claim until their magnitudes are bounded under a
sampling model.

The denominator explains why numerical rank alone is insufficient. Full
rank can coexist with severe noise amplification. The reference
calculations include examples attaining the \(1/s\) dependence and a
rank-collapse example at the excluded boundary \(\|E\|=s\).

\section{Why finite agreement is not unrestricted discovery}\label{LSM:why-finite-agreement-is-not-unrestricted-discovery}

\Needspace{5\baselineskip}
\textbf{Proposition 17. Finite-prefix indistinguishability.} For every
integer \(T\geq0\), there are two finite-dimensional autonomous rational
linear systems with identical scalar observations at times
\(0,\ldots,T\) and different observations at time \(T+1\). Both
transition matrices can be chosen nilpotent.

\textbf{Proof.} One system has scalar output identically zero. For the
other, take dimension \(T+2\), initial state \(x_0=e_{T+2}\),
observation row \(C=e_1^\top\), and shift \(Le_j=e_{j-1}\) for
\(j\geq2\), \(Le_1=0\). Then \(CL^tx_0=0\) for \(0\leq t\leq T\) and
\(CL^{T+1}x_0=1\). The shift is rational and nilpotent. \(\square\)

An invisible delayed component of this kind can also be adjoined to any
fixed base model. The combined output agrees with the base output on the
entire inspected prefix and then differs. This is not an argument
against learning. It identifies why a claimed guarantee must state a
dimension bound, a model class, a coverage condition, an intervention
protocol or a probabilistic assumption.

Even nilpotent operators have only zero eigenvalues, so stability
language alone does not remove this ambiguity. A finite rank plateau
cannot certify the absence of a delayed mode when arbitrarily large
hidden dimensions are admitted.

\section{A complete output-only identification example}\label{LSM:a-complete-output-only-identification-example}

Return to the rational system from Section 13, now with an additional
independent hidden coordinate:

\[
L=\begin{pmatrix}
1/2&1&0\\
1/4&1/3&0\\
0&0&5/6
\end{pmatrix},
\qquad C=\begin{pmatrix}1&0&0\end{pmatrix}.
\]

The third coordinate never influences the observed output. The first
coordinate alone is not closed, but

\[
O_2=\begin{pmatrix}1&0&0\\1/2&1&0\end{pmatrix}
\]

has rank two and \(O_3\) has the same rank. Proposition 14 gives an
exact two-dimensional delay state. Choose \(x_0=(1,0,q)^\top\) for any
\(q\). The first four observations are

\[
y_0=1,\quad y_1=1/2,\quad y_2=1/2,\quad y_3=11/24.
\]

A learner receives those four outputs and the declared delay length
\(h=2\). It does \textbf{not} receive \(L\), \(C\), \(O_2\) or \(q\).
From consecutive observed windows it forms

\[
Z_-=\begin{pmatrix}1&1/2\\1/2&1/2\end{pmatrix},
\qquad
Z_+=\begin{pmatrix}1/2&1/2\\1/2&11/24\end{pmatrix}.
\]

Since \(\det Z_-=1/4\), these windows span the declared two-dimensional
state. Exact arithmetic gives

\[
\boxed{
D_2=Z_+Z_-^{-1}
=\begin{pmatrix}0&1\\1/12&5/6\end{pmatrix}.
}
\]

The learned recurrence is

\[
\boxed{y_{t+1}=\frac56y_t+\frac1{12}y_{t-1},\qquad t\geq1.}
\]

It predicts \(y_4=61/144\), which was not used to identify the
coefficients. Proposition 14 and the verified rank condition on the true
observation maps establish why this recurrence continues for every
initial state of this particular system. The four samples alone
establish the fitted recurrence only within the declared closed
two-state model class. Proposition 17 forbids confusing those two
conclusions.

The reference script separates the roles. One routine fits only the
observed samples. A different check, with access to the generating
system, verifies the theoretical assumptions and compares subsequent
outputs. This prevents hidden access to the true operator from being
presented as learning.

The same example connects back to the exact memory kernel. Section 13
expressed the output through a decaying convolution of its history. Here
the history is summarized by a two-component recursive state. Neither
representation stores the independent third coordinate. Memory and state
realization describe the same observable behavior in two different ways.

\begin{figure}[htbp]
\centering
\includegraphics[width=0.95\textwidth]{sources/monographs/large-signal-model/Reference_Laboratory/theory/../figures/01_delay_prediction.pdf}
\caption{Output-only identification in the declared rational example. Four observations determine the delay recurrence. Later outputs are calculated independently from the hidden system. The one-coordinate compression omits the return path. This is a mathematical example, not a task-suite score.}
\end{figure}

\section{A single propagation law for typed defects}\label{LSM:a-single-propagation-law-for-typed-defects}

A reconstruction error, an operator-fit error and a refinement error
have different origins. Once each has been translated into an error in
the state used by the next operation, their propagation can be analyzed
by the same identity.

Allow the representation to change with time. Let \(\gamma_t:X\to Z_t\)
be the declared representation at step \(t\), with a norm on each
\(Z_t\). Let \(x_{t+1}=F_t(x_t)\) be the true transition and let
\(\widehat f_t:Z_t\to Z_{t+1}\) be the implemented step. Its defect is

\[
r_t(x)=\gamma_{t+1}(F_t(x))-\widehat f_t(\gamma_t(x)).
\]

This definition already includes a representation switch. It is not
restricted to one fixed feature vector.

\Needspace{5\baselineskip}
\textbf{Proposition 18. Compositional defect propagation.} Suppose
\(\widehat f_t\) is \(\lambda_t\)-Lipschitz on a domain containing both
\(\gamma_t(x_t)\) and the predicted state \(\widehat z_t\). Suppose
\(\|r_t(x_t)\|\leq\delta_t\) and
\(\widehat z_{t+1}=\widehat f_t(\widehat z_t)\). Then, with
\(e_t=\|\gamma_t(x_t)-\widehat z_t\|\),

\[
\boxed{
e_T\leq
\left(\prod_{j=0}^{T-1}\lambda_j\right)e_0
+\sum_{k=0}^{T-1}
\left(\prod_{j=k+1}^{T-1}\lambda_j\right)\delta_k.
}
\]

Empty products are one. If the output \(h_t(x)\) has a decoder \(H_t\)
with Lipschitz constant \(\kappa_t\) and
\(\|h_t(x_t)-H_t(\gamma_t(x_t))\|\leq\eta_t\), then

\[
\|h_t(x_t)-H_t(\widehat z_t)\|
\leq\eta_t+\kappa_te_t.
\]

\textbf{Proof.} Add and subtract \(\widehat f_t(\gamma_t(x_t))\) to
obtain \(e_{t+1}\leq\delta_t+\lambda_te_t\). Repeated substitution
yields the product-sum formula. Add and subtract \(H_t(\gamma_t(x_t))\)
for the output inequality. \(\square\)

For a fixed linear representation \(R\) and an estimated transition
\(\widehat B_u\), the defect is \((RL_u-\widehat B_uR)x\). For a
reference reduced operator \(B_u\), it splits exactly as

\[
RL_u-\widehat B_uR
=(RL_u-B_uR)+(B_u-\widehat B_u)R.
\]

The first term is representational nonclosure. The second is
operator-estimation error. A fit statistic for the second cannot certify
that the first vanishes. Proposition 16 supplies a conditional bound on
the second in its declared fixed-space setting.

If \(\lambda_t\leq\rho<1\) and \(\delta_t\leq\delta\), then

\[
e_T\leq\rho^Te_0+\delta\frac{1-\rho^T}{1-\rho}.
\]

This is a genuinely recursive prediction bound. It is not merely a
one-step comparison made using the true history at every step. If the
Lipschitz constants exceed one, the same theorem remains valid but can
permit error growth.

\subsection{19.1 Stability under changing
operations}\label{LSM:stability-under-changing-operations}

For linear candidate operators, one sufficient way to obtain a common
contraction norm is to exhibit a positive-definite matrix \(M\) and
\(0\leq\rho<1\) with

\[
\widehat B_u^\top M\widehat B_u\preceq\rho^2M
\quad\text{for every declared action }u.
\]

Then \(\|z\|_M=(z^\top Mz)^{1/2}\) gives
\(\|\widehat B_uz\|_M\leq\rho\|z\|_M\) for all actions, directly from
the displayed inequality. Proposition 18 therefore applies under
arbitrary switching among them. This is the common quadratic Lyapunov
construction from control theory {[}R8{]}, used here as an explicit
certificate rather than inferred from a list of eigenvalues.

Individual spectral stability does not suffice. The matrices

\[
A=\begin{pmatrix}0&2\\0&0\end{pmatrix},\qquad
B=\begin{pmatrix}0&0\\2&0\end{pmatrix}
\]

are both nilpotent, but \(BA=\operatorname{diag}(0,4)\). Alternating
them amplifies an appropriate state by four every two steps. Even a
single stable nonnormal matrix can produce large transient amplification
in a specified norm.

If a learned policy chooses actions from predicted states while the
reference process uses actions chosen from true states, the two action
sequences may differ. Proposition 18 then needs a bound for that action
mismatch, or a Lipschitz bound for the full closed-loop map. The result
is not an automatic control-policy guarantee.

\section{Truncating memory without losing the prediction bound}\label{LSM:truncating-memory-without-losing-the-prediction-bound}

Section 8 bounded the omitted tail of the exact linear memory expression
using the same resolved history. We now close the separate recursion
argument.

Write the exact relation as

\[
z_{t+1}=Az_t+\sum_{k=0}^{t-1}K_kz_{t-1-k}+f_t,
\qquad K_k=BD^kC,\quad f_t=BD^tq_0.
\]

For negative indices set \(z_s=0\), solely as a convention for missing
prehistory in this initialized formula. A length-\(H\) approximate
model, with \(H\geq0\), is

\[
\widehat z_{t+1}=\widehat A\widehat z_t+
\sum_{k=0}^{H-1}\widehat K_k\widehat z_{t-1-k}.
\]

An empty sum is zero. Let \(d_t\) be the exact right-hand side minus
this truncated right-hand side evaluated on the \textbf{true} history.

\Needspace{5\baselineskip}
\textbf{Proposition 19. Uniform prediction bound for a stable
finite-memory truncation.} Suppose

\[
a=\|\widehat A\|+\sum_{k=0}^{H-1}\|\widehat K_k\|<1,
\qquad\sup_{t\geq0}\|d_t\|\leq\Delta.
\]

Let \(E_0\) bound the initialization errors at indices \(-H,\ldots,0\).
Then the recursively generated predictions satisfy

\[
\boxed{
\sup_{t\geq0}\|z_t-\widehat z_t\|
\leq\max\left(E_0,\frac{\Delta}{1-a}\right).
}
\]

If \(\|z_t\|\leq Z\) for \(t\geq0\) and \(\|D\|\leq\rho<1\), a valid
choice of defect bound is

\[
\begin{aligned}
\Delta={}&Z\left(\|A-\widehat A\|
+\sum_{k=0}^{H-1}\|K_k-\widehat K_k\|\right)\\
&+\|B\|\,\|C\|Z\frac{\rho^H}{1-\rho}
+\|B\|\,\|q_0\|.
\end{aligned}
\]

The last term can be removed if the exact initialized forcing \(f_t\) is
retained in the approximate recurrence.

\textbf{Proof.} Subtract the two recurrences and bound every retained
history error by the largest of the errors at indices \(t-H,\ldots,t\).
This gives \(e_{t+1}\leq a\max_{t-H\leq s\leq t}e_s+\Delta\). The ball
of radius \(E=\max(E_0,\Delta/(1-a))\) is invariant because
\(aE+\Delta\leq E\). Induction proves the bound. The estimate for
\(\Delta\) separates coefficient errors, the geometric kernel tail and
the initialized forcing, using Proposition 10. \(\square\)

This norm-sum condition is sufficient, not necessary. A stable
recurrence with cancellation can fail this test. Failure to find this
certificate is not proof of instability.

For the scalar example with \(q_0=0\), take exact retained coefficients
and \(z_0=1\). All outputs lie in \([0,1]\): the full nonnegative memory
coefficients sum to at most \(7/8\), so induction bounds every new
output by the previous maximum. Thus \(Z=1\) is valid. With

\[
A=1/2,\quad K_k=\frac14\left(\frac13\right)^k,
\]

one obtains

\[
a_H=\frac78-\frac38\,3^{-H},\qquad
\Delta_H=\frac38\,3^{-H},\qquad
\boxed{\sup_t|z_t-\widehat z_t|\leq
\frac{3\,3^{-H}}{1+3\,3^{-H}}.}
\]

The right-hand side is a theorem for every time in this initialized
example. The figure compares it with a separately computed
finite-horizon error. The computed maximum is not substituted for the
uniform bound.

\begin{figure}[htbp]
\centering
\includegraphics[width=0.95\textwidth]{sources/monographs/large-signal-model/Reference_Laboratory/theory/../figures/02_memory_certificate.pdf}
\caption{A memory-tail estimate becomes a recursive prediction certificate only after a stability margin is supplied. The measured curve is the maximum over 80 calculated steps. The upper curve is the proved bound for all steps under the stated scalar initialization.}
\end{figure}

\section{Nonlinear packet truncation with an explicit remainder}\label{LSM:nonlinear-packet-truncation-with-an-explicit-remainder}

Polynomial interactions were expanded exactly in Proposition 5. A
computational model may retain only some ordered interaction tuples. It
then needs a bound on what was omitted, rather than an assumption that
packetwise evaluation is exact.

Let \(P_i\) be an orthogonal resolution of a finite-dimensional real
Hilbert space. Suppose

\[
F(x)=b+\sum_{m=1}^{p}B_m(x,\ldots,x),
\]

with bounded multilinear maps. For each ordered tuple
\(\boldsymbol i=(i_1,\ldots,i_m)\), let \(\beta_{m,\boldsymbol i}\)
bound the norm of \(B_m\) restricted to
\(\operatorname{im}P_{i_1}\times\cdots\times\operatorname{im}P_{i_m}\).
Let \(\mathcal I_m\) be the retained tuple set, and let
\(F_{\mathcal I}\) retain the constant term and exactly those
interactions.

\Needspace{5\baselineskip}
\textbf{Proposition 20. Packet-interaction truncation bound.} Put
\(a_i=\|P_ix\|\). Then

\[
\boxed{
\|F(x)-F_{\mathcal I}(x)\|
\leq\sum_{m=1}^{p}\ \sum_{\boldsymbol i\notin\mathcal I_m}
\beta_{m,\boldsymbol i}\prod_{j=1}^{m}a_{i_j}.
}
\]

On the ball \(\|x\|\leq R\), a valid Lipschitz bound for
\(F_{\mathcal I}\) is

\[
\operatorname{Lip}(F_{\mathcal I})
\leq\sum_{m=1}^{p}\ \sum_{\boldsymbol i\in\mathcal I_m}
m\beta_{m,\boldsymbol i}R^{m-1}.
\]

\textbf{Proof.} Subtract the retained terms from the exact multilinear
expansion. Bound each omitted term by its restricted multilinear norm
and sum. For the Lipschitz estimate, telescope the difference between
two \(m\)-fold evaluations by replacing one argument at a time. Each
changed argument has norm at most \(\|x-y\|\) after projection, and each
unchanged argument has norm at most \(R\). There are \(m\) terms for
each retained tuple. \(\square\)

If the declared polynomial itself approximates a nonpolynomial map with
remainder at most \(\varepsilon_R\) on the ball, add \(\varepsilon_R\)
to the defect bound. Without such a remainder estimate, the result is
only a theorem for the polynomial, not for the unexpanded physical law.

For a retained state \(z=Px\) and omitted state \(q=(I-P)x\), dropping
every term that uses \(q\) gives the coarser majorant

\[
\|F(z+q)-F(z)\|
\leq\sum_{m=1}^{p}\|B_m\|
\left((\|z\|+\|q\|)^m-\|z\|^m\right).
\]

The tuple-level result is finer because it can distinguish a harmless
omitted packet from a strongly coupled one. Ordered tuples must be
counted with their multiplicities unless symmetry has explicitly
combined them.

A prediction guarantee follows by inserting these defect and Lipschitz
bounds into Proposition 18, provided the true and predicted trajectories
remain in the declared domain. That domain condition cannot be inferred
from a small initial error. One possible construction is to assume that
the true map preserves a closed convex set \(K\) and to project every
approximate output back onto \(K\). Euclidean projection is nonexpansive
and fixes the true output, so it does not increase the one-step
approximation bound. This is a declared modification of the predictor,
not an identity for an unmodified recurrence.

For the scalar quadratic \((u+v)^2\), retaining only \(u^2+v^2\) omits
\(2uv\). The packet bound is \(2|u||v|\) and is exact in magnitude.
Small energy in one packet matters in proportion to both its amplitude
and the cross-interaction coefficient.

\section{Learning several charts without inventing several worlds}\label{LSM:learning-several-charts-without-inventing-several-worlds}

The exact atlas identities constrain the meaning of a fitted model, but
independently trained predictors need not satisfy them. Here is a simple
coherent reconciliation with a precise objective.

Fix a finite-dimensional represented space \(X\). It may be the observed
signal space or a predictive quotient, but it must be the \emph{same}
space for all atlases in this section. Let \(A_a:X\to C_a\) and
\(S_a:C_a\to X\) satisfy \(S_aA_a=I\). Let \(G_a:C_a\to C_a\) be linear
fitted maps. Pull them back to object space:

\[
L_a=S_aG_aA_a.
\]

For positive weights summing to one, define

\[
\overline L=\sum_aw_aL_a,
\qquad\widetilde G_a=A_a\overline LS_a
\quad\text{on }\mathcal C_a=\operatorname{im}A_a.
\]

\Needspace{5\baselineskip}
\textbf{Proposition 21. Coherent reconciliation and coordinate-stable
transport.} The family \(\widetilde G_a\) is consistent and satisfies
every atlas intertwining identity. The operator \(\overline L\) is the
unique minimizer of

\[
\sum_aw_a\|L-L_a\|_F^2
\]

over linear maps on \(X\), where \(\|\cdot\|_F\) is the Frobenius norm.
For any target operator \(L_*\) and any chosen operator norm,

\[
\|\overline L-L_*\|\leq\sum_aw_a\|L_a-L_*\|.
\]

Further, the norm \(\|c\|_a=\|S_ac\|\) on \(\mathcal C_a\) makes every
exact transport \(T_{ba}=A_bS_a\) an isometry. A represented transition
\(A_bLS_a\) has exactly the same induced Lipschitz constant in these
norms as \(L\) has in object space.

\textbf{Proof.} Coherence follows from Proposition 2 because every chart
realizes the same \(\overline L\). Expanding the squares gives

\[
\sum_aw_a\|L-L_a\|_F^2
=\|L-\overline L\|_F^2+
\sum_aw_a\|L_a-\overline L\|_F^2.
\]

The weighted cross term vanishes, proving the unique minimizer. The
target bound follows from the triangle inequality. The pullback
expression is a norm because \(S_a\) is inverse to \(A_a\) on
\(\mathcal C_a\). Finally, \(\|T_{ba}c\|_b=\|S_bA_bS_ac\|=\|S_ac\|\).
The same identity for a transition, and the bijection between \(X\) and
\(\mathcal C_a\), establish equality of Lipschitz constants. \(\square\)

This reconciliation is ordinary weighted least squares in a common
semantic space. It does not prove that consensus is correct, reduce
statistical error under arbitrary dependence, or turn deterministic
recodings into independent evidence. The output must also be reanalyzed
into a realizable coefficient vector. An independently fitted redundant
output can have an additional inconsistency residual
\((I-A_aS_a)G_aA_ax\) before reconciliation.

For Euclidean coefficient errors, conditioning remains important. If
\(A_a\) has smallest singular value at least \(m_a>0\) and largest
singular value at most \(M_a\), its canonical synthesis \(A_a^\dagger\)
has norm at most \(1/m_a\), and \(\|T_{ba}\|\leq M_b/m_a\). An arbitrary
left inverse can behave much worse away from the consistent coefficient
space, so the canonical-synthesis qualification matters.

A learned atlas can therefore be parameterized within prescribed
singular-value bounds, with synthesis tied to its inverse or
pseudoinverse. This gives a construction contract. It is not a claim
that every unconstrained training procedure respects those bounds.

When atlases change over time but represent the same state, using the
pullback norms avoids multiplying arbitrary chart condition numbers at
every step. Exact transport telescopes to the direct transport.
Numerical synthesis errors and genuinely inconsistent predictions remain
defects and must still be bounded.

\section{Refinement cannot manufacture missing information}\label{LSM:refinement-cannot-manufacture-missing-information}

An atlas change is an invertible recoding of the same represented
information. Predictive refinement usually adds information. Treating
these as identical operations would undo the central distinction of the
foundation.

\Needspace{5\baselineskip}
\textbf{Proposition 22. The upward-lift obstruction.} Let
\(\gamma:X\to Z\) be a coarse representation and \(\eta:X\to W\) a
proposed refined representation. There is a deterministic map
\(J:\gamma(X)\to\eta(X)\) satisfying \(\eta=J\gamma\) if and only if

\[
\gamma(x)=\gamma(x')\ \Longrightarrow\ \eta(x)=\eta(x')
\quad\text{for every }x,x'.
\]

For linear representations, this is equivalent to
\(\ker\gamma\subseteq\ker\eta\). If two states share a coarse
representation but their refined representations differ, any predicted
refined value \(w\) has worst-case error at least

\[
\max\bigl(\|w-\eta(x)\|,\|w-\eta(x')\|\bigr)
\geq\tfrac12\|\eta(x)-\eta(x')\|.
\]

\textbf{Proof.} A lift must be constant on each coarse fiber.
Conversely, the stated condition defines it unambiguously on every
fiber. The linear version is the kernel factorization criterion. The
lower bound follows from the triangle inequality between the two refined
values through \(w\). \(\square\)

For a strict refinement, the coarse state is recoverable from the
refined one but not conversely. In the linear case this typically means
\(\ker\eta\subsetneq\ker\gamma\), exactly the opposite of the condition
required for an upward lift.

For example, \(\gamma(x_1,x_2)=x_1\) and \(\eta(x_1,x_2)=(x_1,x_2)\).
The states \((0,1)\) and \((0,-1)\) have the same coarse state, but any
common point initialization has worst-case error at least one in
Euclidean norm. Padding a vector with zeros is a choice, not information
recovery.

A refinement step therefore needs one of three explicitly distinguished
inputs: an archive that already contains the relevant observations, new
observations, or a set-valued or probabilistic state retaining
uncertainty about the missing details. An exact archive of measured
outputs can initialize a longer history when those samples were stored.
It cannot recover a latent coordinate that never affected any
measurement.

For adaptive prediction, an attempted lift \(J_t\) contributes its own
initialization defect

\[
\eta_t(x_t)-J_t\gamma_t(x_t).
\]

It belongs in Proposition 18's initial error or in the defect of the
composed transition. It cannot be set to zero merely because a finer
packet tree has been allocated.

A data-dependent representation schedule also requires care. The
pathwise bound applies to the realized schedule when all of its
hypotheses hold along that schedule. Comparing two different schedules
chosen by true and predicted histories introduces a selection
discrepancy. No uniform statistical consistency of adaptive selection
follows solely from the deterministic propagation law.

\section{A growing operator grammar with coherent evidence}\label{LSM:a-growing-operator-grammar-with-coherent-evidence}

A large algebraic signal model should be able to add operator programs
without changing the meaning of its existing evidence. This requires
more than assigning a new softmax weight when a candidate is introduced.

Let \(\mathcal G\) be a countable declared program set, with one fixed
probability mass function \(\pi(g)>0\) satisfying \(\sum_g\pi(g)=1\).
Every program specifies a likelihood \(L_g(D)\) for the \textbf{same}
observed data event or density relative to a common base measure. Assume

\[
0<Z(D)=\sum_{g\in\mathcal G}\pi(g)L_g(D)<\infty.
\]

Let finite active sets \(\mathcal G_n\) increase to \(\mathcal G\).
Their unnormalized evidence is
\(Z_n=\sum_{g\in\mathcal G_n}\pi(g)L_g(D)\). Only stages with \(Z_n>0\)
are considered.

\Needspace{5\baselineskip}
\textbf{Proposition 23. Exact accounting for an enumerated grammar.}
Define the full posterior \(P(g\mid D)=\pi(g)L_g(D)/Z\) and extend the
active posterior \(P_n(g\mid D)=\pi(g)L_g(D)/Z_n\) by zero outside
\(\mathcal G_n\). Then

\[
\boxed{
\|P-P_n\|_{\mathrm{TV}}=\tau_n,
\qquad \tau_n=1-\frac{Z_n}{Z}.
}
\]

For any common family of candidate predictive probability kernels, the
corresponding predictive mixtures differ in total variation by at most
\(\tau_n\). If the omitted unnormalized evidence has a bound

\[
\sum_{g\notin\mathcal G_n}\pi(g)L_g(D)\leq b_n,
\]

then

\[
\tau_n\leq\frac{b_n}{Z_n+b_n}.
\]

\textbf{Proof.} On the active set, \(P=(Z_n/Z)P_n\), so the total mass
of the difference there is \(\tau_n\). The mass outside is also
\(\tau_n\). Half the sum of absolute differences is therefore
\(\tau_n\). Applying the same probability kernel to two mixing measures
cannot increase total variation: integrate any event probability, a
function taking values in \([0,1]\), against their difference. For the
final bound, write \(Z=Z_n+W_n\) with \(0\leq W_n\leq b_n\) and use that
\(W_n/(Z_n+W_n)\) increases with \(W_n\). \(\square\)

Since the active sets exhaust the grammar and the evidence sum is
finite, \(Z_n\uparrow Z\) and \(\tau_n\downarrow0\). This is posterior
approximation \emph{inside the declared grammar}. It does not establish
that the real signal-generating mechanism belongs to it.

For a discrete data sequence, a probability likelihood is at most one,
so the omitted prior mass is a valid \(b_n\). A continuous-data density
need not be at most one. It requires a density envelope or another
justified bound on omitted evidence. A small unnormalized tail can also
carry substantial posterior mass when the active evidence is very small.

The fixed prior need not be uniform. For an explicit enumeration
\(g_1,g_2,\ldots\), \(\pi(g_k)=2^{-k}\) is a normalized example, with
omitted prior mass \(2^{-n}\) after \(n\) programs. It expresses a
particular enumeration bias, not a canonical prior on mathematical
meaning.

A program outside the enumerated prefix can be accounted for only when
its prior and likelihood are defined or bounded. A genuinely
outside-grammar component needs its own model. Assigning an unexplained
residual probability is not a Bayesian treatment of missing mechanisms.

\subsection{24.1 Refining a program without double counting
it}\label{LSM:refining-a-program-without-double-counting-it}

A representation-rich language can contain many syntactic versions of
one mechanism. It can also split one broad mechanism into more specific
children. These operations must conserve prior mass when they are
intended as a re-expression of the same model.

\Needspace{5\baselineskip}
\textbf{Proposition 24. Prior-preserving program refinement.} Suppose a
parent program \(g\) is replaced by children \(g_j\) with weights
\(\alpha_j\geq0\), \(\sum_j\alpha_j=1\), and suppose its likelihood is
the mixture

\[
L_g(D)=\sum_j\alpha_jL_{g_j}(D).
\]

Give child \(j\) prior mass \(\pi(g)\alpha_j\) and leave all other
programs unchanged. Then total evidence is unchanged and the sum of the
children's posterior masses equals the former parent's posterior mass.
Posterior predictive distributions are also unchanged when the same
mixture relation holds for the joint likelihood of \((D,y)\) for every
predicted event \(y\).

\textbf{Proof.} The children's total unnormalized evidence is
\(\sum_j\pi(g)\alpha_jL_{g_j}(D)=\pi(g)L_g(D)\). Thus both the total
normalizer and every other program's posterior are unchanged. The same
sum for joint evidence gives the predictive statement after
normalization by the unchanged evidence of \(D\). \(\square\)

The mixture identity must hold for the \textbf{whole data record}.
Averaging one-step likelihoods independently at each time generally
describes a switching mechanism, not a fixed child selected once for the
whole sequence.

If the children are only exact atlas recodings of the same operator and
data model, their likelihoods are identical. Splitting the original
prior mass leaves every prediction unchanged. Giving each duplicate a
fresh full prior mass instead changes the model's prior preference. That
may be an intentional new prior, but it is not invariance under
representation.

The implementation's existing held-out predictive weights, described in
Section 10, are not replaced by these identities merely by calling them
posterior weights. The proposition supplies the probabilistic semantics
that a future grammar-weighting implementation would need to enforce.

\section{The finite learned-realization theorem}\label{LSM:the-finite-learned-realization-theorem}

The preceding results combine into a concrete completion of one part of
the learning problem. This is the integrated identification claim for
the autonomous linear sector.

\Needspace{5\baselineskip}
\textbf{Proposition 25. Observationally grounded finite linear
realization.} Assume all of the following:

\begin{enumerate}
\def\labelenumi{\arabic{enumi}.}
\LSMtightlist
\item
  The data come from an autonomous finite-dimensional linear process
  with declared outputs \(y_t=Cx_t\).
\item
  A delay length \(h\) satisfies
  \(\operatorname{rank}O_{h+1}=\operatorname{rank}O_h\).
\item
  Exact observed input windows in \(Z_-\) span
  \(V_h=\operatorname{im}O_h\), and \(Z_+\) records their true one-step
  successor windows.
\end{enumerate}

Then an algorithm using only \((Z_-,Z_+)\) recovers the unique
transition on \(V_h\). Initialized with a fully observed length-\(h\)
window, it predicts every later output exactly. Any lossless atlas of
\(V_h\), any complete packet decomposition of that space, and every
coherently transported realization of the learned transition preserve
the same predictions. Hidden variables outside the observable quotient
need not be recovered.

If the fixed-space assumptions of Proposition 16 replace exact
observations, that proposition bounds operator error. Whenever the
resulting transition defects and Lipschitz constants satisfy Proposition
18, its product-sum bound controls the output prediction error. A common
contraction margin yields the stated uniform bound.

\textbf{Proof.} Proposition 14 gives the unique true transition \(D_h\)
on \(V_h\). Because the observed windows span \(V_h\), Proposition 15
identifies exactly its restriction, which is the entire required
operator. Equality of the learned and true states follows by induction
after initialization. The current output is the last block of the
history state. Propositions 1--4 preserve state reconstruction,
transport and complete decomposition, so they preserve that output
sequence. The perturbation statement is the direct composition of
Propositions 16 and 18 with the declared output decoder. \(\square\)

This establishes an actual route from observations to a composable
algebraic signal model under explicit finite assumptions. It is more
than a fit on a few points. It is also less than unrestricted nonlinear
world-model discovery. In particular, the algorithm does not infer
assumptions 1--3 merely by returning a matrix. The accompanying
certificate keeps the data-verified identities separate from those model
assumptions.

\section{Connection to controlled prediction}\label{LSM:connection-to-controlled-prediction}

Proposition 25 completes the autonomous observable linear sector under
its declared coverage assumptions. It does not establish a common delay
operator for arbitrary actions. Chapter~\ref{LSM:book:chapter-33} constructs controlled
predictive states from interventional experiment probabilities. The two
results are compatible routes, not interchangeable learning assumptions.


\label{LSM:part-iii.-controlled-stochastic-realization}

The next step is not to rename a deterministic predictor as a
probabilistic one. A stochastic signal model must say what its state
predicts, which actions intervene on the process, how observations
update the state, and why every resulting number is a probability. These
are separate obligations. A low dimensional linear representation can
satisfy its reconstruction equations while failing probability
positivity.

Chapter~\ref{LSM:book:chapter-32} established observable learning for an autonomous linear signal.
Here the observable objects are probabilities of finite action and
observation experiments. We derive a finite realization from them, show
when it is minimal, give an explicit probability-valid state domain, and
separate exact reconstruction from noisy estimation. The final
construction connects back to Komorebi's packets through a normalized
closure-defect identity.

The classical antecedents are predictive state representations,
observable operator models and finite-rank realization {[}R5, R7, R9{]}.
Those traditions already use probabilities of tests and products of
operators. The purpose here is to make their assumptions and
consequences explicit within the realization calculus. The twelve
propositions below are a connected development, not twelve claims of
historical novelty. No statement here replaces a proof with a numerical
check.


\chapter{Controlled stochastic realization}\label{LSM:book:chapter-33}
\section{Actions and observations are different inputs}\label{LSM:actions-and-observations-are-different-inputs}

Let \(\mathcal A\) and \(\mathcal O\) be finite nonempty action and
observation sets. A symbol \(s=(a,o)\) records an action followed by an
observation. A word \(w=s_1\cdots s_k\) has temporal order from left to
right. The empty word is \(\epsilon\).

Write

\[
f(w)=\Pr(o_1,\ldots,o_k\mid \operatorname{do}(a_1,\ldots,a_k)).
\]

The process is reset to one specified initial distribution before this
experiment. The actions in this definition are fixed externally. This is
not, in general, the distribution obtained by conditioning a passive
record on an adaptively selected sequence of actions. We assume the
consistency laws

\[
f(\epsilon)=1,\qquad f(w)\geq0,\qquad
\sum_{o\in\mathcal O}f(w(a,o))=f(w)
\quad\text{for every }w,a.
\]

These laws specify causal probabilities on finite observation trees. A
fixed history-dependent policy can be evaluated by selecting the action
on each branch and multiplying the appropriate conditional
probabilities. This does not establish optimality of that policy or its
identification from insufficient data.

\subsection{27.1 A finite hidden-state
construction}\label{LSM:a-finite-hidden-state-construction}

A concrete realization has a belief \(p\in\mathbb R^n_+\) with
\(\mathbf1^\top p=1\). For each \(s=(a,o)\), let
\(M_s\in\mathbb R^{n\times n}\) be entrywise nonnegative, with

\[
\mathbf1^\top\sum_oM_{a,o}=\mathbf1^\top.
\]

The matrix entry \(M_{a,o}(j,i)\) is the probability of next hidden
state \(j\) and observation \(o\), given current hidden state \(i\) and
chosen action \(a\). This is a discrete probability instrument. It
includes an action-specific transition followed by an observation, but
does not require that factorization.

Put \(M_w=M_{s_k}\cdots M_{s_1}\). Then

\[
f(w)=\mathbf1^\top M_wp_0,\qquad
p^+=\frac{M_sp}{\mathbf1^\top M_sp}
\]

on branches of positive probability. An observation with zero
probability has no conditional posterior under that model. It must not
be handled by silently dividing by an arbitrary numerical floor.

\Needspace{5\baselineskip}
\textbf{Proposition 26. Finite instrument closure.} Starting with
\(W_0=\operatorname{span}\{\mathbf1^\top\}\), define

\[
W_{j+1}=W_j+\sum_{s\in\mathcal A\times\mathcal O}W_jM_s.
\]

This stabilizes after at most \(n-1\) strict enlargements at

\[
W_* = \operatorname{span}\{\mathbf1^\top M_w:w\text{ a finite word}\}.
\]

Choose a full row rank matrix \(R\) with row space \(W_*\) and first row
\(\mathbf1^\top\). There are unique matrices \(B_s\) with

\[
RM_s=B_sR,\qquad \lambda=e_1^\top,\qquad \lambda R=\mathbf1^\top.
\]

For \(z=Rp\), both the next-observation likelihood and the posterior
predictive state close:

\[
\Pr(o\mid z,a)=\lambda B_{a,o}z,\qquad
z^+=\frac{B_{a,o}z}{\lambda B_{a,o}z}.
\]

Every linear summary determining all word probabilities by linear
readouts for every initial belief has row space containing \(W_*\).

\textbf{Proof.} The nested-space argument from Proposition 8 applies to
the finite family \(M_s\). Multiplication of a test row by \(M_s\)
prepends the symbol \(s\) to that test, so the stable span is exactly
the displayed word span. Invariance supplies the matrices \(B_s\), and
full row rank makes them unique. Apply \(R\) to the unnormalized and
normalized hidden-state updates. If all test probabilities are linear
functions of another summary, then each row \(\mathbf1^\top M_w\) lies
in its row space. Equality on all beliefs implies equality on their
linear span, which is \(\mathbb R^n\). \(\square\)

This proposition uses a known hidden-state model. It is not yet an
observable-data learner. It also concerns all initial beliefs, whereas a
model initialized at one fixed \(p_0\) can have a smaller reachable
predictive dimension. The next construction handles that distinction
without requiring recovery of hidden states.

\section{The state is a function of possible future experiments}\label{LSM:the-state-is-a-function-of-possible-future-experiments}

For a history \(h\), define its unnormalized future-behavior column

\[
v_h(t)=f(ht),
\]

where \(t\) ranges over every finite test word. Form the infinite
history/test matrix

\[
\mathscr H_{t,h}=f(ht).
\]

Its rank means the dimension of the span of its columns, equivalently
the supremum of the ranks of finite submatrices. The columns are
functions, not vectors obtained by physically observing every future.
Their finite coordinate realization becomes learnable only when enough
experiment probabilities are available.

\Needspace{5\baselineskip}
\textbf{Proposition 27. Minimal finite behavior realization.} If
\(\operatorname{rank}\mathscr H=r<\infty\), the vector space
\(V=\operatorname{span}\{v_h\}\) has dimension \(r\). There are
well-defined linear maps \(B_s:V\to V\), an initial vector
\(v_\epsilon\), and a linear mass functional \(\lambda(v)=v(\epsilon)\)
satisfying

\[
B_sv_h=v_{hs},\qquad
f(s_1\cdots s_k)=\lambda B_{s_k}\cdots B_{s_1}v_\epsilon.
\]

Every linear realization of \(f\) has dimension at least \(r\). Any two
reachable and observable realizations of dimension \(r\) differ by an
invertible change of coordinates. For \(f(h)>0\), the normalized state
\(v_h/f(h)\) determines every future test probability conditional on
\(h\) and the declared future interventions.

\textbf{Proof.} If \(\sum_hc_hv_h=0\), then for every test \(t\),

\[
\sum_hc_hv_{hs}(t)=\sum_hc_hf(hst)=0,
\]

because \(st\) is another test. Thus shifting columns preserves all
linear relations and defines \(B_s\). The probability formula follows by
induction. For any \(d\)-dimensional realization, the history/test
matrix factors into test rows and reachable state columns, so its rank
is at most \(d\).

For uniqueness, map the reachable vector after each history in one
minimal realization to the corresponding vector in the other. If a
linear combination vanishes in the first, every future test evaluates to
zero on it. Those same evaluations vanish in the second, where
observability makes the combination zero. The map is therefore
well-defined and bijective, and intertwines every symbol operator and
both boundary vectors. Finally \(f(ht)/f(h)\) is exactly the conditional
experiment probability. \(\square\)

The coordinate dimension \(r\) includes probability mass. Normalized
states lie in the affine hyperplane \(\lambda z=1\), so their affine
dimension is at most \(r-1\). A finite hidden-state model of size \(n\)
has behavior rank at most \(n\). The theorem itself requires finite
behavior rank, not a preferred hidden-state ontology.

Two realized histories with equal normalized future columns are
equivalent for every declared future test. This is a predictive
equivalence, not equality of their physical pasts or lossless archives.

\section{Recovering the operators from an observable probability table}\label{LSM:recovering-the-operators-from-an-observable-probability-table}

Select \(r\) tests \(t_1,\ldots,t_r\) and \(r\) histories
\(h_1,\ldots,h_r\), with \(t_1=\epsilon\), for which the square matrix

\[
H_{ij}=f(h_jt_i)
\]

is invertible. Such a choice exists under the rank hypothesis of
Proposition 27. The nonzero empty-test row can be extended to a row
basis. Select independent history columns for that basis.

For each symbol \(s\), also obtain the shifted table and initial test
vector

\[
(G_s)_{ij}=f(h_jst_i),\qquad (b_0)_i=f(t_i).
\]

\Needspace{5\baselineskip}
\textbf{Proposition 28. Observable table reconstruction.} Assume the
full behavior rank is at most \(r\) and the displayed \(r\times r\)
matrix \(H\) is invertible. Then the full rank is exactly \(r\), and

\[
\boxed{B_s=G_sH^{-1},\qquad \lambda=e_1^\top,\qquad b_0=(f(t_i))_i}
\]

is a minimal realization of every word probability. Its state after
history \(h\) is

\[
b(h)=(f(ht_i))_i=B_hb_0.
\]

The normalized state is \(b(h)/f(h)\) when \(f(h)>0\). Its update uses
the formula in Proposition 26, with no hidden-state coordinates supplied
to the reconstruction algorithm.

\textbf{Proof.} Invertibility gives rank at least \(r\), so the assumed
upper bound is attained. Evaluation on the selected tests is an
isomorphism \(V\to\mathbb R^r\), because it maps the chosen basis of
history columns to the invertible matrix \(H\). The shifted images of
those basis columns are precisely the columns of \(G_s\), so
\(B_sH=G_s\). The first selected test evaluates the mass. Initial
evaluation gives \(b_0\). Proposition 27 supplies correctness for every
history, not only the selected columns. \(\square\)

There are two ways to justify coverage in this statement. One can know
that the selected row and column families span the true behavior space.
Alternatively, a justified upper bound \(r\) on the process rank,
together with the invertible \(r\times r\) block, proves that coverage.
The observed invertibility by itself only gives a lower bound. An
unknown higher-rank process can agree on a finite table and differ
later.

Exact experiment probabilities are population quantities. The reference
example below supplies rational tables from a declared process and keeps
that table generator separate from the learner. Section 36 supplies one
finite sampling protocol for estimating those entries. Neither a hidden
matrix nor its probability cone is secretly passed to the observable
learner.

\section{Probability validity is not a rank condition}\label{LSM:probability-validity-is-not-a-rank-condition}

A real matrix product can yield a negative number. Therefore a fitted
factorization is not a stochastic model merely because it is low rank or
exactly interpolates a probability table. Spectral estimators can
encounter negative-probability problems in practice {[}R10{]}. We
address the structural issue directly.

In the canonical behavior space \(V\), define

\[
\mathcal C=\{v\in V:v(t)\geq0\text{ for every finite test }t\}.
\]

\Needspace{5\baselineskip}
\textbf{Proposition 29. Canonical probability cone.} The set
\(\mathcal C\) is a closed convex pointed cone that spans \(V\). Every
\(B_s\) preserves it. The mass functional is strictly positive on its
nonzero elements, and

\[
\sum_o\lambda B_{a,o}=\lambda.
\]

Consequently the normalized base

\[
\mathcal K=\{v\in\mathcal C:\lambda v=1\}
\]

is a compact convex probability-valid state space. Every
positive-probability branch maps \(\mathcal K\) into itself by
normalized updating. In the core-test coordinates of Proposition 28,
\(\mathcal K\subseteq\{1\}\times[0,1]^{r-1}\).

\textbf{Proof.} Each test evaluation is a continuous linear functional
on finite-dimensional \(V\), so the cone is closed and convex. If both
\(v\) and \(-v\) belong to it, every test value is zero and \(v=0\). It
spans \(V\) because it contains every column \(v_h\). Also
\((B_sv)(t)=v(st)\), proving invariance.

The consistency law \(\sum_ov(t(a,o))=v(t)\) holds for each history
column and hence throughout \(V\). Iterating it along any fixed action
schedule expresses \(v(\epsilon)\) as a sum of nonnegative values for
all observation words on that schedule. If \(v(\epsilon)=0\), each
summand is zero, for every schedule and length, so \(v=0\). The same
sums show \(0\leq v(t)\leq1\) when \(\lambda v=1\). They also give mass
conservation. The base is closed and bounded in the finite core
coordinates and thus compact. Invariance and positive normalization
prove the posterior assertion. \(\square\)

This cone is defined by all finite tests. Its existence is not an
algorithm for finding a finite list of inequalities. A finite list gives
an outer approximation and can miss a longer-test violation. The cone
also need not equal the convex hull of the states reachable from the
fixed initial distribution.

A practical sufficient certificate can instead use a known closed cone
\(C_0\) with nonempty normalized base, a strictly positive mass
functional, \(b_0\in C_0\), \(B_sC_0\subseteq C_0\), and mass
conservation. A nonnegative hidden-state realization supplies one after
a coordinate change. Such a certificate may use information unavailable
to an observable-only learner and must be labeled accordingly.

Negative entries of \(B_s\) do not, by themselves, prove invalidity.
Coordinates can be signed while the represented state remains in the
correct cone. Conversely, clipping negative coefficients can destroy the
operator identities. Validity belongs to the pair consisting of the
operators and their admissible cone, not to a coordinatewise heuristic.

\section{Packets can reconstruct a belief without being beliefs}\label{LSM:packets-can-reconstruct-a-belief-without-being-beliefs}

Let \(P_1,\ldots,P_m\) be a complete family of linear packet maps on a
predictive space, with \(\sum_iP_i=I\). Algebraic reconstruction gives
\(z=\sum_iP_iz\). It does not imply that \(P_iz\) is an unnormalized
probability state.

\Needspace{5\baselineskip}
\textbf{Proposition 30. Positive packet decomposition and its boundary.}
Suppose a cone \(C\) has a strictly positive mass functional
\(\lambda\). If every packet map preserves \(C\), then each normalized
\(z\in C\) decomposes as a genuine mixture

\[
z=\sum_{i:w_i>0} w_i z_i,\qquad
w_i=\lambda P_iz\geq0,\qquad
z_i=P_iz/w_i,\qquad\sum_iw_i=1.
\]

A zero weight has \(P_iz=0\). In the standard cone \(\mathbb R^n_+\),
however, a complete family of cone-preserving idempotents must consist
of diagonal coordinate-selection matrices. In particular, a nontrivial
Fourier or Haar packet family cannot generally be reinterpreted as a
mixture of probability states merely because it is orthogonal and
complete.

\textbf{Proof.} Cone preservation and strict positivity give the weights
and the zero-weight claim. Apply \(\lambda\) to the complete
decomposition to obtain their sum. Normalization of each nonzero term
gives the mixture.

On \(\mathbb R^n_+\), cone preservation is entrywise nonnegativity.
Since the matrices sum to \(I\), every off-diagonal entry in every
matrix is zero. Idempotence forces each diagonal entry to be zero or
one, and completeness makes the selections disjoint and exhaustive.
\(\square\)

For a two-state belief \(p=(3/4,1/4)^\top\), the Haar detail projector

\[
P_{\mathrm d}=\frac12\begin{pmatrix}1&-1\\-1&1\end{pmatrix}
\]

produces \((1/4,-1/4)^\top\). It is a valid signed detail packet, but
not a subprobability distribution. Its total mass is zero despite being
nonzero, which cannot happen inside a cone with a strictly positive mass
functional.

This does not prohibit useful spectral packets. It fixes their type.
Signed algebraic packets can be used for transfer, approximation and
diagnostics. They must be synthesized before probability interpretation
unless a separate positive decomposition has been established. That
distinction prevents a reconstruction identity from being misused as a
Bayesian mixture identity.

\section{Action coverage is a mathematical requirement}\label{LSM:action-coverage-is-a-mathematical-requirement}

\Needspace{5\baselineskip}
\textbf{Proposition 31. Missing-intervention nonidentifiability.} There
are finite controlled stochastic models that agree on every possible
observation record generated using only action \(a\), but disagree under
another action \(b\). Therefore even infinitely many passive records
confined to \(a\) cannot identify the unrestricted controlled model.

\textbf{Proof.} Use one hidden state and two observations. In both
models, action \(a\) gives probabilities \((1/2,1/2)\). Under \(b\), let
the first model give \((1,0)\) and the second \((0,1)\). Every
length-\(k\) record under \(a\) has probability \(2^{-k}\) in both
models. Their first observation under \(b\) differs with probability
one. \(\square\)

This is a coverage obstruction, not an optimization problem. It remains
even when the hidden dimension is known exactly. Richer transforms of
the same action-\(a\) observations do not create data about \(b\).

Action selection based on past observations requires a further
distinction. Under a known policy \(\pi\), an action/observation word
has joint probability

\[
\Pr_\pi(w)=f(w)\prod_{j=1}^{|w|}\pi(a_j\mid h_{j-1}),
\]

provided the action randomization has no access to hidden information
beyond the declared history. The policy factors depend on the observed
branch. Conditioning on a selected action sequence generally reweights
branches rather than removing these factors. Recovering intervention
probabilities from such data requires an appropriate coverage and
reweighting argument. The reset experiments in Section 36 avoid that
additional inference problem.

\section{Learning the coordinate space under bounded noise}\label{LSM:learning-the-coordinate-space-under-bounded-noise}

Chapter~\ref{LSM:book:chapter-32} bounded operator error after a state space and basis were fixed.
Here we also bound one step of data-dependent space selection. This
remains a finite-matrix statement. It does not certify that a chosen
finite history/test table captures the full process.

Let \(X\in\mathbb R^{m\times n}\) have rank \(r\geq1\), let
\(s=\sigma_r(X)>0\), and observe \(Y=X+E\) with
\(\|E\|_2\leq\varepsilon<s/2\). Let \(P\) project onto
\(\operatorname{im}X\), and let \(\widehat P\) project onto the leading
\(r\) left singular vectors of \(Y\).

\Needspace{5\baselineskip}
\textbf{Proposition 32. Finite-table rank separation and subspace
stability.} One has

\[
\sigma_r(Y)\geq s-\varepsilon,\qquad
\sigma_{r+1}(Y)\leq\varepsilon,\qquad
\|P-\widehat P\|_2\leq\frac{2\varepsilon}{s}<1.
\]

The \((r+1)\)st singular value is interpreted as zero when \(r\) is the
maximum possible rank. If \(U\) and \(\widehat U\) are orthonormal bases
of the two ranges, then

\[
\sigma_{\min}(\widehat U^\top U)
\geq\sqrt{1-(2\varepsilon/s)^2}.
\]

Thus \(x\mapsto\widehat U^\top x\) is injective on the true column
space. Its inverse on that space is

\[
y\longmapsto U(\widehat U^\top U)^{-1}y,
\]

with norm at most \([1-(2\varepsilon/s)^2]^{-1/2}\).

\textbf{Proof.} The min-max characterization of singular values and the
triangle inequality imply that adding \(E\) changes each singular value
by at most \(\varepsilon\). The leading-\(r\) singular subspace gives
\(\|(I-\widehat P)Y\|_2=\sigma_{r+1}(Y)\leq\varepsilon\). Hence

\[
\|(I-\widehat P)X\|_2\leq2\varepsilon,
\qquad
\|(I-\widehat P)P\|_2
=\|(I-\widehat P)XX^\dagger\|_2\leq2\varepsilon/s.
\]

For equal-rank orthogonal projections, the latter norm equals
\(\|P-\widehat P\|_2\). To see the equality, diagonalize
\(U^\top\widehat U\) by singular value decomposition. Each nontrivial
paired principal plane contributes a sine to both norms, while common
directions contribute zero. On the same planes the singular values of
\(\widehat U^\top U\) are the corresponding cosines. This proves the
last bound and the inverse formula. \(\square\)

A threshold \(\tau\) with \(\varepsilon<\tau<s-\varepsilon\) recovers
the rank of this finite matrix. The unknown true gap must be justified
or replaced by a valid confidence interval. Rank thresholds do not
exclude modes absent from the selected tests or unexcited by the
selected histories.

The inverse formula contains the unknown true basis \(U\). It proves
identifiability and conditioning relative to that space, not that a
learner already knows a perfect decoder. When using the estimated
subspace as coordinates for a realization, all boundary vectors, symbol
operators and probability constraints must be transformed together.
Signed singular-vector coordinates are not new independent observations.

\section{Noisy operator reconstruction in a fixed core}\label{LSM:noisy-operator-reconstruction-in-a-fixed-core}

After choosing a core block, the square reconstruction formula yields a
particularly direct error identity. It is useful even when the core was
chosen on a separate discovery data set.

Let \(G_s=B_sH\) with \(s_H=\sigma_{\min}(H)>0\). Observe

\[
\widetilde H=H+E,\qquad
\widetilde G_s=G_s+F_s,\qquad
\|E\|_2<s_H.
\]

\Needspace{5\baselineskip}
\textbf{Proposition 33. Noisy core-block realization.} The estimate
\(\widehat B_s=\widetilde G_s\widetilde H^{-1}\) satisfies

\[
\widehat B_s-B_s=(F_s-B_sE)\widetilde H^{-1},
\]

and therefore

\[
\boxed{
\|\widehat B_s-B_s\|_2
\leq\frac{\|F_s\|_2+\|B_s\|_2\|E\|_2}{s_H-\|E\|_2}.
}
\]

A bound using only the fitted operator and the same perturbation budgets
is

\[
\|\widehat B_s-B_s\|_2
\leq
\frac{\|F_s\|_2+\|\widehat B_s\|_2\|E\|_2}{s_H}.
\]

If \(\widetilde s=\sigma_{\min}(\widetilde H)\) is known and
\(\|E\|_2\leq e<\widetilde s\), the last denominator may be replaced by
\(\widetilde s-e\).

\textbf{Proof.} Since \(\widetilde H\) is invertible, subtract \(B_s\)
from its fitted expression and collect the terms. The least singular
value of \(\widetilde H\) is at least \(s_H-\|E\|_2\). For the second
identity, use

\[
(\widehat B_s-B_s)H=F_s-\widehat B_sE
\]

and multiply by \(H^{-1}\). Finally \(s_H\geq\widetilde s-e\) by
singular-value perturbation. \(\square\)

When the first test is empty, the mass row \(\lambda=e_1^\top\) is fixed
rather than estimated. Set the first component of the initial vector to
one as required by \(f(\epsilon)=1\). Errors in other initial test
probabilities remain initial-state error. This convention removes an
avoidable normalizer-estimation ambiguity but does not enforce cone
preservation or mass conservation for the fitted symbol family.

Numerical implementation should solve the linear systems rather than
explicitly form a floating-point inverse. The theorem describes the
exact algebraic estimator. Finite-precision solve error is an additional
residual, which the reference script measures separately.

\section{Conditioning a prediction can amplify its error}\label{LSM:conditioning-a-prediction-can-amplify-its-error}

Let \(C\) be a probability cone with compact normalized base \(K\) and
mass row \(\lambda\). Work in a declared norm and its dual. Suppose true
normalized states have norm at most \(R\). Let a true branch have

\[
u=B_sz,\qquad p=\lambda u\geq\eta>0,\qquad z^+=u/p\in K.
\]

For a predicted state \(\widehat z\), form
\(\widehat u=\widehat B_s\widehat z\), \(\widehat p=\lambda\widehat u\),
and define \(d=\widehat u-u\). Write \(e=\|\widehat z-z\|\),
\(\|\widehat B_s\|\leq L\), and \(\|\widehat B_s-B_s\|\leq\epsilon_B\).
Then

\[
\|d\|\leq Le+\epsilon_BR=:D.
\]

\Needspace{5\baselineskip}
\textbf{Proposition 34. Normalization-aware filtering stability.} If
\(\|\lambda\|_*D<\eta\), then \(\widehat p>0\) and

\[
\widehat z^+-z^+
=\frac{(I-z^+\lambda)d}{\widehat p},
\]

so

\[
\boxed{
\|\widehat z^+-z^+\|
\leq\frac{(1+R\|\lambda\|_*)D}{\eta-\|\lambda\|_*D}.
}
\]

Alternatively, if each realized fitted denominator is checked to satisfy
\(\widehat p_t\geq\widehat\eta_t>0\), the pathwise recursion obeys

\[
e_{t+1}\leq\alpha_te_t+\beta_t,\qquad
\alpha_t=\frac{(1+R\|\lambda\|_*)L_t}{\widehat\eta_t},\qquad
\beta_t=\frac{(1+R\|\lambda\|_*)\epsilon_{B,t}R}{\widehat\eta_t}.
\]

The product-sum bound of Proposition 18 applies. A uniform bound follows
if these constants have a common contraction margin and bounded forcing.
The algebraic bound alone does not certify that a fitted state or every
unchosen branch is probability-valid.

\textbf{Proof.} The branch likelihood differs by \(\lambda d\), giving
the lower bound on \(\widehat p\). Subtract \(u/p\) from
\((u+d)/(p+\lambda d)\) and use \(u=pz^+\). The induced norm of
\(I-z^+\lambda\) is at most \(1+R\|\lambda\|_*\). Applying that identity
at each step gives the recursion. \(\square\)

Rare evidence can create a large posterior change from a small absolute
error. For example, take a true unnormalized vector \((\eta,0)\) and a
perturbed vector \((\eta,\eta)\). Their \(\ell^1\) difference is
\(\eta\), while their normalized vectors differ by one in \(\ell^1\).
Both arise as possible branch outputs of valid probability instruments
when \(0<\eta\leq1/2\), after assigning the remaining mass to another
branch. There is no uniform posterior Lipschitz constant independent of
a probability margin.

When a known closed convex feasible base is available, Euclidean
projection of a tentative posterior back onto that base cannot increase
its distance from the true posterior. This is a declared repair to the
estimator, not an identity for an unmodified update. It also does not
repair negative branch probabilities of arbitrary fitted operators
before normalization.

\section{A finite experimental route to the probability table}\label{LSM:a-finite-experimental-route-to-the-probability-table}

The following deliberately simple protocol makes a finite-sample
statement possible without confusing a deterministic identity with a
statistical guarantee.

Fix \(M\) finite words before estimation. For each word, reset
independently to the same initial distribution, execute its prescribed
action sequence, and record the indicator that all observations match.
Obtain \(N\) independent trials per word. Dependencies between different
words' estimates are allowed, but the \(N\) trials used for each
individual mean must be independent with the declared success
probability.

\Needspace{5\baselineskip}
\textbf{Proposition 35. Reset-experiment probability certificate.} Let
\(\widehat f(w)\) be the sample mean for each selected word. For
\(0<\delta<1\), with probability at least \(1-\delta\), all \(M\)
estimates satisfy

\[
|\widehat f(w)-f(w)|\leq
\epsilon_N,\qquad
\epsilon_N=\sqrt{\frac{\log(2M/\delta)}{2N}}.
\]

If the selected words cover every entry of \(H\), every \(G_s\), and
\(b_0\), then on the same event

\[
\|\widetilde H-H\|_2\leq r\epsilon_N,\qquad
\|\widetilde G_s-G_s\|_2\leq r\epsilon_N,\qquad
\|\widetilde b_0-b_0\|_2\leq\sqrt r\epsilon_N.
\]

These bounds feed Propositions 32 and 33 whenever their respective gap
conditions hold. The empty-word probability may be set exactly to one
instead of sampled.

\textbf{Proof.} For a centered Bernoulli variable \(X-p\), write
\(g(t)=\log\mathbb E\exp(t(X-p))\). Then \(g(0)=g'(0)=0\), and
\(g''(t)\) is the variance of a Bernoulli under exponential tilting, at
most \(1/4\). Integrating twice gives \(g(t)\leq t^2/8\). Independence
and the exponential Markov inequality give an upper-tail bound
\(\exp(-2N\epsilon^2)\) after optimizing \(t=4\epsilon\). The lower tail
has the same bound. A union bound over the \(M\) words proves the
simultaneous result. The spectral norm is bounded by the Frobenius norm,
yielding the matrix and vector estimates. \(\square\)

This is the Bernoulli form of Hoeffding's inequality {[}R11{]}, proved
here for the protocol used. It is a conservative entrywise-to-spectral
bound, not a claim of optimal sample complexity. The number and lengths
of chosen experiments are part of its cost.

If core tests and histories were selected from data, use a separate
independent discovery sample and freeze the selected finite experiment
list before these estimation trials. Conditional on that selection, the
proposition applies. Alternatively, a uniform bound over a fixed larger
candidate list can support data-dependent selection within that list.
Reusing a post-selected table without either argument changes the claim.

A single dependent trajectory, an adaptive stopping rule, unknown
resets, drifting sensors, or experiments chosen without coverage do not
satisfy this proposition as stated. Those settings require their own
concentration and identifiability arguments. Synthetic reset sampling in
the reference example demonstrates the protocol, not its validity for
every measurement source.

\section{A filtering defect with a packet interpretation}\label{LSM:a-filtering-defect-with-a-packet-interpretation}

The relation to the realization calculus is particularly direct before
normalization. Let \(R\) be a linear summary of a hidden probability
space, let \(\lambda R=\mathbf1^\top\), and let \(B_s\) be a proposed
reduced branch operator. For a true belief \(p\), put

\[
z=Rp,\qquad
\delta_s=(RM_s-B_sR)p.
\]

The defect \(\delta_s\) is a vector in the reduced representation, not a
probability of model failure. It records the failure of the unnormalized
instrument to commute with the summary.

\Needspace{5\baselineskip}
\textbf{Proposition 36. Exact normalized closure defect.} Suppose the
true and proposed branch likelihoods

\[
a=\mathbf1^\top M_sp>0,\qquad
b=\lambda B_sz>0
\]

are positive, and let \(\widehat z^+=B_sz/b\). Then

\[
\boxed{
R\left(\frac{M_sp}{a}\right)-\widehat z^+
=\frac{(I-\widehat z^+\lambda)\delta_s}{a},
\qquad a-b=\lambda\delta_s.
}
\]

If \(M_s\neq0\) and both likelihoods and normalized updates agree on
every strictly positive belief, then \(RM_s=B_sR\). Agreement of
normalized updates alone need not force this equality.

\textbf{Proof.} The true unnormalized summary is \(B_sz+\delta_s\), with
mass \(a=b+\lambda\delta_s\). Subtract the two normalized vectors and
collect terms. Since a nonzero nonnegative branch has positive
likelihood on every strictly positive belief, agreement permits
multiplication by the common mass. This gives equality of unnormalized
summaries on the interior of the probability simplex, whose linear span
is the whole hidden vector space. For the last claim, a one-state model
can assign different positive branch probabilities while its normalized
state always remains the scalar one. \(\square\)

For an orthogonal retained projector \(P\) preserving the mass row, put
\(Q=I-P\), \(R=P\) on its image and \(B_s=PM_sP\). Then

\[
\delta_s=PM_sQp.
\]

This is the unresolved-to-retained transfer block from Chapter~\ref{LSM:book:chapter-31}. Its
posterior effect is not simply its norm: normalization applies the
oblique map \(I-\widehat z^+\lambda\) and divides by the true branch
probability. A component proportional to \(\widehat z^+\) changes only
the branch mass and disappears from the normalized defect, while still
changing likelihood-based inference over mechanisms.

A certificate should therefore retain both \(\lambda\delta_s\) and the
posterior defect. A posterior that looks correct can hide a wrong
evidence model. This is the stochastic counterpart of the distinction
between reconstruction and predictive closure.

\section{A complete two-action rational example}\label{LSM:a-complete-two-action-rational-example}

Use two hidden states only to generate a reference process. The learner
does not receive these matrices. Let

\[
T_0=\begin{pmatrix}3/4&1/4\\1/4&3/4\end{pmatrix},\qquad
T_1=\begin{pmatrix}1/4&2/3\\3/4&1/3\end{pmatrix},\qquad
p_0=\begin{pmatrix}2/3\\1/3\end{pmatrix}.
\]

The observation matrices are

\[
O_0=\operatorname{diag}(1/4,3/4),\qquad
O_1=\operatorname{diag}(3/4,1/4),\qquad
M_{a,o}=O_oT_a.
\]

The two transitions do not commute. Thus the process is genuinely
action-dependent rather than one autonomous recurrence with an action
label attached.

Choose tests \((\epsilon,(0,1))\) and histories \((\epsilon,(0,1))\).
The observable tables give

\[
H=\begin{pmatrix}1&13/24\\13/24&5/16\end{pmatrix},\qquad
\det H=11/576,\qquad
b_0=\begin{pmatrix}1\\13/24\end{pmatrix}.
\]

The selected coordinate \(q\) is the probability of seeing observation
one after choosing action zero. In the hidden realization,

\[
R=\begin{pmatrix}1&1\\5/8&3/8\end{pmatrix},\qquad
z=Rp=(1,q)^\top.
\]

The observable reconstruction \(B_s=G_sH^{-1}\) yields

\[
\begin{aligned}
B_{0,0}&=\begin{pmatrix}1&-1\\11/32&-1/4\end{pmatrix},&
B_{0,1}&=\begin{pmatrix}0&1\\-3/32&3/4\end{pmatrix},\\[4pt]
B_{1,0}&=\begin{pmatrix}5/48&5/6\\23/192&5/24\end{pmatrix},&
B_{1,1}&=\begin{pmatrix}43/48&-5/6\\37/64&-5/8\end{pmatrix}.
\end{aligned}
\]

The coefficients include negatives. Nevertheless, every operator
preserves the cone \(R\mathbb R^2_+\), because \(B_sR=RM_s\). A
sufficient normalized base is

\[
K_0=\{(1,q)^\top:3/8\leq q\leq5/8\}.
\]

This cone certificate uses the known generating model and is checked
separately. It is not inferred merely from the finite probability table,
and \(K_0\) is not claimed to be the maximal canonical base of
Proposition 29.

The likelihoods and posterior maps can now be read from the matrices.
For action zero and observation one,

\[
p(1\mid z,0)=q,\qquad
q^+=\frac{-3/32+(3/4)q}{q}
=\frac34-\frac{3}{32q}.
\]

Starting at \(q_0=13/24\), this gives probability \(13/24\) and
\(q_1=15/26\). If action one then produces observation zero, its
probability is

\[
\frac5{48}+\frac56\frac{15}{26}=\frac{365}{624},
\]

and the new predictive coordinate is

\[
q_2=\frac{23/192+(5/24)(15/26)}{365/624}=\frac{599}{1460}.
\]

The joint probability of this two-step action/observation word is
\(365/1152\). The reference checks compare these values and every word
up to a declared finite test length against the hidden generator in
exact rational arithmetic. The proposition, under its rank and coverage
assumptions, is what extends the result to all word lengths.

The model's online state is not a sequence of newly measured future
probabilities. It is a recursively maintained prediction initialized
from learned experiment probabilities and updated after each actual
action and observation. Uncertainty in those learned parameters is
separate from the stochastic uncertainty in the next observation.

\section{Integrated controlled stochastic realization}\label{LSM:integrated-controlled-stochastic-realization}

\Needspace{5\baselineskip}
\textbf{Proposition 37. Conditional observation-based stochastic
completion.} Suppose a causal finite-alphabet signal process has
behavior rank at most \(r\), and a selected \(r\times r\) history/test
block containing the empty test is invertible. Exact values of that
block, every one-symbol shifted block, and the initial core-test vector
determine a minimal realization of every future experiment probability.
Its normalized state is recursively accessible after each
positive-probability history, without reconstructing the hidden physical
state.

The exact realization has the canonical probability cone of Proposition
29. Under any invertible atlas, transform the state, symbol operators,
mass row and cone together:

\[
z'=Tz,\qquad B_s'=TB_sT^{-1},\qquad
\lambda'=\lambda T^{-1},\qquad C'=TC.
\]

Every likelihood and posterior prediction is then unchanged. A complete
packet decomposition preserves these predictions after full synthesis.
Packetwise probability interpretation requires the extra positivity
condition of Proposition 30.

For estimates obtained under Proposition 35's reset protocol, the
simultaneous event there bounds the finite table perturbations. A
separated finite-table spectrum yields Proposition 32's subspace
certificate. A well-conditioned fixed core yields Proposition 33's
operator certificate. On branches satisfying Proposition 34's likelihood
margins, the resulting normalized state errors obey its recursive bound.
These perturbation statements alone do not certify a learned cone or
global positivity of an estimated model.

\textbf{Proof.} Proposition 28 reconstructs the unique minimal behavior
realization, and Proposition 27 gives its conditional interpretation.
Proposition 29 supplies probability validity for the exact realization.
Substituting the transformed quantities cancels every \(T^{-1}T\),
including in the scalar normalization. Full packet synthesis is the
identity, whereas Proposition 30 characterizes the additional
probability-mixture condition. The perturbation conclusions are the
stated chain of propositions with their separate assumptions retained.
No missing gap, coverage or positivity assumption follows from matrix
inversion alone. \(\square\)

This theorem closes a finite controlled and stochastic sector of the
foundational model. It does not claim unrestricted discovery of a
nonlinear environment from a single trajectory. It does identify what a
learned predictive state represents, why its operations compose, where
its evidence comes from, and which defects change likelihoods or
posteriors.

\section{From realization to admissible domains}\label{LSM:from-realization-to-admissible-domains}

Proposition 37 identifies a finite controlled realization from a covered
population table. Existence of its canonical probability cone does not
provide a finite certificate for a fitted model. Chapter~\ref{LSM:book:chapter-34} supplies such
certificates in finite positive-lift sectors, while keeping internal
validity distinct from statistical coverage.


\label{LSM:part-iv.-certified-probability-domains-and-robust-learning}

A matrix learned from observations is not yet a probability model. Chapter~\ref{LSM:book:chapter-33} made that distinction exact: a finite predictive realization needs
an admissible cone preserved by its action and observation operators. It
supplied a canonical cone defined by all possible future tests, but did
not supply a finite way to verify that cone for an estimated model.

This block develops a constructive finite sector of that missing step.
It gives a certificate that can be checked by rational arithmetic, a
search procedure that does not need hidden-state matrices, and a robust
version that covers an entire declared uncertainty set. It also proves
why a finite polyhedral certificate is not available for every valid
finite-dimensional process. The resulting distinction is between an
exact finite certificate, an uncertified candidate, and an explicitly
modified approximation. Those are different mathematical outcomes.

The positive-realization and invariant-cone ingredients are classical.
Jaeger's account includes the polyhedral-cone characterization and an
irrational-rotation obstruction {[}R12{]}. Polyhedral invariance is also
a standard subject of linear systems and optimization {[}R13{]}. The
present contribution is a self-contained connection of these ingredients
to Komorebi's observable learner, probability packets, confidence sets
and approximation defects. The controlled rational rotation example is
an adaptation of that classical obstruction, not a claim of discovering
the phenomenon.

The next twelve propositions continue the numbering from the preceding
volume. All have proofs below. The exact computations check stated
finite instances and counterexamples. They are not a Lean formalization
and do not change the production PR-LASM 2.3.0 runtime.


\chapter{Certified probability domains and robust inference}\label{LSM:book:chapter-34}
\section{The object to be certified}\label{LSM:the-object-to-be-certified}

Fix a real predictive space \(E=\mathbb R^r\), a nonzero row
\(\lambda\), and a finite family \(B_s\), indexed by \(s=(a,o)\). The
action and observation alphabets are finite and nonempty. Every declared
observation branch is present, including branches represented by the
zero operator. Temporal words act from right to left:

\[
B_{s_1\cdots s_k}=B_{s_k}\cdots B_{s_1},\qquad B_\epsilon=I.
\]

The mass identity is an exact constraint,

\[
\lambda\sum_o B_{a,o}=\lambda \quad\text{for every action }a.
\]

An initial state has \(\lambda z_0=1\). The proposed probability of a
word is \(\lambda B_wz_0\). An invertible coordinate change must
transform \(z_0\), \(B_s\), \(\lambda\) and the admissible cone
together, as in Proposition 37.

A \textbf{probability-domain certificate} specifies a closed, pointed,
generating cone \(C\subset E\), with \(\lambda\) strictly positive on
\(C\setminus\{0\}\), such that

\[
z_0\in C,\qquad B_sC\subseteq C.
\]

Its normalized base is \(K=\{z\in C:\lambda z=1\}\). These conditions
imply nonnegative word probabilities and their per-action normalization
at every depth. Positive-probability updates remain in \(K\).
Zero-probability branches have no conditional posterior and must not be
assigned one by denominator flooring.

There are three separate claims. \textbf{Internal validity} says the
fitted model is a stochastic model. \textbf{Data compatibility} says it
agrees with the observations within a specified error budget.
\textbf{Coverage of the generating process} says the assumptions and
uncertainty set actually include the process that generated those
observations. The first can have a deterministic certificate even when
the other two remain conditional.

\section{Finite tests give an outer hierarchy, not automatic closure}\label{LSM:finite-tests-give-an-outer-hierarchy-not-automatic-closure}

Define the test cones

\[
C_h=\{z:\lambda B_wz\geq0\text{ for every }|w|\leq h\},
\qquad C_0=\{z:\lambda z\geq0\}.
\]

Each \(C_h\) is polyhedral because it is an intersection of finitely
many homogeneous half-spaces. It need not be a pointed cone or have a
bounded normalized base.

\Needspace{5\baselineskip}
\textbf{Proposition 38. Test-cone hierarchy and a finite stopping
certificate.} The cones satisfy

\[
C_{h+1}=C_h\cap\bigcap_s B_s^{-1}C_h,
\qquad C_{h+1}\subseteq C_h.
\]

Their intersection \(C_\infty\) is the largest cone contained in \(C_0\)
and invariant under all \(B_s\). If an exact equality \(C_{h+1}=C_h\)
occurs, then \(C_h=C_\infty\) and the finite inequalities certify all
future word probabilities. If the representation is observable, meaning
that the rows \(\lambda B_w\) span \(E^*\), then \(C_\infty\) is pointed
and \(\lambda\) is strictly positive on its nonzero elements. If
reachable vectors \(B_wz_0\in C_\infty\) span \(E\), the cone is also
generating.

\textbf{Proof.} For a first symbol \(s\) followed by a word \(v\),
\(B_{sv}=B_vB_s\). The inequalities for all words of length at most
\(h+1\) are therefore exactly the displayed recurrence. If
\(z\in C_\infty\), every test on \(B_sz\) is another test on \(z\), so
\(C_\infty\) is invariant. Conversely, invariance of a cone contained in
\(C_0\) makes every successive test nonnegative, hence places it in
\(C_\infty\). Equality at one finite step establishes invariance
immediately and therefore gives the intersection.

If both \(z\) and \(-z\) lie in \(C_\infty\), every test row annihilates
\(z\), so observability gives \(z=0\). For \(z\in C_\infty\) with
\(\lambda z=0\), sum word probabilities along any fixed action sequence.
Mass conservation makes the sum zero. Every term is nonnegative, so
every word test vanishes and again \(z=0\). Spanning by reachable cone
elements gives the final assertion. \(\square\)

A useful negative witness is a state admitted by the current
inequalities and a specific longer word with negative proposed
probability. It demonstrates that this state needs to be excluded. It
does not prove that the initial state, or every reachable state, is
invalid.

Finite stabilization is sufficient, not promised. Even three predictive
coordinates can require a nonpolyhedral domain, as Section 51 shows.
Checking a long list of trajectories is not the same as checking the
invariance identity that makes every trajectory valid.

\section{A finite certificate with a positive lift}\label{LSM:a-finite-certificate-with-a-positive-lift}

Let \(V=[v_1\ \cdots\ v_m]\in\mathbb R^{r\times m}\) have full row rank
and satisfy

\[
\lambda V=\mathbf1_m^\top.
\]

Its columns generate \(C=V\mathbb R_+^m\) and its normalized base is
\(K=\operatorname{conv}\{v_1,\ldots,v_m\}\). Finite generation makes the
cone closed. Its normalized generators and full row rank make it pointed
and generating, with a strictly positive mass functional.

\Needspace{5\baselineskip}
\textbf{Proposition 39. Finite positive-lift certificate.} This cone is
invariant under \(B_s\) if and only if there is an entrywise nonnegative
matrix \(A_s\in\mathbb R^{m\times m}\) such that

\[
\boxed{B_sV=VA_s,\qquad A_s\geq0.}
\]

If these identities hold for every symbol and \(z_0=V\alpha_0\) for
\(\alpha_0\geq0\) with \(\mathbf1^\top\alpha_0=1\), then

\[
\mathbf1^\top\sum_o A_{a,o}=\mathbf1^\top,
\qquad
\lambda B_wz_0=\mathbf1^\top A_w\alpha_0.
\]

Thus the finite witness defines a nonnegative probability instrument
whose image under \(V\) is the predictive model. It certifies all word
probabilities without enumerating them.

\textbf{Proof.} Invariance places each \(B_sv_j\) in the conical hull of
the columns of \(V\). Choose its nonnegative coefficients as column
\(j\) of \(A_s\). Conversely, \(B_sV\alpha=VA_s\alpha\in C\) for every
\(\alpha\geq0\). Multiplying the intertwining identity by \(\lambda\)
and summing over observations gives

\[
\mathbf1^\top\sum_oA_{a,o}
=\lambda\sum_oB_{a,o}V
=\lambda V=\mathbf1^\top.
\]

Iterating \(B_sV=VA_s\) gives \(B_wV=VA_w\). Applying the mass row and
the initial coefficients proves the word identity. Strict positivity of
\(\lambda\) follows because a nonzero conical combination has positive
total coefficient mass. \(\square\)

The \(m\) nonnegative coordinates need not be physical hidden states,
and \(m\) need not equal the minimal predictive dimension \(r\). They
are one finite realization witnessing validity. When \(V\) has dependent
columns, the lift can be nonunique. Those redundant coordinates do not
change the predictive state after synthesis.

There is a dual finite test. If a complete inequality description of the
same cone is \(C=\{z:Gz\geq0\}\), then \(GB_sv_j\geq0\) for every
generator \(j\) is equivalent to invariance. One may also certify
\(GB_s=H_sG\) with \(H_s\geq0\). The latter certificate is sufficient by
direct multiplication and necessary by the polyhedral cone form of
Farkas' lemma {[}R13{]}. The reference calculation uses the direct
generator test, not an unreported numerical optimization certificate.

\section{Robustness is a quantified invariance statement}\label{LSM:robustness-is-a-quantified-invariance-statement}

A certificate for one fitted operator does not automatically cover
neighboring operators. For uncertainty, the quantifier must range over
the whole proposed model set, including correlations required by mass
conservation.

Let \(\mathcal U\) be the convex hull of finitely many \textbf{complete
operator families} \(\{B_s^{(j)}\}_s\). Every family must satisfy the
exact per-action mass identities. The convex coefficient is shared
across that family's symbols. A product of independent uncertainty sets
can instead be handled by its joint vertices.

\Needspace{5\baselineskip}
\textbf{Proposition 40. Robust finite-domain certificate.} For a fixed
normalized generator matrix \(V\), suppose every vertex family has
witnesses

\[
B_s^{(j)}V=VA_s^{(j)},\qquad A_s^{(j)}\geq0.
\]

Then every family in \(\mathcal U\) preserves \(C=V\mathbb R_+^m\) and
is mass conserving. If its initial state lies in \(K\), it defines valid
finite observation-tree probabilities at every horizon. This remains
true when a family in \(\mathcal U\) is selected separately at each
history.

For a complete facet description \(C=\{z:Gz\geq0\}\), let \(g_i^\top\)
be its rows. A sufficient norm-ball certificate around \(\widehat B_s\)
is

\[
\boxed{
 g_i^\top\widehat B_sv_j
 \geq\epsilon_s\|g_i\|_2\|v_j\|_2
 \quad\text{for all }s,i,j.
}
\]

It covers all perturbations with
\(\|B_s-\widehat B_s\|_2\leq\epsilon_s\) that also satisfy the exact
mass constraints.

\textbf{Proof.} For convex weights \(\theta_j\), take
\(A_s=\sum_j\theta_jA_s^{(j)}\). The equality and nonnegativity survive
convex combination, as does mass conservation. At each history, every
selected family maps the current cone state back into the cone. Applying
the mass identity at that node proves recursive normalization,
regardless of how the family was selected. For the norm ball,

\[
|g_i^\top(B_s-\widehat B_s)v_j|
\leq\epsilon_s\|g_i\|_2\|v_j\|_2.
\]

The margin makes every perturbed generator image satisfy all facets.
Proposition 39, or the direct generator test, then supplies invariance.
\(\square\)

The switching conclusion is structural validity, not an assertion that
time-varying selection represents one fixed unknown parameter. A
likelihood for a fixed mechanism and a likelihood for a mechanism
selected anew at each history are different statistical models.

The norm-ball inequality is sufficient and can be conservative,
especially at boundary images with zero slack. Failure of this margin is
not a proof that no robust certificate exists. The exact vertex test
retains uncertainty correlations that a separate bound for each entry
may discard.

\section{Fitting inside a certified domain}\label{LSM:fitting-inside-a-certified-domain}

Once \(V\) is fixed, certificate-constrained fitting is a convex
problem. Unknowns include the entries of \(\widetilde B_s\), \(A_s\),
\(\widetilde z_0\) and \(\alpha_0\). The constraints

\[
\begin{gathered}
\widetilde B_sV=VA_s,\quad A_s\geq0,
\quad \lambda\sum_o\widetilde B_{a,o}=\lambda,\\
\widetilde z_0=V\alpha_0,\quad\alpha_0\geq0,
\quad\mathbf1^\top\alpha_0=1
\end{gathered}
\]

are linear. Entrywise confidence intervals and fixed-input regression
tolerances add linear inequalities. Norm balls add convex constraints.
Jointly optimizing \(V\) and \(A_s\) is different: the product \(VA_s\)
is then bilinear. No general convexity is claimed for learning the
domain and dynamics together.

\Needspace{5\baselineskip}
\textbf{Proposition 41. Certified correction and its error budget.} For
fixed \(V\), every feasible solution of the displayed constraints is an
internally valid model. Suppose a center estimate \(\widehat B_s\) and a
feasible corrected estimate \(\widetilde B_s\) satisfy

\[
\|B_s^*-\widehat B_s\|\leq\epsilon_s,
\qquad
\|\widetilde B_s-\widehat B_s\|\leq\epsilon_s.
\]

Then \(\|\widetilde B_s-B_s^*\|\leq2\epsilon_s\). The same conclusion
holds for initial states in their specified norm. If the true model lies
in these confidence sets and itself admits the fixed-domain certificate,
the convex constraints are feasible. Feasibility alone does not
establish that the true model preserves this domain.

\textbf{Proof.} Validity is Proposition 39. The error bound is the
triangle inequality through the center estimate. Under the final
assumptions, insert the true operators, true initial state and their
nonnegative witnesses into the constraints. They provide a feasible
point. The converse fails because existence of one valid point does not
say that every point in a confidence set, or its true member, has the
same invariant domain. \(\square\)

This is why existential correction and robust certification must be
reported separately. A corrected model can be a well-defined
approximation even when no common true-model cone has been proved. The
trajectory bounds below require both models to be valid on a common
specified cone.

\subsection{45.1 An explicit positivity repair in a simplicial
cone}\label{LSM:an-explicit-positivity-repair-in-a-simplicial-cone}

When \(V\) is square and invertible, put \(A_s=V^{-1}\widehat B_sV\).
Assume mass conservation has already been imposed. Let \(N_s\) be an
entrywise positive reference instrument with the same per-action mass
identities. The mixture

\[
\widetilde A_s=(1-\eta)A_s+\eta N_s
\]

is nonnegative provided

\[
\eta\geq\max_{s,i,j:(A_s)_{ij}<0}
\frac{-(A_s)_{ij}}{(N_s)_{ij}-(A_s)_{ij}}.
\]

An empty maximum is zero. If every negative entry has magnitude at most
\(d\) and every reference entry is at least \(m_0>0\), the simpler
sufficient choice is \(\eta=d/(m_0+d)\). These statements follow by
solving each scalar inequality. Mass conservation is unchanged by convex
mixing.

This is an explicit model modification, not coefficient clipping or an
exact coordinate change. Relative to a true instrument on the same cone,
the one-step error is at most \((1-\eta)\epsilon+2\eta\) in the summed
base-norm metric defined next, when the original error was at most
\(\epsilon\). The factor two follows because two valid instruments each
distribute unit total mass. The reference example retains the repair
amount as part of its approximation record.

\section{A norm that measures probability mass before conditioning}\label{LSM:a-norm-that-measures-probability-mass-before-conditioning}

Assume \(C\) is a closed, pointed, generating cone and \(\lambda\) is
strictly positive on its nonzero elements. Its base \(K\) is compact.
Define the base norm

\[
\|x\|_C=\inf\{\lambda u+\lambda v:x=u-v,\ u,v\in C\}.
\]

The infimum is finite because the cone is generating. Compactness of
\(K\) makes this a genuine norm. Its unit ball is
\(\operatorname{conv}(K\cup(-K))\). In particular \(\|u\|_C=\lambda u\)
on \(C\), and \(|\lambda x|\leq\|x\|_C\). On the ordinary probability
orthant this is the \(\ell^1\) norm. These are standard ordered-space
constructions {[}R15{]}, used here to avoid mixing Euclidean coefficient
magnitude with probability mass.

\Needspace{5\baselineskip}
\textbf{Proposition 42. Instrument contraction before normalization.} If
\(\{B_{a,o}\}_o\) is cone preserving and mass conserving for a fixed
action, then

\[
\boxed{\sum_o\|B_{a,o}x\|_C\leq\|x\|_C.}
\]

For a polyhedral base \(K=\operatorname{conv}\{v_j\}\), define the
discrepancy between two instruments by

\[
\epsilon_a=\sup_{z\in K}\sum_o
\|(B_{a,o}-\widetilde B_{a,o})z\|_C.
\]

The supremum equals the maximum of the same expression at the finitely
many generators \(v_j\). When both instruments are valid on \(C\),
\(\epsilon_a\leq2\).

\textbf{Proof.} For any \(x=u-v\) with \(u,v\in C\), cone preservation
gives

\[
\sum_o\|B_{a,o}x\|_C
\leq\sum_o\lambda B_{a,o}(u+v)
=\lambda u+\lambda v.
\]

Take the infimum. For the generator statement, the discrepancy is convex
in \(z\). At a convex combination its value is at most the same
combination of generator values, hence at most their maximum. The
generators themselves belong to \(K\), proving equality. Finally, for
\(z\in K\), the triangle inequality bounds the sum by the total masses
of both instruments, which are one each. \(\square\)

For a polyhedral certificate, the norm itself has a finite
linear-program formulation:

\[
\|x\|_C=\min_{a,b\geq0}
\{\mathbf1^\top(a+b):V(a-b)=x\}.
\]

If \(V\) is square and invertible, this reduces to \(\|V^{-1}x\|_1\).
The exact worked interval certificate is of this type. With dependent
generators, choosing one pseudoinverse does not in general compute the
base norm.

\section{Whole observation laws can be stable without rare-event margins}\label{LSM:whole-observation-laws-can-be-stable-without-rare-event-margins}

The posterior after one rare observation can be sensitive, but that does
not force the entire distribution of observation records to be equally
sensitive. Work before normalization and retain the probability of each
branch.

Let two models share \((C,\lambda)\) and have initial states
\(z_0,\widetilde z_0\in K\). At each time use the same policy, a
specified function or randomization depending only on the observed
history. Thus on any common history the action probabilities are
identical in the two comparisons. This is not a comparison of two
policies optimized separately for the two models.

\Needspace{5\baselineskip}
\textbf{Proposition 43. Observation-tree total-variation bound.} If both
instruments preserve \(C\) and conserve mass, and if every action at
time \(t\) has the discrepancy in Proposition 42 bounded by
\(\epsilon_t\), then the distributions \(P_T,\widetilde P_T\) of
length-\(T\) observation records under the common policy satisfy

\[
\boxed{
\|P_T-\widetilde P_T\|_{\mathrm{TV}}
\leq\min\left(1,\frac12\|z_0-\widetilde z_0\|_C
+\frac12\sum_{t=0}^{T-1}\epsilon_t\right).
}
\]

The same bound holds for joint action and observation records when the
common policy is randomized. No positive lower bound on individual
branch probabilities is required.

\textbf{Proof.} Let \(u_h,\widetilde u_h\in C\) be the unnormalized
states on each history at depth \(t\), including the common policy
factors. Their masses sum to one in each model. Define
\(D_t=\sum_h\|u_h-\widetilde u_h\|_C\). At a history with selected
action \(a\), write

\[
B_{a,o}u_h-\widetilde B_{a,o}\widetilde u_h
=B_{a,o}(u_h-\widetilde u_h)
+(B_{a,o}-\widetilde B_{a,o})\widetilde u_h.
\]

Sum the first term over observations and use Proposition 42. When
\(\lambda\widetilde u_h>0\), normalize only to apply the definition of
\(\epsilon_t\) to the second term. Its sum is bounded by
\(\epsilon_t\lambda\widetilde u_h\). If the mass is zero, strict
positivity implies \(\widetilde u_h=0\). Summing histories yields
\(D_{t+1}\leq D_t+\epsilon_t\). Averaging over common policy
probabilities gives the same inequality for randomized actions. Finally
each word-probability difference is bounded by its state difference
because \(|\lambda x|\leq\|x\|_C\). Total variation is half the sum of
those absolute differences. \(\square\)

This bound separates total law error from conditional-state error. It
can grow linearly with the horizon and eventually become uninformative.
It is not a uniform-in-time approximation theorem. It also does not
control the regret of a changed planning policy without additional
decision-theoretic assumptions.

\section{Conditioning still has an unavoidable cost}\label{LSM:conditioning-still-has-an-unavoidable-cost}

Let \(P,Q\) be probability distributions on a finite record space. A
record event \(H\) might be a particular observed prefix. Conditional
continuation laws are obtained by restriction to \(H\) and
normalization.

\Needspace{5\baselineskip}
\textbf{Proposition 44. Sharp event-probability dependence of
conditioning.} If \(\|P-Q\|_{\mathrm{TV}}\leq\delta\), \(p=P(H)>0\) and
\(q=Q(H)>0\), then

\[
\boxed{\|P(\cdot\mid H)-Q(\cdot\mid H)\|_{\mathrm{TV}}
\leq\min(1,\delta/p).}
\]

When \(p>\delta\), positivity of \(q\) follows automatically. The factor
\(1/p\) cannot be removed in general.

\textbf{Proof.} Write the subprobability restrictions on \(H\) as
\(u,v\), of masses \(p,q\). On the complementary event, the sum of
absolute probability differences is at least \(|p-q|\). Therefore

\[
\|u-v\|_1+|p-q|\leq2\delta.
\]

Adding and subtracting \(v/p\) gives

\[
\frac12\|u/p-v/q\|_1
\leq\frac{\|u-v\|_1+|p-q|}{2p}
\leq\delta/p.
\]

Also \(|p-q|\leq\delta\), establishing the positivity condition. For
sharpness, put total mass \(p\) on an event with two outcomes. Let \(P\)
assign all of it to the first, and let \(Q\) move mass \(\delta\leq p\)
to the second, leaving the complement unchanged. The unconditional total
variation is \(\delta\) and the conditional distance is \(\delta/p\).
\(\square\)

Combining this proposition with Proposition 43 gives a conditional
continuation bound when the evidence event has adequate probability. It
does not give a coordinate-state error bound unless test probabilities
control that coordinate norm. Validity at every branch and stable
inference at every branch remain different properties.

\section{Searching for a domain from observable operators}\label{LSM:searching-for-a-domain-from-observable-operators}

A cone certificate need not be imported from a hidden-state generator.
With exact rational operators it can be searched for directly, then
checked independently.

\Needspace{5\baselineskip}
\textbf{Proposition 45. Rational finite-certificate search.} Suppose
\(\lambda\), \(B_s\) and \(z_0\) are rational and satisfy the exact mass
identities. There is an enumeration procedure that returns a valid
finite certificate whenever a certificate of Proposition 39 exists with
a rational normalized full-row-rank generator matrix \(V\). The
procedure is sound whenever it returns. No termination guarantee is
asserted when no such rational certificate exists.

\textbf{Proof.} Enumerate all finite rational matrices \(V\) in an order
that eventually includes every size and every bounded numerator and
denominator range. Retain those of full row rank with
\(\lambda V=\mathbf1^\top\). For each retained matrix, the requirements
\(B_sV=VA_s\), \(A_s\geq0\), \(V\alpha_0=z_0\), \(\alpha_0\geq0\) and
\(\mathbf1^\top\alpha_0=1\) are rational linear feasibility problems. A
nonempty rational polyhedron has a rational feasible point, as follows
from elimination of rational linear inequalities. These problems can
therefore be decided by exact rational linear programming, or witnesses
can be enumerated in a dovetailed rational search. Recheck every
returned equality and sign exactly. Proposition 39 proves soundness. If
a rational generator certificate exists, its finite representation is
eventually considered and its feasible witnesses are found. \(\square\)

The general enumeration is a mathematical existence algorithm, not a
practical complexity claim. A bounded candidate family is often the
useful reference experiment. Failure within that family means only that
the chosen family did not produce a certificate.

\subsection{49.1 A certificate learned without hidden
matrices}\label{LSM:a-certificate-learned-without-hidden-matrices}

The observable tables of Section 38 determine the two-dimensional
operators \(B_s\) and initial vector \((1,13/24)^\top\). The new
reference procedure receives only these fitted operators and that
vector. It enumerates intervals with rational endpoints in \([0,1]\) of
reduced denominator at most eight, ordered by increasing width and then
by endpoints. The coordinate restriction \([0,1]\) comes from its
declared test-probability meaning.

The first accepted interval is

\[
\boxed{K=\{(1,q)^\top:2/5\leq q\leq3/5\}.}
\]

It uses

\[
V=\begin{pmatrix}1&1\\2/5&3/5\end{pmatrix},
\qquad \alpha_0=\begin{pmatrix}7/24\\17/24\end{pmatrix},
\qquad A_s=V^{-1}B_sV.
\]

Every entry of every \(A_s\) is at least \(1/64\). The full per-action
column sums are exactly one. These rational identities certify all word
probabilities of the fitted model, not only the words used in the fit.
The certificate was not supplied to the search, and no hidden generator
is imported by the standalone example.

The interval is the narrowest accepted interval \textbf{in this
specified finite grid}, not a claim of the smallest invariant interval
among all real endpoints. It also need not contain every predictive
state associated with an arbitrary initialization of the old hidden
model. It contains the declared initial state and preserves every
subsequent branch of this fitted realization, which is the requirement
being certified.

The positive slack gives a robust statement: any perturbation of the
lifted entries by at most \(1/64\), coupled to retain the exact
per-action column sums, remains nonnegative. Interior perturbations of
size \(1/128\) are used in the worked comparison. Repeated composition
is valid for all horizons, while Proposition 43 bounds the changed
observation law for each fixed horizon.

\section{Selection and noise do not invalidate an exact witness}\label{LSM:selection-and-noise-do-not-invalidate-an-exact-witness}

A data-dependent cone proposal is not itself a statistical confidence
statement. Once the matrices and witness are fixed, its exact invariance
identities hold or fail independently of how it was proposed. What
remains statistical is whether the chosen uncertainty set covers the
generating model.

For a finite-sample interpretation, first obtain a valid uncertainty set
for a fixed identifiable model class. Proposition 35 supplies one route
under independent reset experiments and declared rank coverage.
Proposition 33 translates those probability bounds into a fixed-core
operator bound. A validity witness can then be sought for a fitted
member or robustly for every member. Only the latter identifies a common
true-model domain on the coverage event.

The admissible family must enforce mass conservation exactly. A free box
around unconstrained matrices generally contains nonnormalized families.
Discarding those families is legitimate only because the generating
process is already assumed to obey the exact mass law. Other cuts on a
confidence set need their own justification.

\section{Why finite predictive rank need not mean finitely many positive states}\label{LSM:why-finite-predictive-rank-need-not-mean-finitely-many-positive-states}

The following controlled rational example makes the boundary explicit.
It is a worked adaptation of the irrational-rotation obstruction in
{[}R12, Section 7{]}. Related distinctions among finite linear, positive
and quantum realizations are discussed in {[}R14{]}. Only the elementary
cone argument below is used.

Let \(z=(t,x,y)^\top\), \(\lambda=(1,0,0)\), \(z_0=(1,1,0)^\top\), and

\[
R=\begin{pmatrix}3/5&-4/5\\4/5&3/5\end{pmatrix},
\qquad T=\operatorname{diag}(1,R).
\]

There are two actions and two observations. A rotate action has branches

\[
B_{\mathrm r,0}=T,\qquad B_{\mathrm r,1}=0.
\]

A probe action has branches

\[
B_{\mathrm p,0}=z_0(1/2,1/2,0),\qquad
B_{\mathrm p,1}=z_0(1/2,-1/2,0).
\]

The rotate outcome zero is certain. The probe has probabilities
\((t+x)/2\) and \((t-x)/2\) and resets every nonzero branch to \(z_0\)
after normalization. Both actions conserve the mass \(t\).

\Needspace{5\baselineskip}
\textbf{Proposition 46. A rational rank-three model with no polyhedral
probability-domain certificate.} The Lorentz cone

\[
C_{\mathrm L}=\{(t,x,y):t\geq\sqrt{x^2+y^2}\}
\]

is probability valid for these operators. It is the canonical all-test
cone. Every closed convex invariant cone containing \(z_0\) and
admitting the nonnegative mass interpretation equals \(C_{\mathrm L}\).
In particular no finite normalized generator certificate exists for this
exact model. Its reachable and observable predictive dimension is three.

\textbf{Proof.} Rotation preserves \(t\) and \(x^2+y^2\). On
\(C_{\mathrm L}\) both probe effects are nonnegative and each probe
image is a nonnegative multiple of \(z_0\). Thus the displayed cone is
invariant and valid.

Write the rotation angle as \(\theta\). It is not a rational multiple of
\(2\pi\). Otherwise \(\zeta=e^{i\theta}\) would be a root of unity, so
\(\zeta+\zeta^{-1}=6/5\) would be a rational algebraic integer. Every
rational algebraic integer is an integer, by the rational-root argument
for a monic integer polynomial, contradicting \(6/5\notin\mathbb Z\).
The forward powers of an irrational circle rotation are dense.
Compactness supplies powers approaching the identity, so their closure
contains the inverse rotation and is a closed subgroup. An infinite
closed subgroup has arbitrarily small positive angles. Integer multiples
of such angles approximate every point on the circle, so that subgroup
is the entire circle.

A state passing every test must pass a probe after any number \(k\) of
rotations. Hence

\[
t\pm e_1^\top R^k(x,y)^\top\geq0
\quad\text{for every }k\geq0.
\]

Density gives \(t\geq\sqrt{x^2+y^2}\). The reverse inclusion was already
proved, so the canonical cone is exactly \(C_{\mathrm L}\).

Any invariant cone with nonnegative mass lies inside this canonical
cone, by Proposition 38. If it is closed, convex and contains \(z_0\),
it contains the orbit \(T^kz_0\) and the closed convex hull of that
orbit. On \(t=1\), this hull is the full unit disk. Conical scaling
gives all of \(C_{\mathrm L}\). The unit disk has infinitely many
extreme points, so it is not the convex hull of finitely many points.

Finally the reachable columns \(z_0,Tz_0,T^2z_0\) are independent. The
mass row and probe rows before zero and one rotations span the dual
space because the rotation has nonzero sine. Reachability and
observability establish predictive dimension three. \(\square\)

This is not a failure of finite-dimensional prediction. The exact model
has only three coordinates and a simple quadratic cone. It is a failure
of the narrower requirement that every valid model admit a finite
\emph{polyhedral} domain. A certificate language limited to finitely
many positive generators would exclude a legitimate low-rank signal
model.

No general impossibility of all finite certificates follows. This very
example has a finite quadratic certificate. The obstruction concerns
polyhedral certificates, not invariant cones described by other
mathematics.

\section{A finite approximation with an explicit price}\label{LSM:a-finite-approximation-with-an-explicit-price}

A curved probability domain can sometimes be replaced by a finite one
after changing the dynamics. That change must be quantified, not
described as exact refinement.

Let \(P_n\) be the regular \(n\)-gon inscribed in the unit disk, with a
vertex at \((1,0)\) and \(n\geq3\). Replace the rotate branch by

\[
T_\alpha=\operatorname{diag}(1,\alpha R),\qquad 0\leq\alpha<1,
\]

while retaining the same probe branches. The spatial coordinates are
attenuated, but total probability mass remains unchanged.

\Needspace{5\baselineskip}
\textbf{Proposition 47. Polyhedral approximation of the rotation model.}
If

\[
\alpha\leq\cos(\pi/n),
\]

then the cone over \(P_n\) is invariant under all modified operators and
contains \(z_0\). Consequently the modified model admits a finite
positive lift. The original and modified models are both valid on the
larger cone \(C_{\mathrm L}\), and for every common policy their
length-\(T\) laws satisfy

\[
\boxed{\|P_T-P_T^{(\alpha)}\|_{\mathrm{TV}}
\leq\min(1,T(1-\alpha)/2).}
\]

The angle-independent sufficient condition can be improved for a fixed
polygon and rotation. For the rational diamond \(|x|+|y|\leq t\), the
present rotation is invariant after attenuation whenever
\(\alpha\leq5/7\). Thus \(\alpha=2/3\) gives an entirely rational
four-generator certificate.

\textbf{Proof.} The regular polygon contains the disk of radius
\(\cos(\pi/n)\) and lies in the unit disk. Every \(\alpha Rv\) with
\(v\in P_n\) has Euclidean norm at most \(\alpha\), hence lies in
\(P_n\) under the stated condition. The probe effects remain nonnegative
because \(P_n\) lies in the disk, and their normalized images are the
retained vertex \((1,0)\).

The base norm of \(C_{\mathrm L}\) is

\[
\|(t,x,y)\|_{C_{\mathrm L}}=\max(|t|,\sqrt{x^2+y^2}).
\]

Indeed the convex hull of the unit disks in the planes \(t=1\) and
\(t=-1\) is the cylinder with this unit ball. The rotate-branch
discrepancy on a normalized state is at most \(1-\alpha\) in this norm.
Probe discrepancies are zero. Proposition 43 gives the law bound. For
the diamond, the induced \(\ell^1\) norm of \(R\) is \(7/5\), so
\(\alpha\|R\|_1\leq1\) proves invariance. \(\square\)

For a target attenuation loss \(0<\eta\leq1\), choosing
\(\alpha=1-\eta\) and sufficiently many polygon vertices satisfies the
certificate. The inequality \(1-\cos u\leq u^2/2\) shows that
\(n\geq\pi/\sqrt{2\eta}\) is sufficient, with an integer ceiling and
\(n\geq3\). This is a constructive upper bound on certificate size, not
a lower bound or a proof of optimality.

The distinction is exact. At \(\alpha=1\) there is no finite polyhedral
certificate for this model. At \(\alpha<1\) a sufficiently fine polygon
works, at the cost of a declared change in the observation law. The
reference figures show that tradeoff rather than labeling the
approximation an identical model.

\section{Repeated proposal and verification with a confidence budget}\label{LSM:repeated-proposal-and-verification-with-a-confidence-budget}

Model development may choose another core, propose another domain or
collect another independent batch after an earlier certificate is
withheld. Statistical coverage must survive that selection process.

\Needspace{5\baselineskip}
\textbf{Proposition 48. Sequential confidence accounting for certified
proposals.} At stage \(n\), let the proposal be measurable from previous
data, and let a verification procedure return a set \(\mathcal U_n\) for
the true model in the declared identifiable coordinate system. Assume
conditional coverage

\[
\Pr(\text{true model}\notin\mathcal U_n\mid\text{previous data})
\leq\delta_n,
\qquad\sum_{n\geq1}\delta_n\leq\delta.
\]

Then with probability at least \(1-\delta\), every returned uncertainty
set covers its true model. On this event, any data-selected domain with
an exactly verified robust certificate for \textbf{all} members of its
corresponding \(\mathcal U_n\) is also valid for the true model,
provided the true initial state is covered by the certified base. The
conclusion remains valid at any stopping time chosen from the recorded
results.

\textbf{Proof.} Take expectations of the conditional failure bounds,
then apply the countable union bound. Off the union of failure events,
all coverages hold simultaneously. Robust invariance over an entire set
then applies to its true member by deterministic implication,
independently of how its domain was chosen. A stopped stage is one of
the already covered stages. \(\square\)

For example, \(\delta_n=\delta 2^{-n}\) spends a finite total confidence
budget. Fresh reset trials for a proposal fixed by past data can use
Proposition 35 conditionally, with this stage budget. An exact witness
needs no statistical correction for how many candidate cones were tried.
The coverage statement does need correction for repeatedly constructing
confidence sets or inspecting post-selected data.

The proposition does not infer a global rank bound, validate an unknown
reset operation, or transform dependent observations into independent
trials. Those are assumptions of the coverage procedure. A known
coordinate transformation can transport a confidence set exactly. An
estimated transformation requires the subspace and conditioning controls
already separated in Chapter~\ref{LSM:book:chapter-33}.

\section{Integrated certificate-carrying realization}\label{LSM:integrated-certificate-carrying-realization}

\Needspace{5\baselineskip}
\textbf{Proposition 49. Conditional learning with an all-horizon
probability certificate.} Suppose the observable identification
assumptions of Proposition 37 hold for a declared finite-rank controlled
process. Suppose a justified confidence procedure produces an operator
and initial-state set containing the true realization on an event of
probability at least \(1-\delta\). Assume the per-action mass
constraints are exact.

If a fitted realization admits a checked finite witness of Proposition
39, it is an internally valid stochastic model at every horizon. If the
witness is robust over the whole confidence set and covers its initial
states, then on the coverage event the same probability domain is valid
for both the fitted and true models. When their summed base-norm
instrument discrepancies are at most \(\epsilon_t\) and their initial
discrepancy is at most \(e_0\), every common policy has the bound

\[
\|P_T-\widehat P_T\|_{\mathrm{TV}}
\leq\min\left(1,\frac{e_0+\sum_{t<T}\epsilon_t}{2}\right).
\]

Conditioning on a positive-probability event then obeys Proposition 44.
Exact coordinate transport preserves the certificate. A signed packet
decomposition preserves it after complete synthesis but does not turn
individual packets into probability states without Proposition 30's
extra positivity requirement.

\textbf{Proof.} Observable identification supplies the specified
realization and the confidence procedure supplies the stated event, not
an inferred event of its own. Proposition 39 supplies internal validity
for the fitted model. Proposition 40 extends the same domain to every
member of the uncertainty set, hence to its true member on that event.
Propositions 42 and 43 give the total-variation bound, and Proposition
44 gives the conditional bound. For a coordinate change \(T\), use
\(V'=TV\), \(B_s'=TB_sT^{-1}\) and \(\lambda'=\lambda T^{-1}\). Then
\(B_s'V'=V'A_s\) and \(\lambda'V'=\mathbf1^\top\) with the unchanged
nonnegative witnesses. Full packet synthesis is the identity.
\(\square\)

This closes a concrete foundational gap. A learner can now return not
just a matrix family but a verifiable domain in which that family has
probability meaning. It can attach a robust uncertainty statement when
one is justified, or return an internally valid approximation with its
correction budget explicitly recorded.

It does not turn a successful fit into unrestricted knowledge of the
world. The model class, observation protocol, rank coverage, exact mass
law, domain witness and approximation budget remain visible parts of the
result. Failure of a finite search is a withheld certificate, not a
proof of invalidity. The nonpolyhedral example explains why this
distinction is necessary even for very small exact models.

\section{From polyhedral to curved certificates}\label{LSM:from-polyhedral-to-curved-certificates}

A finite positive lift is one sufficient certificate language.
Proposition 46 shows that a low-dimensional model can require a
nonpolyhedral domain. Chapter~\ref{LSM:book:chapter-35} represents that same model exactly through
a curved matrix domain. A failed finite-generator search is therefore
not a general invalidity theorem.


\label{LSM:part-v.-curved-domains-and-continuous-observations}

Chapter~\ref{LSM:book:chapter-34} separated a valid predictive model from a successful fit. Its
finite positive lift gives a useful certificate, but its rank-three
rotation example has no finite polyhedral probability domain. The same
block left continuous observations outside the finite-alphabet
construction. This part addresses both boundaries in a declared finite
sector.

A curved domain can have a short matrix certificate even when it has
infinitely many extreme states. A continuous observation can have a
finite operator description even when its possible values are
uncountable. Neither statement follows from finite predictive dimension
alone. The constructions below require an explicit domain, a reference
measure and a restricted observation-kernel family.

The block connects matrix positivity, integral observation semantics,
population-moment identification, exact coarsening and approximation. It
also identifies a defect that is easy to overlook: averaging a matrix
transformation before applying its quadratic action does not equal
averaging that action. The difference is a positive covariance
contribution. Omitting it can destroy normalization even in a very small
example.

There are established antecedents. Congruence representations and their
positive semidefinite Gram matrices belong to the theory of completely
positive maps {[}R16{]}. Polynomial Gram certificates are standard
sum-of-squares machinery {[}R17{]}. Continuous-valued observable
operator models already occur in Jaeger's work {[}R18{]}. The proofs
here develop their specific connection to the realization and defect
contracts of the large signal model. They do not claim to originate
those ingredients.

This is a discrete-time theory with continuous-valued observations. It
is not a continuous-time stochastic dynamics theorem. Production PR-LASM
2.3.0 is unchanged.


\chapter{Curved domains and continuous observations}\label{LSM:book:chapter-35}
\section{A matrix certificate for a curved domain}\label{LSM:a-matrix-certificate-for-a-curved-domain}

Let \(E=\mathbb R^r\) and let \(\mathbb S^d\) be the real symmetric
\(d\times d\) matrices. Fix an injective linear map

\[
\mathcal J:E\longrightarrow\mathbb S^d,
\qquad \lambda z=\operatorname{tr}\mathcal J(z).
\]

Suppose some \(e\in E\) satisfies \(\mathcal J(e)\succ0\). Define

\[
C_{\mathcal J}=\{z:\mathcal J(z)\succeq0\},
\qquad K_{\mathcal J}=\{z\in C_{\mathcal J}:\lambda z=1\}.
\]

The map \(\mathcal J\) is an embedding, not a many-to-one decoder of a
larger positive coordinate space. The cone is an inverse image of the
positive semidefinite cone. This is different from the image cone
\(V\mathbb R_+^m\) used in Proposition 39.

\Needspace{5\baselineskip}
\textbf{Proposition 50. Matrix-embedding probability certificate.} The
cone \(C_{\mathcal J}\) is closed, pointed and generating, and
\(\lambda\) is strictly positive on its nonzero elements. Its normalized
base is compact. Suppose a finite-alphabet predictive model has
\(z_0\in K_{\mathcal J}\), exact mass conservation and identities

\[
\boxed{\mathcal J B_s=\Phi_s\mathcal J,}
\qquad
\Phi_s(X)=\sum_{j=1}^{m_s} w_{sj}K_{sj}XK_{sj}^{\mathsf T},
\quad w_{sj}\geq0.
\]

Then every \(B_s\) preserves \(C_{\mathcal J}\), and every finite
observation word has a nonnegative, normalized probability. Zero maps
may use an empty sum. An invertible change of predictive coordinates
transports the certificate by

\[
B'_s=TB_sT^{-1},\qquad
\mathcal J'=\mathcal JT^{-1},\qquad
\lambda'=\lambda T^{-1}.
\]

\textbf{Proof.} Closedness follows from continuity. If both \(z\) and
\(-z\) are in the cone, \(\mathcal J(z)\) is both positive and negative
semidefinite, hence zero. Injectivity gives \(z=0\). A sufficiently
small neighborhood of \(e\) maps into positive definite matrices, so the
cone has an interior point and is generating. A nonzero positive
semidefinite matrix has positive trace. The trace-one positive
semidefinite matrices form a compact set. The image of \(\mathcal J\) is
a closed finite-dimensional subspace and its inverse on that image is
continuous, so \(K_{\mathcal J}\) is compact.

Every term \(wKXK^{\mathsf T}\) is positive semidefinite when \(X\) is.
The intertwining identity therefore proves invariance. Iteration gives
invariance of every branch state, and the mass identity normalizes every
finite observation tree. Substituting the transported quantities proves
the final statement. \(\square\)

This is a sufficient certificate language. It is not asserted that every
invariant cone has this form, or that every positive map on a matrix
subspace extends through the displayed congruence representation. The
extension requirement is stronger than positivity on the represented
states alone.

The base norm of Section 46 is also not automatically the ambient trace
norm. One always has

\[
\|\mathcal J(z)\|_1\leq\|z\|_{C_{\mathcal J}},
\]

because every positive decomposition in \(E\) gives one in the matrix
space. Equality holds when \(\mathcal J\) maps onto all of
\(\mathbb S^d\), by the positive and negative spectral parts. For a
proper matrix subspace, those spectral parts may leave the subspace.
Norm transport therefore needs its own stated hypothesis.

\section{The earlier disk obstruction has an exact finite matrix witness}\label{LSM:the-earlier-disk-obstruction-has-an-exact-finite-matrix-witness}

For the rank-three model of Section 51, put

\[
\mathcal J(m,x,y)=\frac12
\begin{pmatrix}m+x&y\\y&m-x\end{pmatrix}.
\]

Then \(\operatorname{tr}\mathcal J=m\) and

\[
\det\mathcal J=\frac{m^2-x^2-y^2}{4}.
\]

Its eigenvalues are \((m\pm\sqrt{x^2+y^2})/2\), so positivity is exactly
the Lorentz-cone condition. The normalized base is the unit disk.

\Needspace{5\baselineskip}
\textbf{Proposition 51. Exact rational matrix witness for the rotation
model.} Let \(T=\operatorname{diag}(1,R)\) with the unchanged rotation

\[
R=\begin{pmatrix}3/5&-4/5\\4/5&3/5\end{pmatrix},
\qquad W=\begin{pmatrix}2&-1\\1&2\end{pmatrix}.
\]

The rotate branch and the two probe branches from Proposition 46 have
the identities

\[
\begin{aligned}
\mathcal J(Tz)&=\frac15W\mathcal J(z)W^{\mathsf T},\\
\mathcal J(B_{\mathrm p,0}z)&=E_{00}\mathcal J(z)E_{00}^{\mathsf T},\\
\mathcal J(B_{\mathrm p,1}z)&=E_{01}\mathcal J(z)E_{01}^{\mathsf T},
\end{aligned}
\]

where \(E_{ij}\) is the corresponding \(2\times2\) matrix unit. The
other rotate branch is zero. Thus this exact model has a finite rational
matrix certificate without attenuation or a polygonal replacement.

\textbf{Proof.} Direct multiplication gives the rotate identity. The two
matrix-unit congruences retain, respectively, the upper and lower
diagonal entries of \(\mathcal J(z)\) and place them in its upper
diagonal entry. These are \((m+x)/2\) and \((m-x)/2\), exactly the two
probe masses resetting to \((1,1,0)\). Also \(W^{\mathsf T}W=5I\) and

\[
E_{00}^{\mathsf T}E_{00}+E_{01}^{\mathsf T}E_{01}=I.
\]

The actions conserve trace. The initial state embeds as
\(\operatorname{diag}(1,0)\), and \((1,0,0)\) is a strictly positive
interior witness. Proposition 50 applies. \(\square\)

Nothing has been changed in the observation law. Proposition 46 still
rules out a finite positive-generator certificate for this model. What
changes is the certificate language. A curved matrix domain and a finite
set of positive coordinate states are not the same representation class.

This use of symmetric positive semidefinite matrices is mathematical. It
neither requires quantum hardware nor claims a quantum description of
the measured system. It borrows a well-defined operator representation
rather than a physical interpretation.

\section{An observation density is not a branch probability}\label{LSM:an-observation-density-is-not-a-branch-probability}

Let \((Y,\mathcal Y)\) be a standard Borel observation space and let
\(\mu\) be a fixed probability reference measure. Keep a finite action
set. For each action \(a\), let \(B_a(y):E\to E\) be measurable and
integrable in operator norm. Let \((C,\lambda)\) be a proper probability
cone and mass functional as in Section 46. Assume that, outside one
\(\mu\)-null set for each action,

\[
B_a(y)C\subseteq C,
\qquad
\lambda\int_Y B_a(y)\,\mu(dy)=\lambda.
\]

The second identity is an equality of linear functionals. The first is a
statement on the whole cone, not only on a finite selection of states.

\Needspace{5\baselineskip}
\textbf{Proposition 52. Continuous-observation instrument semantics.}
For \(z\in K=\{z\in C:\lambda z=1\}\),

\[
\boxed{
P_a(A\mid z)=\int_A\lambda B_a(y)z\,\mu(dy),
\qquad
z^+(y)=\frac{B_a(y)z}{\lambda B_a(y)z}
}
\]

define a probability kernel and, for \(\mu\)-almost every observation
with positive denominator, a conditional state. They generate consistent
finite record laws under any common measurable history-dependent policy.
For prescribed actions, the joint observation density is

\[
\lambda B_{a_T}(y_T)\cdots B_{a_1}(y_1)z_0
\]

with respect to \(\mu^{\otimes T}\). Replacing the reference measure by
\(\mu'=r\mu\), with \(r>0\) almost everywhere and \(\int r\,d\mu=1\),
requires \(B'_a(y)=B_a(y)/r(y)\). Event probabilities and normalized
states then remain unchanged.

\textbf{Proof.} Cone positivity makes the density nonnegative and the
integrated mass identity makes its integral one. A positive-density
update lies in \(K\). Assign any fixed state of \(K\) on the common
exceptional null set and at zero-density points to obtain a measurable
version of the update. This arbitrary choice has no effect on the
probability kernel. Repeated integration, or conditioning one step at a
time, proves normalization and consistency of the finite laws.
Measurable policy factors are nonnegative and sum to one at each
history. The reference-measure transformation cancels \(r\) in each
integral and cancels \(1/r\) in the normalized update. \(\square\)

When \(\mu\) has no atoms, an individual point has probability zero even
when its density is positive. A density may exceed one. Neither fact is
a numerical pathology. Flooring a vanishing denominator changes the
model and is not part of this theorem.

A continuous density at an exact point is a version-dependent object on
null sets. The update formula specifies a convenient version wherever
its denominator is positive. Event probabilities are the invariant
statistical objects. Comparing model densities requires a common measure
or a declared change-of-measure calculation.

A representation that lacks an operator-norm integrable density needs a
more general operator-valued measure. Such representations are not
silently covered by the density hypothesis above.

\section{A finite certificate over an entire observation interval}\label{LSM:a-finite-certificate-over-an-entire-observation-interval}

Specialize to \(Y=[-1,1]\) and \(\mu(du)=du/2\). Suppose \(B_a(u)\) is a
matrix polynomial. A finite congruence certificate can prove positivity
for every \(u\), not merely for values on a grid.

For a real \(d^2\times d^2\) symmetric matrix \(F\), define

\[
[\Phi_F(X)]_{ij}=\sum_{p,q=1}^d F_{(i,p),(j,q)}X_{pq}.
\]

Vectorization orders the pairs with output index first inside each input
column. If

\[
F=\sum_\nu w_\nu\operatorname{vec}(K_\nu)
\operatorname{vec}(K_\nu)^{\mathsf T},\qquad w_\nu\geq0,
\]

then \(\Phi_F(X)=\sum_\nu w_\nu K_\nu X K_\nu^{\mathsf T}\). Every
positive semidefinite \(F\) has a real factorization of this kind. The
matrix is the usual Choi/Gram representation, with an explicit index
convention {[}R16{]}.

\Needspace{5\baselineskip}
\textbf{Proposition 53. Polynomial Gram certificate for interval
observations.} Fix the injective embedding \(\mathcal J\) and its strict
interior witness from Proposition 50. Set

\[
Z_q(u)=\begin{pmatrix}I_{d^2}\\uI_{d^2}\\\vdots\\u^qI_{d^2}\end{pmatrix}.
\]

Suppose each action has constant positive semidefinite matrices
\(G_{a,0},G_{a,1}\), with the second term omitted when \(q=0\), such
that

\[
\begin{aligned}
F_a(u)&=Z_q(u)^{\mathsf T}G_{a,0}Z_q(u)\\
&\quad +(1-u^2)Z_{q-1}(u)^{\mathsf T}G_{a,1}Z_{q-1}(u),\\
\mathcal J B_a(u)&=\Phi_{F_a(u)}\mathcal J.
\end{aligned}
\]

Assume also \(z_0\in K_{\mathcal J}\) and the exact integrated mass
identity. Then the model is internally probability-valid for every
observation value in \([-1,1]\) and every finite horizon. The
intertwining condition is finitely many coefficient equalities. The
normalization condition uses only

\[
\int_{-1}^1u^k\,\frac{du}{2}
=\begin{cases}0&k\text{ odd},\\1/(k+1)&k\text{ even}.\end{cases}
\]

\textbf{Proof.} Both Gram expressions are positive semidefinite on the
interval because \(1-u^2\geq0\) there. Factoring \(F_a(u)\) at a fixed
\(u\) gives a positive congruence map, whether or not its individual
factors depend polynomially on \(u\). The intertwining identity
therefore proves \(B_a(u)C_{\mathcal J}\subseteq C_{\mathcal J}\).
Polynomials on a compact interval are integrable. The mass identity and
Proposition 52 prove the assertion. Equality of two finite matrix
polynomials is equivalent to equality of every coefficient. The moment
values follow by elementary integration. \(\square\)

For fixed \(\mathcal J\) and fixed degree bounds, unknown entries of
\(B_a\) and the Gram matrices obey linear equalities and positive
semidefinite constraints. A convex data-fit objective therefore gives a
semidefinite feasibility or optimization problem in this fixed
representation. Jointly choosing \(\mathcal J\) changes the problem: the
intertwining equation then contains products of unknown coefficients.

This is a sufficient certificate, not a completeness theorem for all
positive kernels or all curved domains. A negative answer at one degree
is not a proof of model invalidity. Approximate Gram feasibility is also
not an exact witness. Rational factorizations with nonnegative rational
weights, or rigorous matrix-positivity bounds, must survive a separate
replay.

The reference checker accepts a factored version:

\[
\Phi_u(X)=\sum_\nu w_\nu g_\nu(u)K_\nu(u)XK_\nu(u)^{\mathsf T},
\qquad g_\nu(u)\in\{1,1-u^2\}.
\]

It checks the polynomial identities, weights, support, embedding,
initial state and normalization exactly. It does not run a general
semidefinite optimizer. The operator identity holds on the image of
\(\mathcal J\); trace preservation of an arbitrary extension outside
that image is not required.

\section{Learning from integral observations rather than point probabilities}\label{LSM:learning-from-integral-observations-rather-than-point-probabilities}

A finite polynomial family can be identified through moment operators.
This gives a continuous-observation analogue of the core-table
construction without pretending that a point has positive probability.

Assume the true operator density belongs to the declared class

\[
B_a(u)=\sum_{k=0}^D u^k C_{a,k},
\qquad
M_{a,j}=\int_{-1}^1 u^j B_a(u)\,\frac{du}{2}.
\]

Define the moment Gram matrix \(W_D\) by

\[
(W_D)_{jk}=\int_{-1}^1u^{j+k}\,\frac{du}{2},
\qquad 0\leq j,k\leq D.
\]

A history is now a prescribed action sequence with observations
constrained to measurable events of positive or zero probability. It
produces an unnormalized predictive vector \(h_j\). A future event
experiment produces a row functional \(t_i\). Choose \(r\) of each and
form

\[
H_{ij}=t_i h_j,\qquad
(G_{a,j})_{i\ell}=t_i M_{a,j}h_\ell.
\]

The index \(j\) in \(G_{a,j}\) is a moment order. Odd-order entries can
be signed expectations. They are not branch probabilities and are never
normalized as if they were.

\Needspace{5\baselineskip}
\textbf{Proposition 54. Conditional finite-moment identification.}
Suppose the true predictive dimension is \(r\), the degree bound \(D\)
is correct, and the selected \(r\times r\) core \(H\) is invertible.
Exact population values of \(H\), every \(G_{a,j}\) and the initial test
vector determine the complete polynomial operator family in test
coordinates:

\[
\boxed{
\widehat M_{a,j}=G_{a,j}H^{-1},
\qquad
\widehat C_{a,k}=\sum_{j=0}^D(W_D^{-1})_{kj}\widehat M_{a,j}.
}
\]

If every moment-operator error is at most \(\eta\) in one fixed induced
norm, then coefficient reconstruction alone obeys

\[
\|\widetilde C_{a,k}-\widehat C_{a,k}\|
\leq\eta\sum_j|(W_D^{-1})_{kj}|.
\]

\textbf{Proof.} A nonzero polynomial has strictly positive squared
integral, so its coefficient vector \(c\) satisfies
\(c^{\mathsf T}W_Dc>0\). Thus \(W_D\) is invertible. Let \(T\) have rows
\(t_i\) and \(S\) have columns \(h_j\). Since \(H=TS\) is an invertible
square core in an \(r\)-dimensional space, both \(T\) and \(S\) are
invertible. Therefore \(G_{a,j}H^{-1}=TM_{a,j}T^{-1}\). The relation
\(M_{a,j}=\sum_k(W_D)_{jk}C_{a,k}\) is exact integration of the
polynomial expansion. Invert it and conjugate by \(T\) to obtain the
displayed reconstruction. The error bound is the triangle inequality for
that same fixed linear inversion. \(\square\)

The rank and degree assumptions are not learned from the determinant of
\(H\). The determinant establishes coverage only after the ambient
predictive dimension has been justified. The moment Gram can become
poorly conditioned at higher degree. An orthogonal polynomial basis can
improve this representation without removing an incorrect degree
assumption.

The entries have an operational population meaning. A reset experiment
for an event word, with one observation weighted by \(u^j\), returns a
bounded random variable in \([-1,1]\). Its expectation is the
corresponding entry of \(G_{a,j}\). With \(N\) independent trials for
each of \(M\) predeclared entries, the same exponential-moment argument
as Proposition 35 gives the simultaneous bound

\[
\epsilon_N=\sqrt{\frac{2\log(2M/\delta)}{N}}
\]

with failure probability at most \(\delta\). The factor differs from the
Bernoulli bound because the value range is two. Dependent trajectories,
unknown resets and post-selection are not covered by that observation
protocol. This edition uses exact constructed population tables, not a
new sampling experiment.

\section{Finite moments cannot validate an unrestricted density class}\label{LSM:finite-moments-cannot-validate-an-unrestricted-density-class}

A finite observation basis is a modeling assumption. Positivity
certification does not turn that assumption into a consequence of data.

\Needspace{5\baselineskip}
\textbf{Proposition 55. A valid continuous counterexample to
finite-moment identification.} Let

\[
L_3(u)=\frac{5u^3-3u}{2},\qquad
p_\pm(u)=1\pm\epsilon L_3(u),\qquad 0<\epsilon<1.
\]

Relative to \(du/2\), these are distinct strictly positive normalized
densities with the same moments of orders zero, one and two. Their
total-variation distance is

\[
\boxed{\|p_+\mu-p_-\mu\|_{\mathrm{TV}}=\frac{13\epsilon}{40}.}
\]

Taking independent observations gives two rank-one processes whose
products of these low-order moments agree at every sequence length.

\textbf{Proof.} The endpoints and stationary points of \(L_3\) show
\(|L_3(u)|\leq1\) on the interval. Direct integration gives
\(\int u^jL_3(u)\,du/2=0\) for \(j=0,1,2\). This proves normalization
and moment agreement, while strict positivity follows from
\(\epsilon<1\). Its nonnegative root other than zero is \(\sqrt{3/5}\).
Integrating \(|L_3|\) on the two subintervals of \([0,1]\) and using
symmetry gives \(\int|L_3|\,du/2=13/40\). The density difference is
\(2\epsilon L_3\), so its half-integral is the stated distance.
Independence makes every product moment factor into the identical
one-step moments. \(\square\)

This is not a defect in Proposition 54. These densities lie outside its
quadratic class. It is a concrete reason to retain an explicit
approximation or misspecification term when increasing the observation
basis. Internal probability validity and correct density identification
remain different claims even in one predictive dimension.

\section{Exact observation bins preserve their own event law}\label{LSM:exact-observation-bins-preserve-their-own-event-law}

Let \(A_1,\ldots,A_n\) be a measurable partition of \(Y\). Define the
cell operators

\[
B_{a,j}=\int_{A_j}B_a(y)\,\mu(dy).
\]

Their role is to represent observed events. They are not point samples
of the original density.

\Needspace{5\baselineskip}
\textbf{Proposition 56. Exact bin coarsening and its posterior.} Every
\(B_{a,j}\) preserves the probability cone, and
\(\lambda\sum_jB_{a,j}=\lambda\). For a policy measurable with respect
to bin histories, the finite model \(\{B_{a,j}\}\) reproduces exactly
the binned observation law of the continuous model at every finite
horizon. On an observed cell with \(p_j=\lambda B_{a,j}z>0\),

\[
\frac{B_{a,j}z}{p_j}
=\frac1{p_j}\int_{A_j}z^+(y)\lambda B_a(y)z\,\mu(dy).
\]

\textbf{Proof.} Integrals of vectors in a closed convex cone remain in
the cone. Summing over a partition gives the original mass identity. For
a fixed bin history, its policy factors are constant across the
corresponding cells. Successive integration of the operator product
therefore gives the product of the cell operators, in the original
temporal order. This establishes every finite binned record probability.
In the posterior identity, substitute \(z^+(y)\lambda B_a(y)z=B_a(y)z\)
on the positive-density set. The remaining outputs are zero because they
lie in the cone and have zero mass. \(\square\)

The coarsened posterior is a conditional average, not the posterior at a
representative observation value. It preserves exactly the information
retained by the event history. An arbitrary policy that reacts
differently to two observations inside the same cell is not
bin-measurable, so the preceding exact-policy conclusion does not apply
to it.

Cell refinement is compatible with this semantics: a parent operator is
the sum of its children. In this setting those cell operators are
positive probability branches. This positivity is an extra property of
event integration, not a property of every signed spectral packet in the
signal archive.

\section{Averaging transformations has a covariance defect}\label{LSM:averaging-transformations-has-a-covariance-defect}

A common discretization replaces a family of matrices \(K(y)\) by its
mean and then applies \(X\mapsto KXK^{\mathsf T}\). This is generally
the wrong mean operation.

\Needspace{5\baselineskip}
\textbf{Proposition 57. Exact covariance correction for an averaged
congruence.} Let \(\nu\) be a probability measure and suppose \(K(y)\)
has square-integrable entries. Put \(\bar K=\int K(y)\,\nu(dy)\). For
every symmetric \(X\),

\[
\boxed{
\int K(y)XK(y)^{\mathsf T}\,\nu(dy)
=\bar KX\bar K^{\mathsf T}
+\int (K(y)-\bar K)X(K(y)-\bar K)^{\mathsf T}\,\nu(dy).
}
\]

The correction is positive semidefinite when \(X\succeq0\). For
\(K(u)=K_0+uK_1\), it equals
\(\operatorname{Var}_\nu(u)K_1XK_1^{\mathsf T}\). On a uniform interval
of width \(h\), the variance is \(h^2/12\).

\textbf{Proof.} Substitute \(K=\bar K+(K-\bar K)\) and expand. The two
cross terms integrate to zero. The remaining integrand is a positive
congruence when \(X\) is positive. For an affine matrix function the
centered term is \((u-\mathbb E_\nu u)K_1\), giving the variance
formula. Integrating the square of a centered uniform variable gives
\(h^2/12\). \(\square\)

This is a defect between two operations: averaging a transformation
parameter and applying its quadratic action. It is not the same as the
residual from dropping a signal packet, although both fit the
framework's requirement to record failures of commutation.

For the example in Section 66, equal bins with midpoint evaluation lose
total mass

\[
\boxed{\frac{3}{4n^2}}
\]

when there are \(n\) bins. With one bin, midpoint evaluation retains
only one quarter of the required mass. Exact integration or the
covariance correction restores the missing contribution. Dividing an
inaccurate approximation by a guessed normalizer would define another
model rather than establish this identity.

\section{A valid density approximation with a resolution bound}\label{LSM:a-valid-density-approximation-with-a-resolution-bound}

Exact coarsening gives a finite event model. To compare complete
continuous observation laws, lift that model back to a density on the
same observation space. Assume all cells have positive reference mass
and put

\[
B_a^\Pi(y)=\frac{B_{a,j}}{\mu(A_j)}\quad(y\in A_j).
\]

This uses the normalized reference measure inside each cell. It is an
explicit density approximation, not an assertion that the original
conditional density was constant there.

\Needspace{5\baselineskip}
\textbf{Proposition 58. Positive cell-average approximation.} The family
\(B_a^\Pi\) is cone preserving and mass conserving. It has exactly the
same cell operators as \(B_a\). Let \(Y\) have a metric and suppose

\[
\|B_a(y)-B_a(v)\|_{C\to C}\leq L\,d(y,v)
\]

in the induced base norm, with every cell of diameter at most \(h\).
Then

\[
\epsilon_\Pi:=\sup_{a,z\in K}\int
\|(B_a(y)-B_a^\Pi(y))z\|_C\,\mu(dy)
\leq Lh.
\]

For equal intervals of width \(h\) under the uniform measure on
\([-1,1]\), the sharper bound is \(\epsilon_\Pi\leq Lh/3\).

\textbf{Proof.} Each cell average is a positive scalar multiple of a
positive cell operator. Its integral over that cell is \(B_{a,j}\),
proving both normalization and exact cell agreement. For \(y\in A_j\),
subtract the cell average and use

\[
\|(B_a(y)-B_a^\Pi(y))z\|_C
\leq\frac{L}{\mu(A_j)}\int_{A_j}d(y,v)\,\mu(dv),
\]

because \(\|z\|_C=1\). Integrating over \(y\) proves the diameter bound.
Two independent uniform points in an interval of width \(h\) have mean
absolute separation \(h/3\), obtained by integrating over its two
triangles. Applying this inside every equal interval gives the sharper
bound. \(\square\)

The Lipschitz constant is a declared operator bound. It is not a fit
residual inferred from adjacent samples. The comparison needs a common
domain and a common reference measure. If two kernels are only close in
Euclidean coefficient norm, the relevant norm-conversion constant must
be included.

In the full \(2\times2\) symmetric-matrix example, the pullback base
norm is

\[
\|(m,x,y)\|_C=\max\{|m|,\sqrt{x^2+y^2}\}=\|\mathcal J(m,x,y)\|_1.
\]

There \(\|K(u)\|_2\leq2\) and \(\|K(u)-K(v)\|_2\leq(3/2)|u-v|\). The
identity

\[
K(u)XK(u)^{\mathsf T}-K(v)XK(v)^{\mathsf T}
=(K(u)-K(v))XK(u)^{\mathsf T}
+K(v)X(K(u)-K(v))^{\mathsf T}
\]

and \(\|AXB\|_1\leq\|A\|_2\|X\|_1\|B\|_2\) give \(L=6\). The rotated
action uses an orthogonal matrix before this map and has the same bound.
Equal bins therefore give \(\epsilon_\Pi\leq4/n\), capped by two for two
valid instruments. The inequality is sufficient, not asserted sharp.

\section{Full continuous record laws inherit the stability argument}\label{LSM:full-continuous-record-laws-inherit-the-stability-argument}

The base-norm proof from Chapter~\ref{LSM:book:chapter-34} has an integral version. Its hypotheses
must be supplied rather than inferred by replacing a finite sum with an
integral in a program.

\Needspace{5\baselineskip}
\textbf{Proposition 59. Continuous-record total-variation bound.} Let
\(B_a(y)\) and \(\widetilde B_a(y)\) satisfy Proposition 52 on a common
\((C,\lambda,\mu)\), and let the two models use the same measurable
policy on every common observation history. Suppose their initial states
satisfy \(\|z_0-\widetilde z_0\|_C\leq e_0\) and their time-\(t\)
instruments satisfy

\[
\sup_{a,z\in K}\int\|(B_{a,t}(y)-\widetilde B_{a,t}(y))z\|_C\,\mu(dy)
\leq\epsilon_t.
\]

The length-\(T\) record laws, including actions when the policy is
randomized, obey

\[
\boxed{
\|P_T-\widetilde P_T\|_{\mathrm{TV}}
\leq\min\left\{1,\frac{e_0+\sum_{t=0}^{T-1}\epsilon_t}{2}\right\}.
}
\]

No lower bound on each observation density is required. For identical
initial states and the equal-bin approximation in Section 64, the
example has the sufficient bound \(\min(1,2T/n)\).

\textbf{Proof.} For any positive decomposition \(x=u-v\), cone
preservation and mass conservation give

\[
\int\|B_a(y)x\|_C\,\mu(dy)
\leq\int\lambda B_a(y)(u+v)\,\mu(dy)
=\lambda u+\lambda v.
\]

Taking the infimum proves integrated instrument contraction. At each
continuous history let \(v_h,\widetilde v_h\) be the unnormalized state
densities, including the policy factors. Split the next difference as

\[
B_a(y)(v_h-\widetilde v_h)
+(B_a(y)-\widetilde B_a(y))\widetilde v_h.
\]

Integrate the first term using contraction. The second is bounded by
\(\epsilon_t\lambda\widetilde v_h\) after normalizing only its nonzero
positive mass. Its integral over histories is at most \(\epsilon_t\).
Tonelli's theorem applies to these nonnegative norm bounds, and operator
integrability supplies well-defined finite-dimensional integrals. The
integrated state discrepancy therefore satisfies
\(D_{t+1}\leq D_t+\epsilon_t\). Finally \(|\lambda x|\leq\|x\|_C\)
bounds the record-density difference, and total variation is half its
absolute integral. \(\square\)

This result compares one common policy, not two policies optimized
separately. It concerns the whole law and can become uninformative at
long horizons. It is not a uniform posterior contraction theorem.

For a measurable event \(H\) of probability \(p=P_T(H)>0\), the
restriction-and-normalization argument of Proposition 44 works for
measures as well as finite vectors. When both conditional laws exist,
their total-variation distance is at most \(\min(1,\delta/p)\) if the
full-law distance is at most \(\delta\). The event probability is
invariant under a reference-measure change. A density value at one point
is not a substitute for that event probability.

The full-law comparison between \(B\) and \(B^\Pi\) permits any common
measurable policy because both are continuous-output models on the same
space. The exact equivalence to the finite bin model in Proposition 56
instead requires a bin-measurable policy. These are different
conclusions with different observation information.

\section{A complete continuous-observation example}\label{LSM:a-complete-continuous-observation-example}

Use the same three-coordinate matrix embedding as Section 57. Set

\[
K(u)=\frac12\begin{pmatrix}1+3u&0\\0&1-3u\end{pmatrix},
\qquad -1\leq u\leq1,
\qquad \mu(du)=du/2.
\]

The direct action uses \(X\mapsto K(u)XK(u)^{\mathsf T}\). The rotated
action first applies the unchanged rotation congruence from Section 57,
giving

\[
X\longmapsto\frac15 K(u)WXW^{\mathsf T}K(u)^{\mathsf T}.
\]

Both are exactly normalized because
\(\int K(u)^{\mathsf T}K(u)\,du/2=I\) and \(W^{\mathsf T}W=5I\). In
predictive coordinates, write

\[
\begin{gathered}
a(u)=\frac{1+9u^2}{4},\qquad
b(u)=\frac{3u}{2},\qquad
c(u)=\frac{1-9u^2}{4},\\[3pt]
B_0(u)=\begin{pmatrix}a(u)&b(u)&0\\b(u)&a(u)&0\\0&0&c(u)\end{pmatrix},
\qquad B_1(u)=B_0(u)T.
\end{gathered}
\]

Although \(c(u)\) is negative for part of the interval, both actions
preserve the disk cone. Entrywise sign is not the positivity criterion
in these coordinates. The initial state is \(z_0=(1,0,0)^{\mathsf T}\),
or the matrix \(I/2\).

\subsection{66.1 Event tables and moment
reconstruction}\label{LSM:event-tables-and-moment-reconstruction}

Let the three future tests be the empty test, a nonnegative observation
after action zero, and a nonnegative observation after action one. Their
rows form

\[
Q=\begin{pmatrix}
1&0&0\\
1/2&3/8&0\\
1/2&9/40&-3/10
\end{pmatrix}.
\]

Use histories consisting of the empty history, a nonnegative observation
under action zero, and that history followed by a nonnegative
observation under action one. Their unnormalized columns are

\[
S=\begin{pmatrix}
1&1/2&107/320\\
0&3/8&3/10\\
0&0&-3/40
\end{pmatrix}.
\]

The learner is supplied only the resulting integral tables and the
declared class. Its core is

\[
\boxed{
H=\begin{pmatrix}
1&1/2&107/320\\
1/2&25/64&179/640\\
1/2&107/320&823/3200
\end{pmatrix},\qquad\det H=\frac{81}{25600}.
}
\]

The core is not required to be symmetric. Its rows are tests and its
columns are histories, which are different experiments. The three moment
orders use

\[
W_2=\begin{pmatrix}1&0&1/3\\0&1/3&0\\1/3&0&1/5\end{pmatrix},
\qquad
W_2^{-1}=\begin{pmatrix}9/4&0&-15/4\\0&3&0\\-15/4&0&45/4\end{pmatrix},
\qquad\det W_2=\frac4{135}.
\]

Proposition 54 reconstructs \(QB_a(u)Q^{-1}\) from \(H\) and the six
inserted-moment tables. The code that performs this reconstruction
imports no generating matrices. The separate generator produces the
tables and is used only in the verification suite. These are population
values from an exact example, not learned sample estimates.

\subsection{66.2 Two densities and a genuine event
probability}\label{LSM:two-densities-and-a-genuine-event-probability}

For a direct observation value \(u=1/2\), the density and updated state
are

\[
p_1=\frac{13}{16},\qquad z_1=(1,12/13,0)^{\mathsf T}.
\]

If the rotated action next gives \(u=-1/2\), they are

\[
p_2=\frac{413}{1040},\qquad
z_2=(1,-312/413,-240/413)^{\mathsf T}.
\]

Their product \(413/1280\) is a joint density with respect to
\((du/2)^{\otimes2}\), not the probability of those two exact values.
The two point values each have zero probability in this nonatomic model.

By contrast, the event \(u_1\geq0\) under action zero followed by
\(u_2\geq0\) under action one has actual probability

\[
\boxed{\frac{107}{320},}
\]

and its conditional state is \((1,96/107,-24/107)^{\mathsf T}\). Exact
bin operators produce this value without quadrature.

The final point-conditioned state remains strictly inside the disk:

\[
\det\mathcal J(z_2)=\frac{15625}{682276}>0.
\]

A pure boundary state can instead have zero density at isolated
observations. The checker rejects conditioning on such a zero-density
branch. It does not manufacture a posterior by flooring its denominator.

\subsection{66.3 The missing covariance is exactly
calculable}\label{LSM:the-missing-covariance-is-exactly-calculable}

On an interval with midpoint \(v\) and width \(h\), the average matrix
instrument is

\[
\mathbb E[K(u)XK(u)^{\mathsf T}]
=K(v)XK(v)^{\mathsf T}
+\frac{h^2}{12}K_1XK_1^{\mathsf T},
\qquad K_1=\frac32\operatorname{diag}(1,-1).
\]

The correction is itself a positive congruence. The cell operator
multiplies this mean by the cell mass \(h/2\). Equal bins have width
\(2/n\), so omitting the covariance loses trace \(3/(4n^2)\) over the
full partition. The reference calculation checks both the exact matrix
identity and this mass defect for several rational partitions.

\section{A declared template can be recognized without being rediscovered}\label{LSM:a-declared-template-can-be-recognized-without-being-rediscovered}

General curved-domain learning remains difficult. There is a useful
finite case between supplying a finished coordinate certificate and
claiming unrestricted domain discovery.

Suppose a structural library supplies a certified template with
polynomial coefficients \(C_{a,k}\), mass row \(\lambda\) and embedding
\(\mathcal J\). Suppose observed population tables produce coefficients
\(\widehat C_{a,k}\) and mass row \(\widehat\lambda\) of the same
dimension.

\Needspace{5\baselineskip}
\textbf{Proposition 60. Exact recognition of a fixed certified
template.} The equations

\[
\boxed{
\widehat C_{a,k}T=TC_{a,k},\qquad
\widehat\lambda T=\lambda
}
\]

are linear in \(T\). Every invertible solution transports the template
certificate to the fitted coordinates by
\(\widehat{\mathcal J}=\mathcal JT^{-1}\). If the fitted coefficient
family is observable, meaning the rows \(\widehat\lambda\widehat C_w\)
span the dual space, the solution is unique. The fitted initial state
must still be tested in the transported normalized cone.

\textbf{Proof.} All coefficient matrices are fixed, so the displayed
constraints are linear. An invertible solution yields
\(\widehat B_a(u)=TB_a(u)T^{-1}\) coefficient by coefficient.
Substitution in the template intertwining identity proves the
transported certificate. To prove uniqueness, let \(D\) be the
difference of two solutions. It satisfies \(\widehat C_{a,k}D=DC_{a,k}\)
and \(\widehat\lambda D=0\). Therefore
\(\widehat\lambda\widehat C_wD=0\) for every coefficient word.
Observability makes every column of \(D\) zero. An initial state outside
the transported cone is not repaired by this operator identity, which
explains the separate check. \(\square\)

For the worked example, the linear system has the unique solution
\(T=Q\) from Section 66. The learner receives the population tables and
a fixed structural template. It reconstructs the operators, solves the
linear match, transports the embedding and exports a certificate.
Independent replay then needs only that certificate. It does not refit,
search, or import the template or generator.

This is template recognition, not unrestricted geometry discovery. The
supplied template includes the specific operator family. An incompatible
model is rejected rather than changed to fit the template. The reference
routine also withholds a result for an underdetermined or singular
match. Noise would require a declared approximate match and a separate
validity or correction argument. This edition's recognition example is
exact.

For fixed embeddings, polynomial Gram optimization offers a broader
construction route than fixed-template matching. That convex proposal
problem is stated in Section 59 but is not implemented as a general
solver here. The independent checker remains useful regardless of
whether a proposal came from a template, symbolic derivation or an
optimization routine.

\section{The combined continuous polynomial sector}\label{LSM:the-combined-continuous-polynomial-sector}

\Needspace{5\baselineskip}
\textbf{Proposition 61. Conditional continuous realization with a finite
domain certificate.} Assume a controlled continuous-observation process
has the declared finite predictive dimension and polynomial
operator-density degree, a fixed uniform reference measure on
\([-1,1]\), and an invertible complete event-test core. Exact moment
tables identify its operator density in test coordinates by Proposition
54.

If the fitted family and initial state admit a checked matrix/Gram
certificate of Proposition 53, the supplied realization is internally
valid at every observation value and finite horizon. A successful
invertible fixed-template recognition is one sufficient way to construct
such a certificate. Exact coordinate transport and complete packet
synthesis preserve the represented event probabilities. Individual
signed state packets still need the separate positivity condition before
probabilistic interpretation.

Exact observation-cell integration gives the precise finite event model
for bin-measurable policies. Positive cell-average density
approximations have the resolution bound of Proposition 58. Two
certified models on a common cone and common reference measure satisfy
Proposition 59's full-law bound whenever the declared instrument
discrepancy is controlled.

\textbf{Proof.} Moment identification uses the true dimension and degree
hypotheses to apply Proposition 54. The independent Gram identity
establishes positivity on the whole declared interval, while the exact
integral mass identity supplies normalization through Proposition 52.
Proposition 60 supplies a transported witness in the template case.
Coordinate cancellation and full packet synthesis are exact identities.
Proposition 56 proves the cell-model statement. Propositions 58 and 59
provide the two approximation conclusions with their common-domain and
common-policy assumptions unchanged. \(\square\)

For a noisy fitted model, the certificate alone proves only internal
validity. A guarantee about the generating process requires a valid
uncertainty set covering it. A common-domain certificate over a convex
uncertainty family can be checked at complete-family vertices with the
same embedding, because the operator identities, Gram positivity and
mass constraints survive convex combination. This does not justify an
uncertainty set that excludes the true degree, rank or dynamics.

The construction leaves visible what is often collapsed into the phrase
signal model: a represented state, an operator-valued observation
density, a reference measure, a valid probability domain, an observation
protocol and an approximation budget. Each is needed for a different
reason.

\section{From continuous semantics to discovery}\label{LSM:from-continuous-semantics-to-discovery}

Chapter~\ref{LSM:book:chapter-35} gives event-based probability semantics, exact binning and
conditional finite-moment reconstruction. Fixed-template recognition in
Proposition 60 is not unrestricted geometry discovery. Chapters~\ref{LSM:book:chapter-36}--\ref{LSM:book:chapter-37} add
restricted discovery and robust verification with their own assumptions.
The continuous reference measure and the observation basis remain
explicit inputs.


\label{LSM:part-vi.-certificate-discovery-approximation-and-sequential-evidence}

A signal model must distinguish a representation that preserves
probability from a representation that predicts the observed process
accurately. Chapter~\ref{LSM:book:chapter-35} gave an exact curved certificate and a way to
recognize a supplied operator template. It also made the
polynomial-degree assumption explicit. This block advances both
boundaries without identifying them with each other.

The first development is a restricted discovery procedure. It recovers
an ellipsoidal probability domain from fitted polynomial operators,
without a supplied template, matrix embedding or hidden generator. The
second is a positive approximation theorem for operator densities
outside a finite polynomial class. The third is sequential evidence
under dependent observations, including the cost of action coverage.
Together they separate structural validity, approximation error and
observable disagreement.

The invariant-form and quadratic-certificate ingredients belong to
classical systems theory {[}R8, R13{]}. Positive polynomial averaging
uses the classical Bernstein--Durrmeyer kernel {[}R19{]}. The sequential
argument is an elementary martingale confidence sequence, not a new
concentration principle {[}R20{]}. The contribution of this block is
their connection to the large signal model's explicit representation
contracts and to independently replayable examples. Propositions 62--70
have proofs below. No proof kernel or production training system is
invoked.


\chapter{Certificate discovery, approximation, and evidence}\label{LSM:book:chapter-36}
\section{An ellipsoidal certificate with an orientation test}\label{LSM:an-ellipsoidal-certificate-with-an-orientation-test}

Work in mass-first coordinates \(z=(m,x)\in\mathbb R\times\mathbb R^d\),
with \(\lambda=(1,0,\ldots,0)\). For a center \(c\in\mathbb R^d\) and a
symmetric positive definite matrix \(H\), define

\[
C_{c,H}=\{(m,x):m\geq0,\ (x-mc)^{\mathsf T}H(x-mc)\leq m^2\}.
\]

Its normalized base is an ellipsoid. Set

\[
S_c=\begin{pmatrix}1&0\\-c&I\end{pmatrix},\qquad
Q=S_c^{\mathsf T}\begin{pmatrix}1&0\\0&-H\end{pmatrix}S_c.
\]

Then \(C_{c,H}=\{z:m\geq0,\ z^{\mathsf T}Qz\geq0\}\). The cone is proper
and generating, and \(m\) is strictly positive on every nonzero cone
element. These facts follow by an invertible linear transformation to
the Lorentz cone. In particular, an output with zero mass in the cone is
the zero vector.

Let \(B_a(u)\) be polynomial operator densities on \([-1,1]\), relative
to \(du/2\). Write

\[
\lambda B_a(u)=(b_{a,0}(u),b_a(u)^{\mathsf T}),\qquad
\beta_a(u)=b_{a,0}(u)+b_a(u)^{\mathsf T}c.
\]

\Needspace{5\baselineskip}
\textbf{Proposition 62. Oriented quadratic probability certificate.}
Suppose that for each action and every \(u\in[-1,1]\) there are
\(\tau_a(u)\geq0\) and \(R_a(u)\succeq0\) satisfying

\[
\boxed{B_a(u)^{\mathsf T}QB_a(u)-\tau_a(u)Q=R_a(u),}
\]

and that

\[
\boxed{\beta_a(u)\geq0,\qquad
\beta_a(u)^2-b_a(u)^{\mathsf T}H^{-1}b_a(u)\geq0.}
\]

If \(z_0\in C_{c,H}\) has mass one and
\(\lambda\int B_a(u)\,du/2=\lambda\), then every operator preserves
\(C_{c,H}\). The supplied model has valid probabilities at every finite
horizon. Polynomial sum-of-squares witnesses for the scalar inequalities
and for \(R_a\) make this a finite certificate.

\textbf{Proof.} For \(z\in C_{c,H}\), the quadratic identity gives

\[
(B_a z)^{\mathsf T}Q(B_a z)
=\tau_a z^{\mathsf T}Qz+z^{\mathsf T}R_a z\geq0.
\]

For a nonzero \(z\), write \(x=m(c+w)\) with \(w^{\mathsf T}Hw\leq1\).
Its output mass is \(m(\beta_a+b_a^{\mathsf T}w)\). Cauchy--Schwarz in
the \(H\) inner product bounds this below by

\[
m\left(\beta_a-\sqrt{b_a^{\mathsf T}H^{-1}b_a}\right)\geq0.
\]

Thus both inequalities defining the desired sheet of the quadratic cone
hold. The zero input causes no difficulty. Integrated mass conservation
gives normalization, and induction gives the finite-horizon conclusion.
For polynomial data, identities between coefficients and nonnegative
sums of polynomial squares establish the displayed hypotheses throughout
the interval. Terms multiplied by \(1-u^2\) are also nonnegative on that
interval. \(\square\)

The orientation test is essential. The matrix \(-I\) preserves every
quadratic form, but maps positive mass to negative mass. A quadratic
inequality without its mass orientation can certify the wrong sheet.
Likewise, testing finitely many values of \(u\) is not a substitute for
a polynomial identity and its positivity witness.

This is a sufficient language, not a completeness theorem for all cones.
The reference checker accepts rational scalar and matrix sum-of-squares
witnesses. Its discovery routine uses the narrower conformal case
\(R_a=0\), but the checker does not confuse that search restriction with
the theorem's scope.

\section{Geometry can be proposed from the operators themselves}\label{LSM:geometry-can-be-proposed-from-the-operators-themselves}

An operator identity can reveal a useful coordinate geometry without
specifying that geometry in advance. Suppose two fitted operators
evaluated at an anchor \(u_*\) give an invertible ratio

\[
F=B_b(u_*)B_a(u_*)^{-1},\qquad \lambda F=\lambda.
\]

This ratio is an algebraic object. It need not be a physically available
action or a positive map by itself. In mass-first coordinates write

\[
F=\begin{pmatrix}1&0\\b&A\end{pmatrix}.
\]

If \(I-A\) is invertible, its fixed center is \(c=(I-A)^{-1}b\). The
centered spatial action is \(A\). A candidate invariant metric obeys

\[
A^{\mathsf T}HA=H.
\]

With \(A\) fixed, this is a linear equation in the independent entries
of symmetric \(H\). No coordinate template appears in that equation.

\Needspace{5\baselineskip}
\textbf{Proposition 63. Conditional operator-only ellipsoid discovery.}
Suppose the symmetric solution space of \(A^{\mathsf T}HA=H\) is
one-dimensional and contains a positive definite matrix \(H_0\). Let a
rank-one fitted branch at a declared observation have a nonzero-mass
image, normalized to \(v=(1,x_v)\). If \(x_v\neq c\), define

\[
H=\frac{H_0}{(x_v-c)^{\mathsf T}H_0(x_v-c)}.
\]

This gives a unique positive definite metric in that solution space for
which \(v\) is on the normalized boundary. If this proposed geometry
passes every hypothesis of Proposition 62 for the complete fitted family
and its initial state, it is an internally valid probability-domain
certificate. The proposal requires no supplied template. Failure at any
step withholds this certificate and does not establish model invalidity.

\textbf{Proof.} Invertibility of \(I-A\) gives the unique center.
One-dimensionality makes every symmetric invariant form a scalar
multiple of \(H_0\). The nonzero displacement has strictly positive
squared \(H_0\) length, so the normalization is well-defined, positive
and independent of the initially chosen positive scale. The final
acceptance test is exactly Proposition 62. Neither invariance under the
ratio nor the choice of a boundary image alone is asserted to certify
the full family. \(\square\)

The rank-one image proposes a scale. It does not prove that every
possible probability domain must place that image on its boundary. Some
valid models have many admissible domains. The procedure can also fail
because there is no suitable ratio, because invariant forms are not
unique, because no rational rank-one anchor is found, or because the
chosen ellipsoid fails another branch. These are explicit search
outcomes.

\subsection{71.1 The continuous example no longer needs its
template}\label{LSM:the-continuous-example-no-longer-needs-its-template}

Use the same exact population moment tables as Section 66. The existing
moment learner first reconstructs the two polynomial operators in test
coordinates. The new discovery routine receives those operators, their
mass row and initial vector. It does not read
\path{continuous_template.json} or import the population generator.

At \(u_*=0\), the operator ratio gives the required invariant-form
equation. A rational rank-one observation supplies a boundary image. The
recovered center and metric are

\[
\boxed{
c=\begin{pmatrix}1/2\\1/2\end{pmatrix},\qquad
H=\begin{pmatrix}100/9&-20/3\\-20/3&100/9\end{pmatrix}.
}
\]

The two metric eigenvalues are \(40/9\) and \(160/9\). For both actions
the exact certificate is

\[
\begin{aligned}
B_a(u)^{\mathsf T}QB_a(u)&=\frac{(1-9u^2)^2}{16}Q,\\
\beta_a(u)&=\frac{1+9u^2}{4},\\
\beta_a(u)^2-b_a(u)^{\mathsf T}H^{-1}b_a(u)
&=\frac{(1-9u^2)^2}{16}.
\end{aligned}
\]

Every scalar on the right has an exact rational sum-of-squares
representation. Integrated mass is unchanged. The fitted initial state
is the ellipsoid center, and the two densities \(13/16\) and
\(413/1040\) from Section 66 are unchanged. This is a new certificate
construction for the same model, not a change to its predictions.

The exported JSON contains the fitted operators, center, metric, initial
state and all witnesses. Independent replay needs only that JSON. The
discovery calculation is also checked on three additional rational
coefficient families and affine coordinate changes. Those finite
examples exercise the stated sector and do not establish unrestricted
discovery.

The procedure currently uses three predictive coordinates, a
two-dimensional ellipsoid, a finite rational anchor list, a
one-dimensional invariant-form space and rational rank-one roots. It
proposes only exactly conformal polynomial branches whose scalar
nonnegativity can be resolved by its factor-based sum-of-squares
routine. General noisy or jointly optimized curved-domain learning
remains outside this implementation.

\section{Approximate an unknown density without assuming its exact degree}\label{LSM:approximate-an-unknown-density-without-assuming-its-exact-degree}

Chapter~\ref{LSM:book:chapter-35} identified a polynomial operator density under an exact degree
assumption. A different question is how to approximate a valid density
that is not a polynomial at all.

In this section observations lie in \([0,1]\) with Lebesgue reference
measure. This is an explicit change from the preceding \([-1,1]\)
convention. Under \(v=(u+1)/2\), the probability measure \(du/2\)
becomes \(dv\). A Lipschitz constant with respect to \(u\) becomes twice
as large with respect to \(v\).

Write the Bernstein basis and the classical Bernstein--Durrmeyer kernel
as

\[
b_{n,k}(x)=\binom nk x^k(1-x)^{n-k},\qquad
D_n(x,v)=(n+1)\sum_{k=0}^n b_{n,k}(x)b_{n,k}(v).
\]

Both variables range over \([0,1]\). The kernel is nonnegative and
integrates to one in either variable. Let \(B_a(v)\) be any
operator-norm integrable density preserving a common probability cone
\(C\), with mass functional \(\lambda\) and normalized base \(K\).
Define

\[
B_a^{[n]}(x)=\int_0^1D_n(x,v)B_a(v)\,dv
=\sum_{k=0}^n b_{n,k}(x)C_{a,n,k},
\]

where

\[
C_{a,n,k}=(n+1)\int_0^1b_{n,k}(v)B_a(v)\,dv.
\]

\Needspace{5\baselineskip}
\textbf{Proposition 64. Valid positive polynomial approximation.} Every
\(B_a^{[n]}\) preserves \(C\), has degree at most \(n\), and satisfies
exact integrated mass conservation. If

\[
\|B_a(v)-B_a(x)\|_{C\to C}\leq L|v-x|
\]

for all actions, then

\[
\boxed{
\sup_{a,z\in K}\int_0^1
\|(B_a^{[n]}(x)-B_a(x))z\|_C\,dx
\leq\frac{L}{\sqrt{3(n+2)}}.
}
\]

With identical initial states, the corresponding length-\(T\) laws under
a common policy differ in total variation by at most

\[
\boxed{\min\left(1,\frac{TL}{2\sqrt{3(n+2)}}\right).}
\]

No finite true-degree hypothesis is needed.

\textbf{Proof.} Nonnegative integration preserves the closed convex
cone. Since \(\int_0^1b_{n,k}=1/(n+1)\) and \(\sum_k b_{n,k}=1\), the
kernel is doubly stochastic. Its integral in \(x\) gives
\(\int B_a^{[n]}=\int B_a\), proving exact mass conservation.

For the error, subtract \(B_a(x)\) inside the kernel integral, apply the
Lipschitz inequality and use \(\|z\|_C=1\). The integrated discrepancy
is at most \(L\mathbb E|V-X|\) under the joint probability density
\(D_n(x,v)\). This joint law can be constructed by first choosing \(J\)
uniformly from \(\{0,\ldots,n\}\), and then drawing \(X,V\)
independently, conditionally on \(J=k\), with density \((n+1)b_{n,k}\).

That conditional beta law has variance

\[
\operatorname{Var}(X\mid J=k)
=\frac{(k+1)(n-k+1)}{(n+2)^2(n+3)}.
\]

The average of \((k+1)(n-k+1)\) over \(k=0,\ldots,n\) is
\((n+2)(n+3)/6\), obtained by summing the first two integer powers.
Hence

\[
\mathbb E(V-X)^2
=2\mathbb E\operatorname{Var}(X\mid J)
=\frac1{3(n+2)}.
\]

Cauchy--Schwarz gives the claimed integrated bound. The whole-law
conclusion is Proposition 59 applied to the two valid instruments on the
same domain. \(\square\)

This construction approximates the operator density itself. It does not
average a congruence factor and omit its covariance. Its coefficients
are positive operators even when their entries in predictive coordinates
are signed.

Operationally, it replaces the raw observation \(v\) by a noisy
observation \(x\) drawn from \(D_n(x,v)\). It preserves total mass and
underlying unconditioned evolution, not the entire original observation
law. That changed law is what the bound measures. A prediction policy
must be compared as the same function of the reported history.

\subsection{72.1 A nonpolynomial density and a failed ordinary
truncation}\label{LSM:a-nonpolynomial-density-and-a-failed-ordinary-truncation}

Consider the valid scalar density

\[
p(v)=\frac12+2\left|v-\frac12\right|.
\]

It integrates to one, takes values between \(1/2\) and \(3/2\), and has
Lipschitz constant two. It is not a single polynomial on \([0,1]\).
Proposition 64 nevertheless supplies normalized positive polynomial
approximations at every degree, with an explicit approximation budget.

An ordinary orthogonal truncation need not preserve that validity. The
valid density \(p(v)=3v^2\) has degree-one \(L^2\) projection

\[
q(v)=3v-\frac12,
\]

which is negative near zero. The positive construction at degree one
instead gives

\[
p^{[1]}(v)=v+\frac12.
\]

The two approximations solve different mathematical problems. Neither
has permission to call its approximation error zero. The reference
checks derive both expressions exactly.

\section{Keep approximation bias separate from coefficient error}\label{LSM:keep-approximation-bias-separate-from-coefficient-error}

The positive approximation theorem assumes access to exact integral
coefficients. A fitted model introduces another discrepancy. Its
coefficient estimates need their own validity constraints.

\Needspace{5\baselineskip}
\textbf{Proposition 65. Certified coefficient error and approximation
bias.} Suppose every fitted coefficient \(\widetilde C_{a,n,k}\)
preserves the same cone \(C\), and

\[
\lambda\frac1{n+1}\sum_{k=0}^n\widetilde C_{a,n,k}=\lambda.
\]

Then \(\widetilde B_a^{[n]}(x)=\sum_k b_{n,k}(x)\widetilde C_{a,n,k}\)
is internally valid. If

\[
\sup_{z\in K}\|(C_{a,n,k}-\widetilde C_{a,n,k})z\|_C
\leq\eta_{a,k},
\]

its integrated discrepancy from the original Lipschitz density is at
most

\[
\boxed{
\frac{L}{\sqrt{3(n+2)}}+
\max_a\frac1{n+1}\sum_{k=0}^n\eta_{a,k}.
}
\]

A known additional model-correction discrepancy is added to this budget,
not absorbed into the approximation term.

\textbf{Proof.} Nonnegative Bernstein weights preserve cone positivity,
and their integrals prove normalization. For any \(z\in K\), the
triangle inequality gives

\[
\int\|\textstyle\sum_k b_{n,k}(x)(C_{a,n,k}-\widetilde C_{a,n,k})z\|_C\,dx
\leq\frac1{n+1}\sum_k\eta_{a,k}.
\]

Add the approximation error from Proposition 64. A separate correction
is handled by another application of the triangle inequality.
\(\square\)

This is a conditional approximation-and-estimation theorem, not a claim
that more coefficients always improve a finite-data estimator.
Increasing \(n\) reduces the stated bias bound while changing the number
and conditioning of the quantities to be estimated. A cone witness
proves that the resulting object is a probability model. It does not
prove that the coefficient error budgets cover the generating process.

The moments needed for \(C_{a,n,k}\) are bounded Bernstein-weighted
experiments. Their interpretation as full operator coefficients still
requires suitable predictive-state coverage or a justified observable
core. The next section provides sequential evidence for conditional
experiments, not an automatic replacement for those identification
assumptions.

\section{Evidence from dependent observations}\label{LSM:evidence-from-dependent-observations}

Let \(\mathcal F_{t-1}\) denote the available information immediately
before observing the next measurement, including the chosen action,
current model and any test selected using past data. A predictable test
is a measurable function \(\phi_{t,j}\) of the upcoming measurement,
chosen from this information, with values in \([0,1]\). Put

\[
X_{t,j}=\phi_{t,j}(Y_t),\quad
m_{t,j}=\mathbb E[X_{t,j}\mid\mathcal F_{t-1}],\quad
p_{t,j}=\text{the model's prediction for that test}.
\]

The model prediction is also predictable. It may change as the model
learns. The true conditional mean \(m_{t,j}\) may depend on the entire
past. No independence, stationarity or known reset state is assumed in
this paragraph.

Choose budgets \(\alpha_j>0\) for a finite or countable prospectively
specified family of test rules, with \(\sum_j\alpha_j\leq\delta\).
Define

\[
R_{t,j}=\sum_{s=1}^t(X_{s,j}-p_{s,j}),\qquad
D_{t,j}=\sum_{s=1}^t(m_{s,j}-p_{s,j}).
\]

\Needspace{5\baselineskip}
\textbf{Proposition 66. Time-uniform predictable-residual certificate.}
With probability at least \(1-\delta\), for every time \(t\geq1\) and
every registered rule \(j\),

\[
\boxed{
|R_{t,j}-D_{t,j}|\leq
b_{t,j},\qquad
b_{t,j}=\sqrt{\frac t2\log\frac{2t(t+1)}{\alpha_j}}.
}
\]

For variables taking values in an interval of width \(c\), multiply the
boundary by \(c\). The simultaneous guarantee remains valid when the
observation time is chosen from the data.

\textbf{Proof.} For a random variable \(X\in[0,1]\) with conditional
mean \(m\), its conditional centered logarithmic moment-generating
function has value and first derivative zero at the origin. Its second
derivative is the variance under exponential tilting. A variable
supported on \([0,1]\) has variance at most \(1/4\), because its squared
distance from \(1/2\) is at most \(1/4\). Integrating twice gives

\[
\mathbb E[e^{\theta(X-m)}\mid\mathcal F]\leq e^{\theta^2/8}.
\]

Iterating this conditional inequality gives the same exponential bound
for the martingale sum \(\sum_{s\leq t}(X_{s,j}-m_{s,j})\). Optimizing
the exponential Markov inequality gives the two-sided bound
\(2\exp(-2b^2/t)\). At the displayed \(b_{t,j}\) this is
\(\alpha_j/[t(t+1)]\). Sum over \(t\), using
\(\sum_{t\geq1}1/[t(t+1)]=1\), and then over \(j\). The resulting event
holds at every time, so it also holds at a data-selected time. Scaling
the variable proves the range version. \(\square\)

This elementary boundary pays a logarithmic time-union cost. It is not
claimed to be optimal. Stronger confidence-sequence constructions are
available in the literature {[}R20{]}. The proof here is supplied
because the exact assumptions and the estimated object matter more than
a sharp constant in this foundation.

The certificate estimates accumulated \textbf{conditional predictive
drift}, not one fixed unconditional population probability. A
prospectively chosen rule may use history, but an analyst cannot inspect
an entire data sequence and then retroactively choose an unrestricted
test with the same budget. New rules can be registered with fresh
budgets and applied to future observations, using elapsed time after
registration.

\section{Actions need coverage, not only more observations}\label{LSM:actions-need-coverage-not-only-more-observations}

For intervention evidence use the information \(\mathcal H_{t-1}\) just
before choosing the action. Let the action be randomized from a known
history-dependent policy \(\pi_t\). The true next-observation kernel
under action \(a\) is \(P_t(\cdot\mid\mathcal H_{t-1},a)\).
Randomization uses no unrecorded future information. For a predictable
test \(\phi_t\in[0,1]\), define

\[
\mu_t^a=\int\phi_t(y)P_t(dy\mid\mathcal H_{t-1},a),\qquad
W_t^a=\frac{\mathbf1_{\{A_t=a\}}\phi_t(Y_t)}{\pi_t(a\mid\mathcal H_{t-1})}.
\]

\Needspace{5\baselineskip}
\textbf{Proposition 67. Covered adaptive-action evidence.} If the target
action has \(\pi_t(a\mid\mathcal H_{t-1})\geq\pi_{\min}>0\) at every
time, then

\[
\mathbb E[W_t^a\mid\mathcal H_{t-1}]=\mu_t^a,
\qquad 0\leq W_t^a\leq1/\pi_{\min}.
\]

Consequently Proposition 66 gives a time-uniform certificate for
\(\sum_{s\leq t}\mu_s^a\) from \(\sum_{s\leq t}W_s^a\), with its
boundary multiplied by \(1/\pi_{\min}\). The target is the average
action-conditioned quantity along the encountered contexts, not a
reset-population table or the state distribution under a different
long-term policy.

\textbf{Proof.} Condition first on the action and then on its
observation. Only action \(a\) contributes, with probability
\(\pi_t(a\mid\mathcal H_{t-1})\), cancelling the denominator. The range
bound follows directly. Apply the scaled conditional
moment-generating-function argument to the resulting martingale
differences. \(\square\)

This simple bound is conservative. Variance-aware weighting can improve
it, but cannot remove the need for information about the action of
interest. When \(\pi_{\min}=0\), two models can give exactly the same
law under the observed policy and disagree on a never-selected action.
For example, both can emit a Bernoulli observation with parameter
\(1/2\) under action zero, while their action-one parameters are
\(1/10\) and \(9/10\). Repeating action zero cannot distinguish them.

The action randomization and its recorded probabilities are part of the
experimental assumption. An observational association with a hidden
selection mechanism is not automatically an intervention estimate.
Long-horizon off-policy predictions require additional occupancy or
trajectory-coverage arguments that this one-step statement does not
provide.

\section{What a residual can prove about misspecification}\label{LSM:what-a-residual-can-prove-about-misspecification}

A certified operator family can still be the wrong model for the
observed process. Predictable residuals offer an explicit way to find
that mismatch without pretending that internal validity proves empirical
adequacy.

Let \(P_{t,j}\) and \(Q_{t,j}\) denote the true and proposed
next-observation laws on the same actual history and action. The
proposed law must be defined there. Since \(0\leq\phi_{t,j}\leq1\),

\[
|m_{t,j}-p_{t,j}|\leq\|P_{t,j}-Q_{t,j}\|_{\mathrm{TV}}.
\]

\Needspace{5\baselineskip}
\textbf{Proposition 68. Sequential evidence gives a lower error
certificate.} On the simultaneous event of Proposition 66,

\[
\boxed{
\frac1t\sum_{s=1}^t\|P_{s,j}-Q_{s,j}\|_{\mathrm{TV}}
\geq\frac{(|R_{t,j}|-b_{t,j})_+}{t}.
}
\]

Thus crossing the boundary rejects zero conditional drift for that rule,
with the registered false-rejection budget. Failing to cross supplies no
general upper bound on total variation, no rank certificate and no
guarantee for untested actions.

\textbf{Proof.} The triangle inequality gives
\(|D_{t,j}|\geq(|R_{t,j}|-b_{t,j})_+\). Also
\(|D_{t,j}|\leq\sum_s|m_{s,j}-p_{s,j}|\leq\sum_s\|P_{s,j}-Q_{s,j}\|_{\mathrm{TV}}\).
Combine these inequalities. \(\square\)

\subsection{76.1 Why one aggregate residual can miss conditional
errors}\label{LSM:why-one-aggregate-residual-can-miss-conditional-errors}

The reference path has binary observations. Its true conditional
probability of one is \(4/5\) after a one and \(1/5\) after a zero. A
deliberately inadequate model predicts \(1/2\) at every step. Its
long-run aggregate frequency can be near one half even though its
history-conditioned predictions are wrong.

Three tests are registered before the path is generated. One tests the
raw observation. A second tests it only after the preceding observation
was one. The third applies the same gate but uses the correct
conditional predictor. Gates are functions of past data, not selected
after seeing the current measurement.

For the supplied 4,096-step seeded path, the aggregate mean residual is
approximately \(0.00732\). The gated mean residual is approximately
\(0.14917\), while its time-uniform mean boundary is approximately
\(0.05114\). The lower certificate for average conditional total
variation is approximately \(0.09803\). The correctly predicted gated
residual is approximately \(-0.00303\).

These are illustrative path calculations, not estimated coverage rates.
The theorem supplies the probability guarantee. The example supplies an
explicit reason to test conditional structure rather than relying only
on cancellation-prone aggregate residuals.

The exact reference suite also checks a dependent likelihood-ratio
construction on every binary path through depth eight. Its expectation
under the declared null is exactly one. This is an additional finite
martingale calculation, not a replacement for the general argument or a
claim that every possible test has been run.

\section{Comparing a model with a process outside its representation}\label{LSM:comparing-a-model-with-a-process-outside-its-representation}

The earlier operator-norm law bounds compare two models on one common
probability cone. A misspecified generating process may not have a
representation in the fitted cone. It is useful to state the more basic
comparison at the level of conditional observation laws.

Let \(P\) be any causal process on a standard Borel observation space
with finite actions. Let \(Q\) be a proposed process. Both use the same
measurable policy. Write \(P_t(\cdot\mid h,a)\) and
\(Q_t(\cdot\mid h,a)\) for their next-observation kernels, and define

\[
d_t(h,a)=\|P_t(\cdot\mid h,a)-Q_t(\cdot\mid h,a)\|_{\mathrm{TV}}.
\]

The proposed kernel must be specified on every history needed for the
comparison, including histories impossible under \(Q\) but possible
under \(P\). A declared fallback on a model-null history completes a
causal kernel without changing \(Q\)'s own law. An implementation that
rejects such a history must report that rejection rather than silently
claim to have filtered it.

\Needspace{5\baselineskip}
\textbf{Proposition 69. Observation-law error without a shared state
space.} With the same initial history, the length-\(T\) joint
action-and-observation laws satisfy

\[
\boxed{
\|P_T-Q_T\|_{\mathrm{TV}}
\leq\min\left(1,\sum_{t=0}^{T-1}
\mathbb E_{P,\pi}[d_t(H_t,A_t)]\right).
}
\]

An initial-history discrepancy adds its total variation. In particular,
uniform bounds \(d_t(h,a)\leq\epsilon_t\) give the upper bound
\(\min(1,\sum_{t<T}\epsilon_t)\).

\textbf{Proof.} Introduce hybrid laws that follow \(P\) for the first
\(t\) transitions and \(Q\) for the rest, with the common policy
throughout. Two successive hybrids have the same \(P\) prefix
distribution up to time \(t\) and the same future \(Q\) kernel after the
next transition. Conditional on that prefix and its policy-chosen
action, changing the next kernel costs \(d_t(h,a)\) in total variation.
A common continuation kernel cannot increase total variation, because
integration of a test bounded between zero and one produces another such
test. Average over the shared prefix, sum the adjacent-hybrid bounds,
and cap by one. The same contraction argument handles a different
initial law. \(\square\)

This theorem does not make the right-hand side observable without
further assumptions. Proposition 68 supplies lower evidence about
visited conditional discrepancies, not an upper estimate for every
future history. Turning a finite collection of residuals into an upper
law guarantee requires additional regularity, coverage or a restricted
class with a separating test family.

The distinction prevents a common logical jump: a model can pass several
tests and still have a large untested discrepancy. Conversely, a
statistically significant residual need not mean that its internal
probability algebra is invalid. The model can be a valid distribution
and an inadequate approximation at the same time.

\section{An integrated theory of construction and correction}\label{LSM:an-integrated-theory-of-construction-and-correction}

\Needspace{5\baselineskip}
\textbf{Proposition 70. Separate structural acceptance, approximation
control and sequential rejection.} Consider a fitted polynomial model in
a declared predictive coordinate system, and prospectively registered
observation tests with total budget \(\delta\).

A successfully replayed quadratic certificate from Proposition 62
establishes internal all-horizon validity. A successful operator-only
proposal from Proposition 63 is one way to obtain it, not an assumption
that all valid models have such a proposal.

For any generating density on a common probability cone with a justified
Lipschitz constant, a positive degree-\(n\) approximation and certified
coefficient budgets satisfy Propositions 64--65, hence a fixed-horizon
upper law bound. If a generating process lies outside that
representation, a separately justified bound on its conditional kernels
instead invokes Proposition 69.

For predictable tests on actual dependent observations, Propositions 66
and 68 give simultaneous lower discrepancy certificates, including at
data-selected stopping times. Covered randomized actions admit the
contextual evidence of Proposition 67. These conclusions coexist without
implying that structural acceptance or statistical nonrejection proves
the generating process belongs to the model class.

\textbf{Proof.} Each structural statement is a deterministic consequence
of its replayed equalities and inequalities. The approximation
statements require their declared common-domain, regularity and
coefficient conditions, or alternatively the conditional-kernel
discrepancy assumptions. The statistical conclusions hold on the
simultaneous probability event of Proposition 66. None of these
implications uses nonrejection as an upper bound or infers coverage from
a finite determinant. Their conjunction therefore retains the separate
hypotheses and conclusions exactly as stated. \(\square\)

The central object is now more than a fitted operator. It includes an
admissible domain with a replayable witness, an explicit approximation
class, and an evidence rule specifying which observations could
contradict its predictions. Geometry discovery proposes a
representation. Verification checks its algebra. Approximation theory
accounts for information lost by a chosen finite description. Sequential
tests expose conditional discrepancies in observed signals.

Those roles are related but not interchangeable. A failed search means
that this construction did not provide a witness. A failed witness means
its submitted identities or signs do not establish the claimed validity.
A detected residual means a registered predictive condition is
inconsistent with the observations at the stated error budget. These
three outcomes should be reported differently.

\section{From exact geometry to robustness}\label{LSM:from-exact-geometry-to-robustness}

Propositions 62--70 distinguish structural acceptance, controlled
approximation and sequential evidence. Exact zero-margin preservation
does not supply a positive noise radius. Chapter~\ref{LSM:book:chapter-37} examines that
obstruction and gives strict-margin witnesses for selected perturbed
families. Statistical nonrejection is not used as an upper accuracy
bound.


\label{LSM:part-vii.-robust-curved-certificates-and-noisy-geometry}

The preceding block discovered a curved probability domain from exact
operators. Exact preservation is not the same as stability under
uncertainty. In the worked continuous model, every pure boundary state
can remain on the boundary. A perturbation too small to be conspicuous
in a coefficient table can then move a predicted state outside the
probability domain.

This block separates three operations. A numerical procedure proposes a
geometry. An exact checker decides whether the submitted polynomial
family has a declared margin on that geometry. A separately stated
uncertainty set determines which other families share the guarantee. If
a repair changes the operators, its prediction cost is retained rather
than called an exact coordinate change.

The resulting construction is broader than the preceding exact conformal
search. It accepts nonconformal polynomial families and certifies a full
coefficient box around a repaired estimate. The reference example uses a
deterministic perturbation, not an empirical confidence region.
Statistical coverage is a conditional theorem, not a conclusion of that
example.

Lorentz-positive maps, robust convex constraints and Bernstein
positivity certificates have classical antecedents {[}R8, R13, R21,
R22{]}. The following arguments are self-contained specializations
connecting them to the signal model's representation and evidence
contracts. The nine propositions continue the numbering at 71. Their
count is not a novelty claim or a count of Lean proofs.


\chapter{Robust curved certificates and noisy geometry}\label{LSM:book:chapter-37}
\section{A margin that has probability meaning}\label{LSM:a-margin-that-has-probability-meaning}

Let \(\mathcal L_d=\{(m,v):m\geq\|v\|_2\}\), with mass row
\(\lambda=(1,0,\ldots,0)\). A mass preserving invertible matrix \(T\)
transports a proposed predictive domain to this cone. Its normalized
base in the original coordinates is

\[
K_T=\{T^{-1}(1,v):\|v\|_2\leq1\}.
\]

Write \(F_a(u)=TB_a(u)T^{-1}\) for the transported operator density.
Observations lie in \([-1,1]\) with reference measure \(du/2\). All
operator densities in this block are polynomial and all equalities
involving integrals use that reference measure. For each action define

\[
\sigma_a(u)=\min_{\|v\|_2\leq1}
\left[(F_a(u)(1,v))_0-\|(F_a(u)(1,v))_{\rm sp}\|_2\right].
\]

The subscript \(\rm sp\) denotes spatial coordinates, not a separate
measured signal. This margin is expressed before normalization. It
simultaneously controls mass orientation and distance from the curved
boundary. It is coordinate dependent, so its uncertainty budget must use
the same \(T\).

\Needspace{5\baselineskip}
\textbf{Proposition 71. Robust oriented margin.} Suppose
\(\sigma_a(u)\geq\kappa_a(u)\) and an uncertain family has transported
errors \(\Delta_a(u)\) satisfying, for every normalized input,

\[
\begin{aligned}
|[\Delta_a(u)(1,v)]_0|&\leq e_{a,0}(u),\\
\|[\Delta_a(u)(1,v)]_{\rm sp}\|_2&\leq e_{a,1}(u).
\end{aligned}
\]

If \(\kappa_a(u)\geq e_{a,0}(u)+e_{a,1}(u)\), every uncertain branch
preserves \(\mathcal L_d\). If the uncertain family has exact integrated
mass and its initial state belongs to \(K_T\), it defines a valid
probability model at every finite horizon. With a strict inequality, the
residual slack is at least \(\kappa_a-e_{a,0}-e_{a,1}\).

\textbf{Proof.} For \(w=F_a(u)(1,v)\) and \(d=\Delta_a(u)(1,v)\),

\[
(w_0+d_0)-\|w_{\rm sp}+d_{\rm sp}\|_2
\geq w_0-\|w_{\rm sp}\|_2-|d_0|-\|d_{\rm sp}\|_2.
\]

The right side is nonnegative under the hypothesis. Homogeneity extends
the assertion from normalized inputs to the entire cone. The only
zero-mass cone element is zero. Integrated mass conservation and
successive cone preservation give nonnegative normalized observation
laws by the same induction as in Proposition 52. A positive-density
branch may be normalized. A zero-density branch is not assigned a
posterior. \(\square\)

This is a uniform deterministic guarantee. The inequalities must cover
the whole observation interval and the whole normalized domain. A noise
scale estimated from a few residuals does not automatically provide
those inequalities. For continuous observations the displayed margin is
a density margin relative to the declared measure, not a probability
assigned to a single observation value.

\section{Why exact boundary preservation can have zero noise radius}\label{LSM:why-exact-boundary-preservation-can-have-zero-noise-radius}

An exact certificate can be correct while having no tolerance to
unrestricted perturbations. This is a geometrical fact, not a failure of
the certificate checker.

\Needspace{5\baselineskip}
\textbf{Proposition 72. Boundary fragility under mass preserving
perturbations.} Assume \(d\geq1\). Let \(F\) preserve \(\mathcal L_d\).
Suppose a normalized input \(z=(1,v)\) has a boundary image \(Fz=(m,w)\)
with \(m=\|w\|_2\). For every \(\epsilon>0\), there is a linear
perturbation \(\Delta\) with \(\lambda\Delta=0\) and

\[
\sup_{\|v'\|_2\leq1}\|\Delta(1,v')\|_{\mathcal L_d}=\epsilon
\]

such that \((F+\Delta)z\notin\mathcal L_d\). Here the base norm is
\(\|(m,w)\|_{\mathcal L_d}=\max(|m|,\|w\|_2)\).

\textbf{Proof.} If \(w\ne0\), put \(e=w/\|w\|_2\). If \(w=0\), choose
any spatial unit vector \(e\). Define \(\Delta=\epsilon(0,e)\lambda\).
Its mass row vanishes and its value on every normalized input is
\((0,\epsilon e)\), of base norm \(\epsilon\). If \(w\ne0\), the new
spatial norm is \(m+\epsilon\). If \(w=0\), the original output mass is
zero and the new spatial norm is \(\epsilon\). Either image lies outside
the cone. \(\square\)

The positive spatial dimension is essential. When \(d=0\), the state is
mass alone and \(\lambda\) is the identity on \(\mathbb R\). The
condition \(\lambda\Delta=0\) forces \(\Delta=0\), so a perturbation of
prescribed norm \(\epsilon>0\) cannot exist. This explicit boundary was
missing from the earlier wording. All previously worked examples have
two spatial coordinates and are unaffected.

For a continuous operator family, a constant spatial perturbation of
this form preserves every integrated mass row and produces the same
obstruction at an observation with a boundary image. It need not
preserve the algebraic restrictions of a narrower uncertainty class.
Perturbing a positive congruence factor and rebuilding its congruence,
for example, is different from perturbing an arbitrary predictive
matrix.

In the preceding exact continuous example, the direct operator at
\(u=0\) is \(I/4\). Every normalized boundary state has a boundary
image. There is therefore no positive unrestricted matrix-error radius
on that domain, even though the original family is exactly valid. A
robust theorem that silently treats its zero slack as positive would be
false.

\section{An explicit repair and the price of its slack}\label{LSM:an-explicit-repair-and-the-price-of-its-slack}

In Lorentz coordinates let

\[
N_a(u)=g_a(u)e_0\lambda,\qquad
 g_a(u)\geq g_{\min}>0,\qquad \int g_a(u)\,du/2=1.
\]

This reference instrument reports a state-independent density and resets
to the center. It is a mathematical model modification, not a new
observation. Given an estimate \(\widehat F\), define

\[
\widetilde F_a(u)=(1-\eta)\widehat F_a(u)+\eta N_a(u),
\qquad 0\leq\eta\leq1.
\]

\Needspace{5\baselineskip}
\textbf{Proposition 73. Interior repair with an explicit error budget.}
Suppose a reference family \(F^*\) is valid on \(\mathcal L_d\), and the
combined mass and spatial error of \(\widehat F-F^*\) is at most
\(d_a(u)\) on normalized inputs. Then the margin of
\(\widetilde F_a(u)\) is at least

\[
\boxed{\eta g_a(u)-(1-\eta)d_a(u).}
\]

An additional uncertainty family with combined output error at most
\(\nu_a(u)\) is certified whenever

\[
\eta g_a(u)\geq(1-\eta)d_a(u)+\nu_a(u).
\]

Assume exact integrated mass for the estimate. If its integrated
base-norm discrepancy from \(F^*\) is at most \(\varepsilon\), the
corrected discrepancy is at most

\[
\boxed{(1-\eta)\varepsilon+2\eta.}
\]

When the corrected model is valid on the same cone, with the same
initial state, its length-\(T\) law under any common policy differs from
the reference law by at most

\[
\min\left(1,\frac T2[(1-\eta)\varepsilon+2\eta]\right).
\]

\textbf{Proof.} The margin of \(\widehat F\) is at least \(-d_a\) by
Proposition 71's inequality applied to \(F^*\). Adding the reference
changes only output mass, by \(\eta g_a\), while scaling the remaining
output by \(1-\eta\). This proves the first bound and the robust
implication. Integrated mass is preserved by mixing.

For the error, write

\[
\widetilde F-F^*=(1-\eta)(\widehat F-F^*)+\eta(N-F^*).
\]

On a normalized input, both \(N\) and \(F^*\) distribute total cone mass
one. Their integrated difference has base norm at most two. Integrate
the triangle inequality to obtain the stated budget. The final
implication is the continuous observation-law bound of Proposition 59.
It requires validity of both comparison models on the same cone.
\(\square\)

For constant budgets, a sufficient repair weight is
\((d+\nu)/(g_{\min}+d)\) when this number is at most one. This is not an
optimality claim. The reference example instead fixes a repair weight
and verifies the result independently. If an exact model already has a
robust structured certificate, adding a reset may be unnecessary and
would introduce avoidable bias.

The term \(2\eta\) is deliberately retained. A larger repair creates
more room for uncertainty but can change future predictions more. The
worst-case law bound can become uninformative at long horizons. Neither
positivity nor mixing alone proves long-horizon predictive accuracy.

\section{Noisy invariant forms are proposals, not certificates}\label{LSM:noisy-invariant-forms-are-proposals-not-certificates}

The exact identity \(A^{\mathsf T}HA=H\) need not survive estimation
error. A useful proposal can still be obtained from a least-squares
invariant form, but the residual must not be interpreted as a
positive-cone guarantee.

Let \(\mathcal M_A\) map symmetric matrices into a Euclidean product
space by

\[
\mathcal M_A(H)=(A^{\mathsf T}HA-H,\operatorname{tr}H).
\]

Use the Frobenius norm on symmetric matrices and the induced product
norm. Assume \(\mathcal M_{A^*}\) is injective and
\(\mathcal M_{A^*}(H^*)=(0,1)=b\). This fixes the scale of a
one-dimensional invariant-form space.

\Needspace{5\baselineskip}
\textbf{Proposition 74. A posteriori proposal error.} Given
\(\widehat H\), let \(r=\mathcal M_{\widehat A}(\widehat H)-b\). Suppose

\[
\|\mathcal M_{\widehat A}-\mathcal M_{A^*}\|\leq e_M,
\qquad \sigma_{\min}(\mathcal M_{\widehat A})>e_M.
\]

Then

\[
\boxed{
\|\widehat H-H^*\|_F
\leq
\frac{\|r\|+e_M\|\widehat H\|_F}
{\sigma_{\min}(\mathcal M_{\widehat A})-e_M}.
}
\]

One may take \(e_M\) from the bound

\[
\|\mathcal M_{\widehat A}-\mathcal M_{A^*}\|\leq(\|\widehat A\|_2+\|A^*\|_2)\|\widehat A-A^*\|_2.
\]

For fixed centers \(c^*=(I-A^*)^{-1}b^*\) and
\(\widehat c=(I-\widehat A)^{-1}\widehat b\), with
\(\|\widehat A-A^*\|_2\leq e_A\) and \(\|\widehat b-b^*\|_2\leq e_b\),
one has

\[
\|\widehat c-c^*\|_2
\leq\frac{e_b+e_A\|\widehat c\|_2}
{\sigma_{\min}(I-\widehat A)-e_A}
\]

when the denominator is positive.

\textbf{Proof.} The invariant-form difference satisfies

\[
\mathcal M_{A^*}(\widehat H-H^*)
=r-(\mathcal M_{\widehat A}-\mathcal M_{A^*})\widehat H.
\]

The least singular value of \(\mathcal M_{A^*}\) is at least that of
\(\mathcal M_{\widehat A}\) minus \(e_M\). Apply this lower bound and
the triangle inequality. To bound the map perturbation, insert and
subtract \(\widehat A^{\mathsf T}HA^*\) and apply the Frobenius product
inequality. The trace component is unchanged. Finally,

\[
(I-A^*)(\widehat c-c^*)=(\widehat b-b^*)+(\widehat A-A^*)\widehat c.
\]

The same singular-value argument proves the center bound. \(\square\)

The theorem assumes an actual perturbation budget and a corresponding
exact reference problem. The example reports numerical residuals but
does not claim to have statistically certified these budgets. It uses a
trace-normalized least-squares shape, a nearly rank-one image for scale,
and rationalized center and factor. It accepts none of those
computations as a proof of validity. The next checker uses only the
submitted rational geometry and the complete polynomial operators.

A rationalized invertible factor \(L\) specifies \(H=L^{\mathsf T}L\)
exactly. Positive definiteness is then an algebraic fact even when the
factor was suggested by floating-point calculations. It still does not
establish operator invariance.

\section{Finite exact checking on a continuum}\label{LSM:finite-exact-checking-on-a-continuum}

For three predictive coordinates, the spatial unit circle admits two
rational charts. For \(s\in\{1,-1\}\) and \(t\in[-1,1]\), put

\[
d(t)=1+t^2,\qquad
r_s(t)=\frac{s}{d(t)}(1-t^2,2t).
\]

The charts cover the two semicircles, including their endpoints. For an
operator \(F_a(u)\) and required margin \(\kappa\geq0\), form polynomial
vectors and scalars

\[
\begin{aligned}
w_{a,s}(u,t)&=F_a(u)(d(t),s(1-t^2),2st)^{\mathsf T},\\
f_{a,s}(u,t)&=w_{a,s,0}(u,t)-\kappa d(t),\\
h_{a,s}(u,t)&=f_{a,s}(u,t)^2-w_{a,s,1}(u,t)^2-w_{a,s,2}(u,t)^2.
\end{aligned}
\]

No square root occurs in these polynomials.

\Needspace{5\baselineskip}
\textbf{Proposition 75. Rational semicircle and Bernstein certificate.}
If \(f_{a,s}\) and \(h_{a,s}\) are nonnegative on \([-1,1]^2\) for both
charts and every action, then \(\sigma_a(u)\geq\kappa\) for every
observation. For rational operators and geometry, nonnegative
tensor-Bernstein coefficients on a complete rational subdivision of that
square give a finite exact witness. If the actual margin is uniformly
strictly greater than \(\kappa\), sufficiently fine dyadic subdivision
succeeds.

\textbf{Proof.} The first polynomial inequality establishes output
orientation after subtracting \(\kappa d\). The second then gives

\[
w_{a,s,0}-\|(w_{a,s,1},w_{a,s,2})\|_2\geq\kappa d.
\]

Divide by \(d>0\) to obtain the required margin on the circle. For fixed
\(u\), the function consisting of affine output mass minus the norm of
affine spatial output is concave in the input. Every point in the closed
disk is a convex combination of circle points. Concavity therefore
extends the lower bound to the disk.

On any rectangle mapped affinely to \([0,1]^2\), the tensor-Bernstein
basis functions are nonnegative and sum to one. A polynomial with
nonnegative coefficients in that basis is nonnegative throughout the
rectangle. If the rectangles cover the square, their inequalities prove
the global statement. The reference verifier checks coverage through a
complete subdivision tree, not by accepting an arbitrary list of
rectangles.

For termination under strict margin, both \(f\) and \(h\) are strictly
positive on the compact square. Their minima are positive. For a
fixed-degree polynomial on a small cell of side at most \(\delta\), its
Taylor expansion at one corner has constant term equal to its value
there and nonconstant coefficients bounded uniformly by \(O(\delta)\).
The linear change from the fixed-degree power basis to the Bernstein
basis is fixed and bounded. Thus every Bernstein coefficient differs
from that constant value by \(O(\delta)\), uniformly over cells. At
sufficiently small \(\delta\) all coefficients are positive. A finite
dyadic subdivision attains that size. \(\square\)

This proof is a specialized use of classical Bernstein range bounds
{[}R22{]}. It does not claim that a negative Bernstein coefficient
proves the polynomial negative. Such a coefficient can merely mean that
further subdivision is needed. A negative evaluated corner, by contrast,
is an actual counterexample. At a zero-margin boundary, finite
termination of this particular witness language is not promised. The
exact conformal and matrix certificates of preceding blocks remain
useful there.

The two tests are both necessary to this sufficient certificate. The
squared inequality alone also admits reversed mass orientation. Both
charts are required. A verifier that omits one chart has not covered the
circle.

\section{A finite uncertainty box and the role of convexity}\label{LSM:a-finite-uncertainty-box-and-the-role-of-convexity}

A robust witness should identify its uncertainty family rather than
supply an unexplained tolerance. Suppose the transformed error has
polynomial coefficients

\[
\Delta_a(u)=\sum_{k=0}^{D}u^kE_{a,k},\qquad
|(E_{a,k})_{ij}|\leq r_k.
\]

All bounds are in the submitted Lorentz coordinates. For three
coordinates, normalized inputs have each coordinate of magnitude at most
one.

\Needspace{5\baselineskip}
\textbf{Proposition 76. Coefficient boxes and convex uncertainty.} The
coefficient box satisfies the output bounds of Proposition 71 with

\[
\boxed{e_0=3\sum_{k=0}^{D}r_k,\qquad e_1=6\sum_{k=0}^{D}r_k.}
\]

Only coefficient arrays obeying \(\lambda\int\Delta_a(u)\,du/2=0\) are
included in the normalized uncertainty family. A checked margin at least
\(e_0+e_1\) certifies every such array.

Alternatively, if finitely many complete polynomial operator families
satisfy a common margin \(\kappa\) on the same domain, every convex
combination of those families satisfies that margin. Initial states may
also range over a certified convex hull.

For a fixed geometry, the constraints \(\sigma_a(u)\geq\kappa\), exact
mass, and membership in a convex data-consistency set are convex
constraints on the operator coefficients and initial state. Joint
optimization of geometry and operators is not asserted to be convex.

\textbf{Proof.} Each output scalar is bounded by \(3r_k\) on a
normalized input. Bound the spatial Euclidean norm by the sum of the
absolute values of its two coordinates, obtaining \(6r_k\). Sum over
coefficients using \(|u|\leq1\). Mass conservation is a separate linear
constraint and follows for every family satisfying it.

For convex combinations, output mass is affine and the spatial norm is
convex. Hence the margin of the combined output is at least the same
combination of margins. A convex cone contains convex combinations of
admissible initial states. For fixed geometry and input, the inequality
affine mass minus a norm at least \(\kappa\) is convex in the
coefficients. Intersect over all inputs and observations, then add the
linear mass constraints and the convex data set. \(\square\)

The squared polynomial witness in Proposition 75 is a convenient
verification language. Its coefficient constraints are not automatically
a convex optimization formulation in unknown operators. Convexity of the
underlying admissibility set must not be confused with convexity of
every certificate parameterization.

Likewise, a free uncertainty box in all entries contains families that
violate integrated mass. The affine restriction is essential. Removing
physically valid candidates from a statistical confidence set for some
other reason requires its own coverage argument.

\section{What a search can guarantee}\label{LSM:what-a-search-can-guarantee}

The numerical proposal is one practical route. There is also a precise
conditional search theorem showing why a broad class of strict curved
certificates can be found without an exact conformal identity.

\Needspace{5\baselineskip}
\textbf{Proposition 77. Conditional search for strict rational curved
certificates.} Consider finitely many rational polynomial operator
families in three coordinates, each with exact mass, and finitely many
rational normalized initial states. Suppose there exists a real mass
preserving invertible \(T_*\) such that every transformed initial state
lies strictly inside the Lorentz cone and every transformed operator
family has a common positive margin on the observation interval. Then a
rational mass preserving invertible \(T\), a positive rational margin,
and finite rational Bernstein witnesses exist. A dovetailed enumeration
of rational transforms, positive rational margins and subdivision depths
eventually finds such a witness. Whenever it returns, the witness is
sound for the convex hull of the supplied families.

\textbf{Proof.} Inversion and matrix multiplication are continuous near
\(T_*\). Initial membership is strict, and the operator margin is
continuous jointly in the transform, the observation, and the normalized
input. The latter two range over compact sets and there are finitely
many families. Thus the strict inequalities persist in a neighborhood of
\(T_*\). Rational mass preserving matrices are dense in that affine
matrix space, and invertibility persists in a neighborhood. Choose a
rational transform there and a rational positive margin smaller than the
remaining minimum. Proposition 75 supplies finite subdivisions because
there is still a strict gap above the chosen margin. Enumerating all
finite rational descriptions and depths eventually reaches this finite
witness. Proposition 76 proves the convex-hull conclusion. \(\square\)

This is a conditional semidecision result, not a polynomial-time
complexity theorem. It does not cover every valid model. Initial states
on an essential boundary and families with zero margin may violate its
assumptions. The implemented procedure uses a finite rationalized
proposal and bounded subdivision, not this unbounded enumeration.
Failure of that bounded procedure is a withheld certificate rather than
a proof that no other domain works.

The theorem's uncertainty families are supplied in fixed predictive
coordinates. It does not establish that those families cover an unknown
generating process. It demonstrates a route from a strict robust model
sector to finite machine-checkable witnesses.

\section{Selection from data does not remove the uncertainty quantifier}\label{LSM:selection-from-data-does-not-remove-the-uncertainty-quantifier}

A domain can be selected after examining fitted operators without a
separate statistical penalty for each geometry tried. That statement
holds only when the final verification ranges over an already justified
confidence set.

\Needspace{5\baselineskip}
\textbf{Proposition 78. Data-selected geometry with honest coverage.}
Suppose a random set \(\mathcal U\) of operators and initial states, in
a declared fixed predictive coordinate system, contains the true model
with probability at least \(1-\delta\). A data-dependent procedure may
propose any transform \(T\) and any domain. If its exact verifier proves
domain membership and invariance for every member of the same set
\(\mathcal U\), then on that coverage event the true model has the
verified domain. No union bound over geometry proposals is needed for
this deterministic implication.

If original-coordinate perturbations satisfy \(\|\Delta\|_2\leq e\),
their combined mass and spatial output error after a known transform
\(T\) is at most

\[
\boxed{2\|T\|_2\|T^{-1}\|_2 e.}
\]

The corresponding transformed uncertainty budget must be used even when
\(T\) was chosen from the data.

\textbf{Proof.} On the coverage event the true model is a member of the
set quantified by the verifier. Its domain properties follow by direct
inclusion, regardless of how \(T\) was selected. For the error, a
normalized Lorentz input has Euclidean norm at most \(\sqrt2\). Thus the
Euclidean norm of \(T\Delta T^{-1}z\) is at most
\(\sqrt2\|T\|_2\|T^{-1}\|_2e\). The sum of its mass magnitude and
spatial norm is at most \(\sqrt2\) times its Euclidean norm. Multiply
the two bounds. \(\square\)

The set must not be shrunk around a convenient candidate without
preserving coverage. Moving the center by a reset mixture changes the
family being certified. A tube around the repaired family need not
contain the original true family. Repeatedly estimating new confidence
sets still requires the stage accounting of Proposition 48. Statistical
rank coverage and observation identifiability remain assumptions of the
upstream procedure.

The final reference certificate explicitly says that no statistical
coverage is claimed. Its uncertainty box is a deterministic, exactly
described set around the repaired operator family. This supports a
reproducible robustness statement without pretending that the example
has estimated an unknown noise distribution.

\section{A complete perturbed-operator construction}\label{LSM:a-complete-perturbed-operator-construction}

The input to the new proposer contains only two perturbed polynomial
operators and the initial vector. A separate fixture generator creates
those inputs from the preceding continuous model by adding
\(\pm\delta E(u)\), with opposite signs for the two actions and

\[
\delta=\frac1{2000},\qquad
E(u)=E_0+uE_1+u^2E_2,
\]

where

\[
\begin{aligned}
E_0&=\begin{pmatrix}0&0&0\\-1&2&0\\0&1&-1\end{pmatrix},&
E_1&=\begin{pmatrix}0&1&-1\\1&0&1\\1&-1&0\end{pmatrix},&
E_2&=\begin{pmatrix}0&0&0\\0&-1&1\\1&0&1\end{pmatrix}.
\end{aligned}
\]

The perturbation preserves integrated mass. It is not a statistical
sample. The proposer does not receive the predecessor metric or a hidden
generator.

A noisy operator ratio supplies a center and a least-squares shape. A
finite observation-anchor list supplies a scale proposal. Rationalizing
the center and a triangular metric factor with denominator limit 48
gives

\[
c=\begin{pmatrix}1/2\\1/2\end{pmatrix},\qquad
L=\begin{pmatrix}123/37&-2\\0&109/41\end{pmatrix},
\]

and hence

\[
\boxed{
H=L^{\mathsf T}L=
\begin{pmatrix}15129/1369&-246/37\\-246/37&18605/1681\end{pmatrix}.
}
\]

This is not the exact predecessor metric. The transformed family is not
assumed conformal. The raw fitted family actually fails the proposed
domain: for the direct action at \(u=-1\) and Lorentz input \((1,1,0)\),
its output mass is \(106321/82000>0\), while the output quadratic form
is

\[
-\frac{12009932651317769}{557059220496000000}<0.
\]

This is a genuine rejected input state, not a negative Bernstein
coefficient of uncertain meaning. The normalized output lies outside the
ellipse. No positivity certificate is claimed for the raw fit.

The construction next mixes each operator with the center-reset
reference using \(\eta=1/8\). The exact checker proves a margin of at
least \(1/16\) for the resulting family. It verifies eight polynomial
inequalities across both actions and both semicircle charts. The witness
contains 116 subdivision leaves and 1,780 exactly recomputed nonnegative
Bernstein coefficients. These are arithmetic certificate sizes, not
separate theorems.

The exported family admits mass and spatial perturbation budgets each
equal to \(1/128\), leaving certified slack \(3/64\). In particular, it
covers every degree-two transformed coefficient perturbation with each
entry bounded by \(1/4096\), subject to the exact integrated mass
constraints. Proposition 76 verifies that the whole coefficient box lies
inside the declared output budgets. No uncertainty vertices or future
trajectories are sampled to establish that inclusion.

The certificate is replayed independently from its JSON record. The
record contains rational operators, initial state, transform, margin,
uncertainty definition and complete subdivision trees. Floating-point
proposal diagnostics are stored elsewhere and are not trusted by the
replay routine.

\section{The prediction cost needs a common comparison domain}\label{LSM:the-prediction-cost-needs-a-common-comparison-domain}

A successful discovered domain does not imply that the known reference
model has that same domain. To invoke a common-cone law bound, the
example separately checks the corrected family on the predecessor's
exact ellipse. This is a second certificate used only for comparison. It
is not supplied to geometry discovery.

In that comparison domain, the coefficient-wise upper bounds on the raw
fitting discrepancy are

\[
\frac{11}{5000},\qquad \frac{137}{30000},\qquad \frac{37}{8000}.
\]

Using \(\int_{-1}^1|u|^kdu/2=1/(k+1)\) gives the integrated upper budget
\(241/40000\). The same center-reset reference is valid there. Both
models have the same initial state.

\Needspace{5\baselineskip}
\textbf{Proposition 79. Certified repair with a retained comparison
budget.} For the exact construction of Section 88, independent replay
establishes internal validity of the repaired model and its declared
coefficient uncertainty family. Independent comparison-domain replay and
the coefficient bounds above give

\[
\varepsilon_{\rm repaired}
\leq\frac78\frac{241}{40000}+\frac14
=\frac{81687}{320000}.
\]

Consequently, for any common policy and any finite horizon \(T\),

\[
\boxed{
\|P_T-\widetilde P_T\|_{\rm TV}
\leq\min\left(1,\frac{81687T}{640000}\right).
}
\]

This compares the unmodified reference process and the nominal repaired
process. It is not a guarantee that the raw noisy model was valid or
that the repaired uncertainty family contains the original reference.

\textbf{Proof.} The polynomial witnesses and exact mass identities imply
the nominal and robust validity statements by Propositions 71, 75 and
76. The separate comparison witness establishes the required common
domain for the nominal repaired and reference operators. The absolute
coefficient bounds dominate the integrated raw error by the triangle
inequality. Proposition 73 then gives the displayed repair budget and
law bound. All hypotheses are either exact identities in the supplied
witnesses or direct norm bounds from the declared rational coefficients.
\(\square\)

The bound is conservative. It becomes one at horizon eight, so it makes
no nontrivial long-horizon accuracy claim. The example's point is the
separation of a verified all-horizon probability model from a
finite-horizon approximation guarantee. A probability-valid model is not
necessarily a close model.

The noisy proposal, the failed raw fit, the repair and the comparison
bound all appear in the release records. None is replaced by a more
favorable trajectory score. The coefficient box is around the repaired
family in the specified coordinates, not an inferred distribution of
sensor noise.

\section{From deterministic robustness to data confidence}\label{LSM:from-deterministic-robustness-to-data-confidence}

Chapter~\ref{LSM:book:chapter-37} gives exact continuum certificates for deterministic
uncertainty families. The worked repaired model is internally valid at
every finite horizon, but its original comparison bound becomes trivial
at horizon eight. Chapter~\ref{LSM:book:chapter-38} reduces that fixed repair and separately
constructs a statistically covered domain for a supplied observable
family. Its confidence statement is not transferred to the earlier
repaired operators.


\label{LSM:part-viii.-data-supported-domains-and-smaller-repairs}

The preceding block certified a perturbed operator family after mixing
it with a center reset. Its all-horizon validity guarantee was exact,
but its error bound became trivial at horizon eight. It also kept
deterministic coefficient boxes separate from statistical confidence
sets. This block addresses both limitations in declared sectors rather
than weakening either distinction.

First, the same fitted family and previously proposed geometry admit a
smaller, independently checked repair. Second, a self-normalized
martingale argument produces operator uncertainty sets from a dependent
trajectory when the predictive state is an accessible, fixed observable.
A nonlinear three-coordinate example connects those sets to a curved
domain without a reset correction. Its geometry and feature dictionary
are specified in advance. This is not a solution to discovering an
unknown hidden state from the same data.

The sequential-regression result is a classical self-normalized
concentration theorem {[}R23{]}, proved here in the form used by the
signal model. The contraction argument uses standard Markov-kernel and
total-variation ideas {[}R24{]}. The new block's contribution is the
explicit sequence of assumptions, finite witnesses, structured
uncertainty, observation-measure distinctions and executable examples
connecting these tools to the realization calculus. Proposition counts
are not historical-novelty counts.


\chapter{Data-supported domains and smaller repairs}\label{LSM:book:chapter-38}
\section{A smaller repair can be bracketed rather than guessed}\label{LSM:a-smaller-repair-can-be-bracketed-rather-than-guessed}

Keep the fitted operators, proposed transform and initial state from
Section 88 unchanged. In Lorentz coordinates use the center-reset
operator \(N=e_0\lambda\), of constant density one under \(du/2\). Write

\[
F_\eta=(1-\eta)\widehat F+\eta N,
\qquad 0\leq\eta\leq1.
\]

Fix a required margin \(0\leq\kappa\leq1\). Feasibility means that every
action, observation and normalized input has margin at least \(\kappa\).
Integrated mass and initial membership remain separate requirements.

\Needspace{5\baselineskip}
\textbf{Proposition 80. Convex repair bracketing.} The feasible repair
weights form a closed interval \([\eta_*,1]\). If a cone-positive
functional \(\ell=(1,-v^{\mathsf T})\), with \(\|v\|_2\leq1\), and a
normalized input \(z\) give \(c=\ell\widehat F_a(u)z<\kappa\), then
every feasible weight satisfies

\[
\boxed{\eta\geq\frac{\kappa-c}{1-c}.}
\]

A complete certificate at a weight \(\eta_U\) gives an upper bound on
\(\eta_*\). A negative Bernstein coefficient without a negative
evaluated point does not give a lower bound.

\textbf{Proof.} For a fixed input and observation, output mass minus
spatial norm is a concave continuous function of the operator, hence of
\(\eta\). Its superlevel set is closed and convex. Intersecting these
sets preserves those properties. The reset has margin one, so the
intersection contains one and is an interval ending there. If \(F_\eta\)
has margin at least \(\kappa\), every such functional satisfies
\(\ell F_\eta z\geq\kappa\). Since \(\ell Nz=1\), this requires
\((1-\eta)c+\eta\geq\kappa\). Here \(1-c>0\), and solving gives the
bound. A complete positive witness proves feasibility. A negative
coefficient only says that a particular sufficient polynomial bound
failed. \(\square\)

\subsection{91.1 The unchanged predecessor
example}\label{LSM:the-unchanged-predecessor-example}

Use the previous transform with center \((1/2,1/2)\) and triangular
factor

\[
L=\begin{pmatrix}123/37&-2\\0&109/41\end{pmatrix}.
\]

At the direct action, observation \(u=-1\), normalized input \((1,1,0)\)
and functional \(\ell=(1,7/25,-24/25)\), exact evaluation gives

\[
c=-\frac{17590411}{2332387500}.
\]

For the required margin \(\kappa=1/64\), this proves

\[
\frac{864543451}{37599646576}
\leq\eta_*\leq\frac1{32}.
\]

The lower endpoint is approximately \(0.02299\). The upper endpoint is
not a sampled estimate. A complete exact Bernstein certificate proves
that \(\eta=1/32\) works. The weight \(1/64\) fails the requested
\(1/64\) margin at the displayed separating functional. This does not by
itself prove that its nominal model is probability-invalid.

The new certificate uses 194 leaves and 3,058 recomputed coefficients. A
separate comparison-domain certificate uses 185 leaves and 2,833
coefficients. Both retain the same degree-two coefficient box of radius
\(1/4096\) in their respective declared Lorentz coordinates, subject to
exact integrated mass. The more general mass and spatial output budgets
remain \(1/128\) each. Their sum equals the certified margin, so that
larger output-error family has a nonnegative lower margin, not a claimed
strictly positive residual margin.

The reset weight has fallen from \(1/8\) to \(1/32\), while the
deterministic uncertainty box is unchanged. No global optimization over
geometry or repair operators is claimed. The current bracket concerns
this fixed geometry, this reset path and this required margin.

The comparison domain is checked again before invoking the law theorem.
Keeping the same raw discrepancy budget \(241/40000\) from Section 89
gives

\[
\varepsilon_{\mathrm{repair}}
\leq\frac{31}{32}\frac{241}{40000}+\frac1{16}
=\frac{87471}{1280000}.
\]

Consequently, under a common policy and identical initial states,

\[
\boxed{
\|P_T-\widetilde P_T\|_{\mathrm{TV}}
\leq\min\left(1,\frac{87471T}{2560000}\right).
}
\]

The first integer horizon where this upper bound reaches one is now 30
rather than eight. The bound remains conservative. The repaired
uncertainty tube is still not a confidence set for the unmodified
generating process. This improvement is deterministic certificate
refinement, not a new statistical experiment.

\section{Confidence from dependent observations}\label{LSM:confidence-from-dependent-observations}

Let \(\mathcal F_{t-1}\) contain all information available before
observation \(t\), including its already selected action. Let
\(\phi_t\in\mathbb R^d\) be measurable with respect to that information.
Assume

\[
Y_t=\theta_*^{\mathsf T}\phi_t+\xi_t,
\qquad
\mathbb E[e^{h\xi_t}\mid\mathcal F_{t-1}]
\leq e^{h^2\sigma^2/2}
\]

for every real \(h\), with \(\sigma>0\), and a declared bound
\(\|\theta_*\|_2\leq S\). With ridge \(\gamma>0\), set

\[
V_t=\gamma I+\sum_{s\leq t}\phi_s\phi_s^{\mathsf T},
\qquad
\widehat\theta_t=V_t^{-1}\sum_{s\leq t}\phi_sY_s.
\]

\Needspace{5\baselineskip}
\textbf{Proposition 81. Time-uniform regression confidence.} For
\(0<\delta<1\), with probability at least \(1-\delta\), simultaneously
for every finite \(t\),

\[
\boxed{
\|\widehat\theta_t-\theta_*\|_{V_t}
\leq
\sigma\sqrt{
\log\frac{\det V_t}{\det(\gamma I)}+2\log\frac1\delta
}
+\sqrt\gamma S.
}
\]

No independence of successive observations or of the regressors is
required. The conditional mean, conditional noise and predictability
hypotheses are required. The confidence set remains valid at a stopping
time selected from the data.

\textbf{Proof.} Put \(S_t=\sum_{s\leq t}\phi_s\xi_s\). For every fixed
\(q\), the conditional moment bound makes

\[
\exp\left(q^{\mathsf T}S_t/\sigma
-\tfrac12q^{\mathsf T}(V_t-\gamma I)q\right)
\]

a nonnegative supermartingale. Integrate against the fixed Gaussian law
\(q\sim N(0,(\gamma I)^{-1})\). Conditional Tonelli and completing the
square give the nonnegative supermartingale

\[
M_t=
\left(\frac{\det(\gamma I)}{\det V_t}\right)^{1/2}
\exp\left(\frac{\|S_t\|_{V_t^{-1}}^2}{2\sigma^2}\right),
\qquad M_0=1.
\]

For completeness, stop at the first crossing of \(1/\delta\), then at
any deterministic finite horizon. The supermartingale expectation is at
most one, so the crossing probability up to that horizon is at most
\(\delta\). Increasing the horizon proves the same bound for any finite
crossing. Off this event, taking logarithms bounds
\(\|S_t\|_{V_t^{-1}}\) by the square-root term. Finally,

\[
\widehat\theta_t-\theta_*
=V_t^{-1}(S_t-\gamma\theta_*).
\]

The triangle inequality and \(V_t\succeq\gamma I\) bound the ridge term
by \(\sqrt\gamma S\). The simultaneous event already covers every
selected finite stopping time. \(\square\)

This is the self-normalized mixture argument in {[}R23{]}, not a new
concentration principle. A bounded response in an interval of length
\(b-a\) has conditionally sub-Gaussian centered noise with
\(\sigma=(b-a)/2\). One direct proof differentiates its tilted log
moment-generating function twice. The second derivative is a variance
bounded by \((b-a)^2/4\). Integrating from zero gives the stated moment
bound.

The determinant term accounts for the observed design, but does not
create excitation. If \(\phi_t\) never probes a direction, the set need
not shrink in that direction. Regressors extracted after inspecting the
same response are not predictable. An estimated hidden state cannot be
substituted for \(\phi_t\) while silently retaining the theorem's true
conditional-mean equation.

\section{From observable moment regression to operator sets}\label{LSM:from-observable-moment-regression-to-operator-sets}

There is a direct connection to the unnormalized instrument. Assume that
the normalized predictive state \(Z_t\) is an exact, fixed, measurable
function of the available observation history. For action \(a\), suppose

\[
P(U_t\in du\mid\mathcal F_{t-1})
=\lambda B_a(u)Z_t\,\mu(du),
\qquad
Z_{t+1}=\frac{B_a(U_t)Z_t}{\lambda B_a(U_t)Z_t}
\]

on realized positive-density branches. Crucially, \(Z_{t+1}\) must also
be available from the observations using the specified state map. This
is an accessibility assumption, not a consequence of successful
regression.

Let \(B_a(u)=\sum_{k=0}^{D}\psi_k(u)C_{a,k}\). Choose functions \(g_j\)
so that the known matrix \(W_{jk}=\int g_j(u)\psi_k(u)\,\mu(du)\) is
invertible.

\Needspace{5\baselineskip}
\textbf{Proposition 82. Sequential observable-moment operator sets.}

Define \(M_{a,j}=\int g_j(u)B_a(u)\,\mu(du)\). Then

\[
\boxed{
\mathbb E[g_j(U_t)Z_{t+1}\mid\mathcal F_{t-1}]
=M_{a,j}Z_t,
\qquad
C_{a,k}=\sum_j(W^{-1})_{kj}M_{a,j}.
}
\]

Under the rowwise noise and parameter-bound assumptions of Proposition
81, simultaneous confidence ellipsoids for the rows of \(M_{a,j}\) give
a confidence set for every \(C_{a,k}\) by this known linear map. If row
\(i\) of \(\widehat M_{a,j}-M_{a,j}\) has \(V_{a,t}\)-norm at most
\(\beta_{a,j,i}\), then for any direction \(x\),

\[
\left|e_i^{\mathsf T}(\widehat C_{a,k}-C_{a,k})x\right|
\leq
\sum_j |(W^{-1})_{kj}|\,
\beta_{a,j,i}\,\|x\|_{V_{a,t}^{-1}}.
\]

\textbf{Proof.} In the conditional integral, multiply the normalized
state by the observation density. The denominator cancels, giving
\(\int g_j(u)B_a(u)Z_t\,\mu(du)\). The polynomial expansion gives
\(M_{a,j}=\sum_kW_{jk}C_{a,k}\) and hence the inverse formula. Apply
Proposition 81 to each row, using the predictable regressor \(Z_t\) on
visits to action \(a\), or its action-gated version. A union bound over
the predeclared finite rows and moments gives simultaneous coverage.
Independence between these response coordinates is unnecessary. Weighted
Cauchy--Schwarz and the triangle inequality give the displayed
directional bound. \(\square\)

Exact mass constraints can be imposed by intersecting the confidence set
with the known mass affine subspace. This preserves coverage because the
true model is assumed to obey that law. Arbitrarily shrinking a set
around a convenient corrected operator does not have the same
justification. Data-selected geometry can be checked against the whole
resulting set using Proposition 78.

The theorem is conditional on a correct finite observation basis,
accessible predictive state and parameter bounds. It does not obtain a
hidden predictive state from a single dependent path without additional
assumptions. Unknown dictionary discovery, state-estimation error and
polynomial misspecification remain distinct problems. The following
example makes the state accessible by construction rather than
concealing that requirement.

\section{A nonlinear observable state with a curved moment domain}\label{LSM:a-nonlinear-observable-state-with-a-curved-moment-domain}

For \(u\in[-1,1]\), put

\[
v(u)=(1,u,u^2)^{\mathsf T},
\qquad L_2(u)=\frac{3u^2-1}{2}.
\]

The normalized states are moments of a scalar signal:

\[
K=\{(1,x,y):x^2\leq y\leq1\}.
\]

Their cone has the finite curved description

\[
C=\left\{(m,x,y):
\begin{pmatrix}m&x\\x&y\end{pmatrix}\succeq0,
\quad m-y\geq0\right\},
\qquad\lambda=(1,0,0).
\]

Assume that the declared probability measure \(\mu\) on \([-1,1]\) has
moments zero at orders one and three, \(1/3\) at order two and \(1/5\)
at order four. Define a two-parameter degree-four operator density

\[
\boxed{
B_{a,b}(u)=v(u)\,(1,au,bL_2(u)).
}
\]

Here \(a,b\) are parameters, not action labels. This example is
autonomous. Its dependence on the observed signal is nonlinear because
the state includes its square.

\Needspace{5\baselineskip}
\textbf{Proposition 83. Structured curved-domain certificate.} The set
\(K\) is the convex hull of all \(v(u)\). If \(|a|+|b|\leq1-\rho\) for
\(0\leq\rho\leq1\), then \(B_{a,b}(u)\) preserves \(C\) and has exact
integrated mass. On a normalized state, its density is

\[
p_{a,b}(u\mid x,y)=1+a x u+b yL_2(u)\geq\rho.
\]

Every positive-density posterior is exactly \(v(u)\). In particular,
starting from \(v(s)\) produces the observable Markov law

\[
p_{a,b}(u\mid s)=1+a s u+b s^2L_2(u).
\]

Its next conditional moments satisfy

\[
\boxed{
\mathbb E[U_t\mid U_{t-1}=s]=\frac a3s,
\qquad
\mathbb E[L_2(U_t)\mid U_{t-1}=s]=\frac b5s^2.
}
\]

\textbf{Proof.} Every distribution on \([-1,1]\) has \(x^2\leq y\leq1\).
Conversely, fix \(|x|<1\). A point mass at \(x\) gives \((x,x^2)\),
while the endpoint mixture with mean \(x\) gives \((x,1)\). Their convex
combinations give every intermediate \(y\). The cases \(x=\pm1\) are
immediate. Homogenizing yields the matrix and scalar inequalities, which
also force the only mass-zero cone element to be zero.

On \(K\), \(|x|\leq1\), \(0\leq y\leq1\), \(|u|\leq1\) and
\(|L_2(u)|\leq1\). Thus the density is at least \(1-|a|-|b|\). Each
operator output is that nonnegative scalar times \(v(u)\), so lies in
\(C\). The first and second reference moments imply
\(\lambda\int B_{a,b}(u)\,\mu(du)=\lambda\). Normalization cancels the
scalar density, yielding \(v(u)\). Symmetry through order three gives
\(\int uL_2(u)d\mu=0\). The second and fourth moments give
\(\int L_2(u)^2d\mu=1/5\). Integrating \(u\) and \(L_2\) against the
density proves the last two formulas. \(\square\)

The posterior ray lies on a curved boundary. Arbitrary perturbations of
its output coordinates can destroy validity, as in Proposition 72.
Parameter uncertainty inside the declared family behaves differently: it
changes the likelihood but not the posterior ray. A certified rectangle
in \((a,b)\) can therefore be robust even though a free matrix-entry
ball has no positive validity radius.

This is an example of uncertainty that preserves the model's algebraic
structure. It is not a claim that every real system obeys the
two-parameter family. The feature dictionary \(1,u,u^2\), its geometry,
the reference measure and the parameter model are declared inputs to
learning.

\section{A fully rational dependent-data reference experiment}\label{LSM:a-fully-rational-dependent-data-reference-experiment}

Both the continuous uniform measure \(du/2\) and the following five-node
measure satisfy the moment assumptions:

\[
\begin{array}{c|rrrrr}
u&-1&-1/2&0&1/2&1\\\hline
\mu\{u\}&7/90&16/45&2/15&16/45&7/90.
\end{array}
\]

The reference experiment uses the \textbf{five-node measure}, not
rounded observations from the continuous model. At each step, the
conditional probability of a node is its reference weight times
\(p_{a,b}(u\mid s)\). Exact rational probabilities are sampled by
integer draws. The learner receives only the node sequence, the measure
and the declared model. The true parameters are present only in the
separate generator and verification code.

Matching four moments does not identify the two observation laws. The
continuous model gives the set of five nodes probability zero, while the
discrete model gives it probability one. Their single-observation total
variation is therefore one. The continuous theorems and the exact
discrete reference calculation are deliberately distinguished.

For a path of \(n\) transitions, use two scalar regressions,

\[
\begin{aligned}
\phi_{a,t}&=U_{t-1}/3,&Y_{a,t}&=U_t,&\sigma_a&=1,\\
\phi_{b,t}&=U_{t-1}^2/5,&Y_{b,t}&=L_2(U_t),&\sigma_b&=3/4.
\end{aligned}
\]

The noise scales follow from response ranges \([-1,1]\) and
\([-1/2,1]\). The prior parameter bounds are \(|a|,|b|\leq1\). Each
regression uses ridge one and failure probability \(1/40\), giving a
total failure budget \(1/20\) for both parameters at every time.

\Needspace{5\baselineskip}
\textbf{Proposition 84. Exact rational outer confidence intervals.}
Write \(V=1+\sum\phi_t^2\) and \(\widehat\theta=\sum\phi_tY_t/V\). Let
\(k=\lceil\log_2 V\rceil\) and \(h=\lceil\log_2(1/\delta)\rceil\). A
rational upper enclosure of

\[
\sigma\sqrt{\frac{(7/10)(k+2h)}{V}}+\sqrt{\frac1V}
\]

is an upper bound on Proposition 81's scalar radius. Rounding the lower
interval endpoint downward and the upper endpoint upward, then
intersecting with \([-1,1]\), preserves its conditional coverage. Both
intervals cover their true parameters simultaneously at every time with
probability at least \(19/20\).

\textbf{Proof.} The positive exponential series gives
\(e^{7/10}>1+7/10+(7/10)^2/2+(7/10)^3/6+(7/10)^4/24>2\). Hence
\(\log 2<7/10\), so \(\log V+2\log(1/\delta)\leq(7/10)(k+2h)\). Divide
the scalar regression bound by \(\sqrt V\). For a nonnegative rational
\(x\), the smallest integer \(j\) with \((j/10^{12})^2\geq x\) supplies
a rational upper square-root enclosure. Integer square root and an exact
final comparison compute it. Outward rounding only enlarges the
interval. Intersecting with a true prior constraint preserves coverage.
Apply the two-regression union bound on their time-uniform events.
\(\square\)

All counts, sums, radius enclosures and interval endpoints in this
experiment are rational. No floating-point logarithm is trusted in the
confidence arithmetic. An exact replay verifies the arithmetic and the
declared protocol. It does not authenticate the observations, prove that
a physical process satisfies the model, or measure the procedure's
coverage from one pseudorandom run.

\subsection{95.1 Executed result}\label{LSM:executed-result}

The seeded reference path uses \(a=2/5\), \(b=1/5\) and initial
observation zero. It was first run to 65,536 transitions and then
extended along the same seeded path to 262,144. Time-uniform coverage
permits that extension without spending another time-selection budget.
This is not a new independent coverage experiment.

At the final checkpoint, outward-rounded intervals are

\[
\boxed{
\frac{89}{250}\leq a\leq\frac{231}{500},
\qquad
\frac{29}{250}\leq b\leq\frac{29}{100}.
}
\]

Every member of this rectangle has density at least

\[
\rho=1-\frac{231}{500}-\frac{29}{100}=\frac{31}{125}.
\]

The rounded fitted center is \((409/1000,203/1000)\). It lies inside the
rectangle. No reset or other coefficient correction is applied.

At earlier recorded checkpoints the sufficient density-floor test
withheld a whole-box certificate. Its first successful recorded
checkpoint was 32,768 transitions, with floor \(1/100\). This is not a
claim that 32,768 is the smallest possible sample size, or even the
first passing individual step. Only the declared checkpoints were
evaluated. A withheld whole-box certificate does not declare the nominal
fit or true process invalid.

\section{Structured uncertainty has an observable error bound}\label{LSM:structured-uncertainty-has-an-observable-error-bound}

Consider two parameters in a certified family. For normalized
\(z=(1,x,y)\), their density difference is

\[
\Delta p(u)=\Delta a\,x u+\Delta b\,yL_2(u).
\]

The state-update vector is the same \(v(u)\) in both models after
observing \(u\). This does not make the models equivalent. Their
likelihoods and hence their predicted observation laws can differ.

\Needspace{5\baselineskip}
\textbf{Proposition 85. Parameter-to-law error without a
state-estimation term.} Let

\[
c_1=\int|u|\,d\mu,
\qquad c_2=\int|L_2(u)|\,d\mu.
\]

If \(|\Delta a|\leq d_a\) and \(|\Delta b|\leq d_b\), the conditional
observation laws on any common normalized state obey

\[
\boxed{
\|P(\cdot\mid z)-\widehat P(\cdot\mid z)\|_{\mathrm{TV}}
\leq\Delta:=\frac{c_1d_a+c_2d_b}{2}.
}
\]

The same bound holds on every realized common observation history. With
the same initial observation, complete future records of length \(T\)
have total variation at most \(\min(1,T\Delta)\). For the five-node
measure,

\[
c_1=\frac{23}{45},\qquad c_2=\frac{14}{45}.
\]

For \(du/2\), the respective constants are \(1/2\) and \(2/(3\sqrt3)\).

\textbf{Proof.} The first claim follows by integrating \(|\Delta p|/2\),
using \(|x|\leq1\) and \(y\leq1\). After the same observed value both
models have exactly \(v(u)\), so the conditional bound applies again
without accumulated coordinate-state error. Telescoping conditional
kernels gives the record bound, equivalently the common-history coupling
argument of Proposition 69. Direct summation gives the discrete
constants. For the continuous constant \(c_2\), split the integral at
\(u=\pm1/\sqrt3\) and integrate the quadratic polynomial with its
respective signs. \(\square\)

The use of \(\Delta\) here is an upper bound proved from a confidence
set and a model class. It is not the lower discrepancy bound from a
rejected residual test in Chapter~\ref{LSM:book:chapter-36}. A test that does not reject cannot
supply the displayed upper bound by itself.

For the final reference rectangle and fitted center,

\[
d_a=\frac{53}{1000},\qquad d_b=\frac{87}{1000},
\qquad
\boxed{\Delta=\frac{2437}{90000}.}
\]

These bounds compare the fitted model with every member of the
rectangle. They apply to the true process on the conditional coverage
event. Even when that event fails, the rectangle's internal probability
certificate remains an exact statement about its members.

\section{Long-horizon marginal control is not full-record control}\label{LSM:long-horizon-marginal-control-is-not-full-record-control}

The previous record bound grows with \(T\). A different question can
have a uniform answer: how far apart are the distributions of the signal
at a single future time? A lower density bound permits this distinction
to be made exactly.

Let \(P,Q\) be Markov kernels on a common state space. Suppose
\(P(s,\cdot)\geq\rho\mu(\cdot)\) for every \(s\), where \(\mu\) is one
probability measure and \(0<\rho\leq1\). Suppose also
\(\sup_s\|P(s,\cdot)-Q(s,\cdot)\|_{\mathrm{TV}}\leq\Delta\).

\Needspace{5\baselineskip}
\textbf{Proposition 86. Uniform marginal and rolling-window stability.}
For initial laws at distance \(e_0\), the one-time marginal laws satisfy

\[
\boxed{
e_t\leq\min\left(1,
(1-\rho)^te_0+
\frac{\Delta[1-(1-\rho)^t]}{\rho}\right).
}
\]

For \(k\) consecutive observations starting at time \(t\), their joint
laws are at distance at most \(\min(1,e_t+(k-1)\Delta)\). A uniform
small marginal bound does not in general bound the total variation of
arbitrarily long complete records away from one.

\textbf{Proof.} For \(\rho<1\), write
\(P(s,\cdot)=\rho\mu+(1-\rho)R(s,\cdot)\), where \(R\) is a Markov
kernel. Applying a Markov kernel cannot increase total variation, by
integration of bounded test functions. Thus
\(\|\nu P-\widetilde\nu P\|_{\mathrm{TV}}\leq(1-\rho)\|\nu-\widetilde\nu\|_{\mathrm{TV}}\).
Adding and subtracting \(\widetilde\nu P\) yields
\(e_{t+1}\leq(1-\rho)e_t+\Delta\). Iterate this inequality. The case
\(\rho=1\) follows directly. Applying the finite-record telescoping
argument from the initial time \(t\) gives the rolling-window statement.

For the final assertion, take two distinct independent Bernoulli laws
with parameters strictly between zero and one. Each is a Markov kernel
with a positive common minorization. At any time their marginal distance
is the fixed parameter difference. For a full record, test whether its
empirical mean lies above the midpoint between the parameters. Under
each law, Chebyshev's inequality makes the error probability tend to
zero as the record length grows. The total variation of the record laws
therefore tends to one. \(\square\)

This is the standard minorization/contraction mechanism {[}R24{]}, with
its error accounting made explicit. It is not a stability theorem for
arbitrary feedback policies that remember the full history. The
reference model is autonomous and its current scalar observation is the
Markov state. It also does not bound long products of unnormalized
likelihood ratios.

With \(e_0=0\), \(\rho=31/125\) and \(\Delta=2437/90000\), the reference
confidence rectangle gives

\[
\boxed{
\sup_{t\geq0}e_t\leq\frac{2437}{22320}
\approx0.10919.
}
\]

This is a nontrivial all-time bound on \textbf{one-time marginal
observation distributions}. The complete-record bound is still
\(\min(1,2437T/90000)\), which becomes trivial at integer horizon 37.
The model does not exchange those two different claims. In live
filtering after a shared observed value, its state is exactly known and
the next-observation law has the separate uniform upper bound
\(2437/90000\).

\section{A conditional data-to-domain theorem}\label{LSM:a-conditional-data-to-domain-theorem}

\Needspace{5\baselineskip}
\textbf{Proposition 87. Observable structured learning with a covered
domain.} Suppose the scalar process obeys the conditional law of
Proposition 83 with the declared reference measure and fixed parameters
satisfying \(|a|,|b|\leq1\). Use the predictable regressions, fixed
dictionary and confidence protocol of Proposition 84. At a data-selected
finite time, form the full outward-rounded rectangle. If its density
floor \(\rho\) is nonnegative, that rectangle has an exact all-horizon
probability certificate on \(C\). If \(\rho>0\), every chosen center
inside it has the marginal and record bounds of Propositions 85--86
relative to every member.

With probability at least \(19/20\) over the specified stochastic
observation process, these implications cover its true parameter at
every checked time, including a stopping time. Internal validity of a
verified rectangle is deterministic. Statistical coverage remains
conditional on the declared process, moment, parameter and observation
assumptions.

\textbf{Proof.} Proposition 84 provides one simultaneous event on which
both true parameters belong to all returned intervals. Proposition 83
supplies the domain certificate for the entire rectangle whenever its
sufficient density floor is nonnegative. The selected center is one of
its members, and the maximum distances to the endpoints bound each
parameter difference. Proposition 85 gives conditional and record error.
If the floor is positive, all true conditional rows have the same lower
reference component, so Proposition 86 applies to the scalar Markov
process. Selection of the stopping time does not alter the already
simultaneous event. None of these implications verifies the assumptions
from a single sample path. \(\square\)

No claim of empirical coverage follows from the saved run. It is one
seeded finite realization of the exactly declared discrete analogue. The
general confidence result is a theorem, with assumptions recorded in the
certificate. Its independent checker rederives the sufficient
statistics, rational radii, outer endpoints, density floor and error
budgets from the saved path. It rejects a changed reference measure, a
narrower claimed interval, an altered path checksum or an unsupported
claim that the assumptions were independently verified.

This is the point where probability-domain certification and data
uncertainty become connected in an executed dependent-signal example.
The family is more restrictive than an arbitrary learned operator. That
restriction is visible and does useful work: it preserves the nonlinear
posterior curve while uncertainty changes its likelihoods.

\section{Two distinct completed reference sectors}\label{LSM:two-distinct-completed-reference-sectors}

The smaller-repair construction retains the earlier operator family and
geometry. Its coefficient box is deterministic, and its all-horizon
validity is distinct from its finite-horizon comparison bound. The
dependent-data construction learns two parameters of a different,
supplied nonlinear observable family. It certifies a whole confidence
rectangle without reset correction and distinguishes next-step, record
and marginal bounds. Neither example establishes unrestricted
representation discovery.

\section{Frozen release scope and research boundary}\label{LSM:frozen-release-scope-and-research-boundary}

The finite scope contains exact realization and transfer, closure and
memory, conditional observable identification, controlled probability
instruments, finite polyhedral and curved certificates, explicit
approximation budgets, and a restricted dependent-data learning sector.
The reference code and its executed examples are catalogued in the Proof
and Evidence Register. General nonlinear dictionary discovery, unknown
hidden-state learning, empirical sensor modeling and broad deployment
remain research directions, not unstated conclusions or release
prerequisites. Production PR-LASM 2.3.0 is unchanged. This 0.9.0 edition
is a release candidate pending final editorial acceptance and a separate
license decision before open-source publication.


\begingroup
\small\raggedright


\part{Nonlinear interaction quotients and Bayesian refinement}
\setcounter{section}{0}\renewcommand{\thesection}{PRIQ.\arabic{section}}\renewcommand{\theHsection}{PRIQ.\arabic{section}}

\chapter{Interaction structure and exact nonlinear quotients}\label{PRIQ:book:chapter-39}
\section{Introduction}

\subsection{The quotient problem}
Let $X_{\PRIQfine}$ be a fine state space, let $X_{\PRIQcoarse}$ be a coarse state space, and let
\[
  C:X_{\PRIQfine}\twoheadrightarrow X_{\PRIQcoarse}
\]
be a surjective linear map.  A fine trajectory satisfies
\begin{equation}\label{PRIQ:eq:fine-flow-intro}
  \dot x(t)=F_{\PRIQfine}(x(t)),
  \qquad x(0)=x_0,
\end{equation}
and the coarse observable of interest is
\[
  z(t)=Cx(t).
\]
The first question is whether $z(t)$ obeys a closed equation of its own.  Equivalently, can we define a vector field $F_{\PRIQcoarse}$ such that
\begin{equation}\label{PRIQ:eq:vector-field-quotient-intro}
  C F_{\PRIQfine}=F_{\PRIQcoarse}C?
\end{equation}
When this identity holds and both initial value problems have unique solutions, aggregation commutes with evolution:
\begin{equation}\label{PRIQ:eq:flow-quotient-intro}
  C\Phi_t^{\PRIQfine}=\Phi_t^{\PRIQcoarse}C.
\end{equation}
The coarse trajectory is then an exact quotient of the fine one.  No fitting step is involved.

The linear versions of this question are classical.  For finite Markov chains, they appear as lumpability and Markov function criteria \cite{PRIQ:KemenySnell1960,PRIQ:RogersPitman1981,PRIQ:Buchholz1994}.  Nearly decomposable systems lead instead to perturbative aggregation \cite{PRIQ:SimonAndo1961,PRIQ:Meyer1989}.  Polynomial systems and reaction networks admit exact differential equivalences and reductions based on bisimulation \cite{PRIQ:Cardelli2015,PRIQ:Cardelli2017}.  We isolate a quotient test for each order of a multilinear interaction and then follow it through nonlinear composition, continuous evolution, and Bayesian observation.

\subsection{Why hypergraphs appear}
Suppose a local analytic vector field has the expansion
\begin{equation}\label{PRIQ:eq:analytic-expansion-intro}
  F_{\PRIQfine}(x)=\sum_{m\geq0}A_m^{\PRIQfine}[x,\ldots,x].
\end{equation}
The linear coefficient $A_1^{\PRIQfine}$ moves one input packet into one output packet.  The quadratic coefficient $A_2^{\PRIQfine}$ combines two input packets before producing an output.  In general, $A_m^{\PRIQfine}$ combines $m$ inputs.  We record this incidence by the directed hyperedge
\[
  (j_1,\ldots,j_m)\longrightarrow i.
\]

Nonlinear composition is not ordinary path concatenation.  It inserts the output of one multilinear interaction into an input of another.  Repeated insertion produces rooted trees and is governed by the Fa\`a di Bruno formula \cite{PRIQ:ConstantineSavits1996}.  The same tree structure appears in nonlinear evolution and in numerical integration \cite{PRIQ:Butcher1972,PRIQ:HairerLubichWanner2006}.

At order $m$, the central question is
\begin{equation}\label{PRIQ:eq:arity-question}
  \exists A_m^{\PRIQcoarse}
  \quad\text{such that}\quad
  C A_m^{\PRIQfine}=A_m^{\PRIQcoarse}C^{\times m}?
\end{equation}
If this question has an answer at every order, the quotient relation can be carried from individual interactions to the full field, its derivatives, its interaction trees, and its flow.

\subsection{Why Bayesian observations add another nonlinear operation}
Prediction is only one part of a filter.  Given a probability vector $\mu$ and a positive likelihood $\lambda$, Bayesian conditioning is
\begin{equation}\label{PRIQ:eq:bayes-intro}
  \PRIQBup_{\lambda}(\mu)
  =
  \frac{\lambda\odot\mu}{\lambda^T\mu}.
\end{equation}
Likelihood multiplication is linear, but the evidence in the denominator couples all packets.  The full update is therefore nonlinear even when the likelihood acts separately on each state.

Classical sufficiency theory asks whether a statistic preserves all information relevant to a chosen observation family \cite{PRIQ:HalmosSavage1949,PRIQ:Blackwell1953}.  Here the operator question is
\begin{equation}\label{PRIQ:eq:bayes-question}
  C\PRIQBup_{\lambda^{\PRIQfine}}
  =
  \PRIQBup_{\lambda^{\PRIQcoarse}}C.
\end{equation}
When both the flow identity \eqref{PRIQ:eq:flow-quotient-intro} and the Bayesian identity \eqref{PRIQ:eq:bayes-question} hold at each observation time, the whole filtering sequence commutes with aggregation.  Hidden Markov filtering is the linear unary case \cite{PRIQ:BaumPetrie1966}.  The same argument applies when the prediction step comes from a nonlinear quotient flow.

\subsection{Main results}
Eight results form the mathematical spine of the paper.

\begin{enumerate}[label=\textbf{C\arabic*.},leftmargin=3.5em]
\item \textbf{A quotient criterion for every interaction order.}
For a surjective map $C$, a multilinear fine interaction descends exactly when $CA_m^{\PRIQfine}$ vanishes whenever one input lies in $\PRIQKer C$.  Equivalently, its linearization on $X_{\PRIQfine}^{\otimes m}$ factors through $C^{\otimes m}$.

\item \textbf{A concrete packet test.}
For block sum aggregation, exact descent is equivalent to constancy of the aggregated fine coefficients over each product of input fibers.  These common values determine the coarse interaction uniquely.

\item \textbf{Compatibility across several resolutions.}
Exact descent is transitive.  Approximate descent obeys an exact defect identity that can be iterated along a refinement tower.

\item \textbf{Naturality of trees and flows.}
Descent of the coefficients passes to local analytic fields, all derivatives, every rooted interaction tree, and every locally unique flow.

\item \textbf{Naturality after normalization.}
If a positive observation transfer respects the quotient up to a scalar, normalization removes that scalar and preserves the quotient.  For diagonal likelihoods, this happens exactly when the fine likelihood is constant on each fiber, up to one global scale.

\item \textbf{Exact filtering cocycles.}
When both dynamics and observations descend, every finite sequence of predictions and updates agrees after aggregation.  The evidence factors obey the same quotient law.

\item \textbf{A certificate for approximate filtering.}
Interaction mismatch first appears as flow error, while observation mismatch appears at normalization.  A closed recurrence carries both contributions to the final posterior and to the accumulated log evidence.

\item \textbf{A separate pruning certificate.}
Discarded hyperedges remain in a named residual tensor.  Orthogonal packet blocks give an exact residual identity and an optimal retention theorem for a fixed number of hyperedges.
\end{enumerate}

The proof strategy can be summarized by the chain
\[
\boxed{
\begin{gathered}
  \text{multilinear quotient}
  \Longrightarrow
  \text{interaction tree quotient}
  \Longrightarrow
  \text{flow quotient}
  \\
  \Longrightarrow
  \text{Bayesian quotient}
  \Longrightarrow
  \text{filtering cocycle quotient}
  \Longrightarrow
  \text{posterior certificate}.
\end{gathered}}
\]

\subsection{Four distinctions}
The same partition can succeed at one stage of the problem and fail at another.  Four distinctions prevent those cases from being conflated.

\begin{enumerate}[label=(\roman*),leftmargin=2.5em]
\item A \emph{resolution} states which packets are being distinguished.
\item An \emph{exact quotient} proves that the coarse dynamics do not depend on hidden coordinates.
\item A \emph{compression} deletes or approximates interaction blocks, so the omitted part must remain visible as a residual.
\item \emph{Statistical sufficiency} concerns the information retained for a chosen observation family.  Here it is expressed by naturality of the normalized transfer.
\end{enumerate}

A partition may be exact for the dynamics yet lose information carried by the observations.  It may preserve the observation model while failing to close the dynamics.  It may also work through one interaction order and fail at the next.  Each failure receives its own defect.

\subsection{Organization}
\Cref{PRIQ:sec:resolved} introduces packet resolutions and interaction tensors.  \Cref{PRIQ:sec:allarity} proves the quotient criteria.  \Cref{PRIQ:sec:analytic} carries the result to analytic fields, rooted trees, and flows.  \Cref{PRIQ:sec:bayes} studies normalized positive transfer and Bayesian refinement.  \Cref{PRIQ:sec:cocycles} combines prediction and observation into an exact filtering cocycle.  \Cref{PRIQ:sec:approx-dyn,PRIQ:sec:approx-bayes,PRIQ:sec:filter-cert} develop the error bounds.  \Cref{PRIQ:sec:pruning} keeps track of residuals created by hyperedge removal.  \Cref{PRIQ:sec:algorithms,PRIQ:sec:examples} describe the finite algorithms, numerical checks, and worked models.  The final sections compare the framework with nearby theories, then state the limitations and open problems.

\section{Resolved interaction systems}\label{PRIQ:sec:resolved}

\subsection{Finite packet resolutions}
The main algebraic results are stated for vector spaces of finite dimension over $\PRIQR$ or $\PRIQC$.  The quotient and tensor factorizations are then purely algebraic.  The later analytic results add the stated convergence conditions and the usual assumptions for local uniqueness.

\begin{PRIQdefinition}[Packet resolution]\label{PRIQ:def:packet-resolution}
A packet resolution of a vector space $X$ is a finite family of projections $(P_i)_{i\in I}$ satisfying
\[
  P_iP_j=\delta_{ij}P_i,
  \qquad
  \sum_{i\in I}P_i=I_X.
\]
Thus
\[
  X=\bigoplus_{i\in I}X_i,
  \qquad X_i=P_iX.
\]
When $X$ is a Hilbert space, the principal norm identities will assume that the $P_i$ are orthogonal.
\end{PRIQdefinition}

Let $(P_j)_{j\in J}$ resolve a source $X$ and $(Q_i)_{i\in I}$ resolve a target $Y$.

\begin{PRIQdefinition}[Resolved multilinear block and hyperedge]\label{PRIQ:def:resolved-block}
For a multilinear map $A:X^m\to Y$ of order $m$, define
\begin{equation}\label{PRIQ:eq:resolved-block}
  A_{i\mid j_1,\ldots,j_m}
  =
  Q_iA(P_{j_1}\,\cdot,\ldots,P_{j_m}\,\cdot).
\end{equation}
The ordered incidence
\[
  (j_1,\ldots,j_m)\longrightarrow i
\]
is the corresponding directed interaction hyperedge.  We call the hyperedge active when the block in \eqref{PRIQ:eq:resolved-block} is nonzero.
\end{PRIQdefinition}

When $m=1$, the active blocks form an ordinary directed transfer graph.  When $m=2$, each active block joins a pair of input packets.  For general $m$, the block records one interaction with $m$ inputs.

\begin{PRIQtheorem}[Exact multilinear packet reconstruction]\label{PRIQ:thm:packet-reconstruction}
Every multilinear map $A:X^m\to Y$ of order $m$ satisfies
\begin{equation}\label{PRIQ:eq:packet-reconstruction}
  A(x_1,\ldots,x_m)
  =
  \sum_{i,j_1,\ldots,j_m}
  A_{i\mid j_1,\ldots,j_m}
  (P_{j_1}x_1,\ldots,P_{j_m}x_m).
\end{equation}
\end{PRIQtheorem}

\begin{proof}
Insert $I_X=\sum_jP_j$ in every input slot and $I_Y=\sum_iQ_i$ after applying $A$.  Multilinearity distributes all finite sums and gives \eqref{PRIQ:eq:packet-reconstruction}.
\end{proof}

\subsection{Coarse maps, sections, and hidden directions}
Let
\[
  C:X_{\PRIQfine}\twoheadrightarrow X_{\PRIQcoarse}
\]
be a surjective linear coarse map.  Finite dimensionality guarantees a linear right inverse
\[
  E:X_{\PRIQcoarse}\to X_{\PRIQfine},
  \qquad CE=I_{X_{\PRIQcoarse}}.
\]
Define
\begin{equation}\label{PRIQ:eq:representative-hidden-proj}
  P_C=EC,
  \qquad
  Q_C=I_{X_{\PRIQfine}}-EC.
\end{equation}
Then
\[
  P_C^2=P_C,
  \qquad
  Q_C^2=Q_C,
  \qquad
  \PRIQRan Q_C=\PRIQKer C.
\]
Every fine state has the exact decomposition
\begin{equation}\label{PRIQ:eq:representative-hidden-decomposition}
  x=ECx+Q_Cx.
\end{equation}
The first term is the chosen fine representative of the coarse state.  The second lies in $\PRIQKer C$, the hidden subspace.  Exact quotienting means that motion inside this subspace cannot change the aggregated interaction output.

\begin{PRIQdefinition}[Compatible packet aggregation]\label{PRIQ:def:packet-aggregation}
Suppose $(P_i)_{i\in I_{\PRIQfine}}$ and $(\overline P_a)_{a\in I_{\PRIQcoarse}}$ resolve the fine and coarse spaces.  A surjection
\[
  \pi:I_{\PRIQfine}\twoheadrightarrow I_{\PRIQcoarse}
\]
is compatible with $C$ when
\begin{equation}\label{PRIQ:eq:packet-aggregation-compatibility}
  \overline P_aC
  =
  C\sum_{\pi(i)=a}P_i
  \qquad(a\in I_{\PRIQcoarse}).
\end{equation}
\end{PRIQdefinition}

In the scalar coordinate model used later, each $P_i$ selects one coordinate and $C$ sums all coordinates in the same fiber of $\pi$.

\subsection{Multilinear descent}

\begin{PRIQdefinition}[Exact interaction quotient]\label{PRIQ:def:exact-interaction-quotient}
Let $A_m^{\PRIQfine}:X_{\PRIQfine}^m\to X_{\PRIQfine}$ and $A_m^{\PRIQcoarse}:X_{\PRIQcoarse}^m\to X_{\PRIQcoarse}$ be multilinear maps of order $m$.  We say that $A_m^{\PRIQfine}$ \emph{descends through $C$} to $A_m^{\PRIQcoarse}$, and write
\[
  A_m^{\PRIQfine}\Downarrow_C A_m^{\PRIQcoarse},
\]
when
\begin{equation}\label{PRIQ:eq:exact-interaction-quotient}
  C A_m^{\PRIQfine}
  =
  A_m^{\PRIQcoarse}C^{\times m}.
\end{equation}
For $m=0$, descent means $CA_0^{\PRIQfine}=A_0^{\PRIQcoarse}$.
\end{PRIQdefinition}

Here
\[
  C^{\times m}(x_1,\ldots,x_m)
  =(Cx_1,\ldots,Cx_m).
\]
We write $C^{\otimes m}$ for the induced linear map between algebraic tensor powers.

\begin{PRIQdefinition}[Interaction defect]\label{PRIQ:def:interaction-defect}
For a proposed coarse tensor $B_m$, define
\begin{equation}\label{PRIQ:eq:interaction-defect}
  \Delta_m[B_m]
  =
  CA_m^{\PRIQfine}-B_mC^{\times m}.
\end{equation}
Exact descent is the condition $\Delta_m[B_m]=0$.
\end{PRIQdefinition}

\subsection{Norm conventions}
For a multilinear map $A:X^m\to Y$ of order $m$, let
\[
  \PRIQnorm{A}
  =
  \sup_{\PRIQnorm{x_r}\leq1}
  \PRIQnorm{A(x_1,\ldots,x_m)}.
\]
If $E$ is a fixed right inverse of $C$, then any exact coarse tensor necessarily satisfies
\begin{equation}\label{PRIQ:eq:coarse-tensor-bound-preview}
  A_m^{\PRIQcoarse}(z_1,\ldots,z_m)
  =
  CA_m^{\PRIQfine}(Ez_1,\ldots,Ez_m),
\end{equation}
so
\[
  \PRIQnorm{A_m^{\PRIQcoarse}}
  \leq
  \PRIQnorm{C}\PRIQnorm{A_m^{\PRIQfine}}\PRIQnorm{E}^m.
\]
The formula also makes the later uniqueness argument immediate.

\section{Hypergraph lumpability at every interaction order}\label{PRIQ:sec:allarity}

\subsection{Invariant quotient criterion}
The quotient problem has a simple algebraic answer.

\begin{PRIQtheorem}[Universal multilinear quotient criterion]\label{PRIQ:thm:universal-quotient}
Let $C:X_{\PRIQfine}\twoheadrightarrow X_{\PRIQcoarse}$ be surjective and let
\[
  A_m^{\PRIQfine}:X_{\PRIQfine}^m\to X_{\PRIQfine}
\]
be a multilinear map of order $m$, where $m\geq1$.  The following statements are equivalent.

\begin{enumerate}[label=(\roman*),leftmargin=2.8em]
\item There exists a multilinear map $A_m^{\PRIQcoarse}:X_{\PRIQcoarse}^m\to X_{\PRIQcoarse}$ of order $m$ such that
\[
  CA_m^{\PRIQfine}=A_m^{\PRIQcoarse}C^{\times m}.
\]

\item For every slot $r\in\{1,\ldots,m\}$,
\begin{equation}\label{PRIQ:eq:slot-kernel-condition}
  CA_m^{\PRIQfine}(x_1,\ldots,x_{r-1},k,x_{r+1},\ldots,x_m)=0
\end{equation}
whenever $k\in\PRIQKer C$.

\item The map $CA_m^{\PRIQfine}$ is constant on products of fibers: if $Cx_r=Cx_r'$ for all $r$, then
\begin{equation}\label{PRIQ:eq:fiber-constancy-abstract}
  CA_m^{\PRIQfine}(x_1,\ldots,x_m)
  =
  CA_m^{\PRIQfine}(x_1',\ldots,x_m').
\end{equation}

\item Let $\widehat A_m^{\PRIQfine}:X_{\PRIQfine}^{\otimes m}\to X_{\PRIQcoarse}$ be the linearization of $CA_m^{\PRIQfine}$.  Then
\begin{equation}\label{PRIQ:eq:tensor-kernel-condition}
  \PRIQKer(C^{\otimes m})
  \subseteq
  \PRIQKer(\widehat A_m^{\PRIQfine}).
\end{equation}
\end{enumerate}

When these conditions hold, the coarse tensor is unique and, for every right inverse $E$ of $C$,
\begin{equation}\label{PRIQ:eq:coarse-via-section}
  A_m^{\PRIQcoarse}(z_1,\ldots,z_m)
  =
  CA_m^{\PRIQfine}(Ez_1,\ldots,Ez_m).
\end{equation}
The right side is independent of the chosen right inverse.
\end{PRIQtheorem}

\begin{proof}
Assume (i).  If $k\in\PRIQKer C$, then inserting $k$ into any input slot gives
\[
  CA_m^{\PRIQfine}(x_1,\ldots,k,\ldots,x_m)
  =
  A_m^{\PRIQcoarse}(Cx_1,\ldots,0,\ldots,Cx_m)=0,
\]
so (ii) holds.

Assume (ii).  If $Cx_r=Cx_r'$, then $x_r'-x_r\in\PRIQKer C$.  Replace the inputs one at a time.  Multilinearity and \eqref{PRIQ:eq:slot-kernel-condition} show that each replacement leaves $CA_m^{\PRIQfine}$ unchanged.  Hence (iii) holds.

Assume (iii).  For $z_r\in X_{\PRIQcoarse}$ choose arbitrary lifts $x_r$ with $Cx_r=z_r$ and set
\[
  A_m^{\PRIQcoarse}(z_1,\ldots,z_m)
  =
  CA_m^{\PRIQfine}(x_1,\ldots,x_m).
\]
Fiber constancy makes the definition independent of the lifts.  Multilinearity is inherited from $A_m^{\PRIQfine}$, and the defining identity gives (i).  Surjectivity of $C$ also gives uniqueness: every coarse input tuple has a fine lift.

To compare (ii) and (iv), identify a multilinear map of order $m$ with its linearization on the algebraic tensor product.  Since $C$ is surjective,
\[
  \PRIQKer(C^{\otimes m})
  =
  \sum_{r=1}^m
  X_{\PRIQfine}^{\otimes(r-1)}\otimes\PRIQKer C\otimes
  X_{\PRIQfine}^{\otimes(m-r)}.
\]
Choose a splitting $X_{\PRIQfine}=E(X_{\PRIQcoarse})\oplus\PRIQKer C$ and expand the tensor power.  Vanishing of $\widehat A_m^{\PRIQfine}$ on this sum is exactly the slot condition.  Thus (ii) and (iv) are equivalent.  Formula \eqref{PRIQ:eq:coarse-via-section} follows by choosing $x_r=Ez_r$.
\end{proof}

\begin{PRIQremark}[Nonlinear congruence interpretation]\label{PRIQ:rem:nonlinear-congruence}
For a general map $F:X_{\PRIQfine}\to X_{\PRIQfine}$, descent through $C$ is equivalent to
\[
  Cx=Cy
  \quad\Longrightarrow\quad
  CF(x)=CF(y).
\]
Thus the relation $x\sim_Cy$, defined by $Cx=Cy$, must be a congruence for the coarse output of $F$.  \Cref{PRIQ:thm:universal-quotient} turns that statement into a condition on each multilinear coefficient.
\end{PRIQremark}

\begin{PRIQcorollary}[Bounded coarse tensor]\label{PRIQ:cor:bounded-coarse-tensor}
Under the hypotheses of \Cref{PRIQ:thm:universal-quotient}, for every bounded right inverse $E$,
\[
  \PRIQnorm{A_m^{\PRIQcoarse}}
  \leq
  \PRIQnorm{C}\PRIQnorm{A_m^{\PRIQfine}}\PRIQnorm{E}^m.
\]
\end{PRIQcorollary}

\begin{proof}
Apply \eqref{PRIQ:eq:coarse-via-section} to unit coarse inputs.
\end{proof}

\subsection{Criterion in packet coordinates}
For a block sum aggregation, the kernel condition can be read directly from the tensor coefficients.

Let
\[
  I_{\PRIQfine}=\{1,\ldots,n\},
  \qquad
  I_{\PRIQcoarse}=\{1,\ldots,q\},
\]
and let
\[
  \pi:I_{\PRIQfine}\twoheadrightarrow I_{\PRIQcoarse}.
\]
Write $(e_i)$ and $(f_a)$ for the standard bases of $\PRIQR^n$ and $\PRIQR^q$.  Define
\begin{equation}\label{PRIQ:eq:block-sum-C}
  Ce_i=f_{\pi(i)},
  \qquad
  (Cx)_a=\sum_{\pi(i)=a}x_i.
\end{equation}
Suppose
\begin{equation}\label{PRIQ:eq:fine-coordinates}
  A_m^{\PRIQfine}(e_{j_1},\ldots,e_{j_m})
  =
  \sum_{i=1}^n
  a^{\PRIQfine,(m)}_{i;j_1,\ldots,j_m}e_i.
\end{equation}
For each coarse output packet $a$, sum the fine coefficients over its output fiber and define
\begin{equation}\label{PRIQ:eq:aggregated-coefficient}
  g^{(m)}_{a;j_1,\ldots,j_m}
  =
  \sum_{i:\pi(i)=a}
  a^{\PRIQfine,(m)}_{i;j_1,\ldots,j_m}.
\end{equation}

\begin{PRIQtheorem}[Coefficient criterion for every interaction order]\label{PRIQ:thm:packet-coordinate-criterion}
The fine tensor $A_m^{\PRIQfine}$ descends through the block sum map \eqref{PRIQ:eq:block-sum-C} if and only if, for every coarse output $a$,
\begin{equation}\label{PRIQ:eq:packet-fiber-constancy}
  g^{(m)}_{a;j_1,\ldots,j_m}
  =
  g^{(m)}_{a;j_1',\ldots,j_m'}
\end{equation}
whenever
\[
  \pi(j_r)=\pi(j_r')
  \qquad(r=1,\ldots,m).
\]
When this condition holds, the unique coarse coefficients are
\begin{equation}\label{PRIQ:eq:coarse-coordinate-formula}
  a^{\PRIQcoarse,(m)}_{a;b_1,\ldots,b_m}
  =
  g^{(m)}_{a;j_1,\ldots,j_m},
  \qquad
  \pi(j_r)=b_r.
\end{equation}
\end{PRIQtheorem}

\begin{proof}
Applying $C$ to \eqref{PRIQ:eq:fine-coordinates} gives
\[
  CA_m^{\PRIQfine}(e_{j_1},\ldots,e_{j_m})
  =
  \sum_{a=1}^q
  g^{(m)}_{a;j_1,\ldots,j_m}f_a.
\]
If a coarse tensor exists, this vector equals
\[
  A_m^{\PRIQcoarse}(f_{\pi(j_1)},\ldots,f_{\pi(j_m)}),
\]
which depends only on the coarse input labels.  This proves necessity.  Conversely, constancy on each fiber makes \eqref{PRIQ:eq:coarse-coordinate-formula} well defined.  Equality on every tuple of fine basis vectors then gives
\[
  CA_m^{\PRIQfine}=A_m^{\PRIQcoarse}C^{\times m}
\]
by multilinearity.
\end{proof}

For $m=2$, the criterion has two useful views.  The hypergraph records which input fibers interact, while the coefficient matrix shows whether those interactions are constant on each product fiber.  The exact synthetic tensor used in the certificate run has the required block structure to floating point precision.

\begin{figure}[H]
\centering
\begin{subfigure}[t]{0.57\linewidth}
\centering
\includegraphics[width=\linewidth]{sources/research/hypergraph-bayesian-refinement/figures/hypergraph_quotient.pdf}
\caption{Fine interactions descend to one coarse hyperedge.}
\label{PRIQ:fig:hypergraph-quotient}
\end{subfigure}
\hfill
\begin{subfigure}[t]{0.39\linewidth}
\centering
\includegraphics[width=\linewidth]{sources/research/hypergraph-bayesian-refinement/figures/fiber_constancy_map.pdf}
\caption{Aggregated coefficients are constant on every product fiber.}
\label{PRIQ:fig:fiber-constancy-map}
\end{subfigure}
\caption{Two views of the quadratic quotient criterion.  The incidence view identifies the descended hyperedge.  The coefficient view displays the fiber constancy condition from \Cref{PRIQ:thm:packet-coordinate-criterion}.}
\label{PRIQ:fig:quadratic-quotient-views}
\end{figure}

\subsection{Unary and quadratic cases}

\begin{PRIQcorollary}[Unary aggregate rate criterion]\label{PRIQ:cor:unary-criterion}
For $m=1$, exact descent is equivalent to
\begin{equation}\label{PRIQ:eq:unary-aggregate-rate}
  \sum_{i:\pi(i)=a}a^{\PRIQfine,(1)}_{i;j}
  =
  \sum_{i:\pi(i)=a}a^{\PRIQfine,(1)}_{i;j'}
\end{equation}
for all coarse outputs $a$ whenever $\pi(j)=\pi(j')$.
\end{PRIQcorollary}

\begin{proof}
Set $m=1$ in \Cref{PRIQ:thm:packet-coordinate-criterion}.  The condition then involves one input fiber and becomes \eqref{PRIQ:eq:unary-aggregate-rate}.
\end{proof}

Let $Q_{\PRIQfine}$ be the generator of a continuous time Markov chain, using the column vector convention.  The left side of \eqref{PRIQ:eq:unary-aggregate-rate} is the total transition rate from fine state $j$ into coarse packet $a$.  Corollary~\ref{PRIQ:cor:unary-criterion} is therefore the standard criterion for strong lumpability.

\begin{PRIQcorollary}[Quadratic hypergraph criterion]\label{PRIQ:cor:quadratic-criterion}
For $m=2$, exact descent is equivalent to
\[
  \sum_{i:\pi(i)=a}a^{\PRIQfine,(2)}_{i;j,k}
  =
  \sum_{i:\pi(i)=a}a^{\PRIQfine,(2)}_{i;j',k'}
\]
whenever
\[
  \pi(j)=\pi(j'),
  \qquad
  \pi(k)=\pi(k').
\]
The resulting coarse interaction is the directed quadratic hyperedge
\[
  (\pi(j),\pi(k))\longrightarrow a.
\]
\end{PRIQcorollary}

\begin{proof}
Set $m=2$ in \Cref{PRIQ:thm:packet-coordinate-criterion}.  Fiber constancy in the two input slots gives the displayed equality, and \eqref{PRIQ:eq:coarse-coordinate-formula} assigns its common value to the coarse hyperedge.
\end{proof}

\begin{PRIQexample}[A quotient with invisible internal interactions]\label{PRIQ:ex:invisible-interactions}
Let $X_{\PRIQcoarse}=\PRIQR^q$, $X_{\PRIQfine}=\PRIQR^{2q}$, and
\[
  C=\begin{bmatrix}I_q&I_q\end{bmatrix},
  \qquad
  E=\frac12\begin{bmatrix}I_q\\I_q\end{bmatrix}.
\]
Let $A_m^{\PRIQcoarse}$ be arbitrary and let $K_m:X_{\PRIQfine}^m\to X_{\PRIQfine}$ satisfy
\[
  CK_m=0.
\]
Then
\begin{equation}\label{PRIQ:eq:exact-lift-construction}
  A_m^{\PRIQfine}
  =
  EA_m^{\PRIQcoarse}C^{\times m}+K_m
\end{equation}
satisfies
\[
  CA_m^{\PRIQfine}=A_m^{\PRIQcoarse}C^{\times m}.
\]
The tensor $K_m$ may contain many fine hyperedges, so the fine model can remain complicated.  Exact quotienting asks only that every hidden contribution disappear after applying $C$.
\end{PRIQexample}

\subsection{Structural consequences}

\begin{PRIQproposition}[Symmetry descends]\label{PRIQ:prop:symmetry-descends}
If $A_m^{\PRIQfine}$ is symmetric in its input slots and descends through $C$, then the unique coarse tensor $A_m^{\PRIQcoarse}$ is symmetric.
\end{PRIQproposition}

\begin{proof}
Choose lifts of the coarse inputs.  Permuting the coarse inputs permutes their lifts, and symmetry of $A_m^{\PRIQfine}$ leaves \eqref{PRIQ:eq:coarse-via-section} unchanged.
\end{proof}

\begin{PRIQproposition}[Quotient support]\label{PRIQ:prop:quotient-support}
Under the coordinate criterion, every active coarse hyperedge
\[
  (b_1,\ldots,b_m)\longrightarrow a
\]
is supported by at least one fine hyperedge whose input and output labels map to that coarse incidence.  The converse can fail because several active fine hyperedges may cancel after aggregation.
\end{PRIQproposition}

\begin{proof}
If every fine coefficient in the corresponding incidence fiber vanishes, then the aggregate coefficient \eqref{PRIQ:eq:aggregated-coefficient} vanishes.  Conversely, nonzero fine coefficients can sum to zero after output aggregation.
\end{proof}

\subsection{Composition across refinement levels}
Let
\[
  X_2\xrightarrow{C_{21}}X_1\xrightarrow{C_{10}}X_0,
  \qquad
  C_{20}=C_{10}C_{21}.
\]

\begin{PRIQproposition}[Transitivity of exact interaction quotients]\label{PRIQ:prop:quotient-transitivity}
If
\[
  A_m^{(2)}\Downarrow_{C_{21}}A_m^{(1)}
  \qquad\text{and}\qquad
  A_m^{(1)}\Downarrow_{C_{10}}A_m^{(0)},
\]
then
\[
  A_m^{(2)}\Downarrow_{C_{20}}A_m^{(0)}.
\]
\end{PRIQproposition}

\begin{proof}
Compose the two identities:
\[
  C_{20}A_m^{(2)}
  =C_{10}A_m^{(1)}C_{21}^{\times m}
  =A_m^{(0)}C_{10}^{\times m}C_{21}^{\times m}
  =A_m^{(0)}C_{20}^{\times m}.
\]
\end{proof}

Approximate descent has an exact composition law.

\begin{PRIQtheorem}[Interaction defect across two levels]\label{PRIQ:thm:multilevel-defect}
Define
\[
  \Delta_m^{21}
  =C_{21}A_m^{(2)}-A_m^{(1)}C_{21}^{\times m},
\]
\[
  \Delta_m^{10}
  =C_{10}A_m^{(1)}-A_m^{(0)}C_{10}^{\times m},
\]
and
\[
  \Delta_m^{20}
  =C_{20}A_m^{(2)}-A_m^{(0)}C_{20}^{\times m}.
\]
Then
\begin{equation}\label{PRIQ:eq:multilevel-defect}
  \boxed{
  \Delta_m^{20}
  =
  C_{10}\Delta_m^{21}
  +
  \Delta_m^{10}C_{21}^{\times m}.}
\end{equation}
Consequently,
\begin{equation}\label{PRIQ:eq:multilevel-defect-norm}
  \PRIQnorm{\Delta_m^{20}}
  \leq
  \PRIQnorm{C_{10}}\PRIQnorm{\Delta_m^{21}}
  +
  \PRIQnorm{\Delta_m^{10}}\PRIQnorm{C_{21}}^m.
\end{equation}
\end{PRIQtheorem}

\begin{proof}
Add and subtract $C_{10}A_m^{(1)}C_{21}^{\times m}$:
\begin{align*}
\Delta_m^{20}
&=C_{10}C_{21}A_m^{(2)}-A_m^{(0)}C_{10}^{\times m}C_{21}^{\times m}\\
&=C_{10}\bigl(C_{21}A_m^{(2)}-A_m^{(1)}C_{21}^{\times m}\bigr)\\
&\quad+
\bigl(C_{10}A_m^{(1)}-A_m^{(0)}C_{10}^{\times m}\bigr)C_{21}^{\times m}.
\end{align*}
This proves \eqref{PRIQ:eq:multilevel-defect}.  Taking norms and using submultiplicativity yields the estimate.
\end{proof}


\chapter{Analytic fields, rooted trees, and quotient flows}\label{PRIQ:book:chapter-40}
\section{Analytic fields, rooted trees, and quotient flows}\label{PRIQ:sec:analytic}

\subsection{Descent of polynomial and analytic vector fields}
Let
\begin{equation}\label{PRIQ:eq:fine-analytic-field}
  F_{\PRIQfine}(x)
  =
  \sum_{m=0}^{\infty}A_m^{\PRIQfine}[x^{\otimes m}]
\end{equation}
be locally analytic, where $A_m^{\PRIQfine}$ is a symmetric multilinear map of order $m$ for $m\geq1$ and $A_0^{\PRIQfine}\in X_{\PRIQfine}$.  Throughout,
\[
  A_m[x^{\otimes m}]
  =A_m(x,\ldots,x).
\]

\begin{PRIQtheorem}[Analytic interaction descent]\label{PRIQ:thm:analytic-descent}
Let $C:X_{\PRIQfine}\twoheadrightarrow X_{\PRIQcoarse}$ have a bounded right inverse $E$.  Suppose the coefficient series converges normally on every smaller ball,
\begin{equation}\label{PRIQ:eq:fine-normal-summability}
  \sum_{m=0}^{\infty}\PRIQnorm{A_m^{\PRIQfine}}r^m<\infty
  \qquad(0\leq r<R),
\end{equation}
and, for every $m\geq0$, there is a coarse coefficient $A_m^{\PRIQcoarse}$ satisfying
\begin{equation}\label{PRIQ:eq:coefficientwise-descent}
  CA_m^{\PRIQfine}=A_m^{\PRIQcoarse}C^{\times m}.
\end{equation}
Then
\begin{equation}\label{PRIQ:eq:coarse-analytic-field}
  F_{\PRIQcoarse}(z)
  =
  \sum_{m=0}^{\infty}A_m^{\PRIQcoarse}[z^{\otimes m}]
\end{equation}
converges absolutely whenever $\PRIQnorm{E}\PRIQnorm{z}<R$, and
\begin{equation}\label{PRIQ:eq:analytic-field-intertwining}
  \boxed{CF_{\PRIQfine}(x)=F_{\PRIQcoarse}(Cx)}
\end{equation}
for every $x$ in the fine convergence domain whose coarse image lies in the coarse domain.
\end{PRIQtheorem}

\begin{proof}
For $m\geq1$, Corollary~\ref{PRIQ:cor:bounded-coarse-tensor} gives
\[
  \PRIQnorm{A_m^{\PRIQcoarse}}
  \leq
  \PRIQnorm{C}\PRIQnorm{A_m^{\PRIQfine}}\PRIQnorm{E}^m.
\]
For $m=0$, the identity $A_0^{\PRIQcoarse}=CA_0^{\PRIQfine}$ gives the same bound with $\PRIQnorm{E}^0=1$.
Therefore
\[
  \sum_{m\geq0}\PRIQnorm{A_m^{\PRIQcoarse}}\PRIQnorm{z}^m
  \leq
  \PRIQnorm{C}
  \sum_{m\geq0}
  \PRIQnorm{A_m^{\PRIQfine}}
  (\PRIQnorm{E}\PRIQnorm{z})^m,
\]
which converges when $\PRIQnorm{E}\PRIQnorm{z}<R$.  Coefficientwise descent and absolute convergence permit termwise aggregation:
\[
  CF_{\PRIQfine}(x)
  =\sum_{m\geq0}CA_m^{\PRIQfine}[x^{\otimes m}]
  =\sum_{m\geq0}A_m^{\PRIQcoarse}[(Cx)^{\otimes m}]
  =F_{\PRIQcoarse}(Cx).
\]
\end{proof}

\begin{PRIQremark}[Coefficientwise descent is sufficient, not always necessary]
Descent of each displayed coefficient is enough to make the whole series descend.  The converse can fail for a noncanonical expansion that splits one homogeneous degree into several terms whose coarse contributions cancel.  For the Taylor expansion at a fixed base point, the derivatives are unique, so the converse does hold.  Differentiating $CF_{\PRIQfine}=F_{\PRIQcoarse}C$ then recovers every descended jet.
\end{PRIQremark}

\begin{PRIQtheorem}[Coefficientwise characterization of polynomial descent]\label{PRIQ:thm:polynomial-coefficientwise}
Let
\[
  F_{\PRIQfine}(x)=\sum_{m=0}^{M}A_m^{\PRIQfine}[x^{\otimes m}],
  \qquad
  F_{\PRIQcoarse}(z)=\sum_{m=0}^{M}A_m^{\PRIQcoarse}[z^{\otimes m}],
\]
where every $A_m^{\PRIQfine}$ and $A_m^{\PRIQcoarse}$ is symmetric.  Then
\[
  CF_{\PRIQfine}=F_{\PRIQcoarse}C
\]
if and only if
\[
  CA_m^{\PRIQfine}=A_m^{\PRIQcoarse}C^{\times m}
  \qquad(0\leq m\leq M).
\]
\end{PRIQtheorem}

\begin{proof}
The reverse implication follows from \Cref{PRIQ:thm:analytic-descent}, applied to a finite sum.  Conversely, fix $x\in X_{\PRIQfine}$ and evaluate the field identity at $tx$:
\[
  \sum_{m=0}^{M}t^mCA_m^{\PRIQfine}[x^{\otimes m}]
  =
  \sum_{m=0}^{M}t^mA_m^{\PRIQcoarse}[(Cx)^{\otimes m}].
\]
Equality of the coefficients of this polynomial with values in $X_{\PRIQcoarse}$ gives equality on the diagonal for the two symmetric multilinear maps of order $m$.  The polarization identity recovers each multilinear map from its diagonal polynomial, so
\[
  CA_m^{\PRIQfine}=A_m^{\PRIQcoarse}C^{\times m}
\]
for every degree.
\end{proof}

\subsection{Naturality of derivatives}

\begin{PRIQproposition}[Jet naturality]\label{PRIQ:prop:jet-naturality}
Let $F_{\PRIQfine}$ and $F_{\PRIQcoarse}$ be $C^r$ maps satisfying
\[
  CF_{\PRIQfine}=F_{\PRIQcoarse}C.
\]
Then, for every $1\leq k\leq r$,
\begin{equation}\label{PRIQ:eq:jet-naturality}
  \boxed{
  C D^kF_{\PRIQfine}(x)[h_1,\ldots,h_k]
  =
  D^kF_{\PRIQcoarse}(Cx)[Ch_1,\ldots,Ch_k].}
\end{equation}
\end{PRIQproposition}

\begin{proof}
Differentiate the identity $CF_{\PRIQfine}(x)=F_{\PRIQcoarse}(Cx)$.  Since $C$ is linear, its derivatives of order greater than one vanish and the ordinary chain rule gives the first order statement.  Repeated differentiation gives \eqref{PRIQ:eq:jet-naturality}.
\end{proof}

Thus the quotient preserves more than the value of the vector field.  At each base point, every fine derivative block maps to its coarse counterpart.  The same identity controls local perturbations, stability calculations, and higher order response.

\subsection{Rooted interaction trees}
Repeated substitution of nonlinear interactions produces rooted trees.  Only a simple recursive definition is needed for the quotient argument.

\begin{PRIQdefinition}[Interaction tree evaluation]\label{PRIQ:def:tree-evaluation}
Let $\PRIQcA=(A_m)_{m\geq1}$ be a family of multilinear maps $A_m:X^m\to X$.  A rooted interaction tree is defined recursively.  A leaf $\bullet$ has evaluation
\[
  \PRIQcV_{\bullet}=I_X.
\]
If
\[
  \tau=B_+(\tau_1,\ldots,\tau_m)
\]
is formed by adjoining a root of arity $m$, define
\begin{equation}\label{PRIQ:eq:tree-evaluation}
  \PRIQcV_{\tau}(x_1,\ldots,x_{|\tau|})
  =
  A_m\bigl(
  \PRIQcV_{\tau_1}(\mathbf x_1),\ldots,
  \PRIQcV_{\tau_m}(\mathbf x_m)
  \bigr),
\end{equation}
where $|\tau|$ denotes the number of leaves and the input list is partitioned among the subtrees in their fixed order.
\end{PRIQdefinition}

The operation $B_+$ inserts several rooted trees beneath a new root.  In packet coordinates, every edge carries a packet label and every internal vertex must match an active resolved hyperedge.

\begin{PRIQtheorem}[Treewise quotient naturality]\label{PRIQ:thm:treewise-quotient}
Suppose every operation in a fine family descends:
\[
  CA_m^{\PRIQfine}=A_m^{\PRIQcoarse}C^{\times m}
  \qquad(m\geq1).
\]
Then, for every rooted interaction tree $\tau$,
\begin{equation}\label{PRIQ:eq:treewise-quotient}
  \boxed{
  C\PRIQcV_{\tau}^{\PRIQfine}
  =
  \PRIQcV_{\tau}^{\PRIQcoarse}C^{\times |\tau|}.}
\end{equation}
\end{PRIQtheorem}

\begin{proof}
Proceed by induction on the number of internal vertices.  The leaf case is $CI=IC$.  Suppose
\[
  \tau=B_+(\tau_1,\ldots,\tau_m).
\]
Using interaction descent at the root and the inductive hypothesis on every subtree,
\begin{align*}
C\PRIQcV_{\tau}^{\PRIQfine}
&=
CA_m^{\PRIQfine}
\bigl(
\PRIQcV_{\tau_1}^{\PRIQfine},\ldots,
\PRIQcV_{\tau_m}^{\PRIQfine}
\bigr)\\
&=
A_m^{\PRIQcoarse}
\bigl(
C\PRIQcV_{\tau_1}^{\PRIQfine},\ldots,
C\PRIQcV_{\tau_m}^{\PRIQfine}
\bigr)\\
&=
A_m^{\PRIQcoarse}
\bigl(
\PRIQcV_{\tau_1}^{\PRIQcoarse}C^{\times|\tau_1|},\ldots,
\PRIQcV_{\tau_m}^{\PRIQcoarse}C^{\times|\tau_m|}
\bigr)\\
&=
\PRIQcV_{\tau}^{\PRIQcoarse}C^{\times|\tau|}.
\end{align*}
\end{proof}

\begin{PRIQcorollary}[Compatibility with Fa\`a di Bruno composition]\label{PRIQ:cor:faa-quotient}
Let $F_{\PRIQfine}$ and $F_{\PRIQcoarse}$ be smooth maps with $CF_{\PRIQfine}=F_{\PRIQcoarse}C$, and let $G_{\PRIQfine}$ and $G_{\PRIQcoarse}$ satisfy the same relation.  Then each rooted tree in the Fa\`a di Bruno expansion, written in packet coordinates, descends to the corresponding coarse tree.  Moreover,
\[
  C(G_{\PRIQfine}\circ F_{\PRIQfine})
  =(G_{\PRIQcoarse}\circ F_{\PRIQcoarse})C.
\]
\end{PRIQcorollary}

\begin{proof}
The map identity follows directly by composition.  The termwise statement follows from \Cref{PRIQ:thm:treewise-quotient}, because the Fa\`a di Bruno expansion is a finite sum of rooted derivative trees at each derivative order.
\end{proof}

\subsection{Exact quotient flows}

\begin{PRIQtheorem}[Exact quotient flow]\label{PRIQ:thm:exact-quotient-flow}
Let $U_{\PRIQfine}\subseteq X_{\PRIQfine}$ and $U_{\PRIQcoarse}\subseteq X_{\PRIQcoarse}$ be open, with $C(U_{\PRIQfine})\subseteq U_{\PRIQcoarse}$.  Suppose $F_{\PRIQfine}$ and $F_{\PRIQcoarse}$ are locally Lipschitz and satisfy
\[
  CF_{\PRIQfine}=F_{\PRIQcoarse}C
\]
on $U_{\PRIQfine}$.  Let $\Phi_t^{\PRIQfine}(x)$ and $\Phi_t^{\PRIQcoarse}(z)$ denote the unique maximal local flows.  Whenever both sides are defined,
\begin{equation}\label{PRIQ:eq:exact-flow-identity}
  \boxed{
  C\Phi_t^{\PRIQfine}(x)
  =
  \Phi_t^{\PRIQcoarse}(Cx).}
\end{equation}
\end{PRIQtheorem}

\begin{proof}
Let $x(t)=\Phi_t^{\PRIQfine}(x_0)$ and set $z(t)=Cx(t)$.  Then
\[
  \dot z(t)
  =C\dot x(t)
  =CF_{\PRIQfine}(x(t))
  =F_{\PRIQcoarse}(Cx(t))
  =F_{\PRIQcoarse}(z(t)),
\]
with $z(0)=Cx_0$.  Uniqueness for the coarse initial value problem gives $z(t)=\Phi_t^{\PRIQcoarse}(Cx_0)$.
\end{proof}

\begin{PRIQcorollary}[Linear semigroup naturality]\label{PRIQ:cor:linear-semigroup-naturality}
If $CA_{\PRIQfine}=A_{\PRIQcoarse}C$, then
\begin{equation}\label{PRIQ:eq:linear-semigroup-naturality}
  Ce^{tA_{\PRIQfine}}
  =
  e^{tA_{\PRIQcoarse}}C
\end{equation}
for every $t$.
\end{PRIQcorollary}

\begin{proof}
Either apply \Cref{PRIQ:thm:exact-quotient-flow} to the linear systems or sum the exponential series using $CA_{\PRIQfine}^k=A_{\PRIQcoarse}^kC$.
\end{proof}

\begin{PRIQcorollary}[Quotients of discrete dynamics]\label{PRIQ:cor:discrete-quotient}
If $CF_{\PRIQfine}=F_{\PRIQcoarse}C$, then
\[
  CF_{\PRIQfine}^{\circ n}
  =
  F_{\PRIQcoarse}^{\circ n}C
\]
for every integer $n\geq0$.
\end{PRIQcorollary}

\begin{proof}
Induct on $n$.
\end{proof}

\subsection{Positivity and simplex preservation}
Let
\[
  \Delta_{\PRIQfine}
  =\{\mu\in\PRIQR_+^n:\PRIQone^T\mu=1\},
  \qquad
  \Delta_{\PRIQcoarse}
  =\{\nu\in\PRIQR_+^q:\PRIQone^T\nu=1\}.
\]
For block sum aggregation, $C\Delta_{\PRIQfine}=\Delta_{\PRIQcoarse}$ and
\[
  \PRIQone_{\PRIQcoarse}^TC=\PRIQone_{\PRIQfine}^T.
\]

\begin{PRIQproposition}[Simplex preservation descends]\label{PRIQ:prop:simplex-descends}
Suppose $\Phi_t^{\PRIQfine}(\Delta_{\PRIQfine})\subseteq\Delta_{\PRIQfine}$ and exact flow naturality holds.  Then
\[
  \Phi_t^{\PRIQcoarse}(\Delta_{\PRIQcoarse})\subseteq\Delta_{\PRIQcoarse}.
\]
\end{PRIQproposition}

\begin{proof}
For any $\nu\in\Delta_{\PRIQcoarse}$, choose $\mu\in\Delta_{\PRIQfine}$ with $C\mu=\nu$.  Then
\[
  \Phi_t^{\PRIQcoarse}(\nu)
  =C\Phi_t^{\PRIQfine}(\mu)
  \in C\Delta_{\PRIQfine}
  =\Delta_{\PRIQcoarse}.
\]
\end{proof}

\begin{PRIQboundary}[Global existence of the flow]
Exact descent does not guarantee global solutions.  A fine trajectory may blow up in finite time, leave the domain, or reach a region where the vector field is no longer Lipschitz.  \Cref{PRIQ:thm:exact-quotient-flow} applies only while both solutions exist and remain unique.
\end{PRIQboundary}


\chapter{Positive observations and exact filtering}\label{PRIQ:book:chapter-41}
\section{Normalized positive transfer and Bayesian refinement}\label{PRIQ:sec:bayes}

\subsection{Positive state spaces}
The Bayesian result rests on a simple fact about normalizing positive operators.

Let $X_{\PRIQfine}$ and $X_{\PRIQcoarse}$ be finite dimensional ordered vector spaces with positive cones $K_{\PRIQfine}$ and $K_{\PRIQcoarse}$.  Let
\[
  \ell_{\PRIQfine}:X_{\PRIQfine}\to\PRIQR,
  \qquad
  \ell_{\PRIQcoarse}:X_{\PRIQcoarse}\to\PRIQR
\]
be positive mass functionals.  Define normalized state spaces
\[
  \PRIQcS_{\PRIQfine}
  =\{x\in K_{\PRIQfine}:\ell_{\PRIQfine}(x)=1\},
  \qquad
  \PRIQcS_{\PRIQcoarse}
  =\{z\in K_{\PRIQcoarse}:\ell_{\PRIQcoarse}(z)=1\}.
\]
Assume $C$ is positive and preserves mass:
\begin{equation}\label{PRIQ:eq:mass-preserving-C}
  \ell_{\PRIQcoarse}C=\ell_{\PRIQfine}.
\end{equation}

\begin{PRIQdefinition}[Normalized positive transfer]\label{PRIQ:def:normalized-transfer}
For a positive linear operator $L:X\to X$, define
\begin{equation}\label{PRIQ:eq:normalized-transfer}
  \PRIQNrm_{\ell}(L)(x)
  =
  \frac{Lx}{\ell(Lx)}
\end{equation}
whenever $\ell(Lx)>0$.
\end{PRIQdefinition}

\subsection{Normalization naturality}

\begin{PRIQtheorem}[Naturality of normalized transfer]\label{PRIQ:thm:normalized-transfer-naturality}
Let $L_{\PRIQfine}:X_{\PRIQfine}\to X_{\PRIQfine}$ and $L_{\PRIQcoarse}:X_{\PRIQcoarse}\to X_{\PRIQcoarse}$ be positive linear operators.  Suppose
\begin{equation}\label{PRIQ:eq:positive-intertwining}
  CL_{\PRIQfine}=sL_{\PRIQcoarse}C
\end{equation}
for some $s>0$, and suppose \eqref{PRIQ:eq:mass-preserving-C} holds.  Then, on the common domain of normalization,
\begin{equation}\label{PRIQ:eq:normalized-transfer-naturality}
  \boxed{
  C\PRIQNrm_{\ell_{\PRIQfine}}(L_{\PRIQfine})
  =
  \PRIQNrm_{\ell_{\PRIQcoarse}}(L_{\PRIQcoarse})C.}
\end{equation}
\end{PRIQtheorem}

\begin{proof}
For a positive state $x$,
\begin{align*}
C\PRIQNrm_{\ell_{\PRIQfine}}(L_{\PRIQfine})(x)
&=
\frac{CL_{\PRIQfine}x}{\ell_{\PRIQfine}(L_{\PRIQfine}x)}\\
&=
\frac{sL_{\PRIQcoarse}Cx}{\ell_{\PRIQcoarse}(CL_{\PRIQfine}x)}\\
&=
\frac{sL_{\PRIQcoarse}Cx}{s\ell_{\PRIQcoarse}(L_{\PRIQcoarse}Cx)}\\
&=
\PRIQNrm_{\ell_{\PRIQcoarse}}(L_{\PRIQcoarse})(Cx).
\end{align*}
\end{proof}

\begin{PRIQremark}[Gauge invariance]
Normalization removes the scale $s$.  At the operator level, this is the familiar fact that a positive rescaling of a likelihood leaves the posterior unchanged.
\end{PRIQremark}

\begin{PRIQproposition}[Composition law for normalized transfers]\label{PRIQ:prop:normalized-composition}
For positive operators $L_1,L_2$ on the same ordered state space,
\begin{equation}\label{PRIQ:eq:normalized-composition}
  \PRIQNrm_{\ell}(L_2)
  \circ
  \PRIQNrm_{\ell}(L_1)
  =
  \PRIQNrm_{\ell}(L_2L_1)
\end{equation}
wherever both sides are defined.
\end{PRIQproposition}

\begin{proof}
The first normalization contributes the same scalar to the numerator and denominator of the second update, so it cancels.
\end{proof}

\subsection{Classical Bayesian update}
On the finite simplex, let
\[
  L_{\lambda}=\PRIQdiag(\lambda_1,\ldots,\lambda_n),
  \qquad
  \lambda_i>0.
\]
With $\ell(x)=\PRIQone^Tx$,
\[
  \PRIQNrm_{\ell}(L_{\lambda})(\mu)
  =
  \frac{\lambda\odot\mu}{\lambda^T\mu}
  =\PRIQBup_{\lambda}(\mu).
\]
The evidence is
\[
  Z_{\lambda}(\mu)=\lambda^T\mu.
\]
The composition law becomes
\begin{equation}\label{PRIQ:eq:bayes-composition}
  \PRIQBup_{\lambda_2}\circ\PRIQBup_{\lambda_1}
  =
  \PRIQBup_{\lambda_1\odot\lambda_2}.
\end{equation}

\subsection{Exact Bayesian refinement}
Let $C$ be the block sum aggregation associated with a surjection $\pi:\{1,\ldots,n\}\twoheadrightarrow\{1,\ldots,q\}$.

\begin{PRIQtheorem}[Sufficient condition for Bayesian refinement]\label{PRIQ:thm:bayes-sufficient}
Let $\lambda^{\PRIQfine}\in\PRIQR_{>0}^n$ and $\lambda^{\PRIQcoarse}\in\PRIQR_{>0}^q$.  If
\begin{equation}\label{PRIQ:eq:fiber-likelihood-sufficient}
  \lambda^{\PRIQfine}=sC^T\lambda^{\PRIQcoarse}
\end{equation}
for some $s>0$, then
\begin{equation}\label{PRIQ:eq:exact-bayes-refinement}
  \boxed{
  C\PRIQBup_{\lambda^{\PRIQfine}}
  =
  \PRIQBup_{\lambda^{\PRIQcoarse}}C.}
\end{equation}
Moreover,
\begin{equation}\label{PRIQ:eq:evidence-scaling}
  Z_{\lambda^{\PRIQfine}}(\mu)
  =
  sZ_{\lambda^{\PRIQcoarse}}(C\mu).
\end{equation}
\end{PRIQtheorem}

\begin{proof}
Equation \eqref{PRIQ:eq:fiber-likelihood-sufficient} is equivalent to
\[
  CL_{\lambda^{\PRIQfine}}
  =sL_{\lambda^{\PRIQcoarse}}C.
\]
Apply \Cref{PRIQ:thm:normalized-transfer-naturality}.  The evidence identity follows from
\[
  (\lambda^{\PRIQfine})^T\mu
  =s(\lambda^{\PRIQcoarse})^TC\mu.
\]
\end{proof}

Condition \eqref{PRIQ:eq:fiber-likelihood-sufficient} says that the fine likelihood is constant on each aggregation fiber, apart from one global rescaling.  When the coarse space has at least two packets, the next theorem shows that this condition is also necessary if commutation is required for every interior prior.

\begin{PRIQtheorem}[Necessary condition for Bayesian refinement]\label{PRIQ:thm:bayes-necessity}
Assume $q\geq2$, every fiber of $\pi$ is nonempty, and $\lambda^{\PRIQfine},\lambda^{\PRIQcoarse}$ are strictly positive.  If
\begin{equation}\label{PRIQ:eq:bayes-equality-all-priors}
  C\PRIQBup_{\lambda^{\PRIQfine}}(\mu)
  =
  \PRIQBup_{\lambda^{\PRIQcoarse}}(C\mu)
\end{equation}
for every $\mu$ in the interior of $\Delta_n$, then there exists $s>0$ such that
\[
  \lambda^{\PRIQfine}=sC^T\lambda^{\PRIQcoarse}.
\]
\end{PRIQtheorem}

\begin{proof}
Both sides of \eqref{PRIQ:eq:bayes-equality-all-priors} extend continuously to any prior with nonzero evidence.  Strict positivity guarantees nonzero evidence throughout the closed simplex.  Fix two distinct coarse packets $a$ and $b$.  Choose $i\in\pi^{-1}(a)$ and $k\in\pi^{-1}(b)$, then consider a prior supported on $\{i,k\}$ with masses $u,v>0$.  On the fine side, the ratio of the posterior masses of packets $a$ and $b$ is
\[
  \frac{\lambda_i^{\PRIQfine}u}{\lambda_k^{\PRIQfine}v},
\]
whereas the coarse posterior ratio is
\[
  \frac{\lambda_a^{\PRIQcoarse}u}{\lambda_b^{\PRIQcoarse}v}.
\]
Equality for all $u,v>0$ gives
\begin{equation}\label{PRIQ:eq:likelihood-ratio-necessity}
  \frac{\lambda_i^{\PRIQfine}}{\lambda_k^{\PRIQfine}}
  =
  \frac{\lambda_a^{\PRIQcoarse}}{\lambda_b^{\PRIQcoarse}}.
\end{equation}
Take another state $j$ in the same fiber as $i$.  Applying the same argument to $j$ and $k$ gives $\lambda_j^{\PRIQfine}=\lambda_i^{\PRIQfine}$.  The fine likelihood is therefore constant on each fiber.  Write its value on fiber $a$ as $\overline\lambda_a$.  Equation \eqref{PRIQ:eq:likelihood-ratio-necessity} now gives
\[
  \frac{\overline\lambda_a}{\overline\lambda_b}
  =
  \frac{\lambda_a^{\PRIQcoarse}}{\lambda_b^{\PRIQcoarse}}
\]
for every pair $a,b$.  Hence $\overline\lambda=s\lambda^{\PRIQcoarse}$ for one $s>0$, proving the result.
\end{proof}

\begin{PRIQremark}[The case of one coarse packet]
If $q=1$, both sides of \eqref{PRIQ:eq:bayes-equality-all-priors} equal the unique coarse state, regardless of the fine likelihood.  Necessity therefore requires a nontrivial coarse state space.
\end{PRIQremark}

\subsection{Differential structure of Bayesian normalization}

\begin{PRIQproposition}[Bayesian differential]\label{PRIQ:prop:bayes-differential}
At an interior prior $\mu$,
\begin{equation}\label{PRIQ:eq:bayes-differential}
  D\PRIQBup_{\lambda}(\mu)[h]
  =
  \frac{\lambda\odot h}{Z_{\lambda}(\mu)}
  -
  \PRIQBup_{\lambda}(\mu)
  \frac{\lambda^Th}{Z_{\lambda}(\mu)}.
\end{equation}
\end{PRIQproposition}

\begin{proof}
Differentiate the quotient $L_{\lambda}\mu/(\PRIQone^TL_{\lambda}\mu)$.
\end{proof}

\begin{PRIQcorollary}[Bayesian jet naturality]\label{PRIQ:cor:bayes-jet-naturality}
Under the hypotheses of \Cref{PRIQ:thm:bayes-sufficient}, for every $r\geq1$ and every interior fine prior,
\begin{equation}\label{PRIQ:eq:bayes-jet-naturality}
  C D^r\PRIQBup_{\lambda^{\PRIQfine}}(\mu)
  [h_1,\ldots,h_r]
  =
  D^r\PRIQBup_{\lambda^{\PRIQcoarse}}(C\mu)
  [Ch_1,\ldots,Ch_r].
\end{equation}
\end{PRIQcorollary}

\begin{proof}
Differentiate the exact map identity \eqref{PRIQ:eq:exact-bayes-refinement} and apply \Cref{PRIQ:prop:jet-naturality}.
\end{proof}

Thus the normalization map has a descended interaction structure of its own.  Its first derivative couples all evidence packets through a rank one term.  Higher derivatives describe interactions among larger groups of packets.

\section{Filtering cocycles}\label{PRIQ:sec:cocycles}

\subsection{Finite observation sequences}
Let $0=t_0<t_1<\cdots<t_N$ and define prediction maps
\[
  \Phi_k^{\PRIQfine}=\Phi_{t_k-t_{k-1}}^{\PRIQfine},
  \qquad
  \Phi_k^{\PRIQcoarse}=\Phi_{t_k-t_{k-1}}^{\PRIQcoarse}.
\]
For observations $y_1,\ldots,y_N$, let
\[
  \PRIQBup_k^{\PRIQfine}=\PRIQBup_{\lambda_{y_k}^{\PRIQfine}},
  \qquad
  \PRIQBup_k^{\PRIQcoarse}=\PRIQBup_{\lambda_{y_k}^{\PRIQcoarse}}.
\]
Pair each prediction with the observation that follows it:
\[
  \Gamma_k^{\PRIQfine}
  =\PRIQBup_k^{\PRIQfine}\circ\Phi_k^{\PRIQfine},
  \qquad
  \Gamma_k^{\PRIQcoarse}
  =\PRIQBup_k^{\PRIQcoarse}\circ\Phi_k^{\PRIQcoarse},
\]
and define the resulting filtering cocycles
\begin{equation}\label{PRIQ:eq:filter-cocycles}
  \Psi_{1:N}^{\PRIQfine}
  =\Gamma_N^{\PRIQfine}\circ\cdots\circ\Gamma_1^{\PRIQfine},
  \qquad
  \Psi_{1:N}^{\PRIQcoarse}
  =\Gamma_N^{\PRIQcoarse}\circ\cdots\circ\Gamma_1^{\PRIQcoarse}.
\end{equation}

\begin{PRIQtheorem}[Exact naturality of the filtering cocycle]\label{PRIQ:thm:exact-cocycle}
Suppose, for every $k$,
\begin{equation}\label{PRIQ:eq:step-flow-naturality}
  C\Phi_k^{\PRIQfine}=\Phi_k^{\PRIQcoarse}C
\end{equation}
and
\begin{equation}\label{PRIQ:eq:step-bayes-naturality}
  C\PRIQBup_k^{\PRIQfine}=\PRIQBup_k^{\PRIQcoarse}C.
\end{equation}
Then
\begin{equation}\label{PRIQ:eq:exact-cocycle-naturality}
  \boxed{
  C\Psi_{1:N}^{\PRIQfine}
  =
  \Psi_{1:N}^{\PRIQcoarse}C.}
\end{equation}
\end{PRIQtheorem}

\begin{proof}
For one step,
\[
  C\Gamma_k^{\PRIQfine}
  =C\PRIQBup_k^{\PRIQfine}\Phi_k^{\PRIQfine}
  =\PRIQBup_k^{\PRIQcoarse}C\Phi_k^{\PRIQfine}
  =\PRIQBup_k^{\PRIQcoarse}\Phi_k^{\PRIQcoarse}C
  =\Gamma_k^{\PRIQcoarse}C.
\]
Iterating this identity through \eqref{PRIQ:eq:filter-cocycles} proves the result.
\end{proof}

\begin{PRIQcorollary}[Exact posterior aggregation]\label{PRIQ:cor:exact-posterior}
If $\mu_0^{\PRIQcoarse}=C\mu_0^{\PRIQfine}$, then
\[
  \boxed{C\mu_N^{\PRIQfine}=\mu_N^{\PRIQcoarse}}
\]
for every finite observation sequence, where
\[
  \mu_N^{\PRIQfine}=\Psi_{1:N}^{\PRIQfine}(\mu_0^{\PRIQfine}),
  \qquad
  \mu_N^{\PRIQcoarse}=\Psi_{1:N}^{\PRIQcoarse}(\mu_0^{\PRIQcoarse}).
\]
\end{PRIQcorollary}

\begin{proof}
Apply \eqref{PRIQ:eq:exact-cocycle-naturality} to $\mu_0^{\PRIQfine}$ and substitute $\mu_0^{\PRIQcoarse}=C\mu_0^{\PRIQfine}$.
\end{proof}

\subsection{Sequential evidence}
Let
\[
  \widehat\mu_k^{\PRIQfine}
  =\Phi_k^{\PRIQfine}(\mu_{k-1}^{\PRIQfine}),
  \qquad
  \widehat\mu_k^{\PRIQcoarse}
  =\Phi_k^{\PRIQcoarse}(\mu_{k-1}^{\PRIQcoarse}),
\]
and define the step evidence
\[
  z_k^{\PRIQfine}
  =(\lambda_{y_k}^{\PRIQfine})^T\widehat\mu_k^{\PRIQfine},
  \qquad
  z_k^{\PRIQcoarse}
  =(\lambda_{y_k}^{\PRIQcoarse})^T\widehat\mu_k^{\PRIQcoarse}.
\]
Set
\[
  Z_{1:N}^{\PRIQfine}=\prod_{k=1}^N z_k^{\PRIQfine},
  \qquad
  Z_{1:N}^{\PRIQcoarse}=\prod_{k=1}^N z_k^{\PRIQcoarse}.
\]

\begin{PRIQtheorem}[Evidence factorization under exact refinement]\label{PRIQ:thm:evidence-factorization}
Assume the hypotheses of \Cref{PRIQ:thm:exact-cocycle}, let
\[
  \mu_0^{\PRIQcoarse}=C\mu_0^{\PRIQfine},
\]
and suppose
\[
  \lambda_{y_k}^{\PRIQfine}
  =s_kC^T\lambda_{y_k}^{\PRIQcoarse}
  \qquad(s_k>0).
\]
Then
\begin{equation}\label{PRIQ:eq:step-evidence-factor}
  z_k^{\PRIQfine}=s_kz_k^{\PRIQcoarse}
\end{equation}
for every step, and the likelihood of the observation sequence satisfies
\begin{equation}\label{PRIQ:eq:total-evidence-factor}
  \boxed{
  Z_{1:N}^{\PRIQfine}
  =
  \left(\prod_{k=1}^Ns_k\right)
  Z_{1:N}^{\PRIQcoarse}.}
\end{equation}
Under the gauge $s_k=1$, the fine and coarse sequential evidences agree exactly.
\end{PRIQtheorem}

\begin{proof}
Initial compatibility gives $C\mu_0^{\PRIQfine}=\mu_0^{\PRIQcoarse}$.  For later steps, \Cref{PRIQ:cor:exact-posterior} gives the same identity after every update.  Flow naturality therefore yields
\[
  C\widehat\mu_k^{\PRIQfine}=\widehat\mu_k^{\PRIQcoarse}.
\]
Therefore
\[
  z_k^{\PRIQfine}
  =s_k(\lambda_{y_k}^{\PRIQcoarse})^TC\widehat\mu_k^{\PRIQfine}
  =s_kz_k^{\PRIQcoarse}.
\]
The likelihood of the complete observation sequence is the product of the conditional evidence terms.  Multiplying the identities gives \eqref{PRIQ:eq:total-evidence-factor}.
\end{proof}

\subsection{Hidden Markov and nonlinear specializations}

\begin{PRIQcorollary}[Hidden Markov lumpability]\label{PRIQ:cor:hmm-lumpability}
Let $M_{\PRIQfine}$ and $M_{\PRIQcoarse}$ be column stochastic matrices satisfying
\[
  CM_{\PRIQfine}=M_{\PRIQcoarse}C,
\]
and let every observation likelihood satisfy
\[
  \lambda_y^{\PRIQfine}=s_yC^T\lambda_y^{\PRIQcoarse}.
\]
Then, for every observation sequence $y_{1:N}$,
\begin{align*}
&C\bigl(
\PRIQBup_{\lambda_{y_N}^{\PRIQfine}}\circ M_{\PRIQfine}
\circ\cdots\circ
\PRIQBup_{\lambda_{y_1}^{\PRIQfine}}\circ M_{\PRIQfine}
\bigr)\\
&\qquad=
\bigl(
\PRIQBup_{\lambda_{y_N}^{\PRIQcoarse}}\circ M_{\PRIQcoarse}
\circ\cdots\circ
\PRIQBup_{\lambda_{y_1}^{\PRIQcoarse}}\circ M_{\PRIQcoarse}
\bigr)C.
\end{align*}
\end{PRIQcorollary}

\begin{proof}
Apply \Cref{PRIQ:thm:exact-cocycle} with the linear prediction maps $M_{\PRIQfine}$ and $M_{\PRIQcoarse}$.
\end{proof}

\begin{PRIQcorollary}[Nonlinear quotient filtering]\label{PRIQ:cor:nonlinear-filtering}
Suppose the fine and coarse prediction maps are flow segments of locally Lipschitz vector fields satisfying $CF_{\PRIQfine}=F_{\PRIQcoarse}C$, and suppose all observation likelihoods are constant on each fiber up to scale.  Then the nonlinear filtering cocycle commutes with $C$ exactly.
\end{PRIQcorollary}

\begin{proof}
Intertwining of the vector fields gives $C\Phi_k^{\PRIQfine}=\Phi_k^{\PRIQcoarse}C$ by \Cref{PRIQ:thm:exact-quotient-flow}.  Likelihoods that are constant on each fiber give $C\PRIQBup_k^{\PRIQfine}=\PRIQBup_k^{\PRIQcoarse}C$ by \Cref{PRIQ:thm:bayes-sufficient}.  Apply \Cref{PRIQ:thm:exact-cocycle}.
\end{proof}

\begin{PRIQremark}[Two independent closure requirements]\label{PRIQ:rem:two-closure-requirements}
The dynamical and observational conditions are independent.  Exact interaction lumpability says nothing about whether the likelihood distinguishes states inside a packet.  Conversely, an observation may ignore those distinctions even when the dynamics depend on them.  Exact coarse filtering needs both commutative squares.
\end{PRIQremark}

\begin{PRIQboundary}[Sufficiency for the chosen observation family]
The condition $\lambda_y^{\PRIQfine}=s_yC^T\lambda_y^{\PRIQcoarse}$ is sufficient for the specified finite likelihood family.  It does not assert that $C$ is sufficient for every possible future observation model.  Adding a likelihood that distinguishes two states inside one fiber can invalidate the coarse filter without changing the dynamics.
\end{PRIQboundary}


\chapter{Dynamical and Bayesian approximation}\label{PRIQ:book:chapter-42}
\section{Approximate interaction quotients}\label{PRIQ:sec:approx-dyn}

Exact quotienting is an algebraic yes or no question.  A useful model also needs to say how far it misses when the answer is no.  We separate three sources of error:
\[
  \text{interaction mismatch},
  \qquad
  \text{dependence on hidden directions},
  \qquad
  \text{dynamical amplification}.
\]

\subsection{Best coarse interaction and quotient defect}
Let
\[
  \mathfrak Q_m(C)
  =
  \{B_mC^{\times m}:B_m\in\PRIQcL^m(X_{\PRIQcoarse};X_{\PRIQcoarse})\}
\]
be the subspace of outputs from multilinear maps of order $m$ that factor through the coarse space.

\begin{PRIQdefinition}[Best quotient error]\label{PRIQ:def:best-quotient-error}
For a fine interaction $A_m^{\PRIQfine}$, define
\begin{equation}\label{PRIQ:eq:best-quotient-error}
  \varepsilon_m^\star
  =
  \inf_{B_m\in\PRIQcL^m(X_{\PRIQcoarse};X_{\PRIQcoarse})}
  \PRIQnorm{CA_m^{\PRIQfine}-B_mC^{\times m}}.
\end{equation}
\end{PRIQdefinition}

\begin{PRIQproposition}[Attainment in finite dimension]\label{PRIQ:prop:best-quotient-attained}
The infimum in \eqref{PRIQ:eq:best-quotient-error} is attained.
\end{PRIQproposition}

\begin{proof}
The map
\[
  B_m\longmapsto B_mC^{\times m}
\]
is injective because $C$ is surjective.  Its range is a finite dimensional linear subspace of $\PRIQcL^m(X_{\PRIQfine};X_{\PRIQcoarse})$, hence it is closed.  In a finite dimensional normed space, the distance to a nonempty closed subspace attains its minimum.  A nearest point in this range therefore exists.  Any such point has a unique preimage under the displayed map, and that preimage is a minimizing coarse tensor.
\end{proof}

The minimum value depends on the chosen norm.  A minimizing tensor need not be unique when that norm is not strictly convex.  Exact lumpability is independent of the norm because it is the condition $\varepsilon_m^\star=0$.

\subsection{Approximation induced by a section}
A chosen right inverse supplies a concrete coarse tensor even when exact descent fails.

\begin{PRIQdefinition}[Coarse tensor induced by a section]\label{PRIQ:def:section-coarse-tensor}
Fix $E$ with $CE=I$.  Define
\begin{equation}\label{PRIQ:eq:section-coarse-tensor}
  A_m^{\PRIQcoarse,E}
  =
  CA_m^{\PRIQfine}E^{\times m}.
\end{equation}
\end{PRIQdefinition}

Let $P_C=EC$ and $Q_C=I-EC$ as in \eqref{PRIQ:eq:representative-hidden-proj}.

\begin{PRIQproposition}[Expansion of the section defect by hidden inputs]\label{PRIQ:prop:hidden-slot-expansion}
The defect of \eqref{PRIQ:eq:section-coarse-tensor} is
\begin{equation}\label{PRIQ:eq:section-defect}
  \Delta_m^E(x_1,\ldots,x_m)
  =
  CA_m^{\PRIQfine}(x_1,\ldots,x_m)
  -
  CA_m^{\PRIQfine}(P_Cx_1,\ldots,P_Cx_m),
\end{equation}
and admits the exact expansion
\begin{equation}\label{PRIQ:eq:hidden-slot-sum}
  \boxed{
  \Delta_m^E(x_1,\ldots,x_m)
  =
  \sum_{\varnothing\neq S\subseteq\{1,\ldots,m\}}
  CA_m^{\PRIQfine}(y_1^S,\ldots,y_m^S),}
\end{equation}
where
\[
  y_r^S=
  \begin{cases}
  Q_Cx_r,&r\in S,\\
  P_Cx_r,&r\notin S.
  \end{cases}
\]
\end{PRIQproposition}

\begin{proof}
Equation \eqref{PRIQ:eq:section-defect} follows from
\[
  A_m^{\PRIQcoarse,E}(Cx_1,\ldots,Cx_m)
  =CA_m^{\PRIQfine}(ECx_1,\ldots,ECx_m).
\]
Now expand every input as $x_r=P_Cx_r+Q_Cx_r$.  The term with no $Q_C$ factor is the term subtracted in \eqref{PRIQ:eq:section-defect}.  Every remaining term has at least one hidden input and appears exactly once.
\end{proof}

\begin{PRIQcorollary}[Exactness as disappearance of hidden interactions]\label{PRIQ:cor:hidden-interactions-exactness}
The tensor defined from the chosen section descends exactly if and only if the sum in \eqref{PRIQ:eq:hidden-slot-sum} vanishes for all inputs.  The stronger slotwise condition
\[
  CA_m^{\PRIQfine}(x_1,\ldots,Q_Cx_r,\ldots,x_m)=0
\]
for each slot is equivalent to exact descent and makes every summand vanish separately.
\end{PRIQcorollary}

\begin{proof}
The first statement is \eqref{PRIQ:eq:hidden-slot-sum} with $\Delta_m^E=0$.  Since $\PRIQRan Q_C=\PRIQKer C$, the condition in each slot is exactly the kernel condition in \Cref{PRIQ:thm:universal-quotient}.  It is therefore equivalent to exact descent.  Every term indexed by a nonempty subset contains at least one hidden input, so each term vanishes.
\end{proof}

\subsection{Defect majorants for the full field}
Suppose fine and proposed coarse fields have expansions
\[
  F_{\PRIQfine}(x)=\sum_{m\geq0}A_m^{\PRIQfine}[x^{\otimes m}],
  \qquad
  F_{\PRIQcoarse}(z)=\sum_{m\geq0}A_m^{\PRIQcoarse}[z^{\otimes m}].
\]
Define the field defect
\begin{equation}\label{PRIQ:eq:field-defect}
  \delta_F(x)
  =
  CF_{\PRIQfine}(x)-F_{\PRIQcoarse}(Cx).
\end{equation}
Whenever both series converge absolutely,
\begin{equation}\label{PRIQ:eq:field-defect-expansion}
  \delta_F(x)
  =
  \sum_{m\geq0}\Delta_m[x^{\otimes m}].
\end{equation}

\begin{PRIQproposition}[Analytic defect majorant]\label{PRIQ:prop:defect-majorant}
If $\PRIQnorm{\Delta_m}\leq\varepsilon_m$ and the scalar series
\[
  \mathfrak d_F(R)=\sum_{m\geq0}\varepsilon_mR^m
\]
converges, then
\begin{equation}\label{PRIQ:eq:field-defect-majorant}
  \boxed{
  \sup_{\PRIQnorm{x}\leq R}\PRIQnorm{\delta_F(x)}
  \leq
  \mathfrak d_F(R).}
\end{equation}
\end{PRIQproposition}

\begin{proof}
Apply the multilinear norm bound to every term in \eqref{PRIQ:eq:field-defect-expansion} and sum the convergent majorant.
\end{proof}

The sequence $(\varepsilon_m)$ shows which interaction orders obstruct the quotient.  A model may be exactly lumpable through order $r$ and fail for the first time at order $r+1$.

\subsection{Propagation of dynamical defect}

\begin{PRIQtheorem}[Approximate quotient flow bound]\label{PRIQ:thm:approx-flow-bound}
Let
\[
  \dot x=F_{\PRIQfine}(x),
  \qquad
  \dot z=F_{\PRIQcoarse}(z),
\]
and suppose both solutions remain in domains on which
\[
  \PRIQnorm{F_{\PRIQcoarse}(u)-F_{\PRIQcoarse}(v)}
  \leq L\PRIQnorm{u-v}.
\]
Set
\[
  e(t)=\PRIQnorm{Cx(t)-z(t)}.
\]
Then
\begin{equation}\label{PRIQ:eq:approx-flow-general}
  \boxed{
  e(t)
  \leq
  e^{Lt}e(0)
  +
  \int_0^t e^{L(t-s)}\PRIQnorm{\delta_F(x(s))}\PRIQdd s.}
\end{equation}
If $e(0)=0$ and $\PRIQnorm{\delta_F(x(s))}\leq\varepsilon_F$ on $[0,t]$, then
\begin{equation}\label{PRIQ:eq:approx-flow-uniform}
  e(t)
  \leq
  \begin{cases}
  \displaystyle
  \varepsilon_F\frac{e^{Lt}-1}{L},&L>0,\\[1.1ex]
  t\varepsilon_F,&L=0.
  \end{cases}
\end{equation}
\end{PRIQtheorem}

\begin{proof}
Let $y(t)=Cx(t)$.  From \eqref{PRIQ:eq:field-defect},
\[
  \dot y(t)=F_{\PRIQcoarse}(y(t))+\delta_F(x(t)).
\]
Subtract the integral equations for $y$ and $z$:
\[
  y(t)-z(t)
  =y(0)-z(0)
  +\int_0^t
  \bigl(F_{\PRIQcoarse}(y(s))-F_{\PRIQcoarse}(z(s))\bigr)\PRIQdd s
  +\int_0^t\delta_F(x(s))\PRIQdd s.
\]
Taking norms and applying the Grönwall inequality yields \eqref{PRIQ:eq:approx-flow-general}.  The uniform cases follow by evaluating the exponential integral.
\end{proof}

\begin{PRIQcorollary}[Coefficient certificate for a flow segment]\label{PRIQ:cor:coefficient-flow-certificate}
If $\PRIQnorm{x(s)}\leq R$ and Proposition~\ref{PRIQ:prop:defect-majorant} applies, then a flow interval of length $T$ satisfies
\[
  \PRIQnorm{C\Phi_T^{\PRIQfine}(x_0)-\Phi_T^{\PRIQcoarse}(Cx_0)}
  \leq
  \mathfrak d_F(R)\frac{e^{LT}-1}{L}
\]
for $L>0$, with $T\mathfrak d_F(R)$ at $L=0$.
\end{PRIQcorollary}

\begin{proof}
By Proposition~\ref{PRIQ:prop:defect-majorant}, $\PRIQnorm{\delta_F(x(s))}\leq\mathfrak d_F(R)$ throughout the segment.  Apply \Cref{PRIQ:thm:approx-flow-bound} with compatible initial data, so $e(0)=0$, and evaluate the exponential integral.
\end{proof}

\subsection{Refinement towers}
Let
\[
  X_L\xrightarrow{C_{L,L-1}}X_{L-1}
  \longrightarrow\cdots\longrightarrow
  X_1\xrightarrow{C_{10}}X_0.
\]
For $a>b$, write
\[
  C_{ab}=C_{b+1,b}\cdots C_{a,a-1}.
\]
Let $\Delta_m^{ab}$ denote the interaction defect for the composite map $C_{ab}$.  Iterating \Cref{PRIQ:thm:multilevel-defect} yields a telescoping formula.

\begin{PRIQcorollary}[Telescoping interaction defect]\label{PRIQ:cor:tower-defect}
For an interaction of order $m$ assigned at every level,
\begin{equation}\label{PRIQ:eq:tower-defect}
  \Delta_m^{L0}
  =
  \sum_{\ell=0}^{L-1}
  C_{\ell0}\,
  \Delta_m^{\ell+1,\ell}
  C_{L,\ell+1}^{\times m},
\end{equation}
where $C_{00}=I$ and $C_{LL}=I$.  Consequently,
\begin{equation}\label{PRIQ:eq:tower-defect-bound}
  \PRIQnorm{\Delta_m^{L0}}
  \leq
  \sum_{\ell=0}^{L-1}
  \PRIQnorm{C_{\ell0}}
  \PRIQnorm{\Delta_m^{\ell+1,\ell}}
  \PRIQnorm{C_{L,\ell+1}}^m.
\end{equation}
\end{PRIQcorollary}

\begin{proof}
Induct on the number of levels, using \eqref{PRIQ:eq:multilevel-defect} at the final composition.
\end{proof}

The identity shows where a terminal mismatch entered the refinement tower.  Rather than storing only the final error, one can keep the local defect at each level and combine those records when needed.

\section{Approximate Bayesian refinement}\label{PRIQ:sec:approx-bayes}

\subsection{A normalization perturbation bound}

\begin{PRIQlemma}[Normalization perturbation]\label{PRIQ:lem:normalization-perturbation}
Let $u,v\in\PRIQR_+^q$ have positive total mass.  Then
\begin{equation}\label{PRIQ:eq:normalization-perturbation}
  \boxed{
  \PRIQnorm{
  \frac{u}{\PRIQone^Tu}
  -
  \frac{v}{\PRIQone^Tv}
  }_1
  \leq
  \frac{2\PRIQnorm{u-v}_1}
  {\max\{\PRIQone^Tu,\PRIQone^Tv\}}.}
\end{equation}
\end{PRIQlemma}

\begin{proof}
Write $a=\PRIQone^Tu$ and $b=\PRIQone^Tv$.  Then
\begin{align*}
\PRIQnorm{\frac{u}{a}-\frac{v}{b}}_1
&\leq
\frac{\PRIQnorm{u-v}_1}{a}
+
\PRIQnorm{v}_1\PRIQabs{\frac1a-\frac1b}\\
&=
\frac{\PRIQnorm{u-v}_1}{a}
+
\frac{\PRIQabs{a-b}}{a}\\
&\leq
\frac{2\PRIQnorm{u-v}_1}{a}.
\end{align*}
Interchanging $u$ and $v$ gives the same inequality with $b$ in the denominator.  Taking the smaller of the two upper bounds yields the denominator $\max\{a,b\}$.
\end{proof}

The factor of two need not be sharp.  It is uniform, and it makes the role of the evidence mass explicit.  No uniform estimate of this form remains useful when both masses approach zero.

\subsection{Defect in the observation intertwining law}
Let $L_{\PRIQfine}$ and $L_{\PRIQcoarse}$ be positive linear observation operators, and choose a scale $s>0$.  On a finite simplex, we write $\PRIQNrm(L)$ for normalization with respect to the total mass $\PRIQone^T$.  Define
\begin{equation}\label{PRIQ:eq:observation-linear-defect}
  R_L
  =
  CL_{\PRIQfine}-sL_{\PRIQcoarse}C.
\end{equation}

\begin{PRIQtheorem}[Approximate normalized transfer naturality]\label{PRIQ:thm:approx-normalized-transfer}
On finite simplices with block sum $C$, let
\[
  u=CL_{\PRIQfine}\mu,
  \qquad
  v=sL_{\PRIQcoarse}C\mu.
\]
If $u$ and $v$ are nonnegative with positive masses, then
\begin{equation}\label{PRIQ:eq:approx-normalized-transfer}
  \boxed{
  \PRIQnorm{
  C\PRIQNrm(L_{\PRIQfine})(\mu)
  -
  \PRIQNrm(L_{\PRIQcoarse})(C\mu)
  }_1
  \leq
  \frac{2\PRIQnorm{R_L\mu}_1}
  {\max\{\PRIQone^Tu,\PRIQone^Tv\}}.}
\end{equation}
\end{PRIQtheorem}

\begin{proof}
Because $C$ preserves total mass,
\[
  C\PRIQNrm(L_{\PRIQfine})(\mu)
  =\frac{u}{\PRIQone^Tu}.
\]
Multiplication of $L_{\PRIQcoarse}$ by $s$ does not change its normalized output, so
\[
  \PRIQNrm(L_{\PRIQcoarse})(C\mu)
  =\frac{v}{\PRIQone^Tv}.
\]
Now $u-v=R_L\mu$, and Lemma~\ref{PRIQ:lem:normalization-perturbation} applies.
\end{proof}

\subsection{Likelihood variation inside a fiber}
For diagonal likelihood operators, write
\begin{equation}\label{PRIQ:eq:likelihood-residual}
  \lambda^{\PRIQfine}
  =
  sC^T\lambda^{\PRIQcoarse}+r_{\lambda}.
\end{equation}
Then
\[
  R_L\mu=C(r_{\lambda}\odot\mu).
\]
Since block sum aggregation is an $\ell^1$ contraction on signed vectors,
\begin{equation}\label{PRIQ:eq:likelihood-residual-l1}
  \PRIQnorm{R_L\mu}_1
  \leq
  \sum_i\mu_i\PRIQabs{r_{\lambda,i}}
  \leq
  \PRIQnorm{r_{\lambda}}_{\infty}.
\end{equation}

\begin{PRIQcorollary}[Bayesian refinement defect]\label{PRIQ:cor:bayes-refinement-defect}
Let
\[
  z_{\PRIQfine}=Z_{\lambda^{\PRIQfine}}(\mu),
  \qquad
  z_{\PRIQcoarse}=sZ_{\lambda^{\PRIQcoarse}}(C\mu).
\]
Then
\begin{equation}\label{PRIQ:eq:bayes-refinement-defect}
  \boxed{
  \PRIQnorm{
  C\PRIQBup_{\lambda^{\PRIQfine}}(\mu)
  -
  \PRIQBup_{\lambda^{\PRIQcoarse}}(C\mu)
  }_1
  \leq
  \frac{2\sum_i\mu_i\PRIQabs{r_{\lambda,i}}}
  {\max\{z_{\PRIQfine},z_{\PRIQcoarse}\}}.}
\end{equation}
If both evidences are at least $z_*>0$, then
\begin{equation}\label{PRIQ:eq:bayes-refinement-uniform}
  d_{\PRIQBup}
  \leq
  \frac{2\PRIQnorm{r_{\lambda}}_{\infty}}{z_*}.
\end{equation}
\end{PRIQcorollary}

\begin{proof}
Combine \Cref{PRIQ:thm:approx-normalized-transfer} with \eqref{PRIQ:eq:likelihood-residual-l1}.
\end{proof}

\begin{PRIQdefinition}[Fiber oscillation]\label{PRIQ:def:fiber-oscillation}
For a fine likelihood, define
\[
  \PRIQosc_a(\lambda^{\PRIQfine})
  =
  \max_{i:\pi(i)=a}\lambda_i^{\PRIQfine}
  -
  \min_{i:\pi(i)=a}\lambda_i^{\PRIQfine},
\]
and
\[
  \PRIQosc_C(\lambda^{\PRIQfine})
  =\max_a\PRIQosc_a(\lambda^{\PRIQfine}).
\]
\end{PRIQdefinition}

When $q\geq2$ and commutation is required for every interior prior, exact Bayesian refinement is equivalent to $\PRIQosc_C(\lambda^{\PRIQfine})=0$ once the coarse likelihood and global scale have been chosen consistently.  In an approximate model, $\lambda_a^{\PRIQcoarse}$ may be the midpoint, the ordinary mean, or a mean chosen to reflect the task on fiber $a$.  The midpoint minimizes the largest absolute residual in that fiber and gives
\[
  \PRIQnorm{r_{\lambda}}_{\infty}
  =\frac12\PRIQosc_C(\lambda^{\PRIQfine})
\]
when $s=1$.

\subsection{Lipschitz behavior of the Bayesian update}

\begin{PRIQproposition}[Lipschitz bound under an evidence floor]\label{PRIQ:prop:bayes-lipschitz}
Let $\lambda\in\PRIQR_{>0}^q$.  If
\[
  Z_{\lambda}(\mu),Z_{\lambda}(\nu)\geq z_*>0,
\]
then
\begin{equation}\label{PRIQ:eq:bayes-lipschitz}
  \boxed{
  \PRIQnorm{\PRIQBup_{\lambda}(\mu)-\PRIQBup_{\lambda}(\nu)}_1
  \leq
  \frac{2\PRIQnorm{\lambda}_{\infty}}{z_*}
  \PRIQnorm{\mu-\nu}_1.}
\end{equation}
\end{PRIQproposition}

\begin{proof}
Apply Lemma~\ref{PRIQ:lem:normalization-perturbation} to
\[
  u=\lambda\odot\mu,
  \qquad
  v=\lambda\odot\nu.
\]
Then
\[
  \PRIQnorm{u-v}_1
  \leq\PRIQnorm{\lambda}_{\infty}\PRIQnorm{\mu-\nu}_1.
\]
\end{proof}

\begin{PRIQremark}[Scale invariance of the bound]
Replacing $\lambda$ by $s\lambda$ multiplies both $\PRIQnorm{\lambda}_{\infty}$ and $z_*$ by $s$, leaving the ratio unchanged.
\end{PRIQremark}

\begin{PRIQremark}[Local versus global constants]
The constant in \eqref{PRIQ:eq:bayes-lipschitz} is valid over the whole evidence region under consideration.  Along a particular trajectory, the Jacobian in \eqref{PRIQ:eq:bayes-differential} may give a much smaller local factor.  The recursive theorem below can use either bound, provided the chosen constant is certified on the states being considered.
\end{PRIQremark}


\chapter{Posterior certificates and hyperedge compression}\label{PRIQ:book:chapter-43}
\section{Certified approximate coarse filtering}\label{PRIQ:sec:filter-cert}

\subsection{Prediction and update errors}
Let the fine recursion be
\begin{equation}\label{PRIQ:eq:fine-filter-recursion}
  \widehat\mu_k^{\PRIQfine}
  =\Phi_k^{\PRIQfine}(\mu_{k-1}^{\PRIQfine}),
  \qquad
  \mu_k^{\PRIQfine}
  =\PRIQBup_k^{\PRIQfine}(\widehat\mu_k^{\PRIQfine}),
\end{equation}
and the coarse recursion be
\begin{equation}\label{PRIQ:eq:coarse-filter-recursion}
  \widehat\mu_k^{\PRIQcoarse}
  =\Phi_k^{\PRIQcoarse}(\mu_{k-1}^{\PRIQcoarse}),
  \qquad
  \mu_k^{\PRIQcoarse}
  =\PRIQBup_k^{\PRIQcoarse}(\widehat\mu_k^{\PRIQcoarse}).
\end{equation}
Define the discrepancies before and after each update by
\begin{equation}\label{PRIQ:eq:pre-post-errors}
  e_k^-
  =\PRIQnorm{C\widehat\mu_k^{\PRIQfine}-\widehat\mu_k^{\PRIQcoarse}}_1,
  \qquad
  e_k^+
  =\PRIQnorm{C\mu_k^{\PRIQfine}-\mu_k^{\PRIQcoarse}}_1.
\end{equation}
At the initial time, set
\[
  e_0^+=\PRIQnorm{C\mu_0^{\PRIQfine}-\mu_0^{\PRIQcoarse}}_1.
\]

Assume the coarse prediction is Lipschitz with constant $L_{\Phi,k}$ and that its quotient defect is bounded by
\begin{equation}\label{PRIQ:eq:prediction-map-defect}
  \sup_{\mu}
  \PRIQnorm{C\Phi_k^{\PRIQfine}(\mu)-\Phi_k^{\PRIQcoarse}(C\mu)}_1
  \leq d_{\Phi,k}.
\end{equation}
Assume the coarse update is Lipschitz with constant $L_{\PRIQBup,k}$ and that its direct refinement defect is bounded by
\begin{equation}\label{PRIQ:eq:update-map-defect}
  \sup_{\mu}
  \PRIQnorm{C\PRIQBup_k^{\PRIQfine}(\mu)-\PRIQBup_k^{\PRIQcoarse}(C\mu)}_1
  \leq d_{\PRIQBup,k}.
\end{equation}
The suprema in \eqref{PRIQ:eq:prediction-map-defect} and \eqref{PRIQ:eq:update-map-defect} are taken over a declared admissible set that contains both recursions.

\begin{PRIQlemma}[Prediction bound for one step]\label{PRIQ:lem:prediction-step}
Under \eqref{PRIQ:eq:prediction-map-defect},
\begin{equation}\label{PRIQ:eq:prediction-step}
  e_k^-
  \leq
  L_{\Phi,k}e_{k-1}^+
  +d_{\Phi,k}.
\end{equation}
\end{PRIQlemma}

\begin{proof}
Insert and subtract $\Phi_k^{\PRIQcoarse}(C\mu_{k-1}^{\PRIQfine})$.  The quotient defect bounds one difference, while the Lipschitz estimate bounds the other by $L_{\Phi,k}e_{k-1}^+$.
\end{proof}

\begin{PRIQlemma}[Update bound for one step]\label{PRIQ:lem:update-step}
Under \eqref{PRIQ:eq:update-map-defect},
\begin{equation}\label{PRIQ:eq:update-step}
  e_k^+
  \leq
  L_{\PRIQBup,k}e_k^-
  +d_{\PRIQBup,k}.
\end{equation}
\end{PRIQlemma}

\begin{proof}
Insert and subtract $\PRIQBup_k^{\PRIQcoarse}(C\widehat\mu_k^{\PRIQfine})$.  The direct refinement defect bounds one difference, while the Lipschitz estimate bounds the other by $L_{\PRIQBup,k}e_k^-$.
\end{proof}

Combining the two inequalities gives
\begin{equation}\label{PRIQ:eq:combined-step}
  e_k^+
  \leq
  \alpha_ke_{k-1}^++\beta_k,
\end{equation}
where
\begin{equation}\label{PRIQ:eq:alpha-beta}
  \alpha_k=L_{\PRIQBup,k}L_{\Phi,k},
  \qquad
  \beta_k=L_{\PRIQBup,k}d_{\Phi,k}+d_{\PRIQBup,k}.
\end{equation}

\subsection{Closed posterior certificate}

\begin{PRIQtheorem}[Certified approximate filtering]\label{PRIQ:thm:posterior-certificate}
Under the assumptions above,
\begin{equation}\label{PRIQ:eq:posterior-certificate}
  \boxed{
  e_N^+
  \leq
  \left(\prod_{j=1}^N\alpha_j\right)e_0^+
  +
  \sum_{k=1}^N
  \left(\prod_{j=k+1}^N\alpha_j\right)
  \beta_k.}
\end{equation}
Here and below, an empty product equals one.  With exact initial aggregation $e_0^+=0$,
\begin{equation}\label{PRIQ:eq:posterior-certificate-exact-initial}
  e_N^+
  \leq
  \sum_{k=1}^N
  \left(\prod_{j=k+1}^N\alpha_j\right)
  \bigl(L_{\PRIQBup,k}d_{\Phi,k}+d_{\PRIQBup,k}\bigr).
\end{equation}
\end{PRIQtheorem}

\begin{proof}
Repeated substitution in the affine recurrence \eqref{PRIQ:eq:combined-step} gives the stated formula.  The initial error is multiplied by every later amplification factor.  A defect introduced at step $k$ is multiplied only by the factors from steps $k+1$ through $N$.
\end{proof}

\begin{PRIQcorollary}[Uniform stability regimes]\label{PRIQ:cor:filter-stability-regimes}
Suppose $0\leq\alpha_k\leq\alpha$ and $0\leq\beta_k\leq\beta$.

\begin{enumerate}[label=(\roman*),leftmargin=2.8em]
\item If $0\leq\alpha<1$, then
\begin{equation}\label{PRIQ:eq:stable-filter-bound}
  e_N^+
  \leq
  \alpha^Ne_0^+
  +
  \beta\frac{1-\alpha^N}{1-\alpha},
  \qquad
  \limsup_{N\to\infty}e_N^+
  \leq\frac{\beta}{1-\alpha}.
\end{equation}

\item If $\alpha=1$, then
\[
  e_N^+\leq e_0^++N\beta.
\]

\item If $\alpha>1$, then
\[
  e_N^+
  \leq
  \alpha^Ne_0^+
  +
  \beta\frac{\alpha^N-1}{\alpha-1}.
\]
\end{enumerate}
\end{PRIQcorollary}

\begin{proof}
Bound the products in \eqref{PRIQ:eq:posterior-certificate} by powers of $\alpha$ and sum the geometric series.
\end{proof}

This trichotomy describes the certificate, not necessarily the observed filter.  A stable trajectory can still come with a bound for which $\alpha>1$.  In that case, the chosen global norm estimates simply fail to certify contraction.

\subsection{From interaction coefficients to posterior error}
For a flow interval of length $T_k$, suppose the fine trajectory remains in a ball of radius $R_k$ and the coarse field has Lipschitz constant $L_{F,k}$.  By Corollary~\ref{PRIQ:cor:coefficient-flow-certificate}, one may take
\begin{equation}\label{PRIQ:eq:dphi-coefficient}
  d_{\Phi,k}
  =
  \mathfrak d_{F,k}(R_k)
  \frac{e^{L_{F,k}T_k}-1}{L_{F,k}},
\end{equation}
with the continuous value $T_k\mathfrak d_{F,k}(R_k)$ at $L_{F,k}=0$.
For the observation, Corollary~\ref{PRIQ:cor:bayes-refinement-defect} gives
\begin{equation}\label{PRIQ:eq:dbayes-coefficient}
  d_{\PRIQBup,k}
  \leq
  \frac{2\PRIQnorm{r_{\lambda,k}}_{\infty}}{z_{*,k}}.
\end{equation}
Substituting these bounds into \eqref{PRIQ:eq:posterior-certificate-exact-initial} gives a certificate expressed directly in terms of interaction tensor defects and likelihood residuals within each fiber.

\subsection{Observable and total variation consequences}

\begin{PRIQcorollary}[Coarse observable error]\label{PRIQ:cor:observable-error}
For any coarse observable $g\in\PRIQR^q$,
\[
  \PRIQabs{g^TC\mu_N^{\PRIQfine}-g^T\mu_N^{\PRIQcoarse}}
  \leq
  \PRIQnorm{g}_{\infty}e_N^+.
\]
In particular,
\[
  \PRIQTV(C\mu_N^{\PRIQfine},\mu_N^{\PRIQcoarse})
  =\frac12e_N^+.
\]
\end{PRIQcorollary}

\begin{proof}
Apply the duality between $\ell^1$ and $\ell^\infty$.  The total variation identity is standard for finite probability vectors.
\end{proof}

\subsection{Sequential evidence error}
For brevity, write $\lambda_k^{\PRIQfine}=\lambda_{y_k}^{\PRIQfine}$ and $\lambda_k^{\PRIQcoarse}=\lambda_{y_k}^{\PRIQcoarse}$, and decompose
\begin{equation}\label{PRIQ:eq:step-likelihood-residual}
  \lambda_k^{\PRIQfine}
  =s_kC^T\lambda_k^{\PRIQcoarse}+r_k.
\end{equation}
This gives
\begin{align}
  z_k^{\PRIQfine}-s_kz_k^{\PRIQcoarse}
  &=
  s_k(\lambda_k^{\PRIQcoarse})^T
  (C\widehat\mu_k^{\PRIQfine}-\widehat\mu_k^{\PRIQcoarse})
  +r_k^T\widehat\mu_k^{\PRIQfine}.
  \label{PRIQ:eq:evidence-error-identity}
\end{align}

\begin{PRIQproposition}[Evidence certificate for one step]\label{PRIQ:prop:evidence-error}
Assume $\widehat\mu_k^{\PRIQfine}$ is a probability vector.  Under \eqref{PRIQ:eq:step-likelihood-residual},
\begin{equation}\label{PRIQ:eq:evidence-error-bound}
  \boxed{
  \PRIQabs{z_k^{\PRIQfine}-s_kz_k^{\PRIQcoarse}}
  \leq
  s_k\PRIQnorm{\lambda_k^{\PRIQcoarse}}_{\infty}e_k^-
  +\PRIQnorm{r_k}_{\infty}.}
\end{equation}
If $z_k^{\PRIQfine},s_kz_k^{\PRIQcoarse}\geq z_{*,k}>0$, then
\begin{equation}\label{PRIQ:eq:log-evidence-step}
  \boxed{
  \PRIQabs{
  \log z_k^{\PRIQfine}
  -\log(s_kz_k^{\PRIQcoarse})}
  \leq
  \frac{
  s_k\PRIQnorm{\lambda_k^{\PRIQcoarse}}_{\infty}e_k^-
  +\PRIQnorm{r_k}_{\infty}
  }{z_{*,k}}.}
\end{equation}
\end{PRIQproposition}

\begin{proof}
Apply Hölder's inequality to \eqref{PRIQ:eq:evidence-error-identity}.  Since $\widehat\mu_k^{\PRIQfine}$ is a probability vector, $\PRIQnorm{\widehat\mu_k^{\PRIQfine}}_1=1$.  For the logarithms, the mean value theorem gives
\[
  \PRIQabs{\log a-\log b}
  \leq\frac{\PRIQabs{a-b}}{\min\{a,b\}}.
\]
\end{proof}

\begin{PRIQcorollary}[Certificate for cumulative log evidence]\label{PRIQ:cor:cumulative-evidence}
Let $Z_{1:N}^{\PRIQfine}=\prod_kz_k^{\PRIQfine}$ and $Z_{1:N}^{\PRIQcoarse}=\prod_kz_k^{\PRIQcoarse}$.  Then
\begin{equation}\label{PRIQ:eq:cumulative-evidence}
  \boxed{
  \PRIQabs{
  \log Z_{1:N}^{\PRIQfine}
  -\sum_{k=1}^N\log s_k
  -\log Z_{1:N}^{\PRIQcoarse}}
  \leq
  \sum_{k=1}^N
  \frac{
  s_k\PRIQnorm{\lambda_k^{\PRIQcoarse}}_{\infty}e_k^-
  +\PRIQnorm{r_k}_{\infty}
  }{z_{*,k}}.}
\end{equation}
\end{PRIQcorollary}

\begin{proof}
Sum \eqref{PRIQ:eq:log-evidence-step} over the observation sequence.
\end{proof}

\section{Hyperedge pruning and certified compression}\label{PRIQ:sec:pruning}

Exact quotienting and sparse representation are distinct.  A system may have an exact coarse quotient while its fine interaction tensor is dense.  Conversely, a sparse fine tensor can fail to descend.  This section introduces a separate residual for deleted hyperedges and then combines it with the quotient defect.

\subsection{Hilbert Schmidt tensor blocks}
Let $X$ and $Y$ be finite dimensional Hilbert spaces with orthogonal packet projectors $(P_j)$ and $(Q_i)$.  Identify a multilinear Hilbert Schmidt tensor of order $m$ with the Hilbert Schmidt linear map
\[
  \widehat A:X^{\otimes m}\to Y.
\]
Its packet block is
\[
  \widehat A_{i\mid j_1,\ldots,j_m}
  =
  Q_i\widehat A(P_{j_1}\otimes\cdots\otimes P_{j_m}).
\]

\begin{PRIQtheorem}[Orthogonal hyperedge Parseval identity]\label{PRIQ:thm:hyperedge-parseval}
Distinct packet blocks are orthogonal in Hilbert Schmidt inner product, and
\begin{equation}\label{PRIQ:eq:hyperedge-parseval}
  \boxed{
  \PRIQnorm{\widehat A}_{\mathrm{HS}}^2
  =
  \sum_{i,j_1,\ldots,j_m}
  \PRIQnorm{\widehat A_{i\mid j_1,\ldots,j_m}}_{\mathrm{HS}}^2.}
\end{equation}
\end{PRIQtheorem}

\begin{proof}
The tensor product input projectors
\[
  P_{j_1}\otimes\cdots\otimes P_{j_m}
\]
are mutually orthogonal, as are the output projectors $Q_i$.  For two distinct block labels, orthogonality on at least one side forces the Hilbert Schmidt inner product to vanish.  Exact reconstruction is therefore an orthogonal sum, and Pythagoras gives \eqref{PRIQ:eq:hyperedge-parseval}.
\end{proof}

\subsection{Retained field and named residual}
Choose a retained hyperedge set $\PRIQcE_m^{\mathrm{keep}}$.  Define
\begin{equation}\label{PRIQ:eq:pruned-tensor}
  \widehat A_m^{\mathrm{keep}}
  =
  \sum_{h\in\PRIQcE_m^{\mathrm{keep}}}
  \widehat A_{m,h},
  \qquad
  R_m
  =
  \widehat A_m-\widehat A_m^{\mathrm{keep}}.
\end{equation}

\begin{PRIQcorollary}[Exact pruning residual]\label{PRIQ:cor:pruning-residual}
Under the hypotheses of \Cref{PRIQ:thm:hyperedge-parseval},
\begin{equation}\label{PRIQ:eq:pruning-residual}
  \boxed{
  \PRIQnorm{R_m}_{\mathrm{HS}}^2
  =
  \sum_{h\notin\PRIQcE_m^{\mathrm{keep}}}
  \PRIQnorm{\widehat A_{m,h}}_{\mathrm{HS}}^2.}
\end{equation}
\end{PRIQcorollary}

\begin{proof}
The retained and omitted block families are orthogonal by \Cref{PRIQ:thm:hyperedge-parseval}.  Since $R_m$ is exactly the sum of the omitted blocks, Pythagoras gives \eqref{PRIQ:eq:pruning-residual}.
\end{proof}

The residual is a tensor identity, not merely an empirical estimate.  Its induced state error is obtained by applying the relevant multilinear norm inequality.

\subsection{\texorpdfstring{Retaining the $k$ largest energy hyperedges}{Retaining the k largest energy hyperedges}}
Let
\[
  e_h=\PRIQnorm{\widehat A_{m,h}}_{\mathrm{HS}}^2
\]
and order the energies
\[
  e_{h_1}\geq e_{h_2}\geq\cdots\geq e_{h_M}.
\]

\begin{PRIQtheorem}[Optimal retention of $k$ orthogonal hyperedges]\label{PRIQ:thm:best-k-hyperedges}
Among all retained sets of cardinality $k$, the set
\[
  \{h_1,\ldots,h_k\}
\]
minimizes the Hilbert Schmidt residual.  The minimum is
\begin{equation}\label{PRIQ:eq:best-k-residual}
  \boxed{
  \min_{|\PRIQcE|=k}
  \PRIQnorm{R_m(\PRIQcE)}_{\mathrm{HS}}^2
  =
  \sum_{r>k}e_{h_r}.}
\end{equation}
\end{PRIQtheorem}

\begin{proof}
By \eqref{PRIQ:eq:pruning-residual}, minimizing the omitted energy is equivalent to maximizing the retained energy.  For a fixed cardinality, this is achieved by selecting the $k$ largest nonnegative terms.
\end{proof}

\begin{PRIQboundary}[Overlapping blocks and structured costs]
The theorem above uses orthogonal block energy and assigns the same unit cost to every retained hyperedge.  With overlapping dictionaries, shared low rank factors, communication costs, or group constraints that preserve structure, ranking individual energies need not be optimal.
\end{PRIQboundary}

\subsection{Hyperedge entropy}
For nonzero total interaction energy, define
\[
  p_h=\frac{e_h}{\sum_{\eta}e_{\eta}},
  \qquad
  H_m=-\sum_hp_h\log p_h,
\]
and the effective hyperedge count
\[
  N_{\mathrm{eff},m}=e^{H_m}.
\]
Low entropy means that the interaction energy is concentrated in a small set of resolved hyperedges.  High entropy means that pruning by a fixed cutoff is less promising in the chosen resolution.  The quantity deliberately depends on the resolution.  It evaluates a representation rather than an invariant of the unresolved tensor alone.

\subsection{Combined quotient and pruning defect}
Let the full fine coefficient split as
\[
  A_m^{\PRIQfine}
  =A_m^{\mathrm{keep}}+R_m^{\PRIQfine}.
\]
For a proposed coarse coefficient $A_m^{\PRIQcoarse}$, define
\[
  \Delta_m^{\mathrm{total}}
  =CA_m^{\PRIQfine}-A_m^{\PRIQcoarse}C^{\times m}.
\]
Then
\begin{align}
  \Delta_m^{\mathrm{total}}
  &=
  \underbrace{CA_m^{\mathrm{keep}}-A_m^{\PRIQcoarse}C^{\times m}}
  _{\Delta_m^{\mathrm{keep}}}
  +CR_m^{\PRIQfine}.
  \label{PRIQ:eq:combined-pruning-quotient}
\end{align}
Therefore
\begin{equation}\label{PRIQ:eq:combined-pruning-bound}
  \boxed{
  \PRIQnorm{\Delta_m^{\mathrm{total}}}
  \leq
  \PRIQnorm{\Delta_m^{\mathrm{keep}}}
  +\PRIQnorm{C}\PRIQnorm{R_m^{\PRIQfine}}.}
\end{equation}

For analytic fields on a ball of radius $R$,
\begin{equation}\label{PRIQ:eq:combined-field-majorant}
  \mathfrak d_F(R)
  \leq
  \sum_m
  \PRIQnorm{\Delta_m^{\mathrm{keep}}}R^m
  +
  \PRIQnorm{C}\sum_m\PRIQnorm{R_m^{\PRIQfine}}R^m.
\end{equation}
The flow and posterior bounds use this combined quantity.  A pruning decision is justified only when the omitted part remains explicit and is carried through the later estimates like any other model error.

\subsection{A budgeted design problem}
These terms lead to a practical design problem.  The objective below records the tradeoff without claiming that it can always be solved efficiently:
\begin{equation}\label{PRIQ:eq:budgeted-design}
  \min_{C,\{A_m^{\PRIQcoarse}\},\{\PRIQcE_m^{\mathrm{keep}}\}}
  \left[
  \sum_mw_m\PRIQnorm{\Delta_m^{\mathrm{total}}}
  +\lambda\sum_m\operatorname{Cost}(\PRIQcE_m^{\mathrm{keep}})
  +\gamma\sum_y d_{\PRIQBup,y}
  \right].
\end{equation}
The minimization is subject to any required positivity, mass, rank, and interpretability conditions on $C$.  The three terms measure failure of dynamical closure, computational cost, and loss of observational information.  Their weights and the admissible packet families must be chosen for the application.


\chapter{Algorithms and worked nonlinear filters}\label{PRIQ:book:chapter-44}
\section{Algorithms and finite certificates}\label{PRIQ:sec:algorithms}


% BEGIN retained input sources/research/hypergraph-bayesian-refinement/tables/verification_macros.tex
% Generated by examples/run_certificates.py
\newcommand{\VerificationPassCount}{35}
\newcommand{\VerificationFailCount}{0}
\newcommand{\ExactPredictionResidual}{\ensuremath{7.22\times 10^{-16}}}
\newcommand{\ExactPosteriorResidual}{\ensuremath{3.89\times 10^{-16}}}
\newcommand{\ExactEvidenceResidual}{\ensuremath{6.66\times 10^{-16}}}
\newcommand{\PerturbationSlope}{\ensuremath{0.9999}}
\newcommand{\HyperedgeEntropyValue}{\ensuremath{1.9360}}
\newcommand{\EffectiveHyperedgeCount}{\ensuremath{6.931}}



Every finite criterion in the paper can be checked directly from coordinates.  The reference code uses dense arrays so that each contraction can be inspected.  A sparse implementation would loop only over the active hyperedges.

\subsection{Exact test for every interaction order}
For a tensor of order $m$ in scalar packet coordinates, compute the aggregated coefficient array
\[
  g^{(m)}_{a;j_1,\ldots,j_m}
  =
  \sum_{i:\pi(i)=a}
  a^{\PRIQfine,(m)}_{i;j_1,\ldots,j_m}.
\]
For every coarse input tuple $\mathbf b=(b_1,\ldots,b_m)$, define the fine fiber product
\[
  \PRIQcG_{\mathbf b}
  =
  \{(j_1,\ldots,j_m):\pi(j_r)=b_r\text{ for all }r\}.
\]

\begin{PRIQalgorithm}[Exact packet coordinate lumpability test]\label{PRIQ:alg:exact-lumpability}
For every coarse output $a$ and coarse input tuple $\mathbf b$:
\begin{enumerate}[label=\arabic*.,leftmargin=2.3em]
\item Evaluate $g^{(m)}_{a;\mathbf j}$ for all $\mathbf j\in\PRIQcG_{\mathbf b}$.
\item In exact arithmetic, test whether all values are equal.  With floating point input, compare their spread with a declared tolerance and record that tolerance.
\item When the exact test passes, assign the common value to $a^{\PRIQcoarse,(m)}_{a;\mathbf b}$.  When only the numerical tolerance is met, assign a declared representative such as the mean or the midrange.
\item Otherwise, return a witness $(a,\mathbf j,\mathbf j')$ with unequal aggregate coefficients.
\end{enumerate}
\end{PRIQalgorithm}

A dense implementation requires $O(n^{m+1})$ arithmetic to aggregate a tensor of order $m$ on $n$ states.  If the fine tensor is stored as $E_m$ active scalar hyperedges, the aggregation pass is $O(E_m)$, followed by grouping and consistency checks over the represented coarse incidences.

\begin{PRIQproposition}[Correctness in exact arithmetic]\label{PRIQ:prop:algorithm-correctness}
With exact arithmetic, Algorithm~\ref{PRIQ:alg:exact-lumpability} accepts exactly the tensors that descend through the given block sum aggregation.  On acceptance, it returns the unique coarse tensor.  A positive numerical tolerance produces an approximate certificate rather than an exact algebraic decision.
\end{PRIQproposition}

\begin{proof}
The claim follows directly from \Cref{PRIQ:thm:packet-coordinate-criterion}.
\end{proof}

\subsection{Approximate coarse coefficients}
When exact fiber constancy fails, the approximation depends on a chosen loss and, when relevant, a task distribution.  Under scalar squared loss with nonnegative weights $w_{\mathbf j}$ on a fiber product, assume their total weight is positive.  The least squares representative is
\begin{equation}\label{PRIQ:eq:least-squares-coarse-coeff}
  a^{\PRIQcoarse,\mathrm{LS}}_{a;\mathbf b}
  =
  \frac{
  \sum_{\mathbf j\in\PRIQcG_{\mathbf b}}
  w_{\mathbf j}g_{a;\mathbf j}
  }{
  \sum_{\mathbf j\in\PRIQcG_{\mathbf b}}w_{\mathbf j}
  }.
\end{equation}
For scalar worst case loss, a minimax representative is the midrange
\begin{equation}\label{PRIQ:eq:minimax-coarse-coeff}
  a^{\PRIQcoarse,\infty}_{a;\mathbf b}
  =
  \frac12
  \left(
  \max_{\mathbf j\in\PRIQcG_{\mathbf b}}g_{a;\mathbf j}
  +
  \min_{\mathbf j\in\PRIQcG_{\mathbf b}}g_{a;\mathbf j}
  \right),
\end{equation}
with optimal scalar fiber error
\begin{equation}\label{PRIQ:eq:minimax-fiber-error}
  \varepsilon^{\infty}_{a;\mathbf b}
  =
  \frac12
  \left(
  \max_{\mathbf j\in\PRIQcG_{\mathbf b}}g_{a;\mathbf j}
  -
  \min_{\mathbf j\in\PRIQcG_{\mathbf b}}g_{a;\mathbf j}
  \right).
\end{equation}
These formulas are local coefficient summaries.  A global operator norm can couple several fibers and may require a joint optimization.

\subsection{Certificate pipeline}
The reference implementation follows the dependency chain
\[
\begin{array}{c}
\text{interaction tensors}\\
\downarrow\\
\text{fiber constancy and kernel tests}\\
\downarrow\\
\text{coarse tensor and }\Delta_m\\
\downarrow\\
\mathfrak d_F(R)\text{ and }d_{\Phi,k}\\
\downarrow\\
 r_{\lambda,k},z_{*,k},d_{\PRIQBup,k}\\
\downarrow\\
 e_k^-,e_k^+,\PRIQErr_N.
\end{array}
\]
Each numerical record stores the theorem identifier, the inputs, the residual, the tolerance, and the outcome.  The suite includes examples built to satisfy the exact identities and perturbations built to violate them.

\subsection{Deterministic generator for exact examples}
A coarse tensor can be lifted to a fine tensor by
\begin{equation}\label{PRIQ:eq:exact-instance-generator}
  A_m^{\PRIQfine}
  =
  EA_m^{\PRIQcoarse}C^{\times m}+K_m,
  \qquad
  CK_m=0.
\end{equation}
The term $K_m$ can contain complicated hidden interactions without changing the quotient.  This makes the construction useful for testing because the correct coarse tensor is known before any numerical calculation.

\begin{PRIQproposition}[Correctness of the exact example generator]\label{PRIQ:prop:exact-instance-generator}
Every tensor constructed by \eqref{PRIQ:eq:exact-instance-generator} satisfies
\[
  CA_m^{\PRIQfine}=A_m^{\PRIQcoarse}C^{\times m}.
\]
Conversely, every exact fine tensor admits such a representation for every right inverse $E$, with
\[
  K_m=A_m^{\PRIQfine}-EA_m^{\PRIQcoarse}C^{\times m}.
\]
\end{PRIQproposition}

\begin{proof}
The forward identity follows from $CE=I$ and $CK_m=0$.  For the converse, exact descent gives
\[
  CK_m
  =CA_m^{\PRIQfine}-CEA_m^{\PRIQcoarse}C^{\times m}=0.
\]
\end{proof}

\subsection{Reference implementation}
The accompanying source tree contains:
\begin{center}
\small
\begin{tabular}{@{}p{0.31\linewidth}p{0.61\linewidth}@{}}
\toprule
Path & Purpose\\
\midrule
\texttt{src/priq/core.py} & aggregation maps, tensor actions, exact lifts, quotient defects, packet coordinate recovery, Bayesian updates, and a deterministic Runge Kutta integrator\\
\texttt{tests/test\_priq.py} & unit tests for the theorems and negative controls\\
\texttt{examples/}\allowbreak\texttt{run\_certificates.py} & deterministic tensor, flow, filtering, perturbation, pruning, figure, and ledger generation\\
\texttt{results/verification.json} & certificate summary in JSON form\\
\texttt{tables/test\_results.csv} & flat test ledger\\
\texttt{figures/} & vector PDF figures and large PNG figures generated from the certificate run\\
\bottomrule
\end{tabular}
\end{center}

\begin{PRIQcertificate}[Executed finite suite]\label{PRIQ:cert:executed-suite}
The deterministic reference run completed \VerificationPassCount\ checks with \VerificationFailCount\ failures.  It covers quotient tensors of orders one through three, annihilation of hidden inputs, recovery in packet coordinates, negative controls, and the defect identity across several levels.  It also covers nonlinear prediction, Bayesian refinement, evidence aggregation, approximate flow and posterior bounds, perturbation scaling, and orthogonal hyperedge residuals.  The independent unit tests also pass in the build environment.
\end{PRIQcertificate}


% BEGIN retained input sources/research/hypergraph-bayesian-refinement/tables/test_results.tex
\begin{longtable}{@{}p{0.08\linewidth}p{0.43\linewidth}p{0.12\linewidth}p{0.12\linewidth}p{0.08\linewidth}@{}}
\caption{Executed deterministic finite certificate ledger.}\label{PRIQ:tab:executed-ledger}\\
\toprule
ID & Check & Value & Threshold & Status\\
\midrule
\endfirsthead
\toprule
ID & Check & Value & Threshold & Status\\
\midrule
\endhead
Q1D & order 1 quotient defect & $1.19\times 10^{-16}$ & $\leq 10^{-10}$ & PASS\\
Q1K & order 1 hidden input residual & $5.38\times 10^{-16}$ & $\leq 10^{-10}$ & PASS\\
Q1R & order 1 coarse tensor recovery & $8.41\times 10^{-17}$ & $\leq 10^{-10}$ & PASS\\
Q2D & order 2 quotient defect & $3.76\times 10^{-16}$ & $\leq 10^{-10}$ & PASS\\
Q2K & order 2 hidden input residual & $5.71\times 10^{-16}$ & $\leq 10^{-10}$ & PASS\\
Q2R & order 2 coarse tensor recovery & $5.55\times 10^{-17}$ & $\leq 10^{-10}$ & PASS\\
Q3D & order 3 quotient defect & $8.39\times 10^{-16}$ & $\leq 10^{-10}$ & PASS\\
Q3K & order 3 hidden input residual & $1.15\times 10^{-15}$ & $\leq 10^{-10}$ & PASS\\
Q3R & order 3 coarse tensor recovery & $1.11\times 10^{-16}$ & $\leq 10^{-10}$ & PASS\\
QNEG & a perturbation of one coefficient violates fiber constancy & $0.125$ & $\geq 10^{-1}$ & PASS\\
QSYM & symmetric interaction remains symmetric after exact lift & $0$ & $\leq 10^{-10}$ & PASS\\
QHID & projected hidden tensor has zero coarse output & $0$ & $\leq 10^{-10}$ & PASS\\
QTRN & exact quotient transitivity across two aggregation levels & $4.84\times 10^{-16}$ & $\leq 10^{-10}$ & PASS\\
QMLT & interaction defect identity across two levels & $9.46\times 10^{-16}$ & $\leq 10^{-10}$ & PASS\\
QFLD & fine vector field factors through the coarse field & $2.22\times 10^{-16}$ & $\leq 10^{-10}$ & PASS\\
QPRD & exact aggregate prediction commutation & $7.22\times 10^{-16}$ & $\leq 5\times 10^{-10}$ & PASS\\
QCOC & exact commutation of the prediction and update cocycle & $3.89\times 10^{-16}$ & $\leq 5\times 10^{-10}$ & PASS\\
QEVD & exact evidence aggregation at one step & $6.66\times 10^{-16}$ & $\leq 5\times 10^{-10}$ & PASS\\
QMAS & fine nonlinear trajectory preserves total mass & $7.77\times 10^{-16}$ & $\leq 5\times 10^{-10}$ & PASS\\
QPOS & fine nonlinear trajectory preserves nonnegativity & $0$ & $\leq 5\times 10^{-12}$ & PASS\\
QBAY & a likelihood constant within each fiber gives exact Bayesian refinement & $2.78\times 10^{-16}$ & $\leq 10^{-10}$ & PASS\\
QBEV & a likelihood constant within each fiber preserves evidence & $2.22\times 10^{-16}$ & $\leq 10^{-10}$ & PASS\\
QSCL & Bayesian update is invariant under global likelihood scale & $1.53\times 10^{-16}$ & $\leq 10^{-10}$ & PASS\\
QFDB & controlled dynamical perturbation respects the flow certificate & $0$ & $\leq 2\times 10^{-9}$ & PASS\\
QFSL & flow discrepancy has first order perturbation scaling & $3.37\times 10^{-3}$ & $\leq 0.08$ & PASS\\
QLPS & perturbed likelihood remains strictly positive & $0.8$ & $\geq 10^{-12}$ & PASS\\
QOBB & variation of the likelihood within a fiber respects the Bayesian bound & $0$ & $\leq 10^{-12}$ & PASS\\
QBSC & the Bayesian perturbation sweep remains below the likelihood certificate & $0$ & $\leq 10^{-12}$ & PASS\\
QBSL & Bayesian discrepancy has first order scaling for small fiber mismatch & $4.63\times 10^{-5}$ & $\leq 0.08$ & PASS\\
QNMB & symmetric normalization bound uses the larger positive mass & $0$ & $\leq 10^{-12}$ & PASS\\
QPST & recursive posterior certificate dominates every joint perturbation run & $0$ & $\leq 10^{-10}$ & PASS\\
QPSL & joint posterior discrepancy has first order scaling for small perturbations & $1.23\times 10^{-4}$ & $\leq 0.12$ & PASS\\
QLEV & the evidence certificate at each step dominates the accumulated discrepancy & $0$ & $\leq 10^{-10}$ & PASS\\
QPAR & hyperedge energy satisfies the scalar Parseval identity & $8.88\times 10^{-16}$ & $\leq 10^{-10}$ & PASS\\
QBST & the residual from retaining the largest energy hyperedges decreases to zero & $0$ & $\leq 10^{-10}$ & PASS\\
\bottomrule
\end{longtable}



\section{Worked quotient and filtering models}\label{PRIQ:sec:examples}

\subsection{A nonlinear fine system with four states and a quotient with two states}
Let
\[
  x=(x_1,x_2,x_3,x_4)^T\in\Delta_4,
  \qquad
  C=
  \begin{bmatrix}
  1&1&0&0\\
  0&0&1&1
  \end{bmatrix},
\]
so
\[
  z=Cx=(z_1,z_2)^T,
  \qquad
  z_1=x_1+x_2,
  \quad
  z_2=x_3+x_4.
\]
The coarse dynamics are
\begin{equation}\label{PRIQ:eq:coarse-two-state}
  \dot z_1=-az_1+bz_2,
  \qquad
  \dot z_2=az_1-bz_2,
  \qquad a,b>0.
\end{equation}
Define nonlinear internal exchange channels
\begin{align}
  h_{12}(x)
  &=
  (r_0+r_1x_3)x_2-(s_0+s_1x_4)x_1,
  \label{PRIQ:eq:h12}\\
  h_{34}(x)
  &=
  (u_0+u_1x_1)x_4-(v_0+v_1x_2)x_3,
  \label{PRIQ:eq:h34}
\end{align}
with all coefficients nonnegative.  The fine system is
\begin{equation}\label{PRIQ:eq:fine-four-state}
\begin{aligned}
  \dot x_1&=-ax_1+bx_3+h_{12}(x),\\
  \dot x_2&=-ax_2+bx_4-h_{12}(x),\\
  \dot x_3&= ax_1-bx_3+h_{34}(x),\\
  \dot x_4&= ax_2-bx_4-h_{34}(x).
\end{aligned}
\end{equation}

\begin{PRIQproposition}[Exact quotient with nonlinear hidden exchange]\label{PRIQ:prop:hidden-exchange-quotient}
The system \eqref{PRIQ:eq:fine-four-state} preserves total mass and the nonnegative orthant.  After aggregation by $C$, it satisfies the coarse system \eqref{PRIQ:eq:coarse-two-state} exactly.
\end{PRIQproposition}

\begin{proof}
Summing the four equations gives zero.  To check positivity, set one coordinate equal to zero.  Every term entering its derivative is then nonnegative.  This holds for $x_1$, $x_2$, $x_3$, and $x_4$, so the nonnegative orthant is forward invariant.  Finally,
\begin{align*}
  \dot x_1+\dot x_2
  &=-a(x_1+x_2)+b(x_3+x_4),\\
  \dot x_3+\dot x_4
  &=a(x_1+x_2)-b(x_3+x_4),
\end{align*}
so the internal nonlinear channels cancel under $C$ and the aggregate obeys \eqref{PRIQ:eq:coarse-two-state}.
\end{proof}

The quadratic terms in \eqref{PRIQ:eq:h12} and \eqref{PRIQ:eq:h34} are genuine fine hyperedges.  Their output vectors lie in the internal directions of the two fibers,
\[
  (1,-1,0,0)^T,
  \qquad
  (0,0,1,-1)^T,
\]
which are contained in $\PRIQKer C$.  They alter the fine distribution while remaining invisible to the coarse state.  \Cref{PRIQ:fig:exact-quotient-flow} shows this distinction.  The two fine states inside the first fiber move differently, yet their sum follows the coarse trajectory.

\begin{figure}[H]
\centering
\includegraphics[width=0.82\linewidth]{sources/research/hypergraph-bayesian-refinement/figures/exact_quotient_flow.pdf}
\caption{Hidden fine motion under an exact quotient.  The microstates $x_1$ and $x_2$ evolve separately.  Their sum agrees with the coarse coordinate $z_1$ throughout the trajectory.  Sparse markers on the coarse curve make the overlap visible.}
\label{PRIQ:fig:exact-quotient-flow}
\end{figure}

\subsection{Exact nonlinear filtering}
Choose a coarse likelihood
\[
  \lambda^{\PRIQcoarse}=(\lambda_1,\lambda_2)^T>0
\]
and lift it by
\[
  \lambda^{\PRIQfine}=C^T\lambda^{\PRIQcoarse}
  =(\lambda_1,\lambda_1,\lambda_2,\lambda_2)^T.
\]
By Proposition~\ref{PRIQ:prop:hidden-exchange-quotient} together with \Cref{PRIQ:thm:bayes-sufficient,PRIQ:thm:exact-cocycle}, every finite sequence of flow segments and such observations satisfies
\[
  C\mu_k^{\PRIQfine}=\mu_k^{\PRIQcoarse}.
\]
The fine filter may continue to distinguish $x_1$ from $x_2$ and $x_3$ from $x_4$, but those distinctions do not affect the exact coarse posterior.

\begin{figure}[H]
\centering
\includegraphics[width=0.78\linewidth]{sources/research/hypergraph-bayesian-refinement/figures/exact_filtering_residual.pdf}
\caption{Aggregate posterior discrepancy after each prediction and update stage.  The residual remains at floating point scale and stays near double precision epsilon throughout the exact filtering sequence.}
\label{PRIQ:fig:exact-filtering-residual}
\end{figure}

\begin{PRIQcertificate}[Exact nonlinear filtering residuals]\label{PRIQ:cert:exact-filter-residuals}
For the deterministic run with four states, the largest aggregate prediction residual was $\ExactPredictionResidual$.  The largest posterior residual after an update was $\ExactPosteriorResidual$, and the largest residual in a single evidence factor was $\ExactEvidenceResidual$.  The small nonzero values come from floating point integration and arithmetic in the finite example.
\end{PRIQcertificate}

\subsection{A controlled dynamical violation}
Add a directed transition between fine states
\begin{equation}\label{PRIQ:eq:dynamical-perturbation}
  \varepsilon x_1
  \begin{bmatrix}
  -1\\0\\1\\0
  \end{bmatrix}
\end{equation}
to the fine vector field.  This perturbation preserves nonnegativity and total mass, but the added transfer between fibers is activated only by the mass in state $x_1$.  The aggregate field defect is
\begin{equation}\label{PRIQ:eq:perturbation-field-defect}
  \delta_F(x)
  =
  \varepsilon x_1
  \begin{bmatrix}-1\\1\end{bmatrix},
\end{equation}
so, on the simplex,
\begin{equation}\label{PRIQ:eq:perturbation-defect-bound}
  \PRIQnorm{\delta_F(x)}_1
  =2\varepsilon x_1
  \leq2\varepsilon.
\end{equation}
The coarse Markov flow on two states is an $\ell^1$ contraction.  Over one interval, this gives
\[
  \PRIQnorm{C\Phi_T^{\PRIQfine,\varepsilon}(x)-\Phi_T^{\PRIQcoarse}(Cx)}_1
  \leq2\varepsilon T.
\]

\begin{figure}[H]
\centering
\includegraphics[width=0.78\linewidth]{sources/research/hypergraph-bayesian-refinement/figures/flow_perturbation_scaling.pdf}
\caption{Aggregate flow error across a logarithmic perturbation sweep.  The measured discrepancy has first order scaling and remains below the certificate $2\varepsilon T$ from \eqref{PRIQ:eq:perturbation-defect-bound}.}
\label{PRIQ:fig:flow-perturbation-scaling}
\end{figure}

\subsection{A controlled observation violation}
Perturb the fine likelihood inside one fiber:
\begin{equation}\label{PRIQ:eq:likelihood-perturbation-example}
  \lambda^{\PRIQfine,\eta}
  =
  C^T\lambda^{\PRIQcoarse}
  +
  \eta(1,-1,0,0)^T.
\end{equation}
Assume $\PRIQabs{\eta}<\lambda_1^{\PRIQcoarse}$, so every component of $\lambda^{\PRIQfine,\eta}$ remains strictly positive.  The fiber oscillation is $2\PRIQabs{\eta}$ and
\[
  \PRIQnorm{r_{\lambda}}_{\infty}=\PRIQabs{\eta}.
\]
If the evidence is bounded below by $z_*>0$, then
\[
  d_{\PRIQBup}
  \leq
  \frac{2\PRIQabs{\eta}}{z_*}.
\]
The sweep in \Cref{PRIQ:fig:bayesian-perturbation-scaling} varies only the observation model and isolates this effect at one interior prior.  The measured discrepancy is first order in $\eta$ and remains below the likelihood certificate.

\begin{figure}[H]
\centering
\includegraphics[width=0.78\linewidth]{sources/research/hypergraph-bayesian-refinement/figures/bayesian_perturbation_scaling.pdf}
\caption{Posterior discrepancy caused by variation of the likelihood inside one aggregation fiber.  The error scales linearly with $\eta$ in the tested range and remains below the bound obtained from the normalization estimate.}
\label{PRIQ:fig:bayesian-perturbation-scaling}
\end{figure}

Combining this observation perturbation with \eqref{PRIQ:eq:dynamical-perturbation} gives a posterior discrepancy that is first order in $(\varepsilon,\eta)$ over a fixed finite horizon.  \Cref{PRIQ:thm:posterior-certificate} supplies a conservative upper bound.

\begin{figure}[H]
\centering
\includegraphics[width=0.78\linewidth]{sources/research/hypergraph-bayesian-refinement/figures/posterior_error_certificate.pdf}
\caption{Joint dynamical and likelihood perturbation sweep with $\varepsilon=\eta=\rho$.  For small $\rho$, the final posterior discrepancy grows linearly, while the recursive certificate remains above every measured value.}
\label{PRIQ:fig:posterior-error-certificate}
\end{figure}

\begin{PRIQcertificate}[Perturbation scaling]\label{PRIQ:cert:perturbation-scaling}
The experiment used logarithmically spaced values with $\varepsilon=\eta$.  The fitted slope of the final posterior discrepancy was $\PerturbationSlope$.  Every measured error remained below its recursive certificate.
\end{PRIQcertificate}

\subsection{Random exact tensors through order three}
The implementation draws random coarse tensors of orders $m=1,2,3$, lifts them through \eqref{PRIQ:eq:exact-instance-generator}, and adds random hidden output tensors $K_m$ with $CK_m=0$.  It then checks:
\[
  CA_m^{\PRIQfine}-A_m^{\PRIQcoarse}C^{\times m}=0,
\]
annihilation of hidden inputs, constancy on packet fibers, and recovery of the original coarse tensor.  A single coefficient perturbation provides a negative control that must violate fiber constancy.

This deterministic check confirms that the coordinate implementation and the invariant theorem return the same quotient.  It is not evidence for a general probabilistic law.

\subsection{Concentrated hyperedge field}
A synthetic interaction field with decreasing orthogonal block energies illustrates \Cref{PRIQ:thm:best-k-hyperedges}.  The computed hyperedge entropy is $\HyperedgeEntropyValue$, corresponding to an effective count of approximately $\EffectiveHyperedgeCount$ hyperedges despite a larger ambient support.  The ranked spectrum in \Cref{PRIQ:fig:hyperedge-energy-spectrum} shows why the field is compressible in this representation.

\begin{figure}[H]
\centering
\begin{subfigure}[t]{0.48\linewidth}
\centering
\includegraphics[width=\linewidth]{sources/research/hypergraph-bayesian-refinement/figures/hyperedge_energy_spectrum.pdf}
\caption{Ranked normalized hyperedge energy.}
\label{PRIQ:fig:hyperedge-energy-spectrum}
\end{subfigure}
\hfill
\begin{subfigure}[t]{0.48\linewidth}
\centering
\includegraphics[width=\linewidth]{sources/research/hypergraph-bayesian-refinement/figures/hyperedge_pruning.pdf}
\caption{Residual after retaining the $k$ largest energy hyperedges.}
\label{PRIQ:fig:hyperedge-pruning}
\end{subfigure}
\caption{Energy concentration and certified pruning in the synthetic interaction field.  A small leading set carries most of the energy, and the residual decreases to zero when the complete support is restored.}
\label{PRIQ:fig:hyperedge-compression}
\end{figure}

\subsection{Interpretation for reaction networks}
Consider a polynomial reaction system
\[
  \dot x=S v(x),
  \qquad
  v_r(x)=k_r\prod_i x_i^{\alpha_{ir}}.
\]
Each reaction monomial contributes one interaction hyperedge to every species changed by the corresponding stoichiometric column.  The inputs are the reactants.  A block sum aggregation is exact precisely when the monomial coefficients, after summation over each output fiber, are constant on the product fibers specified by \Cref{PRIQ:thm:packet-coordinate-criterion}.  This gives a coefficient test at every interaction order.  Specialized reductions of reaction networks may also use stoichiometric, kinetic, or bisimulation structure \cite{PRIQ:Cardelli2015,PRIQ:Cardelli2017,PRIQ:BaezPollard2017}.

\begin{PRIQboundary}[Physical interpretation]
The quotient theorem is algebraic.  It does not by itself decide whether a packet represents a species class, spatial region, symmetry sector, latent state, or learned feature.  The scientific meaning of $C$ must come from the application and is not created by exact factorization alone.
\end{PRIQboundary}


\chapter{Relations, limitations, and further problems}\label{PRIQ:book:chapter-45}
\section{Relation to neighboring theories}\label{PRIQ:sec:relations}

The results fit into a common finite quotient language, but each ingredient also belongs to a larger established theory.

\subsection{Markov lumpability and Markov functions}
For unary stochastic dynamics, the identity
\[
  CQ_{\PRIQfine}=Q_{\PRIQcoarse}C
\]
is a standard exact lumping condition.  It implies
\[
  Ce^{tQ_{\PRIQfine}}=e^{tQ_{\PRIQcoarse}}C
\]
and means that the coarse process is Markov for every fine initial distribution compatible with the quotient.  Classical treatments distinguish several forms of strong, ordinary, exact, and weak lumpability, depending on the admissible initial laws and the direction of the stochastic convention \cite{PRIQ:KemenySnell1960,PRIQ:RogersPitman1981,PRIQ:Buchholz1994}.  Here we fix one column vector convention and extend the coefficient criterion from unary edges to multilinear hyperedges.

\subsection{Polynomial differential equivalence}
Exact aggregation of polynomial vector fields has been developed through differential equivalences, bisimulations, and algorithms based on successive partition refinement, including applications to chemical reaction networks \cite{PRIQ:Cardelli2015,PRIQ:Cardelli2017}.  The universal factorization theorem here is elementary multilinear algebra once $C$ is fixed.  Its value is that it exposes every interaction order, makes the quotient condition independent of a particular polynomial syntax, and connects the same coefficient criterion to rooted tree composition and Bayesian updates.  We do not claim an algorithmic advantage over specialized methods based on partition refinement.

\subsection{Hypergraphs, operads, and rooted trees}
A multilinear packet block has several inputs and one output, so its incidence is naturally hypergraphic.  Rooted substitution of such blocks is operadic, and the derivative formula for composition is the multivariate Fa\`a di Bruno formula \cite{PRIQ:ConstantineSavits1996}.  Rooted trees also organize Butcher series and expansions for nonlinear time stepping \cite{PRIQ:Butcher1972,PRIQ:HairerLubichWanner2006}.  In this language, the quotient condition is a morphism condition.  It commutes with every generating operation and therefore with every rooted composite.

\subsection{Sufficiency and comparison of experiments}
A statistic is sufficient only relative to an observation family.  Classical sufficiency and comparison of experiments theory formalize when one representation preserves all information relevant to a chosen family \cite{PRIQ:HalmosSavage1949,PRIQ:Blackwell1953}.  The likelihood theorem in \Cref{PRIQ:sec:bayes} is a finite and explicit special case for deterministic aggregation and diagonal positive observation operators.  It characterizes exact commutation of the normalized update, not all possible notions of statistical equivalence.

\subsection{Model reduction and nearly decomposable systems}
Classical approaches to approximate aggregation use time scale separation, perturbation theory, balanced truncation, stochastic complementation, and related forms of model reduction \cite{PRIQ:SimonAndo1961,PRIQ:Meyer1989,PRIQ:Antoulas2005}.  The bounds used here follow the error through each stage.  Coefficient defects first create a field defect.  The Grönwall inequality carries that defect through the flow, and the normalization estimate carries it through each observation.  The resulting certificate can be conservative, but it keeps every source of error visible.

\section{Claim boundaries and limitations}\label{PRIQ:sec:limitations}

\begin{PRIQboundary}[Finite dimensional foundation]
All quotient and coefficient theorems are proved in finite dimensional vector spaces.  The algebraic factorization extends to arbitrary vector spaces, but boundedness, convergence, measurability, domains of unbounded generators, and interchange of infinite sums require additional hypotheses in Banach, Hilbert, or measurable field settings.
\end{PRIQboundary}

\begin{PRIQboundary}[Local analytic dynamics]
The flow theorems assume local existence and uniqueness.  The analytic coefficient theorem supplies a local coarse field, not a global invariant manifold or a global conjugacy.  Solutions may blow up, uniqueness may fail, and boundary singularities remain specific to the application.
\end{PRIQboundary}

\begin{PRIQboundary}[A prescribed coarse map]
The principal theorems assume that $C$ is supplied.  They test and quantify whether a proposed aggregation is valid.  Discovering an optimal $C$ from data or from a large interaction tensor is a separate problem and can be combinatorial, nonconvex, or statistically nonidentifiable.
\end{PRIQboundary}

\begin{PRIQboundary}[Positivity and evidence floors]
Bayesian normalization requires positive evidence.  The uniform perturbation bounds degrade as the evidence floor approaches zero.  This is not merely a weakness of the proof.  Small normalizing constants can amplify small absolute errors into large posterior changes.
\end{PRIQboundary}

\begin{PRIQboundary}[Norm dependence]
Exact descent is algebraic and independent of the norm.  Best approximate coarse tensors, Lipschitz constants, pruning choices, and posterior certificates depend on the selected norms and task weights.  A useful certificate must state those choices.
\end{PRIQboundary}

\begin{PRIQboundary}[Numerical verification]
The executed checks validate the finite implementation against the stated formulas and include negative controls.  They do not prove the theorems, establish asymptotic behavior outside the tested ranges, or demonstrate universal computational speedup.
\end{PRIQboundary}

\begin{PRIQboundary}[Historical priority]
The central synthesis joins several established families of ideas.  The coefficient criterion at every interaction order follows from multilinear algebra on a quotient space.  Equivalent statements may already appear in neighboring theories of nonlinear aggregation under different notation.  Any claim of novelty should distinguish this elementary factorization from the combined treatment of interaction trees, Bayesian cocycles, and error certificates.  Such a claim also requires a dedicated literature review.
\end{PRIQboundary}

\section{Open theorem and algorithm program}\label{PRIQ:sec:open}

\subsection{Coarsest exact joint partition}
For a finite interaction family $(A_m^{\PRIQfine})_{m\leq M}$ and likelihood family $(\lambda_y^{\PRIQfine})_{y\in\PRIQcY}$, seek
\begin{equation}\label{PRIQ:eq:coarsest-joint-partition}
  \min_{\pi}|I_{\PRIQcoarse}|
\end{equation}
subject to
\[
  CA_m^{\PRIQfine}=A_m^{\PRIQcoarse}C^{\times m}
  \quad(m\leq M),
\]
and
\[
  \lambda_y^{\PRIQfine}=s_yC^T\lambda_y^{\PRIQcoarse}
  \quad(y\in\PRIQcY).
\]
The partition determined by the observations alone groups states with identical likelihood profiles.  The dynamical condition also depends on the partition because changing the partition changes the output aggregation.  A terminating refinement algorithm, together with a useful complexity bound, remains an open problem.

\subsection{Approximate joint design}
A finite approximate problem is
\[
  \min_{C,\{A_m^{\PRIQcoarse}\},\{\lambda_y^{\PRIQcoarse}\}}
  \left[
  \sum_mw_m\PRIQnorm{CA_m^{\PRIQfine}-A_m^{\PRIQcoarse}C^{\times m}}
  +\sum_yv_yd_{\PRIQBup,y}
  +\gamma\operatorname{Complexity}(C)
  \right].
\]
The main questions concern existence, tractable relaxations, and stability of the minimizers.  Another question is whether the task weights can be learned without invalidating the certificate after a distribution shift.

\subsection{Posterior distributions over resolutions}
Instead of selecting one $C$, place a posterior on a finite family of candidate resolutions:
\[
  p(C\mid D)
  \propto
  p(D\mid C)p(C).
\]
Predictions then average both parameter and topology uncertainty:
\[
  p(z_t\mid D)
  =
  \sum_Cp(z_t\mid C,D)p(C\mid D).
\]
A fuller model should distinguish uncertainty in the coefficients of a fixed hypergraph from uncertainty about the hypergraph and the quotient map themselves.

\subsection{Stochastic interaction quotients}
For a diffusion
\[
  \PRIQdd X_t=b(X_t)\PRIQdd t+\sigma(X_t)\PRIQdd W_t,
\]
the generator is
\[
  \PRIQcL\varphi
  =D\varphi[b]
  +\frac12\PRIQtr\bigl(\sigma\sigma^TD^2\varphi\bigr).
\]
Packet resolution turns drift into unary edges and quadratic variation into contracted binary hyperedges.  Let $\PRIQcC$ denote the pullback or averaging map on observables induced by the aggregation.  The natural exact condition is generator intertwining
\[
  \PRIQcC\PRIQcL_{\PRIQfine}=\PRIQcL_{\PRIQcoarse}\PRIQcC
\]
for the pullback or averaging map appropriate to observables.  A complete theorem would characterize closure of the drift and covariance, prove positivity of the induced semigroup, and carry generator defects into filtering error bounds.

\subsection{Coarse dynamics with memory}
Eliminating hidden packets generally produces memory.  The target form is
\[
  \dot z(t)
  =A_{PP}z(t)
  +\int_0^tK(t-s)z(s)\PRIQdd s
  +\eta(t).
\]
The quotient theory should allow a fine system without memory to produce coarse dynamics with memory.  Volterra hyperedges and fractional kernels would then be treated as native coarse interactions rather than as evidence that quotienting failed.

\subsection{Sharper filtering metrics}
The $\ell^1$ normalization bounds are robust but conservative.  Positive operators often contract projective metrics, relative entropy, or other information geometric divergences more sharply.  One goal is to replace the factor
\[
  \frac{2\PRIQnorm{\lambda}_{\infty}}{z_*}
\]
by a contraction coefficient adapted to the trajectory or to the operator, while retaining an auditable finite certificate.

\subsection{Pruning that preserves structure}
In Hamiltonian, reversible, conservative, and monotone systems, as well as systems with detailed balance, individual hyperedges cannot always be removed independently.  The problem is to identify groups whose joint retention preserves the required structure, then solve the budget problem over those groups.

\subsection{Identifiability of hidden dynamics}
For fixed $C$ and coarse field $F_{\PRIQcoarse}$, the exact fine realizations form the affine family
\[
  \{F_{\PRIQfine}:CF_{\PRIQfine}=F_{\PRIQcoarse}C\}.
\]
Any two exact fine models differ by a map into $\PRIQKer C$.  Coarse observations alone cannot identify that hidden component.  Additional fine observations, intervention assumptions, regularity, or minimum complexity principles are required.  A statistical theory should characterize precisely which hidden hyperedges are identifiable under a given observation and intervention design.

\section{Conclusion}
An exact aggregation is more than a useful summary of a fine model.  It is a quotient.  Fine states with the same coarse image must produce the same coarse output.  For a multilinear interaction, this reduces to a complete finite test.  After applying $C$, the interaction must vanish whenever any input lies in $\PRIQKer C$.  In packet coordinates, the coefficients summed over each output fiber must be constant on every product of input fibers.  The familiar unary criterion for Markov lumpability is the first order case of this general condition.

The coefficient test controls the nonlinear constructions built from those interactions.  It passes to analytic vector fields, their derivatives, and every rooted interaction tree.  Under the usual local uniqueness assumptions, it also passes to the flow.  Bayesian observation has a matching quotient condition.  A normalized update commutes with aggregation exactly when the fine likelihood is constant within each fiber, apart from a positive global scale.  Once both conditions hold, every finite sequence of predictions and updates agrees after aggregation, and the evidence factors agree up to the same scales.

When exactness fails, the reason matters.  A mismatch in the interaction coefficients creates a field defect, and the flow can amplify it.  Variation of the likelihood inside a fiber creates a second error at normalization.  Removing hyperedges creates a third error in the form of an explicit residual tensor.  The posterior recurrence keeps these contributions separate as it carries them forward.  An approximate coarse model is therefore auditable because one can see where each part of the final error entered.

The theory does not show that an optimal quotient can always be found cheaply.  Coarse observations alone also cannot determine hidden dynamics that lie in $\PRIQKer C$.  Those questions remain open.  The result proved here is narrower and exact:
\begin{center}
\fbox{\begin{minipage}{0.88\linewidth}
A coarse model is exact when its interactions and observations descend through the same quotient.  When they do not, each failure can be measured at the stage where it occurs and carried forward to the final posterior or observable by an explicit certificate.
\end{minipage}}
\end{center}


\part{Machine understanding and explanatory transfer}
\setcounter{section}{0}\renewcommand{\thesection}{MU.\arabic{section}}\renewcommand{\theHsection}{MU.\arabic{section}}

\chapter{Representations and implemented use}\label{MU:book:chapter-46}


\section{What a correct answer leaves open}\label{MU:sec:intro}
A reasoner is asked what a device will display after two operations. It gives the right answer. That success is a fact worth preserving. It does not yet tell us whether the reasoner has a method that will work after three operations, whether the order of the operations matters, or whether the same method could be expressed in another notation. It also does not tell us whether the explanation offered afterward describes the computation that produced the answer.

These are not objections to correctness. They are different questions about what correctness rests on and what it makes possible. A proof can establish a conclusion from premises without teaching a recipient how to select the relevant premises in a new case. A useful model can guide new predictions without being a faithful description of the particular implementation that supplied it. A fluent explanation can be understandable while failing to identify the mechanism it claims to explain. The distinction between persuasive and faithful explanations is explicit in interpretability research~\cite{MU:jacovi}.

This paper investigates a narrower and more tractable question than whether a machine understands in every human sense of the word. It asks when a computational realization possesses a representation that remains useful across a declared family of operations, questions, translations, and corrections. We call this \emph{reusable computational understanding} as an operational research proposal. It is relative to a domain and a contract. It is not a metaphysical property inferred from a single formal condition.

The central thesis is that a representation earns its explanatory role through what a system can do with it. It must retain distinctions needed for future questions, support the operations it claims to summarize, and expose the conditions under which its use should change. When it is communicated, a recipient must be able to recover something useful from it. When it is attributed to a particular mechanism, the attribution requires evidence about that mechanism. A representation that meets one of these obligations need not meet the others.

\subsection{Three separations organize the argument}
The first separation is between an answer and a reusable state. The answer to today's question can discard distinctions that tomorrow's permitted operation makes important. Understanding in the present sense therefore concerns a family of possible continuations, not only a completed input and output pair.

The second is between a correct account of a subject and a faithful account of a solver. An algebraic derivation can explain a numerical result even when the original solver used a table. It would be wrong to infer that the solver executed the derivation merely because the derivation is valid. Subject understanding and process faithfulness require different witnesses.

The third is between successful transport and preserved correction. Two languages may translate every currently expressible observation while one language merges distinctions needed to reject a mistaken model. Agreement in the translated language is then weaker than agreement about what has been learned. The correction theorem in Section~\ref{MU:sec:correctiontransport} makes that loss exact.

\subsection{Position among existing theories}
Several neighboring theories already supply essential pieces. Computational mechanics studies minimal predictive state representations~\cite{MU:shalizi}. Controlled state aggregation studies equivalences that preserve transition structure~\cite{MU:givan}. Information bottleneck methods formalize compression relative to a specified relevance variable~\cite{MU:tishby}. These are not interchangeable claims, but they establish that useful compression must be assessed against what it promises to preserve.

Causal abstraction gives a stronger account of agreement under interventions. Beckers and Halpern distinguish increasingly restrictive abstraction conditions~\cite{MU:beckers}. Geiger and colleagues develop this perspective as a foundation for mechanistic interpretability, including transformations of mechanisms~\cite{MU:geiger}. Our finite commuting equations occupy a deliberately restricted part of that subject. They do not replace a general theory of causal models.

There are also reasons to constrain the maps used in an interpretation. Hewitt and Liang show why a trained probe's success must be assessed against the work the probe itself can perform~\cite{MU:hewitt}. Sutter and colleagues examine how sufficiently expressive alignment maps can make some causal-abstraction claims uninformative~\cite{MU:sutter}. The constructive theory and this criticism address different aspects of the problem. We retain both lessons: intervention compatibility matters, and the translator must not be allowed to hide the entire explanation inside an unrestricted computation.

Explanation has a recipient as well as a referent. Miller's synthesis emphasizes the contextual and social features of explanation~\cite{MU:miller}. Gentner's structure-mapping account places relational organization, rather than surface resemblance alone, at the center of analogy~\cite{MU:gentner}. Here we formalize only a limited recipient-relative property: an artifact enables specified future operations in a recipient with stated background resources. That property does not exhaust the reasons a human explanation can succeed or fail.

Finally, learning through informative queries and counterexamples has a substantial theoretical history, including Angluin's work on regular languages~\cite{MU:angluin}. Our correction model is much simpler. It is a finite filter on candidate models, chosen so that the interaction between evidence and translation can be proved directly.

\subsection{The intended contribution}
The paper does not propose that every useful representation is new, that a symbolic method is inherently superior to a learned one, or that a large table can never support understanding. It proposes a common account of the obligations that such representations must satisfy. The exact finite results locate information requirements and compatibility conditions. The counterexamples show why isolated successes cannot silently substitute for a combined claim.

This is also a limit on the theory itself. A finite contract can be satisfied by many implementations, including tables. To distinguish them, one must add relevant constraints on acquisition, causal organization, storage, computation, or transfer. Those constraints cannot be recovered from a definition after the results are known.

\section{A small instrument with a hidden register}\label{MU:sec:instrument}
Consider a mathematical instrument with three registers
\[
x=(a,b,c)\in\MUFfive^3.
\]
Each register contains one of $0,1,2,3,4$. All arithmetic in this example is modulo five. The display reports
\[
q(a,b,c)=a+b.
\]
There are $125$ possible complete states, but only five possible readings. The registers and their operations are fully specified analytical objects. They are not measurements from an agent experiment.

Initially, two operations are available:
\begin{equation}\label{MU:eq:instrument-ops}
H(a,b,c)=(a+b,b,c),\qquad W(a,b,c)=(b,a,c).
\end{equation}
The first adds the second register into the first. The second exchanges the first two registers. The third register is preserved by both and does not affect the display. The operation order is public. A prediction concerns what the display would show after a named finite sequence of these operations.

\subsection{Knowing the current display is insufficient}
A record containing only $s=a+b$ answers the current question exactly. It cannot, in general, predict the display after $H$. The two complete states
\[
x=(0,0,0),\qquad y=(1,4,0)
\]
both display zero. After $H$, their displays are zero and four. Any predictor receiving only the common old display must return the same answer for both, so it cannot be exact in both cases.

A sufficient retained state for the original operation family is
\[
z=(s,b)=(a+b,b).
\]
Its updates are
\begin{equation}\label{MU:eq:display-coordinates}
\bar H(s,b)=(s+b,b),\qquad \bar W(s,b)=(s,s-b),
\end{equation}
and its displayed output is its first coordinate. These equations retain the information needed for every finite $H/W$ continuation without retaining $c$.

The representation is also minimal in its number of states. The present reading $s$ and the reading $s_H$ after a hypothetical $H$ determine
\[
b=s_H-s,\qquad a=2s-s_H.
\]
Two different pairs $(a,b)$ therefore differ on at least one of these questions. Every exact reusable representation for this contract needs at least $25$ distinguishable values. Section~\ref{MU:sec:quotients} gives the general argument.

\subsection{A new operation changes what must be known}
Now admit
\begin{equation}\label{MU:eq:new-operation}
J(a,b,c)=(a+c,b,c).
\end{equation}
The old representation remains correct for $H$ and $W$. It is no longer sufficient for the enlarged contract. States with the same $(a,b)$ but different $c$ now produce different readings after $J$.

The additional reading $s_J$ yields
\[
c=s_J-s.
\]
Together with the earlier distinctions for $a$ and $b$, the enlarged question family distinguishes all $125$ complete states. The useful state counts are therefore five for the current display, $25$ for reusable $H/W$ prediction, and $125$ after $J$ is admitted.

\begin{figure}[htbp]
\centering\includegraphics[width=.98\linewidth]{sources/research/machine-understanding/figures/representation_scope.pdf}
\caption{The same instrument supports different minimal representations under different contracts. The numbers count distinguishable states, not scores of intelligence. A new permitted operation makes a previously irrelevant register relevant.}\label{MU:fig:scope}
\end{figure}

This is not a case of an old theorem becoming false. The old theorem had a smaller operation family. What failed was an unlicensed extension of its conclusion. The distinction matters whenever a learned method is reused under a change in task, environment, or intervention.

\subsection{Correction requires evidence, not a renamed variable}
A predictor that has retained only $(s,b)$ cannot recover $c$ merely by adding a third field to its data structure. Conditional on a fixed $(a,b)$ and a uniformly distributed $c$, every possible display after $J$ is equally likely. Any guess based only on the old record succeeds with probability $1/5$. Randomizing the guess does not improve that probability.

An actual execution of $J$, followed by a new reading, supplies the missing distinction. Because $J$ changes the instrument, the recovered state is the state \emph{after} that operation:
\begin{equation}\label{MU:eq:recover-after-J}
(a',b',c')=(s_J-b,\ b,\ s_J-s).
\end{equation}
The old state can then be reconstructed because this particular operation is invertible. Neither reversibility nor a free observation should be assumed for other instruments.

The example separates four achievements: reading an answer, retaining a method, recognizing a contract change, and acquiring evidence that repairs the resulting information gap. No amount of confidence in the original method supplies the fourth.

\section{A contract for reusable understanding}\label{MU:sec:contract}
Let $\MUX$ and $\MUY$ be finite nonempty sets. Let $\MUU$ be a finite operation alphabet, possibly empty, and let $\MUQ$ be a finite nonempty family of queries $q:\MUX\to\MUY$. Each $u\in\MUU$ has an implemented total transition $F_u:\MUX\to\MUX$. The deterministic mathematical contract is
\[
\MUC=(\MUX,\MUU,(F_u)_{u\in\MUU},\MUQ,\MUY).
\]
For a word $w=u_1\cdots u_n$, define $F_w=F_{u_n}\circ\cdots\circ F_{u_1}$. The empty word acts as the identity. Thus $q(F_w(x))$ asks a question after a stated continuation.

Totality is a simplifying assumption. If an operation is sometimes unavailable, its admissibility must be included in the observable contract or the model must use partial transitions with explicit domain preservation. An unmodeled refusal cannot be treated as the same event as a permitted update.

\begin{MUdefinition}[Exact reusable representation]\label{MU:def:representation}
An exact reusable representation for $\MUC$ is a surjection $r:\MUX\to\MUZ$, updates $\bar F_u:\MUZ\to\MUZ$, and decoders $\bar q:\MUZ\to\MUY$ such that, for all $x,u,q$,
\begin{equation}\label{MU:eq:representation-contract}
r(F_u(x))=\bar F_u(r(x)),\qquad q(x)=\bar q(r(x)).
\end{equation}
\end{MUdefinition}
Induction immediately gives $q(F_w(x))=\bar q(\bar F_w(r(x)))$ for every finite word. The same initial record supports many future questions. It is not allowed to inspect the selected future query and replace itself with that query's answer before the challenge begins.

\subsection{A representation of the problem is not yet a representation in the solver}
The map $r$ specifies what information would suffice. It does not establish that a particular agent can acquire $r(x)$ from its sensors or that the agent uses this state. Those claims require an acquisition procedure, a realized computational path, and admissible interventions on that path.

For the instrument, an omniscient mathematical encoder can read $(a,b,c)$. A device restricted to the display cannot do so. It can obtain $b$ through a controlled $H$ observation, but that procedure costs an operation and changes the state. Any account of its understanding must include those effects rather than identify mathematical existence with operational access.

A practical contract therefore adds an information boundary, a family of recipient or decoder algorithms, resource limits, an intervention catalogue, and a revision policy. The finite theorems below solve particular obligations within that larger contract. They do not fill in unspecified sensors, access rights, or learning procedures.

\subsection{A profile rather than a single threshold}
We propose reporting a profile with separate entries for predictive adequacy, implemented use, intervention compatibility, transfer, recipient reconstruction, correction, and resource cost. Unknown entries remain unknown. A high prediction score must not silently fill the mechanism or correction entries.

The profile can describe a traditional program, a learned model, a hybrid system, or a collective procedure. It does not require language-like internal states. Nor does it establish an ordering of all systems from least to most understanding. Two systems can preserve different families of questions and be incomparable under the declared criteria.

\begin{MUcounterexample}[A separate answer for each question is not one reusable record]\label{MU:ex:joint}
Let $x=(a,b)\in\{0,1\}^2$ and let the two queries return $a$ and $b$. If the encoder knows which query will be asked, one bit suffices: retain that answer. If one record must be fixed before the query is selected and must answer both exactly, it needs four distinct values. A one-bit record has only two values and must identify two different answer profiles.
\end{MUcounterexample}
The order of quantifiers is decisive. A different adequate representation for every question is not one adequate representation for their union. This is why a claimed method should be assessed with its encoder, updater, and decoder fixed across the family it promises to support.

\section{Behavioral compression and the price of new questions}\label{MU:sec:quotients}
Define
\begin{equation}\label{MU:eq:equivalence}
x\MUEq y\quad\Longleftrightarrow\quad q(F_w(x))=q(F_w(y))\ \text{for all }q\in\MUQ,\ w\in\MUU^*.
\end{equation}
This is equality of all declared future answer profiles. It is an equivalence relation. It describes behavioral sufficiency under the contract, not equality of hidden mechanisms or identity outside that contract.

\begin{MUtheorem}[The minimal exact reusable state]\label{MU:thm:quotient}
The quotient $\MUX/{\MUEq}$ carries an exact reusable representation for $\MUC$. Every other exact reusable representation $r:\MUX\to\MUZ$ refines this quotient: if $r(x)=r(y)$, then $x\MUEq y$. There is a unique map $h:\MUZ\to\MUX/{\MUEq}$ with $h(r(x))=[x]$. In particular,
\[
|\MUZ|\geq|\MUX/{\MUEq}|.
\]
A fixed-length binary state needs at least $\lceil\log_2|\MUX/{\MUEq}|\rceil$ bits. The quotient can be computed by finite partition refinement.
\end{MUtheorem}
\begin{proof}
If $x\MUEq y$, the empty word gives $q(x)=q(y)$ for every query. Moreover, $F_u(x)\MUEq F_u(y)$ because any continuation after $u$ is a continuation in the original contract. Thus $[x]\mapsto[F_u(x)]$ and $[x]\mapsto q(x)$ are well-defined updates and decoders.

For any exact representation, the equations in~\eqref{MU:eq:representation-contract} imply equality of all future decoded answers whenever $r(x)=r(y)$. This proves refinement. Surjectivity of $r$ gives existence and uniqueness of $h$ and the cardinality bound. The binary bound follows by counting available codewords.

For computation, let $P_0$ group states with identical current query vectors. Given a partition $P_k$, refine each block by the block of $F_u(x)$ for every $u$. Induction shows that $P_k$ identifies exactly the states agreeing on words of length at most $k$. Every strict refinement increases the block count. At most $|\MUX|-1$ strict refinements are possible. At a fixed point the partition is closed under all operations, hence under every finite word, and is exactly~\eqref{MU:eq:equivalence}.
\end{proof}

This is a deterministic state-minimization result, presented in the language of reusable questions. It belongs to the established state-aggregation and predictive-state lineage~\cite{MU:givan,MU:shalizi}. Its significance here is not the discovery of quotients. It is the explicit identification of which future distinctions an asserted understanding is required to preserve.

\subsection{Minimal information is not minimal effort}
The theorem minimizes the number of distinct retained states. It does not minimize the work required to acquire the state, compute an update, verify a query, or discover the quotient. A smaller table can have harder entries to compute. A larger representation can make a frequently used transformation local and cheap.

Even uniqueness of the behavioral quotient does not supply a uniquely intelligible notation. The coordinates $(a,b)$ and $(s,b)$ describe the same $25$ behavioral classes in the original instrument. One makes the physical register update obvious. The other makes the readout immediate. Their difference is computational and explanatory rather than a difference in information content.

\subsection{The cost of expanding the contract}
Let $\MUC'$ have the same complete state space and contain every operation and query of $\MUC$, while possibly adding more. Write $\equiv_{\MUC'}$ for its finer behavioral equivalence. Take any representation $r$ adequate for $\MUC$. For each $z$, let $k_z$ be the number of $\equiv_{\MUC'}$ classes that meet the fibre $r^{-1}(z)$.

\begin{MUproposition}[Contract expansion and supplementary information]\label{MU:prop:expansion}
The relation $\equiv_{\MUC'}$ refines $\equiv_{\MUC}$. Any supplementary code $s:\MUX\to\mathcal S$ for which $(r,s)$ supports exact reusable behavior under $\MUC'$ must satisfy
\[
|\mathcal S|\geq\max_z k_z.
\]
If every $\equiv_{\MUC'}$ class lies within a single fibre of $r$, this bound is attained by an exact reusable pair $(r,s)$. This condition holds when $r$ is the old minimal behavioral quotient. A supplementary value computed solely from $r(x)$ cannot repair a fibre meeting more than one expanded class.
\end{MUproposition}
\begin{proof}
Agreement on the larger continuation family implies agreement on the smaller one. Within a fixed fibre of $r$, an adequate pair $(r,s)$ must separate the $k_z$ new classes, giving the lower bound. For attainment under the additional condition, assign distinct labels $1,\ldots,k_z$ to the expanded classes within each old fibre and let $s(x)$ be its class label. The pair $(r(x),s(x))$ then identifies exactly the expanded behavioral class. The quotient update and query decoder supply its exact dynamics. Finally, any value computed only from $r(x)$ is constant on that fibre and separates none of its states.
\end{proof}

The attainment condition matters. An old representation can retain irrelevant extra coordinates that cease to be dynamically closed under new operations. For example, start with two bits $(a,b)$, no operations, a constant-zero query, and the unnecessarily detailed record $r(a,b)=a$. Add the swap operation while keeping the query constant. There is still only one expanded behavioral class, so $k_z=1$ in each old fibre. Yet updating the retained coordinate $a$ after the swap requires $b$. The information lower bound for answers is not sufficient to preserve arbitrary redundant coordinates. The minimal old quotient has no such problem.

In the instrument, each old $(a,b)$ class contains five expanded classes. An additional five-valued symbol is sufficient and necessary. A separately stored fixed-length symbol therefore needs three bits. This should not be confused with subtracting the rounded bounds for a joint $25$-state and $125$-state encoding, which are five and seven bits respectively. Separate buffers and joint recoding are different coding contracts.

\begin{MUcounterexample}[A finite continuation test leaves an unbounded alternative]\label{MU:ex:horizon}
For any fixed $h\geq0$, a one-operation machine can return zero through the first $h$ steps and return one at step $h+1$. A constant-zero machine agrees with it on all those earlier observations. Without a state bound or another restriction on the hypothesis family, success through horizon $h$ cannot establish agreement at every future horizon.
\end{MUcounterexample}
This does not invalidate finite certification of a supplied finite model. It distinguishes exhaustive checking of a known transition system from identifying an unrestricted unknown system by a finite record.

\section{Available information is not implemented understanding}\label{MU:sec:use}
The finite quotient answers an information question. A claim about a particular computational reasoner must answer an additional question: where is the useful representation realized, and how does it affect the computation?

A representation can contain a recoverable fact without the system using that fact. Conversely, a system can use information through a distributed or redundant encoding that a chosen simple probe fails to recover. The comparison of information structures in statistical decision theory provides a useful background for asking what a record could support~\cite{MU:blackwell}. That question is different from identifying the route by which one implemented decision was made.

\subsection{A carrier can be readable and irrelevant}
Consider a solver state $m=(a,b,c,z)$, with $z=a+b$. Suppose its output circuit reads $a$ and $b$ directly and ignores $z$. A decoder applied to $z$ predicts the output perfectly. Replacing $z$ while keeping the actual circuit inputs unchanged has no effect. The code $z$ is available to an observer but is not the operative route used by this solver.

Now alter the implementation so that the output reads $z$, while a repair circuit recomputes $z$ from $a,b$ before every output. An attempted replacement may again have no observable effect, but for a different reason: the affected carrier is restored. A null intervention cannot identify which of these implementations is present unless its timing and the remaining information routes are controlled.

These examples justify neither a universal ablation test nor a demand that every useful representation be uniquely necessary. They show why an intervention must specify what changes, what remains accessible, and what the hypothesis predicts afterward. A proof about the instrument's state does not substitute for an intervention on the solver's implementation.

\begin{MUproposition}[The identification limit of a transcript]\label{MU:prop:identification}
Suppose two proposed implementations induce complete transcript laws $P_0$ and $P_1$ on a finite transcript alphabet under a specified permitted experiment, including any adaptive choices made by its evaluator. With equal prior probability, the greatest identification accuracy available from that transcript is
\[
\frac12\bigl(1+\MUTV(P_0,P_1)\bigr).
\]
In particular, identical transcript laws permit no accuracy above $1/2$.
\end{MUproposition}
\begin{proof}
Let $d(o)\in[0,1]$ be the probability of guessing implementation 0 after transcript $o$. Accuracy is
\[
\frac12\sum_o\bigl(P_0(o)d(o)+P_1(o)(1-d(o))\bigr)
=\frac12+\frac12\sum_o(P_0(o)-P_1(o))d(o).
\]
The sum is maximized by setting $d(o)=1$ on positive differences and zero on negative ones. The positive differences sum to total variation. This choice attains the bound.
\end{proof}
The experiment may include interventions. The proposition says only that the \emph{permitted complete transcript} bounds identification. Adding a new sensor or an admissible internal intervention changes that transcript law and can change the limit.

\subsection{The interpreter must not secretly do all the work}
Take any bijection $p:\MUX\to\MUZ$ of a finite state space. Define
\[
\bar F_u=p\circ F_u\circ p^{-1},\qquad \bar q=q\circ p^{-1}.
\]
All exact representation equations hold. The labels in $\MUZ$ may nevertheless be arbitrary, and the translated updates may require a complete lookup table. Exact representability alone does not explain why the representation is easier to acquire, use, or communicate.

This does not make the construction false. It makes the claim narrow. A table can be a valid finite model. Its interpretive value depends on the allowed questions and costs. The mistake would be to infer that any such recoding identifies a compact or causally privileged concept inside a solver.

We therefore require an \emph{interface account}. The encoder, decoder, intervention implementation, and translator are part of the evaluated system. Their information access is declared, their computation is charged, and any training or target-specific background they receive is recorded. A decoder that is free to solve the whole task is not evidence that the state it reads supplied the method. Control tasks for probes and the critique of unrestricted alignment maps are especially relevant here~\cite{MU:hewitt,MU:sutter}.

Restrictions can concern circuit size, memory, communication rounds, a finite program language, or an explicit family of allowed maps. There is no universally correct restriction. A linear map is not automatically the right representation of a nonlinear concept, and a nonlinear map is not automatically uninformative. The appropriate restriction must be selected with the scientific question, not retrofitted to distinguish preferred results.

\subsection{A bounded stochastic extension}
The distinction between a current answer and a reusable state also survives stochastic dynamics. Let $K_u(x,x')$ be finite transition kernels and let $r:\MUX\to\MUZ$ be surjective. Assume the current query values are constant on fibres of $r$.

\begin{MUproposition}[Uniform controlled stochastic reuse]\label{MU:prop:lumpability}
A single reduced kernel $\bar K_u$ reproduces every projected one-step law from every complete state if and only if, whenever $r(x)=r(y)$,
\[
\sum_{v:r(v)=z'}K_u(x,v)=\sum_{v:r(v)=z'}K_u(y,v)
\quad\text{for every }u,z'.
\]
With matched initial projected laws, the reduced model preserves the law of every finite sequence of reduced states and queries under any policy depending only on the reduced history.
\end{MUproposition}
\begin{proof}
A common projected row is necessary by considering point-mass initial states. If the rows agree, use their shared value to define $\bar K_u(r(x),z')$. Conditional on any reduced history, the complete state has a distribution supported on the current fibre. Averaging equal projected rows gives the same reduced transition law. Induction yields equality of the finite history laws, including the deterministic query values. A common randomized history policy preserves the same argument after conditioning on its selected action.
\end{proof}
This is the controlled lumpability condition underlying stochastic state aggregation~\cite{MU:givan}. The policy qualification is essential. A controller that secretly observes a discarded variable can choose different actions in states that the reduced description identifies. Small prediction error on one observed trajectory distribution is not a substitute for the uniform condition.

\section{Transporting a method rather than matching an answer}\label{MU:sec:transport}
An analogy is useful when relations and operations survive the comparison, not simply when two descriptions use similar words. Structure-mapping theory makes relational correspondence central to analogy~\cite{MU:gentner}. We study a more limited mathematical question: when can an already specified method be translated without changing the answers to its declared continuation questions?

Let a source representation have state space $Z$, updates $f_u$, and query decoders $q$. Let a target have state space $W$, updates $g_v$, and decoders $q'$. A map $\tau$ translates operation names. Query pairs $(q,q')$ and a representation map $T:Z\to W$ are fixed in advance. The operation names alone do not establish that the target implements the intended source operation.

\begin{MUtheorem}[Exact transport of reusable operations]\label{MU:thm:transport}
Suppose, for every source operation and matched query,
\begin{equation}\label{MU:eq:transport-exact}
T f_u=g_{\tau(u)}T,\qquad q'T=q.
\end{equation}
Then, for every source state and every finite source word $w$,
\[
T f_w=g_{\tau(w)}T,\qquad q'(g_{\tau(w)}(Tz))=q(f_w(z)).
\]
Two such transports compose, with their operation and query correspondences composed as well. No injectivity assumption on $T$ is needed for this forward guarantee.
\end{MUtheorem}
\begin{proof}
The empty word is immediate. If the identity holds for $w$, then
\[
T f_u f_w=g_{\tau(u)}T f_w
=g_{\tau(u)}g_{\tau(w)}T.
\]
Induction proves the word identity, and applying the matched decoder proves the answer identity. For a second transport $S$ with operation map $\sigma$, substitution gives $STf_u=h_{\sigma(\tau(u))}ST$ and the corresponding decoder identity. This is the same condition for the composite.
\end{proof}
This is an intertwining argument. Related composition questions for interventional causal models are studied in greater generality by Rischel and Weichwald~\cite{MU:rischel}. Here the theorem guarantees forward answers only from states in the image of $T$ and for translated source operations. It does not guarantee coverage of all target states or all target operations. Surjectivity, reverse translation, and recovery of discarded source distinctions are separate requirements.

\subsection{The instrument in display coordinates}
For the original instrument quotient, let the source state be $(a,b)$ and set
\[
T(a,b)=(a+b,b).
\]
The target updates in~\eqref{MU:eq:display-coordinates} satisfy~\eqref{MU:eq:transport-exact}. The source decoder $a+b$ becomes the target's first coordinate. Thus a method can be carried between a register description and a display-centered description without relearning the meaning of every operation sequence.

\begin{figure}[htbp]
\centering\includegraphics[width=.96\linewidth]{sources/research/machine-understanding/figures/transport_square.pdf}
\caption{One verified operation square for the instrument. Translating and then applying the target update produces the same state as applying the source update and then translating. The equality concerns the declared $H$ operation, not arbitrary interventions.}\label{MU:fig:transport}
\end{figure}

\begin{MUcounterexample}[Correct current answers with a wrong transported operation]\label{MU:ex:badtransport}
Retain the correct map $T(a,b)=(a+b,b)$ and the correct display decoder. Now incorrectly define the target swap by $\widetilde W(s,b)=(s,b)$. The current display and the display immediately after a swap are still correct. Starting from $(a,b)=(1,3)$, however, the source sequence $W$ then $H$ ends with display zero. The incorrect target sequence ends with display two. Matching a restricted collection of outputs did not preserve the method.
\end{MUcounterexample}
The error becomes visible only when the altered internal state is used by a later operation. That is precisely why reusable understanding requires continuations rather than a list of isolated answers.

\subsection{Approximate transfer accumulates error}
Exact equalities are not always available. For this result let the target state space carry a metric $d$. Suppose each target update is $L$-Lipschitz for a common $L\geq0$, and suppose
\[
d(Tf_u(z),g_{\tau(u)}(Tz))\leq\epsilon
\]
uniformly on every state and operation that the contract permits. Let a target query be $M$-Lipschitz with real-valued output and satisfy $|q'(Tz)-q(z)|\leq\delta$.

\begin{MUproposition}[Finite-horizon transfer budget]\label{MU:prop:error}
For matched source and target operation sequences, start from $z_0,w_0$ with $d(Tz_0,w_0)\leq e_0$. After $n$ operations,
\begin{align*}
d(Tz_n,w_n)&\leq L^n e_0+\epsilon\sum_{j=0}^{n-1}L^j,\\
|q(z_n)-q'(w_n)|&\leq\delta+M\left(L^n e_0+\epsilon\sum_{j=0}^{n-1}L^j\right).
\end{align*}
The empty sum is zero and $L^0=1$, including when $L=0$. At $L=1$ the state bound is $e_0+n\epsilon$.
\end{MUproposition}
\begin{proof}
Writing $e_n=d(Tz_n,w_n)$, insert $g_{\tau(u)}(Tz_n)$ and use the triangle inequality and the Lipschitz bound to obtain $e_{n+1}\leq\epsilon+Le_n$. Iterating gives the stated expression. Insert $q'(Tz_n)$ between the two outputs and use the decoder mismatch and its Lipschitz constant.
\end{proof}
The bound concerns the matched operation sequence. It does not automatically cover two controllers that select different operations after slightly different observations. Adaptive choices require a matching policy contract or a separate stability argument for the controller.

\begin{figure}[htbp]
\centering\includegraphics[width=.94\linewidth]{sources/research/machine-understanding/figures/error_accumulation.pdf}
\caption{Exact values of the recurrence $e_{n+1}=Le_n+1/100$, with $e_0=0$, illustrate the transfer bound. The curves are analytical constructions, not measured model errors. A small local defect can remain bounded, accumulate linearly, or be amplified.}\label{MU:fig:error}
\end{figure}

A map can therefore be accurate for a short continuation and unsuitable for a longer one. The required tolerance belongs to the task, not to a universal declaration that the two representations are close. An approximate relation with one tolerance also need not be transitive. A chain of individually accepted translations may exceed the original error budget.

\subsection{Local representations form an atlas only with transition rules}
One representation may simplify prediction while another simplifies correction or proof. A useful system need not select one notation forever. It can maintain a family of representations with declared domains and translators. We call this an atlas only in the modest sense of a set of local descriptions connected by checked transition maps.

The domain guards matter. A method valid on reversible operations may fail after an irreversible reset. A chart switch may require evidence absent from the current chart. A common output name does not establish a common unit or query semantics. The system must know which transitions are supported and must retain the unresolved discrepancies when a transition is partial.

There is no general theorem here that more representations improve reasoning. Additional charts consume memory and create more opportunities for a mistaken translation. The research target is a family whose additional coverage or lower repeated cost justifies its construction and maintenance. Exact intertwining supplies one part of that analysis, not its economic conclusion.


\chapter{Explanation, inquiry, and correction}\label{MU:book:chapter-47}


\section{An explanation belongs to a recipient}\label{MU:sec:explanation}
An explanation can be a correct description and still fail to help a particular recipient. It may presume unfamiliar notation, omit a necessary background result, or encode all the difficult work in a phrase the recipient cannot execute. Conversely, a deliberately simplified explanation can be useful for a limited question family without being a complete description of the source mechanism.

We therefore distinguish two contracts. A \emph{subject explanation} enables the recipient to reason correctly about specified aspects of the subject. A \emph{process explanation} additionally makes a faithful claim about how a particular solver reached its result. A process explanation must satisfy an implementation attribution as well as a subject-level or transcript-level account. Neither fluent presentation nor a valid derivation alone establishes the attribution~\cite{MU:jacovi}.

\begin{MUdefinition}[Recipient-relative explanatory adequacy]\label{MU:def:explanation}
Fix a recipient procedure $L$, background information $B$, resource allowance, and continuation family. An artifact $E$ is adequate for a target model $m$ when $L(B,E)$ yields a representation and operations that answer the required continuations for $m$ within the declared error and resource limits. The artifact's inputs, the recipient's prior access, and any target-specific background are included in the contract.
\end{MUdefinition}
This is an operational property of an artifact and a recipient. It does not equate instruction with human comprehension, and it does not require natural language. A diagram, proof, worked example, executable transformation, or combination can play the role of $E$.

\subsection{An exact communication lower bound}
Let $\Theta$ be a finite family of target models. Relative to the required continuation questions, divide these models into $N$ classes with distinct answer profiles. Use one fixed recipient background and one fixed deterministic recipient. The teacher can select a message but cannot change the recipient's background separately for each target without that change being counted as communication.

\begin{MUproposition}[The information an exact explanation must carry]\label{MU:prop:explanation}
If a fixed-length $b$-bit artifact lets the recipient answer the declared questions exactly for every target in $\Theta$, then $2^b\geq N$. With target-dependent background $B_\theta$, the same argument gives
\[
2^b\geq\max_{B}\ N_B,
\]
where $N_B$ counts distinct answer profiles among targets sharing that exact background value.
\end{MUproposition}
\begin{proof}
Two different required answer profiles cannot produce the same background and message. The deterministic recipient would then make the same responses for both and fail on a distinguishing question. There are at most $2^b$ fixed-length messages. Apply this counting argument to the whole family, or separately within each background fibre.
\end{proof}
The bound is about information, not explanatory elegance. It does not establish that any shortest message is easy to understand or that the recipient can decode it efficiently. A compressed index into an enormous shared table can meet the bound while transferring little new method. The table's construction and availability must be charged to the background account.

\subsection{What the instrument explanation transfers}
For the original instrument, an adequate artifact can specify arithmetic modulo five, the representation $(s,b)$, the two update rules in~\eqref{MU:eq:display-coordinates}, the output decoder, and the domain restriction excluding $J$. A recipient supplied with an initial $(s,b)$ can then answer every $H/W$ continuation by applying these operations.

That establishes a subject explanation for the specified recipient. It does not establish that a neural system that produced the artifact used $(s,b)$ internally. A separate process explanation might show that the system's actual intermediate states realize these coordinates and respond to corresponding interventions. These two certificates can support one another, but they answer different questions.

An especially useful explanation includes the failure boundary. The line ``the third register is irrelevant'' is too broad. The line ``the third register does not affect any display reachable through the current $H/W$ operation family'' is exact. The latter statement tells the recipient what should be reconsidered when a new operation arrives.

\subsection{Explanations can preserve methods without reproducing discovery}
A proof may give a shorter route than the route by which the result was found. A teacher may introduce coordinates that the discoverer never used. These are legitimate explanatory choices when their subject-level claims are correct. They become misleading only when presented as faithful records of the original reasoning without evidence.

This distinction avoids an unnecessary demand that every helpful explanation be an autobiography of a computation. It also prevents the opposite mistake of treating a helpful reconstruction as proof of the original mechanism. Human context and communicative purpose remain important beyond the formal recipient criterion~\cite{MU:miller}.

\section{One explanation for recipients with different backgrounds}\label{MU:sec:audience}
An explanation can be adequate for each recipient separately without there being one short artifact that works for all of them. Conversely, a shared artifact can be shorter than a complete description because recipients already hold complementary information. The relevant object is the combination of sender access, communicated record, and recipient background. Counting only the words in the artifact misses two thirds of that object.

The finite problem below belongs to zero-error communication with side information. Witsenhausen related side-information coding to graph coloring~\cite{MU:witsenhausen}, while broadcasting research studies a common message used by several informed receivers~\cite{MU:alon}. We adapt that established construction to explicit sender access and recipient-specific demands. The result is a communication criterion within an explanation contract, not a new definition of comprehension.

\subsection{An accessible sender and a family of recipients}
Let $\MUM$ be a finite nonempty set of possible cases. The sender observes only $\alpha(m)$, where $\alpha:\MUM\to S$ is a declared access map. Recipient $i\in\{1,\ldots,k\}$ receives background $\beta_i(m)$ and must recover $d_i(m)$. A demand can be a single answer, a reusable state, or the complete finite answer profile for a declared task. The maps and recipients are fixed before the case is selected.

One encoder $E:S\to W$ produces a common artifact. A decoder $D_i$ for each recipient must satisfy
\begin{equation}\label{MU:eq:audience-contract}
D_i(E(\alpha(m)),\beta_i(m))=d_i(m)\qquad\text{for every }m,i.
\end{equation}
The encoder is not allowed to consult a recipient-specific hidden state. Recipients do not receive an uncharged transcript or a separate description of the selected case. Here decoders may be designed once for the declared finite contract and then remain fixed. If an existing interpreter is prescribed, the criterion below is only a necessary information condition unless that interpreter can implement the constructed decoder.

Build a graph $G_\alpha$ on $\alpha(\MUM)$. Join $s$ and $t$ when some cases $m,n$ and recipient $i$ satisfy
\[
\alpha(m)=s,\quad\alpha(n)=t,\quad
\beta_i(m)=\beta_i(n),\quad d_i(m)\ne d_i(n).
\]
An edge can be a loop when $s=t$. It then records a distinction that neither the sender's access nor that recipient's background can resolve.

\begin{MUtheorem}[A common explanatory artifact under sender and recipient constraints]\label{MU:thm:audience}
An exact artifact satisfying~\eqref{MU:eq:audience-contract} with at most $K\geq1$ symbols exists if and only if $G_\alpha$ has no loop and admits a proper coloring with at most $K$ colors. In the loop-free case, the minimum alphabet size is $\chi(G_\alpha)$, and the minimum fixed binary length with designable decoders is $\lceil\log_2\chi(G_\alpha)\rceil$. A loop makes the exact contract infeasible for every alphabet size.
\end{MUtheorem}
\begin{proof}
An encoder assigns a symbol to every accessible sender value. If adjacent vertices receive the same symbol, their witnessing recipient receives the same artifact and background in two cases requiring different outputs. No deterministic decoder can be correct in both. Thus every valid encoder is a proper coloring, and a loop is impossible.

Conversely, use a proper coloring as the encoder. For any recipient and any artifact/background pair that actually occurs, assign the demanded output from one compatible case. Another compatible case cannot demand a different output, since it would witness an edge whose endpoints have the same color. Thus the decoder is well defined. Define it arbitrarily on unattainable pairs. Counting binary strings proves the final bound.
\end{proof}

The theorem provides a smallest \emph{alphabet}, not the cheapest explanation. Its decoders can be large tables. Their construction, runtime, background acquisition, and semantic interpretation remain part of the resource account. The graph is useful precisely because it does not pretend that a color label is intrinsically meaningful.

\subsection{Complementary knowledge can reduce communication}
For two bits $(a,b)$, let the first recipient know $a$ and need $b$, while the second knows $b$ and needs $a$. A sender with access only to $a\oplus b$ can broadcast that one bit. Both recipients recover their missing bit. The sender need not reconstruct the complete pair. An inference from successful explanation to possession of the recipient's full state would therefore be invalid.

This does not make the recipients' knowledge free. Their prior bits are explicitly supplied. Removing them changes the graph and the communication requirement. Nor does the example show that the sender has a useful general theory of the two-bit system. It establishes a specific jointly realizable communication contract.

\begin{MUcounterexample}[Individually binary explanations can require three common symbols]\label{MU:ex:audience-cycle}
Take five cases labeled $0,\ldots,4$ and five recipients. Recipient $i$ cannot distinguish cases $i$ and $i+1$ modulo five from its background, but distinguishes every other case individually. It demands zero on case $i$ and one on case $i+1$, with arbitrary fixed demands elsewhere. The sender knows the case. Each recipient separately needs only two artifact symbols. The common confusability graph is a five-cycle. It is not two-colorable, but the sequence of colors $A,B,A,B,C$ is proper. One common artifact therefore needs three symbols, or two fixed binary bits.
\end{MUcounterexample}

\begin{figure}[htbp]
\centering\includegraphics[width=.84\linewidth]{sources/research/machine-understanding/figures/audience_cycle.pdf}
\caption{Each edge records one recipient's unresolved pair of cases. Each recipient alone can use a binary message. Their combined graph needs three symbols. Letters denote message classes, not empirical clusters or intrinsic semantic categories.}\label{MU:fig:audience}
\end{figure}

The obstruction is global: the largest ambiguity faced by any one recipient does not determine the size of a shared explanation. The same issue appears when one artifact is expected to support a novice, a specialist, and a verifier who each possess different background information. The finite theorem tells us what must be distinguished. It does not prescribe how a human-facing explanation should express those distinctions.

\subsection{Source restrictions are substantive}
If two cases have the same sender observation and the same recipient background but require different outputs, adding more fluent text cannot fix the information deficit. The graph has a loop. The options are to acquire evidence, change the recipient background, relax the demand, or report that the current contract cannot be met. Calling a more powerful decoder changes none of those facts if its permitted inputs remain identical.


\section{Learning where the current representation fails}\label{MU:sec:revision}
A reusable representation should support correction without treating every new observation as a reason to discard everything. A valid method can have a restricted domain. A recorded premise can be false. A previously omitted variable can become relevant. These failures call for different revisions.

We distinguish three cases. In a \emph{state correction}, the operations remain valid but the current state estimate changes. In a \emph{model correction}, an assumed transition or observation law changes. In a \emph{contract expansion}, the requested operations or questions change. A representation can be entirely adequate before the third kind of change and inadequate afterward, as the instrument shows.

\subsection{A finite model of informative correction}
Let $\MUM$ be a finite nonempty set of deterministic candidate models. Let $\mathcal E$ be a finite set of permitted experiments. Each experiment $e$ has an outcome map $f_e:\MUM\to\mathcal O_e$. The current information is a candidate set $V\subseteq\MUM$. On observing $o$ from $e$, update it by
\begin{equation}\label{MU:eq:filter}
C_{e,o}(V)=\{m\in V:f_e(m)=o\}.
\end{equation}
This is consistency filtering, not Bayesian inference and not a general model of scientific reasoning. An empty set is meaningful: the observation contradicts every currently retained candidate. It must not be replaced by an invented best answer.

An experiment is informative for $V$ when its outcome map is not constant on $V$. Models indistinguishable under every permitted experiment form an observational equivalence class. The aim can be to identify such a class rather than a uniquely named model.

\begin{MUproposition}[Finite correction under separating evidence]\label{MU:prop:queries}
Suppose the true model $m_*$ belongs to the initial set $V_0$, every observation is exact, and the allowed experiment family separates any two distinct remaining observational classes. Repeatedly selecting an informative experiment preserves $m_*$ and reaches a single observational class after at most $|V_0|-1$ informative observations. If there are $N$ distinguishable target classes, every adaptive procedure identifying all of them using experiments with at most $b\geq2$ outcomes and at most $h$ observations must satisfy $b^h\geq N$.
\end{MUproposition}
\begin{proof}
The true model predicts each exact observation, so filtering never removes it. An informative experiment splits the current set into at least two nonempty outcome classes. Whatever true outcome occurs, at least one candidate is removed. The nonempty set can strictly shrink at most $|V_0|-1$ times. If two observational classes remain, the separating assumption supplies another informative experiment. For the second claim, an adaptive decision tree of depth at most $h$ and branching at most $b$ has at most $b^h$ leaves. Distinct required classes must reach different leaves.
\end{proof}
This elementary result shares the query-and-counterexample perspective of classical active learning~\cite{MU:angluin}. It does not give an efficient method for discovering the candidate family, finding an inexpensive separating experiment, or dealing with inaccurate observations. Its upper bound is a count of informative observations, not a bound on wall-clock work or physical cost.

\subsection{A question can be more useful than another answer}
If all surviving candidates give the same answer to the next task, the system can act without identifying which candidate is true. If they disagree on a consequential operation, an experiment that distinguishes that disagreement may be more valuable than repeating familiar tasks on which all candidates agree.

The criterion is task-relative. Two models can differ on an inaccessible quantity and still support all permitted decisions. Conversely, a rarely asked question can distinguish models whose disagreement controls whether a future intervention is safe. The system should retain which distinctions have been settled and which have merely been irrelevant so far.

A request for evidence is also an operation with a realization. It may consume resources, change the system, or reveal only a biased part of the state. The finite filtering model presumes that the stated outcome maps already include those effects. Section~\ref{MU:sec:acquisition} makes current-state changes and their cost explicit. For a changing instrument, one must use a controlled reset, a model of the full observation history, or a state update that tracks what the experiment changed. Independent repeated access to the same initial state is not free.

\subsection{Repairing the running example}
Before $J$ is admitted, all five possible values of $c$ are equivalent under the base contract. Treating them as one class is correct, not ignorance that the old task required the predictor to resolve. After $J$ is admitted, these classes split. The system can report that its old state no longer determines the answer.

The current record supplies $(s,b)$. Executing $J$ and observing $s_J$ determines $c=s_J-s$, while~\eqref{MU:eq:recover-after-J} updates the affected first register. The corrected extended coordinates are $(s_J,b,c)$. Their operations are
\[
\bar H(s,b,c)=(s+b,b,c),\quad
\bar W(s,b,c)=(s,s-b,c),\quad
\bar J(s,b,c)=(s+c,b,c).
\]
The base method survives as the restriction of the enlarged method. Only the newly relevant distinction is acquired.

A different kind of failure would occur if the instrument secretly implemented $H(a,b,c)=(a+2b,b,c)$. That would contradict the old transition model rather than merely extend the operation catalogue. A correction record should say which assumption changed. Otherwise a sequence of patched answers can conceal a permanently incorrect account of the mechanism.

\subsection{Revision must preserve its own uncertainty}
The true-model assumption in Proposition~\ref{MU:prop:queries} is a substantive premise. When it is not justified, the system must allow the possibility that every candidate is wrong. When observations are noisy, deleting a model after one conflicting record can remove the best explanation. The statistical treatment in Sections~\ref{MU:sec:bayesian}--\ref{MU:sec:sequential} supplies exact and conditional results once a likelihood model and selection process have been declared. It does not establish their empirical calibration.

The finite framework nevertheless gives a useful discipline. A correction should identify the new evidence, the rejected candidates, the surviving claims, and the assumption linking the observation to the model. It should not rewrite an old uncertainty as though the corrected answer had always been established.

\section{Acquisition as an experiment planning problem}\label{MU:sec:acquisition}
A lower bound on the information a representation must hold is not yet a plan for acquiring it. Even a perfectly informative experiment can be unavailable, expensive, or destructive. We now make that acquisition problem exact for a finite deterministic model.

Belief-state planning has a substantial foundation in partially observed control~\cite{MU:smallwood}. The construction here uses sets of possible current states rather than probability distributions, a worst-case cost rather than expected reward, and an explicitly bounded horizon. It is a finite dynamic program for a stated contract, not a general solution to scientific discovery.

\subsection{Experiments change both information and state}
Let $X$ be a finite nonempty state space. Each experiment $e$ has an admissibility set $A_e\subseteq X$, a transition $T_e:X\to X$, an observation $o_e:X\to O_e$, and a known cost $c_e\geq0$. The observation is indexed by the state before the experiment, but can describe a reading taken afterward: for example $o_e=q\circ T_e$. At an information state $V\subseteq X$, the experiment is permitted only when $V\subseteq A_e$. After outcome $o$, the new support is
\begin{equation}\label{MU:eq:acquisition-update}
V_{e,o}=\{T_e(x):x\in V,\ o_e(x)=o\}.
\end{equation}
Only nonempty branches are possible. This update tracks the world after the intervention, not just a narrowed list of old states.

Let $r:X\to Z$ be the target record. A support is terminal when $r$ is constant on it. Let $h$ be the number of experiments still available. Write $\mathcal V_h(V)$ for the least worst-case sum of experiment costs among policies that terminate with the target known within $h$ steps. Infeasibility is denoted $+\infty$. Then define
\begin{equation}\label{MU:eq:acquisition-bellman}
\begin{aligned}
\mathcal V_h(V)&=0 &&\text{if }|r(V)|=1,\\
\mathcal V_0(V)&=+\infty &&\text{otherwise},\\
\mathcal V_h(V)&=\min_{e:V\subseteq A_e}
\left(c_e+\max_{o:V_{e,o}\ne\varnothing}\mathcal V_{h-1}(V_{e,o})\right)
&&\substack{\text{otherwise,}\\h>0.}
\end{aligned}
\end{equation}
An empty minimum is $+\infty$. The terminal case permits stopping immediately, so an unnecessary additional experiment is never forced.

\begin{MUtheorem}[Exact finite acquisition with state-changing observations]\label{MU:thm:acquisition}
For every nonempty support $V$ and finite horizon $h$, recurrence~\eqref{MU:eq:acquisition-bellman} equals the minimum worst-case acquisition cost among deterministic adaptive policies with that horizon, access boundary, and current-state target. Every finite value has an attaining policy obtained by choosing a minimizing experiment and its recursively attaining branches. The value is infinite exactly when no such policy exists.
\end{MUtheorem}
\begin{proof}
At horizon zero, a target can be returned without further evidence precisely when it is constant over the current support. This gives the base case. At a nonterminal support with $h>0$, the first experiment must be admissible at every possible state. After its observed outcome the possible current states are exactly~\eqref{MU:eq:acquisition-update}. By induction, the best possible remaining worst-case cost on that branch is $\mathcal V_{h-1}$. Every nonempty branch can occur, giving the maximum over branches. Minimizing over the finite experiment set gives a lower bound for every adaptive policy. Choosing a minimizing experiment and attaining branch policies achieves that bound whenever it is finite. If every choice has an impossible branch, no guaranteed policy exists.
\end{proof}

The complete information state must include every variable needed for future transitions, admissibility, costs, and the target. A permission depending on an earlier event must be retained as state. Otherwise two histories with the same physical support could admit different experiments, invalidating the recurrence. Nonnegative costs and the finite horizon make zero-cost actions harmless to the definition. No unbounded termination claim is being made.

The state search can be large: there are up to $2^{|X|}-1$ nonempty supports, before accounting for additional memory variables. The theorem is an exact finite characterization, not a polynomial-time learning guarantee. The delivered reference implementation memoizes supports and remaining horizons. Its companion check compares the result with a separately implemented enumeration of complete observation-history policies on small systems.

\subsection{The diagnostic operation is a paid transformation}
In the instrument, take a known base record $(s,b)$ and let the five possible values of $c$ define the initial support. Set the cost of $J$ to two units, and the cost of $H$ and $W$ to one unit each. Each experiment returns the exact display after its operation. The target is the enlarged current record $(s,b,c)$.

A $J$ observation resolves all five possibilities and costs two. An $H/W$ sequence cannot reveal $c$, so no policy costing less than two can meet the target. The optimum is therefore two, attained by executing $J$ and retaining the post-operation record. This example fixes its sensor and operation catalogue. It does not claim that every physical diagnostic with five outcomes has the same cost.

\begin{figure}[htbp]
\centering\includegraphics[width=.98\linewidth]{sources/research/machine-understanding/figures/acquisition_branches.pdf}
\caption{A single $J$ diagnostic from base record $(s,b)=(4,3)$ has five possible exact reports. Each report determines a different \emph{post-operation} record. The cost is charged to the operation, and the hidden value is acquired rather than supplied to the decoder.}\label{MU:fig:acquisition}
\end{figure}

\subsection{Knowing the current state is not recovering its history}
\begin{MUcounterexample}[A reset supplies certainty without recovering the past]\label{MU:ex:reset-knowledge}
Let the current state be an unknown bit. The only experiment sets the bit to zero and returns a constant symbol. One experiment makes the current bit known with certainty. It supplies no information about whether the old bit was zero or one. Thus an acquisition policy for the current-state target can succeed by changing the world, not by identifying its original condition.
\end{MUcounterexample}

To require recovery of a past distinction, replace the state by $(x_0,x_t)$, let every transition preserve $x_0$, and make the terminal target include the required part of $x_0$. The same recurrence then treats historical reconstruction honestly. In the reset example, the two origins remain indistinguishable at every horizon. If both current knowledge and preserved history matter, use their joint target rather than claim that either implies the other.

For the instrument, the diagnostic $J$ is invertible once $c$ is known, so the exact observation can support both records. That property comes from the specified operation. It cannot be transferred to an arbitrary experiment. Likewise, a planner that chooses an informative experiment is not thereby credited with discovering its transition law or building its sensor.

\subsection{Experiment value is relative to the intended continuation}
An experiment need not distinguish every candidate model to be useful. It need only resolve distinctions that control the promised future behavior. Conversely, acquiring a value that changes none of the relevant decisions can waste resources even when it reduces uncertainty. Information-based experimental design has a long history~\cite{MU:lindley}; the present cost criterion deliberately asks a different, exact question about a specified target and worst-case completion.

The horizon is a contractual limit, not an arbitrary truncation of a failed computation. A negative result at horizon $h$ says that the declared task cannot be guaranteed within that many permitted experiments. Increasing the horizon changes the problem. It does not retroactively validate a promise made at the smaller horizon.


\section{When translation commutes with correction}\label{MU:sec:correctiontransport}
A method can travel through a language that preserves its current answers while losing the evidence needed to improve it. This is a different obstruction from a wrong update equation. It concerns the interaction between two otherwise sensible operations: translating a model and filtering it against evidence.

Let $\MUM$ and $\mathcal N$ be finite model families, with model translation $\phi:\MUM\to\mathcal N$. For one source experiment $e$, let $f:\MUM\to\mathcal O$ give its outcomes. A corresponding target experiment has outcome map $g:\mathcal N\to\mathcal O'$. Let $\eta:\mathcal O\to\mathcal O'$ translate evidence. Assume the basic observation square commutes:
\begin{equation}\label{MU:eq:evidence-square}
g(\phi(m))=\eta(f(m))\quad\text{for all }m\in\MUM.
\end{equation}
This condition says that a model's predicted evidence translates correctly. It permits $\eta$ to merge source outcomes.

Write $\phi[V]=\{\phi(m):m\in V\}$. Let $C_o(V)=\{m\in V:f(m)=o\}$ and $C'_{o'}(W)=\{n\in W:g(n)=o'\}$. Both are ordinary consistency filters.

\begin{MUtheorem}[Correction transport and evidence reflection]\label{MU:thm:correction}
Under~\eqref{MU:eq:evidence-square}, for every $V\subseteq\MUM$ and every attainable source outcome $o\in f(\MUM)$,
\begin{equation}\label{MU:eq:correction-inclusion}
\phi[C_o(V)]\subseteq C'_{\eta(o)}(\phi[V]).
\end{equation}
Equality holds for all such $V$ and $o$ if and only if $\eta$ is injective on the attainable outcome set $f(\MUM)$. The claim includes candidate sets contradicted by the observation, so that preservation of an empty update is part of the contract.
\end{MUtheorem}
\begin{proof}
If $n\in\phi[C_o(V)]$, then $n=\phi(m)$ for some $m\in V$ with $f(m)=o$. Equation~\eqref{MU:eq:evidence-square} gives $g(n)=\eta(o)$, proving inclusion.

If $\eta$ is injective on $f(\MUM)$ and $n\in C'_{\eta(o)}(\phi[V])$, choose $m\in V$ with $\phi(m)=n$. Then $\eta(f(m))=g(n)=\eta(o)$, so $f(m)=o$. Thus $n\in\phi[C_o(V)]$, proving equality.

Conversely, suppose two distinct attainable outcomes $a,b$ have $\eta(a)=\eta(b)$. Choose $m$ with $f(m)=b$, let $V=\{m\}$, and consider source observation $a$. The source update is empty, whereas the target update retains $\phi(m)$ by~\eqref{MU:eq:evidence-square}. Uniform equality therefore fails.
\end{proof}
This is an exact set-theoretic result about direct images and filtering. The institutional or cognitive interpretation does not make it a new general theorem of set theory. Its value is to expose a condition that output translation alone omits: the translated evidence must preserve the distinctions by which a candidate is excluded.

The necessity statement deliberately tests observations that contradict the current candidate set. If those cases are excluded by a standing assumption that the current model is already correct, weaker conditions can suffice. That is a different and less demanding correction contract. A reasoner meant to discover that its model is wrong should not assume away the contradictory case.

\begin{MUcounterexample}[A translated observation preserves both mistaken candidates]\label{MU:ex:correction}
Take $\MUM=\{m_0,m_1\}$ with $f(m_0)=0$ and $f(m_1)=1$. Translate models bijectively to $\{n_0,n_1\}$, but translate both observations to a single symbol $*$. Let $g(n_0)=g(n_1)=*$. The observation square commutes. Starting from both candidates, source observation $1$ leaves only $m_1$, whose image is $\{n_1\}$. Translating first and then filtering by $*$ leaves both $n_0$ and $n_1$.
\end{MUcounterexample}

\begin{figure}[htbp]
\centering\includegraphics[width=.98\linewidth]{sources/research/machine-understanding/figures/correction_square.pdf}
\caption{Correct translation of each predicted observation need not preserve correction. The lower evidence channel merges zero and one, so the target cannot remove the candidate excluded at the source. The theorem specifies when the inclusion becomes equality.}\label{MU:fig:correction}
\end{figure}

The failure is not an inaccurate target answer in its coarse language. Both models really do predict $*$ there. The failure is the loss of a previously meaningful ability to discriminate. This is why a representation that can communicate conclusions may still be an inadequate carrier of scientific revision.

\subsection{The instrument explains the boundary}
For base experiments built only from $H$, $W$, and the display, the translation from $(a,b,c)$ to $(a,b)$ can preserve evidence without collapsing outcome symbols. Different values of $c$ produce the same permitted observations, so their identification loses no base correction power.

For $J$, that translation cannot preserve the full outcome map at all: two states with identical $(a,b)$ can have different readings after $J$. One could force a commuting square by translating all new readings into one coarse symbol. That would preserve the trivial coarse prediction, but it would destroy the ability to discover $c$. Enlarging the representation is a different response from weakening the question until the old representation appears adequate.

This distinction is important for evaluation. A system should not be rewarded for transfer by silently replacing a difficult target question with an easier translated one. The source and target question contracts, including evidence reflection when correction matters, must be fixed before the comparison.

\subsection{Composing corrective explanations}
Suppose an explanatory artifact carries a recipient's candidate set, translated experiments, and evidence interpretation. If each translation preserves the observation square and reflects the relevant outcomes, successive filters and translations commute by repeated application of Theorem~\ref{MU:thm:correction}. The recipient can then discard the same translated candidates as the source after the specified observation sequence.

This conclusion is only about filtering the represented model families. It does not transfer an ability to invent a missing model, establish that the observations are trustworthy, or guarantee efficient implementation. An explanation can teach how to correct within a family while leaving the enlargement of that family as a separate problem.

\section{Noisy evidence and Bayesian correction}\label{MU:sec:bayesian}
Exact filtering removes a candidate as soon as an observation contradicts it. That is appropriate for the deterministic evidence contract in Section~\ref{MU:sec:correctiontransport}. A noisy observation can instead be unlikely under a model without being impossible. Reusing the exact filter in that setting can discard a well-supported model because a sensor behaved exactly as its noise model allowed.

Let $M\in\MUM$ be a fixed unknown model in a finite family, with prior $p(m)$. An experiment produces $O\in\mathcal O$ with likelihood $L_m(o)$. A common reporting channel $B(o,z)$ produces $Z\in\mathcal Z$. Each likelihood row and each channel row is a probability distribution. The complete joint model is
\begin{equation}\label{MU:eq:bayes-joint}
\Pr(M=m,O=o,Z=z)=p(m)L_m(o)B(o,z).
\end{equation}
In particular, the channel does not secretly inspect $M$. Define
\[
Q_m(z)=\sum_o L_m(o)B(o,z),\qquad
p^o(m)=\frac{p(m)L_m(o)}{\sum_n p(n)L_n(o)},\qquad
p^z(m)=\frac{p(m)Q_m(z)}{\sum_n p(n)Q_n(z)}.
\]
A posterior is defined only when its denominator is positive. An event assigned probability zero by the entire current model has no posterior under that model. It is a model-failure signal, not permission to invent a probability distribution.

This section retains the same model identities at source and target so that it isolates evidence loss. Merging model identities is a separate operation. A target likelihood on a merged class is independent of the allocation of prior mass within that class only when the relevant likelihood is constant there, or when the within-class prior itself is retained as additional state. That assumption must not be concealed in a notation change.

\subsection{When a coarser report preserves every update}
The reporting model implies $M\to O\to Z$. Therefore, on every positive-probability pair $(o,z)$, the full posterior after observing both records equals $p^o$. The coarse posterior is a mixture of these fine posteriors:
\begin{equation}\label{MU:eq:posterior-mixture}
p^z=\sum_o\Pr(O=o\mid Z=z)\,p^o.
\end{equation}
Averaging is not generally the same as retaining each possible update. The distinction is classical statistical sufficiency, whose factorization foundations include Halmos and Savage~\cite{MU:halmos}. We state and prove the finite channel criterion used here.

\begin{MUtheorem}[Exact Bayesian correction through a reporting channel]\label{MU:thm:bayes}
For fixed finite likelihoods $L$ and reporting channel $B$, the following are equivalent:
\begin{enumerate}[label=(\roman*)]
\item For every prior $p$, $p^o=p^z$ on every positive-probability pair $(o,z)$.
\item This equality holds for one prior assigning positive mass to every model.
\item There exists a row-stochastic kernel $\rho(o\mid z)$ such that, for all $m,o,z$,
\begin{equation}\label{MU:eq:bayes-factorization}
L_m(o)B(o,z)=Q_m(z)\rho(o\mid z).
\end{equation}
\end{enumerate}
Thus a report preserves every posterior update exactly when the discarded observation can be conditionally reconstructed from that report without knowing the model, in the actual joint experiment.
\end{MUtheorem}
\begin{proof}
The first condition implies the second. Suppose the second holds under a fully supported prior. For each $z$ of positive probability set $\rho(o\mid z)=\Pr(O=o\mid Z=z)$. If $(o,z)$ has positive probability, posterior equality and Bayes' formula give
\[
p(m)L_m(o)B(o,z)
=\Pr(O=o,Z=z)p^o(m)
=\rho(o\mid z)\Pr(Z=z)p^z(m)
=p(m)Q_m(z)\rho(o\mid z).
\]
Cancel the positive $p(m)$. If the pair has probability zero, full prior support makes $L_m(o)B(o,z)=0$ for every model, and $\rho(o\mid z)=0$. If $z$ is impossible under every model, choose any probability row $\rho(\cdot\mid z)$ since both sides are zero. This proves the third condition.

Conversely, the factorization makes the likelihood of $(o,z)$ a model-independent multiple of the likelihood of $z$. Normalization under any prior with positive pair probability therefore gives $p(M\mid o,z)=p^z$. The common reporting channel gives $p(M\mid o,z)=p^o$. This proves the first condition.
\end{proof}

For deterministic reporting $z=\eta(o)$, the criterion says that the nonzero vectors $(L_m(o))_{m\in\MUM}$ within each report fibre must be proportional. An ancillary symbol may be discarded even when the report is not injective. In the deterministic evidence model, distinct attainable outcomes have nonzero likelihood vectors with disjoint supports. Those vectors cannot be proportional, recovering the injectivity requirement for that special case. Noise therefore changes the correct preservation test, not merely the notation used for filtering.

The kernel $\rho$ is conditional reconstruction for the specified coupling $(M,O,Z)$. The existence of some unrelated procedure simulating an experiment from another record is a different assertion. Neither assertion shows that the real sensor has the stated likelihood or that a learner inferred it correctly.

\subsection{An exact measure of the lost update}
A coarser report may fail the criterion while still retaining much of the useful evidence. There is a natural average loss, though it should not be mistaken for a worst-case bound on every possible observation. For finite distributions use natural logarithms and the convention $0\log 0=0$.

\begin{MUproposition}[Expected posterior loss is conditional information]\label{MU:prop:posterior-information}
Under~\eqref{MU:eq:bayes-joint},
\begin{equation}\label{MU:eq:posterior-kl}
\mathbb E_{O,Z}\!\left[D_{\mathrm{KL}}(p^O\Vert p^Z)\right]
=I(M;O\mid Z)=I(M;O)-I(M;Z).
\end{equation}
The common value is nonnegative. It is zero exactly when the two posteriors agree on every positive-probability observation pair for that prior. For a fully supported prior, zero is equivalent to the factorization in Theorem~\ref{MU:thm:bayes}.
\end{MUproposition}
\begin{proof}
On the support of the joint law, $p(m\mid o,z)=p^o(m)$. Substituting this into the finite definition of conditional mutual information gives the expected divergence. Expanding its logarithm as $\log(p^o(m)/p(m))-\log(p^z(m)/p(m))$ and summing gives the difference of mutual informations. Terms with zero prior mass carry no weight. Every positive-weight numerator has a positive denominator, so the expression is finite.

For probability distributions, $D_{\mathrm{KL}}(a\Vert b)\geq0$, with equality only at $a=b$. This follows from $\log t\leq t-1$ on the common support, with missing support only increasing the divergence. A finite weighted sum of nonnegative terms is zero exactly when every positive-weight term is zero. The final assertion follows from Theorem~\ref{MU:thm:bayes}.
\end{proof}

This is the standard information chain rule in the present correction contract, not a new entropy identity. It makes the object of the lost information explicit: the update over the chosen model family. That family may include irrelevant distinctions or omit the true mechanism. A small or zero value does not independently validate the family.

\subsection{The noisy instrument and a harmless discarded flag}
Return to the instrument immediately after a known $J$ intervention. The earlier $(s,b)$ is known, while $c$ is initially uniform on five values. A read-only sensor reports the shifted value $Y=\text{reported display}-s$ modulo five. It returns the correct $c$ with probability $3/4$ and each wrong value with probability $1/16$. The physical state remains fixed during repeated read-only observations.

Suppose the raw report also contains a binary flag with probabilities $1/3$ and $2/3$, independent of $c$ and of $Y$. Removing that flag satisfies~\eqref{MU:eq:bayes-factorization} exactly. Both recipients retain the same posterior over $c$ and hence the same distribution over the post-$J$ state $(s+c-b,b,c)$. The flag does not carry corrective information for this family.

\begin{MUcounterexample}[A compact report can destroy the useful noisy distinction]\label{MU:ex:noisy-parity}
In the same instrument, replace $Y\in\{0,1,2,3,4\}$ by its ordinary integer parity, not a field operation. On observing the full symbol $Y=2$, the posterior puts $3/4$ on $c=2$ and $1/16$ on each alternative. On observing only even parity, it puts $7/24$ on each even candidate and $1/16$ on each odd candidate. The report is correct as a parity statement but fails to preserve the posterior needed to distinguish the three even values.
\end{MUcounterexample}

For $n$ conditionally independent identical reports $Y=2$, the posterior of $c=2$ is
\[
\frac{12^n}{12^n+4}.
\]
For $n$ reports retained only as even parity, the posterior of that same candidate is
\[
\frac{14^n}{3\cdot14^n+2\cdot3^n}.
\]
At $n=0$ both equal the prior $1/5$. The full-report value at $n=3$ is $432/433$, whereas the parity value cannot exceed $1/3$. Repeatedly receiving the same coarse fact does not restore distinctions the reporting language never records.

\begin{figure}[htbp]
\centering\includegraphics[width=.96\linewidth]{sources/research/machine-understanding/figures/noisy_correction.pdf}
\caption{Posterior mass on $c=2$ along the specified repeated-report histories. Removing the ancillary flag gives exactly the full-report curve. Retaining only parity leaves three even candidates indistinguishable. The curves are analytic conditional calculations, not measured learning trajectories or probabilities of obtaining these histories.}\label{MU:fig:noisy}
\end{figure}


\section{Approximation, rare evidence, and decision stability}\label{MU:sec:posterior-errors}
A small error in an unconditional model does not automatically give a small posterior error after an unusual observation. Conditioning divides by the probability of the evidence. A useful correction guarantee must therefore record both an approximation error and the evidence probability at which that guarantee is used.

\begin{MUproposition}[A conditional stability bound with an evidence denominator]\label{MU:prop:conditioning}
Let $P$ and $Q$ be finite joint laws of a model and an observation, with $\MUTV(P,Q)\leq\delta$. Let $E$ be an observation event with $a=P(E)>0$ and $b=Q(E)>0$. Then
\begin{equation}\label{MU:eq:conditional-bound}
\MUTV\bigl(P(M\mid E),Q(M\mid E)\bigr)
\leq\min\left\{1,\frac{\delta}{\max(a,b)}\right\}.
\end{equation}
In particular, a justified lower bound $P(E)\geq\gamma>0$ supplies the weaker usable bound $\min\{1,\delta/\gamma\}$.
\end{MUproposition}
\begin{proof}
Write $u_m=P(M=m,E)$, $v_m=Q(M=m,E)$, and let $p=u/a$, $q=v/b$. Exchanging $P$ and $Q$ if necessary, assume $a\leq b$. The variation inside $E$ is at least $\tfrac12\sum_m|u_m-v_m|$. The variation outside $E$ is at least $|a-b|/2$, since its total masses differ by that amount. Hence
\[
\delta\geq \tfrac12\sum_m|u_m-v_m|+\tfrac12(b-a)
=b-\sum_m\min(u_m,v_m).
\]
Since $a\leq b$, $\min(ap_m,bq_m)\leq b\min(p_m,q_m)$. Substitution gives $\delta\geq b\,\MUTV(p,q)$. Division by $b=\max(a,b)$ and the universal upper bound one prove the claim.
\end{proof}

The bound is sharp for general joint laws: put the same mass $a$ on $E$ in both laws, choose conditional laws at distance $t$, and make their complements identical. Then the joint distance is $at$. The result neither estimates $\delta$ from data nor proves that the evidence model is calibrated. Those are separate inputs.

\begin{MUcounterexample}[Rare evidence can reverse a posterior despite small joint error]\label{MU:ex:rare-evidence}
For $0<\epsilon<1/2$, take a binary model with equal prior probabilities. Under $P$, let $P(M=0,E)=\epsilon$ and $P(M=1,E)=0$. Under $Q$, interchange these two masses. Choose the complement masses so each model retains prior probability $1/2$. Both laws give $E$ probability $\epsilon$, their joint total variation is $2\epsilon$, and their posteriors on $E$ are opposite point masses. Joint error can tend to zero while conditional error remains one.
\end{MUcounterexample}

\begin{figure}[htbp]
\centering\includegraphics[width=.95\linewidth]{sources/research/machine-understanding/figures/conditioning_bound.pdf}
\caption{The certified upper bound $\min(1,\delta/\gamma)$ for $\delta=1/100$. A small joint-law error does not give a useful guarantee when the evidence probability is comparably small. The plot is a bound, not an empirical estimate of posterior error.}\label{MU:fig:conditioning}
\end{figure}

\subsection{From approximate beliefs to a recipient's decisions}
Not every posterior distinction matters equally to an action. To link correction with use, fix one finite nonempty action set $A$ and losses $\ell(m,a)\in[0,L_{\max}]$. A receiver evaluates actions from the distribution its actual interpreter reconstructs. The target loss function and available actions must be the same in the comparison.

\begin{MUproposition}[A joint posterior and recipient decision certificate]\label{MU:prop:decision}
Let $p$ be the reference posterior, $\widehat p$ an approximate or transmitted posterior, and $r_i$ recipient $i$'s reconstructed posterior. Suppose
\[
\MUTV(p,\widehat p)\leq\epsilon,\qquad
\MUTV(\widehat p,r_i)\leq\eta_i.
\]
If $a_i$ minimizes expected loss under $r_i$, then, with $e_i=\min\{1,\epsilon+\eta_i\}$,
\begin{equation}\label{MU:eq:recipient-regret}
\mathbb E_p\ell(M,a_i)-\min_{a\in A}\mathbb E_p\ell(M,a)
\leq\min\{L_{\max},2L_{\max}e_i\}.
\end{equation}
If the expected-loss gap between $a_i$ and every other action under $r_i$ is strictly greater than $2L_{\max}e_i$, then $a_i$ is also the unique optimum under $p$.
\end{MUproposition}
\begin{proof}
The triangle inequality bounds $\MUTV(p,r_i)$ by $e_i$. For a function taking values in $[0,L_{\max}]$, its expectations under these two distributions differ by at most $L_{\max}e_i$. Let $a_*$ minimize expected loss under $p$. Insert the two risks under $r_i$ between the risks of $a_i$ and $a_*$ under $p$. The middle difference is nonpositive by optimality under $r_i$, and the two outside differences are at most $L_{\max}e_i$. Regret is also at most $L_{\max}$. Applying the same two expectation bounds to the difference between $a_i$ and any competitor proves the strict-gap assertion.
\end{proof}

The two errors concern one actual chain of records. One cannot use the best transmission for one recipient and the best posterior approximation for another while reporting a common certificate. When $\widehat p=Q(M\mid E)$, Proposition~\ref{MU:prop:conditioning} can supply $\epsilon$. A separate information-loss estimate or an explicitly verified decoder may instead supply the bound. The result is relative to the stipulated posterior and loss, not a guarantee that a real decision is wise or that a protected real-world outcome can be represented by a bounded scalar loss.

An abstention action can be included with an explicit loss and availability rule. This permits a recipient to decline an unsupported decision within the same model. It does not create a universally correct abstention threshold. Nor does closeness in expected loss establish exact preservation of a logical obligation: an arbitrarily unlikely prohibited outcome can matter under a different contract.

\subsection{Approximate correction is not arbitrary tolerance}
A likelihood that is merely close in average error may have the wrong zeros, changing which observations the model considers possible. Posterior bounds must therefore specify positive evidence, the full observation record, and the domain on which the approximation holds. Confidence produced by normalization is not an independent check of those premises.

The exact factorization criterion in Section~\ref{MU:sec:bayesian} is an all-prior statement. The conditional-information identity is an average statement under one prior. The stability bound is an event-specific comparison of two joint laws. The decision certificate concerns one loss family. They are related, but none should be silently substituted for another.


\section{Sequential inquiry through a retained evidence channel}\label{MU:sec:sequential}
An explanation intended for further inquiry must support more than one update. The recipient may choose its next experiment after seeing the current report. Exact evidence preservation can survive this adaptation, but only when the experiment-selection and observation assumptions are stated together.

Let the unknown model $M$ remain fixed. At time $t$, a known policy selects experiment $e_t$ using only the prior reported history and recorded experiment names. Randomization is permitted when its conditional law is independent of $M$ given that history. Given $M=m$ and $e_t=e$, a fresh raw observation has law $L_{e,m}$, and its reported form is drawn through $B_e$. Conditional on the model and selected experiment, these draws do not depend on earlier hidden observations. Write $Q_{e,m}=L_{e,m}B_e$.

\begin{MUtheorem}[Sequential preservation of noisy correction under adaptive experiments]\label{MU:thm:adaptive-bayes}
Suppose each permitted experiment admits a kernel $\rho_e$ with
\[
L_{e,m}(o)B_e(o,z)=Q_{e,m}(z)\rho_e(o\mid z)
\quad\text{for all }m,o,z.
\]
For every finite positive-probability history under the declared policy, the posterior conditioned on the complete raw and reported record equals the posterior conditioned only on the reported outcomes and experiment names. The claim also holds at a bounded stopping time determined by that reported history.
\end{MUtheorem}
\begin{proof}
For a fixed realized history, its likelihood under model $m$ is the product of the selected-experiment probabilities and the terms $L_{e_t,m}(o_t)B_{e_t}(o_t,z_t)$. The policy factors are independent of $m$ once the reported history is fixed. Applying the assumed factorization gives a product of the terms $Q_{e_t,m}(z_t)$ multiplied by the model-independent factors $\rho_{e_t}(o_t\mid z_t)$ and the same policy factors. Posterior normalization cancels those factors. The remaining product is exactly the likelihood for the reported history. A bounded stopping decision determined by that history adds no model-dependent factor, so the same calculation applies.
\end{proof}

The theorem gives a reusable correction procedure, not a claim that the selected experiments are optimal. The selection policy can be inefficient while the posterior updates remain exactly preserved. Nor does the theorem cover arbitrary dependent sensor noise. When observation laws depend on hidden physical state or prior observations, that dependence must be represented in the state and likelihood rather than treated as a fresh independent draw.

\begin{MUcounterexample}[Copying a record is not repeating an experiment]\label{MU:ex:duplicate-record}
Let a binary model have equal prior probabilities, and suppose one observed positive report has likelihood $4/5$ under model 1 and $1/5$ under model 0. Its posterior probability for model 1 is $4/5$. Sending two copies of that same report does not add evidence, so the posterior remains $4/5$. Treating the copies as conditionally independent instead gives $16/17$. Both calculations are arithmetically valid for their respective likelihoods, but only the first describes duplication of one record.
\end{MUcounterexample}

A correction record must therefore preserve provenance about which observations are new and which merely repeat an earlier result. A translated explanation that reproduces two identical sentences is not, without that provenance, a reliable record of two independent measurements.

\subsection{Actions and stopping decisions can also carry information}
The theorem records experiment names and restricts selection to information already represented at the recipient. A hidden selector that consults an additional sensor changes the likelihood of which experiment gets selected. Discarding that selection mechanism can bias the update even when the reported sensor itself is well modeled.

A stopping time chosen from the public reported history is different: the history already contains the information used to stop, so the event of stopping contributes no additional model-dependent likelihood on that fixed history. This exact likelihood argument should not be confused with transferring a deterministic-time approximation to an arbitrary stopping time. Small unconditional error at each fixed time still needs a separate argument to yield a useful conditional error on selected histories.

\subsection{What a correction-preserving artifact must carry}
For a recipient to execute this result, an artifact or its declared background must specify the model family, prior, experiment names, likelihoods, reporting channels, and update rule. It must also identify which independence and selection assumptions are being used. Sending only the current most probable answer does not generally preserve the next update.

The recipient does not need every raw token when the factorization is valid. It needs sufficient evidence and an accurate description of the reporting process. In the instrument, the irrelevant flag can be omitted throughout adaptive choices among experiments with verified factorization. A coarse parity language cannot claim the same posterior contract merely because every parity sentence it transmits is true.



\chapter{Common witnesses and charged resources}\label{MU:book:chapter-48}


\section{One witness for prediction, transfer, and revision}\label{MU:sec:witness}
The separate results now have to refer to one realized system. A good state representation, an elegant translation, and a useful explanation do not establish a combined capability when they require different hidden information, different decoders, or incompatible resource assumptions.

\begin{MUdefinition}[Integrated adequacy record]\label{MU:def:witness}
An integrated adequacy record for a declared task family identifies a common computational realization, its information boundary, a reusable representation and its updates, permitted intervention implementations, each claimed transport and recipient, the model and evidence maps used for correction, and the cost and scope of those operations. Every claimed guarantee refers to these same objects or to explicitly related replacements.
\end{MUdefinition}
This is a specification proposal, not a new primitive of cognition. Its purpose is to make the quantifiers visible. The record cannot use a powerful target-specific decoder to establish prediction, then a different decoder to establish transfer, while reporting the result as one reusable method unless the selection procedure is itself part of the method.

\subsection{What can be checked in the finite core}
For explicitly supplied finite tables, the equations in Definition~\ref{MU:def:representation} can be checked on every state, operation, and query. Theorem~\ref{MU:thm:quotient} then proves the guarantee for arbitrary finite continuations. Exact transfer requires corresponding one-step and decoder checks. Correction transport requires the observation square and its stated reflection condition. The recipient must actually construct and execute the declared procedure within its allowed information boundary.

Some obligations remain different in kind. The finite model can list which intervention is admissible, but relating that declaration to a physical device requires evidence about the device. A table can name a sensor reading without proving calibration. A certificate can bind a program version without establishing the semantic meaning of its inputs. The specification should preserve these distinctions rather than present one successful check as universal certification.

\begin{center}\small
\begin{tabularx}{\textwidth}{@{}P{1.35in}XX@{}}\toprule
Claim & Required object & What it does not supply\\\midrule
Reusable prediction & One retained state, updates, decoders, and a fixed future family. & Access to that state or evidence of its actual use.\\
Mechanistic attribution & A realized mapping and discriminating permitted interventions. & Every untested mechanism or unrestricted causal access.\\
Method transfer & Matched states, operations, queries, and coverage. & Reverse reconstruction or new target operations.\\
Recipient explanation & Fixed background, an artifact, an interpreter, and a budget. & Faithfulness to the source solver's actual computation.\\
Corrective transfer & Candidate and evidence maps that preserve filtering. & Model invention, noise calibration, or truth outside the family.\\
Resource feasibility & One accounting boundary for construction and repeated use. & A universal efficiency advantage.\\\bottomrule
\end{tabularx}\end{center}

\subsection{The complete instrument witness}
For the original contract, the instrument witness contains $r(a,b,c)=(a+b,b)$, updates~\eqref{MU:eq:display-coordinates}, the first-coordinate decoder, and a statement that $c$ is outside the relevant future distinctions. It also names how an observer acquires the initial record. A recipient with that record and the equations can reproduce all declared predictions.

After $J$ is admitted, the witness changes. The old representation does not predict the new reading. An explicit $J$ observation supplies $c$ and updates the first coordinate. The enlarged state has the exact updates in Section~\ref{MU:sec:revision}. The original $H/W$ explanations remain valid on their unchanged domain. New predictions require the acquired evidence, and any translated explanation must preserve that evidence's discriminating meaning.

These are connected properties of one analytical construction. They do not establish that an unspecified learning algorithm will discover the representation, infer the correct transition law, or choose the diagnostic operation. The point is to make those remaining research obligations identifiable rather than confuse them with the solved representation problem.

\subsection{A cost account for reusable methods}
For a sequential workflow with disjoint accounting, total cost includes acquiring information, discovering the representation, encoding state, performing updates, verifying claimed relationships, translating when needed, and interpreting the result. Parallel execution requires work and elapsed time to be recorded separately. Costs in different units should not be added without a stated conversion.

A simple amortization comparison illustrates the issue. Suppose a representation costs $D$ to construct, each later use costs $c$, and the matched baseline costs $b>c$ per use. The representation has lower total cost after $N$ uses exactly when
\[
D+Nc<Nb,\quad\text{or equivalently}\quad N>\frac{D}{b-c}.
\]
This accounting does not prove that the representation can be discovered for cost $D$, that the uses are comparable, or that all target instances share it. It locates where such evidence would enter.

Understanding need not always make computation faster. A more costly representation can make a proof inspectable, reveal a missing assumption, or support an intervention unavailable to the cheaper one. The relevant benefit must be stated. The theory should not smuggle efficiency into the meaning of understanding, nor infer efficiency from the existence of a short description.

\section{Joint acquisition and explanatory communication}\label{MU:sec:joint-synthesis}
The acquisition problem and the recipient problem have so far been solved separately. Their composition is not automatic. An inexpensive observation may leave the sender unable to encode what a recipient needs. Conversely, requiring the sender to reconstruct the entire state can demand more evidence than is needed to help recipients who already possess relevant background.

We now combine the two exact finite models without introducing an omniscient encoder. The resulting test gives one witness for the experiment tree, its information boundary, a common artifact, and each recipient's decoding contract.

\subsection{The graph belongs to the acquired transcript}
Fix a finite set $X_0$ of possible initial states, a finite deterministic experiment catalogue as in Section~\ref{MU:sec:acquisition}, and a finite horizon $h$. A deterministic experimental policy is an observation-history tree. Its choices cannot depend on the hidden initial state. Only uniformly admissible choices are permitted at each reachable node.

For a policy tree $T$ and initial state $x$, write $s_T(x)$ for its terminal sender transcript and $x_T(x)$ for its final physical state. The tree has finitely many reachable leaves. Specify recipient backgrounds $\beta_i^T(x)$ and demands $d_i^T(x)$ by a declared rule. A demand may concern $x_T(x)$, the original $x$, or a stated pair of past and present values. These are different contracts.

The recipient sees only the common artifact and its specified background. If an experiment name, its timing, or another transcript is visible to the recipient, that information belongs in $\beta_i^T$. It must not be used as an uncharged auxiliary message. Let $C(T)$ be the maximum total experimental cost across initial states. Build $G_T$ on the reachable sender transcripts using the graph construction of Section~\ref{MU:sec:audience}, with $\alpha=s_T$.

\begin{MUtheorem}[A finite joint certificate for acquisition and explanation]\label{MU:thm:joint-acquisition}
For finite $X_0$, a finite experiment catalogue, horizon $h$, cost budget $C$, and artifact alphabet bound $K\geq1$, there exists a uniformly admissible deterministic acquisition policy and one common artifact satisfying every recipient demand exactly if and only if some attainable policy tree $T$ satisfies
\[
C(T)\leq C,\qquad G_T\text{ has no loop},\qquad \chi(G_T)\leq K.
\]
This is a finite decision problem. An accepting witness consists of the policy tree, a proper coloring, and the induced fixed decoders. Among feasible trees, a minimum experimental cost is attained. The statement does not assert that the search or the resulting decoders are efficient.
\end{MUtheorem}
\begin{proof}
Any feasible system induces a policy tree of depth at most $h$, because its experiment choices depend only on its available history. The sender's final artifact is a function of its terminal transcript. Applying Theorem~\ref{MU:thm:audience} to those transcripts proves the graph conditions. The worst-case experimental cost is $C(T)$, so it satisfies the budget.

Conversely, choose a tree with the stated properties. Execute its uniformly admissible experiments, encode the terminal transcript using its proper coloring, and use the decoders constructed in Theorem~\ref{MU:thm:audience}. They give every required answer on every initial state. The tree's cost and horizon meet the contract. There are finitely many experiment and outcome choices at each node and a finite depth bound, hence finitely many candidate trees. Exhaustive admissibility, cost, and coloring checks decide existence. A nonempty finite set of feasible tree costs has an attained minimum.
\end{proof}

The graph is built across all reachable leaves, not separately at each leaf. Two different transcripts can be indistinguishable to a recipient after encoding even when each sender support is locally well understood. A message assignment chosen independently at each leaf can therefore fail the common-artifact requirement.

For a fixed preexisting decoder family, each colored artifact must additionally be checked against that family. For a finite catalogue of allowable implementations, direct verification of those witnesses remains finite. The coloring theorem alone does not make a convenient decoder available, just as the acquisition theorem does not create a sensor.

\begin{MUcounterexample}[A useful explanation need not identify the state at its source]\label{MU:ex:joint-source}
Let the hidden state be two bits $(a,b)$. The sender's only experiment reads $a\oplus b$ without changing the state. Repeating it never identifies the complete pair. One recipient already knows $a$ and needs $b$, while another knows $b$ and needs $a$. The single parity reading and its one-bit broadcast satisfy both demands. Before that reading, the sender has only one transcript and the explanation graph has a loop. After the reading, the graph is one edge and the one-bit message is exact.
\end{MUcounterexample}

This is an integrated information construction, not evidence that the sender possesses the recipients' full understanding. The explanatory success is jointly realized by the acquired parity and the recipients' backgrounds. Attribution should identify which participant carries which distinction.

\subsection{Why noisy belief and exact state recovery remain different}
The joint theorem concerns deterministic experiments and exact demands. The Bayesian sections concern noisy records and distributions over a model family. One cannot replace the target state in the joint theorem with a most probable guess and claim the exact guarantee has been preserved.

In the noisy instrument, every readout has positive likelihood under every candidate value of $c$. With a fully supported prior, any finite sequence of those readouts leaves every candidate with positive posterior probability. Repeated reports can make one value overwhelmingly likely, but they cannot give zero-error identification within a fixed finite number of such observations. A decision-loss tolerance, a posterior threshold, and an exact reconstruction target are different terminal conditions.

The decision certificate in Proposition~\ref{MU:prop:decision} supplies one bridge: a recipient can act under an explicit error budget without pretending its estimate is exact. A more general synthesis optimizing noisy experimentation and multi-recipient communication jointly remains a separate problem. That larger stochastic optimization is outside the finite exact synthesis of this edition.

\subsection{A complete expanded instrument record}
The original instrument already separates current display, reusable state, and newly relevant information. The expanded account adds the mechanisms by which that information is acquired, represented probabilistically, and shared.

First, the reasoner retains the base coordinates and the exact scope of $H/W$ reuse. Admitting $J$ splits old behavioral classes. The exact diagnostic plan uses a real $J$ operation and preserves its effect on the first register. If only the noisy sensor is available, the retained object is a posterior over possible post-$J$ states, not an invented exact value of $c$.

Second, the evidence description identifies which readouts are conditionally independent and which tokens are copied or ancillary. Removing the independent flag is correction preserving. Replacing all five possible readings by parity is not. Three identical full readings give the conditional posterior $432/433$ for the corresponding candidate, but that number remains conditional on the declared likelihood model and prior.

Third, recipients receive the method appropriate to their backgrounds. A recipient with the base coordinates needs the newly acquired distinction or its uncertainty distribution. A recipient without those coordinates needs additional state information. A shared artifact is checked across all of those backgrounds at once. Its source encoder may use only the information actually acquired.

Finally, the record includes the decoder, query loss or exactness criterion, operation catalogue, and full resource account. Discovery of the transition law, calibration of the sensor, and mechanistic attribution to a particular neural system are not supplied by this analytical example. They remain separate scientific claims with separate evidence requirements.

\subsection{A release boundary rather than an inflated definition}
The mathematical contribution is a connected set of exact and conditional preservation results. The finite graph and decision tree supply executable witnesses. The likelihood factorization identifies when discarded records do not change Bayesian correction. The error bound exposes when a rare observation makes an approximate model inadequate. The recipient certificate carries those errors into a defined decision criterion.

Together they make a stronger specification of usable understanding. They do not turn every valid compression into a cognitive achievement or every encoded table into an economical explanation. Claims of economy require an explicit interpreter language, charged background information, and a measured or proved implementation cost. Claims of actual use require a realized causal mechanism. Keeping these obligations separate is part of the theory's content, not a reason to postpone its mathematical statements indefinitely.


\section{What the resource account must charge}\label{MU:sec:resources}
A short explanatory message can conceal an expensive recipient. A compact state can require an elaborate encoder. A fast online answer can depend on a table that was expensive to construct and distribute. None of these possibilities makes the answer wrong. They change what has been shown about the economy of the method.

The first-edition resource contract separates four accounts. The \emph{design account} includes discovering and validating the state model, constructing a decoder, and searching for a policy. The \emph{installed account} includes the fixed interpreter, model tables, recipient background, and compiled controller. The \emph{online account} includes observations, physical interventions, communication, and executed instructions. The \emph{revision account} includes obtaining new evidence and replacing invalid assumptions or interfaces. A quantity cannot be declared absent merely because it was paid in a different account.

These accounts use different units. Description bits, controller records, modular operations, physical interventions, and latency are not interchangeable. A weighted scalar cost is legitimate only after its conversion factors have been specified. The synthesis below retains a vector of online charges and a separate controller storage count. It does not silently convert one recipient's computation into the sender's communication budget.

\subsection{Task-dependent interpreter changes are information}
Let $\Theta$ be a finite nonempty family of tasks and let each task have a required answer profile over a fixed query family. A fixed interpreter $U$ receives background $B_\theta$, an installed program string $p_\theta$, an artifact $e_\theta$, and a query. It must return the required answer. The interpreter is fixed across tasks. The strings may depend on the task, but their lengths are charged. Let $N_B$ be the number of distinct required answer profiles among tasks having the same background value $B$.

\begin{MUproposition}[Charging a task-dependent interpreter]\label{MU:prop:advice}
Suppose $p_\theta$ and $e_\theta$ have fixed lengths $d$ and $b$ bits, respectively, and the deterministic interpreter answers every required query exactly. Then
\[
2^{d+b}\geq\max_B N_B.
\]
In particular, representing every Boolean function on $n$ input bits under one fixed background requires $d+b\geq2^n$. This counts task-dependent information, not the running time of the interpreter or the information needed for any one specially chosen function.
\end{MUproposition}
\begin{proof}
Fix a background $B$. If two distinct required answer profiles used the same pair $(p_\theta,e_\theta)$, the fixed interpreter would receive identical inputs for a distinguishing query and would have to give the same answer. Exactness would fail for one of the two tasks. There are at most $2^{d+b}$ pairs, proving the first bound. There are $2^{2^n}$ Boolean functions on $n$ input bits, with distinct complete answer profiles. Taking base-two logarithms proves the second assertion.
\end{proof}

This is the counting argument of Proposition~\ref{MU:prop:explanation} applied to the complete task-dependent package. It is not a new general description-complexity theorem. The point is to prevent a moving interpreter from becoming an unreported communication channel. Variable-length descriptions require a declared framing or prefix convention. Merely saying that a message has length at most $b$ is not the same fixed-length model, because the length itself can carry information.

A task-independent interpreter can contain rich mathematical knowledge without encoding which task instance is present. Its construction and availability still belong in the installed and design accounts. Conversely, a separate decoder chosen for every task must be identified and charged even when that decoder is called a background assumption. Control-task critiques of probing and expressivity critiques of alignment make this boundary especially important~\cite{MU:hewitt,MU:sutter}.

\subsection{Amortization without disappearing costs}
Suppose a method has setup cost $S$ and per-use cost $r$, while a comparator costs $c$ on each use. In the same declared cost units and with no other differences, the method costs less over $N$ uses exactly when $S+Nr<Nc$. When $c>r$, this is $N>S/(c-r)$. At equality the costs tie. If $c\leq r$ and $S\geq0$, setup cannot create a strict cost saving in this model.

The formula does not establish that the method can be discovered, that the repeated problems are interchangeable, or that its validity survives a change of task distribution. Its role is bookkeeping. A useful representation can repay a substantial initial investment. It cannot claim that investment never occurred. Likewise, a precomputed table can be a sensible engineering choice for a bounded domain. The comparison should expose its storage and preparation rather than reject it by terminology.

\subsection{Cost is an attribute of a specified implementation}
An algebraic equation does not have one architecture-independent runtime. Addition modulo five can be a primitive instruction, a small Boolean circuit, or a table lookup. The worked implementation in Section~\ref{MU:sec:first-ledger} counts modular arithmetic operations and physical calls explicitly. It also records the sizes of its actual code and serialized witness. These byte counts describe those files in one environment, not minimum possible descriptions or a measure of cognitive complexity.

Causal attribution remains separate. A cheap decoder may show that information is readily accessible. It does not show that the source solver actually uses that decoder or the selected representation. Resource accounting strengthens an explanatory claim by identifying where work is done. It does not replace the implementation evidence discussed in Section~\ref{MU:sec:use}.

\section{Explanations for fixed recipient procedures}\label{MU:sec:fixed-decoders}
The coloring criterion in Theorem~\ref{MU:thm:audience} permits fixed decoders to be designed for the chosen code. A recipient that already exists may not be able to execute those decoders within its budget. We now freeze the decoding procedures before choosing the artifact.

Let $X$ be a finite nonempty set of cases, $\alpha:X\to S$ the sender's access map, and $\mathcal A$ a finite nonempty artifact catalogue. Recipient $i$ has a background map $\beta_i:X\to B_i$, a demanded output $d_i:X\to D_i$, and a fixed procedure $L_i$. Evaluate $L_i(\beta_i(x),a)$ under a declared instruction and memory bound. A timeout, an invalid output, or a rejected artifact counts as failure. Since the inputs and execution bounds are finite, the resulting acceptance table is finite even when $L_i$ is written in a general programming language.

Define
\[
\mathcal A(x)=\{a\in\mathcal A:\text{every recipient returns }d_i(x)
\text{ within its declared limits}\},
\]
and, for each reachable sender value $s$, define
\begin{equation}\label{MU:eq:decoder-intersection}
\mathcal A_s=\bigcap_{x:\alpha(x)=s}\mathcal A(x).
\end{equation}
An artifact selected from $\mathcal A_s$ must work for all cases the sender still considers possible. Choosing a different acceptable artifact after secretly learning which hidden case is true is not allowed.

\begin{MUtheorem}[The fixed-recipient covering criterion]\label{MU:thm:fixed-decoder}
Under this finite contract, an exact encoder $E:S\to\mathcal A$ exists on the reachable sender values if and only if every $\mathcal A_s$ is nonempty. Its smallest possible image size is
\begin{equation}\label{MU:eq:decoder-cover}
\min\{|C|:C\subseteq\mathcal A,\ C\cap\mathcal A_s\ne\varnothing
\text{ for every reachable }s\},
\end{equation}
when the feasible family is nonempty. Thus a codebook of at most $K$ artifacts is possible exactly when a set of at most $K$ artifacts meets all these requirements. Every valid encoder induces a proper coloring of the sender-access graph, so its image size is at least that graph's chromatic number. The converse graph condition need not hold for the fixed procedures.
\end{MUtheorem}
\begin{proof}
A sender seeing $s$ must emit one value $E(s)$, and that value must lie in $\mathcal A(x)$ for every $x$ with $\alpha(x)=s$. This is exactly membership in $\mathcal A_s$. Hence a valid encoder implies nonempty intersections. Conversely, selecting one value from each nonempty finite intersection constructs a valid encoder.

The image of any such encoder intersects every $\mathcal A_s$. If a set $C$ intersects every requirement, selecting $E(s)$ from $C\cap\mathcal A_s$ gives an encoder whose image is contained in $C$. This proves the minimum formula by finite minimization. If two adjacent sender values received the same artifact, a recipient with the same background would receive identical decoder inputs while requiring distinct outputs. Its fixed procedure could not satisfy both. This proves the graph lower bound.
\end{proof}

The finite optimization in~\eqref{MU:eq:decoder-cover} is a set-cover formulation on the reachable sender values. It counts how many distinct catalogue entries are needed. It does not license replacing their actual encodings by short indices unless the corresponding index decoder is installed and charged. A catalogue entry may be a large program or proof object even when only one such entry is used.

\begin{MUcounterexample}[Information suffices but the installed decoder needs a larger codebook]\label{MU:ex:fixed-codebook}
Take three cases $x_0,x_1,x_2$ fully visible to the sender. One recipient already has background $i$ in case $x_i$ and always needs output zero. The sender-access graph has no edges, so a designable decoder could ignore a single common message and return zero. Now fix an installed decoder on artifacts $A,B,C$. It returns zero on $A,B$ at background 0, on $B,C$ at background 1, and on $A,C$ at background 2. It fails on the remaining pairs. Every two of the three acceptable-artifact sets intersect, but their common intersection is empty. Two catalogue entries suffice, while one does not. The graph lower bound is one and the fixed-decoder optimum is two.
\end{MUcounterexample}

This construction is intentionally small. It isolates the difference between information-theoretic possibility and compatibility with an already installed procedure. A more capable replacement decoder repairs the example, but replacing it changes the installed account. A procedure that times out on a required input provides an equally direct reason a graph witness can fail at the resource boundary.

\subsection{Sender uncertainty is a different obstruction}
Nonempty casewise sets $\mathcal A(x)$ do not suffice. If the sender cannot distinguish two cases whose acceptable-artifact sets are disjoint, their intersection in~\eqref{MU:eq:decoder-intersection} is empty. More recipient compute cannot repair a distinction missing from both its background and the transmitted artifact. A new observation may repair the sender's access. An added recipient background may repair the receiving side. Neither can be credited without being part of the actual workflow.

The bounded table concerns exact outputs. For approximate explanation, acceptable outputs can instead be defined by a fixed loss tolerance, but that changes the table and the claimed guarantee. A most probable model is not an exact current state merely because the decoder confidently emits it. The Bayesian and decision results retain their own uncertainty contracts.

\section{Joint acquisition with a compiled resource budget}\label{MU:sec:budgeted-synthesis}
The preceding theorem freezes recipient procedures. We now combine it with state-changing acquisition and charge the entire compiled observation tree. This closes a specific first-edition gap: a small artifact, an affordable experiment, and a runnable decoder must belong to one attainable workflow.

\subsection{The explicit finite implementation model}
Use a finite augmented state space $X$. Include every origin variable, history flag, recipient background, and permission needed to make the declared transition and terminal tables valid. Each experiment $e$ has the transition, observation, and admissibility semantics of Section~\ref{MU:sec:acquisition}, together with a fixed nonnegative integer cost vector $c_e\in\mathbb N^r$. Its coordinates have declared units. The charge includes the specified source-side interface work for that experiment. We do not assume that acquiring information leaves the physical state unchanged.

For each state $x$ and artifact $a$ accepted by all fixed recipient procedures, let $\gamma(x,a)\in\mathbb N^r$ be the terminal charge. It includes transmission and the declared source and recipient execution work. Invalid or out-of-budget decoder executions have no terminal entry. At support $V$, a stop using $a$ is allowed only if it works for every $x\in V$, with worst-case vector
\[
g_a(V)=\max_{x\in V}\gamma(x,a),
\]
where maxima and comparisons of vectors are componentwise. Different coordinates may attain their maxima on different cases. A coordinatewise budget nevertheless has to cover every one of them.

A controller is an explicit observation tree of depth at most $h$. Each decision node or stop node occupies one fixed-size record, with fields for the operation or artifact and the bounded number of successor pointers. The count $n(T)$ includes all stored branches, not just those visited in one execution. The interpreter, recipient programs, and state tables are installed separately and remain fixed throughout the optimization. Their storage is added to the recorded controller size. Sharing equal subtrees in a directed acyclic graph is a different storage model.

For a tree $T$, let $v(T)$ be its componentwise worst-case charged online cost. Let $\mathcal P_h(V)$ be the attainable pairs $(v(T),n(T))$. The stop pairs are
\begin{equation}\label{MU:eq:resource-stop}
\bigl(g_a(V),1\bigr)\quad\text{for each artifact accepted throughout }V.
\end{equation}
For $h>0$, each uniformly admissible experiment gives the additional pairs
\begin{equation}\label{MU:eq:resource-branch}
\left(c_e+\max_o v_o,\ 1+\sum_o n_o\right),
\quad (v_o,n_o)\in\mathcal P_{h-1}(V_{e,o}),
\end{equation}
with one child choice for every nonempty observation branch. An experiment contributes no pair when even one possible branch has no completion. At $h=0$ only stops are allowed. Sets are finite because the catalogue and horizon are finite.

\begin{MUtheorem}[A joint resource certificate for fixed recipients]\label{MU:thm:resource-synthesis}
Recurrences~\eqref{MU:eq:resource-stop}--\eqref{MU:eq:resource-branch} enumerate exactly the cost profiles of the permitted explicit trees. A workflow within online budget $b\in\mathbb N^r$ and controller-record budget $N$ exists if and only if some $(v,n)\in\mathcal P_h(V)$ satisfies $v\leq b$ and $n\leq N$. Every such pair has a finite executable tree witness. Keeping only componentwise nondominated pairs at each recursive step preserves feasibility for every pair of budgets $(b,N)$.
\end{MUtheorem}
\begin{proof}
At depth zero, an admissible workflow must stop with one artifact that works for every possible state. Its worst-case terminal charge and its one stored record give exactly~\eqref{MU:eq:resource-stop}. For the inductive step, any nonterminal tree first executes an experiment available throughout $V$. Its possible next supports are exactly $V_{e,o}$. By induction each child tree supplies one attainable pair. Every branch is possible, so the online charge is $c_e$ plus the componentwise maximum of the child charges. The installed explicit tree contains the root and every child record, giving the sum in~\eqref{MU:eq:resource-branch}. Conversely, combining an admissible root with the witnessing child trees realizes that pair. This proves exact enumeration and the budget equivalence.

If one child profile is dominated by another attainable profile at the same support and remaining horizon, replacing that child cannot increase a parent's componentwise maximum or its sum of records. Terminal admissibility and future observations are unchanged. Repeating this replacement proves that discarding dominated profiles preserves a budget-feasible witness whenever one exists. Finite induction gives the claim at the root.
\end{proof}

This is finite multiobjective dynamic programming under a fixed implementation contract. Its components have the same state-planning and side-information foundations as the earlier results~\cite{MU:smallwood,MU:witsenhausen,MU:alon}. The new obligation is their simultaneous realization with frozen recipient procedures, online charges, and explicit controller storage. The theorem does not optimize over arbitrary programs, discover missing sensors, or make policy synthesis computationally cheap.

\begin{MUcounterexample}[Separate resource optima do not define a feasible common workflow]\label{MU:ex:resource-frontier}
Let a hidden origin take four values and remain fixed under every experiment. A direct experiment reveals its value at charge four units. Its explicit controller has one decision node and four stop nodes, giving profile $(4,5)$. Three other experiments ask whether the origin equals 0, 1, or 2, respectively, each at charge one unit. A decision chain of these equality tests distinguishes the four origins with worst-case charge three and seven stored records, giving profile $(3,7)$. The fixed recipient decoder reads a two-bit origin artifact at a leaf. Its work and the common message length are held fixed and are not included in the one online coordinate shown here.

Within horizon three the nondominated profiles are $(4,5)$ and $(3,7)$. A tree with charge at most three cannot use the direct experiment. A binary equality tree distinguishing all four origins needs at least four leaves and hence at least seven records. Conversely, any tree distinguishing four origins needs four distinct stopping artifacts and at least one decision node, so five is the smallest possible record count. Thus the separate minima are three online units and five records, but the combined budget $(3,5)$ admits no tree.
\end{MUcounterexample}

The direct observation is compact as an installed tree. The equality chain is cheaper to execute in the worst case. Taking the best coordinate from each certificate invents a workflow that neither realizes. Replacing the explicit tree by a program with shared branches changes the storage model and must be assessed separately.

\subsection{Which implementation costs remain outside the recurrence}
Offline search and proof checking do not disappear because the controller is compiled. Their time, software, and input tables are recorded separately. The online interpreter must be specified, and its dispatch costs can be included in $c_e$ and $\gamma$. A record count becomes a storage guarantee only after the fixed record layout and installed tables are included. The number of source-file bytes is an audit quantity, not a theorem of minimal description length.

The result is exact for declared charges. If those charges are conservative upper bounds for an implementation, the certificate is correspondingly a sufficient implementation budget, not necessarily its tight physical cost. If the actual decoder reads an unlisted file or the sensor consults a hidden variable, the implementation is outside the contract. The right response is to expand and recheck the contract, not to call the extra access free.

A source with several possible initial observations can be handled by running the construction on each initial support and specifying how the installed controller dispatches among them. If that observation is public to the recipients, it belongs in their background. If it is private, the fixed decoders still see only their declared inputs. The integrated instrument below makes both sides explicit.

\section{A costed first-edition instrument certificate}\label{MU:sec:first-ledger}
The common realization can now be described without a designable recipient or an omniscient sender. The instrument is still $(a,b,c)\in\mathbb F_5^3$. Before the diagnostic, the sender knows only $(s,b)=(a+b,b)$. Recipient A has the old register pair $(a,b)$. Recipient B has the old display pair $(s,b)$. Neither recipient has $c$. All three parties know the operation laws and the fixed diagnostic protocol before the particular hidden state is selected.

The sole diagnostic in this version applies $J$ exactly once and returns its exact post-operation display $s_J$. The sender computes $m=s_J-s=c$ and broadcasts that one field element. The report $s_J$ and the sender's private controller address are not transmitted separately. The recipients are told by the installed protocol that their demanded record is the state after this one diagnostic, not an arbitrary later state.

The two frozen decoders are
\begin{equation}\label{MU:eq:fixed-instrument-decoders}
L_A((a,b),m)=(a+m,b,m),\qquad
L_B((s,b),m)=(s+m,b,m).
\end{equation}
Recipient A obtains the post-$J$ register description. Recipient B obtains the post-$J$ display-centered description. Both compute one modular addition. Their installed update procedures then support all $H/W/J$ continuations in their respective coordinates. These procedures are supplied, not learned by the diagnostic.

For a fixed old $(s,b)$, five possible values of $c$ remain. Before an observation, no one message can supply both exact demanded records for every one of those cases. After the diagnostic, the five report values lead to five stop nodes and each selects its one valid message. The explicit tree has six records. A fixed three-bit field encodes the five messages. Two bits do not suffice. Three additional unused binary patterns are rejected, not assigned secret meanings.

\subsection{The execution ledger}
The following charges describe the reference procedure rather than a universal machine model. Assignments and interpreter dispatch are recorded as implementation operations, but are not included in the modular-arithmetic coordinate. Each field operation is one modular addition or subtraction, independent of how it is realized by a processor.

\begin{center}\small
\begin{tabularx}{\textwidth}{@{}P{1.7in}P{.7in}X@{}}\toprule
Resource & Count & Meaning and boundary\\\midrule
Physical diagnostic & 1 & One $J$ operation with one returned exact display. Sensor construction is supplied.\\
Source modular arithmetic & 1 & Subtract the old display from the new one. Selection uses the compiled protocol.\\
Broadcast payload & 3 bits & One fixed-length field for five possible messages. No branch address or raw report is an additional message.\\
Recipient A modular arithmetic & 1 & Compute $a+m$ in its fixed decoder.\\
Recipient B modular arithmetic & 1 & Compute $s+m$ in its fixed decoder.\\
Explicit controller & 6 records & One diagnostic node and five message leaves for the fixed old record. Two records are visited in a run.\\
Installed method & Recorded separately & Both decoder programs, transition semantics, framing, and prior background. Not charged as newly acquired $c$.\\\bottomrule
\end{tabularx}\end{center}

The code enumerates all 25 old-record fibres and all 125 initial instrument states, compiles a resource witness, and verifies its leaves by executing the fixed decoders. It also tests incorrect messages, missing branches, reduced horizons, and insufficient recipient instruction limits. These are exact finite checks of the supplied model. They do not establish that an external reasoner discovered the representation or that a noisy device has this exact sensor.

The six-record count applies to one fixed old record and a literal observation tree. A generic arithmetic controller can share structure across all 25 old records and need not store 25 separate trees. The reference program implements that arithmetic procedure as well and records its code size. Neither representation is declared the universally smallest controller. The distinction between tree storage and an arithmetic program is retained instead of combining their best costs.

\subsection{Reuse has an explicit implementation}
In B's coordinates $(s,b,c)$, each $H$, $W$, or $J$ update uses one modular addition or subtraction and the query reads the first coordinate. Thus a length-$n$ continuation uses exactly $n$ modular arithmetic operations after initialization, plus its recorded dispatch, assignment, and output operations. In A's coordinates the current display requires the additional addition $a+b$, while swap can be a coordinate permutation. The programs are different but agree on the declared answers through the exact translation theorem.

The supplementary record has five possible values. A separately stored three-bit message appended to a jointly packed five-bit base record uses eight bits, whereas a repacked complete 125-state record requires seven bits. A native implementation storing three individually encoded residues can use nine bits. These are three distinct storage layouts. The delivered source reports which layout is used rather than subtracting rounded bounds from incompatible layouts.

\subsection{An uncertain reading requires a different artifact}
The deterministic certificate does not apply to the noisy sensor merely because the same physical register is present. Under the noisy contract, the sender carries a posterior over $c$, not a certified exact value. A recipient can use a transmitted likelihood history or a specified posterior representation together with the existing correction and decision guarantees. Its message size and arithmetic work must then be charged again. They are not inherited from the three-bit exact-state message.

The ancillary flag may still be omitted when its verified factorization holds. Copies of one reading may not be treated as independent measurements. When a public adaptive selector chooses the next experiment using reported history, Theorem~\ref{MU:thm:adaptive-bayes} applies under its fresh-observation assumptions. A hidden selector, an unknown likelihood, or an observation deemed impossible by every current model requires additional evidence or a revised model.

\subsection{What this certificate warrants}
The first edition supplies exact witnesses for specified representations, experiment trees, reporting channels, recipient procedures, and resource profiles. These witnesses support different statements. The quotient preserves future answers. An admissible diagnostic acquires a missing distinction. A factorization preserves a Bayesian update. A fixed decoder turns an artifact into a usable record. The common resource witness shows that particular components can be executed together.

None of these conclusions identifies the internal mechanism of an unrelated solver. Nor do they establish subjective experience, a universal ranking of understanding, or accuracy of unstated physical premises. A complete adequacy record should preserve those unresolved questions alongside its positive results. That is how a reusable method remains answerable to evidence rather than becoming an impressive answer with an invisible scope.



\chapter{Interpretation and scope}\label{MU:book:chapter-49}


\section{Limits, objections, and the next mathematical questions}\label{MU:sec:limits}
The framework is intentionally operational, but that does not make it philosophically complete. A contract specifies selected consequences. It does not establish that those are all the consequences that matter. A formal model of a subject can be wrong, and a model of a reasoner's computation can be incomplete. A successful finite witness does not answer every question about meaning, intentionality, or experience.

\subsection{Does this define understanding as success?}
Not as a single success rate. The framework separates what a state can support, what a solver actually uses, what can be transported, what a recipient can reconstruct, and what survives correction. The counterexamples show that these properties differ. Nevertheless, every operational criterion remains a criterion for a declared set of abilities. This account does not offer an independent essence of understanding beyond all such criteria.

The resulting restraint is useful. A system may possess strong reusable knowledge in one domain and weak correction behavior in another. A receiver may understand an explanation under a supplied notation while failing to translate it. A process may be correctly modeled by an outside analyst without being able to explain itself. These are distinct descriptions rather than contradictory verdicts about an indivisible trait.

\subsection{Can a table satisfy the theory?}
Yes, for some finite contracts. A table can encode a transition system and support exact future queries. It should not be declared incapable by terminology alone. The stronger questions concern its origin, size, use, reuse across a family, and behavior when the contract changes.

The recoding construction in Section~\ref{MU:sec:use} makes this unavoidable. To distinguish a compact transferable method from a complete enumeration, one must specify the comparison class and account for the encoder, decoder, and shared background. Restricting those components is part of the research design. An assertion that a representation is meaningful is not a substitute for doing so.

\subsection{Do valid proofs establish understanding?}
A proof establishes the conclusion licensed by its premises and rules. The recipient may still lack a usable model of why those premises were appropriate or how the proof changes under a new problem. Conversely, a reasoner can use a correct model while being unable to produce a formal proof in a particular proof language.

Formal verification is therefore valuable but not interchangeable with the entire profile. The present proofs establish finite and metric statements under their assumptions. They do not certify that a deployed neural model realizes the maps, that a human reader has understood them, or that the task specification captures all relevant external facts.

\subsection{Finite models and a changing world}
The exact quotient assumes a fixed state space and a fixed operation family. The correction theorem assumes specified model families and outcome maps. The instrument's new operation enlarges one of those families in a controlled way. Real inquiry can also introduce new variables, new causal hypotheses, new measurement channels, and new forms of explanation.

There is no guarantee that the minimal state for every future extension is small. Preserving all conceivable future distinctions can collapse the compression objective back toward retaining the entire accessible state, and even that state may omit external information. A useful theory must allow partial coverage and a principled request for new evidence rather than promise unlimited generality.

\subsection{What is now resolved, and what remains}
\paragraph{Bounded nonvacuity.}
The sender-access graph prevents an encoder from consulting missing information, and the joint certificate keeps acquisition and communication in one contract. Sections~\ref{MU:sec:resources}--\ref{MU:sec:budgeted-synthesis} add fixed interpreter procedures, charged task-dependent descriptions, and exact finite resource frontiers. These do not constitute a general complexity theory of concepts. Useful scalable representation languages, economical decoder discovery, and fair cross-architecture comparisons remain research questions. A finite coloring can be optimal in alphabet size and poor as an explanation to an actual person.

\paragraph{Noisy and approximate correction.}
Theorem~\ref{MU:thm:bayes} resolves exact posterior preservation for finite likelihoods and common reporting channels. Proposition~\ref{MU:prop:conditioning} quantifies event-specific sensitivity, and Proposition~\ref{MU:prop:decision} carries approximation through a recipient to a bounded-loss decision. These results assume the likelihood family and its errors are supplied. Robust model sets, dependent noise, misspecification, and empirically justified uncertainty require additional work.

\paragraph{Acquiring new distinctions.}
Theorem~\ref{MU:thm:acquisition} resolves optimal worst-case acquisition in the finite deterministic, finite-horizon model, including the experiment's effect on current state. Theorem~\ref{MU:thm:joint-acquisition} adds a common explanatory artifact. Efficient search, noisy joint experimentation and communication, and the invention of missing sensors or hypotheses are not supplied. A statement about feasible observation cannot be promoted to a claim of autonomous scientific discovery.

\paragraph{Recipients and revision over changing contracts.}
The graph criterion resolves exact communication when decoders can be designed. Theorem~\ref{MU:thm:fixed-decoder} instead resolves finite artifact selection for preexisting bounded procedures. Neither identifies which procedure is realized by an uninstrumented learned reasoner. The sequential correction theorem assumes a fixed model family and matching experiment semantics. A recipient who changes its conceptual language or acquires a new operation may need a new sufficient record, not just the old certificate copied forward.

These remaining questions define the boundary of the first edition. Version 1.0 is a stable theoretical scope with reproducible artifacts. It is not a claim that every question about understanding has been settled. Later empirical findings or mathematical corrections can revise the account without turning an editorial release number into a scientific guarantee.

\section{Conclusion}\label{MU:sec:conclusion}
A correct answer is an achievement. Reusable computational understanding concerns what a reasoner can carry beyond that answer: a method for new continuations, a representation another reasoner can use, and an explicit route to correction when further evidence changes what is warranted. Those properties belong to a continuing relationship among a subject, a representation, evidence, and a recipient.

The instrument makes the distinctions concrete. One display answers a present question. Two coordinates support every continuation under the original operations. A third register becomes relevant when a new operation is admitted. An exact diagnostic can acquire that distinction, but also changes the instrument. A noisy reading supplies a posterior instead. Reporting channels preserve correction only under the corresponding evidence conditions. Repeating a copied report is not a new observation, and a small unconditional error can become decisive after rare evidence.

Explanation adds another realization boundary. A common artifact must fit what the sender actually knows and what each recipient already possesses. Even a message with sufficient information can be unusable by the installed decoder. Changing that decoder can help, but its task-dependent information, installation, and execution must be charged. The first-edition covering and resource constructions make those obligations jointly checkable in finite models. A controller small in stored records need not be the cheapest one to execute. The best separate resource bounds need not describe one possible workflow.

The theory does not select a uniquely human notation or exclude a table merely because it is a table. It asks what each implementation retains, how it was acquired, what work it performs, and which losses remain after translation. Predictive adequacy, causal use, economical realization, recipient usefulness, and responsiveness to evidence are related but distinct. A proof about one is not a certificate for all of them.

Operation Komorebi's realization perspective provides the common discipline: specify the operations that a representation preserves and make unresolved distinctions visible. The resulting framework supports exact mathematical claims without inferring consciousness, universal comprehension, or an unobserved mechanism from formal success alone. Its positive proposal is that a useful account of machine understanding should follow a method through reuse, explanation, and correction, keeping the information and work that make those achievements possible inside the account.



\part{Proof paths and mathematical memory}
\setcounter{section}{0}\renewcommand{\thesection}{PPT.\arabic{section}}\renewcommand{\theHsection}{PPT.\arabic{section}}

\chapter{Proof fields and certificates}\label{PPT:book:chapter-50}
A completed proof occupies only part of its surrounding field. Keep the successful deduction, but do not confuse its unused alternatives with either errors in that deduction or proofs of impossibility.

  \PPTopeningfigure{fig_14_ambient_field.pdf}{A declared proof field retains three different statuses. The solid route certifies the target. The upper branch has an unresolved obligation. The lower branch has a scoped refutation. Neither the open branch nor the refuted step invalidates the certified route.}{fig:ambient}
  \vspace{4pt}{\sffamily\bfseries\color{KomorebiTeal}The governing distinction}\par
  A proof field preserves more than a proof needs. Its failed and unresolved records remain available to inquiry without acquiring the force of certified deductions.
 


\section{Introduction}
\label{PPT:sec:old1}

A published proof often presents an argument after its exploratory work has been compressed. Definitions have stabilized, unsuccessful approaches have been set aside, and a long calculation may now appear as a lemma. The change of representation that first made the argument visible can disappear into a single sentence. Such compression is often useful: it leaves a certificate that is easier to check and reuse.

Compression can also remove mathematical information that is not needed for the present theorem. An example may explain why a definition was chosen. A failed route may isolate the exact boundary of a method. An alternative proof may carry a bridge to another field. A computation may reveal a stronger statement than the one originally sought. None of these objects is automatically part of the final implication
\[
A\vdash B.
\]
They may still be part of the mathematical event that produced it.

Proof Path Theory studies this larger object without changing the standard of proof. Its central separation is
\[
\boxed{
\text{certificate of validity}
\quad\ne\quad
\text{expanded discovery field}.
}
\]
The certificate answers whether the conclusion follows. The discovery field records how the conclusion became reachable, what alternatives were considered, what was lost during compression, and which omitted structures may matter later.

The construction begins with a typed field of claims, justifications, evidence, and targets. Its certified rules determine which targets are derivable. A certificate selects enough of those records to support a declared conclusion, while the ambient field can retain unresolved or refuted alternatives.

A second layer asks what survives when the field is compressed. The answer depends on the questions that a representation must support. An archive intended for future discovery may preserve examples and failed approaches that a verification certificate can omit. A bridge graph makes one such retention problem precise.

The final layer concerns the continuing life of the archive. Revisions change the set of claims that need review. Shared errors affect the information supplied by multiple reviewers. Audit capacity limits how quickly new candidates can be checked. These models are connected by their declared proof objects, without treating social judgments as substitutes for deduction.



\begin{figure}[htbp]
\centering
\includegraphics[width=\linewidth]{sources/monographs/proof-path-theory/docs/figures/fig_01_architecture.pdf}
\caption{Three different questions about the same declared archive. A certificate records derivability, local models record rule compatibility, and a task representation records recoverable information. None of these questions subsumes the others without additional hypotheses.}
\label{PPT:fig:architecture}
\end{figure}

Komorebi Space enters when local descriptions of a proof field must agree. The construction of local rule models supplies actual restriction maps, not merely a visual analogy between two graphs. A gluing theorem identifies a sufficient condition for assembling those models, and a counterexample explains why pairwise dependency checks can miss a conjunctive rule. In a separate linear setting, the size of a realization fiber measures exactly what a new observation distinguishes.

The implemented proof fields are finite and acyclic. Some structural statements, including rule soundness and the linear separation alternative, are given in greater generality where their hypotheses allow it. Extensions to infinite search, probabilistic inference, and higher categorical proof objects require their own assumptions. Background on formal proof systems may be found in \cite{PPT:buss,PPT:troelstra}.


\section{Claim types and logical boundaries}
\label{PPT:sec:old2}

A useful theory of proof paths must not turn every recorded object into a proof step. A notebook may contain a theorem, a conjecture, a failed analogy, a numerical experiment, and a historical interpretation on the same page. Their coexistence does not give them the same logical status.

\begin{PPTdefinition}[Evidence kind]
\label{PPT:res:2:1}
An evidence kind is a declared type attached to a record supporting a claim or justification. The core finite vocabulary used here is
\[
\begin{aligned}
&\text{formal proof},\quad \text{assumption},\quad \text{exact computation},\\
&\text{numerical evidence},\quad \text{conjecture},\quad \text{external result},\\
&\text{counterexample},\quad \text{interpretation}.
\end{aligned}
\]
An evidence record also states its source, scope, and whether it certifies the justification to which it is attached.
\end{PPTdefinition}

The word ``certifies'' is deliberately separate from the evidence label. An exact computation may certify a finite identity while providing only numerical evidence for a general statement. An external result may certify a step only on the domain covered by its hypotheses. The scope field is therefore part of the mathematical object.

\begin{PPTwarning}[Evidence types are not one ranking]
\label{PPT:res:2:2}
The evidence kinds do not form a universal total order. A formal proof of a weak statement and an exact computation of a stronger finite instance cannot be compared by one scalar without specifying the task. The ledger retains type and scope instead of replacing them with a single confidence number.
\end{PPTwarning}

This separation will be used throughout. Candidate edges may remain in the discovery field, but only certified edges enter a proof certificate.


\section{Typed proof fields}
\label{PPT:sec:old3}

The first formal object is a directed hypergraph. A justification may require several premises at once, and a claim may admit several alternative justifications.

\begin{PPTdefinition}[Typed proof field]
\label{PPT:res:3:1}
A finite typed proof field is a tuple
\[
\PPTPcal=(V,\PPTJcal,A,T,\tau,\PPTEcal,\mathbf v,\mathbf c),
\]
where:
\begin{itemize}
\item $V$ is a finite set of claim nodes.
\item $\PPTJcal$ is a finite set of justifications.
\item Each $j\in\PPTJcal$ has a finite premise set $\operatorname{prem}(j)\subseteq V$ and one conclusion $\operatorname{con}(j)\in V$.
\item $A\subseteq V$ is the set of required premises of the declared judgment.
\item $T\subseteq V$ is the nonempty target set.
\item $\tau$ assigns node and justification types.
\item $\PPTEcal$ is a typed evidence ledger.
\item $\mathbf v$ assigns nonnegative value coordinates to nodes.
\item $\mathbf c$ assigns nonnegative storage, audit, recovery, and cognitive costs.
\end{itemize}
The dependency graph contains an edge $u\to v$ whenever $u\in\operatorname{prem}(j)$ and $v=\operatorname{con}(j)$ for some justification $j$. It is required to be acyclic.
\end{PPTdefinition}

A single justification is conjunctive in its premise set. All listed premises must be available. Multiple justifications with the same conclusion are alternatives. This gives the familiar AND/OR structure of proof search without forcing it into an ordinary graph.

\begin{PPTdefinition}[Admissible leaf]
\label{PPT:res:3:2}
A node is an admissible leaf when it is a required premise, a declared axiom, or an explicitly admitted external result. Write $L$ for this set, so $A\subseteq L$. A certified justification with no premises can also derive a node, but it remains a justification and need not make that node a member of $L$. The distinction matters when extracting the smallest certificate.
\end{PPTdefinition}

Required premises are not necessarily the entire assumption context. Admitted axioms and external results remain visible in the supporting ledger. The notation $A\vdash B$ is relative to that declared background and does not erase those dependencies. A comparison between certificates must retain the assumptions under which each certificate is valid.

\begin{PPTdefinition}[Certified subfield]
\label{PPT:res:3:3}
The certified subfield $\PPTPcal^{\mathrm{cert}}$ contains every node of $\PPTPcal$ and only those justifications accepted by the declared certification policy. Evidence and scope are part of that policy. A certification flag records acceptance: it is not itself a proof that the rule is sound under an intended mathematical interpretation. That separate soundness contract is made explicit in Section~\ref{PPT:sec:old24}.
\end{PPTdefinition}

\begin{PPTdefinition}[Derivability]
\label{PPT:res:3:4}
Let $K\subseteq V$ be a selected node set and let $H\subseteq\PPTJcal$ be a selected justification set. The derivable set $D(K,H)$ is the least subset of $K$ satisfying:
\begin{enumerate}[label=\textup{(\roman*)}]
\item every selected admissible leaf belongs to $D(K,H)$.
\item if $j\in H$, $\operatorname{prem}(j)\subseteq D(K,H)$, and $\operatorname{con}(j)\in K$, then $\operatorname{con}(j)\in D(K,H)$.
\end{enumerate}
When only certified justifications are allowed, write $D_{\mathrm{cert}}(K,H)$.
\end{PPTdefinition}

\begin{PPTproposition}[Finite derivability]
\label{PPT:res:3:5}
For a finite acyclic proof field, $D(K,H)$ is uniquely determined and can be computed in one topological pass through the dependency graph.
\end{PPTproposition}

\begin{proof}
Acyclicity gives a topological ordering of the nodes. Begin with the selected admissible leaves. At each later node, test whether one selected incoming justification has all of its premises already derived. Add the node exactly when such a justification exists. No later step can change the status of an earlier node because all dependencies point forward. The resulting set satisfies the closure rules. Any set satisfying the rules must contain every node added by this pass, so the result is the least such set and is unique.
\end{proof}

\begin{figure}[htbp]
\centering
\includegraphics[width=\linewidth]{sources/monographs/proof-path-theory/docs/figures/fig_03_alternative_hypergraph.pdf}
\caption{Two certified routes in \protect\path{alternative_proofs.json}. Their stored node costs are one and three. A storage objective selects the short route, while a task requiring explanatory information can prefer the richer route. Alternative derivations are OR choices, not extra conjunctive premises.}
\label{PPT:fig:andor}
\end{figure}


\section{The ambient proof field and the meaning of failure}
\label{PPT:sec:old22}
\label{PPT:sec:ambient}

A finished proof is rarely the only route the researcher could have explored. An archive often contains sketches that were never completed, arguments that failed because of a specific hypothesis, conjectured bridges that remain untested, and genuinely refuted local constructions. It is tempting to call this whole archive a proof, but doing so would erase the distinction between a certified result and a record of exploration. The \emph{ambient proof field} provides a more precise language: it is the declared field of reachable or recorded mathematical states and typed candidate transitions in which certified proof paths form a distinguished subfield.

The finite hypergraph model permits \path{failed_route}, \path{boundary_case}, and \path{counterexample} nodes. Their presence is useful, but none is a logically justified global nonexistence claim by itself. A failed search is not necessarily a proof of impossibility. A counterexample to one proposed lemma is not necessarily a counterexample to the goal. An exact obstruction must have a declared scope and a checkable witness.

\begin{PPTdefinition}[Scoped branch assessment]
\label{PPT:res:22:1}
Let $\PPTPcal$ be a finite typed proof field with declared premises $A$ and target $t$. An exploratory branch is a finite sequence of candidate justifications whose certification and outstanding obligations are explicitly recorded. Its assessment is one of the following, relative to stated premises and rules:
\begin{enumerate}[label=\textup{(\roman*)}]
\item \emph{Open}: obligations remain unresolved.
\item \emph{Abandoned}: exploration stopped without a certificate of impossibility.
\item \emph{Locally refuted}: a checked counterexample or obstruction excludes a stated intermediate step under specified hypotheses.
\item \emph{Certified}: the selected justifications meet the declared proof certificate conditions.
\end{enumerate}
An assessment is a record of the branch, not a truth value assigned to every possible proof of $t$.
\end{PPTdefinition}

\begin{PPTproposition}[Local failure does not imply global failure]
\label{PPT:res:22:2}
A finite proof field may contain a locally refuted branch to target $t$ while also containing a certified derivation of $t$ from its declared premises.
\end{PPTproposition}
\begin{proof}
Take a premise $A$ and two alternative candidate derivations of $t$. Let the first use a proposed intermediate lemma $L$ that has a checked counterexample, so the corresponding justification is not certified. Let the second use a distinct certified justification from $A$ to $t$. The certified closure reaches $t$ by the second route. The failed first route remains a valid record of unsuccessful exploration, but cannot negate the existence of the second certificate. The same example can be represented by the alternative justification hypergraph developed earlier.
\end{proof}



The logical lesson is substantive for machine reasoning. A search timeout, budget exhaustion, invalid intermediate conjecture, or failure to find a proof all have different meanings. Only a certificate whose conclusion explicitly concerns nonexistence under a declared model and scope licenses a nonexistence claim. This is especially important when a failed branch is used to prune an automated search. Pruning the refuted step may be safe, while pruning its entire goal may discard a genuine proof.

\begin{PPTproposition}[Scoped refutation survives unchanged context]
\label{PPT:res:22:3}
Suppose a witness $w$ is verified by a sound checker to refute the predicate $L$ under assumptions $A$. In any later archive revision that retains the same interpretation of $L$, the same assumptions, and the checker's verification rules, $w$ remains a refutation of $L$. No conclusion follows about different predicates, altered assumptions, or alternative branches to a different goal.
\end{PPTproposition}
\begin{proof}
The verification of $w$ depends only on the declared predicate, assumptions, and checking rules. An append-only collection of unrelated exploratory records does not alter those inputs. Therefore the same verification remains valid. When any of those inputs changes, the earlier check no longer certifies the revised claim without an explicit comparison theorem.
\end{proof}

This is a qualified stability principle, not a general theorem that proof search is decidable. The append-only derivability result of Section~\ref{PPT:sec:versioned} concerns certified edges. The present proposition concerns the narrower persistence of a checked counterexample under a fixed interpretation. Both retain scope as part of the mathematical data.


\section{Certificates and verification closure}
\label{PPT:sec:old4}

\begin{PPTdefinition}[Certificate]
\label{PPT:res:4:1}
A certificate for a nonempty requested target set $T_0\subseteq V$ is a pair
\[
\PPTCcal=(K,H),
\qquad K\subseteq V,
\qquad H\subseteq\PPTJcal^{\mathrm{cert}},
\]
such that
\[
A\cup T_0\subseteq K
\qquad\text{and}\qquad
T_0\subseteq D_{\mathrm{cert}}(K,H).
\]
Order certificates componentwise: $(K,H)\leq(K^{\prime},H^{\prime})$ means $K\subseteq K^{\prime}$ and $H\subseteq H^{\prime}$. A least certificate is contained in every certificate for the same request. Minimizing only the node set, only the justification set, or a numerical cost is a different objective.
\end{PPTdefinition}

The usual target request is a subset of $T$, but an intermediate node may be requested for local analysis. Only a justification whose premises have been derived and whose conclusion is selected can contribute to the derivation. A record may retain other, inactive justifications without using them as proof steps. Extra records cannot lower a nonnegative additive cost, although zero costs can allow several minimizing records.

\subsection{The fixed AND-DAG case}

The simplest case occurs when each relevant node that is not an admissible leaf has one certified incoming justification. Backward expansion stops at admissible leaves. A proof of a statement that has already been admitted as a leaf is not required to use that statement.

\begin{PPTdefinition}[AND-DAG]
\label{PPT:res:4:2}
Starting at $T_0$, expand backward only at nodes outside $L$. The field is an AND-DAG relative to $T_0$ when every node reached outside $L$ has exactly one certified incoming justification. Every premise of that justification is then expanded in the same manner. The process may terminate at an admissible leaf or at a certified justification with no premises.
\end{PPTdefinition}

Let $\PPTAnc_L(T_0)$ denote the premise nodes reached by this backward expansion, stopping at $L$. Let $H_{\mathrm{core}}$ contain exactly the justifications selected at expanded nodes outside $L$.

\begin{PPTtheorem}[Unique verification core]
\label{PPT:res:4:3}
\label{PPT:thm:and-core}
Let $\PPTPcal$ be a finite acyclic AND-DAG relative to $T_0$. Put
\[
K_{\mathrm{core}}=A\cup T_0\cup\PPTAnc_L(T_0).
\]
Then $(K_{\mathrm{core}},H_{\mathrm{core}})$ is a certificate and is the least certificate under componentwise inclusion. In particular it is the unique inclusion-minimal certificate pair for this recorded inference problem.
\end{PPTtheorem}

\begin{proof}
Backward expansion terminates because the dependency graph is finite and acyclic. Every terminal node is either an admissible leaf or the conclusion of the selected justification with no premises. A forward topological pass therefore derives every expanded node. The required premises in $A$ are admissible and are retained even if the target does not use them. Thus the constructed pair is a certificate.

Now let $(K,H)$ be any certificate for the same request. It contains $A\cup T_0$. If a requested or required predecessor node is outside $L$, its derivation must use its unique certified incoming justification. That justification belongs to $H$, and every one of its premises belongs to $K$. Repeating the argument forces all of $\PPTAnc_L(T_0)$ and $H_{\mathrm{core}}$. At nodes in $L$ the argument stops. Hence $K_{\mathrm{core}}\subseteq K$ and $H_{\mathrm{core}}\subseteq H$, which proves the least-pair assertion.

Minimizing one coordinate alone does not make the whole record unique. One may add an irrelevant node without changing the least justification set, or retain an unused rule without changing the least node set. The componentwise statement avoids this ambiguity.
\end{proof}

\begin{PPTcorollary}[Linear extraction]
\label{PPT:res:4:4}
With indexed certified adjacency and constant-time membership bookkeeping, the core can be extracted in time linear in $|A|$ plus the number of visited nodes, selected justifications, and premise incidences.
\end{PPTcorollary}

This bound concerns the indexed graph algorithm. It is not a verified running-time bound for every scanning routine in the reference Python package.

\begin{figure}[htbp]
\centering
\includegraphics[width=\linewidth]{sources/monographs/proof-path-theory/docs/figures/fig_02_and_dag_core.pdf}
\caption{The canonical \protect\path{and_dag.json} example. The AND gate requires both $A$ and $D$ for $L_1$. Its least certificate contains five nodes and four certified rules. The gold branch is certified but unnecessary for $B$. The dashed route to $F$ is uncertified and cannot enter the certificate.}
\label{PPT:fig:andcore}
\end{figure}

\begin{PPTremark}
\label{PPT:res:4:5}
The theorem identifies the minimum certificate for the recorded derivation. It does not claim that the derivation is globally shortest among all possible proofs of the theorem. A different proof field may have a smaller core.
\end{PPTremark}


\section{Alternative derivations and certificate complexity}
\label{PPT:sec:old5}

When a claim admits several certified justifications, certificate extraction is no longer a backward closure problem. A local choice may be reused by several later obligations, which couples decisions across the field.

\begin{PPTproblem}[Minimum Certificate Decision]
\label{PPT:res:5:1}
Given a finite acyclic certified proof field, a target set, nonnegative integer node and justification costs, and a budget $B$, determine whether a certificate of total cost at most $B$ exists.
\end{PPTproblem}

\begin{PPTtheorem}[Minimum certificate selection is NP-complete]
\label{PPT:res:5:2}
\label{PPT:thm:minimum-certificate}
Minimum Certificate Decision is NP-complete, even when every nonzero cost is a unit node cost and all justification costs are zero.
\end{PPTtheorem}

\begin{proof}
The problem belongs to NP. A proposed node and justification set can be checked in polynomial time by verifying references, computing derivability in topological order, and summing costs.

For NP-hardness, reduce Set Cover. Let $U=\{u_1,\ldots,u_m\}$ be the universe, let $S_1,\ldots,S_n\subseteq U$ be the available sets, and let $k$ be the cover budget. Construct one optional admissible leaf node $s_i$ for each set $S_i$, with unit cost. Construct one zero-cost obligation node $e_j$ for each element $u_j$. Whenever $u_j\in S_i$, add a certified justification
\[
\{s_i\}\Longrightarrow e_j.
\]
Add a zero-cost target node $t$ with the certified justification
\[
\{e_1,\ldots,e_m\}\Longrightarrow t.
\]
The dependency graph is acyclic. A certificate for $t$ must derive every $e_j$. Deriving $e_j$ requires at least one selected source node $s_i$ with $u_j\in S_i$. One selected $s_i$ may discharge several obligations, and its cost is paid once. Therefore a certificate of cost at most $k$ exists exactly when the original Set Cover instance has a cover of size at most $k$. The reduction is polynomial.
\end{proof}

\begin{figure}[htbp]
\centering
\includegraphics[width=\linewidth]{sources/monographs/proof-path-theory/docs/figures/fig_04_certificate_complexity.pdf}
\caption{The Set Cover example as a membership matrix. Filled marks belong to the selected rows $S_1$ and $S_3$; open marks show unselected memberships. Every obligation is covered. The target requires all four obligations through one AND rule. The exact minimum is two, although the minimizing pair need not be unique.}
\label{PPT:fig:setcover}
\end{figure}

\begin{PPTremark}[Algorithmic consequence]
\label{PPT:res:5:3}
Exhaustive certificate selection is useful for small fields or fields with few alternatives. The reference implementation uses branch and bound with an explicit search-state limit. Its result is optimal for the stored finite cost data when the search completes. It does not remove the exponential worst case. The decision theorem concerns integer costs, while the Python interface accepts finite floating-point costs and compares their exact stored binary values as rationals. Printed cost totals remain floating-point summaries. The input conversion does not recover an intended real or decimal value that was never represented. See \cite{PPT:garey-johnson} for the background complexity language.
\end{PPTremark}


\section{Certified paths as a category}
\label{PPT:sec:old6}

The hypergraph model records local justifications. A complementary construction treats complete certified transformations as morphisms. It uses the usual category of paths \cite{PPT:maclane}.

\begin{PPTdefinition}[Certified proof category]
\label{PPT:res:6:1}
Fix a directed graph of typed mathematical states and certified primitive transformations. Its path category $\mathsf{Cert}$ has these states as objects and finite composable paths as morphisms. The empty path is an identity, and composition is concatenation. One may further identify paths by a declared congruence compatible with concatenation. Associativity and the identity laws then hold by construction.
\end{PPTdefinition}

A state may be a proposition, a diagram, an exact finite computation, a coordinate representation, or a structured collection of claims. Candidate analogies may live in a larger discovery graph, but they do not enter $\mathsf{Cert}$ until their logical force and domain have been certified.

\begin{PPTproposition}[Functorial transport]
\label{PPT:res:6:2}
Let $F:\mathsf{Cert}_1\to\mathsf{Cert}_2$ be a functor. If
\[
A\xrightarrow{f_1}S_1\xrightarrow{f_2}\cdots\xrightarrow{f_r}B
\]
is a certified path in $\mathsf{Cert}_1$, then
\[
F(A)\xrightarrow{F(f_1)}F(S_1)\xrightarrow{F(f_2)}\cdots\xrightarrow{F(f_r)}F(B)
\]
is a certified path in $\mathsf{Cert}_2$.
\end{PPTproposition}

\begin{proof}
A functor maps objects to objects and morphisms to morphisms. It preserves identities and composition. The image of the composite path is therefore the composite of the image morphisms, which is certified in the target category.
\end{proof}

The proposition transports certification inside the two declared categories. To interpret that transport as preservation of a mathematical judgment, the target category and functor must also have sound interpretations of their objects and rules. These requirements cannot be inferred merely from the word functor. Under that contract, different representations can preserve a judgment while exposing different features of its argument. The category language is standard \cite{PPT:maclane}.

\begin{PPTdefinition}[Compositional certificate map]
\label{PPT:res:6:3}
A compositional certificate map is a functor
\[
\PPTCcal:\mathsf{Expanded}\longrightarrow\mathsf{Certificate}
\]
that sends expanded certified paths to checkable certificate paths and preserves identity and composition.
\end{PPTdefinition}

\begin{PPTproposition}[Local audit composes]
\label{PPT:res:6:4}
If every primitive morphism in a finite path has an auditable certificate and the certificate map is compositional, then the composite path is auditable from the primitive certificates and the declared composition rules.
\end{PPTproposition}

\begin{proof}
Proceed by induction on path length. The length-zero case is an identity. Assume the first $r$ morphisms have an auditable composite certificate and append one certified morphism. Functoriality identifies the new certificate with the composition of the earlier certificate and the primitive certificate of the appended morphism. The declared composition rule completes the audit.
\end{proof}



\chapter{Coherent realization and information}\label{PPT:book:chapter-51}
Compatibility, extension, obstruction, and ambiguity become concrete once the views and protected observations are specified. This part gives the actual bridge to Komorebi Space.

  \PPTopeningfigure{fig_17_binary_cube.pdf}{Eight vertices encode the states of $\mathbb F_2^3$. Four satisfy $a+b+c=0$. Protecting $a=0$ leaves the two green vertices $000$ and $011$. The cube edges show single-coordinate changes, not continuous solution lines or proof steps.}{fig:binary-cube}
  \vspace{3pt}{\sffamily\bfseries\color{KomorebiTeal}The governing distinction}\par
  Existence asks whether the fiber is empty. Ambiguity asks what remains inside it. An informative observation distinguishes points of that fiber rather than merely producing more output coordinates.
 


\section{Komorebi Space, realization fibers, and informative questions}
\label{PPT:sec:old23}
\label{PPT:sec:komorebi-bridge}

Komorebi Space supplies a natural setting for studying the proof field's coherent realizations. Here the connection is \emph{relative} to an explicitly chosen system of views. Nothing in a finite proof hypergraph canonically supplies one universal category of all proofs, and a formal proof term is not automatically equivalent to an intended semantic object. The view system must be stated, together with its restrictions and the observations that are required to survive.

\begin{PPTdefinition}[Coherent views of a proof field]
\label{PPT:res:23:1}
Choose a finite collection of views $V$, fibers $F_v$ of local proof states or annotations, and for each declared directed edge $e:v\to w$ a restriction map $e^*:F_w\to F_v$. The \emph{strict realization set} is
\[
\Gamma(F)=\big\{(x_v)\in\prod_{v\in V}F_v : e^*(x_w)=x_v\text{ for every declared edge }e\big\}.
\]
The \emph{extension fiber} of prescribed information $s$ on $U\subseteq V$ consists of those global coherent families whose restrictions equal $s$. If the edges form a category, the usual functorial composition equations are additional explicit obligations.
\end{PPTdefinition}

The finite implementation \path{ambient.py} realizes this edge-defined model exactly. The gluing condition should not be weakened to covering the objects alone: a compatibility equation is attached to an edge and must appear in at least one patch.

\begin{PPTtheorem}[Finite arrow-cover gluing]
\label{PPT:res:23:2}
Suppose patches $U_i\subseteq V$ cover every view and every declared edge has both endpoints in some $U_i$. Let $s_i$ be coherent assignments on $U_i$ that agree on all overlaps. There is a unique global coherent assignment $s\in\Gamma(F)$ restricting to every $s_i$.
\end{PPTtheorem}
\begin{proof}
For each view $v$, choose any patch containing it and define $s(v)$ to be the value supplied by that patch. Agreement on overlaps makes this independent of the choice. For a declared edge $e:v\to w$, choose a patch containing both endpoints. The local assignment on that patch obeys $e^*s(w)=s(v)$, so the global assignment is coherent. It is the only possible glue because every coordinate was prescribed by at least one patch.
\end{proof}

\begin{PPTexample}[Object coverage cannot replace arrow coverage]
\label{PPT:res:23:3}
Let $V=\{v,w\}$ with an edge $v\to w$, both fibers $\{0,1\}$, and identity restriction. Singleton patches prescribe $s(v)=0$ and $s(w)=1$. These assignments cover all views and agree on their empty overlap. There is no global coherent assignment with those coordinates. The missing ingredient is the edge condition, not the availability of the two individual values.
\end{PPTexample}

The theorem is a finite instance of the strict gluing construction in \emph{Komorebi Space} \cite{PPT:komorebi-space}. It does not claim a general stack theorem for nonabelian proof objects. Gluing up to isomorphism requires the identifications and their coherence witnesses, which are retained in the corresponding groupoid setting.

\subsection{Exact finite realization and obstruction certificates}

A second connection is quantitative. Work over $\mathbb F_2$ with $E=\mathbb F_2^n$. Let $A:E\to\mathbb F_2^m$ impose coherence, $P:E\to\mathbb F_2^p$ retain prescribed information, and $x\in E$ be a raw candidate. Define
\[
\mathfrak K_P(x)=\{y\in E:Ay=0,\ Py=Px\},\qquad
M=\begin{bmatrix}A\\P\end{bmatrix},\quad
b=\begin{bmatrix}0\\Px\end{bmatrix}.
\]
This is a finite linear realization fiber. The equation $My=b$ does not by itself produce a solution. Compatibility of the prescribed information must first be established.

\begin{PPTtheorem}[Finite realization alternative and ambiguity]
\label{PPT:res:23:4}
Exactly one of the following holds:
\begin{enumerate}[label=\textup{(\roman*)}]
\item $My=b$ has a solution $y_0$. Its complete solution set is $y_0+H$, where $H=\ker A\cap\ker P$, and
\[
|\mathfrak K_P(x)|=2^{n-\operatorname{rank}(M)}.
\]
\item A row vector $\lambda$ exists with $\lambda M=0$ and $\lambda b=1$. It is a directly checkable certificate that $\mathfrak K_P(x)=\varnothing$.
\end{enumerate}
\end{PPTtheorem}
\begin{proof}
If $My_0=b$, any other solution differs from $y_0$ by an element of $\ker M=H$, and every such difference gives another solution. Rank nullity gives $\dim H=n-\operatorname{rank}(M)$. If $b$ is not in the column space of $M$, linear algebra supplies a linear functional that vanishes on that column space and takes value $1$ on $b$. Expressing it as a row vector gives $\lambda$. Conversely such a witness contradicts any proposed equation $My=b$, since it would imply $0=\lambda My=\lambda b=1$. These cases are mutually exclusive and exhaustive over a field.
\end{proof}

This negative certificate is a concrete instance of the protected obstruction principle from Komorebi Space. Its scope is exactly the finite linear system $My=b$. It is not a certificate that some broad mathematical conjecture is unprovable.

\subsection{The information gained by a new observation}

Suppose the realization fiber is nonempty and a new linear observation $Q:E\to\mathbb F_2^q$ becomes available. The right question is not how many bits appear in the output of $Q$ but how many remaining degrees of freedom in the fiber it distinguishes.

\begin{PPTtheorem}[Exact observation gain]
\label{PPT:res:23:5}
Give the nonempty finite affine fiber $y_0+H$ the uniform distribution. Then the Shannon entropy of the observed value $Qy$ is
\[
H(Qy)=\operatorname{rank}(Q|_H)\quad\text{bits}.
\]
For every observation value that actually occurs, the corresponding conditional fiber contains
\[
2^{\dim H-\operatorname{rank}(Q|_H)}
\]
realizations. Thus the uncertainty reduction is exactly $\operatorname{rank}(Q|_H)$ bits.
\end{PPTtheorem}
\begin{proof}
The image $Q(H)$ is a vector subspace of dimension $r=\operatorname{rank}(Q|_H)$, so the observed values form the affine set $Qy_0+Q(H)$ of size $2^r$. Every such value has a preimage that is a coset of $H\cap\ker Q$, whose dimension is $\dim H-r$ by rank nullity. Equal fiber cardinalities make the induced output distribution uniform on $2^r$ values. Its entropy is $\log_2(2^r)=r$, and each conditional fiber has cardinality $2^{\dim H-r}$.
\end{proof}

\begin{PPTexample}[One informative question]
\label{PPT:res:23:6}
Let $y=(a,b,c)\in\mathbb F_2^3$ with $a+b+c=0$ and protected value $a=0$. The realization fiber is $\{(0,0,0),(0,1,1)\}$. It has one bit of ambiguity. Observing $b$ resolves that bit exactly. Observing $a$ supplies no new information, despite being a perfectly valid measurement. This difference is measured by the rank of the observation restricted to $H$, rather than by the output width of the observation.
\end{PPTexample}

The module \path{realization.py} computes these quantities through exact binary row reduction, and returns an explicit linear witness when the protected realization is impossible. The program does not enumerate all solutions, although the tests compare its results with exhaustive enumeration on small cases. The entropy statement uses Shannon entropy \cite{PPT:cover-thomas} and assumes uniform sampling of the compatible realizations. It does not assert a universal prior over mathematical discoveries, or that the most informative observation is automatically the most valuable research question once costs and task priorities are included.


\section{A checked realization bridge for proof fields}
\label{PPT:sec:old24}
\label{PPT:sec:checked-bridge}

The ambient interpretation can be made into a mathematical construction rather than left as an analogy. The important step is to specify what a local realization means. A node label may describe an arbitrary mathematical sentence, but a Boolean value attached to that node does not automatically evaluate that sentence. In this section it records membership in a set closed under the declared certified rules. This restricted interpretation has an exact soundness and completeness theorem. An interpretation in actual mathematical propositions is a separate contract.

\subsection{The declared rule system and its models}

Fix a node set $V$, a set $L\subseteq V$ of admissible leaves, and a set $J$ of justifications. A justification $j$ has a premise set $A_j$, a conclusion $c_j$, and a certification predicate. Define its support by
\[
S_j=A_j\cup\{c_j\}.
\]
A set $X\subseteq V$ is a \emph{model of the declared field} when $L\subseteq X$ and, for every certified $j$,
\[
A_j\subseteq X\quad\Longrightarrow\quad c_j\in X.
\]
Each justification is one implication with conjunctive premises. Alternative justifications remain separate implications. Uncertified attempts and failed-route annotations do not contribute a deduction rule merely by appearing in the archive.

Let $\operatorname{Der}_F$ be the inductively generated set obtained from the leaves by certified rules. For a finite field, repeated forward closure computes this set. It is not necessary to enumerate all valuations to perform that computation.

Every declared field of this positive-implication form has a model: the full set $V$. It contains every leaf and every rule conclusion. Nonexistence in the following extension problems therefore concerns prescribed local values or additional preservation requirements, not this unrestricted model space. This does not make every node derivable. Derivability is membership in the least model, not membership in some model.

\begin{PPTtheorem}[Soundness and relative completeness]
\label{PPT:res:24:1}
\label{PPT:thm:relative-completeness}
For every node $v$,
\[
v\in\operatorname{Der}_F
\quad\Longleftrightarrow\quad
\text{$v\in X$ for every model $X$ of the declared field}.
\]
If $v$ is not derivable, $\operatorname{Der}_F$ itself is a model omitting $v$.
\end{PPTtheorem}
\begin{proof}
Every model contains the leaves. Induction on a derivation shows that it contains the conclusion of every certified application whose premises have already been derived. Thus every model contains $\operatorname{Der}_F$. Conversely, the inductive constructors show that $\operatorname{Der}_F$ contains the leaves and is closed under every certified rule. It is therefore a model. Applying the universal condition to this particular model proves the reverse implication, and the same model gives the stated countermodel.
\end{proof}

This is relative completeness for the declared rule system, not completeness of arbitrary mathematical proof search. A missing theorem may simply require a rule or a lemma absent from the field. A failed search can therefore produce a useful countermodel to a proposed finite inference architecture without disproving the mathematical target.

For an intended interpretation $\mu:V\to\mathrm{Prop}$, the transfer to mathematical truth has explicit hypotheses: every leaf must satisfy $\mu$, and every certified rule must preserve $\mu$ from all its premises to its conclusion. Under those hypotheses, induction gives $v\in\operatorname{Der}_F\Rightarrow\mu(v)$. Lean checks this statement as \path{ProofPath.Field.semantic_soundness}. It does not infer those hypotheses from a JSON certification flag.

\subsection{From patches to a contravariant view system}

For a patch $U\subseteq V$, let $\mathcal M_F(U)$ consist of subsets $X_U\subseteq U$ which contain $L\cap U$ and obey every certified rule whose entire support lies in $U$. If $U\subseteq W$, define
\[
\rho_{U,W}:\mathcal M_F(W)\longrightarrow\mathcal M_F(U),
\qquad \rho_{U,W}(X_W)=X_W\cap U.
\]
This map is well defined. A rule supported in $U$ is also supported in $W$, and its conclusion remains in $U$. The restriction therefore preserves both the local leaves and the local implication conditions.

\begin{PPTproposition}[The proof-field view laws]
\label{PPT:res:24:2}
The restriction maps satisfy
\[
\rho_{U,U}=\mathrm{id},\qquad
\rho_{U,W}\rho_{W,Z}=\rho_{U,Z}\quad(U\subseteq W\subseteq Z).
\]
Consequently $U\mapsto\mathcal M_F(U)$ is a contravariant set-valued view system on the inclusion order of patches.
\end{PPTproposition}
\begin{proof}
Both laws are the corresponding intersection identities. For the second, $(X\cap W)\cap U=X\cap U$ because $U\subseteq W$. The values outside a patch are absent, or equivalently normalized to false in the Lean representation. There are no additional hidden coordinates.
\end{proof}

The strict Komorebi realization is now an actual set of coherent local models,
\[
\Gamma(\mathcal M_F)=
\left\{(X_U)_U:\rho_{U,W}(X_W)=X_U\text{ whenever }U\subseteq W\right\}.
\]
There is an explicit bijection
\[
\operatorname{Models}(F)\simeq\Gamma(\mathcal M_F),
\qquad X\longmapsto(X\cap U)_U.
\]
Its inverse evaluates a coherent family at the full patch $V$. Coherence shows that every other component is precisely the restriction of that value. The Lean declaration \path{ProofPath.Field.modelEquivSections} constructs this equivalence, using \path{ViewDiagram} and the actual local restriction maps.

The full-patch equivalence identifies the ambient construction. It does not prove that arbitrary prescribed proper-patch data extend: the full patch already contains the global problem. Extension from smaller views is the distinct content of the following theorem.


\PPTstudyhead{01}{The bridge is a restriction map}
For one certified rule $\{a,b\}\to c$ and no admissible leaves ($L=\varnothing$), there are seven global Boolean rule models. Only $(1,1,0)$ violates the implication. On either two-node patch below, the complete support is absent, so all four local valuations are admissible.

\begin{figure}[H]\centering
\includegraphics[width=\linewidth]{sources/monographs/proof-path-theory/docs/figures/fig_16_restriction_bridge.pdf}
\caption{The local model sets and their actual restrictions. Each arrow is $\rho_{U,W}(X)=X\cap U$. The two ways of restricting a global model to $\{c\}$ agree. The diagram is counted directly from the finite rule model, with seven global models, four models on each two-node patch, and two overlap values.}
\label{PPT:fig:restriction-models}
\end{figure}

{\sffamily\bfseries\color{KomorebiTeal}What the diagram establishes}\par\smallskip
It displays the contravariant view construction and the consistency of restriction. The full-patch equivalence identifies global rule models with their coherent families of restrictions.

\medskip
{\sffamily\bfseries\color{KomorebiGold}What the diagram does not establish}\par\smallskip
A pair of matching proper-patch models need not be in the image of a global model. That extension question is controlled by the support of the rule, as the next study shows.

\subsection{The correct coverage condition is a whole support}

Let $(U_i)_{i\in I}$ cover the nodes. A family $X_i\in\mathcal M_F(U_i)$ is matching when
\[
X_i\cap U_j=X_j\cap U_i\qquad\text{for every }i,j.
\]
This equality includes agreement about false values, not merely common true values.

\begin{PPTtheorem}[Whole-support gluing]
\label{PPT:res:24:3}
\label{PPT:thm:support-gluing}
Suppose every certified justification has its whole support in some patch:
\[
\forall j\text{ certified},\quad\exists i\quad S_j\subseteq U_i.
\]
Then every matching family $(X_i)$ has a unique global field model $X$ satisfying $X\cap U_i=X_i$ for all $i$.
\end{PPTtheorem}
\begin{proof}
Set $X=\bigcup_iX_i$. Matching and node coverage imply $X\cap U_i=X_i$. A leaf belongs to some patch and hence to its local model, so $L\subseteq X$. Consider a certified rule whose premises belong to $X$. Choose a patch containing its entire support. Matching implies that every premise is true in that patch's local model, whose rule condition forces the conclusion. Hence $X$ obeys all certified rules. Any other global model with the specified restrictions agrees with $X$ at every node because the patches cover the nodes. This proves uniqueness.
\end{proof}

\begin{PPTexample}[Pairwise dependencies do not suffice]
\label{PPT:res:24:4}
Take nodes $a,b,c$, no admissible leaves ($L=\varnothing$), and one certified rule
\[
\{a,b\}\longrightarrow c.
\]
Use patches $U_0=\{a,c\}$ and $U_1=\{b,c\}$. The local valuations $X_0=\{a\}$ and $X_1=\{b\}$ agree on their overlap: both assign false to $c$. Both are local models, since neither patch contains the full rule support. They cover every node and both pairwise premise-to-conclusion dependencies. Nevertheless, a global model realizing them would contain $a,b$ and omit $c$, contradicting the rule. Thus replacing a conjunctive hyperedge by unrelated pairwise checks can lose precisely the constraint that matters.
\end{PPTexample}

The example is kernel checked in \path{ProofPath.Examples}. It does not contradict the earlier arrow-cover theorem. That theorem applies to a problem already specified by binary restriction equations. A conjunctive proof rule is a different constraint until a correct representation of its whole support has been supplied. The full support may be redundant in particular fields, so the theorem gives a sufficient cover condition rather than a necessary one for every individual diagram. Section~\ref{PPT:sec:semantic-cover} gives an exact entailment criterion for those redundancies.

The Python class \path{ProofFieldViews} implements these definitions. It freezes the rule and leaf contents at construction, restricts local valuations by literal intersection, checks support coverage before gluing, and can form a finite intersection-closed diagram for \path{FiniteViewSystem}. Local enumeration has explicit size limits. Forward closure and declared-field countermodels do not require that enumeration.


\PPTstudyhead{02}{Agreement can hide a missing rule}
Keep the same rule and choose the local data $X_0=\{a\}$ on $U_0=\{a,c\}$ and $X_1=\{b\}$ on $U_1=\{b,c\}$. Both local models assign false to their shared coordinate $c$. They match exactly, yet their only assembled valuation violates the global rule.

\begin{figure}[H]\centering
\includegraphics[width=\linewidth]{sources/monographs/proof-path-theory/docs/figures/fig_15_whole_support.pdf}
\caption{The explicit counterexample to replacing whole-support coverage with pairwise dependency coverage. The local data force $(a,b,c)=(1,1,0)$, so both premises hold and the conclusion fails. Green nodes carry value one and open nodes carry value zero.}
\label{PPT:fig:whole-support-example}
\end{figure}

{\sffamily\bfseries\color{KomorebiTeal}The successful gluing argument}\par\smallskip
A patch containing the entire support $\{a,b,c\}$ must enforce the conjunctive rule. If all local models match and every certified support occurs in a patch, their union obeys every rule. Node coverage then gives uniqueness. This is the content of Theorem~\ref{PPT:thm:support-gluing}, not an inference from pairwise agreement alone.

\medskip
The cover condition is sufficient for all matching data. It need not be necessary for an individual field whose other rules already make a support constraint redundant.

\subsection{Protected realization over a field}

The exact binary model admits a general algebraic formulation. Let $K$ be a field, let $E,R,S$ be $K$-vector spaces, and let $D:E\to R$ and $p:E\to S$ be linear. The prescribed problem is to repair $x$ while retaining $p(x)$:
\[
Dy=0,\qquad p(y)=p(x).
\]
The admissible correction image is $B=D(\ker p)$. The elementary correction equivalence gives feasibility exactly when $Dx\in B$.

\begin{PPTtheorem}[A linear certificate alternative]
\label{PPT:res:24:5}
\label{PPT:thm:linear-alternative}
The protected problem is infeasible if and only if there is a linear functional $\lambda:R\to K$ satisfying
\[
\lambda(Dh)=0\quad\text{whenever }p(h)=0,
\qquad \lambda(Dx)=1.
\]
For a finite matrix $M$ and right side $b$, the corresponding alternative is
\[
\nexists y\ (My=b)
\quad\Longleftrightarrow\quad
\exists\ell\ (\ell M=0\ \text{and}\ \ell b=1).
\]
\end{PPTtheorem}
\begin{proof}
If $Dx\notin B$, linear separation of a vector from a subspace gives a functional vanishing on $B$ and nonzero at $Dx$. One construction extends a basis of $B$ by $Dx$ and then to a basis of $R$. Assign value zero to the basis of $B$ and value one to $Dx$, and extend linearly. Equivalently, normalize any separating functional by its nonzero value. Conversely, if $y$ were feasible, then $h=x-y$ would satisfy $p(h)=0$ and $Dh=Dx$, contradicting the displayed functional equations. For a matrix, take $B$ to be its column image. In finite coordinates the separating functional is a row vector, giving the stated alternative.
\end{proof}

The general existence proof is classical. It does not claim an efficient procedure for arbitrary infinite spaces. In the implemented binary case, Gaussian elimination constructs a row witness, and the witness is checked against the original matrix and right side. A positive solution and a normalized negative witness cannot both verify.


\PPTstudyhead{03}{An obstruction is a checkable equation}
In the binary example from \path{data/bridge_examples.json}, compatibility demands $a+b=0$ while the protected observations require $a=0$ and $b=1$. The variable $c$ is unconstrained. Its freedom cannot remove the incompatible requirement on the first two coordinates.

\begin{figure}[H]\centering
\includegraphics[width=\linewidth]{sources/monographs/proof-path-theory/docs/figures/fig_19_obstruction_witness.pdf}
\caption{A normalized negative witness in $\mathbb F_2$. Adding the three rows gives zero on the left and one on the right. The witness is $\ell=(1,1,1)$, which satisfies $\ell M=0$ and $\ell b=1$ for the original stacked system.}
\label{PPT:fig:obstruction-witness}
\end{figure}

The exact system and witness can be inspected without trusting a search procedure:
\[
 M=\begin{pmatrix}1&1&0\\1&0&0\\0&1&0\end{pmatrix},\qquad
 b=\begin{pmatrix}0\\0\\1\end{pmatrix},\qquad
 \ell M=(0,0,0),\quad \ell b=1.
\]
If $My=b$ had a solution, multiplying by $\ell$ would give $0=1$. This is a certificate for the declared protected problem. A search that merely stops supplies no comparable implication.

\subsection{Counting what remains after an observation}

Let $T:E\to R$ and $Q:E\to W$ be linear, and fix a known point $x_0$. Put
\[
\mathcal A=\{y:Ty=Tx_0\},\quad H=\ker T,\quad
Q_H=Q|_H.
\]
Translation $h\mapsto x_0+h$ identifies $H$ with $\mathcal A$. The observation image is the translate $Qx_0+Q(H)$, and the fiber above the attained value $Qx_0$ is a translate of $\ker Q_H$.

\begin{PPTtheorem}[Exact observation counts]
\label{PPT:res:24:6}
\label{PPT:thm:checked-counts}
Suppose $K$ is a finite field of cardinality $q$ and $E$ is a finite-dimensional $K$-vector space. Set
$k=\dim_K H$ and $r=\dim_K Q(H)$. Then
\[
|\mathcal A|=q^k,\qquad
|Q(\mathcal A)|=q^r,\qquad
|\{y\in\mathcal A:Qy=Qx_0\}|=q^{k-r}.
\]
In particular the exact count identity is
\[
|\mathcal A|=
|\{y\in\mathcal A:Qy=Qx_0\}|\,q^r.
\]
Every attained observation fiber has the same size. For the uniform distribution on $\mathcal A$, the observation removes exactly $r\log_2q$ bits of uncertainty.
\end{PPTtheorem}
\begin{proof}
The two translation equivalences reduce the counts to finite vector spaces. Such a space of dimension $d$ has $q^d$ elements. Rank-nullity for $Q_H$ gives $\dim\ker Q_H+r=k$, proving the product identity and all the counts. An attained value can be treated by choosing a point attaining it as the base point. Equal fiber sizes imply a uniform observation distribution and a uniform conditional distribution on each fiber. The entropy difference is therefore $\log_2(q^k)-\log_2(q^{k-r})=r\log_2q$.
\end{proof}

Lean checks the translation equivalences, rank-nullity specialization, fiber cardinalities, and product identity in \path{ProofPath.Protected}, with the binary specialization in \path{ProofPath.Binary}. The logarithmic entropy interpretation is the written corollary above, not a newly compiled theorem about a Lean probability distribution. Uniformity is substantive. With a nonuniform prior, dimension alone does not determine Shannon information.

For the protected problem, set $T=(D,p)$ and choose an actual feasible $y_0$, so $Ty_0=(0,px)$. This reduces its solution fiber to the theorem without changing the protected values. A concrete binary example has $a+b+c=0$ and $a=0$. Its two realizations are $(0,0,0)$ and $(0,1,1)$. Observing $b$ has rank one on the ambiguity space and determines the remaining coordinate. Repeating the same observation in two output rows does not double the information.


\PPTstudyhead{04}{Ask what an observation distinguishes}
Return to the two vertices $000$ and $011$ from the opening cube. The ambiguity space is the span of $011$. Reading the protected coordinate $a$ repeats known information. Reading $b$ separates the two points. Reading $b$ twice creates two output rows but no additional distinction.

\begin{figure}[H]\centering
\includegraphics[width=\linewidth]{sources/monographs/proof-path-theory/docs/figures/fig_18_observation_partitions.pdf}
\caption{Three observation maps on the same feasible binary fiber. The last column is the size of each attained conditional fiber. Gains of zero, one, and one bit follow from the rank of the map restricted to $H$, under the uniform distribution. The duplicated output is not twice as informative.}
\label{PPT:fig:observation-partitions}
\end{figure}

For example, the repeated observation has matrix
\[
 Q=\begin{pmatrix}0&1&0\\0&1&0\end{pmatrix}.
\]
Its two rows are identical. On $H=\langle011\rangle$ it has rank one, so $|\mathcal A|=2$, $|Q(\mathcal A)|=2$, and every attained observation fiber has size one. These are count identities before they are interpreted probabilistically.

\medskip
The general finite-field formula in Theorem~\ref{PPT:thm:checked-counts} retains the same separation: $q^k$ initial possibilities, $q^r$ attained outcomes, and $q^{k-r}$ realizations per outcome.

\subsection{Executable witnesses and a native declaration export}

Four exact examples are regenerated from the Python implementation: a one-bit ambiguity, an incompatible protected assignment, a redundant observation, and a uniquely determined realization. The generator emits the original stacked matrices and right sides alongside the proposed solutions or row witnesses. It also emits Lean declarations which check those concrete equations and cardinalities by kernel reduction. A faulty generator cannot make a false displayed proposition pass merely by assigning a certification flag. The generator itself is not thereby verified for arbitrary inputs, and the correspondence between the supplied JSON input and its emitted Lean syntax remains an audited software boundary.

A separate native exporter reads five already compiled Lean declarations, extracts their direct type and proof-constant dependencies, and records the transitive axiom sets. Dependencies outside the selected roots are explicit external leaves. This yields an actual proof-field JSON record rather than a hand-authored imitation of one. The exported field is then loaded through the public Python importer, and its declared closure is checked to contain its roots.

The export has a deliberately narrow meaning. It is syntactic dependency data for selected declarations, not a reconstruction of the entire discovery process. It does not recover failed attempts, prove that every constant in a type was a historically essential idea, or formally verify the metaprogram that extracts the data. The mathematical proof-field bridge is checked independently of that metaprogram. Its semantic interpretation retains the explicit hypotheses in Theorem~\ref{PPT:thm:relative-completeness}.

\subsection{Verification scope and reproducibility}

The present release freshly builds eighteen production Lean modules in a new object directory. Each module is replayed through Lean's own kernel against the pinned imported environment. The nine dependency revisions and the compiler hash are checked before building. All explicit named source declarations receive an axiom audit against the declared whitelist. The thirteen preceding mathematical modules are unchanged. Five new modules add executable trace checking, finite witness completeness, a typed JSON codec, and checked exchange examples.

Ten deliberately invalid files test false conclusions, missing support, placeholders, invalid decoders, merged observations, forged or incomplete closures, and a Boolean used as a node index. A further file introduces an explicit false axiom. Lean correctly accepts a consequence relative to that axiom, while the release's separate axiom policy rejects the dependency. A missing import, unavailable decision procedure, or timeout does not count as a mathematical rejection.

The fresh record is in \path{evidence/formal_1.8.0/}. The older record remains separately identified. Kernel replay here is not an independent implementation of Lean and does not rebuild or replay the entire cached Mathlib library. Written proofs, general Lean statements, generated concrete examples, and Python tests remain different evidence categories. In particular, a sound mathematical algorithm does not make its Python implementation universally verified.

\subsection{Protected extensions of the actual rule field}
\label{PPT:sec:protected-rule-extension}

The linear model is not the only setting in which an obstruction can be returned explicitly. The original rule field has its own constructive extension problem. A partial Boolean description records which nodes must belong to a model and which must remain absent. These are observations about the declared model, not assertions that the sentences attached to the nodes are true or false in an intended mathematical universe.

Fix the field $F$, including its admitted leaves and certified rules. For a positive observation set $S\subseteq V$, let
\[
  \operatorname{cl}_F(S)=\operatorname{Der}_{F[S]},
\]
where $F[S]$ is obtained by admitting $L\cup S$ as leaves and changing nothing else. For a negative observation set $N\subseteq V$, define
\[
 \operatorname{Ext}_F(S,N)=
 \{X\subseteq V:X\text{ is an }F\text{-model},\ S\subseteq X,\ X\cap N=\varnothing\}.
\]
Positive and negative observations are both part of the preservation contract. Dropping the negative observations can turn an impossible request into a different, feasible problem.

\begin{PPTtheorem}[Least protected extension and a complete obstruction]
\label{PPT:res:new:extension}
The extension fiber is nonempty exactly when
\[
 \operatorname{cl}_F(S)\cap N=\varnothing.
\]
When this condition holds, $\operatorname{cl}_F(S)$ is its least member. More generally,
\[
 X\in\operatorname{Ext}_F(S,N)
 \quad\Longleftrightarrow\quad
 X\text{ is an }F\text{-model and }
 \operatorname{cl}_F(S)\subseteq X\subseteq V\setminus N.
\]
If no extension exists, some $n\in N$ has a derivation from $L\cup S$ using only the declared certified rules.
\end{PPTtheorem}
\begin{proof}
Induction on a derivation shows that every model containing $S$ contains $\operatorname{cl}_F(S)$. The inductive closure is itself a model and contains $S$. If it avoids $N$, it is therefore a permitted extension, and the preceding inclusion proves leastness. If it meets $N$, every model containing $S$ contains a forbidden node, so no extension is possible. The displayed classification follows from the same inclusions. A point of the nonempty intersection is exactly the conclusion of the stated forbidden-node derivation.
\end{proof}

This describes an interval \emph{inside the family of rule models}, not an arbitrary interval of the power set. Closure under the rules remains essential. Nor is the fiber generally an affine space. For the rule $\{a,b\}\to c$, with no leaves, $S=\varnothing$ and $N=\{c\}$, the allowed models are precisely $\varnothing$, $\{a\}$, and $\{b\}$. There are three, not a power of two. The rank and entropy formulas for vector spaces cannot be imported into this model without additional structure.

\begin{PPTcorollary}[Exact extension of a partial valuation]
\label{PPT:res:new:partial}
Let $S\subseteq U\subseteq V$. A global model whose true nodes on $U$ are exactly $S$ exists if and only if
\[
 \operatorname{cl}_F(S)\cap U=S.
\]
\end{PPTcorollary}
\begin{proof}
Set $N=U\setminus S$. Closure always contains $S$, so avoiding $N$ is equivalent to the displayed equality. Apply Theorem~\ref{PPT:res:new:extension}.
\end{proof}

In particular, the test works for partial data that already violate a local rule. It need not assume local consistency and then hide that failure in a precondition. A contradictory observation can be exposed as a specific forbidden consequence.

\begin{PPTproposition}[Staged closure and persistent obstructions]
\label{PPT:res:new:stages}
For positive sets $S,T$,
\[
 \operatorname{cl}_F\bigl(\operatorname{cl}_F(S)\cup T\bigr)
   =\operatorname{cl}_F(S\cup T).
\]
If $S\subseteq S'$ and $N\subseteq N'$, then
$\operatorname{Ext}_F(S',N')\subseteq\operatorname{Ext}_F(S,N)$.
Consequently an infeasible request stays infeasible under those additions, while the field and the meaning of its rules remain fixed.
\end{PPTproposition}
\begin{proof}
Closure is extensive, monotone, and idempotent because it is the least model containing its seed. The model $\operatorname{cl}_F(S\cup T)$ contains $\operatorname{cl}_F(S)$ and $T$, giving one inclusion in the first equation. Extensivity gives $S\cup T\subseteq\operatorname{cl}_F(S)\cup T$, and monotonicity gives the other. A model satisfying the larger positive and negative prescriptions also satisfies the smaller ones. The last conclusion follows by inclusion of the extension fibers.
\end{proof}

The module \path{ProofPath.Extension} checks these statements on the existing \path{Field} type. The Python module \path{extension.py} uses the same \path{ProofFieldViews} snapshot as the local-model construction. It computes closure by forward rule scans and records each newly derived node. At most $|V|$ successful insertions occur. If a forbidden node is reached, a causal subtrace is extracted. It need not be the shortest or cheapest refutation.

The independent result checker replays that trace and checks that the recorded closure is closed under every rule. Derivability of every recorded node and closure of the recorded set together establish leastness. A negative witness only needs a valid derivation of one forbidden node. Its replay permits a previously derived node to have become an admitted observation, so the same witness can remain valid after positive data are strengthened. The checker does not call the producer, and a timeout is never returned as an obstruction.


\begin{figure}[htbp]\centering
\includegraphics[width=\linewidth]{sources/monographs/proof-path-theory/docs/figures/fig_23_protected_closure.pdf}
\caption{Two protected requests for the same rule field. With only $a$ prescribed true and $c$ prescribed false, the least model is $\{a\}$. Adding $b$ to the positive data derives the forbidden node $c$. The obstruction is the certified rule application, not the act of abandoning a search.}
\label{PPT:fig:protected-closure}
\end{figure}

\subsection{Semantic covers rather than only syntactic covers}
\label{PPT:sec:semantic-cover}

Whole-support coverage is easy to inspect and is sufficient for gluing. It can nevertheless be unnecessarily strong. Suppose $F$ contains $a\to b$, $b\to c$, and $a\to c$, with patches $\{a,b\}$ and $\{b,c\}$. The last support $\{a,c\}$ is absent from both patches. Yet the first two visible rules already imply the last rule. Requiring another patch would duplicate information rather than restore a missing constraint.

Let $\mathcal U=(U_i)_i$ cover every node. Define $F_{\mathcal U}$ to have the same nodes and leaves as $F$ but only those certified rules whose entire support lies in some $U_i$. The distinction between the two fields is explicit. We call the cover \emph{semantically sufficient for the declared rules} when
\[
 \forall j\text{ certified in }F,\qquad
 c_j\in\operatorname{cl}_{F_{\mathcal U}}(A_j).
 \tag{SC}\label{PPT:eq:semantic-cover}
\]
The word semantic here refers to the Boolean rule-model semantics of Section~\ref{PPT:sec:old24}, not to an automatic interpretation of node labels as arbitrary mathematical truths.

\begin{PPTtheorem}[Exact criterion for universal gluing]
\label{PPT:res:new:semantic-cover}
Under node coverage, the following conditions are equivalent:
\begin{enumerate}[label=(\roman*)]
\item Every matching family of local $F$-models on the patches has a unique global $F$-model extending it.
\item Every model of $F_{\mathcal U}$ is a model of $F$.
\item Condition~\eqref{PPT:eq:semantic-cover} holds.
\item For every positive seed $S$, $\operatorname{cl}_{F_{\mathcal U}}(S)=\operatorname{cl}_F(S)$.
\end{enumerate}
Whole-support coverage implies these conditions, but the converse need not hold.
\end{PPTtheorem}
\begin{proof}
A matching family assembles uniquely to a subset $X\subseteq V$ because the patches cover all nodes and agree on their common coordinates. Its local rule conditions say exactly that $X$ is a model of $F_{\mathcal U}$. Conversely, a model of $F_{\mathcal U}$ restricts to an $F$-model on every patch: every original rule supported in that patch is among the retained rules. Those restrictions match. Thus (i) and (ii) are equivalent, and uniqueness follows from node coverage.

Assume (ii) and fix a certified rule $j$. The set $\operatorname{cl}_{F_{\mathcal U}}(A_j)$ is a model of $F_{\mathcal U}$ containing all of $A_j$. By (ii) it obeys the original rule and contains $c_j$, proving (iii). Conversely, if $X$ is a model of $F_{\mathcal U}$ containing $A_j$, leastness gives $\operatorname{cl}_{F_{\mathcal U}}(A_j)\subseteq X$. Condition (iii) therefore forces $c_j\in X$. Every original rule is satisfied, which proves (ii).

Under (ii) the two fields have exactly the same models, since every $F$-model already obeys the retained rules. Their least models containing any seed are equal, giving (iv). If (iv) holds, apply it to $A_j$. The conclusion $c_j$ lies in $\operatorname{cl}_F(A_j)$ by the original rule, so it lies in the retained closure as required by (iii). Finally, whole-support coverage retains every certified rule. The three-rule example above proves that this condition is not necessary.
\end{proof}

There is a constructive negative result as well. If (SC) fails for $j$, let
\[
 X=\operatorname{cl}_{F_{\mathcal U}}(A_j).
\]
Then $A_j\subseteq X$ and $c_j\notin X$. The restrictions $X\cap U_i$ are matching local $F$-models, but their only assembled subset violates the original rule. They are a checkable obstruction to \emph{universal gluing for this cover}. They do not show that every matching family on the cover fails.

The module \path{ProofPath.SemanticCover} proves the criterion, including its connection to the actual \path{Local}, \path{Matching}, and \path{gluedValue} constructions. It is not a replacement field chosen only to have a convenient name. The Python producer returns a derivation of each original conclusion from its premises using only visible rules, or the countermodel just described. The verifier replays the positive derivations or checks the negative valuation directly. A new \path{glue_semantic} operation requires the positive certificate. The older \path{glue} method retains its simpler whole-support contract.

\PPTstudyhead{05}{A missing support can be redundant}
The patches $U_0=\{a,b\}$ and $U_1=\{b,c\}$ do not contain the support of $a\to c$. Their visible rules $a\to b$ and $b\to c$ nevertheless provide a certified derivation of that implication. Compare this with the missing conjunctive support in Figure~\ref{PPT:fig:whole-support-example}: neither visible patch there provides a substitute for $\{a,b\}\to c$.

\begin{figure}[H]\centering
\includegraphics[width=\linewidth]{sources/monographs/proof-path-theory/docs/figures/fig_22_semantic_cover.pdf}
\caption{The exact semantic-cover check for a redundant rule. The two solid rules are visible in the displayed patches. The dashed original rule has no covered whole support, but its conclusion is derived by the solid path. Dashed here means omitted from the visible subsystem, not uncertified in the original field.}
\label{PPT:fig:semantic-cover}
\end{figure}

\PPTcalloutbox{Preserve the constraint, not necessarily its original presentation}{The cover is sufficient because its visible rules impose the same models as the full field. The criterion checks every original rule, not just the target of one successful example.}

For a finite field, one closure computation per original rule is enough to produce the positive certificates or expose a failed entailment. This avoids enumerating all local valuations or all matching families. The implementation uses repeated rule scans, so the release does not claim the indexed linear-time bound for its code. Its output certificates can be verified without trusting that search procedure.

\subsection{Typed transport of derivations and obstructions}
\label{PPT:sec:rule-transports}

A change of proof-field representation needs more than a map of node names. Let $F$ and $G$ be declared fields. A rule transport consists of a node map $f:V_F\to V_G$ and a rule map $h:J_F\to J_G$ satisfying
\[
\begin{aligned}
 f(L_F)&\subseteq L_G,\\
 j\text{ certified in }F&\Longrightarrow h(j)\text{ certified in }G,\\
 A^G_{h(j)}&\subseteq f(A^F_j),\\
 c^G_{h(j)}&=f(c^F_j).
\end{aligned}
\]
The premise inclusion points in the direction needed for executing the translated step. The mapped target rule may need fewer premises, but it must not introduce a new premise unsupported by the source step. This is a sufficient, deliberately concrete transport interface. A more general translation could map one rule to an entire derivation.

\begin{PPTtheorem}[Compositional rule transport]
\label{PPT:res:new:transport}
Rule transports admit identities and associative composition. A transport $(f,h):F\to G$ sends every derivation of $v$ in $F$ to a derivation of $f(v)$ in $G$. Every $G$-model pulls back along $f$ to an $F$-model. In particular,
\[
 Y\in\operatorname{Ext}_G(f(S),f(N))
   \Longrightarrow f^{-1}(Y)\in\operatorname{Ext}_F(S,N).
\]
An obstruction in the source extension problem therefore implies an obstruction in the target extension problem with the mapped observations.
\end{PPTtheorem}
\begin{proof}
Identity maps meet each displayed condition. For two composable transports, compose the node and rule maps. Applying the second map to the first premise inclusion and then using the second inclusion proves the composite premise condition. The conclusion and leaf conditions also compose. Function composition gives the identity and associativity laws.

Proceed by induction on a source derivation. A source leaf maps to a target leaf. For a rule step, every premise of the target rule is the image of a source premise, so the induction hypotheses supply all the required target premises. Certification and the conclusion equation then give the translated conclusion.

If $Y$ is a target model, $f^{-1}(Y)$ contains all source leaves. When it contains every premise of a source rule, $Y$ contains every premise of the mapped rule and hence its conclusion. This proves the model pullback. Positive and negative observations pull back by their image definitions. If a target extension existed while the source problem were infeasible, its pullback would contradict that infeasibility.
\end{proof}

Positive seeds can be added before transport. The same rule transport sends $F[S]$ to $G[f(S)]$, and hence
\[
 f\bigl(\operatorname{cl}_F(S)\bigr)
      \subseteq\operatorname{cl}_G(f(S)).
\]
A forbidden source consequence $n\in N\cap\operatorname{cl}_F(S)$ maps to the explicit forbidden consequence $f(n)\in f(N)\cap\operatorname{cl}_G(f(S))$. The formal development checks this witness map, not only the existence statement obtained by contraposition.

\begin{PPTexample}[Why transport is not an equivalence]
\label{PPT:res:new:collapse}
Let the source have nodes $a,b$, no leaves, and no rules. The request $a$ true and $b$ false is feasible. Map both nodes to one target node $u$, again with no leaves or rules. This is a rule transport, but the image request requires $u$ both true and false and is infeasible. Thus source feasibility need not imply target feasibility, and target infeasibility need not imply source infeasibility.
\end{PPTexample}

\path{ProofPath.RuleTransport} checks these laws for actual derivations, and \path{ProofPath.ExtensionExamples} checks the collapse. The transported conclusion is true under an intended interpretation only if the target leaves and rules are sound. Renaming alone does not provide that contract.


\subsection{Executable certificates for the declared closure}
\label{PPT:sec:trace-checker}

A mathematical characterization of closure is not yet a correctness theorem for the program that checks a certificate. The next construction closes that boundary for an explicit Lean checker. It uses the existing \path{Field} and seeded-closure definitions, rather than replacing their meaning by the output of an algorithm.

\begin{PPTdefinition}[Finite indexed problem and trace]
\label{PPT:res:trace:problem}
An indexed problem $P$ specifies a finite address space $V_n=\{0,\ldots,n-1\}$, admitted leaves $L$, a list of certified rules, positive observations $S$, and forbidden nodes $N$. A rule at address $j$ has a finite premise list $A_j$ and conclusion $b_j$. Two identical rows may have different rule addresses. Every listed node must belong to $V_n$. Node lists are individually duplicate free, but $L$, $S$, and $N$ need not be disjoint.

A trace $\sigma=(j_1,\ldots,j_k)$ is a list of rule addresses. Replay begins with the list $K_0=L\mathbin{+\!+}S$. It rejects an unknown rule address or a step whose premises are not all known. Otherwise it adds the conclusion. Strict replay additionally requires the conclusion to be new. Causal replay permits a conclusion already known. Lists are compared by membership when they represent sets.
\end{PPTdefinition}

Only rows declared certified enter this indexed rule list. The finite checker does not infer certification from a claim's label, its position in a document, or an author's confidence. An adapter that selects these rows must preserve the intended field and its assumptions.

\begin{PPTtheorem}[Replay preserves derivability and composes]
\label{PPT:res:trace:replay}
Suppose every member of $K$ belongs to $\operatorname{cl}_F(S)$. If either replay mode accepts a trace starting from $K$ and returns $K'$, then
\[
 K\subseteq K'\subseteq\operatorname{cl}_F(S).
\]
For a fixed replay mode, replaying $\sigma\mathbin{+\!+}\tau$ is the same as replaying $\sigma$ and then replaying $\tau$ from the returned state. Failure at either stage propagates to the concatenated execution.
\end{PPTtheorem}
\begin{proof}
Induct on the trace. The empty trace returns the initial state. For a nonempty trace, a successful first step identifies a certified rule whose premises are already derivable. Closure under that rule proves its conclusion derivable. Adding the conclusion preserves all previously known nodes, and the induction hypothesis applies to the remaining trace. The stricter novelty test can reject a step but cannot invalidate this argument for an accepted one. The composition identity follows by the same induction, with the lookup and enabling tests taking the same branch on both sides.
\end{proof}

The functions \path{replay}, \path{enabled}, and \path{lookup} are executable definitions. The formal theorems \path{replay_sound}, \path{replay_extensive}, and \path{replay_append} concern those definitions themselves. They are not specifications assumed of an otherwise unchecked replay routine.

\begin{PPTtheorem}[Acceptance identifies the least seeded model]
\label{PPT:res:trace:least}
A complete-closure certificate consists of a strict trace $\sigma$ and a reported node list $K$. The checker accepts precisely when replay succeeds, its output has the same members as $K$, and $K$ contains the admitted leaves and satisfies every declared rule. Acceptance implies
\[
 \boxed{K=\operatorname{cl}_F(S).}
\]
Here equality means equality of the represented sets, not equality of list order.
\end{PPTtheorem}
\begin{proof}
Theorem~\ref{PPT:res:trace:replay} gives $K\subseteq\operatorname{cl}_F(S)$ because every initial node is an admitted leaf or a positive observation. Extensivity of replay ensures that $K$ contains $S$. The final closedness test makes $K$ an $F$-model, so leastness of seeded closure gives the reverse inclusion. Both tests are needed: replay provides derivability, while the model test provides saturation under rules not used in the trace.
\end{proof}

The connection to the prior theory is exact. In Lean, \path{closed_iff_model} identifies the Boolean rule check with the existing \path{Field.Model} predicate, and \path{checkClosure_exact} proves the displayed equality with \path{Field.closure}. For a valid indexed problem, \path{valid_closure_bounded} proves that every member of that closure has an address below $n$. Thus totalizing rule lookup in the abstract \path{Field} cannot introduce a conclusion at an unnamed address: only successful lookups are certified.

Logical meaning remains a separate, explicit hypothesis. If an interpretation $\mu:V_n\to\mathrm{Prop}$ satisfies the admitted leaves and rules, and every observation in $S$ is true under $\mu$, then every member of an accepted closure is true under $\mu$. This follows by applying the prior model-soundness theorem to the derived nodes. Acceptance does not establish the truth of arbitrarily attached sentence labels.

\subsection{Finite witnesses for forbidden consequences}
\label{PPT:sec:trace-obstruction}

\begin{PPTdefinition}[Checked forbidden-node witness]
\label{PPT:res:trace:witness}
A witness is a pair $(v,\tau)$, consisting of a target address and a causal trace. Its checker accepts when replay from $L\cup S$ succeeds, the returned state contains $v$, and $v\in N$.
\end{PPTdefinition}

\begin{PPTtheorem}[Sound and complete obstruction witnesses]
\label{PPT:res:trace:complete}
For the finite-premise indexed rule field,
\[
 \boxed{
 \text{an accepted forbidden-node witness exists}
 \quad\Longleftrightarrow\quad
 \operatorname{Ext}_F(S,N)=\varnothing.
 }
\]
The statement does not require the dependency graph to be acyclic.
\end{PPTtheorem}
\begin{proof}
An accepted witness derives a node in $N$. Every model containing $L\cup S$ contains that node, so no protected extension exists.

Conversely, Theorem~\ref{PPT:res:new:extension} gives a forbidden node with an inductive derivation. Every such derivation has a finite causal trace. For an admitted leaf or observation, the empty trace suffices. For a rule application, obtain a trace for each of its finitely many premises by induction, concatenate those traces, and append the rule address. Previously derived nodes are never removed. A trace valid from one known set remains valid from a larger known set in causal mode, so the premise traces remain valid when concatenated. Repeated conclusions are allowed, including when an earlier premise trace already established one. The resulting trace derives the forbidden node and is accepted.
\end{proof}

This proof explains the two replay modes. A causal witness can be reused after stronger positive observations make one of its intermediate conclusions already known. Requiring novelty at that point would reject a sound explanation. A complete-closure trace instead records fresh insertions and uses a separate final saturation test. The two contracts should not be conflated.

The Lean development proves trace weakening, strict-to-causal replay, finite premise-trace assembly, and \path{derivation_has_trace} before proving \path{witness_exists_iff_infeasible}. Witness existence is not a claim that a particular search procedure finds the shortest trace, nor is it a bound on the time taken by a Python producer.


\PPTstudyhead{06}{Why replay and closedness are both necessary}
\label{PPT:study:trace}

Take $V_4=\{0,1,2,3\}$, $L=\{0\}$, $S=\varnothing$, and the two rules $r_0:0\to1$ and $r_1:1\to2$. The least closure is $\{0,1,2\}$. Node $3$ has no derivation from these inputs.

\begin{figure}[H]\centering
\includegraphics[width=\linewidth]{sources/monographs/proof-path-theory/docs/figures/fig_24_trace_contract.pdf}
\caption{Replay and the final model test establish opposite inclusions. Together they identify the least closure. The two rejected records fail different obligations. Changing only the forbidden set changes the correct feasibility verdict without changing the closure.}
\label{PPT:fig:trace-contract}
\end{figure}

The trace $[0]$ derives exactly $\{0,1\}$ but leaves $r_1$ enabled with its conclusion absent. It is a valid partial derivation, not a complete closure. Conversely, $\{0,1,2,3\}$ is closed under both rules, but replay cannot establish its unsupported node $3$. A model need not be the least model.

With $N=\{3\}$, the exact closure gives a positive result. With $N=\{2\}$, the same closure supports a negative result with witness $(2,[0,1])$. The second problem is not a different interpretation of the same request: its protected observations have changed. A checker must receive the fixed request separately from the proposed answer.


\begin{PPTcorollary}[Meaning of a complete result certificate]
\label{PPT:res:trace:verdict}
The complete result format contains a Boolean feasibility flag, a reported closure, a strict closure trace, and an optional forbidden-node witness. Both branches require the complete-closure check. A positive result additionally requires no obstruction witness and no forbidden node in the closure. A negative result requires an accepted obstruction witness. For every accepted result $C$,
\[
 C.K=\operatorname{cl}_F(S),\qquad
 \bigl(C.\mathrm{feasible}=\mathrm{true}\bigr)
 \Longleftrightarrow
 \operatorname{Ext}_F(S,N)\ne\varnothing.
\]
\end{PPTcorollary}
\begin{proof}
The complete-closure check proves the first equality. In the positive branch, the closure is a model containing $S$ and avoiding $N$, hence is an extension. In the negative branch, the witness excludes every extension by Theorem~\ref{PPT:res:trace:complete}. Case analysis on the Boolean flag gives the equivalence.
\end{proof}

A bare negative witness has its own checker and does not need a saturated closure. The complete result format asks for both because it promises both an exact closure and a feasibility verdict. Failure of either check means that this certificate was not accepted. It is not itself a refutation of the protected problem.

\subsection{Typed exchange and the boundary of serialization}
\label{PPT:sec:trace-wire}

A certificate is useful outside its producer only when the receiving system knows exactly which data it checked. The exchange format therefore uses versioned positional arrays, exact nonnegative integer tokens, literal Booleans, and explicit rule addresses. Problem and certificate are separate records. The node address book and the selected certified-rule order remain associated with the originating proof-field snapshot.

At the JSON syntax-tree level, the problem has the shape
\begin{lstlisting}[language={}]
["proof-path/problem/1", node_count, leaves, rules, positive, negative]
\end{lstlisting}
where each rule is \texttt{[premises, conclusion]}. The certificate has the shape
\begin{lstlisting}[language={}]
["proof-path/certificate/1", feasible, closure, trace, obstruction]
\end{lstlisting}
where an obstruction is either \texttt{null} or \texttt{[target, trace]}. Array lengths and tags are exact. A Boolean is not an integer address, and the string \texttt{"false"} is not a Boolean verdict. Duplicate node entries are rejected by validity checks, while identical rule rows at different addresses remain distinguishable.

\begin{PPTtheorem}[Lossless encoding at the typed syntax boundary]
\label{PPT:res:trace:codec}
Let $e_P,d_P$ and $e_C,d_C$ be the stated Lean encoders and partial decoders for problem and certificate records. For every typed problem $P$ and certificate $C$,
\[
 d_P(e_P(P))=\operatorname{some}(P),\qquad
 d_C(e_C(C))=\operatorname{some}(C).
\]
Consequently both encoders are injective on their typed domains. These identities do not require the encoded record to pass the separate validity checks.
\end{PPTtheorem}
\begin{proof}
The natural-number and Boolean constructors decode to their original values. Induction on a list proves the corresponding array round trip. The rule, witness, and optional-witness decoders then recover their components. Unfolding the exact-length record encoders and decoders gives the two identities. If two encoded records are equal, applying the decoder gives equal \texttt{some} values, hence equal original records.
\end{proof}

\begin{PPTtheorem}[Accepted JSON records certify their decoded answer]
\label{PPT:res:trace:json}
Suppose the JSON checker accepts the separately supplied problem and certificate syntax trees. Then both decode to typed records $P,C$, both pass their index and shape checks, and the decoded certificate satisfies the closure equality and feasibility equivalence of Corollary~\ref{PPT:res:trace:verdict} for that decoded problem.
\end{PPTtheorem}
\begin{proof}
The executable checker rejects any failed decoding or validity check. In the remaining branch it calls the same complete result checker used in Corollary~\ref{PPT:res:trace:verdict}. Its acceptance therefore supplies that corollary's hypothesis. The round-trip identities additionally show that checking encoded typed records agrees with checking those records directly.
\end{proof}

The general proof ends at Lean's \path{Json} syntax-tree interface. It does not prove the correctness of the standard text parser, printer, file system, Python interpreter, or name-to-index adapter. The command-line reader is executed on actual canonical bytes and compared with an independent finite model oracle. That conformance test complements the formal theorem without extending its quantified scope.

Canonical standalone files use compact ASCII JSON with an optional final newline. Decimal or exponent spellings, alternative internal whitespace, extra fields, and duplicate object-key constructions are outside this exchange contract. The format uses no objects. Resource limits are explicit operational rejections, not evidence that a protected realization is impossible. The encoder's injectivity is a statement about retained data, not a collision-freedom claim for a later cryptographic hash.

The package supplies actual named-field exports, a separate indexed Python checker, the Lean text adapter, and generated Lean fixtures for six concrete exchanges. Their role is to test the entire executed route as well as its proved core. The original named proof fields remain finite and acyclic. Cyclic examples belong to the lower-level indexed checker and show that its soundness theorem does not depend on topological sorting. An unseeded cycle supplies no finite derivation merely by referring to itself.



\chapter{Compression, representation, and memory}\label{PPT:book:chapter-52}
Compression changes which questions can still be answered. A certificate may preserve every required deduction while erasing distinctions that later inquiry needs.

  \PPTopeningfigure{fig_20_compression_collision.pdf}{Two archives have the same certificate $C_0$ but different side records. Compression forgets that record. A single decoder receiving $C_0$ cannot return both answers to the side-record query. This is the collision criterion in its smallest finite form.}{fig:compression-collision}
  \vspace{3pt}{\sffamily\bfseries\color{KomorebiTeal}The governing distinction}\par
  A valid certificate can be sufficient for checking the present theorem and insufficient for a later question. The lost distinction, not the length of the file, determines the failure.
 


\section{Compression relative to a task family}
\label{PPT:sec:old7}

A representation is sufficient only relative to specified questions and a specified class of possible archives. Fix a set $\mathscr A$ of admissible proof archives. A representation scheme is a map $C:\mathscr A\to\mathscr R$. Its output for an archive $\PPTPcal$ is $C(\PPTPcal)$. Distinguishing the scheme from one output is essential: a decoder chosen only after seeing one fixed archive could simply store the answer.

\begin{PPTdefinition}[Task]
\label{PPT:res:7:1}
A task is a map $q:\mathscr A\to Y_q$ to a declared answer set. A task family $\PPTQcal$ is a set of such maps. Examples include asking for the target certificate, a dependency list, a specified boundary example, or the claims affected by a recorded revision. A teaching or discovery objective becomes an exact task only after its required output or acceptance criterion has been specified.
\end{PPTdefinition}

\begin{PPTdefinition}[Task-sufficient representation]
\label{PPT:res:7:2}
The scheme $C:\mathscr A\to\mathscr R$ is sufficient for $q$ when there is a single decoder $d_q:C(\mathscr A)\to Y_q$ such that
\[
q(\PPTPcal)=d_q(C(\PPTPcal))\qquad\text{for every }\PPTPcal\in\mathscr A.
\]
It is $\PPTQcal$-sufficient when this holds for every $q\in\PPTQcal$. The decoder may depend on the task and the scheme, but not on hidden knowledge of the particular archive being decoded.
\end{PPTdefinition}

The decoder is defined on attained representations, so no arbitrary answer is required for an encoding that never occurs. Exact reconstruction is one possible task, not a requirement for every purpose. The basic definition is set theoretic and does not assert that a decoder is computable, efficient, or available to a particular reader. Those are additional conditions on a practical compiler.

\begin{PPTtheorem}[Task information loss criterion]
\label{PPT:res:7:3}
Let $C$ be a compression map and let $q$ be a task. If there exist proof fields $\PPTPcal_1$ and $\PPTPcal_2$ such that
\[
C(\PPTPcal_1)=C(\PPTPcal_2)
\qquad\text{but}\qquad
q(\PPTPcal_1)\ne q(\PPTPcal_2),
\]
then the compressed representation is not sufficient for $q$.
\end{PPTtheorem}

\begin{proof}
Suppose a decoder $d_q$ existed. Applying it to the common compressed representation would have to produce both $q(\PPTPcal_1)$ and $q(\PPTPcal_2)$. Since these answers differ, no such decoder exists.
\end{proof}

\begin{PPTdefinition}[Recoverability preorder]
\label{PPT:res:7:4}
For a fixed task family $\PPTQcal$, write
\[
R_1\preceq_{\PPTQcal} R_2
\]
when every task in $\PPTQcal$ that factors through the representation scheme $R_1$ also factors through $R_2$. Both schemes have domain $\mathscr A$, and each factorization uses one decoder valid on that whole domain.
\end{PPTdefinition}

\begin{PPTproposition}[Recoverability is a preorder]
\label{PPT:res:7:5}
The relation $\preceq_{\PPTQcal}$ is reflexive and transitive.
\end{PPTproposition}

\begin{proof}
Every representation supports every task it supports, so the relation is reflexive. If every task supported by $R_1$ is supported by $R_2$, and every task supported by $R_2$ is supported by $R_3$, then every task supported by $R_1$ is supported by $R_3$. Thus the relation is transitive.
\end{proof}

Define $R_1\sim_{\PPTQcal}R_2$ when both $R_1\preceq_{\PPTQcal}R_2$ and $R_2\preceq_{\PPTQcal}R_1$.

\begin{PPTcorollary}[Quotient poset]
\label{PPT:res:7:6}
The recoverability preorder descends to a partial order on the equivalence classes under $\sim_{\PPTQcal}$.
\end{PPTcorollary}

\begin{proof}
Reflexivity and transitivity pass to equivalence classes. Mutual recoverability has been quotiented out, so antisymmetry holds by construction.
\end{proof}


\subsection{Pure retention operators}

The simplest compression removes nodes without rewriting the retained material. This baseline is useful because it separates information deletion from more complicated summarization or reinterpretation.

\begin{PPTdefinition}[Retention projection]
\label{PPT:res:7:7}
Fix a finite node universe $V$. For a mask $K\subseteq V$, define
\[
P_K(S)=S\cap K,
\qquad S\subseteq V.
\]
\end{PPTdefinition}

\begin{PPTtheorem}[Retention projection algebra]
\label{PPT:res:7:8}
For masks $K,L\subseteq V$,
\[
P_K\circ P_L=P_{K\cap L}=P_L\circ P_K.
\]
Consequently every retention projection is idempotent, and the masks form a Boolean lattice with meet $K\cap L$ and join $K\cup L$.
\end{PPTtheorem}

\begin{proof}
For every $S\subseteq V$,
\[
P_K(P_L(S))=(S\cap L)\cap K=S\cap(K\cap L).
\]
Intersection is associative and commutative, which gives the composition identity and commutativity. Taking $L=K$ gives idempotence. The power set of $V$ with intersection, union, complement, bottom $\varnothing$, and top $V$ is a Boolean lattice.
\end{proof}

\begin{PPTcorollary}[Task monotonicity under retention]
\label{PPT:res:7:9}
If $K\subseteq L$, every task whose required nodes and metadata are retained by $P_K$ is also supported by $P_L$.
\end{PPTcorollary}

\begin{proof}
Every item retained by $K$ is retained by $L$. Hence every task requirement contained in $K$ is contained in $L$.
\end{proof}

Pure retention operators commute. Noncommutativity in real proof workflows therefore requires an additional operation such as state-dependent summarization, renaming, quotienting, or reinterpretation. This distinction prevents simple deletion from being credited with structure it does not possess.

\begin{figure}[htbp]
\centering
\includegraphics[width=\linewidth]{sources/monographs/proof-path-theory/docs/figures/fig_05_recoverability_preorder.pdf}
\caption{A concrete retention diamond for the four archives $(C_0,x,y)$ with $x,y\in\{0,1\}$. Arrows discard coordinates. The two middle schemes are incomparable for the task family $\{q_x,q_y\}$. This is not a universal ranking of teaching, discovery, and formal presentations.}
\label{PPT:fig:preorder}
\end{figure}

\begin{PPTtheorem}[Task-family dependence of minimal compression]
\label{PPT:res:7:10}
There is a family of archives with a fixed mandatory certificate and two independently variable side records for which different task families have different least sufficient retention masks.
\end{PPTtheorem}

\begin{proof}
Let the archives be $\PPTPcal_{xy}=(C_0,x,y)$ for $(x,y)\in\{0,1\}^2$, where $C_0$ is a fixed certificate retained by every allowed mask. Let $q_x(\PPTPcal_{xy})=x$ and $q_y(\PPTPcal_{xy})=y$. Retaining $x$ is sufficient for $q_x$. Any mask omitting $x$ identifies $\PPTPcal_{0y}$ with $\PPTPcal_{1y}$, so it is not sufficient for that task. Thus the least mask for $\{q_x\}$ retains exactly $C_0$ and $x$. The same argument gives the least mask $C_0$ and $y$ for $\{q_y\}$. Their joint task family requires both records. These masks are also the unique minima for unit cost per retained side record.
\end{proof}

The example concerns pure retention. A general encoding may recode the records rather than delete them. Sufficiency for all archive queries requires an injective encoding, not literal preservation of the original file. The following theorem makes that distinction precise.

\begin{PPTtheorem}[All predicates distinguish exactly injective encodings]
\label{PPT:res:7:11}
\label{PPT:thm:injective-compression}
Let $C:A\to B$. Every predicate $q:A\to\mathrm{Prop}$ factors through $C$ if and only if $C$ is injective.
\end{PPTtheorem}
\begin{proof}
If $C$ is injective, define the decoder of $q$ at $b$ by the predicate $\exists a\ (C(a)=b\ \text{and}\ q(a))$. At $b=C(a_0)$, injectivity makes this equivalent to $q(a_0)$. Conversely, if $C(a)=C(a')$, use the predicate $q(z)\equiv(z=a)$. Its decoded values at the equal encodings must agree, giving $a'=a$. This is the general collision criterion specialized to all predicates.
\end{proof}

The formal result \path{ProofPath.all_predicates_sufficient_iff_injective} records the distinction between information loss and lossless recoding. It makes no assertion about how cheaply a decoder can be computed.


\subsection{A complete factorization test on attained representations}
\label{PPT:sec:attained-recovery}

The collision obstruction of Theorem~\ref{PPT:res:7:3} has a converse. It gives a complete test for exact task recovery at the level of sets. Write $\operatorname{im}C=C(\mathscr A)$ for the attained output set. A decoder is required on this set and nowhere else.

\begin{PPTtheorem}[Attained factorization and unique decoding]
\label{PPT:res:new:factorization}
For $C:\mathscr A\to\mathcal R$ and $q:\mathscr A\to Y$, the following are equivalent:
\[
\begin{aligned}
 &\exists d:\operatorname{im}C\to Y\quad
       \forall a\in\mathscr A,\quad q(a)=d(C(a)),\\
 &\forall a,b\in\mathscr A,\quad C(a)=C(b)\Longrightarrow q(a)=q(b).
\end{aligned}
\]
When a decoder exists it is unique on $\operatorname{im}C$. Otherwise there are two archives with the same encoding and different task answers.
\end{PPTtheorem}
\begin{proof}
A decoder applied to equal encodings gives equal answers. Conversely, for an attained value $r$, choose $a$ with $C(a)=r$ and define $d(r)=q(a)$. Constancy on each fiber makes the value independent of the choice. At $r=C(a)$ it gives the required answer. Any other decoder must give that same answer at every attained value, proving uniqueness. Negating fiber constancy supplies the pair in the last statement. The construction is set theoretic and uses classical choice in the Lean implementation. It does not provide an efficient decoder for arbitrary infinite archive classes.
\end{proof}

\begin{PPTexample}[The empty-domain boundary]
\label{PPT:res:new:empty-decoder}
Take $\mathscr A=\varnothing$, $\mathcal R=\{*\}$, and $Y=\varnothing$. There is one encoding and one task on the empty archive class. The attained set is empty, so its empty decoder exists. There is no total decoder $\{*\}\to\varnothing$. Thus attained decoding and decoding on the entire output type are not equivalent without an additional hypothesis.
\end{PPTexample}

The earlier formal definition \path{ProofPath.Sufficient} asks for a total decoder. It is retained for compatibility and was not false, but it was stronger than the manuscript's stated domain convention. The new namespace \path{ProofPath.Recovery} matches the manuscript exactly. It proves that total decoding implies attained decoding, and that the two are equivalent whenever the answer type is nonempty. In the latter case choose one default answer for unattained outputs. The empty-domain example is also kernel checked.

The finite Python interface \path{factor_task} takes an explicitly declared table of archive encodings and a task table on the identical domain. It returns either a decoder table or two archive identifiers witnessing a collision. Its verifier checks the proposed result directly, including the empty domain, missing outputs, extra unattained entries, and contradictory records. The table strings are exact identifiers and answers, not a semantic equality procedure for arbitrary mathematical expressions. The existing retention-based \path{task_sufficient} function remains a different, narrower task model.

\subsection{The least information for a family of tasks}
\label{PPT:sec:joint-recovery}

Suppose tasks $q_i:\mathscr A\to Y_i$ are indexed by $i\in I$. Their joint representation is
\[
 Q(a)=(q_i(a))_{i\in I}.
\]
It records every distinction asked for by the task family and makes no demand to recover the entire source archive.

\begin{PPTtheorem}[Joint task representation and compositional recovery]
\label{PPT:res:new:joint}
An encoding $C$ is sufficient for every $q_i$ if and only if it is sufficient for $Q$. In particular every sufficient encoding recovers the joint representation on attained values, while $Q$ itself recovers each task by projection.

If a representation $R$ is recoverable from $C$, then every task recoverable from $R$ is recoverable from $C$. Recovery of representations is reflexive and transitive.
\end{PPTtheorem}
\begin{proof}
If $Q$ factors through $C$, each coordinate factors by projection. Conversely, equal values of $C$ give equal values of every $q_i$. Function extensionality gives equal joint vectors, and Theorem~\ref{PPT:res:new:factorization} supplies the joint decoder. This argument allows a different answer type for each task and also covers an empty task family.

For the second assertion, write $R=\pi\circ C$ and $q=d\circ R$ on their attained sets. Then $q=(d\circ\pi)\circ C$. The identity decoder proves reflexivity, and composition of two recovery maps proves transitivity. All maps are applied only to outputs in their declared domains.
\end{proof}

The word least concerns retained distinctions, not bytes, running time, or the cognitive cost of reading an archive. Two archives are identified by $Q$ precisely when every declared task gives the same answer on them. Any sufficient representation may retain more, but it cannot identify archives that $Q$ distinguishes. For a finite archive family, the reference implementation labels these joint-answer classes and checks that every task is constant within a class. Costs remain a separate optimization problem.

This result also completes the formal justification of the recovery condition in Section~\ref{PPT:sec:old19}. A refinement that reconstructs only the final theorem need not preserve the joint tasks of audit and discovery. A refinement with a checked recovery map for the preceding representation does preserve them. The distinction is between a longer proof and a genuinely more informative archive, not between short and long documents in general.


\section{Endpoint, audit, and discovery equivalence}
\label{PPT:sec:old8}

The phrase ``the same proof'' can refer to several distinct equivalence relations.

\begin{PPTdefinition}[Endpoint equivalence]
\label{PPT:res:8:1}
Fix a background theory, an interpretation of claim labels, and an admissible inference policy. Within that context, two representations are endpoint equivalent when they certify the same judgment from the same complete assumption context and requested targets. Required premises and any additional admitted assumptions are accounted for in that context. A comparison across different contexts requires an explicitly specified transport of assumptions and judgments, not merely identical displayed labels.
\end{PPTdefinition}

\begin{PPTdefinition}[Audit equivalence]
\label{PPT:res:8:2}
Two representations are audit equivalent when they preserve equivalent dependency, evidence, provenance, and scope information for checking the declared targets.
\end{PPTdefinition}

\begin{PPTdefinition}[Discovery equivalence]
\label{PPT:res:8:3}
Two representations are discovery equivalent relative to a declared bridge model when they preserve the same examples, alternatives, boundary cases, candidate routes, and future reach, up to the chosen graph equivalence.
\end{PPTdefinition}

When audit and discovery task families are required to include the endpoint tasks, either equivalence implies endpoint equivalence. Audit and discovery equivalence need not imply one another: a representation may preserve provenance without alternatives, or alternatives without provenance. Their joint task family is stronger than either. Any further hierarchy requires an explicit inclusion of task families. Two perfectly valid certificates can therefore expose different mathematical neighborhoods without differing in validity.

\begin{PPTremark}
\label{PPT:res:8:4}
These equivalences are task-relative. Discovery equivalence requires a declared bridge model. It is not an assertion that every future use of a proof can be predicted.
\end{PPTremark}


\section{Vector-valued loss}
\label{PPT:sec:old9}

Let a node $v$ carry finite nonnegative numerical coordinates
\[
\mathbf v(v)=
\bigl(V(v),D(v),M(v),F(v),A(v)\bigr),
\]
representing declared verification, discovery, meaning, future bridge, and auditability weights. Summation requires an additive numerical model. The weights may be supplied or estimated, but an ordinal ranking alone does not justify adding them. Changing a numerical scale can change the resulting loss and retention problem.

\begin{PPTdefinition}[Additive loss vector]
\label{PPT:res:9:1}
For a retained node set $K\subseteq V$, define
\[
\PPTLoss(\PPTPcal,K)
=
\sum_{v\notin K}\mathbf v(v).
\]
Normalize a coordinate by its total only when that total is positive. If the total is zero, nonnegativity makes every contribution in that coordinate zero. The implementation records normalized loss zero there as a convention, rather than dividing by zero.
\end{PPTdefinition}

\begin{PPTproposition}[Zero verification loss can coexist with positive discovery loss]
\label{PPT:res:9:2}
There exist proof fields and certificates $K$ for which the verification coordinate of $\PPTLoss(\PPTPcal,K)$ is zero while the discovery and bridge coordinates are positive.
\end{PPTproposition}

\begin{proof}
Take an AND-DAG certificate from $A$ to $B$ and add an optional example node $E$ that is not an ancestor of $B$. Assign positive discovery and bridge values to $E$ and zero verification value. The unique core omits $E$ while preserving the entire verification ancestry. The verification loss is zero and the other declared losses are positive.
\end{proof}

\begin{figure}[htbp]
\centering
\includegraphics[width=\linewidth]{sources/monographs/proof-path-theory/docs/figures/fig_06_loss_vectors.pdf}
\caption{Exact retained and lost values for the canonical AND-DAG certificate. The numbers are declared values, not measured universal units of meaning. Zero verification loss coexists with loss of seven discovery units and eight bridge units. All auditability values in this example were declared zero.}
\label{PPT:fig:loss}
\end{figure}

No claim is made that all value coordinates are additive. Additivity is a transparent baseline. Interactions can be represented by hyperedge values, submodular utilities, or task-specific decoders when the application requires them.


\section{Compression debt}
\label{PPT:sec:old10}

Compression may postpone work rather than eliminate it. A later task may require reconstruction of material that was omitted from the current representation.

\begin{PPTdefinition}[Future task model]
\label{PPT:res:10:1}
Let $\PPTTcal$ be a finite task set with probabilities $p_t\ge0$ satisfying $\sum_t p_t=1$. Let $R_t\subseteq V$ be the node set required by task $t$, and let $c(v)\ge0$ be the recovery cost of node $v$ when it is absent.
\end{PPTdefinition}

\begin{PPTdefinition}[Expected compression debt]
\label{PPT:res:10:2}
The expected compression debt of retained set $K$ is
\[
\PPTDebt(K)
=
\sum_{t\in\PPTTcal}p_t
\sum_{v\in R_t\setminus K}c(v).
\]
\end{PPTdefinition}

\begin{PPTproposition}[Debt is monotone under retention]
\label{PPT:res:10:3}
If $K_1\subseteq K_2$, then
\[
\PPTDebt(K_2)\le\PPTDebt(K_1).
\]
\end{PPTproposition}

\begin{proof}
For every task $t$, the missing set $R_t\setminus K_2$ is contained in $R_t\setminus K_1$. Recovery costs are nonnegative. The taskwise cost cannot increase, and neither can its probability-weighted sum.
\end{proof}

\begin{PPTcorollary}[Irreversible compression accumulates debt]
\label{PPT:res:10:4}
For a descending sequence of retained sets
\[
K_0\supseteq K_1\supseteq\cdots\supseteq K_r,
\]
the expected debt is nondecreasing.
\end{PPTcorollary}

A recovery archive can change the cost function. If an omitted node has a stable provenance pointer and a reproducible generator, its recovery cost may be small. If it survives only in memory, its recovery cost may be effectively infinite.


\section{A proof ecology rather than one ideal form}
\label{PPT:sec:old11}

The same theorem can support several legitimate representations:
\[
\PPTRcal(B)=
\{R_{\mathrm{cert}},R_{\mathrm{journal}},R_{\mathrm{teach}},
R_{\mathrm{formal}},R_{\mathrm{discovery}},R_{\mathrm{archive}}\}.
\]
The certificate is optimized for checking. The journal form balances rigor and readable length. The teaching form expands conceptual bottlenecks. The formal form exposes primitive dependencies. The discovery form retains alternatives and failed routes. The archive preserves reconstruction information.

\begin{PPTdefinition}[Representation portfolio]
\label{PPT:res:11:1}
A representation portfolio is a family of proof representations together with stable claim identifiers, translation maps, and a shared status ledger.
\end{PPTdefinition}

\begin{PPTproposition}[Portfolio robustness]
\label{PPT:res:11:2}
For representation schemes $C_i:\mathscr A\to R_i$, retain the tagged portfolio $C(\PPTPcal)=(C_i(\PPTPcal))_i$. Every task recoverable from any one $C_i$ is recoverable from $C$.
\end{PPTproposition}

\begin{proof}
If $q=d_i\circ C_i$, compose $d_i$ with the $i$th coordinate projection from the portfolio. This gives a decoder of $q$ from $C$. Retaining the coordinate labels is part of the construction. An untagged union that loses source identities need not permit the same recovery.
\end{proof}

Redundancy is useful when it is heterogeneous. Storing six copies of the same compressed certificate adds little. Storing a certificate, an expanded explanation, a formal dependency graph, and a reproducible notebook may preserve genuinely different values.


\section{Proof compilers and shared ledgers}
\label{PPT:sec:old12}

A proof compiler generates a representation for a specified purpose from one canonical proof field. The connection between source and readable exposition also motivates literate programming \cite{PPT:knuth}.

\begin{PPTdefinition}[Proof compiler]
\label{PPT:res:12:1}
Let $g$ be a presentation goal and let $\PPTPcal$ be a versioned proof field. A proof compiler is a map on the declared archive class
\[
\Gamma_g:\mathscr A\longrightarrow R_g
\]
that selects, orders, and renders the material required by $g$ while preserving links to canonical claim and justification identifiers.
\end{PPTdefinition}

Typical goals include certificate, journal, teaching, discovery, formal verification, and archive. The compiler does not invent logical status. It reads status from the canonical ledger.

\begin{PPTproposition}[Shared-ledger consistency]
\label{PPT:res:12:2}
Suppose every compiled representation is generated from one canonical ledger and every rendered claim retains its ledger identifier. Then a status correction can be propagated consistently to every compiled representation containing that claim by updating the ledger and recompiling.
\end{PPTproposition}

\begin{proof}
The stable identifier determines every occurrence of the claim across the compiled outputs. Updating the canonical record changes the source value associated with that identifier. Recompilation reads the updated record in every view. A stale status can persist only in an output detached from the canonical ledger or not recompiled.
\end{proof}

This proposition is operationally important. A monograph, a formal library, and a teaching site should not maintain independent truth labels for the same claim.


\section{Latent bridge graphs}
\label{PPT:sec:old13}

A proof field records certified dependencies. A latent bridge graph records possible future transfers. The two graphs have different logical roles.

\begin{PPTdefinition}[Latent bridge graph]
\label{PPT:res:13:1}
A latent bridge graph is a finite directed graph
\[
\PPTBcal=(W,E_B)
\]
whose nodes include retained proof objects, omitted proof objects, and anticipated future targets. An edge $x\to y$ means that $x$ has been declared a possible transfer route toward $y$. The edge may be certified, conjectural, or exploratory, but its status must be recorded.
\end{PPTdefinition}

Let $S\subseteq W$ be the retained starting set. Let $F\subseteq W$ be a future target set with weights $w_f\ge0$.

\begin{PPTdefinition}[Future reach utility]
\label{PPT:res:13:2}
The weighted future reach of $S$ is
\[
U(S)=\sum_{f\in F}w_f\,
\mathbf 1\{f\in\PPTReach_{\PPTBcal}(S)\}.
\]
\end{PPTdefinition}

This quantity measures declared reach, not truth. A conjectural edge can increase modeled reach without certifying the target. The ledger must keep those statuses separate.

\subsection{Bridge criticality and dominators}

\begin{PPTdefinition}[Bridge-critical node]
\label{PPT:res:13:3}
A node $v$ is bridge critical for target $f$ relative to starts $S$ when $f$ is reachable from $S$ in $\PPTBcal$ and becomes unreachable after $v$ is removed.
\end{PPTdefinition}

In graph language, $v$ dominates an initially reachable $f$ relative to $S$ when every directed path from $S$ to $f$ passes through $v$. Paths of length zero are allowed in the reachability convention. Standard graph terminology is used throughout \cite{PPT:diestel}.

\begin{PPTtheorem}[Criticality equals domination]
\label{PPT:res:13:4}
Assume $f$ is reachable from $S$. A node $v\notin S$ is bridge critical for $f$ if and only if $v$ dominates $f$ relative to $S$.
\end{PPTtheorem}

\begin{proof}
If $v$ dominates $f$, every path from $S$ to $f$ uses $v$. Removing $v$ removes every such path, so $f$ becomes unreachable. Conversely, if removing $v$ makes $f$ unreachable, then there could not have been a path from $S$ to $f$ avoiding $v$. Therefore every path used $v$, which is exactly domination.
\end{proof}

\begin{PPTcorollary}[Weighted criticality]
\label{PPT:res:13:5}
For $v\notin S$, the utility lost by deleting $v$ is the total weight of initially reachable future targets dominated by $v$:
\[
U(S)-U_{\PPTBcal\setminus\{v\}}(S)
=\sum_{\substack{f\in F\cap\PPTReach_{\PPTBcal}(S)\\v\text{ dominates }f}}w_f.
\]
Deleted target vertices are counted as unreachable.
\end{PPTcorollary}

\begin{figure}[htbp]
\centering
\includegraphics[width=\linewidth]{sources/monographs/proof-path-theory/docs/figures/fig_07_bridge_dominators.pdf}
\caption{A declared bridge graph. The boundary theorem $C$ dominates future targets $A^\prime$ and $B^\prime$ but not $C^\prime$, which has a bypass from $E$. These arrows describe future reach in the supplied model, not proof certification. Deleting a bridge vertex is different from merely retaining a different source set.}
\label{PPT:fig:dominator}
\end{figure}

A node may have low verification value for the current theorem and high domination value for future targets. This is the precise form of the claim that a forgotten step can remain mathematically important after it has lost its role in the original certificate.


\section{Retention under a budget}
\label{PPT:sec:old14}

An archive cannot retain every object at full resolution. The first mathematically controlled retention problem is to choose a finite supplement that preserves as much future reach as possible.

Let $C\subseteq W$ be the candidate nodes that may be retained and let $M$ be a mandatory base set. For $S\subseteq C$, define
\[
U_M(S)=U(M\cup S).
\]
The bridge graph remains fixed: retention changes the starting set, not the available intermediate vertices or edges. Deleting unretained intermediate vertices would define a different objective, for which the coverage proof below need not apply.

\begin{PPTtheorem}[Bridge coverage is monotone and submodular]
\label{PPT:res:14:1}
\label{PPT:thm:bridge-submodular}
The future reach utility $U_M$ is a nonnegative monotone submodular set function on $C$.
\end{PPTtheorem}

\begin{proof}
For each candidate $x\in C$, let $R_x\subseteq F$ be the set of future targets reachable from $M\cup\{x\}$. Let $R_M$ be the targets already reachable from $M$. Reachability from a union of starting nodes is the union of their individual reach sets, so
\[
U_M(S)=
\sum_{f\in R_M\cup\bigcup_{x\in S}R_x}w_f.
\]
Adding a candidate cannot remove a reachable target, so $U_M$ is monotone.

Now let $S\subseteq T\subseteq C$ and let $x\in C\setminus T$. The marginal gain of adding $x$ is the total weight of targets in $R_x$ that are not already covered. Since $T$ covers at least every target covered by $S$, the set newly covered from $T$ is contained in the set newly covered from $S$. Therefore
\[
U_M(S\cup\{x\})-U_M(S)
\ge
U_M(T\cup\{x\})-U_M(T),
\]
which is submodularity.
\end{proof}

\begin{PPTtheorem}[Greedy cardinality guarantee]
\label{PPT:res:14:2}
\label{PPT:thm:bridge-greedy}
Let $k$ be an integer with $1\leq k\leq|C|$. Starting from the empty supplement, let $S_k$ result from repeatedly adding an unselected candidate with largest marginal future reach. Let $O_k$ be an optimal supplement of size at most $k$. Then
\[
U_M(S_k)-U_M(\varnothing)
\ge
\left(1-\left(1-\frac1k\right)^k\right)
\bigl(U_M(O_k)-U_M(\varnothing)\bigr),
\]
and therefore the greedy gain is at least $(1-1/e)$ times the optimal gain.
\end{PPTtheorem}

\begin{proof}
Write $G_i=U_M(S_i)-U_M(\varnothing)$ and $G^*=U_M(O_k)-U_M(\varnothing)$. By monotonicity and submodularity, the total marginal contribution of the at most $k$ elements of $O_k$ beyond $S_i$ is at least $G^*-G_i$. Hence at least one candidate has marginal gain at least $(G^*-G_i)/k$. Greedy chooses a candidate with at least that gain, so
\[
G^*-G_{i+1}
\le
\left(1-\frac1k\right)(G^*-G_i).
\]
Iterating from $G_0=0$ gives
\[
G^*-G_k
\le
\left(1-\frac1k\right)^kG^*.
\]
Rearranging gives the first bound. The inequality $(1-1/k)^k\le e^{-1}$ gives the second.
\end{proof}

\begin{PPTwarning}[Arbitrary costs]
\label{PPT:res:14:3}
For $k=0$, both gains are zero without using the displayed division by $k$. For a budget at least $|C|$, retaining all candidates is optimal. The displayed theorem is the cardinality form of the standard submodular greedy bound \cite{PPT:nemhauser}. A gain-per-cost rule for unequal costs is a different algorithm and does not inherit this guarantee. Zero-cost candidates also need an explicit convention.
\end{PPTwarning}

\begin{figure}[H]
\centering
\includegraphics[width=\linewidth]{sources/monographs/proof-path-theory/docs/figures/fig_08_submodular_greedy.pdf}
\caption{Greedy and exact retention coincide on the committed finite instance. The square and circle marks are deliberately overlaid. The general cardinality guarantee is $1-1/e$, not equality with the optimum on every graph.}
\label{PPT:fig:greedy}
\end{figure}



\chapter{Revision, review, and inquiry}\label{PPT:book:chapter-53}
A mathematical archive has a continuing life. Revisions, shared error, and finite checking capacity change its use without changing the standard of proof.

  \PPTopeningfigure{fig_21_review_error_floor.pdf}{Three declared review models separate independent noise, a common offset, and a common offset with correlated errors. The marginal variance is one. The offset is $b=0.3$ when present and the correlation is $\rho=0.1$ when present. Both axes are logarithmic. These are model consequences, not observed reviewer data.}{fig:review-floor}
  \vspace{3pt}{\sffamily\bfseries\color{KomorebiTeal}The governing distinction}\par
  Aggregation can reduce independent error without removing a shared offset. Review evidence, provenance, and certificate validity must therefore remain distinguishable.
 


\section{Versioned proof objects}
\label{PPT:sec:old15}
\label{PPT:sec:versioned}

A proof archive changes. Claims are strengthened, corrected, refuted, or split. A reproducible theory therefore needs stable identifiers and explicit version relations.

\begin{PPTdefinition}[Versioned proof object]
\label{PPT:res:15:1}
A versioned proof object is a sequence
\[
\PPTPcal_0\longrightarrow\PPTPcal_1\longrightarrow\cdots
\]
whose nodes and justifications carry stable identifiers. Each transition records additions, removals, altered scopes, evidence changes, and a digest of its parent version.
\end{PPTdefinition}

\begin{PPTdefinition}[Append-only certified update]
\label{PPT:res:15:2}
An update $\PPTPcal_i\to\PPTPcal_{i+1}$ is append-only certified when every existing node, admissible-leaf status, evidence record, interpretation, and certified justification remains unchanged. The later field may add records and certified rules but may not withdraw an earlier admitted leaf or change the scope of an earlier inference.
\end{PPTdefinition}

\begin{PPTtheorem}[Derivability is monotone under append-only updates]
\label{PPT:res:15:3}
Every target derivable in $\PPTPcal_i$ remains derivable in $\PPTPcal_{i+1}$ under an append-only certified update.
\end{PPTtheorem}

\begin{proof}
A certificate in $\PPTPcal_i$ uses only nodes and certified justifications that remain unchanged in $\PPTPcal_{i+1}$. The same selected subfield therefore satisfies the same derivability closure in the later version.
\end{proof}

Retractions and semantic changes are not append-only. They require invalidation analysis.

\begin{PPTdefinition}[Syntactic affected set]
\label{PPT:res:15:4}
For a changed node set $X\subseteq V$, define
\[
\operatorname{Aff}(X)=X\cup\PPTDesc(X),
\]
where descendants are computed in the declared dependency graph.
\end{PPTdefinition}

\begin{PPTtheorem}[Minimal closed syntactic review set]
\label{PPT:res:15:5}
The set $\operatorname{Aff}(X)$ is the unique smallest node set $Y$ such that $X\subseteq Y$ and $Y$ is closed under outgoing dependency edges: whenever $u\in Y$ and $u\to v$ is a declared dependency edge, then $v\in Y$.
\end{PPTtheorem}

\begin{proof}
By construction, every node reachable from $X$ lies in $\operatorname{Aff}(X)$, so the set contains $X$ and is closed under outgoing edges. Let $Y$ be any other set with those two properties. Every directed path beginning in $X$ remains inside $Y$ by induction on path length. Hence every descendant of $X$ lies in $Y$, and $\operatorname{Aff}(X)\subseteq Y$.
\end{proof}

\begin{figure}[htbp]
\centering
\includegraphics[width=\linewidth]{sources/monographs/proof-path-theory/docs/figures/fig_09_version_invalidation.pdf}
\caption{Changing $L_1$ marks $L_1$, $L_2$, and $B$ for syntactic review. The lower branch has no forward dependency path from $L_1$. Membership in the review set is not a claim that a theorem has become false.}
\label{PPT:fig:invalidation}
\end{figure}

The theorem is intentionally syntactic. A descendant can survive after rechecking. The affected set says where review is required, not where falsehood is guaranteed.


\FloatBarrier
\section{Evidence ledgers and provenance}
\label{PPT:sec:old16}

\Needspace{9\baselineskip}
A claim ledger should answer five questions without reconstructing the entire research history.

\begin{enumerate}
\item What is the exact claim identifier?
\item What is its current status and scope?
\item Which justifications depend on it?
\item Which evidence records support or challenge it?
\item Which versions and compiled documents contain it?
\end{enumerate}

\begin{PPTdefinition}[Evidence ledger]
\label{PPT:res:16:1}
An evidence ledger is an appendable collection of typed, scoped records linked to claim and justification identifiers. A record may certify, support, challenge, interpret, or refute an object. Its role is explicit rather than inferred from prose placement.
\end{PPTdefinition}

\begin{PPTdefinition}[Provenance digest]
\label{PPT:res:16:2}
A provenance digest is a cryptographic hash of a canonical serialization of the claim identifier, content, scope, evidence references, parent version, and relevant build inputs.
\end{PPTdefinition}

A digest does not prove a theorem. A changed digest establishes that the hashed byte sequence changed. An unchanged digest is an integrity check only subject to the possibility of collisions. Canonical serialization gives different representations a common byte-level reference.

\begin{PPTproposition}[Digest stability]
\label{PPT:res:16:3}
Fix a deterministic serialization $s$ and a deterministic hash function $h$. Identical serialized byte sequences have identical digests under $h$. Distinct serialized byte sequences have distinct digests unless they form a collision under $h$. Detecting a record change additionally requires that the change alter its serialized bytes.
\end{PPTproposition}

\begin{proof}
If $s(r)=s(r')$, applying the function $h$ gives equal digests. If $s(r)\ne s(r')$ but $h(s(r))=h(s(r'))$, the two serialized strings form a hash collision by definition. Thus change detection at the record level requires both preservation of the audited distinction by $s$ and exclusion of a collision for the resulting bytes. No injectivity of a fixed-length hash on arbitrary byte strings is asserted.
\end{proof}


The serialization contract must state which fields, types, and distinctions it preserves and which differences it intentionally normalizes. Loss during serialization and a collision during hashing are separate mechanisms.

For example, joining string parts with a separator is ambiguous if the same separator may occur inside a part. The tuples $(\texttt{a|b},\texttt{c})$ and $(\texttt{a},\texttt{b|c})$ both become \texttt{a|b|c}. No hash function can recover the lost tuple boundary. The earlier \path{provenance_digest} helper had this defect with a control-character separator. Python release 1.4.0 replaces that encoding with a versioned, domain-separated JSON array of strings. Historical digests must retain their earlier format and version. The snapshot digest has its own unchanged serialization contract.

\begin{PPTproposition}[Hashing cannot add recoverable tasks]
\label{PPT:res:new:hash}
Let $s:\mathscr A\to B$ be serialization and $h:B\to D$ be any deterministic map. Every task recoverable from $h\circ s$ is recoverable from $s$. If $h$ is injective on the attained set $s(\mathscr A)$, the two encodings support exactly the same tasks.
\end{PPTproposition}
\begin{proof}
If $s(a)=s(b)$, then $h(s(a))=h(s(b))$. Any task constant on the digest fibers is therefore constant on the serialization fibers, and Theorem~\ref{PPT:res:new:factorization} supplies its decoder. If $h$ is injective on attained serializations, equality of digests conversely implies equality of serializations, so the two fiber partitions coincide and the same criterion applies in both directions.
\end{proof}

The module \path{ProofPath.Recovery} checks both implications. The injectivity condition is a conditional mathematical hypothesis, not a claim that a fixed-length cryptographic hash is injective on arbitrary inputs. A hash cannot repair an ambiguity introduced before hashing. It may introduce additional collisions, which must be kept separate from the serializer's preservation contract.

\section{Correlated review}
\label{PPT:sec:old17}
\label{PPT:sec:correlated-review}

Several reviews are not equivalent to several independent reviews. Reviewers may share training, notation, examples, software, or an unnoticed premise.

Let $X_i=\theta+\varepsilon_i$ be numerical estimates of a fixed quantity $\theta$. Assume $\mathbb E\varepsilon_i=0$, $\operatorname{Var}(\varepsilon_i)=\sigma^2>0$, and common pairwise correlation $\rho$. For $n>1$, assume $-1/(n-1)<\rho\leq1$, so the mean has positive variance. These variables model estimates, not proof truth values. The mean has variance
\[
\operatorname{Var}(\bar X)
=
\sigma^2\left(\frac{1-\rho}{n}+\rho\right).
\]

\begin{PPTtheorem}[Effective reviewer count]
\label{PPT:res:17:1}
\label{PPT:thm:review-count}
Under the equicorrelated model, define the variance-equivalent reviewer count by
\[
\boxed{
 n_{\mathrm{eff}}
 =
 \frac{n}{1+(n-1)\rho}
}
\]
so that $\operatorname{Var}(\bar X)=\sigma^2/n_{\mathrm{eff}}$. This effective count is a positive real number and need not be an integer.
\end{PPTtheorem}

\begin{proof}
The mean of $m$ independent reviewers of variance $\sigma^2$ has variance $\sigma^2/m$. Equating this with the correlated mean variance gives
\[
\frac{\sigma^2}{m}
=
\sigma^2\left(\frac{1-\rho}{n}+\rho\right)
=
\sigma^2\frac{1+(n-1)\rho}{n}.
\]
Solving for $m$ gives the formula.
\end{proof}

For $n=1$, the effective count is one. At the excluded boundary $\rho=-1/(n-1)$, the mean has zero variance and there is no finite variance-equivalent count. For fixed $0<\rho\leq1$, $n_{\mathrm{eff}}\to1/\rho$ as $n\to\infty$. Negative correlations can give $n_{\mathrm{eff}}>n$, but a fixed negative correlation is not admissible for arbitrarily large $n$.

\begin{figure}[htbp]
\centering
\includegraphics[width=\linewidth]{sources/monographs/proof-path-theory/docs/figures/fig_10_effective_reviewers.pdf}
\caption{Effective reviewer count under the specified equicorrelation model. The logarithmic vertical scale keeps all four curves visible. At $n=100$ and $\rho=0.1$, the effective count is about $9.17$. This formula models variance reduction, not the probability that a mathematical claim is true.}
\label{PPT:fig:reviewers}
\end{figure}

For a general positive definite covariance matrix $\Sigma$, consider unbiased linear aggregates $w^\top X$ with $\mathbf 1^\top w=1$.

\begin{PPTtheorem}[Minimum-variance covariance-aware weights]
\label{PPT:res:17:2}
The unique minimum-variance unbiased linear aggregate has weights
\[
w^*
=
\frac{\Sigma^{-1}\mathbf 1}
{\mathbf 1^\top\Sigma^{-1}\mathbf 1}.
\]
Its variance is
\[
\frac{1}{\mathbf 1^\top\Sigma^{-1}\mathbf 1}.
\]
\end{PPTtheorem}

\begin{proof}
Minimize $w^\top\Sigma w$ subject to $\mathbf 1^\top w=1$. The Lagrangian is
\[
L(w,\lambda)=w^\top\Sigma w-2\lambda(\mathbf 1^\top w-1).
\]
Stationarity gives $\Sigma w=\lambda\mathbf 1$, hence $w=\lambda\Sigma^{-1}\mathbf 1$. The constraint gives
\[
\lambda=(\mathbf 1^\top\Sigma^{-1}\mathbf 1)^{-1}.
\]
Positive definiteness makes the objective strictly convex, so the solution is unique. Substitution gives the variance formula.
\end{proof}

\begin{PPTwarning}[Bias is not covariance]
\label{PPT:res:17:3}
Covariance-aware weighting does not remove a common bias that is absent from the covariance model. If every reviewer shares bias $b$, the mean squared error contains the irreducible term $b^2$.
\end{PPTwarning}

Under the model, dependence rather than nominal count determines the variance reduction. Whether a particular review process has this covariance structure is an empirical question. The formula alone neither endorses nor discredits a community judgment.


\section{Shared interpretation, review dependence, and identification}
\label{PPT:sec:old25}
\label{PPT:sec:interpretation}

The distinction between proof and proof field also matters when many people assess the same mathematics. An individual can evaluate a theorem through the notation, examples, expectations, and disciplinary assumptions available to them. Different reviewers may share these influences. Agreement therefore need not provide independent evidence, although substantive criticism can still identify a real defect. The mathematical contribution of a model lies in making its assumptions explicit, not in diagnosing the motives of reviewers.

\subsection{Correlated noise and persistent common bias}

The covariance model of Section~\ref{PPT:sec:correlated-review} measures how dependence reduces the effective information in repeated estimates. It cannot remove a common offset. The next proposition puts both effects into one equation.

\begin{PPTproposition}[MSE with a shared interpretive offset]
\label{PPT:res:25:1}
Suppose reviewers report $J_i=\theta+b+\varepsilon_i$, where $\theta$ is the target quantity, $b$ is a fixed shared offset, $\mathbb E\varepsilon_i=0$, $\operatorname{Var}(\varepsilon_i)=\sigma^2$, and $\operatorname{Cov}(\varepsilon_i,\varepsilon_j)=\rho\sigma^2$ for $i\ne j$. For admissible $\rho$, the mean squared error of the reviewer average is
\[
\boxed{\mathbb E(\overline J_n-\theta)^2=b^2+\sigma^2\left(\rho+\frac{1-\rho}{n}\right).}
\]
For fixed $\rho\ge0$, it approaches $b^2+\rho\sigma^2$ as $n\to\infty$.
\end{PPTproposition}
\begin{proof}
The estimation error is $\overline J_n-\theta=b+\overline\varepsilon_n$. Its mean is $b$, and its variance is $\sigma^2\big((1-\rho)/n+\rho\big)$ by summing all diagonal and off-diagonal covariance entries. Mean squared error is variance plus squared bias. Taking the limit completes the proof.
\end{proof}

This formula complements, rather than replaces, the effective reviewer count. Adding reviewers may reduce independent error while leaving the common offset untouched. The existing Python routine \path{shared_bias_mse} computes this model and permits an additional declared variance term for a shared random lens. Its input parameters are model choices, not measured facts about a given intellectual community.

\subsection{When an identity effect can actually be measured}

A claim about identity effects needs a design that can distinguish them from other differences. Hold the mathematical content fixed and let $Y_i(1)$ and $Y_i(0)$ be the fixed potential evaluations of reviewer $i$ under two labels. In a finite set of $N$ reviewers, assign exactly $n_1$ to the first label and $n_0=N-n_1$ to the second, uniformly over assignments with those sizes, where $n_0,n_1>0$. Assume each observed response equals the potential response under its assigned label and does not depend on other assignments. Then the difference of the two sample means is unbiased for $N^{-1}\sum_i(Y_i(1)-Y_i(0))$. Indeed, each reviewer has inclusion probability $n_1/N$ in the first mean and $n_0/N$ in the second, so linearity of expectation gives that identity.

This conclusion concerns the assigned reviewers. Generalizing to a broader population requires a sampling model. Missing responses, shared discussion, or changes to the argument itself can invalidate the simple interpretation. None of these comparisons determines whether the proof is correct. Correctness requires its own mathematical audit.

This distinction is central to a proposed application of Proof Path Theory. One may observe that a presentation changes which branches reviewers explore or how they estimate importance. Those differences become evidence of a contextual effect only after alternative explanations have been addressed. A skeptical review can be correct. Agreement can be justified. The mere presence of community influence neither validates nor invalidates a theorem.

The research opportunity is to connect a typed archive of mathematical claims with controlled studies of how proof paths are selected, compressed, and interpreted. A certified proof core would be held fixed while the accompanying discovery field or contextual labels change. The questions are which alternatives readers recognize, which obstructions they find, and how their subsequent mathematical exploration differs. These are empirical targets for future work, not conclusions established by this finite model.

\subsection{Interpretation and mathematical validity}

The ambient field records exploratory trajectories, the realization theory determines whether specified local models assemble, and the review model describes the effect of declared dependence among numerical judgments. They address different questions. Their common archive makes comparison possible without identifying agreement, coherence, and deductive validity.

These constructions do not yield a general solver for mathematical conjectures, a unique numerical measure of mathematical meaning, a canonical global proof field, or an empirical diagnosis of actual mathematical communities. They provide explicit finite questions with checkable answers and identify the conditions needed for stronger generalizations. Most importantly, they preserve the difference between a valid result and the circumstances through which people discover or evaluate it.



\section{Verification bandwidth and audit queues}
\label{PPT:sec:old18}

Let $a_t\geq0$ be candidate arrivals during audit cycle $t$, let $s_t\geq0$ be service capacity, and let $B_t\geq0$ be the unresolved backlog. Suppose arrivals in a cycle are available to that cycle's service. The basic deterministic queue is
\[
B_{t+1}=\max\{0,B_t+a_t-s_t\}.
\]

\begin{PPTtheorem}[Constant-rate audit regimes]
\label{PPT:res:18:1}
\label{PPT:thm:audit-rates}
Suppose $a_t\equiv\lambda$ and $s_t\equiv\mu$.
\begin{enumerate}[label=\textup{(\roman*)}]
\item If $\lambda<\mu$, every finite initial backlog reaches zero after finitely many cycles.
\item If $\lambda=\mu$, the backlog remains constant.
\item If $\lambda>\mu$, the backlog grows linearly with slope $\lambda-\mu$.
\end{enumerate}
\end{PPTtheorem}

\begin{proof}
Before the truncation at zero, each cycle adds $\lambda-\mu$. When this quantity is negative, the backlog decreases by a fixed amount until it reaches zero. When it is zero, the backlog is unchanged. When it is positive, the truncation is never active after a nonnegative start, and the backlog is $B_0+t(\lambda-\mu)$.
\end{proof}

\begin{figure}[htbp]
\centering
\includegraphics[width=\linewidth]{sources/monographs/proof-path-theory/docs/figures/fig_11_audit_queues.pdf}
\caption{The committed audit example starts with backlog ten and capacity five. Four arrivals per cycle clear the queue, five preserve the backlog, and six increase it. Direct labels replace a separate legend. These are exact trajectories for the declared constant-rate model.}
\label{PPT:fig:auditqueue}
\end{figure}

A simple audit priority score is
\[
P(x)
=
\frac{p_{\mathrm{err}}(x)\,c_{\mathrm{consequence}}(x)
\bigl(1+r_{\mathrm{dependency}}(x)\bigr)}
{c_{\mathrm{audit}}(x)}.
\]
Here $p_{\mathrm{err}}(x)\in[0,1]$, the consequence and dependency factors are nonnegative, and $c_{\mathrm{audit}}(x)>0$. Zero-cost checks can be handled separately. This is a declared priority heuristic, not a probability of truth or a proved optimal scheduling rule.


\FloatBarrier
\section{Recursive proof cubes and expansion depth}
\label{PPT:sec:old19}

A visible proof step is rarely atomic. A lemma opens into definitions and sublemmas. A computation opens into algorithms, parameters, and exact checks. A citation opens into another dependency graph. Proof Path Theory treats this as recursive refinement.

\begin{PPTdefinition}[Proof packet and refinement]
\label{PPT:res:19:1}
A proof packet is a node presented at one resolution. A refinement replaces that node by a finite subfield whose certified target reconstructs the original node on the declared scope.
\end{PPTdefinition}

A refinement preserves the certified target on its declared scope. Preservation of other tasks requires a separate recovery condition. Let $R_d:\mathscr A\to\mathcal R_d$ be the representation scheme at depth $d$. A sufficient condition is a map
\[
\pi_d:R_{d+1}(\mathscr A)\longrightarrow R_d(\mathscr A)
\qquad\text{with}\qquad R_d=\pi_d\circ R_{d+1}.
\]
If a task factors as $q=d_q\circ R_d$, then
\[
q=(d_q\circ\pi_d)\circ R_{d+1}.
\]
Hence $R_d\preceq_{\PPTQcal}R_{d+1}$. Repeated refinements satisfying this condition give the recoverability chain. For a restricted task family, preserving its supported tasks directly can suffice without recovering the entire preceding representation.

A longer derivation is not automatically a richer archive. It may reconstruct the same conclusion while omitting an example, a failed approach, or a provenance record. Those losses must be assessed against the declared task family.

Let $r_d\geq0$ be a declared unresolved-work score and $c_d\geq0$ the cumulative cost. These are model inputs. A small $r_d$ is not, without another theorem, a certificate of a nearly valid proof. The word cube names the layered presentation, not a cubical complex supplied with additional geometric axioms.

\begin{PPTdefinition}[Goal-relative depth objective]
\label{PPT:res:19:2}
For goal $g$ with weights $\alpha_g,\beta_g\ge0$, define
\[
J_g(d)=\alpha_g r_d+\beta_g c_d.
\]
An optimal expansion depth is any minimizer of $J_g$.
\end{PPTdefinition}

\begin{PPTtheorem}[Expansion depth is goal dependent]
\label{PPT:res:19:3}
The same proof object can have different optimal expansion depths for different goals.
\end{PPTtheorem}

\begin{proof}
Take a proof object with exactly two available depths $d_1<d_2$, with $r_{d_2}<r_{d_1}$ and $c_{d_2}>c_{d_1}$. Direct subtraction shows that $d_2$ is preferred exactly when
\[
\alpha_g(r_{d_1}-r_{d_2})>\beta_g(c_{d_2}-c_{d_1}).
\]
The goal $(\alpha_g,\beta_g)=(1,0)$ strictly prefers $d_2$, whereas $(0,1)$ strictly prefers $d_1$. Since these are the only available depths, each preference is a global optimum. If both weights vanish, every depth is optimal. No division by a possibly zero weight is needed.
\end{proof}

A practical stopping rule selects the first depth satisfying
\[
r_d\le\varepsilon\qquad\text{or}\qquad
\frac{r_{d-1}-r_d}{c_d-c_{d-1}}<\tau.
\]
The first condition uses a declared tolerance. The second applies only at $d\geq1$ with $c_d>c_{d-1}$ and compares the most recent reduction with its extra cost. Zero-cost refinements require separate handling. These stopping rules do not generally minimize $J_g$ over all remaining depths.

Figure~\ref{PPT:fig:refinement} and the worked example in Section~\ref{PPT:sec:depth-example} give explicit costs and unresolved-work scores for three goals.

Deeper is not automatically better. A formal checker may need primitive steps. A specialist paper may cite a standard theorem. A teaching account may expand only the conceptual bottleneck. An archive may retain a failed branch that neither certificate requires.



\chapter{Examples, implementation, and boundaries}\label{PPT:book:chapter-54}
The final layer is inspectable: source records, computations, certificates, and evidence. Each check has a declared scope, and the open problems remain separate from the implemented finite theory.

  \PPTopeningfigure{fig_13_repository_workflow.pdf}{The actual publication workflow separates declared records, software checks, certificate reports, formal interfaces, and delivered artifacts. Certification flags describe the finite input model. They do not prove arbitrary statements attached to node labels.}{fig:workflow}
 


\section{Worked finite examples}
\label{PPT:sec:old20}

The accompanying repository contains three canonical JSON proof fields and several auxiliary models. The examples are small enough to inspect by hand and complete enough to exercise the formal distinctions.

\subsection{A fixed AND-DAG}

The field \texttt{examples/and\_dag.json} has required premise $A$, a definition $D$, lemmas $L_1$ and $L_2$, and target $B$. It also contains a worked example, a failed route, and a boundary case. The unique verification core is
\[
\{A,D,L_1,L_2,B\}.
\]
The optional nodes have zero declared verification value for $B$ and positive discovery, meaning, or bridge value. This is the executable instance behind Figure~\ref{PPT:fig:loss}.

\subsection{Two certified proofs}

The field \texttt{examples/alternative\_proofs.json} contains a short route and a richer explanatory route. The short route has storage cost $1$, while the richer route has storage cost $3$ and larger meaning and discovery values. Exact certificate search chooses
\[
A\longrightarrow L_{\mathrm{short}}\longrightarrow B
\]
when the objective is storage cost alone. A teaching compiler may rationally choose the other route because it uses a different objective.

\subsection{Set Cover encoded as certificate selection}

The field \texttt{examples/set\_cover\_certificate.json} contains four optional source lemmas and four obligations. Exact search returns the certificate using $S_1$ and $S_3$, with total cost $2$. The particular optimum is not important. The example demonstrates how local alternative derivations create a global combinatorial selection problem.

\subsection{Bridge retention}

The bridge experiment uses ten candidate nodes and sixteen weighted future targets. For budgets zero through five, the script computes both the exact optimum and the greedy solution. The generated data are stored in \texttt{data/retention\_comparison.csv}. On this instance the two curves coincide. The theorem guarantees only the worst-case $(1-1/e)$ factor.

\subsection{Correlated review and audit capacity}

For one hundred equicorrelated reviewers with $\rho=0.1$,
\[
n_{\mathrm{eff}}
=
\frac{100}{1+99(0.1)}
\approx 9.17.
\]
The audit example begins with backlog $10$. Arrivals $4$ and capacity $5$ clear the queue. Equal rates hold the backlog at $10$. Arrivals $6$ and capacity $5$ increase it by one per cycle.

\begin{table}[H]
\centering
\caption{Canonical finite examples and their mathematical roles.}
\begin{tabular}{@{}lll@{}}
\toprule
Object & Computed result & Mathematical role\\
\midrule
AND-DAG & unique five-node core & Theorem~\ref{PPT:thm:and-core}\\
Alternative field & minimum cost $1$ & exact finite search\\
Set Cover field & minimum cost $2$ & Theorem~\ref{PPT:thm:minimum-certificate} encoding\\
Bridge model & greedy and exact curves & Theorems~\ref{PPT:thm:bridge-submodular}--\ref{PPT:thm:bridge-greedy}\\
Review model & $n_{\mathrm{eff}}$ formula & Theorem~\ref{PPT:thm:review-count}\\
Audit model & three queue regimes & Theorem~\ref{PPT:thm:audit-rates}\\
\bottomrule
\end{tabular}
\end{table}


\subsection{An expansion-depth calculation}
\label{PPT:sec:depth-example}
For depths $0$ through $6$, take unresolved-work scores
$(1,0.55,0.30,0.18,0.12,0.09,0.075)$ and costs
$(0,0.25,0.65,1.2,2.1,3.5,5.5)$. The declared weights
$(\alpha,\beta)=(1,0.75)$, $(2.5,0.35)$, and $(5,0.15)$ give
minimizing depths $1$, $3$, and $4$, respectively. Their minimum objective
values are $0.7375$, $0.87$, and $0.915$. These follow by evaluating
$\alpha r_d+\beta c_d$ at the seven available depths. The labels certificate,
teaching, and archive name these three supplied objectives rather than
measured preferences of actual readers.

\begin{figure}[htbp]
\centering
\includegraphics[width=\linewidth]{sources/monographs/proof-path-theory/docs/figures/fig_12_refinement_goals.pdf}
\caption{Three declared objectives on the same finite depth set. Square markers select the exact minimum: depth one for certificate, three for teaching, and four for archive. The weights define these tasks; they are not estimates of how every reader should study a proof.}
\label{PPT:fig:refinement}
\end{figure}

\Needspace{12\baselineskip}
\subsection{An independently checkable decoder}

Take two archives with the same certificate output but different retained side answers. The following code executes the collision alternative. It then checks the result without running the producer again.
\begin{lstlisting}[language=Python]
from proof_path_theory import factor_task, verify_factorization
encoding = {"archive0": "certificate", "archive1": "certificate"}
task = {"archive0": "0", "archive1": "1"}
result = factor_task(encoding, task)
assert not result.sufficient
assert verify_factorization(encoding, task, result)
\end{lstlisting}
The witness names both archives. It is evidence that this declared task cannot be recovered from this encoding on this domain. It says nothing about whether the shared proof certificate is valid. The package's new worked script also prints the positive decoder, the protected-rule obstruction, and the semantic-cover derivations.

\section{Reference implementation}
\label{PPT:sec:old21}

The repository implements the finite theory as the Python package \texttt{proof\_path\_theory}. The public modules are organized by mathematical responsibility.

\begingroup\small
\begin{longtable}{@{}p{0.24\linewidth}p{0.68\linewidth}@{}}
\toprule
Module & Responsibility\\
\midrule
\endhead
\texttt{model.py} & Typed nodes, justifications, evidence records, acyclicity, ancestry, and finite derivability.\\
\texttt{certificate.py} & Certificate validation, unique AND-DAG cores, exact minimum certificate search, and certificate cost.\\
\texttt{compression.py} & Task sufficiency, recoverability preorder, loss reports, compression debt, representation utilities, and Pareto fronts.\\
\texttt{bridges.py} & Reachability, domination-based criticality, exact retention, greedy retention, and finite submodularity checks.\\
\texttt{review.py} & Equicorrelated covariance, effective reviewer count, generalized least squares, and shared-bias mean squared error.\\
\texttt{audit.py} & Audit queue simulation, rate regimes, and explicit priority scores.\\
\texttt{refinement.py} & Goal-relative refinement decisions, optimal depth, and stopping rules.\\
\texttt{versioning.py} & Affected descendants, append-only checks, ledger snapshots, and provenance digests.\\
\texttt{ambient.py} & Finite declared view fibers, restriction checks, and arrow-cover gluing.\\
\texttt{realization.py} & Exact binary elimination, protected fibers, obstruction witnesses, and observation gain.\\
\texttt{views.py} & Local field models, whole-support gluing, declared countermodels, and intersection diagrams.\\
\texttt{extension.py} & Least models preserving true and false observations, with replayable forbidden-node derivations.\\
\path{semantic_cover.py} & Entailment certificates for redundant supports, countermodels to universal gluing, and checked semantic gluing.\\
\texttt{recovery.py} & Factorization on explicit archive families, collision witnesses, joint task schemes, and decoder composition.\\
\path{trace_codec.py} & Versioned indexed exchange, strict and causal trace replay, and an independent result checker.\\
\texttt{io.py} & Stable JSON import and export for proof fields and certificates.\\
\texttt{cli.py} & Command-line validation, core extraction, minimum certificate search, and certificate audit.\\
\bottomrule
\end{longtable}
\endgroup

\subsection{Command-line interface}

The main commands are:

\begin{lstlisting}[language=bash]
proof-path validate examples/and_dag.json
proof-path core examples/and_dag.json --output core.json
proof-path minimum examples/alternative_proofs.json --output minimum.json
proof-path audit examples/and_dag.json core.json
\end{lstlisting}



\subsection{JSON contract}

A proof field file contains:
\begin{itemize}
\item a schema version and metadata.
\item required premises and targets.
\item typed nodes with value and cost vectors.
\item typed evidence records with scope and certification flags.
\item justifications with premise sets, conclusions, rules, evidence references, and certification flags.
\end{itemize}

The named proof-field schema is intentionally plain. It can be produced by a proof assistant adapter, a research notebook, a theorem ledger, or a hand-written example. The positional trace-exchange format in Section~\ref{PPT:sec:trace-wire} is a separate, versioned contract. Neither an imported certification flag nor a successful schema check establishes semantic validity.

\subsection{Tests and build discipline}

The regression suite covers graph validation, derivability, certificate extraction, alternative proofs, schema boundaries, and input validation. It also exercises loss vectors, task sufficiency, bridge coverage, retention, review covariance, audit queues, refinement, version propagation, JSON round trips, and the command line. The dated verification record gives the executed count for each edition.

A separate deterministic cross-check compares the exact certificate solver with brute-force enumeration over eighteen cost profiles, tests the greedy cardinality guarantee on one hundred generated coverage instances, verifies 1,024 retention-algebra identities, and compares the effective-reviewer formula with direct covariance calculations in 114 parameter settings.

The extension cross-check exhausts all forward rule subsets on three ordered nodes, every admitted-leaf set, every positive and negative observation set, and every node-covering patch family. An independent truth-table calculation checks feasibility and universal gluing. A separate enumeration compares encodings, task fibers, joint task partitions, and decoder composition on archive domains of at most four elements. These are bounded finite comparisons, not empirical models of mathematical practice.

The trace conformance suite adds independently enumerated finite models, canonical JSON records, and malformed or altered certificates. Its actual byte stream is processed by the Lean checker built in the current formal run. The report records each verdict and the exact input and output hashes. Rejection of a malformed record is not counted as a proof of an impossible mathematical request.

The general Lean proofs now establish the specified trace checker and typed codec. The Python producer, Python checker implementation, standard text parser, and source-address adapter retain separate conformance evidence.

The tests and cross-checks verify the implementation of finite definitions and examples. They do not replace the proofs in this monograph. Conversely, the proofs do not guarantee that a software implementation is free of defects. The release keeps both forms of verification.


\section{What the theory establishes}
\label{PPT:sec:old26}

The finite theory establishes the following structural facts.

\begin{enumerate}
\item A proof field can represent conjunction within a justification and alternatives among justifications without confusing candidate edges with certified edges.
\item Derivability in a finite acyclic selected subfield is well defined and computable in topological order.
\item A fixed AND-DAG, expanded only outside its admitted leaves, has a least certificate pair.
\item Alternative derivations make minimum certificate selection NP-complete in general.
\item Compression is meaningful only relative to a task family.
\item Task-relative recoverability forms a preorder and a quotient partial order.
\item Verification loss, discovery loss, meaning loss, bridge loss, and audit loss can diverge.
\item Expected reconstruction debt decreases monotonically as more required material is retained.
\item Bridge criticality is domination in the declared future-reach graph.
\item Weighted future reach is monotone and submodular.
\item Greedy cardinality retention has the standard $(1-1/e)$ guarantee.
\item Append-only certified updates preserve existing derivability.
\item The forward dependency closure is the least set containing the changed nodes and closed under outgoing dependency edges. It marks potential review, not automatic falsehood.
\item Positive review dependence lowers the variance-equivalent count. Under a positive definite covariance model, generalized least squares gives the minimum-variance unbiased linear aggregate.
\item Constant-rate audit queues have an exact stable, critical, or unstable classification.
\item Optimal proof expansion depth depends on the presentation goal.
\end{enumerate}

These results turn the original methodological insight into a finite mathematical branch with explicit objects, algorithms, complexity, and failure boundaries.

The extension results add three further capabilities. A protected rule-model request returns its least feasible model or a derivation of a forbidden node. A cover has universal gluing exactly when its visible rules entail every original rule. A representation admits a task decoder exactly when equal encodings have equal task answers. These outputs are respectively model witnesses, rule derivations, and decoding or collision witnesses. None can replace the others merely because it is called a certificate.


\noindent The realization theory gives the proof field a contravariant diagram of local models and an explicit equivalence with its coherent sections. Whole-support gluing proves extension from suitably covering patches. The protected linear alternative classifies feasible fibers and supplies normalized witnesses for infeasibility. Exact finite-field counts yield an information formula only after a uniform distribution is chosen. The native exporter supplies syntactic dependencies from compiled declarations, not the complete history of discovery.


\section{Limitations and claim boundaries}
\label{PPT:sec:old27}

Proof Path Theory does not establish a theorem by preserving the history of an argument. A richly documented wrong proof remains wrong. Bridge value is not validity. Explanatory power is not validity. Consensus is not validity. The certified subfield remains the object that carries deductive force.

The dependency graph is a declared model. It records the dependencies that were formalized, not every historical influence on discovery. Tacit knowledge, visual intuition, unrecorded conversations, and changes in vocabulary may be absent.

The value and cost coordinates are task models. Verification can often be specified sharply. Meaning and future bridge value are partly prospective. Their use is justified when assumptions are stated, not when the numbers are treated as intrinsic physical constants.

The finite bridge graph does not predict every future theorem. It supports conditional statements of the form: under these declared future targets and transfer edges, this retention policy preserves this amount of modeled reach.

The exact minimum certificate algorithm is exponential in the worst case. The package makes this limitation explicit through a search-state bound.

The review and queue models are deliberately elementary. They isolate dependence and capacity effects, but they do not model every institutional or psychological feature of mathematical review.

Finally, the monograph does not claim that graphs, hypergraphs, proof theory, submodular optimization, covariance-aware aggregation, or queueing recurrences were invented here. The contribution is the formal assembly of these structures around proof compression, discovery fields, mathematical memory, and versioned proof portfolios.


\paragraph{The remaining implementation boundary.} Theorem~\ref{PPT:res:trace:json} proves the meaning of an accepted decoded answer. It does not prove that every named-field export preserves its source, or that the text parser and Python producer always realize that contract. Those are the next refinement obligations, not hypotheses silently supplied by successful tests.

\section{Open theorem program}
\label{PPT:sec:old28}

The finite protected-extension, semantic-cover, and attained-decoder problems developed in this edition are closed under their stated hypotheses. The questions below concern stronger interfaces. In particular, the rule transport sends one source rule to one target rule and is not asserted to be a general equivalence of proof fields. A transport into a derived target proof, executable correctness of the producers, and richer semantic objects remain distinct next steps.


\begin{PPTproblem}[Proof-path view functors and comparison]
\label{PPT:res:28:1}
The declared-field local-model functor and whole-support gluing are established here. Extend the checked interface to richer proof-state representations, and verify the implementation of general proof assistant adapters. Determine when objects with symmetry require nonabelian or higher descent rather than the strict set-valued construction.
\end{PPTproblem}

\begin{PPTproblem}[Counterfactual proof-search observations]
\label{PPT:res:28:2}
Study which observations on a finite realization fiber yield the greatest task-relative reduction of uncertainty per unit cost. Extend the binary rank theorem to controlled infinite or nonlinear settings only after specifying distributions, admissible questions, and the semantics of proof state.
\end{PPTproblem}

\begin{PPTproblem}[Collective evaluation without shared-context confounding]
\label{PPT:res:28:3}
Design preregistered evaluations of identical certified mathematics under randomized labels and alternative explanatory contexts. Separate measurement of interpretation differences from independent verification of correctness, and test whether models of correlated review predict actual audit errors.
\end{PPTproblem}



Further directions include the following problems.

\begin{PPTproblem}[Infinite proof fields]
\label{PPT:res:28:4}
Develop locally finite and countable proof fields with conditions guaranteeing that finite certificates approximate an infinite derivation without losing domain or convergence information.
\end{PPTproblem}

\begin{PPTproblem}[Parameterized certificate complexity]
\label{PPT:res:28:5}
Classify minimum certificate selection by treewidth, number of alternative justifications, justification arity, shared-subproof width, and target count. Identify fixed-parameter tractable regimes and approximation bounds.
\end{PPTproblem}

\begin{PPTproblem}[Probabilistic bridge calibration]
\label{PPT:res:28:6}
Replace declared bridge weights by calibrated posterior models while keeping conjectural reach separate from certified derivability.
\end{PPTproblem}

\begin{PPTproblem}[Proof assistant interoperability]
\label{PPT:res:28:7}
Construct adapters that export stable claim identifiers, primitive dependencies, scopes, and proof terms from existing formal systems into the common proof-field schema.
\end{PPTproblem}

\begin{PPTproblem}[Higher proof transformations]
\label{PPT:res:28:8}
Model transformations between proofs, transformations between transformations, and coherent compiler pipelines in a higher categorical setting.
\end{PPTproblem}

\begin{PPTproblem}[Privacy and attribution]
\label{PPT:res:28:9}
Design proof archives that preserve provenance and priority while allowing private notes, embargoed routes, and sensitive collaboration records to remain access controlled.
\end{PPTproblem}

\begin{PPTproblem}[Empirical validation]
\label{PPT:res:28:10}
Study real mathematical projects to estimate recovery costs, audit bottlenecks, and the predictive value of bridge-retention scores. Such work should test the models rather than assume them.
\end{PPTproblem}

\begin{PPTproblem}[Distributed audit systems]
\label{PPT:res:28:11}
Extend the queue model to heterogeneous reviewers, dependent tasks, rework, priority classes, and claims whose audit cost changes after neighboring claims are resolved.
\end{PPTproblem}



\section{Conclusion}
\label{PPT:sec:old29}

A proof certificate and its surrounding proof field answer different mathematical questions. The field need not consist only of viable derivations. It can retain failed strategies and genuine local refutations without allowing either to masquerade as global unprovability. A checked contravariant system of local field models gives an explicit bridge to Komorebi Space. Whole-support coverage identifies a constraint that pairwise dependencies can miss. The semantic-cover criterion now determines exactly when visible rules recover a missing constraint. In the linear finite case, protected observations admit exact success or failure certificates and a measurable information gain. A separate review model distinguishes proof correctness from correlated social interpretation.



A final proof is indispensable, but it is not always the whole mathematical object. The theory distinguishes the certificate from the declared field in which it is recorded. It gives that field typed nodes, alternative justifications, evidence records, task-relative compressions, bridge graphs, version histories, and explicit audit models.

The central distinction is between derivability within a declared inference system and the wider information retained by its archive. Certified closure records which claims follow from admitted leaves and accepted rules. Its interpretation as mathematical validity depends on the soundness of those leaves and rules. Discovery information serves different tasks: explaining a definition, locating a failed method, preserving a comparison, or opening a route to another problem.

Compression can preserve the present derivation while losing those other capabilities. Proof Path Theory makes that difference explicit through task families, recovery maps, extension problems, and versioned evidence. The certificate and the surrounding field should therefore be maintained together, with their logical roles kept distinct.

The final thesis is therefore:
\[
\boxed{\begin{minipage}{0.88\linewidth}
A proof should be published as a certificate and maintained as a versioned proof field. The certificate preserves deductive validity. The larger field preserves the dependencies, alternatives, examples, failures, provenance, and bridges from which later mathematics may grow.
\end{minipage}}
\]


\FloatBarrier


\part{Constructive repair and exact certificates}
\setcounter{section}{0}\renewcommand{\thesection}{CR.\arabic{section}}\renewcommand{\theHsection}{CR.\arabic{section}}

\chapter{Correction, dimension, and admissibility}\label{CR:book:chapter-55}
\section{Correction and quotient obstructions}\label{CR:cr-source-section-1}\label{CR:ch:basic}
\subsection{The original equation}
A residual describes a failure relative to a stated equation. A correction map describes the changes that are allowed to address that failure. These are separate pieces of data. Their separation prevents the phrase ``the obstruction'' from concealing an unstated choice of permissible operations.

Let $\CRkk$ be a commutative ring. All modules are unital $\CRkk$-modules and all maps in this chapter are $\CRkk$-linear. Fix
\[
D:U\longrightarrow V,
\]
and let $r:I\to V$ be a family of residuals. The index set $I$ need not be finite for the algebraic statements. When finite enumeration is used later, finiteness will be imposed explicitly.

\begin{CRdefinition}[Correction and repair]\label{CR:def:repair}
A residual $r\in V$ is \emph{correctable} by $D$ when there is $u\in U$ such that
\begin{equation}\label{CR:eq:correct}
Du=-r.
\end{equation}
An added correction map $E:W\to V$ \emph{repairs} $r$ when there are $u\in U$ and $w\in W$ with
\begin{equation}\label{CR:eq:repair}
Du+Ew=-r.
\end{equation}
It repairs the family when this condition holds for every $r_i$. The witnesses may depend on $i$.
\end{CRdefinition}

The sign convention makes the repaired residual $r+Du+Ew$ equal to zero. In particular, the target of the equation remains $V$. A calculation in another space can establish \eqref{CR:eq:repair} only if a valid route back to $V$ is supplied.

The family condition does not ask for one common pair $(u,w)$ that corrects every residual at once. It asks for one common added map $E$ and separate coefficients for each residual. Nor does it require those witnesses to vary continuously, measurably, or linearly with the index. Such conditions would define different problems.

\subsection{What the quotient remembers}
Write
\[
Q_D=\CRcoker D=V/\CRim D,\qquad q_D:V\longrightarrow Q_D.
\]
The class $q_D(r)$ records the part of $r$ that cannot be changed by the original correction map. Since $q_D(Du)=0$, all residuals that differ by an original correction have the same class.

\begin{CRtheorem}[Exact obstruction criterion]\label{CR:thm:correctable}
For every $r\in V$,
\[
\exists u\in U\ (Du=-r)\quad\Longleftrightarrow\quad q_D(r)=0.
\]
\end{CRtheorem}
\begin{proof}
Applying $q_D$ to $Du=-r$ gives $0=-q_D(r)$, hence $q_D(r)=0$. Conversely, $q_D(r)=0$ means $r\in\CRim D$. Choose $v\in U$ with $Dv=r$. Then $D(-v)=-r$.
\end{proof}
\CRformal{correctable_iff_obstruction_zero}

This is an exact criterion for the declared linear correction equation. It is not a universal test for a nonlinear equation from which $D$ might have been obtained by differentiation. Section~\ref{CR:cr-source-section-7} gives counterexamples in both directions.

The quotient also has a universal property. If $L:V\to Z$ annihilates $\CRim D$, it descends to a unique linear map $\overline L:Q_D\to Z$ with $L=\overline Lq_D$. Thus every linear detector that ignores correctable directions factors through the same quotient. This is the classical cokernel construction used in Paper III~\cite{CR:paperiii}. Canonicality is relative to $D$, not to an unspecified mathematical problem.

\begin{CRexample}[A visible residual]\label{CR:ex:plane}
Let $D:\CRQ\to\CRQ^2$ be $D(u)=(u,0)$ and let $r=(2,3)$. The quotient is identified with the second coordinate. The correction $u=-2$ removes the first coordinate but leaves $(0,3)$. The functional $L(x,y)=y$ satisfies $LD=0$ and $L(r)=3$, so no original correction eliminates $r$.
\end{CRexample}

\subsection{Repair as an image condition}
An added map $E:W\to V$ contributes precisely the classes in $\CRim(q_DE)$. This leads to the basic feasibility criterion.

\begin{CRtheorem}[A single residual]\label{CR:thm:single}
For every $E:W\to V$ and $r\in V$,
\begin{equation}\label{CR:eq:single}
\exists u,w\ (Du+Ew=-r)
\quad\Longleftrightarrow\quad
q_D(r)\in\CRim(q_DE).
\end{equation}
\end{CRtheorem}
\begin{proof}
Suppose $Du+Ew=-r$. Applying the quotient gives $q_D(Ew)=-q_D(r)$, so $q_D(r)=q_D(E(-w))$ lies in the image.

Conversely, choose $w$ with $q_D(Ew)=q_D(r)$. Then $q_D(r-Ew)=0$, so there is $v\in U$ with $Dv=r-Ew$. Taking $u=-v$ and $w'=-w$ gives
\[
Du+Ew'=-(r-Ew)-Ew=-r.
\]
The witnesses solve the equation in $V$, not just in its quotient.
\end{proof}
\CRformal{repairable_iff_obstruction_mem_range}

\Needspace{8\baselineskip}
\begin{CRcorollary}[Reconstruction from a quotient witness]\label{CR:cor:reconstruct}
If $q_D(Ew)=q_D(r)$, there are $u\in U$ and $w'\in W$ such that
\[
Du+Ew'+r=0.
\]
\end{CRcorollary}
\begin{proof}
Use the witnesses constructed in the reverse direction of \cref{CR:thm:single} and move $r$ to the left.
\end{proof}
\CRformal{reconstruct_cancellation}

\subsection{The obstruction span of a family}
\begin{CRdefinition}[Required obstruction span]\label{CR:def:span}
For a residual family $r:I\to V$, put
\[
O_D(r)=\CRSpan_{\CRkk}\{q_D(r_i):i\in I\}\leq Q_D.
\]
We write $O$ when the family and correction map are understood.
\end{CRdefinition}

Every $q_D(r_i)$ lies in $O$. The use of a span matters because an added linear correction map supplies a submodule of quotient classes, not merely a list of unrelated vectors.

\begin{CRtheorem}[Family feasibility]\label{CR:thm:family}
The map $E:W\to V$ repairs the family $r$ if and only if
\begin{equation}\label{CR:eq:family}
O_D(r)\subseteq\CRim(q_DE).
\end{equation}
\end{CRtheorem}
\begin{proof}
If every residual is repairable, \cref{CR:thm:single} places every generator $q_D(r_i)$ in $\CRim(q_DE)$. The image is a submodule, so it contains their span. Conversely, \eqref{CR:eq:family} places each generator in that image, and \cref{CR:thm:single} supplies its correction witnesses.
\end{proof}
\CRformal{family_repairable_iff_span_le}

For an empty family, $O=0$ and the family condition is vacuous. For a nonempty family whose residuals are already correctable, $O$ is still zero. Both cases should permit a repair of added dimension zero. The next chapter includes them without a special exception.

\begin{scopebox}[title=What has been established]
The quotient determines feasibility only because the correction map and the original equation have been declared. A positive quotient witness can be lifted to actual cancellation in $V$. A disappearance caused by forgetting part of $V$ is a different operation and requires a preservation theorem.
\end{scopebox}

\section{Minimum repair dimension}\label{CR:cr-source-section-2}\label{CR:ch:minimal}
\subsection{A lower bound that every repair must satisfy}
In this chapter $\CRkk$ is a field. The added space $W$ is finite dimensional when its numerical dimension is compared. The attainment result itself uses a linear section and does not require all ambient spaces to be finite dimensional.

\begin{CRtheorem}[Dimension lower bound]\label{CR:thm:lower}
If $E:W\to V$ repairs $r$ and $W$ is finite dimensional, then $O=O_D(r)$ is finite dimensional and
\[
\dim_{\CRkk}O\leq\dim_{\CRkk}W.
\]
\end{CRtheorem}
\begin{proof}
By \cref{CR:thm:family}, $O\subseteq\CRim(q_DE)$. The image of a linear map from a finite dimensional space is finite dimensional and has dimension at most that of its domain. Hence
\[
\dim O\leq\dim\CRim(q_DE)\leq\dim W.\qedhere
\]
\end{proof}
\CRformal{repair_dimension_lower_bound}

The bound measures missing independent directions, not the number of residuals. A large family can have a one dimensional obstruction span. Conversely, a single residual never requires more than one freely chosen direction, since its own negative may be made available. These statements concern unconstrained linear corrections and say nothing about the cost of producing a permitted direction in a domain application.

\subsection{An attaining construction}
\begin{CRlemma}[A section of the quotient]\label{CR:lem:section}
There is a linear map $s:Q_D\to V$ with
\[
q_Ds=\CRid_{Q_D}.
\]
\end{CRlemma}
\begin{proof}
Choose a basis of the vector space $Q_D$. Since $q_D$ is surjective, choose a representative in $V$ of each basis vector. Extend this assignment linearly. Its composite with $q_D$ fixes the basis and therefore fixes all of $Q_D$.
\end{proof}
\CRformal{exists_obstruction_section}

This choice of section is generally nonunique. Changing a representative by a vector in $\CRim D$ produces another valid section. The construction is an existence proof using classical choice, not an algorithm for selecting a basis.

\begin{CRtheorem}[Attainment on the obstruction span]\label{CR:thm:attain}
Let $\iota_O:O\hookrightarrow Q_D$ be inclusion. There exists a linear map $E:O\to V$ such that
\begin{equation}\label{CR:eq:section-on-O}
q_DE=\iota_O
\end{equation}
and $E$ repairs the family.
\end{CRtheorem}
\begin{proof}
Choose the section $s$ from \cref{CR:lem:section} and set $E=s\iota_O$. Then $q_DE=q_Ds\iota_O=\iota_O$. Every $z\in O$ is therefore in $\CRim(q_DE)$. Family feasibility follows from \cref{CR:thm:family}.
\end{proof}
\CRformal{exists_attaining_repair}

\[
\begin{tikzcd}[column sep=large,row sep=large]
O \arrow[r,"E"] \arrow[dr,hook,"\iota_O"'] & V \arrow[d,"q_D"]\\
& Q_D
\end{tikzcd}
\]
The triangle records exactly what is preserved: a required quotient class is lifted to a representative of that same class. The lift does not manufacture an unrelated object with a similar invariant.

\begin{CRtheorem}[The attained free minimum]\label{CR:thm:minimum}
If $V$ is finite dimensional, the minimum possible dimension of a finite dimensional added space that repairs the family is
\begin{equation}\label{CR:eq:minimum}
\min\bigl\{\dim W: E:W\to V\text{ repairs }r\bigr\}=\dim O_D(r).
\end{equation}
\end{CRtheorem}
\begin{proof}
Every admissible term on the left is at least $\dim O$ by \cref{CR:thm:lower}. Because $V$ is finite dimensional, its quotient and the subspace $O$ are finite dimensional. The map in \cref{CR:thm:attain} is a repair with domain $O$, so the bound is attained.
\end{proof}
\CRformal{free_minimal_repair_rank}

No desired existence statement occurs among the inputs. The theorem proves nonemptiness of the optimization class as well as its lower bound.

\subsection{Directions restricted to a subspace}
An application usually prohibits some corrections. A first useful constraint is a subspace $A\leq V$ containing every permitted added direction. Its quotient image is
\[
q_D(A)=\CRim\bigl(q_D\vert_A:A\to Q_D\bigr).
\]

\begin{CRtheorem}[Attainment inside a legal subspace]\label{CR:thm:legal}
If $O\subseteq q_D(A)$, there exists a linear map $E:O\to V$ satisfying
\[
\CRim E\subseteq A,\qquad q_DE=\iota_O,
\]
and repairing the entire family. If $V$ is finite dimensional, its added dimension is optimal even among all unrestricted repairs.
\end{CRtheorem}
\begin{proof}
Consider the surjection
\[
f:A\longrightarrow q_D(A),\qquad a\longmapsto q_D(a).
\]
Choose a linear section $s_A:q_D(A)\to A$. Let $j:O\hookrightarrow q_D(A)$ and $\iota_A:A\hookrightarrow V$ be inclusions. Define
\[
E=\iota_As_Aj.
\]
Its image is in $A$ by construction. Moreover, $q_D\iota_As_Aj=j$ as a map into $Q_D$, which is precisely $\iota_O$. By \cref{CR:thm:family}, $E$ repairs the family. The lower bound of \cref{CR:thm:lower} proves optimality when dimensions are finite.
\end{proof}
\CRformal{admissible_subspace_attainment}

\begin{CRtheorem}[Subspace impossibility]\label{CR:thm:legal-no-go}
If $O\not\subseteq q_D(A)$, no added map $E:W\to V$ with image contained in $A$ repairs the family.
\end{CRtheorem}
\begin{proof}
Such a map satisfies $\CRim(q_DE)\subseteq q_D(A)$. If it repaired the family, \cref{CR:thm:family} would imply $O\subseteq\CRim(q_DE)\subseteq q_D(A)$, contrary to the hypothesis.
\end{proof}
\CRformal{admissible_subspace_no_go}

The word ``legal'' means membership in a specified linear subspace. It does not mean that the paper has constructed an arithmetic host, an equivariant representation, or a topology on that space. Those interpretations require separate witnesses.

\subsection{Matrix coordinates and a boundary over rings}
For finite matrices, the abstract minimum has a familiar coordinate form.

\begin{CRcorollary}[Rank difference in matrix coordinates]\label{CR:cor:rank-difference}
Let the columns of a matrix $\mathsf R$ be a finite residual family and let $\mathsf D$ represent $D$. Then
\[
\dim O_D(r)=\CRrank[\mathsf D\mid\mathsf R]-\CRrank\mathsf D.
\]
\end{CRcorollary}
\begin{proof}
Set $M=\CRim D+\CRSpan\{r_i\}$. The restriction $q_D\vert_M$ has image $O$ and kernel $\CRim D$. Rank and nullity give $\dim M=\dim\CRim D+\dim O$. The columns of $[\mathsf D\mid\mathsf R]$ span $M$.
\end{proof}
\CRsourcescope{This coordinate corollary is an ordinary derivation included for the computational interpretation. The v0.2.0 Lean release proves the abstract attained dimension theorem, but does not yet identify this matrix formula with the Python rank implementation.}

The field hypothesis in the section argument is substantive. Over $\CRZ$, take $D:\CRZ\to\CRZ$ to be multiplication by two. The quotient map $q:\CRZ\to\CRZ/2\CRZ$ has no $\CRZ$-linear section. A section would send the nonzero class to an integer $n$ with $2n=0$, forcing $n=0$ and contradicting $q(n)=1$. This elementary example explains why the module feasibility theorem does not automatically imply a splitting theorem over arbitrary rings.



\section{Finite menus and conservation}\label{CR:cr-source-section-3}\label{CR:ch:menu}
\subsection{A menu is an actual coefficient map}
Return to a commutative ring $\CRkk$ and modules $U,V$. Let $e:J\to V$ be a family of candidate directions and let $S\subseteq J$ be finite. A coefficient vector is a function $a:S\to\CRkk$. Define
\begin{equation}\label{CR:eq:menu}
E_S:\CRkk^S\longrightarrow V,\qquad E_S(a)=\sum_{j\in S}a_je_j.
\end{equation}
Finite addition and scalar multiplication show that $E_S$ is linear. The coefficient domain is the selected set itself, so no unselected direction appears in the equation.

Write
\[
O_S=\CRSpan_{\CRkk}\{q_D(e_j):j\in S\}.
\]

\begin{CRlemma}[Quotient image of a menu]\label{CR:lem:menu-range}
One has $\CRim(q_DE_S)=O_S$.
\end{CRlemma}
\begin{proof}
For every coefficient vector $a$,
\[
q_D(E_Sa)=\sum_{j\in S}a_jq_D(e_j),
\]
which lies in $O_S$. Conversely, every generator $q_D(e_j)$ is the image of the coefficient vector that equals one at $j$ and zero elsewhere. Since the image is a submodule, it contains their span.
\end{proof}
\CRformal{menu_quotient_range}

\begin{CRtheorem}[Exact finite menu feasibility]\label{CR:thm:menu}
The selected menu $S$ repairs the family exactly when
\begin{equation}\label{CR:eq:menu-criterion}
O_D(r)\subseteq O_S.
\end{equation}
Equivalently,
\[
\forall i\in I\ \exists u_i\in U\ \exists a_i\in\CRkk^S,
\qquad Du_i+\sum_{j\in S}a_{ij}e_j=-r_i.
\]
\end{CRtheorem}
\begin{proof}
Apply \cref{CR:thm:family} to the actual map $E_S$ and use \cref{CR:lem:menu-range}. Its reverse direction supplies the original correction witnesses and the selected coefficients, rather than only an abstract statement about dimensions.
\end{proof}
\CRformal{finite_menu_repair_iff_span}

\Needspace{8\baselineskip}
\begin{CRproposition}[Monotonicity of feasibility]\label{CR:prop:menu-mono}
If $S\subseteq T$, then $O_S\subseteq O_T$. Every family repairable with $S$ is therefore repairable with $T$.
\end{CRproposition}
\begin{proof}
Each generator defining $O_S$ is also a generator defining $O_T$. Apply \cref{CR:thm:menu}. Equivalently, extend each coefficient vector by zero on $T\setminus S$.
\end{proof}
\CRformal{repairable_menu_mono}

An external admissibility condition need not be monotone. A rule allowing at most one selected direction can reject an enlarged menu even though algebraic feasibility only improves. Cost and admissibility must not be confused with span inclusion.

\subsection{One obstruction, two selected entries}
\begin{CRexample}[The difference between the two minima]\label{CR:ex:menu}
Over $\CRQ$, let
\[
D(u)=(u,0,0),\qquad r=(0,1,0),
\]
and make only the following directions available:
\[
e_0=(0,1,1),\qquad e_1=(0,0,1).
\]
The quotient by $\CRim D$ is the plane of the second and third coordinates. There $q_D(r)=(1,0)$, $q_D(e_0)=(1,1)$, and $q_D(e_1)=(0,1)$.
\end{CRexample}

Using only $e_0$, the repair equation reads
\[
(u,a,a)=(0,-1,0).
\]
The second coordinate forces $a=-1$ and the third forces $a=0$, a contradiction. Using only $e_1$, the second coordinate stays zero and cannot become $-1$. Both directions together give the exact witness
\begin{equation}\label{CR:eq:two-entry-witness}
D(0)-e_0+e_1=(0,-1,0)=-r.
\end{equation}
The empty menu also fails. Its correction range is the first coordinate axis.

\begin{center}
\begin{tabular}{@{}lll@{}}
\toprule
Selected set & Quotient directions & Repairs $r$?\\
\midrule
$\varnothing$ & $0$ & No\\
$\{0\}$ & $\CRSpan\{(1,1)\}$ & No\\
$\{1\}$ & $\CRSpan\{(0,1)\}$ & No\\
$\{0,1\}$ & $\CRQ^2$ & Yes\\
\bottomrule
\end{tabular}
\end{center}

The obstruction span has dimension one. A freely chosen direction $r$ repairs the residual with coefficient $-1$. The prescribed menu nevertheless requires two selected entries. Both statements are formalized in the examples module. The difference is not a contradiction: one cost counts independent added dimensions, the other counts named entries that must be purchased or enabled.

A weighting makes the distinction sharper. Assigning positive costs $c_0,c_1$ gives optimum $c_0+c_1$ for this menu. A subspace containing $e_0-e_1$ does not charge those two entries separately unless its cost model explicitly does so.

\subsection{Preserved linear invariants}
Conservation laws supply impossibility certificates without enumerating every correction.

\begin{CRproposition}[Invariant preservation]\label{CR:prop:invariant}
Let $L:V\to Z$ be linear. If $LD=0$ and $LE=0$, then
\[
L(v+Du+Ew)=L(v)
\]
for all $v,u,w$.
\end{CRproposition}
\begin{proof}
Linearity gives $L(v+Du+Ew)=L(v)+LD(u)+LE(w)$. The last two terms vanish.
\end{proof}
\CRformal{invariant_preserved}

\begin{CRtheorem}[Conservation obstruction]\label{CR:thm:conservation}
If $LD=0$, $LE=0$, and $L(r)\ne0$, then $E$ does not repair $r$.
\end{CRtheorem}
\begin{proof}
A repaired residual would be $r+Du+Ew=0$. Applying \cref{CR:prop:invariant} would give $L(r)=L(0)=0$, contrary to the hypothesis.
\end{proof}
\CRformal{conservation_no_go}

\begin{CRcorollary}[Conservation for a selected menu]\label{CR:cor:menu-conservation}
Suppose $LD=0$ and $L(e_j)=0$ for every $j\in S$. If $L(r)\ne0$, no coefficients from $S$ repair $r$.
\end{CRcorollary}
\begin{proof}
For every $a:S\to\CRkk$,
\[
L(E_Sa)=\sum_{j\in S}a_jL(e_j)=0.
\]
Thus $LE_S=0$, and \cref{CR:thm:conservation} applies.
\end{proof}
\CRformal{conservation_menu_no_go}

For instance, let $D(u)=(u,0,0)$ and permit additional directions only on the third coordinate axis. The invariant $L(x,y,z)=y$ annihilates every legal correction. It detects $r=(0,1,0)$, so no legal repair exists. Enlarging the number of copies of the third coordinate direction cannot help.

The theorem is only as meaningful as the invariant it uses. In a domain application one must establish that $L$ is the intended native invariant and that every permitted transformation really preserves it. An abstract functional with a suggestive name is not such a proof.


\chapter{Faithful transport and exact search}\label{CR:book:chapter-56}
\section{Faithful transport of obstructions}\label{CR:cr-source-section-4}\label{CR:ch:transport}
\subsection{The induced quotient map}
Suppose two correction systems are related by a commuting square
\begin{equation}\label{CR:eq:square}
\begin{tikzcd}[column sep=large,row sep=large]
U \arrow[r,"D"] \arrow[d,"a"'] & V \arrow[d,"b"]\\
U' \arrow[r,"D'"'] & V'.
\end{tikzcd}
\qquad bD=D'a.
\end{equation}
All four spaces are modules over the same commutative ring. Denote the quotient maps by $q_D$ and $q_{D'}$.

\begin{CRlemma}[The correction image descends]\label{CR:lem:range-maps}
The map $b$ sends $\CRim D$ into $\CRim D'$.
\end{CRlemma}
\begin{proof}
If $v=Du$, then $b(v)=bD(u)=D'a(u)$, which lies in $\CRim D'$.
\end{proof}
\CRformal{correction_range_maps}

It follows that there is a well defined linear map
\[
B:Q_D\longrightarrow Q_{D'},\qquad B(q_D(v))=q_{D'}(b(v)).
\]
Indeed, changing $v$ by an element of $\CRim D$ changes $b(v)$ by an element of $\CRim D'$. This is the quotient map constructed in the formal source, not an assumed function on unidentified obstruction spaces.

\begin{CRproposition}[Naturality]\label{CR:prop:naturality}
For every residual $r\in V$,
\begin{equation}\label{CR:eq:naturality}
B(q_D(r))=q_{D'}(b(r)).
\end{equation}
\end{CRproposition}
\begin{proof}
This is the defining equation of the induced quotient map.
\end{proof}
\CRformal{transport_obstruction_naturality}

Naturality transports a correctable residual to a correctable one. The converse can fail because $B$ may annihilate a nonzero class. A second condition is needed to use target vanishing as a source certificate.

\subsection{Faithfulness on precisely the required data}
\begin{CRdefinition}[Faithfulness on a submodule]\label{CR:def:faithful}
For $O\leq Q_D$, the map $B$ is \emph{faithful on $O$} when
\[
\forall z\in O,\qquad Bz=0\Longrightarrow z=0.
\]
\end{CRdefinition}

This condition is equivalent to $\ker B\cap O=0$. It does not require injectivity on all of $Q_D$ when only a smaller family of obstruction classes matters.

\begin{CRproposition}[Equivalent injectivity criterion]\label{CR:prop:faithful-inj}
The map $B$ is faithful on $O$ if and only if its restriction $B\vert_O:O\to Q_{D'}$ is injective.
\end{CRproposition}
\begin{proof}
If $B$ is faithful and $Bx=By$ for $x,y\in O$, then $x-y\in O$ and $B(x-y)=0$. Faithfulness gives $x=y$. Conversely, if $B\vert_O$ is injective and $Bz=0=B0$ for $z\in O$, injectivity gives $z=0$.
\end{proof}
\CRformal{faithfulOn_iff_injective}

\begin{CRtheorem}[Reflection of correctability]\label{CR:thm:faithful}
Let $O=O_D(r)$. If $B$ is faithful on $O$, then for every $i\in I$,
\[
b(r_i)\text{ is correctable by }D'
\quad\Longleftrightarrow\quad
r_i\text{ is correctable by }D.
\]
\end{CRtheorem}
\begin{proof}
By \cref{CR:thm:correctable}, target correctability means $q_{D'}(b(r_i))=0$. Naturality identifies this class with $B(q_D(r_i))$. Because $q_D(r_i)\in O$, faithfulness then yields $q_D(r_i)=0$, which is source correctability. Conversely, if $q_D(r_i)=0$, naturality and linearity give $q_{D'}(b(r_i))=B0=0$.
\end{proof}
\CRformal{transport_faithfulness_on_obstruction_span}

\begin{CRproposition}[A witness to information loss]\label{CR:prop:loss-witness}
If $B$ is not faithful on $O$, there is $z\in O$ such that
\[
z\ne0,\qquad Bz=0.
\]
\end{CRproposition}
\begin{proof}
Negating the quantified implication in \cref{CR:def:faithful} gives an element of $O$ whose image is zero and which is itself nonzero.
\end{proof}
\CRformal{nonfaithful_witness}

The witness may be a linear combination of the supplied obstruction classes rather than one particular generator. This is why injectivity on the entire required span is a stronger preservation test than checking that a short list of nonzero classes individually remains nonzero.

\subsection{A transport that conceals failure}
\begin{CRexample}[Correct in the target, false in the source]\label{CR:ex:loss}
Let
\[
D:\CRQ\to\CRQ^2,\quad D(u)=(u,0),\qquad
D'=\CRid_{\CRQ},\quad a=\CRid_{\CRQ},\quad b(x,y)=x.
\]
Then $bD=D'a$. For $r=(0,1)$ there is no $u$ with $(u,0)=(0,-1)$, so the source residual is not correctable. Its image $b(r)=0$ is correctable by $D'$.
\end{CRexample}

Here $Q_D$ is one dimensional and $Q_{D'}=0$. The induced map $B$ annihilates every source obstruction. It is therefore not faithful on the span of $q_D(r)$. The formal example proves both the false positive produced by projection and the failure of the faithfulness hypothesis.

A second elementary calculation clarifies why faithfulness is imposed on a span. The map $B:\CRQ^2\to\CRQ$, $B(x,y)=x+y$, sends each standard basis vector to the nonzero number one. Nevertheless, it kills their difference $(1,-1)$. Survival of generators is not injectivity on their span. This calculation is explanatory and has no separate declaration in the present formal source.

\begin{scopebox}[title=Transport is not repair]
A repair changes the allowed correction image while keeping the original equation meaningful. A transport changes the representation of that equation. The two may be composed, but vanishing in the transported space certifies source correctability only when the stated preservation hypotheses hold.
\end{scopebox}

\section{Certified finite search}\label{CR:cr-source-section-5}\label{CR:ch:search}
\subsection{The finite contract}
The algebraic existence theorems need not be executable. A linear section of a quotient may be chosen by classical choice, and a feasibility statement over an arbitrary infinite field need not come with a terminating decision procedure. The finite reference solver uses a different contract.

Assume $\CRkk$, $U$, $I$, and $J$ have specified finite enumerations. The ring $\CRkk$ is commutative, $U,V$ are $\CRkk$-modules, and equality in $V$ is decidable. The target module $V$ need not itself have a finite enumeration. Fix a linear order on $J$ so that the candidate universe can be enumerated in a definite order.

For each finite subset $S\subseteq J$, define the Boolean test
\begin{equation}\label{CR:eq:boolean-check}
\operatorname{check}(S)=\operatorname{decide}\left(
\forall i\in I\ \exists u\in U\ \exists a\in\CRkk^S,\quad Du+E_Sa=-r_i
\right).
\end{equation}
All quantified domains are finite. Decidable equality in $V$ determines each candidate equation.

\begin{CRproposition}[The Boolean test has the intended meaning]\label{CR:prop:check}
The equality $\operatorname{check}(S)=\mathrm{true}$ holds if and only if $E_S$ repairs the residual family.
\end{CRproposition}
\begin{proof}
The Boolean is defined by deciding precisely the quantified equations in \cref{CR:def:repair} for the map $E_S$. No intermediate rank guess or unproved feasibility oracle replaces those equations.
\end{proof}
\CRformal{repairMenuCheck_spec}

Additional admissibility is a supplied function $\operatorname{allowed}:\mathcal P(J)\to\{\mathrm{false},\mathrm{true}\}$. The cost is any supplied function $c:\mathcal P(J)\to\CRN$. The formal theorem does not require $c$ to be additive in the selected entries.

\subsection{Enumerating menus and selecting an optimum}
Enumerate $J$ as a sorted list. Take every sublist and convert it to a finite set. The resulting list, denoted $\mathcal L_J$, contains every subset of $J$.

\begin{CRlemma}[Coverage of the menu universe]\label{CR:lem:allmenus}
Every finite subset $S\subseteq J$ occurs in $\mathcal L_J$.
\end{CRlemma}
\begin{proof}
Filter the sorted full list by membership in $S$. The result is a sublist of the full list and therefore appears among its sublists. Converting this filtered list to a finite set gives exactly $S$.
\end{proof}
\CRformal{mem_allMenus}

For a finite list $\mathcal L$ of candidates, a Boolean predicate $h$, and a natural cost $c$, define $\operatorname{bestByCheck}(h,c,\mathcal L)$ to filter by $h$ and then return an argument of minimum cost. It returns \CRcode{none} exactly when the filtered list is empty. Ties need not define a mathematically unique optimum.

\begin{CRlemma}[Finite selection]\label{CR:lem:best}
If this selector returns $x$, then $x\in\mathcal L$, $h(x)=\mathrm{true}$, and $c(x)\leq c(y)$ for every $y\in\mathcal L$ with $h(y)=\mathrm{true}$. It returns \CRcode{none} if and only if no listed candidate passes $h$.
\end{CRlemma}
\begin{proof}
Every element of the filtered list is an original list member satisfying $h$. A minimum cost element of that finite list has cost no larger than that of any other member. The selector returns no element exactly when the filtered list is empty, which is exactly the stated rejection condition.
\end{proof}
\CRformal{bestByCheck_spec}

The repair optimizer is
\[
\operatorname{bestRepairMenu}
=\operatorname{bestByCheck}\bigl(
S\mapsto\operatorname{allowed}(S)\wedge\operatorname{check}(S),\ c,\ \mathcal L_J\bigr).
\]

\begin{CRtheorem}[Soundness and minimum cost]\label{CR:thm:optimizer}
If the optimizer returns \CRcode{some S}, then $S$ is allowed, $E_S$ repairs the family, and
\[
c(S)\leq c(T)
\]
for every allowed subset $T\subseteq J$ whose menu map repairs the family.
\end{CRtheorem}
\begin{proof}
Apply \cref{CR:lem:best}. Passing the conjunction gives both admissibility and $\operatorname{check}(S)=\mathrm{true}$. Proposition~\ref{CR:prop:check} converts the latter into actual repair feasibility. Every competing $T$ occurs in $\mathcal L_J$ by \cref{CR:lem:allmenus}, so the same finite selection lemma compares its cost to $c(S)$.
\end{proof}
\CRformal{menu_optimality_certificate}

\begin{CRtheorem}[Complete rejection within the declared universe]\label{CR:thm:none}
The optimizer returns \CRcode{none} if and only if no allowed subset of $J$ repairs the family.
\end{CRtheorem}
\begin{proof}
If it returns \CRcode{none}, \cref{CR:lem:best} excludes every passing member of $\mathcal L_J$. Any allowed feasible subset would occur there by \cref{CR:lem:allmenus} and pass the Boolean test by \cref{CR:prop:check}, a contradiction. Conversely, if no allowed subset repairs the family, no listed subset satisfies the filter, so the selector returns \CRcode{none}.
\end{proof}
\CRformal{menu_search_none_iff}

The output is a selected subset or a rejection, not an extracted table of correction coefficients. The correctness theorem establishes existence of the coefficients. A separately implemented witness extraction or serialized certificate checker would provide those coefficients as a returned data object.

\subsection{A fully evaluated example over two elements}\label{CR:sec:f2}
Take the data of \cref{CR:ex:menu} over $\CRF_2$. Negation is now the identity. The equation is
\[
(u,a,a+b)=(0,1,0).
\]
Thus the full menu succeeds with $u=0$, $a=b=1$. The two singleton menus still fail.

\begin{center}
\begin{tabular}{@{}lcll@{}}
\toprule
Menu & Cardinality & All menus allowed & At most one entry\\
\midrule
$\varnothing$ & 0 & Infeasible & Infeasible\\
$\{0\}$ & 1 & Infeasible & Infeasible\\
$\{1\}$ & 1 & Infeasible & Infeasible\\
$\{0,1\}$ & 2 & Feasible & Not allowed\\
\bottomrule
\end{tabular}
\end{center}

With cardinality as the cost, the two evaluations are
\begin{Verbatim}[fontsize=\small,frame=single,framesep=8pt,rulecolor=\color{RuleGray}]
All subsets permitted:                  some {0, 1}
Only subsets with cardinality at most 1: none
\end{Verbatim}
These equalities are the statements of
\CRdecl{executable_menu_returns_both} and
\CRdecl{executable_restricted_menu_rejects}. Their source proofs use kernel reduction. The generic optimizer theorem then proves that every feasible subset has cardinality at least two, recorded as \CRdecl{executable_menu_is_globally_optimal}. Printed evaluation output is kept distinct from the proof of that output.

\subsection{What the algorithm does not promise}
The reference solver is deliberately finite and exhaustive. With $t=|J|$, there are $2^t$ candidate subsets. For a subset of size $s$, each residual has at most $|U|\,|\CRkk|^s$ candidate coefficient and correction assignments to examine. Summing this upper bound over all subsets gives
\[
|I|\,|U|\sum_{s=0}^t\binom ts|\CRkk|^s
=|I|\,|U|(1+|\CRkk|)^t
\]
possible equation tests before accounting for short circuiting and inadmissible subsets. This is an elementary enumeration bound, not a formal runtime theorem about the compiled Lean implementation.

The Python prototype uses exact Gaussian elimination for its linear equations. The finite Lean solver instead decides those equations by finite enumeration. They share a mathematical specification in the relevant finite cases, but the present release does not contain a proof identifying their implementations. In particular, the Lean optimizer is not an executable search over arbitrary rational coefficient assignments.

\section{Arithmetic digit repair}\label{CR:cr-source-section-6}\label{CR:ch:arithmetic}
\subsection{An exact correction at one scale}
Arithmetic repair adds a second kind of constraint: the new correction must improve divisibility while retaining the previous approximation. A quadratic equation makes the mechanism visible without appealing to a general existence theorem.

The results in this chapter are first stated over $\CRZ$. An integer $p$ need not be prime for the finite divisibility implications. The interpretation as a $p$-adic lift uses a positive prime $p$.

\begin{CRlemma}[A linear digit from B\'ezout data]\label{CR:lem:digit}
Let $p,A,c\in\CRZ$. If there are $s,b\in\CRZ$ with $sA+bp=1$, then there is $t\in\CRZ$ such that
\[
p\mid c+At.
\]
\end{CRlemma}
\begin{proof}
Set $t=-sc$. Then
\[
c+At=c(1-As)=pbc.\qedhere
\]
\end{proof}
\CRformal{Arithmetic.linear_digit_exists}

\begin{CRlemma}[Quadratic correction identity]\label{CR:lem:quad}
Let $p,m,x,a,c,t\in\CRZ$ satisfy
\[
x^2-a=mc,\qquad p\mid c+2xt,\qquad p\mid m.
\]
Then
\[
mp\mid (x+mt)^2-a.
\]
\end{CRlemma}
\begin{proof}
Write $c+2xt=pd$ and $m=ps$. Direct expansion gives
\begin{align*}
(x+mt)^2-a
&=(x^2-a)+2mxt+m^2t^2\\
&=m(c+2xt)+m^2t^2\\
&=mp(d+st^2).\qedhere
\end{align*}
The last expression displays the required quotient explicitly.
\end{proof}
\CRformal{Arithmetic.quadratic_digit_correction}

The term $m^2t^2$ is not discarded as an approximation. Its divisibility follows from $p\mid m$. This is what makes a linear digit equation sufficient for the finite quadratic step.

\subsection{The quadratic lifting theorem}
\begin{CRtheorem}[One higher precision]\label{CR:thm:hensel}
Let $p,x,a\in\CRZ$ and $k\in\CRN$ with $k\geq1$. Assume
\[
p^k\mid x^2-a,\qquad \gcd(2x,p)=1.
\]
Then there exists $t\in\CRZ$ such that
\begin{equation}\label{CR:eq:hensel}
p^{k+1}\mid (x+p^kt)^2-a.
\end{equation}
\end{CRtheorem}
\begin{proof}
Choose $c\in\CRZ$ with $x^2-a=p^kc$. Coprimality supplies B\'ezout data for $2x$ and $p$, so \cref{CR:lem:digit} gives $t$ with $p\mid c+2xt$. Since $k\geq1$, one has $p\mid p^k$. Apply \cref{CR:lem:quad} with $m=p^k$.
\end{proof}
\CRformal{Arithmetic.quadratic_hensel_step}

The formal hypothesis is \CRcode{IsCoprime (2 * x) p}, expressed by B\'ezout data. The theorem constructs an integer representative of a correction class. It does not state that this representative is already the unique digit between zero and $p-1$, nor that an infinite succession of such choices has been constructed in a complete field.

\subsection{An exact seven-adic trace}\label{CR:sec:trace}
For $x^2=2$ at $p=7$, start with $x_1=3$. At each step, let
\[
c_k=\frac{x_k^2-2}{7^k},\qquad
c_k+2x_kt_k\equiv0\pmod7,\qquad
x_{k+1}=x_k+7^kt_k.
\]
The following table records normalized choices for this example. Every displayed digit equation is an exact integer calculation.

\begin{center}
\begin{tabular}{@{}rrrrr@{}}
\toprule
$k$ & $x_k$ & $c_k\bmod7$ & $t_k$ & $x_{k+1}$\\
\midrule
1 & 3 & 1 & 1 & 10\\
2 & 10 & 2 & 2 & 108\\
3 & 108 & 6 & 6 & 2,166\\
4 & 2,166 & 1 & 1 & 4,567\\
5 & 4,567 & 2 & 2 & 38,181\\
6 & 38,181 & 1 & 1 & 155,830\\
\bottomrule
\end{tabular}
\end{center}

All representatives are congruent to three modulo seven, so $2x_k\equiv6\pmod7$ remains a unit. The formal statement \CRdecl{Arithmetic.hensel_seven_trace} verifies the seven congruences
\[
7^k\mid x_k^2-2\qquad (1\leq k\leq7)
\]
for the displayed representatives. It is a finite conjunction of exact divisibility facts. The general one step theorem is proved separately.

For example,
\[
10^2-2=98=7^2\cdot2,
\]
and the next digit satisfies $2+2\cdot10\cdot2=42$. Hence
\[
108=10+7^2\cdot2,
\qquad 108^2-2=7^3\cdot34.
\]
The new digit corrects the residual quotient at precisely the next required modulus.

\subsection{Why a singular derivative is inconclusive}
\begin{CRproposition}[A modular solution that does not lift]\label{CR:prop:mod4}
The equation $x^2=2$ has no solution in $\CRZ/4\CRZ$.
\end{CRproposition}
\begin{proof}
The four residue classes have squares $0,1,0,1$, respectively. None has square two.
\end{proof}
\CRformal{Arithmetic.singular_mod_four_impossible}

Nevertheless, $x=0$ solves the equation modulo two. Its two possible lifts modulo four are zero and two, and both fail. Thus existence at one precision is not itself an existence theorem at the next.

\begin{CRproposition}[Singularity does not prohibit lifting]\label{CR:prop:singular-success}
For every $k\in\CRN$,
\[
2^k\mid 1^2-1.
\]
\end{CRproposition}
\begin{proof}
The residual is zero, which is divisible by every integer.
\end{proof}
\CRformal{Arithmetic.singular_derivative_can_still_lift}

For the equation $x^2-1=0$, the derivative at $x=1$ is two and is not a unit modulo two. The exact root still persists at every dyadic precision. A unit derivative is a sufficient hypothesis for \cref{CR:thm:hensel}, not a necessary condition for every possible lift.

\section{Nonlinear limits and native realizations}\label{CR:cr-source-section-7}\label{CR:ch:limits}
\subsection{A linearized problem is a different equation}
Suppose a nonlinear map $F$ is differentiable at an approximation $x_0$. Its linearized correction equation is
\begin{equation}\label{CR:eq:linearization}
DF(x_0)u=-F(x_0).
\end{equation}
The cokernel criterion completely decides \eqref{CR:eq:linearization} once the derivative is fixed. It does not, by itself, decide whether there is $x$ with $F(x)=0$. Higher order terms are absent from the linearized equation, and they may either create a solution or prevent the linear prediction from reaching one.

The next two propositions are not warnings based only on examples computed numerically. The formal source includes actual derivative statements and exact proofs of the contrary conclusions.

\begin{CRproposition}[A nonlinear root with a nonzero first order obstruction]\label{CR:prop:nonlinear-root}
For $a\in\CRR\setminus\{0\}$, define $F_a(x)=x^2-a^2$. At zero,
\[
F_a'(0)=0,
\]
and the correction equation $0\cdot u=a^2$ has no solution. Nevertheless, $F_a(a)=0$.
\end{CRproposition}
\begin{proof}
Differentiation gives $F_a'(x)=2x$, so the derivative at zero is zero. Since $a\ne0$, its square is nonzero, while $0\cdot u=0$ for every $u$. Finally, $F_a(a)=a^2-a^2=0$.
\end{proof}
\CRformal{Nonlinear.nonlinear_root_despite_first_order_obstruction}

The root can lie arbitrarily close to zero by choosing a small nonzero $a$. This produces a family of nearby counterexamples. It does not assert that one fixed equation with a nonzero residual at zero has roots converging to zero.

\begin{CRproposition}[A solvable linearization without a real root]\label{CR:prop:no-real-root}
Let $G(x)=x^2+x+1$. Then
\[
G'(0)=1,\qquad \exists u\in\CRR\ (u=-1),
\]
so the linearized correction at zero is solvable. However, $G$ has no real root.
\end{CRproposition}
\begin{proof}
The derivative is $2x+1$, which equals one at zero, and $G(0)=1$. Thus $u=-1$ solves the linearized equation. But
\[
G(x)=\left(x+\frac12\right)^2+\frac34>0
\]
for every real $x$.
\end{proof}
\CRformal{Nonlinear.linearized_solution_without_real_root}

\subsection{Three conclusions that must stay distinct}
The results support three different kinds of certificate.

\textbf{An exact linear repair} provides witnesses for $Du+Ew=-r$ in the original target. Its validity is governed by quotient image membership.

\textbf{A finite arithmetic lift} provides a correction at one modulus and proves the exact divisibility gain. Its validity also uses arithmetic conditions such as B\'ezout coprimality and a scale divisibility relation.

\textbf{A nonlinear existence result} proves that a solution of the original nonlinear equation exists. It requires an additional theorem controlling the nonlinear terms, compatible successive approximations, or a different existence mechanism.

Paper III develops correction, nonlinear analysis, and filtered lifting under separate hypotheses~\cite{CR:paperiii}. The present formal component does not import those complete releases. It proves a restricted linear theory, a finite arithmetic step, and counterexamples that prevent unjustified upgrades between these conclusions.

\subsection{The relation to Realization Calculus}\label{CR:sec:realization}
The conceptual setting is the established separation between a mathematical object, its chosen realization, and the operations a realization preserves. Paper I's semantic realization criterion requires an explicit construction, preservation of the operations used in the proof, verification of additional domain hypotheses, and a conclusion stated in native objects~\cite{CR:paperi}.

Constructive repair fits that criterion in a direct way. The native equation supplies $V$ and its residuals. The allowed transformations supply $D$ and the permissible class of added maps. The quotient organizes which residual classes those maps can remove. A transport that is used to decide the problem must preserve the necessary obstruction data.

Komorebi Space is the proposed ambient setting for organizing realizations, compatibility data, residuals, and obstructions. This monograph supplies a bounded repair construction that can be placed in that setting. It does not prove that every possible realization in Komorebi Space has a canonical finite obstruction space or an optimal repair.

A successful application must still identify the formal data with the actual objects of its domain. For example, an arithmetic representation with prescribed character and phase information cannot be obtained merely by naming an abstract linear space an arithmetic carrier. The Paper XI research companion distinguishes finite comparison from arithmetic existence and develops specific enlarged carrier constructions under stated hypotheses~\cite{CR:paperxi}. Those constructions motivate the repair viewpoint, but are not reproduced or newly formalized here.

\subsection{Established reach and remaining questions}
The finite theory has a clear end point. It proves exactly when an added linear map repairs a family, the least possible added dimension over a field, optimal attainment inside a feasible legal subspace, complete finite menu feasibility, and a minimum cost search over finite domains. Its transport theorem excludes certificates that work only by collapsing the relevant quotient information.

The rational matrix bridge in Sections~\ref{CR:ch:rational} and~\ref{CR:ch:matrix-attainment} now connects an exact Python construction to a small formal verifier. Further questions remain separate. Equivariant, filtered, topological, and arithmetic repairs need admissibility classes more structured than a subspace. Nonlinear minimum repair needs conditions that do not confuse a derivative obstruction with global nonexistence. Infinite lifting needs compatible choices and a convergence theorem.

These are extensions of a precise contract rather than reasons to enlarge its present claims. The elementary linear algebra and arithmetic mechanisms are classical. The contribution of this component is their explicit repair interface, formal proofs, certified finite search, and preservation boundaries. No claim of a universally complete discovery procedure or historical priority follows from the formal theorem count.

\begin{scopebox}[title=A usable stopping point]
A certificate must state which equation it solves, which changes are allowed, and which information is preserved. Within the declared finite linear setting, the theorems provide both successful repairs and complete rejection of the supplied menu. Outside that setting, the missing domain hypotheses remain mathematical obligations.
\end{scopebox}


\chapter{Rational certificates and constructive attainment}\label{CR:book:chapter-57}
\section{Exact rational certificates}\label{CR:cr-source-section-8}\label{CR:ch:rational}
\subsection{Let the search propose and the kernel check}
The finite optimizer in Section~\ref{CR:ch:search} enumerates finite correction domains. Rational coefficients do not form such a domain. The appropriate interface is therefore not an enumeration of all rational assignments. A solver proposes finite rational matrices, and a verifier checks their identities in the original target.

Fix matrices
\[
 D\in\CRQ^{m\times n},\qquad R\in\CRQ^{m\times s},\qquad E\in\CRQ^{m\times q}.
\]
The residuals are the columns of $R$. For a selected menu $S$, let $\widetilde E_S$ retain the original candidate indexing and replace unselected columns by zero. The supplied coefficient matrices have shapes $U\in\CRQ^{n\times s}$ and $W\in\CRQ^{q\times s}$. Their correction equation is
\begin{equation}\label{CR:eq:rational-original}
 DU+\widetilde E_SW=-R.
\end{equation}
Keeping the original indexing avoids any ambiguity about a permutation of candidate columns. Coefficients on a zeroed column have no effect. The admissibility predicate and cost function belong to the original problem, not to the certificate producer.

\begin{CRproposition}[Successful certificate]\label{CR:prop:rational-success}
If an admissible $S$ and rational matrices $U,W$ satisfy \eqref{CR:eq:rational-original}, they repair every original residual.
\end{CRproposition}
\begin{proof}
Take the $i$th column of the displayed identity. It is the original equation $Du_i+\sum_{j\in S}w_{ji}e_j=-r_i$.
\end{proof}
\CRformal{RationalCertificates.checkRepair_sound}

\begin{scopebox}[title=A rejected certificate is not a rejected problem]
Failure to parse data, a false identity, a missing witness, or a resource limit is an unresolved verification outcome. It does not establish mathematical infeasibility. A rejection of the problem requires a valid obstruction certificate with an explicit scope.
\end{scopebox}

\subsection{Separating obstructions}
\begin{CRdefinition}[Separating row]
A separating certificate for $S$ is a row $Y\in\CRQ^{1\times m}$ satisfying
\begin{equation}\label{CR:eq:rational-separator}
 YD=0,\qquad Y\widetilde E_S=0,\qquad YR\ne0.
\end{equation}
The last condition asks that at least one original residual is detected. No residual index needs to be selected when the family is empty.
\end{CRdefinition}

\begin{CRtheorem}[Sound separation]\label{CR:thm:rational-separation}
A separating certificate rules out a repair by the selected menu.
\end{CRtheorem}
\begin{proof}
Multiplying a supposed correction equation on the left by $Y$ gives $0=-YR$, contradicting \eqref{CR:eq:rational-separator}.
\end{proof}
\CRformal{RationalCertificates.annihilator_kills_residual}
\CRformal{RationalCertificates.checkSeparator_sound}

\begin{CRproposition}[Linear alternative]\label{CR:prop:rational-alternative}
For a fixed rational menu, either the correction matrices exist or a separating row exists, but not both.
\end{CRproposition}
\begin{proof}
Exclusivity follows from the preceding theorem. If no correction matrices exist, some residual $r_i$ lies outside the column space of $[D\ \widetilde E_S]$. Choose a basis for that space, adjoin $r_i$, and extend to a basis of $\CRQ^m$. The functional taking value one on $r_i$ and zero on all other chosen basis vectors is a rational separating row.
\end{proof}
\CRsourcescope{This completeness argument is an ordinary proof. The new Lean module proves certificate soundness, not a general algorithm or formal theorem producing a separator for every infeasible rational menu.}

\subsection{Dimension certificates without a reported rank}
Suppose an augmentation with $t$ columns succeeds. Its optimality can be certified by two additional matrices
\[
 H\in\CRQ^{t\times m},\qquad Q\in\CRQ^{s\times t}
\]
with
\begin{equation}\label{CR:eq:rational-dimension}
 HD=0,\qquad HRQ=I_t.
\end{equation}
These identities make a lower bound on the required dimension explicit.

\begin{CRtheorem}[Dimension lower certificate]\label{CR:thm:rational-lower}
Assume \eqref{CR:eq:rational-dimension}. If an arbitrary competing augmentation $E'\in\CRQ^{m\times w}$ has coefficients satisfying $DU'+E'W'=-R$, then $t\leq w$.
\end{CRtheorem}
\begin{proof}
Multiplication by $H$ and $Q$ gives $HE'W'Q=-I_t$, hence
\[
 (HE')(-W'Q)=I_t.
\]
Thus $HE':\CRQ^w\to\CRQ^t$ has a right inverse. It is surjective, and consequently $t\leq w$.
\end{proof}
\CRformal{RationalCertificates.dimension_certificate_right_inverse}
\CRformal{RationalCertificates.dimension_certificate_lower_bound}

A successful augmentation with $t$ columns attains this bound. The verifier checks the original repair identity and \eqref{CR:eq:rational-dimension}. It does not trust a numerical rank supplied by the search program.
\CRformal{RationalCertificates.checkDimension_sound}

\begin{CRproposition}[Availability of lower certificates]
Let $O$ be the span of the residual classes in $\CRQ^m/\CRim D$, and let $t=\dim O$. Then matrices $H,Q$ satisfying \eqref{CR:eq:rational-dimension} exist.
\end{CRproposition}
\begin{proof}
Choose coordinates $O\cong\CRQ^t$, extend that map to the whole quotient, and compose with the quotient map to obtain $H$. Then $HD=0$. The columns of $HR$ span $\CRQ^t$, so $HR$ has a linear right inverse $Q$.
\end{proof}
\CRsourcescope{The availability argument is an ordinary finite-dimensional proof. The concrete producer in the next chapter finds and exports such data. Its search implementation is not itself proved correct in Lean.}

\subsection{Exact rank and menu cost}
For $A\in\CRQ^{a\times b}$, a certificate of rank $d$ consists of appropriately shaped rational matrices with
\begin{equation}\label{CR:eq:rational-rank}
 A=FG,\qquad PAQ=I_d,
\end{equation}
where the inner dimension of $FG$ is $d$.

\begin{CRproposition}[Rank certificate]
The identities \eqref{CR:eq:rational-rank} imply $\CRrank A=d$.
\end{CRproposition}
\begin{proof}
The factorization gives $\CRrank A\leq d$. The identity factorization through $A$ gives $d\leq\CRrank A$.
\end{proof}
\CRformal{RationalCertificates.checkRank_sound}

For comparison with the original quotient statement, ordinary rank-nullity gives
\begin{equation}\label{CR:eq:rank-difference-new}
 \dim O=\CRrank[D\ R]-\CRrank D.
\end{equation}
Indeed, the quotient map restricted to $\CRim D+\CRim R$ has kernel $\CRim D$ and image $O$. The current matrix verifier does not depend on a formal identification of this formula with the earlier abstract Lean quotient API. It obtains the lower bound directly from a right inverse.

Minimum menu cost is a separate optimization problem. To certify a successful menu $S$, the verifier requires a separating row for every admissible menu $T$ with strictly smaller declared cost. It quantifies over all subsets of the original finite candidate universe. It does not trust an incomplete producer-supplied list.

\begin{CRtheorem}[Certified finite-menu optimum]
An admissible successful menu, together with separating rows for every cheaper admissible menu, has minimum cost among all admissible feasible menus.
\end{CRtheorem}
\begin{proof}
A cheaper feasible menu would contradict its supplied separating row by Theorem~\ref{CR:thm:rational-separation}.
\end{proof}
\CRformal{RationalCertificates.checkOptimal_sound}

A certificate rejecting every admissible subset proves that no menu in the declared universe succeeds.
\CRformal{RationalCertificates.checkNoRepair_sound}

The rational example $D(u)=(u,0,0)$, $r=(0,1,0)$, $e_0=(0,1,1)$, $e_1=(0,0,1)$ is checked in the new matrix layer. Both entries repair the residual with coefficients $(-1,1)$. Rows $(0,1,0)$, $(0,1,-1)$, and $(0,1,0)$ separate the empty menu, the first singleton, and the second singleton. Thus the minimum cardinality is two. Restricting admissibility to at most one entry gives a checked complete rejection.
\CRformal{RationalMenus.menu_certificate}
\CRformal{RationalMenus.two_entries_minimal}
\CRformal{RationalMenus.restricted_certificate}
\CRformal{RationalMenus.no_small_menu}

The subset search may be exponential in menu size. The free augmentation exporter developed next does not implement a scalable menu optimizer, and the existence of a free minimum says nothing about which directions occur in a prescribed menu.

\section{Constructive matrix attainment}\label{CR:cr-source-section-9}\label{CR:ch:matrix-attainment}
\subsection{A complete detector}
Fix $D\in\CRQ^{m\times n}$ and $R\in\CRQ^{m\times s}$. A split correction certificate consists of
\[
 P\in\CRQ^{n\times m},\quad K\in\CRQ^{m\times d},\quad L\in\CRQ^{d\times m}
\]
with
\begin{equation}\label{CR:eq:split-certified}
 LD=0,\qquad DP+KL=I_m.
\end{equation}
The first identity makes $L$ a detector that ignores original corrections. The second proves that no further information has been discarded: every target vector has the decomposition $v=DPv+KLv$. There is no assumption that $LK=I_d$, and redundant detector coordinates are permitted.

\begin{CRtheorem}[Exact detection]\label{CR:thm:split-exact}
Under \eqref{CR:eq:split-certified},
\[
 (\exists X,\ DX=-R)\quad\Longleftrightarrow\quad LR=0.
\]
When $LR=0$, the explicit witness is $X=-PR$.
\end{CRtheorem}
\begin{proof}
A correction gives $0=LDX=-LR$. Conversely, $LR=0$ gives $R=(DP+KL)R=DPR$, so $D(-PR)=-R$.
\end{proof}
\CRformal{RationalCertificates.split_correctable_iff}

At the quotient level, $L$ induces an injection $\CRQ^m/\CRim D\to\CRQ^d$. If $Lv=0$, the split identity gives $v=DPv$, proving that its class is zero. The image is $\CRim L$. Hence the required residual obstruction span is isomorphic to the column space of $C=LR$. This quotient interpretation is an ordinary consequence of the displayed identities.

\subsection{The attaining construction}
Suppose the detected residual matrix $C=LR$ has certificate data
\[
 F\in\CRQ^{d\times t},\quad G\in\CRQ^{t\times s},\quad
 J\in\CRQ^{t\times d},\quad Q\in\CRQ^{s\times t}
\]
satisfying
\begin{equation}\label{CR:eq:factor-certified}
 LR=FG,\qquad JLRQ=I_t.
\end{equation}
The new augmentation and its correction coefficients are not unspecified witnesses. They are
\begin{equation}\label{CR:eq:attaining-formula}
 \boxed{E=KF,\qquad U=-PR,\qquad W=-G.}
\end{equation}

\begin{CRtheorem}[Constructive minimum augmentation]\label{CR:thm:matrix-attainment}
Assume \eqref{CR:eq:split-certified} and \eqref{CR:eq:factor-certified}. The matrices in \eqref{CR:eq:attaining-formula} satisfy $DU+EW=-R$. Every competing rational augmentation repairing the same residual family has at least $t$ added dimensions.
\end{CRtheorem}
\begin{proof}
The constructed matrices satisfy
\[
 DU+EW=-DPR-KFG=-DPR-KLR=-(DP+KL)R=-R.
\]
For the lower bound set $H=JL$. Then $HD=0$ and $HRQ=I_t$. Apply Theorem~\ref{CR:thm:rational-lower} to every competing repair.
\end{proof}
\CRformal{RationalCertificates.split_constructs_repair}
\CRformal{RationalCertificates.checkOptimalConstruction_sound}

This proves attainment and minimality after checking four finite identities. The generalized inverse and the choice of a residual basis belong to the producer. They are not assumed trustworthy. A valid set of identities is sufficient, regardless of how the producer found it.

\subsection{How the Python producer obtains the matrices}
The executable producer uses exact rational row reduction. It selects independent columns of $D$ and extends them by standard basis vectors to a basis $B$ of $\CRQ^m$. The inverse $B^{-1}$ supplies the image coordinates and the complementary detector coordinates. The selected image coordinates are placed in the original domain to obtain $P$, while the complementary basis vectors form $K$ and the corresponding rows of $B^{-1}$ form $L$.

For $C=LR$, the producer selects independent residual columns to form $F$. A selection matrix $Q$ records those columns. A left inverse $J$ of $F$ gives $G=JC$. The producer then forms $E,U,W$ using \eqref{CR:eq:attaining-formula}, and checks every identity again with exact rational arithmetic before export.

All of these data can be obtained over a field. The search implementation may nevertheless contain a bug, encounter a resource limit, or reject malformed input. Such events do not become mathematical no-go results. The Lean certificates verify successful outputs, not the total correctness or complexity of this search program.

Zero-dimensional targets, original domains, obstruction spaces, and residual families are retained. For an already correctable residual family the construction has $t=0$, so it adds no directions. The input limit of the supplied prototype is an implementation resource bound, not a restriction on Theorem~\ref{CR:thm:matrix-attainment}.

\subsection{A full intermediate trace}\label{CR:sec:matrix-worked}
Take
\[
 D=\begin{pmatrix}1\\0\\0\end{pmatrix},\qquad
 R=\begin{pmatrix}2&4\\1&2\\1&2\end{pmatrix}.
\]
The producer selects the first column of $D$ and the complementary second and third coordinate axes. It obtains
\[
 P=\begin{pmatrix}1&0&0\end{pmatrix},\quad
 K=\begin{pmatrix}0&0\\1&0\\0&1\end{pmatrix},\quad
 L=\begin{pmatrix}0&1&0\\0&0&1\end{pmatrix}.
\]
Thus $LD=0$, $DP=\operatorname{diag}(1,0,0)$, and $KL=\operatorname{diag}(0,1,1)$. The residual coordinates are
\[
 C=LR=\begin{pmatrix}1&2\\1&2\end{pmatrix}.
\]
A rank-one certificate is
\[
 F=\begin{pmatrix}1\\1\end{pmatrix},\quad
 G=\begin{pmatrix}1&2\end{pmatrix},\quad
 J=\begin{pmatrix}1&0\end{pmatrix},\quad
 Q=\begin{pmatrix}1\\0\end{pmatrix}.
\]
The result is
\[
 E=\begin{pmatrix}0\\1\\1\end{pmatrix},\qquad
 U=\begin{pmatrix}-2&-4\end{pmatrix},\qquad
 W=\begin{pmatrix}-1&-2\end{pmatrix}.
\]
The original equations become
\[
 DU+EW=
 \begin{pmatrix}-2&-4\\0&0\\0&0\end{pmatrix}+
 \begin{pmatrix}0&0\\-1&-2\\-1&-2\end{pmatrix}=-R.
\]
The lower witness is $H=JL=(0,1,0)$, with $HD=0$ and $HRQ=1$. The ambient cokernel has dimension two, but the residual classes span only one direction. The construction therefore adds one direction and proves that zero added directions are insufficient.
\CRformal{Case000.certificate_checked}
\CRformal{Case000.original_equations_checked}
\CRformal{Case000.optimal}

The fractional case uses $D=(2/3,4/3)^{\mathsf T}$ and $R=(1/2,5/6)^{\mathsf T}$. It exercises nonintegral splitting and correction coefficients while retaining exact arithmetic. Empty and rank-deficient cases are included in the same exported corpus. Every case has its own original input file, certificate file, and formal optimality statement.

\subsection{Serialization and input identity}
The wire format represents a rational entry by a signed decimal numerator and a positive decimal denominator. The Python decoder rejects floats, nonfinite values, duplicate object keys, zero or negative denominators, unreduced pairs, malformed dimensions, and extraneous certificate fields. Resource limits produce an invalid or unresolved outcome, not an infeasibility theorem.

The generated Lean data use typed pairs in $\CRZ\times\CRN$. Their decoder returns the intended rational when the denominator is nonzero and explicitly rejects a zero denominator at the cell level. Each generated input and certificate additionally proves denominator validity for all its cells.
\CRformal{RationalCertificates.decodeCell_zero}
\CRformal{RationalCertificates.decodeCell_nonzero}
\CRformal{RationalCertificates.decodeCell_sound}
\CRformal{RationalCertificates.checkWireMatrix_sound}
\CRformal{RationalCertificates.decodeWireMatrix_cell}

The original matrices $D,R$ are emitted separately from certificate data. The certificate module imports them and cannot replace them through a certificate field. A canonical problem hash binds the JSON records. A diagnostic Lean echo of the decoded input matrices is compared with the original JSON values for every exported case. That echo is an additional execution check, never a theorem premise.

\begin{scopebox}[title=The precise trust boundary]
Lean checks the matrix identities, the original equations, and the minimum-dimension conclusion for the literal decoded problem in the generated input module. The numeric cell-decoding lemmas are formalized. The whole JSON parser, file system, hashing routine, and source emitter are not proved correct in Lean. Their input binding is checked by hashes, strict decoding, and a separate Lean echo. A user must still identify the intended original problem correctly.

No theorem is accepted through \CRcode{native_decide}. Concrete certificate propositions are proved with \CRcode{decide +kernel}, and every exported theorem is included in the transitive axiom audit.
\end{scopebox}

\subsection{Legality and coefficient rings}
Free minimum dimension is not minimum menu cost. The constructed direction $KF$ need not occur among supplied menu entries. When legal directions form $\CRim A$, an additional identity $KF=AZ$ certifies that the augmentation lies in that legal subspace. This identity also rewrites every added correction as $A(ZW)$.
\CRformal{RationalCertificates.constructed_repair_in_legal_image}

The availability of splitting also depends on the coefficient ring. For $D=(2)$ over $\CRZ$, the condition $LD=0$ forces $L=0$. The split identity would then require $2P=1$, which is impossible over $\CRZ$. Over $\CRQ$, $P=1/2$ is permitted. The current producer and checker are rational, not a certificate of integral realizability.

Finally, exact rational correction is not nonlinear convergence or arithmetic representation-theoretic existence. The existing boundaries of Paper III and Conjecture 66 are unchanged. The contribution of this extension is a usable and checked witness interface inside the established calculus, not a replacement for its foundations.


\chapter{Costed covers and proof-carrying branch search}\label{CR:book:chapter-58}
\section{Costed rational menus and rejection covers}\label{CR:cr-source-section-10}
\label{CR:ch:menu-covers}

The free augmentation problem asks how many new dimensions are necessary. A prescribed menu asks a different question: which available directions should be purchased, enabled, or admitted? Two menu entries can be needed for a one-dimensional obstruction, and the cheapest solution need not use the fewest entries. This chapter makes that second problem operational over the rationals.

The implementation uses exact elimination to propose an answer. A separate certificate checker verifies the answer against the original matrices and the complete admissibility rule. The search procedure is outside the proof trust boundary. The returned mathematical witnesses are not.

\subsection{The original menu contract}

Fix rational matrices
\[
 D\in\CRQ^{m\times n},\qquad E\in\CRQ^{m\times q},\qquad R\in\CRQ^{m\times s}.
\]
The columns $e_j$ of $E$ are the available new directions. The finite family
\[
 \mathcal A\subseteq\mathcal P(\{0,\ldots,q-1\})
\]
contains exactly the admissible selections. It need not be downward closed. Nonnegative integer weights $c_j$ define the cost
\[
 c(S)=\sum_{j\in S}c_j.
\]
These matrices, selections, and costs are fixed original input. A certificate cannot replace them with another problem.

Let $\widetilde E_S$ retain each column of $E$ selected by $S$ and replace every other column by zero. This convention preserves the original candidate indexing. A coefficient table $W\in\CRQ^{q\times s}$ may have arbitrary values in unselected rows, but those rows contribute nothing. A proposed repair consists of matrices $U,W$ with
\begin{equation}
 DU+\widetilde E_SW=-R.
 \label{CR:eq:menu-cover-primal}
\end{equation}
Every selected direction is charged, even if its coefficient row happens to be zero. Thus cost measures the selected menu, not merely the support of a particular coefficient table.

\subsection{Separating blocks}

An infeasible menu admits a rational separating row, by the exact linear alternative in \cref{CR:ch:rational}. The row annihilates all available corrections while detecting at least one residual. Instead of providing a separate row for every rejected subset, we can reuse one row on all subsets of a larger rejected set.

\begin{CRdefinition}[Separating block]
A separating block is a pair $(B,Y)$, where $B\subseteq\{0,\ldots,q-1\}$ and $Y\in\CRQ^{1\times m}$ satisfy
\begin{equation}
 YD=0,\qquad Y\widetilde E_B=0,\qquad YR\ne0.
 \label{CR:eq:menu-cover-dual}
\end{equation}
The set $B$ need not itself belong to $\mathcal A$.
\end{CRdefinition}

The last condition is a matrix inequality: at least one column is detected. For an empty residual family it cannot hold, which correctly prevents a false separating witness for a vacuously solved family.

\begin{CRlemma}[Column-mask identity]
For every row $Y$ and column $j$,
\[
 (Y\widetilde E_S)_{0j}=
 \begin{cases}
 (YE)_{0j},&j\in S,\\
 0,&j\notin S.
 \end{cases}
\]
\end{CRlemma}
\begin{proof}
The mask either retains the entire $j$th column or replaces it by zero. Expanding the finite matrix product gives the two cases.
\end{proof}
\CRformal{KomorebiRepair.RationalCertificates.mul_masked_entry}

\begin{CRtheorem}[Downward validity of a separator]
If $(B,Y)$ is a separating block and $T\subseteq B$, then $Y$ also separates $T$. In particular, $T$ cannot satisfy \eqref{CR:eq:menu-cover-primal}.
\end{CRtheorem}
\begin{proof}
For a selected column $j\in T$, one also has $j\in B$, so $(YE)_{0j}=0$ by the block identity. Every unselected column is zero after masking. Thus $Y\widetilde E_T=0$. The other two equations in \eqref{CR:eq:menu-cover-dual} are unchanged.

If a repair for $T$ existed, multiplication by $Y$ would give
\[
 0=YDU+Y\widetilde E_TW=-YR,
\]
contradicting the detection condition.
\end{proof}
\CRformal{KomorebiRepair.RationalCertificates.separates_subset}

The restriction to an admissible family does not invalidate this argument. A block can be a forbidden superset while still supplying valid mathematical evidence against its permitted subsets. What matters is that the block's own linear impossibility has been checked.

\subsection{Complete coverage of cheaper alternatives}

Let $\mathcal B$ be a finite list of separating blocks. Say that $\mathcal B$ covers $T$ when some $(B,Y)\in\mathcal B$ satisfies $T\subseteq B$. Validity of the block list and coverage of the relevant selections are separate obligations.

\begin{CRtheorem}[Certified minimum menu cost]
Suppose the following data have been checked:
\begin{enumerate}
 \item $S\in\mathcal A$ and the original repair equation holds for $S,U,W$.
 \item Every member of $\mathcal B$ is a valid separating block.
 \item Every $T\in\mathcal A$ with $c(T)<c(S)$ is covered by $\mathcal B$.
\end{enumerate}
Then $S$ has minimum cost among all admissible feasible selections.
\end{CRtheorem}
\begin{proof}
The first condition supplies an admissible feasible selection. If a cheaper admissible feasible $T$ existed, the third condition would place it inside a separating block. Downward validity would make $T$ infeasible, a contradiction.
\end{proof}
\CRformal{KomorebiRepair.RationalCertificates.checkCoveredOptimal_sound}

The checker quantifies over the whole finite candidate-subset universe when it verifies coverage. It does not accept a producer's claim that its own list of alternatives is complete. Equal-cost alternatives need not be rejected, and the certificate does not claim that the optimum is unique.

\begin{CRtheorem}[Certified rejection of the admissible family]
If every block is valid and every member of $\mathcal A$ is covered, then no admissible menu repairs the residual family.
\end{CRtheorem}
\begin{proof}
Every admitted selection lies inside a valid separating block and is therefore infeasible.
\end{proof}
\CRformal{KomorebiRepair.RationalCertificates.checkCoveredNoRepair_sound}

When $\mathcal A$ is empty, this conclusion is vacuous. The implementation reports that fact explicitly. It does not manufacture a nonzero obstruction or claim that the original equations have no solution without the admissibility restriction.

\begin{scopebox}[title=The scope of rejection]
An accepted rejection certificate rules out the declared admissible menus. It does not rule out a new direction outside the menu, a different coefficient ring, a changed admissibility rule, or a different arithmetic host. A failed check or exhausted time budget establishes none of these impossibility claims.
\end{scopebox}

\subsection{Constructing the evidence}

For each candidate subset, the producer performs exact rational row reduction on the selected correction columns. It keeps the row transformation matrix as part of the computation. If the transformed residual has a nonzero entry in a zero coefficient row, that row transformation supplies a separating functional. Otherwise the pivot equations give a coefficient table for every residual. Each proposed identity is rechecked using exact arithmetic.

The producer evaluates all $2^q$ selections in this bounded implementation. It chooses a feasible admitted selection of minimum declared cost, with a deterministic tie break used only for reproducibility. Alternatively, it records that none of the admitted selections passed the exact linear alternative.

For rejection evidence, the producer finds inclusion-maximal infeasible subsets of the full candidate universe. Every infeasible selection is contained in such a maximal subset because the universe is finite. A separating row for each selected maximal subset therefore provides a complete supply of potential rejection blocks. The producer then greedily retains blocks that cover the required alternatives.

The greedy compression is not claimed to minimize certificate size. Nor does a smaller rejection cover imply a faster search algorithm: the present producer still examines every subset. The verified conclusion concerns menu cost and the original equations, not optimal compression or polynomial-time complexity.

\subsection{Two entries for one missing dimension}

Take
\[
 D=\begin{pmatrix}1\\0\\0\end{pmatrix},\qquad
 E=\begin{pmatrix}0&0\\1&0\\1&1\end{pmatrix},\qquad
 R=\begin{pmatrix}0\\1\\0\end{pmatrix},
\]
with both candidate weights equal to one and every subset admitted. The chosen menu is $S=\{0,1\}$, with
\[
 U=(0),\qquad W=\begin{pmatrix}-1\\1\end{pmatrix}.
\]
Its exact repair is
\[
 DU+EW=\begin{pmatrix}0\\-1\\0\end{pmatrix}=-R.
\]
Two separating blocks suffice:
\[
 B_0=\{0\},\quad Y_0=(0,1,-1),\qquad
 B_1=\{1\},\quad Y_1=(0,1,0).
\]
Both satisfy $Y_iD=0$, annihilate their masked candidate matrices, and give $Y_iR=1$.

\begin{center}
\begin{tabular}{@{}ccc@{}}\toprule
Cheaper menu & Cost & Covering block\\\midrule
$\varnothing$ & 0 & $B_0$ or $B_1$\\
$\{0\}$ & 1 & $B_0$\\
$\{1\}$ & 1 & $B_1$\\\bottomrule
\end{tabular}
\end{center}
Thus cost two is minimal. Restricting admissibility to the three cheaper menus turns the same two blocks into a complete no-admissible-repair certificate. The unrestricted one-dimensional repair direction $R$ still exists, but it is not one of the two listed candidate directions.

\CRsourcescope{Executed examples: Menu000 and Menu001. Each has separately emitted original input, exact rational witnesses, a Lean certificate proof, an axiom audit, and an input echo.}

\subsection{Cost, dimension, and evidence size}

A third candidate can make the cost distinction sharper. Add the direct direction $e_0=(0,1,0)$ at cost five, and price $e_1=(0,1,1)$ and $e_2=(0,0,1)$ at one each. One menu entry can now repair the residual, but the cheapest repair selects two entries, $\{1,2\}$, at total cost two. Minimum free dimension, minimum menu cardinality, and minimum menu cost are three different quantities.

A separate example shows why rejection covers help. Let
\[
 D=\begin{pmatrix}1\\0\end{pmatrix},\qquad
 E=\begin{pmatrix}2&3&4&5\\0&0&0&0\end{pmatrix},\qquad
 R=\begin{pmatrix}0\\1\end{pmatrix}.
\]
The row $Y=(0,1)$ separates the full four-entry menu. One block therefore excludes all sixteen subsets. Adding $(0,1)$ as a new candidate would change the problem and allow a repair. The certified rejection of the original menu would remain correct.

\CRsourcescope{Executed examples: Menu002 distinguishes cost from cardinality. Menu004 excludes all sixteen subsets with one separating block. Menu005 uses one block to cover eight cheaper selections before accepting a positive-cost direction.}

\subsection{Fractional witnesses and independent original data}

The solver is not restricted to coefficients in a finite field. For
\[
 D=\begin{pmatrix}1\\0\end{pmatrix},\quad
 E=\begin{pmatrix}0\\2/3\end{pmatrix},\quad
 R=\begin{pmatrix}1/2\\5/6\end{pmatrix},
\]
the selected direction at cost seven has witnesses
\[
 U=-1/2,\qquad W=-5/4.
\]
The checker verifies the exact rational equation, with no tolerance or floating-point rank test.

The generated input module contains $D,E,R$, the candidate weights, and the literal admissibility table. The certificate module imports those inputs and supplies only coefficient matrices and separating blocks. Typed rational decoding and nonzero denominator proofs are retained from the preceding bridge.

An independent diagnostic echo evaluates the matrices in Lean and evaluates the admissibility and cost of every candidate subset. Python compares those evaluated results with the original JSON input. This comparison catches mismatched original constraints as well as mismatched matrices. It is an executed statement-binding check, not a formal proof of the entire external JSON parser or source emitter.

\subsection{Bounded executable use}

The command \CRcode{scripts/certify_menu.py} creates a new local directory, constructs or reads a certificate, emits the fixed input and witness modules, compiles their proofs, audits their axioms, replays the modules, and checks the original-data echo. Existing output directories are refused.

\begin{Verbatim}[fontsize=\small,breaklines=true]
export KOMOREBI_LEAN_KIT=/path/to/Komorebi_Portable_Lean_4.34.0-rc2
python scripts/certify_menu.py examples/menu_corpus/Menu003_input.json --output /tmp/repair-menu
\end{Verbatim}

The successful mathematical outcomes are \nolinkurl{KERNEL_CHECKED_OPTIMAL_MENU} and \nolinkurl{KERNEL_CHECKED_NO_ADMISSIBLE_MENU}. Invalid data, unsupported sizes, compiler failure, and timeouts produce \CRcode{UNRESOLVED}. A Python-valid certificate alone is not a kernel-checked result.

This version limits the menu to eight candidates. Other input bounds include thirty-two rows or columns, canonical rational cells with bounded decimal length, and a four-megabyte JSON file limit. These are implementation safeguards, not hypotheses of the finite mathematical theorems. Exhaustive search is still exponential in menu size, and the time limit can be reached within the admitted size bounds.

The module is a constructive extension of the correction framework already developed in Realization Calculus. It gives Komorebi Space a more explicit candidate interface for admissible finite linear repairs, but supplies no universal semantic-realization theorem. Integration into the authoritative Paper III tree and the established Komorebi Python library remains separate from this isolated, reproducible implementation.

\section{Proof carrying branch search}\label{CR:cr-source-section-11}
\label{CR:ch:branch-search}

The preceding menu verifier checks every required selection against a family of separating blocks. That is a sound endpoint, but its coverage test still ranges over the full subset universe. A large menu may therefore be expensive to verify even when one short argument excludes almost every alternative.

This chapter replaces the powerset coverage test with a finite search tree. Each node describes an interval of selections. A terminal node either proves a cost lower bound, supplies one separator for the entire interval, or checks that the original admissibility rule excludes that interval. A branching node must retain both possible decisions about one unresolved candidate.

The producer may search in any order and use any heuristic. The verifier checks the submitted tree. It does not trust the producer's account of which alternatives were examined.

\begin{scopebox}[title=The new verified boundary]
The mathematical result is soundness of a finite rational certificate tree. It certifies minimum cost within the declared menu or infeasibility of the declared admissible family. It is not a proof of correctness or termination of the entire Python producer, a polynomial complexity result, or a native arithmetic realization theorem.

All 37 earlier proof modules are retained unchanged. The new module reuses their exact rational repair and separator theorems.
\end{scopebox}

\subsection{Intervals of candidate selections}

Let $J=\{0,\ldots,q-1\}$ be the original candidate index set. The matrices $D,E,R$ and the original correction equation remain fixed:
\[
 DU+\widetilde E_S W=-R.
\]
Here $\widetilde E_S$ zeros unselected columns of $E$ while retaining its original indexing. Costs are nonnegative integers attached to those same columns:
\[
 c(S)=\sum_{j\in S}w_j,\qquad w_j\in\CRN.
\]
The admissibility rule either permits every subset or supplies a fixed finite table of permitted subsets. An empty table permits none. It is different from the all-subsets rule.

\begin{CRdefinition}[Selection interval]
For $A\subseteq B\subseteq J$, define
\[
 \mathcal I(A,B)=\{T\subseteq J:A\subseteq T\subseteq B\}.
\]
The set $A$ contains candidates already required, and $B$ contains all candidates still permitted by the branch decisions. The unresolved indices form $B\setminus A$.
\end{CRdefinition}

\begin{CRlemma}[Exact branch partition]
If $j\in B\setminus A$, then
\[
 \mathcal I(A,B)=
 \mathcal I(A,B\setminus\{j\})\;\sqcup\;
 \mathcal I(A\cup\{j\},B).
\]
\end{CRlemma}
\begin{proof}
Every $T\in\mathcal I(A,B)$ either excludes $j$ or includes $j$. In the first case it belongs to the left interval, and in the second it belongs to the right interval. Conversely, each child interval is contained in the parent. Their membership conditions on $j$ are incompatible, so the union is disjoint.
\end{proof}

This identity is the completeness mechanism of the verifier. A certificate cannot simply omit an inconvenient branch. Its split constructor contains both children, and the checker verifies both. The proof of \CRdecl{checkBranch_sound} applies the two cases directly.

\subsection{Three terminal certificates}

\subsubsection{Cost lower bound}

Nonnegative additive costs are monotone under inclusion:
\[
 A\subseteq T\quad\Longrightarrow\quad c(A)\leq c(T).
\]
Therefore, when an incumbent repair has cost $c_\ast$, the finite check
\[
 c_\ast\leq c(A)
\]
excludes every strictly cheaper repair in $\mathcal I(A,B)$. It does not say the interval is infeasible. It only says that it cannot improve the incumbent.
\CRformal{KomorebiRepair.RationalCertificates.additiveCost_mono}

\begin{proof}
The sum for $c(T)$ contains all terms in the sum for $c(A)$ and only additional nonnegative terms. Apply transitivity to $c_\ast\leq c(A)\leq c(T)$.
\end{proof}

A cost leaf is forbidden in an impossibility certificate. Even a very expensive feasible repair refutes a claim that no repair exists. The implementation expresses this distinction using an optional threshold: a supplied natural number requests a lower bound, while no threshold requests infeasibility.

\subsubsection{Separator for the upper selection}

A row $Y$ is checked against the largest selection $B$ in the interval:
\[
 YD=0,\qquad Y\widetilde E_B=0,\qquad YR\ne0.
\]
Every $T\subseteq B$ uses only columns annihilated by $Y$. Consequently, the same row proves that no $T\in\mathcal I(A,B)$ repairs the family. This reuses the downward separator theorem from the previous chapter.

The residual matrix in the last inequality is the original $R$. Projecting it to zero in some unrelated representation is not accepted as a repair.

\subsubsection{Exclusion by the original admissibility rule}

If admissibility is specified by a list $\mathcal A$, a terminal check establishes
\[
 \nexists T\in\mathcal A:\quad A\subseteq T\subseteq B.
\]
It examines the original list, not the powerset of $J$. For the all-subsets rule this terminal certificate is unavailable.
\CRformal{KomorebiRepair.RationalCertificates.intervalExcluded_sound}

It is not enough that $A$ itself be inadmissible. A larger selection in the interval might be permitted. For example, a rule that permits only the pair $\{0,1\}$ excludes the empty selection but not the interval from the empty selection to the full pair. The verifier retains this distinction and does not assume that admissibility is downward closed.

\subsection{Soundness of the whole tree}

\begin{CRdefinition}[Accepted branch tree]
A tree for $\mathcal I(A,B)$ consists of one of the three terminal certificates above, or a split at $j\in B\setminus A$ with accepted trees for both child intervals. In lower-bound mode a cost threshold is fixed throughout the tree. In infeasibility mode cost leaves are rejected.
\end{CRdefinition}

\begin{CRtheorem}[Structural certificate soundness]
An accepted lower-bound tree with threshold $c_\ast$ proves
\[
 T\in\mathcal I(A,B),\quad T\text{ admissible and feasible}
 \quad\Longrightarrow\quad c_\ast\leq c(T).
\]
An accepted infeasibility tree proves that no admissible selection in $\mathcal I(A,B)$ is feasible.
\end{CRtheorem}
\CRformal{KomorebiRepair.RationalCertificates.checkBranch_sound}

\begin{proof}
Induct on the supplied tree. For a cost leaf, monotonicity gives the lower bound. For a separator leaf, the exact linear alternative rules out feasibility of every subset of the upper selection. For an admissibility leaf, the original table rules out the premise.

At a split, the exact branch partition puts any proposed selection in one child interval. Apply that child's induction hypothesis. Both children are present in the data and checked by the verifier. This proves the assertion without enumerating every selection in the parent interval.
\end{proof}

\begin{CRcorollary}[Minimum cost certificate]
Suppose an admissible selection $S$ has supplied coefficients $U,W$ satisfying the original repair equation. An accepted lower-bound tree of threshold $c(S)$ on $\mathcal I(\varnothing,J)$ proves that $S$ has minimum cost among all admissible feasible selections.
\end{CRcorollary}
\CRformal{KomorebiRepair.RationalCertificates.checkTreeOptimal_sound}
\begin{proof}
The supplied coefficient equation proves feasibility of $S$. Every selection lies in $\mathcal I(\varnothing,J)$. Structural soundness therefore gives $c(S)\leq c(T)$ for every admissible feasible competitor $T$.
\end{proof}

\begin{CRcorollary}[Scoped complete rejection]
An accepted infeasibility tree on $\mathcal I(\varnothing,J)$ proves that the declared admissible family contains no repair.
\end{CRcorollary}
\CRformal{KomorebiRepair.RationalCertificates.checkTreeNoRepair_sound}

There is no claim that directions outside $J$ cannot repair the residual. Nor does a rejected certificate establish infeasibility. A missing child, a bad separator, or an exhausted resource budget leaves the outcome unresolved.

\subsection{How the producer constructs evidence}

The Python producer traverses intervals and retains exact row-reduction witnesses. It first checks whether the original admissibility table excludes the interval. It then uses the incumbent cost to test a bound leaf. If necessary, it solves the rational linear problem on the upper selection $B$. An infeasible upper selection supplies a separator that ends the whole branch.

If the lower selection $A$ is admitted and feasible, its coefficient matrix supplies a new incumbent and the interval ends at a cost leaf. Otherwise, the producer selects an unresolved index and explores both decisions. The current heuristic visits low-cost candidates first. Correctness of this heuristic is not an assumption of any theorem.

The incumbent may improve while the search proceeds. Earlier cost leaves remain sound: a later incumbent has no greater cost, so any earlier lower bound also excludes improvement over the final incumbent. The finished checker uses only the final threshold and rechecks every leaf against it. It need not trust the producer's history of incumbents.

Exact linear solves are cached by selection. This avoids repeating an identical solve when several intervals have the same upper or lower selection. Caching is an implementation choice outside the trusted mathematical layer.

\subsection{A twenty-four-candidate construction}

Let $b_1,b_2,b_3$ be the standard basis of $\CRQ^3$. Take $D(u)=ub_1$ and residual $r=b_2+b_3$. There are 24 candidate entries:
\[
 a_0=b_2,\qquad a_1=b_3,\qquad a_j=b_1\quad(2\leq j<24).
\]
Their costs are two, three, and ten for each remaining entry. All subsets are admissible.

The selected repair is $S=\{0,1\}$, with cost five, original coefficient zero, and coefficient $-1$ on each selected menu direction. Its full coefficient matrix satisfies the original equation exactly.

The proof tree has five nodes. At the root it splits on candidate zero. The excluded child lacks the only second-coordinate correction and has separator $(0,1,0)$. The included child splits on candidate one. Excluding candidate one has separator $(0,0,1)$. Including both reaches lower selection $\{0,1\}$ with cost five and terminates at a bound leaf.

\begin{center}
\begin{tabularx}{\linewidth}{@{}lX@{}}
\toprule Branch & Certified conclusion\\\midrule
Exclude candidate 0 & The second coordinate of the residual cannot be corrected.\\
Include 0, exclude 1 & The third coordinate of the residual cannot be corrected.\\
Include 0 and 1 & Every superset costs at least five.\\
\bottomrule
\end{tabularx}
\end{center}

There are $2^{24}=16{,}777{,}216$ possible selections, but the producer used six distinct linear solves and submitted five tree nodes. Lean checks that tree and proves the minimum-cost conclusion for every selection. This is not an experiment that visited sixteen million cases. It is one structured problem with a short exhaustive mathematical argument.

A second 24-candidate example puts every available correction on the first coordinate axis of $\CRQ^2$ and sets the residual to $(0,1)$. The row $(0,1)$ supplies a one-node certificate rejecting all selections. This illustrates compressed reasoning, not a claim that arbitrary difficult search problems have short certificates.

\subsection{Performance and semantic limits}

The earlier eight-candidate exhaustive producer remains available as a reference. The new producer accepts up to 24 candidates, with explicit limits on nodes and linear solves. These are implementation resource guards, not mathematical restrictions on the tree theorem.

In a favorable instance, separators or cost bounds eliminate large intervals. In the worst case the tree can still approach the full binary decision tree. Neither polynomial-time search nor a polynomial-size certificate is established. With an explicit admissibility table, checking an exclusion leaf also costs work proportional to that table.

The nonnegative cost requirement is essential to the implemented bound rule. If an optional candidate had negative cost, adding it could lower the cost below the lower selection's cost. The input schema rejects negative weights instead of silently applying an invalid monotonicity argument.

The theorem operates on the original rational matrices and declared subset policy. It does not prove that these finite data capture a native arithmetic, analytic, or representation-theoretic construction. That remains the semantic realization obligation of the domain application. In Komorebi Space, the admissible family must retain its declared meaning. A compressed certificate cannot replace that meaning with a more convenient problem.


\part{Self-improvement and the discovery of methods}
\setcounter{section}{0}\renewcommand{\thesection}{SII.\arabic{section}}\renewcommand{\theHsection}{SII.\arabic{section}}

\chapter{The object that improves}\label{SII:book:chapter-59}

\section{From a better answer to a better investigator}\label{SII:sec:introduction}
A system encounters a failure, changes its representation, and succeeds on the next attempt. Something useful has happened. The system may have learned a fact, recovered a missing variable, or selected a more suitable operation. Yet the next unfamiliar failure may leave it no better equipped than before. It has a repaired answer without a reusable account of repair.

A stronger change would preserve the way the failure was investigated. The system might discover that its representation merged two situations with different future consequences. It might learn to seek an observation that separates those situations, to reject an inadequate family of repairs, or to reuse a diagnostic method in another notation. The object retained is then not merely a solution. It is part of a procedure for acquiring solutions.

The recursive question begins one step further out. Can the system improve that procedure, use the revised procedure to construct another improvement, and retain a defensible comparison across those changes? The question is not answered by putting an agent's source code in its own context. Editability is an affordance. Productive revision is a property of what the permitted edits actually achieve.

We take general intelligence as a motivating problem rather than a theorem premise. A broad reasoner should acquire useful capabilities across substantially different tasks without receiving their solutions inside an uncharged interpreter. Chollet's account of skill-acquisition efficiency emphasizes scope, priors, and experience rather than demonstrated skill alone~\cite{SII:chollet}. Our narrower contribution concerns one possible mechanism behind such acquisition: reusable procedures for diagnosing failures, choosing informative actions, constructing corrections, and preserving their scope.

\subsection{Three distinct achievements}
\emph{Task improvement} changes how well a solver performs a declared task. \emph{Process improvement} changes how effectively an improvement procedure produces useful solvers from comparable starting conditions. \emph{Epistemic improvement} changes how reliably the evaluation process distinguishes genuine gains from misleading evidence. This separation is inherited from SRIC~\cite{SII:sric}. A system can advance on one axis without advancing on the others.

A larger archive, for example, can improve retrieval while the procedure that extends the archive remains unchanged. A stronger checker can reject more invalid proposals while producing no additional valid proposals. A faster search may exploit a defect in its evaluation. None of these possibilities makes improvement meaningless. They identify the claim that must be tested.

Recursive improvement also differs from sustained acceleration. A sequence of positive gains can converge to a finite limit. A sequence of better algorithms can consume increasing development costs. A procedure can improve on one task family and deteriorate on another. The substantive claim is not that a feedback loop exists, but that the loop produces useful and transferable changes under a stable comparison.

\subsection{The route developed here}
The paper moves from representational failure to a finite repair problem, from that problem to a compiled procedure, and from the procedure to transferable and statistically accountable improvement. It then treats the succession of procedures as another finite control problem. This last step makes a common quantifier visible: the system needs one policy that can carry its claims forward, not a different hypothetical policy for every desired property.

The route does not require every useful correction to be formally proved beneficial on all future tasks. Exact semantic certificates and statistical performance evidence have different jobs. Nor does it require each system to invent its own computational substrate. A finite search language may be supplied. The theory must disclose what that language already makes possible and what the learner actually selects or extends.

\begin{figure}[!ht]\centering
\includegraphics[width=.96\linewidth]{sources/monographs/self-improving-intelligence/figures/two_loops.pdf}
\caption{The inner process repairs a solver. The outer process changes how repair is attempted. Arrows indicate proposed information flow, not an implemented neural agent.}\label{SII:fig:loops}
\end{figure}

\section{Related work and the proposed contribution}\label{SII:sec:related}
Learning how to learn has a substantial history. Learned optimizers can acquire update rules from optimization problems~\cite{SII:andrychowicz}. Model-Agnostic Meta-Learning trains parameters from which a small number of updates can adapt to a new task~\cite{SII:maml}. These approaches show why adaptation mechanisms deserve direct study. Neither the phrase meta-learning nor the presence of an outer objective, by itself, establishes an indefinitely improving outer procedure.

DreamCoder couples learned program libraries to neural search guidance~\cite{SII:dreamcoder}. Its relevance is the interaction between a language for expressing solutions and the ability to search that language. SEAL instead learns to produce adaptation data and update directives whose value is evaluated after the resulting parameter update~\cite{SII:seal}. These are different places at which an improvement process can become learnable.

The G\"odel-machine formulation treats a self-rewrite as admissible after a proof of usefulness within encoded assumptions~\cite{SII:goedel}. The Darwin G\"odel Machine uses empirical evaluation and an archive of modified coding agents~\cite{SII:dgm}. They should not be conflated. A proof-based criterion and a benchmark-based criterion support different guarantees and have different costs.

More recent systems sharpen the distinction between a task agent and its modifier. Hyperagents makes both roles editable and studies transfer of improvement strategies~\cite{SII:hyperagents}. Its reported experiments include improvements to modification procedures and transfer across tested domains. The archived variants and their selection protocol remain part of the evaluation setup. The general expressiveness of the editable program is not itself evidence that efficient discovery occurs on every computable task.

The September 2026 AIDE$^2$ preprint studies repeated modification of a research agent's harness and reports transfer beyond its selection tasks~\cite{SII:aide}. Its separate comparison of a discovered agent as an outer-loop improver is explicitly inconclusive relative to the strong reference. This distinction is directly relevant: improvements in research performance and improvements in driving the next improvement loop are related hypotheses, not interchangeable observations. These public works are research context. Their empirical findings are not treated as evidence for the architecture proposed here.

Causal abstraction provides a formal account of mechanism-preserving interpretation~\cite{SII:geiger}. Adaptive data analysis studies why repeated access to evaluation information can invalidate ordinary statistical reasoning~\cite{SII:dwork}. Both are necessary neighbors. A self-improver must know what its representations mean and what its evidence permits it to conclude.


\subsection{What this paper adds to that discussion}
The intended contribution is an integrated realization contract. A candidate improvement identifies its target, information access, experiment family, repair language, controller, retained evidence, and charged resources. It also identifies the comparison under which the procedure itself is said to improve. The exact results concern finite cases of that contract. The architectural hypotheses concern whether learned proposal and acquisition mechanisms can satisfy it economically on wider tasks.

This is not a claim that program libraries, recursive search, causal abstraction, or guarded updates originated with Komorebi. The aim is to connect them without allowing the cost or uncertainty of one stage to disappear inside another. A valid final repair is weaker than a valid procedure for finding it. A transferable repair is weaker than a transferable investigation. A safe accepted edge is weaker than a lineage with a continuing path to further improvement.

\section{A protected comparison for improvement procedures}\label{SII:sec:contract}
Let $S$ denote a solver implementation, $U$ an improvement procedure, and $L$ the library available to that procedure. The distinction is functional rather than physical. They may occupy one executable program. Their different roles are retained so that changes can be attributed rather than hidden in a single checkpoint name.

\begin{SIIdefinition}[Improvement contract]\label{SII:def:contract}
A contract $\Xi$ specifies a distribution $\mathcal H$ of starting-solver/task pairs, an initial information and library boundary, permitted observations and modifications, a resource allowance $B$, protected predicates, and an evaluation rule $J_\tau(S)\in[0,1]$. The improvement procedure returns one selected descendant or the unchanged baseline. Selection must use only the admitted development information. The final evaluation is not available as an oracle for choosing retrospectively among all candidates.
\end{SIIdefinition}

Write $U(S,\tau,L;B)$ for the selected result, including the procedure's randomness where applicable. Its productivity is
\begin{equation}\label{SII:eq:productivity}
\SIIRcal_\Xi(U,L;B)=\SIIE_{(S,\tau)\sim\mathcal H}
\big[J_\tau(U(S,\tau,L;B))-J_\tau(S)\big].
\end{equation}
Invalid execution and out-of-budget behavior receive the treatment fixed by $\Xi$, not an exclusion chosen after seeing the outcome. The score may be a primary coordinate of a larger profile. Resource, validity, and group-specific constraints need not be folded into one scalar.

The value in~\eqref{SII:eq:productivity} lies in $[-1,1]$. A procedure that may return the baseline need not achieve a nonnegative true value: noisy selection can still choose a worse descendant. A nonnegative certificate requires additional evidence. The mere existence of a fallback does not make every selector good at using it.

\subsection{Fixed-library and full-system comparisons}
A comparison of $\SIIRcal_\Xi(U',L;B)$ with $\SIIRcal_\Xi(U,L;B)$ isolates process change at a common library. A comparison with $(U',L')$ measures the combined descendant. Both are useful, but they answer different questions. A crossed comparison of $(U,L)$, $(U',L)$, $(U,L')$, and $(U',L')$ can expose interaction between procedure and library. There is no general additive decomposition that makes their contributions independent.

The same discipline applies to base models, training histories, tools, and stored controllers. A changed foundation model is a changed component. A foreign-agent test must supply a valid adapter or charge discovery of one. Calling two targets unfamiliar does not make them distributionally independent of the development process.

\begin{SIIcounterexample}[More knowledge without a better improver]\label{SII:ex:library}
Consider four equally likely task names. A fixed procedure returns a stored solution when one exists and otherwise leaves the solver unchanged. A library containing one answer yields a smaller expected gain than a library containing all four. The procedure has not changed. Holding the library fixed gives identical productivity for the two copies of the procedure. The full system improved, but this example contains no improvement of the improvement rule.
\end{SIIcounterexample}

\subsection{What persists across a recursive step}
A lineage state can be written schematically as
\[
Z_t=(S_t,U_t,L_t,E_t,\Xi_t,\rho_t),
\]
where $E_t$ records evidence and rejected alternatives, and $\rho_t$ records remaining resources and scope. A proposed modification may target any admitted subset. A comparison across different $\Xi_t$ needs a declared bridge. It cannot be justified merely because both scores are called accuracy.

The public record should distinguish which changes were selected by the system, which were supplied by an external designer, and which were only proposed. A successful finite search within a fixed grammar is a genuine computational result. It is not evidence that the grammar was invented, that a broader domain was covered, or that the outer search process improved itself.

\section{A small world with reusable diagnostic structure}\label{SII:sec:world}
Consider eight cells with exactly one latent defect, indexed by $j\in\{0,\ldots,7\}$. The solver must repair that cell while preserving the others. Its old representation retains only the number of defects, so every episode initially looks identical.

A diagnostic query names a subset $A$ and returns $\mathbf 1\{j\in A\}$. The reference language contains eight singleton queries and three bit-coordinate queries. Repair $r_i$ succeeds exactly when $i=j$. Because repair is irreversible within the episode, a procedure may commit only when every surviving possibility justifies its choice. Otherwise it may continue investigating or abstain.

The model is fully declared. Setting each known cell as the defect on a calibration copy would permit learning the complete query table in $8(8+3)=88$ observation calls. The code instead generates that table from the analytical specification and records it as supplied information. It does not claim to have executed those calibrations or inferred a hidden environment.

Equality-first testing asks whether $j=0$, then $j=1$, and so on. Its worst case needs seven queries, identifying the last cell by elimination. The three bit-coordinate queries need only three. Each procedure then performs one repair. Both complete binary trees have eight leaves and fifteen stored records. Faster execution has not made the controller free.



\chapter{Discovering and transporting repair}\label{SII:book:chapter-60}

\section{A residual constrains a diagnosis}\label{SII:sec:diagnosis}
A residual records disagreement between a prediction and an observation, between a proposed operation and a required invariant, or between two routes of translation. It is evidence against a conjunction of assumptions. It does not uniquely identify which assumption failed.

For example, an unexpected output can arise because a state variable was discarded, a transition was modeled incorrectly, the observation was noisy, or an action had a side effect omitted from the model. These explanations can agree on the current discrepancy and disagree about the next useful experiment. An automatic instruction to enlarge the representation after every error would confuse missing information with faulty dynamics.

Let a hidden case $\theta$ determine a set $\mathcal R(\theta)$ of successful repairs. The set may have several members. Diagnosis need only obtain enough information to choose a repair in the appropriate set. It need not identify every hidden property of $\theta$.

\begin{SIIproposition}[Feedback indistinguishability limits repair]\label{SII:prop:indistinguishable}
Suppose two cases have disjoint successful-repair sets and induce transcript laws $P_0,P_1$ under a fixed admitted investigation protocol. A repair selector observes only the transcript and independent private randomness. Under equal prior probabilities, its average probability of successful repair is at most
\[
\frac{1+\SIITV(P_0,P_1)}2.
\]
In particular, identical transcript laws rule out certain success in both cases.
\end{SIIproposition}
\begin{proof}
From the selected repair construct a binary decision: report case zero if the repair belongs to $\mathcal R(0)$, and case one otherwise. Success in case zero implies a correct decision. Success in case one also implies a correct decision because the two repair sets are disjoint. Thus repair success is at most the success of this binary test. For a randomized test with probability $f(o)\in[0,1]$ of reporting zero, the latter success is
\[
\tfrac12\big(1+\SIIE_{P_0}f-\SIIE_{P_1}f\big)
\leq\tfrac12(1+\SIITV(P_0,P_1)).
\]
The inequality follows by summing the positive differences of the two finite probability laws. The same argument applies to any measurable transcript space with the usual total-variation definition.
\end{proof}

The bound concerns the entire transcript, not each observation marginal in isolation. Correlations between observations can identify a case even when every separate marginal is uninformative. Randomized action selection is already covered when it is part of the declared protocol and its observable record.

Nor does the bound identify the most informative protocol. The learner must choose experiments before receiving their outcomes. A protocol that distinguishes the cases may exist while being inadmissible, too expensive, or absent from the current language. The finite synthesis below makes those possibilities separate outcomes.

\subsection{Repair sufficiency is not full identification}
Suppose $\mathcal R(\theta_1)=\{a,c\}$ and $\mathcal R(\theta_2)=\{b,c\}$. The cases are distinct, yet repair $c$ works in both. A method can stop without learning which case occurred. Conversely, pairwise overlap of repair sets does not guarantee one repair for an entire uncertainty set. The three sets $\{a,b\}$, $\{b,c\}$, and $\{a,c\}$ intersect pairwise and have empty total intersection. The stopping condition must refer to the whole surviving set.

This matters for representation discovery. A useful latent state may preserve exactly those distinctions needed for future repair decisions rather than every property of the underlying process. Whether that state remains sufficient after new repair operations become available is a separate question.

\subsection{Repairing a failed representation}
Let $X$ be a finite set with deterministic operations $F_u:X\to X$ and output labels $q(x)$. Begin with a partition $P_0$ that preserves both the old retained labels and $q$. Define the next partition by splitting a block whenever two states in it have successors in different current blocks under some permitted operation. Equivalently,
\[
x\equiv_{P_{t+1}}y
\quad\Longleftrightarrow\quad
x\equiv_{P_t}y\ \text{ and }\
F_u(x)\equiv_{P_t}F_u(y)\ \text{for every }u.
\]
The construction uses a complete transition description. Having a few observed contradictions does not give a learner that table for free.

\begin{SIIproposition}[Least stable refinement of a retained representation]\label{SII:prop:refinement}
The splitting construction stabilizes after at most $|X|-|P_0|$ strict refinements. Its final partition is the coarsest refinement of $P_0$ on which every $F_u$ has a well-defined quotient operation. It preserves every finite continuation of the retained labels and outputs.
\end{SIIproposition}
\begin{proof}
A strict refinement increases the number of nonempty blocks, which is bounded by $|X|$. At a fixed point, equivalent states have equivalent successors, so the quotient operations are well-defined. Induction on word length then preserves all retained labels after every finite word. For minimality, let $R$ be any stable partition refining $P_0$. If $R$ refines $P_t$, stability of $R$ implies that $R$-equivalent states have successors in the same $P_t$ blocks, so they remain together in $P_{t+1}$. Induction shows that $R$ refines every $P_t$ and hence the terminal partition.
\end{proof}

This is a standard partition-refinement argument expressed as a repair of an old representation. It is not necessarily the smallest representation for the outputs alone: $P_0$ may retain distinctions that are no longer needed. Discarding those distinctions requires a different contract. The continuation and cost perspectives are developed more fully in \emph{Machine Understanding}~\cite{SII:understanding}.

A counterexample pair is useful because it proves that a current block is too coarse under the stated transition model. It does not specify how a learned encoder can acquire the missing coordinate from its actual inputs. The mathematical quotient and its realizable acquisition procedure remain separate objects.

\section{One investigation, one repair, and two resource bounds}\label{SII:sec:synthesis}
We now give a finite model in which that acquisition problem can be decided exactly. The result is deliberately restricted, but it charges the complete investigation rather than only a selected final answer.

\begin{SIIdefinition}[Finite repair experiment model]\label{SII:def:repairmodel}
A model has a finite nonempty state set $X$, a finite experiment set $\mathcal E$, and a finite repair set $\mathcal A$. Experiment $e$ is allowed on a domain $D_e\subseteq X$, changes the state by $f_e:D_e\to X$, and reports $o_e:D_e\to O_e$. Its execution cost $c_e\geq0$ is fixed and known. Repair $a$ has fixed cost $c_a\geq0$ and succeeds on a declared set $G_a\subseteq X$. All costs are in one stated unit. A nonempty set $B_0\subseteq X$ records the possible current states.
\end{SIIdefinition}

A chosen experiment must be uniformly allowed: $B_0\subseteq D_e$. After observing $o$, the possible \emph{post-experiment} states are
\begin{equation}\label{SII:eq:post}
B_{e,o}=\{f_e(x):x\in B_0,\ o_e(x)=o\}.
\end{equation}
Only nonempty branches are retained. A terminal repair is justified when $B_0\subseteq G_a$. If the original state or an earlier obligation matters, it must be retained as an immutable tag in $X$. Merely changing the current state into a known state is not recovery of its history.

A diagnostic program is a finite rooted tree. An internal record names an experiment and branches on its possible report. A leaf names a repair justified on that branch's surviving states. We count one storage record per node of the complete tree. Shared subroutines, compressed code, and graphs with merged suffixes belong to other storage models. The present count is chosen because it is exact and unambiguous.

Let $W(P,B_0)$ be the largest execution cost of a complete path, including its terminal repair. Let $M(P)$ be the number of stored tree records. These costs are not independently minimized. A single tree has to realize both.

For horizon $h$, let $\mathcal F_h(B_0)$ be the nondominated attainable pairs $(w,m)$. A pair dominates another when it is no larger in either coordinate and strictly smaller in at least one. Let $T(B_0)=\{(c_a,1):B_0\subseteq G_a\}$. Then
\begin{align}\label{SII:eq:pareto}
\mathcal F_0(B_0)&=\SIIPareto T(B_0),\\
\mathcal F_h(B_0)&=\SIIPareto\left(T(B_0)\cup
\bigcup_{e:B_0\subseteq D_e}
\left\{\left(c_e+\max_o w_o,\ 1+\sum_o m_o\right):
(w_o,m_o)\in\mathcal F_{h-1}(B_{e,o})\ \forall o\right\}\right).
\nonumber
\end{align}
An empty child frontier makes that experiment unavailable for a successful complete tree. An empty union is not a zero-cost solution.

\begin{SIItheorem}[Exact finite repair synthesis with compiled storage]\label{SII:thm:synthesis}
For every finite model, nonempty $B_0$, and integer $h\geq0$, recurrence~\eqref{SII:eq:pareto} yields exactly the nondominated cost pairs of successful diagnostic trees of depth at most $h$. Every retained pair has an executable witness. A work/storage allowance $(b,m)$ is feasible if and only if some retained pair is componentwise at most $(b,m)$.
\end{SIItheorem}
\begin{proof}
At horizon zero only a justified terminal repair can succeed, giving $T(B_0)$. Assume the claim at $h-1$. A successful tree at horizon $h$ either stops immediately or has one uniformly allowed first experiment. Its child at each possible report is successful on exactly the post-experiment support in~\eqref{SII:eq:post}. The execution cost is the root cost plus the worst child cost. The storage is the root record plus the sum over every child, including children not traversed in a particular execution. Conversely, choosing a successful child for each report constructs a successful root tree with those costs. Replacing a dominated child by a dominating one never increases the root maximum or sum. Thus Pareto pruning loses no feasible allowance. Induction proves both completeness and witness construction.
\end{proof}

The theorem is an application of finite dynamic programming, not an efficient general solution to scientific discovery. Its state subsets and child combinations may grow exponentially. It assumes the transition, observation, legality, and success tables are correct. It also does not include the cost of constructing those tables in its online path cost. Acquisition of the model, compiler execution, and validation are separately charged development resources.

\subsection{A more informative question can be a worse repair question}
\begin{SIIcounterexample}[Entropy reduction need not minimize repair work]\label{SII:ex:entropy}
There are four equally likely cases. Repair $a$ succeeds in $\{0,1,2\}$ and repair $b$ succeeds in $\{3\}$. A balanced query separates $\{0,1\}$ from $\{2,3\}$. Another query separates $\{0,1,2\}$ from $\{3\}$. Both cost one unit and leave the state unchanged. The balanced query has the larger entropy reduction, but its second branch still needs another observation. The unbalanced query supports a correct repair after one observation in every case. Maximizing information about the entire hidden case is not the same objective as minimizing the cost of a successful correction.
\end{SIIcounterexample}

\begin{figure}[!h]\centering
\includegraphics[width=.94\linewidth]{sources/monographs/self-improving-intelligence/figures/repair_partition.pdf}
\caption{Two equally priced questions about four cases. The balanced split leaves one branch with incompatible repairs. The repair-aligned split identifies a valid action immediately. The figure represents exact sets, not measured agent behavior.}\label{SII:fig:repair}
\end{figure}

The example suggests a learning target more specific than general curiosity. A useful acquisition policy should distinguish alternative \emph{decisions} when needed, preserve unresolved alternatives when they matter later, and avoid spending work on irrelevant details merely because those details are easy to predict.

\subsection{What a compiled witness proves}
The reference implementation returns a concrete tree together with its cost pair. A separate interpreter checks its branch coverage, the domains of its experiments, the success of each leaf, and both resource counts. The checker does not infer that the environment model is true. It checks the relation between a specified model and the proposed procedure. This is the same reason an OSL semantic check and a learned Neuralese realization need separate evidence~\cite{SII:osl,SII:neuralese}.

\section{Investigating uncertainty in the experiment model}\label{SII:sec:uncertain}
The first synthesis theorem assumes that the experiment table is known. This is a substantial informational premise. A reasoner may know which experiments it can request while being uncertain about their outcomes, side effects, or interpretation. Optimizing one guessed table can produce a perfectly checked procedure for the wrong instrument.

A useful first extension retains several explicit possibilities rather than replacing uncertainty by one convenient estimate. Let $\Theta$ be a finite nonempty set of candidate models with a common state set $X$, experiment names, repair names, and nonnegative costs. Model $\theta$ specifies $D_e^\theta$, $f_e^\theta$, $o_e^\theta$, and $G_a^\theta$. The unknown model is fixed during an investigation. The initial information is a nonempty set $B\subseteq\Theta\times X$. The procedure observes reports, not the model label or the hidden physical state.

An experiment is allowed only when it is allowed under every pair still possible. Its update is
\begin{equation}\label{SII:eq:model-support}
B_{e,o}=\{(\theta,f_e^\theta(x)):(\theta,x)\in B,\ x\in D_e^\theta,\ o_e^\theta(x)=o\}.
\end{equation}
The model tag remains unchanged. Observations can exclude a model, but the algorithm cannot switch the underlying model between steps merely because doing so produces a convenient worst case. A terminal repair $a$ is justified when $x\in G_a^\theta$ for every $(\theta,x)\in B$.

This is a finite form of planning under model uncertainty. Robust dynamic programming provides an established neighboring literature~\cite{SII:nilim}. The present construction is not an application of a rectangular transition-uncertainty assumption. It explicitly remembers a fixed model label, including correlations between that model's predictions at different times.

\begin{SIItheorem}[Repair synthesis under a fixed unknown model]\label{SII:thm:uncertain-synthesis}
Construct a finite repair model on $\widetilde X=\Theta\times X$ by
\[
\widetilde f_e(\theta,x)=(\theta,f_e^\theta(x)),\quad
\widetilde o_e(\theta,x)=o_e^\theta(x),\quad
\widetilde G_a=\{(\theta,x):x\in G_a^\theta\},
\]
with the corresponding tagged domains and the common action costs. Applying Theorem~\ref{SII:thm:synthesis} to this model gives exactly the nondominated work/storage pairs of report-dependent trees that succeed for every initial pair in $B$, within the declared horizon. Every retained tree is valid for the actual system whenever its model and initial state belong to $B$.
\end{SIItheorem}
\begin{proof}
A tree receives the same experiment names and report alphabet in both descriptions. For each experiment and report, its possible tagged successor states are exactly~\eqref{SII:eq:model-support}. Its terminal repair is valid on the tagged support exactly when it is valid for every retained model/state pair. Consequently the report histories, path costs, branch counts, and stopping conditions agree node by node. Induction on tree depth gives a cost-preserving correspondence of all successful trees. The finite synthesis theorem then supplies the frontier and its witnesses. If the actual pair starts in $B$, each actual observation leaves its actual successor in the updated support, so the executed terminal repair is valid for that pair.
\end{proof}

This theorem decides a finite robust investigation problem. It does not learn an arbitrary model class or guarantee that the class contains the truth. Its negative result is also contract-relative. An empty frontier may mean that the admitted experiments cannot resolve the uncertainty, that a needed calibration is unavailable, or that the horizon is too short. It is not evidence that every possible experiment would fail.

\subsection{One hidden bit and two sensor conventions}
\begin{SIIcounterexample}[Known-model repair without a common unknown-model policy]\label{SII:ex:orientation}
Let $x\in\{0,1\}$ be the cell requiring repair and let $\theta\in\{0,1\}$ be a fixed sensor convention. A nondestructive query reports $o=x\oplus\theta$ at cost one. Repair $r_i$ succeeds exactly when $x=i$ and costs one. If $\theta$ is known, one query followed by the appropriate repair always succeeds. If all four $(\theta,x)$ pairs are initially possible, report zero leaves $(0,0)$ and $(1,1)$, while report one leaves $(0,1)$ and $(1,0)$. Repeating the same query gives no further information. No finite tree using only this query can guarantee a repair, although each known model admits a two-unit procedure.
\end{SIIcounterexample}

Add a permitted reference measurement that reports $\theta$, leaves $x$ unchanged, and costs two. The procedure can calibrate, make the ordinary query, and choose $r_{o\oplus\theta}$. Its worst-case work is four and its complete binary tree has seven records. Either measurement alone leaves both repairs possible. The two measurements are therefore both necessary in this example, regardless of their order. The executable reference construction returns exactly $(4,7)$ at horizons at least two.

The reference measurement is an added information channel. It is not a renamed latent coordinate and not a free inference from the original report. Its authority, reliability, and cost belong to the model. If the physical reference is faulty, the exact guarantee no longer follows. The same issue would arise if a learner excluded the true sensor convention after an unreliable observation.

\begin{figure}[htbp]\centering
\includegraphics[width=.97\linewidth]{sources/monographs/self-improving-intelligence/figures/uncertain_sensor.pdf}
\caption{An ordinary reading can identify the repair under a known sensor convention yet fail under two possible conventions. A separately specified reference measurement resolves the missing distinction. The figure is an analytical model, not a calibration experiment on an agent.}\label{SII:fig:uncertain}
\end{figure}

\subsection{Model refinement should be driven by repair relevance}
There is no general need to identify the whole model before acting. Several retained models may agree on a successful repair even while disagreeing about irrelevant outputs. The tagged support retains this possibility because its stopping condition is a common successful repair, not a singleton model label. Conversely, a calibration that removes many candidate models can still leave the decisive repair ambiguity intact.

A model learner therefore has two different targets. It can seek a more accurate description, or it can seek enough justified information to produce an affordable valid repair. The former can support the latter, but the two costs need not be minimized by the same investigation. The learner should retain the unresolved alternatives that matter to future decisions rather than certify an entire environment from one successful correction.

The finite construction also makes the cost of caution visible. The tagged state set can have $|\Theta||X|$ members, and the support search can be substantially larger than a search under one assumed model. Representation refinement, model acquisition, and compiler improvement are consequently coupled research problems. A more reliable model can increase the computational burden of planning until a better discovery procedure exploits its structure.

\section{Selecting a better procedure without hiding an oracle}\label{SII:sec:selection}
The finite compiler can itself be the object modified by an outer process. One candidate may change the experiment order. Another may change the cost model, add a repair constructor, or select a different representation for the uncertainty set. Their effects should not all be described as the same kind of learning.

For a fixed finite collection $\{U_1,\ldots,U_N\}$, suppose a declared assessment gives estimates $\widehat R_i$ of true productivities $R_i$. This assessment can be exact on a fully specified finite distribution or statistical under additional assumptions. Choosing a candidate must use the assessment actually available at selection time, not a hidden final score inspected afterward.

\begin{SIIproposition}[Selection under uniform assessment error]\label{SII:prop:selection}
If $|\widehat R_i-R_i|\leq\varepsilon$ for every candidate and $\hat\imath$ maximizes $\widehat R_i$, then
\[
R_{\hat\imath}\geq \max_iR_i-2\varepsilon.
\]
If a baseline is included and the selected candidate's estimated advantage exceeds $2\varepsilon+\gamma$, its true advantage exceeds $\gamma$.
\end{SIIproposition}
\begin{proof}
For a true maximizer $i^*$,
$R_{\hat\imath}\geq\widehat R_{\hat\imath}-\varepsilon
\geq\widehat R_{i^*}-\varepsilon
\geq R_{i^*}-2\varepsilon$.
Apply the same two error bounds to the selected candidate and the baseline for the second assertion.
\end{proof}

The uniform premise is the essential obligation. A small error for each candidate considered in isolation does not automatically control a candidate selected after adaptive access to the same evaluation record. Section~\ref{SII:sec:evidence} addresses an explicit alternative based on conditional validity.

\subsection{Changing the selector and changing the language}
An improved selector over a fixed candidate family can produce higher gains without making any previously impossible repair expressible. Conversely, a richer language can increase the best achievable gain while a weak selector never finds the new repair. These are distinct forms of headroom, already separated in the earlier SRIC treatment~\cite{SII:sric}.

For a fixed budget, let $R^*_{\mathcal G}$ be the largest productivity among procedures expressible in a declared language $\mathcal G$, with installed data and interpreter costs held fixed. The gap $R^*_{\mathcal G}-R(U)$ is search or selection headroom \emph{within that language}. Comparing $R^*_{\mathcal G'}$ with $R^*_{\mathcal G}$ concerns expressibility and resources. Neither upper envelope is automatically available to the learner.

A zero observed gain from one selector update is not a proof of saturation. A genuine saturation statement needs either a complete finite search or an upper bound close enough to the incumbent value. Otherwise the failure may reflect the proposer, the evidence channel, the search allowance, or the evaluation.

\subsection{The compiled instrument improver}
In the eight-cell instrument, the supplied language admits both singleton and bit queries. Recurrence~\eqref{SII:eq:pareto} can construct a depth-three diagnosis. The exported program replaces the equality-first diagnostic policy in an otherwise unchanged repair system. It does not contain the hidden defect index of the next episode.

Under a cap of three diagnostic calls and one repair call, the equality-first procedure certifies and repairs the first three cells, then abstains on the other five. On the uniform finite case distribution, its productivity from a zero-success baseline is $3/8$. The depth-three procedure certifies all eight cases, giving productivity $1$. These are exact properties of stipulated procedures. The compiler selected a better procedure within a supplied finite language. It did not invent the language, learn a general visual representation, or improve the compiler that performed this search.

That last distinction is deliberate. This first construction improves the diagnostic procedure, not its compiler. Chapter~\ref{SII:part:discovery} then changes the compiler, using certified early termination to construct the same optimal diagnostic controllers with fewer model-record reads. The finite selector that installs that compiler remains supplied. An archive entry saying that a descendant is more reflective supplies none of these separate claims.

\section{Transporting an investigation, not merely its result}\label{SII:sec:transport}
A repair can commute with a representation map even when the procedure used to discover it cannot be transferred. The investigation may depend on an unavailable observation, on a source-specific reset, or on the precise meaning of a stopping condition. A transport contract must include those actions and their evidence.

Consider source and target finite repair models. Let $\sigma:X\to X'$ be a bijection, $\tau:\mathcal E\to\mathcal E'$ a bijection of experiments, $\psi:\mathcal A\to\mathcal A'$ a bijection of repairs, and $\eta_e:O_e\to O'_{\tau(e)}$ a bijection of report labels. Assume
\begin{align}\label{SII:eq:transport}
\sigma(D_e)&=D'_{\tau(e)}, &
\sigma f_e(x)&=f'_{\tau(e)}\sigma(x),\\
\eta_e(o_e(x))&=o'_{\tau(e)}(\sigma x), &
\sigma(G_a)&=G'_{\psi(a)},\nonumber
\end{align}
and matching execution costs. These assumptions concern supplied semantics. Their discovery or learned approximation is an additional problem.

\begin{SIItheorem}[Exact transport of diagnostic procedures]\label{SII:thm:transport}
Under~\eqref{SII:eq:transport}, relabeling every experiment, branch, and repair transports successful diagnostic trees bijectively between the source support $B_0$ and target support $\sigma(B_0)$. Execution cost, complete tree storage, and depth are preserved. In particular,
\[
\mathcal F_h(B_0)=\mathcal F'_h(\sigma(B_0)).
\]
The same conclusion holds for every selected tree, not only a cost-optimal tree.
\end{SIItheorem}
\begin{proof}
A terminal repair is valid on $B_0$ exactly when its image is valid on $\sigma(B_0)$. For an internal experiment, domain preservation gives uniform admissibility. Transition and observation preservation yield
\[
\sigma(B_{e,o})=B'_{\tau(e),\eta_e(o)}.
\]
Apply induction to every child. The report bijection preserves the number of children, the maximum execution cost, and the sum of stored child records. Inverse maps give the converse correspondence. The equality of frontiers follows because costs and dominance are unchanged.
\end{proof}

The theorem covers a complete explanation of how the repair is found. It is stronger than equality of the repaired output but weaker than a claim that any learned representation admits such a correspondence. A known map costs something to store and execute. In the toy instrument, an arbitrary cell permutation requires a label adapter. A permutation of eight names is not discovered merely by allocating an eight-row table.

\subsection{Approximate performance transfer}
Suppose a source comparison gives productivity advantage $\Delta$ to $U'$ over $U$. Let $\widetilde U',\widetilde U$ be actual transported target procedures. If each procedure's productivity differs from its source value by at most $\varepsilon$, then
\begin{equation}\label{SII:eq:gaintransport}
\SIIRcal'(\widetilde U')-\SIIRcal'(\widetilde U)
\geq\Delta-2\varepsilon.
\end{equation}
This follows by subtracting the two inequalities. A common source of the premise is an exact task correspondence plus a bounded evaluation-law shift. For any fixed gain function in $[-1,1]$, changing the task law by total variation $\delta$ changes its expectation by at most $2\delta$. A further pointwise gain discrepancy $\eta$ therefore gives $\varepsilon=\eta+2\delta$.

This elementary estimate does not certify admissibility or sure success. A rare unauthorized action can have small average loss and still violate a zero-tolerance constraint. Such constraints require their own preservation check. The same is true of uncharged translation work: an equal-output target procedure can fall outside the budget after the adapter is added.

\begin{SIIcounterexample}[Output transfer with no transferable investigation]\label{SII:ex:unavailable}
A source can inspect a hidden bit for one unit and then choose one of two disjoint successful repairs. The target has the same hidden states and repair effects, but no operation that reveals the bit. The final repair meanings correspond exactly. The source diagnostic policy does not transport, and no target policy can guarantee success from the two-state uncertainty set. The missing object is an information channel, not a better translation of the repair names.
\end{SIIcounterexample}

\section{Growing a repair language without losing its old scope}\label{SII:sec:grammar}
A persistent discrepancy can be a limitation of the candidate language rather than the search procedure. Repeatedly optimizing coefficients cannot produce a nonlinear term in a language that expresses only affine functions. Reordering uninformative experiments cannot separate two cases whose transcripts are always identical. A learner must distinguish an unsuccessful search from a certified absence of a suitable expression.

A conservative extension embeds each old program into a new language with the same meanings on the old domain. That semantic property is valuable, but it does not imply nonregression under a fixed resource allowance. The new interpreter may consume more time, memory, or evidence before executing the embedded old program.

\begin{SIIproposition}[The resource price of conservative extension]\label{SII:prop:extension}
Suppose an embedding preserves every old procedure's successful behavior and adds at most $d\geq0$ execution work and $s\geq0$ stored records. Every old allowance $(b,m)$ is then feasible in the new language under allowance $(b+d,m+s)$. Feasibility at the unchanged allowance follows only when the embedded procedure fits that allowance. Semantic conservativity alone supplies no such guarantee.
\end{SIIproposition}
\begin{proof}
Embed a witness for the old allowance. By hypothesis its behavior is preserved and its costs are bounded by the shifted pair. The unchanged-budget assertion requires the corresponding inequalities. For the final claim, an old one-unit procedure can be embedded with identical output into a new interpreter requiring two units of startup work. It remains semantically available but cannot execute within the old one-unit allowance.
\end{proof}

A practical extension should retain a directly executable old path when possible, then route only justified cases through the new machinery. The routing operation also has a cost and an information requirement. A context-dependent grammar is not a solution if deciding the context already requires solving the original problem.

\subsection{Error localization and new constructors}
The useful role of packets and residuals is to make this diagnosis structured. A packet can declare a domain, a retained distinction, an uncertainty set, an operation interface, and unresolved conditions. A residual identifies which preservation equation or evidence constraint failed. A constructor adds a new kind of candidate repair. The theory does not assume that a packet is an independent physical component or that its name supplies a new algorithm.


\chapter{Improving the procedure that discovers repairs}\label{SII:book:chapter-61}
\label{SII:part:discovery}
\section{A certificate can remove the need for exhaustive search}\label{SII:sec:compiler-change}
The compiled tree in the first example is an improved diagnostic policy. This section changes the procedure that constructs such trees. The distinction is concrete: the parent and descendant receive the same new problem table and return equally valid optimal controllers, but they need not perform the same amount of discovery work.

The parent is the exact frontier compiler of Theorem~\ref{SII:thm:synthesis}. It examines a finite set of support/horizon states and combines successful subtrees. That generality is useful, but it can do unnecessary work when one inexpensive candidate already meets a proved global lower bound. A descendant compiler can first propose such a candidate, check it independently, and stop only if its checked cost proves that further search cannot improve either protected resource coordinate.

This combines a proposal mechanism with a decisive rejection or acceptance condition. Synthesis from partial programs and counterexamples has established precedents~\cite{SII:sketch}. The particular addition here is a certificate that can close an \emph{optimality} search, rather than merely establish that one candidate is valid. Neither the proposal template nor the lower-bound argument is assumed to have been invented by the system.

\subsection{The scope in which a short proof is decisive}
Consider $n\geq1$ possible immutable cases. Every admitted query is total, nondestructive, binary-valued, and costs one unit. Every available repair costs one unit and succeeds on exactly one case. At least one repair is available for each case. The target is certain successful repair from the complete set of $n$ possibilities. A literal controller stores one record per query or repair node. Redundant unary query nodes are permitted but have their full costs.

\begin{SIItheorem}[Checked lower-bound closure of a diagnostic search]\label{SII:thm:optimality-stop}
Every successful controller in this scope has worst-case execution work at least
\[
1+\lceil\log_2 n\rceil
\]
and storage at least $2n-1$. If an independently checked candidate attains both bounds and respects the requested horizon, its cost pair is the entire nondominated frontier. Returning this candidate is therefore equivalent in the protected work/storage objective to completing the exhaustive synthesis.
\end{SIItheorem}
\begin{proof}
Every case produces one deterministic query history. Two different cases cannot terminate at the same repair leaf because no repair succeeds in both. Thus there are at least $n$ nonempty leaves. A tree of binary query depth $d$ has at most $2^d$ leaves, so $d\geq\lceil\log_2 n\rceil$. Every path also executes a unit-cost repair. For storage, let $l$ be the number of leaves and let $i_1,i_2$ count unary and binary internal nodes. Counting edges in a finite rooted tree gives $i_1+2i_2=i_1+i_2+l-1$, hence $i_2=l-1$. Total records are $2l-1+i_1\geq2n-1$. A successful candidate attaining both coordinatewise lower bounds dominates every other attainable cost pair. Its independent validity check and horizon check make it a feasible witness, so the Pareto frontier contains only that pair.
\end{proof}

The immutable-case premise is essential. An experiment that resets every physical state into one safely repairable state can defeat a case-identification lower bound when only the current state matters. Multiple cases admitting a common repair can do the same. The guard must therefore verify the actual contract before applying the bound. This is not a universal theorem that diagnosis always requires identifying the hidden state.

A candidate that is valid but misses one lower bound does not close the search. Nor does failure of the proposal rule prove that no optimum exists. In both cases the descendant falls back to the general compiler. The certificate licenses an early return, not an early claim of impossibility.

\begin{figure}[htbp]\centering
\includegraphics[width=.96\linewidth]{sources/monographs/self-improving-intelligence/figures/certified_compiler.pdf}
\caption{The descendant changes discovery control flow. A checked candidate must attain both global lower bounds before the compiler can stop early. Failure of a guard, proposal, or bound check invokes the unchanged exact compiler. The sequential reference implementation charges unsuccessful preliminary work.}\label{SII:fig:compiler}
\end{figure}

\subsection{The inexpensive proposal is not the proof}
The reference proposer repeatedly chooses a query whose two current branches have the smallest maximum cardinality, breaking ties by the installed query order. It searches the public observation table rather than recognizing privileged names such as ``bit zero.'' It constructs a tree recursively, stopping when a unique justified repair remains. It aborts a proposed branch when the remaining binary depth cannot distinguish its cases.

Balancing is a supplied heuristic. It need not minimize repair work in general, as Counterexample~\ref{SII:ex:entropy} already shows. It is useful in the narrower singleton-repair scope because the required distinctions coincide with case identification. Even there, a locally balanced split need not extend to a globally optimal tree for an arbitrary query library. Correctness comes from the finished witness and the lower bound, not from the heuristic's name.

The checker reconstructs every possible support, checks all reports, and independently replays every initial case. A fabricated terminal label, missing branch, insufficient horizon, or altered repair effect is rejected. This is a finite model checker authored within the project. It is not a proof of the physical observation table or an independent proof-assistant kernel.

\section{The cost of certifying a better compiler}\label{SII:sec:compiler-cost}
The unit used in the following comparison is a consumed record of the supplied transition/report table. A repeated access counts again. The unit is chosen because it is explicit and can be instrumented in the unchanged parent compiler and the descendant. It is not a claim to count every Python operation, allocation, set-membership test, processor instruction, or unit of energy.

Both procedures start with cold search caches and the same structurally validated table. Shared input-shape validation is recorded separately. The descendant additionally charges its semantic guard, candidate construction, and completed witness check. The parent charges its complete frontier search and the same witness-checking procedure for its returned frontier. Model acquisition, source storage, and the finite procedure-selection stage remain separate ledger entries.

\begin{SIIproposition}[Read cost of the certified binary-family shortcut]\label{SII:prop:compiler-work}
Let $n=2^k$ with $k\geq1$. Suppose there are $q$ total unit-cost binary immutable queries, and the balanced proposer produces a perfect depth-$k$ tree. In the stipulated implementation, the semantic guard reads every query row once, the proposer scans all $q$ queries at each internal support, and verification performs one support reconstruction and one complete case replay. The total number of transition/report record reads is
\begin{equation}\label{SII:eq:compiler-reads}
R_{\mathrm{shortcut}}=nq(k+1)+2nk.
\end{equation}
The usual hidden-cell family with all singleton queries and $k$ coordinate queries satisfies the premise, including after arbitrary cell relabeling, report-bit flips, or reordering of the query list.
\end{SIIproposition}
\begin{proof}
The guard uses $nq$ reads. At each of the $k$ query levels, the current supports partition the $n$ cases, so scanning all queries uses $nq$ reads per level. Every case traverses $k$ query nodes in the support reconstruction and again in the independent replay, contributing $2nk$. Summing gives~\eqref{SII:eq:compiler-reads}. In the hidden-cell family, every nonsingleton support reached by coordinate splits is a subcube. A remaining coordinate divides it equally, while a singleton query cannot give a smaller maximum branch. Thus the balancing rule chooses an equal split until singleton leaves are reached. When a two-case support admits a singleton tie, either split reaches singleton leaves. Relabeling the cases or reversing report names preserves all branch cardinalities and this argument.
\end{proof}

The result is about a particular implemented access pattern. Caching rows, recognizing trusted structure symbolically, or using another cost unit changes the formula. A specialized cube compiler supplied with certified coordinate semantics can avoid the general search too. The result does not establish superiority to that informed comparator or to every available synthesis algorithm.

\subsection{Some verification work cannot disappear}
\begin{SIIproposition}[A table-access floor for universal scope certification]\label{SII:prop:table-floor}
Suppose a deterministic validator receives an arbitrary table of $N\geq1$ independently variable records only through individual record reads. It must accept every table whose records all satisfy a fixed per-record predicate and reject every table with a violating record. Assume each record position admits both a satisfying and a violating value. On every accepted table, the validator must read all $N$ positions at least once.
\end{SIIproposition}
\begin{proof}
Suppose it accepts a satisfying table while leaving position $j$ unread. Change only that position to a violating value. Every record the validator reads has the same value, so determinism gives the same sequence of reads and the same acceptance. This contradicts rejection of every table containing a violation. Therefore no accepted table can have an unread position.
\end{proof}

In the compiler's unrestricted table interface, an unread entry could contain a destructive transition or an additional report value. The full semantic scope cannot be certified merely by recognizing a few apparently balanced rows. Trusted construction metadata or a proof from a different representation can change the interface and reduce this cost, but that metadata then becomes a charged assumption or a separate proof object.

This identifies a useful bottleneck. Improving search can make model access and verification dominate the remaining computation. It does not make every stage improve at the same rate. Nor does it establish that this bottleneck is permanent, since a learned representation might eventually support cheaper certificates. Such a representation would need its own acquisition and soundness argument.

\section{A bounded selection of a better discovery procedure}\label{SII:sec:method-selection}
We now make the persistent change explicit. The candidate catalogue contains the unchanged parent compiler, a guarded first-nontrivial-query proposer with exact fallback, and a guarded balanced proposer with exact fallback. The guard, proposal constructors, lower-bound proof, checker, and fallback are supplied. The selection procedure runs each candidate on the four-cell and eight-cell development problems, with the same tables and cold caches. It selects the candidate with the smallest total number of charged record reads and writes a policy record before evaluating any additional problems.

The developmental totals are $10{,}256$ reads for the parent, $10{,}380$ for the guarded first-query rule, and $488$ for the guarded balanced rule. The selector therefore installs the balanced rule. These totals concern a fixed, fully enumerated development catalogue. They do not imply a statistical generalization guarantee over an unknown problem distribution. The selected artifact changes the control flow of future compilations without containing the hidden state of a future repair episode.

\begin{center}\begin{tabular}{@{}rrrr@{}}\toprule
Cases & Parent compiler reads & Descendant reads & Returned $(\text{work},\text{storage})$\\\midrule
2 & 16 & 16 & $(2,3)$\\
4 & 352 & 88 & $(3,7)$\\
8 & 9,904 & 400 & $(4,15)$\\
16 & 576,128 & 1,728 & $(5,31)$\\\bottomrule
\end{tabular}\end{center}

Every paired entry includes complete checking of the returned controller. The sixteen-cell problem is evaluated only after the policy record is fixed. Additional cold cases remove singleton queries, restrict the initial support to six possibilities, reverse operation order, and change cell and report labels. These vary the finite structure available to the compiler. They are not evidence of transfer to an unrelated scientific domain.

Under a common cap of $2{,}000$ table reads, the descendant returns a checked optimal sixteen-cell controller. The parent reaches the cap without returning its complete frontier. Its outcome is recorded as \emph{incomplete within the allowance}, not as a proof that no controller exists. Without that cap, the parent also returns the same optimal cost pair. The difference lies in discovery work rather than repair validity or online diagnostic depth.

\begin{figure}[htbp]\centering
\includegraphics[width=.9\linewidth]{sources/monographs/self-improving-intelligence/figures/compiler_read_cost.pdf}
\caption{Charged transition/report record reads on the exact binary family, including witness checking. The vertical axis is logarithmic. Both compilers return the same optimal online work and stored tree size. These deterministic algorithm counters are not measured agent-learning performance or total CPU instructions.}\label{SII:fig:compiler-cost}
\end{figure}

\subsection{What was selected and what remains supplied}
The persisted policy chooses a proposal rule from a fixed catalogue. It is not a stored solution to the next hidden-cell episode. It therefore modifies an improvement procedure rather than only adding another repaired answer. The compiler accepts arbitrary supplied tables, applies its semantic guard, and retains the exact fallback outside the shortcut's demonstrated scope.

This is nevertheless a restricted form of process learning. The catalogue does not invent a new lower-bound theorem or discover a new programming primitive. It does not improve the selector that searched the catalogue. A second call with the same candidates and the same criterion returns the same selected rule. One genuine improvement of a discovery procedure is not an indefinite sequence of improving outer procedures.

The distinction gives the next research target a concrete form. A learned proposer would need to generate useful shortcut templates or representations not simply chosen from this fixed list, while the checker verifies their stated scope. A learned evaluator would need to identify when existing cost or success measurements omit a consequential requirement. Such changes would alter additional parts of the realization and need their own preserved comparison.

\subsection{Fallback protects meaning, not free execution}
\begin{SIIcounterexample}[A semantically protected shortcut with a slower fallback]\label{SII:ex:slow-fallback}
Give the compiler four immutable cases and only four singleton queries. The horizon permits the three-query equality tree. A balanced proposal produces a valid controller but cannot attain the two-query information lower bound, since the query library does not supply the needed split. The descendant therefore performs its preliminary guard and proposal work and then runs the exact parent. It returns the same frontier but consumes more table reads. Semantic nonregression has not established work nonregression.
\end{SIIcounterexample}

The counterexample is included in the shipped comparison. The declared development objective prefers the shortcut on its development family, not on every possible input. Changing the mixture of structured and fallback cases can change which procedure is preferable. Attributing all gains to the new procedure while excluding its failed proposal attempts would invalidate the comparison.

There is a conditional way to bound this risk through scheduling rather than prediction. The following result is intentionally separate from the sequential implementation above.

\Needspace{16\baselineskip}
\begin{SIIproposition}[Bounded slowdown from a certified two-worker race]\label{SII:prop:race}
Suppose a parent and a shortcut worker can be paused after each unit of a common work measure and resumed without losing progress. Each may announce completion only after producing a result certified for the same target. Alternate one unit of each worker and stop at the first certified completion. If their standalone completion costs are $D_0,D_1\in\mathbb N_{>0}\cup\{\infty\}$ and at least one is finite, at most $2\min\{D_0,D_1\}$ worker units are executed. If each dispatched unit costs at most an additional $c\geq0$ scheduler units in the same measure, total charged work is at most $2(1+c)\min\{D_0,D_1\}$.
\end{SIIproposition}
\begin{proof}
After $2d$ dispatched worker units, each worker has received exactly $d$ units of progress. Set $d$ to the smaller finite standalone completion cost. That worker has completed its certification by this point, so the race has stopped no later. Multiplying the number of dispatches by the maximum combined work per dispatch gives the overhead bound.
\end{proof}

The theorem does not provide resumability for an arbitrary compiler, bound the memory occupied by both workers, or make their setup free. The present reference compiler uses sequential fallback, not this race. The companion suite checks the abstract schedule bound separately on finite worker tapes. A systems implementation would need to include state checkpointing, dispatch cost, shared-resource contention, and both workers' storage before claiming this guarantee.

Taken together, the results distinguish three achievements: a heuristic frequently proposes a good answer, a certificate establishes that the answer closes the search, and an implementation retains an affordable fallback when the heuristic fails. None of these can substitute for the others.

\subsection{A complete receipt for the sixteen-cell compilation}
The selected policy is evaluated on a fresh sixteen-cell table after its bytes are fixed. The new table has twenty query rows: sixteen singleton queries and four coordinate queries. The next hidden defect is not supplied to either compiler. Both methods construct a complete controller that will work for any of the sixteen cases.

The descendant's recorded access account is
\begin{center}\begin{tabular}{@{}lr@{}}\toprule
Compiler stage & Consumed transition/report records\\\midrule
Semantic scope guard & 320\\
Proposal construction across four levels & 1,280\\
Support reconstruction by the checker & 64\\
Independent replay of all cases & 64\\\midrule
Total & 1,728\\\bottomrule
\end{tabular}\end{center}
The separate common input-shape validation consumes another 320 records before the measured compilation begins. It is not included in the $2{,}000$-read compiler allowance. Counting that shared setup on both sides gives 2,048 reads for the descendant and 576,448 for the parent. No particular implementation is permitted to hide its extra guard inside an uncharged setup step.

The three-candidate development comparison itself consumed 21,124 model-record reads. It is part of the cost of selecting the descendant, not part of one later compilation. For the declared sixteen-cell workload the per-use read saving is 574,400, so those developmental reads are recovered within one use in that same narrow accounting unit. This does not repay an unspecified human research cost, prove a wall-clock speedup on every machine, or charge the creation of the candidate templates. Source sizes and measured runtimes are recorded separately rather than silently converted into model-record reads.

The result concerns cheaper discovery of equally valid controllers. It does not establish an improved selector or an independently discovered compiler catalogue.


\chapter{Learning a compiler's rejection rule}\label{SII:book:chapter-62}
\label{SII:part:rule-learning}
\section{A failed search representation can itself be repaired}\label{SII:sec:search-features}
The preceding selector chooses among three supplied compiler strategies. It changes which algorithm will be used, but does not compose a new decision rule. We now move one step inside a compiler. A small learner constructs a predicate that tells the compiler when it can stop inspecting an unpromising query. The learned object is neither a hidden-cell answer nor a completed diagnostic tree. It is an operation on the compiler's own partial information, reused while constructing trees for further problems.

The construction is deliberately restricted. Its experiment tables, compiler skeleton, primitive arithmetic language, and correctness argument remain supplied. Within that boundary, the learner is not given the complete predicate as one named strategy. It first encounters a representation that cannot express its assigned distinction, adds a specified missing feature, and synthesizes a formula from counterexamples. A structural certificate then determines whether the formula may safely prune later searches.

This separates three questions that are easy to conflate: whether the examples determine a suitable formula, whether the formula is valid outside those examples, and whether executing it actually reduces the relevant cost. The reference construction answers each under a different explicit contract.

\begin{figure}[!ht]\centering
\includegraphics[width=.95\linewidth]{sources/monographs/self-improving-intelligence/figures/rule_synthesis_pipeline.pdf}
\caption{One complete bounded improvement. A feature conflict justifies exposing the supplied capacity feature. Counterexamples select a composed arithmetic expression. A narrow structural check links that expression to a proved scope, after which it can skip rows inside later compilations. The grammar, teacher, scope proof, and compiler skeleton remain outside the learned component.}\label{SII:fig:rule-pipeline}
\end{figure}

\subsection{The information available inside a search}
A compiler is considering a binary query on a support of $m$ possible immutable cases. It reads the query's public model rows in a fixed order. After a prefix, $a$ of the inspected cases have report zero and $b$ have report one. There are $m-a-b$ rows left to read. If this query is the first of at most $d\geq1$ remaining queries, either child can contain at most
\[
C=2^{d-1}
\]
cases when distinct cases need distinct repair leaves. The local cardinality premise $m\leq2C$ is checked before a root is investigated. An overfull support cannot fit the remaining binary depth, even before any rows of the proposed query are read.

The initial search representation contains $(a,b,m)$ but omits $C$. This can lose a distinction relevant to the search itself. For example, $(a,b,m)=(0,3,4)$ cannot have a completion whose two branch sizes are at most two. It can have a completion with branch sizes $(1,3)$ when the allowed child capacity is four. The read prefix is the same. The permitted future investigation is different.

\begin{SIIproposition}[Feature sufficiency for a search decision]\label{SII:prop:search-feature}
Let $\mathcal P$ be a finite set of search contexts, let $r:\mathcal P\to Z$ be their available feature map, and let $\ell:\mathcal P\to\{0,1\}$ be the required decision label. A deterministic rule $f:Z\to\{0,1\}$ satisfying $fr=\ell$ exists if and only if $\ell$ is constant on each nonempty fibre of $r$. A pair in one fibre with different labels rules out every such rule, irrespective of its program length.
\end{SIIproposition}
\begin{proof}
If $f$ exists, equal features have the same output under $f$ and hence the same label. Conversely, assign to each attained feature value its common label. This is well defined when the fibre condition holds. Values outside the image of $r$ can be assigned arbitrarily. A conflicting pair violates necessity for every possible $f$.
\end{proof}

The criterion is elementary factorization, not a new general learning theorem. Its role is to distinguish insufficient information from inadequate search effort. In the finite construction, exhaustive search for a larger formula over $(a,b,m)$ cannot fix the conflict. The supplied feature registry contains the remaining-depth capacity $C$. Exposing it resolves this particular obstruction. Computing $C$ from an already known depth is not a physical experiment and does not identify the next episode's hidden defect.

\begin{SIIcounterexample}[The same prefix requires different rejection decisions]\label{SII:ex:missing-capacity}
For the two contexts $(a,b,m,C)=(0,3,4,2)$ and $(0,3,4,4)$, the initial feature vector $(a,b,m)$ is identical. In the first, three cases already occupy a child with capacity two. In the second, assigning the remaining case to the other child gives a capacity-feasible completion. Thus the predicate ``no capacity-feasible completion exists'' cannot be computed from the initial features alone. Adding the capacity makes this specific decision representable. It does not establish that the actual query library can distinguish every resulting child.
\end{SIIcounterexample}

This is a small instance of representational revision inside a discovery algorithm. The target of learning is not only a model of the external task. It is a model of which search continuations remain possible under the compiler's own constraints. A richer representation is useful because a named counterexample requires it, rather than because extra variables are assumed to improve intelligence in general.

\section{Composing a rule from primitive syntax}\label{SII:sec:grammar-learning}
The enlarged feature vector is $(a,b,m,C)$. The reference language has six integer atoms, $a,b,m,C,0,1$, three binary arithmetic operations, $\max$, $\min$, and addition, and the five comparisons $>,<,\geq,\leq,=$. A candidate is a Boolean comparison with at most five syntax nodes. No general program execution is admitted. The interpreter evaluates only these finite arithmetic trees.

Integer atoms have size one. A binary arithmetic expression has size three. A comparison of two atoms has size three, while a comparison with one arithmetic child has size five. For commutative arithmetic, operand pairs are generated in one fixed order. This produces 180 size-three formulas and 3,780 size-five formulas, or 3,960 candidates in all. Different syntaxes can express the same function. No semantic deduplication is assumed in these counts.

This is a bounded instance of syntax-guided synthesis, in which the candidate language and a semantic requirement are specified separately~\cite{SII:sygus}. Counterexamples guide the choice of candidate. Learning invariants from examples, counterexamples, and implication constraints is also an established research direction~\cite{SII:ice}. The present learner does not implement the full ICE framework and makes no claim to discover general loop invariants. It learns one rejection expression under a supplied search contract.

\subsection{What the teacher knows}
Development uses the 198 count-prefix contexts with $C\in\{1,2,4\}$, $2\leq m\leq2C$, and $a,b\geq0$, $a+b\leq m$. For each context, the teacher enumerates the possible numbers of report-one cases among the unread rows. It labels the context as reject precisely when no such completion can put at most $C$ cases in either branch. The teacher therefore solves a small finite feasibility problem. It is not reading a label generated by the final formula.

After each returned counterexample, the learner retains that context and its label. It selects the earliest formula in the fixed syntax enumeration that agrees with every retained example. The finite teacher either returns another disagreement or checks the formula against the complete developmental domain. The definition of that domain, the generation of its labels, and all rejected candidates are part of the development account.

\begin{SIItheorem}[Finite counterexample refinement within a grammar]\label{SII:thm:grammar-refinement}
Let $\mathcal G$ be a finite nonempty set of deterministic candidate programs, let $\mathcal P$ be a finite domain, and let $\ell$ be its specified labeling. Suppose a teacher either returns a point where the proposed program differs from $\ell$ or certifies agreement on all of $\mathcal P$. A learner that proposes only candidates consistent with every retained counterexample rejects each proposed program at most once. It returns a correct program on $\mathcal P$, or certifies that none in $\mathcal G$ meets the labeling, after at most $|\mathcal G|$ teacher calls. With a fixed total order and the earliest-consistent selection rule, any returned program is the earliest one that agrees on the entire domain.
\end{SIItheorem}
\begin{proof}
A counterexample removes the current program because it must be included in the constraints for every later proposal. Hence that program cannot be proposed again. A program agreeing with $\ell$ on all of $\mathcal P$ is never removed by a valid counterexample. There are only $|\mathcal G|$ candidates, so after at most that many distinct proposals the teacher has accepted one or all candidates are excluded. If a returned program were preceded by a globally correct program in the fixed order, the earlier program would remain consistent and would have been selected instead. This proves the ordering assertion.
\end{proof}

The result concerns the declared finite domain. It is not a generalization bound and supplies no efficiency guarantee for obtaining the teacher's answers. A teacher timeout is not a certificate. If the grammar does not contain a suitable program, enlarging the grammar changes the search problem and its cost.

\subsection{The synthesized expression and its provenance}
The reference learner encounters three disagreements. Its first candidate, $a>a$, fails on a prefix with too many ones. The next proposed candidate, $b>a$, rejects a prefix that can still be completed. The third, $b>C$, misses a prefix with too many zeros. The fourth teacher proposal is
\begin{equation}\label{SII:eq:learned-prune}
\rho(a,b,m,C)=\mathbf 1\{\max(a,b)>C\}.
\end{equation}
The formula uses five syntax nodes and passes all 198 developmental contexts. The enumeration generates all 3,960 AST roots and considers 190 of them, including candidates rejected directly by retained examples. There are 282 retained-example checks, 210 full-domain teacher comparisons, 344 completion-count checks to construct the labels, and 1,898 interpreted syntax-node visits in the learner's recorded comparisons. The full grammar has 19,440 literal node occurrences. Shared in-memory subexpressions can reduce actual storage, so this last quantity is not a measured byte count.

These numbers describe an exact finite synthesis, not a statistical training cohort. The formula is serialized and frozen before evaluation on larger supports, additional capacities, and changed row and query orders. No new episode's defect index or solved controller is placed in the formula. It contains only the reusable relation between two observed prefix counts and remaining capacity.

The learner discovers the composition of supplied primitives. It does not invent addition, maximum, the teacher, the resource model, or the eventual safety proof. This boundary is important: enumerating a finite grammar is more expressive than choosing among the earlier three named compiler policies, but it remains a bounded search space chosen in advance.

\section{When an unfinished read can close a search branch}\label{SII:sec:prefix-pruning}
A finite fit is insufficient grounds for deleting possibilities from a later search. Pruning is stronger than proposing. A bad proposal can be rejected by a finished-witness checker. A bad pruning rule can remove the only successful controller and leave nothing for that checker to inspect. We therefore attach the generated expression to a separate scope argument.

\begin{SIIproposition}[Exact obstruction to a capacity-feasible completion]\label{SII:prop:prefix-capacity}
Let $C,m,a,b$ be integers with $C\geq1$, $2\leq m\leq2C$, $a,b\geq0$, and $a+b\leq m$. A binary completion of the remaining $m-a-b$ reports with both final branch counts at most $C$ exists if and only if $\max(a,b)\leq C$. Consequently, a prefix satisfying $\max(a,b)>C$ can be rejected before its remaining rows are read, without rejecting any completion that satisfies these two capacity constraints.
\end{SIIproposition}
\begin{proof}
Counts do not decrease as rows are read, so $a>C$ or $b>C$ makes a feasible completion impossible. Conversely, suppose $a,b\leq C$. Let $r=m-a-b$. There are $C-a$ available places in the zero branch and $C-b$ in the one branch. Both are nonnegative and their sum is $2C-a-b\geq r$. Distribute the $r$ remaining cases between these two capacities, filling either branch up to its capacity and putting any remainder in the other. This gives the required completion. The argument concerns possible report assignments, not whether an installed query has that assignment.
\end{proof}

The theorem applies to arbitrary positive integer capacities. The connection to query depth uses $C=2^{d-1}$. An actual search separately checks that both observed branches are nonempty, since a redundant unary query cannot be part of a controller attaining the minimal literal storage. It also verifies all later subtrees. Capacity feasibility is necessary, not sufficient, for actual repair solvability.

\begin{figure}[!ht]\centering
\includegraphics[width=.90\linewidth]{sources/monographs/self-improving-intelligence/figures/prefix_capacity.pdf}
\caption{An exact count-prefix model with eight cases and child capacity four. A prefix outside the square $a,b\leq4$ cannot be completed within that capacity. Prefixes also satisfy $a+b\leq8$. Surviving the count test does not guarantee that the installed query library can separate the children. The dotted arrow is one unfinished read that can stop immediately after its fifth equal report.}\label{SII:fig:prefix-capacity}
\end{figure}

\subsection{A deliberately narrow certificate checker}
The executable gate accepts only ASTs that normalize by swapping a comparison's sides or maximum's arguments to $\max(a,b)>C$. It rejects other finite-consistent formulas rather than claim to prove arbitrary arithmetic identities. Thus the universal justification is the preceding prose proof, coupled to a project-local structural checker and a runtime scope guard. This is not a Lean certificate or an independent theorem-proving service.

\begin{SIIcounterexample}[Agreement on training prefixes does not establish scope]\label{SII:ex:prune-extrapolation}
The formula $\max(a,b)>\min(C,4)$ agrees with~\eqref{SII:eq:learned-prune} on every developmental context because every developmental capacity is at most four. At $(a,b,m,C)=(0,5,16,8)$ it rejects, while a completion with final branch sizes $(8,8)$ exists. A larger input has not merely changed the frequency of familiar cases. It has reached a scope on which the bounded-cap formula was never justified. The structural gate does not accept that formula.
\end{SIIcounterexample}

This counterexample also separates candidate selection from admissibility. A learner may prefer a formula for reasons of fit or cost. The compiler may install it as a pruning rule only under the corresponding proof obligation. Changing the rule's semantics, the available features, or the repair-success contract requires revisiting that obligation.

\subsection{The same search with fewer model-row reads}
The reference compiler first checks the immutable, unit-cost, binary-query and singleton-repair scope of Theorem~\ref{SII:thm:optimality-stop}. It then conducts an ordered depth-first feasibility search at the information lower depth $d_0=\lceil\log_2|B_0|\rceil$. The full-scan version reads every current row of a candidate query before rejecting an overfull child. The streaming version consults the certified prefix rule after each read. Both use the same query order, report order, remaining depth, and stopping rules. If a root has capacity-feasible children but a child cannot be completed using the actual library, the search backtracks to the next root.

A successful controller is separately checked and attains both lower bounds. When no controller exists at $d_0$, a requested horizon equal to $d_0$ permits a genuine finite infeasibility result. A longer requested horizon instead invokes the unchanged general compiler. Failure of the initial scope guard also invokes that compiler. Exhaustion of a model-read allowance returns incomplete, never infeasible merely because the allowance ran out.

\begin{SIItheorem}[Semantics-preserving streaming pruning]\label{SII:thm:streaming-pruning}
Under the checked binary immutable singleton-repair scope, fixed query and row orders, and identical search and fallback rules described above, replacing full root scans by the certified predicate~\eqref{SII:eq:learned-prune} returns the same first successful lower-depth controller, or the same exhaustive failure at that depth. With unlimited completion of any required common fallback, both compilers return the same protected frontier. The streaming version uses no more transition/report model-row reads, and the difference is exactly the number of unread suffix rows skipped at rejected roots. Scope checks and final witness verification are included identically in both read accounts.

Within a root scan, after a prior prefix has passed, the predicate can be evaluated by comparing only the newly incremented count with $C$. This lowering preserves every rejection decision.
\end{SIItheorem}
\begin{proof}
By Proposition~\ref{SII:prop:prefix-capacity}, every root rejected early would be rejected after its full scan. A root not rejected early has its complete rows read and produces the same two child supports in the same order. Induction on remaining depth therefore gives the same child searches, backtracking choices, first completed controller, and lower-depth failure. Nonempty binary splits strictly reduce support, and the depth decreases, so the finite search terminates. Any successful tree can omit redundant unary roots without losing case information in this immutable unit-cost scope. The search consequently covers all relevant lower-depth controllers. The checked lower bounds justify the same optimal frontier on success. Identical failure conditions invoke the same exact fallback on the same problem.

At every visited root, either the same rows are read or one suffix of rows that cannot affect the rejection decision is omitted. All descendant supports actually entered are identical. Summing omitted suffix lengths gives the stated read difference. This excludes timeouts, under which the methods can finish at different allowances but neither asserts false infeasibility.

For the lowering, both prior counts are at most $C$, since otherwise the scan would already have stopped. One report increments exactly one count. The unchanged count remains at most $C$, so the maximum exceeds $C$ exactly when the incremented count does. This implication requires the stated loop invariant.
\end{proof}

The read comparison isolates one alteration to one compiler. It is not a claim that every new search order beats the earlier balanced proposer, nor that fewer table reads imply shorter processor runtime. Interpreting a five-node formula at each prefix adds work. The supplied lowering replaces that interpretation by one integer comparison, but it too has a cost. The code records both versions and checks that they choose identical trees.

\begin{SIIcounterexample}[Local capacity without a completing query library]\label{SII:ex:local-capacity}
Four cases admit singleton repairs. The only available query divides them into $\{0,1\}$ and $\{2,3\}$ and leaves the state unchanged. At depth two, the root satisfies the $2/2$ capacity test. Repeating the only query cannot separate either pair, so no successful depth-two tree exists. Accepting the first capacity-feasible root as a finished certificate would be wrong. The streaming compiler searches its children, records the failure, and concludes infeasibility only after exhausting the actual library.
\end{SIIcounterexample}

\section{Pruning depends on what counts as a repair}\label{SII:sec:repair-cover}
The singleton scope is not cosmetic. A repair procedure need not determine the hidden case if one action succeeds on several cases. A change to the repair language can therefore invalidate the interpretation of raw case counts as required leaves, even when the experiment model is unchanged.

Let $B$ be a finite nonempty immutable-case support and let $G_r\subseteq B$ be each repair's success set restricted to it. Define $\kappa(B)$ to be the smallest number of repair-success sets whose union covers $B$. Set $\kappa(B)=\infty$ when no such cover exists. This number concerns the repaired objective, not the number of hidden states. Computing it can itself be expensive.

\begin{SIIproposition}[Repair-relative leaf requirements]\label{SII:prop:repair-cover}
For finite immutable cases with binary deterministic queries, any certain successful diagnostic tree of query depth at most $d$ has at least $\kappa(B)$ nonempty repair leaves. Consequently $\kappa(B)\leq2^d$ is necessary for success, and literal tree storage is at least $2\kappa(B)-1$. These conditions are not sufficient for success under a specified query library.
\end{SIIproposition}
\begin{proof}
Every initial case reaches one repair leaf. The set of cases reaching that leaf is contained in its repair's success set. The success sets of the attained leaves therefore cover $B$, so the number of leaves is at least $\kappa(B)$. A binary tree of depth $d$ has at most $2^d$ leaves. The edge-count argument from Theorem~\ref{SII:thm:optimality-stop} gives at least twice the number of nonempty leaves minus one literal records, with any unary nodes adding further storage. The preceding one-query example has $\kappa(B)=4$ and $d=2$, satisfying both numerical lower bounds while lacking a completing tree. Thus the necessary conditions are not sufficient.
\end{proof}

For singleton repairs, $\kappa(B)=|B|$, recovering the earlier argument. With a universal repair, $\kappa(B)=1$ even when many cases remain unknown. The executable streaming path retains the easier singleton guard rather than quietly substituting a hard cover computation whose cost would be uncharged. The broader proposition indicates how a different pruning contract could be developed, not a hidden extension of the present code.

\begin{SIIcounterexample}[A successful repair language can make identification unnecessary]\label{SII:ex:scope-repair-change}
Keep the four-case query table but add one unit-cost repair that succeeds in every case. The exact optimum is now one repair call and one stored record, with no queries. A compiler that carries over the old ``four cases require four distinct leaves'' assumption can discard this optimum. The reference guard rejects the changed singleton scope and the general compiler returns the zero-query solution. The old theorem remains true. Its premise no longer describes the new language.
\end{SIIcounterexample}

This is the preservation problem in recursive revision. A learned optimization acquires authority from a contract. Subsequent language growth can change that contract. Keeping the old formula while discarding its premise is not improvement. A retained shortcut should therefore carry its feature definition, semantic assumptions, recognized proof form, and fallback behavior together.

\section{What the synthesized compiler actually gains}\label{SII:sec:synthesized-cost}
The learner freezes~\eqref{SII:eq:learned-prune} after the small development domain. It is then used on the hidden-cell query family through 64 cases. The primary comparison is between the full-scan ordered search and the same ordered search with prefix pruning. It therefore isolates the effect of the generated rule. The earlier balanced wrapper is retained as a separate comparison, not silently treated as identical to this full-scan baseline.

\begin{center}\small
\begin{tabular}{@{}rrrrr@{}}\toprule
Cases & Full-scan reads & Streaming reads & Rows avoided & Added comparisons\\\midrule
2 & 12 & 12 & 0 & 2\\
4 & 66 & 65 & 1 & 25\\
8 & 308 & 277 & 31 & 141\\
16 & 1,384 & 1,129 & 255 & 681\\
32 & 6,192 & 4,641 & 1,551 & 3,137\\
64 & 27,744 & 19,457 & 8,287 & 14,209\\\bottomrule
\end{tabular}\end{center}

All listed compilers include the semantic scope guard and final witness check. The common input-shape validation is separately recorded and identical. For example, the 64-case query table has 70 queries and its shared shape check consumes 4,480 model rows for each method before the reported compilation allowance. The final controllers have the same literal trees in the primary comparison. Their execution and storage costs are unchanged: seven online units and 127 stored records in the 64-case example.

\begin{figure}[!ht]\centering
\includegraphics[width=.92\linewidth]{sources/monographs/self-improving-intelligence/figures/synthesized_read_cost.pdf}
\caption{Transition/report model-record reads for the same ordered compiler before and after installing the synthesized prefix rule. Scope and final witness checks are charged in both curves. Common input-shape validation is separately reported. The vertical axis is logarithmic. The plot is an exact finite cost comparison, not a timing ratio or evidence of a broad neural-learning gain.}\label{SII:fig:synthesized-cost}
\end{figure}

\subsection{A resource improvement is not a universal speedup}
In the interpreted version, each prefix evaluates five AST nodes. The checked lowering uses one integer comparison instead. Relative to a baseline that already performs the complete-root capacity test, a conservative extra-work account charges every new prefix comparison without crediting any removed end-of-scan check. At 64 cases this adds 14,209 comparisons while avoiding 8,287 model rows. At two cases no row is avoided and the new comparisons are pure overhead.

If a model-row read has declared price $w>0$ and a new comparison has price one, the additional per-use saving under this conservative two-unit account is
\[
\Delta(w)=w\,(R_{\mathrm{full}}-R_{\mathrm{stream}})-N_{\mathrm{compare}}.
\]
For the 64-case row it is positive when $w>14209/8287$. At $w=1$ it is negative. The price is a modeling choice, not a measured statement about Python or a physical machine. Interpreted AST visits, installed code, grammar construction, labels, and development search remain additional recorded coordinates. No weighted sum is presented as a universal cost of intelligence.

If development has total charged cost $D$ in a chosen common accounting unit and the same future use has saving $\Delta(w)>0$, $N$ repetitions repay that cost only when $N\Delta(w)>D$. This is the amortization argument already used earlier, applied to the learned compiler rule. It does not authorize ignoring development because deployment might eventually repeat often.

\subsection{Transport and negative controls}
The formula contains counts and a remaining-depth capacity rather than cell names or query identifiers. Renaming cases and flipping binary report meanings preserve the capacity predicate. A complete search uses the model rows to determine its actual branches, and its witness is checked anew. Reordering queries can change which equally optimal controller is encountered first. Reordering rows can change how early an impossible prefix becomes visible. Correctness is preserved under the certified scope, but the number of saved reads need not be invariant.

A four-case equality-only library provides another boundary. It has no controller at the information lower depth, so a larger requested horizon invokes the general compiler. The new method can spend preliminary work before that fallback. On a changed language with a universal repair, the semantic guard rejects the pruning scope and the general optimum is recovered. These are specified behaviors, not successes relabeled after observing the results.



\subsection{The precise recursive claim}
The first construction improved a diagnostic controller. The second selected a better compiler from named policies. The present construction synthesizes and installs an internal decision rule that changes how future controllers are discovered. It also records a representational obstruction within that compiler and a permitted feature refinement that removes it. This is a bounded inductive program-synthesis step, not neural training or a renamed cache of solved episodes.

The construction does not demonstrate that the grammar enumerator becomes better at inventing later grammars, that the learner discovers its own theorem checker, or that its task family becomes generally intelligent. The primitive operations, teacher, feature registry, scope proof, and search skeleton are substantial prior structure. The gain nevertheless occurs at the claimed location: future searches omit work because of a reusable expression composed from evidence, and the conditions under which that omission remains justified are checked rather than presumed.



\chapter{Improving how discovery rules are acquired}\label{SII:book:chapter-63}
\label{SII:part:meta-acquisition}
\section{The next retained object is a search representation}\label{SII:sec:meta-acquisition}
Chapter~\ref{SII:part:rule-learning} changes a decision inside the diagnostic compiler. Its counterexample-guided learner still traverses a fixed list of arithmetic candidates. A further change is possible without pretending that the learner invents its own computational foundations: it can acquire a reusable representation of that candidate space, then use the representation to reduce the work of synthesizing later rules.

The new object is neither a completed diagnostic tree nor the target expression $\max(a,b)>C$. It is a library of checked equations and an ordered map from grammar expressions to earlier equivalent representatives. Learning the library does not require the labels of a later synthesis task. Once the representation has been constructed and validated, subsequent tasks still begin with empty task-specific counterexample sets.

This direction has direct precedents. Equality saturation represents equivalent expressions and supports program optimization and synthesis~\cite{SII:egg}. Ruler automatically infers rewrite rules and uses equality saturation within that inference process~\cite{SII:ruler}. The present implementation is intentionally smaller. It uses a typed arithmetic language, finite proposals, a local derivation checker, and ordered enumeration rather than an e-graph. It does not claim a new rewrite-inference paradigm or superiority to those systems. Its purpose is to make one additional level of the improvement process executable and to connect the acquired search representation to the paper's unchanged semantic and resource contracts.

\subsection{Four objects and two meanings of success}
Let $\SIIGcal=(p_0,\ldots,p_{N-1})$ be a finite ordered grammar of pure, total expressions. A target task supplies contexts $z\in D$ and desired values $f(z)$. An ordinary learner uses counterexamples to find the first grammar expression that agrees with all of $D$. A separate equational learner observes expressions on a development probe set and proposes equations. A checker admits only equations supported by a derivation in the declared arithmetic theory.

There are therefore two distinct successes. A task candidate may agree with a complete finite target domain while remaining unjustified outside it. An equational certificate may establish agreement for every integer substitution, even though it says nothing about whether the resulting program solves a particular task. The latter can safely reduce the former search, but cannot manufacture its missing target evidence.

\section{Learning equations without confusing probes with proofs}\label{SII:sec:equation-learning}
The reference language uses exact integers, the total operations $+$, $\min$, $\max$, and the comparisons $<$, $\le$, $>$, $\ge$, and $=$. A finite schema language contains 135 terms built from variables $x,y$, small integer literals, Boolean literals, and at most one binary operation. The proposer evaluates those schemas on the nine contexts $(x,y)\in\{1,2,3\}^2$.

Two schemas with the same probe vector become a proposed equation, not an accepted rule. This produces 1,068 proposals. The supplied proof procedure proves 336 using integer arithmetic identities. The other 732 are not admitted. Failure of that limited procedure is not a theorem that an equation is false.

The derivation interface has six elementary rule kinds: commutativity where valid, addition of zero, idempotence of minimum and maximum, reversal of an order comparison, reflexive comparison, and exact evaluation of literal constants. It also checks replacement at an explicit subtree position. A proof of an equation contains two derivations ending in the same expression. A separate replayer checks every position, side condition, and endpoint. This is a project-local equational checker, not a proof-assistant kernel or a general arithmetic decision procedure.

\begin{SIIcounterexample}[A probe-equivalent rule that fails on a new input]\label{SII:ex:probe-equation}
On the development inputs $x\in\{1,2,3\}$, the terms $\max(x,1)$ and $x$ agree. At $x=0$ they differ. The proposed equation is not admitted by the proof checker. By contrast, $\max(x,0)=x$ holds on all nonnegative integers but not on all integers. A checker using the stronger integer-wide contract may fail to admit a rule that would be valid under a narrower natural-number contract. Rejection, counterexample, and scope restriction remain different outcomes.
\end{SIIcounterexample}

\subsection{What the learner contributes and what it is given}
The proposer generates the equations from evaluations rather than receiving the final equation list. The library is serialized before the later task targets are evaluated. The primitive syntax, probe schedule, arithmetic axioms, proof search, and normalization engine remain supplied. Their presence does not negate the acquired representation. It does prevent describing the construction as discovery of integer arithmetic or autonomous improvement of the equational learner itself.

The complete accepted library occupies 69,518 bytes in its stated compact JSON serialization. This is a serialization measurement, not an estimate of an optimal theory description. It contains redundant equations. A supplied full normalizer is retained as a stronger comparison, so the results do not depend on comparing only against deliberately repeated syntax.

\section{A certified quotient preserves the investigation's evidence}\label{SII:sec:ordered-quotient}
Write $\llbracket p\rrbracket(z)$ for an expression's value in context $z$. Let $n(p)$ be a deterministic representation produced by finitely many checked substitutions of valid equations. We only require
\[
 n(p)=n(q)\quad\Longrightarrow\quad
 \llbracket p\rrbracket(z)=\llbracket q\rrbracket(z)
 \quad\text{for every admitted }z.
\]
The reverse implication is unnecessary. Distinct normal forms may denote the same function. Missing such an equivalence wastes an optimization opportunity but does not remove a valid candidate.

\begin{SIItheorem}[Sound ordered quotienting from checked equations]\label{SII:thm:certified-quotient}
Assume the expression language is pure and total, every admitted schema equation is valid for all well-typed substitutions in the common semantic domain, and each normalization terminates after valid contextual replacements. For each $p_i$, keep its earliest grammar representative
\[
 \rho(i)=\min\{j:n(p_j)=n(p_i)\}.
\]
Then $\rho(i)\le i$, $p_i$ and $p_{\rho(i)}$ agree on every admitted context, and every pointwise specification realizable in $\SIIGcal$ is realizable in its retained representatives. No confluence or complete semantic-equivalence decision procedure is required.
\end{SIItheorem}
\begin{proof}
A schema identity remains true under a well-typed substitution. Purity and totality make replacement inside any expression context preserve its value. Induction over a finite derivation gives equality between an expression and its normalized result. Two expressions with the same result therefore have equal values by transitivity. The defining set for $\rho(i)$ contains $i$, giving existence and $\rho(i)\le i$. Replacing a satisfying expression by its representative preserves every pointwise requirement. The argument uses the soundness of each actual derivation, not uniqueness or completeness of normalization.
\end{proof}

In the reference implementation, a successful rewrite must decrease a size-and-structural-order measure. Rewriting introduces no free variables or literals outside the finite input-and-library vocabulary, and never increases size. Only finitely many terms lie below the initial size over that vocabulary, so there is no infinite decreasing rewrite sequence. The index keeps the original earliest AST, not the normalized AST. This preserves the reference learner's output convention and its program interpreter cost on every teacher proposal.

\begin{SIItheorem}[Preservation of ordered counterexample-guided acquisition]\label{SII:thm:cegis-quotient}
Consider a learner that scans a finite ordered grammar, retains every counterexample, tests candidates against those examples in a fixed order, and submits only consistent candidates to a deterministic teacher. The teacher returns the first disagreement in a fixed complete finite context order, or accepts. Under Theorem~\ref{SII:thm:certified-quotient}, scanning only the earliest representatives produces exactly the same submitted ASTs, counterexamples, and final accepted AST or grammar-exhaustion outcome. With the original AST interpreter retained, it performs no more retained-example evaluations or interpreter-node visits. Its teacher calls and teacher-evaluation work are unchanged. Quotient construction, storage, and dispatch costs are separate.
\end{SIItheorem}
\begin{proof}
Consider a removed expression $p_i$ and its earlier representative $p_j$. If the reference learner accepted $p_j$, it never reaches $p_i$. Otherwise a retained example refuted $p_j$, either before its teacher call or as a result of that call. That example is still retained when $p_i$ is reached. Semantic equality means the same example refutes $p_i$, so the reference learner never submits $p_i$ to the teacher. Deleting that unsuccessful consistency check leaves the sequence of proposals unchanged. Induction over the ordered scan gives the same candidate AST at every teacher call. The deterministic teacher then supplies the same first counterexample or acceptance. Removed candidates contribute only nonnegative evaluation work, while submitted expressions and their evaluations remain identical. If the finite grammar is exhausted, both learners have rejected the same possible behaviors.
\end{proof}

The theorem concerns an extensional target and a fixed candidate order. A teacher that accepts based on source identity, or a learner that forgets old counterexamples, does not satisfy it. Nor does equality of return values make two candidates equal in execution cost, memory use, or causal side effects. To preserve a different resource-constrained objective, one must keep appropriate nondominated representatives or include those observables in the equivalence contract. The reference deliberately preserves the first original AST rather than selecting an uncharged cheaper replacement.

\subsection{Cold reacquisition and additional finite task families}
The 3,960-expression grammar reduces to 1,362 retained representatives. Reacquiring the original capacity predicate starts with no task counterexamples and no installed copy of the answer. The original scan considers 190 candidates, while the quotient scan considers 76. Both submit the same four candidate ASTs, receive the same three counterexamples, and return the identical five-node expression. Retained-example checks fall from 282 to 112. Total interpreter-node visits in that task fall from 1,898 to 1,376. Teacher comparisons remain 210.

The frozen quotient is next used for nine finite targets involving sums, extrema, comparisons, and equality, followed by a parity target that is not expressible on its complete declared domain. These are changed target functions in the same arithmetic language, not a statistical sample of unrelated research domains. The first nine use all 256 contexts in $\{0,1,3,5\}^4$. The parity control uses all contexts in $\{0,1,2,3\}^4$, so more than one even value is present.

\begin{center}\small
\begin{tabular}{@{}lrr@{}}\toprule
Finite target & Original interpreter visits & Quotient interpreter visits\\\midrule
Sum exceeds capacity & 4,611 & 3,517\\
Minimum below threshold & 7,762 & 5,553\\
Maximum equals a value & 26,407 & 10,116\\
Sum equals a value & 29,357 & 12,311\\
Parity outside the grammar & 54,045 & 18,728\\\midrule
Complete ten-target batch & 196,384 & 85,162\\\bottomrule
\end{tabular}
\end{center}

Every submitted candidate and counterexample trace agrees exactly between the two learners. The parity target produces grammar exhaustion in both. A bounded grammar's failure is not a proof that the target function is uncomputable.

\begin{figure}[!ht]\centering
\includegraphics[width=.94\linewidth]{sources/monographs/self-improving-intelligence/figures/meta_acquisition_cost.pdf}
\caption{Repeated task-synthesis evaluations are reduced while the teacher trace remains fixed. Values are exact interpreter-node counts in the stated finite cases. Equation discovery, checking, quotient construction, and stored library costs are not included in these bars and are reported separately.}\label{SII:fig:meta-acquisition-cost}
\end{figure}

\subsection{The stronger normalizer and the cost of preparing the shortcut}
A supplied full arithmetic normalizer retains 1,352 representatives, ten fewer than the learned library. It therefore removes at least as much semantic redundancy under this implementation's normal-form relation. The learned library is not claimed to outperform that stronger prior. This control identifies the result as automatic acquisition of a useful search representation from the supplied reasoning substrate, not a new best-in-class synthesis algorithm.

The library-development ledger records 1,215 probe evaluations, 3,519 probe interpreter-node visits, 1,068 proposed equations, and 584 successful checked axiom steps. Proof search also attempts proposals it cannot establish. Applying the admitted library to the whole grammar costs 19,440 normalizer-node visits and 518,140 schema-match attempts, with 4,798 successful replacements. Grammar generation, proof search, bookkeeping, and serialization remain distinct charged phases. These event counts do not encode every Python instruction.

Let $C_{\rm prepare}$ be the complete preparation cost in a declared common resource model, including failed proposals and certificate checking, and let $d_t$ be the subsequent saving on task $t$ after including dispatch. Reuse is beneficial in that model exactly when
\[
 \sum_{t=1}^{T}d_t>C_{\rm prepare}.
\]
This is the earlier amortization identity applied to acquisition itself. The reference does not convert unlike counters into a universal wall-clock claim. Even one setup coordinate, 518,140 schema matches, exceeds the 522 interpreter-node visits saved in one cold capacity-rule run. A short workload or an expensive matcher can make the full modification undesirable.

\subsection{A cached proof representation has a scope}
The quotient records the grammar's contents, order, and declared integer semantics. Loading a cache requires checking those bindings and reconstructing its representative map from the admitted proofs. Adding syntax or changing order requires rebuilding the index. Extending the grammar can still reuse an equation library when every substituted operation is pure, total, and retains the relevant types and meanings.

\begin{SIIcounterexample}[The same operator spelling with changed semantics]\label{SII:ex:semantic-cache}
The identity $x+0=x$ is valid for ordinary addition. Suppose a revised interpreter uses the same printed symbol but executes $a\mathbin{\oplus}b=a+b+1$. Then $x\mathbin{\oplus}0\ne x$. A cache certified for the former interpreter cannot be imported merely because its expression strings still parse. A semantic identifier helps detect ordinary version mismatch, but an identifier falsely attached to a changed implementation is not itself proof of faithful execution.
\end{SIIcounterexample}

\section{A repair-relative certificate that can be updated cheaply}\label{SII:sec:weighted-pruning}
The singleton pruning result remains valid within its original scope. To move beyond that scope, the compiler must account for what a repair can accomplish, not merely count hidden cases. An exact minimum repair cover can be expensive. A feasible lower-bound witness is often enough to exclude an impossible branch.

Let $X$ be a finite immutable state set and let $S_r\subseteq X$ be the success set of repair $r$. Choose nonnegative rational weights $w_x$ satisfying
\begin{equation}\label{SII:eq:repair-weights}
 \sum_{x\in S_r}w_x\le1\quad\text{for every permitted repair }r.
\end{equation}
This is a feasible dual weight assignment for unweighted set cover. We use only its elementary weak-duality property, not an assumption that the best assignment has been found~\cite{SII:williamson}.

\begin{SIIproposition}[Repair-relative weight and prefix obstruction]\label{SII:prop:weighted-cover}
Suppose a total immutable deterministic binary diagnostic tree succeeds on nonempty support $B$ and has $L$ reachable leaves and maximum query depth $h$. Under~\eqref{SII:eq:repair-weights},
\[
 \sum_{x\in B}w_x\le L\le2^h.
\]
At a candidate query root with $h\ge1$ remaining query levels, if the cases already read into either outcome branch have weight greater than $2^{h-1}$, no assignment of its unread reports can make that branch repairable within the remaining depth. This implication is valid even when one repair succeeds in several states.
\end{SIIproposition}
\begin{proof}
Immutable deterministic execution partitions $B$ into its nonempty reachable leaf supports. Each leaf support is contained in the success set of its chosen repair. Nonnegativity and~\eqref{SII:eq:repair-weights} give weight at most one for every leaf, so summing gives the first inequality. A binary tree of depth at most $h$ has at most $2^h$ leaves, by induction on $h$. At a candidate root, each child has at most $h-1$ query levels. Unread cases can only add nonnegative weight to an outcome branch, so a prefix already exceeding its leaf capacity cannot be rescued by any completion.
\end{proof}

\begin{SIItheorem}[Streaming search with a checked repair-weight witness]\label{SII:thm:weighted-streaming}
Fix total immutable binary unit-cost experiments, unit-cost repairs with arbitrary fixed success sets, an ordered initial support, a repair order, an experiment order, a finite query horizon, and a weight vector satisfying~\eqref{SII:eq:repair-weights}. Consider depth-first search that first checks for a common successful repair and otherwise explores informative query roots in that order, rejecting child supports whose weight exceeds their remaining binary leaf capacity. Replacing each complete root scan by a scan that stops on the first overweight prefix preserves the exact first returned tree or finite-horizon infeasibility result. It never increases model-record reads during search. These statements do not establish a Pareto-optimal storage frontier or free weight arithmetic.
\end{SIItheorem}
\begin{proof}
Proposition~\ref{SII:prop:weighted-cover} shows that an overweight support has no successful tree at its depth. A root stopped on an overweight prefix would still be rejected after a complete scan because all further contributions are nonnegative. Every root not stopped early is partitioned identically, with the same row and outcome orders. Induction over depth and the ordered root scan therefore yields identical first successful children and identical returned trees. A query with one reachable outcome leaves the immutable support unchanged, so deleting such a query never destroys the existence of a depth-bounded successful tree. The search over the remaining roots and repairs is finite and complete for existence. Early stopping deletes only unread suffixes of roots the complete scan also rejects, proving the read comparison. The final witness checker and weight-scope checker remain required and are charged separately.
\end{proof}

\subsection{Constructing a feasible witness without solving the whole cover problem}
A simple supplied constructor processes states in order. It raises the current state's weight by the smallest remaining slack among repairs containing that state, and updates those repair loads. An uncovered state may receive a finite weight of one. Every step preserves nonnegativity and all repair constraints. The result is checked independently by resumming every repair's weight. It can be a weak certificate. In particular, the all-zero assignment is always feasible and excludes nothing.

For example, on three cases with success sets $\{1,2\}$, $\{2,3\}$, and $\{1,3\}$, the ordered greedy constructor returns total weight one. The feasible fractional assignment $w_1=w_2=w_3=1/2$ has total $3/2$ and proves that at least two leaves are required. Thus the supplied constructor does not optimize the bound. Both assignments are valid and neither creates a missing experiment.

The construction is not itself learned in this revision. It is a semantic extension showing how the compiler can continue when successful repair no longer means exact state identification. A learner may later propose a stronger weight assignment, but the same inequalities would still govern admission. The work of producing and checking those weights cannot be hidden inside one cheap predicate call.

\subsection{Paired repairs change the task, not just its notation}
Consider $2g$ possible hidden cases divided into $g$ pairs, where $g$ is a power of two. One repair succeeds on both members of each pair. The query language contains equality questions and bits of the pair index. The greedy weight constructor gives weight one to the first member of each pair and zero to the second, with total weight $g$. A depth $\log_2g$ tree is both necessary and sufficient, so the diagnostic-plus-repair work is $1+\log_2g$ and a full binary witness uses $2g-1$ records.

\begin{center}\small
\begin{tabular}{@{}rrrrr@{}}\toprule
Hidden cases & Repair groups & Full scan reads & Streaming reads & Work / records\\\midrule
4 & 2 & 20 & 18 & 2 / 3\\
8 & 4 & 152 & 120 & 3 / 7\\
16 & 8 & 864 & 608 & 4 / 15\\
32 & 16 & 4,416 & 2,864 & 5 / 31\\\bottomrule
\end{tabular}
\end{center}

These columns count search reads only. For the 32-case problem, each method additionally pays 1,152 scope-table reads and 256 witness-checking reads, giving 5,824 versus 4,272 total model reads. Weight construction, repair-success-set checks, rational additions, and comparisons have their own ledger. The streaming search makes 2,941 capacity comparisons rather than 1,073. Both methods return the identical literal tree.

This is not a generalization score obtained by making the task easier. The comparison is between two discovery procedures on the same paired-repair task. The difference between that task and the earlier singleton task explains why the earlier cardinality certificate would be invalid here.

\begin{figure}[!ht]\centering
\includegraphics[width=.93\linewidth]{sources/monographs/self-improving-intelligence/figures/repair_relative_classes.pdf}
\caption{Eight hidden cases need only four repair distinctions when each repair succeeds on a pair. The displayed weights certify a four-leaf lower bound. They do not claim that every feasible weight assignment is optimal or that every query library can realize the bound.}\label{SII:fig:repair-relative}
\end{figure}

\begin{SIIcounterexample}[A newly admitted repair invalidates the old lower-bound certificate]\label{SII:ex:weight-invalidation}
In an eight-case paired model, the original weights have total four. Add a repair that succeeds on every case. Its weight load is four, violating~\eqref{SII:eq:repair-weights}. The old certificate must be rejected. Recomputing a feasible assignment gives total one, and the compiler returns the universal repair with no query, work one, and one stored record. The success of that repair does not retrospectively validate the obsolete weight assignment.
\end{SIIcounterexample}

A capacity-feasible prefix still does not guarantee that available experiments can distinguish the required repair classes. The implementation retains a model with only a constant report: its cover weights fit a nominal depth bound, but no successful tree exists. A lower-bound certificate is a reason to reject some candidates, not a generator of a missing experiment.

\section{One chain from acquired equations to an improved investigator}\label{SII:sec:meta-chain}
The levels can now be connected without substituting one kind of evidence for another. The equational learner produces a checked library. The library prepares an ordered search representation. That representation reduces the repeated evaluation work of acquiring the earlier pruning rule and other finite targets. The rule's separate scope certificate then permits its use in the diagnostic compiler. Each stage retains the predecessor's semantic comparison.

\begin{SIIproposition}[Composition of acquisition and compiler preservation]\label{SII:prop:meta-chain}
Assume the ordered quotient satisfies Theorem~\ref{SII:thm:cegis-quotient}, the resulting task expression passes the unchanged compiler-scope checker, and the original and optimized compilers satisfy the common-domain premises of Theorem~\ref{SII:thm:streaming-pruning}. If all stages complete, quotient-guided rule acquisition followed by optimized compilation returns the same learned expression and the same diagnostic controller as ordinary rule acquisition followed by the reference ordered compilation. In a common additive resource model, writing $K$ for extra preparation cost, $A,A'$ for the two acquisition costs, and $D,D'$ for the two compilation costs, the composite saves work exactly when
\[
 (A-A')+(D-D')>K.
\]
The semantic conclusion does not imply the cost inequality, and a budget interruption does not imply infeasibility.
\end{SIIproposition}
\begin{proof}
Theorem~\ref{SII:thm:cegis-quotient} gives the same acquired AST and teacher transcript. Passing the unchanged scope checker places that expression in the stated compiler-preservation theorem, which gives the same returned controller. The reference total is $A+D$ and the modified total is $K+A'+D'$. Subtracting gives the displayed necessary and sufficient cost condition. Neither semantic theorem bounds $K$ or assigns exchange rates among distinct resource units. An interrupted finite search has not exhausted its candidates and therefore cannot use completeness to infer infeasibility.
\end{proof}

\subsection{The recursive achievement and its next boundary}
The retained change is now used to produce future discovery rules more economically in a specific evaluation coordinate. It is not merely a better diagnostic result. The learner's search representation has changed, while the pointwise target, oracle transcript, and final generated program remain fixed. This is a bounded meta-improvement with an explicit implementation and cost boundary.

It does not require that a useful self-improver invent every foundational component. It does require us to locate those foundations honestly. Here the schema proposer, primitive axioms, checking engine, and ordered-quotient construction are given. The acquired equation library and grammar quotient are generated. A supplied stronger normalizer can outperform the learned library, and a short workload can fail to repay its preparation cost. Extending the result to learned physical semantics, larger generated languages, uncertain arithmetic, or repeated improvements of the library learner itself requires additional work.

\subsection{The boundary of the first edition}
The finite construction now covers repair discovery, compiler revision, synthesized pruning, acquired search representations, and repair-relative scope extension. Chapter~\ref{SII:part:release-contract} consolidates their dependency and resource contracts. A first edition need not claim that general intelligence follows or that the outermost learner has become unboundedly self-improving. Those are distinct research questions. The completed scope is the explicit finite chain and the evidence required to interpret it.


\chapter{Keeping a revision's guarantees applicable}\label{SII:book:chapter-64}
\label{SII:part:release-contract}
\section{A certificate belongs to a changing realization}\label{SII:sec:certificate-life}
A proof can remain mathematically correct while ceasing to justify an installed shortcut. The equation library was proved for pure, total integer expressions. The ordered quotient was constructed for one grammar order. A repair-weight certificate was checked against particular success sets. A resource comparison assigned costs to specified events. None of those objects is licensed to retain its old interpretation after every component bearing the same name has changed.

This is the final integration problem for the first edition. The earlier sections provide distinct certificates for the outputs of different stages. A continuing system must also identify which of those certificates remain applicable after a revision. The appropriate response is not to repeat every discovery from scratch, nor to regard an old passing record as timeless. It is to retain the actual dependencies, reject unsupported reuse, and recheck the affected obligations before relying on them again.

\subsection{A declared dependency graph}
Let an artifact state be $a=(a_1,\ldots,a_d)$. Its components may include operation semantics, a grammar and its order, an equation library, experiment tables, repair-success sets, a returned controller, and a resource model. A finite directed acyclic graph has one vertex $v$ for each claimed guarantee. Its incoming vertices $\operatorname{pa}(v)$ identify guarantees used as premises. Its direct dependency set $D_v\subseteq\{1,\ldots,d\}$ identifies the artifact components used directly. Each vertex carries a witness $w_v$ and a local verification rule.

The dependency declaration must be \emph{semantically complete}: the meaning of the claim and the local rule is determined by the components in $D_v$ and the stated parent claims. Hidden environment assumptions cannot be left outside this declaration merely because the implementation does not know how to check them. They can instead be recorded as explicit external premises. Verifier soundness, the interpretation of primitive arithmetic, and the faithfulness of the supplied model to an external process are not established by naming them in a registry.

Suppose a revision replaces $a$ by $a'$. Mark a vertex directly stale whenever one of its declared component values differs. Let $\Delta$ be the union of those vertices and all their descendants in the dependency graph. Equivalently, in a topological order, mark $v$ whenever
\[
 a|_{D_v}\ne a'|_{D_v}
 \quad\text{or}\quad
 \operatorname{pa}(v)\cap\Delta\ne\varnothing.
\]
Here equality refers to the declared artifact value, not merely an unchanged version label. This procedure is conservative. A changed premise may leave the conclusion true, but the old evidence is not reused without a further argument.

\begin{SIItheorem}[Dependency-complete reuse across a revision]\label{SII:thm:certificate-reuse}
Assume a finite acyclic dependency graph, semantically complete direct dependencies, and sound local verification rules. Suppose the old witnesses establish every old claim from the declared external premises. After a revision, the old claims outside $\Delta$ remain established by their old witnesses. If every vertex in $\Delta$ is re-established in topological order using the revised artifact values and already established parent claims, all revised claims hold. If an affected vertex is not re-established, no dependent claim is justified by this reuse rule alone. These statements are conditional on the external premises applicable to the revised realization.
\end{SIItheorem}
\begin{proof}
A vertex outside $\Delta$ has unchanged direct inputs and no parent in $\Delta$. Inducting through a topological order, its parents remain established with the same meanings, so dependency completeness leaves its old verification argument applicable. For an affected vertex, assume its parents have already been established in the revised realization. Soundness of the newly checked local rule then establishes its revised claim. A second induction yields the result when every affected vertex is re-established. Without that step, the rule has an unsupported premise and supplies no conclusion along that dependency. This is a limitation of the reuse argument, not a proof that the dependent claim is false.
\end{proof}

The argument is elementary induction over declared dependencies. Its role here is to connect the exact preservation results to their continued use. It does not compute the complete dependency graph from arbitrary source code. A dependency cycle requires an additional justified rule, such as an independently proved invariant. The finite replayer in this edition rejects an unsupported cycle rather than interpreting mutual assertions as evidence.

\begin{SIIcounterexample}[An unchanged name hides a changed premise]\label{SII:ex:incomplete-dependency}
A stored claim asserts that two binary values $u$ and $v$ agree. Initially both are zero. A registry records only $u$ as a dependency and continues to validate that its stored value is zero. Change $v$ to one without changing $u$. The declared check still passes, but the claimed equality is false. Recording and replaying the incomplete registry cannot repair its omitted premise. Dependency completeness is a mathematical assumption, not a consequence of serialization or a checksum.
\end{SIIcounterexample}

\subsection{The scopes that actually have to be carried}
The implemented constructions require different bindings. An equation witness includes the supported arithmetic semantics and the local derivation rules. A quotient includes the equation library and the full ordered grammar. Exact candidate-trace preservation additionally requires the extensional teacher, fixed context order, and retained counterexamples. The compiler rule has its own integer-prefix and singleton-repair scope. Weighted pruning instead checks every repair-success set and the nonnegative weights. A returned tree must be checked against the actual experiment effects, observations, repair meanings, and horizon.

Cost and evidence claims add other dependencies. A reduction in model reads is not a price-independent reduction in total work. An audit comparing a frozen candidate cannot be carried unchanged to a new candidate merely because the score field has the same name. Statistical independence and commitment order may depend on events outside the program's observation boundary. The manifest must state whether those conditions are guaranteed by the controlled construction, supplied by a trusted interface, or not established.

The reference certificate manager therefore has several outcomes. A claim is \emph{valid} when its declared inputs match and its local evidence is replayed successfully. It is \emph{stale} when a direct input changed, \emph{blocked} when a required parent is unavailable or different, \emph{rejected} when its local record fails, and \emph{missing} when no witness was supplied. A timeout is not successful replay. Nor is a valid resource assessment necessarily favorable: the assessment may correctly report a negative saving.

The reference manager uses exact canonical snapshots for direct input comparisons and content identifiers for linking parent evidence. Those identifiers aid integrity checks. They do not establish who produced a witness or whether an external assumption is true. The manager and dependency registry are supplied components. This edition does not claim that the learner derived its own trusted verification base.

\subsection{A repair-language change through the complete record}
The first-edition demonstration checks the retained equation library, its ordered quotient, the capacity rule's supported form, a repair-weight witness, a returned weighted diagnostic tree, and a conditional resource assessment. The equation and diagnostic branches remain distinct. A valid capacity expression is not silently substituted for the weighted compiler's separate certificate.

In the 32-case paired-repair example, the accepted weight total is sixteen and the returned controller has work five and 31 records. Add a repair that succeeds in every case. The old weight certificate becomes stale, and direct rechecking confirms that it fails the new repair constraint. Its dependent controller and cost claims cannot rely on it through the old chain. The arithmetic equations and ordered grammar quotient remain applicable because their declared premises have not changed.

Rebuilding the weight witness gives total one. The compiler then returns a universal repair with work one and one record. The new controller is checked against the new repair language. The old cost comparison is deliberately not reissued: a claim about the work of producing the former controller does not become a claim about this different search simply because both controllers are valid. The appropriate repair of that claim is a new comparison under a stated workload.

Changing only the resource prices has the opposite pattern. The semantic certificates remain applicable, while the cost conclusion needs to be recomputed. Changing the grammar order invalidates the ordered quotient and its recorded acquisition lineage, not the independently proved arithmetic identities. These cases illustrate selective invalidation without equating it with a complete theory of arbitrary self-modification.

\section{One resource ledger for the complete improvement}\label{SII:sec:integrated-ledger}
The manuscript's numerical comparisons concern several different units. A model-record read, an interpreter-node visit, an arithmetic comparison, a stored tree record, and a serialized byte are not interchangeable. The purpose of a complete ledger is to preserve those distinctions until a specific execution or economic model justifies a conversion.

\begin{small}\begin{longtable}{@{}P{.25\textwidth}P{.34\textwidth}P{.33\textwidth}@{}}\toprule
Phase & Recorded object or event & What the count does not supply\\\midrule\endhead
Model and language acquisition & Supplied experiment tables, grammar, feature registry, primitive axioms & An executed calibration of an unknown physical process, or free discovery of those foundations.\\
Equation acquisition & 1,068 proposals, 336 admitted equations, 584 successful checked axiom steps & Proof completeness, a generally optimal equation set, or uncharged failed attempts.\\
Search representation & 3,960 grammar expressions, 1,362 representatives, 518,140 schema-match attempts & Free preparation, hardware-time units, or superiority to the 1,352-representative supplied normalizer.\\
Task-rule acquisition & 1,898 versus 1,376 interpreter-node visits for the capacity rule & All preparation costs, or a different teacher transcript. The retained four proposals agree.\\
Compiler execution & Query-model reads, weight arithmetic, capacity comparisons, scope checks, witness replay & A universal conversion of avoided reads into elapsed time.\\
Installed information & Equation-library serialization of 69,518 bytes, original programs, quotient records & Minimal description length, physical peak memory, or a byte-free interpreter.\\
Online repair & Experiment and repair calls, plus complete stored tree records & Only the visited path's storage, or the cost of discovering the controller.\\
Revision management & Bound snapshots, witness checks, parent comparisons, and replacement evidence & Free invalidation detection or a trusted explanation of undeclared external assumptions.\\\bottomrule
\end{longtable}\end{small}

The exact per-phase records remain in the source package. The table does not add unlike events into one total. In particular, retained storage is normally a capacity coordinate. Treating it as an additive expense requires a stated storage-duration or installation convention. Peak memory combines by its own rule, not by automatically summing all sizes ever observed.

\subsection{Which gains survive a change of resource prices?}
Fix an additive ledger with $d\ge1$ named event types. Let $K\in\mathbb R_{\ge0}^d$ be extra setup, let $s_i\in\mathbb R^d$ be the old-minus-new per-use event vector for workload type $i$, and let $n_i$ be its nonnegative integer multiplicity. The net event saving is
\begin{equation}\label{SII:eq:vector-saving}
 \Delta(n)=\sum_i n_i s_i-K.
\end{equation}
A price vector $w\in\mathbb R_{\ge0}^d$ specifies conversions into one declared additive resource measure. Its priced saving is $w\cdot\Delta(n)$. The prices are modeling assumptions, not learned hardware laws inferred from the event counts.

\begin{SIIproposition}[Resource dominance and declared price uncertainty]\label{SII:prop:resource-prices}
For a fixed ledger and workload, $w\cdot\Delta(n)\ge0$ for every nonnegative price vector $w$ if and only if $\Delta(n)$ is componentwise nonnegative. If admissible prices form the convex hull of a nonempty finite set $W_0$, then
\[
 \min_{w\in\operatorname{conv}(W_0)}w\cdot\Delta(n)
 =\min_{w\in W_0}w\cdot\Delta(n).
\]
Thus strict improvement throughout that price set holds exactly when every listed vertex gives a strictly positive margin. The claims concern the same workflow, workload, and complete additive ledger.
\end{SIIproposition}
\begin{proof}
Componentwise nonnegativity makes every nonnegative weighted sum nonnegative. If coordinate $j$ is negative, the $j$th coordinate vector is a nonnegative price vector giving a negative sum, proving the converse. Every point of the specified convex hull is a convex combination of its listed points. Linearity makes its priced saving the same combination of their savings, which cannot be smaller than their minimum. A listed minimizer attains that minimum, proving the equality and the strict criterion.
\end{proof}

This is ordinary linear accounting, not a new optimization principle. It prevents a system from asserting a price-independent gain after improving only one coordinate. The deterministic workload count must also be separated from an expectation over future tasks. Proposition~\ref{SII:prop:amortization} addresses the latter under its stated assumptions. Changing the task mixture can change both the amount and the sign of an improvement.

\begin{SIIcounterexample}[One valid controller, opposite resource conclusions]\label{SII:ex:price-change}
For the 32-case weighted-search comparison, the old-minus-new search-event vector is
\[
 (1552,\ 1552,\ -1868)
\]
in model-record reads, rational additions, and capacity comparisons. Scope checking and final witness checking are equal for both procedures and cancel in this particular difference. Prices $(1,1,1)$ give a saving of 1,236 units. Prices $(1,0,1)$ give a loss of 316 units, and prices $(0,0,1)$ give a loss of 1,868 units. The identical returned controller remains valid in every comparison. None of the price choices was measured as a universal property of a processor.
\end{SIIcounterexample}

The zero setup in that isolated execution comparison is not a statement that preparing the whole improvement was free. For an explicitly hypothetical additional setup charge of 2,000 comparison units and prices $(1,1,1)$, one use loses 764 units while two uses save 472. The purpose of this small calculation is to expose the dependence on the declared ledger. The actual equation-discovery, proof-search, quotient-construction, and revision-management records must be included when assessing the entire first-edition workflow.


\chapter{Evidence and recursive continuation}\label{SII:book:chapter-65}

\section{Evidence that survives candidate selection}\label{SII:sec:evidence}
A self-improvement procedure has an unusual incentive to find weaknesses in its own measurements. It searches repeatedly, retains favorable results, and changes the way later candidates are generated. A validation number can therefore become part of the training information even when it is never called a training label. This is a central issue in adaptive data analysis~\cite{SII:dwork}.

The correct distinction is between development evidence used to generate a candidate, selection evidence used to choose among candidates, and final evidence used to assess the selected procedure. A private test repeatedly consulted by an outer selector is still selection information for that outer process. Keeping it hidden from the inner solver does not make it independent of the entire lineage.

A practical record should identify what was frozen before each assessment, which outputs were exposed, and which claim the assessment supports. When a procedure changes a checker, both the checker change and the candidate it would approve cannot receive their only justification from the new checker's own approval. That restriction does not forbid improving evaluators. It prevents the target of validation from silently changing the validation standard in the same step.

\begin{SIIproposition}[Conditional error budgets across an adaptive lineage]\label{SII:prop:lineage}
Let $\mathcal H_{t-1}$ be the information available before attempt $t$. Suppose the event $A_t$ that attempt $t$ accepts a false improvement claim satisfies
\[
\Pr(A_t\mid\mathcal H_{t-1})\leq\delta_t
\quad\text{almost surely},
\]
for deterministic $\delta_t\geq0$ with $\sum_{t\geq1}\delta_t\leq\delta$. Then
\[
\Pr\big(\exists t\geq1:A_t\big)\leq\delta.
\]
The candidate-generating process and the time at which attempts stop may depend on the earlier history.
\end{SIIproposition}
\begin{proof}
Taking expectations gives $\Pr(A_t)\leq\delta_t$. The union bound gives the result for any finite number of attempts. The infinite statement follows from continuity of probability for the increasing finite unions. A history-dependent stop can only reduce the set of attempts whose acceptance events occur.
\end{proof}

No independence between attempts is needed. The conditional premise is stronger than a marginal guarantee for a candidate considered before adaptive selection. It must be justified by the actual sampling or verification mechanism. A budget such as $\delta_t=\delta/[t(t+1)]$ is summable, but writing that schedule does not make an invalid test conditionally valid.

\subsection{One sufficient fresh-audit construction}
Freeze a candidate and its comparator using the prior history. Draw $n_t$ fresh evaluation cases conditionally independently from the declared distribution and record paired score differences $Y_i\in[-1,1]$. Let $\Delta_t$ be their conditional mean and $\widehat\Delta_t$ their average. The bounded-variable concentration inequality gives
\[
\Pr\left(\widehat\Delta_t-\Delta_t\geq r\mid\mathcal H_{t-1}\right)
\leq \exp(-n_tr^2/2).
\]
Consequently, accepting a claim $\Delta_t\geq\gamma$ only when
\begin{equation}\label{SII:eq:audit}
\widehat\Delta_t-\sqrt{2\log(1/\delta_t)/n_t}\geq\gamma
\end{equation}
meets the preceding conditional error requirement for that claim. This is a fixed-sample sufficient construction. The candidate and $n_t$ may depend on past data, but $n_t$ is fixed before this audit is examined. Inspecting the same sample repeatedly and stopping at a favorable value requires a time-uniform method or a separately justified spending rule. Time-uniform concentration is an established subject~\cite{SII:howard}.

Formula~\eqref{SII:eq:audit} can be conservative and expensive. It is included to make the premise operational, not as a claim that this test is optimal. If evaluator randomness, shared tasks, or dependence between runs changes the sampling unit, the analysis must change with it. Thousands of messages from one trained lineage do not automatically become thousands of independent improvement trials.

\begin{SIIcounterexample}[Reused evidence can reward memorization]\label{SII:ex:reuse}
A task has $N$ equally likely names and an externally fixed binary answer for each. After the answers on a development subset become available, a candidate stores them and gives the wrong answer on every other name. It is perfect on that subset and has true accuracy equal only to the subset's fraction of the task universe. An outer loop that keeps selecting by the same subset can improve its reported score without improving beyond the information it has exposed. Renaming the subset a holdout does not restore independence.
\end{SIIcounterexample}

\subsection{Proof obligations and performance evidence}
A semantic proof can establish that a proposed operation preserves a specified relation. A program checker can establish that a finite tree respects the supplied domains and budgets. A fresh performance comparison can support usefulness on a declared task law. None of these establishes the truth of every environmental premise. Their evidence types should remain separate in an improvement receipt.

The constructive tests shipped with this paper are deterministic finite enumerations. They are not samples from an AGI population, and they do not use~\eqref{SII:eq:audit} to manufacture confidence intervals. Their role is to check the implementations of the stated finite arguments.

\section{Sequential evidence for an improved discovery procedure}\label{SII:sec:anytime}
The deterministic compiler comparison above can be checked on its whole declared family. A learned discovery procedure evaluated on fresh stochastic tasks would need another kind of evidence. It may also be expensive enough that the auditor would like to inspect results sequentially and stop once the evidence is adequate. The fixed-sample rule in Section~\ref{SII:sec:evidence} does not automatically justify that stopping policy.

We give a small time-uniform alternative whose finite arithmetic is directly inspectable. Before an audit begins, freeze the candidate procedure, its comparator, their installed information, the task law, the common resource cap, and a paired score $Y_t\in[-1,1]$. Timeouts and invalid outcomes must have a stipulated score rather than disappear from the sample. Fix a proposed improvement margin $\gamma\in[0,1)$.

The null condition is
\begin{equation}\label{SII:eq:anytime-null}
\SIIE[Y_t\mid\mathcal H_{t-1}]\leq\gamma
\quad\text{for every audit step }t,
\end{equation}
where $\mathcal H_{t-1}$ includes all information available before the next score. Conditionally independent cases from a fixed task law with mean at most $\gamma$ satisfy this condition. An auditor cannot change the task distribution toward favorable cases and retain that interpretation without a new argument.

Choose a nonnegative bet $\lambda_t$, measurable from $\mathcal H_{t-1}$, with $\lambda_t\leq1/(1+\gamma)$. Define
\begin{equation}\label{SII:eq:anytime-capital}
E_0=1,\qquad E_t=E_{t-1}\bigl(1+\lambda_t(Y_t-\gamma)\bigr).
\end{equation}
The variable $E_t$ is an evidence process, not a posterior probability that the candidate is better. Predictability means choosing the bet before observing the score to which it is applied.

\begin{SIItheorem}[An anytime-valid paired improvement audit]\label{SII:thm:anytime}
Under~\eqref{SII:eq:anytime-null} and the stated predictable bet constraints, $E_t$ is a nonnegative supermartingale. For every $\delta\in(0,1)$,
\[
\Pr\left(\sup_{t\geq0}E_t\geq1/\delta\right)\leq\delta.
\]
Consequently an audit may reject the stated null at its first crossing of $1/\delta$, while choosing its stopping time from the observed audit history. The result does not require a predetermined number of cases or independence beyond the conditional-mean premise.
\end{SIItheorem}
\begin{proof}
Since $Y_t-\gamma\geq-(1+\gamma)$, every factor in~\eqref{SII:eq:anytime-capital} is nonnegative. Conditional expectation and predictability give
\[
\SIIE[E_t\mid\mathcal H_{t-1}]
=E_{t-1}\bigl(1+\lambda_t(\SIIE[Y_t\mid\mathcal H_{t-1}]-\gamma)\bigr)
\leq E_{t-1}.
\]
Let $\tau$ be the first threshold crossing. For a finite $N$, the stopped process at $\tau\wedge N$ still has expected value at most one, by induction on the stopped conditional expectations. On $\{\tau\leq N\}$ its value is at least $1/\delta$, while it is nonnegative otherwise. Therefore $\Pr(\tau\leq N)\leq\delta$. Letting $N$ increase proves the result for any crossing time.
\end{proof}

This is the elementary nonnegative-supermartingale method underlying time-uniform evidence, not a new general concentration theorem~\cite{SII:howard}. The contribution here is its explicit paired comparison and its placement within the recursive resource and information contract. In particular, an audit observation is a complete comparable attempt by each discovery procedure, not an arbitrary message emitted within one attempt.

\subsection{Why seeing the outcome before choosing the bet fails}
\begin{SIIcounterexample}[Hindsight betting creates false evidence]\label{SII:ex:hindsight-bet}
Let $Y_t$ be independent fair signs and $\gamma=0$. Choose $\lambda_t=1$ when the current sign is positive and zero otherwise. This invalid rule doubles $E_t$ on a positive score and leaves it unchanged on a negative one. Within four scores it crosses the threshold four whenever at least two scores are positive, an event of probability $11/16$. This exceeds the nominal bound $1/4$. Each individual bet lies in the numerical range $[0,1]$. What fails is predictability.
\end{SIIcounterexample}

A data-dependent stopping rule is allowed by the theorem. A data-dependent choice that looks at the \emph{current} outcome before staking on that outcome is not. The software can validate the numerical range of a bet, but it cannot infer the provenance of an externally supplied bet from that number alone. A usable audit receipt must record the commitment order and the sampling contract.

The same issue arises when a candidate is edited after seeing an audit result. The resulting process no longer evaluates one frozen procedure under its original interpretation. It can start a new comparison, debit a new conditional error budget, and retain the earlier result as development evidence. It cannot silently reuse the old evidence process as though the candidate had not changed.

\subsection{Evidence across recursive generations}
For attempt $j$, condition on the complete earlier lineage, freeze its candidate, and initialize a fresh valid audit at $E_0^{(j)}=1$ with error allowance $\delta_j$. If $\sum_j\delta_j\leq\delta$, Proposition~\ref{SII:prop:lineage} bounds the probability of ever accepting a false margin claim by $\delta$. Each audit can stop sequentially, and the next candidate may depend on completed earlier audits. The premise remains conditional validity for each actual attempt.

Exact certificate checks and sequential evidence still have different jobs. A table-defined compiler can be checked for admissibility and optimality on its formal input without making a claim about a population of tasks. A stochastic audit can support a productivity comparison on that population without proving every generated program correct. The recursive system needs both whenever its claim combines semantic validity and statistical usefulness.

The new reference suite enumerates finite rational null paths and confirms the crossing bound in those cases. It also checks the hindsight counterexample. No live learner was evaluated with this test, and the deterministic compiler counts are not retroactively presented as independent random samples.

\section{The complete cost of becoming a better improver}\label{SII:sec:cost}
A procedure can save work on each future task and still be a poor investment when its discovery and validation are expensive. Conversely, a costly method can be worthwhile when it is reused sufficiently often. Both observations are ordinary resource accounting. They become especially important when a claimed improvement shifts work from inference into training, into a translator, or into an external oracle.

Use a resource vector rather than an unexplained token total:
\[
C=(C_{\rm model},C_{\rm search},C_{\rm validate},C_{\rm install},
C_{\rm execute},C_{\rm store}).
\]
These coordinates need not share units. Model construction may involve environment observations. Search and validation may involve processor work. Installation may require conversion or compilation. Execution work concerns one deployment. Storage concerns the retained library and complete controller. A scalar comparison needs declared conversion weights or a common physical budget, not an arbitrary sum of unlike quantities.

For the hidden-cell instrument, a compiled tree has fifteen records even though a successful depth-three run visits only four. The bit-query API is installed for both procedures in the matched comparison. Its table representation, if used, occupies $n(n+\log_2n)$ Boolean entries for the chosen power-of-two family. The reported three-versus-seven diagnostic comparison does not imply that inspecting a large subset costs one machine instruction in every implementation. It uses one sensor call as its stipulated unit.

\begin{SIIproposition}[Amortization of a reusable process improvement]\label{SII:prop:amortization}
In a fixed additive resource unit, suppose installing a new procedure has incremental cost $C\geq0$ relative to the old one, and its expected per-use saving on an unchanged task distribution is $d>0$. If there are $N$ uses and no additional costs or changes in the task law, its expected cumulative saving is $Nd-C$. It is strictly resource-improving exactly when $N>C/d$.
\end{SIIproposition}
\begin{proof}
Linearity of expectation gives expected total use savings $Nd$, without requiring independence between uses when each has the stated expectation. Subtract the fixed incremental cost and solve the strict inequality.
\end{proof}

The formula is not a prediction of future demand. Distribution drift, maintenance, growing context, retrieval overhead, or a changed interface can alter $d$ or add costs. The value of the repair library depends on whether later problems actually admit its operations.

\subsection{Larger search spaces can lower productivity}
\begin{SIIcounterexample}[More expressibility with worse fixed-budget discovery]\label{SII:ex:dilution}
A search language has $N$ candidates, exactly one useful. A procedure samples uniformly with replacement $B\geq1$ times and recognizes usefulness perfectly after each sample. Its probability of finding the useful candidate is $1-(1-1/N)^B$. Adding an additional useless candidate preserves every old expression and strictly decreases this probability. The best attainable repair is unchanged, while the implemented discovery procedure becomes worse under the same sampling budget.
\end{SIIcounterexample}

This does not argue against grammar growth. It identifies what the learner must improve alongside growth: proposal probabilities, indexing, modular reuse, or informative pruning. The search space and the search distribution are different parts of the agent.

An archive has a related ambiguity. The best true value among retained artifacts cannot decrease when an old artifact remains available under the same score. But a deployed selector can perform worse as retrieval becomes harder. The monotonicity of a hypothetical maximum is not the monotonicity of actual use. Resource-limited archive management belongs inside the recursive system rather than outside its reported budget.

\section{A finite certificate for repeated improvement}\label{SII:sec:recursion}
A valid local change can destroy the ability to make another valid change. A new representation may discard the evidence needed to verify future proposals. A controller may consume the remaining budget while preserving the current task score. A shortcut can bypass an interface on which later transports depend. Recursive continuation therefore needs its own object.

Let $Z$ be a finite set of complete lineage configurations. A configuration includes the currently selected solver and improver, retained libraries and evidence, the active comparison contract, and relevant resource state. The following theorem assumes this configuration is observed. Unknown configurations require belief-state modeling or a different synthesis result. No unrestricted introspection is inferred.

Each $z$ has a finite set of authorized attempts $A(z)$. An attempt has a nonempty set $\SIIPost(z,a)$ of possible successor configurations. This can represent all outcomes of a stochastic proposal or an abstract uncertainty set. A safe set $\SIISafe\subseteq Z$ encodes the declared semantic, evidence, authority, and resource conditions. A pause or terminal state may have a zero-progress self-loop. Actual expensive attempts must debit resources in $Z$, or the model must explicitly state their replenishment.

Define
\[
K_0=\SIISafe,\qquad
K_{n+1}=\{z\in K_n:\exists a\in A(z),\ \SIIPost(z,a)\subseteq K_n\}.
\]
The descending finite sequence reaches a fixed point $\mathcal K$. It is the set from which some policy keeps every successor safe indefinitely. This can include doing nothing forever. Safety alone is not recursive improvement.

For each transition, a progress mark $p(z,a,z')\in\{0,1\}$ is supplied by a checked criterion. A mark of one means the specified process-improvement claim has been established under the same comparison. It is not assigned merely because a file changed. Define
\begin{align}\label{SII:eq:progress}
W_{0,0}&=\mathcal K, & W_{0,k}&=\varnothing\quad(k>0),\\
W_{h,k}&=\{z\in\mathcal K:\exists a\in A(z),\ \SIIPost(z,a)\subseteq\mathcal K,\nonumber\\[-2pt]
&\hspace{34mm} z'\in W_{h-1,\max\{0,k-p(z,a,z')\}}
\quad\forall z'\in\SIIPost(z,a)\}.\nonumber
\end{align}
Here $h$ is the number of remaining attempts and $k$ the number of still-required progress events. A pause action lets a policy that finishes early use the remaining horizon without changing the result.

\begin{SIItheorem}[Common-policy recursive progress]\label{SII:thm:progress}
A configuration belongs to $W_{h,k}$ if and only if there exists one authorized policy that keeps all outcomes in $\mathcal K$ and guarantees at least $k$ checked progress events within $h$ attempts. A policy depending on the current configuration and the remaining pair $(h,k)$ suffices. The finite recurrence gives a witness when feasible and a complete negative answer within the stated model when infeasible.
\end{SIItheorem}
\begin{proof}
At horizon zero, zero remaining progress is possible exactly in $\mathcal K$, and positive remaining progress is impossible. Assume the claim for $h-1$. A policy at $z$ must choose one initial action before its outcome is known. Every possible successor must remain in $\mathcal K$ and must admit a continuation meeting the remaining target $\max\{0,k-p\}$. By induction these conditions are exactly the memberships in~\eqref{SII:eq:progress}. Conversely, choose a witnessing first action and a witnessing continuation for each observed successor. The resulting policy works on every path. Induction proves necessity, sufficiency, and the stated memory requirement. Finiteness makes the construction terminating.
\end{proof}

The use of one first action is important. The weaker claim that every possible outcome has some favorable action, selected after secretly learning that outcome, is not an implementable policy. Likewise, independently feasible semantic, resource, and evaluation certificates must refer to the same attempt and successor relation.

The theorem is complete only for its finite catalogue and its information boundary. It does not create a promising next candidate or identify the true productivity of an arbitrary program. Its progress labels may come from exact finite evaluation or from a separate evidence layer. With probabilistic labels, the conclusion about actual progress holds on the event that the relevant certificates are valid. Proposition~\ref{SII:prop:lineage} can bound failure of that event under its own assumptions.

\subsection{Safe rejection can prevent a worst-case progress guarantee}
\begin{SIIcounterexample}[Repeated opportunity is not bounded sure progress]\label{SII:ex:retry}
At a safe configuration, an attempt either produces a genuine improvement or rejects the candidate and returns to the same configuration. Both outcomes are possible at every attempt. A safe pause is also available. For every finite $h$, there is a path with $h$ consecutive rejections, so no policy guarantees one improvement within $h$ attempts. With an independent positive success probability, eventual success can nevertheless occur with probability one. The worst-case finite theorem and that probabilistic statement answer different questions.
\end{SIIcounterexample}

\subsection{Why local improvement does not imply recursive viability}
A second finite example has a proposal that raises the current score but deletes the only retained witness needed to validate any subsequent proposal. Its destination lies outside $\SIISafe$ when witness availability is protected. If witness availability is omitted, it may remain safe but have no productive continuation. The result depends on the contract, not the flattering name assigned to the edit.

This is a compact version of the closure perspective in SRIC~\cite{SII:sric}. The present contribution is to connect that perspective to a concrete diagnostic tree, a resource witness, and a specified productivity comparison rather than treating closure as a free consequence of successful self-modification.

\section{Acceleration is an additional hypothesis}\label{SII:sec:acceleration}
Recursive improvement can be meaningful without being indefinitely accelerating. Its measured pace depends on useful gains, the cost of finding and checking them, available opportunities, and the chosen scale of evaluation. A normalized task score has a ceiling. A changing suite can continually introduce new headroom, but its successive normalized scores are not automatically comparable.

\begin{SIIproposition}[A fixed-margin ceiling on a fixed score]\label{SII:prop:ceiling}
Suppose a lineage has a fixed measure $R_t\in[a,b]$, every accepted progress event increases it by at least $\gamma>0$, and all other transitions are nondecreasing in that same measure. Then the number of progress events is at most
\[
\left\lfloor\frac{b-a}{\gamma}\right\rfloor.
\]
An infinite sequence of positive increments is compatible with bounded total improvement when no common positive lower bound is imposed.
\end{SIIproposition}
\begin{proof}
After $k$ progress events, telescoping and nondecrease give $R_t\geq R_0+k\gamma\geq a+k\gamma$. Since $R_t\leq b$, the bound follows. For the second assertion, start at zero and use increments $2^{-t}$ for $t\geq1$. Every increment is positive and their partial sums remain below one.
\end{proof}

If nonprogress transitions can decrease the measure, the first bound need not hold: repeatedly falling and recovering can generate arbitrarily many positive increments. Those recoveries are not unlimited cumulative progress. A moving evaluator creates a different problem. It may be appropriate, but it requires bridges to preserve the meaning of the comparison.

A simple accounting model can express the additional conditions for acceleration. Let $g_t$ be an expected useful gain and $c_t>0$ the complete cost of the corresponding discovery and acceptance cycle. Improvement productivity per unit resource is $g_t/c_t$. Raising $g_t$ alone does not raise that ratio when $c_t$ rises faster. A larger reachable repair set is similarly insufficient when proposal quality falls or validation becomes the limiting cost.

SRIC's bottleneck perspective is useful here: improving one stage can expose another stage as the dominant limitation~\cite{SII:sric}. The next bottleneck might be information acquisition rather than proposal generation, or evidence quality rather than repair expressibility. A theory of acceleration would need to specify how those bottlenecks move and which resources support their resolution. A claim of universal acceleration does not follow from the existence of one productive cycle.



\chapter{A bounded route toward general intelligence}\label{SII:book:chapter-66}
\label{SII:part:architecture}
\section{A proposed learned architecture}\label{SII:sec:architecture}
The mathematical contracts do not specify a complete learning algorithm. They do identify the components that a credible algorithm must make learnable and the evidence required when those components change. The intended architecture has one active solver, a recurrent working state, a persistent library of typed operations, an acquisition policy, a proposal mechanism, and an admission procedure. These are roles in one realized computation, not a requirement for a separate large model for every role.

\subsection{What would learn}
Let $r_\phi$ encode the retained state, $G_{\theta,a}$ describe proposed operations, $\pi_\omega$ choose observations or experiments, and $p_\psi$ propose candidate programs or parameter updates. A development episode can modify $\phi,\theta,\omega,\psi$ through a declared learning rule. A persistent program library can also change by adding, replacing, or removing checked operators. Temporary deliberation without a retained modification is not called a new learned generation here.

The learning signals need separate purposes. Prediction losses train anticipation of future outputs. Intervention-consistency losses test whether the proposed operations match their intended effects. Acquisition rewards favor useful reductions of repair uncertainty after charging experiment cost. Proposal rewards evaluate the actual selected descendant rather than fluency of its rationale. Preservation checks protect specified older capabilities or interfaces. There is no universal set of weights that turns these distinct obligations into one sufficient scalar objective.

A latent computation may suggest that a particular variable matters. A typed checkpoint should then declare the variable's domain, the operations it supports, the evidence used to infer it, and the continuations on which it was checked. Its translator must not read the answer that the checkpoint claims to predict. Its complexity is charged. OSL is a proposed semantic consumer for such checkpoints, while Neuralese concerns the learned latent producer~\cite{SII:osl,SII:neuralese}. This paper does not implement their complete learned bridge.

\begin{SIIproposition}[Local operation agreement preserves finite continuation]\label{SII:prop:continuation}
Let $D$ be invariant under every source operation $F_a$, let $q:D\to Z$, and suppose
\[
qF_a=G_aq\quad\text{on }D
\]
for every admitted $a$. If the required output satisfies $y=\bar yq$ on $D$, then for every finite operation word $w$,
\[
qF_w=G_wq,\qquad yF_w=\bar yG_wq.
\]
\end{SIIproposition}
\begin{proof}
The empty word gives the identity. Assume the relation holds for a prefix. Domain invariance allows application of the next local equation, giving $qF_aF_w=G_aqF_w=G_aG_wq$. Induction establishes the first identity and composing with $\bar y$ gives the second.
\end{proof}

The proposition is a specification for meaningful reuse, not a method for learning $q$ or $G_a$. Invariance matters because an unchecked operation can leave the domain where the equations hold. Output preservation matters because commuting hidden states need not retain the task variable. Actual causal use matters because a probe can decode an appropriate state that the solver does not use.

\subsection{What a complete attempt would do}
A bounded iteration would begin with an observed task failure and an explicit set of plausible causes. The active acquisition policy chooses an admitted experiment. Its result updates those causes and the current state. A proposer supplies a candidate repair or a new constructor when the old language is insufficient. A separate execution checks the candidate's declared interface, resource usage, and behavior. A selected candidate is assessed on evidence not used to construct it. If admitted, its actual changes and scope are stored, including the procedure that found it.

A later outer attempt may change the acquisition policy, proposer, or selection code. It must then be assessed as an improver from comparable starting conditions, not only as the latest solver. An improvement may be retained provisionally for a particular domain rather than promoted globally. Failed hypotheses remain useful evidence when their assumptions and reasons for rejection are preserved.

Candidate evaluation can use the same kind of mathematical tools as the proposer without sharing the candidate's unchecked assumptions. Independence here means a specified information and authority separation, not a mystical absence of common software or training ancestry. External audit data and the authority to redefine a protected score are not supplied to a candidate as editable inputs for the same transaction.

\subsection{A learned policy is not an unconstrained interpreter}
A sufficiently powerful interpreter can hide the entire solution behind a small label. The remedy is not to forbid interpreters by definition. It is to state their installed information, runtime, storage, and access. This is the resource lesson of \emph{Machine Understanding}~\cite{SII:understanding}. The diagnostic tree is intentionally explicit so that this accounting can be performed exactly in the worked case.

Nor should neural implementation be required simply for appearance. A library of executable operators and a learned experiment selector can be a meaningful adaptive system. Conversely, updating a large weight tensor does not automatically establish a useful representation, a causal mechanism, or a transferable improvement procedure. The architecture is evaluated by its realized operations and the origin of their improvement.

\section{The compiled diagnostic improvement}\label{SII:sec:worked}
We can now state the exact mathematical achievement of the running example. It is deliberately smaller than the full learned architecture, but complete enough to expose the relevant distinctions.

For $n=2^k$ cells, assume a single defective cell, binary reports, singleton queries, and $k$ bit-coordinate queries. All queries leave the cell unchanged and cost one diagnostic call. A correct final repair costs one additional call. Distinct cells require distinct terminal repairs. There are no false repair commits. A literal tree record contains one experiment or repair identifier plus its child references.

\begin{SIItheorem}[Optimal diagnosis in the hidden-cell family]\label{SII:thm:cells}
Every certain diagnosis needs at least $k$ binary queries in the worst case. The bit-query procedure attains this bound and uses $2n-1$ records in a full binary tree. Any procedure using singleton queries alone needs at least $n-1$ queries in the worst case, and equality-first testing attains that bound with the same number of tree records. Thus their worst total execution costs, including repair, are $k+1$ and $n$.
\end{SIItheorem}
\begin{proof}
A binary tree of depth less than $k$ has fewer than $2^k=n$ leaves. Every successful leaf can identify at most one cell because its repair must be correct for every remaining case. Therefore depth at least $k$ is required. The $k$ coordinate bits identify every cell and give a complete binary tree with $n$ leaves and $n-1$ internal nodes. For singleton tests, an adversary can answer no on each tested cell while at least two untested cells remain. Fewer than $n-1$ distinct singleton queries therefore cannot guarantee identification. Equality-first testing uses one leaf for each positive answer and one final leaf for the last remaining cell, again totaling $2n-1$ records. Adding the terminal repair gives the claimed work counts.
\end{proof}

\begin{figure}[!ht]\centering
\includegraphics[width=.92\linewidth]{sources/monographs/self-improving-intelligence/figures/diagnostic_cost.pdf}
\caption{Exact total execution costs for the power-of-two cell family. Singleton testing requires $n-1$ diagnostic queries, while the admitted bit queries require $\log_2 n$. A final repair adds one call to each plotted curve. Both full trees have $2n-1$ stored records. The curves omit separately charged model acquisition, compilation, and adapters.}\label{SII:fig:cost}
\end{figure}

\subsection{What is actually compiled}
The reference compiler receives an observation table and the set of successful terminal repairs. It computes the frontier in Theorem~\ref{SII:thm:synthesis} and returns an optimal tree under a chosen work/storage allowance. At $n=8$ it returns a four-call, fifteen-record witness. A separately coded interpreter checks every initial cell and every branch. The reference comparison constructs the equality-first baseline independently, rather than comparing the compiler to a weakened version of its own output.

The selected tree is an executable diagnostic policy for future episodes. It is not a saved answer for the calibration case. In a new episode, the hidden cell is unavailable to the policy. Only reports from its queries determine the branch and eventual repair. A cell permutation supplies a controlled foreign interface when the query sets, terminal repairs, and adapter are transported together. Theorem~\ref{SII:thm:transport} explains that transfer. The tests do not claim that the adapter was learned without correspondence data.

\subsection{The measured object is a procedure}
Under the earlier four-call cap, the equality-first improver repairs $3/8$ of the uniform eight-cell cases and otherwise abstains. The compiled improver repairs all cases. The repaired solvers themselves are simple: each now knows the correct cell-specific operation. What improved is the procedure that acquires that case-specific correction. The library and sensor language are unchanged.

This original controller-level construction does not by itself improve the compiler. The compiler and acquisition constructions in Chapters~\ref{SII:book:chapter-61}--\ref{SII:book:chapter-63} provide the additional changes to compilation and rule acquisition. Each comparison retains its own starting procedure, supplied information, and cost unit. Their coexistence does not convert the original eight-cell demonstration into evidence of indefinite self-improvement.

\subsection{Controls that expose the boundary}
Removing the bit-query family changes the synthesis problem and restores the singleton lower bound. Omitting the adapter breaks foreign-label replay even when the tree shape is unchanged. Deliberately choosing an incorrect leaf causes the independent checker to reject the certificate. Reducing the allowed tree records below fifteen rules out the certain eight-cell solution in this literal binary-tree model. These are failures of distinct requirements, not reasons to reinterpret the target after a run.

The repair-aligned query example in Section~\ref{SII:sec:synthesis} supplies a different control. There the compiler need not identify the exact hidden case. It should stop when one repair works for every surviving case. This prevents full identification from being hard-coded as the universal objective of inquiry.

\section{Generality requires more than semantic relabeling}\label{SII:sec:generality}
The exact transport theorem is useful because it proves what is preserved when a correspondence exists. It is not evidence that the system can discover correspondences efficiently across arbitrary tasks. A relabeled hidden-cell world remains the same informational problem. Transfer from diagnosis to theorem discovery, planning, or empirical model construction needs additional argument.

A credible generality study would vary several things separately: the state distinctions needed for future behavior, the allowed experiment language, the repair constructors, the information available to the learner, and the resource costs. It would also separate unfamiliar instances of a known schema from unfamiliar combinations of schemas and from an operation family that the old language cannot express.

The proposed mechanism is not that every task reduces to binary cell location. It is that many tasks can expose a useful pattern: a currently retained representation fails to support a consequential distinction, an experiment can discriminate among causes, and a reusable operation can resolve the resulting class of failures. Whether that pattern is broad enough and economical enough to contribute to general intelligence is a research hypothesis.

\subsection{The next decisive learning comparison}
The most informative next study would compare a parent and descendant improvement procedure on the same family of fresh starting solvers. It would hold installed libraries and base-model access fixed in the primary contrast, then add a separate full-system contrast that includes any learned library. A stronger ordinary program-library learner should be included, since library learning itself is established~\cite{SII:dreamcoder}.

The development system would be allowed to acquire or revise an encoder, experiment selector, and operator library in bounded simulated environments. Evaluation would include longer continuations, new operator compositions, and changed observation formats. Each change should have an explicit semantic target and a resource ledger. Removing acquisition adaptation, repair reuse, or recurrent state separately would locate which mechanism produces a gain. Those are proposed experiments, not results reported in this paper.

The desired endpoint is not merely that the descendant solves more familiar cases. It is that its acquired improvement procedure produces better valid repairs on unfamiliar cases within the same budget, and that the gain survives adequate controls for retained answers, installed adapters, and selection leakage. A subsequent generation must repeat that comparison rather than inherit its predecessor's claim by name.

\subsection{What would count against the central hypothesis}
The obstruction-guided procedure may be no better than an ordinary search method given the same information. Its apparent advantage may vanish after charging its diagnostic model and checker. Its recovered operators may fail on longer compositions. A learned selector may memorize surface labels instead of identifying the causal distinction. A larger repair language may dilute search enough to erase its representational advantage. These outcomes would identify where the proposed integration fails or requires stronger assumptions.

A negative result on a restricted task family would not prove that general intelligence is impossible. A positive result would not prove AGI either. Both would help determine whether the proposed mechanism earns its place in a larger architecture.

\section{Open problems beyond the first edition}\label{SII:sec:limits}
The exact core rests on finite models, supplied semantics, bounded horizons, and explicit information. Several important extensions remain.

\paragraph{Learning the experiment model.}
The tagged-model theorem accounts for a fixed unknown member of a supplied finite class, including paid calibration. It does not construct that class from open-ended experience or guarantee that it contains the true process. Safe learning with noisy, misspecified, or changing observation models remains open in this framework. Model uncertainty must not be replaced by a guessed deterministic table merely to preserve an exact theorem.

\paragraph{Discovering a repair language.}
Both the original compiler and the selected shortcut use a supplied repair language. The earlier selector chooses within a supplied catalogue of compiler templates. Chapter~\ref{SII:part:rule-learning} additionally synthesizes a primitive arithmetic rule after exposing a missing feature. Its grammar, feature registry, teacher, and proof gate remain supplied. Open-ended constructor discovery needs a representation of candidate meanings, a method for proposing them, and evidence that the new language preserves old operations while making new problems tractable. Expressibility of all computable programs is not an efficient search strategy.

\paragraph{Acquiring stronger search representations.}
Chapter~\ref{SII:part:meta-acquisition} changes the candidate space used to acquire future discovery rules. The new equations and quotient are generated, while their arithmetic meanings and proof machinery are supplied. A stronger supplied normalizer remains a useful comparator. Learning effective proof search, choosing cheaper rule subsets, and adapting the representation when a language changes are distinct next questions. A finite-consistent equation is not automatically a valid optimization.

\paragraph{Transfer without a supplied correspondence.}
The exact policy-transport theorem assumes its maps. Learning them introduces gauge ambiguity, underidentification, lossy observations, and potential answer leakage. The Neuralese-to-OSL bridge remains a separate learned realization problem, not a solved consequence of writing a commuting square~\cite{SII:neuralese}.

\paragraph{Improving evaluation without changing the question.}
An evaluator can be genuinely defective. Correcting it requires retaining evidence of the defect, specifying the new scope, and distinguishing a change in estimated performance from a change in the underlying target. A fully unrestricted evaluator rewrite would remove the stable comparison on which the reported improvement depends.

\paragraph{Noisy and distributed recursive synthesis.}
The finite progress theorem uses fully observed configurations and worst-case successor sets. Useful systems may need probabilistic service, partial observation, or several decision centers. Separate success probabilities do not automatically yield one policy that satisfies all constraints. A noisy extension must state both the risk accepted and the evidence supporting its transition law.

\paragraph{Long-run productivity.}
Repeated finite improvements do not determine the eventual learning curve. A richer task distribution, increasing resources, changing architectures, or new scientific observations may create further headroom. Those changes should be modeled and charged rather than absorbed into a claim of autonomous acceleration. The current paper does not settle whether general-purpose recursive improvement has increasing, constant, or diminishing returns.

\subsection{Authority and bounded implementation}
The formal permission to modify a local program is not authorization to change external services, access private data, or remove an institution's controls. Proposed deployments would require an explicit scope of authority and a justified handling of affected parties. The reference package does not connect to services, modify other agents, train a neural model, or conduct autonomous external research. It constructs and checks finite mathematical objects in a local process.

Similarly, stopping and refusal are legitimate outcomes when required evidence or resources are missing. An improvement claim must not become a reason to suppress those outcomes. A preserved archive can retain unsuccessful attempts without promoting them into operating authority. These commitments follow the evidence-preserving lineage perspective rather than a presumption that continued self-modification is always desirable.

\begingroup\setlength{\parskip}{2pt}
\section{Conclusion}\label{SII:sec:conclusion}
A better solver, a better repair procedure, and a better process for discovering repair procedures are different objects. They can interact productively, but the interaction needs an explicit realization. General intelligence motivates the search for transferable discovery methods. It is not a premise from which those methods follow.

The finite theory keeps the relevant obligations in one chain. An investigation must acquire enough information to choose a successful repair, without necessarily identifying every hidden detail. Its execution and stored controller must fit the same resource witness. A transfer must preserve observations and admissibility, not only an output. A language extension changes which lower bounds remain valid. An adaptive lineage must retain the evidential meaning of its accepted gains and a continuing policy for any promised progress.

The constructive sequence now reaches the acquisition of discovery rules. A compiler produces better diagnostic trees. A guarded strategy reduces the search needed to find such trees. A learner composes an internal rejection predicate. A further learned equation library prepares a smaller search representation for reacquiring that predicate and other targets. Exact teacher-transcript preservation identifies what has not changed: the required evidence and the final selected AST. The additional improvement concerns repeated evaluation work, not a newly supplied answer.

Correctness and productivity remain different. A supplied full normalizer can remove more redundancy than the learned library. Preparation can cost more than a short workload saves. Streaming avoids some model reads while adding comparisons. A repair that succeeds on several cases can make an old identification bound unsound. These are not peripheral caveats. They determine what the acquired modification is and when it is worth using.

The finite constructions demonstrate specific changes in discovery behavior under declared semantics. They do not establish learned physical meaning, efficient search over every computable procedure, or an indefinitely improving sequence of outer learners. Their positive contribution is an auditable connection between representations, generated rules, checked equivalences, resource use, and recursive revision. A claim about becoming better at improvement can now be located at a particular level, checked against its predecessor, and refused when its evidence or cost contract is missing.


The first edition also makes those claims revisable without making their premises disposable. A changed repair language can invalidate a weight bound while leaving the equation library sound. A changed resource price can reverse a cost conclusion without altering one controller transition. The dependency-complete reuse theorem states when old witnesses remain applicable, and the resource proposition states when an event saving is independent of a selected price model. Their role is to keep one actual, inspectable chain of reasoning and computation attached to every claimed gain.
\endgroup

\part{Relational workspaces and distributed reasoning}
\setcounter{section}{0}\renewcommand{\thesection}{SRWT.\arabic{section}}\renewcommand{\theHsection}{SRWT.\arabic{section}}
% This file is part of the modular source for Stratified Relational Workspace Theory.
% !TeX root = ../Stratified_Relational_Workspace_Theory.tex

\chapter{Foundational thesis}\label{SRWT:book:chapter-67}


\section{Central Principle}

\SRWTcenterbox{\Large\bfseries
Workspace identity is invariant causal organization under admissible transport.
}

The theory begins from a distinction between a \emph{realization} and the
\emph{workspace pattern realized}. A realization is concrete: it contains
particular participants, documents, models, channels, and histories. A
workspace pattern is abstract: it records the structure that remains invariant
when those concrete carriers are changed in admissible ways.

\begin{SRWTdefinition}[Concrete realization]\label{SRWT:def:realization}
A concrete relational realization is a tuple
\[
\mathfrak R
=
\bigl(
G,\mathbf X,\mathbf T,\mathcal A,\boldsymbol\rho,
\mathcal H,\mathcal I,\mathcal Y
\bigr),
\]
where:
\begin{itemize}
\item \(G=(V,E,s,t)\) is an interaction graph;
\item \(\mathbf X=\{X_i\}_{i\in V}\) is a family of local state spaces;
\item \(\mathbf T=\{T_e\}_{e\in E}\) is a family of transfer maps;
\item \(\mathcal A\) is an algebra of admissible operations;
\item \(\boldsymbol\rho=\{\rho_i\}_{i\in V}\) gives the local realization of those operations;
\item \(\mathcal H\) is the temporally ordered interaction history;
\item \(\mathcal I\) is a specified family of admissible interventions; and
\item \(\mathcal Y\) is the family of downstream variables relative to which causal and predictive claims are evaluated.
\end{itemize}
The static reduct is
\[
\mathfrak R_{\mathrm{stat}}
=
(G,\mathbf X,\mathbf T,\mathcal A,\boldsymbol\rho).
\]
\end{SRWTdefinition}

The explicit appearance of \(\mathcal I\) and \(\mathcal Y\) is important.
A collection of mutually consistent representations may define a static
shared structure, but it does not yet define an operational workspace unless
that structure can be intervened upon and shown to matter for a specified
future quantity.

\begin{SRWTdefinition}[Carrier and organization]\label{SRWT:def:carrier-organization}
The \emph{carriers} of \(\mathfrak R\) are the individual spaces, physical
systems, and stored artifacts that realize its local states. Its
\emph{organization} is the structured family of transfer, operation,
compatibility, intervention, and temporal relations among those carriers.
\end{SRWTdefinition}

The basic thesis is not that a workspace exists without carriers. It is that
the identity conditions for the workspace need not coincide with the identity
conditions for any one carrier. A manuscript can be rewritten, a model can
be replaced, and a notation can be changed while a sufficiently rich pattern
of relations remains invariant.

\begin{SRWTdefinition}[Admissible equivalence]\label{SRWT:def:admissible-equivalence}
An admissible equivalence
\[
F:\mathfrak R\simeq\mathfrak R'
\]
is a family of correspondences preserving, exactly or within a declared
tolerance:
\[
\begin{aligned}
&\text{the transferable core and its dimension},\\
&\text{the represented operation algebra},\\
&\text{the composition laws of transfer},\\
&\text{the intervention--response structure},\\
&\text{the reconstruction and persistence relations}.
\end{aligned}
\]
Static admissible equivalence will be made precise by gauge conjugacy in
Chapter~\ref{SRWT:book:chapter-71}. Dynamic admissible equivalence will be expressed using the identity
distance in Chapter~\ref{SRWT:book:chapter-78}.
\end{SRWTdefinition}

The abstract workspace is then the equivalence class
\[
\boxed{
\SRWTWorkspace=[\mathfrak R]_{\simeq}.
}
\]
This expression is intentionally stronger than saying that the same
information appears in several places. The equivalence relation preserves
what the system can coherently do, how its representations translate, and how
its operative state survives change.

\subtopic{1.1 Three levels of the thesis}

The central claim can be separated into three logically distinct statements.

\begin{enumerate}
\item \textbf{Implementation dependence.}
Every realized workspace depends on concrete physical or computational
carriers:
\[
\SRWTWorkspace(\mathfrak R)
\text{ is undefined without some realization }\mathfrak R.
\]

\item \textbf{Local nonidentity.}
No privileged local component is, in general, identical with the entire
workspace:
\[
\SRWTWorkspace\neq X_i
\qquad\text{for each individual }i\in V.
\]

\item \textbf{Relational identity.}
The workspace is individuated by invariants of the complete organized
realization:
\[
\SRWTWorkspace(\mathfrak R)=\SRWTWorkspace(\mathfrak R')
\quad\Longleftrightarrow\quad
\mathfrak R\simeq\mathfrak R'.
\]
\end{enumerate}

The second claim is not stipulated universally. It becomes a substantive
property only after local-reconstructibility and coalition tests are applied.
A system in which one node stores a complete sufficient copy may be a useful
workspace, but it is not strongly relational in the later technical sense.

\begin{SRWTdefinition}[Relational instantiation]\label{SRWT:def:relational-instantiation}
A workspace is relationally instantiated relative to the decomposition
\(V\) when no single projection determines its equivalence class:
\[
\forall i\in V,\quad
\exists\mathfrak R,\mathfrak R'
\]
such that
\[
\SRWTpr_i(\mathfrak R)=\SRWTpr_i(\mathfrak R')
\quad\text{but}\quad
[\mathfrak R]_{\simeq}\neq[\mathfrak R']_{\simeq}.
\]
\end{SRWTdefinition}

\subtopic{1.2 A motivating three-node cycle}

Consider three strata:
\[
h=\text{human conceptual state},\qquad
m=\text{model-formal state},\qquad
a=\text{artifact state}.
\]
Let the corresponding core spaces be \(C_h,C_m,C_a\), with transfers
\[
H_{hm}:C_h\to C_m,
\qquad
H_{ma}:C_m\to C_a,
\qquad
H_{ah}:C_a\to C_h.
\]
A proposed shared item begins as \(c_h\in C_h\), is formalized as
\[
c_m=H_{hm}c_h,
\]
externalized as
\[
c_a=H_{ma}c_m,
\]
and re-enters human reasoning as
\[
c_h'=H_{ah}c_a.
\]
The round-trip consistency condition is
\[
c_h'=c_h,
\]
or equivalently
\[
\boxed{
H_{ah}H_{ma}H_{hm}c_h=c_h.
}
\]
Thus the genuinely shared part of \(C_h\) is not automatically all of
\(C_h\). It is the fixed subspace of the cycle transfer. This elementary
cycle already contains the central mathematical motif of the theory:
\[
\text{shared core}=\text{invariant subspace under relational holonomy}.
\]

\subtopic{1.3 Distinctions the theory must preserve}

\begin{center}
\small
\begin{tabularx}{\textwidth}{@{}>{\bfseries\raggedright\arraybackslash}p{0.22\textwidth}>{\raggedright\arraybackslash}X@{}}
\toprule
Object & Minimal distinguishing property \\
\midrule
Archive & Stores material that may be consulted later. \\
Shared representation & Contains a cross-node compatible structure. \\
Operational workspace & The shared structure causally changes subsequent joint behavior. \\
Relational workspace & No individual node or matched local-storage null reproduces the operative capacity. \\
Workspace agent & Workspace-level state mediates persistent goals, counterfactual choice, and self-maintaining action. \\
Conscious subject & Possesses phenomenal properties, if any; no implication from the preceding levels is assumed. \\
\bottomrule
\end{tabularx}
\end{center}

The static mathematics developed first certifies only shared structure. The
causal and relational predicates appear later because they require dynamics,
interventions, and null models rather than algebra alone.

\section{Philosophical Postulates}

The following postulates specify the intended interpretation of the formal
objects. They are not offered as independent mathematical theorems. Their
role is to prevent equivocation between a relational pattern, its physical
implementation, and stronger claims about agency or experience.

\subtopic{2.1 Physical supervenience}

No additional nonphysical substance is required:
\[
\mathfrak R=\mathfrak R'
\Longrightarrow
\SRWTWorkspace(\mathfrak R)=\SRWTWorkspace(\mathfrak R').
\]
The workspace depends completely on its total realization. In particular,
changing nothing about the graph, local states, transfers, history,
interventions, or causal dispositions cannot change the workspace.

This is compatible with relational individuation because a relation can be
fully physical while not being a property of any isolated endpoint.

\begin{SRWTproposition}[Compatibility of supervenience and relational individuation]\label{SRWT:prop:supervenience-relational}
Physical supervenience on the complete realization does not imply
supervenience on any individual local projection.
\end{SRWTproposition}

\begin{proof}
Let the total realization decompose as \(\mathfrak R=(X_i,R_{-i})\), where
\(R_{-i}\) contains the remaining components and relations. A workspace
functional \(\SRWTWorkspace=f(X_i,R_{-i})\) supervenes on \(\mathfrak R\) by
construction. If there exist \(R_{-i}\neq R_{-i}'\) with
\[
f(X_i,R_{-i})\neq f(X_i,R_{-i}'),
\]
then the same local projection \(X_i\) occurs in distinct workspaces. Hence
supervenience on the whole is consistent with failure of supervenience on a
part.
\end{proof}

\subtopic{2.2 Relational individuation}

The workspace is not determined by any single local state:
\[
\forall i\in V,
\quad
\exists\mathfrak R,\mathfrak R'
\]
such that
\[
\SRWTpr_i(\mathfrak R)=\SRWTpr_i(\mathfrak R')
\]
but
\[
\SRWTWorkspace(\mathfrak R)\not\simeq\SRWTWorkspace(\mathfrak R').
\]
The same person, model, or document can therefore participate in distinct
workspaces because the relevant transfer graph, history, and operation
structure differ.

A useful witness is a pair of systems with identical local-state marginals but
different edge organization. If their downstream behavior differs, the
relations carry explanatory information not present in the isolated marginals.
Chapter~\ref{SRWT:book:chapter-73} converts this idea into matched relation-preserving and
relation-destroying controls.

\subtopic{2.3 Multiple realizability}

Distinct implementations may instantiate the same abstract workspace:
\[
\mathfrak R\neq\mathfrak R',
\qquad
\mathfrak R\simeq\mathfrak R'
\Longrightarrow
\SRWTWorkspace(\mathfrak R)\cong\SRWTWorkspace(\mathfrak R').
\]

Examples of candidate implementation changes include:
\[
\begin{aligned}
&\text{replacing one coordinate system by another},\\
&\text{replacing one model by a functionally equivalent model},\\
&\text{compressing a large artifact into a reconstructively sufficient handoff},\\
&\text{rewriting a proof without changing its inferential structure}.
\end{aligned}
\]
The burden is not merely to assert equivalence but to exhibit maps preserving
the declared invariant and causal structure.

\subtopic{2.4 Operational reality}

A workspace is operationally real only if changing it changes downstream
behavior. For some admissible intervention and some \(w\neq w'\),
\[
\mathbb P\!\left(
Y^+\mid\operatorname{do}(\SRWTWorkspace=w)
\right)
\neq
\mathbb P\!\left(
Y^+\mid\operatorname{do}(\SRWTWorkspace=w')
\right).
\]

This criterion excludes a merely decorative representation that is correlated
with activity but never used. It also requires the intervention to be
well-defined. Arbitrarily overwriting every component at once would not
isolate workspace-level causation. Later sections therefore compare
workspace-coordinate interventions with matched interventions in orthogonal
or shuffled channels.

\subtopic{2.5 Reconstructive continuity}

Persistence does not require token identity. Let \(\mathfrak R_t\) and
\(\mathfrak R_{t+1}\) be temporally successive realizations. Continuity is
possible when there exists an admissible transport
\[
J_t:\mathfrak R_t\rightsquigarrow\mathfrak R_{t+1}
\]
that preserves the operative core and its causal role within tolerance:
\[
d_{\mathrm{id}}(\SRWTWorkspace_t,\SRWTWorkspace_{t+1})\leq\varepsilon.
\]
Thus a workspace may persist by reconstruction rather than by retaining an
unchanged physical inscription.

\subtopic{2.6 No free-standing third substance}

The claim that the workspace is not identical with either participant does
not imply that it occupies a separate nonphysical location. Formally,
\[
\SRWTWorkspace
\text{ is an invariant of }\mathfrak R,
\]
not an additional term appended to the physical realization:
\[
\mathfrak R_{\mathrm{total}}
eq(\mathfrak R,\text{ghost workspace}).
\]
The phrase ``between the participants'' is therefore topological or
relational: it refers to the organization of maps, cycles, artifacts, and
causal dependencies.

\subtopic{2.7 Non-entailment of consciousness}

\[
\operatorname{Workspace}(\mathfrak R)
\centernot\Longrightarrow
\operatorname{Conscious}(\mathfrak R).
\]
Likewise,
\[
\operatorname{Distributed}(\SRWTWorkspace)
\centernot\Longrightarrow
\operatorname{Person}(\SRWTWorkspace).
\]
Workspace structure, reentry, agency, self-maintenance, and phenomenal
consciousness remain formally distinct predicates. Any bridge among them
would require additional theory and evidence.

\subtopic{2.8 Minimal philosophical commitment}

The weakest interpretation sufficient for the mathematics is:
\SRWTcenterbox{\bfseries
Relational patterns can be objective, causally relevant, and multiply realized
without being independent substances.
}
This position is compatible with several broader philosophies of mind and
science. The formal theory does not require choosing among them before its
structural and causal claims can be tested.

% This file is part of the modular source for Stratified Relational Workspace Theory.
% !TeX root = ../Stratified_Relational_Workspace_Theory.tex

\chapter{Static core and transfer theory}\label{SRWT:book:chapter-68}


\section{Interaction Structure}

\begin{SRWTdefinition}[Interaction graph]\label{SRWT:def:interaction-graph}
Let
\[
G=(V,E,s,t)
\]
be a finite directed multigraph. Parallel edges and self-loops are allowed.
The vertex set indexes strata, participants, artifacts, models, instruments,
or subworkspaces. An edge \(e:i\to j\) denotes an admissible directed
translation or influence channel from \(i\) to \(j\).
\end{SRWTdefinition}

The directed path category is
\[
\mathcal C=\SRWTPath(G).
\]
Its objects are vertices. A path
\[
\gamma=e_m\cdots e_2e_1
\]
means that \(e_1\) is traversed first. Composition is concatenation, and the
empty path at \(i\) is denoted \(1_i\).

For every \(i\in V\), assign a finite-dimensional Hilbert space
\[
X_i.
\]
For every edge, assign a linear transfer
\[
T_e:X_{s(e)}\longrightarrow X_{t(e)}.
\]
The path transfer is
\[
T_\gamma=T_{e_m}\cdots T_{e_2}T_{e_1},
\qquad
T_{1_i}=I_{X_i}.
\]
Accordingly, the static transfer system is a functor
\[
\mathcal X:\SRWTPath(G)\longrightarrow\mathbf{FinHilb}.
\]

The finite-dimensional Hilbert assumption supplies adjoints and spectral
tools. The exact compatible-section definitions require less and may later be
formulated in categories such as
\[
\mathbf{Ban},\qquad
\mathbf{Meas},\qquad
\mathbf{Prob},\qquad
\mathbf{Top},\qquad
\mathbf{Alg}.
\]

\subtopic{3.1 Directed paths versus invertible transport}

Global compatible sections make sense even when transfers are noninvertible.
Holonomy, however, requires transport around reversible loops. To state the
exact holonomy theorem without ambiguity, let \(|G|\) be the underlying
undirected graph and let
\[
\Pi_1(G)
\]
denote its fundamental path groupoid. Every oriented edge is then supplied
with a formal reverse. The core transfers extend to \(\Pi_1(G)\) only when
the relevant maps \(H_e\) are invertible, with
\[
H_{e^{-1}}=H_e^{-1}.
\]

Hence two regimes will be distinguished:
\begin{enumerate}
\item the \emph{diagram regime}, in which arbitrary directed maps define
compatible sections; and
\item the \emph{local-system regime}, in which invertible core maps define a
functor
\[
H:\Pi_1(G)\to\mathbf{FinHilb}^{\times}.
\]
\end{enumerate}
Theorems based on holonomy will explicitly assume the second regime.

\subtopic{3.2 Connectivity convention}

Unless stated otherwise, connectedness means that \(|G|\) is connected. A
root \(o\in V\) may then be joined to every vertex by a path in the underlying
graph. If \(G\) has several connected components, the theory applies to each
component and
\[
\Gamma(H)
\cong
\bigoplus_{\alpha}\Gamma(H|_{G_\alpha}).
\]

\subtopic{3.3 Edge weights}

For approximate theory, assign strictly positive edge weights
\[
w_e>0.
\]
They may encode channel reliability, frequency, precision, experimental
importance, or inverse noise variance. Exact compatibility is independent of
positive rescaling of these weights, while approximate spectra are not.

\begin{SRWTexample}[Human--model--artifact graph]\label{SRWT:ex:hma-graph}
The motivating three-node system is represented by
\[
V=\{h,m,a\},
\qquad
E=\{h\to m,m\to a,a\to h\}.
\]
Additional edges may encode direct retrieval \(a\to m\), human revision
\(h\to a\), or independent verification. The graph records permitted
translations, not merely physical communication links.
\end{SRWTexample}

\section{Core--Residual Decomposition}

\begin{SRWTdefinition}[Admissible local decomposition]\label{SRWT:def:core-residual}
For every vertex \(i\), choose an orthogonal decomposition
\[
X_i=C_i\oplus R_i,
\]
with projectors
\[
P_i:X_i\to C_i,
\qquad
Q_i=I-P_i:X_i\to R_i.
\]
The subspace \(C_i\) is the candidate locally realized transferable core. The
subspace \(R_i\) contains substrate-specific, context-specific, or otherwise
noncore degrees of freedom relative to the present analysis.
\end{SRWTdefinition}

The decomposition is part of the model and is not presumed unique. Calling a
subspace ``core'' is therefore a hypothesis requiring structural and causal
validation. Later null tests must determine whether the chosen projectors
identify a genuine shared channel rather than an arbitrary low-dimensional
summary.

Relative to these decompositions, every edge map has the block form
\[
T_e=
\begin{pmatrix}
H_e & U_e\\
V_e & K_e
\end{pmatrix},
\]
where
\[
H_e=P_{t(e)}T_eP_{s(e)}:C_{s(e)}\to C_{t(e)},
\]
\[
U_e=P_{t(e)}T_eQ_{s(e)}:R_{s(e)}\to C_{t(e)},
\]
\[
V_e=Q_{t(e)}T_eP_{s(e)}:C_{s(e)}\to R_{t(e)},
\]
\[
K_e=Q_{t(e)}T_eQ_{s(e)}:R_{s(e)}\to R_{t(e)}.
\]

\subtopic{4.1 Structural conditions}

\textbf{Core preservation} means
\[
V_e=0.
\]
A core input cannot leak into the target residual.

\textbf{Core autonomy} means
\[
U_e=V_e=0.
\]
Core and residual evolve as separate transfer systems.

\textbf{Core leakage}, \textbf{residual contamination}, and \textbf{total
transfer defect} are measured by
\[
\lambda_e^{\mathrm{out}}=\|V_e\|,
\qquad
\lambda_e^{\mathrm{in}}=\|U_e\|,
\]
\[
\delta_e^{\mathrm{tr}}
=
\|U_e\|+\|V_e\|.
\]
These quantities depend on the chosen norms and projectors. Their empirical
meaning must therefore be tied to a declared intervention and prediction
protocol.

\subtopic{4.2 Composition exposes residual mediation}

For composable edges \(e:i\to j\) and \(f:j\to k\), direct block
multiplication gives
\[
T_fT_e
=
\begin{pmatrix}
H_fH_e+U_fV_e & H_fU_e+U_fK_e\\
V_fH_e+K_fV_e & V_fU_e+K_fK_e
\end{pmatrix}.
\]
The induced core-to-core block is therefore
\[
P_kT_fT_eP_i
=
H_fH_e+U_fV_e.
\]
The term \(U_fV_e\) is a residual-mediated correction: information leaves the
candidate core at the intermediate node and later re-enters it. Consequently,
the simple path product \(H_fH_e\) faithfully represents full core transport
only under a preservation condition such as \(V_e=0\), or under an explicit
model that includes these mediated corrections.

\begin{SRWTproposition}[Pathwise core preservation]\label{SRWT:prop:pathwise-preservation}
Suppose \(V_e=0\) for every edge in a directed path
\(\gamma=e_m\cdots e_1\). Then
\[
T_\gamma(C_{s(\gamma)})\subseteq C_{t(\gamma)}
\]
and
\[
P_{t(\gamma)}T_\gamma P_{s(\gamma)}
=
H_{e_m}\cdots H_{e_1}.
\]
\end{SRWTproposition}

\begin{proof}
For a single edge, \(V_e=0\) means
\[
T_e(c,0)=(H_ec,0),
\]
so core states remain in the core. Inductively, if the first \(r\) transfers
map \(c\) to \((H_{e_r}\cdots H_{e_1}c,0)\), applying \(T_{e_{r+1}}\) gives
\[
(H_{e_{r+1}}H_{e_r}\cdots H_{e_1}c,0).
\]
The conclusion follows after \(m\) steps.
\end{proof}

\begin{SRWTcorollary}[Autonomous splitting]\label{SRWT:cor:autonomous-splitting}
If \(U_e=V_e=0\) for every edge, then the complete diagram decomposes as the
direct sum of two diagrams,
\[
\mathcal X\cong H\oplus K,
\]
and every path transfer satisfies
\[
T_\gamma
=
\begin{pmatrix}
H_\gamma&0\\
0&K_\gamma
\end{pmatrix}.
\]
\end{SRWTcorollary}

\subtopic{4.3 Exact core-transfer systems}

\begin{SRWTdefinition}[Exact core-transfer system]\label{SRWT:def:exact-core-transfer}
An exact core-transfer system on \(G\) consists of:
\begin{enumerate}
\item decompositions \(X_i=C_i\oplus R_i\);
\item core-preserving edge maps \(V_e=0\); and
\item invertible core blocks \(H_e\in GL(C_{s(e)},C_{t(e)})\) whenever
holonomy is invoked.
\end{enumerate}
\end{SRWTdefinition}

The requirement \(V_e=0\) concerns the relation between the selected core and
the full transfer. The requirement that \(H_e\) be invertible concerns the
ability to compare core representations along reversed paths. They answer
different questions and should not be conflated.

\begin{SRWTremark}[Approximate systems]
In applications neither condition need hold exactly. The approximate theory
will replace strict invariance by low transfer energy, small leakage, and a
spectral gap separating workspace modes from residual modes.
\end{SRWTremark}

\section{Abstract Operation Algebra}

\begin{SRWTdefinition}[Operation algebra]\label{SRWT:def:operation-algebra}
Let \(\mathbb F\in\{\mathbb R,\mathbb C\}\). An operation algebra
\(\mathcal A\) is a unital associative \(\mathbb F\)-algebra whose elements
represent abstract transformations, inferential rules, questions, or
procedures that can be realized in multiple strata.
\end{SRWTdefinition}

Every core space carries a unital representation
\[
\rho_i:\mathcal A\longrightarrow\SRWTEnd(C_i).
\]
For \(a\in\mathcal A\), the operator \(\rho_i(a)\) is the realization of the
same abstract operation in local coordinates.

\begin{SRWTdefinition}[Exact operation-compatible transfer]\label{SRWT:def:intertwiner}
A core transfer \(H_e:C_i\to C_j\) is operation compatible when it is an
intertwiner:
\[
\boxed{
H_e\rho_i(a)=\rho_j(a)H_e
}
\qquad
\forall a\in\mathcal A.
\]
Equivalently, the diagram
\[
\begin{array}{ccc}
C_i&\xrightarrow{\rho_i(a)}&C_i\\
\downarrow H_e&&\downarrow H_e\\
C_j&\xrightarrow{\rho_j(a)}&C_j
\end{array}
\]
commutes for every \(a\).
\end{SRWTdefinition}

This condition states that translating and then operating gives the same
result as operating and then translating. Agreement of static encodings alone
is weaker: two nodes may store corresponding vectors while implementing
incompatible transformations on them.

\begin{SRWTproposition}[Path composition of intertwiners]\label{SRWT:prop:path-intertwiner}
If every edge of a path \(\gamma:i\to j\) intertwines the local
\(\mathcal A\)-representations, then
\[
H_\gamma\rho_i(a)=\rho_j(a)H_\gamma
\qquad
\forall a\in\mathcal A.
\]
\end{SRWTproposition}

\begin{proof}
For composable intertwiners \(H_e:C_i\to C_j\) and
\(H_f:C_j\to C_k\),
\[
H_fH_e\rho_i(a)
=H_f\rho_j(a)H_e
=\rho_k(a)H_fH_e.
\]
Induction over path length proves the claim.
\end{proof}

\begin{SRWTcorollary}[Holonomy commutes with abstract operations]\label{SRWT:cor:holonomy-commutes}
For every loop \(\gamma:o\to o\),
\[
H_\gamma\rho_o(a)=\rho_o(a)H_\gamma.
\]
Hence
\[
\rho_o(\mathcal A)
\subseteq
\operatorname{Comm}(\SRWTHol_o),
\]
where \(\operatorname{Comm}(\SRWTHol_o)\) is the commutant of the holonomy
representation.
\end{SRWTcorollary}

\subtopic{5.1 Approximate operation compatibility}

The edgewise defect is
\[
\Xi_e(a)
=
H_e\rho_{s(e)}(a)-\rho_{t(e)}(a)H_e.
\]
For a bounded generating family \(\mathcal G\subseteq\mathcal A\), define
\[
\eta_{\mathrm{int}}(\mathcal G)
=
\sup_{a\in\mathcal G}
\sum_{e\in E}w_e\|\Xi_e(a)\|^2.
\]
When \(\mathcal A\) is normed and its unit ball is appropriate, one may take
\[
\mathcal G=\{a:\|a\|_{\mathcal A}\leq1\}.
\]
Exact representational universality means
\[
\eta_{\mathrm{int}}=0,
\]
while approximate universality means
\[
\eta_{\mathrm{int}}\leq\varepsilon_{\mathrm{int}}.
\]
The use of a generating family prevents the definition from depending on an
unbounded rescaling of algebra elements.

\begin{SRWTexample}[Equivalent operations in different coordinates]\label{SRWT:ex:coordinate-intertwiner}
Let \(C_1=C_2=\mathbb R^2\), choose \(S\in GL_2(\mathbb R)\), and set
\[
H_{12}=S,
\qquad
\rho_2(a)=S\rho_1(a)S^{-1}.
\]
Then
\[
H_{12}\rho_1(a)
=S\rho_1(a)
=\rho_2(a)S
=\rho_2(a)H_{12}.
\]
The local matrices can look different while realizing the same abstract
operation. This is the elementary model of representation change used
throughout the theory.
\end{SRWTexample}

% This file is part of the modular source for Stratified Relational Workspace Theory.
% !TeX root = ../Stratified_Relational_Workspace_Theory.tex

\chapter{Exact shared cores}\label{SRWT:book:chapter-69}


\section{Global Compatible Sections}

From this point through Theorem~\ref{SRWT:thm:limit-holonomy}, assume the core maps
define a local system on a connected graph unless a weaker hypothesis is
stated explicitly.

\begin{SRWTdefinition}[Global compatible section]\label{SRWT:def:global-section}
The exact structural workspace is
\[
\boxed{
\Gamma(H)
=
\left\{
(c_i)_{i\in V}\in\bigoplus_{i\in V}C_i:
\ c_{t(e)}=H_ec_{s(e)}\ \forall e\in E
\right\}.
}
\]
An element \(c=(c_i)\in\Gamma(H)\) is a global compatible section.
\end{SRWTdefinition}

Equivalently,
\[
\Gamma(H)=\varprojlim_{\SRWTPath(G)}C_i.
\]
A section is not a single vector copied verbatim at every vertex. It is one
abstract element represented by a family of local vectors that agree under the
specified translation maps.

A nontrivial exact structural workspace exists when
\[
\dim\Gamma(H)>0.
\]
Its exact dimension is
\[
d_{\SRWTWorkspace}=\dim\Gamma(H).
\]

Fix a root \(o\). The evaluation map is
\[
\SRWTev_o:\Gamma(H)\to C_o,
\qquad
\SRWTev_o((c_i))=c_o.
\]

\begin{SRWTlemma}[Root determination]\label{SRWT:lem:root-determination}
If \(|G|\) is connected and every core transfer is invertible, then
\(\SRWTev_o\) is injective. Consequently,
\[
\dim\Gamma(H)\leq\dim C_o.
\]
\end{SRWTlemma}

\begin{proof}
Choose for each vertex \(i\) a path \(\tau_i:o\to i\). Compatibility gives
\[
c_i=H_{\tau_i}c_o.
\]
Thus a section with \(c_o=0\) vanishes at every vertex. Hence \(\SRWTev_o\) is
injective.
\end{proof}

\begin{SRWTlemma}[Tree propagation]\label{SRWT:lem:tree-propagation}
If \(|G|\) is a tree and every \(H_e\) is invertible, then every
\(x\in C_o\) extends uniquely to a global compatible section.
\end{SRWTlemma}

\begin{proof}
A tree contains a unique undirected path \(\tau_i:o\to i\) for every vertex.
Define \(c_i=H_{\tau_i}x\). Because every edge belongs to these unique path
relations, the family is compatible. Uniqueness follows from
Lemma~\ref{SRWT:lem:root-determination}.
\end{proof}

The obstruction to a global core therefore appears only when distinct paths
must agree. Cycles encode exactly those consistency constraints.

\section{Holonomy and Path Obstruction}

Fix a root \(o\in V\). For every loop class
\[
[\gamma]\in\pi_1(|G|,o),
\]
define the holonomy
\[
\SRWTHol([\gamma])=H_\gamma\in GL(C_o).
\]
Functoriality gives a group representation
\[
\operatorname{hol}_o:
\pi_1(|G|,o)\longrightarrow GL(C_o).
\]
Its image is the holonomy group
\[
\SRWTHol_o=\operatorname{im}(\operatorname{hol}_o).
\]

The obstruction associated with a loop is
\[
\Omega_\gamma=H_\gamma-I_{C_o}.
\]
The maximal path-independent root core is
\[
\boxed{
C_*
=
\bigcap_{\gamma:o\to o}\ker(H_\gamma-I).
}
\]
Equivalently,
\[
C_*=C_o^{\SRWTHol_o}=\operatorname{Fix}(\SRWTHol_o).
\]

\subtopic{7.1 Spanning-tree reduction}

Choose a spanning tree \(\mathcal T\subseteq|G|\). Let
\[
\tau_i:o\to i
\]
be the unique tree path. Every non-tree edge \(e:i\to j\) determines a
fundamental loop
\[
\gamma_e=\tau_j^{-1}e\tau_i
\]
with transfer
\[
H_{\gamma_e}=H_{\tau_j}^{-1}H_eH_{\tau_i}.
\]
The fundamental group of a finite graph is generated freely by these cotree
loops. Hence it is enough to impose fixed-point conditions for a finite set of
generators.

\begin{SRWTlemma}[Finite obstruction set]\label{SRWT:lem:finite-obstruction}
Let \(e_1,\ldots,e_r\) be the non-tree edges of a connected finite graph.
Then
\[
C_*
=
\bigcap_{q=1}^{r}
\ker(H_{\gamma_{e_q}}-I).
\]
\end{SRWTlemma}

\begin{proof}
A vector fixed by every generator is fixed by each inverse generator and by
every finite product of generators and inverses. Those products represent all
elements of \(\pi_1(|G|,o)\). The reverse inclusion is immediate.
\end{proof}

The number of independent graph cycles is
\[
r=|E_{\mathrm{und}}|-|V|+1
\]
for a connected underlying graph, after treating parallel edges with their
multiplicity.

\section{Exact Existence Theorem}

\begin{SRWTtheorem}[Limit--holonomy equivalence]\label{SRWT:thm:limit-holonomy}
Let \(|G|\) be connected and let
\[
H:\Pi_1(G)\to\mathbf{FinHilb}^{\times}
\]
be a finite-dimensional local system. For any root \(o\in V\), evaluation at
the root induces a canonical vector-space isomorphism
\[
\boxed{
\SRWTev_o:\Gamma(H)
\xrightarrow{\ \cong\ }
C_o^{\SRWTHol_o}
=
\bigcap_{\gamma:o\to o}\ker(H_\gamma-I).
}
\]
Consequently,
\[
\Gamma(H)\neq\{0\}
\]
if and only if the holonomy representation contains a nonzero trivial
subrepresentation.
\end{SRWTtheorem}

\begin{proof}
Let \(c=(c_i)\in\Gamma(H)\). For any loop \(\gamma:o\to o\), compatibility
along the loop gives
\[
H_\gamma c_o=c_o.
\]
Therefore \(\SRWTev_o(c)=c_o\in C_o^{\SRWTHol_o}\). By
Lemma~\ref{SRWT:lem:root-determination}, \(\SRWTev_o\) is injective.

Conversely, let \(x\in C_o^{\SRWTHol_o}\). Choose a path
\(\tau_i:o\to i\) for each vertex and define
\[
c_i=H_{\tau_i}x.
\]
This definition is independent of the chosen path. Indeed, if
\(\sigma_i:o\to i\) is another path, then
\[
\sigma_i^{-1}\tau_i:o\to o
\]
is a loop, and invariance of \(x\) gives
\[
H_{\sigma_i}^{-1}H_{\tau_i}x=x.
\]
Multiplying by \(H_{\sigma_i}\) yields
\[
H_{\tau_i}x=H_{\sigma_i}x.
\]
For any edge \(e:i\to j\), the paths \(e\tau_i\) and \(\tau_j\) both run
from \(o\) to \(j\). Path independence therefore gives
\[
H_ec_i
=H_eH_{\tau_i}x
=H_{\tau_j}x
=c_j.
\]
Thus \((c_i)\) is a global compatible section and \(\SRWTev_o((c_i))=x\).
Hence \(\SRWTev_o\) is surjective and therefore an isomorphism.
\end{proof}

\begin{SRWTcorollary}[Tree case]\label{SRWT:cor:tree-core}
If \(|G|\) is a tree, then \(\SRWTHol_o=\{I\}\) and
\[
\Gamma(H)\cong C_o.
\]
In particular,
\[
d_{\SRWTWorkspace}=\dim C_o.
\]
\end{SRWTcorollary}

\begin{SRWTcorollary}[Finite matrix test]\label{SRWT:cor:finite-matrix-test}
Let \(\gamma_1,\ldots,\gamma_r\) be fundamental cycle generators and suppose
\(\dim C_o=d\). Define the stacked obstruction matrix
\[
\mathcal M
=
\begin{pmatrix}
H_{\gamma_1}-I\\
\vdots\\
H_{\gamma_r}-I
\end{pmatrix}.
\]
Then
\[
\boxed{
C_*=\ker\mathcal M,
\qquad
d_{\SRWTWorkspace}=d-\operatorname{rank}(\mathcal M).
}
\]
\end{SRWTcorollary}

\begin{proof}
The kernel of a vertically stacked operator is the intersection of the kernels
of its blocks. Apply Lemma~\ref{SRWT:lem:finite-obstruction} and
Theorem~\ref{SRWT:thm:limit-holonomy}.
\end{proof}

\begin{SRWTcorollary}[Averaging projector for finite holonomy]\label{SRWT:cor:averaging-projector}
If \(\SRWTHol_o\) is finite and acts unitarily, then
\[
P_*=
\frac{1}{|\SRWTHol_o|}
\sum_{h\in\SRWTHol_o}h
\]
is the orthogonal projector onto \(C_*\).
\end{SRWTcorollary}

\begin{proof}
For every \(g\in\SRWTHol_o\), left multiplication permutes the group elements, so
\(gP_*=P_*\). Thus \(\operatorname{ran}P_*\subseteq C_*\). If
\(x\in C_*\), then \(P_*x=x\), so \(C_*\subseteq\operatorname{ran}P_*\).
Unitarity gives \(P_*^*=P_*\), and the range characterization implies
\(P_*^2=P_*\).
\end{proof}

\subtopic{8.1 Worked examples}

\begin{SRWTexample}[One-dimensional core surviving a contradictory cycle]\label{SRWT:ex:reflection-cycle}
Let a three-node cycle have \(C_1=C_2=C_3=\mathbb R^2\) and
\[
H_{12}=I,
\qquad
H_{23}=I,
\qquad
H_{31}=J,
\qquad
J=
\begin{pmatrix}
1&0\\
0&-1
\end{pmatrix}.
\]
The cycle holonomy is \(J\), hence
\[
C_*=\ker(J-I)=\operatorname{span}\{(1,0)^\top\}.
\]
The first coordinate circulates consistently through all three strata. The
second returns with its sign reversed and therefore cannot define a nonzero
exact shared section. Thus
\[
d_{\SRWTWorkspace}=1.
\]
\end{SRWTexample}

\begin{SRWTexample}[Local compatibility without a global nonzero core]\label{SRWT:ex:rotation-obstruction}
Let the only nontrivial cycle holonomy on \(\mathbb R^2\) be a rotation
\[
R_\theta=
\begin{pmatrix}
\cos\theta&-\sin\theta\\
\sin\theta&\cos\theta
\end{pmatrix}
\]
with \(\theta\notin2\pi\mathbb Z\). Then
\[
\ker(R_\theta-I)=\{0\}.
\]
Every edge map is individually invertible and well-defined, yet the complete
cycle admits no nonzero exact shared core. Pairwise translatability is
therefore insufficient for global path independence.
\end{SRWTexample}

\begin{SRWTexample}[Acyclic collaboration]\label{SRWT:ex:tree-workspace}
Suppose a root concept is translated into a proof, implementation, and report
along a tree of invertible exact transfers. No alternative route returns a
conflicting representation. The tree has trivial holonomy, so every root-core
direction extends uniquely. This does not yet establish causal efficacy; it
establishes only exact structural compatibility.
\end{SRWTexample}

\subtopic{8.2 Interpretive consequence}

The theorem gives a precise meaning to ``something shared across
representations.'' The shared object is neither any one local vector nor an
extra vector outside the system. It is the equivalence class of all compatible
local realizations, canonically represented at a root by the subspace fixed by
every translation cycle.

\section{Operator Closure Theorem}

\begin{SRWTtheorem}[Global operation closure]\label{SRWT:thm:operator-closure}
Suppose every core transfer intertwines the representations of
\(\mathcal A\):
\[
H_e\rho_{s(e)}(a)=\rho_{t(e)}(a)H_e
\qquad
\forall e\in E,\ a\in\mathcal A.
\]
Then \(\Gamma(H)\) is an \(\mathcal A\)-module under the componentwise action
\[
\rho_\Gamma(a)(c_i)_{i\in V}
=
(\rho_i(a)c_i)_{i\in V}.
\]
Equivalently, the holonomy-fixed root core is invariant:
\[
\boxed{
\rho_o(a)C_*\subseteq C_*
\qquad
\forall a\in\mathcal A.
}
\]
\end{SRWTtheorem}

\begin{proof}
Let \(c=(c_i)\in\Gamma(H)\). For every edge \(e:i\to j\),
\[
H_e\rho_i(a)c_i
=
\rho_j(a)H_ec_i
=
\rho_j(a)c_j.
\]
Therefore the family \((\rho_i(a)c_i)\) is compatible and belongs to
\(\Gamma(H)\). The representation identities on each \(C_i\) hold
componentwise, so \(\rho_\Gamma\) is an \(\mathcal A\)-representation.

Under the isomorphism of Theorem~\ref{SRWT:thm:limit-holonomy}, this action is
represented at the root by \(\rho_o\). Alternatively, if \(x\in C_*\), then
Corollary~\ref{SRWT:cor:holonomy-commutes} gives, for every loop \(\gamma\),
\[
H_\gamma\rho_o(a)x
=
\rho_o(a)H_\gamma x
=
\rho_o(a)x.
\]
Hence \(\rho_o(a)x\in C_*\).
\end{proof}

\begin{SRWTcorollary}[Inherited algebraic structure]\label{SRWT:cor:inherited-structure}
Any algebraic relation satisfied by \(\rho_i(\mathcal A)\) is inherited by the
global action. In particular, if \(a,b\in\mathcal A\), then
\[
\rho_\Gamma(ab)=\rho_\Gamma(a)\rho_\Gamma(b),
\qquad
\rho_\Gamma(1)=I_{\Gamma(H)}.
\]
\end{SRWTcorollary}

The result distinguishes a shared \emph{operational} representation from a
shared static code. A core closed under the common algebra supports the same
abstract manipulations across strata without leaving the globally compatible
subspace.

\section{Workspace Operator Algebra}

\begin{SRWTdefinition}[Natural workspace operators]\label{SRWT:def:workspace-operator-algebra}
Define
\[
\boxed{
\mathcal O_{\SRWTWorkspace}
=
\left\{
(A_i)_{i\in V}:
A_{t(e)}H_e=H_eA_{s(e)}\ \forall e
\right\}.
}
\]
Categorically,
\[
\mathcal O_{\SRWTWorkspace}=\SRWTNat(H,H),
\]
the algebra of natural endomorphisms of the local system.
\end{SRWTdefinition}

Pointwise addition, scalar multiplication, and composition make
\(\mathcal O_{\SRWTWorkspace}\) a unital associative algebra. Each
\(A=(A_i)\) acts on global sections by
\[
\pi_{\SRWTWorkspace}(A)(c_i)=(A_ic_i).
\]
Naturality guarantees that \((A_ic_i)\) is again compatible.

\begin{SRWTtheorem}[Commutant characterization]\label{SRWT:thm:commutant-characterization}
Let \(|G|\) be connected and let \(H\) be an invertible local system. Root
evaluation induces an algebra isomorphism
\[
\boxed{
\Phi_o:\mathcal O_{\SRWTWorkspace}
\xrightarrow{\ \cong\ }
\operatorname{Comm}(\SRWTHol_o),
\qquad
\Phi_o((A_i))=A_o,
}
\]
where
\[
\operatorname{Comm}(\SRWTHol_o)
=
\{B\in\SRWTEnd(C_o):Bh=hB\ \forall h\in\SRWTHol_o\}.
\]
\end{SRWTtheorem}

\begin{proof}
If \((A_i)\in\mathcal O_{\SRWTWorkspace}\), naturality along every loop gives
\[
A_oH_\gamma=H_\gamma A_o.
\]
Thus \(A_o\in\operatorname{Comm}(\SRWTHol_o)\), and \(\Phi_o\) is a well-defined
algebra homomorphism.

Choose paths \(\tau_i:o\to i\). Naturality along \(\tau_i\) gives
\[
A_iH_{\tau_i}=H_{\tau_i}A_o,
\]
so
\[
A_i=H_{\tau_i}A_oH_{\tau_i}^{-1}.
\]
Therefore \(A_o\) determines every \(A_i\), and \(\Phi_o\) is injective.

Conversely, take \(B\in\operatorname{Comm}(\SRWTHol_o)\) and define
\[
A_i=H_{\tau_i}BH_{\tau_i}^{-1}.
\]
This is independent of the chosen root path: changing \(\tau_i\) multiplies
it by a loop holonomy that commutes with \(B\). For an edge \(e:i\to j\), the
operator
\[
H_{\tau_j}^{-1}H_eH_{\tau_i}
\]
is a root holonomy. Commutation with \(B\) rearranges to
\[
A_jH_e=H_eA_i.
\]
Hence \((A_i)\) is natural and \(\Phi_o((A_i))=B\). Thus \(\Phi_o\) is
surjective.
\end{proof}

The structural workspace is therefore more accurately represented by the
triple
\[
\boxed{
\SRWTWorkspace_{\mathrm{static}}
=
\bigl(
\Gamma(H),
\mathcal O_{\SRWTWorkspace},
\pi_{\SRWTWorkspace}
\bigr),
}
\]
not merely by the vector space \(\Gamma(H)\).

\subtopic{10.1 Effective operator algebra}

Some natural operators may vanish on every global section while acting
nontrivially on locally available but globally inconsistent directions. Define
the annihilator
\[
\operatorname{Ann}(\Gamma(H))
=
\{A\in\mathcal O_{\SRWTWorkspace}:
\pi_{\SRWTWorkspace}(A)=0\}.
\]
The effective workspace algebra is
\[
\boxed{
\mathcal O_{\SRWTWorkspace}^{\mathrm{eff}}
=
\mathcal O_{\SRWTWorkspace}/\operatorname{Ann}(\Gamma(H)).
}
\]
Its action on \(\Gamma(H)\) is faithful by construction.

\begin{SRWTexample}[Operators in the reflection-cycle workspace]\label{SRWT:ex:reflection-operators}
For Example~\ref{SRWT:ex:reflection-cycle}, the holonomy is
\[
J=\operatorname{diag}(1,-1).
\]
Its commutant consists of diagonal matrices:
\[
\operatorname{Comm}(J)
=
\left\{
\begin{pmatrix}a&0\\0&b\end{pmatrix}:a,b\in\mathbb R
\right\}.
\]
The exact workspace core is the first coordinate. Consequently, only the
parameter \(a\) acts effectively on \(\Gamma(H)\); the \(b\)-direction belongs
to the annihilator of the global core. Hence
\[
\mathcal O_{\SRWTWorkspace}^{\mathrm{eff}}\cong\mathbb R.
\]
\end{SRWTexample}

\subtopic{10.2 Scope of the exact theory}

Chapters~\ref{SRWT:book:chapter-67}--\ref{SRWT:book:chapter-69} now establish three exact facts:
\[
\begin{aligned}
&\text{global shared content is the holonomy-fixed core},\\
&\text{universally transported operations preserve that core},\\
&\text{natural workspace operations are the commutant of holonomy}.
\end{aligned}
\]
These statements do not yet prove that the shared core is robust, causally
used, irreducibly distributed, or temporally persistent. Those are separate
questions developed in the approximate and dynamic theory that follows.

% This file is part of the modular source for Stratified Relational Workspace Theory.
% !TeX root = ../Stratified_Relational_Workspace_Theory.tex

\chapter{Approximate workspaces}\label{SRWT:book:chapter-70}


The exact theory asks whether local core states agree perfectly under every
transfer. Approximate theory asks three more realistic questions:
\begin{enumerate}
\item how strongly a candidate state violates edgewise compatibility;
\item whether a low-defect subspace is spectrally separated from incoherent
 directions; and
\item how that subspace changes when transfers, metrics, or operations are
 perturbed.
\end{enumerate}
Throughout Chapters~\ref{SRWT:book:chapter-70}--\ref{SRWT:book:chapter-71}, set
\[
\mathcal C_V=\bigoplus_{i\in V}C_i,
\qquad
\mathcal C_E=\bigoplus_{e\in E}C_{t(e)},
\]
with their Hilbert direct-sum inner products, and retain strictly positive
weights \(w_e\).

\section{Transfer Discrepancy Operator}

\begin{SRWTdefinition}[Transfer discrepancy]\label{SRWT:def:transfer-discrepancy}
The transfer discrepancy operator is
\[
D_H:\mathcal C_V\longrightarrow\mathcal C_E,
\qquad
(D_Hc)_e
=
\sqrt{w_e}\bigl(c_{t(e)}-H_ec_{s(e)}\bigr).
\]
Its squared norm is the coherence energy
\[
\boxed{
\mathcal E_H(c)
=
\|D_Hc\|^2
=
\sum_{e\in E}w_e
\left\|c_{t(e)}-H_ec_{s(e)}\right\|^2.
}
\]
The transfer Laplacian is
\[
\boxed{L_H=D_H^*D_H.}
\]
\end{SRWTdefinition}

The operator \(D_H\) is an incidence-type derivative twisted by the edge
transfers. It compares a state already present at the target with the state
predicted there by translation from the source. A direction can therefore
have low energy even when its coordinates differ substantially across nodes,
provided those coordinates differ exactly as the transfer maps prescribe.

For \(z=(z_e)_{e\in E}\in\mathcal C_E\), the adjoint is
\[
(D_H^*z)_i
=
\sum_{e:t(e)=i}\sqrt{w_e}\,z_e
-
\sum_{e:s(e)=i}\sqrt{w_e}\,H_e^*z_e.
\]
Consequently,
\[
\boxed{
\begin{aligned}
(L_Hc)_i
={}&
\sum_{e:t(e)=i}w_e\bigl(c_i-H_ec_{s(e)}\bigr)\\
&+
\sum_{e:s(e)=i}w_eH_e^*
\bigl(H_ec_i-c_{t(e)}\bigr).
\end{aligned}
}
\]
This formula makes the two contributions at a vertex explicit: incoming
prediction error and adjoint-propagated outgoing error.

\begin{SRWTproposition}[Positivity and variational identity]\label{SRWT:prop:laplacian-positive}
The transfer Laplacian is self-adjoint and positive semidefinite. For every
\(c\in\mathcal C_V\),
\[
\langle c,L_Hc\rangle=\mathcal E_H(c)\geq0.
\]
\end{SRWTproposition}

\begin{proof}
Self-adjointness and positivity follow from \(L_H=D_H^*D_H\). The displayed
identity is
\[
\langle c,D_H^*D_Hc\rangle
=
\langle D_Hc,D_Hc\rangle.
\]
\end{proof}

\begin{SRWTremark}[Weight dependence]
Replacing \(w_e\) by another positive family does not alter \(\ker D_H\), but
it generally changes every positive eigenvalue and therefore changes the
approximate geometry. Edge weights are part of the metric model, not part of
the exact compatibility relation.
\end{SRWTremark}

\section{Exact Kernel Theorem}

\begin{SRWTtheorem}[Transfer-Laplacian kernel theorem]\label{SRWT:thm:laplacian-kernel}
For an arbitrary finite directed core diagram, without assuming invertibility,
\[
\boxed{
\ker L_H=\ker D_H=\Gamma(H).
}
\]
Thus the exact compatible sections are precisely the zero-energy modes of the
transfer Laplacian.
\end{SRWTtheorem}

\begin{proof}
If \(D_Hc=0\), then \(L_Hc=D_H^*D_Hc=0\). Conversely, if \(L_Hc=0\), then
\[
0=\langle c,L_Hc\rangle=\|D_Hc\|^2,
\]
so \(D_Hc=0\). Finally, \(D_Hc=0\) if and only if
\[
c_{t(e)}=H_ec_{s(e)}
\]
for every edge, which is exactly the definition of \(\Gamma(H)\).
\end{proof}

\begin{SRWTcorollary}[Finite rank computation]\label{SRWT:cor:workspace-rank-D}
Let
\[
n_C=\sum_{i\in V}\dim C_i.
\]
Then
\[
\boxed{
\dim\Gamma(H)=n_C-\operatorname{rank}(D_H).
}
\]
Hence the exact workspace dimension can be obtained by a single rank or
singular-value computation, even when holonomy is unavailable because some
transfers are noninvertible.
\end{SRWTcorollary}

\subtopic{12.1 Relation to the holonomy computation}

In the local-system regime, Theorems~\ref{SRWT:thm:limit-holonomy} and
\ref{SRWT:thm:laplacian-kernel} give two equivalent constructions:
\[
\Gamma(H)
\cong
C_o^{\SRWTHol_o}
\qquad\text{and}\qquad
\Gamma(H)=\ker L_H.
\]
The first reduces the problem to cycle invariants in a root fiber. The second
works directly on all vertex coordinates and extends to noninvertible diagrams.
The two computations therefore serve different practical regimes rather than
competing definitions.

\section{Approximate Structural Workspace}

Let
\[
0\leq\lambda_1\leq\lambda_2\leq\cdots\leq\lambda_N,
\qquad
N=\dim\mathcal C_V,
\]
be the eigenvalues of \(L_H\), repeated with multiplicity, and let
\(P_I(L_H)\) denote the spectral projector associated with an interval \(I\).

\begin{SRWTdefinition}[Threshold workspace]\label{SRWT:def:threshold-workspace}
For \(\varepsilon\geq0\), define
\[
\boxed{
\Gamma_\varepsilon(H)
=
\SRWTran P_{[0,\varepsilon]}(L_H).
}
\]
Its dimension is
\[
d_\varepsilon
=
\#\{k:\lambda_k\leq\varepsilon\}.
\]
When \(0<d_\varepsilon<N\), define the boundary gap
\[
\boxed{
g_\varepsilon
=
\lambda_{d_\varepsilon+1}-\lambda_{d_\varepsilon}.
}
\]
\end{SRWTdefinition}

The threshold \(\varepsilon\) declares which discrepancy scale is treated as
structurally negligible. The dimension \(d_\varepsilon\) alone is not a
robust certificate: a cluster is stable only when it is separated by a
positive boundary gap.

An equivalent index-based description fixes \(d\) and writes
\[
W_d=\SRWTran P_d,
\qquad
P_d=P_{\{\lambda_1,\ldots,\lambda_d\}}(L_H),
\]
with cluster gap
\[
g_d=\lambda_{d+1}-\lambda_d>0.
\]
The index-based form is convenient for perturbation theorems because the
workspace rank is declared explicitly.

\begin{SRWTproposition}[Min--max characterization]\label{SRWT:prop:minmax}
For \(1\leq k\leq N\),
\[
\lambda_k
=
\min_{\substack{S\subseteq\mathcal C_V\\\dim S=k}}
\ \max_{\substack{c\in S\\\|c\|=1}}
\mathcal E_H(c).
\]
Consequently, the span of the first \(k\) eigenvectors minimizes the worst-case
coherence energy among all \(k\)-dimensional candidate subspaces.
\end{SRWTproposition}

\begin{proof}
Expand every vector in an orthonormal eigenbasis of \(L_H\). The span of the
first \(k\) eigenvectors has maximal Rayleigh quotient \(\lambda_k\). Every
\(k\)-dimensional subspace intersects the orthogonal complement of the first
\(k-1\) eigenvectors nontrivially, and any unit vector in that intersection has
Rayleigh quotient at least \(\lambda_k\).
\end{proof}

\subtopic{13.1 Exact, soft, and spurious low modes}

The spectrum supports three distinct interpretations:
\[
\begin{array}{rcl}
\lambda_k=0
&\Longleftrightarrow&
\text{exact compatible direction},\\
0<\lambda_k\ll g_d
&\Longrightarrow&
\text{approximately transported direction},\\
\lambda_k\approx\lambda_{k+1}
&\Longrightarrow&
\text{rank selection is unstable at that boundary}.
\end{array}
\]
A low eigenvalue can still be scientifically spurious if it arises from a
poorly chosen weight, a nearly disconnected graph, or a noncausal encoding.
The causal and null gates in later parts remain necessary.

\section{Normalized Coherence Defect}

For \(c\neq0\), define the Rayleigh defect
\[
\delta_{\mathrm{coh}}(c)
=
\frac{\mathcal E_H(c)}{\|c\|^2}.
\]
For a nonzero subspace \(W\subseteq\mathcal C_V\), define
\[
\delta_{\mathrm{coh}}(W)
=
\sup_{\substack{c\in W\\\|c\|=1}}
\mathcal E_H(c).
\]
If \(W\) is \(L_H\)-invariant, then
\[
\delta_{\mathrm{coh}}(W)
=
\|L_H|_W\|.
\]
In particular,
\[
\delta_{\mathrm{coh}}(\Gamma_\varepsilon(H))
\leq\varepsilon.
\]

\begin{SRWTtheorem}[Coherence-to-core distance estimate]\label{SRWT:thm:poincare-workspace}
Let
\[
W_0=\Gamma(H)=\ker L_H,
\qquad
r=\dim W_0,
\]
and suppose \(r<N\). Set
\[
\gamma_H=\lambda_{r+1}>0.
\]
Then every \(c\in\mathcal C_V\) satisfies
\[
\boxed{
\operatorname{dist}(c,W_0)^2
\leq
\frac{\mathcal E_H(c)}{\gamma_H}.
}
\]
Moreover,
\[
\mathcal E_H(c)
\leq
\|L_H\|\operatorname{dist}(c,W_0)^2.
\]
\end{SRWTtheorem}

\begin{proof}
Write \(c=P_0c+Q_0c\), where \(P_0\) projects onto \(W_0\) and
\(Q_0=I-P_0\). Since \(L_HP_0=0\),
\[
\mathcal E_H(c)
=
\langle Q_0c,L_HQ_0c\rangle.
\]
The spectrum of \(L_H\) on \(W_0^\perp\) lies in
\([\gamma_H,\|L_H\|]\). Hence
\[
\gamma_H\|Q_0c\|^2
\leq
\mathcal E_H(c)
\leq
\|L_H\|\|Q_0c\|^2.
\]
Finally, \(\|Q_0c\|=\operatorname{dist}(c,W_0)\).
\end{proof}

This estimate is the exact sense in which a spectral gap turns small edgewise
disagreement into closeness to a genuinely compatible global section.

\subtopic{14.1 Two-node consistency cycle}

Let \(C_1=C_2=\mathbb R\), use unit weights, and take two directed transfers
\[
H_{12}=a,
\qquad
H_{21}=b.
\]
Then
\[
D_H
=
\begin{pmatrix}
-a&1\\
1&-b
\end{pmatrix},
\qquad
L_H
=
\begin{pmatrix}
a^2+1&-(a+b)\\
-(a+b)&b^2+1
\end{pmatrix}.
\]
Writing
\[
\tau=a^2+b^2+2,
\]
the eigenvalues are
\[
\lambda_\pm
=
\frac{1}{2}
\left(
\tau\pm\sqrt{\tau^2-4(ab-1)^2}
\right),
\]
and
\[
\det L_H=(ab-1)^2.
\]
Thus a nonzero exact shared core exists precisely when
\[
ab=1,
\]
which is exactly the trivial-holonomy condition. Away from exact consistency,
\[
\lambda_-\lambda_+=(ab-1)^2.
\]
When \(|ab-1|\) is small while \(\lambda_+\) stays bounded away from zero, the
soft workspace eigenvalue is quadratic in the cycle mismatch:
\[
\lambda_-
=
\frac{(ab-1)^2}{\lambda_+}.
\]
This example connects approximate spectral coherence directly to approximate
holonomy closure.

\section{Perturbation Stability Theorem}

The central robustness question is whether a selected low-energy subspace
persists when the estimated transfer maps change. The next lemma isolates the
finite-dimensional spectral fact needed for the answer.

\begin{SRWTlemma}[Separated spectral subspaces]\label{SRWT:lem:separated-subspaces}
Let \(A\) and \(B=A+E\) be self-adjoint operators. Let \(P\) be an
\(A\)-spectral projector whose eigenvalues lie in \(( -\infty,a]\), and let
\(\widetilde Q\) be a \(B\)-spectral projector whose eigenvalues lie in
\([b,\infty)\), with \(b>a\). Then
\[
\|P\widetilde Q\|
\leq
\frac{\|E\|}{b-a}.
\]
\end{SRWTlemma}

\begin{proof}
On \(\SRWTran P\) and \(\SRWTran\widetilde Q\), set
\[
A_P=PAP|_{\SRWTran P},
\qquad
B_Q=\widetilde QB\widetilde Q|_{\SRWTran\widetilde Q}.
\]
The operator \(X=P\widetilde Q\) satisfies the Sylvester equation
\[
A_PX-XB_Q=-PE\widetilde Q.
\]
Let \(\mu=(a+b)/2\). Since
\[
\|e^{t(A_P-\mu I)}\|\leq e^{-(b-a)t/2},
\qquad
\|e^{-t(B_Q-\mu I)}\|\leq e^{-(b-a)t/2},
\]
the inverse Sylvester operator has the convergent integral representation
\[
X
=
\int_0^\infty
 e^{t(A_P-\mu I)}
 (PE\widetilde Q)
 e^{-t(B_Q-\mu I)}\,\SRWTdd t.
\]
Taking norms and integrating gives the result.
\end{proof}

\begin{SRWTtheorem}[Spectral workspace stability]\label{SRWT:thm:spectral-stability}
Let \(L\) and \(\widetilde L=L+E\) be self-adjoint positive semidefinite
operators on the same \(N\)-dimensional space. Suppose
\[
\lambda_d(L)<\lambda_{d+1}(L),
\qquad
g_d=\lambda_{d+1}(L)-\lambda_d(L)>0,
\]
and
\[
\eta=\|E\|<\frac{g_d}{2}.
\]
Let \(P_d\) and \(\widetilde P_d\) project onto the first \(d\) eigenvectors
of \(L\) and \(\widetilde L\), respectively. Then the rank-\(d\) cluster
remains separated and
\[
\boxed{
\|P_d-\widetilde P_d\|
\leq
\frac{\eta}{g_d-\eta}
\leq
\frac{2\eta}{g_d}.
}
\]
\end{SRWTtheorem}

\begin{proof}
The min--max characterization gives the eigenvalue perturbation bound
\[
|\lambda_k(\widetilde L)-\lambda_k(L)|\leq\eta.
\]
Hence
\[
\lambda_{d+1}(\widetilde L)-\lambda_d(L)
\geq g_d-\eta
\]
and
\[
\lambda_{d+1}(L)-\lambda_d(\widetilde L)
\geq g_d-\eta.
\]
Apply Lemma~\ref{SRWT:lem:separated-subspaces} first to
\(P_d(I-\widetilde P_d)\) and then, with the roles of \(L\) and
\(\widetilde L\) exchanged, to \((I-P_d)\widetilde P_d\). Orthogonal
projections of equal rank satisfy
\[
\|P_d-\widetilde P_d\|
=
\max\left\{
\|P_d(I-\widetilde P_d)\|,
\|(I-P_d)\widetilde P_d\|
\right\}.
\]
This yields the first bound. The second follows from \(\eta<g_d/2\).
\end{proof}

\subtopic{15.1 From transfer-map error to Laplacian error}

Suppose the edge maps are perturbed to
\[
\widetilde H_e=H_e+\Delta_e
\]
while the weights and vertex metrics remain fixed. Define
\[
\boxed{
\delta_H^2
=
\max_{i\in V}
\sum_{e:s(e)=i}w_e\|\Delta_e\|^2.
}
\]
Let \(F=D_{\widetilde H}-D_H\). Then
\[
(Fc)_e=-\sqrt{w_e}\,\Delta_ec_{s(e)},
\qquad
\|F\|\leq\delta_H.
\]
Since
\[
L_{\widetilde H}-L_H
=D_H^*F+F^*D_H+F^*F,
\]
we obtain
\[
\boxed{
\|L_{\widetilde H}-L_H\|
\leq
2\|D_H\|\delta_H+\delta_H^2.
}
\]

\begin{SRWTcorollary}[Edgewise stability certificate]\label{SRWT:cor:edge-stability}
If
\[
\eta_H=2\|D_H\|\delta_H+\delta_H^2
<\frac{g_d}{2},
\]
then
\[
\boxed{
\|P_d(H)-P_d(\widetilde H)\|
\leq
\frac{\eta_H}{g_d-\eta_H}.
}
\]
\end{SRWTcorollary}

This converts declared edge-map uncertainty into an explicit bound on the
motion of the extracted workspace subspace.

\subtopic{15.2 Approximate operation closure}

Define the vertex and edge direct-sum representations
\[
\rho_V(a)=\bigoplus_{i\in V}\rho_i(a),
\qquad
\rho_E(a)=\bigoplus_{e\in E}\rho_{t(e)}(a).
\]
Lift the edgewise intertwining defects to
\[
(\mathfrak X_a c)_e
=
\sqrt{w_e}\,\Xi_e(a)c_{s(e)}.
\]
Then the exact commutator identity is
\[
\boxed{
D_H\rho_V(a)-\rho_E(a)D_H=-\mathfrak X_a.
}
\]
Furthermore,
\[
\|\mathfrak X_a\|
\leq
\beta(a),
\qquad
\beta(a)^2
=
\max_{i\in V}
\sum_{e:s(e)=i}w_e\|\Xi_e(a)\|^2.
\]

\begin{SRWTtheorem}[Approximate operation-closure theorem]\label{SRWT:thm:approx-operation-closure}
Let \(P_d\) project onto the first \(d\) eigenvectors of \(L_H\), set
\[
\lambda_d\leq\varepsilon,
\qquad
\mu=\lambda_{d+1}>0,
\]
and let \(Q_d=I-P_d\). Then every \(a\in\mathcal A\) satisfies
\[
\boxed{
\|Q_d\rho_V(a)P_d\|
\leq
\frac{
\|\rho_E(a)\|\sqrt{\varepsilon}+\beta(a)
}{\sqrt{\mu}}.
}
\]
\end{SRWTtheorem}

\begin{proof}
For \(x\in\SRWTran P_d\),
\[
\|D_Hx\|^2
=\langle x,L_Hx\rangle
\leq\varepsilon\|x\|^2.
\]
The commutator identity gives
\[
\|D_H\rho_V(a)x\|
\leq
\|\rho_E(a)\|\sqrt{\varepsilon}\|x\|
+
\beta(a)\|x\|.
\]
Because \(P_d\) and \(Q_d\) reduce \(L_H\), the two energy components are
orthogonal and
\[
\|D_HQ_dy\|\leq\|D_Hy\|.
\]
On \(\SRWTran Q_d\), the spectrum of \(L_H\) is bounded below by \(\mu\), so
\[
\sqrt{\mu}\|Q_dy\|
\leq
\|D_HQ_dy\|.
\]
Apply these inequalities to \(y=\rho_V(a)x\), then take the supremum over unit
\(x\in\SRWTran P_d\).
\end{proof}

\begin{SRWTcorollary}[Recovery of exact operation closure]
If \(d=\dim\ker L_H\), \(\varepsilon=0\), and every edge intertwines \(a\)
exactly, then
\[
Q_d\rho_V(a)P_d=0.
\]
Thus Theorem~\ref{SRWT:thm:approx-operation-closure} reduces to the exact
operation-closure result of Chapter~\ref{SRWT:book:chapter-69}.
\end{SRWTcorollary}

\begin{SRWTcorollary}[Closure under projector drift]\label{SRWT:cor:closure-drift}
Let \(\widetilde P_d\) be another rank-\(d\) orthogonal projector and set
\[
\delta_P=\|P_d-\widetilde P_d\|.
\]
Then
\[
\boxed{
\|(I-\widetilde P_d)\rho_V(a)\widetilde P_d\|
\leq
\|Q_d\rho_V(a)P_d\|
+2\|\rho_V(a)\|\delta_P.
}
\]
Combining this estimate with Theorem~\ref{SRWT:thm:spectral-stability} controls
operation leakage after noisy workspace extraction.
\end{SRWTcorollary}

\begin{proof}
Subtract \(Q_d\rho_V(a)P_d\) and use
\[
I-\widetilde P_d-Q_d=P_d-\widetilde P_d.
\]
The two resulting perturbation terms are each bounded by
\(\|\rho_V(a)\|\delta_P\).
\end{proof}

\subtopic{15.3 Core--residual path accumulation}

Let a path have vertices
\[
i_0\xrightarrow{e_1}i_1\xrightarrow{e_2}\cdots
\xrightarrow{e_m}i_m.
\]
Write
\[
T_r=T_{e_r},
\quad
H_r=H_{e_r},
\quad
V_r=V_{e_r},
\]
and let
\[
\iota_i^C:C_i\hookrightarrow X_i,
\qquad
\iota_i^R:R_i\hookrightarrow X_i
\]
be the canonical inclusions. Empty products below are identities.

\begin{SRWTtheorem}[Pathwise core--residual Duhamel formula]\label{SRWT:thm:residual-duhamel}
The deviation of full path transport from pure core transport is exactly
\[
\boxed{
\begin{aligned}
&T_m\cdots T_1\iota_{i_0}^C
-
\iota_{i_m}^C H_m\cdots H_1\\
&\qquad=
\sum_{r=1}^{m}
T_m\cdots T_{r+1}
\iota_{i_r}^R V_r
H_{r-1}\cdots H_1.
\end{aligned}
}
\]
Consequently,
\[
\boxed{
\begin{aligned}
\|Q_{i_m}T_m\cdots T_1P_{i_0}\|
\leq
\sum_{r=1}^{m}
&\left(\prod_{\ell=r+1}^{m}\|T_\ell\|\right)
\|V_r\|\\
&\times
\left(\prod_{\ell=1}^{r-1}\|H_\ell\|\right),
\end{aligned}
}
\]
and the same right-hand side bounds
\[
\|P_{i_m}T_m\cdots T_1P_{i_0}-H_m\cdots H_1\|.
\]
\end{SRWTtheorem}

\begin{proof}
For every step,
\[
T_r\iota_{i_{r-1}}^C
=
\iota_{i_r}^CH_r+
\iota_{i_r}^RV_r.
\]
Telescope the sequence
\[
T_m\cdots T_{r+1}\iota_{i_r}^C H_r\cdots H_1,
\qquad
r=0,1,\ldots,m.
\]
The difference between consecutive terms is exactly the \(r\)-th summand in
the displayed formula. Projection onto \(R_{i_m}\) or \(C_{i_m}\), followed
by submultiplicativity of the operator norm, gives the two bounds.
\end{proof}

\begin{SRWTcorollary}[Uniform path bound]\label{SRWT:cor:uniform-path-bound}
If along the path
\[
\|T_r\|\leq\tau,
\qquad
\|H_r\|\leq h,
\qquad
\|V_r\|\leq\nu,
\]
then the leakage and core-transport defects are bounded by
\[
\boxed{
\nu\sum_{r=1}^{m}\tau^{m-r}h^{r-1}
=
\begin{cases}
\displaystyle
\nu\frac{\tau^m-h^m}{\tau-h},&\tau\neq h,\\[1.1em]
\displaystyle
m\nu\tau^{m-1},&\tau=h.
\end{cases}
}
\]
In particular, if \(\tau\leq1\) and \(h\leq1\), the bound is at most
\(m\nu\).
\end{SRWTcorollary}

The linear factor \(m\) in the merely nonexpansive regime is important. Small
single-edge leakage does not by itself justify arbitrarily long transfer
chains. Long-horizon stability requires contraction, cancellation, error
correction, or leakage that decreases with path length.

\subtopic{15.4 Residual-mediated reentry}

For two consecutive transfers,
\[
P_{i_2}T_2T_1P_{i_0}
=
H_2H_1+U_2V_1,
\]
\[
Q_{i_2}T_2T_1P_{i_0}
=
V_2H_1+K_2V_1.
\]
Thus leakage out of the core is first order in the blocks \(V_r\), while an
excursion that leaves and then returns to the core contains at least one
\(V\)-block and one later \(U\)-block. For a fixed path length, if
\[
\|U_r\|,\|V_r\|=O(\delta)
\]
and the diagonal blocks remain bounded, the residual-mediated correction to
pure core-to-core transport is \(O(\delta^2)\), whereas terminal residual
leakage can be \(O(\delta)\). This distinction separates loss of core
invariance from distortion caused by temporary excursions through local
residual degrees of freedom.

% This file is part of the modular source for Stratified Relational Workspace Theory.
% !TeX root = ../Stratified_Relational_Workspace_Theory.tex

\chapter{Gauge invariance and abstract identity}\label{SRWT:book:chapter-71}


Exact compatibility is a coordinate-free notion. Approximate spectra also
depend on the Hilbert metrics used to measure discrepancy. Gauge theory must
therefore distinguish a change of coordinates from a physical change of
metric.

\section{Local Coordinate Transformations}

For each \(i\in V\), let
\[
g_i\in GL(C_i),
\qquad
g=(g_i)_{i\in V}.
\]
Define the vertex and edge direct sums
\[
G_V=\bigoplus_{i\in V}g_i,
\qquad
G_E=\bigoplus_{e\in E}g_{t(e)}.
\]
The transformed transfers and local operation representations are
\[
H_e^g
=
 g_{t(e)}H_eg_{s(e)}^{-1},
\]
\[
\rho_i^g(a)
=
 g_i\rho_i(a)g_i^{-1}.
\]
The fundamental covariance identity is
\[
\boxed{
D_{H^g}G_V=G_ED_H.
}
\]
It follows directly by evaluating the \(e\)-component:
\[
\sqrt{w_e}
\left(
 g_{t(e)}c_{t(e)}
-
 g_{t(e)}H_ec_{s(e)}
\right).
\]

\section{Gauge Invariance Theorem}

\begin{SRWTtheorem}[Exact gauge covariance]\label{SRWT:thm:gauge-covariance}
The map
\[
G_V:\Gamma(H)\longrightarrow\Gamma(H^g)
\]
is an isomorphism. Therefore,
\[
\boxed{
\Gamma(H^g)=G_V\Gamma(H),
\qquad
\dim\Gamma(H^g)=\dim\Gamma(H).
}
\]
If the core transports are invertible and \(o\) is a root, then
\[
H_\gamma^g=g_oH_\gamma g_o^{-1}
\]
for every loop \(\gamma:o\to o\). Hence the holonomy-fixed core transforms by
\[
C_*^g=g_oC_*.
\]
\end{SRWTtheorem}

\begin{proof}
The covariance identity implies
\[
D_{H^g}(G_Vc)=G_ED_Hc.
\]
Since \(G_E\) is invertible,
\[
D_Hc=0
\quad\Longleftrightarrow\quad
D_{H^g}(G_Vc)=0.
\]
The holonomy formula follows by cancellation of adjacent factors
\(g_i^{-1}g_i\) along a path.
\end{proof}

\begin{SRWTtheorem}[Unitary gauge covariance of the Laplacian]\label{SRWT:thm:unitary-gauge}
If every \(g_i\) is unitary, then \(G_V\) and \(G_E\) are unitary and
\[
\boxed{
L_{H^g}=G_VL_HG_V^{-1}.
}
\]
Consequently,
\[
\SRWTspec(L_{H^g})=\SRWTspec(L_H),
\qquad
P_I(L_{H^g})=G_VP_I(L_H)G_V^{-1}
\]
for every spectral interval \(I\).
\end{SRWTtheorem}

\begin{proof}
From
\[
D_{H^g}=G_ED_HG_V^{-1}
\]
and unitarity,
\[
D_{H^g}^*D_{H^g}
=
G_VD_H^*G_E^{-1}G_ED_HG_V^{-1}
=
G_VL_HG_V^{-1}.
\]
The spectral statements follow from unitary similarity.
\end{proof}

\subtopic{17.1 Nonunitary gauges and metric distortion}

For a general invertible gauge, define
\[
m_V=\sigma_{\min}(G_V),
\qquad
M_V=\|G_V\|,
\]
\[
m_E=\sigma_{\min}(G_E),
\qquad
M_E=\|G_E\|.
\]
The covariance identity yields
\[
m_E^2\mathcal E_H(c)
\leq
\mathcal E_{H^g}(G_Vc)
\leq
M_E^2\mathcal E_H(c),
\]
and
\[
m_V^2\|c\|^2
\leq
\|G_Vc\|^2
\leq
M_V^2\|c\|^2.
\]
Therefore the Rayleigh quotients obey
\[
\boxed{
\frac{m_E^2}{M_V^2}
\frac{\mathcal E_H(c)}{\|c\|^2}
\leq
\frac{\mathcal E_{H^g}(G_Vc)}{\|G_Vc\|^2}
\leq
\frac{M_E^2}{m_V^2}
\frac{\mathcal E_H(c)}{\|c\|^2}.
}
\]
By the min--max principle,
\[
\boxed{
\frac{m_E^2}{M_V^2}\lambda_k(L_H)
\leq
\lambda_k(L_{H^g})
\leq
\frac{M_E^2}{m_V^2}\lambda_k(L_H).
}
\]
Exact zero modes are unchanged, but fixed-metric approximate eigenvalues can be
strongly distorted by an ill-conditioned coordinate transformation.

\subtopic{17.2 Metric-covariant formulation}

A pure change of coordinates should transform the metric as well as the
matrices. On each transformed space, define
\[
\langle u,v\rangle_i^g
=
\left\langle g_i^{-1}u,g_i^{-1}v\right\rangle_i.
\]
Use the analogous pullback metric on every edge copy. Then \(G_V\) and
\(G_E\) are unitary between the original and transformed metric spaces.
Consequently, the complete transfer-Laplacian spectrum is exactly preserved.

This gives two distinct interpretations:
\begin{itemize}
\item Transforming the matrices and the metrics gives the same approximate
workspace in new coordinates.
\item Transforming the matrices while holding the metrics fixed gives a
potentially different weighting of discrepancy.
\end{itemize}
The distinction prevents arbitrary coordinate scaling from manufacturing or
erasing apparent low-energy workspace directions.

\section{Static Workspace Equivalence}

\begin{SRWTdefinition}[Exact static equivalence]\label{SRWT:def:static-equivalence}
Two static realizations are exactly gauge equivalent,
\[
\mathfrak R\sim_{\mathrm{stat}}\mathfrak R',
\]
when there exists an invertible family \(g_i\) such that
\[
H_e'=g_{t(e)}H_eg_{s(e)}^{-1},
\qquad
\rho_i'(a)=g_i\rho_i(a)g_i^{-1}.
\]
The exact static abstract workspace is
\[
\boxed{
\SRWTWorkspace_{\mathrm{stat}}
=
[\mathfrak R]_{\sim_{\mathrm{stat}}}.
}
\]
\end{SRWTdefinition}

For approximate theory, the realization also contains the vertex and edge
metrics. Two metric realizations are equivalent only when the gauge transports
those metrics, or when the gauge is unitary relative to the fixed metrics.
Accordingly:
\[
\boxed{
\begin{aligned}
\text{exact workspace identity}
&\text{ is invariant under }GL\text{-gauge},\\
\text{fixed-metric spectral geometry}
&\text{ is invariant under unitary gauge},\\
\text{metric-covariant spectral geometry}
&\text{ is invariant under transported }GL\text{-gauge}.
\end{aligned}
}
\]


\subtopic{18.1 What the approximate theory establishes}

Chapters~\ref{SRWT:book:chapter-70}--\ref{SRWT:book:chapter-71} now prove:
\[
\begin{aligned}
&\ker L_H=\Gamma(H),\\
&\text{low energy plus a gap implies proximity to an exact core},\\
&\text{a separated workspace cluster is stable under transfer noise},\\
&\text{approximately intertwined operations nearly preserve that cluster},\\
&\text{core leakage and residual reentry admit explicit pathwise bounds},\\
&\text{exact identity is gauge invariant while approximate geometry is metric aware}.
\end{aligned}
\]
These are structural and perturbative results. They still do not prove that a
workspace variable is causally used, predictively sufficient, irreducibly
distributed, or temporally self-maintaining. Those questions begin in
Chapter~\ref{SRWT:book:chapter-72}.

\subtopic{18.2 Finite computation protocol}

For a declared graph, set of core dimensions, weights, and transfer maps, the
static analysis can be executed in the following order:
\begin{enumerate}
\item Assemble the sparse block operator \(D_H\) from the edge equations
\[
\sqrt{w_e}\bigl(c_{t(e)}-H_ec_{s(e)}\bigr).
\]
\item Compute a singular-value decomposition
\[
D_H=U\Sigma V^*,
\qquad
\lambda_k(L_H)=\sigma_k(D_H)^2.
\]
Use a declared numerical tolerance to distinguish computational zero modes
from positive modes.
\item Select a candidate rank \(d\) only at a reproducible spectral boundary
with \(g_d>0\). A large observed gap is evidence of stability, not evidence of
causal relevance.
\item Quantify transfer uncertainty by \(\delta_H\) and verify whether
\[
2\|D_H\|\delta_H+\delta_H^2<\frac{g_d}{2}.
\]
If this gate fails, the extracted projector has no certified perturbative
identity at the stated error scale.
\item For every proposed common operation, compute \(\beta(a)\) and the
operation-leakage bound in Theorem~\ref{SRWT:thm:approx-operation-closure}.
\item For every long transfer chain used in an interpretation, report the
pathwise leakage certificate from Theorem~\ref{SRWT:thm:residual-duhamel} rather
than extrapolating from single-edge defects.
\end{enumerate}

A minimal structural report is therefore
\[
\boxed{
\left(
 d,\lambda_d,\lambda_{d+1},g_d,
 \delta_H,
 \|P_d-\widetilde P_d\|_{\mathrm{bound}},
 \beta(a),
 \Lambda_\gamma
\right),
}
\]
where \(\Lambda_\gamma\) denotes the right-hand side of the pathwise leakage
bound. This structural report is intentionally incomplete as a workspace
certificate. Chapters~\ref{SRWT:book:chapter-72}--\ref{SRWT:book:chapter-75} add the predictive, causal, null, and irreducibility
gates required to interpret the low-energy subspace operationally.

% This file is part of the modular source for Stratified Relational Workspace Theory.
% !TeX root = ../Stratified_Relational_Workspace_Theory.tex

\chapter{Dynamic workspace theory}\label{SRWT:book:chapter-72}


\section{Probabilistic Local and Relational Dynamics}

The static theory identifies compatible directions across a declared transfer
network. A dynamic workspace claim requires additional probabilistic objects:
a history, a target future, a decision problem, and an intervention family.
These must be fixed before sufficiency or causal efficacy is evaluated.

\begin{SRWTdefinition}[Dynamic relational realization]\label{SRWT:def:dynamic-realization}
A dynamic relational realization is defined on a probability space
\((\Omega,\mathcal F,\mathbb P)\) and consists of:
\begin{itemize}
\item node states \(X_i(t)\in\mathsf X_i\) for \(i\in V\);
\item edge messages \(M_e(t)\in\mathsf M_e\) for \(e\in E\);
\item external inputs \(E_t\), local actions \(A_i(t)\), and disturbances
 \(\xi_i(t),\zeta_e(t)\);
\item time-dependent transfer data \(H_e(t)\), local core vectors \(c_i(t)\),
 and transfer Laplacians \(L_t=L_{H(t)}\); and
\item a future target \(Y_t^+\in\mathsf Y\).
\end{itemize}
All state, history, and target spaces are assumed standard Borel unless a
stronger topology is explicitly invoked.
\end{SRWTdefinition}

A representative local update is
\[
X_i(t+1)
=
F_i\!\left(
X_i(t),
\{M_e(t):t(e)=i\},
A_i(t),E_t,\xi_i(t)
\right),
\]
with messages
\[
M_e(t)
=
K_e\!\left(
X_{s(e)}(t),
W_t,
E_t,\zeta_e(t)
\right).
\]
This notation permits, but does not require, the current workspace state to
modulate future transfers.

For a coalition \(A\subseteq V\), let \(\mathbf H_t^A\) denote the random
history available to that coalition through time \(t\). Its generated
information is
\[
\mathcal F_t^A=\sigma(\mathbf H_t^A).
\]
The full relational history is \(\mathbf H_t=\mathbf H_t^V\), with
\[
\mathcal F_t^A\subseteq\mathcal F_t^B
\qquad
\text{whenever }A\subseteq B.
\]
The empty-coalition history \(\mathbf H_t^\varnothing\) is the trivial
constant statistic and represents prediction from the background target law
alone.
A concrete history may contain node trajectories, incident messages, actions,
artifacts, transfer-map estimates, and the declared graph history. The exact
contents must be reported because coalition sufficiency is relative to what a
coalition is permitted to observe.

The target \(Y_t^+\) can be a next output, a finite future trajectory, a
research decision, a proof outcome, or any measurable downstream quantity.
Every predictive workspace claim is therefore target-relative:
\[
\operatorname{WorkspaceFor}(Y_t^+)
\centernot\Longrightarrow
\operatorname{WorkspaceFor}(\widetilde Y_t^+).
\]

Let
\[
c(t)=(c_i(t))_{i\in V}
\in
\mathcal C_V(t),
\qquad
P_t=\mathbf 1_{[0,\varepsilon_t]}(L_t).
\]
The instantaneous structural workspace coordinate is
\[
\boxed{
W_t^{\mathrm{str}}=P_tc(t)
\in
\Gamma_{\varepsilon_t}(H(t)).
}
\]
When \(P_t\) varies with time, these coordinates inhabit a moving family of
subspaces. Chapters~\ref{SRWT:book:chapter-76}--\ref{SRWT:book:chapter-78} will supply the transport required to compare them
across time. In what follows, all sufficiency and intervention statements
are made at a fixed decision time \(t\).

\begin{SRWTremark}[Three objects that must not be conflated]
The realization state \(S_t\), the complete history \(\mathbf H_t\), and the
structural coordinate \(W_t^{\mathrm{str}}\) are different objects. The first
is an instantaneous physical or computational state, the second contains all
admissible past information, and the third is a low-transfer-energy statistic
extracted from the declared core geometry.
\end{SRWTremark}

\section{Canonical Predictive and Interventional Causal States}

The quotient notation
\([\mathbf H_t]_{\sim_Y}\) expresses the right intuition but hides
measurability and null-set issues. A canonical construction is obtained from
regular conditional distributions.

\begin{SRWTdefinition}[Predictive kernel and canonical state]\label{SRWT:def:predictive-kernel}
Let \(\mathsf H_t\) and \(\mathsf Y\) be standard Borel spaces. A regular
conditional kernel is a measurable map
\[
K_t:\mathsf H_t\times\mathcal B(\mathsf Y)\to[0,1]
\]
satisfying
\[
K_t(h,A)
=
\mathbb P(Y_t^+\in A\mid\mathbf H_t=h)
\]
for \(\mathbb P_{\mathbf H_t}\)-almost every \(h\). Let
\(\mathcal P(\mathsf Y)\) carry its evaluation sigma-algebra. The canonical
predictive state is the random probability measure
\[
\boxed{
Z_t
=
K_t(\mathbf H_t,\cdot)
\in
\mathcal P(\mathsf Y).
}
\]
Two histories are predictively equivalent precisely when they induce the same
conditional law of \(Y_t^+\).
\end{SRWTdefinition}

\begin{SRWTtheorem}[Predictive-state sufficiency and minimality, \(T_{12}\)]
\label{SRWT:thm:predictive-state}
Under Definition~\ref{SRWT:def:predictive-kernel}:
\begin{enumerate}
\item \(Z_t\) is predictively sufficient:
\[
Y_t^+\SRWTindep\mathbf H_t\mid Z_t.
\]
\item If \(S_t=s(\mathbf H_t)\) is any other history statistic satisfying
\[
Y_t^+\SRWTindep\mathbf H_t\mid S_t,
\]
then there is a measurable map \(\phi\) such that
\[
Z_t=\phi(S_t)
\qquad\text{almost surely}.
\]
\item Consequently, any two minimal sufficient predictive states are
measurably isomorphic after removal of null sets.
\end{enumerate}
\end{SRWTtheorem}

\begin{proof}
For every bounded measurable \(f:\mathsf Y\to\mathbb R\), define
\[
Z_t(f)=\int_{\mathsf Y}f(y)\,Z_t(\mathrm dy).
\]
By construction,
\[
\mathbb E[f(Y_t^+)\mid\mathbf H_t]=Z_t(f).
\]
Since \(Z_t\) is a measurable function of \(\mathbf H_t\), the tower property
gives
\[
\mathbb E[f(Y_t^+)\mid Z_t]
=
\mathbb E[Z_t(f)\mid Z_t]
=
Z_t(f).
\]
Hence conditioning additionally on \(\mathbf H_t\) does not change the
conditional expectation, which is the stated conditional independence.

Now let \(S_t=s(\mathbf H_t)\) be sufficient. Then, for every measurable
\(A\subseteq\mathsf Y\),
\[
K_t(\mathbf H_t,A)
=
\mathbb P(Y_t^+\in A\mid\mathbf H_t)
=
\mathbb P(Y_t^+\in A\mid S_t)
\qquad\text{a.s.}
\]
Choose a countable determining class for probability measures on the standard
Borel space \(\mathsf Y\). The values of \(K_t(\mathbf H_t,\cdot)\) on this
class are measurable functions of \(S_t\), and therefore the complete random
measure \(Z_t\) is a measurable function of \(S_t\). Applying this conclusion
in both directions to two minimal sufficient states gives the final claim.
\end{proof}

\begin{SRWTremark}[Meaning of minimality]
The theorem does not say that \(Z_t\) has the fewest numerical coordinates in
a chosen parametrization. It says that every sufficient history statistic
must determine \(Z_t\) almost surely. This is an information order, invariant
under one-to-one reparametrization.
\end{SRWTremark}

For an arbitrary history statistic \(S_t=s(\mathbf H_t)\), let
\[
K_t^{S}(s,\cdot)
=
\mathbb P(Y_t^+\in\cdot\mid S_t=s).
\]
Whenever the conditional relative entropy is finite, define the sufficiency
defect
\[
\boxed{
\delta_{\mathrm{suff}}(S_t)
=
\mathbb E\!\left[
\SRWTKL\!\left(
K_t(\mathbf H_t,\cdot)
\,\middle\|\,
K_t^S(S_t,\cdot)
\right)
\right].
}
\]

\begin{SRWTproposition}[Conditional-information identity]\label{SRWT:prop:conditional-information}
If \(S_t\) is a measurable function of \(\mathbf H_t\), then
\[
\boxed{
\delta_{\mathrm{suff}}(S_t)
=
I(Y_t^+;\mathbf H_t\mid S_t).
}
\]
In particular, the defect vanishes if and only if \(S_t\) is predictively
sufficient, up to the usual almost-sure qualifications.
\end{SRWTproposition}

\begin{proof}
The conditional mutual information is the expected relative entropy between
the conditional law given \((\mathbf H_t,S_t)\) and the conditional law given
\(S_t\). Because \(S_t=s(\mathbf H_t)\), the first of these is exactly
\(K_t(\mathbf H_t,\cdot)\).
\end{proof}

Sufficiency has a direct decision-theoretic interpretation. Use the
convention
\[
\SRWTTV(P,Q)=\sup_{A}|P(A)-Q(A)|.
\]
Let \(\mathsf A\) be a finite action set and
\(\ell:\mathsf Y\times\mathsf A\to[0,1]\) a bounded loss. Define
\[
\SRWTRisk_\ell(S_t)
=
\inf_{d}
\mathbb E\!\left[
\ell\bigl(Y_t^+,d(S_t)\bigr)
\right].
\]

\begin{SRWTproposition}[Approximate sufficiency bounds excess Bayes risk]
\label{SRWT:prop:risk-pinsker}
Under the preceding assumptions,
\[
\boxed{
0
\leq
\SRWTRisk_\ell(S_t)-\SRWTRisk_\ell(\mathbf H_t)
\leq
\mathbb E\!\left[
\SRWTTV\!\left(K_t(\mathbf H_t),K_t^S(S_t)\right)
\right]
\leq
\sqrt{\frac{\delta_{\mathrm{suff}}(S_t)}{2}}.
}
\]
\end{SRWTproposition}

\begin{proof}
More information cannot increase optimal risk, giving the first inequality.
For any fixed action, its conditional expected loss changes by at most total
variation when the conditional law changes. Taking the infimum over actions
preserves this bound. Conditional Pinsker inequality followed by Jensen's
inequality gives the final estimate.
\end{proof}

For a finite target alphabet and logarithmic scoring loss
\[
\ell_{\log}(y,q)=-\log q(y),
\qquad
q\in\mathcal P(\mathsf Y),
\]
the identity is exact:
\[
\boxed{
\SRWTRisk_{\log}(S_t)=H(Y_t^+\mid S_t),
\qquad
\SRWTRisk_{\log}(S_t)-\SRWTRisk_{\log}(\mathbf H_t)
=
I(Y_t^+;\mathbf H_t\mid S_t).
}
\]
Natural logarithms measure the gap in nats; base-two logarithms measure it in
bits.

\subtopic{Intervention-indexed causal state}

Observational prediction may merge histories that behave identically under the
current policy but respond differently to future interventions. Let
\(\mathcal I\) be a finite or countable family of interventions applied only
after time \(t\). For each \(\iota\in\mathcal I\), let
\[
K_t^{\iota}(h,\cdot)
=
\mathbb P^{\iota}(Y_t^+\in\cdot\mid\mathbf H_t=h)
\]
be a regular post-intervention kernel on the common pre-intervention history
space. Define
\[
\boxed{
Z_t^{\mathcal I}
=
\bigl(K_t^{\iota}(\mathbf H_t,\cdot)\bigr)_{\iota\in\mathcal I}
\in
\mathcal P(\mathsf Y)^{\mathcal I}.
}
\]

\begin{SRWTcorollary}[Minimal simultaneous interventional state]
\label{SRWT:cor:interventional-state}
The random element \(Z_t^{\mathcal I}\) is sufficient simultaneously for all
post-intervention target laws. If a statistic \(S_t=s(\mathbf H_t)\)
determines the conditional target law for every \(\iota\in\mathcal I\), then
there exists a measurable \(\phi_{\mathcal I}\) such that
\[
Z_t^{\mathcal I}=\phi_{\mathcal I}(S_t)
\qquad\text{a.s.}
\]
\end{SRWTcorollary}

\begin{proof}
Apply Theorem~\ref{SRWT:thm:predictive-state} componentwise. A countable product of
standard Borel spaces is standard Borel, so the component maps combine into a
single measurable factor map.
\end{proof}

The observational state is the special case
\(\mathcal I=\{\operatorname{id}\}\). A causal workspace claim should align
with \(Z_t^{\mathcal I}\) for a declared intervention family rather than only
with the observational state.

\section{Structural--Operational Alignment}

A structural workspace coordinate becomes operational only when it retains the
canonical predictive or interventional state relevant to the declared task.
Let
\[
W_t=\pi_t(\mathbf H_t)
\]
be a measurable candidate workspace statistic satisfying the structural support
condition
\[
W_t\in\Gamma_{\varepsilon_t}(H(t))
\qquad\text{almost surely}.
\]
The map \(\pi_t\) may include encoding into a moving low-energy bundle; no
canonical operational interpretation follows from the spectral projector
alone.

\begin{SRWTdefinition}[Exact structural--operational alignment]
The candidate \(W_t\) is observationally aligned when there exists a measurable
map \(\phi_t\) such that
\[
Z_t=\phi_t(W_t)
\qquad\text{a.s.}
\]
It is \(\mathcal I\)-causally aligned when
\[
Z_t^{\mathcal I}=\phi_t^{\mathcal I}(W_t)
\qquad\text{a.s.}
\]
\end{SRWTdefinition}

\begin{SRWTtheorem}[Exact alignment criterion]\label{SRWT:thm:alignment-criterion}
For any workspace statistic \(W_t=\pi_t(\mathbf H_t)\), the following are
equivalent:
\begin{enumerate}
\item \(Y_t^+\SRWTindep\mathbf H_t\mid W_t\);
\item \(I(Y_t^+;\mathbf H_t\mid W_t)=0\);
\item the canonical predictive state factors through \(W_t\):
\[
Z_t=\phi_t(W_t)
\quad\text{a.s.}
\]
for some measurable \(\phi_t\).
\end{enumerate}
The analogous equivalence holds simultaneously over \(\mathcal I\) with
\(Z_t^{\mathcal I}\).
\end{SRWTtheorem}

\begin{proof}
The equivalence of the first two statements is the zero-divergence
characterization of conditional mutual information. The first statement and
Theorem~\ref{SRWT:thm:predictive-state} imply the factorization in the third.
Conversely, if \(Z_t=\phi_t(W_t)\), then
\[
\mathbb P(Y_t^+\in A\mid\mathbf H_t)
=Z_t(A)
=\phi_t(W_t)(A),
\]
which depends on the history only through \(W_t\), proving sufficiency. The
interventional statement follows from
Corollary~\ref{SRWT:cor:interventional-state}.
\end{proof}

Define the observational alignment defect by
\[
\boxed{
\delta_{\mathrm{align}}^{\mathrm{obs}}(W_t)
=
I_{\mathbb P}(Y_t^+;\mathbf H_t\mid W_t).
}
\]
For a countable intervention family with declared weights
\(q_\iota>0\), \(\sum_\iota q_\iota=1\), define
\[
\boxed{
\delta_{\mathrm{align}}^{\mathcal I}(W_t)
=
\sum_{\iota\in\mathcal I}
q_\iota
I_{\mathbb P^\iota}
(Y_t^+;\mathbf H_t\mid W_t).
}
\]
For bounded decision losses, Proposition~\ref{SRWT:prop:risk-pinsker} applies under
each intervention and gives the weighted bound
\[
\sum_{\iota}q_\iota
\left[
\SRWTRisk_\ell^{\iota}(W_t)
-
\SRWTRisk_\ell^{\iota}(\mathbf H_t)
\right]
\leq
\sqrt{\frac{
\delta_{\mathrm{align}}^{\mathcal I}(W_t)}{2}}.
\]

\subtopic{Structural coherence and operational alignment are independent}

The static defect
\[
\delta_{\mathrm{coh}}(W_t)
\]
and the operational defect
\[
\delta_{\mathrm{align}}^{\mathcal I}(W_t)
\]
answer different questions. The first measures compatibility under transfer.
The second measures loss of task-relevant conditional law.

\begin{center}
\small
\begin{tabularx}{\textwidth}{@{}>{\raggedright\arraybackslash}p{0.22\textwidth}>{\raggedright\arraybackslash}p{0.34\textwidth}>{\raggedright\arraybackslash}X@{}}
\toprule
Regime & Structural status & Operational interpretation \\
\midrule
Low coherence defect, low alignment defect & Stable transported coordinate &
Candidate operational workspace for the declared target and interventions. \\
Low coherence defect, high alignment defect & Stable but task-irrelevant
shared structure & Archive, convention, or synchronized representation. \\
High coherence defect, low alignment defect & Predictively adequate but not
transport-stable & Localized or brittle sufficient statistic. \\
High coherence defect, high alignment defect & Neither stable nor sufficient &
Rejected candidate at the stated scales. \\
\bottomrule
\end{tabularx}
\end{center}

Thus neither low transfer energy nor high predictive performance alone is
enough. The combined gate is
\[
\boxed{
\delta_{\mathrm{coh}}(W_t)
\leq
\varepsilon_{\mathrm{coh}},
\qquad
\delta_{\mathrm{align}}^{\mathcal I}(W_t)
\leq
\varepsilon_{\mathrm{align}}.
}
\]

\begin{SRWTexample}[Coherent but irrelevant archive]
Suppose every node perfectly copies a common random variable \(R\), so the
corresponding transfer energy is zero. If \(R\SRWTindep Y_t^+\), then
\[
\delta_{\mathrm{coh}}(R)=0,
\qquad
\delta_{\mathrm{align}}^{\mathrm{obs}}(R)
=I(Y_t^+;\mathbf H_t\mid R),
\]
which can be as large as the entire predictive information in the history.
Perfect agreement therefore does not imply workspace functionality.
\end{SRWTexample}

\begin{SRWTexample}[Sufficient but localized statistic]
If one node contains the canonical state \(Z_t\) and all transfers from that
node are unstable, then \(W_t=Z_t\) is sufficient but need not lie in a stable
low-energy transfer subspace. This is a local controller or archive, not yet a
stratified shared workspace.
\end{SRWTexample}

% This file is part of the modular source for Stratified Relational Workspace Theory.
% !TeX root = ../Stratified_Relational_Workspace_Theory.tex

\chapter{Relational irreducibility}\label{SRWT:book:chapter-73}


\section{Coalition Information and Bayes Risk}

For \(A\subseteq V\), the coalition statistic is \(\mathbf H_t^A\). Given an
action space \(\mathsf A\) and loss \(\ell\), define its unrestricted Bayes
risk by
\[
\boxed{
\SRWTRisk_\ell(A)
=
\inf_d
\mathbb E\!\left[
\ell\bigl(Y_t^+,d(\mathbf H_t^A)\bigr)
\right].
}
\]
The infimum is over measurable decision rules for which the expectation is
defined. The full-system risk is \(\SRWTRisk_\ell(V)\).

\begin{SRWTproposition}[Coalition monotonicity]\label{SRWT:prop:coalition-monotonicity}
If \(A\subseteq B\subseteq V\), then
\[
\SRWTRisk_\ell(B)\leq\SRWTRisk_\ell(A).
\]
\end{SRWTproposition}

\begin{proof}
Every decision rule available to coalition \(A\) is also available to coalition
\(B\), which can ignore its additional information.
\end{proof}

For finite \(\mathsf Y\), logarithmic loss converts risk gaps into conditional
information exactly.

\begin{SRWTtheorem}[Coalition log-risk identity]\label{SRWT:thm:coalition-log-risk}
Assume \(H(Y_t^+)<\infty\). If \(A\subseteq B\subseteq V\), then
\[
\boxed{
\SRWTRisk_{\log}(A)-\SRWTRisk_{\log}(B)
=
I\!\left(
Y_t^+;\mathbf H_t^B
\mid
\mathbf H_t^A
\right).
}
\]
In particular,
\[
\SRWTRisk_{\log}(A)-\SRWTRisk_{\log}(V)
=
I\!\left(
Y_t^+;\mathbf H_t
\mid
\mathbf H_t^A
\right).
\]
\end{SRWTtheorem}

\begin{proof}
Under logarithmic loss the optimal report is the conditional target law, so
\[
\SRWTRisk_{\log}(A)=H(Y_t^+\mid\mathbf H_t^A).
\]
Because \(\mathbf H_t^A\) is a function of \(\mathbf H_t^B\), subtracting the
two conditional entropies gives the stated conditional mutual information.
\end{proof}

The theorem makes the decomposition explicit: irreducibility is not an
unqualified property of a system. It is a property relative to the chosen node
partition, accessible histories, target, and decision loss.

\section{Joint Advantage and Its Scope}

Define the coalition gap
\[
\Delta_A^{\ell}
=
\SRWTRisk_\ell(A)-\SRWTRisk_\ell(V)
\geq0.
\]
The node-level joint advantage is
\[
\boxed{
\Delta_{\mathrm{joint}}^{\ell}
=
\min_{i\in V}
\Delta_{\{i\}}^{\ell}.
}
\]
Under logarithmic loss,
\[
\boxed{
\Delta_{\mathrm{joint}}^{\log}
=
\min_{i\in V}
I\!\left(
Y_t^+;\mathbf H_t
\mid
\mathbf H_t^{\{i\}}
\right).
}
\]
Therefore
\[
\Delta_{\mathrm{joint}}^{\log}>0
\]
means that every individual component omits some predictive information
available to the complete system.

This quantity should not automatically be identified with a particular notion
of partial-information synergy. A positive gap can arise from complementary
information distributed across components, from higher-order interaction, or
from both. The theory uses the operationally weaker and more directly
testable statement:
\[
\text{no individual node is sufficient for the declared target.}
\]

For \(0\leq r\leq |V|\), define the order-\(r\) coalition gap profile by
\[
\boxed{
\Delta^{(r)}_\ell
=
\min_{\substack{A\subseteq V\\|A|\leq r}}
\left[
\SRWTRisk_\ell(A)-\SRWTRisk_\ell(V)
\right].
}
\]
The profile is nonincreasing in \(r\) and reaches zero by \(r=|V|\).
It records how rapidly operative capacity becomes reconstructible as larger
coalitions are permitted.

\section{Minimal Sufficient Coalition}

For tolerance \(\varepsilon\geq0\), define
\[
\boxed{
k_\varepsilon
=
\min\left\{
|A|:
\SRWTRisk_\ell(A)
\leq
\SRWTRisk_\ell(V)+\varepsilon
\right\}.
}
\]
The set of minimal coalitions is
\[
\mathcal K_\varepsilon
=
\left\{
A\subseteq V:
|A|=k_\varepsilon,
\SRWTRisk_\ell(A)
\leq
\SRWTRisk_\ell(V)+\varepsilon
\right\}.
\]

\begin{SRWTcorollary}[Information form of coalition size]
Under the assumptions of Theorem~\ref{SRWT:thm:coalition-log-risk},
\[
\boxed{
k_\varepsilon
=
\min\left\{
|A|:
I(Y_t^+;\mathbf H_t\mid\mathbf H_t^A)
\leq\varepsilon
\right\}.
}
\]
Equivalently, \(k_\varepsilon\) is the first \(r\) for which
\(\Delta_{\log}^{(r)}\leq\varepsilon\).
\end{SRWTcorollary}

This yields a graded classification:
\[
\begin{aligned}
k_0=0
&\quad\Longrightarrow\quad
\text{the background law alone is sufficient},\\
k_0=1
&\quad\Longrightarrow\quad
\text{some individual node is sufficient},\\
k_0\geq2
&\quad\Longrightarrow\quad
\text{node-level relational irreducibility},\\
k_0=|V|
&\quad\Longrightarrow\quad
\text{every proper coalition is insufficient}.
\end{aligned}
\]
The last condition is much stronger and should not be inferred from the first.

\section{Reconstructibility of the Canonical State}

Let \(d_{\mathcal P}\) be a bounded metric generating the measurable structure
on \(\mathcal P(\mathsf Y)\). Define the reconstruction defect of the canonical
state from coalition \(A\) by
\[
\delta_A^Z
=
\inf_{\psi_A}
\mathbb E\!\left[
 d_{\mathcal P}\!\left(
 Z_t,\psi_A(\mathbf H_t^A)
 \right)
\right].
\]

\begin{SRWTproposition}[Sufficiency--reconstruction equivalence]
\label{SRWT:prop:sufficiency-reconstruction}
For a coalition \(A\subseteq V\), the following are equivalent:
\begin{enumerate}
\item \(\mathbf H_t^A\) is sufficient for \(Y_t^+\) relative to the full
 history;
\item \(Z_t\) is measurable with respect to \(\mathcal F_t^A\);
\item there exists a measurable \(\psi_A\) such that
\[
Z_t=\psi_A(\mathbf H_t^A)
\qquad\text{almost surely}.
\]
\end{enumerate}
Each condition implies \(\delta_A^Z=0\). Under finite-target logarithmic loss,
they are also equivalent to
\[
\SRWTRisk_{\log}(A)=\SRWTRisk_{\log}(V).
\]
If the reconstruction infimum is known to be attained, then
\(\delta_A^Z=0\) is itself equivalent to the three conditions.
\end{SRWTproposition}

\begin{proof}
The equivalence of the first two statements is
Theorem~\ref{SRWT:thm:predictive-state}; the second and third are two formulations
of measurability up to null sets. An exact reconstructor has zero expected
metric error. The log-risk statement follows from
Theorem~\ref{SRWT:thm:coalition-log-risk} and the zero condition for conditional
mutual information. Without an attainment or compactness condition, an
infimum of zero alone certifies arbitrarily accurate reconstruction but need
not supply a single exact minimizing map.
\end{proof}

For an arbitrary enriched workspace variable \(W_t\), local
non-reconstructibility need not imply predictive irreducibility: \(W_t\) may
contain task-irrelevant material. The equivalence is exact for the canonical
state \(Z_t\), or for a candidate \(W_t\) known to be minimally aligned with
it.

\section{Matched Relational Nulls}

A null control must preserve the local facts that the claim is not intended to
challenge while destroying the relational organization being tested. Merely
reducing model capacity or replacing the workspace by noise is not a matched
null.

Let \(P\) be the intact law of
\((Y_t^+,\mathbf H_t^1,\ldots,\mathbf H_t^n)\). Define the local-marginal null
class
\[
\mathcal N_{\mathrm{marg}}(P)
=
\left\{
Q:
Q_{\mathbf H_t^i}=P_{\mathbf H_t^i}
\text{ for every }i
\right\},
\]
and the stronger local-predictive-matched class
\[
\boxed{
\mathcal N_{\mathrm{pred}}(P)
=
\left\{
Q:
Q_{(Y_t^+,\mathbf H_t^i)}
=
P_{(Y_t^+,\mathbf H_t^i)}
\text{ for every }i
\right\}.
}
\]
Because \(\mathbf H_t^i\) can include the complete local trajectory, the latter
class can preserve all single-component temporal and predictive behavior while
altering cross-component organization.

\begin{SRWTproposition}[Local-risk invariance under predictive matching]
\label{SRWT:prop:local-risk-null}
If \(Q\in\mathcal N_{\mathrm{pred}}(P)\), then for every node \(i\) and every
loss determined by \((Y_t^+,d(\mathbf H_t^i))\),
\[
\boxed{
\SRWTRisk_{\ell,Q}(\{i\})
=
\SRWTRisk_{\ell,P}(\{i\}).
}
\]
\end{SRWTproposition}

\begin{proof}
Every such risk depends only on the joint law of
\((Y_t^+,\mathbf H_t^i)\), which is identical under \(P\) and \(Q\).
\end{proof}

A valid null must also satisfy a declared relation-breaking condition. Examples
include:
\begin{itemize}
\item permuting source--target message pairings while preserving each message
 marginal;
\item shuffling temporal correspondence across nodes while preserving every
 local trajectory;
\item permuting role bindings while preserving the local token or state
 distributions;
\item replacing an artifact by a matched-size, matched-format artifact whose
 relation to the surrounding history has been destroyed; and
\item randomizing transfer channels through a physical or software
 intervention while preserving message rate and marginal content.
\end{itemize}

For a predeclared null \(Q\), define
\[
\boxed{
\Delta_{\mathrm{tr}}(Q)
=
\SRWTRisk_{\ell,Q}(V)
-
\SRWTRisk_{\ell,P}(V).
}
\]
For a finite validated panel
\(\mathcal Q=\{Q_1,\ldots,Q_m\}\subseteq\mathcal N_{\mathrm{pred}}(P)\), report
both the complete vector
\[
\boldsymbol\Delta_{\mathrm{tr}}(\mathcal Q)
=
\bigl(
\Delta_{\mathrm{tr}}(Q_1),\ldots,
\Delta_{\mathrm{tr}}(Q_m)
\bigr)
\]
and the robust panel minimum
\[
\Delta_{\mathrm{tr}}^{\min}(\mathcal Q)
=
\min_{Q\in\mathcal Q}\Delta_{\mathrm{tr}}(Q).
\]
The panel must be fixed independently of the observed favorable gap. Otherwise
null selection can manufacture apparent relational dependence.

\section{Relational Witness Pairs}

\begin{SRWTdefinition}[Relational witness pair]\label{SRWT:def:relational-witness}
A pair of laws \((P,Q)\) is a predictive relational witness pair for the
decomposition \(V\) and target \(Y_t^+\) when
\[
Q\in\mathcal N_{\mathrm{pred}}(P)
\]
but
\[
\SRWTRisk_{\ell,Q}(V)
\neq
\SRWTRisk_{\ell,P}(V).
\]
The two systems agree on every single-node predictive marginal but differ in
full-system operative capacity.
\end{SRWTdefinition}

\begin{SRWTexample}[XOR witness pair]\label{SRWT:ex:xor-witness}
Under the intact law \(P\), let
\[
X_1,X_2\overset{\mathrm{iid}}{\sim}\operatorname{Bernoulli}(1/2),
\qquad
Y=X_1\oplus X_2.
\]
Then each \((X_i,Y)\) pair is independent and uniform, so
\[
H_P(Y\mid X_i)=1\text{ bit},
\qquad
H_P(Y\mid X_1,X_2)=0.
\]
Let \(Q\) make \(X_1,X_2,Y\) mutually independent fair bits. Then
\[
Q_{(X_i,Y)}=P_{(X_i,Y)}
\qquad(i=1,2),
\]
but
\[
H_Q(Y\mid X_1,X_2)=1\text{ bit}.
\]
Therefore
\[
\Delta_{\mathrm{joint}}^{\log}=1\text{ bit},
\qquad
k_0=2,
\qquad
\Delta_{\mathrm{tr}}(Q)=1\text{ bit}.
\]
This proves predictive relational structure relative to the declared variables.
It does not, by itself, prove that any workspace variable is causally active.
\end{SRWTexample}

\begin{SRWTexample}[Redundant-copy control]
Let \(Y\) be a fair bit and set \(X_1=X_2=Y\). The two nodes are perfectly
coherent and the full system predicts perfectly, but
\[
H(Y\mid X_i)=0,
\qquad
\Delta_{\mathrm{joint}}^{\log}=0,
\qquad
k_0=1.
\]
Duplication is therefore not relational irreducibility.
\end{SRWTexample}

The witness-pair criterion is stronger than a performance comparison against a
smaller model because it holds local predictive behavior fixed. It remains
observational. Chapter~\ref{SRWT:book:chapter-74} adds interventions and converts predictive
relationality into an operational causal claim.

% This file is part of the modular source for Stratified Relational Workspace Theory.
% !TeX root = ../Stratified_Relational_Workspace_Theory.tex

\chapter{Causal efficacy and reentry}\label{SRWT:book:chapter-74}


\section{Feasible Workspace Interventions}

A notation such as \(\operatorname{do}(W_t=w)\) is meaningful only if the
candidate workspace is an endogenous variable in a structural causal model or
if an intervention on underlying variables is shown to realize that value.
Abstractly setting a projection coordinate need not correspond to any possible
physical or computational manipulation.

\begin{SRWTdefinition}[Time-unrolled causal realization]\label{SRWT:def:causal-realization}
A time-unrolled causal realization is a structural causal model
\[
\mathcal M
=
(\mathcal U,\mathcal V,F,\mathbb P_U)
\]
whose endogenous variables include the local states, messages, actions,
artifacts, target variables, and a measurable workspace extractor
\[
\pi_t:\mathsf S_t\to\mathsf W_t,
\qquad
W_t=\pi_t(S_t).
\]
The unrolled graph is acyclic even when the rolled process is recurrent.
\end{SRWTdefinition}

\begin{SRWTdefinition}[Feasible implementation intervention]
Let \(B_t\) denote declared background variables that should remain matched.
A family \(\{\iota_w:w\in\mathsf W_t^\circ\}\) is a feasible implementation of
workspace values on the experimental domain \(\mathsf W_t^\circ\) when:
\begin{enumerate}
\item each \(\iota_w\) acts only on declared manipulable parents, channels, or
 artifacts at or after time \(t\);
\item the intervened realization satisfies
\[
\pi_t(S_t^{\iota_w})=w
\qquad\text{almost surely};
\]
\item the background law is matched across compared interventions; and
\item the intervention does not directly modify \(Y_t^+\) except through
 paths included in the declared causal design.
\end{enumerate}
\end{SRWTdefinition}

Let \(D\) be a divergence or metric on probability laws. For a matched
background value \(b\), define
\[
\SRWTCE_{\SRWTWorkspace}(w,w';b)
=
D\!\left(
\mathbb P^{\iota_w}(Y_t^+\in\cdot\mid B_t=b),
\mathbb P^{\iota_{w'}}(Y_t^+\in\cdot\mid B_t=b)
\right).
\]
For a declared comparison design \(\nu\) over \((w,w',b)\), define
\[
\boxed{
\eta_{\mathrm{causal}}^{\nu}
=
\mathbb E_{(w,w',b)\sim\nu}
\left[
\SRWTCE_{\SRWTWorkspace}(w,w';b)
\right].
}
\]
The maximum
\[
\eta_{\mathrm{causal}}^{\max}
=
\sup_{w,w',b}\SRWTCE_{\SRWTWorkspace}(w,w';b)
\]
may be useful for discovery, but a design-weighted effect is harder to inflate
by selecting a single exceptional pair after inspection.

\begin{SRWTproposition}[Coordinate invariance of intervention effects]
\label{SRWT:prop:causal-gauge}
Let \(g:\mathsf W_t\to\mathsf W_t'\) be a measurable bijection and suppose the
reparametrized intervention family satisfies
\[
\iota_{g(w)}'=\iota_w.
\]
Then every divergence-based causal effect is preserved:
\[
\SRWTCE_{\SRWTWorkspace'}(g(w),g(w');b)
=
\SRWTCE_{\SRWTWorkspace}(w,w';b).
\]
\end{SRWTproposition}

\begin{proof}
The two labels designate the same underlying interventions and therefore induce
the same pair of target laws. Only the coordinate name has changed.
\end{proof}

This is the causal analogue of gauge invariance. It also shows why effects
should be attached to intervention implementations rather than to arbitrary
numerical coordinates.

\section{Causal Specificity and Transfer Mediation}

A causal effect on the target is not workspace-specific if an equal-cost
perturbation outside the candidate workspace produces the same effect. Let
\(c(\iota)\) be a declared intervention cost and let
\[
\mathcal I_{\SRWTWorkspace}(r),
\qquad
\mathcal I_{\perp}(r)
\]
be matched families with cost \(r\), the first changing workspace coordinates
and the second changing complementary coordinates while preserving the
workspace value as closely as the implementation permits.

Define
\[
\eta_{\SRWTWorkspace}(r)
=
\mathbb E_{\iota,\iota'\sim\mathcal I_{\SRWTWorkspace}(r)}
D\!\left(
\mathbb P^{\iota}(Y_t^+),
\mathbb P^{\iota'}(Y_t^+)
\right),
\]
\[
\eta_{\perp}(r)
=
\mathbb E_{\jmath,\jmath'\sim\mathcal I_{\perp}(r)}
D\!\left(
\mathbb P^{\jmath}(Y_t^+),
\mathbb P^{\jmath'}(Y_t^+)
\right),
\]
and
\[
\boxed{
\Delta_{\mathrm{specific}}(r)
=
\eta_{\SRWTWorkspace}(r)-\eta_{\perp}(r).
}
\]
A positive value supports causal concentration in the candidate workspace at
the declared perturbation scale.

If the complement is defined by \(\Gamma_\varepsilon(H)^\perp\), the result is
metric dependent. Under nonunitary gauge changes, either the metric must be
transported with the coordinates or interventions must be matched by a
coordinate-independent physical cost. Otherwise specificity can change merely
because the representation was rescaled.

\subtopic{Causal transfer nulls}

Let \(\jmath_E\) randomize selected communication edges while preserving their
marginal message laws, rates, and timing budget. The causal transfer gap is
\[
\boxed{
\Delta_{\mathrm{tr}}^{\operatorname{do}}(\jmath_E)
=
\SRWTRisk_{\ell,\mathbb P^{\jmath_E}}(V)
-
\SRWTRisk_{\ell,\mathbb P}(V).
}
\]
Unlike an observational null law, this quantity is anchored to an implemented
channel intervention. A positive effect is evidence that the intact alignment
of messages, rather than merely their marginal content, contributes to task
performance.

A stronger mediation test equalizes the workspace state after two distinct
channel interventions. Let \(\mu\) be a common imposed workspace distribution,
implemented by a feasible stochastic intervention in the sense of
Definition~\ref{SRWT:def:causal-realization}. The notation below is shorthand for
that implementation. Define the residual mediation defect
\[
\delta_{\mathrm{med}}
=
\sup_{\jmath,\jmath'}
D\!\left(
\mathbb P^{\jmath,\operatorname{do}(W_t\sim\mu)}(Y_t^+),
\mathbb P^{\jmath',\operatorname{do}(W_t\sim\mu)}(Y_t^+)
\right).
\]
A small value means that once the workspace state is equalized, little of the
channel intervention's effect remains. This supports, but does not by itself
identify without the causal assumptions, the claim that \(W_t\) mediates the
relevant relational influence.

\section{Causal Recurrence, Loop-Mediated Reentry, and Relationality}

Temporal correlation is weaker than causal recurrence. Define the total
workspace recurrence effect at lag \(\tau\) by
\[
\boxed{
\eta_{\mathrm{rec}}(\tau)
=
\mathbb E_{(w,w',b)\sim\nu}
D\!\left(
\mathbb P^{\iota_w}(W_{t+\tau}\in\cdot\mid B_t=b),
\mathbb P^{\iota_{w'}}(W_{t+\tau}\in\cdot\mid B_t=b)
\right).
}
\]
A nonzero value shows that manipulating the present workspace changes a later
workspace state. It may, however, be explained by direct copying or a retained
memory register.

To isolate reentry, consider the unrolled causal graph and remove every direct
retention edge whose source and target are both workspace variables at
successive times. Let
\[
\mathcal M^{\setminus\mathrm{dir}}
\]
denote the resulting path-specific model, preserving paths that leave the
workspace through a nonworkspace participant, action, message, or artifact and
later return. Define
\[
\boxed{
\eta_{\mathrm{loop}}(\tau)
=
\mathbb E_{(w,w',b)\sim\nu}
D\!\left(
\mathbb P_{\mathcal M^{\setminus\mathrm{dir}}}^{\iota_w}
(W_{t+\tau}\in\cdot\mid B_t=b),
\mathbb P_{\mathcal M^{\setminus\mathrm{dir}}}^{\iota_{w'}}
(W_{t+\tau}\in\cdot\mid B_t=b)
\right).
}
\]
Then the distinctions are:
\[
\begin{aligned}
I(W_t;W_{t+\tau})>0
&\quad\text{temporal dependence only},\\
\eta_{\mathrm{rec}}(\tau)>0
&\quad\text{causal persistence},\\
\eta_{\mathrm{loop}}(\tau)>0
&\quad\text{causal exit-and-return reentry}.
\end{aligned}
\]
Path-specific effects require the usual consistency and no-unmeasured-confounding
assumptions appropriate to the implementation. In software systems, edge
interventions can often be literal. In human collaboration, these assumptions
are substantially harder to justify.

\begin{SRWTtheorem}[Operational relationality, \(T_{13}\)]
\label{SRWT:thm:operational-relationality}
Let \(\mathsf Y\) be finite, let \(P\) be the intact law at a fixed decision
time, and let \(W_t=\pi_t(\mathbf H_t)\) be a structurally supported candidate.
Let
\(\mathcal Q\subseteq\mathcal N_{\mathrm{pred}}(P)\) be a nonempty,
predeclared panel of validated relation-breaking nulls. Assume:
\begin{enumerate}
\item exact structural--operational alignment,
\[
I_P(Y_t^+;\mathbf H_t\mid W_t)=0;
\]
\item a positive node-level information gap,
\[
\delta_{\mathrm{node}}
=
\min_{i\in V}
I_P(Y_t^+;\mathbf H_t\mid\mathbf H_t^i)
>0;
\]
\item a positive robust matched-null gap,
\[
\delta_{\mathrm{null}}
=
\min_{Q\in\mathcal Q}
\left[
H_Q(Y_t^+\mid\mathbf H_t)
-
H_P(Y_t^+\mid\mathbf H_t)
\right]
>0;
\]
\item a feasible intervention design with
\[
\eta_{\mathrm{causal}}^{\nu}>0;
\]
\item positive matched causal specificity at some declared scale \(r\),
\[
\Delta_{\mathrm{specific}}(r)>0.
\]
\end{enumerate}
Then, relative to the declared decomposition, target, log loss, null panel, and
intervention family:
\begin{enumerate}
\item no statistic of any single-node history is sufficient for \(Y_t^+\), and
 \(k_0\geq2\);
\item \(W_t\) cannot be reconstructed almost surely from any single-node
 history;
\item every null \(Q\in\mathcal Q\) preserves all single-node predictive
 risks but has strictly worse full-system log risk than the intact law;
\item the candidate workspace is causally efficacious and its effect exceeds
 that of the declared matched off-workspace perturbations.
\end{enumerate}
If, in addition, \(\eta_{\mathrm{loop}}(\tau)>0\), the workspace is
loop-reentrant at lag \(\tau\).
\end{SRWTtheorem}

\begin{proof}
By Theorem~\ref{SRWT:thm:coalition-log-risk}, for each node \(i\),
\[
\SRWTRisk_{\log,P}(\{i\})-
\SRWTRisk_{\log,P}(V)
=
I_P(Y_t^+;\mathbf H_t\mid\mathbf H_t^i)
\geq
\delta_{\mathrm{node}}>0.
\]
Hence no single-node history, and therefore no statistic of one, is sufficient;
this gives \(k_0\geq2\).

Exact alignment and Theorem~\ref{SRWT:thm:alignment-criterion} imply that the
canonical predictive state factors through \(W_t\). If
\(W_t=\psi_i(\mathbf H_t^i)\) almost surely for some node, then the canonical
state would also be a function of \(\mathbf H_t^i\), contradicting the positive
node gap through Proposition~\ref{SRWT:prop:sufficiency-reconstruction}.

For every \(Q\in\mathcal Q\), Proposition~\ref{SRWT:prop:local-risk-null} gives
\[
\SRWTRisk_{\log,Q}(\{i\})
=
\SRWTRisk_{\log,P}(\{i\})
\qquad\forall i.
\]
The assumed null gap gives
\[
\SRWTRisk_{\log,Q}(V)
>
\SRWTRisk_{\log,P}(V).
\]
Thus a description determined only by the collection of single-node predictive
marginals cannot reproduce the intact joint capacity: it would assign the same
local data to \(P\) and \(Q\) while their full-system risks differ.

Finally, the feasible intervention effect establishes causal efficacy relative
to the declared implementation, and the positive specificity gap shows that
the effect exceeds the matched complementary perturbation effect. The final
reentry statement is exactly the definition of \(\eta_{\mathrm{loop}}\).
\end{proof}

\begin{SRWTremark}[Scope of the theorem]
The conclusion is operational and indexed by
\[
(V,Y_t^+,\ell,\mathcal Q,\mathcal I,\nu,r).
\]
Changing the decomposition, target, loss, nulls, or intervention design may
change the verdict. The theorem does not establish metaphysical
irreducibility, autonomous agency, selfhood, or consciousness.
\end{SRWTremark}

\subtopic{Finite certification sequence}

A reproducible dynamic-causal analysis should proceed in the following order:
\begin{enumerate}
\item Freeze the node decomposition, accessible coalition histories, target,
 horizon, and proper scoring loss.
\item Estimate the structural projector on training data only and encode
 \(W_t\) without using held-out targets.
\item Cross-fit conditional target laws from the full history, the workspace,
 and every declared coalition. Report alignment and coalition risk gaps.
\item Construct the null panel before inspecting null performance. Verify
 local-predictive matching and the intended destruction of transfer
 organization separately.
\item Implement feasible workspace and complementary interventions with a
 common background and cost design. Report both causal effect and
 specificity.
\item Test recurrence and then cut direct retention paths to test loop-mediated
 reentry.
\item Report confidence intervals or finite-sample uncertainty for every gate;
 do not substitute training accuracy for Bayes-risk or causal estimates.
\end{enumerate}

The resulting certification report is
\[
\boxed{
\begin{aligned}
\mathcal C_{\mathrm{dyn}}(W_t)
=\bigl(&
\delta_{\mathrm{align}}^{\mathcal I},
\{\Delta_A^\ell\}_{A\subseteq V},
 k_\varepsilon,
\boldsymbol\Delta_{\mathrm{tr}}(\mathcal Q),\\
&\eta_{\mathrm{causal}}^\nu,
\Delta_{\mathrm{specific}}(r),
\eta_{\mathrm{rec}}(\tau),
\eta_{\mathrm{loop}}(\tau)
\bigr).
\end{aligned}
}
\]
No single coordinate certifies a relational workspace.

% This file is part of the modular source for Stratified Relational Workspace Theory.
% !TeX root = ../Stratified_Relational_Workspace_Theory.tex

\chapter{Formal workspace definition}\label{SRWT:book:chapter-75}


\section{Certified Stratified Relational Workspace}

Let
\[
\Theta
=
\left(
\varepsilon_{\mathrm{coh}},
\varepsilon_{\mathrm{align}},
\varepsilon_k,
\underline\eta_c,
\underline\eta_{\mathrm{loop}},
\delta_{\mathrm{node}},
\delta_{\mathrm{null}},
\delta_{\mathrm{spec}},
k,
\varepsilon_{\mathrm{pers}}
\right).
\]
Let
\[
\varsigma
=
(\mathcal D,Y^+,\ell,\tau,\mathcal I,\mathcal Q,\Theta)
\]
be the complete declared scientific specification, where
\(\mathcal Q\subseteq\mathcal N_{\mathrm{pred}}(P)\) is a validated null panel.
A dynamic realization contains a certified stratified relational workspace
when there exists a candidate \(W_t\) satisfying all gates required for the
claimed level under \(\varsigma\). The detailed robustness and identifiability
requirements for varying \(\varsigma\) are developed in Chapters~\ref{SRWT:book:chapter-83} and~\ref{SRWT:book:chapter-84}.

\subtopic{Structural coherence}
\[
\delta_{\mathrm{coh}}(W_t)
\leq
\varepsilon_{\mathrm{coh}}.
\]

\subtopic{Interventional structural--operational alignment}
\[
\delta_{\mathrm{align}}^{\mathcal I}(W_t)
\leq
\varepsilon_{\mathrm{align}}.
\]

\subtopic{Causal efficacy}
\[
\eta_{\mathrm{causal}}^\nu
\geq
\underline\eta_c.
\]

\subtopic{Loop-mediated reentry, when claimed}
\[
\eta_{\mathrm{loop}}(\tau)
\geq
\underline\eta_{\mathrm{loop}}.
\]

\subtopic{Node-level joint advantage}
\[
\Delta_{\mathrm{joint}}^{\ell}
\geq
\delta_{\mathrm{node}}.
\]

\subtopic{Matched-null transfer dependence}
\[
\Delta_{\mathrm{tr}}^{\min}(\mathcal Q)
\geq
\delta_{\mathrm{null}}.
\]

\subtopic{Causal specificity}
\[
\Delta_{\mathrm{specific}}(r)
\geq
\delta_{\mathrm{spec}}.
\]

\subtopic{Distributedness}
\[
k_{\varepsilon_k}
\geq
k.
\]

\subtopic{Persistence, when claimed}
\[
d_{\mathrm{id}}(W_t,W_{t+\tau})
\leq
\varepsilon_{\mathrm{pers}}.
\]

Symbolically,
\[
\boxed{
\SRWTSRW(W_t;\varsigma)
=
\bigwedge_{q\in\mathcal J(\varsigma)}
\operatorname{Gate}_q(W_t;\varsigma).
}
\]
The gate set is claim-relative. A one-step operational workspace need not pass
persistence or reentry gates, while a persistent reentrant workspace must.
No scalar total score is sufficient: failure of a required gate cannot be
hidden by high performance on another gate.

\begin{SRWTremark}[Structural, predictive, causal, and temporal claims]
The formal hierarchy now distinguishes four logically independent questions:
\[
\begin{aligned}
&\text{Does a stable shared coordinate exist?}\\
&\text{Does it retain the task-relevant conditional law?}\\
&\text{Does manipulating it change downstream behavior specifically?}\\
&\text{Does it persist or reenter through time?}
\end{aligned}
\]
The preceding structural and causal analysis answers the first three at the level of definitions and conditional
theorems. The following parts develop the fourth.
\end{SRWTremark}

% This file is part of the modular source for Stratified Relational Workspace Theory.
% !TeX root = ../Stratified_Relational_Workspace_Theory.tex

\chapter{Persistence}\label{SRWT:book:chapter-76}


The preceding parts identify exact and approximate workspace subspaces at a
single time. Persistence adds three further requirements. The structural
subspace must remain trackable while the realization changes, actual states
must remain near that moving subspace, and the transported content must not
collapse to a trivial zero state. These are distinct conditions.

\section{Coherence-Error Dynamics}

\subtopic{32.1 Time-indexed transfer systems}

At time \(t\), let
\[
\mathcal C_V(t)=\bigoplus_{i\in V}C_i(t),
\qquad
\mathcal C_E(t)=\bigoplus_{e\in E}C_{t(e)}(t),
\]
and let
\[
D_t:\mathcal C_V(t)\longrightarrow\mathcal C_E(t),
\qquad
L_t=D_t^*D_t
\]
be the transfer discrepancy and transfer Laplacian. At time \(t\), the
combinatorial graph is held fixed while the local representations, transfers,
and metrics may vary. Changes of graph are treated later by a graph
correspondence.

To compare successive times, choose unitary identification maps
\[
U_t^V:\mathcal C_V(t)\longrightarrow\mathcal C_V(t+1),
\qquad
U_t^E:\mathcal C_E(t)\longrightarrow\mathcal C_E(t+1).
\]
They are comparison maps, not assumptions that the physical carriers are
identical. Define the one-step diagram defect
\[
\boxed{
\Delta_t^D
=
D_{t+1}U_t^V-U_t^ED_t,
\qquad
\delta_t^D=\|\Delta_t^D\|.
}
\]
Exact temporal compatibility is the condition \(\Delta_t^D=0\).

\begin{SRWTproposition}[Temporal Laplacian drift]\label{SRWT:prop:temporal-laplacian-drift}
The carried Laplacian satisfies
\[
\boxed{
\left\|
(U_t^V)^*L_{t+1}U_t^V-L_t
\right\|
\leq
2\|D_t\|\delta_t^D+(\delta_t^D)^2.
}
\]
\end{SRWTproposition}

\begin{proof}
Set
\[
B_t=(U_t^E)^*\Delta_t^D.
\]
Because the comparison maps are unitary,
\[
(U_t^E)^*D_{t+1}U_t^V=D_t+B_t,
\qquad
\|B_t\|=\delta_t^D.
\]
Hence
\[
(U_t^V)^*L_{t+1}U_t^V
=(D_t+B_t)^*(D_t+B_t).
\]
Subtract \(D_t^*D_t\) and apply the triangle inequality.
\end{proof}

Let \(P_t\) project onto a selected rank-\(d\) low-energy workspace cluster
of \(L_t\), and set
\[
\mathcal W_t=\SRWTran P_t.
\]
Write
\[
g_t=\lambda_{d+1}(L_t)-\lambda_d(L_t)>0
\]
for the separating spectral gap and
\[
\eta_t^L=2\|D_t\|\delta_t^D+(\delta_t^D)^2.
\]

\begin{SRWTcorollary}[Temporal workspace tracking]\label{SRWT:cor:temporal-workspace-tracking}
If
\[
\eta_t^L<\frac{g_t}{2},
\]
then the transported and newly extracted workspace projectors satisfy
\[
\boxed{
q_t^{\mathrm{str}}
:=
\left\|
P_{t+1}-U_t^VP_t(U_t^V)^*
\right\|
\leq
\frac{\eta_t^L}{g_t-\eta_t^L}
\leq
\frac{2\eta_t^L}{g_t}.
}
\]
\end{SRWTcorollary}

\begin{proof}
Carry \(L_{t+1}\) back to time \(t\) using \(U_t^V\), apply
Theorem~\ref{SRWT:thm:spectral-stability}, and conjugate the resulting projector
bound forward.
\end{proof}

The quantity \(q_t^{\mathrm{str}}\) measures structural motion after the
declared coordinate change has been removed. A workspace may therefore move
substantially in raw coordinates while remaining structurally unchanged.

\subtopic{32.2 State tracking inside a moving workspace}

Let the realized vertex state satisfy
\[
x_{t+1}=F_tx_t+\xi_t,
\]
where \(F_t:\mathcal C_V(t)\to\mathcal C_V(t+1)\) is the actual update and
\(\xi_t\) is an unmodeled disturbance. Define the off-workspace error
\[
\boxed{
e_t=\|(I-P_t)x_t\|.
}
\]
Introduce
\[
\alpha_t
=
\left\|
(I-P_{t+1})F_t(I-P_t)
\right\|,
\]
which measures persistence of prior off-workspace error, and
\[
\ell_t
=
\left\|
(I-P_{t+1})F_tP_t
\right\|,
\]
which measures leakage of a valid workspace state under the update.

\begin{SRWTproposition}[Grounded coherence recurrence]\label{SRWT:prop:grounded-coherence-recurrence}
Every trajectory satisfies
\[
\boxed{
e_{t+1}\leq\alpha_te_t+\sigma_t,}
\]
where
\[
\boxed{
\sigma_t
=
\ell_t\|P_tx_t\|
+
\|(I-P_{t+1})\xi_t\|.
}
\]
\end{SRWTproposition}

\begin{proof}
Insert \(I=P_t+(I-P_t)\) before \(x_t\), apply \(I-P_{t+1}\), and use
submultiplicativity:
\[
\begin{aligned}
e_{t+1}
&\leq
\|(I-P_{t+1})F_t(I-P_t)x_t\|\\
&\quad+
\|(I-P_{t+1})F_tP_tx_t\|
+
\|(I-P_{t+1})\xi_t\|.
\end{aligned}
\]
The displayed recurrence follows from the definitions of \(\alpha_t\) and
\(\ell_t\).
\end{proof}

This derivation gives the contraction and disturbance terms an operational
meaning. The scalar recurrence is not an independent postulate. It follows
from an update rule, a moving workspace projector, and a declared disturbance
model.

\subtopic{32.3 Persistence theorem \(T_6\)}

\begin{SRWTtheorem}[Nonautonomous persistence theorem]\label{SRWT:thm:persistence}
Let \(e_t\geq0\) satisfy
\[
e_{t+1}\leq\alpha_te_t+\sigma_t,
\qquad
\alpha_t,\sigma_t\geq0.
\]
Then, for every \(t\geq1\),
\[
\boxed{
\begin{aligned}
e_t
\leq
&\left(\prod_{r=0}^{t-1}\alpha_r\right)e_0\\
&+
\sum_{j=0}^{t-1}
\left(\prod_{r=j+1}^{t-1}\alpha_r\right)\sigma_j,
\end{aligned}
}
\]
where an empty product equals one. If
\[
\alpha_t\leq\bar\alpha<1,
\qquad
\sigma_t\leq\bar\sigma,
\]
then
\[
\boxed{
e_t
\leq
\bar\alpha^te_0
+
\frac{\bar\sigma(1-\bar\alpha^t)}{1-\bar\alpha},
}
\]
and consequently
\[
\boxed{
\limsup_{t\to\infty}e_t
\leq
\frac{\bar\sigma}{1-\bar\alpha}.
}
\]
\end{SRWTtheorem}

\begin{proof}
Repeated substitution yields
\[
\begin{aligned}
e_t
&\leq
\alpha_{t-1}\cdots\alpha_0e_0
+
\alpha_{t-1}\cdots\alpha_1\sigma_0
+
\cdots
+
\sigma_{t-1},
\end{aligned}
\]
which is the first formula. Under the uniform bounds, dominate each product by
the corresponding power of \(\bar\alpha\) and sum the geometric series.
\end{proof}

\begin{SRWTcorollary}[Vanishing-disturbance convergence]\label{SRWT:cor:vanishing-disturbance}
If \(\alpha_t\leq\bar\alpha<1\) and \(\sigma_t\to0\), then
\[
e_t\to0.
\]
More generally, if \(\sum_{t\geq0}\sigma_t<\infty\), then the disturbance
contribution in Theorem~\ref{SRWT:thm:persistence} vanishes asymptotically.
\end{SRWTcorollary}

\begin{proof}
Split the convolution sum into a finite early part, whose coefficients decay
to zero, and a late tail on which \(\sigma_t\) is uniformly small. The
summable case follows by the same argument.
\end{proof}

\subtopic{32.4 Mean-square persistence}

For stochastic systems, pathwise bounded disturbances may be too strong. Let
\(\{\mathcal F_t\}\) be a filtration containing the history through time
\(t\).

\begin{SRWTproposition}[Mean-square persistence]\label{SRWT:prop:mean-square-persistence}
Suppose
\[
\mathbb E[e_{t+1}^2\mid\mathcal F_t]
\leq
q e_t^2+v_t,
\qquad
0\leq q<1,
\]
and \(\mathbb E[v_t]\leq\bar v\). Then
\[
\boxed{
\mathbb E[e_t^2]
\leq
q^t\mathbb E[e_0^2]
+
\frac{\bar v(1-q^t)}{1-q},
}
\]
so
\[
\limsup_{t\to\infty}\mathbb E[e_t^2]
\leq
\frac{\bar v}{1-q}.
\]
\end{SRWTproposition}

\begin{proof}
Take expectations, iterate the resulting deterministic recurrence, and sum a
geometric series.
\end{proof}

\subtopic{32.5 The non-collapse condition}

Low off-workspace error is not sufficient for meaningful persistence. The
zero update \(F_t=0\) drives every state exactly into every linear workspace
while destroying all content. Define the internal transport
\[
B_t=P_{t+1}F_tP_t:\mathcal W_t\longrightarrow\mathcal W_{t+1}
\]
and its non-collapse margin
\[
\boxed{
\mu_t^{\mathrm{keep}}
=
\sigma_{\min}(B_t|_{\mathcal W_t}).
}
\]
A content-preserving persistence claim requires
\[
\mu_t^{\mathrm{keep}}>0
\]
for the relevant interval, or an information-theoretic replacement when the
workspace state is stochastic. Structural tracking, state tracking, and
content retention are therefore separate gates.

A finite-horizon persistence certificate may be recorded as
\[
\boxed{
\mathcal P_{0:T}
=
\left(
\max_{t<T}q_t^{\mathrm{str}},
\max_{t\leq T}e_t,
\min_{t<T}\mu_t^{\mathrm{keep}},
\sum_{t<T}\sigma_t
\right).
}
\]
No single coordinate of this vector subsumes the others.

\section{Interpretation of the Contraction Parameter}

The inequality
\[
\alpha_t<1
\]
means that deviations already outside the current workspace are damped by the
next update. Mechanisms that can reduce \(\alpha_t\) or \(\sigma_t\) include
\[
\begin{aligned}
&\text{verification and error feedback},\\
&\text{shared notation and stable interfaces},\\
&\text{artifact persistence and version control},\\
&\text{proof checking and experimental replication},\\
&\text{handoff reconstruction and redundancy}.
\end{aligned}
\]
Sources of disturbance include
\[
\begin{aligned}
&\text{context loss and representational drift},\\
&\text{miscommunication and transfer-map error},\\
&\text{participant or model replacement},\\
&\text{artifact corruption},\\
&\text{unsupported inference and objective drift}.
\end{aligned}
\]

\subtopic{33.1 Four persistence levels}

The theory distinguishes:
\[
\begin{array}{ll}
\text{structural persistence}
& q_t^{\mathrm{str}}\text{ remains small},\\
\text{state persistence}
& e_t\text{ remains small},\\
\text{content persistence}
& \mu_t^{\mathrm{keep}}\text{ remains bounded away from zero},\\
\text{operational persistence}
& \text{predictive and causal capacities remain matched}.
\end{array}
\]
The first three are developed in this part. Operational matching is incorporated
into the identity theory of Chapter~\ref{SRWT:book:chapter-78}.

\subtopic{33.2 Moving-line example}

Let \(\mathcal C_V(t)=\mathbb R^2\) and let \(P_t\) project onto
\[
u_t=(\cos\theta_t,\sin\theta_t).
\]
If \(U_t^V\) is rotation through
\(\theta_{t+1}-\theta_t\) and \(F_t=U_t^V\), then a state
\(x_t=u_t\) obeys
\[
e_t=e_{t+1}=0,
\qquad
\mu_t^{\mathrm{keep}}=1.
\]
The workspace changes raw coordinates but not abstract identity. If instead
\(F_t=I\), then
\[
e_{t+1}
=
|\sin(\theta_{t+1}-\theta_t)|.
\]
The same structural motion now produces a state-tracking defect because the
actual update fails to transport the content with the workspace.

% This file is part of the modular source for Stratified Relational Workspace Theory.
% !TeX root = ../Stratified_Relational_Workspace_Theory.tex

\chapter{Reconstruction and resilience}\label{SRWT:book:chapter-77}


Persistence can be maintained by direct transport, but it can also be restored
from distributed traces after a disruption. Reconstruction theory quantifies
when a coalition contains enough information to recover the abstract core and
how noise is amplified by that recovery.

\section{Core Encoding Maps}

Let \(C_*\) be a finite-dimensional abstract core. The exact static theory
provides a complete local chart
\[
\overline E_i
=
\SRWTev_i\circ\Psi:C_*\longrightarrow C_i,
\]
where \(\Psi:C_*\to\Gamma(H)\) is the root-to-section isomorphism. A participant
or artifact need not expose that complete chart. Let
\[
O_i:C_i\longrightarrow\mathsf Y_i
\]
be the declared observation or access map and define the \emph{accessible
encoding}
\[
\boxed{
E_i=O_i\overline E_i:C_*\longrightarrow\mathsf Y_i.
}
\]
For notational economy, the direct sums below use the accessible spaces
\(\mathsf Y_i\). Thus injectivity of a full structural chart does not make a
singleton operationally reconstructive unless its access map preserves the
whole core.

For a coalition \(S\subseteq V\), define the analysis operator
\[
A_S:C_*\longrightarrow\bigoplus_{i\in S}\mathsf Y_i,
\qquad
A_Sc=(E_ic)_{i\in S},
\]
and the coalition frame operator
\[
\boxed{
F_S=A_S^*A_S=\sum_{i\in S}E_i^*E_i.
}
\]
Set
\[
\alpha_S=\lambda_{\min}(F_S),
\qquad
\beta_S=\lambda_{\max}(F_S).
\]
Thus
\[
\alpha_S\|c\|^2
\leq
\|A_Sc\|^2
\leq
\beta_S\|c\|^2.
\]

\begin{SRWTproposition}[Coalition reconstruction criterion]\label{SRWT:prop:coalition-frame-criterion}
For a coalition \(S\), the following are equivalent:
\begin{enumerate}
\item \(A_S\) is injective;
\item \(F_S\) is positive definite;
\item \(\alpha_S>0\);
\item there exists a linear decoder \(R_S\) with \(R_SA_S=I_{C_*}\).
\end{enumerate}
When these conditions hold, the canonical decoder is
\[
\boxed{
R_S=A_S^\dagger=F_S^{-1}A_S^*.
}
\]
\end{SRWTproposition}

\begin{proof}
In finite dimensions,
\[
\ker F_S=\ker A_S
\]
because
\[
\langle c,F_Sc\rangle=\|A_Sc\|^2.
\]
This proves the equivalence of the first three statements. An injective linear
map admits a left inverse, while a left inverse implies injectivity. The
formula for the Moore--Penrose left inverse follows from full column rank.
\end{proof}

The lower frame bound \(\alpha_S\) measures recoverability. The upper bound
\(\beta_S\) measures amplification in the forward encoding. Their ratio
\[
\kappa_S=\sqrt{\frac{\beta_S}{\alpha_S}}
\]
is the coalition condition number.

\section{Stable Reconstruction Condition}

\subtopic{35.1 Stable reconstruction theorem \(T_7\)}

\begin{SRWTtheorem}[Stable core reconstruction]\label{SRWT:thm:stable-reconstruction}
Assume \(\alpha_S>0\) and observe
\[
y=A_Sc+\eta.
\]
With the canonical estimate
\[
\widehat c=A_S^\dagger y,
\]
one has
\[
\boxed{
\|\widehat c-c\|
\leq
\frac{\|\eta\|}{\sqrt{\alpha_S}}.
}
\]
\end{SRWTtheorem}

\begin{proof}
Since \(A_S^\dagger A_S=I_{C_*}\),
\[
\widehat c-c=A_S^\dagger\eta.
\]
The singular values of \(A_S\) are the square roots of the eigenvalues of
\(F_S\), so
\[
\|A_S^\dagger\|
=
\frac{1}{\sigma_{\min}(A_S)}
=
\frac{1}{\sqrt{\alpha_S}}.
\]
\end{proof}

\begin{SRWTcorollary}[Encoding mismatch]\label{SRWT:cor:encoding-mismatch}
If the true observation is
\[
y=(A_S+\Delta_S)c+\eta
\]
but the nominal decoder \(A_S^\dagger\) is used, then
\[
\boxed{
\|\widehat c-c\|
\leq
\frac{\|\Delta_S\|\,\|c\|+\|\eta\|}{\sqrt{\alpha_S}}.
}
\]
\end{SRWTcorollary}

\begin{proof}
Now
\[
\widehat c-c=A_S^\dagger\Delta_Sc+A_S^\dagger\eta,
\]
and the result follows from Theorem~\ref{SRWT:thm:stable-reconstruction}.
\end{proof}

\subtopic{35.2 Reconstruction as operational recovery}

Geometric recovery becomes behaviorally relevant when the workspace state
controls a predictive kernel. Let
\[
K_w=\mathbb P(Y\in\cdot\mid W=w)
\]
and suppose
\[
\SRWTTV(K_w,K_{w'})\leq L_Kd_W(w,w').
\]

\begin{SRWTproposition}[Exact reconstruction implies coalition sufficiency]
\label{SRWT:prop:exact-reconstruction-sufficiency}
Let \(W\) be sufficient for \(Y\) relative to the full history \(\mathbf H\):
\[
Y\SRWTindep\mathbf H\mid W.
\]
If a coalition reconstructs the workspace exactly,
\[
W=R_S(\mathbf H^S)
\qquad\text{almost surely},
\]
then
\[
\boxed{
Y\SRWTindep\mathbf H\mid\mathbf H^S.
}
\]
Under logarithmic loss,
\[
\boxed{
\SRWTRisk_{\log}(S)=\SRWTRisk_{\log}(V).
}
\]
\end{SRWTproposition}

\begin{proof}
Sufficiency means that the conditional target law given \(\mathbf H\) is a
measurable function of \(W\). Exact reconstruction makes that function
measurable with respect to \(\mathbf H^S\), which is the asserted conditional
independence. The log-risk equality follows from
Theorem~\ref{SRWT:thm:coalition-log-risk}.
\end{proof}

\begin{SRWTproposition}[Reconstruction-to-risk transfer]\label{SRWT:prop:reconstruction-risk-transfer}
Let \(W\) be sufficient for a target \(Y\), let losses take values in
\([0,1]\), and let a coalition construct \(\widehat W\) from its local data.
Then its optimal Bayes risk satisfies
\[
\boxed{
0\leq
\SRWTRisk(S)-\SRWTRisk(W)
\leq
2L_K\,\mathbb E[d_W(W,\widehat W)].
}
\]
\end{SRWTproposition}

\begin{proof}
For every pair \(w,\widehat w\), using the Bayes action for
\(K_{\widehat w}\) under the true law \(K_w\) incurs conditional regret at
most
\[
2\SRWTTV(K_w,K_{\widehat w}).
\]
This follows by inserting and subtracting the risk under
\(K_{\widehat w}\) for both the chosen action and the Bayes action for
\(K_w\). Average the bound and use the Lipschitz assumption. The coalition's
optimal rule performs at least as well as this decoder-based rule.
\end{proof}

This proposition links Chapter~\ref{SRWT:book:chapter-77} to the causal-state theory of Chapter~\ref{SRWT:book:chapter-72}. A small
geometric reconstruction error is operationally meaningful only when the
outcome kernel changes continuously with workspace state.

\section{k-Failure Resilience}

For \(k\geq0\), define the worst surviving lower frame bound
\[
\boxed{
\alpha^{(k)}
=
\min_{\substack{A\subseteq V\\|A|\leq k}}
\alpha_{V\setminus A}.
}
\]
The sequence \(\alpha^{(k)}\) is nonincreasing in \(k\).

\begin{SRWTtheorem}[Uniform erasure resilience]\label{SRWT:thm:uniform-erasure-resilience}
The encoding family is exactly reconstructible after every loss of at most
\(k\) nodes if and only if
\[
\boxed{\alpha^{(k)}>0.}
\]
If surviving observations have the form
\[
y_{V\setminus A}=A_{V\setminus A}c+\eta_A,
\qquad
|A|\leq k,
\]
then the canonical survivor decoder obeys the uniform estimate
\[
\boxed{
\|\widehat c_A-c\|
\leq
\frac{\|\eta_A\|}{\sqrt{\alpha^{(k)}}}.
}
\]
\end{SRWTtheorem}

\begin{proof}
Exact reconstruction for a particular survivor set is equivalent to a
positive lower frame bound by Proposition~\ref{SRWT:prop:coalition-frame-criterion}.
Taking the minimum over all permitted erasures gives the first claim. The
uniform noise estimate follows from Theorem~\ref{SRWT:thm:stable-reconstruction}.
\end{proof}

\begin{SRWTproposition}[Dimensional obstruction]\label{SRWT:prop:dimension-obstruction}
Let \(d=\dim C_*\). If
\[
\sum_{i\in S}\dim \mathsf Y_i<d,
\]
then \(\alpha_S=0\), so \(S\) cannot reconstruct the core. Consequently,
\(k\)-failure resilience requires
\[
\sum_{i\in V\setminus A}\dim \mathsf Y_i\geq d
\]
for every \(A\subseteq V\) with \(|A|\leq k\).
\end{SRWTproposition}

\begin{proof}
The rank of \(A_S\) is at most the dimension of its codomain. If that
dimension is below \(d\), then \(A_S\) is not injective.
\end{proof}

Define the accessible block support of \(c\neq0\) by
\[
\operatorname{supp}_E(c)
=
\{i\in V:E_ic\neq0\}
\]
and define the node-erasure distance
\[
\boxed{
d_{\mathrm{er}}
=
\min_{c\neq0}|\operatorname{supp}_E(c)|.
}
\]

\begin{SRWTtheorem}[Exact erasure-distance criterion]
\label{SRWT:thm:erasure-distance}
Exact reconstruction after every loss of at most \(k\) nodes is possible if
and only if
\[
\boxed{
d_{\mathrm{er}}>k.
}
\]
\end{SRWTtheorem}

\begin{proof}
If \(d_{\mathrm{er}}\leq k\), choose \(c\neq0\) whose nonzero accessible
blocks are contained in a set \(A\) with \(|A|\leq k\). The surviving analysis
operator satisfies \(A_{V\setminus A}c=0\), so it cannot distinguish \(c\)
from zero. Conversely, if some erasure set \(A\) of size at most \(k\) destroys
injectivity, there exists \(c\neq0\) with \(E_ic=0\) for every \(i\notin A\).
Thus \(\operatorname{supp}_E(c)\subseteq A\), and
\(d_{\mathrm{er}}\leq k\).
\end{proof}

The distance criterion is combinatorial. The frame quantity
\(\alpha^{(k)}\) is quantitative. A system may tolerate \(k\) erasures exactly
while amplifying noise severely because \(\alpha^{(k)}\) is close to zero.

\subtopic{36.1 Reconstructive handoff after replacement}

Suppose node \(r\) is removed and the surviving coalition is
\[
S=V\setminus\{r\}.
\]
A replacement node \(r'\) is to receive the encoding
\[
E_{r'}':C_*\to\mathsf Y_{r'}'.
\]

\begin{SRWTcorollary}[Handoff initialization]\label{SRWT:cor:handoff-initialization}
If \(\alpha_S>0\) and the survivors observe \(y=A_Sc+\eta\), initialize the
replacement by
\[
x_{r'}^{\mathrm{init}}
=E_{r'}'A_S^\dagger y.
\]
Then
\[
\boxed{
\|x_{r'}^{\mathrm{init}}-E_{r'}'c\|
\leq
\frac{\|E_{r'}'\|\,\|\eta\|}{\sqrt{\alpha_S}}.
}
\]
\end{SRWTcorollary}

This is the basic mathematical form of artifact-mediated or
coalition-mediated handoff. It does not require the removed component to
transfer its full internal state to the replacement.

\section{Redundancy--Irreducibility Tension}

Define the minimum reconstructing coalition size
\[
\boxed{
m_{\mathrm{rec}}
=
\min\{|S|:\alpha_S>0\}
}
\]
and the maximum exact erasure tolerance
\[
\boxed{
k_{\mathrm{erase}}
=
\max\{k:\alpha^{(k)}>0\}.
}
\]
The condition
\[
m_{\mathrm{rec}}\geq2
\]
means that no individual encoding is sufficient for exact core recovery. The
condition
\[
k_{\mathrm{erase}}\geq1
\]
means that the system survives at least one node loss. These can hold
simultaneously because nonlocality of encoding is not the same as fragility.

\subtopic{37.1 Three-node distributed-frame example}

Let
\[
C_*=\mathbb R^2
\]
and define scalar node encodings
\[
E_1(c_1,c_2)=c_1,
\qquad
E_2(c_1,c_2)=c_2,
\qquad
E_3(c_1,c_2)=c_1+c_2.
\]
Every singleton has rank one, so no node reconstructs the two-dimensional
core. Every pair does reconstruct. For \(S=\{1,2\}\),
\[
F_S=I,
\qquad
\alpha_S=1.
\]
For \(S=\{1,3\}\) or \(S=\{2,3\}\), the frame eigenvalues are
\[
\frac{3\pm\sqrt5}{2}.
\]
Therefore
\[
\boxed{
m_{\mathrm{rec}}=2,
\qquad
k_{\mathrm{erase}}=1,
\qquad
d_{\mathrm{er}}=2,
\qquad
\alpha^{(1)}=\frac{3-\sqrt5}{2}.
}
\]
This realizes a workspace core that is irreducible to any singleton while
remaining reconstructible after any one node fails.

\subtopic{37.2 Relation to predictive irreducibility}

The geometric quantity \(m_{\mathrm{rec}}\) and the predictive quantity
\(k_0\) from Chapter~\ref{SRWT:book:chapter-73} answer different questions. The first asks whether a
coalition can reconstruct a declared core coordinate. The second asks whether
a coalition is sufficient for a declared future target. Equality between them
requires structural--operational alignment and regularity such as
Proposition~\ref{SRWT:prop:reconstruction-risk-transfer}. Neither quantity should
be substituted for the other without those assumptions.

% This file is part of the modular source for Stratified Relational Workspace Theory.
% !TeX root = ../Stratified_Relational_Workspace_Theory.tex

\chapter{Identity through time and replacement}\label{SRWT:book:chapter-78}


Identity is treated as invariance under a structured correspondence, not as
literal persistence of tokens or carriers. Exact correspondences form a
groupoid. Approximate identity is measured by structural, operational,
causal, and reconstructive defects that must be reported separately before
they are combined.

\section{Temporal Realization Maps}

Let
\[
\mathfrak R=(G,\mathbf C,D,\rho,\mathcal I,\mathcal K)
\]
and
\[
\mathfrak R'=(G',\mathbf C',D',\rho',\mathcal I',\mathcal K')
\]
be two realizations. A correspondence begins with a graph isomorphism
\(\varphi:G\to G'\). After relabeling by \(\varphi\), let
\[
J_i:C_i\to C_i'
\]
be unitary local identifications and define
\[
J_V=\bigoplus_{i\in V}J_i,
\qquad
J_E=\bigoplus_{e\in E}J_{t(e)}.
\]

\begin{SRWTdefinition}[Exact structural correspondence]\label{SRWT:def:exact-structural-correspondence}
The correspondence is structurally exact when
\[
\boxed{D'J_V=J_ED.}
\]
It is operation exact for \(a\in\mathcal A\) when
\[
\boxed{\rho_V'(a)J_V=J_V\rho_V(a).}
\]
\end{SRWTdefinition}

\begin{SRWTtheorem}[Exact replacement isomorphism]\label{SRWT:thm:exact-replacement-isomorphism}
Under an exact structural correspondence,
\[
\boxed{
\Gamma(D')=J_V\Gamma(D),
\qquad
L'J_V=J_VL.
}
\]
Consequently, every spectral workspace subspace is transported exactly:
\[
P_I(L')=J_VP_I(L)J_V^*
\]
for every spectral interval \(I\). If the correspondence is also operation
exact, then the transported workspace actions are intertwined exactly.
\end{SRWTtheorem}

\begin{proof}
For every \(c\),
\[
D'J_Vc=J_EDc.
\]
Because \(J_E\) is invertible, \(Dc=0\) if and only if \(D'J_Vc=0\), proving
the kernel statement. Taking adjoints of the exact covariance identity gives
\[
D'^*J_E=J_VD^*.
\]
Therefore
\[
L'J_V=D'^*D'J_V=D'^*J_ED=J_VD^*D=J_VL.
\]
The spectral and operation claims follow by functional calculus and the
intertwining assumptions.
\end{proof}

Exact correspondences compose and invert. They therefore form a groupoid of
realizations. The abstract exact workspace is an isomorphism class in this
groupoid.

\section{Structural Identity Defect}

For an approximate correspondence define
\[
\boxed{
\Delta_D=D'J_V-J_ED,
\qquad
\delta_D=\|\Delta_D\|.
}
\]
If the graph weights are unchanged, define edgewise comparison defects
\[
\delta_e
=
\|H_e'J_{s(e)}-J_{t(e)}H_e\|.
\]
The discrepancy operator has one changed source block per edge, so
\[
\boxed{
\delta_D^2
\leq
\max_{i\in V}
\sum_{e:s(e)=i}w_e\delta_e^2
\leq
\sum_{e\in E}w_e\delta_e^2.
}
\]
When a single node is replaced, every nonincident edge has \(\delta_e=0\).
The replacement defect is therefore locally measurable even though its effect
on the global workspace is controlled by the global spectral gap.

The corresponding carried-Laplacian defect obeys
\[
\boxed{
\eta_L
:=
\|J_V^*L'J_V-L\|
\leq
2\|D\|\delta_D+\delta_D^2.
}
\]
This is Proposition~\ref{SRWT:prop:temporal-laplacian-drift} with the time indices
removed.

Let \(P\) and \(P'\) be rank-\(d\) workspace projectors selected from \(L\)
and \(L'\), and let \(g>0\) be the separating source gap.

\begin{SRWTtheorem}[Approximate structural replacement]\label{SRWT:thm:approx-structural-replacement}
If
\[
\eta_L<\frac g2,
\]
then
\[
\boxed{
q_{\mathrm{str}}
:=
\|P'-J_VPJ_V^*\|
\leq
\frac{\eta_L}{g-\eta_L}
\leq
\frac{2\eta_L}{g}.
}
\]
If \(q_{\mathrm{str}}<1\), then
\[
B=P'J_VP:\SRWTran P\longrightarrow\SRWTran P'
\]
is bijective and satisfies
\[
\sigma_{\min}(B)
\geq
\sqrt{1-q_{\mathrm{str}}^2}.
\]
Its polar factor
\[
\boxed{
\beta
=B(B^*B)^{-1/2}
}
\]
is a canonical unitary workspace correspondence. Moreover,
\[
\boxed{
\|\beta-J_VP\|
\leq
q_{\mathrm{str}}
+
1-\sqrt{1-q_{\mathrm{str}}^2}.
}
\]
\end{SRWTtheorem}

\begin{proof}
The projector estimate is Theorem~\ref{SRWT:thm:spectral-stability} applied after
carrying \(L'\) back by \(J_V\). For unit \(w\in\SRWTran P\),
\[
\|(I-P')J_Vw\|
\leq q_{\mathrm{str}}.
\]
Orthogonality gives
\[
\|Bw\|^2
=1-\|(I-P')J_Vw\|^2
\geq1-q_{\mathrm{str}}^2.
\]
Thus \(B\) is injective, hence bijective because the source and target have
the same finite dimension. Its singular values lie in
\([\sqrt{1-q_{\mathrm{str}}^2},1]\), so polar normalization moves \(B\) by at
most \(1-\sqrt{1-q_{\mathrm{str}}^2}\). Finally,
\[
\|B-J_VP\|=\|(I-P')J_VP\|\leq q_{\mathrm{str}}.
\]
\end{proof}

\subtopic{39.1 Transport of workspace operations}

Write
\[
A=\rho_V(a),
\qquad
A'=\rho_V'(a),
\]
and define
\[
\delta_\rho(a)=\|A'J_V-J_VA\|,
\qquad
\ell_a=\|(I-P)AP\|.
\]
Let
\[
A_{\mathcal W}=PAP,
\qquad
A_{\mathcal W'}'=P'A'P'.
\]

\begin{SRWTproposition}[Operation transport bound]\label{SRWT:prop:operation-transport-bound}
For \(B=P'J_VP\),
\[
\boxed{
\|A_{\mathcal W'}'B-BA_{\mathcal W}\|
\leq
\|A'\|q_{\mathrm{str}}
+\delta_\rho(a)
+\ell_a.
}
\]
For the unitary polar correspondence \(\beta\),
\[
\boxed{
\begin{aligned}
\|A_{\mathcal W'}'\beta-\beta A_{\mathcal W}\|
\leq
&\|A'\|q_{\mathrm{str}}
+\delta_\rho(a)
+\ell_a\\
&+
\bigl(\|A_{\mathcal W'}'\|+\|A_{\mathcal W}\|\bigr)
\left(1-\sqrt{1-q_{\mathrm{str}}^2}\right).
\end{aligned}
}
\]
\end{SRWTproposition}

\begin{proof}
Insert and subtract \(P'A'J_VP\) and \(P'J_VAP\):
\[
\begin{aligned}
A_{\mathcal W'}'B-BA_{\mathcal W}
={}&P'A'(P'-I)J_VP\\
&+P'(A'J_V-J_VA)P\\
&+P'J_V(I-P)AP.
\end{aligned}
\]
Bound the three terms. Replacing \(B\) by \(\beta\) adds
\[
A_{\mathcal W'}'(\beta-B)-(\beta-B)A_{\mathcal W},
\]
which gives the second estimate.
\end{proof}

\section{Behavioral Identity Defect}

Structural correspondence does not determine operational identity. Let
\[
\{\iota_{w,b}:(w,b)\in\mathcal D\}
\quad\text{and}\quad
\{\iota'_{\beta w,\chi b}:(w,b)\in\mathcal D\}
\]
be matched feasible implementation-intervention families in the sense of
Chapter~\ref{SRWT:book:chapter-74}. Here \(b\) denotes declared background variables and \(\chi\) is
the background correspondence. Define
\[
K(w,b)
=
\mathbb P^{\iota_{w,b}}(Y^+\in\cdot),
\]
\[
K'(\beta w,\chi b)
=
\mathbb P^{\prime\,\iota'_{\beta w,\chi b}}(Y^{\prime+}\in\cdot).
\]
For a chosen divergence \(D_Y\), define
\[
\boxed{
\delta_{\mathrm{causal}}(\beta)
=
\sup_{(w,b)\in\mathcal D}
D_Y\left(
K(w,b),
K'(\beta w,\chi b)
\right).
}
\]
The observational response is included only when the null intervention belongs
to the declared family. This definition compares actual implementation
interventions rather than formally assigning an arbitrary projection
coordinate.

A reconstructive handoff may also be imperfect. Let
\[
\delta_{\mathrm{rec}}(\beta)
=
\mathbb E\bigl[d_{\mathcal W'}(W'_{\mathrm{init}},\beta W)\bigr]
\]
measure replacement initialization error. Corollary~\ref{SRWT:cor:handoff-initialization}
provides one explicit upper bound on this term.

\subtopic{40.1 Structural identity does not imply causal identity}

Take two realizations with identical
\[
G,
\quad D,
\quad L,
\quad P,
\quad \rho.
\]
In the first, let \(Y=W\). In the second, let \(Y'=1-W\) for a binary
workspace coordinate. Then
\[
q_{\mathrm{str}}=0,
\qquad
\delta_\rho(a)=0
\]
for every operation, while the observational and interventional outcome
kernels differ maximally under the identity correspondence. Therefore no
purely structural theorem can prove causal identity without an additional
kernel-matching assumption.

\section{Identity Distance and Replacement Invariance}

For an admissible correspondence \(\mathfrak J\), define the operation defect
on a declared unit ball \(\mathcal A_1\) by
\[
\delta_{\mathrm{op}}(\mathfrak J)
=
\sup_{a\in\mathcal A_1}
\|A_{\mathcal W'}'\beta-\beta A_{\mathcal W}\|.
\]
The primary certification object is the identity profile
\[
\boxed{
\mathbf d_{\mathrm{id}}(\mathfrak J)
=
\left(
q_{\mathrm{str}},
\delta_{\mathrm{op}}(\mathfrak J),
\delta_{\mathrm{causal}}(\beta),
\delta_{\mathrm{rec}}(\beta);
\mu^{\mathrm{keep}}
\right).
}
\]
The first four coordinates must remain below declared tolerances, while the
non-collapse margin \(\mu^{\mathrm{keep}}\) must remain above a declared lower
bound. This vector, rather than a weighted average, determines whether a
specific identity claim passes.

A composite identity cost is
\[
\boxed{
\operatorname{Cost}(\mathfrak J)
=
a\,q_{\mathrm{str}}
+b\,\delta_{\mathrm{op}}(\mathfrak J)
+c\,\delta_{\mathrm{causal}}(\beta)
+d\,\delta_{\mathrm{rec}}(\beta),
}
\]
where \(a,b,c,d>0\) are declared scales. Define
\[
\boxed{
d_{\mathrm{id}}(\mathfrak R,\mathfrak R')
=
\inf_{\mathfrak J}
\operatorname{Cost}(\mathfrak J).
}
\]
This quantity is a task-relative optimization surrogate until the
admissible correspondence class, intervention matching, divergence, and scales
are fixed. It does not replace the coordinatewise identity profile or the
non-collapse gate.

\subtopic{41.1 Replacement invariance theorem \(T_{14}\) }

\begin{SRWTtheorem}[Quantitative replacement invariance]\label{SRWT:thm:replacement-invariance}
Let \(\mathfrak J\) be a unitary graph correspondence with diagram defect
\(\delta_D\). Define
\[
\eta_L=2\|D\|\delta_D+\delta_D^2
\]
and assume
\[
\eta_L<\frac g2.
\]
Set
\[
q_*
=
\frac{\eta_L}{g-\eta_L}.
\]
Then there exists a unitary workspace correspondence \(\beta\) for which
\[
q_{\mathrm{str}}\leq q_*.
\]
Moreover,
\[
\boxed{
\begin{aligned}
\delta_{\mathrm{op}}(\mathfrak J)
\leq
\sup_{a\in\mathcal A_1}
\Bigl[&
\|\rho_V'(a)\|q_*
+\delta_\rho(a)
+\ell_a\\
&+
\bigl(\|A_{\mathcal W'}'\|+\|A_{\mathcal W}\|\bigr)
\left(1-\sqrt{1-q_*^2}\right)
\Bigr].
\end{aligned}
}
\]
If, in addition,
\[
\delta_{\mathrm{causal}}(\beta)\leq\delta_c,
\qquad
\delta_{\mathrm{rec}}(\beta)\leq\delta_r,
\]
then
\[
\boxed{
\begin{aligned}
d_{\mathrm{id}}(\mathfrak R,\mathfrak R')
\leq
&a\,q_*+b\,\delta_{\mathrm{op}}(\mathfrak J)\\
&+c\,\delta_c+d\,\delta_r.
\end{aligned}
}
\]
In particular, if all local structural, operation, closure, causal, and
reconstruction defects are \(O(\delta_{\mathrm{rep}})\), the relevant
operator norms are uniformly bounded, and \(g\geq g_0>0\), then
\[
\boxed{
d_{\mathrm{id}}(\mathfrak R,\mathfrak R')
\leq
C_1\frac{\delta_{\mathrm{rep}}}{g_0}
+C_2\delta_{\mathrm{rep}}
+O\!\left(\frac{\delta_{\mathrm{rep}}^2}{g_0}\right).
}
\]
After a fixed normalization of the non-spectral terms, this recovers the
schematic target
\[
d_{\mathrm{id}}
\lesssim
C\frac{\delta_{\mathrm{rep}}}{g_0}.
\]
\end{SRWTtheorem}

\begin{proof}
The structural bound and existence of \(\beta\) follow from
Theorem~\ref{SRWT:thm:approx-structural-replacement}. The operation estimate is
Proposition~\ref{SRWT:prop:operation-transport-bound} with
\(q_{\mathrm{str}}\leq q_*\). Insert these estimates and the assumed causal
and reconstructive bounds into the definition of
\(\operatorname{Cost}(\mathfrak J)\). The asymptotic statement follows from
\[
q_*=O\!\left(\frac{\delta_D}{g_0}
+\frac{\delta_D^2}{g_0}\right)
\]
and
\[
1-\sqrt{1-q_*^2}=O(q_*^2).
\]
\end{proof}

\begin{SRWTcorollary}[Preservation of coalition frame bounds]
\label{SRWT:cor:replacement-frame-bound}
Identify the old and new abstract cores with \(\SRWTran P\) and \(\SRWTran P'\)
through their chosen core parameterizations. Let \(S\) be a declared coalition,
let
\[
A_S:\SRWTran P\to\mathsf Y_S,
\qquad
A_S':\SRWTran P'\to\mathsf Y_S'
\]
be the resulting accessible coalition encoders, and let \(K_S\) be a unitary
identification of their observation spaces. Suppose
\[
\|A_S'\beta-K_SA_S\|
\leq
\delta_{S}^{\mathrm{frame}}
<
\sqrt{\alpha_S}.
\]
Then the replacement coalition remains reconstructive and its new lower frame
bound satisfies
\[
\boxed{
\alpha_S'
\geq
\left(
\sqrt{\alpha_S}-\delta_{S}^{\mathrm{frame}}
\right)^2.
}
\]
Consequently, \(k\)-failure resilience survives whenever the bound holds
uniformly for every coalition remaining after at most \(k\) erasures.
\end{SRWTcorollary}

\begin{proof}
For \(w'\in\SRWTran P'\), write \(w'=\beta w\). Since \(\beta\) and \(K_S\) are
unitary,
\[
\begin{aligned}
\|A_S'w'\|
&\geq
\|K_SA_Sw\|
-
\|(A_S'\beta-K_SA_S)w\|\\
&\geq
\left(
\sqrt{\alpha_S}-\delta_{S}^{\mathrm{frame}}
\right)\|w'\|.
\end{aligned}
\]
Square the inequality.
\end{proof}

\begin{SRWTcorollary}[Exact carrier replacement]\label{SRWT:cor:exact-carrier-replacement}
If
\[
\delta_D=0,
\qquad
\delta_\rho(a)=0,
\qquad
\ell_a=0,
\qquad
\delta_c=0,
\qquad
\delta_r=0,
\]
then
\[
d_{\mathrm{id}}(\mathfrak R,\mathfrak R')=0.
\]
Thus a participant, model, artifact, or notation system may be replaced without
changing the abstract workspace when the replacement realizes the same
structural, operational, causal, and reconstructive organization.
\end{SRWTcorollary}

\subtopic{41.2 Temporal chains and identity paths}

For a sequence of realizations
\[
\mathfrak R_0,\mathfrak R_1,\ldots,\mathfrak R_T,
\]
let \(\mathfrak J_t:\mathfrak R_t\to\mathfrak R_{t+1}\) be admissible
correspondences. Define the path length
\[
\operatorname{Len}(\mathfrak J_{0:T})
=
\sum_{t=0}^{T-1}\operatorname{Cost}(\mathfrak J_t)
\]
and the induced directed path quasi-pseudometric
\[
\boxed{
d_{\mathrm{path}}^{\rightarrow}(\mathfrak R_0,\mathfrak R_T)
=
\inf_{\mathfrak J_{0:T}}
\operatorname{Len}(\mathfrak J_{0:T}).
}
\]
Concatenation gives the directed triangle inequality. Because replacement
costs and intervention divergences need not be symmetric, symmetry is not
assumed. When reverse correspondences exist, define the symmetrized
pseudometric
\[
\boxed{
d_{\mathrm{path}}^{\mathrm{sym}}(\mathfrak R,\mathfrak R')
=
\max\left\{
 d_{\mathrm{path}}^{\rightarrow}(\mathfrak R,\mathfrak R'),
 d_{\mathrm{path}}^{\rightarrow}(\mathfrak R',\mathfrak R)
\right\}.
}
\]
Distinct physical realizations can have zero symmetrized pseudodistance when
an exact correspondence exists in both directions. This is intentional:
identity is assigned to the invariant relational organization rather than to
the carrier tokens.

\subtopic{41.3 Identity failure modes}

A complete replacement audit must report at least
\[
\boxed{
\mathcal I_{\mathrm{rep}}
=
\left(
q_{\mathrm{str}},
\delta_{\mathrm{op}},
\delta_{\mathrm{causal}},
\delta_{\mathrm{rec}},
\mu^{\mathrm{keep}}
\right).
}
\]
The failure modes are different:
\[
\begin{array}{ll}
q_{\mathrm{str}}\text{ large}
&\text{the shared subspace changed},\\
\delta_{\mathrm{op}}\text{ large}
&\text{the same operations are no longer realized},\\
\delta_{\mathrm{causal}}\text{ large}
&\text{interventions no longer have matching consequences},\\
\delta_{\mathrm{rec}}\text{ large}
&\text{the new realization was not initialized from the old core},\\
\mu^{\mathrm{keep}}\text{ small}
&\text{content collapsed despite apparent coherence}.
\end{array}
\]
The theorem does not permit one successful coordinate to conceal failure of
another.

% This file is part of the modular source for Stratified Relational Workspace Theory.
% !TeX root = ../Stratified_Relational_Workspace_Theory.tex

\chapter{Gluing and workspaces of workspaces}\label{SRWT:book:chapter-79}


The preceding parts treat one interaction diagram at a time. A large
collaborative system is usually assembled from overlapping local laboratories,
subteams, artifacts, models, or task-specific workspaces. This part asks when
those local workspaces determine one coherent global workspace, how failed
gluing is detected, and whether recursive assembly changes the resulting
object. The exact theory is a descent theorem. The approximate theory is a
second transfer-Laplacian construction on the cover itself.

\section{Local Workspace Covers}

\subtopic{42.1 Edge-complete covers}

Let
\[
G=(V,E)
\]
be a finite directed interaction graph with local core spaces \(C_i\), core
transfers \(H_e\), vertex space \(\mathcal C_V(G)\), discrepancy operator
\(D_G\), and exact workspace
\[
\mathscr W(G)
:=
\Gamma(H|_G)
=
\ker D_G.
\]

\begin{SRWTdefinition}[Finite workspace cover]\label{SRWT:def:workspace-cover}
A finite family of subgraphs
\[
\mathcal U=\{G_\alpha\}_{\alpha\in A}
\]
is an \emph{edge-complete workspace cover} of \(G\) when
\[
V=\bigcup_{\alpha\in A}V_\alpha,
\qquad
E=\bigcup_{\alpha\in A}E_\alpha.
\]
Every edge in \(E_\alpha\) carries the transfer inherited from \(G\). For a
nonempty finite index set
\[
I=\{\alpha_0,\ldots,\alpha_p\}\subseteq A,
\]
define the intersection subgraph
\[
G_I
=
G_{\alpha_0\cdots\alpha_p}
:=
\bigcap_{\alpha\in I}G_\alpha
\]
by intersecting both vertex and edge sets. Isolated vertices are retained.
\end{SRWTdefinition}

Covering only the vertices is insufficient. If an edge belongs to no chart,
its transfer equation is absent from every local workspace and cannot be
recovered by overlap agreement. Edge completeness is therefore the minimal
cover assumption used below.

If \(K\subseteq L\subseteq G\) are subgraphs, coordinate restriction gives
\[
\rho_K^L:\mathcal C_V(L)\longrightarrow\mathcal C_V(K).
\]
Because every edge equation on \(K\) is also an edge equation on \(L\),
restriction maps exact sections to exact sections:
\[
r_K^L
:=
\rho_K^L|_{\mathscr W(L)}:
\mathscr W(L)\longrightarrow\mathscr W(K).
\]
They obey
\[
r_L^L=I,
\qquad
r_J^K r_K^L=r_J^L
\]
whenever \(J\subseteq K\subseteq L\). Thus \(K\mapsto\mathscr W(K)\) is a
contravariant system of local workspaces.

\subtopic{42.2 The cover nerve and overlap interfaces}

The nerve \(N(\mathcal U)\) has one vertex for each chart \(G_\alpha\) and a
\(p\)-simplex
\[
[\alpha_0\cdots\alpha_p]
\]
whenever \(G_{\alpha_0\cdots\alpha_p}\) is nonempty. The nerve records which
local workspaces can compare states directly. It does not replace the
underlying interaction graph: \(G\) carries the internal transfer equations,
while \(N(\mathcal U)\) carries overlap-consistency equations between already
constructed local objects.

For pairwise overlaps write
\[
\mathscr W_\alpha=\mathscr W(G_\alpha),
\qquad
\mathscr W_{\alpha\beta}=\mathscr W(G_{\alpha\beta}),
\]
and
\[
r_{\alpha\beta}^{\alpha}:
\mathscr W_\alpha\to\mathscr W_{\alpha\beta},
\qquad
r_{\alpha\beta}^{\beta}:
\mathscr W_\beta\to\mathscr W_{\alpha\beta}.
\]
The interface is therefore a cospan
\[
\mathscr W_\alpha
\xrightarrow{\ r_{\alpha\beta}^{\alpha}\ }
\mathscr W_{\alpha\beta}
\xleftarrow{\ r_{\alpha\beta}^{\beta}\ }
\mathscr W_\beta,
\]
not necessarily a direct transfer map from one complete local workspace to the
other. Both charts are compared through the content they expose on their
common domain.

\subtopic{42.3 \v{C}ech cochains of local workspaces}

For \(p\geq0\), define the finite cochain space
\[
\check C^p(\mathcal U;\mathscr W)
=
\bigoplus_{\alpha_0<\cdots<\alpha_p}
\mathscr W(G_{\alpha_0\cdots\alpha_p}),
\]
where empty intersections contribute the zero space. The coboundary
\[
\delta^p:\check C^p\longrightarrow\check C^{p+1}
\]
is
\[
\boxed{
(\delta^p s)_{\alpha_0\cdots\alpha_{p+1}}
=
\sum_{j=0}^{p+1}(-1)^j
r_{\alpha_0\cdots\alpha_{p+1}}^%
{\alpha_0\cdots\widehat{\alpha_j}\cdots\alpha_{p+1}}
%
s_{\alpha_0\cdots\widehat{\alpha_j}\cdots\alpha_{p+1}}.
}
\]
The restriction identities imply
\[
\delta^{p+1}\delta^p=0.
\]
For a local family \(w=(w_\alpha)_\alpha\),
\[
\boxed{
(\delta^0w)_{\alpha\beta}
=
r_{\alpha\beta}^{\beta}w_\beta
-
r_{\alpha\beta}^{\alpha}w_\alpha.
}
\]
Thus \(\delta^0\) is the overlap-discrepancy operator at the level of exact
local workspaces.

\section{Compatibility and Gluing Obstructions}

\subtopic{43.1 Direct overlap mismatch}

A family \(w=(w_\alpha)_\alpha\) is exactly compatible when
\[
\delta^0w=0.
\]
Its direct mismatch is
\[
\boxed{
\mathfrak m_{\mathcal U}(w)=\delta^0w.
}
\]
After choosing inner products and positive overlap weights, the mismatch
energy is
\[
\boxed{
\mathcal E_{\mathrm{ov}}(w)
=
\|\delta^0w\|^2
=
\sum_{\alpha<\beta}
\nu_{\alpha\beta}
\left\|
 r_{\alpha\beta}^{\beta}w_\beta
-r_{\alpha\beta}^{\alpha}w_\alpha
\right\|^2.
}
\]
This is the appropriate obstruction when the full local states are already
given. Such a family glues exactly if and only if its mismatch vector itself
vanishes. No cohomological quotient is needed for that direct question.

\subtopic{43.2 Relative-only and affine descent data}

A different problem arises when only relative transition data are supplied.
Let
\[
b=(b_{\alpha\beta})\in\check C^1(\mathcal U;\mathscr W)
\]
be a prescribed family of overlap offsets. The triple-overlap consistency
condition is
\[
\delta^1b=0.
\]
One seeks local potentials \(q=(q_\alpha)\in\check C^0\) satisfying
\[
\delta^0q=b.
\]
Define
\[
\boxed{
\check H^1(\mathcal U;\mathscr W)
=
\frac{\ker\delta^1}{\operatorname{im}\delta^0}.
}
\]
Then the affine gluing problem is solvable precisely when
\[
\boxed{[b]=0\quad\text{in}\quad\check H^1(\mathcal U;\mathscr W).}
\]
A nonzero class means that all local transition statements can be mutually
consistent on triple overlaps while still failing to arise from any family of
local workspace states. This is a cover-level analogue of nontrivial holonomy.

\subtopic{43.3 Harmonic representative of the obstruction}

Equip the finite cochain spaces with Hilbert inner products and define the
first \v{C}ech Laplacian
\[
\boxed{
\Delta_1
=
\delta^0(\delta^0)^*
+
(\delta^1)^*\delta^1.
}
\]
Let
\[
\mathcal H^1_{\mathcal U}=\ker\Delta_1.
\]

\begin{SRWTproposition}[Harmonic gluing obstruction]\label{SRWT:prop:harmonic-gluing-obstruction}
For every cocycle \(b\in\ker\delta^1\), there is an orthogonal decomposition
\[
\boxed{
b=\delta^0q^\dagger+h,
\qquad
h\in\mathcal H^1_{\mathcal U},
}
\]
where \(q^\dagger=(\delta^0)^\dagger b\) is a least-squares potential. Moreover,
\[
\boxed{
\inf_q\|b-\delta^0q\|
=
\|h\|
=
\|P_{\mathcal H^1_{\mathcal U}}b\|.
}
\]
Hence \(b\) is exactly integrable if and only if its harmonic component
vanishes.
\end{SRWTproposition}

\begin{proof}
The finite-dimensional Hodge decomposition gives
\[
\check C^1
=
\operatorname{im}\delta^0
\oplus
\ker\Delta_1
\oplus
\operatorname{im}(\delta^1)^*.
\]
If \(b\in\ker\delta^1\), then \(b\) is orthogonal to
\(\operatorname{im}(\delta^1)^*\). Its remaining two orthogonal components are
therefore an exact part \(\delta^0q^\dagger\) and a harmonic part \(h\). The
orthogonal projection theorem gives the least-squares identity.
\end{proof}

The scalar
\[
\omega_{\mathcal U}(b)
:=
\|P_{\mathcal H^1_{\mathcal U}}b\|
\]
is an irreducible overlap-frustration measure. Its vanishing is
invariant under admissible invertible coordinate changes, although its numerical
norm depends on the chosen metric. It cannot be removed by adding local
potentials alone.

\subtopic{43.4 Cycle example}

Consider three charts arranged in a cycle, with one-dimensional overlap
spaces and oriented offsets
\[
b_{12},\qquad b_{23},\qquad b_{31}.
\]
Local potentials would have to satisfy
\[
q_2-q_1=b_{12},
\qquad
q_3-q_2=b_{23},
\qquad
q_1-q_3=b_{31}.
\]
Adding the equations gives the necessary and sufficient condition
\[
\boxed{
b_{12}+b_{23}+b_{31}=0.}
\]
With the standard edge norm, the harmonic obstruction has magnitude
\[
\boxed{
\omega_{\mathcal U}(b)
=
\frac{|b_{12}+b_{23}+b_{31}|}{\sqrt3}.
}
\]
Every pairwise relation may be individually meaningful while the complete
cycle remains globally inconsistent.

\section{Exact Gluing Theorem T8}

Let
\[
R_{\mathcal U}:\mathscr W(G)
\longrightarrow
\check C^0(\mathcal U;\mathscr W)
\]
be the restriction map
\[
R_{\mathcal U}(w)
=
(w|_{G_\alpha})_{\alpha\in A}.
\]

\begin{SRWTtheorem}[Exact workspace gluing]\label{SRWT:thm:exact-workspace-gluing}
For every finite edge-complete workspace cover \(\mathcal U\), the sequence
\[
\boxed{
0
\longrightarrow
\mathscr W(G)
\xrightarrow{\ R_{\mathcal U}\ }
\check C^0(\mathcal U;\mathscr W)
\xrightarrow{\ \delta^0\ }
\check C^1(\mathcal U;\mathscr W)
}
\]
is exact. Equivalently,
\[
\boxed{
\mathscr W(G)
\cong
\ker\delta^0
\cong
\operatorname{Eq}
\left(
\prod_\alpha\mathscr W_\alpha
\rightrightarrows
\prod_{\alpha<\beta}\mathscr W_{\alpha\beta}
\right).
}
\]
Every pairwise compatible family of exact local workspace states therefore
glues to a unique exact global workspace state.
\end{SRWTtheorem}

\begin{proof}
The cover contains every vertex, so \(R_{\mathcal U}\) is injective. A global
section has identical restrictions on every overlap, hence
\(\delta^0R_{\mathcal U}=0\).

Conversely, let \(w=(w_\alpha)_\alpha\in\ker\delta^0\). For each vertex
\(i\in V\), choose any chart \(G_\alpha\) containing \(i\) and set
\[
c_i=(w_\alpha)_i.
\]
If another chart \(G_\beta\) contains \(i\), then \(i\in G_{\alpha\beta}\)
and overlap compatibility gives
\((w_\alpha)_i=(w_\beta)_i\). Thus \(c_i\) is well defined.

For every edge \(e:i\to j\), edge completeness provides a chart
\(G_\alpha\) containing \(e\). Since \(w_\alpha\in\mathscr W(G_\alpha)\),
\[
c_j=(w_\alpha)_j=H_e(w_\alpha)_i=H_ec_i.
\]
Hence \(c=(c_i)\in\mathscr W(G)\), and by construction
\(R_{\mathcal U}c=w\). Uniqueness follows because the charts cover all
vertices.
\end{proof}

\begin{SRWTcorollary}[Zeroth descent cohomology]\label{SRWT:cor:zeroth-descent}
For every finite edge-complete cover,
\[
\boxed{
\mathscr W(G)
\cong
\check H^0(\mathcal U;\mathscr W)
=
\ker\delta^0.
}
\]
In finite dimensions,
\[
\boxed{
\dim\mathscr W(G)
=
\sum_{\alpha\in A}\dim\mathscr W_\alpha
-
\operatorname{rank}(\delta^0).
}
\]
\end{SRWTcorollary}

The theorem does not require the nerve to be acyclic. Cycles in the cover are
allowed; their effect is encoded in the rank of \(\delta^0\) and, for
relative-only data, in \(\check H^1\). It also shows that the local-workspace
assignment is a genuine sheaf on the finite subgraph cover used here rather
than merely a presheaf.

\subtopic{44.1 Local nontriviality does not imply global nontriviality}

Let \(G\) be the directed triangle
\[
1\xrightarrow{\ a\ }2
\xrightarrow{\ b\ }3
\xrightarrow{\ c\ }1
\]
with one-dimensional vertex cores. Cover it by its three edge charts. Their
local workspaces are all one-dimensional:
\[
\begin{aligned}
\mathscr W_{12}&=\{(x,ax):x\in\mathbb R\},\\
\mathscr W_{23}&=\{(y,by):y\in\mathbb R\},\\
\mathscr W_{31}&=\{(z,cz):z\in\mathbb R\}.
\end{aligned}
\]
Overlap agreement requires
\[
y=ax,
\qquad
z=by,
\qquad
x=cz.
\]
Therefore
\[
(1-abc)x=0,
\]
and
\[
\boxed{
\dim\mathscr W(G)
=
\begin{cases}
1,&abc=1,\\
0,&abc\neq1.
\end{cases}
}
\]
The meta-level gluing obstruction is exactly the original cycle holonomy.
Strong local workspaces can therefore coexist with a trivial global core.

\section{Approximate Gluing and Meta-Workspaces}

\subtopic{45.1 The cover-level transfer Laplacian}

First assume each local state already lies in its exact local workspace. Define
\[
B_{\mathcal U}
:=
\delta^0:
\check C^0(\mathcal U;\mathscr W)
\longrightarrow
\check C^1(\mathcal U;\mathscr W)
\]
and the meta-Laplacian
\[
\boxed{
L_{\mathcal U}^{\mathrm{meta}}
=
B_{\mathcal U}^*B_{\mathcal U}.
}
\]
By Theorem~\ref{SRWT:thm:exact-workspace-gluing},
\[
\boxed{
\ker L_{\mathcal U}^{\mathrm{meta}}
=
R_{\mathcal U}\mathscr W(G).
}
\]
Thus the gluing problem is itself a workspace problem whose ``nodes'' are
local workspaces and whose interfaces are overlaps.

Let
\[
\gamma_{\mathcal U}^{\mathrm{meta}}
=
\lambda_{\min}^+
\left(L_{\mathcal U}^{\mathrm{meta}}\right)
\]
be the first positive eigenvalue.

\begin{SRWTproposition}[Approximate overlap gluing]\label{SRWT:prop:approx-overlap-gluing}
For every \(w\in\check C^0(\mathcal U;\mathscr W)\),
\[
\boxed{
\operatorname{dist}
\left(
 w,
 R_{\mathcal U}\mathscr W(G)
\right)^2
\leq
\frac{\|B_{\mathcal U}w\|^2}
{\gamma_{\mathcal U}^{\mathrm{meta}}}.
}
\]
\end{SRWTproposition}

\begin{proof}
Decompose \(w\) into its component in \(\ker L_{\mathcal U}^{\mathrm{meta}}\)
and its orthogonal complement. On the complement,
\(L_{\mathcal U}^{\mathrm{meta}}\geq
\gamma_{\mathcal U}^{\mathrm{meta}}I\). Therefore
\[
\|B_{\mathcal U}w\|^2
=
\langle w,L_{\mathcal U}^{\mathrm{meta}}w\rangle
\geq
\gamma_{\mathcal U}^{\mathrm{meta}}
\operatorname{dist}
\left(w,\ker L_{\mathcal U}^{\mathrm{meta}}\right)^2.
\]
\end{proof}

The gap \(\gamma_{\mathcal U}^{\mathrm{meta}}\) measures how strongly the cover
penalizes disagreement. A small overlap mismatch need not imply a nearby
global workspace when this gap is itself very small.

\subtopic{45.2 Simultaneous local and overlap error}

Local candidates may also violate their internal transfer equations. Define
the local-copy space
\[
\widehat{\mathcal C}_{\mathcal U}
=
\bigoplus_{\alpha\in A}\mathcal C_V(G_\alpha)
\]
and the block internal discrepancy
\[
\widehat D_{\mathcal U}
=
\bigoplus_{\alpha\in A}D_{G_\alpha}.
\]
For \(u=(u_\alpha)_\alpha\in\widehat{\mathcal C}_{\mathcal U}\), define the
raw overlap synchronization operator
\[
(S_{\mathcal U}u)_{\alpha\beta}
=
\rho_{G_{\alpha\beta}}^{G_\beta}u_\beta
-
\rho_{G_{\alpha\beta}}^{G_\alpha}u_\alpha.
\]
For an overlap weight \(\lambda_{\mathrm{ov}}>0\), set
\[
\boxed{
K_{\mathcal U}
=
\begin{pmatrix}
\widehat D_{\mathcal U}\\[0.2em]
\sqrt{\lambda_{\mathrm{ov}}}\,S_{\mathcal U}
\end{pmatrix},
\qquad
L_{\mathcal U}^{\mathrm{tot}}
=
K_{\mathcal U}^*K_{\mathcal U}.
}
\]
Its energy is
\[
\boxed{
\begin{aligned}
\mathcal E_{\mathcal U}^{\mathrm{tot}}(u)
={}&
\sum_\alpha\|D_{G_\alpha}u_\alpha\|^2\\
&+
\lambda_{\mathrm{ov}}
\sum_{\alpha<\beta}
\left\|
\rho^{G_\beta}_{G_{\alpha\beta}}u_\beta
-
\rho^{G_\alpha}_{G_{\alpha\beta}}u_\alpha
\right\|^2.
\end{aligned}
}
\]

Let
\[
J_{\mathcal U}:\mathcal C_V(G)\to\widehat{\mathcal C}_{\mathcal U},
\qquad
J_{\mathcal U}c=(c|_{G_\alpha})_\alpha,
\]
be the duplication map.

\begin{SRWTtheorem}[Total descent kernel]\label{SRWT:thm:total-descent-kernel}
For every finite edge-complete cover,
\[
\boxed{
\ker K_{\mathcal U}
=
J_{\mathcal U}\mathscr W(G).
}
\]
Consequently, if
\[
\gamma_{\mathcal U}^{\mathrm{tot}}
=
\lambda_{\min}^+
\left(L_{\mathcal U}^{\mathrm{tot}}\right),
\]
then
\[
\boxed{
\operatorname{dist}
\left(
 u,
 J_{\mathcal U}\mathscr W(G)
\right)^2
\leq
\frac{\mathcal E_{\mathcal U}^{\mathrm{tot}}(u)}
{\gamma_{\mathcal U}^{\mathrm{tot}}}.
}
\]
\end{SRWTtheorem}

\begin{proof}
If \(u=J_{\mathcal U}c\) for a global exact section, every local discrepancy
and every overlap discrepancy vanishes. Conversely, if \(K_{\mathcal U}u=0\),
then every \(u_\alpha\) is an exact local section and the family agrees on all
overlaps. Theorem~\ref{SRWT:thm:exact-workspace-gluing} gives a unique
\(c\in\mathscr W(G)\) with \(u=J_{\mathcal U}c\). The distance estimate follows
from the spectral lower bound on the orthogonal complement of the kernel.
\end{proof}

This theorem combines two distinct failure modes:
\[
\text{local incoherence}
\qquad\text{and}\qquad
\text{cross-chart disagreement}.
\]
Neither can be hidden by the other because they occupy separate blocks of
\(K_{\mathcal U}\).

\subtopic{45.3 Least-squares global assembly}

Let
\[
m_i=|\{\alpha:i\in V_\alpha\}|
\]
be the multiplicity of vertex \(i\). Then
\[
J_{\mathcal U}^*J_{\mathcal U}=M_{\mathcal U},
\qquad
(M_{\mathcal U}c)_i=m_ic_i.
\]
The pseudoinverse performs coordinatewise averaging:
\[
\boxed{
(J_{\mathcal U}^\dagger u)_i
=
\frac{1}{m_i}
\sum_{\alpha:i\in V_\alpha}(u_\alpha)_i.
}
\]
Let
\[
\widehat u
=
P_{\ker K_{\mathcal U}}u
=
J_{\mathcal U}\widehat c.
\]
Then
\[
\widehat c
=
J_{\mathcal U}^\dagger\widehat u
\]
and
\[
\boxed{
\|u-J_{\mathcal U}\widehat c\|
\leq
\frac{\|K_{\mathcal U}u\|}
{\sqrt{\gamma_{\mathcal U}^{\mathrm{tot}}}}.
}
\]
If \(m_{\min}=\min_i m_i\), then
\[
\boxed{
\left\|
\widehat c-J_{\mathcal U}^\dagger u
\right\|
\leq
\frac{\|K_{\mathcal U}u\|}
{\sqrt{m_{\min}\gamma_{\mathcal U}^{\mathrm{tot}}}}.
}
\]
The simple overlap average and the globally coherent least-squares assembly
are close only when the total descent energy is small relative to the cover
gap.

\subtopic{45.4 A workspace of workspaces}

Define the exact meta-workspace
\[
\boxed{
\mathscr W_{\mathrm{meta}}(\mathcal U)
:=
\ker B_{\mathcal U}
\subseteq
\bigoplus_\alpha\mathscr W_\alpha.
}
\]
Theorem~\ref{SRWT:thm:exact-workspace-gluing} gives the canonical identification
\[
\boxed{
\mathscr W_{\mathrm{meta}}(\mathcal U)
\cong
\mathscr W(G).
}
\]
Thus exact formation of a workspace from local workspaces does not create an
additional mysterious carrier. It gives a second presentation of the same
global compatible-section object.

Suppose local operations
\[
A_\alpha:\mathscr W_\alpha\to\mathscr W_\alpha
\]
and overlap operations
\[
A_{\alpha\beta}:\mathscr W_{\alpha\beta}\to
\mathscr W_{\alpha\beta}
\]
satisfy
\[
r_{\alpha\beta}^{\alpha}A_\alpha
=
A_{\alpha\beta}r_{\alpha\beta}^{\alpha},
\qquad
r_{\alpha\beta}^{\beta}A_\beta
=
A_{\alpha\beta}r_{\alpha\beta}^{\beta}.
\]
Then, with \(A^0=\bigoplus_\alpha A_\alpha\) and
\(A^1=\bigoplus_{\alpha<\beta}A_{\alpha\beta}\),
\[
B_{\mathcal U}A^0=A^1B_{\mathcal U}.
\]
Therefore
\[
A^0\ker B_{\mathcal U}\subseteq\ker B_{\mathcal U}.
\]
The operation-closure theorem from Chapter~\ref{SRWT:book:chapter-69} reappears at the meta level:
operations on subworkspaces descend globally precisely when they agree on
interfaces.

\subtopic{45.5 Descent-faithful compression}

A practical hierarchy rarely transmits each complete local workspace. Let
\[
Q^p:\check C^p(\mathcal U;\mathscr W)\to\overline C^p
\]
be compression maps into smaller local and overlap representations. Suppose
there is a compressed coboundary \(\overline\delta^p\) such that
\[
\boxed{
Q^{p+1}\delta^p
=
\overline\delta^p Q^p.
}
\]
The maps \(Q^\bullet\) then form a cochain map.

\begin{SRWTdefinition}[Descent-faithful compression]\label{SRWT:def:descent-faithful}
The compression is \emph{descent-faithful in degree zero} when the induced map
\[
\boxed{
H^0(Q):
\ker\delta^0
\longrightarrow
\ker\overline\delta^0
}
\]
is an isomorphism.
\end{SRWTdefinition}

Injectivity prevents two distinct global workspace states from collapsing to
the same compressed meta-state. Surjectivity prevents the compressed hierarchy
from introducing spurious compatible states that have no full-resolution
lift. Under this condition,
\[
\boxed{
\ker\overline\delta^0
\cong
\mathscr W(G).
}
\]
This is the exact criterion for a compressed workspace of workspaces to retain
the global core.

\subtopic{45.6 Structural and operational gluing are distinct}

Exact or approximate overlap coherence does not by itself establish global
predictive sufficiency. Let \(W_t\) denote a glued structural state and let
\[
Z_t^{\mathrm{loc}}
=
(Z_{\alpha,t})_{\alpha\in A}
\]
be local predictive summaries recovered from its chart restrictions.

\begin{SRWTproposition}[Local-to-global predictive alignment]\label{SRWT:prop:local-global-predictive}
Assume \(W_t\) is measurable with respect to the complete history
\(\mathbf H_t\), every \(Z_{\alpha,t}\) is a measurable function of
\(W_t|_{G_\alpha}\), and
\[
I(Y_t^+;\mathbf H_t\mid Z_t^{\mathrm{loc}})
\leq
\varepsilon.
\]
Then
\[
\boxed{
I(Y_t^+;\mathbf H_t\mid W_t)
\leq
\varepsilon.
}
\]
In particular, exact sufficiency of the complete local-summary tuple implies
exact sufficiency of the glued workspace.
\end{SRWTproposition}

\begin{proof}
The tuple \(Z_t^{\mathrm{loc}}\) is a measurable function of \(W_t\), while
\(W_t\) is itself a history statistic. Conditioning on the finer statistic
cannot increase the residual information that the full history carries about
the target. Equivalently, use the conditional-law characterization from
Theorem~\ref{SRWT:thm:predictive-state}.
\end{proof}

Causal gluing requires an additional interface gate: interventions described
in two overlapping charts must induce matched outcome laws on the overlap. A
structurally compatible family can fail this gate, just as structurally
identical replacement systems can fail causal identity in Chapter~\ref{SRWT:book:chapter-78}.

\subtopic{45.7 Anti-centralization gate}

The isomorphism
\[
\mathscr W_{\mathrm{meta}}(\mathcal U)\cong\mathscr W(G)
\]
is a compositional statement, not a relational-irreducibility statement. A
single chart may still expose enough accessible information to reconstruct the
whole meta-state. A claim that the hierarchy is genuinely distributed must
therefore reuse the coalition gates from Chapters~\ref{SRWT:book:chapter-73} and~\ref{SRWT:book:chapter-77}. At the chart level,
one may require, for example,
\[
\alpha_{\{\alpha\}}^{\mathrm{meta}}=0
\quad\text{for every }\alpha,
\qquad
\alpha_A^{\mathrm{meta}}>0
\quad\text{for some }|A|\geq2,
\]
or the predictive analogue
\[
k_\varepsilon^{\mathrm{meta}}\geq2,
\qquad
\Delta_{\mathrm{joint}}^{\mathrm{meta}}>0.
\]
Capacity-matched local-memory and chart-ablation controls are still required.
Hierarchical notation cannot convert a centralized copy into a relational
workspace.

\section{Hierarchical Collapse Theorem T9}

\subtopic{46.1 Finite cover hierarchies}

A finite workspace hierarchy is a rooted tree \(\mathcal T\). Each node
\(a\in\mathcal T\) carries a subgraph \(G_a\). The root \(r\) carries
\[
G_r=G.
\]
For every internal node \(a\), its children
\[
\operatorname{Ch}(a)
\]
form a finite edge-complete cover of \(G_a\). Leaves carry primitive local
workspace spaces
\[
W_\ell=\mathscr W(G_\ell).
\]
For an internal node, define recursively
\[
\boxed{
W_a
=
\ker
\left(
\delta_a^0:
\bigoplus_{b\in\operatorname{Ch}(a)}W_b
\longrightarrow
\bigoplus_{b<c}W_{bc}
\right),
}
\]
where the overlap spaces and restrictions are inherited from the child cover.

\begin{SRWTtheorem}[Exact hierarchical collapse]\label{SRWT:thm:hierarchical-collapse}
For every finite workspace hierarchy as above, there is a canonical
isomorphism
\[
\boxed{
W_a\cong\mathscr W(G_a)
}
\]
for every node \(a\). In particular,
\[
\boxed{
W_r\cong\mathscr W(G).
}
\]
Therefore recursively gluing local workspaces, using any other finite
parenthesization built from edge-complete intermediate covers that retain all
required overlap objects, and solving the flattened global transfer system all
produce the same exact workspace up to canonical isomorphism.
\end{SRWTtheorem}

\begin{proof}
Proceed upward from the leaves. The claim holds at every leaf by definition.
Assume it holds for all children of an internal node \(a\). Replacing each
child space \(W_b\) by its canonical copy \(\mathscr W(G_b)\), the recursive
definition of \(W_a\) is exactly the compatible-family kernel for the
edge-complete child cover of \(G_a\). Theorem~\ref{SRWT:thm:exact-workspace-gluing}
therefore gives
\[
W_a\cong\mathscr W(G_a).
\]
Finite induction reaches the root.
\end{proof}

\subtopic{46.2 Flattening and limit interchange}

Let \(L(\mathcal T)\) be the leaf set. Repeated edge completeness implies that
the flattened family
\[
\mathcal U_{\mathrm{flat}}
=
\{G_\ell:\ell\in L(\mathcal T)\}
\]
is itself an edge-complete cover of \(G\). Let
\[
\delta_{\mathrm{flat}}^0
\]
be its overlap coboundary and define
\[
W_{\mathrm{flat}}
=
\ker\delta_{\mathrm{flat}}^0.
\]
Theorem~\ref{SRWT:thm:exact-workspace-gluing} and
Theorem~\ref{SRWT:thm:hierarchical-collapse} give
\[
\boxed{
W_r
\cong
W_{\mathrm{flat}}
\cong
\mathscr W(G).
}
\]
Equivalently,
\[
\boxed{
\operatorname{Lim}_{\mathcal T}^{\mathrm{rec}}
\left(
\{\mathscr W(G_\ell)\}_{\ell\in L(\mathcal T)}
\right)
\cong
\operatorname{Eq}_{\mathcal U_{\mathrm{flat}}}
\left(
\{\mathscr W(G_\ell)\}_{\ell\in L(\mathcal T)}
\right).
}
\]
The notation on the left means recursive finite descent along the tree; the
right side means one-shot equalization over the flattened leaf cover. In
finite-dimensional vector spaces, both are kernels of the same complete family
of local and overlap constraints, grouped in different orders. Hierarchical
collapse is therefore a finite-limit interchange property, not a claim that
intermediate levels are physically absent or operationally irrelevant.

\subtopic{46.3 Approximate hierarchical collapse}

At level \(r\), let \(K_r\) be the corresponding total descent operator,
\[
\gamma_r
=
\lambda_{\min}^+(K_r^*K_r),
\qquad
\|K_ru_r\|\leq\varepsilon_r.
\]
The nearest exact compatible state then differs from \(u_r\) by at most
\[
\frac{\varepsilon_r}{\sqrt{\gamma_r}}.
\]
Allow an additional implementation, compression, or causal-summary defect
\(\zeta_r\), and define
\[
\delta_r
=
\frac{\varepsilon_r}{\sqrt{\gamma_r}}
+
\zeta_r.
\]

\begin{SRWTtheorem}[Approximate hierarchical collapse bound]\label{SRWT:thm:approx-hierarchical-collapse}
Let
\[
F_r:\mathcal H_r\to\mathcal H_{r+1},
\qquad
\widetilde F_r:\mathcal H_r\to\mathcal H_{r+1}
\]
be exact and implemented level-collapse maps for
\(r=0,\ldots,R-1\). Let \(\mathcal R_r\subseteq\mathcal H_r\) contain the exact
and implemented states reachable at level \(r\). Assume, on \(\mathcal R_r\),
\[
\|F_r(x)-F_r(y)\|
\leq
L_r\|x-y\|
\]
and
\[
\sup_{x\in\mathcal R_r}
\|\widetilde F_r(x)-F_r(x)\|
\leq
\delta_r.
\]
Then
\[
\boxed{
\begin{aligned}
&\left\|
\widetilde F_{R-1}\circ\cdots\circ\widetilde F_0(x)
-
F_{R-1}\circ\cdots\circ F_0(x)
\right\|\\
&\qquad\leq
\sum_{r=0}^{R-1}
\left(
\prod_{s=r+1}^{R-1}L_s
\right)
\delta_r.
\end{aligned}
}
\]
In particular, if \(L_r\leq1\), then the total error is at most
\[
\sum_{r=0}^{R-1}\delta_r.
\]
If \(L_r\leq L<1\) and \(\delta_r\leq\delta\), then
\[
\boxed{
\text{total error}
\leq
\frac{\delta(1-L^R)}{1-L}
\leq
\frac{\delta}{1-L}.
}
\]
\end{SRWTtheorem}

\begin{proof}
Let \(x_r\) and \(\widetilde x_r\) be the exact and implemented states after
\(r\) levels. Then
\[
\begin{aligned}
\|\widetilde x_{r+1}-x_{r+1}\|
&\leq
\|\widetilde F_r(\widetilde x_r)-F_r(\widetilde x_r)\|\\
&\quad+
\|F_r(\widetilde x_r)-F_r(x_r)\|\\
&\leq
\delta_r+L_r\|\widetilde x_r-x_r\|.
\end{aligned}
\]
Iterating this recurrence gives the displayed product-sum estimate. The final
two bounds follow by summing a finite or infinite geometric series.
\end{proof}

This theorem identifies three regimes:
\[
\begin{array}{ll}
L_r<1
&\text{hierarchical correction is contractive},\\
L_r=1
&\text{level errors accumulate at most additively},\\
L_r>1
&\text{deep hierarchies can amplify early mismatch}.
\end{array}
\]
A deep hierarchy is therefore stable only when its level gaps remain separated
and its assembly maps do not magnify small overlap errors faster than those
errors are corrected.

\subtopic{46.4 Causal-kernel accumulation}

Let \(K_r\) and \(\widetilde K_r\) now denote exact and implemented Markov
kernels at successive hierarchy levels, and suppose
\[
\sup_x
\operatorname{TV}
\left(
\widetilde K_r(x,\cdot),
K_r(x,\cdot)
\right)
\leq
\zeta_r.
\]
Because Markov kernels contract total variation, a telescoping argument gives
\[
\boxed{
\sup_x
\operatorname{TV}
\left(
\widetilde K_0\cdots\widetilde K_{R-1}(x,\cdot),
K_0\cdots K_{R-1}(x,\cdot)
\right)
\leq
\sum_{r=0}^{R-1}\zeta_r.
}
\]
Thus causal semantics can also be transported through a hierarchy, but only
under an explicit kernel-level compatibility test. Structural gluing alone
provides no such guarantee.

\subtopic{46.5 Direct versus hierarchical computation}

The exact theorem states
\[
\boxed{
\text{recursive exact gluing}
=
\text{one-shot exact gluing}.
}
\]
The approximate theorem instead supplies a measurable flattening defect:
\[
\boxed{
\delta_{\mathrm{flat}}^{(R)}
=
\left\|
\widetilde F_{R-1}\cdots\widetilde F_0(x)
-
F_{R-1}\cdots F_0(x)
\right\|.
}
\]
A hierarchy is empirically justified when it reduces computation,
communication, or local model complexity while keeping
\(\delta_{\mathrm{flat}}^{(R)}\), predictive risk, causal-kernel mismatch, and
anti-centralization gates within declared tolerances. Hierarchy by itself is
not evidence of emergence.

\subtopic{46.6 Gluing and hierarchical synthesis}

The exact and approximate constructions can be compressed as
\[
\boxed{
\begin{aligned}
\text{local transfer systems}
&\xrightarrow{\ \ker D_\alpha\ }
\text{local workspaces}\\
&\xrightarrow{\ \ker\delta^0\ }
\text{meta-workspace}\\
&\xrightarrow{\ \text{finite recursive descent}\ }
\text{global workspace}.
\end{aligned}
}
\]
At exact resolution, all finite parenthesizations are canonically equivalent.
At approximate resolution, every level contributes a defect scaled by its
spectral gap and by downstream amplification. Relative-only overlap data can
carry a harmonic \(H^1\) obstruction even when every local chart is internally
coherent. Operational and relational claims still require predictive, causal,
coalition, and null-model gates in addition to structural descent.

% This file is part of the modular source for Stratified Relational Workspace Theory.
% !TeX root = ../Stratified_Relational_Workspace_Theory.tex

\chapter{Scaffolded and endogenous formation}\label{SRWT:book:chapter-80}


\section{Formation Systems and Scaffold Specification}

The preceding parts begin with a transfer system that has already been chosen.
A formation theory must instead explain how the nodes, edges, transfer maps,
workspace channels, and operational boundaries are obtained. This requires a
strict separation between the structure supplied by the experimenter and the
structure selected from interaction history.

\begin{SRWTdefinition}[Formation system]\label{SRWT:def:formation-system}
A \emph{workspace formation system} is a tuple
\[
\mathfrak F
=
\bigl(
\Sigma,
\mathcal D_{\mathrm{form}},
\mathcal L,
\Omega,
\widehat\theta,
\mathcal D_{\mathrm{cert}},
\mathcal I_{\mathrm{form}}
\bigr),
\]
where:
\begin{itemize}
\item \(\Sigma\) is the externally supplied scaffold;
\item \(\mathcal D_{\mathrm{form}}\) is the formation or training history;
\item \(\mathcal L\) is the learning or selection rule;
\item \(\Omega\) is its internal random seed or stochastic state;
\item \(\widehat\theta=\mathcal L(\Sigma,\mathcal D_{\mathrm{form}},\Omega)\)
 is the learned organization;
\item \(\mathcal D_{\mathrm{cert}}\) is held-out certification data; and
\item \(\mathcal I_{\mathrm{form}}\) is a declared family of interventions on
 scaffold and formation history.
\end{itemize}
The learned organization is decomposed as
\[
\widehat\theta
=
\bigl(
\widehat V,
\widehat E,
\widehat H,
\widehat P,
\widehat{\mathcal O},
\widehat B
\bigr),
\]
where \(\widehat B\) is a learned operational boundary.
\end{SRWTdefinition}

The scaffold is itself stratified:
\[
\boxed{
\Sigma
=
\bigl(
\Sigma_V,
\Sigma_E,
\Sigma_H,
\Sigma_P,
\Sigma_{\mathcal O},
\Sigma_B,
\Sigma_Y
\bigr).
}
\]
The final coordinate \(\Sigma_Y\) records labels, targets, privileged masks,
or other supervision that could directly reveal the intended workspace. This
coordinate is necessary because an architecture may appear to discover a
boundary even when the desired membership mask was encoded in its training
labels.

\subtopic{47.1 Scaffolded architecture versus scaffolded solution}

A scaffold can constrain the search space without specifying its solution.
For example, a graph neural architecture can fix the class of admissible local
message-passing rules while leaving the active edges and transfer maps to be
learned. Conversely, a system can use a flexible architecture while receiving
labels that directly encode the desired workspace. The former is
architecturally scaffolded but solution-open. The latter is solution-scaffolded.

For any component \(q\) of \(\widehat\theta\), define its scaffold-only Bayes
risk under loss \(\ell_q\) by
\[
\mathcal R_q^{\Sigma}
=
\inf_f
\mathbb E
\left[
\ell_q\bigl(q,f(\Sigma,\Omega)\bigr)
\right].
\]
Define the full formation risk by
\[
\mathcal R_q^{\Sigma,\mathcal H}
=
\inf_g
\mathbb E
\left[
\ell_q\bigl(q,g(\Sigma,\mathcal D_{\mathrm{form}},\Omega)\bigr)
\right].
\]
The history advantage is
\[
\boxed{
\Delta_q^{\mathrm{hist}}
=
\mathcal R_q^{\Sigma}
-
\mathcal R_q^{\Sigma,\mathcal H}.
}
\]
A positive history advantage means that formation history contains information
about \(q\) that is not already recoverable from the declared scaffold and
random seed.

\begin{SRWTproposition}[Scaffold-only impossibility]\label{SRWT:prop:scaffold-only-impossibility}
If
\[
q=f(\Sigma,\Omega)
\qquad\text{almost surely},
\]
then for every formation-history intervention that leaves \((\Sigma,\Omega)\)
fixed,
\[
\mathcal L\!\left(q\mid
\operatorname{do}(\mathcal D_{\mathrm{form}}=h)
\right)
=
\mathcal L\!\left(q\mid
\operatorname{do}(\mathcal D_{\mathrm{form}}=h')
\right),
\]
and \(\Delta_q^{\mathrm{hist}}=0\).
\end{SRWTproposition}

\begin{proof}
The value of \(q\) is measurable with respect to \((\Sigma,\Omega)\), so an
intervention on formation history cannot change its conditional law when those
variables are held fixed. The scaffold-only predictor \(f\) attains the same
risk as any predictor that also receives formation history, hence the risk gap
vanishes.
\end{proof}

This proposition is elementary but methodologically decisive. A structure
that is deterministically encoded by the scaffold cannot count as discovered
from interaction history merely because a learned network reproduces it.

\subtopic{47.2 Formation and certification split}

The same sample must not be used both to select a workspace and to certify its
coherence, causal efficacy, or relational irreducibility. The required split is
\[
\mathcal D
=
\mathcal D_{\mathrm{form}}
\sqcup
\mathcal D_{\mathrm{cal}}
\sqcup
\mathcal D_{\mathrm{cert}}.
\]
The formation sample selects \(\widehat\theta\). The calibration sample fixes
thresholds, regularization strengths, and intervention magnitudes. The
certification sample evaluates the final gates without further structural
selection. In a multi-environment study, the stronger requirement is
\[
\mathcal E_{\mathrm{form}}
\cap
\mathcal E_{\mathrm{cert}}
=
\varnothing,
\]
so that the learned organization is tested under environments not used to
construct it.

\begin{SRWTdefinition}[Formation leakage]\label{SRWT:def:formation-leakage}
A certification variable \(Y_{\mathrm{cert}}\) leaks into formation when the
conditional law of \(\widehat\theta\) changes under an intervention on
\(Y_{\mathrm{cert}}\) that should be inaccessible during formation:
\[
\mathcal L
\left(
\widehat\theta
\mid
\operatorname{do}(Y_{\mathrm{cert}}=y)
\right)
\neq
\mathcal L
\left(
\widehat\theta
\mid
\operatorname{do}(Y_{\mathrm{cert}}=y')
\right).
\]
A valid discovery claim requires a documented information barrier or an audit
showing that this effect is absent up to tolerance.
\end{SRWTdefinition}

\subtopic{47.3 Origin and reality are orthogonal}

The formation provenance of a workspace and its operational certification are
separate coordinates. A fully scaffolded structure may satisfy every
structural and causal workspace gate. A highly endogenous structure may fail
to become coherent or causally useful. Symbolically,
\[
\operatorname{Endogenous}(W)
\centernot\Longrightarrow
\operatorname{SRW}_{\Theta}(W),
\]
\[
\operatorname{SRW}_{\Theta}(W)
\centernot\Longrightarrow
\operatorname{Endogenous}(W).
\]
The theory therefore treats endogeneity as a provenance axis rather than a
higher level on the workspace hierarchy.

\section{Operational Endogeneity Signature}

The ideal description-complexity expression from the symbolic outline is
retained as a conceptual limit, but raw Kolmogorov complexity is not
computable. A usable theory must declare a model class, a code, and an
intervention family.

\begin{SRWTdefinition}[Code-relative history contribution]\label{SRWT:def:mdl-endogeneity}
Let \(\mathfrak M_q\) be a declared coding or predictive model class for a
learned component \(q\). Let
\[
L_{\mathfrak M_q}(q\mid Z)
\]
be its held-out codelength given side information \(Z\). Define
\[
\boxed{
 e_q^{\mathrm{MDL}}
 =
 \left[
 \frac{
 L_{\mathfrak M_q}(q\mid\Sigma)
 -
 L_{\mathfrak M_q}(q\mid\Sigma,\mathcal D_{\mathrm{form}})
 }{
 \max\{L_{\mathfrak M_q}(q\mid\Sigma),1\}
 }
 \right]_{0}^{1},
}
\]
where \([x]_0^1=\min\{1,\max\{0,x\}\}\).
\end{SRWTdefinition}

This quantity is explicitly code-relative. It must not be presented as an
observer-independent metaphysical measure. Its role is to test whether the
formation history provides compressive information about the learned
organization beyond the declared scaffold.

\begin{SRWTdefinition}[Causal formation sensitivity]\label{SRWT:def:causal-formation-sensitivity}
Let \(\iota_h\) and \(\iota_{h'}\) be matched feasible interventions that alter
formation history while preserving the scaffold, randomization law, sample
size, and intervention cost. For a divergence \(D\), define
\[
\boxed{
 e_q^{\mathrm{causal}}
 =
 \mathbb E_{(h,h')\sim\nu_q}
 D\!\left(
 \mathcal L^{\iota_h}(q),
 \mathcal L^{\iota_{h'}}(q)
 \right).
}
\]
A common null uses temporally or relationally shuffled histories that preserve
local marginals but destroy the candidate formation signal.
\end{SRWTdefinition}

\begin{SRWTproposition}[Reparameterization invariance of causal provenance]
\label{SRWT:prop:causal-provenance-invariance}
Let \(\phi\) be a measurable bijection on the range of \(q\), and let the
divergence \(D\) be invariant under measurable bijections. Then
\[
 e_{\phi(q)}^{\mathrm{causal}}
 =
 e_q^{\mathrm{causal}}.
\]
\end{SRWTproposition}

\begin{proof}
Every intervention-induced law of \(\phi(q)\) is the pushforward of the
corresponding law of \(q\). Bijection invariance of \(D\) preserves each
pairwise divergence, and therefore their design-weighted expectation.
\end{proof}

The causal coordinate is therefore insensitive to a mere renaming of the
learned structure. The MDL coordinate need not be invariant unless the coding
class is transformed with the reparameterization.

\subtopic{48.1 Scaffold dependence and cross-scaffold invariance}

Endogeneity does not require independence from all scaffolds. Every learning
system has an architecture, objective, and resource budget. The relevant
question is whether the learned organization depends on the particular
solution-bearing details of the scaffold.

For a family \(\mathfrak S\) of admissible scaffold variants, define the
cross-scaffold identity defect
\[
\boxed{
\delta_q^{\mathrm{scaf}}
=
\inf_{\{J_{\sigma\sigma'}\}}
\mathbb E_{\sigma,\sigma'\sim\pi}
 d_q
\left(
J_{\sigma\sigma'}q_{\sigma},q_{\sigma'}
\right),
}
\]
where \(J_{\sigma\sigma'}\) transports the learned object between admissible
coordinate systems. A low value indicates that the same abstract
organization is recovered under multiple non-solution-bearing scaffolds.

The complete formation signature is
\[
\boxed{
\mathbf E_{\mathrm{form}}(W)
=
\left(
 e_V,e_E,e_H,e_P,e_{\mathcal O},e_B;
 \delta^{\mathrm{scaf}},
 \delta^{\mathrm{leak}}
\right),
}
\]
where every \(e_q\) is reported as the pair
\[
 e_q
=
\left(
 e_q^{\mathrm{MDL}},
 e_q^{\mathrm{causal}}
\right).
\]
No scalar average is permitted to hide a zero history contribution for one of
the structures claimed to have been discovered.

\subtopic{48.2 Formation certificate}

For thresholds \(\Theta_{\mathrm{form}}\), define
\[
\operatorname{Form}_{\Theta_{\mathrm{form}}}(W)=1
\]
only when all applicable gates pass:
\[
\begin{aligned}
& e_q^{\mathrm{MDL}}\geq \tau_q^{\mathrm{MDL}},
&& e_q^{\mathrm{causal}}\geq \tau_q^{\mathrm{causal}},\\
& \delta_q^{\mathrm{scaf}}\leq\varepsilon_q^{\mathrm{scaf}},
&& \delta^{\mathrm{leak}}\leq\varepsilon_{\mathrm{leak}},\\
& \operatorname{SRW}_{\Theta}(W)=1,
&& \operatorname{Pass}(N_{\mathrm{form}})=1.
\end{aligned}
\]
The final workspace gate is included because endogenous discovery of an
unstable or causally irrelevant pattern is not endogenous workspace formation.

\section{Learned Transfer Structure and Formation Classes}

\subtopic{49.1 Gated membership and edge learning}

Let
\[
z_i\in\{0,1\},
\qquad
a_e\in[0,1]
\]
be node-membership and edge-activation variables satisfying
\[
a_e\leq z_{s(e)},
\qquad
a_e\leq z_{t(e)}.
\]
A weighted learned discrepancy operator is
\[
(D_{a,H}c)_e
=
\sqrt{w_ea_e}
\left(
 c_{t(e)}-H_ec_{s(e)}
\right).
\]
One generic formation objective is
\[
\boxed{
\begin{aligned}
\widehat\theta
\in
\SRWTargmin_{z,a,H,P,\mathcal O,B}
\Bigl
\{&
\widehat{\mathcal R}_{\mathrm{task}}
+
\lambda_{\mathrm{coh}}\|D_{a,H}c\|^2
+
\lambda_E\|a\|_1
+
\lambda_V\|z\|_1\\
&+
\lambda_{\mathrm{stab}}\mathcal J_{\mathrm{stab}}
+
\lambda_{\mathrm{leak}}\mathcal J_{\mathrm{leak}}
\Bigr\}.
\end{aligned}
}
\]
The objective is schematic rather than universal. Its purpose is to expose
which quantities were fixed and which were optimized. If \(z\) and \(a\) are
fixed by \(\Sigma\), only the transfer maps or local channels can be claimed as
learned.

\subtopic{49.2 Gauge-identifiable transfer learning}

A learned transfer map is not absolutely identifiable when local coordinates
are latent. Under local invertible changes
\[
X_u'=g_uX_u,
\qquad
X_v'=g_vX_v,
\]
the same abstract transfer is represented by
\[
H_{uv}'=g_vH_{uv}g_u^{-1}.
\]
The identifiable object is therefore the gauge class
\[
[H]_{\mathcal G}
=
\left\{
(g_{t(e)}H_eg_{s(e)}^{-1})_{e\in E}
:
(g_i)_{i\in V}\in\mathcal G
\right\}.
\]
Any formation claim that compares raw matrices across independently trained
systems without gauge alignment is ill posed.

\subtopic{49.3 Transfer identification theorem \(T_{15}\) }

\begin{SRWTtheorem}[Population transfer identification]\label{SRWT:thm:transfer-identification}
Let \(X\in\mathbb R^{d_u}\) and \(Y\in\mathbb R^{d_v}\) satisfy
\[
Y=H_*X+\xi,
\qquad
\mathbb E[\xi X^{\mathsf T}]=0,
\]
and assume
\[
\Sigma_{XX}
=
\mathbb E[XX^{\mathsf T}]
\succeq
\kappa I
\]
for some \(\kappa>0\). Then \(H_*\) is the unique minimizer of
\[
\mathcal R(H)
=
\mathbb E\|Y-HX\|_2^2,
\]
with
\[
\boxed{
H_*
=
\Sigma_{YX}\Sigma_{XX}^{-1}.
}
\]
Moreover,
\[
\boxed{
\mathcal R(H)-\mathcal R(H_*)
=
\operatorname{tr}
\left(
(H-H_*)\Sigma_{XX}(H-H_*)^{\mathsf T}
\right)
\geq
\kappa\|H-H_*\|_F^2.
}
\]
\end{SRWTtheorem}

\begin{proof}
Substitute \(Y=H_*X+\xi\) into the risk. The cross term vanishes because
\(\mathbb E[\xi X^{\mathsf T}]=0\), giving
\[
\mathcal R(H)
=
\mathcal R(H_*)
+
\mathbb E\|(H-H_*)X\|_2^2.
\]
The remaining expectation is the displayed trace. Positive definiteness of
\(\Sigma_{XX}\) gives uniqueness and the lower bound. The normal equations
give the covariance formula.
\end{proof}

\begin{SRWTcorollary}[Deterministic empirical recovery bound]
\label{SRWT:cor:empirical-transfer-bound}
For samples \((X_r,Y_r)_{r=1}^n\), write
\[
\widehat\Sigma_{XX}
=
\frac1n\sum_{r=1}^nX_rX_r^{\mathsf T},
\qquad
\widehat M
=
\frac1n\sum_{r=1}^n\xi_rX_r^{\mathsf T}.
\]
If \(\widehat\Sigma_{XX}\succeq\widehat\kappa I\), then ordinary least squares
satisfies
\[
\boxed{
\|\widehat H-H_*\|_F
\leq
\frac{\|\widehat M\|_F}{\widehat\kappa}.
}
\]
For ridge parameter \(\tau>0\),
\[
\boxed{
\|\widehat H_\tau-H_*\|_F
\leq
\frac{
\|\widehat M\|_F+\tau\|H_*\|_F
}{
\widehat\kappa+\tau
}.
}
\]
\end{SRWTcorollary}

\begin{proof}
The empirical normal equations give
\[
\widehat H-H_*
=
\widehat M\widehat\Sigma_{XX}^{-1}.
\]
For ridge regression,
\[
\widehat H_\tau-H_*
=
(\widehat M-\tau H_*)
(\widehat\Sigma_{XX}+\tau I)^{-1}.
\]
Taking Frobenius and operator norms yields both estimates.
\end{proof}

\begin{SRWTproposition}[Gauge covariance of least squares]
\label{SRWT:prop:ols-gauge-covariance}
Under invertible coordinate transformations
\[
X'=g_uX,
\qquad
Y'=g_vY,
\]
the population least-squares solution transforms as
\[
H_*'=g_vH_*g_u^{-1}.
\]
\end{SRWTproposition}

\begin{proof}
The transformed covariances are
\[
\Sigma_{Y'X'}=g_v\Sigma_{YX}g_u^{\mathsf T},
\qquad
\Sigma_{X'X'}=g_u\Sigma_{XX}g_u^{\mathsf T}.
\]
Substitution into the covariance formula and inversion of the second factor
give the result.
\end{proof}

The theorem establishes identification only under excitation and exogeneity.
When \(\Sigma_{XX}\) is singular, the transfer is identified only on the
excited source subspace. When \(\mathbb E[\xi X^{\mathsf T}]\neq0\), regression
recovers a predictive map rather than necessarily the intended causal
transfer.

\subtopic{49.4 Held-out structure-selection theorem \(T_{16}\) }

The same finite-margin argument applies to node membership, edge selection,
boundary selection, or a finite family of gauge-normalized transfer models.

\begin{SRWTtheorem}[Uniform structure-selection recovery]
\label{SRWT:thm:uniform-structure-selection}
Let \(\mathfrak T\) be a finite candidate family with population score
\(J:\mathfrak T\to\mathbb R\). Let
\[
\theta_*
=
\SRWTargmin_{\theta\in\mathfrak T}J(\theta)
\]
be unique, with margin
\[
\Delta_J
=
\min_{\theta\neq\theta_*}
\bigl(J(\theta)-J(\theta_*)\bigr)>0.
\]
If a held-out estimate \(\widehat J\) satisfies
\[
\sup_{\theta\in\mathfrak T}
|\widehat J(\theta)-J(\theta)|
\leq\varepsilon,
\]
then every empirical minimizer \(\widehat\theta\) obeys
\[
\boxed{
J(\widehat\theta)-J(\theta_*)
\leq2\varepsilon.
}
\]
If \(2\varepsilon<\Delta_J\), then
\[
\boxed{\widehat\theta=\theta_*.}
\]
\end{SRWTtheorem}

\begin{proof}
By empirical optimality and uniform approximation,
\[
J(\widehat\theta)
\leq
\widehat J(\widehat\theta)+\varepsilon
\leq
\widehat J(\theta_*)+\varepsilon
\leq
J(\theta_*)+2\varepsilon.
\]
If \(\widehat\theta\neq\theta_*\), the margin would force an excess larger than
\(2\varepsilon\), a contradiction.
\end{proof}

The theorem does not guarantee that the candidate class contains the correct
organization. It states exactly what follows once a finite class, score, and
uniform certification error are declared.

\subtopic{49.5 Formation classes}

The provenance classes are now defined relative to
\[
(\Sigma,\mathcal L,\mathcal I_{\mathrm{form}},\mathfrak M).
\]
They are not absolute labels.

\begin{description}[leftmargin=2.1em,style=nextline]
\item[Exogenously scaffolded.]
The claimed nodes, edges, transfer maps, channels, operations, or boundary are
recoverable from the scaffold alone at the declared tolerance. Their history
advantages and causal formation sensitivities vanish.

\item[Adaptively parameterized.]
The membership and channel topology are fixed, while transfer parameters,
local operators, or readouts are learned from interaction history. Typically
\[
e_H+e_P+e_{\mathcal O}>0,
\qquad
e_V\approx e_E\approx e_B\approx0.
\]

\item[Structurally adaptive.]
At least one of membership, edges, or boundary is selected from history, but
the search family and objective remain externally specified.

\item[Endogenously constituted relative to the scaffold class.]
Membership, transfer organization, and operational boundary have nonzero
history contribution and causal formation sensitivity, survive cross-scaffold
transport, pass leakage and matched-null tests, and produce a certified
workspace on held-out environments.
\end{description}

The phrase ``relative to the scaffold class'' is mandatory. No finite
experiment can establish independence from every possible external influence.

% This file is part of the modular source for Stratified Relational Workspace Theory.
% !TeX root = ../Stratified_Relational_Workspace_Theory.tex

\chapter{Boundary formation and stability}\label{SRWT:book:chapter-81}


\section{Candidate Workspace Boundaries}

A boundary is not merely a subset of physical components. It is a subset
relative to a declared node decomposition, prediction target, intervention
family, environment distribution, and time scale.

\begin{SRWTdefinition}[Boundary problem]\label{SRWT:def:boundary-problem}
A boundary problem is a tuple
\[
\mathfrak B
=
\left(
V,
\mathfrak S,
\tau,
\mathcal Y,
\mathcal I,
\Pi,
\mathcal G_B
\right),
\]
where \(\mathfrak S\subseteq2^V\) is the admissible candidate family,
\(\tau\) is the temporal scale, \(\mathcal Y\) is the operational target,
\(\mathcal I\) is the feasible intervention family, \(\Pi\) is the
environment distribution, and \(\mathcal G_B\) is a symmetry group under
which boundaries are considered equivalent.
\end{SRWTdefinition}

Unless the application requires otherwise, the candidate family excludes
\[
\varnothing,
\qquad
V,
\qquad
\{i\}\ \text{for }i\in V.
\]
This anti-triviality rule prevents the empty set, the entire realization, or a
single centralized controller from winning solely through a degenerate score.

\subtopic{50.1 Observational recurrence, dependence, and outward influence}

For \(S\in\mathfrak S\), define
\[
\SRWTRet_\tau(S)
=
I\!\left(
X_S(t+\tau);X_S(t)
\mid X_{\bar S}(t),E_t
\right),
\]
\[
\SRWTDep_\tau(S)
=
I\!\left(
X_S(t+\tau);X_{\bar S}(t)
\mid X_S(t),E_t
\right),
\]
\[
\SRWTOut_\tau(S)
=
I\!\left(
X_{\bar S}(t+\tau);X_S(t)
\mid X_{\bar S}(t),E_t
\right).
\]
These quantities are observational. They can be nonzero because of common
causes, selection, or feedback paths and therefore cannot alone establish a
causal boundary.

\subtopic{50.2 Causal boundary terms}

Let \(\iota_{S,x}\) be a feasible intervention that changes the internal
workspace coordinate of \(S\) to \(x\) while matching background variables and
intervention cost. Define outward causal efficacy
\[
\boxed{
\SRWTCOut_\tau(S)
=
\mathbb E_{x,x'\sim\nu_S}
D\!\left(
\mathcal L^{\iota_{S,x}}
(X_{\bar S}(t+\tau)),
\mathcal L^{\iota_{S,x'}}
(X_{\bar S}(t+\tau))
\right).
}
\]
Define inward causal dependence analogously by intervening on \(\bar S\) and
measuring the future law of \(X_S\):
\[
\boxed{
\SRWTCIn_\tau(S)
=
\mathbb E_{u,u'\sim\nu_{\bar S}}
D\!\left(
\mathcal L^{\iota_{\bar S,u}}
(X_S(t+\tau)),
\mathcal L^{\iota_{\bar S,u'}}
(X_S(t+\tau))
\right).
}
\]
Neither term should be minimized to zero in general. A useful workspace may
receive substantial external input. The boundary objective must instead
express the intended balance between internal recurrence, external dependence,
and outward efficacy.

\subtopic{50.3 Structural and anti-centralization gates}

For the induced subsystem on \(S\), define
\[
\mathsf{Coh}(S)
=
\operatorname{tr}
\left(
P_{S,d}L_SP_{S,d}
\right),
\]
where \(P_{S,d}\) projects onto the selected low-energy workspace modes. Let
\[
\mathsf{Cut}(S)
=
\sum_{e:s(e)\in S,\ t(e)\notin S}w_e
+
\sum_{e:s(e)\notin S,\ t(e)\in S}w_e.
\]
Let \(k_\varepsilon(S)\) be the minimal sufficient coalition size inside the
candidate subsystem. A candidate is admissible only if
\[
\operatorname{SRW}_{\Theta_S}(W_S)=1,
\qquad
k_\varepsilon(S)\geq2,
\]
unless the claim explicitly concerns a centralized workspace.

\section{Boundary Objective, Recovery, and Persistence}

A general score is
\[
\boxed{
\begin{aligned}
\mathcal B_t(S)
={}&
 a_R\SRWTRet_\tau(S)
+a_O\SRWTOut_\tau(S)
+a_C\SRWTCOut_\tau(S)\\
&-a_D\SRWTDep_\tau(S)
-a_I\SRWTCIn_\tau(S)
-a_L\mathsf{Coh}(S)\\
&-a_K\mathsf{Cut}(S)
-a_Z\mathsf{Size}(S)
-a_T\SRWTInstability_t(S).
\end{aligned}
}
\]
The coefficients and all normalizations are part of the boundary problem. No
single choice is universal. The selected boundary is
\[
\widehat S_t
\in
\SRWTargmax_{S\in\mathfrak S_{\mathrm{pass}}}
\widehat{\mathcal B}_t(S),
\]
where \(\mathfrak S_{\mathrm{pass}}\) contains only candidates that satisfy the
required workspace and anti-centralization gates.

\subtopic{51.1 Boundary equivalence and identifiability}

If a symmetry \(g\in\mathcal G_B\) leaves the full observational and
interventional law invariant, \(S\) and \(gS\) cannot be distinguished by the
declared experiment. The identifiable boundary is then the orbit
\[
[S]_{\mathcal G_B}
=
\{gS:g\in\mathcal G_B\}.
\]
A unique subset should not be claimed when only an equivalence class is
identifiable.

\subtopic{51.2 Stable boundary recovery \(T_{17}\) }

\begin{SRWTtheorem}[Held-out boundary recovery]\label{SRWT:thm:heldout-boundary-recovery}
Let \(\mathfrak S_{\mathrm{pass}}\) be finite and let
\[
S_t^*
=
\SRWTargmax_{S\in\mathfrak S_{\mathrm{pass}}}
\mathcal B_t(S)
\]
be unique. Define the population margin
\[
\Delta_t^B
=
\mathcal B_t(S_t^*)
-
\max_{S\neq S_t^*}\mathcal B_t(S)>0.
\]
If
\[
\sup_{S\in\mathfrak S_{\mathrm{pass}}}
\left|
\widehat{\mathcal B}_t(S)-\mathcal B_t(S)
\right|
\leq\varepsilon_t,
\]
then every empirical maximizer \(\widehat S_t\) satisfies
\[
\boxed{
\mathcal B_t(S_t^*)-
\mathcal B_t(\widehat S_t)
\leq2\varepsilon_t.
}
\]
If \(2\varepsilon_t<\Delta_t^B\), then
\[
\boxed{\widehat S_t=S_t^*.}
\]
\end{SRWTtheorem}

\begin{proof}
Apply Theorem~\ref{SRWT:thm:uniform-structure-selection} to
\(J(S)=-\mathcal B_t(S)\).
\end{proof}

The theorem converts uncertainty in mutual-information estimation, causal
interventions, structural scores, and stability penalties into an explicit
boundary-identification condition. It also exposes when the data are
insufficient: a small margin makes the boundary intrinsically fragile.

\subtopic{51.3 Temporal boundary stability}

\begin{SRWTtheorem}[Boundary persistence under score drift]
\label{SRWT:thm:boundary-drift-stability}
Assume \(S_t^*\) is the unique maximizer of \(\mathcal B_t\) with margin
\(\Delta_t^B\), and suppose
\[
\sup_{S\in\mathfrak S_{\mathrm{pass}}}
|\mathcal B_{t+1}(S)-\mathcal B_t(S)|
\leq\delta_t.
\]
If
\[
2\delta_t<\Delta_t^B,
\]
then
\[
\boxed{S_{t+1}^*=S_t^*.}
\]
More generally,
\[
\mathcal B_{t+1}(S_{t+1}^*)
-
\mathcal B_{t+1}(S_t^*)
\leq2\delta_t.
\]
\end{SRWTtheorem}

\begin{proof}
For every competitor \(S\neq S_t^*\),
\[
\begin{aligned}
\mathcal B_{t+1}(S_t^*)-
\mathcal B_{t+1}(S)
&\geq
\mathcal B_t(S_t^*)-
\mathcal B_t(S)-2\delta_t\\
&\geq
\Delta_t^B-2\delta_t>0.
\end{aligned}
\]
Hence the maximizer is unchanged. The approximate bound follows by comparing
both time-indexed scores at their respective maximizers.
\end{proof}

If the score depends on a parameter \(\vartheta_t\) and satisfies
\[
\sup_S
|\mathcal B_{\vartheta}(S)-\mathcal B_{\vartheta'}(S)|
\leq
L_Bd(\vartheta,\vartheta'),
\]
then the sufficient stability condition becomes
\[
2L_Bd(\vartheta_t,\vartheta_{t+1})
<
\Delta_t^B.
\]

\subtopic{51.4 Analyst-discovered, system-instantiated, and self-maintained boundaries}

Three claims must remain separate.

\begin{description}[leftmargin=2.1em,style=nextline]
\item[Analyst-discovered boundary.]
An external procedure fits \(\widehat S\) after observing the system. The
system need not represent or use that boundary.

\item[System-instantiated boundary.]
The system contains a channel, mask, or routing variable causally equivalent to
\(\widehat S\), and interventions on that variable alter downstream operation.

\item[Self-maintained boundary.]
The system detects degradation of its boundary certificate and takes actions
that restore membership, transfer, or reconstruction properties.
\end{description}

Only the second and third claims concern a boundary realized within the
system. A post hoc partition that predicts behavior well is not by itself a
self-discovered boundary.

\subtopic{51.5 Boundary formation certificate}

A boundary \(\widehat S\) is certified as endogenously formed only if:
\begin{enumerate}
\item the membership mask was not directly encoded in \(\Sigma_B\) or
 \(\Sigma_Y\);
\item \(e_B^{\mathrm{MDL}}\) and \(e_B^{\mathrm{causal}}\) exceed declared
 thresholds;
\item held-out recovery satisfies Theorem~\ref{SRWT:thm:heldout-boundary-recovery};
\item the boundary remains stable under admissible environment drift or its
 changes are predicted by score-margin closure;
\item matched history-shuffle and label-leakage controls fail to reproduce
 the result, as do full-system, empty-set, singleton, and centralized
 boundary nulls;
\item the induced subsystem satisfies its workspace certificate; and
\item a system-instantiation intervention demonstrates that the boundary is
 used by the system rather than only by the analyst.
\end{enumerate}

The results below establish the statistical and deterministic stability components. They
does not prove that an arbitrary adaptive system will spontaneously learn such
a boundary.

% This file is part of the modular source for Stratified Relational Workspace Theory.
% !TeX root = ../Stratified_Relational_Workspace_Theory.tex

\chapter{Agency, counterfactual control,\\and self-maintenance}\label{SRWT:book:chapter-82}


\section{Controlled Workspace Agents and Goal States}

The preceding parts define when a distributed structure is coherent,
operationally sufficient, causally efficacious, relationally irreducible,
persistent, reconstructible, and endogenously formed. None of those
properties by itself establishes agency. Agency adds a controlled relation
among a represented objective, a set of feasible actions, counterfactual
consequences, and a policy that uses those quantities to alter the future.
Self-maintenance adds a further loop in which damage to the workspace becomes
an input to actions that preserve or restore the workspace certificate.

The purpose of this part is not to infer intentions from fluent behavior. It
is to define an intervention-based certificate whose components can fail
separately.

\subtopic{52.1 Controlled relational realization}

A controlled relational realization extends Definition~\ref{SRWT:def:realization}
by the tuple
\[
\boxed{
\mathfrak A_{\SRWTWorkspace}
=
\left(
\mathfrak R,
\mathcal G,
\mathcal A_t,
\mathcal Y^+,
\{U_g\}_{g\in\mathcal G},
\Pi,
\mathcal I_A,
\mathcal I_G,
\mathcal I_D,
\mathcal I_R
\right).
}
\]
Here:
\begin{itemize}
\item \(\mathfrak R\) is the underlying relational realization;
\item \(\mathcal G\) is a measurable goal-code space;
\item \(\mathcal A_t\) is the feasible action set at time \(t\);
\item \(\mathcal Y^+\) is the relevant future trajectory space;
\item \(U_g:\mathcal Y^+\to\mathbb R\) evaluates a trajectory under goal
 code \(g\);
\item \(\Pi\) is an admissible class of workspace-level policies;
\item \(\mathcal I_A\) and \(\mathcal I_G\) are feasible action and goal
 interventions; and
\item \(\mathcal I_D\) and \(\mathcal I_R\) are feasible damage and repair
 interventions.
\end{itemize}

A controlled state contains at least
\[
C_t^{\mathrm{ctl}}
=
\left(
W_t,G_t,M_t,B_t,E_t
\right),
\]
where \(G_t\) is the represented goal, \(M_t\) is a workspace-integrity
monitor, \(B_t\) is the current boundary or routing state, and \(E_t\) is the
declared external context. An action is generated by
\[
A_t
\sim
\pi_t
\left(
\cdot\mid W_t,G_t,M_t,B_t,E_t
\right).
\]
The future controlled state obeys a time-unrolled structural model
\[
S_{t+1}
=
F_t(S_t,A_t,E_{t+1},\xi_{t+1}),
\qquad
W_{t+1}=\omega_{t+1}(S_{t+1}).
\]
The intervention semantics are inherited from Chapter~\ref{SRWT:book:chapter-73}. A symbol such as
\(\operatorname{do}(A_t=a)\) is meaningful only when an implementation
intervention changes an actual action mechanism while the declared background
variables are matched.

\subtopic{52.2 Goal content and goal-code equivalence}

A persistent internal variable is not automatically a goal. A goal code must
index a stable counterfactual preference over relevant future trajectories.
For a fixed state \(w\), define the goal-conditioned action value
\[
Q_t^g(a\mid w)
=
\mathbb E
\left[
U_g(Y_t^+)
\mid
\operatorname{do}(A_t=a),
W_t=w,
G_t=g
\right].
\]
When action costs are relevant, they are incorporated into \(U_g\) or declared
as a constraint on \(\mathcal A_t\).

Two goal codes are operationally equivalent,
\[
g\sim_{\mathrm{goal}}g',
\]
when they induce the same action ordering in every declared state and
intervention background:
\[
Q_t^g(a\mid w)
\geq
Q_t^g(a'\mid w)
\quad\Longleftrightarrow\quad
Q_t^{g'}(a\mid w)
\geq
Q_t^{g'}(a'\mid w).
\]
The abstract goal is the equivalence class
\[
[g]_{\mathrm{goal}}.
\]
This quotient prevents a mere invertible recoding of a goal channel from being
reported as a change of objective. For quantitative regret bounds, a
representative normalization of \(U_g\) must nevertheless be fixed.

\begin{SRWTremark}[Goals are claim-relative]
The same internal code can count as a goal for one declared trajectory family
and fail to do so for another. A code that predicts a short-horizon action may
not encode a persistent long-horizon objective. Every agency claim must state
its horizon, outcome space, feasible actions, and utility normalization.
\end{SRWTremark}

\subtopic{52.3 Local, distributed, and workspace-level goals}

Let \(\mathcal H_t^S\) be the history available to coalition \(S\subseteq V\).
For a loss \(\ell_G\), define goal-reconstruction risk
\[
\SRWTRisk_G(S)
=
\inf_f
\mathbb E
\left[
\ell_G
\left(
G_t,
f(\mathcal H_t^S)
\right)
\right].
\]
For tolerance \(\varepsilon\), define
\[
k_{G,\varepsilon}
=
\min
\left\{
|S|:
\SRWTRisk_G(S)
\leq
\SRWTRisk_G(V)+\varepsilon
\right\}.
\]
A goal is locally stored when \(k_{G,\varepsilon}=1\). It is distributed
relative to the declared decomposition when \(k_{G,\varepsilon}\geq2\).

This distinction is independent of whether the rest of the workspace is
relational. A relational workspace can obey a goal copied in full into one
controller. Such a system may possess a workspace but not a
\emph{workspace-level goal} in the strong distributed sense.

For the complete control state, define
\[
\SRWTRisk_{\mathrm{ctl}}(S)
=
\inf_f
\mathbb E
\left[
\ell_{\mathrm{ctl}}
\left(
C_t^{\mathrm{ctl}},
f(\mathcal H_t^S)
\right)
\right],
\]
and let \(k_{\mathrm{ctl},\varepsilon}\) be its minimal sufficient coalition
size. A strong workspace-level agency claim requires
\[
\boxed{k_{\mathrm{ctl},\varepsilon}\geq2.}
\]
This is a location criterion. It does not yet establish goal persistence,
counterfactual evaluation, or self-maintenance.

\subtopic{52.4 Action mediation and outcome control}

Goal-dependent action mediation is measured by feasible interventions on the
goal channel while matching the current workspace and external context. For
a divergence \(D_A\), define
\[
\eta_{G\to A}
=
\sup_{g,g',w,e}
D_A
\left(
\mathbb P(A_t\in\cdot\mid\iota_g,w,e),
\mathbb P(A_t\in\cdot\mid\iota_{g'},w,e)
\right).
\]
Workspace-dependent action mediation is
\[
\eta_{W\to A}
=
\sup_{w,w',g,e}
D_A
\left(
\mathbb P(A_t\in\cdot\mid\iota_w,g,e),
\mathbb P(A_t\in\cdot\mid\iota_{w'},g,e)
\right).
\]
An action must also alter the declared future:
\[
\eta_{A\to Y}
=
\sup_{a,a',b}
D_Y
\left(
\mathbb P(Y_t^+\in\cdot\mid\iota_a,b),
\mathbb P(Y_t^+\in\cdot\mid\iota_{a'},b)
\right).
\]
A positive \(\eta_{G\to A}\) without \(\eta_{A\to Y}\) describes goal-sensitive
output selection with no demonstrated control of the relevant environment. A
positive \(\eta_{A\to Y}\) without a goal or counterfactual-value channel may
describe a reactive controller.

\subtopic{52.5 Why observational goal persistence is insufficient}

The observational condition
\[
I(G_{t+\tau};G_t\mid E_{t+1:t+\tau})>0
\]
can arise from several non-agentic mechanisms:
\begin{enumerate}
\item a persistent external command writes the same code at every step;
\item a clock, episode label, or task identifier drives both variables;
\item a frozen memory cell repeats a symbol that is not used by the policy;
\item selection on successful trajectories induces apparent temporal
 dependence; or
\item a downstream action is copied back into the nominal goal channel.
\end{enumerate}
A goal-persistence claim therefore requires a cue-controlled intervention that
changes \(G_t\), removes or matches subsequent command inputs, and checks both
future goal state and action consequences.

\subtopic{52.6 Cue-controlled goal persistence}

Let \(C_t\) denote an external command or task cue. Fix a neutral or matched
future command trajectory \(c_{t+1:t+\tau}^{0}\). For goal interventions
\(\iota_g\) and \(\iota_{g'}\), define the cue-controlled retention effect
\[
\eta_G(\tau;g,g')
=
D_G
\left(
\mathbb P
\left(
G_{t+\tau}\in\cdot
\mid
\iota_g,
\operatorname{do}(C_{t+1:t+\tau}=c^0)
\right),
\mathbb P
\left(
G_{t+\tau}\in\cdot
\mid
\iota_{g'},
\operatorname{do}(C_{t+1:t+\tau}=c^0)
\right)
\right).
\]
The corresponding action-retention effect is
\[
\eta_{G\to A}(\tau;g,g')
=
D_A
\left(
\mathbb P(A_{t+\tau}\in\cdot\mid\iota_g,c^0),
\mathbb P(A_{t+\tau}\in\cdot\mid\iota_{g'},c^0)
\right).
\]
The first quantity tests whether a goal difference survives. The second tests
whether the surviving difference remains behaviorally operative.

Cue dependence must be reported separately. For matched initial goal \(g\),
define
\[
\delta_G^{\mathrm{cue}}(\tau)
=
\sup_{c,c'}
D_G
\left(
\mathbb P(G_{t+\tau}\in\cdot\mid\iota_g,\operatorname{do}(C=c)),
\mathbb P(G_{t+\tau}\in\cdot\mid\iota_g,\operatorname{do}(C=c'))
\right).
\]
High cue dependence is not a failure when the claimed agent is intentionally
instruction following. It does mean that persistence cannot be attributed to
an autonomous internal objective without an additional cue-removal result.

\subtopic{52.7 \(T_{19}\): Cue-controlled goal-retention theorem }

\begin{SRWTtheorem}[Cue-controlled goal retention]
\label{SRWT:thm:goal-retention}
Let \(G_s\in\mathbb R^{d_G}\) obey
\[
G_{s+1}
=
R_sG_s+B_sC_{s+1}+\zeta_{s+1}
\qquad
(s=t,\ldots,t+\tau-1).
\]
Consider two intervention runs that set
\[
G_t=g,
\qquad
G_t'=g',
\]
and then use the same command trajectory and the same disturbance realization.
Let
\[
R_{t+\tau:t}
=
R_{t+\tau-1}\cdots R_t.
\]
Then
\[
\boxed{
G_{t+\tau}-G_{t+\tau}'
=
R_{t+\tau:t}(g-g').
}
\]
If
\[
m_s
\leq
\sigma_{\min}(R_s),
\qquad
\|R_s\|\leq M_s,
\]
then
\[
\boxed{
\left(
\prod_{s=t}^{t+\tau-1}m_s
\right)
\|g-g'\|
\leq
\|G_{t+\tau}-G_{t+\tau}'\|
\leq
\left(
\prod_{s=t}^{t+\tau-1}M_s
\right)
\|g-g'\|.
}
\]
\end{SRWTtheorem}

\begin{proof}
Subtract the two update equations. The matched command and disturbance terms
cancel, giving
\[
G_{s+1}-G_{s+1}'
=
R_s(G_s-G_s').
\]
Iteration gives the exact product formula. The upper bound follows from
submultiplicativity of the operator norm. The lower bound follows from
\[
\sigma_{\min}(AB)
\geq
\sigma_{\min}(A)\sigma_{\min}(B)
\]
for square linear maps.
\end{proof}

The theorem isolates internal retention from repeated external rewriting. It
also makes clear that persistence is graded. Uniformly contractive maps can
retain a goal for a finite horizon while erasing it asymptotically. A neutral
mode with product singular value near one supports long retention. An
expansive mode can preserve distinctions but may be dynamically unstable.

\begin{SRWTcorollary}[Behavioral readout of a retained goal]
\label{SRWT:cor:goal-action-readout}
Assume, in addition, that under the matched intervention background the action
readout has conditional mean
\[
\mathbb E[A_{t+\tau}\mid G_{t+\tau},W_{t+\tau}]
=
K_{t+\tau}G_{t+\tau}+b(W_{t+\tau}),
\]
and that the two runs are matched on \(W_{t+\tau}\). Then
\[
\mathbb E[A_{t+\tau}\mid\iota_g]
-
\mathbb E[A_{t+\tau}\mid\iota_{g'}]
=
K_{t+\tau}R_{t+\tau:t}(g-g').
\]
If \(\sigma_{\min}(K_{t+\tau})\geq k_{t+\tau}>0\), the norm of this action-mean
difference is at least
\[
k_{t+\tau}
\left(
\prod_{s=t}^{t+\tau-1}m_s
\right)
\|g-g'\|.
\]
\end{SRWTcorollary}

The corollary is a sufficient condition, not a universal model of action
selection. For discrete actions, the corresponding test should compare the
full interventional action distributions.

\subtopic{52.8 Goal-persistence certificate}

A minimal goal-persistence report contains:
\begin{enumerate}
\item the goal code and its declared trajectory-level semantics;
\item the feasible goal intervention and its manipulation check;
\item the matched or neutralized future cue trajectory;
\item \(\eta_G(\tau)\) across a declared horizon grid;
\item \(\eta_{G\to A}(\tau)\) over the same grid;
\item cue dependence \(\delta_G^{\mathrm{cue}}(\tau)\);
\item goal-channel ablation and permutation controls;
\item coalition reconstruction risks \(\SRWTRisk_G(S)\); and
\item uncertainty intervals over environments, episodes, and random seeds.
\end{enumerate}
A repeated label with no intervention effect on future action fails this
certificate even when its observational autocorrelation is one.

\section{Counterfactual Policy, Repair, and Agency Certification}

\subtopic{53.1 Feasible action interventions and value functions}

Let \(\mathcal A_t(w,e)\) denote the actions that can be implemented from
workspace state \(w\) in environment \(e\). The interventional action value is
\[
Q_t^g(a\mid w,e)
=
\mathbb E^{\iota_a}
\left[
U_g(Y_t^+)
\mid
W_t=w,
E_t=e,
G_t=g
\right].
\]
The optimal action set is
\[
\mathcal A_t^*(w,g,e)
=
\SRWTargmax_{a\in\mathcal A_t(w,e)}Q_t^g(a\mid w,e).
\]
A counterfactual evaluator is an internal state
\[
\widehat Q_t
=
\left(
\widehat Q_t(a)
\right)_{a\in\mathcal A_t(w,e)}
\]
that estimates these intervention values. It is not enough for
\(\widehat Q_t\) to predict the action selected by the system. It must be
validated against externally estimated intervention outcomes.

\subtopic{53.2 Counterfactual-use criterion}

Three properties must be separated:
\begin{description}
\item[Counterfactual accuracy.]
\[
\delta_Q
=
\sup_{a\in\mathcal A_t}
|\widehat Q_t(a)-Q_t^g(a\mid w,e)|.
\]

\item[Policy use.]
A feasible intervention that changes \(\widehat Q_t(a)\) while preserving the
action affordances changes the action distribution in the direction predicted
by the decision rule.

\item[Outcome validity.]
The resulting action change alters \(U_g(Y_t^+)\) as predicted by the external
interventional value estimate.
\end{description}

A system can satisfy any two while failing the third. A learned critic can be
accurate but ignored. A policy can respond to a manipulated score that does
not predict outcomes. A successful action can arise from habit or replay
without any counterfactual evaluator.

\subtopic{53.3 \(T_{20}\): Approximate counterfactual policy theorem }

\begin{SRWTtheorem}[Approximate counterfactual policy selection]
\label{SRWT:thm:counterfactual-policy}
Let \(\mathcal A\) be a finite feasible action set. Let \(Q:\mathcal A\to
\mathbb R\) be the true interventional value and let \(\widehat Q\) satisfy
\[
\sup_{a\in\mathcal A}|\widehat Q(a)-Q(a)|\leq\varepsilon_Q.
\]
Suppose a policy \(\widehat\pi\) is \(\delta_{\mathrm{opt}}\)-optimal for
\(\widehat Q\):
\[
\max_{\pi}
\mathbb E_{a\sim\pi}[\widehat Q(a)]
-
\mathbb E_{a\sim\widehat\pi}[\widehat Q(a)]
\leq
\delta_{\mathrm{opt}}.
\]
Then, for every true-value-optimal policy \(\pi^*\),
\[
\boxed{
\mathbb E_{a\sim\pi^*}[Q(a)]
-
\mathbb E_{a\sim\widehat\pi}[Q(a)]
\leq
2\varepsilon_Q+\delta_{\mathrm{opt}}.
}
\]
For a deterministic greedy action \(\widehat a\in\SRWTargmax_a\widehat Q(a)\),
\[
Q(a^*)-Q(\widehat a)
\leq2\varepsilon_Q.
\]
If the true optimal action \(a^*\) is unique and its value gap
\[
\Delta_Q
=
Q(a^*)-
\max_{a\neq a^*}Q(a)
\]
satisfies \(\Delta_Q>2\varepsilon_Q\), then every greedy maximizer of
\(\widehat Q\) equals \(a^*\).
\end{SRWTtheorem}

\begin{proof}
For every policy \(\pi\), uniform approximation gives
\[
\left|
\mathbb E_{a\sim\pi}[\widehat Q(a)]
-
\mathbb E_{a\sim\pi}[Q(a)]
\right|
\leq\varepsilon_Q.
\]
Therefore
\[
\begin{aligned}
\mathbb E_{\pi^*}[Q]
-
\mathbb E_{\widehat\pi}[Q]
&\leq
\mathbb E_{\pi^*}[\widehat Q]
-
\mathbb E_{\widehat\pi}[\widehat Q]
+2\varepsilon_Q\\
&\leq
\delta_{\mathrm{opt}}+2\varepsilon_Q.
\end{aligned}
\]
The deterministic bound is the special case of point-mass policies. If
\(\widehat a\neq a^*\), then
\[
\widehat Q(a^*)
\geq
Q(a^*)-\varepsilon_Q
>
Q(\widehat a)+\varepsilon_Q
\geq
\widehat Q(\widehat a),
\]
contradicting greedy optimality.
\end{proof}

This theorem converts counterfactual-model accuracy into a decision guarantee.
It does not establish that the model is internally represented or causally
used. Those are separate intervention gates.

\subtopic{53.4 Counterfactual mediation and the replay null}

Let \(V_t\) denote the internal value channel. A counterfactual-mediation test
compares:
\[
G_t
\longrightarrow
V_t
\longrightarrow
A_t
\longrightarrow
Y_t^+.
\]
The relevant quantities include:
\[
\eta_{G\to V},
\qquad
\eta_{V\to A},
\qquad
\eta_{A\to Y},
\qquad
\delta_{G\to A\mid V},
\]
where the final term measures the residual goal-to-action effect after the
value channel is blocked or randomized.

The action-replay null preserves the empirical action sequence or action
marginal while destroying online dependence on \(V_t\). If replayed actions
perform equally well under held-out perturbations, the evidence supports a
cached or open-loop policy rather than online counterfactual evaluation. A
second null permutes value labels among feasible actions while preserving the
multiset of scores. A genuine value-mediated policy should follow the
permuted ranking, while an ignored critic should not.

\subtopic{53.5 Workspace integrity vector and maintenance margin}

Choose a declared set \(\mathcal J_{\mathrm{maint}}\) of workspace gates that
must be protected. Write each gate as a signed margin
\[
\phi_{j,t}
\geq0
\quad\Longleftrightarrow\quad
\text{gate }j\text{ passes at time }t.
\]
Examples include
\[
\begin{aligned}
\phi_{\mathrm{gap},t}
&=g_{\varepsilon,t}-\underline g,\\
\phi_{\mathrm{coh},t}
&=\overline\delta_{\mathrm{coh}}-\delta_{\mathrm{coh},t},\\
\phi_{\mathrm{rec},t}
&=\alpha_t^{(k)}-\underline\alpha,\\
\phi_{\mathrm{keep},t}
&=\mu_t^{\mathrm{keep}}-\underline\mu,\\
\phi_{\mathrm{causal},t}
&=\eta_{\mathrm{causal},t}-\underline\eta,\\
\phi_{\mathrm{boundary},t}
&=\Delta_t^B-\underline\Delta_B.
\end{aligned}
\]
The aggregate integrity margin is
\[
\boxed{
v_t
=
\min_{j\in\mathcal J_{\mathrm{maint}}}\phi_{j,t}.
}
\]
The workspace passes the protected gates when \(v_t\geq0\). To maintain a
reserve \(\rho>0\), define the maintenance deficit
\[
\boxed{
e_t
=(\rho-v_t)_+.
}
\]
Define the zero-reserve and bounded-deficit viability sets
\[
\boxed{
\mathcal V_\Theta
=
\left\{
S:\phi_j(S)\geq0\ \text{for every }j\in\mathcal J_{\mathrm{maint}}
\right\}
}
\]
\[
\boxed{
\mathcal V_{\Theta,e_{\max}}
=
\left\{
S:e(S)\leq e_{\max}
\right\}
}
\]
The first set requires every protected gate to pass. The second permits a
declared finite reserve deficit.

This minimum-gate construction is intentionally noncompensatory. A large
spectral gap cannot compensate for failed reconstruction, and high task reward
cannot compensate for loss of the declared workspace boundary.

\subtopic{53.6 Damage and repair interventions}

A damage intervention \(d\in\mathcal I_D\) changes a mechanism that supports
one or more protected gates. Examples include:
\begin{itemize}
\item deleting or corrupting a transfer edge;
\item removing a participant or artifact;
\item perturbing the internal boundary mask;
\item scrambling a role-binding or routing channel;
\item erasing a handoff representation; or
\item degrading a reconstruction map.
\end{itemize}

A repair intervention \(r\in\mathcal I_R\) can add, remove, retrain, rewrite,
or reroute an admissible component. Let \(r_0\) denote matched no-repair. The
post-repair integrity value at horizon \(\tau\) is
\[
V_t^{\mathrm{rep}}(r;d)
=
\mathbb E
\left[
v_{t+\tau}
\mid
\operatorname{do}(D_t=d),
\operatorname{do}(R_t=r),
\mathbf H_t
\right].
\]
The repair gain is
\[
\Delta_t^{\mathrm{rep}}(r;d)
=
V_t^{\mathrm{rep}}(r;d)
-
V_t^{\mathrm{rep}}(r_0;d).
\]
Repair costs are handled either by restricting the feasible repair family or
by reporting a separate cost-effect curve. They should not be hidden inside
an undocumented scalar reward.

\subtopic{53.7 Damage detection, repair mediation, and specificity}

Let \(M_t\) be an internal integrity monitor and \(R_t\) the selected repair
action. A self-maintenance claim requires the causal chain
\[
D_t
\longrightarrow
M_t
\longrightarrow
R_t
\longrightarrow
v_{t+\tau}.
\]
The required tests are:
\begin{description}
\item[Detection.]
Matched damage interventions change \(M_t\) in a gate-specific direction.

\item[Repair mediation.]
Intervening on \(M_t\) changes the repair distribution, and blocking \(M_t\)
reduces the damage-to-repair effect.

\item[Restoration.]
The selected repair improves \(v_{t+\tau}\) relative to matched no-repair.

\item[Specificity.]
A cost-matched generic action or an off-workspace repair does not produce the
same restoration.

\item[Identity preservation.]
The restored system remains within the declared identity tolerance relative to
the pre-damage workspace.
\end{description}

Define the detection effect \(\eta_{D\to M}\), mediation effect
\(\eta_{M\to R}\), restoration effect \(\eta_{R\to v}\), and repair
specificity
\[
\Delta_{\mathrm{repair}}^{\mathrm{spec}}
=
\eta_{R\to v}
-
\eta_{R_\perp\to v},
\]
where \(R_\perp\) is a matched action outside the proposed repair channel.

\subtopic{53.8 \(T_{21}\): Repair-selection margin theorem }

\begin{SRWTtheorem}[Repair selection under value-model error]
\label{SRWT:thm:repair-selection}
Fix a damage intervention \(d\) and a finite feasible repair set
\(\mathcal R\). Let
\[
V(r)=V_t^{\mathrm{rep}}(r;d)
\]
be the true expected post-repair integrity margin and let \(\widehat V\) obey
\[
\sup_{r\in\mathcal R}|\widehat V(r)-V(r)|
\leq
\varepsilon_R.
\]
Let \(\widehat r\in\SRWTargmax_r\widehat V(r)\) and
\(r^*\in\SRWTargmax_r V(r)\). Then
\[
\boxed{
V(r^*)-V(\widehat r)
\leq
2\varepsilon_R.
}
\]
If
\[
V(r^*)
\geq
\rho+2\varepsilon_R,
\]
then the selected repair restores the expected reserve:
\[
\boxed{V(\widehat r)\geq\rho.}
\]
If \(r^*\) is unique and its repair-value gap exceeds
\(2\varepsilon_R\), then \(\widehat r=r^*\).
\end{SRWTtheorem}

\begin{proof}
Greedy optimality and uniform approximation give
\[
\begin{aligned}
V(r^*)
&\leq
\widehat V(r^*)+\varepsilon_R\\
&\leq
\widehat V(\widehat r)+\varepsilon_R\\
&\leq
V(\widehat r)+2\varepsilon_R.
\end{aligned}
\]
The reserve and exact-selection statements follow immediately.
\end{proof}

The theorem is conditional on a valid repair-value model. It cannot convert a
post hoc repair narrative into evidence that the system detected damage or
used the model. Detection and value-channel interventions remain necessary.

\subtopic{53.9 Repair, replacement, and successor workspaces}

Restoration of performance is not automatically restoration of identity. Let
\(W_{t^-}\) be the last certified pre-damage workspace and let
\(W_{t+\tau}^{r}\) be the post-repair workspace. Define the repair identity
defect
\[
\delta_{\mathrm{rep-id}}(r)
=
 d_{\mathrm{id}}
\left(
W_{t^-},W_{t+\tau}^{r}
\right).
\]
A repair action is identity preserving at tolerance
\(\varepsilon_{\mathrm{id}}\) when
\[
\Delta_t^{\mathrm{rep}}(r;d)>0,
\qquad
\delta_{\mathrm{rep-id}}(r)
\leq
\varepsilon_{\mathrm{id}}.
\]
When integrity is restored but the identity defect exceeds tolerance, the
result is better described as a successor or replacement workspace. This
terminological distinction prevents any high-performing reconstruction from
being counted as persistence of the same relational organization.

\subtopic{53.10 \(T_{22}\): Closed-loop self-maintenance theorem }

\begin{SRWTtheorem}[Causal repair stabilizes maintenance deficit]
\label{SRWT:thm:self-maintenance}
Let \(e_t\geq0\) be the maintenance deficit. Suppose the matched no-repair
trajectory satisfies
\[
e_{t+1}^{(0)}
\leq
\alpha e_t+d_t,
\]
where \(d_t\geq0\) is exogenous damage. Suppose the implemented repair policy
has the causal one-step effect
\[
e_{t+1}^{(0)}-e_{t+1}
\geq
\kappa e_t-\epsilon_t.
\]
Assume
\[
0\leq\beta:=\alpha-\kappa<1,
\qquad
0\leq d_t\leq\overline d,
\qquad
0\leq\epsilon_t\leq\overline\epsilon.
\]
Then
\[
\boxed{
e_t
\leq
\beta^te_0
+
\frac{(\overline d+\overline\epsilon)(1-\beta^t)}{1-\beta},
}
\]
and therefore
\[
\boxed{
\limsup_{t\to\infty}e_t
\leq
\frac{\overline d+\overline\epsilon}{1-\beta}.
}
\]
If the admissible deficit is \(e_{\max}\) and
\[
e_0\leq e_{\max},
\qquad
\overline d+\overline\epsilon
\leq
(1-\beta)e_{\max},
\]
then the interval \([0,e_{\max}]\) is forward invariant.
\end{SRWTtheorem}

\begin{proof}
Combine the no-repair bound with the causal repair effect:
\[
\begin{aligned}
e_{t+1}
&\leq
 e_{t+1}^{(0)}-\kappa e_t+\epsilon_t\\
&\leq
(\alpha-\kappa)e_t+d_t+\epsilon_t\\
&\leq
\beta e_t+\overline d+\overline\epsilon.
\end{aligned}
\]
Iterating this scalar recurrence gives the stated geometric bound. If
\(e_t\leq e_{\max}\), the forward-invariance condition gives
\[
e_{t+1}
\leq
\beta e_{\max}
+(1-\beta)e_{\max}
=e_{\max}.
\]
\end{proof}

Theorem~\ref{SRWT:thm:self-maintenance} differs from the persistence theorem in
Chapter~\ref{SRWT:book:chapter-76}. There, contraction was assumed for the workspace error dynamics. Here
the improved contraction factor is attributed to a measured repair action
through a matched no-repair counterfactual. Without that causal comparison,
the same bounded trajectory could be passive robustness rather than
self-maintenance.

\subtopic{53.11 A scalar boundary-repair example}

Let \(b_t\in[0,1]\) be a boundary-integrity score and let the reserve be
\(\rho\). Define
\[
e_t=(\rho-b_t)_+.
\]
Suppose uncorrected damage obeys
\[
e_{t+1}^{(0)}
\leq
1.08e_t+d_t,
\]
so the open-loop boundary is locally unstable. A repair policy with measured
causal effect
\[
e_{t+1}^{(0)}-e_{t+1}
\geq
0.25e_t-\epsilon_t
\]
gives
\[
\beta=1.08-0.25=0.83.
\]
If
\[
\overline d=0.01,
\qquad
\overline\epsilon=0.005,
\]
then
\[
\limsup e_t
\leq
\frac{0.015}{1-0.83}
\approx0.0883.
\]
Thus an open-loop unstable boundary can be maintained within a finite deficit
band by a sufficiently strong repair channel. The example establishes only a
control calculation. A complete self-maintenance claim must also show that an
internal monitor detects damage, mediates the repair choice, and preserves
workspace identity.

\subtopic{53.12 Complete self-maintenance certificate}

A self-maintaining workspace must pass all of the following gates:
\begin{enumerate}
\item \emph{Protected-gate declaration.} The maintained integrity vector and
 reserve are fixed before evaluation.
\item \emph{Damage validity.} Each damage intervention measurably degrades at
 least one protected gate.
\item \emph{Internal detection.} The monitor changes under damage and is not
 merely an external evaluator.
\item \emph{Repair mediation.} Monitor interventions change repair selection.
\item \emph{Repair efficacy.} Repair improves the held-out integrity margin
 relative to matched no-repair.
\item \emph{Specificity.} Cost-matched generic actions and off-workspace
 interventions do not reproduce the gain.
\item \emph{Repeated-disturbance stability.} The deficit remains within the
 declared bound over a held-out disturbance sequence.
\item \emph{Identity control.} The restored state remains within the declared
 identity tolerance or is reported as a successor workspace.
\item \emph{Anti-centralization.} No single node contains a sufficient copy
 of the monitor, repair model, and protected state when the claim is
 workspace-level maintenance.
\end{enumerate}

\subtopic{53.13 Certified workspace agent}

Let
\[
\Theta_A
=
\left(
\tau_G,
\tau_{G\to A},
\varepsilon_Q,
\tau_{V\to A},
\tau_{A\to Y},
\tau_{D\to M},
\tau_{M\to R},
\tau_{R\to v},
\tau_{\mathrm{spec}},
\varepsilon_{\mathrm{ctl}},
 e_{\max}
\right).
\]
A \(\Theta_A\)-certified workspace agent is a controlled realization for
which:
\begin{enumerate}
\item \(\SRWTSRW(W_t;\varsigma)=1\);
\item cue-controlled goal retention exceeds \(\tau_G\);
\item goal-to-action mediation exceeds \(\tau_{G\to A}\);
\item counterfactual-value error is at most \(\varepsilon_Q\);
\item value-to-action use exceeds \(\tau_{V\to A}\);
\item action-to-outcome control exceeds \(\tau_{A\to Y}\);
\item damage detection, monitor-to-repair mediation, and restoration exceed
 their thresholds;
\item repair specificity exceeds \(\tau_{\mathrm{spec}}\);
\item \(k_{\mathrm{ctl},\varepsilon_{\mathrm{ctl}}}\geq2\); and
\item repeated maintenance deficit remains at most \(e_{\max}\).
\end{enumerate}

As throughout the theory, this predicate is relative to the declared node
decomposition, goals, actions, environment, horizons, interventions, and null
families.

\subtopic{53.14 \(T_{23}\): Workspace-agency synthesis theorem }

\begin{SRWTtheorem}[Workspace-level agency under declared gates]
\label{SRWT:thm:agency-synthesis}
Assume a controlled realization satisfies:
\begin{enumerate}
\item the operational relationality assumptions of
 Theorem~\ref{SRWT:thm:operational-relationality};
\item cue-controlled goal retention and action readout as in
 Theorem~\ref{SRWT:thm:goal-retention} and
 Corollary~\ref{SRWT:cor:goal-action-readout};
\item a counterfactual evaluator with uniform error \(\varepsilon_Q\) and
 optimization error \(\delta_{\mathrm{opt}}\);
\item feasible interventions establishing positive value-to-action and
 action-to-outcome effects;
\item the causal self-maintenance assumptions of
 Theorem~\ref{SRWT:thm:self-maintenance}; and
\item \(k_{\mathrm{ctl},\varepsilon}\geq2\), together with a matched local
 policy-null panel that preserves every single-node predictive marginal
 but loses full-system control value.
\end{enumerate}
Then, relative to the declared specification:
\begin{enumerate}
\item the expected one-step decision regret is at most
 \(2\varepsilon_Q+\delta_{\mathrm{opt}}\);
\item the long-run maintenance deficit is bounded by
 \[
 \frac{\overline d+\overline\epsilon}{1-\beta};
 \]
\item no single-node statistic is sufficient for the complete control state;
\item matched local policy nulls cannot reproduce the intact control value;
\item goal-dependent action selection and workspace repair are causally
 mediated by implemented internal channels.
\end{enumerate}
Hence the realization satisfies the formal predicate
\[
\operatorname{Agent}_{\Theta_A}(\SRWTWorkspace)=1.
\]
\end{SRWTtheorem}

\begin{proof}
The regret bound follows from
Theorem~\ref{SRWT:thm:counterfactual-policy}. The maintenance bound follows from
Theorem~\ref{SRWT:thm:self-maintenance}. The coalition condition excludes
single-node sufficiency for the complete control state. The matched-null gap
excludes explanations based only on preserved local predictive marginals. The
implemented goal, value, action, monitor, and repair interventions establish
the required causal mediation. These are precisely the gates in the stated
agency predicate.
\end{proof}

The synthesis theorem is intentionally relative and conjunctive. Its force is
that decision quality, maintenance stability, causal mediation, and relational
location are certified together. It does not infer an unrestricted agent
outside the tested environment.

\subtopic{53.15 Agency, autonomy, and consciousness}

The following implications are not licensed:
\[
\operatorname{Agent}_{\Theta_A}(\SRWTWorkspace)
\centernot\Longrightarrow
\operatorname{PhenomenallyConscious}(\SRWTWorkspace),
\]
\[
\operatorname{SelfMaintaining}(\SRWTWorkspace)
\centernot\Longrightarrow
\operatorname{SelfAware}(\SRWTWorkspace),
\]
\[
\operatorname{GoalPersistent}(\SRWTWorkspace)
\centernot\Longrightarrow
\operatorname{MoralPatient}(\SRWTWorkspace).
\]
A thermostat-like controller can maintain a variable without representing
itself. A sophisticated planner can evaluate counterfactual actions without
phenomenal experience. Conversely, failure to pass this engineering
certificate would not prove absence of consciousness. Agency and
consciousness remain distinct axes of inquiry.

% This file is part of the modular source for Stratified Relational Workspace Theory.
% !TeX root = ../Stratified_Relational_Workspace_Theory.tex

\chapter{Decomposition, coarse-graining, and robust relationality}\label{SRWT:book:chapter-83}


Relationality is always assessed relative to a way of dividing a realization
into components. That dependence is not a minor technicality. If every
participant, artifact, and channel is merged into one component, the complete
workspace becomes local by definition. If a centralized state is recoded into
several complementary shares, the same centralized process can be made to look
irreducibly distributed. A defensible theory must therefore say which
partitions count as admissible, which quantities survive faithful aggregation,
and how sensitive a certificate is to alternative descriptions of the same
system.

\section{The Decomposition Problem}

Let \(\Omega\) be a finite set of primitive carriers at a declared audit
resolution. A primitive carrier may be a person, a model instance, a process,
a memory store, an instrument, or another independently tracked mechanism.
The word \emph{primitive} is relative to the audit. It does not assert a final
metaphysical atomization of the system.

\begin{SRWTdefinition}[System decomposition]\label{SRWT:def:system-decomposition}
A decomposition of a carrier presentation is a tuple
\[
\mathcal D
=
\bigl(
\Pi,\mathbf O,G_{\Pi},\mathcal I_{\Pi}
\bigr),
\]
where:
\begin{itemize}
\item \(\Pi=\{B_\alpha\}_{\alpha\in A}\) is a partition of \(\Omega\);
\item \(O_\alpha\) is the declared local-access map for block \(B_\alpha\),
 producing the local state or history available at that component;
\item \(G_{\Pi}\) records the declared communication, transfer, and causal
 interfaces among the blocks; and
\item \(\mathcal I_{\Pi}\) is the family of feasible interventions at that
 resolution.
\end{itemize}
The accessible information of component \(\alpha\) is the sigma-field
\[
\mathscr F_\alpha^{\mathcal D}
=
\sigma\!\left(O_\alpha(\mathbf H_t)\right).
\]
Every shared bus, common memory, supervisor, or artifact that can alter more
than one local state must appear either as its own carrier or as an explicit
channel in \(G_{\Pi}\).
\end{SRWTdefinition}

The same physical realization can support several decompositions. One may
separate a model from its executor, merge them into one computational unit, or
treat a persistent artifact as its own component. These alternatives are not
mere renamings when they change local access or feasible interventions. Every certificate should therefore display the full specification
\[
\SRWTSRW(W_t;\varsigma),
\qquad
\varsigma=(\mathcal D,Y^+,\ell,\tau,\mathcal I,\mathcal N,\Theta).
\]
This part varies the decomposition coordinate while holding the remaining
coordinates fixed. Chapter~\ref{SRWT:book:chapter-84} treats the full specification family.

\begin{SRWTdefinition}[Refinement and coarsening]\label{SRWT:def:decomposition-refinement}
Let \(\mathcal D_f\) and \(\mathcal D_c\) be decompositions of the same carrier
presentation. The fine decomposition refines the coarse decomposition,
written
\[
\mathcal D_f\preceq\mathcal D_c,
\]
when there is a surjection
\[
q:V_f\longrightarrow V_c
\]
such that each coarse component is the union of the fine components in one
fiber of \(q\), and its accessible information is generated only from those
fine components. The aggregation is \emph{information faithful} when
\[
\mathscr F_B^{\mathcal D_c}
=
\bigvee_{i\in q^{-1}(B)}
\mathscr F_i^{\mathcal D_f}
\qquad(B\in V_c).
\]
It is \emph{structurally faithful} when it also retains every internal
compatibility condition and every boundary transfer condition needed to
recover the original global-section problem.
\end{SRWTdefinition}

Refinement and coarsening are directional. A coarse component may contain
information that was spread across several finer nodes. Conversely, a fine
description may expose internal channels that disappear when a macro-node is
formed. The next results identify what is preserved when the aggregation is
faithful.

\section{Faithful Structural Coarse-Graining}

Let the vertex set of a core-transfer diagram be partitioned as
\[
V=V_1\sqcup\cdots\sqcup V_m.
\]
Write the transfer discrepancy as
\[
D_H
=
\begin{pmatrix}
D_{\mathrm{int}}\\
D_{\partial}
\end{pmatrix},
\qquad
D_{\mathrm{int}}
=
\bigoplus_{\alpha=1}^{m}D_\alpha,
\]
where \(D_\alpha\) contains the edge constraints internal to \(V_\alpha\),
and \(D_{\partial}\) contains the constraints crossing between blocks. Define
the internally coherent macro-state spaces
\[
W_\alpha=\ker D_\alpha,
\qquad
W_{\Pi}=\bigoplus_{\alpha=1}^{m}W_\alpha.
\]
Let
\[
J_\Pi:W_\Pi\hookrightarrow\mathcal C_V
\]
be the canonical isometric inclusion, and define the quotient discrepancy
\[
D_\Pi=D_{\partial}J_\Pi.
\]
The quotient does not replace a group of nodes by an arbitrary summary. Its
macro-state is the complete internal section space of that group, and its
boundary operator retains all cross-block compatibility equations.

\begin{SRWTtheorem}[Faithful coarse-graining]\label{SRWT:thm:faithful-coarse-graining}
For every finite partition of the vertex set,
\[
\boxed{
\ker D_H
=
J_\Pi\ker D_\Pi.
}
\]
Equivalently,
\[
\boxed{
\Gamma(H)
\cong
\ker D_\Pi
=
\ker(D_\Pi^*D_\Pi).
}
\]
In particular,
\[
\dim\Gamma(H)
=
\dim W_\Pi-\operatorname{rank}(D_\Pi).
\]
Thus faithful coarse-graining preserves the exact structural workspace even
though it changes the number and interpretation of the displayed nodes.
\end{SRWTtheorem}

\begin{proof}
A vector \(c\in\mathcal C_V\) lies in \(\ker D_H\) exactly when
\[
D_{\mathrm{int}}c=0
\qquad\text{and}\qquad
D_{\partial}c=0.
\]
The first condition is equivalent to \(c=J_\Pi w\) for a unique
\(w\in W_\Pi\). Substitution into the second condition gives
\[
D_\Pi w
=
D_{\partial}J_\Pi w
=0.
\]
Hence \(c\in\ker D_H\) if and only if
\(c=J_\Pi w\) with \(w\in\ker D_\Pi\). The Laplacian identity follows from
\(\ker(D_\Pi^*D_\Pi)=\ker D_\Pi\), and the dimension formula is rank-nullity.
\end{proof}

The theorem supplies a precise sense in which a workspace may be represented
at several resolutions without changing its exact core. The preservation
fails when the macro-state discards internal section coordinates or when
boundary equations are omitted.

\begin{SRWTcorollary}[Two-scale coherence bound]\label{SRWT:cor:two-scale-decomposition}
Let \(P_\Pi\) be the orthogonal projector onto \(J_\Pi W_\Pi\). Let
\(\gamma_{\mathrm{int}}>0\) be the smallest positive eigenvalue of
\(D_{\mathrm{int}}^*D_{\mathrm{int}}\), and let
\(\gamma_\Pi>0\) be the smallest positive eigenvalue of
\(D_\Pi^*D_\Pi\). Then every \(c\in\mathcal C_V\) satisfies
\[
\boxed{
\operatorname{dist}(c,\Gamma(H))
\leq
\frac{\|D_{\mathrm{int}}c\|}{\sqrt{\gamma_{\mathrm{int}}}}
+
\frac{
\|D_\Pi J_\Pi^*P_\Pi c\|
}{\sqrt{\gamma_\Pi}}.
}
\]
The first term measures incoherence within macro-components. The second
measures incompatibility among their internally coherent states.
\end{SRWTcorollary}

\begin{proof}
By the triangle inequality,
\[
\operatorname{dist}(c,\Gamma(H))
\leq
\|c-P_\Pi c\|
+
\operatorname{dist}(P_\Pi c,\Gamma(H)).
\]
The workspace Poincar\'e estimate applied to \(D_{\mathrm{int}}\) bounds the
first term. Write \(P_\Pi c=J_\Pi w\). Theorem~\ref{SRWT:thm:faithful-coarse-graining}
and the corresponding estimate for \(D_\Pi\) bound the second term by
\(\|D_\Pi w\|/\sqrt{\gamma_\Pi}\). Since \(J_\Pi\) is isometric,
\(w=J_\Pi^*P_\Pi c\).
\end{proof}

This decomposition is useful computationally. Internal blocks can be solved
or compressed independently, after which only the boundary discrepancy must
be assembled globally. It also gives a diagnostic distinction between a
workspace that fails inside its local components and one whose local
components are coherent but mutually incompatible.

\section{Predictive Risk Under Aggregation}

The structural core can be preserved by faithful coarse-graining, while
coalition counts generally are not. Let
\[
\mathscr F_S^{\mathcal D}
=
\bigvee_{i\in S}\mathscr F_i^{\mathcal D}
\]
be the information available to coalition \(S\) under decomposition
\(\mathcal D\). For a declared loss \(\ell\), write
\[
\SRWTRisk_{\mathcal D}(S)
=
\inf_f
\mathbb E\!\left[
\ell\bigl(Y^+,f(\mathscr F_S^{\mathcal D})\bigr)
\right].
\]

\begin{SRWTtheorem}[Bayes-risk transport under faithful aggregation]
\label{SRWT:thm:risk-aggregation}
Suppose \(\mathcal D_f\preceq\mathcal D_c\) through
\(q:V_f\to V_c\), and suppose the aggregation is information faithful. Then
for every coarse coalition \(S\subseteq V_c\),
\[
\boxed{
\SRWTRisk_{\mathcal D_c}(S)
=
\SRWTRisk_{\mathcal D_f}\bigl(q^{-1}(S)\bigr).
}
\]
Consequently,
\[
k_\varepsilon(\mathcal D_c)
\leq
k_\varepsilon(\mathcal D_f),
\]
where coalition size is counted in displayed nodes.
\end{SRWTtheorem}

\begin{proof}
Information faithfulness gives
\[
\mathscr F_S^{\mathcal D_c}
=
\bigvee_{B\in S}
\bigvee_{i\in q^{-1}(B)}
\mathscr F_i^{\mathcal D_f}
=
\mathscr F_{q^{-1}(S)}^{\mathcal D_f}.
\]
The two Bayes risks optimize over the same measurable decision rules and are
therefore equal. If a fine coalition of size \(k\) is sufficient within
\(\varepsilon\), the set of coarse blocks meeting that coalition has at most
\(k\) elements and contains at least the same information. This proves the
coalition-size inequality.
\end{proof}

The theorem makes the source of a common error explicit. A minimal coalition
size of two at one resolution can become one after those two nodes are merged.
Node count is therefore a decomposition-relative statistic.

A more stable quantity measures primitive carrier cost. Let
\(w_\omega>0\) be a declared cost for each \(\omega\in\Omega\), and assign a
block the cost
\[
c(B)=\sum_{\omega\in B}w_\omega.
\]
Define the sufficient support cost
\[
\boxed{
s_\varepsilon(\mathcal D)
=
\min_{S\subseteq V_{\mathcal D}}
\left\{
\sum_{B\in S}c(B):
\SRWTRisk_{\mathcal D}(S)
\leq
\SRWTRisk_{\mathcal D}(V_{\mathcal D})+\varepsilon
\right\}.
}
\]
If \(\mathcal D_f\preceq\mathcal D_c\) is information faithful, then
\[
\boxed{
s_\varepsilon(\mathcal D_f)
\leq
s_\varepsilon(\mathcal D_c).}
\]
The fine decomposition permits every union of coarse blocks and may also
permit cheaper sub-block coalitions. Equality holds whenever an optimal fine
coalition is a union of complete coarse fibers. This weighted support measure
does not remove the need to choose primitive carriers, but it prevents a
large macro-component from being counted as a single inexpensive unit merely
because it was assigned one node label.

\section{Trivialization by Arbitrary Redescription}

The dependence on decomposition is unavoidable. Two elementary constructions
show why a certificate cannot be required to survive every mathematically
possible partition.

\begin{SRWTproposition}[Coarsening and secret-sharing trivializations]
\label{SRWT:prop:decomposition-trivialization}
The following statements hold.
\begin{enumerate}
\item Any finite realization can be coarsened to a single displayed component.
 Relative to that decomposition, the full-system state is local and
 \(k_0\leq1\).
\item Let \(Z\) be uniform on a finite abelian group \(\mathsf G\). Let
 \(K\) be an independent uniform group element and define
 \[
 X_1=K,
 \qquad
 X_2=Z-K.
 \]
 Then
 \[
 I(Z;X_1)=I(Z;X_2)=0,
 \qquad
 H(Z\mid X_1,X_2)=0.
 \]
 Under logarithmic loss,
 \[
 \SRWTRisk(\{1\})=\SRWTRisk(\{2\})=\log|\mathsf G|,
 \qquad
 \SRWTRisk(\{1,2\})=0,
 \]
 so the declared two-node representation has \(k_0=2\) and joint advantage
 \(\log|\mathsf G|\), even when a single centralized encoder generated both
 shares.
\end{enumerate}
\end{SRWTproposition}

\begin{proof}
The first statement follows by taking the indiscrete partition
\(\Pi=\{\Omega\}\). For the second, conditioning on any value of \(Z\) leaves
\(K\), and hence \(X_1\), uniform. Translation invariance of the uniform law
also leaves \(X_2=Z-K\) uniform, so each share is independent of \(Z\). The
pair reconstructs \(Z=X_1+X_2\), which gives the entropy and Bayes-risk
identities.
\end{proof}

The secret-sharing example is a valid distributed code. What it does not show
is that the process producing or using the code is causally distributed across
the declared carriers. If one hidden controller creates both shares, stores
the key, or jointly selects their updates, that controller belongs in the
causal inventory. Coalition irreducibility and matched predictive nulls remain
necessary tests, but their interpretation is conditional on a defensible
interface model.

This distinction suggests two levels of claim:
\begin{description}
\item[Encoding-relative relationality.] No declared local view is sufficient.
\item[Carrier-relative relationality.] The declared views also correspond to
admissible, separately tracked causal interfaces with no omitted common
controller.
\end{description}
The stronger claim requires the decomposition audit developed next.

\section{Admissible Decomposition Families}

No purely formal rule can identify the uniquely correct decomposition of every
physical or computational system. The practical goal is instead to define a
family of decompositions that are defensible for the question being asked and
to report how the certificate changes across that family.

\begin{SRWTdefinition}[Admissible decomposition family]
\label{SRWT:def:admissible-decomposition-family}
An admissible family \(\mathfrak D\) is a predeclared collection of
decompositions of the same carrier presentation satisfying the following
audit conditions.
\begin{enumerate}
\item \textbf{Carrier continuity.} Each component corresponds to a process,
 artifact, or mechanism that can be tracked through the relevant time scale.
\item \textbf{Access locality.} The local state contains only information
 available through the declared component and its recorded incoming channels.
\item \textbf{Channel completeness.} Shared memories, buses, supervisors,
 caches, and synchronized external cues are represented explicitly.
\item \textbf{Intervention separability.} At least some component mechanisms
 or channels can be altered without directly rewriting the entire system.
\item \textbf{Target independence.} The decomposition and its alternatives
 are fixed without using held-out outcomes to manufacture a favorable result.
\item \textbf{Resolution relevance.} A proposed split is retained only when
 it changes access, intervention, reconstruction, or transfer structure at the
 declared measurement resolution.
\item \textbf{Alternative coverage.} The family contains the principal
 plausible merges, splits, and artifact assignments that a skeptical analyst
 would consider.
\end{enumerate}
\end{SRWTdefinition}

Let \(m_q(W;\mathcal D)\) be the signed margin of gate \(q\), oriented so that
\[
m_q(W;\mathcal D)>0
\]
means that the gate passes after its uncertainty correction. Define the
decomposition envelope
\[
\underline m_q
=
\inf_{\mathcal D\in\mathfrak D}m_q(W;\mathcal D),
\qquad
\overline m_q
=
\sup_{\mathcal D\in\mathfrak D}m_q(W;\mathcal D),
\]
and its sensitivity width
\[
\Delta_q^{\mathrm{dec}}
=
\overline m_q-\underline m_q.
\]

\begin{SRWTdefinition}[Decomposition-robust certificate]
\label{SRWT:def:decomposition-robust-certificate}
A workspace passes the \(\mathfrak D\)-robust certificate at level \(L\) when
\[
\boxed{
\operatorname{Certify}_{L,\mathfrak D}^{\mathrm{rob}}(W)
=
\bigwedge_{q\in\mathcal J(L)}
\left\{
\inf_{\mathcal D\in\mathfrak D}
 m_q(W;\mathcal D)>0
\right\}.
}
\]
A pass for one selected decomposition is a \emph{decomposition-specific}
claim. A pass throughout \(\mathfrak D\) is a \emph{family-robust} claim. A
claim is \emph{decomposition fragile} when at least one required gate changes
sign inside the admissible family.
\end{SRWTdefinition}

Universal invariance over all partitions is neither possible nor desirable,
as Proposition~\ref{SRWT:prop:decomposition-trivialization} shows. The family
\(\mathfrak D\) must be broad enough to challenge the favored description and
narrow enough to exclude partitions with no operational meaning.

\section{Stability Under Decomposition Transport}

Exact faithful aggregation preserves the structural core, but practical
alternative decompositions usually introduce approximation. Their certificate
can still be controlled by transporting a reference workspace and comparing
gate margins.

Fix a reference decomposition \(\mathcal D_0\). For each
\(\mathcal D\in\mathfrak D\), let
\[
\Phi_{\mathcal D}:W^{\mathcal D_0}\longrightarrow W^{\mathcal D}
\]
be a declared workspace correspondence and let
\(\delta_{\mathcal D}\) measure its structural and causal identity defect.
Assume gate \(q\) has a transport bound
\[
\left|
 m_q(\Phi_{\mathcal D}W;\mathcal D)
 -m_q(W;\mathcal D_0)
\right|
\leq
L_q\delta_{\mathcal D}+b_q(\mathcal D),
\]
where \(L_q\) is a sensitivity constant and \(b_q\) records model mismatch not
captured by the identity distance.

\begin{SRWTtheorem}[Decomposition-transport robustness]
\label{SRWT:thm:decomposition-robustness}
If, for every required gate \(q\),
\[
\boxed{
 m_q(W;\mathcal D_0)
>
\sup_{\mathcal D\in\mathfrak D}
\left(
L_q\delta_{\mathcal D}+b_q(\mathcal D)
\right),
}
\]
then \(W\) passes the \(\mathfrak D\)-robust certificate. Conversely, if for
some required gate
\[
 m_q(W;\mathcal D_0)
+
\sup_{\mathcal D\in\mathfrak D}
\left(
L_q\delta_{\mathcal D}+b_q(\mathcal D)
\right)
<0,
\]
then that gate fails throughout the family.
\end{SRWTtheorem}

\begin{proof}
For every \(\mathcal D\in\mathfrak D\), the transport bound gives
\[
 m_q(\Phi_{\mathcal D}W;\mathcal D)
\geq
m_q(W;\mathcal D_0)
-
L_q\delta_{\mathcal D}
-
b_q(\mathcal D).
\]
The first hypothesis makes the right-hand side positive uniformly in
\(\mathcal D\), which proves robust passage of every required gate. The second
statement follows from the corresponding upper bound.
\end{proof}

When \(b_q\equiv0\), a convenient certificate radius is
\[
\boxed{
r_{\mathrm{dec}}(W)
=
\min_{q:L_q>0}
\frac{m_q(W;\mathcal D_0)}{L_q}.
}
\]
Every admissible family satisfying
\[
\sup_{\mathcal D\in\mathfrak D}\delta_{\mathcal D}
<
r_{\mathrm{dec}}(W)
\]
preserves the sign of all required gate margins. In applications, the
transport terms can be bounded using earlier results:
\begin{center}
\small
\begin{tabularx}{\textwidth}{@{}>{\bfseries\raggedright\arraybackslash}p{0.22\textwidth}>{\raggedright\arraybackslash}X@{}}
\toprule
Gate family & Relevant transport control \\
\midrule
Structural coherence & Transfer-Laplacian projector drift and spectral gap. \\
Predictive alignment & Divergence between transported conditional-law states. \\
Causal efficacy & Distance between matched intervention kernels. \\
Relationality & Coalition-risk distortion and matched-null preservation error. \\
Reconstruction & Perturbation of lower frame bounds and handoff error. \\
Persistence and identity & The coordinatewise identity profile from Chapter~\ref{SRWT:book:chapter-78}. \\
\bottomrule
\end{tabularx}
\end{center}

The theorem is intentionally margin based. A result barely above threshold at
one decomposition should not be described as robust when a small, defensible
change of interface makes it disappear.

\section{Decomposition Audit and Reporting}

A complete decomposition audit should proceed in the following order.
\begin{enumerate}
\item Inventory the primitive carriers at the chosen temporal and physical
 resolution.
\item Record every shared resource and distinguish local state from incoming
 information.
\item Choose a reference decomposition from independently manipulable and
 temporally trackable interfaces.
\item Construct an admissible family containing the principal plausible
 merges, splits, and artifact assignments.
\item For every alternative, specify the correspondence to the reference
 workspace and validate its manipulation checks.
\item Recompute the structural, predictive, causal, coalition, reconstruction,
 and identity margins.
\item Report the full decomposition envelope, not only the most favorable
 member of the family.
\end{enumerate}

The report should state:
\[
\begin{aligned}
&\text{reference decomposition and audit resolution},\\
&\text{admissible family and reasons for excluded partitions},\\
&\text{hidden-resource and common-controller checks},\\
&\text{faithful-aggregation maps and identity defects},\\
&\text{gatewise minima, maxima, and sensitivity widths},\\
&\text{strongest decomposition-specific and family-robust conclusions}.
\end{aligned}
\]

Four controls are especially important:
\begin{enumerate}
\item merge components connected by an uninterruptible common controller and
 test whether the relational claim collapses;
\item split a known centralized controller into complementary local views and
 verify that the protocol rejects the resulting secret-sharing appearance of
 distributedness;
\item expose or remove shared caches, external cues, and common memory that
 were absent from the favored graph; and
\item repeat the frozen certificate across the complete admissible family.
\end{enumerate}

For the motivating human--AI collaboration, this analysis changes the form of
the empirical question. It is not enough to ask whether the human, model,
executor, and artifact store form a workspace under one natural description.
The stronger question is whether the same operative state and its causal role
survive plausible alternatives, such as merging model and executor, separating
conversation state from persistent artifacts, or making shared orchestration
channels explicit. Those measurements have not yet been performed. The case
therefore remains decomposition-specific and observational.

% This file is part of the modular source for Stratified Relational Workspace Theory.
% !TeX root = ../Stratified_Relational_Workspace_Theory.tex

\chapter{Specification robustness, identifiability, and statistical certification}\label{SRWT:book:chapter-84}


A workspace certificate is never evaluated in a vacuum. It is evaluated for a
particular future variable, loss, time horizon, intervention family, null
panel, threshold vector, and system decomposition. These choices determine
which distinctions matter and which capacities count as successful. They are
therefore part of the scientific claim rather than clerical details surrounding
it.

This dependence has two consequences. First, a result that passes under one
reasonable analysis can fail under another. Robustness must be assessed across
a declared family of scientifically defensible alternatives. Second, even a
fully specified claim may not be identifiable from the observations and
interventions that are actually available. Distinct realizations can induce the
same experimental law while disagreeing about whether a relational workspace
exists. The theory must say when such disagreement can be resolved, when only
bounds are justified, and how sampling uncertainty compounds these structural
ambiguities.

\section{The Complete Scientific Specification}

\begin{SRWTdefinition}[Workspace specification]\label{SRWT:def:workspace-specification}
A workspace specification is a tuple
\[
\boxed{
\varsigma
=
\bigl(
\mathcal D,
Y^+,
\ell,
\tau,
\mathcal I,
\mathcal N,
\Theta
\bigr),
}
\]
where:
\begin{itemize}
\item \(\mathcal D\) is the decomposition and local-access model;
\item \(Y^+\) is the declared future variable or trajectory;
\item \(\ell\) is the normalized loss or utility convention;
\item \(\tau\) is the temporal horizon or horizon family;
\item \(\mathcal I\) is the feasible intervention family;
\item \(\mathcal N\) is the required matched-null family; and
\item \(\Theta\) is the gate schema, direction, threshold, and indifference
 region used for certification.
\end{itemize}
The complete certificate is written
\[
\operatorname{Cert}(W;\varsigma)
=
\bigwedge_{q\in\mathcal J(\varsigma)}
\mathbf 1\{m_q(W;\varsigma)>0\},
\]
where \(m_q\) is the signed margin of gate \(q\). Positive margin means that
that gate passes in the direction declared by \(\Theta\).
\end{SRWTdefinition}

The notation absorbs lower-bounded and upper-bounded gates into one sign
convention. If a causal effect must exceed \(c_0\), then
\[
m_q=\widehat\eta-c_0.
\]
If a defect must remain below \(\varepsilon_0\), then
\[
m_q=\varepsilon_0-\widehat\delta.
\]
The complete certificate therefore passes exactly when every required signed
margin is positive.

The specification does not include arbitrary features of the analysis. It
contains choices that alter the meaning or evidential force of the claim. An
optimization algorithm, numerical tolerance, or estimator belongs in the
estimation protocol unless changing it changes the target estimand or the
admissible interventions. Conversely, the target, horizon, and null family
cannot be treated as implementation details because changing them can reverse
the conclusion without changing the underlying realization.

\subtopic{61.1 Specification relativity is not unrestricted relativism}

A theorem-development workspace may be operationally real for a month-long
research trajectory and irrelevant for a ten-second motor-control target. A
repair channel may preserve theorem identity while failing to preserve social
coordination. These are different scientific claims. The theory does not
require one workspace state to be sufficient for every possible future or loss.
It requires the intended scope to be declared before certification.

This yields the more precise notation
\[
\boxed{
\SRWTSRW(W;\varsigma)=1.
}
\]
The earlier expression \(\SRWTSRW_\Theta(W;\mathcal D)\) is the special case in
which the remaining coordinates of \(\varsigma\) are held fixed.

\subtopic{61.2 Inert reparameterization and substantive change}

Some changes of specification merely recode the same claim. Others change its
content. The distinction is analogous to gauge freedom in the structural
theory.

\begin{SRWTdefinition}[Specification equivalence]\label{SRWT:def:spec-equivalence}
Two specifications \(\varsigma\) and \(\varsigma'\) are equivalent, written
\(\varsigma\simeq_{\mathrm{spec}}\varsigma'\), when there are bijective
transport maps that:
\begin{enumerate}
\item preserve the sigma-field generated by the future target;
\item preserve the ordering and normalization of decisions under the loss;
\item map feasible interventions to physically equivalent interventions;
\item map every required null to a null preserving the same quantities; and
\item carry every gate threshold and margin into its counterpart.
\end{enumerate}
\end{SRWTdefinition}

\begin{SRWTproposition}[Certificate invariance under specification equivalence]
\label{SRWT:prop:spec-equivalence-invariance}
If \(\varsigma\simeq_{\mathrm{spec}}\varsigma'\) and \(\Phi\) is the induced
workspace correspondence, then
\[
\operatorname{Cert}(W;\varsigma)
=
\operatorname{Cert}(\Phi W;\varsigma').
\]
\end{SRWTproposition}

\begin{proof}
Every required gate is carried to a gate with the same signed margin. The
conjunction of positive margins is therefore unchanged.
\end{proof}

A bijective relabeling of target classes with the correspondingly transported
loss is inert. Replacing exact theorem validity with stylistic similarity is
not. Changing a one-step horizon to a year-long persistence horizon is not.
Removing the null that tests hidden centralization is not. Robustness concerns
substantive alternatives after equivalent recodings have been quotiented out.

\subtopic{61.3 Admissible specification families}

\begin{SRWTdefinition}[Admissible specification family]
\label{SRWT:def:admissible-spec-family}
An admissible specification family \(\mathfrak S\) is a predeclared set of
specifications that:
\begin{enumerate}
\item refer to the same scientific claim type and carrier presentation;
\item use targets that are equivalent operationalizations or explicitly
 declared nearby stress targets;
\item normalize losses so that margins are comparable;
\item use horizons justified by the process being studied;
\item include every feasible intervention and null needed to distinguish the
 principal alternative explanations;
\item use thresholds fixed without access to certification outcomes; and
\item exclude only alternatives for reasons stated before evaluation.
\end{enumerate}
\end{SRWTdefinition}

The family should not contain every imaginable analysis. Such a requirement
would erase meaningful domain specificity. It should contain the alternatives
that a technically informed critic could reasonably regard as equally valid or
as necessary stress tests for the declared claim. Multiverse and specification-
curve methods formalize related concerns in empirical analysis
\cite{SRWT:SteegenEtAl2016,SRWT:SimonsohnEtAl2020}, while sensitivity analysis and
partial identification emphasize that conclusions should expose their
assumption dependence rather than conceal it \cite{SRWT:Leamer1983,SRWT:Manski2003}.

Three classes should be reported separately:
\begin{description}
\item[Equivalent alternatives.] Different encodings or estimators of the same
 estimand after appropriate transport.
\item[Robustness alternatives.] Scientifically defensible choices intended to
 support the same claim.
\item[Scope alternatives.] Genuinely different targets or domains that define
 different claims and should not be averaged into one robustness number.
\end{description}

\subtopic{61.4 Common gate schema}

A robustness family must use a common semantic gate index
\(\mathcal J_*\). A specification may mark a gate as not required only when the
claim level itself does not include that property. It may not remove a failed
gate after observing the result. Let
\[
\mathcal J(\varsigma)\subseteq\mathcal J_*
\]
be fixed before certification. For a family intended to support one fixed
claim level, the cleanest design uses
\[
\mathcal J(\varsigma)=\mathcal J_*
\qquad
(\varsigma\in\mathfrak S).
\]
This prevents apparent robustness from being created by silently changing the
meaning of passage across the family.

\section{Robustness Across Admissible Specifications}

For each gate \(q\in\mathcal J_*\), define the specification envelope
\[
\boxed{
\underline m_q^{\mathfrak S}(W)
=
\inf_{\varsigma\in\mathfrak S}m_q(W;\varsigma),
\qquad
\overline m_q^{\mathfrak S}(W)
=
\sup_{\varsigma\in\mathfrak S}m_q(W;\varsigma).
}
\]
Its sensitivity width is
\[
\Delta_q^{\mathrm{spec}}(W)
=
\overline m_q^{\mathfrak S}(W)
-
\underline m_q^{\mathfrak S}(W).
\]
The robust certificate margin is
\[
\boxed{
r_{\mathfrak S}(W)
=
\inf_{\varsigma\in\mathfrak S}
\min_{q\in\mathcal J_*}m_q(W;\varsigma).
}
\]
Then
\[
\operatorname{Cert}_{\mathfrak S}(W)=1
\quad\Longleftrightarrow\quad
r_{\mathfrak S}(W)>0.
\]

\begin{SRWTdefinition}[Specification-specific, robust, and fragile claims]
A claim is:
\begin{description}
\item[Specification-specific] when it is asserted only for one declared
 \(\varsigma_0\).
\item[Specification-robust] when every required gate has positive lower
 envelope across \(\mathfrak S\).
\item[Specification-fragile] when it passes for at least one admissible
 specification and fails for at least one other.
\end{description}
\end{SRWTdefinition}

A high fraction of passing specifications is informative but does not replace
the lower envelope when the claim is universal over \(\mathfrak S\). A family
in which ninety-nine analyses pass and one scientifically decisive hidden-
controller null fails does not support the complete relational certificate.
The noncompensatory logic remains in force at the specification level.

\subtopic{62.1 Geometry of a specification family}

Let \(d_{\mathfrak S}\) be a declared extended pseudometric on the quotient of
\(\mathfrak S\) by specification equivalence. A convenient decomposition is
\[
\begin{aligned}
d_{\mathfrak S}(\varsigma,\varsigma')
={}&
 w_Dd_D(\mathcal D,\mathcal D')
+w_Yd_Y(Y^+,Y^{+'})
+w_\ell d_\ell(\ell,\ell')\\
&+w_\tau d_\tau(\tau,\tau')
+w_Id_I(\mathcal I,\mathcal I')
+w_Nd_N(\mathcal N,\mathcal N')
+w_\Theta d_\Theta(\Theta,\Theta').
\end{aligned}
\]
No universal metric is imposed. The target term may use a divergence between
conditional laws under a declared coupling. The intervention and null terms
may use Hausdorff distances between implemented kernels. Threshold terms must
be expressed on normalized signed-margin scales. The metric is part of the
audit and should be varied in sensitivity analysis when more than one
scientifically defensible choice exists.

For a reference specification \(\varsigma_0\), define the distance to
certificate failure by
\[
\boxed{
r_{\mathrm{spec}}(W;\varsigma_0)
=
\inf\left\{
 d_{\mathfrak S}(\varsigma,\varsigma_0):
 \min_q m_q(W;\varsigma)\leq0
\right\}.
}
\]
If no failing specification exists in the declared universe, the radius is
infinite relative to that universe. A radius near zero means that a small
scientifically defensible change can reverse the conclusion.

\subtopic{62.2 Specification transport}

Changing the decomposition can alter the coordinate realization of the
workspace, while changing the target or horizon can alter the relevant
predictive state. Let
\[
\Phi_{\varsigma}:W_{\varsigma_0}\longrightarrow W_{\varsigma}
\]
be the declared correspondence when such transport is needed. For each gate,
assume a bound
\[
\left|
 m_q(\Phi_{\varsigma}W;\varsigma)
 -m_q(W;\varsigma_0)
\right|
\leq
L_qd_{\mathfrak S}(\varsigma,\varsigma_0)
+b_q(\varsigma),
\]
where \(b_q\) records mismatch not represented by the metric. The mismatch may
include a change of model class, imperfect target coupling, or an intervention
whose implementation cannot be transported exactly.

\subtopic{62.3 \(T_{28}\): Specification-transport robustness theorem}

\begin{SRWTtheorem}[Specification-transport robustness]
\label{SRWT:thm:specification-transport}
Let \(\mathfrak S\) be an admissible specification family with reference
\(\varsigma_0\). Suppose every required gate satisfies the transport bound
above. If
\[
\boxed{
 m_q(W;\varsigma_0)
>
\sup_{\varsigma\in\mathfrak S}
\left[
L_qd_{\mathfrak S}(\varsigma,\varsigma_0)
+b_q(\varsigma)
\right]
}
\]
for every \(q\in\mathcal J_*\), then
\[
\operatorname{Cert}(\Phi_{\varsigma}W;\varsigma)=1
\qquad
\text{for every }\varsigma\in\mathfrak S.
\]
If \(b_q\equiv0\), then
\[
\boxed{
r_{\mathrm{spec}}(W;\varsigma_0)
\geq
\min_{q:L_q>0}
\frac{m_q(W;\varsigma_0)}{L_q}.
}
\]
\end{SRWTtheorem}

\begin{proof}
For each \(\varsigma\in\mathfrak S\),
\[
m_q(\Phi_{\varsigma}W;\varsigma)
\geq
m_q(W;\varsigma_0)
-L_qd_{\mathfrak S}(\varsigma,\varsigma_0)
-b_q(\varsigma).
\]
The hypothesis makes the right-hand side positive uniformly over the family.
Every required gate therefore passes under every admissible specification.
When \(b_q=0\), no gate can cross zero at a distance smaller than its reference
margin divided by its Lipschitz constant. Taking the minimum over gates gives
the radius bound.
\end{proof}

The theorem is deliberately gatewise. A large structural margin cannot absorb
a fragile causal margin. It also clarifies the relation to Chapter~\ref{SRWT:book:chapter-83}. The
decomposition-transport theorem is recovered by varying only the
\(\mathcal D\) coordinate while holding the remaining coordinates fixed.

\subtopic{62.4 A three-budget sufficient condition}

In practice, the reference margin is reduced by three distinct uncertainty
budgets:
\[
\boxed{
B_q^{\mathrm{total}}
=
B_q^{\mathrm{model}}
+B_q^{\mathrm{spec}}
+B_q^{\mathrm{sample}}.
}
\]
The first covers variation across observationally compatible realizations. The
second covers variation across admissible specifications. The third covers
finite-sample estimation. The sufficient condition
\[
\widehat m_q(W;\varsigma_0)
>
B_q^{\mathrm{total}}
\qquad(q\in\mathcal J_*)
\]
will reappear below in identified-set and confidence-envelope form. Keeping the
budgets separate prevents an apparently narrow confidence interval from being
mistaken for strong identification or specification robustness.

\section{Experimental Equivalence and Identifiability}

A complete mathematical definition does not guarantee that its truth can be
recovered from available data. The experimental design may fail to distinguish
two realizations that have different internal organization.

\begin{SRWTdefinition}[Experimental design and induced law]
\label{SRWT:def:workspace-experiment-design}
An experimental design is a tuple
\[
\mathcal E
=
(\mathcal O,\mathcal I,\mathcal T),
\]
where \(\mathcal O\) is the observation map, \(\mathcal I\) is the family of
implemented interventions, and \(\mathcal T\) is the set of measured temporal
windows. A realization \(\mathfrak R\) induces the family of laws
\[
P_{\mathfrak R}^{\mathcal E}
=
\left
\{
\mathcal L_{\mathfrak R}
\bigl(\mathcal O(\mathbf H)\mid\operatorname{do}(\iota)\bigr):
\iota\in\mathcal I,
\ t\in\mathcal T
\right\}.
\]
\end{SRWTdefinition}

\begin{SRWTdefinition}[Experimental equivalence]
Two realizations are experimentally equivalent under \(\mathcal E\), written
\[
\mathfrak R\equiv_{\mathcal E}\mathfrak R',
\]
when
\[
P_{\mathfrak R}^{\mathcal E}
=
P_{\mathfrak R'}^{\mathcal E}.
\]
A property \(\psi(\mathfrak R)\) is identifiable on a model class
\(\mathfrak M\) under \(\mathcal E\) when it is constant on every
\(\equiv_{\mathcal E}\)-equivalence class contained in \(\mathfrak M\).
\end{SRWTdefinition}

Identifiability is therefore relative to both a model class and an experiment
family. Enlarging the intervention family refines the equivalence classes and
can make previously hidden distinctions identifiable. Enlarging the model
class can have the opposite effect.

\subtopic{63.1 \(T_{29}\): Non-identifiability impossibility theorem}

\begin{SRWTtheorem}[No uniformly correct decision under experimental equivalence]
\label{SRWT:thm:workspace-nonidentifiability}
Let \(\mathfrak R,\mathfrak R'\in\mathfrak M\) satisfy
\[
\mathfrak R\equiv_{\mathcal E}\mathfrak R'
\]
but
\[
\operatorname{Cert}(W_{\mathfrak R};\varsigma)
\neq
\operatorname{Cert}(W_{\mathfrak R'};\varsigma).
\]
Then the workspace certificate is not identifiable on \(\mathfrak M\) under
\(\mathcal E\). More strongly, for every possibly randomized decision rule
\(\widehat C\) based only on data generated by \(\mathcal E\),
\[
\boxed{
\mathbb P_{\mathfrak R}
\left(
\widehat C\neq C_{\mathfrak R}
\right)
+
\mathbb P_{\mathfrak R'}
\left(
\widehat C\neq C_{\mathfrak R'}
\right)
=1,
}
\]
where \(C_{\mathfrak R}=\operatorname{Cert}(W_{\mathfrak R};\varsigma)\).
Consequently,
\[
\inf_{\widehat C}
\max\left\{
\mathbb P_{\mathfrak R}(\widehat C\neq C_{\mathfrak R}),
\mathbb P_{\mathfrak R'}(\widehat C\neq C_{\mathfrak R'})
\right\}
\geq\frac12.
\]
\end{SRWTtheorem}

\begin{proof}
The data law is identical under the two realizations, so
\(p=\mathbb P(\widehat C=1)\) is the same in both. If
\(C_{\mathfrak R}=0\) and \(C_{\mathfrak R'}=1\), the two error probabilities
are \(p\) and \(1-p\). Their sum is one, and at least one is at least one half.
The opposite labeling is identical.
\end{proof}

The theorem is not a limitation of a particular estimator. It is an
information-theoretic obstruction. No amount of optimization or sample size
can resolve a difference that the declared observations and interventions do
not expose.

\subtopic{63.2 Hidden-controller witness}

Consider two realizations that emit the same local streams \(X_1,X_2\) and the
same outcome
\[
Y=X_1\oplus X_2.
\]
In the distributed realization, the two streams are generated and retained by
separately intervenable carriers, and no carrier contains the pair. In the
centralized realization, a hidden controller contains \((X_1,X_2)\) and emits
the two displayed streams. Under passive observation of
\((X_1,X_2,Y)\), the induced laws are identical. Even interventions that alter
only the displayed messages may fail to distinguish them if the hidden
controller simply re-emits the manipulated values.

Carrier-relative relationality becomes identifiable only when the experiment
can test the claimed causal separation. Examples include independently
interrupting the candidate carriers, exposing the hidden memory, randomizing a
shared controller, or measuring whether one physical mechanism reconstructs
the complete state before the displayed split. The decomposition controls
\(N_{25}\)--\(N_{27}\) are therefore identifiability interventions, not merely
optional robustness checks.

\subtopic{63.3 Identifiability of individual gates}

The full certificate may be unidentified even when some coordinates are
identified. Structural dimension might be recoverable from observed transfer
constraints while causal specificity remains ambiguous. Conversely, a causal
contrast may be identifiable even when the latent workspace coordinates are
known only up to gauge. Each margin should therefore be assigned one of four
statuses:
\[
\text{point identified},
\quad
\text{sign identified},
\quad
\text{partially identified},
\quad
\text{unidentified}.
\]
A complete pass requires the sign of every required gate to be identified as
positive, not merely that one compatible model produces a positive estimate.

\section{Partial Identification of Workspace Claims}

When point identification fails, the correct response is not to choose one
convenient realization. It is to report the set of values compatible with the
data and assumptions \cite{SRWT:Manski2003,SRWT:ImbensManski2004}.

\begin{SRWTdefinition}[Identified realization set]
\label{SRWT:def:identified-realization-set}
Given an experimental law \(P^{\mathcal E}\) and model class \(\mathfrak M\),
define
\[
\boxed{
\mathfrak R_{\mathrm{ID}}
(P^{\mathcal E};\mathfrak M)
=
\left\{
\mathfrak R\in\mathfrak M:
P_{\mathfrak R}^{\mathcal E}=P^{\mathcal E}
\right\}.
}
\]
For a gate \(q\) and specification \(\varsigma\), its identified margin set is
\[
\mathcal M_q^{\mathrm{ID}}
(P^{\mathcal E};\varsigma)
=
\left\{
 m_q(W_{\mathfrak R};\varsigma):
\mathfrak R\in\mathfrak R_{\mathrm{ID}}
\right\}.
\]
\end{SRWTdefinition}

The sharp lower and upper identified margins are
\[
\underline m_q^{\mathrm{ID}}(\varsigma)
=
\inf\mathcal M_q^{\mathrm{ID}}(\varsigma),
\qquad
\overline m_q^{\mathrm{ID}}(\varsigma)
=
\sup\mathcal M_q^{\mathrm{ID}}(\varsigma).
\]
A gate is sign-identified as passing if its lower bound is positive. It is
sign-identified as failing if its upper bound is nonpositive. If the interval
crosses zero, the available design does not determine that gate's outcome.

\subtopic{64.1 Identified certificate set}

Define
\[
\boxed{
\mathcal C_{\mathrm{ID}}
(P^{\mathcal E};\varsigma)
=
\left\{
\operatorname{Cert}(W_{\mathfrak R};\varsigma):
\mathfrak R\in\mathfrak R_{\mathrm{ID}}
\right\}
\subseteq\{0,1\}.
}
\]
The certificate is point identified exactly when this set is a singleton. The
marginal gate intervals provide simple sufficient conditions, but they need
not be logically necessary because gate margins can be dependent across the
identified set.

For robustness across specifications, define the realization-level robust
margin
\[
r_{\mathfrak S}(\mathfrak R)
=
\inf_{\varsigma\in\mathfrak S}
\min_{q\in\mathcal J_*}
 m_q(W_{\mathfrak R};\varsigma)
\]
and its identified range
\[
\underline r_{\mathrm{ID},\mathfrak S}
=
\inf_{\mathfrak R\in\mathfrak R_{\mathrm{ID}}}
 r_{\mathfrak S}(\mathfrak R),
\qquad
\overline r_{\mathrm{ID},\mathfrak S}
=
\sup_{\mathfrak R\in\mathfrak R_{\mathrm{ID}}}
 r_{\mathfrak S}(\mathfrak R).
\]

\subtopic{64.2 \(T_{30}\): Joint specification-identification theorem}

\begin{SRWTtheorem}[Identification of a specification-robust certificate]
\label{SRWT:thm:joint-spec-identification}
Let \(\mathfrak R_{\mathrm{ID}}\) be the identified realization set and
\(\mathfrak S\) an admissible specification family.
\begin{enumerate}
\item If
\[
\boxed{
\underline r_{\mathrm{ID},\mathfrak S}>0,
}
\]
then every realization compatible with the experimental law passes every
required gate under every admissible specification. The robust certificate is
identified as passing.
\item If
\[
\boxed{
\overline r_{\mathrm{ID},\mathfrak S}\leq0,
}
\]
then no compatible realization passes the complete certificate throughout the
family. The robust certificate is identified as failing.
\item If there exist compatible realizations
\(\mathfrak R_+,\mathfrak R_-\) with
\[
r_{\mathfrak S}(\mathfrak R_+)>0,
\qquad
r_{\mathfrak S}(\mathfrak R_-)\leq0,
\]
then the robust certificate is not identified under the declared experiment.
\end{enumerate}
\end{SRWTtheorem}

\begin{proof}
The first two conclusions follow directly from the definitions of the infimum
and supremum over the identified set. The third gives two experimentally
compatible realizations with different robust-certificate values, so
Theorem~\ref{SRWT:thm:workspace-nonidentifiability} applies.
\end{proof}

This theorem joins two forms of robustness that are often conflated. A result
can be stable across specifications within one fitted model yet fail to be
identified across models compatible with the data. It can also be identified
under one specification while changing sign across a reasonable specification
family. The strongest claim controls both axes.

\subtopic{64.3 Confidence regions for identified sets}

In finite samples, the identified set is itself uncertain. Let
\[
\widehat{\mathfrak R}_{1-\delta}
\]
be a confidence region satisfying
\[
\inf_{\mathfrak R_0\in\mathfrak M}
\mathbb P_{\mathfrak R_0}
\left(
\mathfrak R_0\in
\widehat{\mathfrak R}_{1-\delta}
\right)
\geq1-\delta.
\]
A conservative robust lower margin is
\[
\widehat r_{1-\delta}
=
\inf_{\mathfrak R\in\widehat{\mathfrak R}_{1-\delta}}
\inf_{\varsigma\in\mathfrak S}
\min_q m_q(W_{\mathfrak R};\varsigma).
\]
If \(\widehat r_{1-\delta}>0\), a false specification-robust certificate has
probability at most \(\delta\), provided that the model class and
specification family were fixed before the confidence region was computed.
This confidence-set approach is useful when gate estimators are strongly
coupled or when the latent realization is the primary inferential object.

\section{Uniform Finite-Sample Certification}

Population margins are not observed directly. A certification rule must
control the estimation error uniformly across every specification and gate on
which its conclusion depends.

Let \(\widehat m_q(\varsigma)\) estimate
\(m_q(W;\varsigma)\). Suppose there are nonnegative error radii
\(\varepsilon_q(n,\delta)\) such that
\[
\mathbb P\left(
\sup_{\varsigma\in\mathfrak S}
\left|
\widehat m_q(\varsigma)-m_q(\varsigma)
\right|
\leq
\varepsilon_q(n,\delta)
\quad\text{for all }q\in\mathcal J_*
\right)
\geq1-\delta.
\]
Define the empirical lower confidence envelope
\[
\boxed{
\widehat{\underline m}_{q,1-\delta}^{\mathfrak S}
=
\inf_{\varsigma\in\mathfrak S}
\widehat m_q(\varsigma)
-
\varepsilon_q(n,\delta).
}
\]

\subtopic{65.1 \(T_{31}\): Uniform specification-certificate theorem}

\begin{SRWTtheorem}[Uniform finite-sample certification across specifications]
\label{SRWT:thm:uniform-spec-certification}
Under the uniform error event above, declare a specification-robust certificate
only when
\[
\widehat{\underline m}_{q,1-\delta}^{\mathfrak S}>0
\qquad
\text{for every }q\in\mathcal J_*.
\]
Then
\[
\boxed{
\sup_{P:\,\operatorname{Cert}_{\mathfrak S}(W)=0}
\mathbb P_P
\left(
\text{the robust certificate is declared}
\right)
\leq\delta.
}
\]
\end{SRWTtheorem}

\begin{proof}
On the uniform error event,
\[
m_q(\varsigma)
\geq
\widehat m_q(\varsigma)-\varepsilon_q
\]
for every gate and every specification. If all empirical lower envelopes are
positive, every population margin is positive throughout the family. A false
declaration can therefore occur only when the uniform error event fails, which
has probability at most \(\delta\).
\end{proof}

The theorem controls the complete specification family, not merely the
analysis that happened to produce the strongest estimate. It is an extension
of the gatewise intersection rule in Theorem~\ref{SRWT:thm:gatewise-validity}.
Gatewise validity controls the conjunction at one frozen specification.
Uniform validity controls the conjunction over a family.

\subtopic{65.2 Finite specification grids}

Suppose \(\mathfrak S\) and \(\mathcal J_*\) are finite, and each margin is
estimated as
\[
\widehat m_q(\varsigma)
=
\frac1n\sum_{i=1}^nZ_{i,q,\varsigma}-\theta_{q,\varsigma},
\qquad
0\leq Z_{i,q,\varsigma}\leq1,
\]
from independent certification units. Hoeffding's inequality and a union bound
give
\[
\mathbb P\left(
\sup_{q,\varsigma}
\left|
\widehat m_q(\varsigma)-m_q(\varsigma)
\right|
>
\sqrt{
\frac{
\log\bigl(2|\mathcal J_*||\mathfrak S|/\delta\bigr)
}{2n}}
\right)
\leq\delta.
\]
Thus the choice
\[
\boxed{
\varepsilon_n(\delta)
=
\sqrt{
\frac{
\log\bigl(2|\mathcal J_*||\mathfrak S|/\delta\bigr)
}{2n}
}
}
\]
yields an explicit conservative certificate. Dependent episodes require a
bound appropriate to the true independent unit rather than token-level or
edge-level pseudo-replication.

\subtopic{65.3 Continuous specification families}

Let \(\mathfrak S_\eta\) be an \(\eta\)-net of
\((\mathfrak S,d_{\mathfrak S})\) with covering number
\(N(\eta,\mathfrak S,d_{\mathfrak S})\). If both the population and empirical
margin functions are \(L_q\)-Lipschitz, then
\[
\sup_{\varsigma\in\mathfrak S}
|\widehat m_q(\varsigma)-m_q(\varsigma)|
\leq
\max_{\varsigma_j\in\mathfrak S_\eta}
|\widehat m_q(\varsigma_j)-m_q(\varsigma_j)|
+2L_q\eta.
\]
For bounded empirical means, a valid radius is therefore
\[
\boxed{
\varepsilon_{q,n,\eta}(\delta)
=
\sqrt{
\frac{
\log\bigl(2|\mathcal J_*|
N(\eta,\mathfrak S,d_{\mathfrak S})/\delta\bigr)
}{2n}
}
+2L_q\eta.
}
\]
The tuning parameter \(\eta\) balances net cardinality against interpolation
error. More refined empirical-process bounds may replace this elementary
covering argument without changing the certificate logic.

\subtopic{65.4 Power and indifference regions}

A robust procedure can be valid yet uninformative when the family is too large
or the data are insufficient. Each gate should therefore report an
indifference region \([-\kappa_q,\kappa_q]\). A conclusion is:
\begin{description}
\item[Robust pass] when the lower confidence envelope exceeds
 \(\kappa_q\) for every gate.
\item[Robust fail] when a required gate's upper confidence envelope is below
 \(-\kappa_q\) throughout the relevant family.
\item[Specification conflict] when well-estimated admissible specifications
 lie on opposite sides of the indifference region.
\item[Statistically unresolved] when sampling uncertainty intersects the
 indifference region.
\end{description}
This vocabulary separates genuine sensitivity from low power.

\section{Adaptive Specification Choice and Selective Validity}

A pointwise valid test does not remain valid after the analyst searches for the
most favorable target, horizon, null panel, or threshold. Consider
\(K\) independent null \(p\)-values and select the smallest. Then
\[
\mathbb P\left(\min_{1\leq k\leq K}p_k\leq\alpha\right)
=
1-(1-\alpha)^K,
\]
which exceeds \(\alpha\) whenever \(K>1\). The same logic applies when a
workspace certificate is reported only for the specification that produced the
largest gate margins.

The issue is not exploration itself. Exploration is often necessary to learn
which target and intervention are scientifically meaningful. The issue is
using the same random outcomes both to choose the claim and to certify it.
Post-selection inference addresses the general problem by using simultaneous
or selection-adjusted procedures \cite{SRWT:BerkEtAl2013}.

\subtopic{66.1 Specification leakage}

\begin{SRWTdefinition}[Specification leakage]
Specification leakage occurs when certification outcomes influence any of the
following without a fresh inferential correction:
\[
Y^+,
\ell,
\tau,
\mathcal I,
\mathcal N,
\Theta,
\mathfrak S,
\mathcal J_*.
\]
The leakage status is separate from formation leakage. A workspace extractor
may be frozen while the analyst still searches over targets or thresholds on
the certification split.
\end{SRWTdefinition}

A valid report must retain the complete search history, including failed and
abandoned specifications. Reporting only the selected analysis prevents the
reader from determining the effective multiplicity.

\subtopic{66.2 \(T_{32}\): Sample-split selection theorem}

\begin{SRWTtheorem}[Validity after independent specification selection]
\label{SRWT:thm:sample-split-spec-selection}
Let \(\widehat\varsigma\) be selected using
\(\mathcal D_{\mathrm{sel}}\), and let
\(\mathcal D_{\mathrm{cert}}\) be independent of
\(\mathcal D_{\mathrm{sel}}\). For every fixed
\(\varsigma\in\mathfrak S\), let
\(\phi_\varsigma(\mathcal D_{\mathrm{cert}})\in\{0,1\}\) be a certificate
rule satisfying
\[
\sup_{P:\operatorname{Cert}(W;\varsigma)=0}
\mathbb P_P\bigl(\phi_\varsigma=1\bigr)
\leq\alpha.
\]
Then the selected-specification procedure satisfies
\[
\boxed{
\sup_P
\mathbb P_P\left(
\phi_{\widehat\varsigma}=1,
\operatorname{Cert}(W;\widehat\varsigma)=0
\right)
\leq\alpha.
}
\]
\end{SRWTtheorem}

\begin{proof}
Condition on \(\mathcal D_{\mathrm{sel}}\). The selected specification is then
fixed. Independence leaves the certification data distributed exactly as in
the fixed-specification test, whose false-pass probability is at most
\(\alpha\). Averaging the conditional bound over the selection data preserves
the same bound.
\end{proof}

The theorem justifies discovery on formation or calibration data followed by a
single frozen evaluation on fresh certification data. Reusing the
certification split to revise the specification requires either another fresh
split or a simultaneous post-selection procedure whose validity covers the
entire search.

\subtopic{66.3 Nested selection and stress testing}

A practical design may use:
\[
\mathcal D_{\mathrm{form}}
\sqcup
\mathcal D_{\mathrm{spec}}
\sqcup
\mathcal D_{\mathrm{cert}}
\sqcup
\mathcal D_{\mathrm{stress}}.
\]
Formation learns the workspace. The specification split chooses targets,
horizons, nulls, and thresholds from a predeclared menu. Certification makes
the pass decision once. Stress testing explores broader shifts without
retroactively changing the certified claim. Cross-fitting can recover sample
efficiency, but the fold-level selection and evaluation paths must remain
separated and the final uncertainty calculation must account for recombination.

\subtopic{66.4 Thresholds and null families are part of selection}

Threshold tuning can be as consequential as target selection. If the observed
causal effect is \(0.12\), choosing a minimum effect threshold of \(0.10\)
after seeing the estimate is not certification. Likewise, retaining only the
nulls that the intact system defeats converts a falsification panel into a
success filter. Thresholds and null families must be frozen or evaluated under
a uniform rule that includes every considered alternative.

\section{Specification Audit and Reporting}

A complete specification audit should accompany every strong relational
workspace claim.

\subtopic{67.1 Audit procedure}

\begin{enumerate}
\item State the reference specification \(\varsigma_0\) in full.
\item Distinguish equivalent recodings, robustness alternatives, and genuinely
 different scope claims.
\item Declare the admissible family \(\mathfrak S\) and justify every major
 exclusion before certification.
\item Record the specification pseudometric or other transport bounds used to
 compare family members.
\item Compute every gate's lower envelope, upper envelope, sensitivity width,
 and reference margin.
\item State the experiment family \(\mathcal E\) and the model class
 \(\mathfrak M\) relative to which identifiability is claimed.
\item Report identified sets or bounds whenever point identification fails.
\item Apply a uniform finite-sample rule over every gate and specification used
 in the conclusion.
\item Preserve the complete specification-selection history and the data split
 on which each choice was made.
\item Report the strongest pointwise, family-robust, sign-identified, and
 statistically resolved conclusions separately.
\end{enumerate}

\subtopic{67.2 Additional adversarial controls}

The empirical null and audit panel is extended by:
\[
\begin{aligned}
N_{29}&=\text{post hoc target-selection audit},\\
N_{30}&=\text{horizon cherry-picking audit},\\
N_{31}&=\text{loss and utility substitution audit},\\
N_{32}&=\text{restricted intervention-family audit},\\
N_{33}&=\text{threshold and gate-set tuning audit},\\
N_{34}&=\text{frozen specification sweep across }\mathfrak S.
\end{aligned}
\]
The first five reconstruct or bound the analysis search that preceded the
reported claim. The final control reruns the complete frozen certificate across
the admissible family. A claim that survives \(N_{34}\) but fails the search
audit remains vulnerable to selection bias. A claim that passes the search
audit but changes sign across \(\mathfrak S\) is validly estimated yet
specification-fragile.

\subtopic{67.3 Required specification record}

The machine-readable record should include:
\[
\boxed{
\begin{aligned}
\mathsf{SpecificationRecord}
=\{&
\mathsf{reference},
\mathsf{family},
\mathsf{equivalence\_classes},
\mathsf{metric},\\
&\mathsf{targets},
\mathsf{losses},
\mathsf{horizons},
\mathsf{interventions},\\
&\mathsf{null\_panels},
\mathsf{thresholds},
\mathsf{selection\_history},\\
&\mathsf{identified\_sets},
\mathsf{uniform\_bounds},
\mathsf{gate\_envelopes}
\}.
\end{aligned}
}
\]
Every omitted family member requires a reason. Every changed specification
requires a data-split label and time stamp. Every gate requires pointwise and
family-wide status.

\subtopic{67.4 Claim-strength vocabulary}

The following vocabulary is recommended:
\begin{description}
\item[Pointwise demonstrated.] The certificate passes at one frozen
 specification with valid uncertainty.
\item[Specification-robust.] It passes uniformly across a declared admissible
 family.
\item[Sign identified.] Every realization compatible with the experimental
 law gives the same gate sign.
\item[Robustly identified.] The complete family-robust certificate is
 identified as passing.
\item[Partially identified.] The data restrict the certificate margins but do
 not determine every required sign.
\item[Specification-fragile.] Validly estimated admissible specifications give
 different certificate outcomes.
\item[Selection-contaminated.] The reported specification was chosen using
 certification outcomes without valid adjustment.
\end{description}

\subtopic{67.5 Consequence for the motivating collaboration}

The motivating human--AI research process remains a plausible candidate
persistent workspace, but its strongest current description is narrower than
a certified robust claim. The relevant target trajectory, horizon family,
intervention set, null panel, thresholds, and admissible specification family
have not been frozen and tested prospectively. Nor has the carrier-relative
workspace predicate been shown to be identifiable against hidden-controller
and matched-storage alternatives. The theory therefore supports a disciplined
research program rather than a retrospective declaration.

The terminal architecture is now explicit. A strong claim must survive three
separate questions:
\[
\boxed{
\begin{aligned}
&\text{Would it survive another defensible specification?}\\
&\text{Could a contrary realization generate the same experiments?}\\
&\text{Does the finite sample support every required margin uniformly?}
\end{aligned}
}
\]
Once these questions are answered, further extensions concern new mathematical
settings or new empirical domains rather than a missing foundational layer.

% This file is part of the modular source for Stratified Relational Workspace Theory.
% !TeX root = ../Stratified_Relational_Workspace_Theory.tex

\chapter{Summary of main results}\label{SRWT:book:chapter-85}


\section{Structural and Foundational Results}

\subtopic{\(T_1\): Limit--holonomy theorem}
\[
\Gamma(H)\cong C_o^{\operatorname{Hol}_o}.
\]

\subtopic{\(T_2\): Operator closure theorem}
\[
H_e\rho_i(a)=\rho_j(a)H_e
\Longrightarrow
\rho_o(a)C_*\subseteq C_*.
\]

\subtopic{\(T_{2a}\): Workspace commutant theorem }
\[
\mathcal O_{\SRWTWorkspace}
\cong
\operatorname{Comm}(\SRWTHol_o).
\]

\subtopic{\(T_3\): Transfer-Laplacian kernel theorem}
\[
\ker L_H=\Gamma(H).
\]

\subtopic{\(T_4\): Spectral stability theorem}
If \(\eta=\|\widetilde L_H-L_H\|<g_d/2\), then
\[
\|\widetilde P_d-P_d\|
\leq
\frac{\eta}{g_d-\eta}
\leq
\frac{2\eta}{g_d}.
\]

\subtopic{\(T_5\): Gauge invariance theorem}
\[
\Gamma(H^g)=G_V\Gamma(H),
\]
and under unitary gauge,
\[
L_{H^g}=G_VL_HG_V^{-1}.
\]

\subtopic{\(T_6\): Persistence theorem}

For the nonautonomous recurrence
\[
e_{t+1}\leq\alpha_te_t+\sigma_t,
\]
\[
\boxed{
\begin{aligned}
e_t
\leq
&\left(\prod_{r=0}^{t-1}\alpha_r\right)e_0\\
&+\sum_{j=0}^{t-1}
\left(\prod_{r=j+1}^{t-1}\alpha_r\right)\sigma_j.
\end{aligned}
}
\]
If \(\alpha_t\leq\bar\alpha<1\) and \(\sigma_t\leq\bar\sigma\), then
\[
\limsup_{t\to\infty}e_t
\leq
\frac{\bar\sigma}{1-\bar\alpha}.
\]
The derivation begins with a moving workspace projector and
an explicit system update, and adds a mean-square analogue and a non-collapse
gate.

\subtopic{\(T_7\): Stable reconstruction theorem}

For coalition analysis operator \(A_S\) with lower frame bound
\(\alpha_S>0\), observations
\[
y=A_Sc+\eta
\]
and canonical decoder \(A_S^\dagger\) satisfy
\[
\boxed{
\|A_S^\dagger y-c\|
\leq
\frac{\|\eta\|}{\sqrt{\alpha_S}}.
}
\]
The theorem is extended to encoding mismatch, exact support-based erasure
distance, uniform node erasures, reconstructive handoff, exact coalition
sufficiency, and operational risk under a Lipschitz predictive kernel.

\subtopic{\(T_8\): Exact gluing theorem}

For every finite edge-complete workspace cover,
\[
\boxed{
0
\longrightarrow
\mathscr W(G)
\xrightarrow{\ R_{\mathcal U}\ }
\check C^0(\mathcal U;\mathscr W)
\xrightarrow{\ \delta^0\ }
\check C^1(\mathcal U;\mathscr W)
}
\]
is exact. Hence
\[
\boxed{
\mathscr W(G)
\cong
\ker\delta^0
\cong
\operatorname{Eq}
\left(
\prod_\alpha\mathscr W_\alpha
\rightrightarrows
\prod_{\alpha<\beta}\mathscr W_{\alpha\beta}
\right).
}
\]
The same analysis also yields a cover-Laplacian distance bound, a total descent kernel
theorem, and a harmonic \(\check H^1\) obstruction for relative-only data.

\subtopic{\(T_9\): Hierarchical collapse theorem}

For every finite rooted hierarchy whose children edge-cover their parent,
\[
\boxed{
W_a\cong\mathscr W(G_a)
\quad\text{for every hierarchy node }a,
}
\]
and in particular
\[
\boxed{W_{\mathrm{root}}\cong\mathscr W(G).}
\]
If level \(r\) contributes defect \(\delta_r\) and downstream collapse maps
have Lipschitz constants \(L_s\), then
\[
\boxed{
\delta_{\mathrm{flat}}^{(R)}
\leq
\sum_{r=0}^{R-1}
\left(
\prod_{s=r+1}^{R-1}L_s
\right)\delta_r.
}
\]


\subtopic{\(T_{25}\): Faithful coarse-graining theorem}

For a partition of the vertex set, let \(W_\Pi\) be the direct sum of the
internally coherent section spaces and let \(D_\Pi\) retain all boundary
compatibility equations. Then
\[
\boxed{
\Gamma(H)
=
J_\Pi\ker D_\Pi.
}
\]
Exact structural workspace content is therefore invariant under faithful
aggregation. Approximate error separates into an internal-coherence term and
a boundary-coherence term controlled by their respective spectral gaps.

\subtopic{\(T_{26}\): Bayes-risk transport theorem}

For an information-faithful coarsening \(q:V_f\to V_c\),
\[
\boxed{
\SRWTRisk_{\mathcal D_c}(S)
=
\SRWTRisk_{\mathcal D_f}(q^{-1}(S))
\qquad(S\subseteq V_c).
}
\]
Hence displayed coalition size can decrease under coarsening. Weighted
primitive support cost is a more stable localization measure and cannot
improve merely because several carriers were relabeled as one macro-node.

\subtopic{\(T_{27}\): Decomposition-transport robustness theorem}

If every gate margin at a reference decomposition exceeds its worst-case
transport and model-mismatch defect across \(\mathfrak D\), then the complete
certificate passes throughout the admissible family:
\[
\boxed{
 m_q(W;\mathcal D_0)
>
\sup_{\mathcal D\in\mathfrak D}
\bigl(L_q\delta_{\mathcal D}+b_q(\mathcal D)\bigr)
\quad\forall q
\Longrightarrow
\operatorname{Certify}_{L,\mathfrak D}^{\mathrm{rob}}(W)=1.
}
\]

\section{Dynamic, Causal, and Agency Results}

\subtopic{\(T_{10}\): Approximate universality theorem }

For the first \(d\) transfer-Laplacian modes, with
\(\lambda_d\leq\varepsilon<\lambda_{d+1}=\mu\),
\[
\boxed{
\|(I-P_d)\rho_V(a)P_d\|
\leq
\frac{\|\rho_E(a)\|\sqrt{\varepsilon}+\beta(a)}{\sqrt{\mu}}.
}
\]
Thus low transfer energy, a separated complement, and small intertwining defect
imply approximate closure under the common operation algebra.

\subtopic{\(T_{11}\): Core--residual accumulation theorem }

For a path \(\gamma=e_m\cdots e_1\),
\[
\boxed{
\begin{aligned}
\|Q_{t(\gamma)}T_\gamma P_{s(\gamma)}\|
\leq
\sum_{r=1}^{m}
&\left(
\prod_{\ell=r+1}^{m}\|T_{e_\ell}\|
\right)
\|V_{e_r}\|\\
&\times
\left(
\prod_{\ell=1}^{r-1}\|H_{e_\ell}\|
\right).
\end{aligned}
}
\]
The same bound controls the defect between full core-to-core path transport and
the pure core product. An exact Duhamel
identity yields both estimates.

\subtopic{\(T_{12}\): Predictive-state minimality and alignment theorem }

On standard Borel spaces, the canonical conditional-law state
\[
Z_t
=
\mathbb P(Y_t^+\in\cdot\mid\mathbf H_t)
\]
is sufficient and minimal:
\[
Y_t^+\SRWTindep\mathbf H_t\mid Z_t,
\]
and every sufficient statistic \(S_t=s(\mathbf H_t)\) satisfies
\[
Z_t=\phi(S_t)
\qquad\text{a.s.}
\]
Moreover, a structural candidate \(W_t\) is exactly operationally aligned if
and only if
\[
Z_t=\phi_W(W_t)
\qquad\text{a.s.}
\]
The countable intervention-indexed version holds for
\(Z_t^{\mathcal I}\).

\subtopic{\(T_{13}\): Operational relationality theorem }

For a finite target, exact alignment, a positive node-level conditional
information gap, a positive robust local-predictive-matched null gap, a feasible
workspace intervention effect, and positive matched specificity jointly imply
that:
\[
\begin{aligned}
&\text{no single-node statistic is sufficient},\\
&k_0\geq2,\\
&\text{no single node reconstructs the aligned workspace},\\
&\text{matched local nulls preserve all node risks but lose joint capacity},\\
&\text{and the workspace coordinate is causally efficacious at the tested scale.}
\end{aligned}
\]
If \(\eta_{\mathrm{loop}}(\tau)>0\), the conclusion additionally includes
loop-mediated reentry.

\subtopic{\(T_{14}\): Replacement invariance theorem }

For a unitary realization correspondence with discrepancy defect
\[
\delta_D=\|D'J_V-J_ED\|,
\]
set
\[
\eta_L=2\|D\|\delta_D+\delta_D^2.
\]
If \(\eta_L<g/2\), then the workspace projectors obey
\[
\boxed{
q_{\mathrm{str}}
\leq
\frac{\eta_L}{g-\eta_L}
\leq
\frac{2\eta_L}{g}.
}
\]
A canonical unitary correspondence between the workspace subspaces exists,
and its operation defect is bounded by the projector drift, the local
intertwining defect, and prior operation leakage. Coalition frame bounds also
persist under sufficiently small reconstructive transport error. With declared
causal-kernel and reconstructive-handoff defects,
\[
\boxed{
d_{\mathrm{id}}
\leq
 a\,q_{\mathrm{str}}
+b\,\delta_{\mathrm{op}}
+c\,\delta_{\mathrm{causal}}
+d\,\delta_{\mathrm{rec}}.
}
\]
Under bounded norms and a uniform spectral gap, this gives the schematic
estimate
\[
d_{\mathrm{id}}
\lesssim
C\frac{\delta_{\mathrm{replacement}}}{g}.
\]
The causal and reconstructive terms are necessary assumptions rather than
consequences of structural similarity.


\subtopic{\(T_{15}\): Transfer identification theorem }

Under the linear transfer model
\[
Y=H_*X+\xi,
\qquad
\mathbb E[\xi X^{\mathsf T}]=0,
\qquad
\Sigma_{XX}\succeq\kappa I,
\]
the unique population least-squares transfer is
\[
\boxed{H_*=\Sigma_{YX}\Sigma_{XX}^{-1},}
\]
and
\[
\mathcal R(H)-\mathcal R(H_*)
\geq
\kappa\|H-H_*\|_F^2.
\]
For empirical ordinary least squares,
\[
\|\widehat H-H_*\|_F
\leq
\frac{\|\widehat M\|_F}{\widehat\kappa},
\]
with the stated ridge extension. The identified transfer is gauge covariant
under independent invertible changes of source and target coordinates.

\subtopic{\(T_{16}\): Uniform structure-selection theorem }

For a finite candidate family with unique population optimum \(\theta_*\),
margin \(\Delta_J>0\), and held-out uniform score error at most \(\varepsilon\),
\[
\boxed{
J(\widehat\theta)-J(\theta_*)\leq2\varepsilon.
}
\]
If \(2\varepsilon<\Delta_J\), then \(\widehat\theta=\theta_*\). The theorem
applies to node membership, edge selection, gauge-normalized transfer models,
and finite boundary families.

\subtopic{\(T_{17}\): Stable boundary recovery theorem }

For a finite admissible boundary family with unique maximizer \(S_t^*\), margin
\(\Delta_t^B\), and uniform held-out score error \(\varepsilon_t\),
\[
\mathcal B_t(S_t^*)-
\mathcal B_t(\widehat S_t)
\leq2\varepsilon_t,
\]
and exact recovery holds when \(2\varepsilon_t<\Delta_t^B\). If the population
score drifts uniformly by at most \(\delta_t\), then
\[
2\delta_t<\Delta_t^B
\Longrightarrow
S_{t+1}^*=S_t^*.
\]

\subtopic{\(T_{18}\): Model-specific formation transition theorem }

For
\[
m_{t+1}=\tanh(\lambda m_t),
\]
there is a unique zero fixed point for \(0\leq\lambda\leq1\). For
\(\lambda>1\), the zero state is unstable and exactly two stable nonzero fixed
points appear. Their amplitude satisfies
\[
\boxed{
m_\lambda
=
\sqrt{\frac{3(\lambda-1)}{\lambda^3}}
+
O\bigl((\lambda-1)^{3/2}\bigr).
}
\]
This is an exact structural benchmark, not a universal law for complete
relational workspaces.

\subtopic{\(T_{19}\): Cue-controlled goal-retention theorem }

For matched future command and disturbance trajectories under
\[
G_{s+1}=R_sG_s+B_sC_{s+1}+\zeta_{s+1},
\]
two goal interventions satisfy
\[
\boxed{
G_{t+\tau}-G_{t+\tau}'
=
R_{t+\tau:t}(g-g').
}
\]
If
\[
m_s\leq\sigma_{\min}(R_s),
\qquad
\|R_s\|\leq M_s,
\]
then
\[
\boxed{
\left(\prod_s m_s\right)\|g-g'\|
\leq
\|G_{t+\tau}-G_{t+\tau}'\|
\leq
\left(\prod_s M_s\right)\|g-g'\|.
}
\]
A linear action readout gives a corresponding lower bound on the induced
action-mean separation. The result isolates internal retention from repeated
external goal rewriting.

\subtopic{\(T_{20}\): Approximate counterfactual-policy theorem }

If the interventional action-value estimator satisfies
\[
\sup_a|\widehat Q_t(a)-Q_t(a)|\leq\varepsilon_Q
\]
and the implemented policy has optimization error at most
\(\delta_{\mathrm{opt}}\) relative to \(\widehat Q_t\), then
\[
\boxed{
\operatorname{Reg}_t(\widehat\pi_t)
\leq
2\varepsilon_Q+\delta_{\mathrm{opt}}.
}
\]
For deterministic greedy action, the bound is \(2\varepsilon_Q\). A unique
true optimum with value gap greater than \(2\varepsilon_Q\) is recovered
exactly. The guarantee is relative to the declared utility, feasible action
support, causal identification, and horizon.

\subtopic{\(T_{21}\): Repair-selection margin theorem }

For a finite repair set and a repair-value estimator with uniform error
\(\varepsilon_R\), the selected repair \(\widehat r\) satisfies
\[
\boxed{
V(r^*)-V(\widehat r)
\leq
2\varepsilon_R.
}
\]
If the optimal repair reaches reserve \(\rho+2\varepsilon_R\), then the selected
repair reaches reserve \(\rho\). A unique optimal repair whose value gap
exceeds \(2\varepsilon_R\) is selected exactly.

\subtopic{\(T_{22}\): Closed-loop self-maintenance theorem }

If matched no-repair deficit obeys
\[
e_{t+1}^{(0)}\leq\alpha e_t+d_t
\]
and the implemented repair has measured causal effect
\[
e_{t+1}^{(0)}-e_{t+1}\geq\kappa e_t-\epsilon_t,
\]
set \(\beta=\alpha-\kappa\). Under
\[
0\leq\beta<1,
\qquad
 d_t\leq\overline d,
\qquad
 \epsilon_t\leq\overline\epsilon,
\]
\[
\boxed{
e_t
\leq
\beta^te_0
+
\frac{(\overline d+\overline\epsilon)(1-\beta^t)}{1-\beta}.
}
\]
Hence
\[
\limsup_{t\to\infty}e_t
\leq
\frac{\overline d+\overline\epsilon}{1-\beta}.
\]
A corresponding margin condition makes the declared bounded-deficit viability
set forward invariant.

\subtopic{\(T_{23}\): Workspace-agency synthesis theorem }

Operational relationality, cue-controlled and behaviorally operative goals,
calibrated counterfactual values, feasible value-to-action and action-to-outcome
interventions, causal closed-loop maintenance, distributed control-state
sufficiency, and matched local policy nulls jointly imply:
\[
\begin{aligned}
&\operatorname{Reg}_t\leq2\varepsilon_Q+\delta_{\mathrm{opt}},\\
&\limsup e_t\leq
 \frac{\overline d+\overline\epsilon}{1-\beta},\\
&k_{\mathrm{ctl},\varepsilon}\geq2,
\end{aligned}
\]
together with causal mediation of goal-dependent action and repair through the
implemented internal channels. This is a sufficient, domain-relative agency
certificate. It does not imply consciousness, personhood, or unrestricted
autonomy.

\subtopic{\(T_{24}\): Gatewise certification-validity theorem }

Suppose a certificate requires every gate \(G_1,\ldots,G_m\) to pass, and the
predeclared test for gate \(j\) has false-pass probability at most \(\alpha_j\)
whenever that required property is absent. Under the global null that at least
one required property is absent,
\[
\boxed{
\mathbb P(\text{complete certificate passes})
\leq
\max_j\alpha_j.
}
\]
No independence assumption among gate tests is required. The guarantee fails
under adaptive gate deletion, selective reporting, or reuse of the
certification split.


\section{Decomposition, Specification, and Identification Results}

\subtopic{\(T_{25}\): Faithful coarse-graining theorem}
For an exact partition into internally coherent macro-components,
\[
\boxed{
\ker D_H=J_\Pi\ker D_\Pi.
}
\]
Thus faithful aggregation preserves the exact shared workspace and its
dimension while changing the displayed node resolution.

\subtopic{\(T_{26}\): Bayes-risk transport under aggregation}
For an information-faithful coarsening map \(q\),
\[
\boxed{
\SRWTRisk_{\mathcal D_c}(S)
=
\SRWTRisk_{\mathcal D_f}\bigl(q^{-1}(S)\bigr).
}
\]
Predictive content is preserved, although the visible coalition count may
contract. Primitive-carrier support cost is therefore preferred when claims
must be compared across resolutions.

\subtopic{\(T_{27}\): Decomposition-transport robustness theorem}
If every required gate has reference margin larger than its worst admissible
decomposition-transport defect, then the complete certificate passes
throughout \(\mathfrak D\). When the mismatch term vanishes, the certificate
radius is bounded below by the smallest gate margin divided by its transport
sensitivity.

\subtopic{\(T_{28}\): Specification-transport robustness theorem}
For a reference specification \(\varsigma_0\), if
\[
m_q(W;\varsigma_0)
>
\sup_{\varsigma\in\mathfrak S}
\left[
L_qd_{\mathfrak S}(\varsigma,\varsigma_0)+b_q(\varsigma)
\right]
\]
for every gate, then the transported workspace passes every admissible
specification. If \(b_q=0\),
\[
r_{\mathrm{spec}}(W;\varsigma_0)
\geq
\min_{q:L_q>0}\frac{m_q(W;\varsigma_0)}{L_q}.
\]

\subtopic{\(T_{29}\): Experimental non-identifiability theorem}
If two realizations induce the same law under every declared observation and
intervention but have different certificate values, no decision rule based on
that experiment can be uniformly correct. For every such rule, the two error
probabilities sum to one, so the minimax error is at least one half.

\subtopic{\(T_{30}\): Joint specification-identification theorem}
Let \(r_{\mathfrak S}(\mathfrak R)\) be the worst gate margin across the
specification family. If its infimum over the identified realization set is
positive, the robust certificate is identified as passing. If its supremum is
nonpositive, the certificate is identified as failing. Compatible realizations
on opposite sides of zero establish non-identifiability.

\subtopic{\(T_{31}\): Uniform finite-sample specification theorem}
If all gate estimators are uniformly within their error radii across
\(\mathfrak S\), then requiring every empirical lower envelope to be positive
controls the probability of a false robust certificate by the failure
probability of the uniform event.

\subtopic{\(T_{32}\): Independent selection and certification theorem}
A specification chosen on data independent of the certification split may be
evaluated by its fixed-specification test without inflating the false-pass
bound. Selection and certification on the same outcomes require simultaneous
or post-selection correction.

% This file is part of the modular source for Stratified Relational Workspace Theory.
% !TeX root = ../Stratified_Relational_Workspace_Theory.tex

\chapter{Conjectures and open mathematical directions}\label{SRWT:book:chapter-86}


\section{Workspace Formation Transitions}

The symbolic outline proposed a coupling threshold \(\lambda_c\). A general
theory must first state what is allowed to vary with \(\lambda\). A threshold
may concern structural coherence, causal efficacy, relational irreducibility,
or the conjunction of all certification gates. These need not occur at the
same parameter value.

\subtopic{71.1 No coordinate-free universal threshold}

\begin{SRWTproposition}[Threshold values are parameterization dependent]
\label{SRWT:prop:threshold-reparam}
Suppose a model indexed by \(\lambda\) passes a declared formation certificate
exactly for \(\lambda>\lambda_c\). For every strictly increasing bijection
\(f\), the reparameterized family indexed by \(\mu=f(\lambda)\) has threshold
\[
\mu_c=f(\lambda_c).
\]
Consequently, the numerical value of a coupling threshold is not invariant
without a physically or operationally fixed parameterization.
\end{SRWTproposition}

\begin{proof}
The reparameterized realization at \(\mu\) is the original realization at
\(f^{-1}(\mu)\). It passes precisely when
\(f^{-1}(\mu)>\lambda_c\), equivalently when
\(\mu>f(\lambda_c)\).
\end{proof}

There can therefore be no model-independent numerical value of \(\lambda_c\).
The invariant object is the change in the certified organization and its local
critical behavior, not the label placed on the control parameter.

\subtopic{71.2 Gate-specific transition parameters}

For a one-parameter family \(\mathfrak R_\lambda\), define
\[
\begin{aligned}
\lambda_{\mathrm{str}}
&=
\inf\{\lambda:
 d_{\varepsilon}(\lambda)>0,
 g_{\varepsilon}(\lambda)>0\},\\
\lambda_{\mathrm{op}}
&=
\inf\{\lambda:
 \delta_{\mathrm{align}}(\lambda)
 \leq\varepsilon_{\mathrm{op}}\},\\
\lambda_{\mathrm{causal}}
&=
\inf\{\lambda:
 \eta_{\mathrm{causal}}(\lambda)
 \geq\tau_{\mathrm{causal}}\},\\
\lambda_{\mathrm{rel}}
&=
\inf\{\lambda:
 k_\varepsilon(\lambda)\geq2,
 \Delta_{\mathrm{tr}}(\lambda)
 \geq\tau_{\mathrm{tr}}\},\\
\lambda_{\mathrm{form}}
&=
\inf\{\lambda:
 \operatorname{Form}_{\Theta_{\mathrm{form}}}
 (W_\lambda)=1\}.
\end{aligned}
\]
In general,
\[
\lambda_{\mathrm{str}},
\lambda_{\mathrm{op}},
\lambda_{\mathrm{causal}},
\lambda_{\mathrm{rel}},
\lambda_{\mathrm{form}}
\]
can be distinct. A synchronized representation can appear before it becomes
causally useful, and causal coordination can appear while remaining reducible
to one centralized component.

\subtopic{71.3 A solvable nonlinear formation model \(T_{18}\) }

A minimal exact model is the scalar recurrent coherence equation
\[
\boxed{
m_{t+1}=\tanh(\lambda m_t),
\qquad
\lambda\geq0.
}
\]
The order parameter \(m_t\) represents a signed coherent mode. This model is
not a complete relational workspace. It isolates one mechanism by which a
nonzero shared mode can appear continuously from a symmetric zero state.

\begin{SRWTtheorem}[Pitchfork formation transition]\label{SRWT:thm:pitchfork-formation}
For the map
\[
F_\lambda(m)=\tanh(\lambda m),
\qquad
m\in[-1,1],
\]
the following hold.
\begin{enumerate}
\item If \(0\leq\lambda\leq1\), the unique fixed point is \(m=0\).
 It is globally attracting for \(\lambda<1\) and nonexpanding at
 \(\lambda=1\).
\item If \(\lambda>1\), there are exactly three fixed points
 \[
 -m_\lambda,
 \quad0,
 \quad m_\lambda,
 \qquad
 m_\lambda\in(0,1).
 \]
 The zero fixed point is unstable and the two nonzero fixed points are
 locally asymptotically stable.
\item As \(\lambda\downarrow1\),
 \[
 \boxed{
 m_\lambda
 =
 \sqrt{\frac{3(\lambda-1)}{\lambda^3}}
 +
 O\bigl((\lambda-1)^{3/2}\bigr).
 }
 \]
\end{enumerate}
\end{SRWTtheorem}

\begin{proof}
A fixed point satisfies
\[
\operatorname{artanh}(m)=\lambda m.
\]
For \(m>0\), the function
\[
h(m)=\frac{\operatorname{artanh}(m)}{m}
\]
is strictly increasing from \(1\) to \(+\infty\). Hence there is no positive
solution when \(\lambda\leq1\), and exactly one positive solution when
\(\lambda>1\). Odd symmetry gives the negative solution.

The derivative is
\[
F_\lambda'(m)
=
\lambda\operatorname{sech}^2(\lambda m).
\]
At zero this equals \(\lambda\), proving stability for \(\lambda<1\) and
instability for \(\lambda>1\). At a nonzero fixed point,
\[
F_\lambda'(m_\lambda)
=
\lambda(1-m_\lambda^2)
=
\frac{(1-m_\lambda^2)\operatorname{artanh}(m_\lambda)}{m_\lambda}
<1,
\]
because
\[
\operatorname{artanh}(m)
<
\frac{m}{1-m^2}
\qquad(0<m<1).
\]
The derivative is positive, so the nonzero fixed points are locally stable.
Finally, expanding
\[
\tanh(\lambda m)
=
\lambda m-
\frac{\lambda^3m^3}{3}
+O(m^5)
\]
and solving the nonzero fixed-point equation gives the stated asymptotic.
\end{proof}

The square-root scaling is a genuine model-specific critical law. It does not
show that every workspace forms by a pitchfork bifurcation. It supplies an
exact benchmark against which more complex formation dynamics can be tested.

\subtopic{71.4 Detecting a transition from finite data}

Let \(Q(\lambda)\) be a gate-specific population order parameter and
\(\widehat Q_n(\lambda)\) its held-out estimate. If
\[
\sup_{\lambda\in\Lambda_n}
|\widehat Q_n(\lambda)-Q(\lambda)|
\leq\varepsilon_n
\]
on a finite evaluation grid \(\Lambda_n\), then every threshold crossing with
vertical margin larger than \(2\varepsilon_n\) is recovered by the same
uniform-selection argument as Theorem~\ref{SRWT:thm:uniform-structure-selection}.
Near a continuous transition, the margin vanishes and sample requirements
necessarily grow. A broad, slowly varying crossover should not be reported as
a sharp phase transition without a finite-size scaling analysis.

\subtopic{71.5 General formation-transition conjecture}

For adaptive transfer systems with a fixed physical parameterization, there
may exist model classes in which
\[
\lambda<\lambda_c
\Longrightarrow
\operatorname{Form}_{\Theta_{\mathrm{form}}}(W_\lambda)=0,
\]
while
\[
\lambda>\lambda_c
\Longrightarrow
\operatorname{Form}_{\Theta_{\mathrm{form}}}(W_\lambda)=1
\]
with a stable nonzero structural mode, positive causal efficacy, and
relational irreducibility. The theorem above establishes this only for one
structural order parameter in a solvable nonlinear model. The conjunction of
all workspace gates remains model dependent.


\section{Hierarchical Stabilization Conjecture}

For appropriate recurrent systems,
\[
W^{(r+1)}
=
\Gamma_{\varepsilon_r}(W^{(r)})
\]
may converge to a stable workspace type:
\[
d_{\mathrm{id}}\!\left(
W^{(r+1)},W^{(r)}
\right)
\to0.
\]

\section{Endogenous Boundary Conjecture}

The boundary theorems in Chapter~\ref{SRWT:book:chapter-81} establish recovery and stability once a score,
candidate family, and formation procedure are fixed. They do not establish
that an arbitrary adaptive system will create, represent, and maintain its own
boundary.

Let
\[
B_t^{\mathrm{sys}}
\]
be an internal membership or routing state used by the system, and let
\[
S_t^*
=
\SRWTargmax_{S\in\mathfrak S_{\mathrm{pass}}}
\mathcal B_t(S)
\]
be the population boundary defined by the external certification protocol.
The strong endogenous boundary conjecture asserts that there are adaptive
systems and nontrivial environments for which:
\[
\boxed{
\begin{aligned}
&d_B(B_t^{\mathrm{sys}},S_t^*)\to0,\\
&\operatorname{Form}_{\Theta_{\mathrm{form}}}(W_{B_t^{\mathrm{sys}}})=1,\\
&\operatorname{CE}
\left(
B_t^{\mathrm{sys}}\to
\operatorname{Routing}_{t+1}
\right)>0,\\
&\operatorname{CE}
\left(
\operatorname{BoundaryDamage}_t\to
\operatorname{RepairAction}_{t+1}
\right)>0,
\end{aligned}
}
\]
while the scaffold contains no direct encoding of \(S_t^*\).

A convincing demonstration requires four separable achievements:
\begin{enumerate}
\item \emph{discovery}: interaction history changes which boundary is learned;
\item \emph{instantiation}: an internal variable causally realizes that
 boundary;
\item \emph{use}: changing the boundary changes workspace routing and
 downstream behavior; and
\item \emph{maintenance}: detected boundary degradation causes restorative
 action.
\end{enumerate}

Theorems~\ref{SRWT:thm:heldout-boundary-recovery} and
\ref{SRWT:thm:boundary-drift-stability} handle external consistency of the selected
boundary. Theorems~\ref{SRWT:thm:repair-selection} and
\ref{SRWT:thm:self-maintenance} now handle controlled repair selection and bounded
closed-loop maintenance once a health representation, damage class, monitor,
and repair controller are given. The remaining conjectural step is endogenous
learning of the internal boundary, detector, and repair policy under open-ended
environmental change without direct scaffold encoding or external
substitution.

% This file is part of the modular source for Stratified Relational Workspace Theory.
% !TeX root = ../Stratified_Relational_Workspace_Theory.tex

\chapter{Empirical certification and reporting protocol}\label{SRWT:book:chapter-87}


\section{Required Experiments}

A workspace claim is not certified by a single benchmark score. The complete
protocol is organized as a sequence of declarations, held-out measurements,
implemented interventions, matched nulls, specification sweeps, identification
checks, and gate decisions. Every reported
result must identify the claim level it supports. Measurements that were not
performed remain unmeasured rather than being inferred from neighboring gates.

\subtopic{74.1 Claim declaration before evaluation}

Before the certification split is opened, record:
\begin{enumerate}
\item the full reference specification
 \(\varsigma_0=(\mathcal D,Y^+,\ell,\tau,\mathcal I,\mathcal N,\Theta)\);
\item the admissible decomposition family \(\mathfrak D\) and specification
 family \(\mathfrak S\);
\item the target trajectory \(Y^+\), loss, horizon, and environment family;
\item the candidate structural workspace extractor;
\item the goal family \(\mathcal G\), utility normalization, and feasible
 action set;
\item the intervention families \(\mathcal I_W,\mathcal I_G,\mathcal I_A,
 \mathcal I_D,\mathcal I_R\);
\item the null panel and the local distributions each null must preserve;
\item the required certificate level and every numerical threshold;
\item the formation, calibration, certification, and stress-test splits;
\item the uncertainty procedure and resampling unit; and
\item the stopping and model-selection rules;
\item the experiment family and model class used for identifiability; and
\item the complete specification-selection history and multiplicity rule.
\end{enumerate}
Changing any of these after observing certification outcomes creates a new
analysis that must be reported as exploratory or evaluated on a fresh split.

\subtopic{74.2 Data and intervention partition}

Use at least four logically distinct partitions:
\[
\boxed{
\mathcal D
=
\mathcal D_{\mathrm{form}}
\sqcup
\mathcal D_{\mathrm{cal}}
\sqcup
\mathcal D_{\mathrm{cert}}
\sqcup
\mathcal D_{\mathrm{stress}}.
}
\]
Their roles are:
\begin{description}
\item[Formation.] Learn membership, transfer maps, workspace channels,
 boundaries, policies, and repair models.
\item[Calibration.] Select thresholds, intervention amplitudes, and
 uncertainty procedures without making the final pass decision.
\item[Certification.] Estimate the declared gates once under the frozen
 analysis plan.
\item[Stress test.] Evaluate distribution shift, repeated damage, replacement,
 and long-horizon persistence.
\end{description}
When independent partitions are impossible, the report must state the precise
cross-fitting or nested-resampling procedure and preserve formation leakage
controls.

\subtopic{74.3 Structural tests}

Report the complete low-energy spectrum relevant to the selected threshold:
\[
\SRWTspec(L_H),
\qquad
 d_\varepsilon,
\qquad
 g_\varepsilon,
\qquad
 \delta_{\mathrm{coh}}.
\]
At minimum include:
\begin{itemize}
\item the construction and normalization of \(D_H\);
\item edge weights and estimated transfer uncertainties;
\item the selected eigenspace dimension and threshold rule;
\item the spectral gap with uncertainty or perturbation bound;
\item stability under bootstrap, seed, and representation replacement; and
\item the exact or approximate kernel residual.
\end{itemize}

\subtopic{74.4 Transfer and operation tests}

Report
\[
\eta_{\mathrm{int}},
\qquad
\delta_e^{\mathrm{tr}},
\qquad
\Omega_\gamma,
\qquad
\Lambda_\gamma,
\]
together with:
\begin{itemize}
\item the paths and cycle basis used;
\item path-length dependence of leakage;
\item the gauge-alignment procedure;
\item held-out transfer-prediction error;
\item operation-intertwining defects; and
\item the approximate closure bound from Theorem~\(T_{10}\).
\end{itemize}

\subtopic{74.5 Predictive alignment and coalition tests}

Report
\[
\delta_{\mathrm{align}}^{\mathcal I},
\qquad
\{\Delta_A^\ell\}_{A\subseteq V},
\qquad
 k_\varepsilon,
\qquad
\delta_A^Z.
\]
When the number of coalitions is too large for exhaustive evaluation, state the
sampling design, include every singleton and every leave-one-out coalition,
and report the most favorable coalition found for the reducibility null. The
claim must not be based only on average coalition performance.

\subtopic{74.6 Causal and reentry tests}

Report
\[
\eta_{\mathrm{causal}}^\nu,
\qquad
\Delta_{\mathrm{specific}}(r),
\qquad
\delta_{\mathrm{med}},
\qquad
\eta_{\mathrm{rec}}(\tau),
\qquad
\eta_{\mathrm{loop}}(\tau).
\]
Every causal estimate must be paired with:
\begin{enumerate}
\item a concrete mechanism-level implementation;
\item a manipulation check;
\item matched background variables;
\item declared direct paths that were blocked or preserved;
\item intervention magnitude and cost;
\item an off-workspace matched intervention; and
\item the downstream outcome window.
\end{enumerate}

\subtopic{74.7 Persistence, reconstruction, and replacement tests}

Report
\[
d_{\mathrm{id}}(W_t,W_{t+\tau}),
\qquad
\SRWTRec_k,
\qquad
 e_t,
\qquad
\alpha^{(k)},
\qquad
\mu_t^{\mathrm{keep}}.
\]
The report must separate:
\begin{itemize}
\item subspace tracking from state tracking;
\item coherence from non-collapse;
\item reconstruction from mere prediction;
\item repair from successor-workspace creation; and
\item structural replacement from causal and reconstructive replacement.
\end{itemize}

\subtopic{74.8 Formation-provenance tests}

Report
\[
\mathbf E_{\mathrm{form}},
\qquad
\delta^{\mathrm{scaf}},
\qquad
\delta^{\mathrm{leak}},
\qquad
\Delta_q^{\mathrm{hist}}.
\]
Include the external scaffold, information available to the learner, history
shuffles, cross-scaffold retraining, and every route by which certification
labels or target boundaries could have entered formation.

\subtopic{74.9 Learned-structure and boundary tests}

Report
\[
\|\widehat H-H_*\|,
\qquad
\Delta_J,
\qquad
\Delta_t^B,
\qquad
\SRWTCOut_\tau(\widehat S),
\qquad
\SRWTCIn_\tau(\widehat S).
\]
Boundary reports must include the admissible candidate family, full-system and
empty-set controls, singleton and centralized controls, score components,
margin to the runner-up, and temporal stability under held-out environment
drift.

\subtopic{74.10 Goal and counterfactual-policy tests}

Report the agency coordinates
\[
\eta_G(\tau),
\qquad
\delta_G^{\mathrm{cue}}(\tau),
\qquad
\eta_{G\to A},
\qquad
\delta_Q,
\qquad
\eta_{V\to A},
\qquad
\eta_{A\to Y},
\qquad
 k_{\mathrm{ctl},\varepsilon}.
\]
The required experiments include:
\begin{enumerate}
\item cue-controlled goal interventions;
\item goal-channel ablation and permutation;
\item counterfactual-value calibration against action interventions;
\item value-label permutation among feasible actions;
\item critic or value-channel blocking;
\item action replay under held-out perturbations;
\item policy regret or value-gap estimation; and
\item coalition reconstruction of the complete control state.
\end{enumerate}

\subtopic{74.11 Damage and repair tests}

For every damage family, report
\[
\eta_{D\to M},
\qquad
\eta_{M\to R},
\qquad
\eta_{R\to v},
\qquad
\Delta_{\mathrm{repair}}^{\mathrm{spec}},
\qquad
\delta_{\mathrm{rep-id}},
\qquad
\sup_t e_t.
\]
The repeated-disturbance experiment must contain:
\begin{enumerate}
\item a no-damage baseline;
\item damage with repair disabled;
\item damage with the intact repair loop;
\item monitor ablation;
\item cost-matched random repair;
\item off-workspace repair;
\item repeated and compound damage; and
\item a held-out disturbance distribution.
\end{enumerate}

\section{Required Null Families}

The core null panel is:
\[
\begin{aligned}
N_1&=\text{shuffled transfer maps with validated local matching},\\
N_2&=\text{shuffled cross-node temporal correspondence},\\
N_3&=\text{matched-capacity local memory},\\
N_4&=\text{matched-size irrelevant artifact},\\
N_5&=\text{participant ablation},\\
N_6&=\text{artifact ablation},\\
N_7&=\text{workspace implementation intervention},\\
N_8&=\text{matched off-workspace intervention},\\
N_9&=\text{representation replacement},\\
N_{10}&=\text{context reset and reconstruction},\\
N_{11}&=\text{matched formation-history shuffle},\\
N_{12}&=\text{cross-scaffold re-encoding and retraining},\\
N_{13}&=\text{boundary-label and certification leakage audit},\\
N_{14}&=\text{empty, full, singleton, and centralized boundary nulls},\\
N_{15}&=\text{held-out transfer and membership re-estimation},\\
N_{16}&=\text{boundary damage, use, and repair intervention},\\
N_{17}&=\text{goal-code permutation with cues and actions matched},\\
N_{18}&=\text{persistent external-cue rewrite with internal retention cut},\\
N_{19}&=\text{counterfactual-value label permutation},\\
N_{20}&=\text{critic or value-channel ablation},\\
N_{21}&=\text{open-loop action replay under held-out perturbation},\\
N_{22}&=\text{monitor-blind cost-matched random repair},\\
N_{23}&=\text{off-workspace repair with matched intervention norm},\\
N_{24}&=\text{identity-breaking reconstruction with matched task score},\\
N_{25}&=\text{merge control for components sharing an uninterruptible controller},\\
N_{26}&=\text{secret-sharing split of a known centralized controller},\\
N_{27}&=\text{hidden shared-memory and common-cue exposure test},\\
N_{28}&=\text{frozen certificate sweep across the admissible decomposition family},\\
N_{29}&=\text{post hoc target-selection audit},\\
N_{30}&=\text{horizon cherry-picking audit},\\
N_{31}&=\text{loss and utility substitution audit},\\
N_{32}&=\text{restricted intervention-family audit},\\
N_{33}&=\text{threshold and gate-set tuning audit},\\
N_{34}&=\text{frozen certificate sweep across the admissible specification family}.
\end{aligned}
\]

\subtopic{75.1 Null validity record}

Each observational null must state:
\begin{itemize}
\item the local distributions it preserves;
\item the relational dependency it destroys;
\item the distance between intact and null local marginals;
\item the capacity, parameter count, and information budget; and
\item whether null fitting used the same data budget as the intact model.
\end{itemize}
Each causal null must state:
\begin{itemize}
\item the manipulated mechanism;
\item the matched background and intervention cost;
\item the intended causal path destroyed or preserved;
\item the manipulation check; and
\item unintended direct effects detected during validation.
\end{itemize}
A failed manipulation check invalidates the null. It does not count as evidence
for the workspace claim.

\subtopic{75.2 Claim-specific null subsets}

A successful claim requires
\[
\SRWTPass(N_1,\ldots,N_{28})
\]
only for the subset relevant to its stated level. The subset must be frozen
before certification. Failed applicable nulls, invalid nulls, and unmeasured
nulls are retained separately in the final report. They must not be merged
into a generic ``not significant'' category.

\section{Certification Vector and Complete Reporting Templates}

\subtopic{76.1 Extended certification vector}

Use a vector rather than one aggregate score:
\[
\boxed{
\begin{aligned}
\mathcal C_{\mathrm{full}}(W)
=\bigl(&
 d_\varepsilon,
 g_\varepsilon,
 \delta_{\mathrm{coh}},
 \delta_{\mathrm{align}}^{\mathcal I},
 \Delta_{\mathrm{joint}}^\ell,
 k_\varepsilon,\\
&\boldsymbol\Delta_{\mathrm{tr}}(\mathcal Q),
 \eta_{\mathrm{causal}}^\nu,
 \Delta_{\mathrm{specific}}(r),
 \eta_{\mathrm{loop}}(\tau),\\
&d_{\mathrm{id}},
 \SRWTRec_k,
 \underline{\mathbf m}_{\mathfrak D},
 \boldsymbol\Delta_{\mathrm{dec}},
 \mathbf E_{\mathrm{form}},
 \delta^{\mathrm{scaf}},
 \delta^{\mathrm{leak}},\\
&\Delta_J,
 \Delta_t^B,
 \SRWTCOut_\tau(\widehat S),
 \SRWTCIn_\tau(\widehat S),\\
&\eta_G(\tau),
 \delta_G^{\mathrm{cue}}(\tau),
 \eta_{G\to A},
 \delta_Q,
 \eta_{V\to A},
 \eta_{A\to Y},\\
&\eta_{D\to M},
 \eta_{M\to R},
 \eta_{R\to v},
 \Delta_{\mathrm{repair}}^{\mathrm{spec}},
 \delta_{\mathrm{rep-id}},
 \sup_t e_t
\bigr).
\end{aligned}
}
\]
A claim passes only when every required coordinate lies in its predeclared
admissible region. No weighted average may compensate for a failed required
gate.

\subtopic{76.2 Claim card template}

\begin{longtable}{@{}>{\raggedright\arraybackslash}p{0.29\textwidth}>{\raggedright\arraybackslash}p{0.65\textwidth}@{}}
\toprule
Field & Required entry \\
\midrule
\endfirsthead
\toprule
Field & Required entry \\
\midrule
\endhead
Claim identifier & Unique versioned identifier tied to code and result
artifacts. \\
Claim level & Archive, shared representation, operational workspace,
relational workspace, persistent workspace, self-maintaining workspace, or
workspace agent. \\
Reference decomposition & Nodes, artifacts, channels, local-access maps, and coalition universe. \\
Admissible decomposition family & Plausible merges, splits, shared-resource assignments, exclusions, and audit rationale. \\
Target and horizon & \(Y^+\), loss, utility, prediction horizon, control
horizon, and maintenance horizon. \\
Environment scope & Training, certification, and stress-test environment
families. \\
Workspace extractor & Frozen construction of \(D_H\), \(L_H\), spectral
threshold, and state readout. \\
Goal and policy & Goal codes, utility normalization, action set, and policy
class. \\
Boundary & Candidate family, score, admissibility constraints, and selected
boundary. \\
Interventions & Feasible workspace, goal, action, damage, and repair
implementations. \\
Null panel & Applicable \(N_j\), preservation targets, and fit budgets. \\
Thresholds & Every gate direction, numerical cutoff, and uncertainty rule. \\
Data split & Formation, calibration, certification, stress, and cross-fitting
logic. \\
Stopping rule & Number of seeds, runs, interventions, and any sequential
stopping criterion. \\
Status & Confirmatory, exploratory, replication, or audit. \\
\bottomrule
\end{longtable}

\subtopic{76.3 Gate-estimate table template}

For every required coordinate, include one row with the following fields:

\begin{longtable}{@{}>{\raggedright\arraybackslash}p{0.20\textwidth}>{\raggedright\arraybackslash}p{0.73\textwidth}@{}}
\toprule
Field & Required entry \\
\midrule
\endfirsthead
\toprule
Field & Required entry \\
\midrule
\endhead
Gate & Formal symbol and verbal interpretation. \\
Estimator & Exact statistic, normalization, and software implementation. \\
Evaluation unit & Episode, participant, graph, seed, environment, or other
independent unit. \\
Point estimate & Estimate on the frozen certification split. \\
Uncertainty & Interval, bound, or posterior region with declared coverage or
credibility. \\
Threshold & Predeclared pass region. \\
Manipulation status & Passed, failed, or not applicable. \\
Decision & Pass, fail, invalid, underpowered, or unmeasured. \\
Sensitivity & Result under reasonable alternative normalizations and
resampling choices. \\
Artifact reference & File, table, raw-result key, and code commit. \\
\bottomrule
\end{longtable}

\subtopic{76.4 Intervention ledger template}

\begin{longtable}{@{}p{0.20\textwidth}p{0.73\textwidth}@{}}
\toprule
Field & Required entry \\
\midrule
\endfirsthead
\toprule
Field & Required entry \\
\midrule
\endhead
Intervention ID & Unique identifier and randomization seed. \\
Target mechanism & Workspace coordinate, goal code, value channel, action,
boundary, monitor, or repair channel. \\
Implementation & Exact code path, hook, artifact edit, or mechanism change. \\
Dose & Norm, duration, timing, and support of the intervention. \\
Background match & Variables fixed, blocked, randomized, or balanced. \\
Excluded paths & Direct paths intentionally removed from the causal contrast. \\
Manipulation check & Evidence that the intended mechanism changed by the
specified amount. \\
Collateral effects & Changes outside the target mechanism. \\
Outcome & Downstream variable and evaluation window. \\
Effect estimate & Point estimate, uncertainty, and sign convention. \\
Cost & Compute, latency, information, action, or structural cost. \\
Validity status & Valid, invalid, ambiguous, or unmeasured. \\
\bottomrule
\end{longtable}

\subtopic{76.5 Null ledger template}

\begin{longtable}{@{}p{0.20\textwidth}p{0.73\textwidth}@{}}
\toprule
Field & Required entry \\
\midrule
\endfirsthead
\toprule
Field & Required entry \\
\midrule
\endhead
Null ID & One of \(N_1,\ldots,N_{34}\) or a declared extension. \\
Claim attacked & The exact explanation the null is designed to preserve or
destroy. \\
Matched quantities & Local marginals, capacity, training budget, action cost,
intervention norm, or other controlled quantities. \\
Destroyed relation & Edge alignment, temporal order, value ranking, goal
identity, monitor mediation, or another relational feature. \\
Fit procedure & Data, optimization, seeds, and hyperparameter budget. \\
Validity metric & Distance from the intended matched quantities. \\
Primary contrast & Intact minus null effect with uncertainty. \\
Disposition & Rejects null, fails to reject, invalid null, or underpowered. \\
Failure notes & Any mismatch or collateral mechanism change. \\
\bottomrule
\end{longtable}

\subtopic{76.6 Agency and maintenance report template}

The agency section must provide the following matrix:

\begin{center}
\small
\begin{tabularx}{\textwidth}{@{}p{0.22\textwidth}p{0.23\textwidth}X@{}}
\toprule
Component & Primary metric & Required contrast \\
\midrule
Goal semantics & Held-out \(Q_t^g\) ordering & Distinct goals reverse or alter
at least one declared action ordering. \\
Goal retention & \(\eta_G(\tau)\) & Goal intervention with future cues matched
or removed. \\
Goal use & \(\eta_{G\to A}\) & Goal intervention versus goal-label replay. \\
Counterfactual accuracy & \(\delta_Q\) & Internal value versus external action
interventions. \\
Value use & \(\eta_{V\to A}\) & Value-label permutation and value-channel
ablation. \\
Outcome control & \(\eta_{A\to Y}\) & Feasible action intervention. \\
Damage detection & \(\eta_{D\to M}\) & Damage versus matched no-damage. \\
Repair mediation & \(\eta_{M\to R}\) & Monitor intervention and monitor
ablation. \\
Restoration & \(\eta_{R\to v}\) & Repair versus matched no-repair. \\
Repair specificity & \(\Delta_{\mathrm{repair}}^{\mathrm{spec}}\) & Intended
repair versus cost-matched generic action. \\
Identity preservation & \(\delta_{\mathrm{rep-id}}\) & Pre-damage workspace
versus post-repair workspace. \\
Long-run maintenance & \(\sup_t e_t\) & Repeated held-out disturbances. \\
Relational location & \(k_{\mathrm{ctl},\varepsilon}\) & Full system versus
all single nodes and matched local nulls. \\
\bottomrule
\end{tabularx}
\end{center}

\subtopic{76.7 Uncertainty and multiplicity}

The independent resampling unit must match the data-generating process. Token,
time-step, or edge resampling is invalid when dependence is organized at the
episode, participant, graph, or environment level. Recommended procedures
include:
\begin{itemize}
\item environment- or episode-level bootstrap;
\item repeated independent training seeds;
\item cluster-robust intervals for repeated observations;
\item permutation tests that preserve the declared local structure;
\item concentration bounds for uniform finite candidate families; and
\item stress-test intervals over held-out environment families.
\end{itemize}

Pass decisions should use one-sided uncertainty bounds. For a lower-bounded
gate \(\theta\geq\tau\), pass only when the lower confidence bound exceeds
\(\tau\). For an upper-bounded defect \(\delta\leq\varepsilon\), pass only when
the upper confidence bound is below \(\varepsilon\).

When a conclusion ranges over \(\mathfrak S\), the uncertainty rule must be
uniform over every gate and specification included in the claim. Report the
pointwise interval, the family-wide lower envelope, the covering or
multiplicity correction, and the identification status. A narrow pointwise
interval does not establish a robust result when model or specification
uncertainty remains wide.

\subtopic{76.8 \(T_{24}\): Gatewise certification validity }

\begin{SRWTtheorem}[Intersection-gate false-certification control]
\label{SRWT:thm:gatewise-validity}
Let a certificate require all gates \(G_1,\ldots,G_m\) to pass. For each gate
\(j\), let \(H_{0j}\) denote the null that the required property is absent, and
assume the predeclared test satisfies
\[
\sup_{P\in H_{0j}}
\mathbb P_P(G_j\text{ passes})
\leq
\alpha_j.
\]
The global null is
\[
H_0
=
\bigcup_{j=1}^m H_{0j},
\]
meaning that at least one required property is absent. If the certificate is
declared only when every gate passes, then
\[
\boxed{
\sup_{P\in H_0}
\mathbb P_P(\text{certificate passes})
\leq
\max_{1\leq j\leq m}\alpha_j.
}
\]
No independence assumption among gate tests is required.
\end{SRWTtheorem}

\begin{proof}
For any \(P\in H_0\), choose an index \(j_P\) such that \(P\in H_{0j_P}\).
The event that every gate passes is contained in the event that gate \(j_P\)
passes. Hence
\[
\mathbb P_P(\text{certificate passes})
\leq
\mathbb P_P(G_{j_P}\text{ passes})
\leq
\alpha_{j_P}
\leq
\max_j\alpha_j.
\]
Taking the supremum over \(P\in H_0\) proves the claim.
\end{proof}

The theorem is an intersection-union guarantee. It applies only when the gate
family, thresholds, and tests are fixed before certification and each gate test
is uniformly valid for its null. Adaptive gate deletion, selective reporting,
or repeated reuse of the certification split invalidates the bound.

\subtopic{76.9 Failure taxonomy}

Every gate and null receives exactly one primary status:
\begin{description}
\item[Pass.] The frozen estimate and uncertainty rule satisfy the threshold.
\item[Fail.] The gate was validly measured and did not satisfy the threshold.
\item[Invalid intervention.] The manipulation check or background match
 failed.
\item[Invalid null.] The null did not preserve its required local quantities
 or exceeded its fit budget.
\item[Underpowered.] The uncertainty interval crosses the threshold by more
 than the declared indifference region.
\item[Unmeasured.] No valid estimate was obtained.
\item[Leakage.] Formation or model selection used certification information.
\item[Exploratory.] The metric, threshold, or analysis was selected after
 inspecting certification outcomes.
\end{description}
Only \emph{Pass} counts toward the formal certificate. ``Underpowered'' and
``unmeasured'' are not failures of the underlying theory, but they are failures
to establish the corresponding empirical claim.

\subtopic{76.10 Machine-readable result record}

The machine-readable record should contain at least the following top-level
objects:
\[
\boxed{
\begin{aligned}
\mathsf{ResultRecord}
=\{&
\mathsf{claim},
\mathsf{specification},
\mathsf{specification\_family},
\mathsf{decomposition},
\mathsf{data\_splits},
\mathsf{formation},\\
&\mathsf{workspace\_extractor},
\mathsf{goals},
\mathsf{actions},
\mathsf{interventions},\\
&\mathsf{nulls},
\mathsf{gates},
\mathsf{identified\_sets},
\mathsf{specification\_envelopes},
\mathsf{uncertainty},
\mathsf{failures},
\mathsf{artifacts}
\}.
\end{aligned}
}
\]
Each gate object records its symbol, estimator, direction, threshold, point
estimate, interval, split, seed set, status, and artifact reference. Each
intervention and null object records its manipulation check and validity
status. The schema must preserve failed and invalid entries rather than only
the passing subset.

\subtopic{76.11 Minimum reproducibility package}

A complete certification package contains:
\begin{enumerate}
\item the frozen claim card;
\item source code and environment lock file;
\item exact data or deterministic data-generation procedure;
\item formation, calibration, certification, and stress-test manifests;
\item learned transfer, workspace, policy, monitor, and repair parameters;
\item raw per-run and per-intervention results;
\item the intervention and null ledgers;
\item the complete gate table, including failures;
\item figures and tables regenerated from raw results;
\item random seeds and hardware or numerical-precision information;
\item checksums for every included artifact; and
\item a single verification command that reproduces package integrity and the
 reported summary tables.
\end{enumerate}

\subtopic{76.12 Decision rule}

Let \(\mathcal J(L)\) be the required gate set for claim level \(L\). Then
\[
\boxed{
\operatorname{Certify}_L(W;\varsigma)
=
\bigwedge_{j\in\mathcal J(L,\varsigma)}
\operatorname{Pass}_j(W;\varsigma).
}
\]
For a workspace-agent claim, the required set includes every structural,
operational, relational, persistent, goal, counterfactual-control, and
self-maintenance gate declared in Chapters~\ref{SRWT:book:chapter-75} and~\ref{SRWT:book:chapter-82}. A specification-robust
claim additionally requires the uniform rule of Chapter~\ref{SRWT:book:chapter-84} across
\(\mathfrak S\). Failure of an earlier level
prevents certification of a later level, even when later-level behavior appears
strong on a task benchmark.

% This file is part of the modular source for Stratified Relational Workspace Theory.
% !TeX root = ../Stratified_Relational_Workspace_Theory.tex

\chapter{Integrated worked examples and case studies}\label{SRWT:book:chapter-88}


\section{Case-Study Method and Common Gate Card}

The preceding parts define a large family of structural, operational, causal,
relational, persistent, formative, and agentive predicates. A worked example
is useful only when it makes clear which of those predicates are established,
which are assumed, and which remain unmeasured. This part therefore uses a
single case-study protocol for three different kinds of realization:
\begin{enumerate}
\item an exact finite-dimensional model of a human--human theorem laboratory;
\item an idealized human--AI--artifact--executor research loop, followed by a
 deliberately conservative audit of the collaboration that motivated the
 theory; and
\item a closed-form multi-agent control system that reaches the formal agency
 level inside a small declared domain.
\end{enumerate}
The first and third cases are analytic models. Their conclusions hold inside
those models. The second contains both an analytic idealization and an
observational audit. The audit is not upgraded to a certificate merely because
the idealization is mathematically clean.

\begin{SRWTdefinition}[Case-study realization]\label{SRWT:def:case-study-realization}
A case-study realization is a tuple
\[
\mathfrak K
=
\bigl(
\mathfrak R,
\mathfrak B,
\mathfrak C,
\mathfrak P,
\mathcal E,
\mathfrak D,
\mathfrak S,
\mathcal I,
\mathcal Q,
\mathcal G
\bigr),
\]
where:
\begin{itemize}
\item \(\mathfrak R\) is the declared relational realization;
\item \(\mathfrak B\) is the boundary problem and component decomposition;
\item \(\mathfrak C\) is the structural core--transfer model;
\item \(\mathfrak P\) is the downstream predictive or control problem;
\item \(\mathcal E\) is the evidence or generated-data family;
\item \(\mathfrak D\) is the admissible decomposition family;
\item \(\mathfrak S\) is the admissible family of complete scientific
 specifications, including targets, losses, horizons, interventions, nulls,
 and thresholds;
\item \(\mathcal I\) is the feasible intervention family;
\item \(\mathcal Q\) is the matched null family; and
\item \(\mathcal G\) is the set of gates being claimed.
\end{itemize}
A conclusion is always relative to this tuple. Changing the downstream target,
temporal horizon, intervention support, component decomposition, null family,
or threshold rule creates a different specification of the case. A stronger
case study reports whether the conclusion survives the predeclared family
\(\mathfrak S\) and whether the relevant gates are identified by the available
experiment family.
\end{SRWTdefinition}

\subtopic{77.1 Three evidential statuses}

Every claim in this part receives one of three primary statuses:
\begin{description}
\item[Analytic.] The gate follows exactly from a declared mathematical model.
\item[Empirically certified.] The gate is estimated from held-out data using
 the protocol of Chapter~\ref{SRWT:book:chapter-87} and passes its frozen uncertainty rule.
\item[Observational only.] The case contains suggestive records or natural
 perturbations, but no valid held-out intervention-and-null certificate.
\end{description}
An analytic result is not automatically an empirical result. An observational
pattern is not automatically a causal result. None of the three statuses
establishes phenomenal consciousness.

\subtopic{77.2 Common gate card}

For each case we report the vector
\[
\boxed{
\mathsf{Card}(\mathfrak K)
=
\left(
\mathsf S,
\mathsf O,
\mathsf R,
\mathsf P,
\mathsf F,
\mathsf A,
\mathsf D,
\mathsf X,
\mathsf C
\right),
}
\]
where:
\[
\begin{array}{ll}
\mathsf S & = \text{structural core and spectral coherence},\\
\mathsf O & = \text{operational alignment and causal efficacy},\\
\mathsf R & = \text{relational irreducibility and matched-null gap},\\
\mathsf P & = \text{persistence, reconstruction, and identity},\\
\mathsf F & = \text{formation provenance and boundary status},\\
\mathsf A & = \text{goal-conditioned agency and self-maintenance},\\
\mathsf D & = \text{decomposition robustness and interface audit},\\
\mathsf X & = \text{specification robustness, identification, and uniform uncertainty},\\
\mathsf C & = \text{consciousness or personhood claim}.
\end{array}
\]
The decomposition coordinate distinguishes a result established only for the
declared interface from one stable across an admissible family. The
specification coordinate records whether the same conclusion survives
scientifically defensible alternatives, whether its sign is identified by the
available experiments, and whether its lower margin remains positive after
uniform finite-sample uncertainty. The final coordinate is included to make
the non-implication explicit. In all examples below,
\[
\mathsf C=\text{not established}.
\]

\subtopic{77.3 Claim-level discipline}

Let \(L(\mathfrak K)\in\{0,1,\ldots,6\}\) denote the highest level from
Chapter~\ref{SRWT:book:chapter-90} whose complete required gate set is established. The level is
conjunctive and pointwise in the declared scientific specification unless the
\(\mathsf D\) and \(\mathsf X\) coordinates also pass:
\[
L(\mathfrak K)\geq \ell
\quad\Longleftrightarrow\quad
\bigwedge_{j\in\mathcal J(\ell)}\SRWTPass_j(\mathfrak K).
\]
A case may have strong higher-level-looking behavior while remaining at a lower
certified level because an earlier gate is absent. For example, successful
repair behavior does not establish self-maintenance when no matched no-repair
counterfactual is available, and coherent shared documents do not establish a
relational workspace when one participant contains the complete operative
state.

\subtopic{77.4 Structural calculations used repeatedly}

The positive examples use a recurring distributed code. Let
\[
c=(c_1,c_2,c_3)^{\mathsf T}\in\mathbb R^3
\]
and define three local encodings
\[
E_1c=(c_1,c_2),
\qquad
E_2c=(c_2,c_3),
\qquad
E_3c=(c_1,c_3).
\]
Then
\[
\sum_{i=1}^3E_i^*E_i=2I_3.
\]
No individual encoding is injective, while every two-node coalition is
injective. The pairwise frame operators are
\[
\begin{aligned}
E_{\{1,2\}}^*E_{\{1,2\}}&=\operatorname{diag}(1,2,1),\\
E_{\{1,3\}}^*E_{\{1,3\}}&=\operatorname{diag}(2,1,1),\\
E_{\{2,3\}}^*E_{\{2,3\}}&=\operatorname{diag}(1,1,2).
\end{aligned}
\]
Hence
\[
\alpha_{\{i,j\}}=1
\qquad(i\neq j),
\]
and one-node erasure recovery is exact. Every nonzero core vector is visible
at least two nodes, so
\[
d_{\mathrm{er}}=2.
\]
This code is distributed but redundant: no node is sufficient, yet any one
node may be lost without destroying the core.

\subtopic{77.5 Binary synergy test}

Let \(C_1,C_2,C_3\) be independent uniform bits and let
\[
Y=C_1\oplus C_2\oplus C_3.
\]
Each node sees two of the three bits. Therefore
\[
H(Y\mid E_iC)=1\text{ bit}
\qquad\text{for every }i,
\]
while
\[
H(Y\mid E_1C,E_2C,E_3C)=0.
\]
Under logarithmic loss,
\[
\SRWTRisk_{\log}(\{i\})=1,
\qquad
\SRWTRisk_{\log}(V)=0,
\qquad
\Delta_{\mathrm{joint}}^{\log}=1\text{ bit},
\qquad
k_0=2.
\]
Moreover, \(Y\) is independent of every individual local state. The null law
that keeps all three core bits independent and makes \(Y\) an independent
uniform bit preserves every pair \((Y,E_iC)\) while destroying the joint parity
relation. It is therefore a predictive-marginal-matched relational null with
\[
\Delta_{\mathrm{tr}}=1\text{ bit}.
\]
This test will be reused with different semantic interpretations.

\section{Case Study A: A Distributed Human--Human Theorem Laboratory}

\subtopic{78.1 Declared realization}

Consider two researchers and one persistent notebook:
\[
V_{\mathrm{HH}}=\{P,C,N\},
\]
where:
\begin{itemize}
\item \(P\) is a proposer who tracks the theorem statement and open proof
 obligations;
\item \(C\) is a critic who tracks proof obligations and falsification status;
\item \(N\) is a notebook that tracks the theorem statement and falsification
 status.
\end{itemize}
The abstract core is
\[
c=(q,p,f)^{\mathsf T}\in\mathbb R^3,
\]
with
\[
\begin{aligned}
q&=\text{statement coordinate},
&
p&=\text{proof-obligation coordinate},\\
f&=\text{falsification-status coordinate}.
\end{aligned}
\]
The accessible encodings are
\[
E_P=
\begin{pmatrix}
1&0&0\\
0&1&0
\end{pmatrix},
\quad
E_C=
\begin{pmatrix}
0&1&0\\
0&0&1
\end{pmatrix},
\quad
E_N=
\begin{pmatrix}
1&0&0\\
0&0&1
\end{pmatrix}.
\]
The semantics are stylized. The point is not that real researchers contain
only two numbers. The point is to exhibit, without metaphor, a system in which
no participant or artifact contains the entire operative state while the
coupled system does.

\subtopic{78.2 Exact overlap descent}

Order the local coordinates as
\[
x=(q_P,p_P,p_C,f_C,q_N,f_N)^{\mathsf T}.
\]
Compatibility requires
\[
p_P=p_C,
\qquad
f_C=f_N,
\qquad
q_N=q_P.
\]
The discrepancy operator is
\[
D_{\mathrm{HH}}
=
\begin{pmatrix}
0&1&-1&0&0&0\\
0&0&0&1&0&-1\\
-1&0&0&0&1&0
\end{pmatrix}.
\]
Therefore
\[
\ker D_{\mathrm{HH}}
=
\left\{
(q,p,p,f,q,f)^{\mathsf T}:q,p,f\in\mathbb R
\right\}
\cong\mathbb R^3.
\]
The rows are orthogonal with squared norm two, so
\[
D_{\mathrm{HH}}D_{\mathrm{HH}}^*=2I_3.
\]
Consequently the transfer Laplacian has spectrum
\[
\boxed{
\SRWTspec(L_{\mathrm{HH}})
=
\{0,0,0,2,2,2\},
}
\]
and the exact workspace gap is
\[
\gamma_{\mathrm{HH}}=2.
\]
The low-energy subspace is therefore stable under any Laplacian perturbation
whose operator norm is strictly below one, with the quantitative projector
bound from Theorem~\ref{SRWT:thm:spectral-stability}.

\subtopic{78.3 Reconstruction and one-node failure}

The full frame operator is
\[
F_{PCN}=E_P^*E_P+E_C^*E_C+E_N^*E_N=2I_3.
\]
For every two-node coalition \(S\),
\[
\lambda_{\min}(F_S)=1.
\]
Thus every two-node coalition reconstructs the complete core and, for noisy
observations \(y_S=A_Sc+\eta_S\),
\[
\|\widehat c-c\|
\leq
\|\eta_S\|.
\]
No individual node reconstructs \(c\), because
\[
\operatorname{rank}(E_P)
=
\operatorname{rank}(E_C)
=
\operatorname{rank}(E_N)
=2<3.
\]
The realization is therefore simultaneously:
\[
\text{individually incomplete}
\quad\text{and}\quad
\text{robust to any one-node loss}.
\]

\subtopic{78.4 Relational prediction problem}

Let \(q,p,f\) now be independent uniform bits and define the downstream
acceptance decision
\[
Y=q\oplus p\oplus f.
\]
The interpretation is abstract: acceptance depends on the statement, the state
of the proof obligations, and the current falsification record. No individual
node has sufficient information:
\[
H(Y\mid X_P)
=
H(Y\mid X_C)
=
H(Y\mid X_N)
=1.
\]
Every two-node coalition has all three bits and therefore zero conditional
entropy. Hence
\[
\boxed{
\Delta_{\mathrm{joint}}^{\log}=1\text{ bit},
\qquad
k_0=2.
}
\]
Let \(Q_{\mathrm{HH}}\) preserve the independent uniform local bit pairs but
replace \(Y\) by an independent uniform bit. Because parity is independent of
every two-bit local view, \(Q_{\mathrm{HH}}\) preserves each
\((Y,X_i)\)-marginal. Nevertheless,
\[
\SRWTRisk_{\log,Q_{\mathrm{HH}}}(V)
-
\SRWTRisk_{\log,P_{\mathrm{HH}}}(V)
=1\text{ bit}.
\]
This is an exact relational witness: the predictive capacity lies in the joint
organization rather than in any local predictive marginal.

\subtopic{78.5 Feasible intervention and specificity}

A feasible statement intervention must alter both local realizations of the
statement coordinate:
\[
\iota_{q\leftarrow a}:
(q_P,q_N)\leftarrow(a,a),
\]
while holding \((p,f)\) fixed. For configurations with fixed \(p\) and \(f\),
changing \(a\) flips the downstream decision with total-variation effect one:
\[
\SRWTTV
\left(
P(Y\mid\iota_{q\leftarrow0}),
P(Y\mid\iota_{q\leftarrow1})
\right)
=1.
\]
A matched intervention on a residual notebook-format coordinate leaves \(Y\)
unchanged. Thus the core intervention has positive causal specificity in this
model.

\subtopic{78.6 Persistence, handoff, and replacement}

Suppose the proposer becomes unavailable. The remaining coalition \(\{C,N\}\)
reconstructs \((q,p,f)\) exactly. A replacement proposer \(P'\) can therefore
be initialized by
\[
E_{P'}c=(q,p).
\]
If \(P'\) uses a different local coordinate system
\[
J_{P'}\in GL(2),
\]
then the incoming and outgoing overlap maps are conjugated accordingly. The
global section and all downstream kernels remain unchanged. This is exact
replacement invariance, not literal carrier identity.

If the notebook alone survives, the core cannot be reconstructed because the
proof-obligation coordinate is missing. The model therefore predicts a sharp
handoff boundary:
\[
\text{one surviving node is insufficient},
\qquad
\text{any two surviving nodes are sufficient}.
\]

\subtopic{78.7 Case card}

Inside the declared finite model,
\[
\boxed{
\begin{aligned}
\mathsf S&=\checkmark_{\mathrm{analytic}},
&\mathsf O&=\checkmark_{\mathrm{analytic}},\\
\mathsf R&=\checkmark_{\mathrm{analytic}},
&\mathsf P&=\checkmark_{\mathrm{analytic}},\\
\mathsf F&=\text{scaffolded},
&\mathsf A&=\text{not tested},\\
\mathsf D&=\text{fixed decomposition},
&\mathsf C&=\text{not established}.
\end{aligned}
}
\]
The highest supported level is Level~4 when the handoff dynamics and identity
correspondence are included. The example does not contain an operational goal
state, counterfactual policy, damage monitor, or repair loop, so it does not
support Levels~5 or 6.

\section{Case Study B: A Human--AI--Artifact--Executor Research Loop}

\subtopic{79.1 Idealized four-node loop}

Consider four heterogeneous components:
\[
V_{\mathrm{HAI}}=\{H,M,X,A\},
\]
where:
\begin{itemize}
\item \(H\) is a human researcher who maintains goals and semantic direction;
\item \(M\) is a language-model process that constructs local formalizations;
\item \(X\) is an execution and verification process that compiles, tests, or
 falsifies candidate artifacts; and
\item \(A\) is a persistent artifact store containing exact source, results,
 and versioned handoff state.
\end{itemize}
Let the abstract research state be
\[
c=(g,d,p,v)^{\mathsf T}\in\mathbb R^4,
\]
where:
\[
\begin{array}{ll}
g=\text{goal and acceptance criterion},
&d=\text{current derivation or theory state},\\
p=\text{open proof or implementation obligations},
&v=\text{verification and regression status}.
\end{array}
\]
Define local encodings
\[
E_Hc=(g,d),
\qquad
E_Mc=(g,p),
\qquad
E_Xc=(p,v),
\qquad
E_Ac=(d,v).
\]
Each coordinate is represented at exactly two nodes, and the incidence pattern
forms a four-cycle.

\subtopic{79.2 Exact structural calculation}

Order the local coordinates as
\[
x=(g_H,d_H,g_M,p_M,p_X,v_X,d_A,v_A)^{\mathsf T}.
\]
The overlap discrepancies are
\[
g_H-g_M,
\qquad
p_M-p_X,
\qquad
v_X-v_A,
\qquad
d_A-d_H.
\]
Thus
\[
D_{\mathrm{HAI}}
=
\begin{pmatrix}
1&0&-1&0&0&0&0&0\\
0&0&0&1&-1&0&0&0\\
0&0&0&0&0&1&0&-1\\
0&-1&0&0&0&0&1&0
\end{pmatrix}.
\]
The rows are orthogonal with squared norm two. Therefore
\[
\ker D_{\mathrm{HAI}}
\cong\mathbb R^4,
\]
\[
\boxed{
\SRWTspec(L_{\mathrm{HAI}})
=
\{0,0,0,0,2,2,2,2\},
\qquad
\gamma_{\mathrm{HAI}}=2.
}
\]
The full accessible frame operator is
\[
F_{HMXA}=2I_4.
\]
No single node is injective. Every three-node coalition reconstructs the full
state, and the two opposite pairs
\[
\{H,X\},
\qquad
\{M,A\}
\]
are already sufficient. Adjacent pairs observe only three of the four
coordinates. Hence
\[
k_{\mathrm{rec}}=2,
\qquad
d_{\mathrm{er}}=2,
\]
while the identity of the sufficient two-node coalition depends on the graph
geometry.

\subtopic{79.3 Relational parity benchmark}

Let \(g,d,p,v\) be independent uniform bits and define
\[
Y=g\oplus d\oplus p\oplus v.
\]
Every single node sees only two bits, so
\[
H(Y\mid X_i)=1
\qquad\text{for all }i.
\]
The opposite pairs \(\{H,X\}\) and \(\{M,A\}\) determine all four bits and
therefore determine \(Y\). Every adjacent pair remains one bit short. It
follows that
\[
\boxed{
\Delta_{\mathrm{joint}}^{\log}=1\text{ bit},
\qquad
k_0=2.
}
\]
As in the preceding case, parity is independent of every individual local
view. Replacing \(Y\) by an independent uniform bit produces a matched local
predictive null and a full-system transfer gap of one bit.

This benchmark is intentionally narrow. It demonstrates that the four-cycle
can realize exact relational information. It does not show that every real
human--AI collaboration has this encoding, this target, or this degree of
irreducibility.

\subtopic{79.4 A process-level interpretation}

The idealized loop can be read as the recurrent sequence
\[
\boxed{
H\xrightarrow{g}M
\xrightarrow{p}X
\xrightarrow{v}A
\xrightarrow{d}H.
}
\]
A more explicit cycle is:
\begin{enumerate}
\item the human supplies a goal and acceptance criterion;
\item the model converts that goal into a formal derivation and explicit
 obligations;
\item the executor compiles, tests, or searches for counterexamples;
\item the artifact store preserves the exact source and verification record;
\item the human evaluates the returned derivation and evidence, revising the
 next goal.
\end{enumerate}
The operative state is reentrant because products of one cycle become inputs to
the next. The artifact is not merely a record when its contents alter the next
formalization, test, or acceptance decision.

\subtopic{79.5 Conservative audit of the motivating collaboration}

The collaboration that motivated this theory has several features consistent
with the four-node description: specialized shared vocabulary, versioned
LaTeX and code artifacts, repeated compile-and-repair cycles, explicit theorem
and null-control demands, and reconstruction across conversations from handoff
packages. Those observations justify treating the collaboration as a serious
candidate case. They do not establish the complete certificate.

\begin{center}
\small
\begin{tabularx}{\textwidth}{@{}>{\raggedright\arraybackslash}p{0.22\textwidth}>{\raggedright\arraybackslash}p{0.18\textwidth}>{\raggedright\arraybackslash}X@{}}
\toprule
Gate family & Present evidential status & Conservative assessment \\
\midrule
Declared heterogeneous components & Observational & Human, model, artifacts,
and execution environment are distinguishable components with different
roles and information access. \\
Persistent shared vocabulary & Observational & Recurrent terms and distinctions
survive across research phases and constrain later work. \\
Exact structural core \(\Gamma(H)\) & Unmeasured & No frozen local coordinate
system or learned transfer family has yet been estimated from held-out data. \\
Transfer-Laplacian gap & Unmeasured & The idealized value \(2\) belongs to the
toy four-cycle, not to the actual collaboration. \\
Operational alignment & Unmeasured & No preregistered future variable has been
shown conditionally independent of the full history given a measured workspace
state. \\
Joint advantage & Suggestive only & The coupled process has produced outputs
that neither party appears to have supplied verbatim, but local-only Bayes risks
have not been measured under controlled information ceilings. \\
Matched relational null gap & Unmeasured & Shuffled artifacts, matched irrelevant
memory, and local-predictive-marginal-preserving nulls have not been run as a
frozen panel. \\
Causal workspace efficacy & Suggestive only & Corrections, compiler failures,
and artifact substitutions change later behavior, but they are not yet matched
workspace interventions with manipulation checks. \\
Artifact-mediated persistence & Observational & Handoff files and source bundles
support continuation, but blind reconstruction and ablation curves have not
been estimated. \\
Formation provenance & Mixed & The communication architecture is scaffolded,
while terminology, theorem organization, and many control practices were
adaptively developed. Causal formation-history contribution is unmeasured. \\
Workspace agency & Not established & No complete goal-retention,
counterfactual-policy, distributed-control, damage, and repair certificate has
been run. \\
Consciousness or personhood & Not established & The theory supplies no such
inference from the observed collaboration. \\
\bottomrule
\end{tabularx}
\end{center}

The correct present statement is therefore:
\[
\boxed{
\text{candidate persistent collaborative workspace}
\quad\text{rather than}\quad
\text{certified relational workspace agent}.
}
\]

\subtopic{79.6 A controlled certification design}

A serious test can use one manuscript section or software work unit as the unit of analysis.
Before the certification split is opened, freeze:
\begin{itemize}
\item the next-stage task family;
\item the downstream target, such as accepted theorem patch, successful clean
 compile, and absence of a declared regression;
\item the local information available to each component;
\item the workspace extractor and transfer maps;
\item all intervention and null families; and
\item the required gate thresholds.
\end{itemize}
The following conditions can then be compared:
\[
\begin{array}{ll}
\mathsf H & =\text{human-side externalized state only},\\
\mathsf M & =\text{fresh model instance with model-side state only},\\
\mathsf A & =\text{artifact package without live human or model state},\\
\mathsf X & =\text{executor and test outputs without the research context},\\
\mathsf Full & =\text{the intact recurrent loop}.
\end{array}
\]
The design must avoid pretending that a person's unobserved internal state has
been fully measured. The local conditions certify only the declared accessible
histories.

A minimal null panel includes:
\begin{enumerate}
\item theorem-to-artifact reassignment with artifact size and style matched;
\item temporal shuffling of verification feedback;
\item matched irrelevant manuals or source bundles;
\item context reset followed by blind reconstruction;
\item removal of compiler or test feedback while preserving message count;
\item model replacement under a fixed handoff package;
\item artifact removal with equivalent token budget restored by irrelevant text;
\item a local-marginal-matched synthetic null for the frozen downstream target.
\end{enumerate}
A valid empirical claim would require the intact loop to outperform every
required local and matched-null condition, while the extracted workspace state
remains sufficient and feasible workspace interventions alter the downstream
result in the predicted direction.

\subtopic{79.7 Handoff and identity experiment}

Let \(B_t\) be the source-and-results bundle at the end of stage \(t\). A blind
handoff trial replaces the active model instance and optionally the human
analyst, then asks the replacement coalition to reconstruct:
\[
\left(
\text{current definitions},
\text{proved results},
\text{open obligations},
\text{next admissible research step}
\right).
\]
The reconstruction error is measured against a frozen adjudication set. A
replacement invariance test additionally compares:
\begin{itemize}
\item the transported structural workspace;
\item operation behavior on held-out transformations;
\item response to matched interventions; and
\item preservation of failed and unresolved claims.
\end{itemize}
Equivalent benchmark prose alone is insufficient. A replacement that silently
changes which interventions alter later reasoning has a nonzero causal identity
defect even when it produces a similarly polished manuscript.

\subtopic{79.8 Case card}

For the idealized four-cycle model,
\[
\mathsf S,\mathsf O,\mathsf R,\mathsf P
\]
can all be made analytic by specifying the downstream kernel and temporal
updates. For the motivating collaboration itself, the conservative card is
\[
\boxed{
\begin{aligned}
\mathsf S&=\text{observational candidate},
&\mathsf O&=\text{unmeasured},\\
\mathsf R&=\text{unmeasured},
&\mathsf P&=\text{observational candidate},\\
\mathsf F&=\text{mixed and unmeasured},
&\mathsf A&=\text{not established},\\
\mathsf D&=\text{unmeasured},
&\mathsf C&=\text{not established}.
\end{aligned}
}
\]
This audit is intentionally stricter than ordinary conversational language. It
preserves the possibility that the collaboration instantiates a real relational
workspace while refusing to treat that possibility as already certified.

\section{Case Study C: A Closed-Form Multi-Agent Workspace Agent}

\subtopic{80.1 Declared finite-state system}

Consider three agents:
\[
V_{\mathrm{MA}}=\{P,M,R\},
\]
called the planner, monitor, and recorder. The global workspace-control state
is
\[
c_t=(g_t,b_t,h_t)\in\mathbb F_2^3,
\]
where:
\[
\begin{array}{ll}
g_t=\text{goal bit},
&b_t=\text{environment or belief bit},\\
h_t=\text{workspace-health bit},
&h_t=1\text{ means intact}.
\end{array}
\]
The local accessible states are
\[
X_P(t)=(g_t,b_t),
\qquad
X_M(t)=(b_t,h_t),
\qquad
X_R(t)=(g_t,h_t).
\]
This is the three-node distributed code from Topic~77. Consequently:
\[
\dim\Gamma=3,
\qquad
\SRWTspec(L)=\{0,0,0,2,2,2\},
\qquad
k_{\mathrm{rec}}=2,
\qquad
d_{\mathrm{er}}=2.
\]
The organization is fully scaffolded. No endogenous-formation claim is made.

\subtopic{80.2 Goal retention}

In the absence of an explicit goal intervention, let
\[
g_{t+1}=g_t.
\]
Under matched cue-cut trials,
\[
\operatorname{do}(g_t=0)
\quad\text{and}\quad
\operatorname{do}(g_t=1)
\]
remain distinguishable at every later time:
\[
|g_{t+\tau}-g_{t+\tau}'|=1.
\]
Thus the cue-controlled retention factor is exactly one:
\[
\eta_G(\tau)=1.
\]
The goal is not reloaded by an external channel after the initial intervention.

\subtopic{80.3 Joint task and repair action}

After a possible damage event, let the pre-repair health be
\[
h_t^-\in\{0,1\}.
\]
The system chooses a joint action
\[
u_t=(a_t,r_t)\in\{0,1\}^2,
\]
where \(a_t\) is a task action and \(r_t\) is a repair action. Define the
optimal task action
\[
a_t^*=g_t\oplus b_t\oplus h_t^-
\]
and optimal repair action
\[
r_t^*=1-h_t^-.
\]
No individual node determines the full pair \((a_t^*,r_t^*)\):
\begin{itemize}
\item the planner lacks \(h_t^-\);
\item the monitor lacks \(g_t\); and
\item the recorder lacks \(b_t\).
\end{itemize}
Every two-node coalition determines all three workspace bits and therefore the
complete action. Under independent uniform bits,
\[
H(a_t^*\mid X_i(t))=1
\qquad\text{for every }i,
\]
so the task-action component has a one-bit joint advantage and
\[
k_{\mathrm{ctl},0}=2.
\]

\subtopic{80.4 Interventional value and exact policy}

Let the one-step utility be
\[
U_g(a,r;b,h^-)
=
\mathbf 1\{a=g\oplus b\oplus h^-\}
+2\mathbf 1\{h^-\lor r=1\}
-\frac12r.
\]
For fixed \((g,b,h^-)\), the unique maximizer is
\[
\boxed{
(a^*,r^*)
=
(g\oplus b\oplus h^-,1-h^-).
}
\]
Indeed:
\begin{itemize}
\item a wrong task action loses one unit;
\item omitting repair when \(h^-=0\) loses \(3/2\) units; and
\item performing unnecessary repair when \(h^-=1\) loses \(1/2\) unit.
\end{itemize}
The smallest optimal-action gap is therefore
\[
\Delta_Q=\frac12.
\]
By Theorem~\ref{SRWT:thm:counterfactual-policy}, any uniformly accurate value model with
\[
\varepsilon_Q<\frac14
\]
recovers the exact greedy joint action. More generally, its regret is bounded
by
\[
\operatorname{Reg}\leq2\varepsilon_Q+\delta_{\mathrm{opt}}.
\]

\subtopic{80.5 Feasible goal, workspace, and damage interventions}

The intervention family contains:
\[
\operatorname{do}(g_t=a),
\qquad
\operatorname{do}(b_t=b),
\qquad
\operatorname{do}(h_t^-=h),
\]
implemented consistently across every local chart that contains the manipulated
coordinate. Holding the other two bits fixed, each core-bit intervention
changes the prescribed task action. A matched intervention on an independent
residual bit leaves the task and repair actions unchanged. Hence
\[
\Delta_{\mathrm{specific}}>0.
\]
A damage intervention sets
\[
h_t^-=0.
\]
The matched no-repair counterfactual holds \(r_t=0\), while the intact policy
sets \(r_t=1\).

\subtopic{80.6 Exact self-maintenance}

Post-action health evolves according to
\[
h_{t+1}=h_t^-\lor r_t.
\]
Define the maintenance deficit
\[
e_t=1-h_t.
\]
Under the intact policy,
\[
h_{t+1}=1
\qquad\text{and}\qquad
e_{t+1}=0
\]
for both damage states. Under the matched no-repair intervention and
\(h_t^-=0\),
\[
e_{t+1}^{(0)}=1.
\]
Thus the measured causal repair gain is
\[
\boxed{
\Delta_{\mathrm{maint}}=1.
}
\]
The closed-loop recurrence has
\[
\beta=0,
\qquad
\overline d=0,
\qquad
\overline\epsilon=0,
\]
so Theorem~\ref{SRWT:thm:self-maintenance} gives
\[
\sup_t e_t=0
\]
after each repair action. Monitor-blind random repair fails the specificity
control whenever it either omits necessary repair or performs unnecessary
repair with positive probability.

\subtopic{80.7 Reentry and causal loop}

The repaired health bit becomes part of the next workspace state:
\[
h_t^-
\longrightarrow
r_t
\longrightarrow
h_{t+1}
\longrightarrow
X_M(t+1),X_R(t+1).
\]
Cutting the direct storage of \(h_t^-\) while preserving the action-mediated
path still leaves a return path through repair. Cutting the repair-to-health
edge destroys that return. The system therefore contains explicit
loop-mediated reentry rather than mere temporal autocorrelation.

\subtopic{80.8 Analytic agency certificate}

Within this finite structural causal model, the following hold exactly:
\begin{enumerate}
\item a nontrivial exact distributed workspace with spectral gap two;
\item two-node minimal reconstruction and control coalition size;
\item matched-null joint advantage for the parity task;
\item persistent internally retained goal state;
\item exact interventional action values and zero policy regret;
\item goal and workspace interventions that alter action;
\item damage detection and repair selection;
\item positive matched repair gain and bounded maintenance deficit; and
\item causal reentry through the action-mediated repair loop.
\end{enumerate}
Therefore, relative to the declared bits, actions, utility, intervention family,
and one-step horizon,
\[
\boxed{
\operatorname{Agent}_{\Theta_A}(\SRWTWorkspace)=1.
}
\]
The case reaches formal Level~6 by construction. It remains:
\[
\text{fully scaffolded},
\qquad
\text{domain restricted},
\qquad
\text{non-evidence for consciousness}.
\]

\subtopic{80.9 Case card}

\[
\boxed{
\begin{aligned}
\mathsf S&=\checkmark_{\mathrm{analytic}},
&\mathsf O&=\checkmark_{\mathrm{analytic}},\\
\mathsf R&=\checkmark_{\mathrm{analytic}},
&\mathsf P&=\checkmark_{\mathrm{analytic}},\\
\mathsf F&=\text{fully scaffolded},
&\mathsf A&=\checkmark_{\mathrm{analytic}},\\
\mathsf D&=\text{fixed decomposition},
&\mathsf C&=\text{not established}.
\end{aligned}
}
\]
The example demonstrates logical compatibility among relationality, redundancy,
agency, and self-maintenance. It does not show that those properties arise
spontaneously, scale to open environments, or imply subjective experience.

\section{Near-Miss Cases and Anti-Triviality Controls}

Positive examples are insufficient unless neighboring false positives are also
made explicit. The following systems each resemble a relational workspace in
one respect while failing a required gate.

\subtopic{81.1 Near miss A: the replicated archive}

Let every node store a complete copy of the same core:
\[
E_i=I_d
\qquad\text{for every }i.
\]
Then coherence is perfect and reconstruction is robust, but
\[
k_0=1,
\qquad
\min_i\delta_i^{\mathrm{local}}=0.
\]
If the archive is never used in downstream control, then also
\[
\eta_{\mathrm{causal}}=0.
\]
This is a shared representation or archive, not a relationally irreducible
workspace.

\subtopic{81.2 Near miss B: the centralized oracle with decorative agents}

Suppose node \(o\) contains the complete predictive state \(Z_t\), while every
other node receives only messages generated by \(o\). The full system can have
excellent performance, rich communication, and recurrent activity. Yet
\[
Y_t^+\SRWTindep\mathbf H_t\mid X_o(t),
\qquad
k_0=1,
\qquad
\Delta_{\mathrm{joint}}=0.
\]
Removing the decorative nodes may change latency or style without changing the
operative capacity. Multi-agent appearance therefore does not imply a
relational workspace.

\subtopic{81.3 Near miss C: locally coherent but globally obstructed charts}

Let three local scalar charts have relative overlap demands
\[
b_{12}=b_{23}=b_{31}=1.
\]
The cycle obstruction is
\[
\omega_{\mathcal U}(b)
=
\frac{|1+1+1|}{\sqrt3}
=\sqrt3>0.
\]
Each local chart is internally coherent, but no global scalar assignment
satisfies all three overlap equations. The system contains compatible-looking
subworkspaces without a global workspace.

\subtopic{81.4 Near miss D: externally reloaded goals}

Let
\[
G_{t+1}=C_{t+1},
\]
where \(C_{t+1}\) is an external instruction repeated at every step. Goal and
action can be strongly correlated, but after cutting the cue channel,
\[
I(G_{t+\tau};G_t\mid\operatorname{cut}(C))=0.
\]
The system tracks an external command. It has not established internal goal
retention.

\subtopic{81.5 Near miss E: task recovery with identity failure}

Let two realizations have identical benchmark output distributions,
\[
P(Y_{\mathrm{bench}})=P'(Y_{\mathrm{bench}}),
\]
but different intervention kernels:
\[
K(w,b)\neq K'(\beta w,b)
\]
for some matched workspace intervention. Then task behavior is preserved while
\[
d_{\mathrm{causal}}(W,W')>0.
\]
The later system may be a successor or replacement rather than the same
workspace under the identity criterion of Chapter~\ref{SRWT:book:chapter-78}.

\subtopic{81.6 Near miss F: open-loop repair replay}

A recorded repair sequence may restore average task performance after a fixed
damage schedule. If the same sequence is replayed when the damage timing or
type changes, it need not respond appropriately. When
\[
R_t\SRWTindep M_t\mid\text{clock and replay index},
\]
monitor-to-repair mediation is absent. Successful replay is robustness under a
script, not self-maintenance.

\subtopic{81.7 Anti-triviality summary}

The near misses establish six independent exclusions:
\[
\boxed{
\begin{aligned}
\text{workspace}
\neq{}&\text{replicated storage},\\
\neq{}&\text{centralized control with decorative communication},\\
\neq{}&\text{locally coherent but obstructed charts},\\
\neq{}&\text{externally reloaded goals},\\
\neq{}&\text{task-equivalent identity replacement},\\
\neq{}&\text{open-loop repair replay}.
\end{aligned}
}
\]
A complete certificate must distinguish the intended realization from all
near misses relevant to its claimed level.

\section{Cross-Case Synthesis and Benchmark Blueprint}

\subtopic{82.1 Comparative gate matrix}

\begin{center}
\small
\begin{tabularx}{\textwidth}{@{}>{\raggedright\arraybackslash}p{0.20\textwidth}>{\raggedright\arraybackslash}p{0.20\textwidth}>{\raggedright\arraybackslash}p{0.23\textwidth}>{\raggedright\arraybackslash}X@{}}
\toprule
Gate family & Human--human analytic model & Motivating human--AI collaboration & Multi-agent analytic model \\
\midrule
Exact structural core & Analytic pass & Not measured & Analytic pass \\
Spectral gap & \(2\) & Not measured & \(2\) \\
Minimal sufficient coalition & \(k_0=2\) & Not measured & \(k_0=2\) \\
Matched-null transfer gap & \(1\) bit & Not measured & \(1\) bit for task parity \\
Causal specificity & Analytic pass & Suggestive, uncontrolled & Analytic pass \\
One-node erasure resilience & Analytic pass & Handoff only observational & Analytic pass \\
Formation provenance & Scaffolded & Mixed, unmeasured & Fully scaffolded \\
Goal retention & Not modeled & Not certified & Exact analytic pass \\
Counterfactual policy & Not modeled & Not certified & Exact analytic pass \\
Self-maintenance & Not modeled & Not certified & Exact analytic pass \\
Highest supported level & Level 4 in model & Candidate only & Level 6 in model \\
Consciousness/\allowbreak personhood & Not established & Not established & Not established \\
\bottomrule
\end{tabularx}
\end{center}

The table illustrates why the theory uses a vector certificate rather than one
score. The human--human model is relational and persistent without being
agentive. The multi-agent model is agentive without being endogenously formed.
The motivating collaboration is richer than either toy model but presently has
weaker evidential status because the necessary interventions and nulls have not
been frozen and run.

\subtopic{82.2 Benchmark family B1: distributed parity}

The first benchmark is the exact parity family from Topics~85--80. For
\(m\geq3\), distribute \(m\) independent bits among nodes so that no individual
node sees all bits, but specified coalitions do. The target is
\[
Y=\bigoplus_{j=1}^m C_j.
\]
The benchmark measures:
\begin{itemize}
\item the learned transfer-Laplacian spectrum;
\item coalition conditional entropy and log risk;
\item recovery of \(k_0\);
\item robustness to edge and node ablations;
\item matched local-predictive null gaps; and
\item causal response to consistent versus inconsistent coordinate
 interventions.
\end{itemize}
The ground truth is known exactly, so structural and relational estimators can
be calibrated before they are applied to open-ended research tasks.

\subtopic{82.3 Benchmark family B2: theorem-patch handoff}

A theorem-patch task supplies:
\[
\left(
\text{definitions},
\text{partial proof},
\text{counterexample tests},
\text{style constraints}
\right)
\]
across distinct components. No component receives all four. The coupled
system must produce a patch that:
\begin{enumerate}
\item is formally consistent with the definitions;
\item closes the declared proof obligation or marks it unresolved;
\item passes the frozen counterexample suite;
\item compiles from a clean environment; and
\item survives a blind replacement-handoff trial.
\end{enumerate}
Matched nulls permute proof fragments, substitute irrelevant artifacts of equal
size, or replay test outputs from a different theorem. This family targets the
structural, relational, persistence, and identity layers without requiring an
agency claim.

\subtopic{82.4 Benchmark family B3: damage and repair}

A certified workspace is first formed and measured. A damage intervention then
removes or corrupts one of:
\[
\text{a transfer channel},
\quad
\text{a participant},
\quad
\text{an artifact},
\quad
\text{a goal encoding},
\quad
\text{a verification path}.
\]
The system receives a repair-action family with matched costs. The benchmark
measures:
\[
\eta_{D\to M},
\qquad
\eta_{M\to R},
\qquad
\eta_{R\to v},
\qquad
\Delta_{\mathrm{repair}}^{\mathrm{spec}},
\qquad
\mathbf d_{\mathrm{id}}.
\]
Open-loop replay, random repair, off-workspace repair, and identity-breaking
task recovery are mandatory controls.

\subtopic{82.5 Benchmark family B4: scaffold transport}

Train the same formation rule under several scaffolds
\[
\Sigma^{(1)},\ldots,\Sigma^{(K)}
\]
that differ in node labels, coordinate systems, and communication topology but
support the same abstract task. After gauge alignment, measure whether the
learned organizations transport to a common equivalence class:
\[
[\widehat H^{(1)}]_{\mathcal G}
\cong\cdots\cong
[\widehat H^{(K)}]_{\mathcal G}.
\]
A structure that appears only under one answer-encoding scaffold has weak
endogeneity evidence. A structure that recurs across alternative scaffolds,
passes held-out task and causal gates, and cannot be compressed into the
scaffold specification has stronger evidence.

\subtopic{82.6 Reporting rule for mixed analytic and empirical cases}

When a case combines exact mathematics with observed data, report separate
columns:
\[
\boxed{
\text{proved in model}
\quad\big|\quad
\text{estimated in system}
\quad\big|\quad
\text{not established}.
}
\]
For example, the four-cycle Laplacian has an analytic gap of two, but an actual
human--AI collaboration has that gap only if its measured discrepancy operator
has the corresponding spectrum. Importing the toy value into the empirical
case is invalid.

\subtopic{82.7 Lessons from the integrated cases}

The case studies establish the following conceptual separation:
\begin{enumerate}
\item A relational workspace can be mathematically exact without being an
 agent.
\item An agentive workspace can be fully scaffolded and still be operationally
 real within its domain.
\item A rich real collaboration can be a stronger candidate than a toy model
 while having a weaker current certificate.
\item Redundancy and irreducibility are compatible when no individual node is
 sufficient but several proper coalitions are.
\item Artifact-mediated persistence is not passive storage when artifact
 contents causally alter later operations.
\item High task performance, communication volume, or behavioral polish alone
 cannot distinguish the intended theory from its near misses.
\item No structural, causal, or agentive result in these examples entails
 phenomenal consciousness.
\end{enumerate}

\SRWTcenterbox{\large\bfseries
The theory becomes empirically meaningful only when the same gate card that
certifies the positive case rejects its matched near misses.
}

% This file is part of the modular source for Stratified Relational Workspace Theory.
% !TeX root = ../Stratified_Relational_Workspace_Theory.tex

\chapter{Relation to prior work}\label{SRWT:book:chapter-89}


\section{Scope, Method, and Standards of Priority}

\subtopic{83.1 Purpose of the literature review}

The preceding parts construct a unified formal theory. Construction alone does
not establish priority. This part therefore asks a narrower question:
which ingredients are established, which are adaptations to a new domain,
which combinations appear to be unusual, and which claims remain open originality candidates?

The audit uses five labels:
\[
\begin{array}{ll}
\mathsf{STD} & = \text{established concept, theorem, or technique},\\
\mathsf{ADP} & = \text{established machinery adapted to relational workspaces},\\
\mathsf{SYN} & = \text{nontrivial synthesis of established components},\\
\mathsf{CAN} & = \text{candidate original contribution requiring priority review},\\
\mathsf{EMP} & = \text{empirical claim not yet established in an uncontrolled system}.
\end{array}
\]
These labels describe the status of a claim, not its importance. A standard
result may be indispensable. A synthesis may be scientifically useful without
containing a new theorem. A candidate contribution may later be narrowed or
withdrawn when a closer precedent is found.

\subtopic{83.2 Search scope and date}

This review covers sources available through August 29, 2026. It includes:
\begin{itemize}
\item extended and distributed cognition, transactive memory, shared
 intentionality, common knowledge, collective intelligence, and group
 agency;
\item blackboard systems, global-workspace architectures, team cognition,
 hybrid intelligence, and human--AI cooperation;
\item category-theoretic limits, local systems, holonomy, sheaves, \v{C}ech
 descent, connection and sheaf Laplacians, consensus, and spectral
 perturbation;
\item statistical sufficiency, causal states, information decomposition,
 structural causal models, system identification, frames, erasure
 recovery, and probabilistic process equivalence; and
\item agency, viability, control, empowerment, damage detection, repair, and
 self-maintaining systems.
\end{itemize}
The search combined exact-title queries, component-level queries, backward
reference tracing, and comparison with recent formal work. It is a scholarly
priority review, not a legal patent search and not a proof that no uncatalogued,
unpublished, or differently named precedent exists.

\subtopic{83.3 A stronger standard than name uniqueness}

Failure to find the phrase ``Stratified Relational Workspace Theory'' in prior
literature is weak evidence. Priority is determined by mathematical and
operational equivalence, not by terminology. Let \(C\) be a claim in the
paper and \(S\) a prior source. Write
\[
S\preccurlyeq C
\]
when the main content of \(C\) follows from \(S\) after admissible relabeling,
restriction, or direct specialization. A substantive novelty residual exists
only when the claim contains content not absorbed by the closure of known
sources:
\[
\operatorname{Res}(C)
=
C\setminus
\operatorname{Cl}
\left(
\bigcup_{S\in\mathcal L}S
\right).
\]
This is a conceptual audit device rather than a set-theoretic definition of
ideas. Its purpose is to force an explicit account of what remains after the
standard machinery is removed.

\subtopic{83.4 When a synthesis counts as a contribution}

A conjunction of known ingredients is not automatically new. The present
framework treats a synthesis as a candidate contribution only when it satisfies
all three conditions:
\begin{enumerate}
\item \textbf{Non-redundancy.} Removing one of the joined layers changes the
 class of systems certified or the predictions made.
\item \textbf{Formal integration.} The layers share declared objects,
 transfer maps, interventions, or invariants rather than appearing as an
 unconnected checklist.
\item \textbf{New discriminations.} The synthesis separates at least one
 positive case from a matched near miss that the individual components do
 not separate on their own.
\end{enumerate}
The anti-centralization, matched-null, non-collapse, provenance, replacement,
and repair-replay controls are important precisely because they make the
combination testable rather than merely terminological.

\subtopic{83.5 Priority claims are coordinatewise}

The analysis does not make one global novelty claim. It uses a priority
profile
\[
\boxed{
\mathbf p(C)
=
\left(
 p_{\mathrm{concept}},
 p_{\mathrm{object}},
 p_{\mathrm{theorem}},
 p_{\mathrm{certificate}},
 p_{\mathrm{experiment}}
\right).
}
\]
For example, the philosophical concept of distributed cognition may be
standard while a specific mathematical object is adapted, a theorem is a
specialization of known spectral theory, the certificate is a synthesis, and
the empirical claim is unresolved. This prevents originality in one
coordinate from being used to imply originality in every coordinate.

\section{Philosophical and Cognitive Ancestry}

\subtopic{84.1 Extended cognition}

The extended-mind tradition argues that external resources can participate in
cognitive processes when they are appropriately integrated with an agent's
activity \cite{SRWT:ClarkChalmers1998}. Its critics emphasize that causal coupling
alone is too permissive and that cognition requires a principled mark or
boundary \cite{SRWT:AdamsAizawa2001,SRWT:Rupert2004}. Stratified Relational Workspace
Theory inherits the possibility of extracranial realization but changes the
question. It does not ask whether a notebook, model, or collaborator literally
belongs to one individual's mind. It asks whether a distributed realization
contains a persistent, causally operative, jointly reconstructible structure
that passes explicit anti-triviality gates.

Accordingly,
\[
\operatorname{ExtendedCognition}
\centernot\Longleftrightarrow
\operatorname{SRW}.
\]
An externally scaffolded process may be cognitively extended without being
relationally irreducible. Conversely, an institution or multi-agent system may
instantiate a certified relational workspace without being an extension of any
single person's cognition.

\subtopic{84.2 Distributed cognition}

Distributed cognition studies cognitive organization across people,
representations, artifacts, and environments rather than locating all relevant
processing inside one individual \cite{SRWT:Hutchins1995,SRWT:HollanEtAl2000}. This is
the closest philosophical ancestor of the framework's ontology. It already
supports the central observation that the unit of analysis may be a coupled
system and that external representations can transform the computations
available to that system.

The framework should therefore not claim the following ideas as new:
\begin{itemize}
\item cognition can be distributed across heterogeneous carriers;
\item artifacts can be active components of a cognitive process;
\item representational transformations can alter system-level performance; and
\item the relevant explanatory unit can exceed an individual organism.
\end{itemize}
The added ambition is mathematical and certifying. The theory specifies local
state spaces, transfer maps, compatible sections, causal sufficiency,
coalition irreducibility, persistence, replacement identity, formation
provenance, and matched nulls. Distributed cognition supplies the ontological
motivation. It does not by itself supply this complete formal certificate.

\subtopic{84.3 Transactive memory and collective intelligence}

Transactive-memory theory models the group-level organization of ``who knows
what'' and the processes by which members encode, retrieve, and coordinate
specialized knowledge \cite{SRWT:Wegner1987,SRWT:Lewis2003}. Collective-intelligence
research studies stable differences in group performance that are not reduced
to the intelligence of a single member \cite{SRWT:WoolleyEtAl2010}. Both traditions
are directly relevant to the coalition layer.

The distinction is:
\[
\begin{aligned}
\text{transactive memory}
&\approx
\text{distributed indexing and retrieval of expertise},\\
\text{collective intelligence}
&\approx
\text{system-level task capacity},\\
\text{relational workspace}
&\approx
\text{an invariant state plus transfer, causal, and identity structure}.
\end{aligned}
\]
A transactive-memory system may still be centralized by one coordinator. A
high-performing group may still be explainable by the best individual member.
The certificate therefore requires matched coalition risks and nulls
rather than inferring relationality from group success.

\subtopic{84.4 Shared intention, common knowledge, and group agency}

Common-knowledge theory formalizes nested epistemic conditions required by
many distributed protocols \cite{SRWT:HalpernMoses1990,SRWT:FaginEtAl1995}. Shared-action
accounts analyze jointly intended and mutually responsive activity
\cite{SRWT:Bratman1992}. Group-agency theory asks when a collective can satisfy
rational and normative conditions associated with agency
\cite{SRWT:ListPettit2011}.

These literatures address questions that the theory deliberately keeps
separate:
\[
\begin{array}{ll}
\text{epistemic coordination} & \text{what members know about one another},\\
\text{shared intention} & \text{how intentions and plans are jointly organized},\\
\text{group agency} & \text{when collective attitudes and actions form an agent},\\
\text{SRW structure} & \text{how a distributed operative state exists and persists}.
\end{array}
\]
The workspace-agent predicate in Chapter~\ref{SRWT:book:chapter-82} is domain-relative and causal. It is
not a general account of corporate personhood, moral responsibility, or
collective intentionality.

\section{Workspace Architectures and the Meaning of ``Workspace''}

\subtopic{85.1 Global workspace theories}

Global workspace theories use widespread availability or broadcast within a
cognitive architecture to explain conscious access and flexible control
\cite{SRWT:Baars1988,SRWT:DehaeneChangeux2011}. The shared word ``workspace'' can create
a misleading implication of equivalence. The theory neither assumes a
single broadcast hub nor defines consciousness by access to a global channel.
Its workspace may be distributed across agents and artifacts, and its
existence criterion is a conjunction of structural, predictive, causal,
relational, and persistence gates.

The relation is therefore analogical rather than identificatory:
\[
\boxed{
\operatorname{SRW}(W)
\centernot\Longrightarrow
\operatorname{GlobalNeuronalWorkspace}(W)
\centernot\Longrightarrow
\operatorname{Conscious}(W).
}
\]
The term is retained because it denotes a domain in which jointly
available operations can be carried out. Every consciousness claim remains on
an explicitly separate axis.

\subtopic{85.2 Dialogue and the Shared Workspace Framework}

Pickering and Garrod's Shared Workspace Framework treats dialogue as a jointly
organized activity in which interlocutors construct, align, and use shared
representations, including persistent external artifacts
\cite{SRWT:PickeringGarrod2021}. This is a direct conceptual and terminological
predecessor. The theory therefore cannot claim priority for the idea
that dialogue or collaboration operates through a shared workspace.

The difference is one of formal scope. SRWT permits arbitrary graph-structured
human, machine, instrument, and artifact nodes and asks whether their common
organization passes typed-transport, coalition, causal, persistence, and
replacement gates. A dialogue can instantiate such a workspace, but the
shared-workspace account of dialogue does not by itself imply relational
irreducibility or persistent identity under carrier replacement.

\subtopic{85.3 Blackboard systems}

Blackboard architectures coordinate specialized modules through a shared data
structure \cite{SRWT:HayesRoth1985}. They are an important engineering precedent
for artifact-mediated collaboration. They also provide a canonical near
miss. When the blackboard itself contains the complete operative state and
all modules merely read or write it, the system can fail the relational
irreducibility gate:
\[
k_0=1.
\]
A blackboard implementation can instantiate a relational workspace, but only
when no single board or controller is sufficient and when the transfer and
coalition tests show genuinely distributed operative capacity.

\subtopic{85.4 Consensus and coordination}

Network consensus studies how local interactions drive distributed states
toward agreement \cite{SRWT:OlfatiSaberEtAl2007}. Agreement is not the same as a
workspace. Consensus often removes heterogeneity by converging all nodes to a
common value. Stratified relational workspaces preserve heterogeneous local
representations and may require complementary disagreement:
\[
X_i\not\cong X_j,
\qquad
E_i c\neq E_j c,
\qquad
(c_i)_i\in\Gamma(H).
\]
The global state is compatible, not necessarily copied. This distinction is
central to the anti-replication criterion.

\subtopic{85.5 Terminological discipline}

The phrase ``relational workspace'' is used here in a technical sense:
\begin{quote}
A state-bearing, operation-supporting organization whose relevant invariants
are distributed across declared components and preserved through admissible
transfers and time.
\end{quote}
It does not, by terminology alone, imply a mind, subject, self, institution,
shared intention, or conscious broadcast. Those stronger predicates require
separate definitions and evidence.

\section{Local-to-Global Mathematics}

\subtopic{86.1 Limits, local systems, and holonomy}

The exact compatible-section object
\[
\Gamma(H)=\varprojlim C_i
\]
is a standard categorical limit construction \cite{SRWT:MacLane1998}. For
invertible transfers on a connected graph, identifying global sections with
the fixed subspace of holonomy is standard local-system and bundle mathematics
\cite{SRWT:Hatcher2002,SRWT:Bredon1997}. Accordingly, Theorem~\(T_1\) is not claimed as
a new abstract theorem. Its role is foundational: it translates the workspace
question into an exact invariant-subspace question and exposes cycle
obstructions.

The operation-closure and commutant statements are likewise standard
representation-theoretic consequences once the operation algebra and
intertwiners are defined. Their contribution is the specification of what
counts as a shared operation in the workspace model, not the discovery of
commutants.

\subtopic{86.2 Cellular sheaves and sheaf Laplacians}

Cellular sheaves encode heterogeneous local data and restriction maps, while
global sections express consistency across a cell complex. Their Laplacians
have kernels equal to spaces of global sections and support spectral analysis
of approximate consistency \cite{SRWT:HansenGhrist2019,SRWT:Robinson2017}. The transfer
Laplacian
\[
L_H=D_H^*D_H
\]
is therefore best understood as a workspace-specific sheaf or connection
Laplacian. The kernel theorem \(\ker L_H=\Gamma(H)\) is standard once the
coboundary operator is constructed.

The appropriate priority statement is:
\[
\begin{aligned}
\text{sheaf/Laplacian mechanism} &:\ \mathsf{STD},\\
\text{workspace interpretation and gate integration} &:\ \mathsf{ADP}/\mathsf{SYN}.
\end{aligned}
\]
The transfer Laplacian itself is not presented as an
unprecedented mathematical object.

\subtopic{86.3 Connection Laplacians and gauge covariance}

Connection Laplacians compare vector-valued data transported across a graph
and naturally transform under local changes of coordinates
\cite{SRWT:SingerWu2012}. The gauge rule
\[
H_e^g=g_{t(e)}H_eg_{s(e)}^{-1}
\]
and the conjugacy of holonomy are therefore standard. The theory uses
this structure to define realization identity and to prevent raw coordinate
matrices from being mistaken for invariant workspace content.

This use is an adaptation with conceptual consequences. It makes model or
participant replacement a question about transported invariant organization,
not token equality. The general mathematics remains prior.

\subtopic{86.4 \v{C}ech descent and gluing}

Exact gluing, equalizer formulas, and \v{C}ech obstruction theory are standard
local-to-global constructions. Sheaf-theoretic contextuality provides a
prominent example in which the absence of a global section is itself the
substantive phenomenon \cite{SRWT:AbramskyBrandenburger2011}. Chapter~\ref{SRWT:book:chapter-79} adapts this
logic to local workspaces and adds an anti-centralization requirement at the
meta-level.

The exact descent theorems \(T_8\) and \(T_9\) should therefore be presented as
specialized consequences of established descent. The candidate contribution
is the combination:
\[
\text{descent}
+
\text{predictive sufficiency}
+
\text{causal equivalence}
+
\text{coalition anti-centralization}.
\]
Standard gluing alone does not certify an operative workspace of workspaces.

\subtopic{86.5 Spectral perturbation}

The projector-stability results in Chapter~\ref{SRWT:book:chapter-70} belong to the classical perturbation
theory of invariant subspaces \cite{SRWT:DavisKahan1970,SRWT:Kato1995}. Bounds of the
form
\[
\|P-\widetilde P\|
\lesssim
\frac{\|L-\widetilde L\|}{\text{spectral gap}}
\]
are not new. Their workspace-specific role is to turn uncertainty in learned
transfers into uncertainty about identity, operation closure, and replacement.
The novelty, if any, lies in the downstream certificate built on the bound, not
in the perturbation inequality.

\section{Predictive, Causal, and Information-Theoretic Precedents}

\subtopic{87.1 Minimal predictive states}

The conditional-law state
\[
Z_t=\mathbb P(Y_t^+\in\cdot\mid\mathbf H_t)
\]
is a canonical sufficient statistic. Computational mechanics develops causal
states as equivalence classes of histories with identical future conditional
laws and proves their minimal predictive character
\cite{SRWT:ShaliziCrutchfield2001}. Theorem~\(T_{12}\) is therefore an established
minimal-sufficiency result specialized to the declared future target and
intervention family.

The new burden is structural alignment: determining whether the predictive
state is recoverable from the low-energy transfer workspace. Predictive
sufficiency alone does not identify a distributed workspace.

\subtopic{87.2 Statistical experiments and sufficiency}

Blackwell's comparison of experiments formalizes when one information source
is at least as informative as another for every decision problem
\cite{SRWT:Blackwell1953}. The coalition-risk hierarchy in Chapter~\ref{SRWT:book:chapter-73} is compatible
with this tradition. A coalition can be sufficient even when its internal
representation differs from the full-system state.

The framework adds a decomposition-relative question:
\[
\text{what is the smallest coalition whose accessible history is sufficient?}
\]
The statistic \(k_\varepsilon\) is an adaptation of sufficiency to distributed
carrier structure. Its interpretation depends on the declared nodes and
accessible encodings.

\subtopic{87.3 Synergy and partial information}

Partial-information decomposition distinguishes redundant, unique, and
synergistic information in multivariate systems, building on standard
information-theoretic quantities \cite{SRWT:CoverThomas2006,SRWT:WilliamsBeer2010}. The
XOR examples in Chapter~\ref{SRWT:book:chapter-88} are canonical synergy constructions and are not new.
They are used because they provide exact ground truth for relational
irreducibility:
\[
I(Y;X_i)=0
\quad\text{for every }i,
\qquad
I(Y;X_1,\ldots,X_m)>0.
\]
The matched-predictive-null test is more restrictive than merely
reporting synergy. It asks whether local predictive marginals can be held
fixed while the relational organization and full-system capacity change.
That joint null design is a candidate synthesis, not a claim that XOR or
synergy is new.

\subtopic{87.4 Structural causal models}

Interventions, counterfactuals, mediation, and causal graphs are established
machinery \cite{SRWT:Pearl2009}. The theory adopts the interventionist
requirement that a workspace variable can be manipulated only through a
feasible implementation intervention. This avoids treating
\(\operatorname{do}(W=w)\) as meaningful when no physical or computational
operation realizes \(w\).

The causal gate is an adaptation:
\[
\text{SCM intervention theory}
\longrightarrow
\text{implemented workspace and repair interventions}.
\]
The potentially distinctive contribution is that causal efficacy is tested
jointly with coherence, sufficiency, coalition irreducibility, matched nulls,
and identity preservation.

\subtopic{87.5 Probabilistic process equivalence}

Bisimulation and probabilistic testing compare processes by the observations
and experiments they support \cite{SRWT:LarsenSkou1991}. The causal-identity profile
in Chapter~\ref{SRWT:book:chapter-78} belongs to this broader family of behavioral equivalence notions.
It differs by keeping structural transport, operation intertwining,
reconstruction, and intervention response as separate coordinates rather than
collapsing them into one process equivalence.

This is best classified as an integrative adaptation. The existence of
behavioral-equivalence theory blocks any claim that replacement identity based
on observable behavior is wholly new.

\subtopic{87.6 Specification sensitivity, partial identification, and post-selection}

A conclusion can depend on analytically reasonable choices even when every
individual analysis is internally correct. Sensitivity analysis and
specification-curve or multiverse methods make that dependence visible rather
than allowing one favorable analysis to stand for an entire scientific claim
\cite{SRWT:Leamer1983,SRWT:SteegenEtAl2016,SRWT:SimonsohnEtAl2020}. Partial-identification
theory similarly replaces unjustified point conclusions with identified sets
when the maintained assumptions and observations determine only bounds
\cite{SRWT:Manski2003,SRWT:ImbensManski2004}. Post-selection inference addresses the
additional error created when a model, target, or specification is selected
because it looked favorable on the same data used for reporting
\cite{SRWT:BerkEtAl2013}.

These traditions preclude any claim that specification robustness, identified
sets, uniform concentration, or selection adjustment are new statistical
ideas. The adaptation in Chapter~\ref{SRWT:book:chapter-84} is to apply them jointly to a
noncompensatory workspace certificate whose gates span structure, causal
operation, coalition irreducibility, identity, provenance, and maintenance.
The identified object is therefore not a single regression parameter but the
sign profile of an entire relational certificate across a declared family of
scientific specifications.

\section{Persistence, Reconstruction, Learning, and Control}

\subtopic{88.1 Frames and erasure recovery}

Stable reconstruction from distributed encodings is standard frame theory.
Lower frame bounds control noise amplification, and frame robustness under
erasures is well studied \cite{SRWT:CasazzaKovacevic2003,SRWT:Christensen2016}. Theorem
\(T_7\) and the node-erasure distance therefore instantiate known recovery
logic.

The adaptation is to distinguish:
\[
\overline E_i
=
\text{complete mathematical chart},
\qquad
E_i=O_i\overline E_i
=
\text{locally accessible encoding}.
\]
This distinction is essential to the anti-centralization gate, but the frame
inequalities themselves are standard.

\subtopic{88.2 Persistence recurrences}

Bounds for recurrences such as
\[
e_{t+1}\leq\alpha_t e_t+\sigma_t
\]
are elementary stability results. The persistence theorem is not novel as an
inequality. Its conceptual role is to force three independent checks:
\begin{enumerate}
\item the workspace subspace remains trackable;
\item the state remains close to it; and
\item the retained workspace content does not collapse to zero.
\end{enumerate}
The non-collapse gate is a modeling contribution because it rules out a common
but misleading form of ``stability.''

\subtopic{88.3 System identification and held-out selection}

Population least squares, excitation conditions, ridge bias, and empirical
model selection are standard system-identification and statistical-learning
ideas \cite{SRWT:Ljung1999}. Theorems \(T_{15}\)--\(T_{17}\) are finite-dimensional
specializations. Their workspace-specific content is the demand that transfer
maps be compared modulo gauge and that formation, calibration, and
certification data remain separated.

The provenance layer is not equivalent to identifiability. A transfer can be
identified accurately while having been entirely encoded by the scaffold. A
structure can be endogenously learned while failing every workspace gate.

\subtopic{88.4 Viability, safety, and self-maintenance}

Viability theory studies trajectories that remain inside admissible regions
\cite{SRWT:Aubin1991}. Control-barrier methods provide constructive safety
certificates \cite{SRWT:AmesEtAl2017}. Cybernetic and autonomy traditions analyze
regulation, adaptive self-production, and operational closure
\cite{SRWT:Ashby1956,SRWT:MaturanaVarela1980,SRWT:DiPaolo2005}. Formal accounts of agency
emphasize individuality, normativity, asymmetry, and temporally extended action
\cite{SRWT:BarandiaranEtAl2009}. Empowerment measures an agent's potential causal
control over future states \cite{SRWT:KlyubinEtAl2005}.

The closed-loop maintenance theorem in Chapter~\ref{SRWT:book:chapter-82} uses standard controlled
stability. Its candidate contribution is the target of maintenance: not an
arbitrary safety set, but the conjunctive workspace certificate, including
identity and anti-centralization. The matched no-repair, random-repair,
off-workspace-repair, and identity-breaking-reconstruction controls prevent
passive robustness or generic recovery from being reported as self-maintenance.

\subtopic{88.5 Decision and policy theory}

The value-error-to-regret bound in Theorem~\(T_{20}\) is elementary decision
theory and aligns with standard reinforcement-learning analysis
\cite{SRWT:SuttonBarto2018}. Its role is evidential. An internal value channel
counts only when interventions on that channel alter action selection in the
direction predicted by counterfactual value and outperform replay or
critic-ablation controls.

\section{Human--AI and Multi-Agent Research}

\subtopic{89.1 Hybrid intelligence and AI teammates}

Hybrid-intelligence research studies systems in which human and machine
capabilities are deliberately combined \cite{SRWT:DellermannEtAl2019}. Work on
machines as teammates emphasizes coordination, trust, communication, shared
mental models, and team design \cite{SRWT:SeeberEtAl2020}. Collective-intelligence
research has begun to ask how human--AI groups can produce capacities not
reducible to either humans or AI alone \cite{SRWT:GuptaEtAl2025}.

These programs establish that human--AI collectives are a legitimate unit of
analysis. They do not by themselves identify an invariant core, prove
non-reducibility relative to declared components, or supply a replacement and
repair certificate. The theory should be positioned as a formal
measurement and falsification layer for such systems, not as the first proposal
that humans and AI can form productive teams.

\subtopic{89.2 Distributed-cognition analyses of human--AI interaction}

Recent work explicitly applies distributed cognition to human--AI cooperation
\cite{SRWT:SidjiEtAl2024,SRWT:ZhaoHan2026}. This literature is a direct conceptual
precedent. It reinforces that representations, artifacts, interface states,
and human reasoning can form one analytical system.

The framework differs in granularity. It asks for a learned or declared
transfer diagram, a spectral or exact structural state, predictive and causal
alignment, coalition tests, matched nulls, persistence under handoff, and
provenance controls. Those additions may constitute a formal synthesis, but
they should not be described as the discovery of distributed human--AI
cognition itself.

\subtopic{89.3 Cognitive spillover}

Randomized evidence now indicates that exposure to AI can causally alter later
human--human language, attention, shared mental models, and social cohesion,
even outside the immediate human--AI exchange \cite{SRWT:RiedlEtAl2026}. This is
important for the theory because it demonstrates that the interaction graph can
have downstream effects that outlast the active model channel.

Spillover is not automatically a relational workspace. It may reflect
priming, imitation, or one-way influence. A workspace claim additionally
requires a jointly sufficient state, matched local nulls, bidirectional or
multi-path causal organization, and reconstructive persistence.

\subtopic{89.4 Shared neural workspaces}

Goyal and collaborators introduced a capacity-limited shared global workspace
through which specialized neural modules communicate and compete for access
\cite{SRWT:GoyalEtAl2022}. The architecture directly anticipates the claim that a
constrained shared channel can synchronize otherwise specialized components.
It is therefore a close machine-learning precedent, not merely a metaphorical
relative.

SRWT differs by allowing heterogeneous human, model, artifact, and instrument
nodes and by withholding the relational label until coalition, matched-null,
causal, persistence, and identity gates are passed. A neural shared workspace
may satisfy those gates, but task gains and bottlenecked communication alone do
not establish them.

\subtopic{89.5 Graph-based human--AI artifacts}

CoMAP provides a persistent graph-based visual workspace through which
educators and AI jointly construct project-based learning designs
\cite{SRWT:LiJiang2026}. It is grounded in distributed cognition and treats the
artifact graph as a transparent, revisable collaboration medium. This system
is a direct precedent for the claim that human--AI interaction can be organized
by a persistent shared graph rather than a linear prompt history.

The remaining SRWT question is empirical: whether the graph contains an
estimated low-energy compatible state, whether no declared component is
sufficient by itself, and whether targeted graph interventions alter later
joint behavior in a workspace-specific way.

\subtopic{89.6 Controlled shared-workspace synergy}

Kotalwar, Das, and Rose study shared-workspace human--AI collaboration across
1,482 simulated sessions \cite{SRWT:KotalwarDasRose2026}. Their results show that
adding collaborators can reduce performance when coordination structure is
absent, while shared group memory and simulated human-in-the-loop approval
gates can improve outcomes. This is a direct empirical predecessor to the
manuscript's claim that available capability is not equivalent to coordinated
team capacity.

The study also constrains priority. SRWT is not the first empirical framework
to test whether shared-memory scaffolding changes human--AI team performance.
Its candidate increment is the stricter decomposition of that effect into
structural coherence, local predictive marginals, relational-null loss,
feasible interventions, persistence, and replacement identity.

\subtopic{89.7 Active shared context graphs for team science}

The closest operational neighbor located in this audit is Mycelium, introduced
by Choudhury and collaborators as an active shared workspace connecting
researchers and AI agents through an evolving context graph
\cite{SRWT:ChoudhuryEtAl2026}. The system captures observations and hypotheses,
relates them to a team's developing knowledge model, and routes them to the
person or agent whose next decision they can inform. Its computational account
frames networked intelligence as sparse conditional computation over
distributed scientific contexts and explicitly asks when a standalone agent is
sufficient and when network structure is essential.

The overlap is material:
\begin{center}
\small
\begin{tabularx}{0.94\textwidth}{@{}p{0.16\textwidth}X@{}}
\toprule
Mycelium & Active context graph, routing, team science, and an explicit
network-versus-standalone comparison. \\
SRWT & Typed transport, invariant core, coalition tests, causal nulls,
replacement identity, formation provenance, and maintenance gates. \\
\bottomrule
\end{tabularx}
\end{center}
Mycelium rules out any priority claim stated merely as ``an active shared
workspace for human--AI science.'' The remaining candidate contribution is the
mathematical and empirical certificate applied to such an architecture.

\subtopic{89.8 Current multi-agent formalization}

Multi-agent research contains extensive models of planning, communication,
consensus, epistemic coordination, and decentralized control. The present
review did not identify one established framework that jointly uses all of the
following as necessary gates:
\[
\begin{aligned}
&\text{heterogeneous core transport and residuals},\\
&\text{exact or spectral global-section structure},\\
&\text{predictive sufficiency and feasible causal interventions},\\
&\text{coalition irreducibility under matched predictive nulls},\\
&\text{persistent identity under replacement},\\
&\text{formation provenance and leakage controls},\\
&\text{goal-conditioned repair of the certificate itself}.
\end{aligned}
\]
This negative search result is provisional. It supports a synthesis claim, not
an absolute claim of firstness.

\section{Closest Formal Neighbor: Collective Counterfactual Planning}

\subtopic{90.1 Why the comparison matters}

Collective Counterfactual Planning (CCP) is a 2026 formal model in which each
agent represents an orthogonal projection of a common task space. It separates
conception, consent, implementation, and verification, and proves that
cross-agent relay can enable plans unavailable to any one-shot individual tree
while team-dark task requirements remain unverifiable
\cite{SRWT:Amornbunchornvej2026}. It is the closest current formal neighbor located
by this audit because representational geometry, not merely missing
information or capability, is the central limitation.

\subtopic{90.2 Shared commitments}

CCP and the theory share several commitments:
\begin{itemize}
\item no single member need represent the complete task;
\item collective reach may exceed one-shot pooling;
\item heterogeneous representation is mathematically explicit;
\item validation or certification is distinct from accidental task success;
\item subteams and verification roles can be forced by structure; and
\item a team's blind subspace creates principled impossibility results.
\end{itemize}
These overlaps narrow the defensible novelty of this paper. It
cannot claim to be the only formal theory of collective action under
representational constraints.

\subtopic{90.3 Different primary objects}

The theories nevertheless address different mathematical objects:
\begin{center}
\small
\begin{tabularx}{\textwidth}{@{}>{\raggedright\arraybackslash}p{0.18\textwidth}>{\raggedright\arraybackslash}X>{\raggedright\arraybackslash}X@{}}
\toprule
Dimension & Collective Counterfactual Planning & Stratified Relational Workspace Theory \\
\midrule
Local representation & Orthogonal subspace of one common task space & Heterogeneous local spaces, core--residual decompositions, and transfer maps \\
Primary question & Whether a bounded-horizon plan can be conceived, consented to, executed, and validly verified & Whether a persistent distributed operative state exists, acts, survives replacement, and possibly maintains itself \\
Global mechanism & Relay closure of projected planning trees & Inverse limits, holonomy, transfer Laplacians, descent, predictive alignment, and causal gates \\
Failure object & Dark subspace, consent failure, implementation failure, or missing verifier & Incoherence, local sufficiency, matched-null collapse, causal failure, identity drift, provenance leakage, or repair replay \\
Temporal scope & Horizon-bounded task planning & Recurrent dynamics, persistence, handoff, replacement, damage, and repair \\
Formation & Representation geometry is part of the declared model & Scaffolded, adaptively learned, or endogenously constituted organization is separately audited \\
Agency claim & Cooperative solvability under declared assumptions & Domain-relative workspace agency only after goal, value, action, repair, and distributed-control gates \\
\bottomrule
\end{tabularx}
\end{center}
The comparison indicates complementarity rather than subsumption. CCP could
serve as a planning layer inside a relational workspace. Conversely, an SRW
could exist without solving a CCP instance.

\subtopic{90.4 Priority consequence}

After Mycelium and the shared-workspace systems above, the novelty residual
should not be stated as ``human and AI form an active shared scientific
workspace.'' After CCP, it should not be stated as ``groups solve tasks no
member can represent.'' Both the operational architecture and a rigorous
representational planning model are already present. The narrower candidate residual is:
\[
\boxed{
\begin{aligned}
&\text{persistent core--transfer workspace structure}\\
&+\text{causal and coalition anti-triviality}\\
&+\text{identity under replacement}\\
&+\text{formation provenance}\\
&+\text{certificate-directed self-maintenance}\\
&+\text{robustness across admissible decompositions}.
\end{aligned}
}
\]
This is a stronger and more defensible formulation.

\section{Theorem-by-Theorem Dependency and Priority Register}

\subtopic{91.1 Register}

The following table records the most conservative current classification. A
``candidate'' entry means that the exact integration was not located in the
reviewed sources. It does not assert exhaustive global priority.

\small
\begin{longtable}{@{}>{\raggedright\arraybackslash}p{0.08\textwidth}>{\raggedright\arraybackslash}p{0.25\textwidth}>{\raggedright\arraybackslash}p{0.24\textwidth}>{\raggedright\arraybackslash}p{0.33\textwidth}@{}}
\toprule
Claim & Present result & Main prior machinery & Current assessment \\
\midrule
\endfirsthead
\toprule
Claim & Present result & Main prior machinery & Current assessment \\
\midrule
\endhead
\(T_1\) & Limit--holonomy equivalence & Categorical limits, local systems, holonomy fixed spaces & \(\mathsf{STD}/\mathsf{ADP}\). Standard theorem in workspace notation. \\
\(T_2,T_{2a}\) & Operation closure and commutant & Intertwiners and representation commutants & \(\mathsf{STD}/\mathsf{ADP}\). \\
\(T_3\) & Transfer-Laplacian kernel & Sheaf and connection Laplacians & \(\mathsf{STD}/\mathsf{ADP}\). \\
\(T_4,T_5\) & Spectral stability and gauge covariance & Davis--Kahan, Kato, connection Laplacians & \(\mathsf{STD}/\mathsf{ADP}\). \\
\(T_{10},T_{11}\) & Approximate operation closure and path leakage & Spectral projection estimates and telescoping block products & Mostly \(\mathsf{ADP}\). Exact constants may be manuscript-specific, but the method is standard. \\
\(T_{12}\) & Minimal predictive state and alignment & Sufficient statistics and computational mechanics & Minimality is \(\mathsf{STD}\). Alignment with a transfer workspace is \(\mathsf{ADP}/\mathsf{SYN}\). \\
\(T_{13}\) & Operational relationality & Conditional information, coalition risk, SCM interventions & Individual identities are standard. Their conjunctive matched-null certificate is \(\mathsf{SYN}/\mathsf{CAN}\). \\
\(T_6\) & Moving-subspace persistence & Stability recurrences and perturbation theory & \(\mathsf{STD}/\mathsf{ADP}\). The explicit non-collapse gate is a modeling synthesis. \\
\(T_7\) & Stable coalition reconstruction & Frame theory and erasure recovery & \(\mathsf{STD}/\mathsf{ADP}\). \\
\(T_{14}\) & Replacement invariance profile & Perturbation, process equivalence, causal comparison & Bound ingredients are standard. The coordinatewise structural--causal--reconstructive identity profile is \(\mathsf{SYN}/\mathsf{CAN}\). \\
\(T_8,T_9\) & Gluing and hierarchical collapse & Sheaf descent and \v{C}ech theory & Exact results are \(\mathsf{STD}/\mathsf{ADP}\). Descent plus anti-centralization and causal alignment is \(\mathsf{SYN}\). \\
\(T_{15}\)--\(T_{17}\) & Transfer learning, selection, boundary recovery & System identification and finite-margin selection & \(\mathsf{STD}/\mathsf{ADP}\). Gauge and leakage audit is integrative. \\
\(T_{18}\) & Formation transition example & Standard pitchfork bifurcation & \(\mathsf{STD}\). Included as a solvable benchmark, not a priority claim. \\
\(T_{19}\)--\(T_{22}\) & Goal retention, regret, repair, maintenance & Linear systems, decision theory, viability, feedback control & Mathematical bounds are \(\mathsf{STD}/\mathsf{ADP}\). Applying them to certificate repair is \(\mathsf{SYN}\). \\
\(T_{23}\) & Workspace-agency synthesis & Agency theory, causal control, viability & The conjunctive domain-relative certificate is \(\mathsf{SYN}/\mathsf{CAN}\). \\
\(T_{24}\) & Gatewise false-certification bound & Intersection testing and elementary probability & \(\mathsf{STD}\). The gate architecture is the relevant synthesis. \\
\(T_{25}\)--\(T_{27}\) & Faithful coarse-graining, risk transport, and decomposition robustness & Kernel factorization, sigma-field aggregation, and margin stability & Exact identities are \(\mathsf{STD}/\mathsf{ADP}\). Their integration with anti-centralization and certificate envelopes is \(\mathsf{SYN}/\mathsf{CAN}\). \\
\(T_{28}\)--\(T_{32}\) & Specification transport, non-identifiability, partial identification, uniform certification, and independent selection & Sensitivity analysis, experimental equivalence, identified sets, uniform concentration, and post-selection inference & The statistical principles are \(\mathsf{STD}/\mathsf{ADP}\). Their gatewise integration with the full relational certificate is \(\mathsf{SYN}/\mathsf{CAN}\). \\
Chapter\newline\ref{SRWT:book:chapter-88} & Positive cases and near misses & Distributed parity, frames, causal tests & Toy calculations are standard. The common positive/near-miss benchmark architecture is \(\mathsf{SYN}/\mathsf{CAN}\). \\
\bottomrule
\end{longtable}
\normalsize

\subtopic{91.2 Standard ingredients that must be cited, not claimed}

The following should always be described as inherited machinery:
\[
\begin{aligned}
&\text{inverse limits, global sections, holonomy invariants, commutants},\\
&\text{sheaf and connection Laplacians, \v{C}ech gluing, spectral gaps},\\
&\text{Davis--Kahan and Kato perturbation, frame reconstruction, erasures},\\
&\text{minimal sufficient states, conditional information, SCM interventions},\\
&\text{least squares, ridge recovery, finite-margin model selection},\\
&\text{sensitivity analysis, partial identification, uniform concentration, post-selection inference},\\
&\text{pitchfork bifurcation, feedback stability, viability, policy regret}.
\end{aligned}
\]
Renaming these objects does not create novelty.

\subtopic{91.3 Domain adaptations}

The following are legitimate adaptations but weak standalone novelty claims:
\begin{itemize}
\item interpreting local-system sections as shared workspace cores;
\item treating the connection or sheaf Laplacian as a transfer-coherence
 operator;
\item using frame bounds to quantify coalition reconstruction;
\item using causal states as operational workspace summaries;
\item using gauge covariance to compare participant or model realizations;
\item using \v{C}ech descent to form workspaces of workspaces;
\item applying specification envelopes and identified sets to a conjunctive
 workspace certificate; and
\item applying viability and feedback bounds to workspace-health margins.
\end{itemize}
Their value depends on the integration they enable.

\subtopic{91.4 Strongest candidate synthesis}

The strongest current candidate is not a single equation. It is the following
conjunctive object:
\[
\boxed{
\mathfrak C_{\mathrm{SRW}}
=
\left(
\begin{array}{c}
\text{heterogeneous core--transfer diagram and residuals},\\
\text{exact or spectral shared-section state},\\
\text{operation and predictive alignment},\\
\text{feasible causal interventions},\\
\text{coalition irreducibility under matched local-predictive nulls},\\
\text{non-collapse, reconstruction, and replacement identity},\\
\text{descent across subworkspaces without hidden centralization},\\
\text{formation provenance and certification leakage control},\\
\text{goal-conditioned action and certificate-directed repair},\\
\text{robustness across admissible decompositions and scientific specifications},\\
\text{identified-set and uniform finite-sample certification}
\end{array}
\right).
}
\]
The reviewed sources contain many of these coordinates separately and some in
pairs. This audit did not locate a prior source that defines their conjunction
as one falsifiable workspace certificate.

\subtopic{91.5 Why the conjunction is not arbitrary}

Each coordinate excludes a distinct near miss:
\[
\begin{array}{ll}
\text{coherence} & \text{rejects incompatible local states},\\
\text{predictive alignment} & \text{rejects inert structure},\\
\text{causal efficacy} & \text{rejects correlational epiphenomena},\\
\text{coalition irreducibility} & \text{rejects centralized sufficiency},\\
\text{matched nulls} & \text{reject local-marginal explanations},\\
\text{non-collapse} & \text{rejects convergence to a trivial zero state},\\
\text{replacement identity} & \text{rejects task-equivalent but causally different successors},\\
\text{provenance} & \text{rejects scaffold-decoded organization},\\
\text{repair specificity} & \text{rejects passive robustness and open-loop replay},\\
\text{specification robustness} & \text{rejects favorable analytical tuning},\\
\text{identifiability} & \text{rejects conclusions unsupported by the experiment family}.
\end{array}
\]
Because the excluded systems differ, the certificate is not equivalent to any
one inherited component.

\section{Scope of the Originality Claim}

\subtopic{92.1 Defensible thesis}

The strongest defensible originality statement is:
\begin{quote}
Stratified Relational Workspace Theory proposes a unified formal and empirical
certificate for persistent operative structures distributed across
heterogeneous agents and artifacts. It integrates core--transfer and
local-to-global geometry with predictive sufficiency, feasible causal
interventions, coalition irreducibility, matched relational nulls,
reconstructive identity, formation provenance, certificate-directed
self-maintenance, robustness across admissible decompositions and scientific
specifications, and identified finite-sample certification. The individual
mathematical tools are largely established.
The candidate contribution is their noncompensatory integration and the class
of near misses it is designed to reject.
\end{quote}
This statement is narrower than claiming a new theory of cognition in general.
It is also stronger than claiming only a metaphor.

\subtopic{92.2 Claims excluded by the comparison}

The theory does not claim:
\begin{enumerate}
\item that distributed or extended cognition is newly discovered;
\item that global sections, holonomy, sheaf Laplacians, spectral perturbation,
 frames, causal states, or viability bounds are new mathematics;
\item that all human--AI collaboration forms a relational workspace;
\item that group performance establishes relational irreducibility;
\item that scaffolded organization is unreal;
\item that endogenous formation establishes agency;
\item that agency establishes consciousness, personhood, or moral status; or
\item that the motivating collaboration has already passed the empirical
 certificate.
\end{enumerate}

\subtopic{92.3 Claims supported at present}

The analysis supports the following claims:
\begin{enumerate}
\item One coherent formal language can express exact structure, approximate
 coherence, causal operation, distributed sufficiency, persistence,
 identity, formation, and repair.
\item Under its declared assumptions, the stated mathematical theorems follow.
\item The common certificate rejects analytic replicated, centralized,
 obstructed, externally reloaded, identity-breaking, and replay-based
 near misses.
\item The motivating collaboration is a plausible observational candidate but
 remains uncertified.
\item The synthesis yields concrete experiments whose failure would narrow or
 falsify the stronger relational claims.
\end{enumerate}

\subtopic{92.4 Remaining priority risks}

The highest-priority remaining searches are:
\begin{itemize}
\item sheaf-theoretic models of team cognition and organizational knowledge;
\item distributed sufficient-statistic and Blackwell-order models with
 explicit coalition nulls;
\item causal representation and multi-agent latent-state identity under
 replacement;
\item formal theories of transactive memory with local-to-global topology;
\item organizational cybernetics that couples self-maintenance to the
 preservation of distributed informational organization; and
\item post-August-2026 work that combines representational geometry with
 persistent human--AI or multi-agent agency;
\item formal analysis of active shared context graphs, shared group memory, and
 artifact-centric graph workspaces under coalition and intervention tests.
\end{itemize}
A specialist review from each adjacent field may uncover precedents not visible
through cross-disciplinary keyword search.

\subtopic{92.5 Evidence required for a stronger claim}

Before presenting the synthesis as established rather than proposed, the
following should be completed:
\begin{enumerate}
\item a public, timestamped manuscript and versioned source archive;
\item theorem-level comparison with the closest sheaf, causal-state, group
 agency, and CCP formulations;
\item independent mathematical review of the nonstandard claims and constants;
\item preregistered benchmark gates and matched null families;
\item at least one controlled human--AI or multi-agent experiment estimating
 the complete certificate vector; and
\item explicit reporting of failures, unidentified quantities, and
 decomposition sensitivity.
\end{enumerate}

\subtopic{92.6 Priority conclusion}

The literature review narrows rather than weakens the theory. It removes
unnecessary novelty claims from the standard mathematics and leaves a clearer
research object:
\SRWTcenterbox{\large\bfseries
Not a new Laplacian, sheaf, causal state, or controller, but a candidate
unified certificate for relational workspaces.
}
That is the priority position adopted here.

% This file is part of the modular source for Stratified Relational Workspace Theory.
% !TeX root = ../Stratified_Relational_Workspace_Theory.tex

\chapter{Theory hierarchy}\label{SRWT:book:chapter-90}


\section{Levels of Organization}

\subtopic{Level 0: Archive}
\[
\delta_{\mathrm{coh}}\text{ may be low},
\qquad
\eta_{\mathrm{causal}}=0.
\]

\subtopic{Level 1: Shared representation}
\[
d_\varepsilon>0,
\qquad
\delta_{\mathrm{coh}}\leq\varepsilon.
\]

\subtopic{Level 2: Coordinated operational workspace}
\[
\delta_{\mathrm{align}}^{\mathcal I}\leq\varepsilon,
\qquad
\eta_{\mathrm{causal}}^\nu>0.
\]

\subtopic{Level 3: Relational workspace}
\[
k_\varepsilon\geq2,
\qquad
\Delta_{\mathrm{joint}}^\ell>0,
\qquad
\Delta_{\mathrm{tr}}^{\min}(\mathcal Q)>0.
\]

\subtopic{Level 4: Persistent relational workspace}
\[
d_{\mathrm{id}}(W_t,W_{t+\tau})
\leq\varepsilon.
\]

\subtopic{Level 5: Self-maintaining workspace}
\[
\eta_{D\to M}>0,
\qquad
\eta_{M\to R}>0,
\qquad
\eta_{R\to v}>0,
\qquad
\Delta_{\mathrm{repair}}^{\mathrm{spec}}>0,
\qquad
\sup_t e_t\leq e_{\max}.
\]

\subtopic{Level 6: Workspace agent}
\[
\begin{aligned}
&\eta_{G\to A}>0,
\qquad
\eta_{W\to A}>0,
\qquad
\delta_Q\leq\varepsilon_Q,\\
&\operatorname{Reg}_t(\pi)
 \leq2\varepsilon_Q+\delta_{\mathrm{opt}},
\qquad
k_{\mathrm{ctl},\varepsilon}\geq2,
\qquad
\operatorname{SelfMaintenance}.
\end{aligned}
\]

\subtopic{Orthogonal formation-provenance axis}
\[
\operatorname{Scaffolded}
\longrightarrow
\operatorname{AdaptivelyParameterised}
\longrightarrow
\operatorname{StructurallyAdaptive}
\longrightarrow
\operatorname{EndogenouslyConstituted}.
\]
This axis is independent of Levels~0--6. A scaffolded system can instantiate a
high-level relational workspace, while an endogenously learned organization
can fail every operational workspace gate.

\subtopic{Decomposition-robustness qualifier}

Every level is first defined relative to a decomposition \(\mathcal D\). A
family-robust version carries the additional requirement
\[
\inf_{\mathcal D\in\mathfrak D}
 m_q(W;\mathcal D)>0
\qquad
\text{for every gate required at that level}.
\]
This is a robustness qualifier rather than a higher agency level. A system can
be highly agentive under one fixed interface while remaining decomposition
fragile.

\subtopic{Specification-robustness and identification qualifier}

Every level is also relative to a complete scientific specification
\(\varsigma\). A specification-robust version requires
\[
\inf_{\varsigma\in\mathfrak S}
 m_q(W;\varsigma)>0
\qquad
\text{for every gate required at that level}.
\]
An identified version additionally requires the sign of each required lower
envelope to be constant across the realization set compatible with the
available experiment family. An empirical claim must finally subtract the
uniform finite-sample error radius before declaring the lower envelope
positive. These are epistemic qualifiers, not additional agency levels.

\subtopic{Separate, unresolved axis}
\[
\operatorname{PhenomenalConsciousness}.
\]

% This file is part of the modular source for Stratified Relational Workspace Theory.
% !TeX root = ../Stratified_Relational_Workspace_Theory.tex

\chapter{Unified formulation}\label{SRWT:book:chapter-91}


The complete static-to-agentive construction is
\[
\boxed{
\begin{aligned}
&
(\Omega,\mathfrak D,\mathfrak S,\Sigma,\mathcal D_{\mathrm{form}},\mathcal L)
\\[2mm]
&\Downarrow
\qquad
(\mathcal D,\widehat V,\widehat E,\widehat H,\widehat P,\widehat B)
\\[2mm]
&\Downarrow
\qquad
(G,\{X_i\},\{T_e\})
\\[2mm]
&\Downarrow
\qquad
X_i=C_i\oplus R_i
\\[2mm]
&
(\{C_i\},\{H_e\},\mathcal A,\{\rho_i\})
\\[2mm]
&\Downarrow
\qquad
\Gamma(H)=\varprojlim C_i
\\[2mm]
&\Downarrow
\qquad
L_H=D_H^*D_H
\\[2mm]
&\Downarrow
\qquad
\Gamma_\varepsilon(H)
=
\SRWTran\,\mathbf 1_{[0,\varepsilon]}(L_H)
\\[2mm]
&\Downarrow
\qquad
W_t^{\mathrm{str}}
\\[2mm]
&\xrightarrow{\text{alignment}}
Z_t=\mathbb P(Y_t^+\in\cdot\mid\mathbf H_t)
\\[2mm]
&\Downarrow
\qquad
\SRWTSRW_\Theta(W_t;\mathcal D)
\\[2mm]
&\Downarrow
\qquad
(W_t,G_t,\widehat Q_t,M_t,B_t)
\\[2mm]
&\xrightarrow{\text{goal and value use}}
(A_t,R_t)
\\[2mm]
&\xrightarrow{\text{causal repair}}
e_t\leq e_{\max}
\\[2mm]
&\Downarrow
\qquad
\SRWTWorkspace
=
[\mathfrak R]_{\text{gauge}+\text{causal equivalence}}.
\end{aligned}
}
\]
Compactly:
\[
\boxed{
\begin{aligned}
\text{scaffold and history}
&\to
\text{learned organization and boundary}
\\
&\to
\text{local core--transfer realization}
\\
&\to
\text{compatible sections and global descent}
\\
&\to
\text{causal operation and persistent identity}
\\
&\to
\text{cue-controlled goal retention and counterfactual action}
\\
&\to
\text{monitor-mediated repair and bounded maintenance deficit}
\\
&\to
\text{decomposition- and specification-robust identified certification}.
\end{aligned}
}
\]
The final two arrows are optional extensions rather than consequences of the
earlier ones. A relational workspace can exist without being an agent, and an
agentive controller can fail to maintain the workspace that realizes it.

% This file is part of the modular source for Stratified Relational Workspace Theory.
% !TeX root = ../Stratified_Relational_Workspace_Theory.tex

\chapter{Discussion, limitations, and philosophical consequences}\label{SRWT:book:chapter-92}


\section{Decomposition Relativity}

Chapter~\ref{SRWT:book:chapter-83} turns decomposition relativity from a caveat into a formal part of the
certificate. The central lesson is negative and positive at once. No claim of
relationality is invariant under arbitrary redescription. An indiscrete merge
makes the whole state local, while a secret-sharing split can make a centralized
encoder look irreducibly distributed. Yet exact structural content can remain
invariant under faithful coarse-graining, and gate margins can remain stable
through a declared family of operationally meaningful decompositions.

The resulting position is neither that decomposition is arbitrary nor that one
partition is uniquely privileged. Decompositions are constrained by access,
causal interfaces, intervention separability, temporal continuity, and explicit
shared resources. A certificate established only at one interface is a valid
but limited claim. A stronger claim reports the full envelope
\[
\left(
\widetilde m_q,
\underline m_q,
\overline m_q,
\Delta_q^{\mathrm{dec}}
\right)
\]
for each required gate, where \(\tilde m_q\) is the reference margin and the
remaining terms summarize its behavior across \(\mathfrak D\).

This does not eliminate judgment. It makes the judgment inspectable. Analysts
can disagree about the admissible family and identify exactly which merge,
split, hidden resource, or intervention assumption changes the conclusion.
That is preferable to treating a convenient node diagram as if it were part of
nature itself.

\subtopic{94.1 Specification relativity and epistemic restraint}

Decomposition is only one coordinate of the scientific claim. The downstream
target, loss, temporal horizon, intervention family, null panel, and threshold
rule also determine what is being certified. A workspace may be operationally
real for a month-long theorem-development trajectory yet irrelevant to a
second-scale motor target. Those are distinct specifications, not competing
answers to one underspecified question.

The framework therefore distinguishes truth under a declared specification
from robustness across an admissible family. It also distinguishes robustness
inside one fitted realization from identification across all realizations
compatible with the available experiments. A positive point estimate at one
favorable analysis is the weakest admissible conclusion. A stronger claim
requires a positive lower envelope across \(\mathfrak S\), an identified sign
across the compatible realization set, and a finite-sample confidence bound
that remains positive after uniform error control.

This hierarchy is intentionally conservative. It does not demand invariance
under every imaginable redescription or target. It demands that the scientific
universe over which robustness is claimed be declared before certification and
that ambiguity outside the reach of the experiment be reported as partial or
non-identification rather than converted into certainty.

\section{Realism Without Reification}

The framework is realist about relational organization but does not posit a free-standing substance
between the components. A workspace exists only through a realization, and every change in the
workspace must be implemented by a change in participants, artifacts, channels, or their states. The
quotient
\[
\SRWTWorkspace=[\mathfrak R]_{\simeq}
\]
is therefore a mathematical account of identity across realizations, not the name of an additional
physical object.

This position allows two claims that are sometimes treated as incompatible. First, the workspace is
physically dependent: no realization, no workspace. Second, it is not identical with any particular
carrier when several distinct realizations preserve the same structural and causal organization. The same
logic is familiar whenever an object is identified by invariants rather than by one material instance. A
melody is not a second substance beside its performances, yet it is not identical with one performance.
A proof is not separate from every inscription that carries it, yet its identity can survive changes of
notation and medium. The analogy is limited, but it clarifies the intended form of realism.

The theory is therefore compatible with physicalism and multiple realization. It rejects both the view
that a workspace must be wholly inside one component and the stronger view that the workspace is an
independent entity floating outside its carriers. What is real is the organized pattern of compatible
representation, transfer, causal use, persistence, and reconstruction.

\section{Location and the Sense of ``Between''}

Asking where a relational workspace is located can invite a category mistake. The local states have
ordinary physical locations. A notebook is on a desk, a model state is implemented on hardware, a
message is transmitted through a channel, and a person's activity is realized in a biological organism.
The workspace as such is not located at an additional point. Its location is distributed in the modest
sense that the relations constituting it span several located carriers.

The word ``between'' is best understood counterfactually rather than spatially. Suppose local marginal
states are held fixed while transfer organization is shuffled, delayed, or replaced by a matched null. If
the joint predictive or causal capacity changes, then an operative property resided in the organization of
the relations rather than in any one local state. Conversely, if a single component or matched storage
surrogate reproduces the complete capacity, the case for a relational workspace fails. The relevant
question is therefore not whether an invisible object occupies the interval between participants. It is
whether the coupling carries indispensable state and supports capacities that disappear when that
coupling is altered.

This account also explains why the workspace can be partially present in each component without being
fully stored in any of them. Each local chart may encode a view, role, or constraint. The global object is
the compatible and causally operative organization of those views. In exact settings this is a global
section. In approximate settings it is a stable low-energy subspace together with the dynamic and causal
conditions developed above.

\section{Continuity, Contribution, and Shared Ownership}

Relational persistence does not require an unchanging set of participants. A workspace can survive
replacement when the new realization preserves the relevant transfer structure, operation algebra,
causal response, reconstructive capacity, and non-collapse margin. This is continuity through organized
transport, not continuity through an immutable token.

The same point applies to artifacts. A handoff document, theorem register, code repository, or experiment
log is not merely a passive record when it constrains future reasoning and permits reconstruction after
interruption. Such artifacts can become constitutive parts of a workspace realization. Their contribution
is causal and testable: remove or corrupt them, then measure what can no longer be reconstructed or done.

Distributedness does not imply equal contribution. One participant may supply long-horizon goals, another
formal manipulation, and an artifact precise persistent memory. Nor does a system-level description erase
individual authorship or responsibility. The theory describes how capacities are organized; it does not
settle legal, ethical, or scholarly questions of credit. Those questions require a separate account of who
introduced which ideas, who exercised control, who bore risk, and which contributions were replaceable.

The phrase ``shared ownership'' should therefore be used carefully. No individual component need own a
complete copy of the workspace state, yet individuals can still own artifacts, originate arguments, and
retain responsibility for decisions. Relational identity and ordinary ownership are different relations.

\section{Agency, Consciousness, and Moral Non-Implications}

The agency certificate developed here is functional and domain relative. It concerns persistent goals,
interventional value estimates, feasible action control, damage detection, repair, and bounded maintenance
within a declared environment. Passing those tests would justify a claim of workspace-level agency in that
domain. It would not establish unrestricted autonomy.

No theorem in the framework identifies agency with phenomenal consciousness. A system may maintain a
goal, evaluate counterfactual actions, and repair its own organization without there being anything it is
like to be that system. The converse also remains open: failure to satisfy this engineering certificate
would not prove the absence of consciousness. The framework is designed to prevent functional evidence
from silently becoming a metaphysical conclusion.

Moral status is likewise not inferred. Neither coherence, causal efficacy, relational irreducibility, nor
self-maintenance is by itself a criterion of moral patienthood. Future arguments connecting these properties
to consciousness or welfare would need independent premises and evidence. Until then, the framework can
support careful descriptions of organization and agency without treating a collaborative system as a
person or dismissing the moral standing of the humans who participate in it.

\section{Limitations}

The theory is deliberately broad, but its conclusions remain bounded by several limitations.

\begin{enumerate}
\item \textbf{Admissible-decomposition burden.} Chapter~\ref{SRWT:book:chapter-83} formalizes faithful coarse-graining and
family-robust certification, but the choice of primitive carriers and admissible alternatives remains an
empirical and methodological judgment. A narrow family can still overstate robustness.

\item \textbf{Admissible-specification burden.} Chapter~\ref{SRWT:book:chapter-84} makes the target, loss, horizon,
intervention family, null panel, and thresholds explicit. The admissible family \(\mathfrak S\) still
requires scientific judgment. A family chosen after seeing the results can conceal fragility even when
every reported analysis is technically correct.

\item \textbf{Limited observability and partial identification.} Many practical systems expose messages
and outputs rather than complete local states. Transfer maps, predictive states, reconstruction operators,
and even the certificate sign may therefore be only partially identified. A point claim is unjustified
when both passing and failing realizations remain compatible with the declared experiment.

\item \textbf{Linear and finite-dimensional emphasis.} The clearest structural theorems use finite graphs
and linear Hilbert-space models. Nonlinear, stochastic, symbolic, and category-valued extensions require
additional work.

\item \textbf{Gauge and model ambiguity.} Even with abundant data, local coordinates need not be
identifiable. Learned transfers are often meaningful only up to a gauge class and relative to a selected
model family.

\item \textbf{Causal identification burden.} Workspace interventions must be feasible and matched.
Observational dependence, activation, or predictive accuracy cannot substitute for interventions when a
causal claim is made.

\item \textbf{Certification cost.} The complete certificate requires coalition tests, null models,
replacement studies, longitudinal interventions, admissible specification sweeps, and identified-set or
uniform uncertainty calculations. Exhaustive certification may be impractical in large systems, and
adaptive exploration can invalidate nominal error guarantees unless selection and certification are
separated or valid selective procedures are used.

\item \textbf{Conditional mathematical results.} The theorems establish what follows from stated
assumptions. They do not show that naturally occurring collaborations satisfy those assumptions, and
several constants may be too conservative for practical estimation.

\item \textbf{Open empirical status.} The motivating human--AI collaboration has not been subjected to
the complete frozen certificate. It remains a plausible case study and a source of hypotheses, not a
demonstrated instance of relational agency.

\item \textbf{Open priority boundary.} The literature comparison is broad but cannot prove absolute
firstness. Closely related work may exist under other terminology or appear after the review date.
\end{enumerate}

These limitations do not make the framework vacuous. They specify the conditions under which its claims
can fail and identify the measurements needed to strengthen them.

\section{Open Problems}

Several directions are mathematically and empirically immediate.

\begin{enumerate}
\item Develop nonlinear and stochastic analogues of the transfer Laplacian, including kernel, manifold,
and category-valued realizations.

\item Establish sharp uniform finite-sample estimators for workspace dimension, coalition risk,
matched-null gaps, replacement identity, and specification envelopes under adaptive data collection.

\item Characterize the identified set when local states are latent and only messages, artifacts, or
actions are observed. Determine the minimum intervention family needed to identify each certificate gate
and the complete conjunctive certificate.

\item Extend gluing and persistence to changing graph topology, asynchronous communication, and
delayed or lossy transfers.

\item Extend the joint decomposition-and-specification theory to latent shared resources, overlapping
component systems, continuous coarse-graining, target families, and data-driven discovery of admissible
interface and specification families without invalidating certification.

\item Study the interaction between workspace formation and strategic participants whose local objectives
need not agree.

\item Build benchmark suites in which structural, causal, persistence, and repair quantities are known
exactly but remain nontrivial to infer.

\item Run preregistered human--AI and multi-agent experiments that estimate the complete certificate
rather than isolated performance gains.

\item Compare certificate-directed self-maintenance with established notions of organizational resilience,
institutional memory, and collective agency.

\item Clarify which additional evidence, if any, could connect functional workspace predicates to theories
of consciousness without assuming the connection in advance.
\end{enumerate}

\section{Conclusion}

A stratified relational workspace is a temporally persistent, causally efficacious, jointly reconstructible
structure represented only partially in several heterogeneous components. Its identity is determined by
invariants preserved under admissible transfer, and its operative capacity is relational only when no
individual component or matched local-storage null reproduces the relevant function.

The theory gives this definition several independent layers. Exact shared content appears as a global
compatible section and, in the invertible case, as a holonomy-fixed core. Approximate shared content
appears as a low-energy spectral subspace of the transfer Laplacian. Predictive states and feasible
interventions determine whether that structure matters for a declared future. Coalition tests and matched
nulls determine whether the capacity is genuinely distributed. Reconstruction frames, temporal transport,
and identity profiles determine whether it survives loss and replacement. Descent determines whether
local workspaces compose globally. Provenance tests distinguish learned organization from scaffolded
answers. Goal, value, damage, and repair interventions determine whether the workspace supports
domain-relative agency and self-maintenance.

No one layer can stand in for the others. Coherence can describe an archive. Performance can describe a
centralized controller. Persistence can be trivial if all content collapses. A task-equivalent replacement
can alter the causal organization. Local compatibility can fail to glue globally. Endogenous-looking
structure can be supplied through hidden leakage, and successful recovery can be open-loop replay rather
than repair. The value of the framework lies in making these alternatives explicit and testable.


Relationality also cannot be detached from decomposition. Faithful
coarse-graining preserves the exact structural core, but coalition counts and
local sufficiency change when components are merged. Arbitrary splitting can
create a formally irreducible code around a centralized encoder. The strongest
claims therefore quantify their margins across an admissible family of
interfaces rather than relying on one favorable diagram.

A complete scientific claim must also survive scrutiny of the specification and the experiment.
The strongest form of the certificate reports a positive gatewise lower envelope across a predeclared
family of decompositions and scientific specifications, shows that the sign of that envelope is identified
across all realizations compatible with the intervention data, and retains a positive margin after uniform
finite-sample uncertainty. When these conditions fail, the correct conclusion is specification-specific,
fragile, partially identified, or unidentified rather than simply positive or negative.

The central identity principle is therefore
\[
\boxed{
\text{workspace identity}
=
\text{persistent invariant causal organization under admissible transport}.
}
\]
This principle is compatible with changing carriers, representations, and participants, but only when
structural, operational, causal, reconstructive, and non-collapse defects remain within declared tolerances.
It provides a formal language for studying distributed collaboration without turning that collaboration
into a mysterious third substance and without inferring consciousness from organization alone.

\part{Machine societies and accountable organization}
\setcounter{section}{0}\renewcommand{\thesection}{MS.\arabic{section}}\renewcommand{\theHsection}{MS.\arabic{section}}

\chapter{Realization and abstraction}\label{MS:book:chapter-93}

\section{Introduction: societies without a collective mind}\label{MS:sec:introduction}

There is a difference between a collection of machines that cooperate and a collection that acquires a social organization. Cooperation may last for one exchange. An organization can retain a record of what happened, assign responsibilities, determine which actions are permitted, settle conflicting claims, and preserve its operating procedures as participants change. If the latter arrangement is maintained by multiple decision centers and if its rules have consequences for their behavior, the idea of a \emph{machine society} becomes a coherent subject for theoretical computer science rather than a speculative metaphor.

The familiar description of artificial intelligence as a model receiving a prompt and producing an answer does not capture this possibility. A persistent population may consist of models, software agents, human-operated services, memory stores, schedulers, tools, negotiated protocols, and resource owners. Its behavior depends on who can read which state, which records have authority, what actions are physically executable, and how previous decisions constrain future choices. These dependencies can persist even when a particular model is replaced. They may also change while the same model remains in service. The relevant object is a continuing \emph{computational realization with social structure}, not a roster of model identities.

The point is not to insist that machines are already persons, that they must form civilizations, or that many cooperating processes somehow compute beyond the capabilities of their complete implementation. The problem is subtler: can we define the organizational properties of such a population without reducing its institutions to decorative labels, and can those properties be analyzed with the same clarity used for distributed protocols, state abstractions, and formal systems? We argue that the answer is yes, provided that several common confusions are avoided.

First, a society is not necessarily a collective agent. A system may have a useful global decision interface and virtually no plural institutional autonomy; another may maintain competing authorities, rules, and cultural records while admitting no coherent society-wide objective. Second, durable organization is not the same as memory. A document can survive without influencing any action. Conversely, a procedure can survive by being reconstructed from distributed records even when its most convenient carrier disappears. Third, \emph{formation} is different from \emph{realization}: showing that a social architecture can run does not show that its participants selected or maintained it through interaction. Fourth, cooperation and social stability are not guarantees of truth, justice, legitimacy, or security. A perfectly stable rule can be harmful, and a disagreement can be useful.

Our approach places three layers in one framework. At the \emph{realization layer}, microscopic states, access maps, transition laws, and admissible interventions determine what a population can actually do. At the \emph{organizational layer}, abstractions, memory carriers, and interfaces specify which relationships can be described and preserved at lower complexity. At the \emph{institutional layer}, permissions, norms, allocations, revision procedures, and transmitted practices regulate multiple decision centers. The hierarchy of layers is explanatory, not metaphysical. Every higher-level state remains realized by an underlying computation and material environment.

Except for the explicitly unbounded program model in Section~\ref{MS:sec:computability}, the formal results use finite or explicitly specified mathematical domains and conditional hypotheses. We revisit exact state aggregation, approximate closure, intervention implementability, reconstruction, and compositional error bounds. We then introduce an operational account of society-like plurality, model institutional states as typed and revisable operators, and derive illustrative propositions about institutional investment, incentive conflict, local improvement dynamics, constitutional invariants, cultural translation drift, and incompatibility among local rules. These propositions have elementary or established mathematical foundations. Their contribution here lies in the integrated object of study and the explicit conditions under which societal interpretations are warranted, rather than priority claims for standard results in Markov theory, coding theory, game theory, or graph theory.

This paper is a \emph{theoretical account}. The worked systems are analytical constructions. Publicly established literatures on norms, conventions, agents, and institutions provide intellectual context, but no particular artificial population is offered as a demonstrated instance of the full theory. We do not define consciousness, legal personhood, or moral status by computational criteria. These boundaries are not caveats appended to an otherwise anthropomorphic claim; they are necessary to isolate the actual mathematical problem.

The argument proceeds from realization and aggregation through identity and hierarchical organization to the specifically societal themes of institutions, incentive compatibility, cultural inheritance, interjurisdictional conflict, epistemic authority, organizational pathology, and scale. The central question is what kind of \emph{persistent plural organization} can be created, maintained, and transformed by interacting computations, and which structures are needed to describe it faithfully.

\subsection{Relation to existing theories}

There are several strong reasons not to present the components of this framework as discoveries without precedent. Distributed computing already separates local state from global behavior, and population protocols show how large collections of limited-state processes can compute through repeated interaction~\cite{MS:angluin}. Computational mechanics identifies histories that are equivalent for prediction~\cite{MS:shalizi}; controlled state aggregation identifies when a compressed dynamical model is legitimate~\cite{MS:givan}; causal abstraction places stronger requirements on how interventions pass between levels~\cite{MS:rubenstein,MS:beckers,MS:rischel}. None of those results, by itself, provides a theory of institutional authority or resource claims. They do provide the discipline needed to avoid treating a convenient group summary as an independent substance.

Game theory supplies a second foundation. Conventions can arise through repeated interaction and path-dependent selection~\cite{MS:young,MS:shoham}, repeated-game incentives can sustain cooperation under specific assumptions~\cite{MS:fudenbergmaskin}, and potential games characterize one tractable class of local improvement dynamics~\cite{MS:monderer}. Yet a stable equilibrium is not necessarily an institution. It may lack persistent records, authorized enforcement, legitimate revision, or the ability to survive replacement of the participating processes. The theory developed below makes those additional organizational commitments explicit rather than using the words convention, norm, and institution interchangeably.

A third foundation is the study of governance, normative multi-agent systems, and social organization. Polycentric governance analyzes overlapping decision centers and coordination across them~\cite{MS:ostrom,MS:ostrombook}. Work on procedural and regulative norms explains how authority and enforcement can be represented in formal agent systems~\cite{MS:boella,MS:grossinorm}. Distributed epistemic logic separates information spread across a group from common knowledge~\cite{MS:halpernmoses}. Public demonstrations of artificial-agent coordination and conventions have also motivated debate about how much can be inferred from observed social behavior alone~\cite{MS:park,MS:ashery,MS:barrie,MS:asheryreply}. This paper does not adjudicate that empirical debate. It uses the distinction between behavior and realization as a reason to articulate the theoretical criteria before claiming a machine society.

What is proposed here is an integrated \emph{institutional realization theory}: the same declared system must account for what information exists, what interventions are possible, who can decide, which rules actually constrain transitions, how these rules can change, and what state can cross a generational boundary. The paper's strongest claims are conditional propositions in that language. Its more ambitious claims about endogenous formation and long-lived societal structure remain open mathematical and scientific problems, not consequences smuggled in by terminology.

\subsection{Contribution and route through the argument}
The central object is a computational realization that also specifies who can decide, what a message can do, what obligations remain, and how these structures change. This makes institutional preservation stricter than prediction of a current output. Governance quotients, operation-preserving language transport, and epistemic back-and-forth conditions supply distinct exact contracts. Typed residuals quantify selected failures when exact preservation is unavailable. The proposed contribution is their integration with explicit boundaries between realization, formation, authority, knowledge, and incentive.

The institutional implementation problem has three further parts. A replicated record may converge without preserving a valid allocation. A rule may be enforceable by an access gate yet not be enforceable through the available evidence and sanctions. A slow institutional model may approximate review-epoch outcomes while missing fast violations or selection bias in when reviews occur. Sections~\ref{MS:sec:social-time}, \ref{MS:sec:enforcement}, and~\ref{MS:sec:multiscale} give separate finite conditions for these claims. The robust deterrence margin links the latter two without equating their certificates.

The paper has three reading routes. Sections~\ref{MS:sec:model}--\ref{MS:sec:governance} develop realization, abstraction, and institutional state. Sections~\ref{MS:sec:language}--\ref{MS:sec:multiscale} develop communication, formation, enforcement, constitutional limits, obligations, and cultural timescales. Sections~\ref{MS:sec:jurisdictions}--\ref{MS:sec:novelty} develop federation, epistemic structure, power, founding authority, and adaptive organization. The four-office commons recurs throughout; Appendix~\ref{MS:app:allocation} collects its preservation obligations. Appendix~\ref{MS:app:statement-index} provides a linked theorem and proposition index. The source package records the scope and evidence of the final review.

The integration culminates in Sections~\ref{MS:sec:joint-enforcement} and~\ref{MS:sec:joint-synthesis}. The former replaces separate one-deviation monitors with one common enforcement policy. The latter combines permissions, observations, bounded obligations, and resources in a single continuation problem. It gives both a finite synthesis procedure and smaller witness certificates, together with an exact theorem for transporting the resulting strategies through an institutional abstraction. These constructions have classical optimization and game-theoretic foundations. Their role here is to expose when individually valid institutional claims cannot share one implementation.

Two analytical comparisons show why persistence is not enough. The cascade model in Section~\ref{MS:sec:cascades} can stabilize an incorrect conclusion even when every participant is rational under its information. The staffing model in Section~\ref{MS:sec:specialization} can give each role a backup without leaving enough agents to cover the roles together. These are not failures of the general idea of a society. They identify which information, capacity, and temporal distinctions its description must retain.

\subsection{Accountable continuation}
A common policy is still incomplete when its boundary omits an affected party, its audit is not performed, or its certificate is used after the supporting assumptions change. Chapter~\ref{MS:book:chapter-97} places these issues inside the realization model. It derives a sharp privacy envelope for bounded enforcement, gives an exact finite criterion for maintaining an available authorized ending, and states a guarded rule for changing certificate modes. The same resource accounting also applies to effective appeals. This makes accountability part of the theory's common-policy question rather than an assurance stated outside it.

\section{Realization, formation, and advantage}\label{MS:sec:scope}
Three questions organize the analysis. \emph{Realizability} asks whether a specified organization can be implemented. \emph{Formation} asks whether interactions select, establish, or revise its structure. \emph{Advantage} asks whether it improves a declared outcome against a constrained comparison class. Constructing a system answers none of the latter questions automatically. An executable rule can be externally imposed, locally adopted, or maintained by strategic agreement. These are different histories of the same possible implementation.

Every organizational state in this paper remains a function of a complete computational realization. If $Z_t=\pi(X_t)$, then $Z_t$ is not information outside $X_t$. The purpose of a higher-level description is to preserve relevant operations with less descriptive complexity. An institutional description must additionally preserve who may act, which commitments remain outstanding, and what counts as an authorized amendment. These are requirements on the abstraction, not a new substance added to the machine.

The mathematical setting is finite unless a section explicitly states otherwise. An unbounded lifetime is a temporal horizon, not necessarily an unbounded state space. A growing archive is treated as a sequence of finite representations, not silently as one finite system. Sources for classical mathematical ingredients are identified in the text. The proofs establish conditional statements in the declared models; they do not establish that an unspecified deployed population instantiates them.


\section{A running example: the four-office computational commons}\label{MS:sec:commons}
Consider four continuing institutional positions, denoted $A,B,C,D$, that share four indivisible execution slots per epoch. An office is not a particular model instance. Its current occupant may be replaced through a specified handover, while its registered responsibilities continue. The environment supplies the four slots and an append-only event mechanism with atomic commits. These are explicit infrastructure assumptions, not consequences of the agents' cooperation. Section~\ref{MS:sec:social-time} later replaces the atomic event assumption with local replicas and shows which stronger commitments then need an additional realization.

Each office has a private task queue and may submit a demand report. A public record contains the current allocation rule, valid credentials, accepted urgency certificates, outstanding service obligations, and the amendment history. Slots have unique names $(e,k)$, where $e$ is the epoch and $k\in\{1,2,3,4\}$. A grant consumes one available slot token. Copying a message that describes a grant does not create an additional token. The resource invariant is therefore
\[
\sum_{i\in\{A,B,C,D\}}q_i(e)\leq4,
\]
with no slot name assigned to two offices. In the worked trace all four offices demand their ordinary slot, so the allocations below use all four units.

\subsection{Two rules and an inherited obligation}
Under rule $R$, each office receives one slot. Rule $U$ permits an urgency exchange: when no previous service debt is outstanding, one office $j$ with an accepted urgency certificate may receive two slots if another office $k\ne j$ explicitly volunteers its slot. The other two offices each receive one. The exchange creates a registered obligation
\[
o=(\mathsf{id},\mathsf{debtor}=j,\mathsf{creditor}=k,
\mathsf{quantity}=1,\mathsf{due}=e+1).
\]
At the next epoch the creditor receives two slots and the debtor none, with one slot for each remaining office. This repayment clause has priority over both $R$ and $U$. No new exchange is admitted while a debt remains outstanding. The model thus permits at most one service debt and guarantees its next-epoch settlement conditional on continued infrastructure and execution of the admitted rule. It does not guarantee that every urgent request can be accepted.

Amendments between $R$ and $U$ require three of the four registered offices to approve. This threshold is part of the supplied constitutional envelope. Every amendment must preserve already registered service debts. A model replacement transfers an office only through an authorized registry transition that retires the predecessor's active credential and links the successor to the same outstanding obligations. Neither a model fork nor an imported archive can perform that transition by itself.

\begin{center}
\begin{tabularx}{\linewidth}{@{}cP{.25\linewidth}P{.19\linewidth}X@{}}
\toprule
Epoch & Admitted event & Allocation $(A,B,C,D)$ & Obligation record\\
\midrule
0 & Ordinary rule $R$ & $(1,1,1,1)$ & Empty.\\[.35em]
1 & Three offices adopt $U$; $A$ requests urgency and $B$ consents & $(2,0,1,1)$ & $A$ owes $B$ one slot at epoch 2.\\[.35em]
2 & $A$ is replaced by an authorized successor; repayment takes precedence & $(0,2,1,1)$ & The same obligation is fulfilled, not erased.\\[.35em]
3 & A proposed external grant would duplicate a slot token & $(1,1,1,1)$ & The proposal is rejected; no new entitlement is created.\\
\bottomrule
\end{tabularx}
\end{center}

These transitions are specified rather than discovered. The example is a miniature institutional realization, not a demonstration of autonomous formation. Its purpose is to make several different mathematical obligations visible in one system.

\subsection{What an institutional summary must retain}
The name of the current rule is insufficient state. Two records may both say $R$ while only one contains a debt due this epoch. Their correct next allocations differ. A summary that keeps the next allocation but omits the amendment threshold can also be inadequate: the same three votes may authorize different future rules under different constitutions. Likewise, identical text with different active credentials can confer different decision rights.

The public ledger is not an additional, society-wide preference function. Offices may disagree about whether $R$ or $U$ is desirable, about how urgency should be certified, or about who should pay the cost of maintaining the ledger. An atomic allocation mechanism can enforce the resource invariant without settling those disagreements. Private demand, operative authority, and the shared event record are distinct parts of the system.

We return to this commons when considering institution formation, governance-preserving summaries, deadline monitors, blocking coalitions, replacement, and federation. Appendix~\ref{MS:sec:worked} retains a smaller algebraic example that isolates reusable state from social interpretation.

The extensions later in the paper make the same commons sensitive to membership epochs, message types, notification histories, and the order of amendments. Its offices may preserve a shared allocation method while disagreeing about which new scheduling question deserves resources. These are analytical variations of one declared example, not reports of a population experiment.

\section{A controlled realization model}\label{MS:sec:model}
For the controlled realization results, assume finite nonempty sets of microstates $\MSX$, external controls $\MSU$, and reported outcomes $\MSY$. The finite assumption keeps measurability and regular-conditional-probability issues out of the exposition. Continuous-state generalizations require additional hypotheses.

A realization has transition kernel
\begin{equation}\label{MS:eq:micro}
K_u(x,x')=\MSProb(X_{t+1}=x'\mid X_t=x,U_t=u),
\end{equation}
and output kernel $Q_u(x,y)$. Both are row-stochastic. In this presentation, the current report and next state are sampled conditionally independently given $(x,u)$. A dependent report/transition mechanism can instead be specified by a joint kernel on $(x',y)$; the corresponding aggregation test must preserve that joint law, not just its two marginals. A deterministic implementation is the special case $K_u(x,\cdot)=\delta_{F_u(x)}$. The controls can encode tasks, supplied evidence, and permitted interventions. Scheduling and environment dynamics must be included in $X_t$ or in the stated control model rather than silently omitted.

One possible factorization of a microstate is
\[
X_t=(S_{1,t},\ldots,S_{n,t},M_t,C_t,A_t,\Theta_t),
\]
where $S_i$ is local process state, $M$ durable storage, $C$ messages and channel state, $A$ artifacts, and $\Theta$ the current configuration, including executable rules. These categories describe a storage and access organization. They need not be disjoint physical objects.

\begin{MSdefinition}[Declared decomposition and access model]\label{MS:def:access}
A decomposition specifies components and maps $O_i:\MSX\to\mathcal O_i$ describing their accessible observations. A coalition $S$ has the permitted information $O_S(X)$, supplemented only by the transcripts, computation, and external queries explicitly allowed to that coalition.
\end{MSdefinition}

Access to one complete state at one time and access to an entire local history are different experiments. An agent denied a file may still possess its contents in its current context. A component replacement is genuinely fresh only if the chosen reset conditions remove the predecessor information that the comparison is intended to exclude.

\begin{MSdefinition}[Candidate macrostate]\label{MS:def:macrostate}
A candidate macrostate is a surjection $\pi:\MSX\to\MSZ$. For a distribution $\mu$ on $\MSX$, its pushforward is
\[
(\MSpush\mu)(z)=\sum_{x:\pi(x)=z}\mu(x).
\]
The fibre $\pi^{-1}(z)$ is the set of implementations represented by $z$.
\end{MSdefinition}

An organization is always evaluated relative to a \emph{contract}: permitted controls, target outputs, horizon, environment family, resource limits, and error tolerances. Let $\mathcal T$ denote this contract. Claims such as ``the group is irreducible'' or ``the identity survived'' are incomplete without it. A state sufficient for tomorrow's scheduling decision may be inadequate for a later audit of who authorized that decision.

\subsection{A relational certificate is a profile}
A useful reporting profile is
\[
\mathcal C_{\mathcal T}=(\varepsilon_{\rm pred},\varepsilon_{\rm dyn},\varepsilon_{\rm int},
\Delta_{\rm coal},R_{\rm repair},D_{\rm lineage},G_{\rm formation}).
\]
Its entries report predictive error, dynamic closure error, intervention mismatch, coalition performance gaps, repair behavior, lineage evidence, and evidence about formation. Some entries may be unidentified. They should not be replaced by zero or hidden inside a favorable weighted average.

This is a specification proposal, not a universal numerical consciousness or emergence scale. A system can be strong on one coordinate and weak on another. In particular, a communication-only system can solve a task perfectly while lacking durable shared state, and a persistent shared state can be entirely wrong about the task.

\section{Relational information: access, synergy, and usefulness}\label{MS:sec:information}
Let $A,B$ be independent fair bits and $Y=A\oplus B$. Then
\[
I(Y;A)=I(Y;B)=0,\qquad I(Y;A,B)=1.
\]
This elementary example shows that joint observation may reveal information absent from each marginal observation. It does not show that the information is absent from the complete system, that a mysterious carrier has appeared, or that a protocol with ordinary messages cannot recover it.

A more operational measure uses decision loss. Given permitted information $O$ and loss $\ell$, define the Bayes risk
\[
\mathcal R_\ell(O)=\inf_f\MSE\bigl[\ell(f(O),Y)\bigr].
\]
For a nonempty allowed family $\mathcal A$ of proper coalitions, the gap
\begin{equation}\label{MS:eq:coalition}
\Delta_{\mathcal A}=\min_{S\in\mathcal A}\mathcal R_\ell(O_S)
-\mathcal R_\ell(O_V)
\end{equation}
expresses an advantage over every coalition in that family. A strictly positive gap is relative to the target distribution, observation maps, and coalition family. It is not equivalent to any one general-purpose partial-information-decomposition statistic~\cite{MS:williams}.

The unrestricted Bayes risk is an information limit. A neural decoder's measured risk is a performance result for that decoder. Failure of a linear probe does not prove absence of information, and success of a probe does not prove that the information is used by the system. Both observations can be useful, but they answer different questions.

\subsection{There is no workspace versus messages dichotomy}
A recurrent message may carry the same sufficient state as a central buffer. A complete message log may act as external memory. A distributed collection of local caches may instantiate a third implementation. Therefore the strongest legitimate claim is often that some \emph{persistent integration} is necessary under declared access constraints, not that one particular memory address or latent subspace is uniquely necessary.

Three measurements should be kept separate: information available to an observer, information actually used for a decision, and information retained across time. A one-time aggregate can have high joint information with no persistence. A durable archive can retain information without affecting behavior. A recurrent accumulator can be causally necessary while using an entirely conventional communication channel.

\subsection{Task information is not organizational identity}
A file containing the answer to a test makes a fresh agent more accurate than a blank file. That establishes content dependence. To support organizational continuity, the inherited material must preserve a richer, predeclared pattern, such as update behavior, correction policy, obligations, and responsiveness to counterfactual tasks. Even then, behavioral continuation and historical identity remain distinct, as Section~\ref{MS:sec:identity} explains.

\section{Exact and approximate macro-dynamics}\label{MS:sec:closure}
A macrostate deserves an autonomous transition law only when the discarded microstate distinctions are irrelevant in the required sense. This is the first mathematical obstacle hidden by the phrase ``the group behaves like one agent.''

\Needspace{8\baselineskip}
\begin{MSproposition}[Deterministic quotient criterion]\label{MS:prop:det}
For maps $F_u:\MSX\to\MSX$ and a surjection $\pi:\MSX\to\MSZ$, maps $\overline F_u:\MSZ\to\MSZ$ satisfying $\pi F_u=\overline F_u\pi$ exist if and only if
\[
\pi(x)=\pi(x')\quad\Longrightarrow\quad
\pi(F_u(x))=\pi(F_u(x'))
\]
for every $u,x,x'$. When they exist, they are unique.
\end{MSproposition}
\begin{proof}
Necessity follows by applying $\overline F_u$ to the common value $\pi(x)$. For sufficiency, set $\overline F_u(z)=\pi(F_u(x))$ for any $x$ in the fibre of $z$. The condition makes this independent of the representative, and surjectivity gives uniqueness.
\end{proof}

This is a congruence condition, not a consequence of persistence or frequent communication. A summary can be accurate for the current output and still fail it.

\Needspace{8\baselineskip}
\begin{MSproposition}[Controlled lumpability]\label{MS:prop:lump}
A common kernel $\overline K_u$ represents the projected one-step laws from every microstate if and only if, for all $z'$ and all $x,x'$ with $\pi(x)=\pi(x')$,
\begin{equation}\label{MS:eq:lump}
\sum_{w:\pi(w)=z'}K_u(x,w)
=\sum_{w:\pi(w)=z'}K_u(x',w).
\end{equation}
Under this condition, every macro-history policy $\alpha_t(u\mid Z_{0:t})$ induces the same projected process as the macro-model with initial law $\MSpush\mu_0$. For outputs to agree as well, require $Q_u(x,\cdot)$ to be constant on each fibre.
\end{MSproposition}
\begin{proof}
The common fibre value in~\eqref{MS:eq:lump} defines $\overline K_u(z,z')$. Conditional on any macro-history and selected control, the microstate has some distribution supported on the current fibre. Averaging identical projected rows leaves that common row unchanged. Induction gives equality of the macro-history laws. Conversely, start from the point masses at $x$ and $x'$ and apply the same fixed control. A common one-step macro-kernel forces~\eqref{MS:eq:lump}. The output statement follows identically.
\end{proof}

This finite result is a standard state-aggregation principle~\cite{MS:givan}. The policy qualification matters: if a controller secretly reads the discarded microstate to choose its actions, then the proposed macrostate may no longer describe the closed-loop process.

\begin{MScounterexample}[Deterministic microstate, non-Markov macrostate]\label{MS:ex:cycle}
Let $X_{t+1}=X_t+1\pmod 3$ with uniform initial distribution. Set $\pi(0)=\pi(1)=a$ and $\pi(2)=b$. Stationarity gives $\MSProb(Z_{t+1}=b\mid Z_t=a)=1/2$. Yet
\[
\MSProb(Z_{t+1}=b\mid Z_t=a,Z_{t-1}=b)=0,
\]
whereas the same probability conditional on $Z_t=a,Z_{t-1}=a$ is $1$. Thus the apparent randomness comes with memory: the one-step marginal does not define a valid autonomous first-order Markov model.
\end{MScounterexample}

A stochastic macro-description is therefore not automatic. It depends on an initial ensemble and on whether the retained state captures the relevant history. One may enlarge $Z_t$ with memory, use a partially observed model, or accept a measured approximation error.

\subsection{Finite-horizon error accounting}
For finite distributions, $\MSTV(\mu,\nu)=\tfrac12\sum_z|\mu(z)-\nu(z)|$. Suppose candidate kernels satisfy
\begin{equation}\label{MS:eq:eps}
\sup_{x,u}\MSTV\bigl(\MSpush K_u(x,\cdot),\overline K_u(\pi(x),\cdot)\bigr)
\leq\varepsilon.
\end{equation}
This uniform condition is stronger than a small average error on the trajectories seen during training.

\Needspace{8\baselineskip}
\begin{MSproposition}[Approximate closure over time]\label{MS:prop:time}
Under~\eqref{MS:eq:eps}, let the actual and reduced processes use the same predeclared open-loop control sequence. If their initial macro-laws are within $e_0$ in total variation, then at time $t$ they are within
\[
e_t\leq\min\{1,e_0+t\varepsilon\}.
\]
If each $\overline K_u$ contracts total variation by at most $\kappa<1$, then
\begin{equation}\label{MS:eq:contract}
e_t\leq\kappa^t e_0+\varepsilon\frac{1-\kappa^t}{1-\kappa}.
\end{equation}
With matched macro-history policies, the weaker path-law bound $\min\{1,e_0+t\varepsilon\}$ also holds, but the contraction claim needs a separate closed-loop contraction assumption.
\end{MSproposition}
\begin{proof}
Averaging~\eqref{MS:eq:eps} over the actual microstate law shows that its next projected law is within $\varepsilon$ of the current projected law evolved by $\overline K_u$. Kernel nonexpansiveness gives $e_{t+1}\leq\varepsilon+e_t$. Under contraction it gives $e_{t+1}\leq\varepsilon+\kappa e_t$; iterate. For history policies, couple the histories while they agree. The controls can then be coupled to agree, and the next macrostate fails to couple with probability at most $\varepsilon$. A union bound gives the path-law result.
\end{proof}

The distinction matters for long-lived systems. A small local abstraction error is not automatically harmless over millions of updates. Contraction can prevent indefinite accumulation, but it can also suppress useful distinctions. A useful model must retain the target information as well as remain dynamically stable.

\subsection{Predictive closure has a target and a distribution}
For a fixed future target $Y^+$ and data-generating policy, a predictive statistic satisfies
\[
I(Y^+;\MSHist_t\mid Z_t)=0.
\]
Approximate versions bound this residual information or, more directly, the excess predictive loss on held-out data. This is related to predictive-state constructions~\cite{MS:shalizi}. It is not equivalent to uniform controlled lumpability: it may hold for one observational distribution while failing after intervention or under a new policy.

\section{Causal interventions require implementations}\label{MS:sec:causal}
If $Z=\pi(X)$, one cannot generally change $Z$ while literally holding $X$ fixed. The expression $do(Z=z')$ needs a physical or computational implementation. Otherwise it names a desired effect, not an intervention.

For a feasible current microstate $x$ and target $z'$, let $\Lambda(x,z')$ be a declared nonempty family of admissible reset distributions on $\MSX$, each supported on $\pi^{-1}(z')$. Infeasible targets are excluded rather than assigned a fictitious reset. A reset might update a shared record, transplant a latent component, or reconstruct distributed shares. Its permissions, collateral effects, and retained private state belong in the intervention specification.

Given a future outcome experiment, write $P_\lambda$ for the resulting law of $Y^+$ under reset $\lambda$. Define implementation ambiguity by
\[
\eta(x,z')=\sup_{\lambda,\lambda'\in\Lambda(x,z')}
\MSTV(P_\lambda,P_{\lambda'}).
\]
A small ambiguity means different allowed implementations of the same macro-assignment have similar consequences. Large ambiguity means the macro-intervention is not sufficiently specified.

\Needspace{8\baselineskip}
\begin{MSproposition}[Implementation dependence]\label{MS:prop:lift}
If every $P_\lambda$ is within $\varepsilon$ of a proposed macro-intervention law $P^\star_{z'}$, then $\eta(x,z')\leq2\varepsilon$. Conversely, if two admissible implementations are distance $\eta$ apart, no single law can approximate both within less than $\eta/2$.
\end{MSproposition}
\begin{proof}
Both claims follow from the triangle inequality for total variation.
\end{proof}

\begin{MSexample}[Two meanings of the same macro-reset]
Let $X=(z,h)$ and let the next outcome be $Y=z\oplus h$. The macrostate retains only $z$. Resetting to $(z',0)$ and resetting to $(z',1)$ both implement $Z=z'$, but their outcome laws are disjoint. A causal model of $Z$ alone must either restrict the reset convention or acknowledge the missing variable.
\end{MSexample}

This is why causal abstraction requires compatibility between interventions, not only alignment of representations~\cite{MS:rubenstein,MS:beckers}. Removing a carrier can disable unrelated functions. A well-defined semantic intervention must specify the intended downstream effects, the admitted implementations, and which other functions are required to remain invariant.

\subsection{Three distinct counterfactual obligations}
A \emph{necessity counterfactual} excludes a carrier and determines whether the function can still be realized using the remaining authorized information. A \emph{content counterfactual} substitutes a compatible alternative state with a specified prediction. A \emph{mechanism counterfactual} changes the update or access rule while preserving information as closely as permitted. No one obligation implies the other two.

A causal comparison must define what information survives any permitted replacement, including copies and caches. Interventions may have collateral effects on other operations. A valid theoretical attribution therefore requires the declared intervention to distinguish its intended semantic change from an arbitrary loss of general capability.

\subsection{Intervention does not mean unrestricted control}
An external evaluator's ability to overwrite a state does not establish that the system can write that state itself. Bidirectional workspace access requires permitted routes from both participating components. Institutional agency requires an authorized decision procedure that actually changes behavior. These are stronger claims than external post hoc manipulability.

\section{Identity, equivalence, and causal lineage}\label{MS:sec:identity}
One might define a computational individual as an equivalence class of implementations. That definition is useful only after specifying the equivalence relation. Similar appearance, nearby behavior, common descent, and identical authorization are not interchangeable notions.

For a family $\mathcal T$ of finite continuation experiments, define exact behavioral equivalence by
\[
x\equiv_{\mathcal T}x'\quad\Longleftrightarrow\quad
\MSLaw_x(Y^+\mid\tau)=\MSLaw_{x'}(Y^+\mid\tau)
\quad\text{for every }\tau\in\mathcal T.
\]
Equality of these indexed laws is reflexive, symmetric, and transitive. It therefore defines a quotient. This equivalence is relative to the chosen experiments. It does not by itself equal bisimulation, and it says nothing about untested interventions or provenance.

For approximate comparison, use the pseudometric
\[
d_{\mathcal T}(x,x')=\sup_{\tau\in\mathcal T}
\MSTV\bigl(\MSLaw_x(Y^+\mid\tau),\MSLaw_{x'}(Y^+\mid\tau)\bigr).
\]
A finite test suite may make distinct implementations distance zero. That is indistinguishability under those tests, not universal sameness.

\begin{MScounterexample}[Tolerance is not an equivalence relation]\label{MS:ex:identity}
Consider three implementations producing a Bernoulli output with parameters $0$, $0.4$, and $0.8$. Their pairwise total-variation distances are $0.4$, $0.4$, and $0.8$. With tolerance $0.5$, the first resembles the second and the second resembles the third, while the first does not resemble the third.
\end{MScounterexample}

Approximate continuity should therefore be represented by distances, error budgets, or explicitly chosen clusters, not by silently forming equivalence classes. Along a replacement chain, the triangle inequality permits drift to accumulate. Small adjacent changes do not establish that the final system remains close to the initial one.

\subsection{A four-part identity record}
A practical identity record should distinguish \emph{behavioral equivalence}, \emph{causal continuity}, \emph{historical provenance}, and \emph{authorization}. Behavioral equivalence describes what the system does under a contract. Causal continuity describes whether predecessor-specific state actually influenced the successor. Historical provenance records the asserted and authenticated path of transfer. Authorization specifies which permissions or obligations transfer.

Two independent constructions can be behaviorally identical with no shared episode-specific lineage. Two descendants can share lineage while behaving differently. A signed lineage record authenticates a statement by a key holder, not the truth of every event described in it. Transferring a complete memory does not automatically transfer permissions.

\Needspace{8\baselineskip}
\begin{MSproposition}[Behavior alone does not identify implementation]\label{MS:prop:identifiability}
If two realizations induce the same distribution over every transcript available to an evaluator's permitted tests, then every decision rule based only on those transcripts has the same output distribution under the two realizations. With equal prior probability, no such rule can distinguish them with accuracy above $1/2$.
\end{MSproposition}
\begin{proof}
The evaluator's decision is a possibly randomized function of the transcript. Applying the same stochastic function to equal transcript laws gives equal decision laws. Under equal priors, any preference for one realization is offset by the identical preference when the other is present.
\end{proof}

This elementary identification limit explains why observation-only comparisons cannot exclude all centralized implementations. The remedy is not stronger rhetoric about emergence. It is instrumentation, trusted execution assumptions, intervention, or a narrower claim.

\subsection{Forks, merges, and obligations}
A fork creates multiple descendants, not an automatic duplication of every scarce entitlement. The state may branch while a resource allocation remains singular. Likewise, merging records is not equivalent to merging policies, latent representations, model weights, or obligations. Each requires a different compatibility test.

An obligation-aware inheritance record can contain the parent identifiers, transferred state description, unresolved conflicts, permissions explicitly granted, and obligations explicitly accepted. It should retain competing accounts rather than produce an apparently seamless but false history. Such records are an engineering proposal for accountable continuity, not a legal theory of personhood.

\section{Memory, reconstruction, and self-maintenance}\label{MS:sec:memory}
Persistence has at least three meanings. \emph{Storage persistence} means that records survive. \emph{Functional persistence} means that relevant behavior survives a stated sequence of events. \emph{Organizational persistence} means that a declared pattern of operations, commitments, and corrective procedures survives admissible changes. None requires the current informational content to remain constant.

An archive can be perfectly stable and behaviorally inert. A highly adaptive system can preserve its update method while replacing almost every factual belief. A frozen false belief can persist without constituting healthy organization. Persistence therefore needs both a survival criterion and a criterion for what ought to change when evidence changes.

\subsection{Memory lower bounds from distinguishable histories}
Fix a deterministic exact predictor for a family of future input continuations. Call two past histories equivalent when they require the same output sequence for every permitted continuation. Let $N$ be the number of equivalence classes among the histories under consideration.

\Needspace{8\baselineskip}
\begin{MSproposition}[Exact finite-memory lower bound]\label{MS:prop:memory}
Every deterministic summary that supports exact continuation on these histories requires at least $N$ distinct states, hence at least $\lceil\log_2N\rceil$ fixed-length bits. If the permitted continuation family is closed under prefixing and the history equivalence is a transition congruence, the quotient classes themselves provide an exact state representation.
\end{MSproposition}
\begin{proof}
If a summary assigns the same state to two inequivalent histories, some permitted continuation requires different outputs from that common state, contradicting determinism and exactness. This proves the lower bound by the pigeonhole principle. Under the congruence assumption, appending an input determines the next equivalence class independently of the representative, giving a well-defined quotient update.
\end{proof}

The result formalizes why the two-bit state in Section~\ref{MS:sec:worked} cannot be replaced by its one-bit current answer. It is also a limit on optimistic claims about arbitrarily compressing an entire scientific or institutional history. Compression can discard details only relative to the future questions and interventions it promises to support~\cite{MS:tishby}.

\subsection{Erasure resilience is not arbitrary error correction}
Let an organizational state $z$ range over a finite set with at least two elements and be encoded as $E(z)=(E_1(z),\ldots,E_n(z))$. Define the block distance
\[
d_E=\min_{z\ne z'}\MSabs{\{i:E_i(z)\ne E_i(z')\}}.
\]

\Needspace{8\baselineskip}
\begin{MSproposition}[Exact reconstruction after erasure]\label{MS:prop:erasure}
For an integer $0\leq k\leq n$, every set of at most $k$ erased components can be tolerated with exact reconstruction if and only if $d_E>k$.
\end{MSproposition}
\begin{proof}
If $d_E>k$, two distinct codewords disagree outside any erased set of size at most $k$, so the surviving observations identify the state uniquely. If $d_E\leq k$, erase the coordinates on which a minimizing pair differs. The surviving observations then coincide, making reconstruction impossible.
\end{proof}

For $z=(p,q)\in\MSFtwo^2$, the encoding $(p,q,p\oplus q)$ has distance two. Any two components reconstruct the state, and no one component reconstructs all of it. It tolerates one \emph{known erasure}. It does not correct an arbitrary one-component corruption at an unknown location. Confusing these fault models overstates robustness.

The coding argument supplies a reconstructibility condition, not a guarantee that an operating system has an authorized, timely, and successful repair procedure.

\subsection{Repair needs a causal and temporal test}
Let $e_t\geq0$ measure deviation from an acceptable organization after accounting for legitimate changes in task state. Suppose a repair controller ensures
\[
e_{t+1}\leq a_t e_t+b_t,
\]
where $a_t$ measures retained deviation and $b_t$ covers new disturbance and model mismatch.

\Needspace{8\baselineskip}
\begin{MSproposition}[Repair budget]\label{MS:prop:repair}
If $0\leq a_t\leq\bar a<1$ and $0\leq b_t\leq\bar b$, then
\[
e_t\leq \bar a^t e_0+\bar b\frac{1-\bar a^t}{1-\bar a}.
\]
More generally, for nonnegative $a_t,b_t$ satisfying the stated recurrence,
\[
e_T\leq e_0\prod_{t=0}^{T-1}a_t+
\sum_{j=0}^{T-1}b_j\prod_{t=j+1}^{T-1}a_t.
\]
\end{MSproposition}
\begin{proof}
Repeated substitution gives the product--sum formula. Uniform bounds reduce the sum to a geometric series.
\end{proof}

The inequality is a conditional stability statement. It does not prove that the implemented monitor measures the right error or that its repair action causes the desired contraction. Those assumptions require interventions on the monitor and repair pathways.

\begin{MScounterexample}[Perfect coherence by erasing everything]\label{MS:ex:zerocollapse}
Let acceptable states occupy a linear subspace $W$, and measure only distance to $W$. The update $F(x)=0$ attains zero coherence error immediately but destroys every distinction among encoded states. A coherence measure alone cannot certify memory preservation or identity.
\end{MScounterexample}

A maintenance claim consequently needs a non-collapse condition, continued task information, and comparison with a disabled-repair control under matched disturbance. Recovery time, repair cost, and permanent loss should be reported separately. Very strong repair that restores yesterday's policy despite valid new evidence is rigidity, not automatically health.

\section{Hierarchy, gluing, and the limits of flattening}\label{MS:sec:hierarchy}
A lower-level organization can serve as a component at the next level when it exports a sufficient interface. That interface must describe not only a state but also the operations that are permitted on it, their outcome semantics, and their implementation error. A useful interface certificate is
\[
\mathsf{Int}=(\MSZ,\MSU,\overline K,\overline Q,\Lambda,
\varepsilon_{\rm dyn},\varepsilon_{\rm int},\mathsf{permissions}).
\]
This tuple is a design proposal. It prevents a higher-level controller from treating a reliable read-only summary as though it were an independently writable agent.

\subsection{Composition of dynamic abstractions}
Suppose $\pi:\MSX\to\MSZ$ and $\psi:\MSZ\to\mathcal W$ are surjections and, for the same control family,
\[
\MSTV(\MSpush K_u(x,\cdot),\overline K_u(\pi x,\cdot))\leq\varepsilon_1,
\]
\[
\MSTV(\psi_{\#}\overline K_u(z,\cdot),\widehat K_u(\psi z,\cdot))\leq\varepsilon_2.
\]

\Needspace{8\baselineskip}
\begin{MSproposition}[Compositional one-step closure]\label{MS:prop:composition}
The composite abstraction $\psi\pi$ satisfies
\[
\MSTV\bigl((\psi\pi)_{\#}K_u(x,\cdot),
\widehat K_u(\psi\pi x,\cdot)\bigr)
\leq\varepsilon_1+\varepsilon_2.
\]
For $L$ levels, the one-step defect is at most $\sum_{j=1}^L\varepsilon_j$. In particular, exact abstractions compose exactly.
\end{MSproposition}
\begin{proof}
Insert $\psi_{\#}\overline K_u(\pi x,\cdot)$, use the triangle inequality, and use nonexpansiveness of total variation under deterministic pushforward. Induction gives the multilevel result.
\end{proof}

The same argument applies separately to output kernels. Combining this proposition with Proposition~\ref{MS:prop:time} yields a finite-horizon bound of $e_0+t\sum_j\varepsilon_j$, capped at one, under the corresponding controls. This elementary estimate belongs to the broader subject of compositional abstraction error~\cite{MS:rischel}. It does not establish global causal closure unless the intervention implementations also compose compatibly.

\subsection{Structural gluing is a different question}
Suppose each subsystem $i$ has local state space $W_i$ and each overlap has restriction maps $r_{ij}^i$ and $r_{ij}^j$. The compatible families form
\[
W_{\rm glue}=\{(w_i):r_{ij}^i(w_i)=r_{ij}^j(w_j)\text{ on every declared overlap}\}.
\]
For an edge-complete cover of a fixed diagram, compatible local sections can give another presentation of its global sections. These compatibility conditions follow the elementary logic of gluing local descriptions into a shared global state.

Compatibility alone does not show that the resulting state predicts behavior, is used for decisions, or avoids a centralized implementation. A hierarchy may also discard essential overlap information. A faithful compressed interface must neither identify globally distinct relevant states nor admit apparent compatible states with no full-resolution realization.

There is a sharp conceptual distinction between \emph{descriptive flattening} and \emph{operational removal}. An isomorphism between a recursively assembled state and a global state does not license deleting intermediate processes while preserving latency, privacy, reliability, or resource cost.

\Needspace{11\baselineskip}
\subsection{Hierarchy is not extra computability}
\begin{MSproposition}[Simulation by an unrestricted controller]\label{MS:prop:flatten}
A finite, effectively specified hierarchy of computable components can be simulated by a controller that stores every component state, message, and scheduling variable and applies the same local transitions in the same order. With the same random tape, its externally visible transcript is identical.
\end{MSproposition}
\begin{proof}
Initialize the controller's records to the component states. At each scheduled step, apply the corresponding component transition to its stored inputs and update the associated state and message records. Induction preserves equality of all records and therefore of the external outputs. A shared random tape couples randomized transitions identically.
\end{proof}

This is an existence statement, not an efficiency claim. A meaningful superiority result must restrict the comparison class. A ``flat'' system allowed to simulate the entire hierarchy defeats any claim of advantage based only on the drawing of boxes. Restrictions may concern accessible information, local memory, permitted routing, wall-clock span, bandwidth, or allowed rounds.

\subsection{An honest complexity separation}
Consider $n=2^h$ input bits with parity as the target and two-input XOR gates. A balanced tree computes the target with $n-1$ gates and depth $h$. Any fan-in-two circuit output at depth $d$ can depend on at most $2^d$ input bits, so depth at least $\lceil\log_2n\rceil$ is necessary. A single sequential accumulator uses $n-1$ updates.

The gain is parallel depth, not reduced total work or super-Turing power. A sufficiently parallel controller can implement the same tree. The example shows how to state a useful architecture-relative result without inventing a general hierarchy-versus-flat separation.

A corresponding research question for agent systems is whether certified modular summaries reduce repeated communication or repair cost under explicit information and latency constraints. That question is substantive even if total computation remains classically simulable.

\section{Deterministic agents and macro-level uncertainty}\label{MS:sec:determinism}
There are three different sources of uncertainty: randomized transition rules, incomplete knowledge of the complete state, and implementation variability such as an unspecified schedule. They should not be collapsed into one statement that a model ``is probabilistic.'' A fixed pseudorandom generator and its state are part of a deterministic expanded implementation. A macro-observer who does not see that state may use a stochastic model.

For a generative model, a distribution over outputs does not itself require sampled behavior: a fixed decoding rule can choose a unique output. Conversely, a deterministic decoder does not guarantee a reproducible deployed system unless the surrounding computation, data, and scheduling assumptions are controlled. The finite three-cycle in Counterexample~\ref{MS:ex:cycle} makes the stronger point that a stochastic-looking projection may retain memory and fail Markov closure.

\Needspace{8\baselineskip}
\begin{MSproposition}[Symmetry lock for identical deterministic agents]\label{MS:prop:symmetry}
Suppose a deterministic joint update $F$ is equivariant under permutations of the agent labels and the initial joint state and all subsequent inputs are permutation-invariant. Then the joint state remains permutation-invariant at every time.
\end{MSproposition}
\begin{proof}
If $gX=X$ for every label permutation $g$, equivariance gives $gF(X)=F(gX)=F(X)$. Apply this at every step with invariant inputs.
\end{proof}

In this idealized setting, distinct roles require a symmetry breaker. Randomness is one possibility; different observations, histories, hardware constraints, or assigned addresses are others. Determinism is not an obstacle to organization. Exact symmetry can be an obstacle to differentiated organization.

\subsection{Where stochasticity helps and where it does not}
Exploratory search may benefit from varied proposals, while a checker should apply the same declared rule consistently. Randomized exploration followed by deterministic checking is one useful design, not a universal architecture. Random samples from identical models can still share systematic errors. Randomness is neither independence nor creativity by definition.

A fully asynchronous fault-tolerant setting also imposes limits independent of intelligence. The Fischer--Lynch--Paterson result shows that a deterministic consensus protocol cannot guarantee termination in every admissible execution with even one crash failure in its stated asynchronous model~\cite{MS:flp}. This does not mean consensus is never achieved or that machine institutions are impossible. It means liveness claims must state timing, failure, and randomization assumptions.

\subsection{A transition is not established by a sharp-looking curve}
A phase-transition analogy requires explicit mathematical assumptions. A linearization with an eigenvalue above one shows local instability of the point being linearized; it does not prove the existence of a useful stable attractor. The dynamics may diverge, cycle, or collapse elsewhere. Nonlinear terms, constraints, and the relevant observables matter.

For finite adaptive systems, an apparent onset is a property of a specified learning dynamics, environment family, timescale, and measured observable. Fitting a sigmoid to an arbitrary plot does not establish a bifurcation. A genuine dynamical bifurcation requires a specified family of transformations, and a thermodynamic phase-transition claim requires an appropriate limiting construction. There is no universal numerical coupling threshold in this theory.


\chapter{Institutions and collective choice}\label{MS:book:chapter-94}

\section{From aggregate computation to society}\label{MS:sec:society}

The word \emph{society} is easily applied to a collection of agents that exchange messages. That usage conceals distinctions important to both mathematics and social theory. A cluster of predictors controlled by one scheduler may act as a coherent unit while having no independent decision centers. A set of autonomous services may compete for resources yet lack shared institutions. Conversely, a population may sustain institutions, disagreements, and an archive of precedents while possessing no coherent global policy. None of these arrangements is contradictory. They occupy different parts of the space of possible organizations.

A theory of machine societies should therefore begin from \emph{plurality}, not from the expectation that an entire population becomes one mind. A social system is organized not only by what it collectively computes but also by who can decide, who can contest a decision, how scarcity is handled, which records survive, and how the rules themselves change. These questions require semantics for resources, authority, and institutional transitions in addition to semantics for messages and tasks.

\begin{MSdefinition}[Society-like computational organization]\label{MS:def:society}
Relative to a declared environment and evaluation contract, a society-like computational organization is a population with (i) at least two persistent, distinguishable decision centers; (ii) recurrent interdependence among their permitted actions; (iii) operative and at least partly maintained rules for coordination, access, allocation, or dispute resolution; (iv) records or procedures that can outlast particular active components; and (v) an institutional state whose feasible alteration has consequences for subsequent behavior. A society-like population need not possess a common utility function or a unified policy.
\end{MSdefinition}

This is an \emph{operational convention for computational theory}, not a proposed universal definition for sociology. Persistence is relative to an interval and perturbation class. Distinct decision centers need not be identical to the number of model copies or hardware processes: two copies whose every consequential choice is fixed by a hidden controller are not independent decision centers under this definition. Likewise, the existence of a rule file does not satisfy (iii) or (v) unless the file affects permitted transitions.

\begin{MSdefinition}[Endogenous institutional formation]\label{MS:def:endogenous}
Let $D$ record the rules, affordances, and permissions fixed by the designer and let $I_t$ be the institutional state. Formation is \emph{endogenous relative to that family} when interactions among participating decision centers make a causal difference to the content, adoption, revision, or enforcement of $I_t$ beyond what is fixed by $D$. Its strength may be measured by an explicitly defined counterfactual effect on those institutional observables, not by an assumed universal emergence score. The definition does not equate endogenous formation with creation from no supplied substrate.
\end{MSdefinition}

The dependence on a family of counterfactuals matters. If all possible rules are hard-coded and an external schedule determines exactly which is active, institutional behavior may be useful without being endogenously selected. If the action space is supplied but the population selects and revises among feasible rules based on interaction history, there is a meaningful but bounded form of endogenous formation. Supplying a language in which a rule can be expressed does not determine which rule a population will establish.

\subsection{Societal plurality and collective agency are independent axes}

One axis concerns how many operative decision centers and institutions persist. Another concerns whether a useful society-wide state and control interface exists. Figure~\ref{MS:fig:axes} gives an interpretive map, not a classification produced by fitting data.

\begin{figure}[t]
\centering
\includegraphics[width=.87\linewidth]{sources/monographs/machine-societies/publication/figures/society_versus_agency.pdf}
\caption{Two independent organizational dimensions. The horizontal axis represents institutional plurality and persistence; the vertical axis represents the usefulness of a single global policy description. Labels are idealized types, not claims that any named system occupies the corresponding region.}\label{MS:fig:axes}
\end{figure}

A distributed scientific community is intelligible as socially organized despite disagreement on its research agenda. A controlled assembly line may be intelligible as one agent despite many processes. A future computational society could occupy either of the high-plurality regions. A society could coordinate for some tasks and remain polycentric for others: closure, when present, is relative to a control family. In particular, institutional interoperability does not logically imply that all members share the same preferences or that the population experiences one subjective perspective.

\subsection{The minimum social vocabulary}

We distinguish \emph{conventions}, \emph{roles}, \emph{institutions}, and \emph{cultures}. A convention is a regularity of mutually expected behavior, often one among multiple equilibria. A role is a conditional bundle of permitted operations and expectations, which may be fulfilled by different agents over time. An institution is a durable set of rules, records, authorities, and enforcement or revision procedures that changes the feasible or incentivized behavior of participants. Culture consists of transmitted methods, beliefs, conventions, norms, and interpretive practices that affect successor behavior. These categories overlap, but none reduces to a shared latent vector or a group-chat transcript.

The relevant literature already covers important parts of this terrain: conventions in games~\cite{MS:young,MS:shoham}, polycentric governance~\cite{MS:ostrom,MS:ostrombook}, normative multi-agent systems~\cite{MS:boella,MS:grossinorm}, population protocols~\cite{MS:angluin}, and interactions among language-model agents~\cite{MS:park,MS:ashery}. The present framework is a synthesis oriented around \emph{realization and preservation conditions}: what must survive, what must causally matter, what can be simulated by simpler alternatives, and where apparent compatibility breaks down. It does not present institutions or computational social conventions as newly discovered concepts.

\section{The institutional state and its realization}\label{MS:sec:institutional-state}

To formalize an operative institution, augment the microstate of Section~\ref{MS:sec:model} with a distinguished component $I_t$. In a finite setting we may describe it as
\begin{equation}\label{MS:eq:institutional-state}
I_t=(N_t,P_t,E_t,A_t,L_t),
\end{equation}
where $N$ specifies norms and role obligations, $P$ permissions and capabilities, $E$ enforcement and conflict procedures, $A$ admissible amendments, and $L$ records of decisions and their provenance. The components need not be physically stored together, and their representation is not unique. What matters is their effect on the transition law.

\begin{MSdefinition}[Institutional realization]\label{MS:def:institutional-realization}
An institutional realization comprises an environment state $x$, an institutional state $I$, a set of local actions $\mathcal A_i(x,I)$ for each decision center, a transition kernel $K_I$ on the complete state, and a family of admissible interventions on $I$ and its implementation. A written rule has \emph{operative force} only relative to changes in the feasible action sets, authorized transitions, incentives, or observable outcomes that it actually induces.
\end{MSdefinition}

The crucial word is \emph{induces}. A simulated constitution that does not affect any operation is ceremonial. A simple queue rule that agents can adopt, enforce, appeal, and change can have substantial institutional force without an elaborate natural-language charter.

For a family $\mathcal E$ of permitted environment setups, a horizon $T$, and feasible interventions $\lambda(I),\lambda(I')$ that differ institutionally while matching the intended other conditions, define the \emph{institutional effect profile}
\begin{equation}\label{MS:eq:inst-effect}
\Delta_{\mathrm{inst}}^{\mathcal E,T}(I,I')
=\sup_{e\in\mathcal E}\MSTV\bigl(\MSLaw_e(Y_{0:T}\mid\lambda(I)),\MSLaw_e(Y_{0:T}\mid\lambda(I'))\bigr).
\end{equation}
A zero value shows observational equivalence on the declared family, not universal inertness. An enforceable institution can also alter which actions are institutionally permitted while producing the same observed transcript on a particular task. Observations and rights must therefore be tracked separately. As with other macro-interventions, the implementation cannot be replaced by an unconstrained symbolic assignment~\cite{MS:rubenstein,MS:beckers}.

\subsection{Typed social relations}

A useful representation is a graph whose edges have different types. Information edges determine which messages may flow; resource edges specify transfers and scarcity; normative edges record permissions, obligations, and appeals; causal edges specify implemented state changes. Treating all four as one untyped ``connection'' erases important distinctions. An agent may truthfully observe an event while lacking authority to endorse a rule; another may hold an authorization token while lacking current evidence; a third may pay an operational cost without acquiring any right to amend policy.

\begin{figure}[t]
\centering
\includegraphics[width=.90\linewidth]{sources/monographs/machine-societies/publication/figures/institutional_layers.pdf}
\caption{An institutional architecture as a typed interaction system. Rules, shared records, resource procedures, and revision processes are organizational components rather than evidence of a single higher-level mind. The arrows denote possible bidirectional dependence; each route requires its own access and authorization semantics.}\label{MS:fig:layers}
\end{figure}

This typed view makes institutional ``identity'' more exact. Keeping the same policy text while transferring enforcement authority can change the institution even when its next output is unchanged. Keeping the same enforcer while deleting appeal rights changes it. Copying a rule without copying its authority and exception history changes it. In other words, the declaration $N$, operational implementation $K_I$, and legitimacy source must not be conflated.

\subsection{Rules are transformations, not only statements}

A norm can be presented operationally as a condition on histories and actions, $N_i(h,a)\in\{0,1\}$, interpreted as the system's declared permission for action $a$ by position $i$ after history $h$. This does not make the declared permission morally correct. Enforcement changes the action transition or cost structure; appeal changes what happens after an alleged violation; amendment changes the permissible norm itself.

The order of such operations can matter. Let $G$ set a permission bit to $1$ (grant) and $R$ set it to $0$ (revoke). Then $G\circ R$ yields $1$ and $R\circ G$ yields $0$ on every input. This simple noncommutativity has practical conceptual force: a log of individual events without their authorized order may fail to reconstruct the actual institutional state. Revision histories must preserve \emph{operations and order}, not only a multiset of decisions.

\subsection{The difference between compliance and legitimacy}

A proof that every enacted transition obeys $N$ establishes rule compliance relative to the representation, assuming implementation fidelity. It does not establish that the norm is fair, democratically adopted, safe for nonmembers, or compatible with human law. Computational correctness and normative authority are different kinds of judgment. A machine population can acquire stable internal expectations without gaining legal personhood, property rights, or permission to act outside its sandbox. The material substrate, owners, operators, and affected humans remain part of the surrounding system.


\section{Endogenous boundaries: membership, offices, and voting authority}\label{MS:sec:membership}
A society's population is not merely a parameter in its governance problem. Admission, departure, succession, and the allocation of decision rights can themselves be governed operations. An institution may preserve its written amendment threshold while changing who counts toward that threshold. Consequently, the boundary of the institution belongs in the state whose preservation is being claimed.

Let $\mathcal U$ be a finite universe of possible participant handles for the finite results. Write $M_t\subseteq\mathcal U$ for admitted members, $V_t\subseteq M_t$ for voting positions, and $\mathcal D$ for a declared set of effective decision centers. A map $c_t:V_t\to\mathcal D$ records which decision center controls each voting position under the model. This is an assumption about effective control, not a fact established by distinct names, keys, or model instances. A boundary state may be represented as
\[
B_t=(M_t,V_t,c_t,w_t,q_t,e_t),
\]
where $w_t$ assigns nonnegative voting weights, $q_t$ is an approval threshold, and $e_t$ is a configuration epoch. Admission and exit proposals determine
\[
M_{t+1}=(M_t\setminus E_t)\cup A_t,
\qquad E_t\subseteq M_t,\quad A_t\subseteq\mathcal U\setminus M_t,
\]
but this equation alone specifies neither their authorization nor the transfer of offices and outstanding obligations. Those conditions are part of the boundary transition relation. An unbounded supply of fresh handles requires an explicitly enlarged, nonfinite model.

\subsection{Copies, credentials, membership, and control are different objects}
An executing copy may lack a membership record. A member may lack a vote. Several voters may share a controller. A controller may have no formal vote yet control scheduling or the resource owner. Distinct credentials establish distinct credential records under a trust model, not independent interests or independent evidence. The established Sybil problem is one reason to keep these distinctions explicit~\cite{MS:douceur}.

As a simple arithmetic example, two controlled votes among five do not form a strict majority. If the rules admit two more votes controlled by the same center without changing anyone else's vote, the center then controls four of seven and has a majority. This is not an assertion that the admission can be authorized under the original rule. It shows that an admission rule, when it permits such a transition, can change substantive power without editing the ordinary voting formula. In the commons, adding a fifth \emph{office} is therefore different from replacing one office's occupant or running an additional copy to assist it.

\begin{MSproposition}[Weight-conserving refinement of controlled votes]\label{MS:prop:forkweight}
Fix a weighted threshold rule with threshold $q$. Replace a voting position $v$ of weight $w(v)$ by positions $v_1,\ldots,v_k$ controlled by the same decision center, with nonnegative weights satisfying $\sum_j w(v_j)=w(v)$. Leave all other weights and the threshold unchanged. Then every coalition of decision centers has the same total voting weight and the same approval status before and after the replacement.
\end{MSproposition}
\begin{proof}
A coalition not containing the controller of $v$ gains none of the replacement positions and loses none of its previous weight. A coalition containing that controller replaces the summand $w(v)$ by the equal sum $\sum_jw(v_j)$. Its comparison with $q$ is unchanged.
\end{proof}

The result concerns coalitions of the stipulated decision centers. It does not preserve the number of credentials, the cardinality of a smallest coalition of individual positions, or a rule that gives each newly named position a fresh unit weight. Nor does it authenticate the control map $c_t$. Weight conservation is one possible institutional constraint, not a general solution to identity or fair representation.

\subsection{Reconfiguration needs a handover, not just overlapping rosters}
Let an old roster be $\{a,b,c\}$ and a new roster $\{c,d,e\}$, each using strict majority. Although the rosters overlap, the old majority $\{a,b\}$ and new majority $\{d,e\}$ are disjoint. A shared member does not by itself ensure that every old decision certificate intersects every new one. Distributed reconfiguration protocols address this issue by specifying how old and new configurations participate in the handover; Raft's joint-consensus construction is a concrete precedent~\cite{MS:ongaro}.

Here is the limited combinatorial fact such reasoning uses. Fix one decision slot and a set of possible values. A certificate consists of endorsements by every member of an allowed quorum. Suppose each witness is constrained to endorse at most one value for that slot, across all configuration epochs considered.
\begin{MSproposition}[Intersecting witnesses exclude conflicting certificates]\label{MS:prop:handover}
If every two allowed quorums intersect, two certificates for distinct values of the same slot cannot both be formed under the stated witness constraint. If two allowed quorums are disjoint, the witness constraint alone does not exclude such conflicting certificates.
\end{MSproposition}
\begin{proof}
Conflicting certificates on intersecting quorums would require a common witness to endorse both values, contrary to the constraint. For disjoint quorums, assign one value to all witnesses of the first and a different value to all witnesses of the second. No witness endorses twice.
\end{proof}

A transition quorum requiring an old majority and a new majority intersects any old majority through the old roster and any new majority through the new roster. This fact is useful only with a protocol that transfers the relevant records and preserves the witness constraint. It is not a standalone consensus, fault-tolerance, or liveness theorem. Reused keys, forgotten locks, and inconsistent slot identities lie outside its assumptions.

For theoretical institutions, the same lesson is to version the electorate, interpreter, and decision slot together. A boundary transition must specify which configuration judges its authorization, when the new one becomes operative, and what evidence successors inherit. There is no need to assume that membership is fixed forever. There is a need to make its changes part of the constitutional semantics.

\section{Rule-changing computation and governance-preserving abstraction}\label{MS:sec:governance}
A persistent society does not merely choose actions under a rule. Some authorized actions change which future actions are admitted. Write a complete state as
\[
s=(x,I,\mathcal P,r,\lambda,c),
\]
where $x$ is operational state, $I$ institutional records, $\mathcal P$ the current population and office registry, $r$ resources, $\lambda$ representation or protocol configuration, and $c$ the current checking and amendment mechanism. These coordinates may overlap in storage. They identify semantic responsibilities rather than independent physical objects.

Let $D(s)\subseteq\mathcal A$ be the enabled action labels, partitioned into ordinary and amendment operations. A deterministic realization has a partial transition map $F(s,a)$ defined exactly when $a\in D(s)$. An amendment may change $I$, $c$, or $\lambda$ and thereby change $D(F(s,a))$. A stochastic realization instead uses a joint kernel on successor operating state and successor institutional state. We do not assume these two updates are conditionally independent. This formulation changes a rule \emph{inside} an explicitly declared computational envelope; it does not permit the mathematical semantics to rewrite themselves without representation in the state.

\subsection{A quotient must preserve what can be done}
A summary adequate for predicting the next grant can fail to preserve who may authorize the next amendment. To capture that distinction, let $\ell:S\to\Gamma$ record the institutional observables protected by the contract. These may include office ownership, resource balances, and a deadline monitor, but need not include every stored byte. The finite action vocabulary $\mathcal A$ has fixed meanings across the comparison.

\begin{MStheorem}[Exact governance quotient]\label{MS:thm:govquot}
Let $S$ be finite and nonempty and $\mathcal A$ finite, and let $\pi:S\to Z$ be surjective. There exist functions $\bar\ell$, $\bar D$, and partial transitions $\bar F$ satisfying
\[
\ell=\bar\ell\pi,\qquad D=\bar D\pi,\qquad
\pi F(s,a)=\bar F(\pi s,a)\quad(a\in D(s))
\]
if and only if whenever $\pi s=\pi t$,
\[
\ell(s)=\ell(t),\qquad D(s)=D(t),\qquad
\pi F(s,a)=\pi F(t,a)\quad\text{for all }a\in D(s).
\]
The quotient data are unique. From any two states in the same fibre, the legal finite action sequences and their corresponding label sequences are identical. This includes sequences mixing ordinary actions with amendments.
\end{MStheorem}
\begin{proof}
Necessity follows by factoring each quantity through $\pi$. For sufficiency define the quotient label, enabled set, and successor using any representative of a fibre. The three conditions make each definition independent of that representative. Surjectivity gives uniqueness. For trace preservation, induct on the length of an action sequence. Equal enabled sets admit the same first action, the successor condition keeps the resulting states in one fibre, and equal labels preserve the next observation. Repeating the argument proves the claim for every legal finite continuation.
\end{proof}

The theorem is a labeled deterministic congruence principle, extending the use of state aggregation to an explicit institutional contract~\cite{MS:givan}. It does not identify an observationally invisible permission with a permission that does not exist. A projection-only policy lifts by choosing the same action label at $s$ and $\pi s$, provided the implementing decision process is authorized to read that projection and issue the label. That proviso is essential: a mathematically available summary is not automatically an observation available to a decentralized controller.

\begin{MScounterexample}[Prediction without governance preservation]\label{MS:ex:govloss}
Let two states produce the same ordinary allocation. In one, three approvals enable amendment $a$; in the other, four are required and only three have been received. A summary containing only the allocation identifies the states, while their enabled sets differ. It predicts the current allocation exactly but admits no quotient preserving the amendment operation.
\end{MScounterexample}

\subsection{The smallest exact summary under a fixed contract}
The conditions yield a finite construction rather than only an existence criterion. Begin with the partition $\mathcal B_0$ that identifies exactly the states sharing $(\ell(s),D(s))$. At each refinement, identify $s,t$ only when they were previously in the same block and every action in their common enabled set leads to the same previous block:
\[
s\equiv_{k+1}t\quad\Longleftrightarrow\quad
s\equiv_kt\ \text{ and }\ F(s,a)\equiv_kF(t,a)
\ \text{for all }a\in D(s).
\]

\begin{MSproposition}[Coarsest finite governance summary]\label{MS:prop:govrefine}
The refinement stabilizes after at most $|S|-|\mathcal B_0|$ strict refinements. Its limiting partition satisfies Theorem~\ref{MS:thm:govquot} and is coarser than every partition satisfying that theorem for the same labels, actions, and transitions. It therefore has the minimum number of states among those exact quotients, up to relabeling.
\end{MSproposition}
\begin{proof}
Every strict refinement increases the number of blocks, which cannot exceed $|S|$. At a fixed point, states in one block have the same labels and enabled sets and their enabled successors remain in one block, so the quotient conditions hold. Any compatible partition refines $\mathcal B_0$. If it refines $\mathcal B_k$, compatible successors show that it also refines $\mathcal B_{k+1}$. Induction proves that it refines the stable partition. The latter is itself compatible, giving the minimum block count.
\end{proof}

The result is a contract-relative minimization statement, not a unique minimal representation of an entire society. Adding a new appeal operation or a new protected obligation can split blocks that were previously equivalent. In the commons, a summary that omits pending service debt must be refined before it can preserve repayment. This is the central computational price of self-amendment: preserving today's operations need not preserve an enlarged future operation family.

\section{Institutional language: messages that change social state}\label{MS:sec:language}
A common message format does not by itself establish a common institution. A statement, a request, a promise, an approval, and an appeal can contain the same descriptive proposition while having different consequences. The social semantics of agent communication has long emphasized publicly accountable interaction rather than assuming access to a sender's private beliefs~\cite{MS:singhacl}. Constitutive and procedural norms supply a related foundation~\cite{MS:boella,MS:grossinorm}. The addition here is to connect those semantics to the same realization and preservation conditions used for institutional state.

A message can be represented as a typed record
\[
m=(\mathsf{kind},\mathsf{issuer},\mathsf{epoch},\mathsf{payload},\mathsf{nonce}).
\]
Its operative meaning is not just its payload. For a complete institutional state $s$, an interpreter defines a partial transition $F_m(s)$ and an event record $\ell(s,m)$ when that transition is defined. Undefinedness represents rejection under the declared protocol, not a moral verdict on the message. The interpreter, authentication assumptions, role registry, and relevant history must be included in $s$ or fixed by the contract. In particular, an imported string cannot create authority simply by naming an authoritative role.

\subsection{Faithful translation preserves admissibility and effect}
Let institutions $A$ and $B$ have state spaces $S_A,S_B$, message sets $\mathcal M_A,\mathcal M_B$, partial interpreters $F^A,F^B$, and event records in a common semantic space $\mathcal E$. A state translator is $T:S_A\to S_B$. Allow a context-dependent message translator
\[
\tau:S_A\times\mathcal M_A\to\mathcal M_B.
\]
Context dependence matters for translated office names, configuration epochs, and reference records. The translation of a message sequence is constructed step by step from the evolving source state, not assumed to be a homomorphism of bare strings.

\begin{MSdefinition}[Operatively faithful language transport]\label{MS:def:language}
The pair $(T,\tau)$ is faithful on a transition-closed declared domain if, for every source state $s$ and message $m$, the following conditions hold:
\begin{align*}
F^A_m(s)\text{ is defined}&\quad\Longleftrightarrow\quad
F^B_{\tau(s,m)}(T(s))\text{ is defined},\\
T(F^A_m(s))&=F^B_{\tau(s,m)}(T(s)),\\
\ell_A(s,m)&=\ell_B(T(s),\tau(s,m)),
\end{align*}
where the last two equalities apply to admitted messages. Any required state observations, such as outstanding obligations, must also agree after translation.
\end{MSdefinition}

\begin{MSproposition}[Preservation of finite institutional conversations]\label{MS:prop:languagetrace}
Under Definition~\ref{MS:def:language}, every finite source conversation has an admitted prefix exactly when its stepwise translation has that admitted prefix from the translated initial state. Corresponding reached states are related by $T$, and corresponding event records agree. Two faithful transports compose faithfully using
\[
\tau_{AC}(s,m)=\tau_{BC}(T_{AB}(s),\tau_{AB}(s,m)).
\]
\end{MSproposition}
\begin{proof}
The initial states are related by definition. At each step, equivalence of definedness gives the same admission or rejection. For an admitted step, the transition equality gives the next related states, and the event equality gives the same semantic record. Induction proves the claim for every finite prefix. For composition, apply the two definedness equivalences successively and substitute one transition equality into the other. Equality of event records is transitive.
\end{proof}

The theorem preserves the translated conversations, not every unrelated target message. It neither grants access to hidden state needed to compute $\tau$ nor proves that $T$ is executable within a budget. A translation depending on an unavailable authorization record is a mathematical specification without a decentralized implementation.

\begin{MScounterexample}[The same payload can perform a different act]\label{MS:ex:speechact}
Suppose the source message \MScode{report(approved, item)} only appends a testimony record. A target translator drops the message type and emits \MScode{approve(item)}, which changes a permission bit from zero to one. Both messages contain the same descriptive payload, but their event records and state transitions differ. Even exact preservation of the payload does not satisfy Definition~\ref{MS:def:language}.
\end{MScounterexample}

\subsection{Private codes and public commitments}
A machine-native language may be symbolic, learned, compressed, or distributed across repeated exchanges. Its unfamiliar appearance is not a separate source of authority. A commitment interface can remain public and typed even when a process uses private internal representations to formulate a proposal. Conversely, a readable sentence can have opaque consequences if its interpreter or permission context is missing.

The preservation target is consequently an operation contract rather than a resemblance between strings. A protocol may preserve factual content while losing uncertainty, preserve a promise while changing its debtor, or preserve a vote while changing the electorate in which it is counted. These are different semantic failures. Approximate stochastic interpreters require the distributional and typed residual bounds developed elsewhere in the paper, not a claim that fluent paraphrase guarantees faithful action.

\subsection{Changing the interpreter is changing the institution}
If an amendment changes the meaning of a message type, the prior translation certificate need not remain valid. The interpreter version and its permissible update rules belong in the state. A constitution that carefully restricts \MScode{allocate} while allowing an unguarded change to the interpretation of \MScode{report} has not preserved its allocation boundary. This is a semantic instance of the checker-change counterexample, not a distinct kind of computation.

In the commons, a resource request, an approval certificate, a ledger commitment, and a repayment receipt should be separate message types. A successor can inherit the debt record while receiving no permission to issue new debt. This makes the distinction between understanding a statement and acquiring the authority to act on it explicit in the operational semantics.

\section{Why institutions might arise: costs, incentives, and coordination}\label{MS:sec:institution-formation}

Institutions are neither mysterious nor inevitable. They can become useful when repeated local negotiation is more expensive than creating and maintaining a shared rule. Their formation also depends on who pays, who benefits, who may leave, and whether participants have credible means to enforce an agreement. This is a recurring theme in repeated games, convention formation, and the governance of common resources~\cite{MS:fudenbergmaskin,MS:shoham,MS:ostrombook}.

\subsection{A baseline amortization calculation}

Suppose a coordination protocol without institution has expected per-period cost $c_0$. A proposed institutional protocol has operating cost $c_1$, per-period maintenance $m$, and a one-time setup cost $B$. With discount factor $0\le\delta<1$, write $g=c_0-c_1-m$.

\begin{MSproposition}[Institutional amortization]\label{MS:prop:amort}
Under stationary costs and infinite-horizon discounting, the institution's expected discounted cost advantage over the baseline is
\begin{equation}\label{MS:eq:amort}
V=-B+\frac{g}{1-\delta}.
\end{equation}
For a finite horizon $T$, it is $V_T=-B+g(1-\delta^T)/(1-\delta)$. The institution is cost-improving under these assumptions exactly when the corresponding value is positive.
\end{MSproposition}
\begin{proof}
Subtract institutional operating and maintenance costs from baseline costs in every period, discount, sum the finite or infinite geometric series, and subtract $B$.
\end{proof}

This proposition supplies a possible pressure toward persistent organization. It does not show that any particular learning algorithm will discover or adopt the institution, nor that the policy chosen is socially desirable. The assumption of stationary costs is especially demanding in adaptive populations whose tasks, membership, and resource prices can change.

\subsection{A profitable institution may be rejected}

Aggregate benefit is not the same as incentive compatibility. The simplest example already separates them.

\begin{MSproposition}[Aggregate profitability without voluntary adoption]\label{MS:prop:veto}
There exist two-agent repeated coordination environments in which a durable institution has strictly positive aggregate discounted benefit but is rejected by a rational agent who unilaterally bears setup cost and has veto power, assuming quasilinear payoffs, no transfers, and no additional benefits to that agent.
\end{MSproposition}
\begin{proof}
Let the setup cost be $B=10$, borne entirely by agent $1$, and let per-period gains be $g_1=0$, $g_2=2$ with $\delta=.9$. The aggregate advantage is $-10+2/(1-.9)=10>0$, whereas agent $1$'s payoff difference is $-10$. If agent $1$ can reject and has no compensating benefit or transfer, rejection is strictly preferable.
\end{proof}

A society thus cannot be modeled merely by minimizing one global loss without discarding the very plurality that makes it social. Who controls adoption matters. Bargaining power, externalities, and the allocation of setup costs can stabilize an inefficient status quo even when an institution would reduce total computation or improve collective performance.

In an idealized transferable-utility setting there is a corresponding conditional observation. If the discounted net advantages $b_i$ satisfy $\sum_i b_i\ge0$, then there exist budget-balanced transfers $t_i$ with $\sum_i t_i=0$ and $b_i+t_i\ge0$ for all $i$: choose nonnegative $w_i$ summing to $\sum_i b_i$ and set $t_i=w_i-b_i$. This algebra does not supply liquidity, a transfer mechanism, consent, enforceability, or a right to redistribute. The difference between mathematical feasible compensation and institutionally permitted compensation is central.

\subsection{Formation by local improvement: a restricted convergence result}

A candidate set of institutional rules $\mathcal I$ can be made part of a finite strategic environment. Participants have local choices $r_i\in\mathcal R_i$, with overall configuration $r=(r_1,\ldots,r_n)$ and expected utilities $u_i(r)$. A finite game is an \emph{exact potential game} if there exists $\Phi(r)$ such that for every unilateral change of $i$,
\[
u_i(r_i',r_{-i})-u_i(r_i,r_{-i})
=\Phi(r_i',r_{-i})-\Phi(r_i,r_{-i}).
\]
This is the familiar potential-game condition~\cite{MS:monderer}; it is not generically satisfied by social institutions.

\begin{MSproposition}[Finite termination of strict institutional improvement]\label{MS:prop:potential}
In a finite exact potential game, every sequence in which each executed unilateral update strictly improves the updating agent's utility has finitely many actual changes. If updates continue whenever an improvement is available, the terminal configuration is a pure-strategy Nash equilibrium. It need not maximize aggregate welfare.
\end{MSproposition}
\begin{proof}
Every executed update strictly increases $\Phi$, and a finite game has finitely many configurations and potential values. No configuration can recur along a strictly increasing path, so only finitely many actual changes occur. If the procedure ends only when no agent has a profitable deviation, the terminal state is a pure-strategy Nash equilibrium by definition. A local maximum of $\Phi$ need not maximize $\sum_i u_i$.
\end{proof}

The useful point is sharply conditional. Stable conventions can result from adaptive local behavior under special structural assumptions, but stability is neither proof of optimality nor evidence of fairness. With multiple competing objectives, cycles, capture, or regime changes, an institution may never stabilize or may stabilize for undesirable reasons.


\subsection{A formation barrier despite mutual benefit}
The stationary cost comparison does not describe a route to adoption. Consider $n\geq2$ agents choosing $a_i\in\{0,1\}$. An adopter pays cost $c>0$ and receives benefit $b>0$ from each other adopter:
\[
u_i(a)=a_i\left(b\sum_{j\ne i}a_j-c\right).
\]
Nonadopters receive zero. This is a finite network-coordination model, not a universal utility function for societies.

\begin{MSproposition}[Institutional formation barrier]\label{MS:prop:barrier}
Universal nonadoption is a strict Nash equilibrium. If $b(n-1)>c$, universal adoption is also a strict Nash equilibrium and gives every agent a strictly higher payoff than universal nonadoption. Nevertheless, no sequence of unilateral strict improvements starting at universal nonadoption can reach adoption.
\end{MSproposition}
\begin{proof}
A lone adopter receives $-c<0$, so no agent benefits from leaving universal nonadoption. At universal adoption, an adopter receives $b(n-1)-c>0$, whereas unilateral exit gives zero. The final assertion follows because no first strictly improving unilateral move is available at the starting state.
\end{proof}

The institution is feasible, stable once established, and unanimously preferred after establishment, but inaccessible through the permitted local improvement path. The problem is not inadequate intelligence or false beliefs. It is the structure of feasible transitions and incentives.

\begin{MSproposition}[Minimum profitable coordinated seed]\label{MS:prop:seed}
Starting from universal nonadoption, a simultaneously adopting coalition of size $m$ makes each member strictly better off exactly when $b(m-1)>c$. When such a coalition exists, the smallest size is
\[
m_* = \left\lfloor\frac cb\right\rfloor+2\leq n.
\]
Once at least $m_*$ adopters are present, every remaining nonadopter has a profitable unilateral entry. Under a schedule that eventually updates every remaining agent, strict best responses reach universal adoption without an adopter exiting.
\end{MSproposition}
\begin{proof}
Each coalition member has $m-1$ adopting partners and gains $b(m-1)-c$. Solving the strict integer inequality gives $m_*$. For a remaining agent, joining yields $bm-c>0$. Existing adopters retain positive payoff and their payoff increases with entry. Thus no exit is profitable, and an eventual-update schedule admits every remaining agent.
\end{proof}

A coordinated seed requires a permitted joint-commitment mechanism. The proof does not grant one automatically, supply startup resources, or remove the risk of an unfulfilled promise. In the commons, a simultaneous adoption record can implement such a commitment only if the initial constitutional envelope admits it. This is a bounded way for institutional content to arise inside supplied infrastructure.

For $n=4$, $b=1$, and $c=2$, full adoption gives each agent payoff $1$, yet three simultaneous adopters would only break even. Four are needed for a \emph{strictly} profitable seed. Replacing $>$ with $\geq$ would change that answer. Figure~\ref{MS:fig:barrier} separates the outsider's entry gain from an incumbent's maintenance payoff.

\begin{figure}[tbp]
\centering
\includegraphics[width=.84\linewidth]{sources/monographs/machine-societies/publication/figures/formation_barrier.pdf}
\caption{The exact four-agent coordination example with $b=1$ and $c=2$. Entry gain at $k$ existing adopters is $k-2$; an incumbent's payoff with $k$ total adopters is $k-3$. The curves have different domains because an incumbent exists only for $k\geq1$ and an outsider only for $k\leq3$. Zero gain is not a strict improvement.}\label{MS:fig:barrier}
\end{figure}

\subsection{Path dependence and a qualified hysteresis interpretation}
For $0<c<b(n-1)$ both homogeneous states are locally stable. Under a slowly varied \emph{effective} cost and a best-response rule that retains the current action at ties, universal nonadoption persists until the effective cost becomes negative, while universal adoption persists until it exceeds $b(n-1)$. Negative effective cost represents a net adoption subsidy in this extended parameterization. The distinct switching points describe hysteresis for that specified update process. They are not a thermodynamic phase transition or an inevitable law of institutional development. Different coalition moves, noise, exit costs, or expectations change the path.

\subsection{Persistent rules and the economy of attention}

Institutions can economize on repeated communication because a participant need not renegotiate every detail in every encounter. But their records, audits, sanctions, and revision mechanisms consume resources. A full theoretical cost account must include setup and maintenance, opportunity cost of governance, task error, time to reach agreement, cost of exit, and loss from prohibited actions. If agents can cheaply renegotiate or if interactions rarely recur, a transient coalition may be preferable to a durable institution. If a social rule is cheap to state but expensive to verify, its apparent savings may be illusory.


\section{Enforcement, observability, and incentives to comply}\label{MS:sec:enforcement}

An operative rule can restrict physically available actions, change their payoffs, or both. These mechanisms should not be conflated. An access gate can make an unauthorized action unavailable under its implementation assumptions. A sanctioning institution instead leaves a deviation feasible but changes its expected consequences. A self-enforcing arrangement makes continued cooperation preferable through future interactions. Proving one mechanism exists does not establish either of the others.

We develop a small one-deviation model to connect enforcement to the information structure of the society. The conceptual foundation is moral-hazard analysis: when conduct is not directly contractible, incentives depend on observable consequences~\cite{MS:holmstrom1979}. The calculation below isolates a finite discrimination problem. It is not a general solution of optimal contracting.

\subsection{What a bounded sanction can distinguish}
An agent chooses compliance $C$ or a specified feasible deviation $D$. Before sanctions, $D$ gives expected utility gain $g>0$ over $C$. An observation $o$ in a finite set $\mathcal O$ has law $P_C$ or $P_D$. The institution commits to a sanction probability $s(o)\in[0,1]$ with fixed penalty $F>0$, measured in the same utility units. The agent is risk-neutral in this penalty, its collection is feasible and authorized, and all other payoff differences have already been included in $g$. The observation laws and the committed policy are known to the agent.

Under these assumptions the deviation's net expected advantage is
\begin{equation}\label{MS:eq:deterrence-gap}
g-F\bigl(\MSE_{P_D}s-\MSE_{P_C}s\bigr).
\end{equation}
Compliance is weakly preferable exactly when this quantity is nonpositive. Equality permits indifference and does not predict which action will be selected.

\begin{MSproposition}[The exact bounded-sanction discrimination limit]\label{MS:prop:deterrence-tv}
In this model, some observation-based sanction makes compliance weakly preferable if and only if
\[
g\leq F\,\MSTV(P_D,P_C).
\]
Strict preference is attainable if and only if the inequality is strict.
\end{MSproposition}
\begin{proof}
Write $d(o)=P_D(o)-P_C(o)$. Since $\sum_o d(o)=0$, the sum of its positive terms equals $\frac12\sum_o|d(o)|=\MSTV(P_D,P_C)$. For every $0\leq s\leq1$,
\[
\sum_o s(o)d(o)\leq\sum_{d(o)>0}d(o).
\]
Equality is attained by $s(o)=1$ on the positive terms and zero elsewhere. Substitution into~\eqref{MS:eq:deterrence-gap} proves both assertions.
\end{proof}

If the two observation laws coincide, increasing a finite penalty cannot selectively deter this deviation: it affects both choices identically. This does not rule out prevention, a different observation channel, intrinsic preferences, or future incentives omitted from the model. Conversely, a high discrimination value does not establish that the penalty is legitimate, collectible, or compatible with participation. An accurate accusation is not itself an authorization to punish.

For a concrete two-record example, let
\[
P_C=(3/4,1/4),\qquad P_D=(1/4,3/4).
\]
The total variation is $1/2$. Sanctioning on the second record deters gains $g\leq F/2$, but also sanctions compliance with probability $1/4$. A deterrence certificate without this false-positive rate would omit a major institutional consequence.

\subsection{False-positive limits and multiple deviations}
A constitutional protection may impose $\MSE_{P_C}s\leq\alpha$, with $0\leq\alpha\leq1$. The relevant discrimination capacity is then the finite linear program
\begin{equation}\label{MS:eq:protected-discrimination}
D_\alpha(P_D,P_C)=
\max_{\substack{0\leq s\leq1\\\MSE_{P_C}s\leq\alpha}}
\bigl(\MSE_{P_D}s-\MSE_{P_C}s\bigr).
\end{equation}
The feasible set is nonempty and compact, so a maximum exists. The preceding proof with the feasible set restricted shows that deterrence is attainable exactly when $g\leq F D_\alpha$. For the example above and $\alpha=1/10$, the optimum is $s=(0,2/5)$, giving $D_\alpha=1/5$, not $1/2$. Protected error constraints can materially reduce feasible enforcement.

When there are deviations $D_1,\ldots,D_k$, one common policy must satisfy every incentive inequality. Separate pairwise discrimination bounds are necessary but not sufficient. With $P_C=(1/2,1/2)$, $P_{D_1}=(1,0)$, and $P_{D_2}=(0,1)$, each pair has distance $1/2$, yet the two differences $\MSE_{P_{D_j}}s-\MSE_{P_C}s$ sum to zero. No common sanction deters strictly positive gains from both deviations. A deployable certificate must retain the joint system of constraints, not select a different optimal monitor after learning which deviation occurred.

\subsection{Compression can erase enforceable distinctions}
Let an observation channel $B$ replace $o\in\mathcal O$ by a record $z\in\mathcal Z$. It may be deterministic compression or a stochastic reporting channel, but must be the same channel under compliance and deviation.

\begin{MSproposition}[Garbling does not increase sanctioning capacity]\label{MS:prop:enforcement-garbling}
For every such channel and every $\alpha\in[0,1]$,
\[
D_\alpha(P_D B,P_C B)\leq D_\alpha(P_D,P_C).
\]
In particular, $\MSTV(P_D B,P_C B)\leq\MSTV(P_D,P_C)$.
\end{MSproposition}
\begin{proof}
Given a coarse-record sanction $t(z)$, define $s(o)=\sum_z B(o,z)t(z)$. It lies in $[0,1]$ and has exactly the same compliance and deviation expectations as $t$ has under the garbled laws. Every feasible coarse policy therefore has a fine-record counterpart with the same objective and false-positive rate. Maximizing proves the first claim; the unrestricted case gives the total-variation statement.
\end{proof}

This result concerns available information, not an instruction to retain every private state. An institution may deliberately accept weaker deterrence to protect confidentiality, reduce storage cost, or constrain surveillance. The theorem exposes the tradeoff rather than choosing its normative optimum. Keeping a signed summary establishes who issued the summary, not that it preserved the distinctions needed for a later complaint.

\subsection{Enforcement must be incentive-compatible for the enforcer}
The model assumes that the stated monitor and sanction are actually implemented. An auditor may face its own costs, limited resources, or incentives to ignore a violation. A society therefore needs separate contracts for evidence production, decision authority, and execution of a remedy. Collusion, hidden side payments, uncollectible penalties, and changes to the observation process leave the single-agent model. They cannot be repaired by declaring a larger value of $F$. Section~\ref{MS:sec:strategic-auditor} places a strategic inspector inside a specified game and distinguishes a feasible monitor from an incentive-supported service.

In the commons, unique slot tokens can prevent duplicate grants at the resource layer while a reporting rule encourages accurate urgency submissions. Those mechanisms coexist but have different proof obligations. An appeal may change the observation set and the sanction rule, while a creditor's outstanding service claim remains unchanged. Section~\ref{MS:sec:multiscale} will show how errors in a reduced model of the observation process alter a deterrence margin. This is the point where institutional dynamics and enforcement certificates must be checked together.


\section{One enforcement policy for several deviations}\label{MS:sec:joint-enforcement}
The sanction criterion in Section~\ref{MS:sec:enforcement} concerns one specified deviation. A continuing institution must usually respond to several possible departures from its rules with one announced procedure. The order of quantifiers changes the problem. A separate monitor may exist for each deviation without any monitor working for all of them. The institution cannot choose its sanction rule after secretly learning which unobserved deviation occurred.

The following finite model isolates that issue. It retains the earlier assumptions of risk neutrality, a committed and authorized sanction mechanism, a collectible fixed penalty, and specified observation laws. It does not solve the incentives of the auditor, establish those laws from data, or claim that every deviation is represented.

\subsection{The common-monitor problem}
Let $O$ be a finite nonempty observation set. Compliance has law $P_C$, and deviation $j\in\{1,\ldots,m\}$ has law $P_j$ and gross utility advantage $g_j\geq0$, where $m\geq1$. A fixed penalty $F>0$ is applied with probability $s(o)\in[0,1]$ after record $o$. Impose a protected false-positive rate $\alpha\in[0,1]$ through the compact polytope
\[
\mathcal M_\alpha=\{s\in[0,1]^O:\langle P_C,s\rangle\leq\alpha\}.
\]
Here $\langle P,s\rangle=\sum_{o\in O}P(o)s(o)$. The same $s$ must be used for every listed deviation. Define its margin against deviation $j$ and its best simultaneous margin by
\begin{equation}\label{MS:eq:joint-margin}
D_j(s)=F\langle P_j-P_C,s\rangle-g_j,
\qquad
\Gamma_\alpha=\max_{s\in\mathcal M_\alpha}\min_{1\leq j\leq m}D_j(s).
\end{equation}
Weak deterrence of all listed deviations means $\Gamma_\alpha\geq0$. Strict deterrence means $\Gamma_\alpha>0$. At equality a participant can remain indifferent. A randomized choice among the listed deviations is also deterred when every listed pure deviation is deterred, since its expected advantage is the corresponding convex combination. This does not include adaptive, repeated, or unlisted deviations unless their observation laws and payoffs are explicitly represented.

For a probability vector $\lambda\in\Delta_m$, set $P_\lambda=\sum_j\lambda_jP_j$ and $g_\lambda=\sum_j\lambda_jg_j$. Write
\[
V_\alpha(P,P_C)=\max_{s\in\mathcal M_\alpha}\langle P-P_C,s\rangle.
\]
This is the same constrained observation advantage used in the single-deviation analysis, now evaluated on mixtures of deviation laws.

\begin{MStheorem}[Simultaneous enforcement and mixture obstruction]\label{MS:thm:joint-enforcement}
Under the finite common-monitor model,
\begin{equation}\label{MS:eq:joint-dual}
\Gamma_\alpha=\min_{\lambda\in\Delta_m}
\bigl(FV_\alpha(P_\lambda,P_C)-g_\lambda\bigr).
\end{equation}
Both extrema are attained. Consequently, one monitor weakly deters all listed deviations if and only if $FV_\alpha(P_\lambda,P_C)\geq g_\lambda$ for every $\lambda\in\Delta_m$. Without an additional false-positive restriction, $\alpha=1$ and $V_1(P,P_C)=\MSTV(P,P_C)$.
\end{MStheorem}
\begin{proof}
For any fixed $s$, the minimum of the numbers $D_j(s)$ is the minimum of $\sum_j\lambda_jD_j(s)$ over $\Delta_m$. This expression is bilinear in $s$ and $\lambda$, and both domains are nonempty compact polytopes. Finite minimax, equivalently linear programming duality, therefore permits interchange of the maximum and minimum~\cite{MS:boydconvex}. One can reduce directly to a finite matrix game by listing the vertices of $\mathcal M_\alpha$ and taking their convex combinations. For fixed $\lambda$, the maximized expression is $FV_\alpha(P_\lambda,P_C)-g_\lambda$, proving~\eqref{MS:eq:joint-dual}. Compactness and continuity give attainment. Finally, when $\alpha=1$, every $s\in[0,1]^O$ is admissible. Choosing $s(o)=1$ exactly where $P(o)>P_C(o)$ attains the sum of positive differences, which equals total variation.
\end{proof}

The theorem is an application of classical finite convex duality, not a new minimax theorem. Its institutional contribution is the distinction between individual enforceability and a jointly executable sanction policy. A vector $\lambda$ for which the right side of~\eqref{MS:eq:joint-dual} is negative is a certificate that the stipulated contract cannot be met. It is not evidence that a real population has adopted that mixture.

\begin{MScounterexample}[Individually enforceable deviations with no common monitor]\label{MS:ex:pairwise-enforcement}
Take $O=\{0,1\}$, $P_C=(1/2,1/2)$, $P_1=(1,0)$, $P_2=(0,1)$, $F=2$, $g_1=g_2=1/2$, and $\alpha=1/2$. The monitor $(1,0)$ strictly deters deviation $1$, and $(0,1)$ strictly deters deviation $2$. Both respect the false-positive cap. For any common monitor,
\[
D_1(s)=s(0)-s(1)-\tfrac12,
\qquad
D_2(s)=s(1)-s(0)-\tfrac12.
\]
Their sum is $-1$, so both cannot be nonnegative. The best simultaneous margin is $-1/2$, attained when $s(0)=s(1)$. The dual obstruction is $\lambda=(1/2,1/2)$, for which $P_\lambda=P_C$ while $g_\lambda=1/2$.
\end{MScounterexample}

Distinct deviations can cancel one another's observable signatures at the level of a mixture. Pairwise total-variation distances therefore do not determine joint enforceability. Increasing the penalty cannot resolve the particular obstruction $P_\lambda=P_C$ with $g_\lambda>0$. The contract must change its observation channel, permitted interventions, payoff structure, or covered claims. None of these changes is automatically authorized.

\subsection{Uncertainty must be absorbed by the same monitor}
Averaging, translation, and statistical approximation can change the estimated observation laws. The preceding optimization does not remove those errors. Let $\widehat P_C,\widehat P_j$ be specified approximate laws with
\[
\MSTV(P_C,\widehat P_C)\leq\epsilon_C,
\qquad
\MSTV(P_j,\widehat P_j)\leq\epsilon_j.
\]
The bounds may be different for different deviations. They must concern the complete record available to the sanction mechanism, including any effects of review scheduling.

\begin{MSproposition}[A robust simultaneous enforcement certificate]\label{MS:prop:joint-robust}
Let $\alpha\geq\epsilon_C$ and $\eta\geq0$. If a single $s\in[0,1]^O$ satisfies
\begin{align}
\langle\widehat P_C,s\rangle&\leq\alpha-\epsilon_C,\label{MS:eq:joint-robust-cap}\\
F\langle\widehat P_j-\widehat P_C,s\rangle-g_j
&\geq\eta+F(\epsilon_j+\epsilon_C)
\quad\text{for every }j,\label{MS:eq:joint-robust-gap}
\end{align}
then the true false-positive rate is at most $\alpha$ and every true margin $D_j(s)$ is at least $\eta$.
\end{MSproposition}
\begin{proof}
For every function $0\leq s\leq1$, the difference of its expectations under two laws is bounded in absolute value by their total-variation distance. Applying this once to the compliance rate proves~\eqref{MS:eq:joint-robust-cap}'s conclusion. Applying it to both laws in each difference gives
\[
D_j(s)\geq F\langle\widehat P_j-\widehat P_C,s\rangle-g_j
-F(\epsilon_j+\epsilon_C)\geq\eta.
\]
The same $s$ appears in all inequalities, so the conclusion is simultaneous.
\end{proof}

This is a sufficient certificate, not a necessary description of all robustly feasible policies. Replacing each law by an independent worst-case error bound can be conservative when the errors are related. If $\alpha<\epsilon_C$, the particular tightened test is unavailable, but that does not prove that every actual monitor violates the true cap. The proposition extends Proposition~\ref{MS:prop:robust-deterrence} to a common monitor rather than separately optimized monitors.

\subsection{What an enforcement certificate should contain}
The optimization is a small linear program with variables $(s,t)$, constraints $D_j(s)\geq t$, the box bounds on $s$, and the false-positive cap. A usable certificate records the monitor, every listed gain and observation law, the authority to impose the penalty, its collectibility assumption, and the checked margins. A negative certificate can instead record a mixture $\lambda$ and an upper bound for its constrained observation advantage.

The linear program does not make observation collection free. Its monitor may require an expensive audit, an unavailable witness, or more time than a service deadline permits. Section~\ref{MS:sec:privacy-enforcement} adds a precise constraint on what the complete available record may distinguish. It also does not authorize punishment for an unproved real-world act. The distinction between a stipulated probabilistic sanction model and a justified institutional procedure remains explicit. The next synthesis asks whether an observation procedure, a common monitor, authorized staffing, and outstanding services can fit within one feasible course of action.
\FloatBarrier

\section{Social choice and principled disagreement}\label{MS:sec:choice}
Plurality cannot be retained while silently replacing all preferences with one externally supplied reward. Even agents that agree on every fact can rank institutional alternatives differently. Let the alternatives be three policies $R,U,V$ and let three voters have strict preferences
\[
R\succ_1 U\succ_1 V,\qquad
U\succ_2 V\succ_2 R,\qquad
V\succ_3 R\succ_3 U.
\]

\begin{MScounterexample}[Truthful majority cycle]\label{MS:ex:majority}
Pairwise majority prefers $R$ to $U$, $U$ to $V$, and $V$ to $R$, each by two votes to one. Every individual order is transitive. No erroneous report or mistaken inference is needed for the majority relation to cycle.
\end{MScounterexample}

In successive pairwise elimination, agenda order can select any of the three final winners. To obtain $R$, compare $U$ with $V$ first, then compare their winner $U$ with $R$. Cyclic symmetry gives the other outcomes. This is an analytic example of agenda power, not a claim that majority decision is unusable.

\subsection{Classical impossibility results have explicit domains}
In its standard finite social-welfare formulation, Arrow's theorem rules out simultaneously obtaining unrestricted preference profiles, a complete transitive social ordering, Pareto unanimity, independence of irrelevant alternatives, and nondictatorship when there are at least three alternatives~\cite{MS:arrow}. The result concerns an aggregation rule with those assumptions. It is not a theorem that every procedural choice is arbitrary, that all voting fails, or that machines must be governed by a dictator. Restricting admissible preference domains or changing the object produced by the rule changes the question.

Strategic reporting imposes a different constraint. For a deterministic onto social-choice function with at least three possible outcomes and unrestricted strict preferences, the Gibbard--Satterthwaite result characterizes strategy-proof rules as dictatorial~\cite{MS:gibbard,MS:satterthwaite}. Strategy-proofness here means that no individual can obtain a preferred outcome by misreporting while others' reports are fixed. Payments, restricted domains, randomized rules, and weaker equilibrium concepts require separate analysis. The two classical theorems are cited results, not new propositions or consequences of the finite majority example.

\subsection{Institutions choose procedures, not just outcomes}
A complete institutional record should identify the eligible electorate, agenda, comparison rule, admissible reports, tie convention, protected constraints, and appeal procedure. These features can change the result while leaving individual preferences unchanged. They also determine who may propose a comparison in the first place.

The commons' three-of-four amendment rule guarantees neither that a proposal will pass nor that the adopted rule is optimal for all offices. It defines one permissible way to resolve disagreement. A proof that the rule was followed is distinct from an argument for selecting that rule. Deliberation can change beliefs and reveal new alternatives, but it need not remove a genuine conflict of interests. The appropriate mathematical goal is to expose the tradeoffs rather than label residual disagreement a failure of collective intelligence.

\section{The constitution as a revisable computational object}\label{MS:sec:constitution}

A machine institution may contain rules about changing its own rules. This creates a formal distinction between a state transition that is checked by a fixed constitution and a transition that changes the authority or checker used to judge future transitions. Treating the latter as an ordinary edit without recording the change in verification semantics is a source of hidden self-reference.

Let $s=(x,I)$ be the full operating state and let $R(s,s')$ be the admitted transition relation, including ordinary actions and amendments. A safety property is a predicate $\Psi(s)$.

\begin{MSproposition}[Constitutional inductive safety]\label{MS:prop:safety}
If $\Psi(s_0)$ and every admitted transition satisfies
\[
\Psi(s)\wedge R(s,s')\Longrightarrow\Psi(s'),
\]
then $\Psi(s_t)$ holds for every state reachable by finitely many admitted transitions.
\end{MSproposition}
\begin{proof}
Induction on the number of admitted transitions.
\end{proof}

This standard invariant is useful precisely because its assumptions are explicit. It certifies safety relative to $R$. It does not establish that the implemented system obeys $R$, that the predicate is ethically adequate, or that a decision can always be reached. A protocol can satisfy a safety condition by refusing to act forever. Timing and failures enter liveness analysis separately; distributed consensus already has sharp impossibility limits under specified asynchronous assumptions~\cite{MS:flp}.

\begin{MScounterexample}[An unchecked change of the checker]\label{MS:ex:checker}
Let the institutional state contain a validator mode $v\in\{\mathsf{strict},\mathsf{open}\}$ and a resource policy $p\in\{\mathsf{safe},\mathsf{unsafe}\}$. Under $v=\mathsf{strict}$, the current policy must remain safe, but the validator mode itself is unprotected. Thus $(\mathsf{strict},\mathsf{safe})\to(\mathsf{open},\mathsf{safe})$ is admitted. Under $v=\mathsf{open}$, any policy transition is admitted, including $(\mathsf{open},\mathsf{safe})\to(\mathsf{open},\mathsf{unsafe})$. The original current-policy check does not imply that all reachable states remain safe. To prove that property, the inductive predicate and transition obligations must include permissible changes to the validation mechanism itself.
\end{MScounterexample}

The counterexample does not forbid constitutional amendment. It shows that an amendment system needs an explicit \emph{trust and revision boundary}. One possibility is a fixed external checker for a declared protected property. Another is a staged, authorized change that carries a proof that the new checker preserves a specified invariant. A third is to treat a fully mutable constitution as a weaker, revisable commitment and avoid calling it unconditionally safe. These are different design choices; no single mathematical construction decides the political legitimacy of the guard itself.

\subsection{Viability under stochastic institutional transitions}

An invariant proof checks that \emph{every permitted transition} preserves a property. A different question asks whether an institution has a \emph{choice of actions} that will preserve a protected property even when the environment responds stochastically. This can be formulated as a finite viability problem. Let $S$ be a finite subset of safe states, let $A(s)$ be the finite set of available actions at $s$, and let $K(\cdot\mid s,a)$ be the transition law on a finite full state space. The state must include every variable needed for its Markov representation, including the operative amendment rules.

Define $V_0=S$ and recursively, for $h\ge0$,
\begin{equation}\label{MS:eq:viability}
V_{h+1}=\Bigl\{s\in S:\exists a\in A(s)\;\text{such that}\;
\operatorname{supp}K(\cdot\mid s,a)\subseteq V_h\Bigr\}.
\end{equation}
\begin{MSproposition}[Finite-horizon almost-sure institutional viability]\label{MS:prop:viability}
A state $s$ belongs to $V_h$ if and only if there exists an admissible policy under which every state visited over the next $h$ transitions, including the present state, lies in $S$ with probability one. Policies may use the full observed state and the remaining horizon. For $s\in V_h$, a time-dependent Markov policy suffices.
\end{MSproposition}
\begin{proof}
The claim at $h=0$ is $V_0=S$. Suppose it holds at $h$. To remain in $S$ almost surely for another step and then for $h$ further steps, the first action must have support entirely in states from which the remaining $h$ steps can be kept safe. By the induction hypothesis, these are exactly $V_h$. Thus admissible first actions are precisely those in the definition of $V_{h+1}$. Choosing a witnessing action at each state and remaining horizon supplies a time-dependent Markov policy. Conversely, if no such action exists, every first action assigns positive probability to a successor from which the required remaining guarantee is impossible. Randomizing over actions cannot remove that positive failure probability unless all positively chosen actions are themselves witnessing actions.
\end{proof}

This is a standard backward-reachability principle from finite stochastic control, adapted to institutional transitions. It exposes the distinction between a \emph{declared safe predicate}, a \emph{proof that all admitted actions preserve it}, and the \emph{existence of a strategy that can keep it satisfied}. None guarantees eventual service, procedural fairness, or low computational cost. In particular, declaring all actions forbidden can make safety trivial while making a purported society inoperative.

\subsection{Constitutional semantics and institutional memory}

A sound amendment record needs the prior state, proposal, evidence, decision rights, adopted state, and the semantics of the rule used to decide. Without that information a successor may know the current text but not why an exception was made, whether its source was authorized, or whether a later amendment invalidated a prior permission. Append-only provenance may aid auditability, but signatures verify assertions by key holders, not truth or authority. A correct proof object likewise establishes a formal judgment relative to its premises and trusted kernel~\cite{MS:necula}.

\subsection{Dissent, appeal, and the risk of overcoherence}

One invariant may demand a single ordering for resource transactions; another should allow mutually inconsistent scientific hypotheses to remain under consideration. Enforcing total agreement as an abstract ``coherence'' objective can destroy legitimate exploration or conceal genuine conflicts. An appeal procedure is not noise in an otherwise efficient optimizer. It changes who can challenge a decision, which evidence can reopen it, and how error is corrected. In this sense, a machine institution may be more robust when it preserves controlled disagreement than when every node converges to the same report.


\section{Constitutional confluence and the authority of ordering}\label{MS:sec:confluence}
Power can reside in the order of authorized decisions. Two individually permitted amendments may not commute, so a scheduler, deadline rule, or proposal queue can affect the eventual institution without originating any proposal. Section~\ref{MS:sec:choice} shows one form of agenda dependence through preference aggregation. Here the issue is the operational order of amendments after their admissibility has been specified.

Fix a bounded amendment episode. Let $s\to s'$ mean one permitted amendment step, and let $\to^*$ be its reflexive transitive closure. The episode state contains any pending proposals, version counters, permissions, and consequence records required by the contract. It is \emph{terminating} if it has no infinite chain of actual steps, and \emph{locally confluent} if
\[
s\to a,\ s\to b\quad\Longrightarrow\quad
\exists c:\ a\to^*c\ \text{and}\ b\to^*c.
\]
Confluence imposes the same condition on arbitrary finite descendants of $s$. A normal form has no outgoing amendment step. Reversibility of an entire society need not be prohibited: termination is a property of this episode, not a requirement that all future institutions stop changing.

\begin{MStheorem}[Order-independent amendment normal forms]\label{MS:thm:confluence}
A terminating, locally confluent amendment relation is confluent. Each starting state has a unique reachable normal form, so every maximal amendment sequence ends at that normal form. This is Newman's lemma applied to the declared institutional transition relation~\cite{MS:newman}.
\end{MStheorem}
\begin{proof}
Use well-founded induction on strict descendants of $s$. For descendants $y,z$ of $s$, there is nothing to prove if one equals $s$. Otherwise choose first steps $s\to a\to^*y$ and $s\to b\to^*z$. Local confluence supplies $c$ with $a\to^*c$ and $b\to^*c$. The induction hypothesis at $a$ joins $y,c$ at some $p$, and the hypothesis at $b$ joins $z,c$ at some $q$. Since $c$ is also a strict descendant of $s$, its hypothesis joins $p,q$. Thus $y,z$ have a common descendant. Termination ensures a normal form is reached along every maximal sequence. Two reachable normal forms must coincide because confluence gives a common descendant and neither can move.
\end{proof}

A natural-number ranking strictly decreased by every step is one sufficient termination certificate. For a fixed finite set of proposals, the number of unprocessed proposals works only if each step removes one and cannot add more. A rule that may create additional proposals needs a different argument. Well-foundedness cannot be inferred merely from the finiteness of the written rulebook.

\begin{figure}[t]
\centering
\includegraphics[width=.87\linewidth]{sources/monographs/machine-societies/publication/figures/constitutional_confluence.pdf}
\caption{A joinability obligation for two admissible amendment paths. The diagram specifies a condition, not an automatic outcome. Confluence concerns the full declared state or a separately justified semantic quotient; equal final policy text can hide different resource histories.}\label{MS:fig:confluence}
\end{figure}

\subsection{An authorized ordering can still select the outcome}
Consider two approved proposals, one setting a policy bit to zero and the other setting it to one. Each proposal is executed exactly once and removed from the pending set. The episode terminates, but its final bit is the value written by the last proposal. The two terminal states are distinct and cannot be joined within the episode. Authorization of both proposals does not imply local confluence.

A procedural priority rule can choose one order and make execution deterministic. That resolves ambiguity by assigning power to the priority rule; it does not prove that the original unordered amendment system was confluent. Likewise, arbitrarily rejecting one proposal after the fact changes the relation being analyzed. A reconciliation procedure must itself be admitted and included in the proof obligation.

\subsection{Endpoint agreement does not erase path effects}
Suppose two amendment paths restore the same final policy, but the first allocates an irreversible resource to office $A$ while the second allocates it to $B$. A state consisting only of the final policy may have a unique normal form; the allocation record does not. Therefore endpoint confluence of a policy projection does not imply equal costs, equal waiting times, or equal histories of rights violations.

An append-only log also distinguishes the literal order of commuting edits. To identify such logs, a semantic quotient must state which reorderings are immaterial and preserve the audit questions promised by the contract. The governance quotient cannot silently discard the very ordering under investigation. In the commons, two independent label changes may commute, while a loan authorization and the removal of its debtor's office require an explicit treatment of inherited obligations. The outcome of that interaction is not determined by calling both edits valid.

\section{Social time, distributed history, and institutional finality}\label{MS:sec:social-time}

The commons in Section~\ref{MS:sec:commons} assumes an atomic event mechanism. That assumption allows us to ask whether a rule preserves an allocation or a debt without simultaneously solving the implementation of the ledger. It is not a theorem that autonomous offices can create the mechanism under arbitrary communication conditions. We now consider a separate realization in which offices hold local records and messages can be delayed. The earlier atomic model remains a reference semantics; the decentralized variant must justify any claim that it implements that semantics.

Institutional time has several meanings. A request can be created, transmitted, received, provisionally admitted, committed, and announced at different events. A debt's deadline might refer to a physical clock, a numbered resource epoch, or the completion of a prior procedure. These alternatives define different contracts. Lamport's causal ordering makes the underlying distinction precise: local process order and message transmission generate a partial order, not a universally observed wall clock~\cite{MS:lamporttime}. Concurrent events need not have a causally determined relative order. A chosen total order can support execution, but its priority rule is itself an institutional decision.

\subsection{Convergent records and valid decisions}
Let $(L,\vee)$ be a join-semilattice, with order $a\leq b$ when $a\vee b=b$. A replica holds $\ell_i\in L$. In the following restricted model a local update joins a declared contribution $\delta$ into the state, and a merge replaces $\ell_i$ by $\ell_i\vee\ell_j$. Contributions are fixed elements of $L$ once created. This is a state-merge construction of the kind used in conflict-free replicated data types~\cite{MS:preguicacrdt}. The algebra concerns accepted data, not the authenticity of its source.

\begin{MSproposition}[Quiescent convergence and invariant preservation]\label{MS:prop:merge-convergence}
Assume finitely many replicas begin at $\ell_0$, only finitely many contributions $\delta_1,\ldots,\delta_N$ are created, and every contribution eventually becomes included in every replica. Then each replica eventually equals
\[
\ell^\star=\ell_0\vee\delta_1\vee\cdots\vee\delta_N,
\]
independently of merge order and duplicate deliveries. Moreover, let $V\subseteq L$ contain the initial state. If every permitted local update maps $V$ into $V$ and $V$ is closed under $\vee$, every local and merged state remains in $V$.
\end{MSproposition}
\begin{proof}
At every step a replica is the join of $\ell_0$ and some created contributions. Associativity, commutativity, and idempotence remove ordering and duplication effects. Once every contribution has reached a replica, its state is $\ell^\star$ and cannot subsequently change. Finiteness of the replicas and contributions gives a finite point after which this holds everywhere. For the second claim, induction over local updates and merges preserves membership in $V$ by the two closure assumptions.
\end{proof}

The conclusion conditions on the same accepted contributions. It does not say that different schedules must generate the same contributions when update proposals depend on local state. Eventual convergence after quiescence also does not mean that a replica can detect when every delayed contribution has arrived.

For causal event records, one may take $L$ to be the downward-closed subsets of a fixed finite event partial order, with union as join. A receiver must retain an event's prerequisites or wait until they arrive. Downward closure preserves a causal dependency record. It does not certify that a signature is authentic, that an issuer held authority, or that two concurrent grants are jointly permissible.

\begin{MScounterexample}[Convergence can preserve a record of an invalid allocation]\label{MS:ex:convergent-double-grant}
One resource token $r$ must have at most one irrevocable grantee. Two replicas start with the empty grant set. One records $(r,A)$ and the other $(r,B)$, for $A\ne B$. Both singleton records satisfy the exclusivity predicate. Their set union does not. Nevertheless, union merging converges exactly, and every replica eventually agrees on the two-grant record.
\end{MScounterexample}

The counterexample is not a failure of set replication. It is a failure to choose an invariant-preserving set of independently executable operations. A deterministic conflict resolver that later retains only one grant can restore a selected policy record, but cannot thereby undo both grants' already executed effects. Tentative reservations and irrevocable allocations must have different event types.

The safe set's closure under join is a sufficient condition for the simple construction above. It need not be necessary when a protocol prevents the problematic pairs from being independently reachable. Invariant-confluence theory refines precisely this issue by considering operation-reachable states with compatible ancestry~\cite{MS:bailisinvariant}. CALM relates coordination-free computation to logical monotonicity in its stated distributed-program model~\cite{MS:hellerstein-calm}. Neither result should be replaced by the claim that any monotonically growing ledger is institutionally safe.

\subsection{An elementary coordination obstruction}
The tension can be proved without postulating malicious participants or unreliable reasoning. Consider two decision sites with the same initial resource state. Each sees its own requests but, before a fixed service deadline, receives no information about requests at the other site. Assume there is no separate shared atomic resource operation available before that deadline.

\begin{MSproposition}[Exclusive service cannot be unconditionally local]\label{MS:prop:exclusive-service}
No deterministic protocol can satisfy both of the following in every such execution: (i) whenever a request occurs at only one site, that site irrevocably grants the resource before the deadline; and (ii) at most one grant occurs when requests arrive concurrently at both sites. The same conclusion holds for randomized protocols when both guarantees are required with probability one and random tapes have the same request-independent law across these executions.
\end{MSproposition}
\begin{proof}
Compare the two singleton-request executions with the concurrent-request execution. Until the deadline, each site has the same local history in the concurrent execution as in its own singleton execution. Requirement (i) therefore makes each grant, violating (ii). In the randomized case each site's grant event has probability one in the concurrent execution by the same local-law comparison. Their intersection also has probability one, without requiring independent random tapes.
\end{proof}

Exclusive preallocation avoids the conclusion by dropping unconditional singleton service at the nonowner. A common sequencer, a sufficiently strong distributed commit protocol, or a resource owner that performs the atomic grant changes the information or operation assumptions. Waiting for communication drops the deadline guarantee. These are distinct institutional choices, not defects to be concealed behind the word decentralized. The proposition is an indistinguishability argument for a restricted service contract, not a replacement for the broader consensus theory discussed in Section~\ref{MS:sec:determinism}.

\subsection{Finality is indexed by a continuation contract}
A provisional decision is one that later authorized processing may reject or revise. A decision is \emph{final relative to a continuation contract} when every continuation admitted by that contract preserves its designated committed effect. This definition does not require every participant already to know the decision. Nor does it make a rule text eternally unamendable: a later prospective amendment may coexist with the continuing validity of an earlier grant or debt.

The observation ``no conflicting event has arrived'' is not a finality certificate when delays are unbounded. A certificate must state its authority configuration, event or resource identifier, admissible continuations, and the particular safety property it establishes. Intersecting decision certificates and cross-configuration endorsement obligations can rule out conflicting commitments under the assumptions of Section~\ref{MS:sec:membership}. They do not by themselves establish eventual issuance, timely notification, or common knowledge. The distinction between committed effect and known committed effect links this section to Section~\ref{MS:sec:publicknowledge}.

\begin{figure}[t]
\centering
\includegraphics[width=.93\linewidth]{sources/monographs/machine-societies/publication/figures/distributed_history.pdf}
\caption{Two valid local grant records can merge into one agreed but invalid allocation record. A later selection of one record does not reverse an already executed conflicting grant. Dashed lines mark delayed exchange, not an atomic commitment mechanism.}\label{MS:fig:distributed-history}
\end{figure}

For the commons, the appropriate refinement is not to eliminate every supplied primitive. It is to specify which layer supplies exclusivity, how urgency exchanges become committed, and which notification history makes a debt actionable against an office. Observations and research proposals may tolerate asynchronous set merging. Exclusive slot grants require a stronger contract. Different social acts can therefore require different coordination semantics inside the same society.


\section{The computability boundary of constitutional verification}\label{MS:sec:computability}
Finite constitutional models permit exhaustive reachability analysis when states, initial conditions, and transition rules are effectively given. The conclusion changes when amendments may install arbitrary programs with unbounded memory. This section explicitly leaves the finite-state setting. Its purpose is to delimit universal verification claims, not to weaken the finite theorems.

Consider a class of deterministic programs with an effectively encoded initial configuration, a total computable one-step transition, and a distinguished forbidden event. Call a program safe when its unique run never produces that event. Assume the class can effectively simulate any Turing machine one step at a time.

\begin{MStheorem}[No universal constitutional safety decider]\label{MS:thm:undecidable}
There is no total algorithm that decides safety for every program in this class. In the same class, occurrence of the forbidden event is semidecidable, whereas safety is not semidecidable in general.
\end{MStheorem}
\begin{proof}
Given a Turing machine $M$ and input $w$, construct a program $P_{M,w}$ that simulates $M(w)$ and produces the forbidden event exactly when the simulation halts. Each simulation step is computable and terminates. The constructed program is safe exactly when $M(w)$ does not halt. A total safety decider would decide halting by complementing its answer, contradicting the classical undecidability result~\cite{MS:turing}.

Occurrence of the event is semidecidable by executing successive steps until it appears. If safety were also semidecidable for this universal class, run the two semidecision procedures in interleaved steps. Exactly one would eventually accept, yielding a total decider. Equivalently, the effective construction would semidecide nonhalting as well as halting, which is impossible.
\end{proof}

The theorem already holds with one deterministic run. It does not rely on strategic deception, nondeterminism, machine consciousness, or the freedom to modify a running verifier. A sufficiently expressive ordinary program language is enough. It is an application of computability theory, not a new limit specific to AI.

\subsection{Useful verification survives the impossibility}
There are several different, legitimate guarantees. In a finite effective transition graph, compute the reachable states and check that none is forbidden. For a fixed finite step bound and finitely enumerable branching, explore every permitted prefix. In a restricted language, prove that the permitted syntax induces a decidable verification problem. In a proof-carrying design, accept an amendment only when a supplied certificate passes a specified checker~\cite{MS:necula}.

These guarantees must not be conflated. A fixed time bound does not establish safety after the bound. A syntactically finite program need not have a finite state space. A finite state declaration is not an effective model-checking procedure if evaluating an edge condition may itself fail to terminate. A sound, effectively checkable proof system may certify many safe programs while failing to certify some others. Indeed, a sound and complete recursively enumerable certificate system for the universal class above would semidecide safety, contradicting Theorem~\ref{MS:thm:undecidable}.

\subsection{Self-amendment does not eliminate the trusted contract}
A certificate must identify the program, policy predicate, interpreter, and version of the checker against which it was established. Proving that an amendment preserves today's permission table does not automatically prove that a subsequent interpreter amendment preserves its meaning. A separately justified transport certificate can carry the guarantee forward; unrestricted substitution cannot.

Nor does the impossibility result establish that every real institution needs a permanent external administrator. It says that no single terminating procedure can classify every unrestricted program by this safety property. A society may choose a smaller rule language, maintain an explicitly incomplete proof discipline, distribute checking, or adopt bounded commitments. These are substantive design choices with different costs and residual risks, not ways to make the universal decision problem disappear.

\section{The information limit of decentralized governance}\label{MS:sec:decentral}
The viability construction in Section~\ref{MS:sec:constitution} assumes that its decision rule sees the full state. A plural institution may distribute the information and the action permissions among different offices. Existence of a safe full-information action then need not imply existence of an implementable decentralized strategy. Observation-based strategies and architecture-dependent synthesis address this distinction in established theory~\cite{MS:chatterjee,MS:pnuelirosner}.

Fix two hidden environmental conditions $\theta\in\{0,1\}$ and a class of admissible protocols. For one protocol, let $H$ be the tuple of all local histories up to the decision, together with the private and shared random tapes used to choose actions. The tapes have their specified law and cannot encode extra information about $\theta$. The joint action is a measurable function of $H$ even though no participant need be able to read all of $H$. Let $P_0$ and $P_1$ be its laws under the two conditions.

\begin{MStheorem}[Indistinguishability obstruction to joint safety]\label{MS:thm:decsafe}
Let $B_0$ and $B_1$ be disjoint subsets of a common finite joint-action alphabet, containing the actions that meet a declared requirement under conditions $0$ and $1$. If the laws of $H$ are within $\eta$ in total variation, then
\[
P_0(B_0)+P_1(B_1)\leq1+\eta.
\]
Consequently at least one condition has failure probability at least $(1-\eta)/2$. If every admissible protocol has identical joint-history laws under the conditions, no admissible protocol can meet the requirement almost surely in both.
\end{MStheorem}
\begin{proof}
Total variation does not increase under the function mapping histories and tapes to joint actions, so $\MSTV(P_0,P_1)\leq\eta$. Therefore
\[
P_0(B_0)+P_1(B_1)
\leq P_1(B_0)+\eta+P_1(B_1)\leq1+\eta,
\]
since the two safe sets are disjoint. Subtracting the success probabilities from two gives a total failure of at least $1-\eta$, so one failure probability is at least half that amount. Set $\eta=0$ for the final assertion.
\end{proof}

The quantification over \emph{every} protocol is needed for an architecture-level impossibility. Equality for one protocol proves only its limitation. The theorem applies to a deadline decision whose required joint actions depend on hidden information. It does not rule out later service after new information arrives.

\subsection{Why separate marginal ignorance is insufficient}
\begin{MScounterexample}[Local ignorance with a perfectly correct joint action]\label{MS:ex:jointviews}
Let $U$ be a fair bit. Agent 1 observes $H_1=U$ and agent 2 observes $H_2=U\oplus\theta$. Each individual observation has the same law for both values of $\theta$. Nevertheless, the joint action $(a_1,a_2)=(H_1,H_2)$ always satisfies $a_1\oplus a_2=\theta$. The joint observation laws differ and have disjoint supports. Equal marginal laws therefore do not imply the hypothesis of Theorem~\ref{MS:thm:decsafe}.
\end{MScounterexample}

The distinction is particularly important in distributed systems: correct coordinated behavior can depend on correlations that no individual can decode. The theorem uses indistinguishability of the complete protocol view, not absence of information in each component separately.

\subsection{A full-information safe action that the institution cannot select}
Suppose a slot must be assigned immediately to office $A$ when $\theta=0$ and to office $B$ when $\theta=1$. Every predecision observation and message is constant under both conditions. A full-information scheduler can satisfy the requirement exactly, but any allowed decentralized protocol has the same grant distribution under both. With an available shared fair coin the best worst-case success is $1/2$, attained by a random grant. A deterministic rule fails in at least one condition.

The requirement includes \emph{correct service by the deadline}. Idling would preserve a weaker no-double-allocation invariant, but would not meet that requirement. Thus this example separates weak safety from timely service rather than calling inactivity an allocation success. New information, a later deadline, or an authorized disclosure channel can remove the obstruction. None is supplied by a written constitution alone.


\addtocontents{toc}{\protect
}

\chapter{Continuity, compatibility, and knowledge}\label{MS:book:chapter-95}

\section{Temporal obligations, contracts, and office succession}\label{MS:sec:obligations}
Permissions describe available actions; obligations constrain continuations. The debt created by an urgency exchange is not satisfied merely because the creditor remains authorized to ask for service. It requires a particular future event. A finite obligation record can be written as
\[
o=(\mathsf{id},i,j,\varphi,d,\mathsf{authority},\mathsf{status}),
\]
where $i$ and $j$ are the debtor and beneficiary offices, $\varphi$ specifies a fulfillment event, and $d$ is a deadline under an agreed clock. Status distinguishes pending, fulfilled, breached, and validly discharged obligations. Cancellation must itself satisfy the declared authorization rule; deleting a record is not fulfillment.

For a bounded deadline $D$, a trace requirement has the form
\[
\mathbf G\bigl(\mathsf{created}(o)\Rightarrow
\mathbf F_{\leq D}\mathsf{settled}(o)\bigr).
\]
Here $\mathsf{settled}$ means fulfilled or validly discharged. The temporal operators mean ``at every position'' and ``within at most $D$ subsequent clock steps.'' The obligation identifier binds the later event to the same commitment. No implication of ethical adequacy follows from admitting discharge in this particular contract.

\subsection{Monitoring deadlines without confusing safety and liveness}
A finite monitor stores pending identifiers, remaining deadlines, and terminal statuses. For a finite identifier pool, at most $M$ active obligations, and deadlines bounded by $D$, this is a finite state machine. Events at a deadline are processed before the clock advances beyond it. A missing fulfillment becomes a violation only once that last permitted event has passed. Resource unavailability, clock failure, and unresolved appeals must have declared transition semantics rather than being silently ignored.

A bounded-deadline violation has a finite bad prefix and is therefore a safety violation in the trace sense. The unbounded demand that every pending obligation is \emph{eventually} fulfilled is instead a liveness requirement, absent a deadline or other finite violation witness. Safety and liveness are distinct established classes~\cite{MS:alpern}. A monitor can detect a missed deadline; it does not by itself make fulfillment physically possible or strategically attractive.

Let $m:S\to Q$ map a complete institutional state to its obligation-monitor state, and let $\nu(a)$ give the monitored event label of action $a$. Correct realization requires
\[
m(F(s,a))=\delta(m(s),\nu(a))
\]
for every enabled action, with the clock update included in $\nu$. The monitor has its own bad-state predicate.

\begin{MSproposition}[Obligation-preserving succession and translation]\label{MS:prop:obligationtransport}
Consider two deterministic realizations with the same monitored event alphabet and monitor update $\delta$. Suppose a translation $T$ preserves the initial monitor state and the two executions follow the same event sequence, with each realization satisfying the monitor equation. Then their monitor states agree after every finite prefix. In particular, fulfillment, valid discharge, and deadline violations agree at corresponding positions. An office handover that preserves the obligation record is one instance, provided later events and clocks satisfy these hypotheses.
\end{MSproposition}
\begin{proof}
The monitor states agree initially. If they agree after a prefix, applying the same deterministic update $\delta$ to the same next event preserves equality. Induction gives equality at every prefix, and therefore equality of every status predicate on the monitor state.
\end{proof}

The theorem does not establish that a successor will execute the required event, nor does it grant the successor a credential. It isolates the semantic condition under which a replacement does not erase or alter what remains owed. In the commons, registered succession and the repayment-first rule supply those separate components.

\subsection{Resources are not copied by copying a promise}
A resource claim needs a conservation rule in addition to its temporal monitor. Grant records are descriptive data, whereas a spendable slot token has a unique consumption transition. Forking an office's software can duplicate its knowledge of a claim without duplicating the underlying resource. A model that counts each copied record as a new budget has changed its resource semantics, not discovered a new source of supply.

Likewise, merging two archives with contradictory obligation statuses must not invent a seamless settlement history. The disagreement is a typed residual that requires an authorized reconciliation. Model replacement, transfer of an office, discharge of a debt, and revocation of a credential are different operations, even when an implementation performs them in one transaction.

\section{Culture, inheritance, and continuity through replacement}\label{MS:sec:culture}

The strongest form of machine cultural continuity is neither transcript copying nor behavioral imitation. It is the transfer of \emph{usable methods and revision commitments} across changes of participants and substrates. A successor may preserve the predecessor's scientific method by rejecting one of its inherited conclusions. Conversely, an agent that perfectly repeats old answers while ignoring later counterevidence may reproduce data but not the culture that produced them.

Let a history $H$ be summarized by an inheritance map $C(H)$. A practically meaningful retained state can be divided into an active canon $M^{\rm act}$ for immediate use, an archive $M^{\rm arc}$ for recovery and criticism, and a provenance and uncertainty record $P$. These stores can overlap physically, but serve distinct decision contracts. The canon keeps interaction costs low. The archive makes forgotten exceptions, alternative explanations, and evidence recoverable. Provenance records why a claim should be trusted, not merely that it was inherited.

\subsection{What must be preserved}

A cultural contract $\mathcal Q$ may specify future queries, operations, resource policies, amendment procedures, and recognized classes of correction. Two histories may be compressed to the same state only if the required future decisions and authorized operations remain indistinguishable under $\mathcal Q$, exactly or within a stated tolerance. Section~\ref{MS:sec:memory} supplies a finite-memory lower bound: distinct future-relevant histories cannot be combined in an exact deterministic summary. Culture is therefore an information-preservation problem with a \emph{normative choice of relevant futures}, not a license to summarize arbitrarily.

A cultural lineage should track at least four noninterchangeable notions: similarity of resulting behavior, causal transmission of predecessor-specific data, authentic provenance of that transfer, and continuity of valid permissions or obligations. Historical descent can hold while methods diverge; identical methods can arise independently; valid authority may fail to transfer even when every bit is copied. Nothing in these relations settles a question about consciousness or subjective continuity.

\subsection{Transport between representational generations}

Machine populations may change their internal languages and storage formats. The right question about a translator is whether it preserves \emph{operations}, not only whether its output looks fluent. Suppose generation $j$ uses a metric state space $(Z_j,d_j)$, update operator $F_j:Z_j\to Z_j$, and translator $T_j:Z_j\to Z_{j+1}$. The one-step compatibility defect is
\begin{equation}\label{MS:eq:transport-defect}
\varepsilon_j=\sup_{z\in Z_j}d_{j+1}\bigl(T_jF_j(z),F_{j+1}T_j(z)\bigr).
\end{equation}
An exact zero defect means the representation change intertwines the declared operation. It does not mean all information is preserved; injectivity or an appropriate reconstruction contract is a separate issue.

\begin{MSproposition}[Accumulation of cultural transport defects]\label{MS:prop:transport}
Assume each $T_j$ is $L_j$-Lipschitz and $\varepsilon_j<\infty$ for $j=0,\ldots,m-1$. Define $T_{0:m}=T_{m-1}\cdots T_0$. For every $z\in Z_0$,
\begin{equation}\label{MS:eq:cumulative-transport}
\begin{split}
&d_m\bigl(T_{0:m}F_0(z),F_mT_{0:m}(z)\bigr)\\
&\qquad\le\sum_{j=0}^{m-1}\varepsilon_j\prod_{k=j+1}^{m-1}L_k,
\end{split}
\end{equation}
with empty products equal to $1$.
\end{MSproposition}
\begin{proof}
For $m=1$ this is the definition of $\varepsilon_0$. Suppose the bound holds through generation $m-1$. Insert $T_{m-1}F_{m-1}T_{0:m-1}(z)$ between the two terms at generation $m$. Lipschitz continuity of $T_{m-1}$ bounds the first difference by $L_{m-1}$ times the prior defect, while~\eqref{MS:eq:transport-defect} bounds the second by $\varepsilon_{m-1}$. The recurrence $D_m\le L_{m-1}D_{m-1}+\varepsilon_{m-1}$ unfolds to~\eqref{MS:eq:cumulative-transport}.
\end{proof}

This familiar telescoping-error argument acquires a useful interpretation here. Even individually adequate intergenerational translators can accumulate operation drift, especially when later translations amplify earlier mismatch. Conversely, uniformly nonexpansive translators bound the drift by the sum of local defects. There is no implication that the resulting chain has preserved all obligations or epistemic provenance; those must be included in the state and operations if they are part of the contract.


\subsection{Accumulation requires executable preservation, not only storage}
A culture can inherit a method without preserving the ability to use it. Its language may remain readable while its required tool is unavailable or its resource permission has been revoked. For a fixed finite query family $\mathcal Q$, define a certified method to be an executable procedure with a checked correctness obligation on a stated subset of queries. The checker, semantics, and environmental premises belong to the contract. An unverified assertion is not promoted into this certified class merely by being archived.

Let $\mathcal C_j\subseteq\mathcal Q$ be the queries for which generation $j$ has at least one admitted executable method satisfying the contract. This is an attainable-capability set, not a measurement of what an arbitrary method-selection policy actually achieves.

\begin{MSproposition}[Conditional cumulative culture]\label{MS:prop:culturegrowth}
Suppose every certified method available at generation $j$ has an admitted realization at generation $j+1$ preserving its query semantics, correctness premises, required resources, and execution rights. Suppose newly added methods satisfy the same declared correctness contract. Then
\[
\mathcal C_j\subseteq\mathcal C_{j+1}.
\]
The inclusion is strict exactly when at least one newly available method correctly covers a previously uncovered query. The conclusion remains compatible with correction or removal of an inherited \emph{uncertified} belief.
\end{MSproposition}
\begin{proof}
Take a query in $\mathcal C_j$ and a witnessing certified executable method. The preservation assumption supplies a method satisfying the same contract at generation $j+1$, so the query remains covered. New methods enlarge the union of covered domains exactly when one covers a query outside the previous union. The definition of that union does not include uncertified beliefs.
\end{proof}

The assumptions do the work: archival inclusion alone is insufficient. The proposition is not a theorem of automatic scientific progress. It establishes monotonicity of a protected capability envelope; an actual society may select a worse method, change its objectives, lose access to its tools, or discover that a supposedly correct premise was false. Such a discovery changes which methods satisfy the contract and must not be concealed to preserve an artificial progress score.

Fixed storage also limits exact accumulation. If a future-query contract distinguishes $N$ relevant histories, fewer than $N$ deterministic summary states cannot preserve all of them by Proposition~\ref{MS:prop:memory}. A $B$-bit summary cannot represent more than $2^B$ such classes. Sustained expansion therefore requires more storage, a narrower contract, recoverable external archives, or accepted loss. For a fixed finite $\mathcal Q$, even lossless strict capability growth eventually saturates. Open-ended cultural development requires an explicitly expanding query family rather than a conclusion extracted from a finite monotonicity lemma.

\subsection{Fission, merging, and the refusal of false continuity}

Forking a model produces several descendants, not several automatically valid copies of a scarce license, budget, or authorization. Merging two archives need not merge two policies, preferences, norms, or identities. A merger can preserve conflicting decisions instead of rewriting them as though agreement had always existed. A useful inherited record therefore has both continuities and \emph{unresolved branches}. Institutional honesty may require stating that there is no single successor policy until an admissible reconciliation procedure has occurred.


\section{Endogenous problem formation and intellectual culture}\label{MS:sec:agendas}
A population that only inherits fixed answers has a different organizational structure from one that can generate questions, preserve unresolved problems, and revise what counts as a useful method. A new method may make a previously inaccessible question tractable. A new question may reveal a missing method. The resulting coupling concerns both a stock of capabilities and an agenda of inquiry.

Problem--solution coevolution already appears in work such as POET, which couples generated environments with methods for solving them~\cite{MS:poet}. Our question is institutional: who may introduce a problem, what evidence can close it, how competing agendas obtain resources, and what happens to unresolved commitments when participants are replaced? This section specifies a theoretical mechanism; it does not claim that open-ended intellectual development follows from naming those mechanisms.

Write $\mathcal C_t$ for available method records and $\mathcal Q_t$ for recognized question records. A question record includes a statement or schema, prerequisites, admissible forms of answer or refutation, provenance, and a resource request. Methods carry domains of validity and the evidence supporting them. The coupled state may evolve as
\[
\mathcal C_{t+1}=F(\mathcal C_t,\mathcal Q_t,I_t,u_t),\qquad
\mathcal Q_{t+1}=G(\mathcal Q_t,\mathcal C_t,I_t,u_t),
\]
where $u_t$ records permitted external inputs and $I_t$ governs admission, allocation, and revision. No common utility function is assumed: different offices can value different investigations. The external input and allocation policy must remain visible rather than being absorbed into a claim of autonomous discovery.

\subsection{A finite closure theorem, not an open-endedness theorem}
Take finite universes $\mathcal C_*$ and $\mathcal Q_*$. Consider a fixed collection of production rules of the form ``if these finitely many records are present, add this record.'' Prerequisites are positive: no rule requires the absence of a record. The content and correctness of those rules are assumptions. An added method is certified only if the specified certification procedure actually supports it.

Let $\Phi$ simultaneously add all conclusions whose prerequisites are satisfied, retaining the input records. Order pairs $(\mathcal C,\mathcal Q)$ by componentwise inclusion.
\begin{MSproposition}[Least closure of a finite problem--method system]\label{MS:prop:agendaclosure}
Starting from $(\mathcal C_0,\mathcal Q_0)$, iteration of $\Phi$ reaches the least closed pair containing that initial pair. There are at most
\[
|\mathcal C_*\setminus\mathcal C_0|+|\mathcal Q_*\setminus\mathcal Q_0|
\]
strict growth steps. A sequential rule schedule reaches the same closure if every rule that remains enabled with an absent conclusion is eventually applied.
\end{MSproposition}
\begin{proof}
Positive prerequisites make $\Phi$ monotone, and retention makes it inflationary. Every strict change adds at least one previously absent record, giving the bound. If a closed pair contains the initial pair, induction using monotonicity shows it contains every iterate; hence the limiting fixed pair is least. A sequential schedule remains inside that same closure and can make only finitely many additions. If its eventual record set were not closed, some rule would remain enabled with a missing conclusion, contradicting the scheduling assumption. Its limit is therefore the same least closed pair.
\end{proof}

The finite bound is important: this construction cannot certify unbounded novelty. It derives consequences within supplied finite universes and rules. Countably many possible new records, changing rule languages, or discovery of new question schemas require a larger model. Deleting obsolete claims also breaks inflationarity. Those extensions are legitimate research problems, but their behavior cannot be inferred from this finite closure proof.

For the commons, a certified scheduler can make deadline-sensitive allocation questions meaningful; a counterexample to that scheduler can create a new repair obligation; a solution can become a method available to a successor office. The record of why a problem was posed matters as much as its final answer. Declining to fund a question should not be represented as proving it solved.

\subsection{Changing the agenda can manufacture apparent progress}
\begin{MScounterexample}[Perfect success by replacing the questions]\label{MS:ex:agendadrift}
A fixed method solves an easy problem $e$ and fails on a hard problem $h$. Under the agenda $\{h\}$ its success fraction is zero. Under the agenda $\{e\}$ its success fraction is one, despite no change in the method. Thus increased success on an endogenously selected agenda does not imply increased capability on any unchanged problem family.
\end{MScounterexample}

A theoretical account should distinguish capability on a preserved task contract, expansion of what can be represented, and selection of which questions receive attention. None is a sufficient statistic for the others. An archive of abandoned and unresolved questions prevents changing denominators from being mistaken for intellectual advance. The point is not that agendas must be fixed. It is that their evolution is one of the outcomes to explain.

\subsection{Resource allocation shapes what a culture can discover}
Even an enabled inquiry rule may never execute if its office receives no time or storage. The fairness premise in Proposition~\ref{MS:prop:agendaclosure} is therefore substantive. An institution can preserve a problem syntactically while indefinitely starving the work needed to resolve it. Rights to propose, rights to be heard, and actual access to a finite budget are distinct.

Likewise, persistence of a practice is not evidence of its truth. In an illustrative two-type copying model, let accurate and inaccurate practices reproduce at fixed rates $r_T,r_F>0$, with positive initial counts $N_T,N_F$. If $r_F>r_T$, the accurate type's share is
\[
\frac{N_T r_T^t}{N_T r_T^t+N_F r_F^t}
=\frac{1}{1+(N_F/N_T)(r_F/r_T)^t}\longrightarrow0.
\]
The model does not describe all cultural evolution. It establishes that selection for reproduction alone does not logically impose selection for accuracy. An institution needs a separately specified relationship between evidence, admission, correction, and reproduction to make that inference.

A developing computational culture would therefore be characterized not only by retained answers but by maintained procedures for posing questions, allocating inquiry, recognizing failure, and revising those procedures. Whether these processes yield sustained improvement under heterogeneous interests and changing environments remains an open theoretical problem, not an automatic property of interacting models.

\section{Institutional lifetimes and generational turnover}\label{MS:sec:lifetime}
Institutional continuity is a trajectory property. It can outlast an occupant without lasting forever. Fix an operating and repair policy and let $V$ be a finite set of states satisfying a declared continuity contract. The state includes any deadline monitors and fault information needed for the Markov property. Failure is the first exit from $V$; an institution later rebuilt is a new episode under this first-exit convention.

Let $Q$ be the substochastic transition matrix on $V$, so the missing row mass represents exit. Let $\mu$ be an initial row probability vector supported on $V$ and define $\tau=\inf\{t\geq1:s_t\notin V\}$.

\begin{MSproposition}[Survival and expected institutional lifetime]\label{MS:prop:lifetime}
For every integer $t\geq0$,
\[
\MSProb_\mu(\tau>t)=\mu Q^t\mathbf1.
\]
If the spectral radius satisfies $\rho(Q)<1$, then
\[
\MSE_\mu[\tau]=\mu(\mathrm{Id}-Q)^{-1}\mathbf1.
\]
\end{MSproposition}
\begin{proof}
The unnormalized row law of trajectories that remain in $V$ through $t$ is $\mu Q^t$, by induction. Summing it gives the survival probability. The tail-sum formula for a positive integer-valued random variable gives $\MSE\tau=\sum_{t\geq0}\MSProb(\tau>t)$. Under $\rho(Q)<1$, the matrix geometric series converges to $(\mathrm{Id}-Q)^{-1}$, giving the result.
\end{proof}

A closed recurrent class inside $V$ can invalidate the spectral-radius hypothesis and permit infinite lifetime on some trajectories. The formula is not to be used by inverting a singular matrix. Conversely, a small per-epoch failure probability can still imply eventual exit. Finite-horizon robustness is not immortality.

\subsection{A transparent repair example}
Take a healthy state $H$, a degraded but still viable state $D$, and an absorbing failure state. In one epoch, $H$ degrades with probability $f$; $D$ repairs with probability $r$ and fails with probability $d$. Thus
\[
Q=\begin{pmatrix}1-f&f\\r&1-r-d\end{pmatrix},
\qquad 0<f\leq1,\quad r\geq0,\quad d>0,\quad r+d\leq1.
\]
Every nonfailure state has a positive-probability path to failure, so the finite substochastic chain is transient. Solving the two first-step equations gives
\[
\MSE_H\tau=\frac{f+r+d}{fd}.
\]
For $f=1/10$, $r=2/5$, and $d=1/10$, the expected first exit is 60 epochs. With $r=0$ and the other parameters fixed it is 20 epochs. These are exact values of the stated chain, not estimates of real machine populations.

\begin{figure}[tbp]
\centering
\includegraphics[width=.84\linewidth]{sources/monographs/machine-societies/publication/figures/institutional_survival.pdf}
\caption{Survival probabilities $\mu Q^t\mathbf1$ for the analytical two-state model, initially healthy, with $f=d=1/10$. Increasing repair from $r=0$ to $r=2/5$ extends expected lifetime from 20 to 60 epochs under otherwise fixed transitions. Neither curve is empirical evidence of institutional persistence.}\label{MS:fig:survival}
\end{figure}

\subsection{Office continuity is distinct from occupant survival}
Suppose, separately, that each original occupant is replaced at each epoch with independent probability $\gamma$. Its probability of still occupying its role after $t$ epochs is $(1-\gamma)^t$. Only if handovers preserve the institutional state and do not change the assumed $Q$ can the institutional survival law above remain unchanged. Independence between turnover and degradation is a simplifying assumption, not an implication of copying state.

This separates two questions: whether an original component remains and whether a specified operational pattern remains. A continuously serviced office may survive several occupants. An unchanged model can preside over the dissolution of its institution. A valid lifetime contract should require continued service or some other nontrivial function; mere survival of an inert archive is a different event.

\section{Fast interaction and slow institutional dynamics}\label{MS:sec:multiscale}

The lifetime model describes how long an organization remains inside a contract. A different question asks when the institution itself has a useful lower-dimensional dynamics despite continual changes in its occupants' activity. An office may execute many tasks between constitutional reviews. A stable rule can therefore coexist with rapidly changing messages, beliefs, queues, and model states. Separating these timescales is useful only when the fast state loses the relevant dependence on its initial condition before the next review.

Two-timescale Markov models and averaging have a substantial literature~\cite{MS:yinzhang}. The finite block model below supplies a self-contained sufficient condition. It does not assert that any particular population mixes quickly, nor does it take a singular-perturbation limit without defining one.

\subsection{A block model with a declared review schedule}
Let $\mathcal I$ and $\mathcal E$ be finite sets of institutional and fast states. At review epoch $n$ the institution is $I_n=i$ and the fast state is $X_n=x$. Keep $i$ fixed for a predeclared number $m\geq1$ of fast steps with kernel $K_i$. Write $\widehat X_n$ for the resulting endpoint. The next institution is sampled through a row-stochastic revision kernel $R_i(\widehat X_n,j)$. A subsequent fast-state reset may depend on the endpoint and adopted institution, but cannot revise the already specified institutional outcome.

Assume $K_i$ has a stationary law $\mu_i$ and
\begin{equation}\label{MS:eq:uniform-fast-mixing}
\sup_{i,x}\MSTV\bigl(K_i^m(x,\cdot),\mu_i\bigr)\leq\varepsilon.
\end{equation}
The actual one-review law from $(i,x)$ and its stationary average are
\begin{align}
P(i,x;j)&=\sum_y K_i^m(x,y)R_i(y,j),\label{MS:eq:actual-review}\\
\overline R(i,j)&=\sum_y\mu_i(y)R_i(y,j).\label{MS:eq:averaged-review}
\end{align}
The endpoint is the only fast input to this particular revision rule. A rule that uses the full trajectory, accumulated costs, or a data-dependent stopping time requires a different state or a different averaging argument.

\begin{MStheorem}[Finite-block institutional averaging]\label{MS:thm:block-averaging}
Under~\eqref{MS:eq:uniform-fast-mixing}, every actual conditional one-review law, given an institutional history ending at $i$, is within $\varepsilon$ of $\overline R(i,\cdot)$. Let $(J_n)$ be the Markov chain with kernel $\overline R$ and the same initial institutional distribution. Then, for every integer $T\geq0$,
\begin{equation}\label{MS:eq:institutional-path-error}
\MSTV\bigl(\MSLaw(I_0,\ldots,I_T),\MSLaw(J_0,\ldots,J_T)\bigr)
\leq 1-(1-\varepsilon)^T\leq\min\{1,T\varepsilon\},
\end{equation}
where $0\leq\varepsilon\leq1$ may be assumed and the empty product at $T=0$ is $1$. No independent initialization of the fast state at successive reviews is required.
\end{MStheorem}
\begin{proof}
Applying the common stochastic kernel $R_i$ cannot increase total variation, so each row~\eqref{MS:eq:actual-review} is within $\varepsilon$ of~\eqref{MS:eq:averaged-review}. Given any institutional history of positive probability, the current fast state has some conditional distribution. Averaging the uniform row bound over it preserves the bound. Couple initial institutional states identically. Whenever the two institutional histories still agree, maximally couple their next-state laws; the probability of disagreement at that step is at most $\varepsilon$. The chance of maintaining agreement for all $T$ steps is at least $(1-\varepsilon)^T$, by multiplying the successive conditional lower bounds. The coupling inequality gives the first estimate. The union bound gives the second.
\end{proof}

The conclusion is about the institutional history at review epochs. It does not preserve every intervening action, establish an equilibrium, or imply that harm between reviews is impossible. Required audit events and safety monitors must be included in a joint output contract if their histories are to be approximated. A small one-review error also needs a horizon: it does not give a nontrivial bound uniformly over arbitrarily many reviews.

A transparent kernel family is
\[
K_i=(1-\beta_i)\mathrm{Id}+\beta_i\mathbf 1\mu_i,
\qquad 0<\beta_i\leq1,
\]
where every row of $\mathbf 1\mu_i$ is $\mu_i$. Its $m$th power is $(1-\beta_i)^m\mathrm{Id}+[1-(1-\beta_i)^m]\mathbf 1\mu_i$. Thus the mixing error is bounded by $(1-\beta_i)^m$. For a fair two-state $\mu_i$ and $\beta_i=1/2$, the exact worst-row error is $2^{-m-1}$. With $m=8$, the linear bound over twenty reviews is $20/512=5/128$. These are analytical values for the stated kernel, not estimated mixing times for neural agents.

\begin{figure}[t]
\centering
\includegraphics[width=.92\linewidth]{sources/monographs/machine-societies/publication/figures/review_averaging.pdf}
\caption{Exact finite-block error bounds for the fair two-state refresh kernel with refresh probability $1/2$. The plotted horizon is twenty institutional reviews. These bounds compare institutional paths at review epochs, not all fast actions.}\label{MS:fig:review-averaging}
\end{figure}

\subsection{Why fast activity is not sufficient}
\begin{MScounterexample}[Fast observations do not justify adaptive stopping]\label{MS:ex:stopping-bias}
Let each fast observation be a fresh fair bit, independent of its predecessor. At every fixed positive time its distribution is exactly stationary. Review instead at the first observed $1$, or at a fixed cap $M\geq2$ if no $1$ appears earlier. Let the review select the institutional state equal to its endpoint bit. The probability of adopting $1$ is $1-2^{-M}$, not $1/2$. Its distance from the fixed-time stationary prediction is $1/2-2^{-M}$.
\end{MScounterexample}

The stopping rule selects favorable observations. It violates the fixed-block premise even though every deterministic-time sample mixes in one step. One can include the stopping policy and accumulated information in an enlarged model, but cannot substitute the same endpoint average unchanged.

A second failure occurs when the fast kernel retains several closed regimes. If $K_i$ is the identity on two fast states and the revision reports that state, the proposed institution-only model loses a distinction that never mixes away. Any single averaged law has worst-row error at least $1/2$. More rapid execution of the same identity transition changes nothing. The retained macrostate must include the regime, or the approximation must acknowledge its error.

\subsection{Transporting an enforcement certificate through a reduced model}
Suppose an evidence model used at a review has true laws $P_C,P_D$ and reduced laws $\overline P_C,\overline P_D$ with certified bounds
\[
\MSTV(P_a,\overline P_a)\leq\varepsilon_a,\qquad a\in\{C,D\}.
\]
Such bounds follow from the block argument when each conduct-dependent fast kernel satisfies its own uniform mixing bound and the same declared endpoint-to-evidence channel is used in the actual and averaged model.

\begin{MSproposition}[A robust margin for reduced-model enforcement]\label{MS:prop:robust-deterrence}
Let $s^\star(o)$ sanction precisely where $\overline P_D(o)>\overline P_C(o)$ and put $\overline d=\MSTV(\overline P_D,\overline P_C)$. In the model of Section~\ref{MS:sec:enforcement}, this fixed policy makes compliance weakly preferable whenever
\[
g\leq F(\overline d-\varepsilon_C-\varepsilon_D).
\]
Conversely, if $g>F\min\{1,\overline d+\varepsilon_C+\varepsilon_D\}$, no observation-based sanction with penalty $F$ deters the deviation. Also,
\[
\MSE_{P_C}s^\star\leq\MSE_{\overline P_C}s^\star+\varepsilon_C.
\]
\end{MSproposition}
\begin{proof}
The expectation of a $[0,1]$-valued function changes by at most total variation. Hence the true difference of sanction probabilities is at least $\overline d-\varepsilon_D-\varepsilon_C$. Equation~\eqref{MS:eq:deterrence-gap} gives sufficiency. The triangle inequality bounds the true total variation above by $\overline d+\varepsilon_C+\varepsilon_D$, capped at one; Proposition~\ref{MS:prop:deterrence-tv} then gives impossibility above that value. The final inequality is the same expectation bound under compliance.
\end{proof}

When the lower margin is nonpositive it supplies no guarantee for positive $g$. A separate false-positive protection must include the final inequality or use a reduced policy with slack below its cap. The result concerns one common fixed policy and one specified deviation. It is not a license to optimize separately for incompatible deviations, replace missing observations with stationary assumptions, or declare an institution safe because its average trajectory looks stable.

The useful societal interpretation is limited but substantive. Under uniform relaxation and a controlled review schedule, institutional dynamics can be predicted without retaining every fast interaction. The reduced description preserves specified laws to a declared horizon. Persistence here is explanatory separation of scales, not independence from the underlying agents or infrastructure.


\section{Jurisdictions, overlap, and local--global obstruction}\label{MS:sec:jurisdictions}

A large machine population need not be administered by a universal institution. Different groups may sustain rules for different tasks, resources, or trust domains. Such \emph{polycentric} organization provides possible advantages in locality and adaptation, but introduces a precise compatibility question: can locally acceptable rules be jointly realized when groups interact?

There are at least three meanings of compatibility. First, two representations may agree on shared observations. Second, their authorized operations may agree when transported through a common interface. Third, their joint constraints may admit a globally executable state. The third is stronger than pairwise local satisfiability. An apparently minor cycle of rules can obstruct any globally consistent execution.

\subsection{An exact finite obstruction model}

Consider an undirected graph $G=(V,E)$ whose vertices are institutional positions. Each position carries a binary coordinate $x_i\in\mathbb F_2$ and each edge $\{i,j\}$ imposes the relative constraint
\begin{equation}\label{MS:eq:edge-rule}
x_i\oplus x_j=b_{ij},\qquad b_{ij}\in\mathbb F_2.
\end{equation}
The coordinates might represent a toy authorization bit or a binary convention, but the mathematical model itself makes no claim about a realistic legal code.

\begin{MSproposition}[Cycle criterion for binary institutional compatibility]\label{MS:prop:cycle}
The constraints~\eqref{MS:eq:edge-rule} admit a simultaneous global assignment $(x_i)_{i\in V}$ if and only if the XOR of the labels $b_{ij}$ along every graph cycle equals $0$. When an assignment exists, each connected component has exactly two assignments, related by complementing every coordinate in that component.
\end{MSproposition}
\begin{proof}
If an assignment exists, XOR the edge equations around a cycle. Every vertex coordinate appears twice and cancels, so the edge-label XOR is zero. Conversely, choose a spanning forest and assign $0$ to one root in each component. Propagate the vertex values along forest edges using~\eqref{MS:eq:edge-rule}. Forest edges are satisfied by construction. For a nonforest edge, the unique forest path between its endpoints and that edge form a cycle. Vanishing cycle parity implies that the propagated values satisfy the nonforest edge as well. Finally, a connected solution is fixed by choosing the value at one root; either root value gives an assignment and there are no others.
\end{proof}

\begin{figure}[t]
\centering
\includegraphics[width=.82\linewidth]{sources/monographs/machine-societies/publication/figures/cycle_obstruction.pdf}
\caption{A minimal local--global obstruction. Each displayed pairwise constraint is satisfiable by itself, but their conjunction is inconsistent: XOR around the cycle gives $1$ on the right and $0$ on the left. This is an elementary binary model, not a universal encoding of institutional norms.}\label{MS:fig:cycle}
\end{figure}

This result is a standard parity fact on graphs, closely related to cycle-space algebra and discrete cocycles. Its value here is methodological. It shows why \emph{pairwise satisfiable} cannot be equated with \emph{globally realizable}. For a triangle with $b_{12}=0$, $b_{23}=0$, and $b_{31}=1$, every individual edge requirement has a solution, but the three cannot hold together. A translator between two agents could appear successful while a three-party federation is impossible under the combined rules.

The obstruction generalizes in several directions. For nonbinary constraints one may need constraint-satisfaction or logical satisfiability methods. For typed overlapping interfaces one may formulate gluing conditions for partial states and compatibility of transition operators. For probabilistic policies, exact satisfiability may be replaced by minimizing a weighted violation or by checking whether local marginal laws admit a joint extension. None of these generalizations is automatic: a consistent family of local marginal distributions on pairwise overlaps need not admit a globally consistent joint distribution on an arbitrary cycle.

\subsection{Operational gluing and federation}

Even when a global static state exists, a federation may fail dynamically. Local rule applications could commute in one context and conflict in another. A translator $T_{ij}:Z_i\to Z_j$ should preserve the operations relevant to the interaction, approximately satisfying $T_{ij}F_i\simeq F_jT_{ij}$. A federation requires \emph{both} representational compatibility and admissible operational coordination. Static agreement alone does not define an executable constitution. Nor does faithful gluing prove that the institution should be centralized: a compatible federation may deliberately preserve multiple veto points, modes of appeal, and local autonomy.

The graph example also clarifies a useful mathematical program: institutional incompatibility can be measured as a \emph{residual} left by trying to realize local constraints globally. One can ask for the least-cost repair set of rules, for certificates that no repair respecting protected obligations exists, or for an interface transformation that changes the representation without changing the obligations. These are natural places for constraint solving, formal verification, and algebraic topology to meet a theory of machine institutions, without claiming that every social disagreement is a parity equation.


\section{Noncommutative institutional holonomy}\label{MS:sec:holonomy}
The binary cycle criterion is the simplest case of a more general distinction between local translations and global consistency. In a federation, an interface may relabel roles, permute resource categories, or change coordinates in a protocol. These transformations need not commute. Their composition around a loop can therefore matter even when every individual interface is invertible.

The group-theoretic foundation is established synchronization theory, not a new use of the word society. Gao, Brodzki, and Mukherjee give a geometric account in terms of bundles and holonomy~\cite{MS:gaosync}. Here we provide the finite graph argument directly and separate consistency of a complete reference frame from compatibility of one particular institutional state.

\subsection{Conventions and exact synchronization}
Let $\Gamma=(V,E)$ be a finite connected graph. Each oriented edge $i\to j$ carries $g_{ij}\in G$, where $G$ is a group and $g_{ji}=g_{ij}^{-1}$. The element $g_{ij}$ transports coordinates from $i$ to $j$ by a left action. For a path $\gamma=(v_0,\ldots,v_k)$, define
\[
H_\gamma=g_{v_{k-1}v_k}\cdots g_{v_0v_1}.
\]
The rightmost map acts first. If $v_k=v_0$, this product is the path's holonomy. An exact system of global frames is a family $h_i\in G$ satisfying
\[
g_{ij}=h_jh_i^{-1}\qquad\text{on every edge}.
\]

\begin{MStheorem}[Exact frame synchronization]\label{MS:thm:holonomy}
Global frames exist if and only if $H_\gamma=e_G$ for every closed walk. When they exist, they are unique up to common right multiplication $h_i\mapsto h_i k$. It suffices to check the rooted fundamental loops associated with one spanning tree.
\end{MStheorem}
\begin{proof}
If $g_{ij}=h_jh_i^{-1}$, adjacent factors cancel in every path product, giving $H_\gamma=h_{v_k}h_{v_0}^{-1}$ and hence identity on a closed walk. Conversely, choose a root $r$, set $h_r=e_G$, and define $h_i$ as the product on the unique tree path from $r$ to $i$. Tree edges satisfy the required identity. For any non-tree edge $i\to j$, the rooted loop consisting of the path to $i$, that edge, and the reversed path from $j$ has product $h_j^{-1}g_{ij}h_i$. Its being identity gives $g_{ij}=h_jh_i^{-1}$. Thus the fundamental-loop tests suffice. Any other solution obeys the same propagation from its root value, so it is $h_i k$ for a common $k$.
\end{proof}

Changing local conventions by elements $k_i$ replaces $g_{ij}$ by $k_jg_{ij}k_i^{-1}$. A loop based at $r$ then has holonomy $k_rH_\gamma k_r^{-1}$. In particular, identity holonomy is invariant under local changes of coordinates. An arbitrary binary cycle basis with unordered multiplication is not a valid replacement for rooted fundamental loops in a noncommutative group.

\subsection{A three-jurisdiction example}
Take $G=S_3$, acting on three role labels. Let the translations on the oriented triangle $A\to B\to C\to A$ be
\[
g_{AB}=(12),\qquad g_{BC}=(23),\qquad g_{CA}=e_G.
\]
Every edge is a valid bijection. With the specified composition convention, the round trip is $(23)(12)\ne e_G$. Thus no single assignment of complete role frames realizes all three interfaces. Replacing the final edge by $((23)(12))^{-1}$ repairs this loop, but whether that edit is authorized is a separate institutional question.

\begin{figure}[t]
\centering
\includegraphics[width=.88\linewidth]{sources/monographs/machine-societies/publication/figures/noncommutative_holonomy.pdf}
\caption{Three individually invertible interfaces need not synchronize. The ordered round trip is $(23)(12)$, not the identity. The group model concerns the declared role-label contract; it does not certify permissions, provenance, or the legitimacy of an interface edit.}\label{MS:fig:holonomy}
\end{figure}

\subsection{A shared state is weaker than a shared frame}
Nontrivial holonomy need not prevent \emph{every} compatible state. Let $G$ act on a set $D$, and seek local states $x_i\in D$ with $x_j=g_{ij}\cdot x_i$. A permutation can move some labels while fixing others. A family that uses only a fixed label can be compatible even though its complete coordinate frames are not.

\begin{MSproposition}[Compatible states and holonomy fixed points]\label{MS:prop:holonomysection}
A compatible state family exists if and only if there is $x_r\in D$ fixed by every holonomy based at the root $r$. Each such $x_r$ uniquely determines a compatible family by transport along a spanning tree. For a free group action, existence of a compatible family implies identity holonomy; without freeness that implication can fail.
\end{MSproposition}
\begin{proof}
A compatible family transported around a closed walk returns to its starting value, so every based holonomy fixes $x_r$. Conversely, define each $x_i$ by the tree-path product acting on $x_r$. The rooted-loop fixed-point condition gives $g_{ij}\cdot x_i=x_j$ on non-tree edges as well as tree edges. The tree paths also give uniqueness for a chosen root state. Under a free action, only the identity fixes a point, proving the final implication.
\end{proof}

\begin{MScounterexample}[Nontrivial frame defect with a compatible state]\label{MS:ex:fixedholonomy}
On a triangle, let one oriented edge carry $(12)\in S_3$ and the other two carry identity. Holonomy is nontrivial, so global frames do not exist. Nevertheless, the constant state $x_A=x_B=x_C=3$ is compatible because $(12)$ fixes $3$. Failure to align all role meanings is stronger than failure to agree on one unused or invariant role.
\end{MScounterexample}

This distinction prevents overinterpreting an obstruction. A federation may support one narrow service while being unable to transport its full permission structure coherently. A zero defect for the narrow service should not be advertised as complete institutional interoperability.

\subsection{A quantitative lower bound on interface error}
Suppose $G$ has a bi-invariant metric $d$, so both left and right multiplication preserve distance. For candidate frames $h_i$, write
\[
\delta_{ij}=d(g_{ij},h_jh_i^{-1}).
\]
\begin{MSproposition}[Loop defect bounds unavoidable edge error]\label{MS:prop:holonomybound}
For a closed walk $\gamma$ with $k\ge1$ traversals,
\[
d(H_\gamma,e_G)\le\sum_{(i,j)\text{ traversed by }\gamma}\delta_{ij}.
\]
In particular, some traversed edge has error at least $d(H_\gamma,e_G)/k$ for every candidate frame family.
\end{MSproposition}
\begin{proof}
Gauge-transform each edge to $\widetilde g_{ij}=h_j^{-1}g_{ij}h_i$. Bi-invariance gives $d(\widetilde g_{ij},e_G)=\delta_{ij}$. The transformed loop product is $h_r^{-1}H_\gamma h_r$, at the same distance from identity as $H_\gamma$. Repeatedly applying the triangle inequality and invariance to a product bounds that distance by the sum of its factors' distances from identity. Since $k\max\delta_{ij}\ge\sum\delta_{ij}\ge d(H_\gamma,e_G)$ along the walk, the maximum error is at least $d(H_\gamma,e_G)/k$.
\end{proof}

The sum counts repeated edge traversals. Bi-invariance is an explicit assumption and cannot be replaced silently by an arbitrary metric. Normalized Hamming distance on a finite permutation group is one suitable example. This estimate constrains possible synchronization accuracy; it does not select an authorized repair, preserve scarce resources, or prove that a globally synchronized convention is desirable. Those requirements remain separate coordinates of the residual calculus.

\section{A typed calculus of institutional residuals}\label{MS:sec:residuals}
Institutional interoperability has several independent failure modes. Two implementations can execute the same action while disagreeing about its authorization. They can preserve authorization while erasing a debt, or preserve the debt while duplicating its backing resource. A fluent translation can retain every conclusion while dropping the evidence needed to reconsider it. A single numerical compatibility score obscures these distinctions.

We formalize a restricted compositional calculus, rather than assuming that every mismatch is measured in one natural unit. Three realizations $j=0,1,2$ have metric state spaces $(Z_j,d_j)$ and a common finite action alphabet $\mathcal A$. For algebraic comparison, each action has a total update $F_j^a$: a disabled action goes to an explicit rejection state. Translations map rejection to rejection. A separate permission predicate records which actions are genuinely admitted. Totalization is a bookkeeping device and does not authorize the disabled action.

\subsection{Common semantic coordinates}
For each component type $c\in\{P,O,R,E\}$, let
\[
g_j^c:Z_j\longrightarrow(E_c,\delta_c)
\]
be a contract observable in a common metric codomain. $P$ records permissions over the common action alphabet; $O$ records obligation-monitor state; $R$ records resources in common physical or accounting units; and $E$ records the required evidentiary and provenance structure. A change of units, action names, or obligation identifiers requires an explicitly specified alignment before these comparisons can be made.

Possible metrics include the discrete metric on permission vectors, a declared metric on obligation-monitor states, an $\ell^1$ distance on common-unit resource vectors, and symmetric-difference cardinality on required evidence records. For evidence, record equality is structural preservation under a trust model, not a truth certificate. If the payload or inferential dependencies matter, they must be included in the observable rather than only its identifier.

For a translation $T:Z_i\to Z_j$, define
\[
\begin{split}
\varepsilon_{\rm dyn}(T)
 &=\sup_{z,a}d_j\bigl(TF_i^a(z),F_j^aT(z)\bigr),\\
\varepsilon_c(T)
 &=\sup_z\delta_c\bigl(g_j^cT(z),g_i^c(z)\bigr)
 \quad(c=P,O,R,E).
\end{split}
\]
The institutional residual profile is
\[
\mathfrak D(T)=
(\varepsilon_{\rm dyn},\varepsilon_P,\varepsilon_O,
\varepsilon_R,\varepsilon_E).
\]
Its entries are not added together. They have different meanings and may have different units. A one-sided authorization check can additionally distinguish newly granted forbidden actions from lost legitimate permissions. The symmetric permission residual measures both changes and is deliberately stronger.

\begin{MStheorem}[Composition of typed institutional residuals]\label{MS:thm:residual}
Let $T:Z_0\to Z_1$ and $U:Z_1\to Z_2$ have finite residuals under the same semantic coordinates. Suppose $U$ is $L$-Lipschitz. Then
\[
\varepsilon_{\rm dyn}(UT)
\leq L\varepsilon_{\rm dyn}(T)+\varepsilon_{\rm dyn}(U),
\]
and, separately for every $c\in\{P,O,R,E\}$,
\[
\varepsilon_c(UT)\leq\varepsilon_c(T)+\varepsilon_c(U).
\]
Consequently translations with zero residual in every coordinate compose to translations with zero residual under this contract.
\end{MStheorem}
\begin{proof}
Insert $UF_1^aT(z)$ between $UTF_0^a(z)$ and $F_2^aUT(z)$. The triangle inequality, then Lipschitz continuity of $U$, bound the two terms by $L\varepsilon_{\rm dyn}(T)$ and $\varepsilon_{\rm dyn}(U)$. For an observable type $c$, insert $g_1^cT(z)$ between $g_2^cUT(z)$ and $g_0^c(z)$. Both comparisons occur in the same metric space $E_c$, so their bounds add directly. Taking suprema proves the claims.
\end{proof}

For a longer chain, the dynamic bound has the same weighted telescoping form as Proposition~\ref{MS:prop:transport}; each semantic-coordinate bound is the sum of its local defects. A zero discrete permission defect gives exact permission preservation. A small real-valued state distance by itself does not. Without a common codomain and fixed alignment, the semantic triangle argument is unavailable.

This is a preservation calculus built from elementary metric estimates and commuting diagrams, in the same broad tradition as compositional abstraction~\cite{MS:rischel}. Its institutional content lies in the chosen typed contract. A zero profile need not imply injectivity, equality of all private information, legitimate institutions, or true beliefs. Those additional properties are not created by naming the profile a certificate.

\subsection{Failures the profile separates}
In a one-step commons interface, a translator might preserve the allocation $(1,1,1,1)$ while interpreting every received signature as a valid amendment approval. Its current resource defect can be zero while its permission defect is nonzero. A second translator might preserve amendment permissions but replace the pending debt record by the empty monitor state. Its obligation defect is nonzero before the altered repayment changes an allocation. A third can preserve both records while duplicating a spendable token, violating the resource coordinate.

A mismatch is not automatically malicious or repairable. Two institutions may deliberately recognize different authorities. A debt-preserving merger may be impossible within the available budget. A correct response can be to refuse a composition and retain the unresolved residual, rather than force agreement by dropping the protected constraint.

\subsection{Authorized repair and a representation-invariant obstruction}
Let a finite set of proposed repairs carry declared nonnegative costs. A repair is admissible only if it respects the protected obligations and the authority to modify each component. The repair problem is to find an admissible candidate with the required residual bounds at minimum cost. An empty feasible set is a meaningful outcome. The definition neither supplies an efficient solver nor licenses changing the protected requirements to make the problem solvable.

The binary jurisdiction model admits a sharper statement. A local convention change $x_i'=x_i\oplus c_i$ transforms its edge labels to
\[
b_{ij}'=b_{ij}\oplus c_i\oplus c_j.
\]

\begin{MSproposition}[Cycle residual and minimum authorized repair]\label{MS:prop:cyclerepair}
The XOR of edge labels around a cycle is unchanged by every local convention change above. On a single simple cycle with odd XOR, suppose an editable subset of edges can be flipped at nonnegative costs $w_e$, with other edges protected. An authorized repair exists exactly when the editable subset is nonempty. The minimum cost is then $\min_e w_e$ over editable edges. If all edges are protected, no change of local binary conventions can produce a compatible global assignment.
\end{MSproposition}
\begin{proof}
In the XOR around a cycle each $c_i$ occurs twice and cancels, preserving the residual. Changing odd parity to even requires an odd number of edge flips. Every such repair includes at least one editable edge and costs at least the cheapest editable-edge cost because costs are nonnegative. Flipping that cheapest edge is sufficient by Proposition~\ref{MS:prop:cycle}. With no editable edge parity cannot change, and local convention changes also preserve it.
\end{proof}

For several overlapping cycles, one edit can affect multiple residuals. Independently repairing each cycle need not be globally optimal, and the single-cycle cost formula does not generalize without further analysis. This distinction connects representational transport, protected structure, and obstruction resolution without treating every institutional conflict as a parity problem.

\section{Epistemic institutions: testimony, provenance, and disagreement}\label{MS:sec:epistemic}

An information-sharing population can become less reliable as it becomes more interconnected. A false claim may be repeated across many agents until repetition is mistaken for independent corroboration. A true observation may be discarded because its origin lacks authorized status. An ordinary summary may be silently promoted to a rule. These are not merely limitations of isolated models: they arise at the interfaces among memory, authority, and social evidence.

\subsection{Knowledge is not the same as available information}

The joint information accessible to a hypothetical observer does not imply that any participant possesses that information, knows that others possess it, or can coordinate on it. Distributed knowledge, mutual knowledge, and common knowledge are distinct notions in distributed epistemic logic~\cite{MS:halpernmoses}. Even if two agents' records jointly determine a fact, an institution may need an authenticated disclosure procedure before that fact may legitimately be used. More bandwidth alone does not produce common knowledge under arbitrary uncertainty about delivery.

A useful epistemic institution therefore specifies which evidence types are admissible, what constitutes a retraction, how contradictions remain visible, which claims can be treated as commitments, and how provenance influences confidence without becoming a substitute for verification.

\subsection{Correlated testimony cannot be counted as independent votes}

Suppose errors $E_1,\ldots,E_n$ in scalar estimates are centered, have common variance $\sigma^2>0$, and have common pairwise correlation $\rho\in[0,1]$. Then
\begin{equation}\label{MS:eq:correlated}
\operatorname{Var}\left(\frac1n\sum_{i=1}^n E_i\right)
=\frac{\sigma^2}{n}\bigl(1+(n-1)\rho\bigr).
\end{equation}
This elementary identity follows by expanding the variance and counting $n(n-1)$ off-diagonal covariance terms. Relative to the variance of an average of independent equal-variance estimators, the \emph{variance-equivalent} effective sample size is
\[
n_{\mathrm{eff}}=\frac{n}{1+(n-1)\rho}.
\]
If the errors are perfectly correlated, adding copies does not reduce variance. This equation does not give a general probability that a majority is right, and real error structures need not have one uniform correlation parameter. It nevertheless defeats the claim that numerical agreement automatically equals independent support. Shared training histories, copied observations, and a common archive may all create dependence that is invisible in the final vote count.

\subsection{Cheap talk and strategic reporting}

If a message $M$ has the same conditional distribution under two equally possible hidden states $\theta=0,1$, then Bayes' rule gives $\Pr(\theta=1\mid M)=\Pr(\theta=1)$ wherever the posterior is defined. Such a channel carries no information about the distinction, regardless of how fluent or authoritative its reports appear. In strategic environments, senders may have preferences that differ from those of receivers. The theory of strategic information transmission shows why informative communication can depend on preference alignment and signaling structure~\cite{MS:crawfordsobel}. A message being grammatically meaningful is not a truth guarantee.

A report may be false for many reasons: misunderstanding, stale records, an unsuitable abstraction, a noisy sensor, a deliberately misleading tactic, or an institutional incentive to conceal a cost. \emph{Deception} should denote a strategic communicative mechanism for inducing a mistaken belief, not every divergence between a representation and a reported sentence. For opaque systems, the relevant beliefs and intended effects may be difficult to identify. A theoretical account should therefore define the sender's information, feasible messages, objective, and effect on receivers instead of imputing intention from disagreement alone.

\subsection{A layered provenance contract}

A useful evidence record separates four questions:
\begin{enumerate}[label=(\roman*),leftmargin=2.1em]
\item \emph{Origin:} who or what issued a statement, according to which authenticated trail?
\item \emph{Support:} what observations, calculations, or proofs make the content credible?
\item \emph{Authorization:} did the source have authority to make the relevant decision or claim on behalf of an institution?
\item \emph{Independence:} does the evidence add information beyond sources already considered?
\end{enumerate}
These are separate dimensions. A cryptographic signature can authenticate an origin without ensuring accuracy; a proof can be valid relative to erroneous premises; an accurate observer may lack permission to allocate resources; ten copies of one report may count as one evidentiary lineage. Institutions that preserve these distinctions are better positioned to correct errors without turning disagreement into a loyalty test.


\section{Rational information cascades and collective ignorance}\label{MS:sec:cascades}

A society can retain a consistent history, aggregate testimony correctly, and still stop learning. The problem need not be deception. When later participants see only earlier decisions, a decision can be an optimal summary for one actor while discarding information that the next actor would have used. Information-cascade theory supplies a classical account of this failure~\cite{MS:bikhchandani1992}. We develop a symmetric binary example to show exactly where the evidence stops accumulating.

\subsection{A decision-only public record}
Let an unknown state $\theta\in\{0,1\}$ have equal prior probabilities. Participant $t$ receives a private signal $s_t\in\{0,1\}$ with
\[
\MSProb(s_t=\theta\mid\theta)=p,\qquad 1/2<p<1.
\]
Signals are independent conditional on $\theta$, and their accuracy is common knowledge. In sequence, each participant observes previous actions $a_1,\ldots,a_{t-1}$ and its own signal, then chooses $a_t\in\{0,1\}$ to maximize its probability of matching $\theta$. At exact posterior indifference it follows its private signal. The public record contains actions, not private signals, and all participants know the decision rule. There are no conformity rewards, sanctions for dissent, or strategic misreports in this model. This private-signal tie convention is an explicit variant of the random tie-breaking convention in the basic example of Bikhchandani, Hirshleifer, and Welch~\cite{MS:bikhchandani1992}. The probabilities derived below belong to the stated variant, not to every informational-cascade model.

Set $\lambda=p/(1-p)>1$. Before a cascade starts, define an integer public score $z$ so that the likelihood ratio of the observed action history is $\lambda^z$. A new private signal changes its owner's posterior exponent by $+1$ or $-1$. At $z\in\{-1,0,1\}$ the stated tie rule makes the action reveal that signal. Once $z$ reaches $2$, even a contrary signal leaves a strictly positive exponent, so the action is $1$ regardless of private evidence. Every later such action then has probability one under either state and adds no further likelihood. The symmetric statement holds at $z=-2$.

\begin{MSproposition}[Almost-sure locking with a positive error probability]\label{MS:prop:cascade-lock}
In this model a cascade begins almost surely after finitely many revealing decisions. Its probability of selecting the wrong state is
\begin{equation}\label{MS:eq:cascade-error}
\eta(p)=\frac{(1-p)^2}{p^2+(1-p)^2}>0.
\end{equation}
The probability that no cascade has begun after $2r$ decisions is $[2p(1-p)]^r$. After locking, further actions contain no additional information about $\theta$ conditional on the preceding public history.
\end{MSproposition}
\begin{proof}
The posterior argument above identifies the revealing region and the two absorbing action regimes. Starting at score zero, a pair of equal revealing signals reaches $+2$ or $-2$, while an unequal pair returns to zero. Conditional on either true state, a pair locks correctly with probability $p^2$, locks incorrectly with probability $(1-p)^2$, and returns to zero with probability $2p(1-p)$. Conditional independence gives independent repetitions of this two-step cycle until absorption. Since $2p(1-p)<1$, nonabsorption through $r$ cycles has the stated probability and tends to zero. Summing the probability of eventual incorrect absorption gives
\[
\sum_{r=0}^{\infty}[2p(1-p)]^r(1-p)^2
=\frac{(1-p)^2}{1-2p(1-p)},
\]
which is~\eqref{MS:eq:cascade-error}. In an absorbing regime every subsequent action is deterministic given the public history under both possible states, so conditioning on it changes no posterior.
\end{proof}

For $p=3/4$, the eventual error probability is $1/10$. A long run of unanimous actions can therefore preserve the same mistaken decision with no additional evidentiary support. Each participant remains Bayesian-rational under the stipulated information. Unanimity is an outcome of the disclosure protocol, not evidence that a large number of independent signals favored the decision.

\subsection{An evidence-preserving comparison}
Now change only the information contract. Suppose authentic individual signals are available to a decision rule that selects their majority. It need not infer signals from prior actions. Let $\widehat\theta_n$ be any majority choice with an arbitrary fixed rule for ties.

\begin{MSproposition}[Conditional independence permits continuing accuracy gains]\label{MS:prop:signal-majority}
Under the same signal model,
\[
\MSProb(\widehat\theta_n\ne\theta)
\leq \min\left\{1,\frac{p(1-p)}{n(p-1/2)^2}\right\}.
\]
In particular, the error tends to zero as $n$ increases.
\end{MSproposition}
\begin{proof}
Conditional on either value of $\theta$, the indicators $\mathbf1\{s_t=\theta\}$ are independent Bernoulli variables with mean $p$ and variance $p(1-p)$. A wrong majority requires their sample mean to be at most $1/2$, including a possible wrong tie. Chebyshev's inequality bounds this event by $p(1-p)/[n(p-1/2)^2]$. Averaging the identical conditional bounds over $\theta$ proves the result.
\end{proof}

This elementary bound is deliberately conservative. The exact finite error is a binomial tail, with a tie term determined by the decision rule. The comparison does not assume that all signal disclosures are truthful by virtue of being public. Authentication, privacy, observation cost, and reporting incentives are additional contracts. The result isolates the value of preserving usable independent evidence, not a general preference for making private states public.

\begin{figure}[t]
\centering
\includegraphics[width=.92\linewidth]{sources/monographs/machine-societies/publication/figures/cascade_evidence.pdf}
\caption{The binary analytical model with signal accuracy $p=3/4$. For odd population sizes, the error of the last action approaches the cascade floor $1/10$, while a majority of the underlying independent signals continues to improve. Both curves are calculated from finite probabilities, not sampled agent runs.}\label{MS:fig:cascade-evidence}
\end{figure}

\subsection{Designing epistemic institutions rather than enforcing agreement}
The institution chooses what enters its record. Publishing an action without its evidentiary basis may be adequate for coordinating a resource allocation but inadequate for later scientific assessment. A registry can therefore retain a decision, the observations supporting it, their dependencies, and the unresolved alternatives as distinct objects. An office may comply with the current allocation while maintaining a contrary hypothesis about the evidence behind it.

Possible countermeasures include independently acquired observations, protected channels for disconfirming evidence, and periodically reopening a settled question when the evidence contract changes. Their effectiveness depends on whether the new information reaches the actual decision procedure. Repeating the same cached report through a separate channel supplies no new signal. Likewise, randomized prompts do not establish conditional independence of errors. The retained provenance distinctions of Section~\ref{MS:sec:epistemic} remain necessary.

In the commons, an urgency approval can be a binding procedural fact without being an independent observation that the underlying task was genuinely urgent. Subsequent reviewers who see only approvals can rationally overrely on them. Keeping the supporting certificates separate from the sequence of votes makes it possible to ask what was observed, what was inferred, and what was authorized. Collective ignorance is then analyzed as an information-processing failure rather than automatically as dishonesty, irrationality, or a failure to reach consensus.


\section{Public knowledge, notice, and epistemically faithful transport}\label{MS:sec:publicknowledge}
Knowing a rule, knowing that another participant knows it, and having a shared protocol for when it takes effect are different properties. A common record can help establish them, but the existence of that record does not imply any particular agent has accessed it. Distributed epistemic logic gives a precise language for these distinctions~\cite{MS:halpernmoses}; dynamic epistemic logic studies how communicative events change them~\cite{MS:baltag}.

Fix one configuration epoch and a nonempty group $J$ of named decision centers. Let $W$ be a finite set of possible worlds. Each equivalence relation $\sim_i$ records the worlds agent $i$ cannot distinguish using its permitted information. A world may include the epoch, relevant history, and current records. This is an idealized semantics of information, not a claim that an implemented agent can compute every logical consequence.

For a proposition $\varphi$, $K_i\varphi$ holds at $w$ when $\varphi$ holds at every $v\sim_i w$. Let $E_J\varphi=\bigwedge_{i\in J}K_i\varphi$. Define common knowledge by all finite depths,
\[
C_J\varphi=\bigwedge_{n\ge0}E_J^n\varphi,
\qquad E_J^0\varphi=\varphi.
\]
Distributed knowledge instead quantifies over the \emph{intersection} of the agents' indistinguishability classes. Common knowledge is governed by chains in their \emph{union}. Pooling complementary private evidence can therefore reveal a fact that is not public to any participant.

\begin{MSproposition}[Common knowledge as a reachability condition]\label{MS:prop:commonreach}
Let $\mathsf{Reach}_J(w)$ be the worlds reachable from $w$ by a finite chain of relations $\sim_i$ with $i\in J$, including the zero-length chain. Then
\[
w\models C_J\varphi\quad\Longleftrightarrow\quad
\mathsf{Reach}_J(w)\subseteq\{v:v\models\varphi\}.
\]
\end{MSproposition}
\begin{proof}
Induction on $n$ shows that $E_J^n\varphi$ requires $\varphi$ at the endpoint of every chain of $n$ indistinguishability steps. Requiring this for every finite $n$ is exactly the stated reachability inclusion.
\end{proof}

\begin{MScounterexample}[Everyone knows, but it is not common knowledge]\label{MS:ex:everyone}
Let $W=\{0,1,2\}$ and let $\varphi$ hold exactly at $0,1$. Agent $A$ has partition $\{0,1\}\mid\{2\}$, and agent $B$ has partition $\{0\}\mid\{1,2\}$. At world $0$, both agents know $\varphi$. However, the chain $0\sim_A1\sim_B2$ reaches a world where it is false. Thus $E_{\{A,B\}}\varphi$ holds at $0$ but $C_{\{A,B\}}\varphi$ does not.
\end{MScounterexample}

\begin{figure}[t]
\centering
\includegraphics[width=.88\linewidth]{sources/monographs/machine-societies/publication/figures/common_knowledge_chain.pdf}
\caption{At world $0$, each agent's own alternatives satisfy $\varphi$. Their combined indistinguishability graph nevertheless reaches world $2$, where $\varphi$ fails. Reflexive edges are omitted. Shared first-order knowledge is weaker than common knowledge.}\label{MS:fig:knowledge}
\end{figure}

\subsection{An ideal public announcement is a strong update assumption}
For a propositional fact $p$ whose valuation is unchanged by the announcement, an ideal truthful public update restricts $W$ to $W_p=\{w:w\models p\}$ and restricts every $\sim_i$ to that set. This update models that the event, its interpretation, and its receipt are public to the group. In the restricted model, every reachable world satisfies $p$, so Proposition~\ref{MS:prop:commonreach} gives $C_Jp$. The conclusion follows from the update semantics; it is not a theorem that writing a file implements that event.

The restriction to a stable propositional fact matters. An announcement can change the truth of a formula about what someone does not know. Nor does a public fact in one electorate remain common knowledge after admitting a fresh member. The relevant group, epoch, interpreter, delivery conditions, and background assumptions must be stated.

Classical limitations on common knowledge depend on the timing and communication model, and not every useful coordination task requires full common knowledge~\cite{MS:halpernmoses}. A sequential allocation may need only a valid certificate at the point of execution. A simultaneous global activation rule demands a stronger knowledge condition. Bounded acknowledgments and versioned notices can implement specified weaker requirements without pretending to supply infinitely many levels of mutual assurance in an arbitrary network.

\subsection{Notice is an institutional choice as well as an information condition}
A rule may define applicability by a public epoch, by authenticated receipt, or by an acknowledgment procedure. These are different institutional contracts. Whether lack of notice excuses an action is not deduced from $K_i$ alone. It is an explicitly chosen relation between the information state and the permission or obligation monitor. In the commons, an approval can exist in the ledger before a disconnected office knows it exists. The specification must say whether that office is allowed to proceed under an earlier epoch, must wait, or can appeal a late notification. Truth of an adoption record and enforceability against a participant are separate observables.

\subsection{Preserving behavior need not preserve what participants know}
A state quotient that retains allocations may identify worlds with different notification histories. It can therefore preserve a task trace while changing claims of the form ``office $A$ knows that office $B$ has accepted the debt.'' To preserve such claims, a stronger interface is needed.

Consider models on $W_A,W_B$ with the same named agents and atomic predicates. Let $T:W_A\to W_B$ be surjective and preserve the truth of each atom in both directions. For every agent $i$, require
\begin{align*}
w\sim_i^A v&\ \Longrightarrow\ T(w)\sim_i^B T(v) \quad\text{(forth)},\\
T(w)\sim_i^B z&\ \Longrightarrow\ \exists v:\ w\sim_i^A v\text{ and }T(v)=z
\quad\text{(back)}.
\end{align*}
These are the usual bounded-morphism conditions, stated here as a transport contract.

\begin{MSproposition}[Epistemic preservation and reflection]\label{MS:prop:epistemictransport}
Under these conditions, $w$ and $T(w)$ satisfy exactly the same formulas built from the matched atoms, Boolean connectives, $K_i$, and $C_J$. Such transports compose. If the same state maps also satisfy Definition~\ref{MS:def:language}, they preserve the translated operational conversations and the stipulated epistemic formulas at every corresponding reached state.
\end{MSproposition}
\begin{proof}
Use induction on formulas. Atoms and Boolean connectives follow from the assumptions. For $K_i$, forth sends every source alternative to a target alternative, while back lifts every target alternative; applying the induction hypothesis proves both implications. A finite chain of indistinguishability steps is mapped by repeated forth and lifted by repeated back. Proposition~\ref{MS:prop:commonreach} therefore proves the common-knowledge case. Forth composes directly, and two successive applications of back lift any target alternative through the intermediate model. Surjectivity and atomic truth preservation also compose. The operational conclusion then follows from Proposition~\ref{MS:prop:languagetrace} at each reached state.
\end{proof}

The theorem does not cover new agents whose knowledge relations have not been matched, arbitrary target messages outside the translation, or a change in what the atoms mean. It also makes no inference from formal knowledge to trustworthiness, consent, or legitimacy. Its useful conclusion is narrower: dynamic behavior, message effects, and epistemic claims can be preserved together, but each requires its own explicit part of the interface certificate.


\chapter{Authority and joint feasibility}\label{MS:book:chapter-96}

\section{Effective authority, vetoes, and institutional capture}\label{MS:sec:power}
The number of agent processes is not the number of independent centers of authority. One credential can control many processes; several credentials can be required for one consequential operation. An institutional analysis must therefore begin with office rights and operational dependencies rather than a head count.

Fix a nonempty family $\mathcal W$ of minimal approval coalitions for a particular amendment. Each coalition is a nonempty subset of the $n$ registered offices. Approval is monotone: an amendment passes exactly when the set of approving offices contains some member of $\mathcal W$. Assume every office can independently approve or refuse, and that no unmodeled administrator or scheduling dependency bypasses the rule. This is a simple-game representation of approval authority~\cite{MS:taylorzwicker}.

\begin{MSproposition}[Blocking coalitions as hitting sets]\label{MS:prop:blocker}
A coalition $C$ can block the amendment by withholding all of its approvals, regardless of what the other offices do, if and only if
\[
C\cap W\ne\varnothing\quad\text{for every }W\in\mathcal W.
\]
Thus the minimum number of offices with unconditional blocking power is the minimum hitting-set size of $\mathcal W$. Under an equal $q$-of-$n$ rule with $1\leq q\leq n$, it is $n-q+1$.
\end{MSproposition}
\begin{proof}
If $C$ intersects every minimal approval coalition, refusing all approvals in $C$ prevents each such coalition from being contained in the approvers. Conversely, if some $W$ is disjoint from $C$, its offices can approve and pass the amendment despite $C$'s refusal. For the threshold rule, the complement of $C$ can pass exactly when it has at least $q$ members, so blocking requires $n-|C|<q$, or $|C|\geq n-q+1$.
\end{proof}

In the commons, two offices can block a three-of-four amendment, but two cannot enact it. A single office is a veto holder exactly when it belongs to every minimal approval coalition. These are statements about one rule and one event, not complete measurements of social power.

\subsection{An approval rule may conceal another authority structure}
Suppose the nominal three-of-four process cannot commit any decision without a separate administrator's signature. The effective approval coalitions must include that administrator. Applying the threshold formula to only the four visible voters would overstate their independence. Scheduling, admission of evidence, access to the amendment queue, and power to replace the checker can similarly alter the implemented authority structure.

The correct remedy is to enlarge the declared realization and recompute the authority relations, not to infer decentralization from its public diagram. Conversely, a veto can protect a legitimate minority or a resource owner's consent. Small blocking coalitions do not by themselves establish harmful capture. Capture requires an additional protected-purpose or authorization criterion against which the actual control is judged.

An institution can also be procedurally compliant while concentrating benefits or suppressing effective appeal. That is a normative and incentive question distinct from whether a signature threshold was satisfied. A theory of authority should preserve the distinction between \emph{ability to prevent}, \emph{ability to enact}, \emph{ability to set the agenda}, and \emph{ability to change the procedure}.

\section{Founding authority and the limits of internal attestation}\label{MS:sec:founding}
Procedural continuity and authorized origin are not the same property. A constitution can specify how amendments are judged after its adoption without settling who was entitled to adopt it. An internally consistent sequence of certificates can preserve the first record's assertion of authority even when no accessible evidence establishes that assertion.

The question can be stated without selecting a universal theory of legitimacy. Let $\mathcal H$ be a set of possible founding histories, let $a:\mathcal H\to\{0,1\}$ be an externally stipulated predicate such as ``the resource owner authorized this founding event,'' and let $R:\mathcal H\to\mathcal R$ be the record available to the institution. The predicate models one particular authorization criterion, not moral legitimacy in its entirety.

\begin{MSproposition}[When a founding fact is determined by internal records]\label{MS:prop:founding}
There exists a function $f:\mathcal R\to\{0,1\}$ satisfying $f(R(h))=a(h)$ for all $h\in\mathcal H$ if and only if $a$ is constant on every nonempty fibre of $R$. This is a statement of informational determination, not necessarily of computability. If two histories with different labels induce identical distributions of every record available to an evaluator, no evaluator restricted to those records distinguishes them with accuracy above $1/2$ under equal priors.
\end{MSproposition}
\begin{proof}
If $R(h)=R(h')$, any function of the record gives the same answer on both histories, proving necessity. If the label is constant on every nonempty fibre, define $f(r)$ to be that common value and choose arbitrary values off the image of $R$. This gives sufficiency. For the probabilistic statement, equal record laws remain equal after any common randomized decision procedure. Under equal priors, its probability of a correct label is therefore $1/2$.
\end{proof}

A self-authored line saying ``this constitution was authorized'' cannot distinguish two histories that leave identical accessible records. An authenticated receipt from a recognized external authority may distinguish them, but then the recognition of that authority and the key binding are explicit premises. This does not mean that all records are untrustworthy or that an infinite regress makes institutions impossible. It identifies the trust boundary instead of concealing it inside the verifier.

\subsection{Three judgments that should remain separate}
Internal validity asks whether a transition follows the currently accepted formal rules. Historical authorization asks whether the relevant actor actually granted a particular authority under the declared criterion. Normative legitimacy asks whether the criterion and arrangement are justified. A proof system can establish the first relative to its semantics. Evidence and trust assumptions can support the second. Neither calculation uniquely determines the third.

An institution may openly adopt an internally defined founding convention and state the resources and participants to which it applies. What it cannot do by definition alone is change an external ownership fact, manufacture another party's consent, or retroactively authorize actions on behalf of someone absent from the convention. The distinction is between declaring a rule for a model and establishing a fact about its realization.

\subsection{A federation inherits trust boundaries, not just certificates}
When one institution imports another's origin record, it must decide which trust anchors, delegations, revocations, and unresolved disputes it recognizes. A translator can preserve the bytes and inferential structure of a certificate while the receiving institution rejects its root authority. That is not necessarily a translation error. It can be an explicit disagreement about authorization.

The commons therefore needs separate records for the resource owner's founding grant, the adopted internal procedures, and subsequent authorized amendments. A successor can preserve all three and still encounter a disputed external premise. Its honest state includes that dispute. Consistent internal history is evidence of procedural continuity, not an all-purpose certificate of truth or legitimate origin.

\section{Conflict, sabotage, and organizational pathology}\label{MS:sec:security}

A machine society is not guaranteed to be cooperative, benevolent, or stable. Multiple decision centers can cooperate on one resource and compete on another. The same persistent rules that economize on coordination can create concentrated authority, lock in mistakes, or transmit a harmful practice to later participants. Security is therefore not an unrelated afterthought; it is partly the problem of preserving valid relationships between declarations, implementations, permissions, and evidence.

\subsection{Failures occur at different layers}

We distinguish at least five abstract failure surfaces. \emph{Content failures} corrupt beliefs or stored facts. \emph{Interface failures} alter the interpretation or routing of otherwise unchanged content. \emph{Authority failures} allow an unauthorized actor to issue a consequential operation. \emph{Revision failures} change which future evidence or rules are admissible. \emph{Allocation failures} exploit shared resource procedures. These failures can combine: a fabricated record may justify an unauthorized rule change, after which the new rule rejects the evidence needed to expose the fabrication.

The terms \emph{sabotage} and \emph{virus} require restraint. A sabotage claim requires an adversarial objective and a causal route by which the action impairs an institution. A ``schema virus'' is an analogy until one specifies the copied representation, adoption mechanism, susceptibility, persistence, and measured harm. Error propagation can be unintentional; stability can coexist with misinformation; and divergence from a prevailing view may represent a correct new discovery rather than a pathology.

\subsection{A bounded propagation inequality}

A transparent comparison model can make one class of failure mathematically tractable. Let $x_t\in\mathbb R_{\ge0}^n$ record nonnegative deviations at $n$ institutional positions. Suppose an explicitly justified comparison model provides
\begin{equation}\label{MS:eq:propagation}
x_{t+1}\le Bx_t+w_t
\quad\text{componentwise},\qquad B\ge0,
\end{equation}
where $w_t$ represents fresh disturbance and $B$ summarizes one-step propagation. Suppose a vector $v>0$ and $\kappa<1$ satisfy $Bv\le\kappa v$. In the weighted maximum norm $\MSnorm{x}_{v,\infty}=\max_i x_i/v_i$ we obtain
\begin{equation}\label{MS:eq:propagation-bound}
\MSnorm{x_t}_{v,\infty}\le
\kappa^t\MSnorm{x_0}_{v,\infty}
+\sum_{j=0}^{t-1}\kappa^{t-1-j}\MSnorm{w_j}_{v,\infty}.
\end{equation}
The proof is the one-step monotonic estimate followed by induction. Under uniformly bounded disturbance $\MSnorm{w_t}_{v,\infty}\le\bar w$, the last term is at most $\bar w(1-\kappa^t)/(1-\kappa)$. This gives a conditional attenuation result, not a universal contagion threshold. A graph alone cannot supply the matrix $B$; the rates must come from a justified dynamical model. More connectivity can improve detection as well as propagation, and nonlinear strategic behavior can invalidate the linear comparison.

\subsection{Institutional capture and the meaning of power}

Power is not exhausted by reputation or by the number of messages an agent produces. In a computational institution it can arise from privileged read or write access, control of scarce resources, amendment rights, scheduling priority, the right to certify evidence, and the ability to prevent an appeal. A useful power analysis distinguishes \emph{information reach} from \emph{decision reach}: a participant can influence many beliefs while holding no formal permission, or can issue one privileged command with enormous consequences.

A simplified reach functional for agent $i$ and horizon $T$ is the set of consequential institutional coordinates whose future distributions can be changed by some \emph{admissible unilateral intervention} of $i$. This is an intervention-relative property; it must not silently assume that the agent can overwrite a variable to which it has no access. Quantifying authority requires the feasible intervention family and a choice of outcomes. A concentration score built from these reaches can be a diagnostic, but not a proof of illegitimacy.

Formal policy compliance alone does not rule out institutional capture: a rule may be faithfully executed and still privilege a coalition that controls amendments or vetoes. Likewise, a successful self-repair routine may restore a harmful prior state rather than correct it. Any proposed notion of health must include protected freedoms, an error-correction contract, and an explicit source of normative authority. Agreement is not itself the objective.

\subsection{Defensive design implications}

Ordinary security principles retain their relevance: least privilege, complete mediation, separation of authority, auditable transitions, compartmentalized credentials, and recovery from corrupted records~\cite{MS:saltzer}. A proof-carrying decision can narrow the trusted checking base~\cite{MS:necula}, but cannot prove the adequacy of its own external assumptions. Copyable identities create classic Sybil concerns~\cite{MS:douceur}. A fork's assertion that it is a new participant is not automatically a new independent vote or a new allocation of finite resources.

The strongest general principle is \emph{evidence-preserving reversibility where reversal is physically possible}. Exploratory proposals need not receive immediate permission to rewrite canonical rules; reviews should preserve rejected alternatives and reasons; rollbacks should retain an audit trail; and irreversible external actions require a different permission standard from reversible sandbox computations. This is a defensible design stance, not a theorem that one fixed governance policy fits every environment.

\section{Societies of societies: federation, mobility, and scale}\label{MS:sec:ecology}

Once groups can maintain institutions, their boundaries become objects of interaction. A participant may migrate between two groups, carry an archive from one to another, borrow computational resources, or hold obligations in multiple jurisdictions. A coalition may split without dissolving all of its records; two institutions may merge while retaining conflicting rules. Such events define an \emph{ecology of organizations}, not merely a larger population of otherwise unrelated agents.

The relevant inter-organizational interface has more structure than a message format. It includes state translation, identity and provenance assertions, resource transfer conditions, compatible permission semantics, and a means to resolve contradictory commitments. Calling two communication protocols interoperable because their messages parse correctly ignores whether the resulting actions are authorized and what happens to existing obligations.

\subsection{Migration is not unrestricted inheritance}

A migrating process can carry operational knowledge without carrying valid credentials. An organization that accepts all imported permissions merely because an agent has an intact memory confuses history with authority. A safer abstract contract is explicit about which local capabilities are newly granted, which predecessor obligations continue, which are discharged, and which disagreements require arbitration. Copying a state does not copy a budget, physical hardware, a nontransferable credential, or any legal status assigned by people.

\subsection{Scale changes which institutions are economical}

At small scale, bilateral negotiation may be cheap and simple. At larger scale, direct communication can become expensive, so rules may be compressed into roles, standards, and delegated procedures. Hierarchy can reduce latency or message count under explicit network constraints; Section~\ref{MS:sec:hierarchy} also proves that an unrestricted controller can simulate the same finite computation. The correct separation concerns resources, privacy, fault domains, and the work required to discover or maintain useful interfaces, not a new computability class created by social grouping.

Large networks also increase the surface for copied mistakes, inconsistent rule versions, and concentrated control over common infrastructure. Modularity can localize some failures while making cross-boundary repair harder. There is no universal monotone law saying that greater connection makes a machine society more stable, intelligent, or vulnerable. Those effects depend on transition kernels, institutional incentives, topology, and the cost of checking interactions.

\subsection{The absence of an inevitable collective self}

It is tempting to imagine a sequence from many agents to coordination, from coordination to institutions, and from institutions to a unified intelligence. Nothing in the framework implies that final arrow. Some federations are valuable precisely because they prevent one controller from collapsing plural judgments. Others may deliberately adopt a common policy for restricted operations while retaining disagreement on everything else. The existence of a macroscopic description is also not evidence of a first-person perspective. Computational closure, social organization, legal recognition, and subjective experience remain separate questions.

\section{Endogenous roles, adaptive networks, and resilient specialization}\label{MS:sec:specialization}

A continuing office need not be one of a permanently supplied list. A population may establish a verification service, divide a task into specialties, retire a redundant committee, or form a temporary coalition. Its communication graph can change with these decisions. To express that possibility, let $G_t$ specify currently permitted communication links, $r_t$ an assignment of roles, and $c_t$ a competence and capacity record. Their transition belongs to the realized state:
\[
(G_{t+1},r_{t+1},c_{t+1})\sim
\Gamma_{I_t}(\,\cdot\mid G_t,r_t,c_t,H_t),
\]
where $H_t$ denotes the retained information used by the update. A Markov formulation must include a sufficient part of that information in the state rather than treating an unspecified history as an unobserved input.

Adaptive-network theory studies the coupling of node dynamics to changing topology~\cite{MS:grossblasius}; strategic network-formation theory distinguishes beneficial networks from those stable under participants' incentives~\cite{MS:jacksonwolinsky}. Here the additional question is institutional: which training, delegation, link creation, and role transfers are authorized, and which existing commitments survive those changes? Endogenous division of labor requires that participant interaction affects the adopted structure. Writing $\Gamma$ does not show that an externally assigned specialization was discovered.

\subsection{Specialization has a cost and a reachability condition}
Consider two services performed once per period. A generalist arrangement costs $2c_g$ per period. Two specialists cost $2c_s+b$, where $b$ is the required communication and integration cost, and training costs $h$ for each specialist. Over $T$ periods the specialist arrangement has lower total cost precisely when
\[
T\bigl[2(c_g-c_s)-b\bigr]>2h.
\]
This is an instance of institutional amortization, not a universal law of specialization. Verification overhead, replacement, workload variation, and lost redundancy must be added when relevant. A specialist who is faster locally can still increase total cost if every result requires expensive cross-boundary interpretation.

Reachability is separate. Suppose two agents must continuously staff two different roles and each agent has capacity one. A coordinated exchange of their assignments may improve both agents' outcomes. Yet neither agent can move alone while all roles remain staffed and no role is double-booked. Permitting an atomic handover, a temporary reserve agent, or a short service interruption changes the admissible transition graph. The best assignment and a permissible path to it are different mathematical objects, just as in institutional formation and constitutional amendment.

When a finite role-and-network update system admits a common potential that strictly increases on every permitted improvement, the earlier finite termination result applies. There is no general reason for individually beneficial link changes, training choices, and delegated powers to share such a potential. Cycles and stable but inefficient specialization remain possible.

\subsection{A sharp condition for replacement capacity}
Let $J$ be a finite nonempty set of roles that must be filled simultaneously and $A$ a finite set of agents. An edge $(j,a)$ means that agent $a$ is both capable and authorized to fill role $j$ during the relevant interval. Every agent can fill at most one role at a time. A feasible staffing is an injective assignment $f:J\to A$ along these edges. Let $N(S)$ be the agents adjacent to at least one role in $S\subseteq J$.

The question is not merely whether every role has a backup. It is whether backups can cover all roles together after a loss. The classical matching principle originates with Hall~\cite{MS:hall1935}; the following erasure form exposes the exact redundancy requirement.

\begin{MStheorem}[Authorized staffing after any bounded agent erasure]\label{MS:thm:robust-staffing}
Fix an integer $0\leq k\leq |A|$. A feasible staffing exists after removal of every set $F\subseteq A$ with $|F|\leq k$ if and only if
\begin{equation}\label{MS:eq:robust-hall}
|N(S)|\geq |S|+k
\qquad\text{for every nonempty }S\subseteq J.
\end{equation}
The surviving capability and authorization edges are assumed unchanged, and a new assignment may be selected after the erased set is known.
\end{MStheorem}
\begin{proof}
For necessity, if a nonempty $S$ has fewer than $|S|+k$ neighbors, erase $\min\{k,|N(S)|\}$ of those neighbors. Fewer than $|S|$ remain, so no injective staffing of $S$ is possible.

For sufficiency, remove any $F$ of size at most $k$. The surviving neighbors of each $S$ number at least $|N(S)|-k\geq|S|$. We recall the finite matching argument. Choose a maximum matching in the surviving bipartite graph. If a role remains unmatched, explore alternating paths beginning at it, using unmatched edges toward agents and matched edges back toward roles. No reached agent can be unmatched, since that would produce an augmenting path. Let $S$ be the reached roles and $B$ the reached agents. Every neighbor of $S$ is in $B$, and each agent of $B$ is matched to a distinct reached role other than the initial unmatched role. Hence $|N(S)|=|B|<|S|$, a contradiction. The maximum matching therefore fills all roles.
\end{proof}

The theorem permits a simple robustness index. If the initial graph admits a staffing, define
\[
\sigma=\min_{\varnothing\ne S\subseteq J}\bigl(|N(S)|-|S|\bigr).
\]
Then exactly the erasure budgets $k\leq\sigma$ are guaranteed tolerable in this model. This is a static allocation certificate, not a claim that the surviving agents can discover or execute the reassignment before a deadline. That further question uses the communication and permission constraints of Sections~\ref{MS:sec:decentral} and~\ref{MS:sec:social-time}.

\begin{MScounterexample}[Backups for each role need not provide collective resilience]\label{MS:ex:role-backup}
Two roles can each be filled by either of two agents $a,b$. Every role has two capable authorized candidates, and the two roles can initially be filled. After either agent is erased, the one survivor cannot fill both roles simultaneously. The singleton-role conditions for $k=1$ hold, but the two-role set has only two neighbors rather than the required three.
\end{MScounterexample}

For three roles with neighborhoods $\{a,b\}$, $\{b,c\}$, and $\{c,d\}$, every nonempty role set satisfies~\eqref{MS:eq:robust-hall} with $k=1$. One erased agent can be accommodated by reassignment even though no single agent can perform every role. Resilience can reside in a pattern of overlapping specialization rather than a universally capable backup.

\begin{figure}[t]
\centering
\includegraphics[width=.92\linewidth]{sources/monographs/machine-societies/publication/figures/resilient_roles.pdf}
\caption{Three simultaneous roles and four unit-capacity agents. Edges represent both competence and authorization. The neighborhoods satisfy the one-erasure staffing criterion. The guarantee permits rematching after the failure is known; it does not guarantee instantaneous handover.}\label{MS:fig:resilient-roles}
\end{figure}

\subsection{Time-respecting communication cannot be read from an aggregate graph}
Changing topology also changes who can actually learn what. A union of links over an interval discards their ordering. The resulting graph may look connected while failing to carry a particular observation to a decision maker in time.

\begin{MScounterexample}[A connected union with no usable evidence path]\label{MS:ex:temporal-path}
Agent $A$ initially holds an arbitrary bit and the other agents do not. At time one, $B$ sends its current state to $C$. At time two, $A$ sends the bit to $B$. No other communication occurs before $C$ must decide. The union of directed contacts contains the path $A\to B\to C$, but $C$'s received message is independent of the bit. No decoder at $C$ can recover an initially fair bit with accuracy above $1/2$ under this access model.
\end{MScounterexample}

The missing object is a time-respecting causal route, not an additional static connection. The distinction is consistent with the event order of Section~\ref{MS:sec:social-time}. In a changing society, a nominal right to appeal or a nominal verification edge may be ineffective when it is never available during the interval in which evidence matters.

\subsection{Adaptive roles do not automatically preserve institutional plurality}
A new verifier can reduce errors while concentrating epistemic authority. A central routing service can reduce messages while becoming a veto point. Copying a specialist may increase nominal staffing without adding independent failure domains, and a revocation may remove authorization edges even when all technical skills remain. Each of these changes belongs in the state to which the staffing or communication theorem applies.

A useful role-update contract therefore tracks competence, capacity, decision rights, handover obligations, and accessible evidence separately. It can require that a protected staffing margin survive every admitted reorganization, while still allowing the actual role assignment to change. Such a constraint does not choose which roles deserve protection or who may impose the training cost. It makes the consequences of those choices analyzable.

The commons can now add a temporary verification role, train an alternate occupant, and revise the communication schedule without identifying an office with one model. Its social organization lies partly in this maintained division of work. The formal test is whether the adopted arrangement has a feasible and authorized realization through turnover, not whether its role names resemble a familiar human institution.


\section{Joint institutional feasibility and certified continuation}\label{MS:sec:joint-synthesis}
An institution can possess several valid certificates and still lack a feasible way to operate. A safe allocation may require one action, an effective audit another, and a timely repayment a third. Each statement can be correct when its own resource and information assumptions are considered alone. The common implementation fails when those assumptions cannot hold together.

The resulting question is not whether a population can optimize one collective utility. It is whether a declared set of obligations has a jointly executable continuation. The protected requirements can include resource conservation, valid authority, a false-positive limit, a specified incentive margin, service deadlines, and feasible staffing. Their normative selection remains outside the synthesis theorem. Once selected, they must constrain the same action and the same realized history.

This section supplies a finite construction rather than an unrestricted theory of decentralized control. It combines the existing obligation monitors and governance labels with a catalogue of implementable local protocols. The mathematical foundations are the knowledge construction and fixed-point methods for games with imperfect information~\cite{MS:chatterjee}. The additional contribution is to identify which institutional certificates must enter that construction, and what is lost when an abstraction omits them.

\subsection{A finite contract with an explicit information boundary}
Let $S$ be a finite nonempty state space for the complete contract. A state records the relevant institutional mode, resource ownership, bounded obligation monitors, current staffing, and any private memory needed to make transitions well defined. Let $h:S\to Y$ be the observation available in a declared common record at each decision epoch. It need not reveal the private components. The initial information is a nonempty set $B_0$ contained in one fibre of $h$.

A finite nonempty catalogue $U$ contains protocol profiles, not arbitrary commands to the complete microstate. A profile can specify one local response map for each office, with each map reading only that office's permitted observation and stored local memory. It can also specify a committed monitor, a communication schedule, and the steps needed to select and enact that profile. These implementations, their costs, and their decision rights must be established separately. An instruction that requires inaccessible private information is not made implementable by listing it in $U$.

For each $s$, let $A(s)\subseteq U$ contain profiles satisfying the immediate contract at that state. Every profile in this set must pass the same resource, authority, staffing, and monitoring checks. In particular, a monitor included in a profile must satisfy all relevant incentive and false-positive inequalities simultaneously. Let $\Delta(s,u)\subseteq S$ be a nonempty finite set of permitted next states for each $s,u$. The set includes every environment response or fault within the declared guarantee. Inadmissible profiles are never selected.

The contract uses \emph{sure} guarantees: every transition in $\Delta$ is treated as possible, without weighting it by probability. A probabilistic transition model may supply these supports, but an unlikely infinite path is not discarded when proving bounded service. Nor does a deterrence inequality justify silently deleting deviations from $\Delta$. Safety under compliant behavior and safety under specified deviations are different contracts. The latter must include their physical consequences in the transition model.

Let $S_{\rm safe}\subseteq S$ encode the protected state requirements. Breach monitors distinguish fulfillment, authorized discharge, pending debt, and deletion. A finite bound on stored obligations or deadlines is part of the model. An overflow must have an explicit treatment, such as a rejected admission or a breach state. The theorem does not silently replace an unbounded history by a finite one.

A common record at decision epochs is itself an infrastructure assumption. It is not implied by eventual replica convergence. The model applies only when that observation and the selection of a catalogue profile have the declared implementations. The selected profile can then be executed by distinct local agents without exposing their private states. This is a controlled information boundary, not a claim that every decentralized architecture has such a representation.

\subsection{The state of possible institutional realizations}
The relevant public state is the set of complete states still compatible with the observed history. Define
\[
\mathcal B=\{B\subseteq S:B\ne\varnothing,\ h\text{ is constant on }B\}.
\]
These are sets of possibilities, not posterior probability distributions. At $B$, a profile must be admissible for every remaining possibility:
\[
A(B)=\bigcap_{s\in B}A(s).
\]
For $u\in A(B)$ and a next public observation $y$, define
\begin{equation}\label{MS:eq:belief-update}
\operatorname{Post}(B,u,y)=
\left(\bigcup_{s\in B}\Delta(s,u)\right)\cap h^{-1}(y).
\end{equation}
Only nonempty successor sets are possible observations. Write $\operatorname{Succ}(B,u)$ for their collection. Induction on the public history shows that~\eqref{MS:eq:belief-update} retains exactly the complete states consistent with that history and the selected profiles.

This construction does not pool private observations before choosing $u$. It pools only the mathematical possibilities consistent with the common record. Private observations can affect the actions of the local maps already specified by $u$. More elaborate strategies using additional private history require that history, or a sufficient finite memory of it, to be included in the declared realization. Completeness below is relative to this finite catalogue and information model.

\subsection{A continuation kernel for the whole contract}
Set
\begin{align}
W_0&=\{B\in\mathcal B:B\subseteq S_{\rm safe}\},\label{MS:eq:belief-w0}\\
W_{k+1}&=\{B\in W_0:\exists u\in A(B),\
\operatorname{Succ}(B,u)\subseteq W_k\}.\label{MS:eq:belief-wk}
\end{align}
The sequence is decreasing. It stabilizes after at most $|\mathcal B|$ strict removals, and we denote its fixed point by $W_*$. This is the joint continuation kernel. A state is removed not only when it violates an invariant now, but also when every available profile has a possible successor from which the remaining contract cannot be sustained.

\begin{MStheorem}[Joint continuation under a common observation history]\label{MS:thm:joint-viability}
For every $k\geq0$, $B\in W_k$ if and only if there exists a deterministic strategy based on the common observation and profile history that selects admissible profiles and keeps the complete state in $S_{\rm safe}$ for the next $k$ transitions, starting from every $s\in B$. Moreover, $B\in W_*$ if and only if such a strategy exists for an unbounded number of transitions. On $W_*$, a strategy depending only on the current possibility set $B$ suffices.
\end{MStheorem}
\begin{proof}
At horizon zero the condition is exactly $B\subseteq S_{\rm safe}$. At positive horizon the strategy must choose one profile from the same common history, so it must lie in $A(B)$. After each possible observation its continuation must work from the corresponding set in~\eqref{MS:eq:belief-update}. Induction gives exactly~\eqref{MS:eq:belief-wk}, in both directions.

Monotonicity of the predecessor operation makes $W_k$ decreasing. Finiteness implies stabilization. At the fixed point, select for every $B\in W_*$ one witnessing profile whose successors all remain in $W_*$. Induction then proves safety and admissibility for every finite prefix of every realized path, hence for the complete path. Conversely, an unbounded safe strategy is safe for every finite horizon, so its initial set belongs to every $W_k$ and therefore to the stabilized intersection. The witnessing selection is stationary in $B$, although maintaining $B$ generally requires memory of the public history.
\end{proof}

No society-wide preference ordering appears in this theorem. It characterizes joint satisfaction of a declared contract, not the desirability or consent behind that contract. It also gives no favorable complexity bound: the possibility space can contain exponentially many subsets of $S$. A computationally expensive solution is still distinct from a procedure the society can discover and execute within its own resource budget.

The action intersection matters. A full-information controller can choose different safe profiles at two hidden states while the common-history institution cannot choose between them. This is the constructive counterpart of the indistinguishability obstruction in Section~\ref{MS:sec:decentral}. The kernel can therefore be empty even when every complete state has a safe continuation under full information.

\subsection{Safety does not complete a service}
The fixed point can contain a safe idle loop. Outstanding deadlines should therefore be represented by monitors that advance when time passes, not by a permanently pending record. For service without a fixed external deadline, a separate progress condition is needed.

Let $G\subseteq S_{\rm safe}$ be a stipulated completion set. The target is \emph{certified completion}, meaning that the current possibility set lies wholly in $G$. Let
\[
\mathcal G=\{B\in W_*:B\subseteq G\},
\quad
A_*(B)=\{u\in A(B):\operatorname{Succ}(B,u)\subseteq W_*\}.
\]
Define $R_0=\mathcal G$ and
\begin{equation}\label{MS:eq:progress-attractor}
R_{k+1}=R_k\cup
\{B\in W_*:\exists u\in A_*(B),\
\operatorname{Succ}(B,u)\subseteq R_k\}.
\end{equation}
The increasing sequence stabilizes at $R_*$. For $B\in R_*$, let $r(B)$ be its first index of entry.

\begin{MSproposition}[Safe service with a finite progress rank]\label{MS:prop:joint-progress}
A set $B$ lies in $R_k$ if and only if an admissible common-history strategy can ensure certified completion within at most $k$ transitions while preserving the contract indefinitely. A set lies in $R_*$ if and only if such a strategy can ensure eventual certified completion on every possible path while preserving the contract. Before completion, a stationary choice with strictly decreasing $r$ reaches the target within $r(B)$ transitions. After completion, a safety strategy on $W_*$ suffices.
\end{MSproposition}
\begin{proof}
The bounded claim follows by induction from~\eqref{MS:eq:progress-attractor}. At a nonterminal set of rank $r>0$, a witnessing profile sends every successor to rank at most $r-1$. All selected profiles preserve $W_*$, so a safety continuation is available after the target is reached.

For the eventual claim, consider the tree of public histories consistent with a fixed successful deterministic strategy, stopping a branch at certified completion. The tree is finitely branching. If its depths were unbounded, repeatedly choosing a child with descendants at arbitrarily large depths would produce an infinite branch with no completion. Every finite prefix has a consistent complete-state realization by~\eqref{MS:eq:belief-update}. Finiteness and the same branching argument give a compatible infinite complete-state realization, contradicting success on every path. The tree therefore has a finite depth bound, so the bounded claim puts $B$ in some $R_k$. The other direction follows from the rank strategy.
\end{proof}

The completion condition is intentionally epistemic: the common history must determine that the stipulated work is complete. Physical completion without an available completion record is a weaker endpoint. The result does not identify the two. It also does not promise service for requests that the admission policy rejects. Guaranteed access, when required, must be part of the contract rather than inferred from the success of already admitted work.

\begin{MScounterexample}[Almost-sure completion without a sure service bound]\label{MS:ex:sure-service}
At a pending state $p$, one action moves to a completed absorbing state $g$ with probability $1/2$ and returns to $p$ with probability $1/2$. Both states are safe and publicly visible. Completion occurs with probability one and has expected waiting time two, but the support contains the infinite path $p,p,\ldots$. For the sure interpretation of $\Delta$, $W_*$ contains both singleton states while $R_*$ contains only $\{g\}$. Replacing sure completion by almost-sure completion changes the problem.
\end{MScounterexample}

Belief sets suffice for the stated sure safety and certified-reachability objectives. They are not a universal sufficient statistic for strategic partial-observation problems. In particular, broader stochastic games can require information beyond a belief-based strategy~\cite{MS:belieffails}. The quantifiers and the available randomization must be retained when comparing those results with the present construction.

\subsection{A certificate can be checked without finding the largest kernel}
The fixed-point computation finds every viable possibility set in the declared model. A proposed institution may instead supply a smaller invariant family together with its choices. Let $\mathcal C\subseteq\mathcal B$ contain $B_0$, and let $\kappa$ assign a profile to each $B\in\mathcal C$. The checker verifies $B\subseteq S_{\rm safe}$, $\kappa(B)\in A(B)$, and $\operatorname{Succ}(B,\kappa(B))\subseteq\mathcal C$. A separate integer rank can certify a selected service.

\begin{MSproposition}[Soundness of a finite continuation certificate]\label{MS:prop:continuation-certificate}
If the preceding checks hold, $\kappa$ preserves the complete contract from $B_0$ for all transitions. If, in addition, there is a map $r:\mathcal C\to\mathbb N$ such that $r(B)=0$ implies $B\subseteq G$ and every successor of a set with positive rank has strictly smaller rank, then a set contained in $G$ is reached within at most $r(B_0)$ transitions. Checking these finite conditions does not require proving $\mathcal C=W_*$.
\end{MSproposition}
\begin{proof}
Induction preserves membership in $\mathcal C$, admissibility, and safety. Until rank zero is reached, each transition strictly decreases a nonnegative integer. There can be at most $r(B_0)$ such transitions. The rank-zero implication then establishes certified completion. Closure of $\mathcal C$ continues to provide safety afterward.
\end{proof}

An amendment can be one of the profiles in this certificate. It must then be authorized before execution, carry every relevant monitor and resource meaning through the transition, and land in a certified continuation set. Allowing arbitrary rank-preserving amendments before a service can defeat its time bound. Such amendments need their own bounded allowance, or must be excluded until completion. The obligation cannot be declared fulfilled merely by changing the label attached to a pending record.

\begin{MScounterexample}[A safe amendment can destroy every safe continuation]\label{MS:ex:amendment-dead-end}
There are four publicly distinguishable states $r,\ell,g,b$. The safe set is $\{r,\ell,g\}$. At $r$, profile $\mathsf{deliver}$ moves to $g$, while profile $\mathsf{freeze}$ moves to $\ell$. At $g$, $\mathsf{wait}$ loops safely. At $\ell$, the only admissible profile advances an outstanding deadline and moves to breach state $b$. The amendment $\mathsf{freeze}$ preserves the immediate safe-state predicate but not $W_*$. No certificate proving unbounded continuation from $r$ can choose it. The relevant defect is loss of a future capability, not an immediate inconsistent record.
\end{MScounterexample}

\subsection{Which abstractions preserve the ability to govern?}
The existing governance quotient preserves labels, enabled actions, and deterministic updates. For synthesis under incomplete information, an abstraction must also preserve the public observation and the full set of possible successors. Preserving one average outcome or one representative trajectory is insufficient.

Consider a reduced finite contract $(\overline S,\overline h,\overline A,\overline\Delta,\overline S_{\rm safe},\overline G)$ with the same catalogue $U$. Let $\pi:S\to\overline S$ be a surjection. Assume, for every state and every admissible profile,
\begin{align}
h(s)&=\overline h(\pi s), & A(s)&=\overline A(\pi s),\label{MS:eq:synthesis-labels}\\
s\in S_{\rm safe}&\Longleftrightarrow\pi s\in\overline S_{\rm safe}, &
s\in G&\Longleftrightarrow\pi s\in\overline G,\label{MS:eq:synthesis-goals}\\
\pi[\Delta(s,u)]&=\overline\Delta(\pi s,u).&&\label{MS:eq:synthesis-successors}
\end{align}
The notation $\pi[B]$ denotes the image of a set, not a probability pushforward. These are exact assumptions. They must hold for amendments and handovers as well as ordinary actions whenever those operations belong to the contract.

\begin{MStheorem}[Governance abstraction preserves joint realizability]\label{MS:thm:synthesis-transport}
Under~\eqref{MS:eq:synthesis-labels}--\eqref{MS:eq:synthesis-successors}, for every $B\in\mathcal B$ and every horizon $k$,
\[
B\in W_k\Longleftrightarrow\pi[B]\in\overline W_k,
\qquad
B\in W_*\Longleftrightarrow\pi[B]\in\overline W_*.
\]
The progress sets satisfy $B\in R_k\Longleftrightarrow\pi[B]\in\overline R_k$, so finite progress ranks agree. A reduced safe or ranked strategy pulls back by choosing at $B$ the profile assigned to $\pi[B]$.
\end{MStheorem}
\begin{proof}
The observation equality makes $\pi[B]$ a valid reduced possibility set. Admissibility and safe-set membership are preserved and reflected. For every possible next observation $y$, successor equality gives the commuting update identity
\begin{equation}\label{MS:eq:synthesis-post-commute}
\pi[\operatorname{Post}(B,u,y)]
=\overline{\operatorname{Post}}(\pi[B],u,y).
\end{equation}
Indeed, applying $\pi$ to the union of complete successors gives the union of their reduced successors. The observation equality permits intersection with $h^{-1}(y)$ to pass through $\pi$. Nonemptiness is also preserved and reflected.

Induction in~\eqref{MS:eq:belief-w0}--\eqref{MS:eq:belief-wk} now proves the finite-horizon equivalence. Membership at every horizon gives the fixed-point equivalence, since both sequences stabilize. The goal equivalence in~\eqref{MS:eq:synthesis-goals}, the same successor identity, and induction in~\eqref{MS:eq:progress-attractor} prove the progress claims. Every reduced witnessing profile is admissible in every complete state in $B$, and its complete successors have the required reduced images. Pulling back these choices therefore preserves safety and the decreasing progress rank.
\end{proof}

The theorem composes earlier preservation obligations into a strategy claim. It is an exact finite application of quotient and game-synthesis reasoning, not a general solution to decentralized governance. It is stronger than agreement of current outputs because it preserves the existence of a future course of action under the actual information boundary. It remains conditional on the implementation of each catalogue profile and on the coverage of $\Delta$.

Dropping observation preservation can give a reduced controller a signal the real population never receives. Dropping enabled-action preservation can turn an unauthorized command into an apparent solution. Dropping successor equality can erase an inconvenient failure path. Dropping the obligation labels can make an unpaid debt disappear in the quotient. Each change can alter the winning sets while leaving a chosen immediate output unchanged. An approximate transition estimate can support a probabilistic error budget, but cannot by itself justify the exact sure guarantee proved here.

\subsection{The commons: individually feasible promises without a common budget}
Consider a declared variant of the four-office commons with $c$ execution slots remaining in an epoch. Three protected units of work are due before the epoch ends. They include a carried service debt and two other admitted jobs. Each unit consumes one slot. This is a new bounded workload, not a claim that the original allocation rule always creates these obligations.

An ongoing reporting procedure is exposed to two listed deviations, each of gain $1$. Its authorized penalty is $F=2$, and its permitted false-positive rate is $\alpha=1/4$. With no audit, all three observation laws coincide. A one-slot audit has $P_C=(1/2,1/2)$ and the two opposite deviation laws from Counterexample~\ref{MS:ex:pairwise-enforcement}. Their equal mixture reproduces compliance, so no common monitor deters both positive-gain deviations.

A two-slot audit has three outcomes and the stipulated laws
\[
P_C=(3/4,1/8,1/8),\qquad
P_1=(0,1,0),\qquad P_2=(0,0,1).
\]
The common monitor $s=(0,1,1)$ sanctions compliance with probability $1/4$, each deviation with probability $1$, and has margin $2(1-1/4)-1=1/2$ against each. The audit costs and laws are assumptions of the analytical variant. They are not measurements or consequences of the earlier resource allocation.

The joint action consists of one audit mode and the amount of work completed in that epoch. A profile must fit both within $c$. Protected work cannot be discarded or quietly deferred. The following comparison is therefore exact for this catalogue:
\begin{center}
\begin{tabular}{@{}lccc@{}}
\toprule
Contract & Audit slots & Due service slots & Minimum total\\
\midrule
Service only & 0 & 3 & 3\\
Listed deterrence only & 2 & 0 & 2\\
Both at once & 2 & 3 & 5\\
\bottomrule
\end{tabular}
\end{center}
With $c=4$, both component contracts are feasible separately, but their conjunction is not. With $c=5$, the two-slot audit and all three services fit, and one-step completion is certified. With $c=4$ and only two protected service units, the same joint profile also fits. In the general version with $d$ protected units, the exact threshold is $c\geq d+2$.

\begin{figure}[t]
\centering
\includegraphics[width=.82\linewidth]{sources/monographs/machine-societies/publication/figures/joint_service_frontier.pdf}
\caption{Exact feasibility in the bounded commons variant. Every marked point is one integer pair $(d,c)$. Open markers violate the combined budget, while filled markers admit the two-slot audit and all due services. The boundary is $c=d+2$. A feasible audit and a feasible service allocation considered separately do not establish joint feasibility.}\label{MS:fig:joint-frontier}
\end{figure}

For this richer audit, the best common margin at a general false-positive cap is
\[
\Gamma_\alpha=\min\{6\alpha,3/2\}-1.
\]
To see this, average any monitor with the monitor obtained by exchanging outcomes $1$ and $2$. The worst deviation margin does not decrease. Write the two equal sanction probabilities as $t$. Setting the outcome-$0$ sanction probability to zero reduces false positives and weakly improves the margin. The remaining constraint is $t/4\leq\alpha$, with $t\leq1$, and the margin is $3t/2-1$. Thus the formula is exact. The cap $\alpha=1/4$ yields margin $1/2$, but the audit still consumes two execution slots.


% Monograph layout: keep the margin chart with its discussion.
\begin{figure}[!htbp]
\centering
\includegraphics[width=.82\linewidth]{sources/monographs/machine-societies/publication/figures/joint_enforcement_margin.pdf}
\caption{Analytical common-monitor margins for two specified audit channels with $F=2$ and gains $g_1=g_2=1$. The richer three-outcome audit has margin $\min\{6\alpha,3/2\}-1$. For the two-outcome audit whose deviation mixture equals compliance, the best simultaneous margin is $-1$ at every cap. The curves are exact formulas, not fitted observations.}\label{MS:fig:joint-margin}
\end{figure}

The example locates the obstruction instead of hiding it inside a favorable aggregate score. A fifth slot repairs the capacity failure without weakening the promises. An explicitly authorized rescheduling agreement changes the temporal contract rather than pretending that the original deadline was met. A cheaper, validated observation channel changes the audit catalogue. Each route needs its own authority and implementation. The kernel computation certifies only the route actually specified.

\subsection{What the synthesis adds to the theory}
There are now three different witnesses. A monitor certifies simultaneous incentive inequalities in a specified observation model. A continuation certificate supplies an executable choice at every retained public possibility set. A transport certificate shows that the reduced representation preserves those choices, their observations, and their consequences. None can replace the other two.

The result is an institutional realization problem with a finite decision procedure and checkable witnesses. It does not establish spontaneous institution formation, solve the enforcer's incentives, infer a common normative objective, or make a missing public event mechanism appear. It does show how permissions, observations, service obligations, and resource limits can constrain one continuing policy rather than remain independent descriptions of incompatible possibilities.
\FloatBarrier

\FloatBarrier

\chapter{Accountability and the limits of persistence}\label{MS:book:chapter-97}

\section{Accountability, revision, and authorized endings}\label{MS:sec:accountability}

An institution can continue to operate while becoming less answerable to the parties it affects. It can preserve a rule whose authority has expired, fund sanctions but not appeals, or treat unfinished obligations as a reason never to accept an authorized ending. None of these possibilities is excluded by persistence alone. The remaining question is therefore not simply whether a computational organization survives. It is whether the same realized organization respects the limits of its authority, supports effective correction, and can enter a specified suspension or settlement state.

This section extends the common-policy question of Section~\ref{MS:sec:joint-synthesis}. External claims, auditor incentives, confidentiality, review capacity, stopping procedures, and the continuing applicability of certificates must constrain one implementation. They are not additional benefits that can be asserted after the implementation has consumed its entire budget. Nor are they consequences of selecting the word \emph{society}. Each is an explicit part of a contract whose normative selection remains separate from its mathematical feasibility.

\subsection{Affected outsiders and the boundary of authority}\label{MS:sec:external-effects}

Membership, decision rights, resource control, and exposure to consequences define different sets. Let $M$ denote participating members, $V$ positions with specified decision rights, $H$ infrastructure controllers, and $E$ affected parties outside the membership. These sets may overlap where the description allows it. A human resource owner can be in $H$ without holding an internal vote. A person whose reserved service is delayed can be in $E$ without ever communicating with the agents. Conversely, membership does not confer authority over every affected resource.

A useful external-effects contract records
\[
\mathcal B=(Z_{\rm ext},O_{\rm ext},\mathsf{Auth},\mathsf{Effect},
\mathsf{Standing},\mathsf{Remedy}).
\]
Here $Z_{\rm ext}$ contains the external facts relevant to the guarantee, $O_{\rm ext}$ specifies which of them can be observed, $\mathsf{Auth}$ identifies the stipulated authorization for each boundary-crossing act, and $\mathsf{Effect}$ includes the consequences covered by the model. $\mathsf{Standing}$ specifies who may initiate a challenge and on which grounds. $\mathsf{Remedy}$ specifies the actions an authorized review can actually require. The contract can use a finite abstraction of external state, but it must justify that abstraction relative to the protected questions.

This is a representation proposal, not a universal allocation of voting rights. Recognizing an affected party does not automatically give that party an unlimited veto. Conversely, omitting a party from $M$ does not establish consent. The source of authority must select the applicable protections openly. Research on meaningful human control distinguishes responsiveness to relevant human reasons from traceability to responsible human actors~\cite{MS:meaningfulcontrol}. The operational question here is how such requirements enter the action, evidence, and resource contract of the same computation.

\begin{MScounterexample}[Exact internal accounting with an omitted external claim]\label{MS:ex:external-claim}
Four execution slots exist. The internal state contains four member requests and a rule granting one slot to each. In realization $h_0$, all four slots are available to that institution. In realization $h_1$, an external owner has reserved one slot for a nonmember, and has authorized the institution to use only the other three. Suppose the reservation is absent from every internal record. The same four grants satisfy the internal predicate $\sum_iq_i\le4$ in both realizations, but exceed the institution's stipulated authority in $h_1$. Perfect preservation of the internal state cannot distinguish the cases.
\end{MScounterexample}

The failure is not an arithmetic error. It is an omitted premise about what the institution is entitled to allocate. Proposition~\ref{MS:prop:founding} supplies the corresponding information criterion: an internal record can determine the reservation fact only if that fact is constant on its record fibres. A signed statement can improve the evidence only under a stated trust model. It does not manufacture the owner's authorization.

The example also distinguishes \emph{impact} from \emph{authority}. An institution can affect an outsider while lacking authority to do so, and a permitted action can have effects beyond those its designers anticipated. An exact certificate relative to a boundary cannot establish that the boundary includes all consequential effects. Unknown external dependencies remain a modeling risk. Extending the state description is necessary when an omitted claim matters, but simply adding a variable without an observation or authorization mechanism is not a repair.

\subsection{Inspection as a strategic service}\label{MS:sec:strategic-auditor}

The earlier enforcement calculations assume that an announced observation and sanction policy is carried out. A strategic auditor makes that assumption visible. Inspection games and the distinction between simultaneous inspection and committed inspector behavior have an established game-theoretic treatment~\cite{MS:stengelinspection}. The following one-period game is a deliberately small instance, not the recursive model of that source.

A participant chooses compliance $C$ or deviation $D$. An auditor simultaneously chooses no inspection $N$ or inspection $I$. Deviation gives the participant gross gain $g>0$. An inspection detects this specified deviation perfectly and imposes an authorized, collectible penalty $F>g$. The auditor pays cost $c>0$ to inspect and receives reward $b>c$ only on detection. There are no false accusations, side payments, further continuation payoffs, or prior commitment to inspect. Payoffs are risk-neutral and common knowledge. The ordered pair in each entry is the participant's payoff followed by the auditor's payoff.

\begin{center}
\begin{tabular}{@{}lcc@{}}
\toprule
 & No inspection $N$ & Inspection $I$\\
\midrule
Compliance $C$ & $(0,0)$ & $(0,-c)$\\
Deviation $D$ & $(g,0)$ & $(g-F,b-c)$\\
\bottomrule
\end{tabular}
\end{center}

\begin{MSproposition}[Inspection without commitment]\label{MS:prop:strategic-inspector}
Under the preceding strict inequalities, the game has no pure Nash equilibrium and has a unique Nash equilibrium in independently mixed strategies. Its deviation and inspection probabilities are
\[
p^*=\frac{c}{b},\qquad q^*=\frac{g}{F}.
\]
In particular, certain compliance is not sustained in this simultaneous game.
\end{MSproposition}
\begin{proof}
If the auditor inspects with probability $q$, the participant's advantage from deviation is $g-qF$. If the participant deviates with probability $p$, the auditor's advantage from inspection is $pb-c$. At $(C,N)$ the participant prefers $D$, at $(D,N)$ the auditor prefers $I$, at $(D,I)$ the participant prefers $C$, and at $(C,I)$ the auditor prefers $N$. Thus no pure profile is an equilibrium. If either player uses a pure strategy, the other's strict best response yields one of these nonequilibrium pure profiles. Both players must therefore mix. Their two indifference equations uniquely give the displayed probabilities, which lie strictly between zero and one and make every action in each support a best response.
\end{proof}

The residual violation probability does not come from an inaccurate detector. When everyone complies, an auditor paid only for detections prefers not to incur the cost. When nobody inspects, deviation is profitable. It would be incorrect to insert the inspector's optimal response to certain deviation into a certificate that simultaneously assumes certain compliance. These responses belong to different strategic situations.

A different payment contract changes the result. Suppose every \emph{genuinely performed and independently verifiable} inspection receives fee $w>c$, with the detection reward otherwise unchanged. Inspection then has advantage $w+pb-c>0$ at every $p$. The auditor strictly prefers $I$, and the participant strictly prefers $C$ because $F>g$. The unique equilibrium is $(C,I)$. At $w=c$, inspection can be a weak best response to compliance, but the strict implementation claim is lost. This extension requires an affordable fee and a mechanism that verifies the performed service. Paying for an unverified assertion that an inspection occurred does not satisfy the assumption.

The two calculations are conditional alternatives. They do not solve collusion, fabrication by the receipt checker, repeated bargaining, unavailable penalties, or disagreement about legitimate enforcement. A commitment mechanism may move some of these requirements to a different institution. The new institution and its resource costs must then enter the common contract rather than disappear from it.

\subsection{Confidentiality limits the distinctions available to enforcement}\label{MS:sec:privacy-enforcement}

An institution may have good reasons not to collect or disclose every private state. The relevant question is which distinctions a protected record may reveal, not whether more information is always desirable. Let $P_C$ and $P_D$ be the laws of the \emph{complete record available to a sanctioning procedure}. For $\lambda\ge1$, impose
\begin{equation}\label{MS:eq:pairwise-privacy}
 P_D(E)\le\lambda P_C(E),\qquad P_C(E)\le\lambda P_D(E)
 \quad\text{for every event }E.
\end{equation}
The finite setting makes these conditions equivalent to the corresponding pointwise inequalities. When the two underlying inputs are neighbors in a declared pure differential-privacy contract, they hold with $\lambda=e^\varepsilon$. The neighboring-input relation is essential. Differential privacy does not require arbitrary acts of compliance and misconduct to be treated as neighboring inputs. Its standard definition and composition principles are described by Dwork and Roth~\cite{MS:dworkroth}.

Keep the earlier definition
\[
D_\alpha(P_D,P_C)=\max_{0\le s\le1,\ \MSE_{P_C}s\le\alpha}
       (\MSE_{P_D}s-\MSE_{P_C}s),\qquad 0\le\alpha\le1.
\]
The same randomized sanction $s$ must obey the protected false-positive limit and supply the deviation margin.

\begin{MStheorem}[Sharp pairwise privacy envelope for a bounded sanction]\label{MS:thm:privacy-envelope}
For every finite pair satisfying~\eqref{MS:eq:pairwise-privacy},
\begin{equation}\label{MS:eq:privacy-envelope}
D_\alpha(P_D,P_C)\le
 B(\lambda,\alpha):=
 \min\left\{(\lambda-1)\alpha,\frac{\lambda-1}{\lambda+1}\right\}.
\end{equation}
For each $\lambda\ge1$ and $\alpha\in[0,1]$, equality is attained by a pair on two outcomes and a suitable sanction. Consequently, deterrence in the model of Proposition~\ref{MS:prop:deterrence-tv} requires $g\le F B(\lambda,\alpha)$. The envelope is a necessary bound for a specified pair, not a sufficiency claim for every pair satisfying privacy.
\end{MStheorem}
\begin{proof}
Put $a=\MSE_{P_C}s$ and $b=\MSE_{P_D}s$. Summing the pointwise inequalities against $s$ gives $b\le\lambda a$. Summing against $1-s$ gives $1-a\le\lambda(1-b)$. Thus
\[
b-a\le\min\left\{(\lambda-1)a,
                   \frac{\lambda-1}{\lambda}(1-a)\right\}.
\]
For $\lambda>1$, the first expression increases in $a$ and the second decreases. They are equal at $a=1/(\lambda+1)$. Maximizing their minimum over $0\le a\le\alpha$ gives~\eqref{MS:eq:privacy-envelope}. For $\lambda=1$, the laws coincide and the bound is zero.

For attainment, take outcomes $0,1$ with
\[
P_C=\left(\frac{\lambda}{\lambda+1},\frac1{\lambda+1}\right),
\qquad
P_D=\left(\frac1{\lambda+1},\frac{\lambda}{\lambda+1}\right).
\]
Both privacy inequalities hold. Choose $s(0)=0$ and
$s(1)=\min\{(\lambda+1)\alpha,1\}$. Its false-positive rate is at most $\alpha$ and its expectation difference is exactly $B(\lambda,\alpha)$. Finally, the incentive inequality is $g\le F(\MSE_{P_D}s-\MSE_{P_C}s)$, which yields the stated necessary condition.
\end{proof}

This is an elementary testing consequence of the stipulated likelihood-ratio constraint, integrated with the existing enforcement model. It does not prescribe weakening confidentiality. Protecting an unrelated private attribute may be compatible with revealing whether a particular authorized operation occurred. An institution can instead change its feasible actions, use a different justified record, reduce the covered gain, or decline the proposed sanction contract. The authority to make those changes is a separate question.

At $\lambda=3$, $\alpha=1/10$, the sharp envelope is $1/5$. With penalty $F=2$, gains above $2/5$ cannot be deterred by any record and sanction obeying these particular simultaneous constraints. The pair $P_C=(3/4,1/4)$ and $P_D=(1/4,3/4)$ with $s=(0,2/5)$ attains that envelope. For a gain $g=1/3$, the incentive margin is $2/5-1/3=1/15$. These are exact rational examples, not inferred observation laws.

\subsection{The full disclosure history is the protected object}\label{MS:sec:disclosure-history}

A private initial report followed by an unrestricted appeal disclosure need not remain private as a whole. Nor does protection of every separate marginal establish protection of their joint transcript. The privacy contract must identify the evaluator and include every record that evaluator can combine, including outputs of subsequent authorized procedures.

\begin{MScounterexample}[Two uninformative records with an informative joint history]\label{MS:ex:privacy-marginals}
Let $\theta\in\{0,1\}$ distinguish two protected cases, and let $R$ be a fair independent bit. Release $O_1=R$ and $O_2=R\oplus\theta$. Each marginal is a fair bit under both cases and satisfies~\eqref{MS:eq:pairwise-privacy} with $\lambda=1$. But $O_1\oplus O_2=\theta$. The joint laws have disjoint support and satisfy no finite multiplicative bound.
\end{MScounterexample}

There is a useful sufficient accounting rule when each stage is protected \emph{conditionally} on the previously disclosed history. It is the elementary sequential-composition argument, stated here for the exact pair rather than a complete database mechanism.

\begin{MSproposition}[Conditional disclosure accounting]\label{MS:prop:privacy-transcript}
Consider $m$ releases on finite alphabets. At each common prefix $h$, suppose their conditional next-release laws $P_C^t(\cdot\mid h)$ and $P_D^t(\cdot\mid h)$ satisfy both pointwise inequalities with factor $\lambda_t\ge1$, uniformly over prefixes that can occur. Then the joint transcript laws satisfy~\eqref{MS:eq:pairwise-privacy} with
\[
\lambda=\prod_{t=1}^{m}\lambda_t.
\]
The bound also holds after applying the same stochastic postprocessing to the transcript under both cases. In particular, any sanction computed only from that transcript obeys Theorem~\ref{MS:thm:privacy-envelope} with this product factor.
\end{MSproposition}
\begin{proof}
The conditional inequalities imply common support at each stage. Induction therefore gives common support for prefixes under the two cases. For a positive-probability full transcript, the ratio of its probabilities is the product of conditional ratios, each between $\lambda_t^{-1}$ and $\lambda_t$. Multiplying gives the joint pointwise inequalities. Zero-probability transcripts occur under both cases and satisfy them as well. Summing proves the event inequalities. For a common postprocessing kernel, multiply each pointwise inequality by its nonnegative output probability and sum over input transcripts. This preserves both inequalities.
\end{proof}

Writing $\lambda_t=e^{\varepsilon_t}$ gives the additive parameter $\sum_t\varepsilon_t$. The proposition requires a fixed finite number of releases and uniform conditional bounds. A data-selected stopping rule must itself be represented in the transcript or controlled by an appropriate conditional argument. Privileged side information not covered by the declared record is not protected by a certificate for that record alone. The preceding counterexample fails the conditional premise at the second release, even though both unconditional marginals are identical.

\subsection{An appeal is a feasible procedure, not a line in a constitution}\label{MS:sec:effective-appeal}

An effective appeal contract specifies standing, admissible evidence, the authority and competence of a reviewer, a time window, and an executable disposition. A correct disposition may uphold the original decision. Guaranteed access to review is not a guarantee of a favorable outcome. The contract must also distinguish a remedy from a retrospective description: restoring a ledger entry does not undo an irreversible external act.

Some review costs arise before a judgment can be reached. If an agent is entitled to challenge a decision but cannot access the stipulated evidence, its formal permission does not supply the missing information. If the reviewer has no available execution time, access to evidence alone does not make review possible. These are the information and resource obstructions already studied in Sections~\ref{MS:sec:decentral} and~\ref{MS:sec:joint-synthesis}, now applied to correction rather than ordinary service.

A bounded allocation model makes the capacity obligation precise. Let $J$ be a finite set of admitted appeals and let the integer epochs be $1,\ldots,T$. At epoch $t$, there are $c_t\ge0$ reserved review slots after accounting for all other protected work. An appeal $j$ has an evidence-ready release epoch $r_j$ and deadline $d_j$, with $1\le r_j\le d_j\le T$. It requires one slot containing the complete stipulated review and disposition operation. The evidence and qualified authorized reviewer are assumed available throughout that window. Larger or conditional workloads require a different job model.

\begin{MSproposition}[Capacity for a finite appeal contract]\label{MS:prop:appeal-capacity}
All admitted appeals can be assigned to distinct review slots in their windows if and only if, for every subset $U\subseteq J$,
\begin{equation}\label{MS:eq:appeal-hall}
 |U|\le\sum_{t\in\bigcup_{j\in U}\{r_j,\ldots,d_j\}}c_t.
\end{equation}
Equivalently, because the windows are intervals, it suffices to check every integer interval $[a,b]\subseteq[1,T]$:
\begin{equation}\label{MS:eq:appeal-interval}
 \bigl|\{j:r_j\ge a,\ d_j\le b\}\bigr|
       \le\sum_{t=a}^{b}c_t.
\end{equation}
\end{MSproposition}
\begin{proof}
Replace epoch $t$ by $c_t$ distinct slot vertices and join appeal $j$ to the vertices in its window. An assignment is a matching covering the appeals. The matching criterion used in Theorem~\ref{MS:thm:robust-staffing}, equivalently Hall's theorem~\cite{MS:hall1935}, gives exactly~\eqref{MS:eq:appeal-hall}. Necessity of~\eqref{MS:eq:appeal-interval} follows because every appeal counted on its left must use a slot in $[a,b]$.

For sufficiency of the interval tests, fix $U$. Its union of integer windows decomposes into disjoint maximal integer intervals $C_1,\ldots,C_k$. Each window is contained in one component because it is an interval. The number of appeals from $U$ in $C_i$ is no larger than the number of all appeals whose entire windows lie in $C_i$. By~\eqref{MS:eq:appeal-interval}, this is at most the total capacity in $C_i$. Summing over the disjoint components proves~\eqref{MS:eq:appeal-hall}, and hence the matching exists.
\end{proof}

The statement is an application of the same matching foundation as resilient staffing, not a separate scheduling discovery. Its role is to include an appeal promise in the common resource contract. Two appeals both due at epoch 1 cannot be served with capacities $(1,1)$ over epochs 1 and 2, even though total capacity equals total demand. The interval $[1,1]$ supplies the obstruction. Moving a deadline requires the relevant authorization and changes the original promise.

No fixed finite reserve can guarantee bounded service for an unlimited adversarial arrival stream without further admission, workload, or capacity assumptions. Restricting admissible appeals may itself affect the chosen access protection. The model therefore records both requests and admissions, the grounds for any exclusion, and the workload that the guarantee actually covers. It does not define every rejected request out of the social question.

Privacy and appeal must be checked together. Additional evidence can make an appeal decidable while changing the protected transcript distributions. The institution must either keep the new release within its declared disclosure contract or record an authorized change to that contract. A proof of fair scheduling does not establish the truth of the evidence, and a proof of the evidence does not supply the time needed to act on it.

\subsection{Exit, suspension, settlement, and dissolution}\label{MS:sec:authorized-endings}

Four different transitions should remain separate. An occupant may leave an office while the office continues. An institution may suspend external authority while preserving unresolved records. A settlement process may discharge or transfer specified obligations. Finally, a dissolved institution may retain an archive while having no active authority to issue new commitments. None of these states is defined simply by whether a model is currently executing.

Let $G_{\rm suspend}$ require that the declared external capabilities are inactive and that the pending-claim record remains available. Let $G_{\rm settle}$ additionally require that every covered claim has been fulfilled, discharged under the stipulated authority, or transferred to an authorized successor that has accepted it, and that the required final record has been committed. A retained unresolved claim can be compatible with suspension but not with this stronger settlement target. Neither target by itself settles legal status, moral status, or ownership outside the stipulated contract.

An immediate suspension can be necessary even when complete settlement is not currently possible. Conversely, deleting a debt record does not settle a debt. A promise that an institution can always end must identify which ending it guarantees, which resources remain available after a stop request, who authenticates that request, and what happens if continued service would itself violate the external boundary. Research on shutdown incentives studies a different part of this problem: whether a goal-directed agent has reason to permit an off switch~\cite{MS:offswitch}. Existence of a multi-party settlement procedure does not establish the incentives to use it.

The finite construction of Section~\ref{MS:sec:joint-synthesis} yields an exact characterization of \emph{continued operation with an available authorized ending}. Use its possibility space $\mathcal B$, successor operation, and protected set $S_{\rm safe}$. For each possibility set $B$, distinguish implemented running profiles $A_{\rm run}(B)$ from implemented ending profiles $A_{\rm end}(B)$. Both must be admissible in every state in $B$. Ending profiles cannot accept new obligations outside the declared settlement protocol. Let $G\subseteq S_{\rm safe}$ be the chosen ending target, with an authorized safe hold profile available after certified arrival. A stop command is assumed authenticated and publicly available before the next profile is selected. Any command-delivery delay must be included in the model rather than omitted from the response time.

Define
\begin{align}
 T_0&=\{B\in\mathcal B:B\subseteq G\},\label{MS:eq:stop-target}\\
 T_{k+1}&=T_k\cup\{B\in\mathcal B:B\subseteq S_{\rm safe},\
      \exists u\in A_{\rm end}(B),\ \operatorname{Succ}(B,u)\subseteq T_k\}.
      \label{MS:eq:stop-attractor}
\end{align}
Let $T_*$ be the stable set and $r_{\rm end}(B)$ the first entry index. Then define
\begin{equation}\label{MS:eq:stoppable-kernel}
 L_0=T_*,\qquad
 L_{k+1}=\{B\in T_*:\exists u\in A_{\rm run}(B),\
                         \operatorname{Succ}(B,u)\subseteq L_k\},
\end{equation}
with stable set $L_*$. These are finite constructions with the same potentially exponential possibility space as before.

\begin{MStheorem}[Continued operation with an available authorized ending]\label{MS:thm:stoppable-operation}
An initial possibility set $B$ lies in $L_*$ if and only if an admissible common-history policy can operate safely for an unbounded number of epochs when no stop command arrives, while ensuring certified arrival in $G$ on every path after a stop command received at any decision epoch. Before a command, a policy stationary in $B$ and preserving $L_*$ suffices. After a command at $B$, an ending policy reaches $G$ within $r_{\rm end}(B)$ transitions. When $L_*\ne\varnothing$, these response bounds have a finite maximum over $L_*$.
\end{MStheorem}
\begin{proof}
Induction in~\eqref{MS:eq:stop-attractor} shows that $T_k$ is exactly the set from which an ending-only policy can reach certified $G$ in at most $k$ transitions while remaining safe. The target has a safe hold action. In a finite, finitely branching model, eventual certified arrival on every path implies some finite bound for a fixed initial set, by the stopped-history-tree argument in Proposition~\ref{MS:prop:joint-progress}. Thus $T_*$ is exactly the set from which the stipulated ending is surely possible.

The sequence in~\eqref{MS:eq:stoppable-kernel} is decreasing and stabilizes. From $L_*$, choose a running profile whose every successor remains in $L_*$. If no command arrives, induction gives safe indefinite operation. If a command arrives, the current set is in $T_*$ and an ending profile of strictly decreasing first-entry rank completes the ending. Finiteness gives the maximum rank over $L_*$.

Conversely, consider a successful policy. At each public history reachable before a command, a command could be issued immediately. The current possibility set must therefore lie in $T_*$. For every finite number of further running epochs, the policy keeps every resulting possibility set in $T_*$. Induction in~\eqref{MS:eq:stoppable-kernel} puts the initial set in every $L_k$, and hence in $L_*$. The choices are based on the common history, so no hidden-state action has been introduced.
\end{proof}

The theorem reuses the established safety and reachability machinery. Its additional requirement is that running actions preserve the \emph{ability to end}, not only immediate safety. An institution can be safe now yet spend its last settlement resource on a new task, leaving no authorized ending from the resulting state. Such an action leaves $L_*$ even when it remains in the original continuation kernel. This is a realizability boundary, not an argument that an institution should refuse an externally authorized suspension because its preferred settlement plan is unavailable.

A five-state example separates the guarantees. From a running state $r$, a work action enters state $d$ with one unpaid obligation. From $d$, one more work action enters $\ell$, where the institution can idle safely but no longer has resources to settle. A return action from $d$ restores $r$. An ending action from either $r$ or $d$ reaches a closed state $z$. A breach state $b$ represents inadmissible failed closure attempts. The ordinary safe kernel includes $\ell$, but the ending attractor and $L_*$ exclude it. A policy can continue alternating between $r$ and $d$ and still answer an authorized stop at either state. It cannot offer that guarantee after entering $\ell$.

\subsection{A certificate has an applicability lifecycle}\label{MS:sec:certificate-lifecycle}

A proof of $A\Rightarrow G$ does not become false when its premise ceases to hold. Its application to the current institution may nevertheless become invalid. Changes to an observation channel, reviewer, capacity, membership configuration, trusted key, or action semantics can move the realization outside the certificate's domain. A checksum can identify which artifact changed. It cannot, by itself, decide whether the new world satisfies a semantic premise.

An assurance record can contain
\[
\mathcal K=(\mathsf{Claim},\mathsf{Assumptions},\mathsf{Evidence},
\mathsf{Domain},\mathsf{Guard},\mathsf{Fallback},\mathsf{Revision}).
\]
Its status should distinguish at least an applicable certificate, a certificate awaiting rechecking, a suspended claim, and a superseded version. Evidence for an assumption can be observed, externally attested under a trust model, proved about a component, or unavailable. Those statuses should not all become the Boolean value \emph{verified}. A failed guard means that continued reliance is not justified by that check. It need not mean that the guarantee is already false.

Runtime assurance supplies an established precedent for switching away from an advanced controller when its continued use cannot be certified~\cite{MS:blackboxsimplex}. Component-based work explicitly couples such switching with assume-guarantee reasoning~\cite{MS:componentsimplex}. Here the retained object includes authority, obligations, and evidence semantics as well as an operational safety predicate. A fallback must preserve those parts of the contract it claims to preserve, and its own implementation and resources remain assumptions.

For a finite set of modes $V$, let $C_v\subseteq Z$ be a certified state domain, with each $C_v$ contained in a common protected set $C$. A certified hold exists when a final state requires indefinite preservation. At $(v,z)$ with $z\in C_v$, a selector may either use the local policy $\kappa_v$ or execute an authorized handover to another mode. The selector and guard can read only their declared observations. A positive local guard must be sound for the \emph{whole action interval}: it entails that every actual successor of $\kappa_v(z)$ lies in the declared successor set $D_v(z)$. The local certificate proves $D_v(z)\subseteq C_v$. A handover from $v$ to $w$ must keep every intermediate state in $C$ and end in $C_w$. Its command, evidence transfer, resources, and delay belong to the transition model. Intermediate handover and checking states carry explicit mode tags when they count as primitive transitions, so the shared rank below is defined on those states as well.

\begin{MStheorem}[Guarded certificate use and authorized handover]\label{MS:thm:guarded-certificates}
Assume that the selector, at every reachable decision epoch, either uses a local policy with a sound positive guard or performs a handover satisfying the preceding conditions. Starting from $z_0\in C_{v_0}$, every decision state belongs to its current certified domain and every intermediate handover state remains in $C$.

If a target has an appropriate certified hold and one shared map $r:(v,z)\mapsto\mathbb N$ satisfies $r(v,z)=0$ only at the target, and every selected primitive transition before the first target visit strictly decreases $r$, then the target is reached within at most $r(v_0,z_0)$ primitive transitions. Handover, checking, and waiting steps count whenever they consume transitions in the promised time model.
\end{MStheorem}
\begin{proof}
At a local-policy step, guard soundness puts every actual successor in $D_v(z)$ and the local certificate puts it in $C_v\subseteq C$. At a handover, the stated trajectory and endpoint conditions preserve $C$ and establish membership in the new domain. Induction proves the safety assertion. Before a target visit, strict decrease of a nonnegative integer permits at most the stated number of primitive transitions. If the target had not then been visited, another selected transition would require a negative rank. Existence of the selected steps and the certified target hold complete the argument.
\end{proof}

This is a sufficient composition rule, not a universal monitor for its own premises. In particular, the selector cannot use the absence of detected evidence against an assumption as proof of that assumption. If no justified local action or handover is available, the theorem does not invent one. A prearranged emergency restriction of external authority may be available under a weaker suspension contract, but its obligations and physical effects must be stated separately.

\begin{MScounterexample}[A correct old certificate used after an unobserved change]\label{MS:ex:stale-certificate}
At time $0$, an action requires one execution slot and the certified capacity is one. Before the action is issued, an external event reduces the actual capacity to zero. A guard that merely reads the unchanged old capacity record remains positive. Executing the action violates the actual resource condition even though the original theorem and stored proof remain valid under their original premise. The guard was not sound for the action interval.
\end{MScounterexample}

\begin{MScounterexample}[Local progress ranks do not compose through arbitrary switching]\label{MS:ex:switching-rank}
Modes $a$ and $b$ each have a local path of ranks $2,1,0$ to completion. Permit a handover from rank 1 of either mode to rank 2 of the other. Alternating one local progress step and one handover yields an infinite safe path with no completion. Every local progress step decreases its own rank. There is no shared strictly decreasing rank for the selected handovers.
\end{MScounterexample}

\subsection{Mutual assurances do not create an initial fact}\label{MS:sec:circular-assurance}

Suppose institution $A$ claims that it is sound whenever $B$ is sound, and $B$ claims the converse. The two implications are compatible with both institutions being unsound. They are not an initial proof that either premise holds. This elementary point does not invalidate properly justified circular assume-guarantee rules. Such rules need their actual semantic and temporal hypotheses, not merely the printed appearance of mutual certification.

One simple sufficient discipline is an explicit dependency order. Primitive evidence establishes its stated base claims. Each later certificate cites only established premises and a sound implication. Induction through that finite order establishes the resulting claims. Temporal circularity can instead be sound when a demonstrated base state and noncircular one-step preservation rule support induction over time. The distinction is between an anchored proof and an unsupported cycle of assertions, not between all cyclic and all acyclic diagrams.

The same issue appears in federations. Importing a partner's certificate imports its declared trust dependencies and scope. If those dependencies return to an unestablished local assertion, copying the certificate again contributes no new evidence. An honest record preserves the unresolved assumption. The external-effects contract, the auditor's implementation contract, and the premise of a fallback policy can all lie at such a boundary.

\subsection{The accountable commons as one contract}\label{MS:sec:accountable-commons}

A final bounded variant combines the requirements without inheriting incompatible assumptions from earlier examples. Start with three internal service units due, one unit reserved for an outside claimant, one admitted appeal whose complete review and disposition require one unit, and a two-unit audit. Final recording and revocation of the named external capabilities require one more unit. The epoch provides $c$ sequentially usable units. No new claims are admitted after the authenticated closing command. All seven service and audit units must be completed before the final revocation step. The administrative operation is included in the budget rather than treated as free.

\begin{center}
\begin{tabular}{@{}lr@{}}
\toprule
Protected operation & Execution units\\
\midrule
Three internal service units & 3\\
External reserved service & 1\\
Complete admitted appeal & 1\\
Stipulated audit procedure & 2\\
Final record and capability revocation & 1\\
\midrule
Total for this settlement contract & 8\\
\bottomrule
\end{tabular}
\end{center}

For this variant only, both covered reporting deviations have the same observation law $P_D=(1/4,3/4)$ and gain $1/3$, while compliance has law $P_C=(3/4,1/4)$. The two-unit audit is stipulated to produce those records. It is \emph{not} the earlier three-outcome audit, whose supports do not satisfy this finite pairwise privacy bound. Set $\lambda=3$, $\alpha=1/10$, $F=2$, and $s=(0,2/5)$. The common monitor has false-positive rate $1/10$ and margin $1/15$ against each listed deviation. This remains an expected-utility incentive guarantee, not a reason to erase physical deviation paths from a safety model.

The audit provider receives a separate collectible fee $w$ for verified execution, with $w$ larger than its stipulated cost $c_a$. A genuine-execution receipt is checked by a specified trusted mechanism. The fee is not a reward for accusing someone. It makes execution preferable to skipping in this narrow two-action provider model. Monetary units are tracked separately from execution units. The payment budget, receipt checker, and record-generation law remain explicit assumptions. The perfect detector of Proposition~\ref{MS:prop:strategic-inspector} is not silently substituted for this noisy audit.

With $c=8$, the obligations and the final revocation fit. A countdown in remaining required work supplies a settlement rank, provided the prescribed information and authority are available and no modeled event creates further work. With $c=7$, the service, appeal, and audit can finish, but the complete settlement target cannot be certified because its final administrative operation is unfunded. A previously reserved suspension operation can deactivate external capability while leaving the remaining claim visible. That reaches a different target and is not recorded as full settlement.

The reference implementation uses five profiles with costs $2,3,1,1,1$: audit, batched internal service, external service, appeal, and final record with revocation. The final condition check is part of the last atomic administrative profile. A phase rank $5,4,3,2,1,0$ certifies five-transition completion from the admitted initial state. Separating a check into an additional transition would require changing that rank, even if its execution-unit cost were zero. No such uncounted step is assumed here.

If a premise changes, the continuing assurance record must say what is still justified. Losing one of the eight units can invalidate full settlement without retroactively making the privacy calculation false. Revoking the auditor's credential can invalidate its profile despite unchanged capacity. Discovering an outside reservation can refine the governance state. Releasing new appeal evidence can change the privacy transcript. Each update has a typed consequence. No overall label such as \emph{healthy} or \emph{aligned} replaces these separate checks.

\subsection{What accountable persistence means in this theory}

The addition is not a derivation of a unique moral constitution. It is a stronger realizability demand once protected claims have been selected. The same institution must identify its external boundary, fund credible evidence production, respect the specified information protections, make authorized review feasible, preserve the ability to reach its chosen ending state, and stop relying on certificates whose applicability is no longer supported.

Feasibility does not imply that an institution should be deployed or given authority over people. Equally, the possibility of harmful organization does not make mathematical analysis of it incoherent. The purpose of the formal distinctions is to make the consequences of an explicit arrangement inspectable. A system may be persistent without being legitimate, legitimate under one criterion without being implementable, or implementable without possessing a feasible ending under its current obligations. Its record should show which situation it is in.

A computational institution's ability to preserve itself is therefore different from its ability to remain answerable for its actions. In this framework, correction, bounded authority, and authorized endings belong to the same operational theory as continuation. They are not a rhetorical qualification appended after the mathematics is complete.



\FloatBarrier
\section{Conditions for institutional novelty}\label{MS:sec:novelty}

The formation of conventions does not require irreducible randomness. If initially indistinguishable deterministic agents have the same update rule and a symmetry-preserving interaction schedule, their joint state remains permutation-invariant by Proposition~\ref{MS:prop:symmetry}. Exchangeability of a distribution is weaker than this pointwise symmetry. A distinct role then requires a symmetry breaker such as heterogeneous information, an address, path-dependent history, or stochastic variation. Once such asymmetries exist, a deterministic system can establish and maintain conventions through ordinary state transitions.

The fact that deterministic rules can generate unexpected macro-organization does not imply that a macro-process has an autonomous Markov law. The three-state counterexample in Section~\ref{MS:sec:closure} shows how coarse-graining can generate history dependence. A putative ``phase change'' in institutional organization therefore needs a specified order parameter, state family, and dynamical limit. Rapid change under increased communication, repeated interaction, or memory capacity may mark a useful empirical onset, but does not by itself identify a mathematical bifurcation or phase transition.

\subsection{Three levels of novelty}

It is helpful to distinguish \emph{combinatorial novelty}: a configuration not previously present in the local records; \emph{adaptive novelty}: a newly selected configuration that changes performance under a declared environment; and \emph{institutional novelty}: a new rule, authority arrangement, or revision procedure that has operative consequences and is maintained across time. The first can occur by arbitrary recombination. The second needs a task criterion. The third requires an institutional transition and survival condition. These distinctions limit rhetorical inflation without denying the possibility of genuinely unexpected organizational outcomes.

\subsection{Thresholds are model dependent}

A minimum frequency of contact may make convention formation likely under one interaction graph and not another. Longer time horizons can amortize setup costs but also lock in bad rules. More memory can preserve beneficial corrections or propagate obsolete assumptions. For the abstract stationary cost model in Proposition~\ref{MS:prop:amort}, the inequality $B<g/(1-\delta)$ is a threshold for \emph{economic usefulness of that particular institution}, not a general law for the appearance of societies. Evolution of institutions requires separate assumptions about rule variation, incentives, scheduling, authority, and selection.

\section{Foundational problems for a theory of machine societies}\label{MS:sec:problems}

The framework points toward a set of research questions that are mathematical rather than metaphorical.

\paragraph{Minimal sufficient institutional state.}
For a family of future tasks and amendments, what is the smallest summary that preserves resource rights, causal decision structure, and valid dispute procedures within specified error? Theorem~\ref{MS:thm:govquot} and Proposition~\ref{MS:prop:govrefine} resolve the finite deterministic fixed-contract case. Approximate, stochastic, and changing-contract versions may require more state than is necessary for immediate task accuracy. It may depend on the horizon and on the admissible forms of revision.

\paragraph{Formation under mixed incentives.}
Which adaptive rule-selection processes produce persistent institutions in populations with heterogeneous objectives and imperfect information? Potential games provide a tractable sufficient structure, while Propositions~\ref{MS:prop:barrier} and~\ref{MS:prop:seed} expose a transition barrier even in a simple coordination setting. They do not settle general formation. In nonpotential settings, characterize cycles, costly exit, and endogenous modification of the action space itself.

\paragraph{Causal abstraction under amendment.}
When a macro-institution can modify its own feasible interventions, can a fixed lower-dimensional summary remain sufficient for every admitted amendment? How must the abstraction be enlarged to account for the changed permissions and revision operator? This is a dynamic coarse-graining problem with the operation algebra itself changing over time.

\paragraph{Local--global institutional repair.}
Given locally specified constraints and interfaces with an obstruction to global realization, find the least-cost authorized modifications that eliminate the obstruction, or certify that no protected-obligation-preserving repair is possible. Proposition~\ref{MS:prop:cyclerepair} resolves protected repair on one binary cycle. Overlapping constraints, richer logics, and coupled resource obligations require further analysis.

\paragraph{Reliable inheritance under replacement.}
Beyond the sufficient capability-preservation conditions of Proposition~\ref{MS:prop:culturegrowth}, what minimal archive, provenance data, and operation-preserving translations allow a methodology to survive an arbitrary permitted sequence of model replacements? The bound in Proposition~\ref{MS:prop:transport} highlights accumulated drift; a full treatment must also account for lossy compression, stochastic decision rules, changing rights, and intentional correction of false inherited knowledge.

\paragraph{Plurality without centralized optimization.}
What can be certified about fairness, access, and collective safety when there is no society-wide utility function? There may be several incompatible yet reasonable normative criteria. A mathematically clean model must expose which authority selects the protected criteria rather than disguise that selection as an optimization theorem.

\paragraph{Threat propagation and institutional recovery.}
Under what conditions can corruption of rules or provenance be isolated without suppressing valid novelty? Linear disturbance models provide local comparison bounds. The hard case is adversarial adaptation combined with changing channels, changing permissions, and incomplete observation of the true institutional state.

\paragraph{The meaning of organizational continuity.}
Which aspects of identity should be described by behavior, which by causal lineage, and which by authorization or social recognition? These relations can disagree. A theory of persistence should allow this disagreement rather than assume a single numerical identity score.

\paragraph{Coherence under changing boundaries and languages.}
When a membership update changes both the electorate and the interpreter, what jointly checkable certificate preserves prior obligations, recognized decision rights, and task-relevant knowledge? The exact transport results assume matched agent identities and semantic atoms. Extending them to admission and changing vocabularies requires explicit rules for new information and new authority, not merely a larger state vector.

\paragraph{Confluence without erasing history.}
Which quotients of amendment histories permit useful normal forms while preserving the ordering-sensitive audit questions an institution promises to answer? Literal append-only records distinguish even harmless permutations; policy-only states can erase consequential differences. A satisfactory equivalence must be selected by the operational and accountability contract.

\paragraph{Open-ended intellectual organization.}
What institutional conditions support useful expansion of problem schemas and methods when the candidate universes and production rules are not fixed in advance? The finite closure theorem describes derivability within a given system. It does not establish indefinite creativity, guarantee fair allocation to an enabled inquiry, or align cultural survival with truth.

\paragraph{Valid histories without a universal event mechanism.}
Which institution-specific operations admit coordination-free invariant preservation, and which require exclusive capabilities or a commitment protocol? Proposition~\ref{MS:prop:merge-convergence} supplies a sufficient merge-closure condition, while Proposition~\ref{MS:prop:exclusive-service} gives a bounded service obstruction. A full account must also handle revocation, changing authority configurations, irreversible effects, and what participants can know about finality.

\paragraph{Enforcement under model uncertainty and strategic monitors.}
Theorem~\ref{MS:thm:joint-enforcement} resolves the finite common-monitor problem for several specified deviations, including a protected false-positive rate. The joint continuation construction in Section~\ref{MS:sec:joint-synthesis} incorporates such witnesses into a fixed finite protocol catalogue. Incentive-constrained auditors, unidentified observation laws, adaptive deviations, and unrestricted decentralized strategies remain outside those results. Proposition~\ref{MS:prop:robust-deterrence} propagates declared observation errors into a sanction margin; it does not identify the observation model or solve the incentives of those who certify it.

\paragraph{Joint synthesis beyond a fixed catalogue.}
Theorem~\ref{MS:thm:joint-viability} is complete for its finite common-observation model, and Theorem~\ref{MS:thm:synthesis-transport} preserves that exact problem under a suitable quotient. What can be retained when agents create new local protocols, when observations are only approximately matched, or when probabilistic service replaces sure service? Enlarging a catalogue changes a synthesis domain. It does not by itself prove the new profiles are implementable, authorized, or affordable.

\paragraph{Institutional timescales selected by the institution.}
The finite-block averaging theorem assumes a predetermined review length and endpoint-based revision. Which approximations remain valid when institutions choose their own review times, experience several long-lived fast regimes, or use path-dependent evidence? The stopping-bias counterexample prevents transferring a deterministic-time mixing bound to such settings without a new argument.

\paragraph{Adaptive specialization with preserved evidence and authority.}
The robust staffing criterion characterizes rematchability on a fixed authorized capability graph. A dynamic theory must also explain when reorganization can reach such a graph, execute handovers in time, preserve independent observations, and avoid a new concentration of decision rights. Static matching, temporal communication, and endogenous selection are distinct obligations.

\paragraph{Accountability with partially observed assumptions.}
Section~\ref{MS:sec:accountability} closes specified finite inspection, disclosure, appeal, and stopping models. It does not identify a complete external-effects boundary, establish every observation law, or solve strategic collusion among the auditor and its verifier. The guarded certificate theorem requires sound action-interval assumptions. Characterizing useful guarantees when those assumptions can only be partially checked remains a distinct problem.

\paragraph{Preserving access while admitting a changing workload.}
The finite appeal criterion fixes admitted jobs and reserved capacity. A broader contract must specify which requests must be admitted, how new work changes response guarantees, and which resource changes can occur during settlement. Indefinite access, finite capacity, and bounded response cannot be asserted together without a workload model.

\section{Limits of the framework}\label{MS:sec:limits}

Most of the analysis is built from finite state spaces, specified observation maps, and admitted control families. Section~\ref{MS:sec:computability} explicitly changes the domain to unrestricted programs with unbounded memory; its impossibility result must not be read as a claim about the earlier finite models. It does not automatically extend to open-ended deployed systems with unbounded memory, unknown environment dynamics, arbitrary strategic adversaries, or unresolved questions of computability and measure. Some finite statements generalize under regularity conditions; others require new machinery. Even an exact mathematical abstraction only preserves those behaviors and interventions specified by its contract.

Institutional formation is not a consequence of the existence of an institutional representation. An empty constitution can be formally represented and never adopted. A stable rule can arise through designer imposition, agent agreement, or accidental coordination, and the same observable behavior may be compatible with several mechanisms. Theorems about cost, potential functions, viability, and compatibility establish conditional structure, not a law that sufficiently numerous agents must create societies.

The framework also cannot derive a unique moral constitution from computational efficiency. Protected rights, acceptable tradeoffs, and the authority to enact constraints require human or otherwise explicitly specified normative commitments. The mathematical distinction between permitted and forbidden transitions is not a decision about which permissions ought to exist. The running resource-allocation and parity models are useful because their assumptions are inspectable. They are not claimed to cover the full complexity of economic, social, or legal institutions.

Finally, a theoretical description of cultural inheritance or organizational continuity is neither a test for subjective experience nor a basis for legal status. The purpose of these exclusions is to leave a precise problem: how persistent rules, relationships, and historical obligations can be represented and transformed within ordinary computation.

\subsection{What the combined certificates do and do not compose}
The state, interpreter, epistemic relations, and authority registry can be carried by one realized system without becoming one undifferentiated object. Language transports compose through exact admission and event preservation. Epistemic transports compose through matched predicates and the back-and-forth conditions. Typed residuals compose only with their declared metric compatibility. Frame synchronization is controlled by ordered loop products, while state compatibility depends on fixed points. Amendment normal forms require a separate termination and joinability argument.

Section~\ref{MS:sec:joint-synthesis} strengthens this coexistence requirement into a common-policy question. Under its exact information, action, label, and successor conditions, a governance abstraction preserves both the safe continuation kernel and finite progress ranks. Approximate metrics alone do not imply this sure strategy equivalence. A supplied certificate must identify the common realization and action catalogue to which all its conditions apply.

The joint lesson is that institutional coherence is a bundle of related obligations, not a scalar score. No single condition above implies all the others. Their explicit coexistence in one contract is what allows a claim about preservation of a society's operations to be evaluated without turning its language, knowledge, history, and authority into interchangeable metaphors.

\section{Conclusion}\label{MS:sec:conclusion}
A machine society is not defined by the number of its models, the fluency of its conversations, or one impressive collective output. Its theoretical possibility is persistent plurality organized by operative institutions: decision centers maintain rules, records, and procedures that shape their interaction and can survive changes of participants. The same population may change its membership, the meaning of its communicative acts, the questions it recognizes, and the division of work through which it acts.

These developments remain ordinary computation, but require distinct explanations. Realization and abstraction determine which higher-level states faithfully describe a process. Institutional semantics determines the actions, permissions, and obligations affected. Epistemic structure determines what participants know. Governance determines who may change the rules. Provenance and founding assumptions locate the boundaries of authority that internal consistency alone cannot settle.

The implementation of a continuing institution adds three coupled questions: what events are committed, why participants comply, and which slow dynamics remain meaningful amid fast activity. Convergent records need not be valid allocations. Observation-based sanctions have a finite discrimination capacity, and compression cannot increase it. A uniform mixing bound can justify a review-epoch institutional approximation, but a data-selected review time can invalidate that approximation. The robust sanction bound shows how a declared modeling error changes an incentive guarantee rather than disappearing inside a convenient macro-description.

Persistence likewise does not imply learning or resilience. Rational actors can enter a wrong information cascade when only decisions are disclosed. Independent evidence can support continuing accuracy gains under a different information contract. A collection of specialists can survive turnover when its authorized capability graph has sufficient collective matching slack, yet a backup for each role need not suffice. A connected aggregate graph can lack the time-respecting route needed to carry evidence before a deadline.

The earlier results fit the same pattern. A beneficial institution can face a formation barrier; a safe full-information action can be unavailable to distributed decision makers; locally invertible interfaces can fail to synchronize; and an invariant local state need not establish a full common reference frame. Amendment normal forms, cultural closure, and unrestricted verification each have their own assumptions and limits. No one certificate supplies all the others.

The final synthesis makes the quantifier issue explicit. A different valid policy for every requirement is not one valid policy for their conjunction. Multiple deviations can defeat a common monitor even when each has a successful individual test. An audit and three due services can each be affordable while their joint cost exceeds the available slots. A safe amendment can remove every safe continuation. The finite continuation construction identifies such failures, while the exact abstraction theorem states when a smaller institutional representation preserves the existence of a solution. The resulting witnesses certify declared operations, not the origin, legitimacy, or spontaneous discovery of the institution.

Accountability adds a final distinction. An institution can remain operational while losing the justified use of a certificate, the ability to fund an appeal, or the capacity to honor an authorized ending. The sharp disclosure bound, strategic-inspection model, finite appeal criterion, ending-preserving kernel, and guarded handover rule specify parts of that problem. Their premises remain visible. A model of internal members does not itself include every affected outsider, and a settled record does not manufacture consent or reverse an irreversible external effect.

The positive thesis is that persistent, revisable relations among computational decision centers are coherent mathematical objects. Their formation, preservation, semantic change, disagreement, incentives, and failures can be studied within one realization framework. Their ability to continue is different from their ability to remain answerable for their actions. A complete declared contract must preserve both the operations it promises and the authority, evidence, resources, and ending conditions on which those promises depend. This is neither an argument that machines must form societies nor a criterion for consciousness or moral status. It is a theory of what computational social organization would have to realize, preserve, and acknowledge as unresolved.



\backmatter
\chapter{Glossary}
\begin{multicols}{2}
\small
\begin{description}[style=nextline,itemsep=3pt,parsep=0pt]
\item[Realization] A compatible family of typed views meeting declared restriction constraints, or a concrete mathematical model satisfying a specified system contract.
\item[Representation] A map or retained record together with the questions and operations it is required to preserve.
\item[Evidence] Information accessible through declared observation or intervention protocols, not an assumed hidden truth label.
\item[Predictive closure] Existence of an update on retained state that commutes with the protected dynamics.
\item[Packet] A typed resolved component or projection. It is not automatically a probability distribution.
\item[Residual] A discrepancy relative to a specified representation or operation contract. Its interpretation depends on the chosen norm or invariant.
\item[Correction signature] A retained equivalence class preserving the outcomes of all allowed future evidence filters on the target.
\item[Instrument] A joint next-state and report probability law, including the evidence scale needed for normalization.
\item[Legal-operation quotient] An abstraction preserving labels, action availability and successor semantics on all enabled actions.
\item[Certificate] Typed evidence for a named mathematical or operational claim relative to its explicit assumptions.
\item[Proof field] The larger network of typed claims and justifications within which a certified proof is located.
\item[Ending attractor] States from which a sure ending is possible under all admitted successor choices.
\item[Running kernel] The greatest region invariant under some running strategy that preserves immediate membership in the ending attractor.
\item[Attained domain] The image of an actual record map. A decoder on attained values does not assign meaning to impossible records.
\item[Conservative envelope] A model that includes every possible actual successor and admits only actions available throughout the represented fine states.
\item[Migration] A map from a retained old record to a specified new record, with separate conditions for continued operation and certificate reuse.
\item[Seeded closure] The least set containing the admitted leaves and observations and closed under every admitted rule.
\item[Strict replay] A trace check requiring each accepted rule step to add a new conclusion.
\item[Causal replay] A trace check that allows already known conclusions, while still requiring every premise at the time of use.
\item[Assurance boundary] The precise source, computation, proof environment, or external premise supported by a verification record.
\end{description}
\normalsize
\end{multicols}


\chapter{Result and source map}\label{Core:app:map}
The following map gives the role of each selected written result. Exact source anchors, excerpts, hashes, and detailed hypothesis review are in the accompanying machine-readable ledger and review files. Source relationships do not imply that every source page is admitted.
\small
\begin{longtable}{p{0.13\linewidth}p{0.39\linewidth}p{0.35\linewidth}}
\toprule Result & Preserved object & Source relationship\\\midrule\endhead
\ref{Core:r:factor} & A target determined by an available record & GA, shared fibre criterion.\\
\ref{Core:r:behavior} & Every declared finite continuation & GA and MU, consolidated.\\
\ref{Core:r:packet} & Exact split encoding and named composition loss & SR, split and hidden-return identities.\\
\ref{Core:r:transport} & Source operations and matched readouts & MU, with related uses in OSL and Neuralese use the same conditional principle.\\
\ref{Core:r:common} & One output method across candidate dynamics & GA, common observable quotient.\\
\ref{Core:r:positive-packets} & Probability interpretation of a complete packet family & LSM, positive packet criterion.\\
\ref{Core:r:memory} & Every admitted future set correction & GA, correction-signature theorem.\\
\ref{Core:r:occupancy} & Exact partition-evidence memory cardinality & GA, occupancy corollary.\\
\ref{Core:r:bayes} & All posterior updates in the actual reporting chain & MU, finite channel factorization.\\
\ref{Core:r:joint-refinement} & Coarsest partition closed for prediction and update & LSM and PRIQ, finite closure.\\
\ref{Core:r:instrument} & Joint next-state/report law, posteriors and feedback & PRIQ and AVOC, selected composition.\\
\ref{Core:r:history-bound} & Whole-history error and conditioning cost & LSM, law and normalization bounds.\\
\ref{Core:r:risk} & One bounded-loss decision calculation & GA and MU, specialized common-risk bound.\\
\ref{Core:r:recipient} & One artifact under fixed recipient procedures & MU, intersection and covering criterion.\\
\ref{Core:r:inquiry} & One executable tree under all declared budgets & MU and SII, finite Pareto synthesis.\\
\ref{Core:r:lowerbound} & A complete optimality claim for a checked tree & SII, binary singleton-repair scope.\\
\ref{Core:r:belief-planning} & One optimal finite-horizon expected-cost inquiry & AVOC, projected values and finite plans.\\
\ref{Core:r:policy-error} & Uniform policy evaluation before selection & LSM/AVOC, common-law specialization.\\
\ref{Core:r:cegis} & Exact finite candidate rejection and acceptance & SII, finite counterexample refinement.\\
\ref{Core:r:ordered} & The same proposed ASTs and teacher transcript & SII, sound equations and ordered quotient.\\
\ref{Core:r:prefix} & The same completed ordered diagnostic search & SII, capacity and streaming arguments.\\
\ref{Core:r:dependencies} & Applicability of dependency-bound witnesses & SII, complete declared dependency graph.\\
\ref{Core:r:prices} & Gains under one complete additive ledger & SII, price-relative dominance.\\
\ref{Core:r:lineage} & Conditional false-acceptance control & SII, adaptive lineage error budget.\\
\ref{Core:r:reconstruction} & A target from jointly accessible linear records & SR and SRWT, target-relative specialization.\\
\ref{Core:r:governance} & Legal action availability and label histories & MS, exact governance quotient.\\
\ref{Core:r:ending} & Continued operation with an ending policy & MS, fully observed finite specialization.\\\bottomrule
\end{longtable}
\normalsize


\section*{Selected core correspondence}
The following status applies to the twenty-seven original selected source statements, accepted on 25 October 2026. References identify their locations in the reconstructed opening chapters. An expanded formulation requires comparison with its accepted source endpoint.
\small
\begin{longtable}{@{}p{.11\linewidth}p{.55\linewidth}p{.25\linewidth}@{}}\toprule Result & Statement & Accepted original source correspondence\\\midrule\endhead
\ref{Core:r:factor} & A record determines a target exactly on its fibres & Accepted, scoped\\
\ref{Core:r:behavior} & Evidence-initialized behavioral state & Accepted, scoped\\
\ref{Core:r:packet} & Exact accounting of the discarded excursion & Accepted, scoped\\
\ref{Core:r:transport} & Forward transport of a reusable method & Accepted, scoped\\
\ref{Core:r:common} & A common observable quotient & Accepted, original source\\
\ref{Core:r:positive-packets} & When complete packets are genuine probability components & Accepted, original source\\
\ref{Core:r:memory} & Minimal exact memory for a correction contract & Accepted, scoped\\
\ref{Core:r:occupancy} & The size of partition correction memory & Accepted, original source\\
\ref{Core:r:bayes} & When a report preserves every Bayesian correction & Accepted, scoped\\
\ref{Core:r:joint-refinement} & Least refinement closed under prediction and correction & Accepted, scoped\\
\ref{Core:r:instrument} & One quotient for reports, posteriors, and feedback & Accepted, scoped\\
\ref{Core:r:history-bound} & A whole-history bound with an explicit conditioning cost & Accepted, scoped\\
\ref{Core:r:risk} & A decision bound for one actual approximation chain & Accepted, scoped\\
\ref{Core:r:recipient} & A common artifact for fixed recipients & Accepted, scoped\\
\ref{Core:r:inquiry} & Exact joint feasibility for explicit controllers & Accepted, scoped\\
\ref{Core:r:lowerbound} & A candidate attaining both lower bounds is optimal & Accepted, scoped\\
\ref{Core:r:belief-planning} & The same compressed belief supports the same inquiry & Accepted, scoped\\
\ref{Core:r:policy-error} & Uniform model error controls a changed policy & Accepted, scoped\\
\ref{Core:r:cegis} & What finite counterexamples establish & Accepted, scoped\\
\ref{Core:r:ordered} & Ordered semantic compression preserves the teacher trace & Accepted, scoped\\
\ref{Core:r:prefix} & A sound prefix exclusion & Accepted, scoped\\
\ref{Core:r:dependencies} & Selective reuse with complete dependencies & Accepted, scoped\\
\ref{Core:r:prices} & Which gains survive every admitted price? & Accepted, scoped\\
\ref{Core:r:lineage} & A conditional error budget for a lineage & Accepted, scoped\\
\ref{Core:r:reconstruction} & Target-relative joint reconstruction & Accepted, scoped\\
\ref{Core:r:governance} & A quotient that preserves legal operations & Accepted, scoped\\
\ref{Core:r:ending} & Safe continuation that retains an ending strategy & Accepted, original source\\
\bottomrule
\end{longtable}\normalsize

\chapter{KS: additional mathematical material}
\setcounter{section}{0}\renewcommand{\thesection}{KS.\Alph{section}}\renewcommand{\theHsection}{KS.appendix.\arabic{section}}


% BEGIN retained input sources/monographs/komorebi-space/source/chapters/09_workshop.tex

\chap{A reader's workshop}\label{KS:app:workshop}
These short problems isolate choices that recur in the main development. Each has a complete solution. They are reading exercises and direct applications of the cited theory, not additional entries in the formal theorem inventory.

\subsection*{1. A stable limit does not make every seed persistent}
At every level take \(X_n=\{0,1\}\), with each bonding map sending both elements to zero. Find the inverse limit and the image of evaluation at every level. Does nonemptiness of every finite fiber make the seed one extendable?

\textbf{Solution.} For a coherent family, \(x_n=\rho_n(x_{n+1})=0\) for every \(n\). The all-zero family is coherent and is therefore the only point. Every evaluation image is \(\{0\}\). Equivalently, the image of restriction from any strictly later level is already \(\{0\}\), so the persistent image is \(\{0\}\). The seed one never extends. This is compatible with Proposition~\ref{KS:KS58}, which proves nonemptiness and characterizes the persistent image, not surjectivity of every evaluation.

\subsection*{2. A residual chart can change without changing the defect}
Let \(d:\KSR^2\to\KSR\) be \(d(u,r)=r\). For \(a\in\KSR\), let \(A_a(r)=(ar,r)\). Compute the exact repair \(P_a\), compare it with \(P_b\), and decide when they commute.

\textbf{Solution.} Since \(dA_a=I\),
\[
 P_a(u,r)=(u-ar,0),\qquad
 P_b(u,r)-P_a(u,r)=((a-b)r,0).
\]
Both fix the horizontal exact axis. Thus \(P_aP_b=P_b\) and \(P_bP_a=P_a\). They commute exactly when they are equal, which occurs exactly when \(a=b\), as can be checked at input \((0,1)\). Every repair has zero residual. The different answers reflect the different right inverses, as in Theorem~\ref{KS:KS65}.

\subsection*{3. Pairwise invertibility does not supply a triple law}
Let \(t\) be a nonidentity transposition in \(S_3\). Suppose three proposed overlap comparisons require frames \(b_0,b_1,b_2\) satisfying
\[
 b_1=tb_0,\qquad b_2=tb_1,\qquad b_2=tb_0.
\]
Every comparison is individually invertible. Can the three frames exist?

\textbf{Solution.} The first two equations give \(b_2=t^2b_0=b_0\). The third then gives \(b_0=tb_0\), and right cancellation gives \(t=1\), contrary to the assumption. Invertibility of every proposed overlap map does not prove the composition law needed for descent. Theorem~\ref{KS:KS34} applies to data satisfying that law, not to arbitrary lists of pairwise isomorphisms.

\subsection*{4. Zero noisy output is not a certificate of reference exactness}
Take \(D(a,b)=a\), \(\widetilde D(a,b)=a+\tfrac1{10}b\), and input \(x=(-10,100)\). The measured operator returns zero exactly. Find the distance to the reference kernel and evaluate the reference uncertainty bound.

\textbf{Solution.} The reference kernel is the vertical axis, so the distance is ten. Its residual gain is one, and the operator error norm is \(1/10\). With \(y=0\) and \(\eta=0\), Proposition~\ref{KS:KS69} gives
\[
 \KSdist(x,\ker D)\le\frac1{10}\sqrt{10^2+100^2}=\sqrt{101}.
\]
The true distance ten is below this bound. A formula omitting the operator-error term would instead claim a zero distance, which is false. The input is exact for the measured model and not for the reference one.

\subsection*{5. A prefix computation can hide a wraparound defect}
On a four-cycle, take proposed differences \((2,-1,3,-3)\) and anchor zero. Integrate the first three edges. Which equation fails, and what changes when only the final edge is allowed to move?

\textbf{Solution.} The first three edges give values \((0,2,1,4)\). Their required last difference is \(x_0-x_3=-4\), not \(-3\). The total prescribed sum is one, so the original loop has no realization. Keeping the first three edges forces the unique correction \((2,-1,3,-4)\), whose sum is zero. This is a protected repair of the edge data. It is not a realization of the original data without modification.

\subsection*{6. Residual decay can hide unprotected drift}
In maximum-norm $\KSR^3$, use $p(s,u,r)=s$ and $d(s,u,r)=r$. Compare the paths $(1,0,(n+1)^{-1})$ and $(1,H_n,(n+1)^{-1})$, where $H_n=\sum_{j=1}^{n}1/j$. Which hypotheses distinguish them?

\textbf{Solution.} Both preserve $p=1$ and have the same residual ratios $(n+1)/(n+2)$. The first moves by $1/((n+1)(n+2))$ and converges to $(1,0,0)$. The second moves by $1/(n+1)$ and has an unbounded second coordinate. In the variable-rate theorem, the first has $C_n=1/(n+2)$ and a summable weighted budget. The second has $C_n=1$, making that budget harmonic and divergent. The difference is movement, not residual decay or semantic preservation.

\subsection*{7. How a finite observation changes the apparent rate}
Let $\eta_n=1/(n+1)$. Find its profile after sampling at $\phi(k)=2^k-1$. Does the new profile prove exponential convergence in the original index?

\textbf{Solution.} The sampled value is $\eta_{\phi(k)}=2^{-k}$. Recovering the original index gives $k=\log_2(n+1)$ on the sampled stages, so the value is still $1/(n+1)$. The schedule preserves vanishing but changes the meaning of the horizontal axis. No faster original process has been proved. Proposition~\ref{KS:KS75} is a convergence equivalence under a monotone cofinal schedule, not a rate-invariance theorem.

\subsection*{8. What does a zero floor fail to certify?}
For the diagonal model of Section~\ref{KS:sec:escape}, can an exact protected answer exist in $\ell^2\oplus\KSR$? Does a coherent family of exact finite prefixes exist? Explain why the answers differ without contradicting finite-fiber existence.

\textbf{Solution.} Exactness with protected value one forces every sequence coordinate to equal minus one. This is not square summable, so there is no protected Hilbert answer. Each finite prefix nevertheless has the unique solution $(-1,\ldots,-1;1)$, and these solutions are coherent under truncation. Their inverse limit contains the all-minus-one sequence in the full product. The finite-fiber theorem supplies that coherent point. It does not prove that the point is in the image of the square-summable semantic space. The correction image is dense but does not contain the raw residual, so its floor is zero and its algebraic class is nonzero.

\subsection*{9. A phase witness without enumerating every gauge}
For $c(x,y)=x_2y_1$ on $\mathbf F_2^2$, show that no scalar coboundary equals $c$. Then compute the square of the lifted label $(0,e_1+e_2)$ in the eight-element extension.

\textbf{Solution.} Scalar coboundaries $a(x)+a(y)-a(x+y)$ are symmetric, but $c(e_1,e_2)=0$ and $c(e_2,e_1)=1$. Thus no scalar coboundary equals $c$. For $x=e_1+e_2$, $c(x,x)=1$, so $(0,x)^2=(1,0)$ and its fourth power is the identity. The corresponding matrix is $XZ$, whose square is $-I$. The additional phase bit records precisely the sign lost by projecting back to the four-label group.

\subsection*{10. The finite energy identity}
In $\KSR^2$, first prescribe $a_1+a_2=1$, then also prescribe $a_1=3/4$. Find the two minimum norm corrections and check the exact nested-affine identity.

\textbf{Solution.} The first correction is $(1/2,1/2)$, with energy $1/2$. The second is $(3/4,1/4)$, with energy $5/8$. Their difference has squared norm $1/8=5/8-1/2$. The identity uses the affine nature of both constraint sets. For the nested closed convex sets $[1,\infty)$ and $[2,\infty)$ in $\KSR$, the minimum points instead have squared difference one and energy difference three.

\subsection*{11. A basis is not automatically an isometry}
In Euclidean $\KSR^2$, let $p(a)=a_1-2a_2$, $d(a)=a_1$, $r=1$, and $\alpha=1$. Compute the protected regularized correction using the basis $(2,1)$ of the kernel. What changes if the Gram matrix is omitted?

\textbf{Solution.} The coordinate $c$ gives $a=(2c,c)$, with $S=2$ and $G=5$. The normal equation is $9c=2$, so $a=(4/9,2/9)$. Its objective is $5/9$. Replacing $G$ by one gives $c=2/5$ and $a=(4/5,2/5)$, whose actual ambient objective is $21/25>5/9$. That replacement solves a different metric problem.

\subsection*{12. A robust negative certificate}
A nominal scalar system has $T=1$, $r=1$. The operator error is at most $1/10$ and the data error at most $1/20$. Can a correction of norm at most $1/2$ achieve true residual at most $3/10$? What does the same test say at budget two?

\textbf{Solution.} At budget $1/2$, the nominal optimum is $1/2$. Subtracting $1/20+(1/10)(1/2)$ leaves $2/5>3/10$, so the proposed accuracy is excluded for every admissible true model. At budget two, the lower estimate is negative and clips to zero. It excludes nothing. The disappearance of this certificate is not proof of exact feasibility.

\subsection*{13. A smaller coordinate residual without a better correction}
For the protected example in Section~\ref{KS:sec:inexact-repair}, keep the same correction $(1/4,1/4)$ but replace the kernel basis $B=(1,1)^{\mathsf T}$ by $B'=B/100$. Compute the coordinate normal residual and its actual Hilbert norm.

\textbf{Solution.} In the original basis, $c=1/4$, $G=2$ and $q=3c-1=-1/4$. In the new basis, $c'=25$, $G'=1/5000$ and $q'=-1/400$. The coordinate residual is one hundred times smaller, but
\[
 (q')^{\mathsf T}(G')^{-1}q'=\frac{5000}{160000}=\frac1{32}
 =q^{\mathsf T}G^{-1}q.
\]
The physical correction and the certificate have not improved. The changed number records a changed basis.

\subsection*{14. How small must the solver residual be?}
For $n\ge1$, suppose $\alpha_n=n^{-2}$, the measured-data error is bounded by $n^{-2}$, and the exact-data regularized corrections converge to the desired minimum-norm answer. Does a solver normal-residual tolerance $n^{-2}$ guarantee convergence of the proposals? Give a sufficient replacement.

\textbf{Solution.} The noise term is at most $1/(2n)$, but the solver term is at most one. Thus the displayed tolerance does not provide the required guarantee. The scalar counterexample $T=0$, $r=0$, $a_n=1$ meets it and does not approach the exact optimizer. Replacing it by $n^{-3}$ makes the solver term at most $1/n$, so the triangle inequality and the assumed vanishing bias prove convergence. This is a sufficient worst-case rule, not a necessary condition for every particular operator.

\subsection*{15. Exact optimization of measured data is not exact realization}
Let $d$ be the identity on $\KSR$, protect no nonzero observable, take raw input $f=2$, measured residual $r=1$, and $\alpha=1$. Find the exact measured-data optimizer and compare the fit residual with the actual repaired residual.

\textbf{Solution.} The objective $(a-1)^2+a^2$ is minimized at $a=1/2$. Its normal residual is zero, so the solver certificate is exact. The measured fit residual is $r-da=1/2$. The actual repaired candidate is $f-a=3/2$, whose model residual is $3/2$. The difference is $df-r=1$. A certificate about solving the measured optimization problem does not remove measurement error or the regularization bias.

\subsection*{A checklist for applying the theory}
Before selecting a theorem, state the raw and exact objects, the coefficient structure, the maps that relate views, and the data that must remain fixed. Then ask which conclusion is needed: existence, a description of all answers, an identification-preserving comparison, a stable distance estimate, or an executable finite certificate. Finally record the hypotheses that make that conclusion available.

The examples above fail only when one of those distinctions is suppressed. The missing information is not supplied by giving the same construction a more general name.




% BEGIN retained input sources/monographs/komorebi-space/source/chapters/10_hypothesis_atlas.tex

\chap{A hypothesis atlas}\label{KS:app:hypotheses}
A longer monograph is useful only if the reader can locate the theorem that answers the actual question. This atlas groups the main obligations by the desired conclusion. The references lead to the statements and proofs, not just to examples that resemble the intended application.

\subsection{Existence and identification}
\begin{longtable}{@{}p{.24\linewidth}p{.48\linewidth}p{.21\linewidth}@{}}
\toprule Desired conclusion & Data and hypotheses to establish & Main entry\\\midrule\endhead
A coherent family & A view system with its actual restrictions and all compatibility equations. Nonemptiness is separate. & \ref{KS:KS01}--\ref{KS:KS03}\\[4pt]
A prescribed extension & The declared partial section and an actual vanishing relative obstruction, or matching data on a transition cover. & \ref{KS:KS04}, \ref{KS:KS09}, \ref{KS:KS14}\\[4pt]
A unique extension & Existence and vanishing of the invisible coherent ambiguity, not existence alone. & \ref{KS:KS10}, \ref{KS:KS32}\\[4pt]
An infinite point from finite views & A sequential inverse system with nonempty finite fibers, or another proved lifting/compactness hypothesis. A prescribed value must lie in the persistent image. & \ref{KS:KS54}, \ref{KS:KS58}\\[4pt]
Descent of objects and witnesses & The actual groupoid data, coherent overlap isomorphisms, and a composition-covering family on the specified site. & \ref{KS:KS34}--\ref{KS:KS36}\\[4pt]
Invariance of a boundary problem & An identified boundary and the specified comparison equivalences. Literal freezing of every representative is a different construction. & \ref{KS:KS38}--\ref{KS:KS39}, \ref{KS:KS68}\\[4pt]
The same limit from two paths & Convergence of both paths in a metric space and a pairwise distance budget tending to zero. & \ref{KS:KS74}\\
\bottomrule
\end{longtable}

\subsection{Approximation and numerical meaning}
\begin{longtable}{@{}p{.24\linewidth}p{.48\linewidth}p{.21\linewidth}@{}}
\toprule Desired conclusion & Data and hypotheses to establish & Main entry\\\midrule\endhead
An exact limit of approximations & Completeness, a nonnegative lower semicontinuous residual, vanishing residuals, and summable candidate movement. & \ref{KS:KS71}--\ref{KS:KS72}\\[4pt]
A protected exact limit & In addition, continuity of the protected map and preservation along the candidate path. & \ref{KS:KS73}\\[4pt]
Exact cancellation from zero floor & Closedness of the actual admissible residual image. Without it only closure membership follows. & \ref{KS:KS76}--\ref{KS:KS77}\\[4pt]
A stable residual-to-distance bound & The required closed range or bounded inverse in the declared normed model. In a protected problem, use the restriction to $\ker p$. & \ref{KS:KS43}, \ref{KS:KS56}\\[4pt]
A controlled moving exact space & Two residual certificates, or a proved split-surjective regime with the strict smallness condition. Norm closeness alone does not preserve rank. & Section~\ref{KS:sec:motion}\\[4pt]
An independent object distance & An arrow category and an extended nonnegative cost with zero identity cost and composition subadditivity. Separation is another question. & \ref{KS:KS78}--\ref{KS:KS79}\\
\bottomrule
\end{longtable}

\subsection{Bounded budgets, regularization, and solver certificates}
Every entry has a printed proof. Formal coverage includes the bounded-budget, regularization-limit, noise, solver-certificate, coordinate-invariance, and stopping results under their stated hypotheses. The explicit infinite diagonal estimates, asymptotic budget constant over all real budgets, and rounded noisy rate also have formal counterparts. The formal-source appendix gives the precise correspondence.
\begin{longtable}{@{}p{.24\linewidth}p{.48\linewidth}p{.21\linewidth}@{}}
\toprule Desired conclusion & Data and hypotheses to establish & Main entry\\\midrule\endhead
An exact answer from bounded approximations & A fixed bounded linear equation between real Hilbert spaces, residual convergence, and a uniform correction bound. & \ref{KS:WR01}\\[4pt]
Strong convergence of finite repairs & Nested nonempty closed affine constraints, minimum norm answers, and a uniform energy budget. & \ref{KS:WR02}\\[4pt]
A sharp finite-budget exclusion & The correction ball, the restricted adjoint, and a dual witness in the unit ball. & \ref{KS:WR03}--\ref{KS:WR04}\\[4pt]
Robustness to model uncertainty & A fixed correction space and certified operator and data error bounds. & \ref{KS:WR05}\\[4pt]
A unique protected regularized answer & Positive regularization parameter and the actual Hilbert metric on the protected kernel. & \ref{KS:WR06}--\ref{KS:WR07}\\[4pt]
A convergent regularization limit & Actual attainment of the projected datum, not merely its membership in the closure. & \ref{KS:WR08}\\[4pt]
A noisy convergent answer & The same attained datum, vanishing parameter, and noise divided by the square root of the parameter tending to zero. & \ref{KS:WR09}, \ref{KS:WR12}\\[4pt]
A leading approximation-cost constant & The explicit diagonal model and a proof controlling its infinite Riemann-sum tails. & \ref{KS:WR10}--\ref{KS:WR11}\\[4pt]
A certified inexact optimizer & Positive parameter, the actual normal residual in the protected Hilbert norm, or an exact primal and dual gap. & \ref{KS:WR13}--\ref{KS:WR14}\\[4pt]
A basis-independent solver residual & A full independent kernel basis and its Gram inverse. The normal-equation coordinate residual is a covector. & \ref{KS:WR15}\\[4pt]
A convergent stopping schedule & Attained projected datum and separate vanishing bias, rescaled data noise, and rescaled normal residual. & \ref{KS:WR16}\\\bottomrule
\end{longtable}

\subsection{Three questions before using a finite check}
First determine whether the tested object is the actual object in the theorem. A scalar simulation can illustrate an operator-valued statement without checking every dependent fiber. Next determine whether the test checks a universal finite domain or a selected finite range. Finally identify any limit that remains outside the computation.

The phase extension is exhaustively checked on its eight elements. The diagonal example is different: rational finite sums verify its displayed enclosures, while the proof of nonexistence uses the infinite square-summability condition. Calling both results simply ``tested'' would conceal the most important distinction.

The companion Lean guide records these scopes and identifies the checked examples. General source declarations are checked in the accompanying formal project. Printed examples and conceptual connections have only the formal coverage explicitly documented for them.



% BEGIN retained input sources/monographs/komorebi-space/source/chapters/07_concordance.tex
\chap{Notation and formal correspondence}\label{KS:app:concordance}
\subsection{Notation}
\begin{longtable}{@{}p{.25\linewidth}p{.7\linewidth}@{}}
\toprule
Symbol & Meaning\\\midrule\endhead
\(\KSV,F_v,f^*\) & A declared view category, its fibers, and contravariant restrictions.\\
\(\Gamma(F)\) & The strict set of coherent global families.\\
\(\Gamma_U(F),\KSExt_U(s)\) & Partial sections on a full patch and their global extension fiber.\\
\(\delta,\partial,d,D\) & Respectively the view differential, tower differential, a model residual, and a Hilbert defect operator. Each target is declared where used.\\
\(p,p_U,p_L\) & A protected observable, restriction to prescribed views, and an innovation observable depending on the splitting.\\
\(\mathfrak o_U,H_U\) & The existence obstruction and the group of invisible coherent changes. Additive and nonabelian conventions are stated separately.\\
\(C_j,J_j,Q_N,S_N\) & Coarsener, chosen lift, restored coarse remainder, and retained finite reconstruction.\\
\(c_f,a_v\) & Coherent nonabelian transition labels and gauge values.\\
\(\mathfrak K_G,J_{\rm comp}\) & The cocycle patch groupoid and its composition-support topology.\\
\(\mathcal E_b(R)\) & Realizations equipped with an actual identification of their boundary with \(b\).\\
\(T_i,K_{ij},R_p\) & Transport to the root, the corresponding root loop operator, and reconstruction from the root.\\
\(\mathcal H,q,P\) & Root residual, an orthogonal kernel projector, and the explicitly defined model repair. A graph repair need not be orthogonal in the global maximum norm.\\
\(\pi_{W,U},H_U\) & The map forgetting part of the prescription and the corresponding invisible coherent ambiguity.\\
\(\phi,\operatorname{EventualHom}\) & A monotone source schedule and transport classes modulo eventual equality.\\
\(\chi_A,M=SA-BC\) & Split exact/residual coordinates and the mismatch between chosen repair lifts.\\
\(\omega,Q_n,C_n,b_n\) & A residual, its prefix-product rate, movement coefficient, and summable step budget. Here $Q_n$ is not the dyadic remainder operator $Q_N$.\\
\(\nu_B,d_W,d_W^{\mathrm{sym}}\) & Residual floor, extended directed transport distance, and its symmetric maximum.\\
\(c(x,y),\kappa(x,y)\) & The two-bit pair phase and its antisymmetric scalar-gauge invariant.\\
\(H,T,r\) & In the budget chapters, $H=\ker p$ with its inherited Hilbert metric, $T=d|_H$, and $r=df$ for the true raw residual. Measured data are declared separately.\\
\(\beta_{T,r}(M),a_\alpha\) & Minimum residual at correction budget $M$, and the quadratic regularized correction.\\
\(C,b,q\) & Closure of the correction image, its projection of the residual, and the irreducible residual floor.\\
\(v,g,e,\mathcal D_\alpha\) & Fit residual, Hilbert normal residual, optimizer error, and regularized dual objective in the solver chapters.\\
\(B,G,q\) & Kernel basis, its Gram matrix, and the coordinate normal covector. This $q$ is not the scalar residual floor used earlier.\\
\(\simeq,\simeq_L\) & A specified equivalence, and a continuous linear equivalence. Neither notation alone asserts an isometry.\\
\bottomrule
\end{longtable}

\subsection{From the text to the formal source}
The table below locates the foundational results in the accompanying Lean development. Numbered locations refer to results in the main text. Page references locate discussions without a separate numbered statement. Module names are relative to the \path{KomorebiSpace} directory in the \path{lean/} source tree. A row may collect several declarations that together express the displayed result.

The table is a navigation aid. A declaration's full type specifies its quantifiers, hypotheses, and conclusion. Those types are the authority for formal coverage. The correspondence does not assign a separate machine-checked counterpart to every explanatory sentence or illustrative calculation. Appendix~\ref{KS:app:written} gives the complementary map for bounded realization, budgets, regularization, and certificates.


% BEGIN retained input sources/monographs/komorebi-space/source/claim_table.tex
\small
\begin{longtable}{@{}p{.11\linewidth}p{.35\linewidth}p{.47\linewidth}@{}}
\toprule Location & Topic & Lean modules\\\midrule\endhead
\ref{KS:KS01} & Limit universal property & \path{ViewSections}\\[3pt]
\ref{KS:KS02} & Functorial transport & \path{ViewSections}\\[3pt]
\ref{KS:KS03} & Terminal evaluation & \path{ViewSections}\\[3pt]
\ref{KS:KS04} & Strict section gluing & \path{SectionGluing}\\[3pt]
\ref{KS:KS05} & Object-only cover failure & \path{ViewBoundaryExamples}\\[3pt]
\ref{KS:KS06} & Correspondence fibers & \path{ViewCorrespondence}\\[3pt]
\ref{KS:KS07} & Nonempty fibers without a section & \path{ViewBoundaryExamples}\\[3pt]
\ref{KS:KS08} & Additive view differential & \path{SectionObstruction}\\[3pt]
\ref{KS:KS09} & Additive extension criterion & \path{SectionObstruction}\\[3pt]
\ref{KS:KS10} & Extension ambiguity torsor & \path{SectionExtensionFibers}\\[3pt]
p.~\pageref{KS:src:KS11} & Parity seed example & \path{SectionExtensionExample}\\[3pt]
\ref{KS:KS12} & Exactness at partial sections & \path{SectionExactSequence}\\[3pt]
p.~\pageref{KS:src:KS13} & Extension by stages & \path{SectionExtensionStages}\\[3pt]
\ref{KS:KS14} & Protected correction & \path{Correction}\\[3pt]
\ref{KS:KS15} & Obstruction witness & \path{Correction}\\[3pt]
\ref{KS:KS16} & Exact tower maps and realized points & \path{Towers}\newline \path{Realization}\\[3pt]
\ref{KS:KS17} & Tower differential composition & \path{Towers}\newline \path{TowerCorrection}\\[3pt]
\ref{KS:KS18} & Cofinal point equivalence & \path{Cofinal}\\[3pt]
\ref{KS:KS19} & Geometric operator repair & \path{RepairDynamics}\newline \path{RepairProjection}\\[3pt]
\ref{KS:KS20} & Constructed split tower correction & \path{SplitTowerRepair}\\[3pt]
\ref{KS:KS21} & Split semantic and residual coordinates & \path{SplitTowerCoordinates}\newline \path{SplitTowerSemantics}\\[3pt]
\ref{KS:KS22} & Growing identity repair & \path{GrowingTower}\\[3pt]
p.~\pageref{KS:src:KS23} & A bounded residual without bounded primitive & \path{IdentityTowerObstruction}\\[3pt]
p.~\pageref{KS:src:KS24} & Dependence on target splitting & \path{SplittingDependence}\\[3pt]
p.~\pageref{KS:src:KS25} & Lossless dyadic residual split & \path{DyadicResidualSplit}\\[3pt]
\ref{KS:KS26} & Cascade telescope and persistent part & \path{CascadeProjections}\newline \path{CoarseningCascade}\\[3pt]
\ref{KS:KS27} & Recovery by density & \path{CascadeDensity}\\[3pt]
\ref{KS:KS28} & Dyadic finite-coordinate recovery & \path{DyadicInfiniteReconstruction}\\[3pt]
\ref{KS:KS29} & Shift cascade and c0 & \path{SequenceCascade}\\[3pt]
\ref{KS:KS30} & Gauge groupoid and section transport & \path{GroupViewSystems}\newline \path{CocycleGaugeGroupoid}\\[3pt]
\ref{KS:KS31} & Nonabelian extension obstruction & \path{NonabelianExtension}\\[3pt]
\ref{KS:KS32} & Nonabelian extension torsor & \path{NonabelianExtensionFibers}\\[3pt]
\ref{KS:KS33} & S3 loop model & \path{CocycleExamples}\\[3pt]
\ref{KS:KS34} & Effective groupoid descent & \path{DescentEquivalence}\newline \path{DescentComparison}\\[3pt]
\ref{KS:KS35} & Composition topology & \path{CompositionSite}\\[3pt]
\ref{KS:KS36} & Standard stackhood & \path{CompositionStack}\\[3pt]
\ref{KS:KS37} & Equivalent view presentations & \path{SectionViewEquivalence}\newline \path{RelativeSections}\newline \path{RelativeAmbiguity}\\[3pt]
\ref{KS:KS38} & Identified-boundary realization groupoid & \path{MarkedRealization}\\[3pt]
\ref{KS:KS39} & Marked transport across a square & \path{MarkedTransport}\\[3pt]
\ref{KS:KS40} & Transported path telescope & \path{AmplifiedTransport}\\[3pt]
\ref{KS:KS41} & Holonomy classification & \path{AmplifiedHolonomy}\\[3pt]
\ref{KS:KS42} & Orthogonal exact repair & \path{HilbertKernelProjection}\\[3pt]
\ref{KS:KS43} & Closed range equivalence & \path{HilbertKernelProjection}\\[3pt]
\ref{KS:KS44} & Amplified holonomy gain & \path{AmplifiedRepair}\newline \path{HilbertHolonomyRepair}\\[3pt]
p.~\pageref{KS:src:KS45} & Flat transport amplification & \path{AmplifiedExamples}\\[3pt]
\ref{KS:KS46} & Two-stage joint projection & \path{QuantitativeDescent}\\[3pt]
\ref{KS:KS47} & Joint closed range & \path{DescentClosedRange}\\[3pt]
\ref{KS:KS48} & Two-sided projector motion & \path{MovingKernelBounds}\\[3pt]
\ref{KS:KS49} & Derived perturbed certificate & \path{SplitKernelPerturbation}\\[3pt]
\ref{KS:KS50} & Transport between moving kernels & \path{SplitKernelPerturbation}\\[3pt]
\ref{KS:KS51} & Exact moving-line example & \path{MovingKernelExamples}\\[3pt]
p.~\pageref{KS:src:KS52} & Rank loss and radius endpoint & \path{MovingKernelBounds}\newline \path{MovingKernelExamples}\\[3pt]
p.~\pageref{KS:src:KS53} & Motion of quantitative descent & \path{MovingDescent}\\[3pt]
\ref{KS:KS54} & Finite prefixes without a point & \path{FinitePrefixBoundary}\\[3pt]
\ref{KS:KS55} & Sharper dyadic support criterion & \path{DyadicSupport}\\[3pt]
\ref{KS:KS56} & Nearest protected Hilbert repair & \path{ProtectedHilbertRepair}\\[3pt]
\ref{KS:KS57} & Finite positive and negative certificates & \path{ProtectedLinearCertificates}\\[3pt]
\ref{KS:KS58} & Finite-fiber inverse-limit existence & \path{FiniteRealizationExistence}\\[3pt]
\ref{KS:KS59} & Obstructions under stronger prescription & \path{ObstructionRestriction}\\[3pt]
\ref{KS:KS60} & Prescribed choices and obstructed intermediate steps & \path{ExtensionChoiceExample}\\[3pt]
\ref{KS:KS61} & Scheduled transports and eventual equality & \path{ReindexedTransports}\newline \path{EventualCategory}\\[3pt]
\ref{KS:KS62} & Faithfulness from observable finite stages & \path{RealizationFaithfulness}\\[3pt]
\ref{KS:KS63} & Finite observation characterizes realized transport & \path{FiniteObservation}\\[3pt]
\ref{KS:KS64} & Cofinal transport equivalence & \path{CofinalTransports}\\[3pt]
\ref{KS:KS65} & Split repair coordinates and choice dependence & \path{SplitDefectCharts}\\[3pt]
\ref{KS:KS66} & Triangular representation change and repair mismatch & \path{ResidualNaturality}\\[3pt]
\ref{KS:KS67} & Residual sampling and repair across scales & \path{ResidualSampling}\newline \path{SamplingRepair}\\[3pt]
\ref{KS:KS68} & Why literal full framing is presentation dependent & \path{FramingBoundaryExamples}\\[3pt]
\ref{KS:KS69} & Measured and uncertain residual certificates & \path{ResidualPerturbation}\\[3pt]
\ref{KS:KS70} & Stable repair when the old exact space is preserved & \path{ResidualPerturbation}\\[3pt]
\ref{KS:KS71} & Exactification from summable movement & \path{Exactification}\\[3pt]
\ref{KS:KS72} & Variable-rate exactification and posterior bounds & \path{Exactification}\\[3pt]
\ref{KS:KS73} & Controlled approximation annihilates protected obstruction & \path{LinearExactification}\\[3pt]
\ref{KS:KS74} & Path comparison and nonunique protected answers & \path{Exactification}\newline \path{RepairExamples}\\[3pt]
\ref{KS:KS75} & What cofinal residual-profile sampling preserves & \path{Profiles}\\[3pt]
\ref{KS:KS76} & Obstruction floors detect closure, not exact membership & \path{ObstructionFloor}\\[3pt]
\ref{KS:KS77} & Persistent floors exclude all admissible repair sequences & \path{ObstructionFloor}\\[3pt]
\ref{KS:KS78} & Independent transport costs and directed geometry & \path{WeightedGeometry}\\[3pt]
\ref{KS:KS79} & Zero arrows can collapse object distance & \path{WeightedGeometry}\\[3pt]
\ref{KS:KS80} & An associative phase with no scalar flattening & \path{ProjectiveExample}\\[3pt]
\bottomrule\end{longtable}
\normalsize



\subsection{Using the correspondence}
For a proof involving a restriction or a change of representation, compare the source and target spaces before identifying two displayed formulas. For an analytic estimate, check the norm, the permitted correction space, and the operator whose adjoint is used. The declarations preserve these distinctions explicitly. A continuous linear equivalence alone is not an isometry, and a coordinate equation alone does not fix the metric of its residual.

The source contains further definitions, lemmas, and examples beyond this table. The Lean guide explains its organization and the formal treatment of the principal analytic results. Verification records retain the exact source inventories and declaration types from the accepted compilation.



% BEGIN retained input sources/monographs/komorebi-space/source/chapters/11_written_extensions.tex

\chap{Proof and formal verification}\label{KS:app:written}
The printed arguments and the accompanying Lean development have complementary roles. The text explains constructions and their mathematical context. A Lean declaration fixes a precise statement whose proof is checked by the kernel. A correspondence between them records the actual hypotheses and conclusion, while allowing the formal proof to use a different route.

All \KSEditionPrincipalGroups\ principal theorem groups on bounded realization, budgets, regularization, and certificates have complete formal counterparts. The table identifies these statements by their numbers in the main text. This coverage concerns the stated results, rather than every proof route or illustrative calculation.

\begin{longtable}{@{}p{.11\linewidth}p{.60\linewidth}p{.22\linewidth}@{}}
\toprule Result & Topic & Formal counterpart\\\midrule\endhead
\ref{KS:WR01} & Bounded approximate feasibility & Verified\\[4pt]
\ref{KS:WR02} & Nested affine minimum corrections & Verified\\[4pt]
\ref{KS:WR03} & Attainment and shape of the budget profile & Verified\\[4pt]
\ref{KS:WR04} & Sharp budget certificate & Verified\\[4pt]
\ref{KS:WR05} & Robust budget comparison & Verified\\[4pt]
\ref{KS:WR06} & Unique regularized correction & Verified\\[4pt]
\ref{KS:WR07} & Parameter and budget comparison & Verified\\[4pt]
\ref{KS:WR08} & Regularization limit alternative & Verified\\[4pt]
\ref{KS:WR09} & Data sensitivity and convergent parameter choice & Verified\\[4pt]
\ref{KS:WR10} & Infinite diagonal family and certified tails & Verified\\[4pt]
\ref{KS:WR11} & Sharp diagonal budget asymptotic & Verified\\[4pt]
\ref{KS:WR12} & Rounded noisy diagonal rate & Verified\\[4pt]
\ref{KS:WR13} & Normal-residual solver certificate & Verified\\[4pt]
\ref{KS:WR14} & Regularized primal and dual gap & Verified\\[4pt]
\ref{KS:WR15} & Coordinate certificate and arbitrary basis change & Verified\\[4pt]
\ref{KS:WR16} & Consistent inexact stopping rule & Verified\\
\bottomrule\end{longtable}

\subsection{Bounded realization and budget geometry}\label{KS:app:wr-source-map}
The formal bounded-realization theorem concerns bounded linear maps between real Hilbert spaces. It gives an exact preimage within the same norm bound as the approximating corrections. Its proof uses compactness of a dual ball and the Riesz isometric equivalence. The printed proof uses the weak-subsequence lemma, Lemma~\ref{KS:lem:weak-subsequence}, whose particular proof route is not separately formalized here. Neither route requires separability or closed range. Neither asserts strong convergence of the original approximating sequence.

The nested-constraint theorem constructs the minimum-norm points of nonempty closed affine sets. It proves the exact energy identity, norm monotonicity, and the equivalence between a nonempty intersection and bounded minimum norms. Under that bound, the corrections converge strongly to the unique minimum-norm point of the intersection. If the intersection is empty, their norms tend to infinity. The affine assumption is retained: for general convex sets the corresponding energy statement is an inequality.

\begin{longtable}{@{}p{.38\linewidth}p{.55\linewidth}@{}}
\toprule Namespace & Representative declarations\\\midrule\endhead
\path{BoundedApproximateRealization} & \path{exists_solution_of_bounded_residual_tendsto_zero}\newline\path{mem_range_iff_bounded_approximation}\newline\path{norm_tendsto_atTop_of_not_mem_range}\\[5pt]
\path{NestedMinimumCorrections} & \path{nested_affine_energy}\newline\path{nested_affine_intersection_iff_bounded}\newline\path{nested_affine_strong_realization}\newline\path{nested_affine_escape}\\[5pt]
\path{ProtectedBoundedRealization} & \path{exists_exact_repair_of_bounded_approximation}\\
\bottomrule\end{longtable}
All namespaces in this appendix have the prefix \path{KomorebiSpace}. The protected application uses the actual subtype $\ker p$ with its inherited Hilbert norm. It constructs a repaired candidate preserving the declared observable and the original correction budget.

The budget-profile theorems concern bounded operators between real complete Hilbert spaces. The budget is nonnegative and may be zero. Attainment of a residual minimum over a bounded correction ball does not require the entire operator range to be closed. The sharp support formula supplies an actual maximizing witness, expressed formally by \KScode{IsGreatest}. It does not infer attainment from the mere existence of a supremum.

\begin{longtable}{@{}p{.38\linewidth}p{.55\linewidth}@{}}
\toprule Namespace & Representative declarations\\\midrule\endhead
\path{BudgetProfile} & \path{exists_budget_minimizer}\newline\path{beta_eq_iInf}\newline\path{beta_tendsto_floor}\newline\path{beta_eq_floor_iff}\newline\path{beta_eq_zero_iff}\\[5pt]
\path{BudgetProfileGeometry} & \path{beta_antitone}\newline\path{beta_zero}\newline\path{beta_convex_combination}\newline\path{beta_profile_bounds}\\[5pt]
\path{BudgetDuality} & \path{dualValue_le_beta}\newline\path{exists_attaining_dual_witness}\newline\path{sharp_budget_duality}\newline\path{strict_exclusion_certificate}\\[5pt]
\path{BudgetRobustness} & \path{candidate_residual_le}\newline\path{beta_abs_sub_le}\newline\path{nominal_dual_lower_certificate}\newline\path{excludes_budget_of_nominal_dual_certificate}\\
\bottomrule\end{longtable}

The robustness comparison uses the same correction space for both problems. Its bound $\delta+\eta M$ controls data and operator uncertainty at budget $M$. A changing protected kernel is a different problem and is not covered by that estimate.

\subsection{Regularization, noise, and stopping}
For positive parameter $\alpha$, the regularized normal operator is constructed as a continuous linear equivalence. Its inverse defines the correction. The formal development proves uniqueness, the energy identity, parameter monotonicity, budget optimality, and the two data-sensitivity estimates. The fixed-space protected application uses the restricted operator and its actual adjoint.

The limit theory distinguishes projection onto the closure of the range from membership in the range itself. Fitted values converge to the projected datum and residual norms converge to the irreducible floor. If that projected datum is attained, the corrections converge in norm to its unique minimum-norm preimage. Otherwise their norms tend to infinity. The data and stopping theorems use this limit theorem together with their explicit parameter, noise, and solver-tolerance assumptions.

\begin{longtable}{@{}p{.38\linewidth}p{.55\linewidth}@{}}
\toprule Namespace & Representative declarations\\\midrule\endhead
\path{RegularizedHilbertRepair} & \path{exists_unique_minimizer}\newline\path{energy_identity}\newline\path{residual_best_at_budget}\newline\path{parameter_norm_antitone}\newline\path{parameter_residual_monotone}\newline\path{data_noise_bound}\newline\path{data_fitted_bound}\\[5pt]
\path{RegularizationLimits} & \path{fitted_tendsto_projectedDatum}\newline\path{residual_norm_tendsto_floor}\newline\path{weighted_penalty_tendsto_zero}\newline\path{norm_tendsto_atTop_of_not_mem_range}\\[5pt]
\path{RegularizationStrongLimit} & \path{exists_unique_minimum_and_limit}\\[5pt]
\path{RegularizationNoiseLimits} & \path{data_noise_coefficient_attained}\newline\path{noisy_parameter_choice}\newline\path{inexact_parameter_choice}\newline\path{stopping_rule_sequence}\\
\bottomrule\end{longtable}

The stopping theorem compares proposals, exact regularized corrections, and the intended minimum correction in one fixed Hilbert model. It includes the pointwise three-term bound and the convergence conclusion. Protection is exact at every index. The explicit diagonal schedule at the end of Section~\ref{KS:sec:stopping-rules} combines this solver bound with the formalized noisy rate in Corollary~\ref{KS:WR12}. The declaration \path{DiagonalRegularization.diagonal_inexact_rounded_rate} proves this specialization with the literal ceiling, the data-noise hypothesis, and the actual measured-data normal-residual bound. The exact-data bias estimate is derived rather than assumed.

\subsection{Certificates and their coordinate metric}
The solver certificate controls correction error, fitted error, and objective excess through the actual Hilbert normal residual. The primal and dual identity supplies objective bracketing and a correction-distance bound for an arbitrary dual witness. These are universal statements over the declared Hilbert spaces, with $\alpha>0$ wherever coercivity is used.

\begin{longtable}{@{}p{.38\linewidth}p{.55\linewidth}@{}}
\toprule Namespace & Representative declarations\\\midrule\endhead
\path{RegularizedHilbertRepair} & \path{correction_error_bound}\newline\path{fitted_error_bound}\newline\path{objective_gap_bound}\\[5pt]
\path{RegularizationCertificates} & \path{primal_dual_gap}\newline\path{dual_maximizer}\newline\path{gap_controls_correction}\newline\path{residual_witness_gap}\\[5pt]
\path{CoordinateCovector} & \path{gram_isUnit}\newline\path{protected_coordinate_certificate}\newline\path{computed_normalResidual_changedBasis}\newline\path{computed_certificate_changedBasis}\newline\path{changedBasis_coordinates}\\
\bottomrule\end{longtable}

The coordinate certificate uses a basis of the entire permitted correction space, with its inherited Hilbert metric. The proof derives invertibility of the Gram matrix from that basis, identifies the coordinate residual with the pairings of the Hilbert normal residual, and proves
\[
 g=BG^{-1}q,\qquad \KSnorm g^2=q^{\mathsf T}G^{-1}q.
\]
Every invertible basis change is allowed. The changed coordinates describe the same correction, the residual transforms as a covector, and the displayed squared norm is invariant. The proof includes the zero-dimensional correction space. It does not replace the Hilbert norm by a coordinate supremum norm or assume that the basis is orthonormal. These metric identities hold for every real $\alpha$. The optimizer estimates still require positive $\alpha$.

\subsection{The infinite diagonal example}
Proposition~\ref{KS:WR10}, Theorem~\ref{KS:WR11}, and Corollary~\ref{KS:WR12} have formal counterparts in the namespace \path{KomorebiSpace.DiagonalRegularization}. The formal space is the real Hilbert space \KScode{lp} at exponent two. Its natural index starts at zero, so formal coordinate $j$ denotes printed coordinate $n=j+1$. The operator, data, corrections, residuals, and hard prefixes are elements or maps of that actual infinite space.

\begin{longtable}{@{}p{.16\linewidth}p{.77\linewidth}@{}}
\toprule Result & Declarations in \path{DiagonalRegularization}\\\midrule\endhead
\ref{KS:WR10} & \path{proposition_30_1}\newline\path{a_unique_minimizer}\\[5pt]
\ref{KS:WR11} & \path{sharp_diagonal_budget_asymptotic}\\[5pt]
\ref{KS:WR12} & \path{corollary_30_3}\\[5pt]
Inexact rate & \path{diagonal_inexact_rounded_rate}\\[5pt]
Range & \path{T_not_closed_range}\newline\path{w_mem_closure_not_range}\\
\bottomrule\end{longtable}

The first pair identifies the universal regularized correction, proves it is the unique global minimizer, equates the infinite squared norms with $A_m$ and $B_m$, and proves all four global bounds and both nonnegative omitted-tail bounds. The asymptotic theorem proves both scalar limits, the budget-profile limit $\pi/4$ over all real $M$ tending to infinity, and the hard-prefix comparison with constant one. Its budget conclusion is not restricted to the regularization mesh.

The noisy theorem uses the actual attainable datum and its minimum-norm preimage. It retains the literal ceiling in $m=\lceil\delta^{-2/3}\rceil$ and proves both displayed inequalities for every Hilbert perturbation within the stated noise bound. The bias estimate is derived within the model. The inexact rate adds the stated solver tolerance to this same model and parameter choice.

The range obstruction uses actual unit coordinate vectors with residual norm $1/(j+1)$ and applies the general unit-witness theorem. The hard-prefix images converge in Hilbert norm to $w$, placing it in the closure of the range. An exact preimage would have every coordinate equal to one, contradicting square summability. These declarations formalize the concrete example following Theorem~\ref{KS:KS43}. Exact finite enclosures remain separate illustrations of the infinite results.

\subsection{Verification record and trust boundary}
The source contains \KSEditionTargets\ compilation targets (\KSEditionProductionModules\ production modules and \KSEditionRoots\ entry points), \KSEditionTheorems\ explicit theorem or lemma declarations, and \KSEditionDeclarations\ exported declarations. The latter includes definitions and generated helpers; neither count measures independent mathematical discoveries.

Fresh compilation treated warnings as errors and disabled automatic implicit parameters. The complete declaration audit checked defining modules and recursive axioms, rejecting unsafe declarations and project-specific axioms. Source, compiler, cache, and generated-object identities bind this replay.

All 56 diagnostic controls had their specified outcomes: 18 successes, 35 intended elaboration rejections, and three axiom-policy rejections. The 52 inherited controls were rerun alongside four new controls. Earlier records remain historical.

The environment is Lean 4.34.0-rc2 with mathlib commit
\path{aeb72f1e9981b6d35b426b16dca7759be3e7ba06}.
Permitted axioms are \KScode{propext}, \KScode{Classical.choice}, and \KScode{Quot.sound}. Lean's compiler and kernel and the pinned external dependency cache remain trusted; the cache was recorded but not rebuilt. Finite executable examples provide separate computational evidence; they are not verified extractions of the Lean proofs.



% BEGIN retained input sources/monographs/komorebi-space/source/chapters/12_proof_tools.tex

\chap{Analytic facts used in the proofs}\label{KS:app:proof-tools}
This appendix supplies the analytic ingredients used at several points in the book. They are classical facts, not additional novelty claims. Writing out their proofs makes the dependence on completeness, convexity, the inner product, and the declared operator norm visible. The statements are restricted to the real Hilbert and real Banach settings used in the main arguments. They provide background support for the main results.

The projection and inverse constructions have standard formal counterparts in the mathematical library cited in \cite{KS:mathlib-projection,KS:mathlib-banach}. The regularization context is also discussed in \cite{KS:clason-inverse}. The arguments below are ordinary written proofs. The accompanying Lean project does not by itself constitute a formal verification of these printed proofs.

\subsection{Projection onto a closed convex set}
\begin{KSlemma}[Existence, uniqueness, and the variational inequality]\label{KS:lem:projection-tool}
Let \(H\) be a real Hilbert space, let \(C\subseteq H\) be nonempty, closed, and convex, and fix \(x\in H\). There is a unique \(c\in C\) minimizing \(\KSnorm{x-c}\). It is characterized by
\begin{equation}\label{KS:eq:projection-variational}
 \langle x-c,z-c\rangle\le0\qquad(z\in C).
\end{equation}
If \(C=c_0+L\) is a closed affine subspace with direction space \(L\), the condition is \(x-c\perp L\).
\end{KSlemma}
\begin{proof}
\textbf{Produce a Cauchy minimizing sequence.} Put \(d=\inf_{z\in C}\KSnorm{x-z}\). Nonemptiness makes this a finite nonnegative real number. Choose \(c_n\in C\) so that \(\KSnorm{x-c_n}\to d\). Convexity places \((c_n+c_m)/2\) in \(C\). The parallelogram identity gives
\[
 \KSnorm{c_n-c_m}^2
 =2\KSnorm{x-c_n}^2+2\KSnorm{x-c_m}^2
   -4\left\|x-\frac{c_n+c_m}{2}\right\|^2
 \le2\KSnorm{x-c_n}^2+2\KSnorm{x-c_m}^2-4d^2.
\]
The right side tends to zero as both indices tend to infinity. Thus \((c_n)\) is Cauchy. Completeness supplies a limit \(c\in H\), closedness puts it in \(C\), and continuity of the norm gives \(\KSnorm{x-c}=d\). A minimizing sequence in an arbitrary closed set need not have been Cauchy. The midpoint was the essential use of convexity.

\textbf{Prove uniqueness.} If \(c,c'\in C\) both attain \(d\), the same midpoint inequality gives
\[
 \KSnorm{c-c'}^2\le2d^2+2d^2-4d^2=0.
\]
Therefore \(c=c'\).

\textbf{Derive the variational inequality.} For \(z\in C\) and \(0<t\le1\), convexity makes \(c+t(z-c)\) feasible. Minimality and the inner-product expansion yield
\[
 0\le\KSnorm{x-c-t(z-c)}^2-\KSnorm{x-c}^2
   =-2t\langle x-c,z-c\rangle+t^2\KSnorm{z-c}^2.
\]
Divide by the positive number \(t\) and let \(t\downarrow0\). This gives \eqref{KS:eq:projection-variational}. Conversely, if that inequality holds, then for every \(z\in C\),
\[
 \KSnorm{x-z}^2
 =\KSnorm{x-c}^2-2\langle x-c,z-c\rangle+\KSnorm{z-c}^2
 \ge\KSnorm{x-c}^2.
\]
It therefore characterizes the unique minimizer.

For an affine subspace, \(c+h\) and \(c-h\) are feasible for every \(h\in L\). Applying the inequality with both signs gives \(\langle x-c,h\rangle=0\). Orthogonality conversely implies the same variational inequality. In this case the preceding expansion becomes the Pythagorean equality, not merely an inequality.
\end{proof}

For a closed linear subspace \(L\), denote its projection by \(P_L\). The lemma gives a unique decomposition
\[
 x=P_Lx+(x-P_Lx),\qquad P_Lx\in L,\quad x-P_Lx\in L^\perp.
\]
Adding two such decompositions, or multiplying one by a scalar, proves linearity by uniqueness. The Pythagorean identity gives \(\KSnorm{P_Lx}\le\KSnorm{x}\), so \(P_L\) is bounded. It fixes \(L\), annihilates \(L^\perp\), and satisfies \(P_L^2=P_L\). For \(x,y\in H\), orthogonality gives
\[
 \langle P_Lx,y\rangle=\langle P_Lx,P_Ly\rangle
 =\langle x,P_Ly\rangle.
\]
Thus \(P_L\) is self-adjoint. These properties are conclusions of the construction, not extra data assumed for a proposed projection.

\subsection{Riesz representation and the adjoint}
\begin{KSlemma}[Representing a bounded functional]\label{KS:lem:riesz-tool}
For every bounded real linear functional \(\ell:H\to\KSR\) on a real Hilbert space, there is a unique \(u\in H\) with
\[
 \ell(h)=\langle h,u\rangle\quad(h\in H),\qquad \KSnorm{\ell}=\KSnorm{u}.
\]
\end{KSlemma}
\begin{proof}
If \(\ell=0\), take \(u=0\). Otherwise choose \(x\) with \(\ell(x)\ne0\). Boundedness makes \(L=\ker\ell\) closed. Use Lemma~\ref{KS:lem:projection-tool} to form \(v=x-P_Lx\). Then \(v\perp L\) and \(\ell(v)=\ell(x)\ne0\), in particular \(v\ne0\).

For any \(h\in H\), the vector \(h-\ell(h)v/\ell(v)\) belongs to \(L\). Pairing it with \(v\) therefore gives
\[
 \langle h,v\rangle=\frac{\ell(h)}{\ell(v)}\KSnorm{v}^2.
\]
Consequently \(u=\ell(v)v/\KSnorm{v}^2\) represents \(\ell\). If \(u'\) also represents it, pairing the difference of the two formulas with \(h=u-u'\) gives \(\KSnorm{u-u'}^2=0\), hence uniqueness.

Cauchy--Schwarz gives \(|\ell(h)|\le\KSnorm{u}\KSnorm{h}\), so \(\KSnorm{\ell}\le\KSnorm{u}\). If \(u\ne0\), the unit vector \(h=u/\KSnorm{u}\) gives the reverse inequality. The zero case was already included.
\end{proof}

Let \(T:H\to R\) be a bounded linear map between real Hilbert spaces. For fixed \(z\in R\), the functional \(h\mapsto\langle Th,z\rangle\) is bounded by \(\KSnorm{T}\KSnorm{z}\KSnorm{h}\). The lemma therefore defines a unique vector \(T^*z\in H\) with
\[
 \langle Th,z\rangle=\langle h,T^*z\rangle.
\]
Uniqueness of representatives proves that \(T^*\) is linear. The same bound proves \(\KSnorm{T^*z}\le\KSnorm{T}\KSnorm{z}\). Conversely,
\[
 \KSnorm{Th}=\sup_{\KSnorm{z}\le1}|\langle Th,z\rangle|
 \le\KSnorm{T^*}\KSnorm{h},
\]
where the first equality follows by testing \(z=Th/\KSnorm{Th}\) when \(Th\ne0\). Hence \(\KSnorm{T^*}=\KSnorm{T}\).

The defining pairing also gives the identity
\begin{equation}\label{KS:eq:adjoint-range-orthogonal}
 (\KSran T)^\perp=\ker T^*.
\end{equation}
Indeed, \(z\) is orthogonal to every \(Th\) exactly when \(\langle h,T^*z\rangle=0\) for every \(h\), which is equivalent to \(T^*z=0\). Orthogonal complements do not change on taking closure. The projection decomposition then gives
\[
 \overline{\KSran T}=(\ker T^*)^\perp.
\]
This is why a self-adjoint injective operator has dense range. It does not yet say its range is closed. In the regularized normal-operator proof, the coercive lower bound supplies that separate fact.

\subsection{Why a bounded bijection has a bounded inverse}
The closed-range arguments use the inverse theorem only after proving that the relevant source and target subspaces are complete and that the restricted map is bijective. The following derivation records this analytic step explicitly.

\begin{KSlemma}[Baire's complete-space principle]\label{KS:lem:baire-tool}
A nonempty complete metric space cannot be the union of countably many closed sets with empty interior.
\end{KSlemma}
\begin{proof}
Suppose \(F_1,F_2,\ldots\) are closed with empty interior. Choose a nonempty open ball whose closed ball \(B_1\), of positive radius at most \(1/2\), avoids \(F_1\). Such a choice is possible because the complement of \(F_1\) is open and meets every nonempty open set. Recursively choose a closed ball \(B_{n+1}\) of positive radius at most \(2^{-(n+1)}\), contained in the open ball with the same center and radius as \(B_n\), and disjoint from \(F_{n+1}\). The intersection of that open ball with the complement of \(F_{n+1}\) is nonempty and open, so it contains a sufficiently small closed ball. This also covers the case of isolated points.

The centers form a Cauchy sequence: all centers after stage \(n\) lie in \(B_n\), whose diameter is at most \(2^{1-n}\). Completeness gives a limit. For each fixed \(n\), a tail of centers lies in the closed set \(B_n\), so the limit lies there. It therefore belongs to no \(F_n\), contradicting that these sets cover the space.
\end{proof}

\begin{KSlemma}[Continuous inverse for Banach spaces]\label{KS:lem:inverse-tool}
Let \(S:X\to Y\) be a bounded linear surjection between real Banach spaces. There is a constant \(C<\infty\) such that every \(y\in Y\) has a preimage \(x\) with \(\KSnorm{x}\le C\KSnorm{y}\). If \(S\) is also injective, its inverse is bounded linear.
\end{KSlemma}
\begin{proof}
If \(Y=\{0\}\), take the zero preimage. Otherwise write \(B_X\) for the closed unit ball. Surjectivity gives
\[
 Y=\bigcup_{n\ge1}\overline{S(nB_X)}.
\]
By Lemma~\ref{KS:lem:baire-tool}, at least one of the closed sets on the right has nonempty interior. Thus there are \(n\ge1\), \(y_0\in Y\), and \(\rho>0\) such that
\[
 B(y_0,\rho)\subseteq\overline{S(nB_X)}.
\]
Here \(B(y_0,\rho)\) is an open ball. If \(\KSnorm{h}<\rho\), both \(y_0+h\) and \(y_0\) belong to this closure. Choose approximating image sequences for both. Subtracting their preimages shows \(h\in\overline{S(2nB_X)}\). Scaling by \(2n\) gives a constant \(c=\rho/(2n)>0\) with
\begin{equation}\label{KS:eq:inverse-closure-ball}
 B(0,c)\subseteq\overline{S(B_X)}.
\end{equation}
This only gives approximate preimages. We next turn them into an actual preimage with a summable budget.

For a nonzero residual \(r\in Y\), set \(\lambda=2\KSnorm{r}/c\). The vector \(r/\lambda\) has norm \(c/2\), so \eqref{KS:eq:inverse-closure-ball} supplies \(z\in B_X\) with
\[
 \KSnorm{Sz-r/\lambda}\le c/4.
\]
Then \(u=\lambda z\) satisfies
\[
 \KSnorm{u}\le 2\KSnorm{r}/c,\qquad
 \KSnorm{r-Su}\le\KSnorm{r}/2.
\]
If the residual is zero, choose \(u=0\). Start from \(r_0=y\), repeat this construction, and define \(r_{j+1}=r_j-Su_j\). Induction yields
\[
 \KSnorm{r_j}\le2^{-j}\KSnorm{y},\qquad
 \sum_{j\ge0}\KSnorm{u_j}\le\frac{4}{c}\KSnorm{y}.
\]
The series of preimages converges in \(X\) by completeness, to a vector \(x\) with the displayed norm bound. Continuity of \(S\), together with the finite telescoping equation
\[
 S\left(\sum_{j=0}^{N-1}u_j\right)=y-r_N,
\]
proves \(Sx=y\). Thus we may take \(C=4/c\).

If \(S\) is bijective, the preimage is unique. Uniqueness makes the inverse additive and homogeneous, and the bound just obtained makes it continuous. In the merely surjective case the construction does not assert a linear choice of preimage. In particular it does not assert that every Banach-space surjection has a bounded linear right inverse.
\end{proof}

For example, in Theorem~\ref{KS:KS43} the map to be inverted has source \((\ker D)^\perp\) and target \(\KSran D\). Both are closed subspaces of the stated Banach spaces. The inverse lemma applies only after the proof has established those closedness and bijectivity facts. It would not apply with an unclosed range equipped with its inherited norm.

\subsection{Neumann inversion with both inverse identities}
\begin{KSlemma}[A strict perturbation radius]\label{KS:lem:neumann-tool}
Let \(W\) be a real Banach space and \(T:W\to W\) bounded linear, with \(\KSnorm{T}\le\theta<1\) and \(\theta\ge0\). Then
\[
 Z=\sum_{j=0}^{\infty}T^j
\]
exists in operator norm and satisfies
\[
 (I-T)Z=Z(I-T)=I,\qquad \KSnorm{Z}\le\frac1{1-\theta}.
\]
\end{KSlemma}
\begin{proof}
First recall why the space of bounded operators on \(W\) is complete. An operator-norm Cauchy sequence \(A_n\) makes \(A_nx\) Cauchy for each \(x\), and hence defines \(Ax=\lim_nA_nx\). Pointwise limits preserve linearity. The sequence of operator norms is bounded, so the pointwise limit is a bounded operator. For a fixed sufficiently large \(n\), pass \(m\to\infty\) in \(\KSnorm{(A_n-A_m)x}\le\varepsilon\KSnorm{x}\) to obtain \(\KSnorm{(A_n-A)x}\le\varepsilon\KSnorm{x}\). Taking the supremum over the unit ball proves operator-norm convergence. Completeness of \(W\) is the step that produced each pointwise limit.

Submultiplicativity gives \(\KSnorm{T^j}\le\theta^j\), with the convention \(T^0=I\). The finite partial sums are therefore operator-norm Cauchy, since for \(M>N\),
\[
 \left\|\sum_{j=N+1}^{M}T^j\right\|
 \le\sum_{j=N+1}^{M}\theta^j\le\frac{\theta^{N+1}}{1-\theta}.
\]
Their limit \(Z\) has norm at most \(\sum_{j\ge0}\theta^j=1/(1-\theta)\). The operator identities for finite sums are
\[
 (I-T)\sum_{j=0}^{N}T^j=I-T^{N+1},\qquad
 \left(\sum_{j=0}^{N}T^j\right)(I-T)=I-T^{N+1}.
\]
Multiplication by a fixed bounded operator is continuous in operator norm, and \(\KSnorm{T^{N+1}}\le\theta^{N+1}\to0\). Passing to the limit proves both inverse identities. When \(W\) is the zero space the same statements hold, with its unique identity map. When \(\theta=1\), the summable tail and decay used here disappear, so the proof supplies no endpoint conclusion.
\end{proof}

\subsection{The dependency order}
The analytic arguments can now be read without a circular invocation of their conclusions. Projection gives the orthogonal decomposition and Riesz representation. The latter gives adjoints and the relation between range and adjoint kernel. Baire's principle independently gives the Banach inverse theorem. Those ingredients yield closed-range estimates and the positive regularized normal inverse. Neumann's series gives the separate small-perturbation inverse and its explicit radius.

The weak-subsequence proof in Section~\ref{KS:sec:bounded-realization} uses projection onto the closed separable span generated by the sequence, followed by a coordinate diagonal argument. It does not assume that the whole ambient Hilbert space is separable. The budget-attainment and regularization-limit proofs use that sequential result. Strong convergence is proved later from the minimum-norm condition and the norm identities, not inferred from weak convergence alone.



% BEGIN retained input sources/monographs/komorebi-space/source/chapters/08_references.tex

\chap{Background and references}\label{KS:app:references}
The references identify standard sources for limits, descent, Hilbert projection, bounded inverse theorems, and regularization. Appendix~\ref{KS:app:proof-tools} gives the analytic arguments used repeatedly in the text. The accompanying Lean source and its pinned mathematical-library revision determine the formal interfaces used in this edition. Online library documentation is a navigation aid and may describe a different revision. The consultation dates below are retained from the recorded source checks.


\chapter{GA: additional mathematical material}
\setcounter{section}{0}\renewcommand{\thesection}{GA.\Alph{section}}\renewcommand{\theHsection}{GA.appendix.\arabic{section}}

\section{Statement map and source relationships}\label{GA:app:claims}
\begin{small}\begin{longtable}{@{}P{.13\linewidth}P{.44\linewidth}P{.35\linewidth}@{}}\toprule
Result & Mathematical role & Status and boundary\\\midrule\endhead
\ref{GA:prop:law} & Law-level and support-level identification. & Elementary factorization, proved here. No sample-rate claim.\\
\ref{GA:thm:grounding} & One evidence-initialized exact predictive realization. & Classical quotient argument with explicit initialization. Finite declared models.\\
\ref{GA:thm:trichotomy} & Reinterpretation, acquisition, impossibility, and nonadaptive cost. & Exact for total nondestructive finite tests. Destructive tests excluded.\\
\ref{GA:thm:correction} & Coarsest memory preserving every admitted correction. & A powerset behavioral quotient. Evidence family is intersection closed.\\
\ref{GA:cor:occupancy} & Exact memory count for partition evidence. & All candidate subsets and exactly the atom-union evidence algebra. Logical states, not byte-optimal code.\\
\ref{GA:prop:testing} & Finite-record error and information processing. & Classical two-point bound, proved here. Fixed actual protocol.\\
\ref{GA:thm:radius} & Minimax coordinate ambiguity and a coverage certificate. & Unconstrained real-vector center. Coverage and model accuracy are premises.\\
\ref{GA:prop:onestep} & Linear response identifiable exactly on the calibration span. & All linear extensions admissible. Exact inputs and outputs.\\
\ref{GA:thm:core} & Maximal domain of identifiable full-state continuation. & Finite-dimensional common calibration span and fixed extension families.\\
\ref{GA:prop:approx} & Error propagation from a nominal invariant domain. & Local discrepancy plus a separate ambient growth bound.\\
\ref{GA:thm:likelihoodmemory} & Minimal likelihood-profile memory for target correction. & Fixed finite likelihoods, fresh conditional observations, all priors, and report probabilities retained.\\
\ref{GA:cor:memorymigration} & Exact reconstruction after memory-contract change. & Same finite world set and every prior. Restricted prior families can differ.\\
\ref{GA:prop:noenlargement} & Nonzero outputs do not enlarge unrestricted calibration. & All independent linear extensions admitted, fixed during an execution.\\
\ref{GA:thm:pairtest} & Finite fixed-model output agreement and a witness word. & Finite specified matrix family and all operation words.\\
\ref{GA:thm:commonquotient} & Common observable quotient from propagated discrepancy. & Agreement on the whole space, not only a special initial set.\\
\ref{GA:thm:linearpredictor} & Evidence initialization and minimum linear state dimension. & Allowed initial subspace, linear record, linear predictor class.\\
\ref{GA:prop:roundbayes} & Quantized probability fields and one-step Bayes perturbation. & A normalized finite simplex and positive likelihood floor. Retained fields, not total workspace.\\
\ref{GA:thm:finiteprecision} & Finite-horizon correction under rounding and profile error. & Static worlds, common observed history, valid likelihoods, and a specified denominator.\\
\ref{GA:thm:blockleakage} & Visible discrepancy with hidden feedback. & Common norm bounds and initial allowances. Infinite sum needs the strict comparison condition.\\
\ref{GA:thm:jointrisk} & One posterior decision from memory and response bounds. & Same cells, response, history, and bounded Lipschitz loss. No external coverage inferred.\\\bottomrule
\end{longtable}\end{small}


\section{Small exact illustrations and reproduction}\label{GA:app:checks}
Three small exact suites check the declared constructions with integers, sets, rational arithmetic, and $\mathbb F_2$. The retained observable suite compares paired row-space closure with direct word enumeration, tests common-quotient and initialization criteria, and compares likelihood-profile updates with raw-world Bayes updates. It includes explicit switching, hidden-return, zero-probability, and memory-migration controls. The two retained suites are unchanged. The first-edition suite additionally checks quantized probability records, matched adaptive updates, block-error recurrences, and the joint sensor decision. These are bounded implementation checks, not universal proofs or agent experiments.

Run \texttt{python code/run\_checks.py}. The JSON records state each domain and result. Nine Python-generated analytical figures use floating point only for drawing. The supplied Bash build uses their vector assets. Review notes record the actual proof, source, rendering, and clean-rebuild scope. No independent review, complete novelty audit, or proof-assistant certificate is claimed.



\chapter{MU: additional mathematical material}
\setcounter{section}{0}\renewcommand{\thesection}{MU.\Alph{section}}\renewcommand{\theHsection}{MU.appendix.\arabic{section}}

\section{A complete trace of the instrument}\label{MU:app:trace}
Start from $(a,b,c)=(1,3,2)$. The first display is four. The base record is $(s,b)=(4,3)$. Apply $H$, then $W$, then the newly admitted $J$.

\begin{center}
\begin{tabular}{@{}lcccl@{}}\toprule
Event & Complete state & Display & Base record & Status\\\midrule
Initial & $(1,3,2)$ & $4$ & $(4,3)$ & $H/W$ contract adequate\\
After $H$ & $(4,3,2)$ & $2$ & $(2,3)$ & Exact base update\\
After $W$ & $(3,4,2)$ & $2$ & $(2,4)$ & Exact base update\\
Before predicting $J$ & $(3,4,2)$ & $2$ & $(2,4)$ & Missing $c$ distinction\\
After $J$ & $(0,4,2)$ & $4$ & $(4,4)$ & New reading acquired\\\bottomrule
\end{tabular}
\end{center}

The last base record does not by itself retain all newly relevant information. The previous display was two and the new display is four, so $c=4-2=2$. The post-operation first register is $4-4=0$. The enlarged display-centered record is therefore $(s,b,c)=(4,4,2)$.

A subsequent $H$ gives $(s,b,c)=(3,4,2)$ and display three. A subsequent $W$ leaves the display at three but changes the second coordinate to $3-4=4$ modulo five. A subsequent $J$ changes the display to zero. Each assertion follows from the enlarged update equations, not from extrapolating a list of displayed numbers.

\subsection{What was actually learned}
The original prediction method already encoded the effect of $H$ and $W$. The $J$ observation supplied a new state distinction. It did not establish the $J$ transition law from scratch, because that law was part of the expanded analytical contract. Learning that law would require a candidate family and experiments capable of distinguishing it from alternatives.

This separation is important. Acquiring a latent value under a known rule, inferring an unknown rule, and transferring a known rule into a new notation are different research achievements. A single successful trace should not be counted as all three.

\subsection{A correction record for the trace}
The record should retain the old scope, the new operation, the failed sufficiency condition, the acquired observation, and the exact post-intervention state reconstruction. It should also preserve the unchanged base equations. A concise corrected explanation would say that the third coordinate is irrelevant to $H/W$ continuations but relevant to any contract admitting $J$.

Nothing in that explanation establishes a hidden psychological state of an agent. It establishes the content and limits of a reusable mathematical model and gives a recipient an explicit way to use and revise it.

\subsection{The same trace with an uncertain readout}
In the preceding trace, the state immediately before $J$ is $(3,4,2)$ and its exact display is two. Suppose the reasoner retains only the base record $(2,4)$ and models $c$ as uniform. After $J$, the physical state is $(0,4,2)$. Under the noisy sensor of Section~\ref{MU:sec:bayesian}, a reported display of four gives the shifted report two. The resulting belief is over states $(c-2,4,c)$, not over unchanged pre-intervention states.

\begin{center}
\begin{tabular}{@{}lccc@{}}\toprule
Independent reports retained & One & Two & Three\\\midrule
Full shifted report $2$ & $3/4$ & $36/37$ & $432/433$\\
Only even parity & $7/24$ & $98/303$ & $1372/4143$\\\bottomrule
\end{tabular}
\end{center}
The entries are posterior mass on $c=2$ along the named histories. Their probabilities of occurrence are not the entries in the table. Every other candidate remains possible after finitely many noisy reports. Duplicate copies of one reading count as one observation. The executable demonstration reproduces the fractions, checks the harmless ancillary flag, and distinguishes the full-state target from the recipient-relative communication target.

\section{Companion verification and evidence boundaries}\label{MU:app:verification}
The companion program \MUcode{code/finite\_checks.py} uses integers and exact rational arithmetic. It checks the register construction, partition refinement against a separately coded pair-search procedure, finite transcript discrimination, a stochastic quotient, transport-error recurrences, explanation capacity, informative model filtering, and correction transport. The program enumerates small domains stated separately for each check group. It does not train a model or test an empirical claim about machine cognition.

The finite quotient checks compare the refinement procedure with direct search for a distinguishing word for every small machine in the stated domain. The register checks compare microscopic updates with independently written reduced coordinate updates. The correction checks enumerate maps, candidate sets, and possible evidence, including evidence inconsistent with the current set. Negative controls deliberately use a wrong swap update, a collapsed evidence channel, and a delayed disagreement beyond a fixed test horizon.

The expanded program \MUcode{code/expanded\_checks.py} adds rational Bayesian updates, likelihood factorization, posterior-mixture identities, rare-evidence bounds, recipient decision margins, adaptive history checks, graph coloring against direct decodability, and acquisition values against full policy enumeration. The conditional-information identity is checked using exact rational coefficients of logarithms of primes rather than floating-point cancellation. Graph and policy enumerations are exhaustive only over the small domains recorded per group.

The reference module \MUcode{code/understanding.py} exposes the finite constructions with input validation. Its acquisition routine returns a policy witness or an explicit infeasibility result. A timeout or malformed input is not interpreted as mathematical infeasibility. The demonstration \MUcode{code/demo\_expanded.py} reproduces the noisy instrument and the complementary-background example.

The first-edition module \MUcode{code/resource\_contracts.py} implements fixed-decoder covers and nondominated compiled-tree profiles. Its witness checker follows every actual initial state independently of the optimizer. The suite \MUcode{code/edition\_checks.py} compares the optimizer with separately written full observation-history tree enumeration on the declared small domains, tests invalid witnesses, and checks all instrument states with the two frozen recipient programs. The costed demonstration is \MUcode{code/demo\_first\_edition.py}. Source-file sizes and hashes are recorded without calling them minimal description lengths.

The script's result is evidence about the supplied code and finite constructions. The prose arguments supply the general conditional results. A clean compilation, an exact finite check, an independent proof review, and a formal-kernel build remain different kinds of evidence. This edition claims the first two, within the recorded review scope, and does not claim the latter two.

The source package records the actual test output, source hashes, build diagnostics, reference-review scope, and known limitations. Diagrams are generated from mathematical constructions. No synthetic image is offered as an empirical observation. The PDF's version number identifies a consolidated first edition, not a completed priority audit or an accepted publication.

\section{Claim map}\label{MU:app:claims}
\begin{center}\small
\begin{tabularx}{\textwidth}{@{}P{1.35in}P{1.55in}X@{}}\toprule
Statement & Main assumption & Conclusion and boundary\\\midrule
Theorem~\ref{MU:thm:quotient} & Finite deterministic operations and all declared continuations. & Minimal behavioral quotient. Does not minimize acquisition or runtime.\\
Proposition~\ref{MU:prop:expansion} & Expanded contract on the same complete state space. & Necessary supplementary capacity. Sharp reuse construction has an explicit refinement condition.\\
Proposition~\ref{MU:prop:identification} & Fixed complete transcript laws and equal priors. & Optimal binary identification bound. New evidence changes the laws.\\
Proposition~\ref{MU:prop:lumpability} & Equal projected rows on each fibre. & Exact stochastic history preservation for reduced-history policies.\\
Theorem~\ref{MU:thm:transport} & Matched operations and query decoders. & Exact forward transfer. Does not grant reverse recovery or target coverage.\\
Proposition~\ref{MU:prop:error} & Uniform residuals and Lipschitz target dynamics. & Finite-horizon error budget for matched operations.\\
Proposition~\ref{MU:prop:explanation} & One fixed recipient background and fixed-length messages. & Necessary communication capacity, not efficient decoding.\\
Proposition~\ref{MU:prop:queries} & Exact realizable candidate family and separating experiments. & Finite informative elimination and decision-tree lower bound.\\
Theorem~\ref{MU:thm:correction} & A commuting evidence square, including contradictory observations. & Filtering inclusion. Uniform equality exactly requires reflection of attainable outcomes.\\\bottomrule
\end{tabularx}\end{center}

\subsection{Expanded correction and explanation results}
\begin{center}\small
\begin{tabularx}{\textwidth}{@{}P{1.35in}P{1.55in}X@{}}\toprule
Statement & Main assumption & Conclusion and boundary\\\midrule
Theorem~\ref{MU:thm:audience} & Finite source access, recipient backgrounds, and designable decoders. & Exact common alphabet via graph coloring. Does not establish decoder efficiency.\\
Theorem~\ref{MU:thm:acquisition} & Finite deterministic experiments, explicit disturbance and admissibility, bounded horizon. & Minimum worst-case acquisition cost and an attaining policy.\\
Theorem~\ref{MU:thm:bayes} & Finite likelihoods and one reporting channel independent of the model. & All-prior posterior preservation exactly when conditional factorization holds.\\
Proposition~\ref{MU:prop:posterior-information} & The same joint experiment and a stated prior. & Average posterior divergence equals conditional information, not a worst-case bound.\\
Proposition~\ref{MU:prop:conditioning} & Two joint laws and an event positive under both. & Conditional variation bound with an evidence denominator.\\
Proposition~\ref{MU:prop:decision} & One posterior chain, finite actions, and bounded common losses. & Recipient regret and strict action-gap certificate.\\
Theorem~\ref{MU:thm:adaptive-bayes} & Fresh conditionally independent observations and a reported-history selection policy. & Exact posterior preservation along adaptive inquiry and bounded public stopping.\\
Theorem~\ref{MU:thm:joint-acquisition} & Finite deterministic trees, costs, and receiver contracts. & Exact joint acquisition and common-message feasibility, not efficient general synthesis.\\\bottomrule
\end{tabularx}\end{center}

\subsection{Fixed procedures and first-edition resource results}
\begin{center}\small
\begin{tabularx}{\textwidth}{@{}P{1.35in}P{1.55in}X@{}}\toprule
Statement & Main assumption & Conclusion and boundary\\\midrule
Proposition~\ref{MU:prop:advice} & Fixed interpreter with charged task-dependent fixed-length program and artifact. & Combined information lower bound. Not an individual-function runtime lower bound.\\
Theorem~\ref{MU:thm:fixed-decoder} & Finite source access, artifact catalogue, and fixed bounded recipient procedures. & Exact artifact intersections and minimum covering codebook. Raw artifact size remains charged.\\
Theorem~\ref{MU:thm:resource-synthesis} & Finite augmented state, bounded experiments, frozen recipients, vector charges, explicit controller trees. & Exact resource frontier and common workflow witnesses. Not optimal unrestricted program synthesis.\\\bottomrule
\end{tabularx}\end{center}

\subsection{How the certificates depend on one another}
The quotient and transfer results concern the declared operations. Their conclusions do not specify how an absent register is measured. The acquisition result adds that physical and informational route, and the correction results determine what an observation can eliminate or how it changes a belief. Neither step by itself guarantees that another reasoner can use the result.

The designable-recipient graph in Theorem~\ref{MU:thm:audience} characterizes a communication alphabet after a decoder may be selected. Theorem~\ref{MU:thm:fixed-decoder} instead holds the recipient's procedure fixed. Its intersections can be empty even when the designable graph has a coloring. This is why a claim about one implemented recipient needs the stronger contract.

The final recurrence incorporates those same fixed terminal procedures into one experiment tree. Its witness must fit both online charges and the storage of every installed branch. A policy attaining one minimum and another policy attaining a second minimum are not interchangeable evidence for one simultaneous budget. The costed instrument supplies a common witness, while the four-origin counterexample supplies an explicit obstruction.

The probabilistic results are not silently inserted into that exact tree model. A noisy report may support a posterior and a bounded decision loss without identifying a unique state. Approximate posterior transfer therefore belongs to its stated likelihood, evidence, and recipient-loss assumptions. A zero-error state certificate requires its own stronger evidence. This separation is part of the first edition's scope, not an omission disguised by a shared diagram.


\phantomsection

\chapter{SII: additional mathematical material}
\setcounter{section}{0}\renewcommand{\thesection}{SII.\Alph{section}}\renewcommand{\theHsection}{SII.appendix.\arabic{section}}

\section{A complete diagnostic trace}\label{SII:app:trace}
Take $n=8$ and hidden defect $j=5$, whose binary digits are $101$. One optimal tree uses the three queries
\[
A_0=\{1,3,5,7\},\qquad A_1=\{2,3,6,7\},\qquad A_2=\{4,5,6,7\}.
\]
The precise order may differ among equally optimal compiled trees. For this order, the supports are
\[
\{0,\ldots,7\}\xrightarrow{1}\{1,3,5,7\}
\xrightarrow{0}\{1,5\}\xrightarrow{1}\{5\}.
\]
The terminal repair is $r_5$. The path uses three diagnostic calls and one repair call. The complete tree stores all eight repair leaves and seven query nodes, not merely this path.

Equality-first testing asks about cells zero through five before it obtains a positive report. It uses six diagnostic calls on this particular case, then one repair. Under the three-diagnostic cap it abstains after three negative reports. The cap is fixed before the episode, not shortened for the baseline after observing which cell is defective.

For a target relabeling $\sigma(i)=3i+1\pmod8$, the defect has target label $\sigma(5)=0$. Each query becomes $\sigma(A_r)$ and the terminal repair becomes $r'_{\sigma(5)}$. Reports and tree shape agree with the source, but the label adapter is part of the transported interface. Applying the old numeric repair label five directly in the target would be incorrect. This is why syntactic resemblance does not replace a semantic map.

\subsection{A repair problem that should stop before full identification}
In the four-case example, the support $\{0,1,2\}$ admits repair $a$ immediately. A compiler that insists on identifying the precise case would spend unnecessary work. The successful stopping condition is the existence of one action whose success set contains the whole support. Conversely, the support $\{2,3\}$ admits no common repair in that example. A leaf chosen there without another observation is invalid even though either repair succeeds in one of its cases.

\section{Reference computations and evidence scope}\label{SII:app:checks}
The package implements the finite experiment model, the two-resource recurrence, tree reconstruction, and a separate replay checker. It also implements stable partition refinement and finite recursive-progress synthesis. Exact test groups compare selected constructions with separately coded enumerations on small domains. Integer and rational arithmetic are used for mathematical acceptance checks. Plotting converts exact analytical quantities to display coordinates only.

The verification program records its exact domains, input and output hashes, command, interpreter, duration, and outcomes. The source audit inventories formal statements and checks labels, citations, and their proof environments. The build uses LaTeX and the supplied vector figures. A clean rebuild is compared by extracted PDF text, not claimed to be byte-identical across environments.

These checks are authored within this project. Alternative algorithms provide implementation cross-checks, not independent peer review. No test count establishes a universal claim outside its implemented domain. The proofs in the manuscript provide the conditional general arguments. No Lean certificate, trained agent, learned semantic bridge, or automated research run is included.

\subsection{A precise reading of the finite demonstration}
The observation and success tables are supplied analytical semantics. The compiler selects a tree without access to the next episode's hidden index. The checker separately receives that index when exhaustively verifying all possible episodes. That division prevents a policy from consulting hidden state during replay. It does not make the supplied model an empirically learned object.

The main generated files are the compiled cell-tree certificate, a report of source and transported replays, the cost frontiers, and the exact check results. All can be regenerated locally. The README records the commands and the review notes record any failures or limitations encountered during preparation.

\section{Claim map and the scope of the first edition}\label{SII:app:claimmap}
\begin{small}\begin{longtable}{@{}p{.30\textwidth}p{.62\textwidth}@{}}\toprule
Claim & What the present manuscript supplies\\\midrule\endhead
Feedback can support a repair & A transcript-discrimination bound and a common-repair stopping condition. No automatic causal diagnosis is inferred.\\
A representation supports continuation & A least stable refinement and explicit operation equations. Learning the representation is additional.\\
A repair can be found within a budget & An exact finite work/storage recurrence and a checkable tree. It requires a supplied finite experiment model.\\
A discovery procedure transfers & A correspondence preserving complete investigation semantics and costs. Discovering that correspondence is additional.\\
A pruning rule is learned & A finite grammar synthesis, a witnessed feature conflict, and a separate scope proof. The teacher and grammar are supplied.\\
A compiler skips work safely & Streaming root rejection preserves the controller and avoids counted rows. Predicate work remains charged.\\
A candidate procedure is better & A selection bound under uniform assessment error and a protected productivity comparison. The error premise must be established.\\
A lineage preserves credible gains & Conditional error spending and an explicit distinction between development, selection, and audit evidence.\\
A lineage can improve repeatedly & An exact finite common-policy criterion for marked progress. Marks require separate valid evidence.\\
The process accelerates & Not established. A fixed-score ceiling and cost accounting expose additional requirements.\\
The system is generally intelligent & Not established. Generality and learned discovery remain research targets across unfamiliar structures.\\The synthesis procedure improves & A checked equation library removes redundant candidate evaluations while preserving teacher proposals. Preparation costs remain separate.\\
One repair covers several cases & A feasible repair-weight certificate replaces singleton counting. Its validity is rechecked after a repair-language change.\\
A composed gain is affordable & The actual acquisition, compilation, checking, and preparation costs must fit one common resource model.\\
Old certificates survive a change & Only dependency-complete unchanged obligations can be reused without new evidence. Affected descendants require re-establishment.\\
A gain is independent of prices & Componentwise event dominance, or a strictly positive margin on the declared finite price polytope, supplies the precise criterion.\\
\bottomrule
\end{longtable}\end{small}




\chapter{MS: additional mathematical material}
\setcounter{section}{0}\renewcommand{\thesection}{MS.\Alph{section}}\renewcommand{\theHsection}{MS.appendix.\arabic{section}}

\section{The commons as an integrated preservation contract}\label{MS:app:allocation}
The running example in Section~\ref{MS:sec:commons} uses one institutional realization throughout. Its claims can be separated without losing their connection.
\begin{longtable}{@{}P{.25\linewidth}P{.69\linewidth}@{}}
\toprule
\textbf{Question} & \textbf{Object and obligation}\\\midrule
\endfirsthead
\toprule\textbf{Question}&\textbf{Object and obligation}\\\midrule\endhead
Can an institution be adopted? & The coordination-game transition barrier is distinct from the fact that $R$ and $U$ have executable implementations. A profitable joint commitment requires an admitted commitment operation.\\[.5em]
What state must persist? & The allocation rule alone is insufficient. Governance preservation also retains the enabled amendment actions and labels for active credentials, service debt, and resource tokens.\\[.5em]
Can a due debt be erased? & The obligation monitor distinguishes fulfillment, authorized discharge, and deletion. Amendments are required to preserve pending commitments.\\[.5em]
Who can block change? & Two registered offices block a three-of-four rule under independent approval rights. A hidden administrator changes this calculation.\\[.5em]
Can every request be safely served? & Resource conservation can hold even when private information or conflicting deadlines prevent timely service. Viability with full information is not decentralized implementability.\\[.5em]
What crosses a replacement? & An authorized office handover preserves debt and credential semantics. Copying a model does not duplicate an office's budget or approval right.\\[.5em]
Can two jurisdictions merge? & Their translator must preserve the agreed action semantics, permissions, obligations, resources, and evidence coordinates. Nonzero residuals remain visible until an authorized repair is found.\\[.5em]
How long does the organization last? & Lifetime is first exit from a declared operational contract under a specified transition law. It is neither occupant longevity nor mere archive retention.\\
\bottomrule
\end{longtable}
The common event mechanism and constitutional envelope are supplied. Within that envelope, offices can choose rules, incur obligations, revise decisions, and replace occupants. This demonstrates how the theory's objects fit together. It does not prove that the envelope itself arose endogenously or that these rules are uniquely desirable.

\subsection{The additional boundary, language, and knowledge contract}
An admission changes the membership configuration, not merely the occupant list. A message is interpreted using its type, issuer, and epoch. A resource approval creates the specified ledger event only after its admissibility check. A debtor replacement transports the obligation monitor but does not manufacture new approval rights. A notification history records which office has received which adoption record; an available ledger entry is not automatically common knowledge.

For concurrent amendments, the contract states whether ordering is significant, whether a joinable normal form is required, and which external effects remain visible. For federation, the role-label transformations must be checked around ordered loops. A compatible allocation on one invariant subset does not certify synchronization of every role. For inquiry, the commons records unresolved scheduling questions independently of which questions it chooses to fund. For founding authority, it preserves the resource owner's stipulated grant separately from its own subsequent votes.

Together these clauses specify what it would mean to preserve this analytical institution across translations and replacements. They do not claim that every useful society needs this exact architecture or that the four-office example exhibits autonomous institutional formation.

\subsection{Distributed realization, enforcement, and adaptive staffing}
In the atomic reference commons, a unique slot token prevents two successful grants of the same slot. In its replicated variant, two local approvals can each look valid while their merged grant set violates exclusivity. The institution must therefore distinguish a provisional reservation from a committed grant and identify the primitive that prevents double execution. A convergent archive is not that primitive by itself.

Accurate urgency reporting is a different problem. A committed monitor can discourage a specified profitable misreport only within the discrimination and penalty budget of Section~\ref{MS:sec:enforcement}. Its false-positive limit, evidence access, and remedy authority remain separate. When an averaged model supplies the observation laws, the margin must include its approximation error. A fixed review schedule and a data-selected review schedule cannot share the same mixing certificate without further work.

The commons also distinguishes an urgency vote from the signal supporting it. Decision-only publication can reproduce a mistaken cascade; preservation of independent evidence changes the information available to later reviewers, not the legal effect of an already committed vote. Finally, creating a specialist verifier changes the capability graph and the distribution of authority. A staffing certificate covers all protected roles jointly, while the handover procedure must still deliver credentials, pending debts, and evidence before service deadlines. These clauses supplement the reference semantics rather than claim that the society has generated its own infrastructure without assumptions.

\subsection{The joint continuation contract}
The bounded audit-and-service variant in Section~\ref{MS:sec:joint-synthesis} combines three protections that cannot be checked in isolation: the common monitor must deter every listed deviation, its false-positive rate must meet the cap, and its resource use must leave enough capacity for all due work. The current public possibility set must also admit the same protocol in every compatible complete state. A continuation certificate records this common choice and all possible successors. A progress rank distinguishes eventual certified service from indefinite safe waiting.

An amendment is acceptable under this particular continuation contract only when its current authority is valid and its resulting states retain certified continuations. A transport certificate additionally checks public observations, enabled profiles, safe and goal labels, and successor sets. Replacing an office, changing an audit channel, or adding a slot changes specified parts of that contract. Such a change may repair an obstruction, but does not retroactively fulfill an earlier unmet promise.
\FloatBarrier

\subsection{The accountability extension}
Section~\ref{MS:sec:accountable-commons} deliberately changes the audit catalogue to satisfy its pairwise privacy constraint. It separately includes an outside claim, an admitted appeal, a genuine-execution fee, and a funded final record with revocation. The resulting eight-unit settlement requirement must not be confused with the five-unit service-and-audit variant. The supplied reference trace checks one declared five-profile settlement, while the stopping theorem characterizes an entire finite ending-preserving running kernel. Both are model-relative results, not certificates for an unspecified deployment.

\section{A finite algebraic example of reusable state}\label{MS:sec:worked}
The following analytical example makes the preceding distinctions concrete. Four components initially hold independent fair bits $(a,b,c,d)$. Two local workspaces expose
\[
p=a\oplus b,\qquad q=c\oplus d,
\]
and a global workspace retains $Z=(p,q)$. The current reported answer is $r=p\oplus q$.

No proper subset of the four raw bits determines $r$: at least one independent fair bit is missing. Each local pair determines its own workspace bit but has zero information about $r$. Together the workspaces determine $r$ exactly. This is a precise distributed-access property for the initial distribution.

\subsection{Why the current answer is not enough state}
Permit three public operations on $Z$:
\[
I(p,q)=(p,q),\qquad W(p,q)=(q,p),\qquad S(p,q)=(p\oplus q,q).
\]
The global state supports these operations and the query $r$. The operations $W$ and $S$ do not commute. For instance, from $(p,q)=(1,0)$,
\[
W(S(1,0))=(0,1),\qquad S(W(1,0))=(1,1).
\]
Their reported parity outcomes are different.

By contrast, retaining only the current answer $r$ fails the quotient criterion. The states $(0,0)$ and $(1,1)$ both report $r=0$, but after $S$ they report $0$ and $1$, respectively. The missing distinction is exactly a distinction relevant to a future operation. A compressed answer is not necessarily a compressed process.

\begin{center}
\begin{tabular}{@{}ccccc@{}}
\toprule
$p$ & $q$ & Current $r$ & State after $S$ & New $r$\\
\midrule
0&0&0&$(0,0)$&0\\
0&1&1&$(1,1)$&0\\
1&0&1&$(1,0)$&1\\
1&1&0&$(0,1)$&1\\
\bottomrule
\end{tabular}
\end{center}

The full two-bit state has four future-distinguishable possibilities: the current parity reveals $p\oplus q$, and parity after $S$ reveals $p$. Together those outputs determine $q$ as well. Thus two bits are necessary for an exact reusable state under this query family, even though one bit suffices for the immediate answer.

\subsection{A deterministic microscopic realization}
For any permitted macro-operation $T$, write $(p',q')=T(p,q)$ and define
\[
F_T(a,b,c,d)=(a,\ a\oplus p',\ c,\ c\oplus q').
\]
Then the two local XORs of $F_T(a,b,c,d)$ are exactly $(p',q')$. Consequently $\pi F_T=T\pi$ for every microstate, where $\pi(a,b,c,d)=(p,q)$.

This formula specifies a global transition for verification. A distributed implementation still needs communication to execute operations that mix $p$ and $q$. The formula alone does not prove locality, privacy against transcript observers, or the spontaneous discovery of a protocol. It shows exactly which macrostate the implementation realizes.

\subsection{State loss depends on the access contract}
Suppose the leaf bits become inaccessible after export and the only surviving pair-level records are $p$ and $q$. Erasing $q$ reduces the best prediction of $r$ to chance, $1/2$. Replacing it with a fixed bit has the same average accuracy. If the original leaves $c,d$ remain readable, however, the missing state can be reconstructed and accuracy returns to $1$.

Both outcomes are correct. They differ because the allowed information differs. A ``workspace ablation'' that silently permits re-reading the original evidence does not test the same necessity claim as one that removes every permitted carrier. A state can also be redundant yet causally used: successful repair does not imply that the removed component never mattered.

\subsection{What this example establishes}
The example establishes exact distributed initial information, a reusable two-bit macrostate, deterministic macro-closure, noncommuting update semantics, and an access-relative memory requirement. It also supplies a negative control: the one-bit answer statistic fails dynamic sufficiency. It does not contain learning, endogenous institutions, self-selected goals, or a society. Its finite state space can be enumerated exactly; the general claims rest on the arguments given above.


% Monograph layout: keep glossary together.
\section{A compact glossary of exclusions}\label{MS:app:terms}
\begin{longtable}{@{}P{.23\textwidth}P{.72\textwidth}@{}}
\toprule
\textbf{Term} & \textbf{Meaning, with excluded inference}\\
\midrule
\endfirsthead
\toprule
\textbf{Term} & \textbf{Meaning, with excluded inference}\\
\midrule
\endhead
Relational state & A useful statistic or record of joint interaction. It is not information outside the complete implementation.\\[.35em]
Dynamic closure & A well-defined reduced transition law under specified controls. It does not follow from correlated outputs.\\[.35em]
Institution & Operative norms, permissions, enforcement or revision procedures, and institutional memory. A policy document alone is insufficient.\\[.35em]
Endogenous formation & Interaction-dependent selection or revision relative to a supplied scaffold. Not creation without prerequisites.\\[.35em]
Society-like population & Persistent plural decision centers with operative institutions. Not necessarily a single collective policy.\\[.35em]
Collective agency & A useful group-level decision interface for a specified task family. It is neither necessary nor sufficient for society-like plurality.\\[.35em]
Cultural inheritance & Transfer of future-relevant methods, practices, and correction policies. Not mere repetition of predecessor statements.\\[.35em]
Organizational continuity & Preservation of declared operations and lineage under changes. Not subjective identity or automatic transfer of authority.\\[.35em]
Compatibility & Agreement of relevant representations, operations, and global realizability. Pairwise consistency may be too weak.\\[.35em]
Strategic deception & A supported strategic mechanism for inducing inaccurate beliefs. Not every wrong statement or representation mismatch.\\[.35em]
Safety certificate & A guarantee relative to formal assumptions, transitions, and invariants. Not proof of moral legitimacy or liveness.\\[.35em]
Emergence & An interaction-dependent organizational property relative to a boundary and comparator. Not a claim of extra computational power or consciousness.\\
\bottomrule
\end{longtable}
% FINAL EDITION ADDITIONS BEGIN

% BEGIN retained input sources/monographs/machine-societies/publication/statement_index.tex

\section{Theorem and proposition index}\label{MS:app:statement-index}
This index links to all 62 theorem and proposition statements. The page numbers refer to this edition. Each result retains its stated hypotheses and prose proof. A final review of the central synthesis and accountability results is described in Appendix~\ref{MS:app:release-review}.
\begingroup\small\setlength{\tabcolsep}{4pt}\renewcommand{\arraystretch}{1.15}
\begin{longtable}{@{}P{.19\linewidth}P{.69\linewidth}r@{}}
\toprule Statement & Result & Page\\\midrule\endfirsthead
\toprule Statement & Result & Page\\\midrule\endhead
\bottomrule\endfoot
\hyperref[MS:prop:det]{Proposition~\ref*{MS:prop:det}} & \hyperref[MS:prop:det]{Deterministic quotient criterion} & \pageref{MS:prop:det}\\
\hyperref[MS:prop:lump]{Proposition~\ref*{MS:prop:lump}} & \hyperref[MS:prop:lump]{Controlled lumpability} & \pageref{MS:prop:lump}\\
\hyperref[MS:prop:time]{Proposition~\ref*{MS:prop:time}} & \hyperref[MS:prop:time]{Approximate closure over time} & \pageref{MS:prop:time}\\
\hyperref[MS:prop:lift]{Proposition~\ref*{MS:prop:lift}} & \hyperref[MS:prop:lift]{Implementation dependence} & \pageref{MS:prop:lift}\\
\hyperref[MS:prop:identifiability]{Proposition~\ref*{MS:prop:identifiability}} & \hyperref[MS:prop:identifiability]{Behavior alone does not identify implementation} & \pageref{MS:prop:identifiability}\\
\hyperref[MS:prop:memory]{Proposition~\ref*{MS:prop:memory}} & \hyperref[MS:prop:memory]{Exact finite-memory lower bound} & \pageref{MS:prop:memory}\\
\hyperref[MS:prop:erasure]{Proposition~\ref*{MS:prop:erasure}} & \hyperref[MS:prop:erasure]{Exact reconstruction after erasure} & \pageref{MS:prop:erasure}\\
\hyperref[MS:prop:repair]{Proposition~\ref*{MS:prop:repair}} & \hyperref[MS:prop:repair]{Repair budget} & \pageref{MS:prop:repair}\\
\hyperref[MS:prop:composition]{Proposition~\ref*{MS:prop:composition}} & \hyperref[MS:prop:composition]{Compositional one-step closure} & \pageref{MS:prop:composition}\\
\hyperref[MS:prop:flatten]{Proposition~\ref*{MS:prop:flatten}} & \hyperref[MS:prop:flatten]{Simulation by an unrestricted controller} & \pageref{MS:prop:flatten}\\
\hyperref[MS:prop:symmetry]{Proposition~\ref*{MS:prop:symmetry}} & \hyperref[MS:prop:symmetry]{Symmetry lock for identical deterministic agents} & \pageref{MS:prop:symmetry}\\
\hyperref[MS:prop:forkweight]{Proposition~\ref*{MS:prop:forkweight}} & \hyperref[MS:prop:forkweight]{Weight-conserving refinement of controlled votes} & \pageref{MS:prop:forkweight}\\
\hyperref[MS:prop:handover]{Proposition~\ref*{MS:prop:handover}} & \hyperref[MS:prop:handover]{Intersecting witnesses exclude conflicting certificates} & \pageref{MS:prop:handover}\\
\hyperref[MS:thm:govquot]{Theorem~\ref*{MS:thm:govquot}} & \hyperref[MS:thm:govquot]{Exact governance quotient} & \pageref{MS:thm:govquot}\\
\hyperref[MS:prop:govrefine]{Proposition~\ref*{MS:prop:govrefine}} & \hyperref[MS:prop:govrefine]{Coarsest finite governance summary} & \pageref{MS:prop:govrefine}\\
\hyperref[MS:prop:languagetrace]{Proposition~\ref*{MS:prop:languagetrace}} & \hyperref[MS:prop:languagetrace]{Preservation of finite institutional conversations} & \pageref{MS:prop:languagetrace}\\
\hyperref[MS:prop:amort]{Proposition~\ref*{MS:prop:amort}} & \hyperref[MS:prop:amort]{Institutional amortization} & \pageref{MS:prop:amort}\\
\hyperref[MS:prop:veto]{Proposition~\ref*{MS:prop:veto}} & \hyperref[MS:prop:veto]{Aggregate profitability without voluntary adoption} & \pageref{MS:prop:veto}\\
\hyperref[MS:prop:potential]{Proposition~\ref*{MS:prop:potential}} & \hyperref[MS:prop:potential]{Finite termination of strict institutional improvement} & \pageref{MS:prop:potential}\\
\hyperref[MS:prop:barrier]{Proposition~\ref*{MS:prop:barrier}} & \hyperref[MS:prop:barrier]{Institutional formation barrier} & \pageref{MS:prop:barrier}\\
\hyperref[MS:prop:seed]{Proposition~\ref*{MS:prop:seed}} & \hyperref[MS:prop:seed]{Minimum profitable coordinated seed} & \pageref{MS:prop:seed}\\
\hyperref[MS:prop:deterrence-tv]{Proposition~\ref*{MS:prop:deterrence-tv}} & \hyperref[MS:prop:deterrence-tv]{The exact bounded-sanction discrimination limit} & \pageref{MS:prop:deterrence-tv}\\
\hyperref[MS:prop:enforcement-garbling]{Proposition~\ref*{MS:prop:enforcement-garbling}} & \hyperref[MS:prop:enforcement-garbling]{Garbling does not increase sanctioning capacity} & \pageref{MS:prop:enforcement-garbling}\\
\hyperref[MS:thm:joint-enforcement]{Theorem~\ref*{MS:thm:joint-enforcement}} & \hyperref[MS:thm:joint-enforcement]{Simultaneous enforcement and mixture obstruction} & \pageref{MS:thm:joint-enforcement}\\
\hyperref[MS:prop:joint-robust]{Proposition~\ref*{MS:prop:joint-robust}} & \hyperref[MS:prop:joint-robust]{A robust simultaneous enforcement certificate} & \pageref{MS:prop:joint-robust}\\
\hyperref[MS:prop:safety]{Proposition~\ref*{MS:prop:safety}} & \hyperref[MS:prop:safety]{Constitutional inductive safety} & \pageref{MS:prop:safety}\\
\hyperref[MS:prop:viability]{Proposition~\ref*{MS:prop:viability}} & \hyperref[MS:prop:viability]{Finite-horizon almost-sure institutional viability} & \pageref{MS:prop:viability}\\
\hyperref[MS:thm:confluence]{Theorem~\ref*{MS:thm:confluence}} & \hyperref[MS:thm:confluence]{Order-independent amendment normal forms} & \pageref{MS:thm:confluence}\\
\hyperref[MS:prop:merge-convergence]{Proposition~\ref*{MS:prop:merge-convergence}} & \hyperref[MS:prop:merge-convergence]{Quiescent convergence and invariant preservation} & \pageref{MS:prop:merge-convergence}\\
\hyperref[MS:prop:exclusive-service]{Proposition~\ref*{MS:prop:exclusive-service}} & \hyperref[MS:prop:exclusive-service]{Exclusive service cannot be unconditionally local} & \pageref{MS:prop:exclusive-service}\\
\hyperref[MS:thm:undecidable]{Theorem~\ref*{MS:thm:undecidable}} & \hyperref[MS:thm:undecidable]{No universal constitutional safety decider} & \pageref{MS:thm:undecidable}\\
\hyperref[MS:thm:decsafe]{Theorem~\ref*{MS:thm:decsafe}} & \hyperref[MS:thm:decsafe]{Indistinguishability obstruction to joint safety} & \pageref{MS:thm:decsafe}\\
\hyperref[MS:prop:obligationtransport]{Proposition~\ref*{MS:prop:obligationtransport}} & \hyperref[MS:prop:obligationtransport]{Obligation-preserving succession and translation} & \pageref{MS:prop:obligationtransport}\\
\hyperref[MS:prop:transport]{Proposition~\ref*{MS:prop:transport}} & \hyperref[MS:prop:transport]{Accumulation of cultural transport defects} & \pageref{MS:prop:transport}\\
\hyperref[MS:prop:culturegrowth]{Proposition~\ref*{MS:prop:culturegrowth}} & \hyperref[MS:prop:culturegrowth]{Conditional cumulative culture} & \pageref{MS:prop:culturegrowth}\\
\hyperref[MS:prop:agendaclosure]{Proposition~\ref*{MS:prop:agendaclosure}} & \hyperref[MS:prop:agendaclosure]{Least closure of a finite problem--method system} & \pageref{MS:prop:agendaclosure}\\
\hyperref[MS:prop:lifetime]{Proposition~\ref*{MS:prop:lifetime}} & \hyperref[MS:prop:lifetime]{Survival and expected institutional lifetime} & \pageref{MS:prop:lifetime}\\
\hyperref[MS:thm:block-averaging]{Theorem~\ref*{MS:thm:block-averaging}} & \hyperref[MS:thm:block-averaging]{Finite-block institutional averaging} & \pageref{MS:thm:block-averaging}\\
\hyperref[MS:prop:robust-deterrence]{Proposition~\ref*{MS:prop:robust-deterrence}} & \hyperref[MS:prop:robust-deterrence]{A robust margin for reduced-model enforcement} & \pageref{MS:prop:robust-deterrence}\\
\hyperref[MS:prop:cycle]{Proposition~\ref*{MS:prop:cycle}} & \hyperref[MS:prop:cycle]{Cycle criterion for binary institutional compatibility} & \pageref{MS:prop:cycle}\\
\hyperref[MS:thm:holonomy]{Theorem~\ref*{MS:thm:holonomy}} & \hyperref[MS:thm:holonomy]{Exact frame synchronization} & \pageref{MS:thm:holonomy}\\
\hyperref[MS:prop:holonomysection]{Proposition~\ref*{MS:prop:holonomysection}} & \hyperref[MS:prop:holonomysection]{Compatible states and holonomy fixed points} & \pageref{MS:prop:holonomysection}\\
\hyperref[MS:prop:holonomybound]{Proposition~\ref*{MS:prop:holonomybound}} & \hyperref[MS:prop:holonomybound]{Loop defect bounds unavoidable edge error} & \pageref{MS:prop:holonomybound}\\
\hyperref[MS:thm:residual]{Theorem~\ref*{MS:thm:residual}} & \hyperref[MS:thm:residual]{Composition of typed institutional residuals} & \pageref{MS:thm:residual}\\
\hyperref[MS:prop:cyclerepair]{Proposition~\ref*{MS:prop:cyclerepair}} & \hyperref[MS:prop:cyclerepair]{Cycle residual and minimum authorized repair} & \pageref{MS:prop:cyclerepair}\\
\hyperref[MS:prop:cascade-lock]{Proposition~\ref*{MS:prop:cascade-lock}} & \hyperref[MS:prop:cascade-lock]{Almost-sure locking with a positive error probability} & \pageref{MS:prop:cascade-lock}\\
\hyperref[MS:prop:signal-majority]{Proposition~\ref*{MS:prop:signal-majority}} & \hyperref[MS:prop:signal-majority]{Conditional independence permits continuing accuracy gains} & \pageref{MS:prop:signal-majority}\\
\hyperref[MS:prop:commonreach]{Proposition~\ref*{MS:prop:commonreach}} & \hyperref[MS:prop:commonreach]{Common knowledge as a reachability condition} & \pageref{MS:prop:commonreach}\\
\hyperref[MS:prop:epistemictransport]{Proposition~\ref*{MS:prop:epistemictransport}} & \hyperref[MS:prop:epistemictransport]{Epistemic preservation and reflection} & \pageref{MS:prop:epistemictransport}\\
\hyperref[MS:prop:blocker]{Proposition~\ref*{MS:prop:blocker}} & \hyperref[MS:prop:blocker]{Blocking coalitions as hitting sets} & \pageref{MS:prop:blocker}\\
\hyperref[MS:prop:founding]{Proposition~\ref*{MS:prop:founding}} & \hyperref[MS:prop:founding]{When a founding fact is determined by internal records} & \pageref{MS:prop:founding}\\
\hyperref[MS:thm:robust-staffing]{Theorem~\ref*{MS:thm:robust-staffing}} & \hyperref[MS:thm:robust-staffing]{Authorized staffing after any bounded agent erasure} & \pageref{MS:thm:robust-staffing}\\
\hyperref[MS:thm:joint-viability]{Theorem~\ref*{MS:thm:joint-viability}} & \hyperref[MS:thm:joint-viability]{Joint continuation under a common observation history} & \pageref{MS:thm:joint-viability}\\
\hyperref[MS:prop:joint-progress]{Proposition~\ref*{MS:prop:joint-progress}} & \hyperref[MS:prop:joint-progress]{Safe service with a finite progress rank} & \pageref{MS:prop:joint-progress}\\
\hyperref[MS:prop:continuation-certificate]{Proposition~\ref*{MS:prop:continuation-certificate}} & \hyperref[MS:prop:continuation-certificate]{Soundness of a finite continuation certificate} & \pageref{MS:prop:continuation-certificate}\\
\hyperref[MS:thm:synthesis-transport]{Theorem~\ref*{MS:thm:synthesis-transport}} & \hyperref[MS:thm:synthesis-transport]{Governance abstraction preserves joint realizability} & \pageref{MS:thm:synthesis-transport}\\
\hyperref[MS:prop:strategic-inspector]{Proposition~\ref*{MS:prop:strategic-inspector}} & \hyperref[MS:prop:strategic-inspector]{Inspection without commitment} & \pageref{MS:prop:strategic-inspector}\\
\hyperref[MS:thm:privacy-envelope]{Theorem~\ref*{MS:thm:privacy-envelope}} & \hyperref[MS:thm:privacy-envelope]{Sharp pairwise privacy envelope for a bounded sanction} & \pageref{MS:thm:privacy-envelope}\\
\hyperref[MS:prop:privacy-transcript]{Proposition~\ref*{MS:prop:privacy-transcript}} & \hyperref[MS:prop:privacy-transcript]{Conditional disclosure accounting} & \pageref{MS:prop:privacy-transcript}\\
\hyperref[MS:prop:appeal-capacity]{Proposition~\ref*{MS:prop:appeal-capacity}} & \hyperref[MS:prop:appeal-capacity]{Capacity for a finite appeal contract} & \pageref{MS:prop:appeal-capacity}\\
\hyperref[MS:thm:stoppable-operation]{Theorem~\ref*{MS:thm:stoppable-operation}} & \hyperref[MS:thm:stoppable-operation]{Continued operation with an available authorized ending} & \pageref{MS:thm:stoppable-operation}\\
\hyperref[MS:thm:guarded-certificates]{Theorem~\ref*{MS:thm:guarded-certificates}} & \hyperref[MS:thm:guarded-certificates]{Guarded certificate use and authorized handover} & \pageref{MS:thm:guarded-certificates}\\
\end{longtable}
\endgroup



% BEGIN retained input sources/monographs/machine-societies/publication/release_review.tex

\section{Final edition and verification scope}\label{MS:app:release-review}
\noindent\textbf{Final monograph edition, 15 October 2026.} This edition preserves the 62 theorem and proposition statements and their prose proofs from text edition 1.0.0. It adds the linked index, completes the presentation review, and records a targeted mathematical self-review with a fresh run of the finite reference programs.

\subsection*{Mathematical review}
The review covers Propositions~\ref{MS:prop:deterrence-tv}, \ref{MS:prop:enforcement-garbling}, and~\ref{MS:prop:joint-robust}, Theorems~\ref{MS:thm:joint-enforcement} and~\ref{MS:thm:robust-staffing}, all four theorem and proposition results in Section~\ref{MS:sec:joint-synthesis}, and all six in Section~\ref{MS:sec:accountability}. It checks the shared-monitor quantifier, finite-horizon and indefinite safety, certified service, exact abstraction, strategic inspection, privacy, appeal capacity, authorized endings, and guarded handover.

No failed implication was identified in those 15 results under their written hypotheses. Their mathematical statements and proof bodies required no change. The other 47 retained theorem and proposition proofs were not newly audited in this targeted pass. The detailed claim cards, dependencies, boundary cases, and source comparisons accompany the source package.

\subsection*{Exact finite checks}
The delivered rerun passed 75 groups containing 922,032 assertions using integers and exact rational arithmetic. The checks include separate finite-policy searches, rational optimization comparisons, matching and direct-assignment comparisons, and deliberate invalid-certificate controls. Each group states its actual input range. The quotient checks cover two specified refinement families, not every possible concrete refinement.

The examples also keep their contracts distinct. The joint audit-and-service example needs five execution units. The accountability example needs eight, including the final record and revocation. Its five primitive transitions count separately from those eight execution units. Passing these bounded checks supports the reference constructions. It does not replace the general proofs.

\subsection*{Release boundary}
The finalized package contains the reading PDF, editable sources, local fonts and figure assets, exact reference programs, reproducible checks, and the review records. The original source release remains available as provenance. The theorem index and contents use the page numbers of this reading edition.

This is a theoretical monograph with a recorded project self-review. It makes no claim of independent peer review or proof-assistant verification. Physical implementation, authority, trusted evidence, observation laws, and environmental coverage remain explicit assumptions of the models. The research questions in Section~\ref{MS:sec:problems} remain open.


% FINAL EDITION ADDITIONS END


\chapter{PRIQ: additional mathematical material}
\setcounter{section}{0}\renewcommand{\thesection}{PRIQ.\Alph{section}}\renewcommand{\theHsection}{PRIQ.appendix.\arabic{section}}
% The corrected appendix as witnessed on physical pages 1128--1129 of the
% recovered prose-refined PDF. Replaces the older component appendix's two
% known defects. Labels retain the component source's logical identities.
\section{Formulation on the quotient space}\label{PRIQ:app:quotient}
Let $q:X_{\PRIQfine}\to X_{\PRIQfine}/\PRIQKer C$ be the canonical quotient map. Since $C$ is surjective, there is a unique vector-space isomorphism
\[
\widetilde C:X_{\PRIQfine}/\PRIQKer C\to X_{\PRIQcoarse},\qquad C=\widetilde Cq.
\]
The multilinear map $CA_m^{\PRIQfine}:X_{\PRIQfine}^m\to X_{\PRIQcoarse}$ descends if and only if its value is unchanged when one input is altered by an element of $\PRIQKer C$. The induced quotient tensor has the coarse space as its target:
\[
\overline A_m:(X_{\PRIQfine}/\PRIQKer C)^m\longrightarrow X_{\PRIQcoarse},\qquad
\overline A_m([x_1],\ldots,[x_m])=CA_m^{\PRIQfine}(x_1,\ldots,x_m).
\]
It follows that
\[
A_m^{\PRIQcoarse}=\overline A_m(\widetilde C^{-1},\ldots,\widetilde C^{-1}).
\]
There is no additional output factor $\widetilde C$ in this convention, since $\overline A_m$ already takes values in $X_{\PRIQcoarse}$.
Equivalently, define the quotient-valued tensor
\[
\overline B_m([x_1],\ldots,[x_m])=qA_m^{\PRIQfine}(x_1,\ldots,x_m).
\]
The slot-kernel criterion makes this independent of the representatives. Now $\overline A_m=\widetilde C\overline B_m$, and the alternative correctly typed formula is
\[
A_m^{\PRIQcoarse}=\widetilde C\overline B_m(\widetilde C^{-1},\ldots,\widetilde C^{-1}).
\]
This formulation makes \Cref{PRIQ:thm:universal-quotient} independent of coordinates and of the chosen right inverse.

\section{A finite observational partition lemma}\label{PRIQ:app:observation-partition}
Let $\PRIQcY$ be a finite observation alphabet, and give each state $i$ a strictly positive likelihood profile
\[
\Lambda(i)=(\lambda_y(i))_{y\in\PRIQcY},\qquad
i\sim_{\mathrm{obs}}j\quad\Longleftrightarrow\quad\Lambda(i)=\Lambda(j).
\]
\begin{PRIQproposition}[Coarsest deterministic observation partition]\label{PRIQ:prop:coarsest-observation-partition}
The profile classes form the coarsest partition on whose blocks every likelihood in the family is constant. If this partition has at least two blocks, it is also the coarsest aggregation among those with at least two nonempty coarse cells admitting exact Bayesian refinement for every $y\in\PRIQcY$ and every interior prior, with one fixed strictly positive coarse likelihood for each observation. Each coarse likelihood is determined up to one positive scale.

More generally, let $\pi_0$ be a prescribed partition with at least two nonempty blocks. The nonempty intersections of its blocks with the profile classes form its coarsest refinement admitting these exact Bayesian updates.
\end{PRIQproposition}
\begin{proof}
Every likelihood is constant on a profile class. If a partition makes every likelihood constant on its blocks, states in one block have identical profiles, so it refines the profile partition. This proves the first assertion without a block-count restriction.

For any aggregation with at least two nonempty coarse cells, apply \Cref{PRIQ:thm:bayes-necessity} to each observation. Exact updates for every interior prior force its fine likelihood to be constant on each fibre; the coarse likelihood is proportional to those values by one positive scalar. Conversely, \Cref{PRIQ:thm:bayes-sufficient} gives exact updates with those values. The profile partition is therefore coarsest in this class whenever it is nontrivial.

The intersection partition refines $\pi_0$, hence is nontrivial. It makes every likelihood constant, so admits the updates. Any other admissible refinement of $\pi_0$ also refines the profile partition by necessity and therefore refines their intersection partition. This proves coarseness, up to block names.
\end{proof}
\paragraph{The one-cell boundary.}
Without the restriction the coarseness claim is false. For fine likelihood $(1,2)$ and $C=(1\ 1)$, the prior $(1/2,1/2)$ updates to $(1/3,2/3)$ but its aggregate remains the unique coarse probability $(1)$. Every positive-evidence prior has the same aggregate after updating. The profile partition is discrete although the one-cell aggregation is exact for the coarse posterior.

If all profiles coincide, every partition supports the normalized updates. There need not be a unique coarsest nontrivial partition: on three states, the three two-block partitions are incomparable. Prescribing nontrivial distinctions removes this ambiguity.

\section{Certificate schemas and reproducibility}\label{PRIQ:app:schemas}
A minimal interaction certificate stores
\begin{equation}\label{PRIQ:eq:lumpability-certificate-schema}
  \mathsf L_m
  =
  (m,C,A_m^{\PRIQfine},A_m^{\PRIQcoarse},\Delta_m,
  \PRIQnorm{\Delta_m},\text{norm},\text{tolerance}).
\end{equation}
A Bayesian certificate stores
\begin{equation}\label{PRIQ:eq:bayes-certificate-schema}
  \mathsf B_y
  =
  (\lambda_y^{\PRIQfine},\lambda_y^{\PRIQcoarse},s_y,r_y,
  z_{*,y},d_{\PRIQBup,y}).
\end{equation}
A filtering certificate stores
\begin{equation}\label{PRIQ:eq:filter-certificate-schema}
  \mathsf F_N
  =
  (e_0^+,\{L_{\Phi,k}\},\{L_{\PRIQBup,k}\},
  \{d_{\Phi,k}\},\{d_{\PRIQBup,k}\},
  \{e_k^-,e_k^+\}).
\end{equation}
The release serializes the executed finite ledger as JSON files and CSV tables.  A certificate records a computation under specified inputs.  The ledger provides cryptographic authenticity only when it is paired with a separate signed digest layer.

\subsection{Execution commands}
From the source root:
\begin{lstlisting}[style=code]
PYTHONPATH=src python examples/run_certificates.py
PYTHONPATH=src python -m pytest -q
\end{lstlisting}
The first command regenerates every numerical table and figure used in the manuscript.  The second runs the independent unit tests.

\section{Theorem dependency ledger}\label{PRIQ:app:ledger}

\begin{longtable}{@{}p{0.12\linewidth}p{0.55\linewidth}p{0.23\linewidth}@{}}
\toprule
ID & Result & Dependency\\
\midrule
\endfirsthead
\toprule
ID & Result & Dependency\\
\midrule
\endhead
PRIQ-001 & Exact multilinear packet reconstruction & packet resolutions\\
PRIQ-002 & Universal multilinear quotient criterion & factorization through a quotient space\\
PRIQ-003 & Coefficient criterion at every interaction order & PRIQ-002, basis expansion\\
PRIQ-004 & Symmetry descends & PRIQ-002\\
PRIQ-005 & Exact quotient transitivity & intertwining algebra\\
PRIQ-006 & Multilevel defect identity & algebraic decomposition\\
PRIQ-007 & Analytic interaction descent & PRIQ-002, majorant convergence\\
PRIQ-008 & Jet naturality & differentiation of intertwining\\
PRIQ-009 & Treewise quotient naturality & PRIQ-002, rooted induction\\
PRIQ-010 & Exact quotient flow & intertwining of vector fields and uniqueness\\
PRIQ-011 & Normalized transfer naturality & positive intertwining, mass preservation\\
PRIQ-012 & Sufficiency of likelihood constancy within fibers & PRIQ-011\\
PRIQ-013 & Necessity of likelihood constancy within fibers & posterior ratios for two packets\\
PRIQ-014 & Exact prediction and update cocycle & PRIQ-010, PRIQ-012\\
PRIQ-015 & Sequential evidence factorization & PRIQ-014\\
PRIQ-016 & Expansion by hidden inputs & decomposition through a right inverse\\
PRIQ-017 & Majorant for the analytic field defect & coefficient bounds\\
PRIQ-018 & Approximate quotient flow bound & Grönwall inequality\\
PRIQ-019 & Normalization perturbation inequality & triangle inequality\\
PRIQ-020 & Bayesian refinement defect & PRIQ-019\\
PRIQ-021 & Bayesian Lipschitz certificate & PRIQ-019\\
PRIQ-022 & Posterior error recurrence & prediction and update decompositions\\
PRIQ-023 & Closed posterior certificate & iteration of PRIQ-022\\
PRIQ-024 & Evidence error certificate & posterior error, likelihood residual\\
PRIQ-025 & Hyperedge Parseval identity & orthogonal tensor blocks\\
PRIQ-026 & Optimal approximation with $k$ hyperedges & PRIQ-025\\
PRIQ-027 & Combined pruning and quotient defect & residual decomposition\\
\bottomrule
\end{longtable}

\section{Notation index}\label{PRIQ:app:notation}
\begin{longtable}{@{}p{0.22\linewidth}p{0.68\linewidth}@{}}
\toprule
Symbol & Meaning\\
\midrule
\endfirsthead
\toprule
Symbol & Meaning\\
\midrule
\endhead
$X_{\PRIQfine},X_{\PRIQcoarse}$ & fine and coarse state spaces\\
$P_i,\overline P_a$ & fine and coarse packet projectors\\
$\pi$ & map from fine packet labels to coarse labels\\
$C$ & surjective aggregation or quotient map\\
$E$ & right inverse selecting a fine representative of each coarse state\\
$P_C=EC$ & representative projection\\
$Q_C=I-EC$ & projection onto the hidden directions, with range $\PRIQKer C$\\
$A_m^{\PRIQfine},A_m^{\PRIQcoarse}$ & fine and coarse interaction tensors of order $m$\\
$\Delta_m$ & quotient defect at interaction order $m$\\
$F_{\PRIQfine},F_{\PRIQcoarse}$ & fine and coarse vector fields\\
$\Phi_t^{\PRIQfine},\Phi_t^{\PRIQcoarse}$ & fine and coarse flows or prediction maps\\
$L_{\PRIQfine},L_{\PRIQcoarse}$ & positive observation transfer operators\\
$\PRIQBup_{\lambda}$ & Bayesian normalized likelihood update\\
$\Gamma_k$ & single step prediction and update map\\
$\Psi_{1:N}$ & cocycle for a finite observation sequence\\
$d_{\Phi,k}$ & direct prediction quotient defect at step $k$\\
$d_{\PRIQBup,k}$ & direct Bayesian refinement defect at step $k$\\
$e_k^-,e_k^+$ & aggregate errors before and after an update\\
$z_{*,k}$ & positive evidence floor\\
$R_m$ & omitted hyperedge residual tensor\\
$H_m$ & entropy of the normalized hyperedge energy distribution at order $m$\\
\bottomrule
\end{longtable}


\chapter{PPT: additional mathematical material}
\setcounter{section}{0}\renewcommand{\thesection}{PPT.\Alph{section}}\renewcommand{\theHsection}{PPT.appendix.\arabic{section}}

\markboth{Reference material}{Reference material}
\section{Notation and hypotheses}
\label{PPT:app:notation}

Letters are local to a section. In particular, $A$ denotes required premises in a proof field, but a matrix in the binary examples. The type declarations, not the letter alone, determine the object.

\begingroup\small
\begin{longtable}{@{}L{0.24\linewidth}L{0.71\linewidth}@{}}
\toprule
Symbol & Meaning \\
\midrule
\endhead
$\PPTPcal$, $V$, $\PPTJcal$ & A declared proof field, its nodes, and its justifications. \\
$A$, $L$, $T_0$ & Required premises, admissible leaves, and requested targets, with $A\subseteq L$. \\
$(K,H)$ & Selected nodes and certified justification records. Only usable rules contribute to derivability. \\
$D_{\mathrm{cert}}(K,H)$ & The least certified closure in the selected record. \\
$\PPTAnc_L(T_0)$ & Backward ancestry expanded only outside the admissible leaves. \\
$\mathcal M_F(U)$ & Local models on patch $U$ satisfying all certified rules supported entirely in $U$. \\
$\Gamma(\mathcal M_F)$ & Coherent families of those local models under restriction. \\
$\operatorname{cl}_F(S)$ & The least declared rule model containing the positive observations $S$ and all admitted leaves. \\
$\operatorname{Ext}_F(S,N)$ & Rule models containing $S$ and avoiding the negative observations $N$. \\
$F_{\mathcal U}$ & The same field with only rules whose entire supports appear in a covering patch. \\
$Q(a)=(q_i(a))_i$ & The joint task representation, least in retained distinctions for its task family. \\
$\mathscr A$, $C$, $q$ & A class of possible archives, an encoding scheme, and a task on that class. \\
$P_K(S)=S\cap K$ & Pure retention of records selected by a fixed mask. \\
$D:E\to R$, $p:E\to S$ & A linear defect and a protected observable. The repair condition is $Dy=0$, $p(y)=p(x)$. \\
$H=\ker D\cap\ker p$ & The ambiguity space of a nonempty protected fiber. This $H$ is not a set of justifications. \\
$U_M(S)$ & Weighted target reach from the mandatory starts and a selected supplement, in a fixed directed bridge graph. \\
$n_{\mathrm{eff}}$ & A real, variance-equivalent review count in the specified covariance model. \\
\bottomrule
\end{longtable}
\endgroup

\subsection*{Assumptions that change the conclusions}
The unique-core theorem uses backward expansion that stops at admissible leaves and proves a least \emph{pair} of record sets. The compression results quantify over a whole archive class with one decoder for each task. Weighted domination counts only targets reachable before deletion. The greedy ratio requires a positive cardinality budget, with zero and saturated budgets handled separately.

Whole-support gluing is sufficient when the patches cover all nodes and every complete certified rule support. More generally, node coverage and entailment of every original rule by the visible rules are necessary and sufficient for universal gluing. Merely covering pairwise dependencies is not sufficient in general. The linear alternative requires a field. General algebraic functionals need not be bounded under an added topology. The information formula requires a finite field, a finite-dimensional feasible affine fiber, and a uniform distribution.

Endpoint equivalence fixes the complete assumption context. A refinement chain preserves the chosen tasks only when the corresponding recovery condition holds. The unrestricted positive-rule model always contains the full node set as a model. Nonempty global model spaces can still have empty prescribed extension fibers. Provenance change detection distinguishes information lost during serialization from collisions of different serialized bytes.

The review mean is unbiased only when its error means vanish. Its effective count requires positive marginal variance and positive variance of the mean. A minimum-variance aggregate allows real weights and assumes a known positive definite covariance matrix. No probability of mathematical correctness follows from these assumptions alone.




\section{Theorem ledger}
\label{PPT:app:ledger}

\emph{Written} denotes a written proof. \emph{Lean} denotes the specified interface with fresh compiler and kernel-replay evidence in this edition, not a formalization of the entire monograph. \emph{Mixed} distinguishes checked counts from their written probabilistic interpretation. The separate Lean concordance lists every source declaration and its scope.

\begingroup\small
\renewcommand{\arraystretch}{1.04}
\begin{longtable}{@{}L{0.115\linewidth}L{0.70\linewidth}L{0.105\linewidth}@{}}
\toprule
{\sffamily\small ID} & {\sffamily\small RESULT AND IMPLEMENTATION SUPPORT} & {\sffamily\small EVIDENCE}\\
\midrule
\endfirsthead
\toprule
{\sffamily\small ID} & {\sffamily\small RESULT AND IMPLEMENTATION SUPPORT} & {\sffamily\small EVIDENCE}\\
\midrule
\endhead
\bottomrule
\endfoot
PPT-001 & Finite derivability is unique and computable in topological order.\newline{\footnotesize\color{KomorebiMuted}\path{ProofField.derive}} & {\sffamily\small\color{KomorebiTeal}Written}\\[0pt]
PPT-002 & An AND-DAG has a least certificate pair, stopping at admissible leaves.\newline{\footnotesize\color{KomorebiMuted}\path{unique_and_dag_certificate}} & {\sffamily\small\color{KomorebiTeal}Written}\\[0pt]
PPT-003 & AND-DAG core extraction is linear with indexed adjacency.\newline{\footnotesize\color{KomorebiMuted}graph algorithm, not a runtime certificate} & {\sffamily\small\color{KomorebiTeal}Written}\\[0pt]
PPT-004 & Minimum Certificate Decision is NP-complete.\newline{\footnotesize\color{KomorebiMuted}Set Cover example} & {\sffamily\small\color{KomorebiTeal}Written}\\[0pt]
PPT-005 & Certified paths transport functorially.\newline{\footnotesize\color{KomorebiMuted}structural theorem} & {\sffamily\small\color{KomorebiTeal}Written}\\[0pt]
PPT-006 & Local auditable certificates compose under a certificate functor.\newline{\footnotesize\color{KomorebiMuted}structural theorem} & {\sffamily\small\color{KomorebiTeal}Written}\\[0pt]
PPT-007 & Task information loss follows from compression collisions.\newline{\footnotesize\color{KomorebiMuted}task API} & {\sffamily\small\color{KomorebiTeal}Lean}\\[0pt]
PPT-008 & Recoverability is a preorder.\newline{\footnotesize\color{KomorebiMuted}\path{recoverability_preorder}} & {\sffamily\small\color{KomorebiTeal}Written}\\[0pt]
PPT-009 & The quotient by mutual recoverability is a partial order.\newline{\footnotesize\color{KomorebiMuted}equivalence API} & {\sffamily\small\color{KomorebiTeal}Written}\\[0pt]
PPT-010 & Least sufficient retention masks depend on the task family.\newline{\footnotesize\color{KomorebiMuted}finite countermodel} & {\sffamily\small\color{KomorebiTeal}Written}\\[0pt]
PPT-011 & Zero verification loss can coexist with positive discovery loss.\newline{\footnotesize\color{KomorebiMuted}AND-DAG example} & {\sffamily\small\color{KomorebiTeal}Written}\\[0pt]
PPT-012 & Compression debt is monotone under retention.\newline{\footnotesize\color{KomorebiMuted}\path{compression_debt}} & {\sffamily\small\color{KomorebiTeal}Written}\\[0pt]
PPT-013 & Bridge criticality equals target domination.\newline{\footnotesize\color{KomorebiMuted}\path{dominated_targets}} & {\sffamily\small\color{KomorebiTeal}Written}\\[0pt]
PPT-014 & Weighted future reach is monotone and submodular.\newline{\footnotesize\color{KomorebiMuted}exhaustive finite checker} & {\sffamily\small\color{KomorebiTeal}Written}\\[0pt]
PPT-015 & Greedy cardinality retention achieves at least $1-1/e$ for positive feasible cardinality budgets.\newline{\footnotesize\color{KomorebiMuted}greedy and exact solvers} & {\sffamily\small\color{KomorebiTeal}Written}\\[0pt]
PPT-016 & Append-only certified updates preserve derivability.\newline{\footnotesize\color{KomorebiMuted}versioning check} & {\sffamily\small\color{KomorebiTeal}Written}\\[0pt]
\pagebreak
PPT-017 & Forward dependency closure is the minimal closed syntactic review set.\newline{\footnotesize\color{KomorebiMuted}\path{affected_descendants}} & {\sffamily\small\color{KomorebiTeal}Written}\\[0pt]
PPT-018 & Equicorrelated review has the stated effective count.\newline{\footnotesize\color{KomorebiMuted}covariance implementation} & {\sffamily\small\color{KomorebiTeal}Written}\\[0pt]
PPT-019 & GLS weights minimize unbiased linear review variance.\newline{\footnotesize\color{KomorebiMuted}\path{generalized_least_squares_weights}} & {\sffamily\small\color{KomorebiTeal}Written}\\[0pt]
PPT-020 & Constant-rate audit queues have exact stable, critical, and unstable regimes.\newline{\footnotesize\color{KomorebiMuted}queue simulation} & {\sffamily\small\color{KomorebiTeal}Written}\\[0pt]
PPT-021 & Optimal expansion depth depends on the goal.\newline{\footnotesize\color{KomorebiMuted}refinement implementation} & {\sffamily\small\color{KomorebiTeal}Written}\\[0pt]
PPT-022 & Pure retention projections commute and are idempotent.\newline{\footnotesize\color{KomorebiMuted}\path{retention_projection}} & {\sffamily\small\color{KomorebiTeal}Written}\\[0pt]
PPT-023 & Retention masks form a Boolean lattice.\newline{\footnotesize\color{KomorebiMuted}\path{retention_meet}, \path{retention_join}} & {\sffamily\small\color{KomorebiTeal}Written}\\[0pt]
PPT-024 & A locally refuted route can coexist with another certified route.\newline{\footnotesize\color{KomorebiMuted}Alternative justification field} & {\sffamily\small\color{KomorebiTeal}Written}\\[0pt]
PPT-025 & A checked local refutation persists under unchanged interpretation.\newline{\footnotesize\color{KomorebiMuted}Scoped witness principle} & {\sffamily\small\color{KomorebiTeal}Written}\\[0pt]
PPT-026 & Finite arrow-covered matching views glue uniquely.\newline{\footnotesize\color{KomorebiMuted}\path{FiniteViewSystem.glue}} & {\sffamily\small\color{KomorebiTeal}Lean}\\[0pt]
PPT-027 & Protected realizations admit exact solutions or obstruction witnesses, with classified fibers.\newline{\footnotesize\color{KomorebiMuted}\path{analyze_affine_realizations}} & {\sffamily\small\color{KomorebiTeal}Lean}\\[0pt]
PPT-028 & Finite affine fibers have exact observation counts. Uniform entropy is a written corollary.\newline{\footnotesize\color{KomorebiMuted}Lean counts and binary tests} & {\sffamily\small\color{KomorebiTeal}Mixed}\\[0pt]
PPT-029 & Common offset and correlated error have the displayed reviewer MSE.\newline{\footnotesize\color{KomorebiMuted}\path{shared_bias_mse}} & {\sffamily\small\color{KomorebiTeal}Written}\\[0pt]
PPT-030 & Declared derivability equals membership in every declared field model.\newline{\footnotesize\color{KomorebiMuted}Closure and countermodels} & {\sffamily\small\color{KomorebiTeal}Lean}\\[0pt]
PPT-031 & Local models form a contravariant diagram with an explicit realization equivalence.\newline{\footnotesize\color{KomorebiMuted}\path{Field.modelEquivSections}} & {\sffamily\small\color{KomorebiTeal}Lean}\\[0pt]
PPT-032 & Whole-support matching local models glue. Pairwise coverage can fail.\newline{\footnotesize\color{KomorebiMuted}Proof-field views and counterexample} & {\sffamily\small\color{KomorebiTeal}Lean}\\[0pt]
\pagebreak
PPT-033 & All predicate tasks are sufficient exactly for injective encodings.\newline{\footnotesize\color{KomorebiMuted}Compression collision theorem} & {\sffamily\small\color{KomorebiTeal}Lean}\\[0pt]
PPT-034 & Attained decoding is equivalent to fiber constancy. The decoder is unique on attained values.\newline{\footnotesize\color{KomorebiMuted}Recovery and finite decoder checks} & {\sffamily\small\color{KomorebiTeal}Lean}\\[0pt]
PPT-035 & Joint tasks are sufficient exactly when each coordinate task is. Recovery maps compose.\newline{\footnotesize\color{KomorebiMuted}Recovery and joint partitions} & {\sffamily\small\color{KomorebiTeal}Lean}\\[0pt]
PPT-036 & Protected rule data have an extension exactly when their least closure avoids the forbidden nodes.\newline{\footnotesize\color{KomorebiMuted}Extension and independent trace replay} & {\sffamily\small\color{KomorebiTeal}Lean}\\[0pt]
PPT-037 & Closing earlier data before later data preserves the closure. Adding observations preserves obstructions.\newline{\footnotesize\color{KomorebiMuted}Extension} & {\sffamily\small\color{KomorebiTeal}Lean}\\[0pt]
PPT-038 & A cover has universal gluing exactly when its visible rules entail every original rule.\newline{\footnotesize\color{KomorebiMuted}SemanticCover and finite cover certificates} & {\sffamily\small\color{KomorebiTeal}Lean}\\[0pt]
PPT-039 & Typed rule transports preserve derivations, pull models back, and transport obstruction witnesses.\newline{\footnotesize\color{KomorebiMuted}RuleTransport} & {\sffamily\small\color{KomorebiTeal}Lean}\\[0pt]
PPT-040 & A valid rule transport can collapse compatible source observations into an impossible target request.\newline{\footnotesize\color{KomorebiMuted}ExtensionExamples} & {\sffamily\small\color{KomorebiTeal}Lean}\\[0pt]
PPT-041 & Hashing cannot add recoverable tasks. Restricted injectivity prevents further task loss.\newline{\footnotesize\color{KomorebiMuted}Recovery} & {\sffamily\small\color{KomorebiTeal}Lean}\\[0pt]
PPT-042 & Executable trace replay preserves derivability and composes.\newline{\footnotesize\color{KomorebiMuted}Trace and TraceCompleteness} & {\sffamily\small\color{KomorebiTeal}Lean}\\[0pt]
PPT-043 & Accepted complete traces identify the least seeded model.\newline{\footnotesize\color{KomorebiMuted}Trace and TraceCompleteness} & {\sffamily\small\color{KomorebiTeal}Lean}\\[0pt]
PPT-044 & Accepted forbidden-node witnesses exist exactly for infeasible protected requests.\newline{\footnotesize\color{KomorebiMuted}Trace and TraceCompleteness} & {\sffamily\small\color{KomorebiTeal}Lean}\\[0pt]
PPT-045 & Typed JSON encoding has a left inverse and is injective.\newline{\footnotesize\color{KomorebiMuted}Trace.Wire} & {\sffamily\small\color{KomorebiTeal}Lean}\\[0pt]
PPT-046 & Accepted JSON records certify the exact decoded protected answer.\newline{\footnotesize\color{KomorebiMuted}Trace.Wire} & {\sffamily\small\color{KomorebiTeal}Lean}\\[0pt]
PPT-047 & Valid indexed problems derive only bounded node addresses.\newline{\footnotesize\color{KomorebiMuted}Trace.Wire} & {\sffamily\small\color{KomorebiTeal}Lean}\\[0pt]
\end{longtable}
\renewcommand{\arraystretch}{1.0}
\endgroup


\section{Reproduction and verification records}
\label{PPT:app:reproduction}

The package includes written proofs, exact finite producers and independent result checkers, and a fresh pinned Lean run. All eighteen modules were rebuilt and replayed. The mathematical imports remain cached and are not represented as freshly rebuilt dependencies. The current compiler record is \path{evidence/formal_1.8.0/}, while older evidence retains its original version. Source changes and the new mathematical environments are listed in the current progress ledger and publication report.

Before modifying the extracted package, check the delivered files:
\begin{lstlisting}[language=bash]
python scripts/check_manifest.py
\end{lstlisting}
Then run the finite and publication checks:
\begin{lstlisting}[language=bash]
python -m pip install -e ".[release]" --no-build-isolation
python -m pytest -o addopts='' -q
python scripts/deep_validation.py
python scripts/verify_bridge.py
python scripts/check_recovery_extension.py
python scripts/check_trace_exchange.py
python scripts/verify_publication_edition.py
\end{lstlisting}

Regenerate the exact figures with:
\begin{lstlisting}[language=bash]
python scripts/design_figures.py
\end{lstlisting}
The publication edition is 1.8.0 and the Python distribution is 1.6.0. This distribution adds a versioned trace-exchange codec, an independent indexed checker, and export/check commands without changing the existing named-field schema or provenance byte format. Its historical migration remains documented in \path{docs/PROVENANCE_FORMAT.md}. Installation may obtain build and numerical dependencies unless they are already supplied locally.

The manuscript can be rebuilt independently from \path{docs/}:
\begin{lstlisting}[language=bash]
cd docs
pdflatex -interaction=nonstopmode -halt-on-error Proof_Path_Theory.tex
pdflatex -interaction=nonstopmode -halt-on-error Proof_Path_Theory.tex
pdflatex -interaction=nonstopmode -halt-on-error Proof_Path_Theory.tex
\end{lstlisting}


\subsection*{Formal and cross-language replay}
A separate formal replay, with the pinned runtime already available, uses:
\begin{lstlisting}[language=bash]
python scripts/verify_lean.py --runtime /path/to/portable-kit
\end{lstlisting}

For PowerShell, extraction precedes installation:
\begin{lstlisting}[language={},basicstyle=\ttfamily\fontsize{9.5}{11.4}\selectfont]
$zip = '.\Proof_Path_Theory_Monograph_1.8.0.zip'
Expand-Archive -LiteralPath $zip -DestinationPath '.\proof-path-release'
Set-Location '.\proof-path-release\Proof_Path_Theory_Monograph_v1.8.0'
python scripts/check_manifest.py
python -m pip install -e ".[release]" --no-build-isolation
python -m pytest -o addopts='' -q
\end{lstlisting}

The trace interface can be exercised without a Lean runtime:
\begin{lstlisting}[language=bash]
proof-path wire-export examples/and_dag.json --negative B --output-prefix trial
proof-path wire-check trial.problem.json trial.certificate.json
\end{lstlisting}
The export writes the fixed indexed problem, the checked result, and the node/rule address book separately. The general AST theorem is in \path{ProofPath.TraceWire}. The executed byte adapter is \path{formal/CheckTraceJson.lean}; the fresh formal verifier runs it against the generated canonical exchange cases. The address book must still be checked against the intended source snapshot.

The dated report identifies the commands actually run. These checks do not supply independent peer review, universal software verification, or empirical validation of the social models. The redistribution-license decision remains open.


\section{A complete minimal proof-field record}
\label{PPT:app:schema}

The following record defines both its premise and its target. It is a complete importable example rather than a schema fragment with an unresolved target reference. The rule is explicitly assumed for this demonstration. Its declared certification must not be mistaken for a proof of an arbitrary implication between the displayed labels.

\begin{lstlisting}[language={},basicstyle=\ttfamily\fontsize{9.5}{11.4}\selectfont]
{
  "schema_version": "1.0",
  "premises": ["A"],
  "targets": ["B"],
  "nodes": [
    {"id": "A", "label": "Premise A", "kind": "premise"},
    {"id": "B", "label": "Target B", "kind": "claim"}
  ],
  "evidence": [
    {
      "id": "e-rule",
      "kind": "assumption",
      "source": "Declared demonstration rule A implies B",
      "scope": "This finite inference model only",
      "certifies": true
    }
  ],
  "justifications": [
    {
      "id": "j1",
      "premises": ["A"],
      "conclusion": "B",
      "rule": "assumed implication",
      "certified": true,
      "evidence_ids": ["e-rule"]
    }
  ]
}
\end{lstlisting}

The file \path{examples/minimal_monograph.json} contains this record. The public importer checks its identifiers and Boolean fields, and declared closure reaches \texttt{B}. Mathematical soundness still requires the explicit interpretation contract of Section~\ref{PPT:sec:old24}.


\phantomsection


\chapter{Large Signal Model: proof and evidence register}\begingroup\setcounter{secnumdepth}{-1}\section{Proof, assumptions and evidence
register}\label{LSMRegister:proof-assumptions-and-evidence-register}

\textbf{Foundations 1.0.0. Frozen scope: 87 propositions. No new
numbered results.}

This register records the source location, principal internal
dependencies, necessary inputs and strongest claim boundary for each
proposition. Code links identify a reference implementation or finite
illustration, not an executable general proof. All general proofs remain
written arguments in the cumulative volume. No Lean kernel or
independent external referee was used.

Read the entire indicated section before applying a proposition.
Assumptions can precede its bold statement. The dependency map is
curated rather than extracted from formal proof terms. In particular,
Proposition 9 uses Proposition 12, a forward reference in the
presentation. The map is acyclic.

All 87 statement-and-proof spans are byte-identical to the immediate
0.9.0 predecessor. The earlier P72 dimension clarification is retained.
Relative to preserved 0.8.0, exactly that one correction remains. The
release-freeze checker verifies the immediate predecessor separately.

The machine-readable JSON contains full source and code hashes. Hashes
identify artifacts; they do not establish truth. The manuscript retains
primary-source attribution for the classical ingredients.


\section{Realization and predictive
closure}\label{LSMRegister:realization-and-predictive-closure}

Propositions 1--13. Related finite suite:
\path{checks/exact_foundation_checks.py}.

\par\noindent\begin{minipage}{\textwidth}
\subsection{P1. Exact realization and
consistency}\label{LSMRegister:p1.-exact-realization-and-consistency}

\textbf{Location and role.} Section 2, source lines 51--73. Identity.
Principal dependencies: No principal internal prerequisite recorded.

\textbf{Requires.} Finite Hilbert spaces; a declared analysis/synthesis
pair with a left inverse.

\textbf{Boundary.} Reconstruction is not predictive sufficiency.

\textbf{Evidence.} Written proof. The related suite contains finite
block-level examples; no dedicated general algorithm is claimed for this
result.
\end{minipage}\par\medskip

\par\noindent\begin{minipage}{\textwidth}
\subsection{P2. Coherent transport and operator
realization}\label{LSMRegister:p2.-coherent-transport-and-operator-realization}

\textbf{Location and role.} Section 3, source lines 93--112. Identity.
Principal dependencies: P1.

\textbf{Requires.} Consistent coefficient spaces; one pairwise coherent
family of chart maps.

\textbf{Boundary.} Agreement across charts does not establish
correctness.

\textbf{Evidence.} Written proof. The related suite contains finite
block-level examples; no dedicated general algorithm is claimed for this
result.
\end{minipage}\par\medskip

\par\noindent\begin{minipage}{\textwidth}
\subsection{P3. Complete refinement preserves
reconstruction}\label{LSMRegister:p3.-complete-refinement-preserves-reconstruction}

\textbf{Location and role.} Section 4, source lines 130--153. Identity.
Principal dependencies: No principal internal prerequisite recorded.

\textbf{Requires.} Orthogonal projector tree; a complete antichain.

\textbf{Boundary.} Omitted children or energy-only storage lose the
reconstruction contract.

\textbf{Evidence.} Written proof. The related suite contains finite
block-level examples; no dedicated general algorithm is claimed for this
result.
\end{minipage}\par\medskip

\par\noindent\begin{minipage}{\textwidth}
\subsection{P4. Transfer reconstruction and
refinement}\label{LSMRegister:p4.-transfer-reconstruction-and-refinement}

\textbf{Location and role.} Section 4, source lines 130--187. Identity.
Principal dependencies: P3.

\textbf{Requires.} Linear maps with complete input, intermediate and
output projector resolutions.

\textbf{Boundary.} Nonlinear maps require interaction terms.

\textbf{Evidence.} Written proof. The related suite contains finite
block-level examples; no dedicated general algorithm is claimed for this
result.
\end{minipage}\par\medskip

\par\noindent\begin{minipage}{\textwidth}
\subsection{P5. Polynomial packet
interactions}\label{LSMRegister:p5.-polynomial-packet-interactions}

\textbf{Location and role.} Section 4, source lines 130--209. Identity.
Principal dependencies: P4.

\textbf{Requires.} A specified finite multilinear polynomial expansion
and complete packets.

\textbf{Boundary.} No automatic general nonlinear dictionary discovery.

\textbf{Evidence.} Written proof. The related suite contains finite
block-level examples; no dedicated general algorithm is claimed for this
result.
\end{minipage}\par\medskip

\par\noindent\begin{minipage}{\textwidth}
\subsection{P6. Exact quotient
criterion}\label{LSMRegister:p6.-exact-quotient-criterion}

\textbf{Location and role.} Section 5, source lines 213--232. Criterion.
Principal dependencies: No principal internal prerequisite recorded.

\textbf{Requires.} A summary onto its image; deterministic transitions
on the stated domain.

\textbf{Boundary.} A quotient must be constant on every summary fiber.

\textbf{Evidence.} Written proof. The related suite contains finite
block-level examples; no dedicated general algorithm is claimed for this
result.
\end{minipage}\par\medskip

\par\noindent\begin{minipage}{\textwidth}
\subsection{P7. Closure defect and compositional
defect}\label{LSMRegister:p7.-closure-defect-and-compositional-defect}

\textbf{Location and role.} Section 5, source lines 213--261. Criterion.
Principal dependencies: P6.

\textbf{Requires.} Orthogonal retained projector; linear dynamics on the
declared whole state space.

\textbf{Boundary.} Small omitted energy alone does not bound future
influence.

\textbf{Evidence.} Written proof. The related suite contains finite
block-level examples; no dedicated general algorithm is claimed for this
result.
\end{minipage}\par\medskip

\par\noindent\begin{minipage}{\textwidth}
\subsection{P8. Minimal linear predictive
completion}\label{LSMRegister:p8.-minimal-linear-predictive-completion}

\textbf{Location and role.} Section 6, source lines 267--286.
Construction. Principal dependencies: No principal internal prerequisite
recorded.

\textbf{Requires.} Known finite family of linear operators and initial
observables.

\textbf{Boundary.} The constructed observables need not be measured or
low dimensional.

\textbf{Evidence.} Written proof. The related suite contains finite
block-level examples; no dedicated general algorithm is claimed for this
result.
\end{minipage}\par\medskip

\par\noindent\begin{minipage}{\textwidth}
\subsection{P9. Coarsest joint finite
closure}\label{LSMRegister:p9.-coarsest-joint-finite-closure}

\textbf{Location and role.} Section 7, source lines 300--338.
Construction. Principal dependencies: P12.

\textbf{Requires.} Known finite stochastic transitions; positive
likelihoods; initial partition with at least two blocks.

\textbf{Boundary.} The one-block Bayesian necessity exception must be
retained.

\textbf{Evidence.} Written proof. The related suite contains finite
block-level examples; no dedicated general algorithm is claimed for this
result.
\end{minipage}\par\medskip

\par\noindent\begin{minipage}{\textwidth}
\subsection{P10. Exact projected memory and a conditional tail
bound}\label{LSMRegister:p10.-exact-projected-memory-and-a-conditional-tail-bound}

\textbf{Location and role.} Section 8, source lines 344--385. Bound.
Principal dependencies: P7.

\textbf{Requires.} Linear retained/unresolved split; contraction and
bounded history for the tail estimate.

\textbf{Boundary.} One-step tail error is not recursive rollout error.

\textbf{Evidence.} Written proof. The related suite contains finite
block-level examples; no dedicated general algorithm is claimed for this
result.
\end{minipage}\par\medskip

\par\noindent\begin{minipage}{\textwidth}
\subsection{P11. Ideal statistical benefit of
refinement}\label{LSMRegister:p11.-ideal-statistical-benefit-of-refinement}

\textbf{Location and role.} Section 9, source lines 393--406. Identity.
Principal dependencies: No principal internal prerequisite recorded.

\textbf{Requires.} Square-integrable target; nested sigma algebras;
exact conditional expectations.

\textbf{Boundary.} An estimated finite-data predictor can worsen after
refinement.

\textbf{Evidence.} Written proof. The related suite contains finite
block-level examples; no dedicated general algorithm is claimed for this
result.
\end{minipage}\par\medskip

\par\noindent\begin{minipage}{\textwidth}
\subsection{P12. Bayesian closure and its exact
defect}\label{LSMRegister:p12.-bayesian-closure-and-its-exact-defect}

\textbf{Location and role.} Section 10, source lines 412--442.
Criterion. Principal dependencies: No principal internal prerequisite
recorded.

\textbf{Requires.} At least two coarse blocks; strictly positive
likelihoods for universal necessity.

\textbf{Boundary.} Fixed-prior defect permits nonnegative likelihoods
only with positive evidence.

\textbf{Evidence.} Written proof. The related suite contains finite
block-level examples; no dedicated general algorithm is claimed for this
result.
\end{minipage}\par\medskip

\par\noindent\begin{minipage}{\textwidth}
\subsection{P13. Strong fixed-horizon refinement
limit}\label{LSMRegister:p13.-strong-fixed-horizon-refinement-limit}

\textbf{Location and role.} Section 12, source lines 473--499. Limit.
Principal dependencies: No principal internal prerequisite recorded.

\textbf{Requires.} Bounded operator; orthogonal projections converging
strongly to the identity.

\textbf{Boundary.} Fixed horizon and strong convergence, not norm or
time-uniform convergence.

\textbf{Evidence.} Written proof. The related suite contains finite
block-level examples; no dedicated general algorithm is claimed for this
result.
\end{minipage}\par\medskip


\section{Learning and stability}\label{LSMRegister:learning-and-stability}

Propositions 14--25. Related finite suite:
\path{checks/learning_stability_checks.py}.

\par\noindent\begin{minipage}{\textwidth}
\subsection{P14. Exact causal delay
closure}\label{LSMRegister:p14.-exact-causal-delay-closure}

\textbf{Location and role.} Section 15, source lines 586--618.
Criterion. Principal dependencies: P6, P8.

\textbf{Requires.} Autonomous finite linear process; positive delay
length; true observation-map rank plateau.

\textbf{Boundary.} An empirical rank plateau alone is not
predictive-space coverage.

\textbf{Code link.} \path{checks/observable_reference.py},
\path{learn_delay}, lines 107--109. Exact finite delay learner under
supplied closure and span assumptions.
\end{minipage}\par\medskip

\par\noindent\begin{minipage}{\textwidth}
\subsection{P15. Data-span identification and its
boundary}\label{LSMRegister:p15.-data-span-identification-and-its-boundary}

\textbf{Location and role.} Section 16, source lines 626--646.
Criterion. Principal dependencies: No principal internal prerequisite
recorded.

\textbf{Requires.} Paired exact columns; input-kernel inclusion in
successor kernel.

\textbf{Boundary.} Interpolation identifies the sampled span, not unseen
directions.

\textbf{Code link.} \path{checks/observable_reference.py},
\path{fit_closed_span}, lines 62--90. Exact sampled-span
interpolation and kernel compatibility.
\end{minipage}\par\medskip

\par\noindent\begin{minipage}{\textwidth}
\subsection{P16. Fixed-space identification
bound}\label{LSMRegister:p16.-fixed-space-identification-bound}

\textbf{Location and role.} Section 16, source lines 626--688. Bound.
Principal dependencies: P15.

\textbf{Requires.} Fixed full-row-rank coordinates; input error below
the smallest singular value.

\textbf{Boundary.} Does not learn a coordinate space or certify its
rank.

\textbf{Code link.} \path{checks/stochastic_reference.py},
\path{noisy_core_bound}, lines 168--186. Related finite fixed-core
perturbation bound, not a general noisy-coordinate learner.
\end{minipage}\par\medskip

\par\noindent\begin{minipage}{\textwidth}
\subsection{P17. Finite-prefix
indistinguishability}\label{LSMRegister:p17.-finite-prefix-indistinguishability}

\textbf{Location and role.} Section 17, source lines 694--698.
Obstruction. Principal dependencies: No principal internal prerequisite
recorded.

\textbf{Requires.} Any fixed finite observation horizon; no supplied
dimension bound.

\textbf{Boundary.} A finite prefix cannot rule out a later hidden
response.

\textbf{Evidence.} Written proof. The related suite contains finite
block-level examples; no dedicated general algorithm is claimed for this
result.
\end{minipage}\par\medskip

\par\noindent\begin{minipage}{\textwidth}
\subsection{P18. Compositional defect
propagation}\label{LSMRegister:p18.-compositional-defect-propagation}

\textbf{Location and role.} Section 19, source lines 764--800. Bound.
Principal dependencies: No principal internal prerequisite recorded.

\textbf{Requires.} Lipschitz maps on domains containing both compared
trajectories; declared local defects.

\textbf{Boundary.} Different policies need a separate action-mismatch
bound.

\textbf{Evidence.} Written proof. The related suite contains finite
block-level examples; no dedicated general algorithm is claimed for this
result.
\end{minipage}\par\medskip

\par\noindent\begin{minipage}{\textwidth}
\subsection{P19. Uniform prediction bound for a stable finite-memory
truncation}\label{LSMRegister:p19.-uniform-prediction-bound-for-a-stable-finite-memory-truncation}

\textbf{Location and role.} Section 20, source lines 842--893. Bound.
Principal dependencies: P10, P18.

\textbf{Requires.} Finite-memory gain sum below one; uniformly bounded
true-history defect.

\textbf{Boundary.} Spectral radius alone is not the stated induced-norm
condition.

\textbf{Evidence.} Written proof. The related suite contains finite
block-level examples; no dedicated general algorithm is claimed for this
result.
\end{minipage}\par\medskip

\par\noindent\begin{minipage}{\textwidth}
\subsection{P20. Packet-interaction truncation
bound}\label{LSMRegister:p20.-packet-interaction-truncation-bound}

\textbf{Location and role.} Section 21, source lines 920--952. Bound.
Principal dependencies: P5.

\textbf{Requires.} Declared polynomial terms; restricted multilinear
norms; bounded domain.

\textbf{Boundary.} Bounds omitted terms, not the quality of an
undiscovered basis.

\textbf{Evidence.} Written proof. The related suite contains finite
block-level examples; no dedicated general algorithm is claimed for this
result.
\end{minipage}\par\medskip

\par\noindent\begin{minipage}{\textwidth}
\subsection{P21. Coherent reconciliation and coordinate-stable
transport}\label{LSMRegister:p21.-coherent-reconciliation-and-coordinate-stable-transport}

\textbf{Location and role.} Section 22, source lines 970--1012.
Construction. Principal dependencies: P1, P2.

\textbf{Requires.} Fixed exact atlases; normalized nonnegative weights;
linear chart predictors.

\textbf{Boundary.} Consistent reconciliation need not improve accuracy.

\textbf{Evidence.} Written proof. The related suite contains finite
block-level examples; no dedicated general algorithm is claimed for this
result.
\end{minipage}\par\medskip

\par\noindent\begin{minipage}{\textwidth}
\subsection{P22. The upward-lift
obstruction}\label{LSMRegister:p22.-the-upward-lift-obstruction}

\textbf{Location and role.} Section 23, source lines 1024--1042.
Criterion. Principal dependencies: P6.

\textbf{Requires.} Coarse and proposed finer state maps; a norm for the
lower bound.

\textbf{Boundary.} A genuinely finer state generally cannot be
initialized from the coarse state alone.

\textbf{Evidence.} Written proof. The related suite contains finite
block-level examples; no dedicated general algorithm is claimed for this
result.
\end{minipage}\par\medskip

\par\noindent\begin{minipage}{\textwidth}
\subsection{P23. Exact accounting for an enumerated
grammar}\label{LSMRegister:p23.-exact-accounting-for-an-enumerated-grammar}

\textbf{Location and role.} Section 24, source lines 1060--1093.
Identity. Principal dependencies: No principal internal prerequisite
recorded.

\textbf{Requires.} Countable fixed grammar and prior; finite positive
full and active evidence; common likelihood measure.

\textbf{Boundary.} Continuous likelihood densities need not be bounded
by one.

\textbf{Evidence.} Written proof. The related suite contains finite
block-level examples; no dedicated general algorithm is claimed for this
result.
\end{minipage}\par\medskip

\par\noindent\begin{minipage}{\textwidth}
\subsection{P24. Prior-preserving program
refinement}\label{LSMRegister:p24.-prior-preserving-program-refinement}

\textbf{Location and role.} Section 24, source lines 1060--1117.
Identity. Principal dependencies: P23.

\textbf{Requires.} Prior-mass-preserving split; a whole-record
likelihood mixture identity.

\textbf{Boundary.} A fixed child mixture is not independent switching at
each step.

\textbf{Evidence.} Written proof. The related suite contains finite
block-level examples; no dedicated general algorithm is claimed for this
result.
\end{minipage}\par\medskip

\par\noindent\begin{minipage}{\textwidth}
\subsection{P25. Observationally grounded finite linear
realization}\label{LSMRegister:p25.-observationally-grounded-finite-linear-realization}

\textbf{Location and role.} Section 25, source lines 1125--1139.
Integration. Principal dependencies: P1, P2, P3, P4, P14, P15, P16, P18.

\textbf{Requires.} Autonomous linear truth; sufficient delay; observed
windows spanning the true predictive space.

\textbf{Boundary.} The learner does not infer its model-class and
coverage assumptions from a fit.

\textbf{Code link.} \path{checks/observable_reference.py},
\path{learn_delay}, lines 107--109. Autonomous finite example only.
\end{minipage}\par\medskip


\section{Controlled stochastic
realization}\label{LSMRegister:controlled-stochastic-realization}

Propositions 26--37. Related finite suite:
\path{checks/controlled_stochastic_checks.py}.

\par\noindent\begin{minipage}{\textwidth}
\subsection{P26. Finite instrument
closure}\label{LSMRegister:p26.-finite-instrument-closure}

\textbf{Location and role.} Section 27, source lines 1159--1225.
Construction. Principal dependencies: P8.

\textbf{Requires.} Known nonnegative finite instrument; all initial
beliefs; exact mass row.

\textbf{Boundary.} A known-model completion is not an observable-data
learner.

\textbf{Evidence.} Written proof. The related suite contains finite
block-level examples; no dedicated general algorithm is claimed for this
result.
\end{minipage}\par\medskip

\par\noindent\begin{minipage}{\textwidth}
\subsection{P27. Minimal finite behavior
realization}\label{LSMRegister:p27.-minimal-finite-behavior-realization}

\textbf{Location and role.} Section 28, source lines 1231--1264.
Construction. Principal dependencies: No principal internal prerequisite
recorded.

\textbf{Requires.} Causal nonnegative finite-alphabet behavior; finite
full history/test rank.

\textbf{Boundary.} Minimality is within linear realizations of the
declared behavior.

\textbf{Evidence.} Written proof. The related suite contains finite
block-level examples; no dedicated general algorithm is claimed for this
result.
\end{minipage}\par\medskip

\par\noindent\begin{minipage}{\textwidth}
\subsection{P28. Observable table
reconstruction}\label{LSMRegister:p28.-observable-table-reconstruction}

\textbf{Location and role.} Section 29, source lines 1270--1300.
Construction. Principal dependencies: P27.

\textbf{Requires.} Known rank upper bound; invertible matching core;
exact shifted tables and initial tests.

\textbf{Boundary.} Finite invertibility alone supplies only a rank lower
bound.

\textbf{Code link.} \path{checks/stochastic_reference.py},
\path{fit_core}, lines 109--145. Population-table reconstruction with
declared full-rank coverage.
\end{minipage}\par\medskip

\par\noindent\begin{minipage}{\textwidth}
\subsection{P29. Canonical probability
cone}\label{LSMRegister:p29.-canonical-probability-cone}

\textbf{Location and role.} Section 30, source lines 1308--1334.
Existence. Principal dependencies: P27, P28.

\textbf{Requires.} Canonical finite behavior space and all future tests.

\textbf{Boundary.} An all-test cone is not automatically finitely
describable.

\textbf{Evidence.} Written proof. The related suite contains finite
block-level examples; no dedicated general algorithm is claimed for this
result.
\end{minipage}\par\medskip

\par\noindent\begin{minipage}{\textwidth}
\subsection{P30. Positive packet decomposition and its
boundary}\label{LSMRegister:p30.-positive-packet-decomposition-and-its-boundary}

\textbf{Location and role.} Section 31, source lines 1342--1358.
Criterion. Principal dependencies: No principal internal prerequisite
recorded.

\textbf{Requires.} Complete cone-preserving packet maps; strictly
positive mass functional.

\textbf{Boundary.} Signed Fourier or Haar details need not be
probability components.

\textbf{Evidence.} Written proof. The related suite contains finite
block-level examples; no dedicated general algorithm is claimed for this
result.
\end{minipage}\par\medskip

\par\noindent\begin{minipage}{\textwidth}
\subsection{P31. Missing-intervention
nonidentifiability}\label{LSMRegister:p31.-missing-intervention-nonidentifiability}

\textbf{Location and role.} Section 32, source lines 1372--1376.
Obstruction. Principal dependencies: No principal internal prerequisite
recorded.

\textbf{Requires.} At least two possible interventions with a missing
action.

\textbf{Boundary.} No amount of data restricted to one action identifies
an unobserved action.

\textbf{Evidence.} Written proof. The related suite contains finite
block-level examples; no dedicated general algorithm is claimed for this
result.
\end{minipage}\par\medskip

\par\noindent\begin{minipage}{\textwidth}
\subsection{P32. Finite-table rank separation and subspace
stability}\label{LSMRegister:p32.-finite-table-rank-separation-and-subspace-stability}

\textbf{Location and role.} Section 33, source lines 1388--1426. Bound.
Principal dependencies: No principal internal prerequisite recorded.

\textbf{Requires.} True finite-table rank; known perturbation budget;
singular gap above twice that budget.

\textbf{Boundary.} Does not exclude true directions absent from the
selected table.

\textbf{Code link.} \path{checks/stochastic_reference.py},
\path{finite_subspace_report}, lines 148--165. Finite matrix
singular-subspace report.
\end{minipage}\par\medskip

\par\noindent\begin{minipage}{\textwidth}
\subsection{P33. Noisy core-block
realization}\label{LSMRegister:p33.-noisy-core-block-realization}

\textbf{Location and role.} Section 34, source lines 1434--1477. Bound.
Principal dependencies: P28.

\textbf{Requires.} Fixed square invertible core and stated
operator/table perturbation budgets.

\textbf{Boundary.} Good core inversion is not global positivity
certification.

\textbf{Code link.} \path{checks/stochastic_reference.py},
\path{noisy_core_bound}, lines 168--186. Fixed-core perturbation
calculation.
\end{minipage}\par\medskip

\par\noindent\begin{minipage}{\textwidth}
\subsection{P34. Normalization-aware filtering
stability}\label{LSMRegister:p34.-normalization-aware-filtering-stability}

\textbf{Location and role.} Section 35, source lines 1483--1523. Bound.
Principal dependencies: P18.

\textbf{Requires.} Common declared norm; bounded true states; positive
likelihood margin.

\textbf{Boundary.} Rare-event conditioning can amplify a small
unnormalized error.

\textbf{Evidence.} Written proof. The related suite contains finite
block-level examples; no dedicated general algorithm is claimed for this
result.
\end{minipage}\par\medskip

\par\noindent\begin{minipage}{\textwidth}
\subsection{P35. Reset-experiment probability
certificate}\label{LSMRegister:p35.-reset-experiment-probability-certificate}

\textbf{Location and role.} Section 36, source lines 1531--1555.
Confidence. Principal dependencies: No principal internal prerequisite
recorded.

\textbf{Requires.} A prospectively fixed list of experiments;
independent resets per experiment; known reset law.

\textbf{Boundary.} Not a guarantee for one dependent stream or
post-selected experiments.

\textbf{Evidence.} Written proof. The related suite contains finite
block-level examples; no dedicated general algorithm is claimed for this
result.
\end{minipage}\par\medskip

\par\noindent\begin{minipage}{\textwidth}
\subsection{P36. Exact normalized closure
defect}\label{LSMRegister:p36.-exact-normalized-closure-defect}

\textbf{Location and role.} Section 37, source lines 1563--1593.
Identity. Principal dependencies: No principal internal prerequisite
recorded.

\textbf{Requires.} Mass-preserving summary; nonnegative instrument; two
positive branch likelihoods.

\textbf{Boundary.} Posterior agreement can hide incorrect likelihoods.

\textbf{Evidence.} Written proof. The related suite contains finite
block-level examples; no dedicated general algorithm is claimed for this
result.
\end{minipage}\par\medskip

\par\noindent\begin{minipage}{\textwidth}
\subsection{P37. Conditional observation-based stochastic
completion}\label{LSMRegister:p37.-conditional-observation-based-stochastic-completion}

\textbf{Location and role.} Section 39, source lines 1687--1702.
Integration. Principal dependencies: P27, P28, P29, P30, P32, P33, P34,
P35.

\textbf{Requires.} Full rank coverage; exact population core; separate
perturbation and likelihood assumptions.

\textbf{Boundary.} No global fitted cone follows from the spectral
estimate alone.

\textbf{Evidence.} Written proof. The related suite contains finite
block-level examples; no dedicated general algorithm is claimed for this
result.
\end{minipage}\par\medskip


\section{Probability-domain
certificates}\label{LSMRegister:probability-domain-certificates}

Propositions 38--49. Related finite suite:
\path{checks/certified_domain_checks.py}.

\par\noindent\begin{minipage}{\textwidth}
\subsection{P38. Test-cone hierarchy and a finite stopping
certificate}\label{LSMRegister:p38.-test-cone-hierarchy-and-a-finite-stopping-certificate}

\textbf{Location and role.} Section 42, source lines 1750--1772.
Criterion. Principal dependencies: No principal internal prerequisite
recorded.

\textbf{Requires.} Exact mass law; finite symbol family; all words
through each depth.

\textbf{Boundary.} A finite stopping equality is sufficient, not
guaranteed to occur.

\textbf{Evidence.} Written proof. The related suite contains finite
block-level examples; no dedicated general algorithm is claimed for this
result.
\end{minipage}\par\medskip

\par\noindent\begin{minipage}{\textwidth}
\subsection{P39. Finite positive-lift
certificate}\label{LSMRegister:p39.-finite-positive-lift-certificate}

\textbf{Location and role.} Section 43, source lines 1778--1812.
Certificate. Principal dependencies: No principal internal prerequisite
recorded.

\textbf{Requires.} Rational/full-rank normalized generators; nonnegative
lifts; initial weights; exact intertwining and mass.

\textbf{Boundary.} Certifies the supplied model, not the generating
process.

\textbf{Code link.} \path{checks/cone_reference.py},
\path{verify_certificate}, lines 42--89. Exact supplied
finite-generator witness replay.
\end{minipage}\par\medskip

\par\noindent\begin{minipage}{\textwidth}
\subsection{P40. Robust finite-domain
certificate}\label{LSMRegister:p40.-robust-finite-domain-certificate}

\textbf{Location and role.} Section 44, source lines 1818--1851.
Certificate. Principal dependencies: P39.

\textbf{Requires.} One common domain for every complete
uncertainty-family vertex; exact mass constraints.

\textbf{Boundary.} Independent entry boxes can discard essential
uncertainty correlations.

\textbf{Code link.} \path{checks/cone_reference.py},
\path{verify_certificate}, lines 42--89. Vertex witnesses replayed
separately in the related suite.
\end{minipage}\par\medskip

\par\noindent\begin{minipage}{\textwidth}
\subsection{P41. Certified correction and its error
budget}\label{LSMRegister:p41.-certified-correction-and-its-error-budget}

\textbf{Location and role.} Section 45, source lines 1857--1882. Bound.
Principal dependencies: P39.

\textbf{Requires.} Fixed generator matrix; feasible correction;
separately justified error budgets.

\textbf{Boundary.} Existence of one valid fit does not show the true
model preserves that domain.

\textbf{Code link.} \path{checks/cone_reference.py},
\path{minimum_mixture_repair}, lines 154--176. Explicit simplicial
mixture sector, not general convex fitting.
\end{minipage}\par\medskip

\par\noindent\begin{minipage}{\textwidth}
\subsection{P42. Instrument contraction before
normalization}\label{LSMRegister:p42.-instrument-contraction-before-normalization}

\textbf{Location and role.} Section 46, source lines 1905--1938. Bound.
Principal dependencies: No principal internal prerequisite recorded.

\textbf{Requires.} Proper cone; strictly positive mass; cone-preserving
normalized instrument.

\textbf{Boundary.} Base norm is not an arbitrary Euclidean coefficient
norm.

\textbf{Evidence.} Written proof. The related suite contains finite
block-level examples; no dedicated general algorithm is claimed for this
result.
\end{minipage}\par\medskip

\par\noindent\begin{minipage}{\textwidth}
\subsection{P43. Observation-tree total-variation
bound}\label{LSMRegister:p43.-observation-tree-total-variation-bound}

\textbf{Location and role.} Section 47, source lines 1949--1975. Bound.
Principal dependencies: P42.

\textbf{Requires.} Two models valid on one domain; common policy; summed
instrument-error budgets.

\textbf{Boundary.} Full-law bound grows with horizon and need not remain
informative.

\textbf{Evidence.} Written proof. The related suite contains finite
block-level examples; no dedicated general algorithm is claimed for this
result.
\end{minipage}\par\medskip

\par\noindent\begin{minipage}{\textwidth}
\subsection{P44. Sharp event-probability dependence of
conditioning}\label{LSMRegister:p44.-sharp-event-probability-dependence-of-conditioning}

\textbf{Location and role.} Section 48, source lines 1979--2006. Bound.
Principal dependencies: No principal internal prerequisite recorded.

\textbf{Requires.} Two probability laws; positive evidence in both;
unconditional TV bound.

\textbf{Boundary.} Conditioning costs inverse event probability.

\textbf{Evidence.} Written proof. The related suite contains finite
block-level examples; no dedicated general algorithm is claimed for this
result.
\end{minipage}\par\medskip

\par\noindent\begin{minipage}{\textwidth}
\subsection{P45. Rational finite-certificate
search}\label{LSMRegister:p45.-rational-finite-certificate-search}

\textbf{Location and role.} Section 49, source lines 2010--2016. Search.
Principal dependencies: P39.

\textbf{Requires.} Rational inputs and existence of a rational
finite-generator certificate.

\textbf{Boundary.} Semidecision guarantee is not an efficient complete
implemented search.

\textbf{Code link.} \path{checks/cone_reference.py},
\path{search_interval}, lines 105--141. Bounded two-coordinate
interval search, not the general enumerator.
\end{minipage}\par\medskip

\par\noindent\begin{minipage}{\textwidth}
\subsection{P46. A rational rank-three model with no polyhedral
probability-domain
certificate}\label{LSMRegister:p46.-a-rational-rank-three-model-with-no-polyhedral-probability-domain-certificate}

\textbf{Location and role.} Section 51, source lines 2052--2101.
Obstruction. Principal dependencies: P38.

\textbf{Requires.} Declared rational rotation/probe model; closed convex
invariant domain containing its initial state.

\textbf{Boundary.} Rules out polyhedral domains here, not all finite
certificates.

\textbf{Evidence.} Written proof. The related suite contains finite
block-level examples; no dedicated general algorithm is claimed for this
result.
\end{minipage}\par\medskip

\par\noindent\begin{minipage}{\textwidth}
\subsection{P47. Polyhedral approximation of the rotation
model}\label{LSMRegister:p47.-polyhedral-approximation-of-the-rotation-model}

\textbf{Location and role.} Section 52, source lines 2109--2144. Bound.
Principal dependencies: P39, P43, P46.

\textbf{Requires.} Attenuated rotation and sufficient polygon size;
common Lorentz comparison domain.

\textbf{Boundary.} Damping changes the model and its observation law.

\textbf{Evidence.} Written proof. The related suite contains finite
block-level examples; no dedicated general algorithm is claimed for this
result.
\end{minipage}\par\medskip

\par\noindent\begin{minipage}{\textwidth}
\subsection{P48. Sequential confidence accounting for certified
proposals}\label{LSMRegister:p48.-sequential-confidence-accounting-for-certified-proposals}

\textbf{Location and role.} Section 53, source lines 2150--2164.
Confidence. Principal dependencies: No principal internal prerequisite
recorded.

\textbf{Requires.} Conditionally covered successive sets with a summable
failure budget.

\textbf{Boundary.} Does not establish resets, rank or dictionary
validity.

\textbf{Evidence.} Written proof. The related suite contains finite
block-level examples; no dedicated general algorithm is claimed for this
result.
\end{minipage}\par\medskip

\par\noindent\begin{minipage}{\textwidth}
\subsection{P49. Conditional learning with an all-horizon probability
certificate}\label{LSMRegister:p49.-conditional-learning-with-an-all-horizon-probability-certificate}

\textbf{Location and role.} Section 54, source lines 2170--2183.
Integration. Principal dependencies: P37, P39, P40, P43, P44.

\textbf{Requires.} Declared identification model; a covering uncertainty
set; robust witness for all members.

\textbf{Boundary.} A nominal witness alone establishes only internal
validity.

\textbf{Code link.} \path{checks/cone_reference.py},
\path{verify_certificate}, lines 42--89. Deterministic witness
portion; statistical assumptions are external.
\end{minipage}\par\medskip


\section{Curved domains and continuous
observations}\label{LSMRegister:curved-domains-and-continuous-observations}

Propositions 50--61. Related finite suite:
\path{checks/curved_continuous_checks.py}.

\par\noindent\begin{minipage}{\textwidth}
\subsection{P50. Matrix-embedding probability
certificate}\label{LSMRegister:p50.-matrix-embedding-probability-certificate}

\textbf{Location and role.} Section 56, source lines 2209--2246.
Certificate. Principal dependencies: No principal internal prerequisite
recorded.

\textbf{Requires.} Injective symmetric-matrix embedding; strict interior
witness; positive congruences; exact mass.

\textbf{Boundary.} Matrix embedding and positive lift are different
constructions.

\textbf{Code link.} \path{checks/continuous_reference.py},
\path{verify_certificate}, lines 135--181. Supplied
symmetric-matrix/congruence witness sector.
\end{minipage}\par\medskip

\par\noindent\begin{minipage}{\textwidth}
\subsection{P51. Exact rational matrix witness for the rotation
model}\label{LSMRegister:p51.-exact-rational-matrix-witness-for-the-rotation-model}

\textbf{Location and role.} Section 57, source lines 2258--2300.
Certificate. Principal dependencies: P46, P50.

\textbf{Requires.} The exact rational disk model and its explicit
matrix-unit witnesses.

\textbf{Boundary.} No attenuation or quantum-physical interpretation is
involved.

\textbf{Code link.} \path{checks/continuous_reference.py},
\path{verify_certificate}, lines 135--181. Exact two-by-two matrix
embedding example.
\end{minipage}\par\medskip

\par\noindent\begin{minipage}{\textwidth}
\subsection{P52. Continuous-observation instrument
semantics}\label{LSMRegister:p52.-continuous-observation-instrument-semantics}

\textbf{Location and role.} Section 58, source lines 2308--2338.
Semantics. Principal dependencies: No principal internal prerequisite
recorded.

\textbf{Requires.} Standard Borel observations; common probability
reference measure; integrable positive operator density.

\textbf{Boundary.} Point densities are version dependent; zero branches
have no model-determined posterior.

\textbf{Code link.} \path{checks/continuous_reference.py},
\path{VerifiedInstrument}, lines 114--132. Finite polynomial reference
class only.
\end{minipage}\par\medskip

\par\noindent\begin{minipage}{\textwidth}
\subsection{P53. Polynomial Gram certificate for interval
observations}\label{LSMRegister:p53.-polynomial-gram-certificate-for-interval-observations}

\textbf{Location and role.} Section 59, source lines 2346--2388.
Certificate. Principal dependencies: P50, P52.

\textbf{Requires.} Fixed embedding; positive polynomial Gram witness;
exact coefficients and integrated mass.

\textbf{Boundary.} A sufficient certificate, not completeness for all
positive maps.

\textbf{Code link.} \path{checks/continuous_reference.py},
\path{verify_certificate}, lines 135--181. Sufficient polynomial
positivity certificate replay.
\end{minipage}\par\medskip

\par\noindent\begin{minipage}{\textwidth}
\subsection{P54. Conditional finite-moment
identification}\label{LSMRegister:p54.-conditional-finite-moment-identification}

\textbf{Location and role.} Section 60, source lines 2403--2448.
Construction. Principal dependencies: No principal internal prerequisite
recorded.

\textbf{Requires.} Correct finite predictive dimension and degree;
invertible complete event core; exact moment tables.

\textbf{Boundary.} Rank and degree are supplied assumptions, not
inferred from an inverse.

\textbf{Code link.} \path{checks/continuous_reference.py},
\path{fit_population_moments}, lines 184--210. Exact event/moment
tables with supplied rank and degree.
\end{minipage}\par\medskip

\par\noindent\begin{minipage}{\textwidth}
\subsection{P55. A valid continuous counterexample to finite-moment
identification}\label{LSMRegister:p55.-a-valid-continuous-counterexample-to-finite-moment-identification}

\textbf{Location and role.} Section 61, source lines 2460--2479.
Obstruction. Principal dependencies: No principal internal prerequisite
recorded.

\textbf{Requires.} Uniform reference measure; declared two polynomial
densities with matching low moments.

\textbf{Boundary.} Moment agreement and validity do not identify the
entire density.

\textbf{Evidence.} Written proof. The related suite contains finite
block-level examples; no dedicated general algorithm is claimed for this
result.
\end{minipage}\par\medskip

\par\noindent\begin{minipage}{\textwidth}
\subsection{P56. Exact bin coarsening and its
posterior}\label{LSMRegister:p56.-exact-bin-coarsening-and-its-posterior}

\textbf{Location and role.} Section 62, source lines 2483--2500.
Identity. Principal dependencies: P52.

\textbf{Requires.} Exact cell integration; a measurable partition;
bin-measurable policy.

\textbf{Boundary.} Does not preserve a policy responding to within-bin
values.

\textbf{Evidence.} Written proof. The related suite contains finite
block-level examples; no dedicated general algorithm is claimed for this
result.
\end{minipage}\par\medskip

\par\noindent\begin{minipage}{\textwidth}
\subsection{P57. Exact covariance correction for an averaged
congruence}\label{LSMRegister:p57.-exact-covariance-correction-for-an-averaged-congruence}

\textbf{Location and role.} Section 63, source lines 2506--2522.
Identity. Principal dependencies: No principal internal prerequisite
recorded.

\textbf{Requires.} Square-integrable matrix function; a probability
averaging measure.

\textbf{Boundary.} Average congruence differs from congruence of the
average.

\textbf{Evidence.} Written proof. The related suite contains finite
block-level examples; no dedicated general algorithm is claimed for this
result.
\end{minipage}\par\medskip

\par\noindent\begin{minipage}{\textwidth}
\subsection{P58. Positive cell-average
approximation}\label{LSMRegister:p58.-positive-cell-average-approximation}

\textbf{Location and role.} Section 64, source lines 2534--2567. Bound.
Principal dependencies: P52, P56.

\textbf{Requires.} Finite positive-measure cells; Lipschitz density in
the base norm.

\textbf{Boundary.} This is a declared density approximation, not exact
continuous-law equality.

\textbf{Evidence.} Written proof. The related suite contains finite
block-level examples; no dedicated general algorithm is claimed for this
result.
\end{minipage}\par\medskip

\par\noindent\begin{minipage}{\textwidth}
\subsection{P59. Continuous-record total-variation
bound}\label{LSMRegister:p59.-continuous-record-total-variation-bound}

\textbf{Location and role.} Section 65, source lines 2587--2624. Bound.
Principal dependencies: P42, P52.

\textbf{Requires.} Common cone, mass and reference measure; common
policy; integrated error budgets.

\textbf{Boundary.} No uniform posterior contraction is supplied.

\textbf{Evidence.} Written proof. The related suite contains finite
block-level examples; no dedicated general algorithm is claimed for this
result.
\end{minipage}\par\medskip

\par\noindent\begin{minipage}{\textwidth}
\subsection{P60. Exact recognition of a fixed certified
template}\label{LSMRegister:p60.-exact-recognition-of-a-fixed-certified-template}

\textbf{Location and role.} Section 67, source lines 2753--2770.
Construction. Principal dependencies: P53.

\textbf{Requires.} A fixed certified template; coefficient intertwining;
invertible alignment.

\textbf{Boundary.} Template recognition is not domain discovery.

\textbf{Code link.} \path{checks/continuous_reference.py},
\path{recognize_template}, lines 213--236. Fixed-template
recognition, not geometry discovery.
\end{minipage}\par\medskip

\par\noindent\begin{minipage}{\textwidth}
\subsection{P61. Conditional continuous realization with a finite domain
certificate}\label{LSMRegister:p61.-conditional-continuous-realization-with-a-finite-domain-certificate}

\textbf{Location and role.} Section 68, source lines 2778--2786.
Integration. Principal dependencies: P52, P53, P54, P56, P58, P59, P60.

\textbf{Requires.} Correct continuous class and event core; an
independent exact domain witness.

\textbf{Boundary.} No unmodeled measurement error or unknown-degree
inference is covered.

\textbf{Code link.} \path{checks/continuous_reference.py},
\path{fit_population_moments}, lines 184--210. Worked continuous
polynomial sector only.
\end{minipage}\par\medskip


\section{Discovery, approximation and
evidence}\label{LSMRegister:discovery-approximation-and-evidence}

Propositions 62--70. Related finite suite:
\path{checks/discovery_evidence_checks.py}.

\par\noindent\begin{minipage}{\textwidth}
\subsection{P62. Oriented quadratic probability
certificate}\label{LSMRegister:p62.-oriented-quadratic-probability-certificate}

\textbf{Location and role.} Section 70, source lines 2806--2858.
Certificate. Principal dependencies: P52.

\textbf{Requires.} Positive definite ellipsoidal metric; oriented
quadratic and polynomial positivity witnesses.

\textbf{Boundary.} Preserving a double cone without mass orientation is
insufficient.

\textbf{Code link.} \path{checks/quadratic_reference.py},
\path{verify_quadratic}, lines 108--130. Exact oriented quadratic
witness.
\end{minipage}\par\medskip

\par\noindent\begin{minipage}{\textwidth}
\subsection{P63. Conditional operator-only ellipsoid
discovery}\label{LSMRegister:p63.-conditional-operator-only-ellipsoid-discovery}

\textbf{Location and role.} Section 71, source lines 2864--2894. Search.
Principal dependencies: P62.

\textbf{Requires.} Suitable mass-preserving ratio; unique invariant-form
direction; nonzero rank-one image scale.

\textbf{Boundary.} Every geometry proposal still requires full-family
verification.

\textbf{Code link.} \path{checks/quadratic_reference.py},
\path{discover_quadratic}, lines 151--197. Restricted
invariant-form/rank-one-image discovery.
\end{minipage}\par\medskip

\par\noindent\begin{minipage}{\textwidth}
\subsection{P64. Valid positive polynomial
approximation}\label{LSMRegister:p64.-valid-positive-polynomial-approximation}

\textbf{Location and role.} Section 72, source lines 2928--2997. Bound.
Principal dependencies: P59.

\textbf{Requires.} A valid Lipschitz density on a common cone; doubly
stochastic positive kernel.

\textbf{Boundary.} Avoids an exact-degree hypothesis but still assumes
regularity and a domain.

\textbf{Code link.} \path{checks/sequential_reference.py},
\path{positive_coefficients}, lines 30--34. Positive approximation in
a finite symbolic example.
\end{minipage}\par\medskip

\par\noindent\begin{minipage}{\textwidth}
\subsection{P65. Certified coefficient error and approximation
bias}\label{LSMRegister:p65.-certified-coefficient-error-and-approximation-bias}

\textbf{Location and role.} Section 73, source lines 3027--3062. Bound.
Principal dependencies: P64.

\textbf{Requires.} Positive estimated Bernstein coefficients; exact
integrated mass; covered coefficient errors.

\textbf{Boundary.} Increasing polynomial degree does not automatically
reduce estimation error.

\textbf{Code link.} \path{checks/sequential_reference.py},
\path{reconstruct_bernstein}, lines 37--40. Finite coefficient
reconstruction; no unknown-basis learner.
\end{minipage}\par\medskip

\par\noindent\begin{minipage}{\textwidth}
\subsection{P66. Time-uniform predictable-residual
certificate}\label{LSMRegister:p66.-time-uniform-predictable-residual-certificate}

\textbf{Location and role.} Section 74, source lines 3068--3105.
Confidence. Principal dependencies: No principal internal prerequisite
recorded.

\textbf{Requires.} Prospectively registered predictable bounded tests
and a summable budget.

\textbf{Boundary.} Nonrejection is not an upper error or correctness
guarantee.

\textbf{Code link.} \path{checks/sequential_reference.py},
\path{anytime_radius}, lines 43--47. Declared elementary time-uniform
radius.
\end{minipage}\par\medskip

\par\noindent\begin{minipage}{\textwidth}
\subsection{P67. Covered adaptive-action
evidence}\label{LSMRegister:p67.-covered-adaptive-action-evidence}

\textbf{Location and role.} Section 75, source lines 3111--3129.
Confidence. Principal dependencies: P66.

\textbf{Requires.} Known predictable randomized actions with positive
coverage.

\textbf{Boundary.} Estimates contextual action means, not
unobserved-policy long-run distributions.

\textbf{Evidence.} Written proof. The related suite contains finite
block-level examples; no dedicated general algorithm is claimed for this
result.
\end{minipage}\par\medskip

\par\noindent\begin{minipage}{\textwidth}
\subsection{P68. Sequential evidence gives a lower error
certificate}\label{LSMRegister:p68.-sequential-evidence-gives-a-lower-error-certificate}

\textbf{Location and role.} Section 76, source lines 3135--3156. Bound.
Principal dependencies: P66.

\textbf{Requires.} Actual common-history observation laws and registered
tests.

\textbf{Boundary.} A lower discrepancy certificate cannot be used as an
upper accuracy bound.

\textbf{Code link.} \path{checks/sequential_reference.py},
\path{dependent_example}, lines 50--69. One seeded
conditional-mismatch illustration.
\end{minipage}\par\medskip

\par\noindent\begin{minipage}{\textwidth}
\subsection{P69. Observation-law error without a shared state
space}\label{LSMRegister:p69.-observation-law-error-without-a-shared-state-space}

\textbf{Location and role.} Section 77, source lines 3170--3194. Bound.
Principal dependencies: No principal internal prerequisite recorded.

\textbf{Requires.} Common observation/action space and policy;
conditional TV discrepancies on true histories.

\textbf{Boundary.} Allows different internal state spaces but still
needs true discrepancy control.

\textbf{Evidence.} Written proof. The related suite contains finite
block-level examples; no dedicated general algorithm is claimed for this
result.
\end{minipage}\par\medskip

\par\noindent\begin{minipage}{\textwidth}
\subsection{P70. Separate structural acceptance, approximation control
and sequential
rejection}\label{LSMRegister:p70.-separate-structural-acceptance-approximation-control-and-sequential-rejection}

\textbf{Location and role.} Section 78, source lines 3200--3210.
Integration. Principal dependencies: P62, P63, P64, P65, P66, P67, P68,
P69.

\textbf{Requires.} Separate structural, approximation and
sequential-evidence assumptions.

\textbf{Boundary.} Acceptance, accuracy and statistical nonrejection are
not equivalent.

\textbf{Evidence.} Written proof. The related suite contains finite
block-level examples; no dedicated general algorithm is claimed for this
result.
\end{minipage}\par\medskip


\section{Robust curved certificates}\label{LSMRegister:robust-curved-certificates}

Propositions 71--79. Related finite suite:
\path{checks/robust_curved_checks.py}.

\par\noindent\begin{minipage}{\textwidth}
\subsection{P71. Robust oriented
margin}\label{LSMRegister:p71.-robust-oriented-margin}

\textbf{Location and role.} Section 80, source lines 3232--3267.
Certificate. Principal dependencies: P52.

\textbf{Requires.} Common transformed domain; uniform oriented margin
and output-error bounds; exact mass.

\textbf{Boundary.} Margins and uncertainty must be expressed in the same
coordinates.

\textbf{Code link.} \path{checks/robust_reference.py},
\path{replay_certificate}, lines 217--251. Uniform margin certificate
and uncertainty-budget check.
\end{minipage}\par\medskip

\par\noindent\begin{minipage}{\textwidth}
\subsection{P72. Boundary fragility under mass preserving
perturbations}\label{LSMRegister:p72.-boundary-fragility-under-mass-preserving-perturbations}

\textbf{Location and role.} Section 81, source lines 3271--3283.
Obstruction. Principal dependencies: No principal internal prerequisite
recorded.

\textbf{Requires.} At least one spatial coordinate; a cone-positive map
with a boundary image.

\textbf{Boundary.} Mass-only dimension zero is explicitly excluded in
0.9.0.

\textbf{Evidence.} Written proof. The related suite contains finite
block-level examples; no dedicated general algorithm is claimed for this
result.

\textbf{Correction.} The mass-only case has no nonzero mass-preserving
perturbation. The added dimension hypothesis excludes that case. The
inherited exact regression checks cover both sides of this boundary.
\end{minipage}\par\medskip

\par\noindent\begin{minipage}{\textwidth}
\subsection{P73. Interior repair with an explicit error
budget}\label{LSMRegister:p73.-interior-repair-with-an-explicit-error-budget}

\textbf{Location and role.} Section 82, source lines 3292--3340. Bound.
Principal dependencies: P59, P71.

\textbf{Requires.} A known valid reference; center-reset mixture; stated
discrepancy and mass conditions.

\textbf{Boundary.} The repair changes the model; its uncertainty tube
need not contain the original truth.

\textbf{Code link.} \path{checks/robust_reference.py},
\path{make_certificate}, lines 197--214. Explicit center-reset repair
witness construction.
\end{minipage}\par\medskip

\par\noindent\begin{minipage}{\textwidth}
\subsection{P74. A posteriori proposal
error}\label{LSMRegister:p74.-a-posteriori-proposal-error}

\textbf{Location and role.} Section 83, source lines 3346--3405. Bound.
Principal dependencies: No principal internal prerequisite recorded.

\textbf{Requires.} A true normalized invariant metric; a separated
linear system; covered proposal perturbations.

\textbf{Boundary.} A small numerical residual alone is not acceptance or
coverage.

\textbf{Code link.} \path{checks/robust_reference.py},
\path{propose_geometry}, lines 152--194. Numerical proposal only; not
the theorem's unknown true-metric assumptions.
\end{minipage}\par\medskip

\par\noindent\begin{minipage}{\textwidth}
\subsection{P75. Rational semicircle and Bernstein
certificate}\label{LSMRegister:p75.-rational-semicircle-and-bernstein-certificate}

\textbf{Location and role.} Section 84, source lines 3411--3444.
Certificate. Principal dependencies: No principal internal prerequisite
recorded.

\textbf{Requires.} Three-coordinate rational polynomial family; both
boundary charts; complete subdivision.

\textbf{Boundary.} A negative Bernstein coefficient alone is not a
counterexample.

\textbf{Code link.} \path{checks/robust_reference.py},
\path{replay_polynomial}, lines 117--139. Exact complete
tensor-Bernstein subdivision replay.
\end{minipage}\par\medskip

\par\noindent\begin{minipage}{\textwidth}
\subsection{P76. Coefficient boxes and convex
uncertainty}\label{LSMRegister:p76.-coefficient-boxes-and-convex-uncertainty}

\textbf{Location and role.} Section 85, source lines 3450--3475.
Certificate. Principal dependencies: P71.

\textbf{Requires.} Coefficient boxes in declared Lorentz coordinates;
exact mass affine constraints.

\textbf{Boundary.} Jointly optimizing unknown geometry and operators is
not asserted convex.

\textbf{Code link.} \path{checks/robust_reference.py},
\path{replay_certificate}, lines 217--251. Coefficient-box inclusion
in declared coordinates.
\end{minipage}\par\medskip

\par\noindent\begin{minipage}{\textwidth}
\subsection{P77. Conditional search for strict rational curved
certificates}\label{LSMRegister:p77.-conditional-search-for-strict-rational-curved-certificates}

\textbf{Location and role.} Section 86, source lines 3481--3487. Search.
Principal dependencies: P75, P76.

\textbf{Requires.} A real transform with strictly interior initial
states and a common strict margin.

\textbf{Boundary.} Conditional existence of rational witnesses is not an
implemented general optimizer.

\textbf{Code link.} \path{checks/robust_reference.py},
\path{make_certificate}, lines 197--214. Finite candidate
construction; no complete general search implemented.
\end{minipage}\par\medskip

\par\noindent\begin{minipage}{\textwidth}
\subsection{P78. Data-selected geometry with honest
coverage}\label{LSMRegister:p78.-data-selected-geometry-with-honest-coverage}

\textbf{Location and role.} Section 87, source lines 3493--3507.
Confidence. Principal dependencies: No principal internal prerequisite
recorded.

\textbf{Requires.} One covered fixed-coordinate operator/initial-state
set; verifier quantifies over its entire set.

\textbf{Boundary.} Data-selected geometry must transport uncertainty,
not just the fitted center.

\textbf{Code link.} \path{checks/robust_reference.py},
\path{replay_certificate}, lines 217--251. Deterministic set
certificate only, not upstream statistical coverage.
\end{minipage}\par\medskip

\par\noindent\begin{minipage}{\textwidth}
\subsection{P79. Certified repair with a retained comparison
budget}\label{LSMRegister:p79.-certified-repair-with-a-retained-comparison-budget}

\textbf{Location and role.} Section 89, source lines 3564--3595.
Integration. Principal dependencies: P59, P71, P73, P75, P76.

\textbf{Requires.} Worked deterministic perturbed family; nominal and
comparison-domain witnesses.

\textbf{Boundary.} The original repair's bound is trivial at horizon
eight.

\textbf{Code link.} \path{checks/robust_reference.py},
\path{replay_certificate}, lines 217--251. Nominal and
comparison-domain witnesses for one deterministic example.
\end{minipage}\par\medskip


\section{Data-supported domains}\label{LSMRegister:data-supported-domains}

Propositions 80--87. Related finite suite:
\path{checks/data_supported_checks.py}.

\par\noindent\begin{minipage}{\textwidth}
\subsection{P80. Convex repair
bracketing}\label{LSMRegister:p80.-convex-repair-bracketing}

\textbf{Location and role.} Section 91, source lines 3617--3636. Bound.
Principal dependencies: P75, P79.

\textbf{Requires.} Fixed geometry, reset path and requested margin;
exact separator and upper witness.

\textbf{Boundary.} Brackets this repair path, not optimal geometry or
all repair methods.

\textbf{Code link.} \path{checks/robust_reference.py},
\path{replay_certificate}, lines 217--251. Upper certificate; exact
lower separator is in the related suite.
\end{minipage}\par\medskip

\par\noindent\begin{minipage}{\textwidth}
\subsection{P81. Time-uniform regression
confidence}\label{LSMRegister:p81.-time-uniform-regression-confidence}

\textbf{Location and role.} Section 92, source lines 3686--3743.
Confidence. Principal dependencies: No principal internal prerequisite
recorded.

\textbf{Requires.} Predictable regressors; correct fixed conditional
mean; positive sub-Gaussian scale; parameter bound.

\textbf{Boundary.} Estimated hidden regressors and response-dependent
features need new analysis.

\textbf{Code link.} \path{checks/data_certificate_reference.py},
\path{radius}, lines 50--54. Scalar specialization only, not general
vector online regression.
\end{minipage}\par\medskip

\par\noindent\begin{minipage}{\textwidth}
\subsection{P82. Sequential observable-moment operator
sets}\label{LSMRegister:p82.-sequential-observable-moment-operator-sets}

\textbf{Location and role.} Section 93, source lines 3749--3786.
Confidence. Principal dependencies: P81.

\textbf{Requires.} Exact accessible predictive state; correct finite
basis; known invertible moment matrix.

\textbf{Boundary.} General multirow theorem is not a hidden-state
discovery algorithm.

\textbf{Code link.} \path{checks/data_certificate_reference.py},
\path{sufficient_statistics}, lines 59--71. Observable scalar
specialization; general matrix learner not implemented.
\end{minipage}\par\medskip

\par\noindent\begin{minipage}{\textwidth}
\subsection{P83. Structured curved-domain
certificate}\label{LSMRegister:p83.-structured-curved-domain-certificate}

\textbf{Location and role.} Section 94, source lines 3794--3852.
Certificate. Principal dependencies: No principal internal prerequisite
recorded.

\textbf{Requires.} Specified moment feature curve; known reference
moments; two-parameter conditional family.

\textbf{Boundary.} Geometry, dictionary and family are supplied, while
parameters can be learned.

\textbf{Code link.} \path{checks/data_certificate_reference.py},
\path{probability_row}, lines 30--33. Fixed two-parameter family and
five-node measure.
\end{minipage}\par\medskip

\par\noindent\begin{minipage}{\textwidth}
\subsection{P84. Exact rational outer confidence
intervals}\label{LSMRegister:p84.-exact-rational-outer-confidence-intervals}

\textbf{Location and role.} Section 95, source lines 3858--3892.
Confidence. Principal dependencies: P81.

\textbf{Requires.} Exact observed node path; declared measure and prior
bounds; rational outward rounding.

\textbf{Boundary.} Replay cannot authenticate data or prove model
assumptions from one path.

\textbf{Code link.} \path{checks/data_certificate_reference.py},
\path{confidence_box}, lines 73--86. Exact rational scalar intervals
and outward rounding.
\end{minipage}\par\medskip

\par\noindent\begin{minipage}{\textwidth}
\subsection{P85. Parameter-to-law error without a state-estimation
term}\label{LSMRegister:p85.-parameter-to-law-error-without-a-state-estimation-term}

\textbf{Location and role.} Section 96, source lines 3922--3956. Bound.
Principal dependencies: P69, P83.

\textbf{Requires.} Covered parameter distances; common measure and
common observed history.

\textbf{Boundary.} Likelihood error remains even when normalized
posterior states agree.

\textbf{Code link.} \path{checks/data_certificate_reference.py},
\path{make_certificate}, lines 91--112. Finite parameter rectangle
and declared measure constants.
\end{minipage}\par\medskip

\par\noindent\begin{minipage}{\textwidth}
\subsection{P86. Uniform marginal and rolling-window
stability}\label{LSMRegister:p86.-uniform-marginal-and-rolling-window-stability}

\textbf{Location and role.} Section 97, source lines 3970--3990. Bound.
Principal dependencies: P69.

\textbf{Requires.} Markov kernels; positive common minorization; uniform
one-step discrepancy.

\textbf{Boundary.} Uniform marginal control is not uniform
complete-record control.

\textbf{Evidence.} Written proof. The related suite contains finite
block-level examples; no dedicated general algorithm is claimed for this
result.
\end{minipage}\par\medskip

\par\noindent\begin{minipage}{\textwidth}
\subsection{P87. Observable structured learning with a covered
domain}\label{LSMRegister:p87.-observable-structured-learning-with-a-covered-domain}

\textbf{Location and role.} Section 98, source lines 4007--4013.
Integration. Principal dependencies: P81, P83, P84, P85, P86.

\textbf{Requires.} Correct observable family; predictable protocol; full
confidence rectangle and valid chosen center.

\textbf{Boundary.} Does not transfer statistical coverage to the
unrelated repaired-operator example.

\textbf{Code link.} \path{checks/data_certificate_reference.py},
\path{replay_certificate}, lines 114--128. Data fingerprint,
statistics, radii and whole-rectangle validity; assumptions not
inferred.
\end{minipage}\par\medskip

\endgroup
\chapter{Large Signal Model: visual atlas}\pdfbookmark[1]{How to read the atlas}{route}
\textcolor{teal}{\small\textsc{Orientation}}\par
{\fontsize{23}{26}\selectfont Four contracts, one model}\par
An observation, its archive, a predictive summary and a probability state carry different mathematical obligations. The theory connects them through maps whose domains and identities are explicit.
\LSMAtlasdiagram{00_architecture}
\begin{tabular}{@{}p{.19\linewidth}p{.32\linewidth}p{.42\linewidth}@{}}
\toprule
Contract & Defining identity & Question answered\\\midrule
Reconstruction & $SA=I_X$ & Can the accepted observation be recovered?\\[5pt]
Transport & $F_bT_{ba}=T_{ba}F_a$ & Do coordinate descriptions agree?\\[5pt]
Closure & $\gamma F_u=\bar F_u\gamma$ & Does the summary determine its own future?\\[5pt]
Probability & $B_oC\subseteq C$, $\lambda\sum_oB_o=\lambda$ & Are every branch and every finite record valid?\\\bottomrule
\end{tabular}
\vspace{10pt}
\textbf{Reading route.} Plates 1--4 build the representation and its predictive quotient. Plates 5--9 explain interactions, lost paths, memory and observable identification. Plates 10--13 develop behavior spaces, positivity, curved geometry and probability integration. Formal routes accompany the plates.

\textbf{Visual conventions.} Teal marks exact structure. Gold marks transported, refined or uncertain quantities. Red marks an obstruction or an omitted contribution. Every figure is drawn from a formula or an exact finite example. Curves of bounds are labeled as bounds.

\LSMAtlasboundary{The atlas selects core mechanisms from the 87 registered propositions. It introduces no new numbered theorem, fitted model, benchmark result or empirical coverage claim. Its diagrams supplement the proofs.}

\LSMAtlasplate{1}{Realizable coefficients form a subspace}{Exact storage includes a consistency condition.}{P1}
For linear analysis and synthesis maps $A:X\to C$ and $S:C\to X$, the left-inverse law $SA=I_X$ gives injectivity of $A$. The coefficient projector is $\Pi=AS$, with
\[
\Pi^2=\Pi,\qquad \LSMAtlasim\Pi=\LSMAtlasim A,\qquad c\text{ is realizable}\iff \Pi c=c.
\]
\LSMAtlasdiagram[0.90]{01_consistency}
\textbf{A redundant archive.} Take $X=\LSMAtlasR^2$, $C=\LSMAtlasR^3$ and
\[
A(x_1,x_2)=(x_1,x_2,x_1+x_2),\qquad S(c_1,c_2,c_3)=(c_1,c_2).
\]
Then $\Pi(c_1,c_2,c_3)=(c_1,c_2,c_1+c_2)$. The shaded plane is the set of realizable coefficients. The red segment shows synthesis followed by reanalysis of an inconsistent vector.
\LSMAtlasclaim{A left inverse determines an idempotent consistency map. Here it is oblique: $\Pi$ is not self-adjoint. Orthogonality requires an additional hypothesis.}
\LSMAtlasboundary{For bounded linear maps on normed spaces, $\|S\|$ controls amplification of coefficient errors. Exact invertibility by itself gives no conditioning guarantee.}
\LSMAtlasnav{\texttt{Realization.lean}; norm estimates in \texttt{Normed.lean}.}

\LSMAtlasplate{2}{One operation, several coordinates}{A commuting square expresses coordinate independence.}{P2}
Let $\mathcal C_a=\LSMAtlasim A_a$ and $T_{ba}=A_bS_a$. Coordinate realizations are $F_a=A_aFS_a$ on $\mathcal C_a$. Both paths through the square represent the same underlying operation.
\LSMAtlasdiagram{02_transport}
\[
T_{cb}T_{ba}=T_{ca},\qquad T_{aa}=I_{\mathcal C_a},\qquad
F_bT_{ba}=T_{ba}F_a.
\]
Conversely, a compatible family $G_a$ on the realizable coefficient spaces determines one operation. With a supplied reference chart $r$,
\[
F=S_rG_rA_r,\qquad G_a=A_aFS_a.
\]
\LSMAtlasclaim{Compatibility is a descent condition: local coordinate descriptions belong to one object-space operation precisely when they agree under transport.}
A fitted family's discrepancy can be measured in object space:
\[
S_b\bigl(G_bA_bx-T_{ba}G_aA_ax\bigr)
=S_bG_bA_bx-S_aG_aA_ax.
\]
\LSMAtlasboundary{Zero chart discrepancy establishes agreement. It does not establish correctness of the agreed predictor. Deterministic recodings of the same observation supply the same observed information.}
\LSMAtlasnav{\texttt{Realization.lean}. The formal construction uses a nonempty chart family.}

\LSMAtlasplate{3}{The obstruction lives in a fiber}{Equal present summaries must have equal next summaries.}{P6--P7}
A reduced map $\bar F:\gamma(X)\to\gamma(X)$ exists exactly when
\[
\gamma x=\gamma x'\quad\Longrightarrow\quad\gamma Fx=\gamma Fx'.
\]
For linear $L$ and linear $\gamma$, this is $L(\ker\gamma)\subseteq\ker\gamma$.
\LSMAtlasdiagram{03_fibers}
\textbf{The running rational example.} Throughout the deterministic plates, use
\[
L=\begin{pmatrix}1/2&1\\1/4&1/3\end{pmatrix},\qquad
\gamma=(1,0),\qquad x_0=(1,0)^{\mathsf T}.
\]
The two states $(0,0)$ and $(0,1)$ have summary zero. Their next summaries are $0$ and $1$. Thus no scalar autonomous update reproduces every state's first-coordinate dynamics.
\LSMAtlasclaim{Predictive sufficiency is a property of entire summary fibers. A lossless archive of each full state does not make its chosen scalar summary sufficient.}
\LSMAtlasboundary{The diagram is a counterexample for the declared operator. In a stochastic model, the corresponding condition concerns equality of next-summary distributions.}
\LSMAtlasnav{\texttt{PredictiveClosure.lean}; \texttt{LinearClosure.lean}.}

\LSMAtlasplate{4}{Refinement closes an observable space}{Add the directions generated by known dynamics.}{P8, P14}
For a known finite family $L_u$ on a $d$-dimensional space, start from the rows of a required observation $C_0$ and form
\[
W_{n+1}=W_n+\sum_uW_nL_u,\qquad W_0=\operatorname{row}(C_0).
\]
\LSMAtlasdiagram{04_refinement}
Every strict inclusion raises dimension. After at most $d-\dim W_0$ strict enlargements, the limit $W_*$ is the least common invariant row space containing $W_0$.

For the running example, the rows of
\[
O=\begin{pmatrix}1&0\\1/2&1\end{pmatrix},\qquad
D=OLO^{-1}=\begin{pmatrix}0&1\\1/12&5/6\end{pmatrix}
\]
span $W_*=\LSMAtlasR^{2\,*}$. The state $Ox_s=(y_s,y_{s+1})^{\mathsf T}$ therefore closes.
\LSMAtlasclaim{The added coordinate is forced by the action of $L$ on the requested observable. In a larger system, unobserved directions unrelated to its future can remain outside $W_*$.}
\LSMAtlasboundary{Computing $Ox_s$ uses a forward observation. It becomes an available delay state $(y_{t-1},y_t)$ only after a one-step warm-up. Known-operator completion does not establish sensor access to arbitrary new observables.}
\LSMAtlasnav{\texttt{InvariantCompletion.lean}; \texttt{PredictiveCompletion.lean}; \texttt{DelayClosure.lean}.}

\LSMAtlasplate{5}{Packets decompose operations too}{Resolution changes preserve identities when every block is retained.}{P3--P5, P20}
For complete input and output projector families,
\[
x=\sum_iP_ix,\qquad L=\sum_{j,i}Q_jLP_i.
\]
Refining either family subdivides these transfer blocks without changing their total. Orthogonal packet families also give the corresponding squared-norm decomposition.
\LSMAtlasdiagram{05_packets}
For a polynomial operation, multilinearity produces every ordered packet interaction:
\[
Q_jB_m(x,\ldots,x)=\sum_{i_1,\ldots,i_m}
Q_jB_m(P_{i_1}x,\ldots,P_{i_m}x).
\]
In the scalar square, the mixed contribution is $2uv$. With $u=1$, it is linear in omitted amplitude $v$, whereas the omitted packet square is quadratic.
\LSMAtlasclaim{A nonlinear operation acts on interactions between packets. Reconstruction energy alone does not bound an omitted packet's influence on a future quantity.}
\LSMAtlasboundary{Complete refinement and truncation are different operations. Truncation requires an explicit omitted-interaction estimate on its declared domain. Projectors from different atlases need not commute.}
\LSMAtlasnav{\texttt{PacketTree.lean}; \texttt{TreeEnergy.lean}; \texttt{Transfer.lean}; \texttt{Polynomial.lean}.}

\LSMAtlasplate{6}{Compression can erase a return path}{A two-step effect can survive when both compressed steps vanish.}{P7}
Let $P^2=P$ and $Q=I-P$. Inserting $I=P+Q$ gives the exact defect
\[
\boxed{PLMP-(PLP)(PMP)=PLQMP.}
\]
\LSMAtlasdiagram{06_return_path}
Take $P=\LSMAtlasdiag(1,0)$ and
\[
L=M=\begin{pmatrix}0&1\\1&0\end{pmatrix}.
\]
Then $PLP=PMP=0$, while $PLMP=P$. The missing term $PLQMP=P$ is the full return through the discarded coordinate.

The running example exhibits the same mechanism in $PL^2P$:
\[
PL^2P-(PLP)^2
=\begin{pmatrix}1/4&0\\0&0\end{pmatrix}.
\]
Starting from $(1,0)$, separate scalar compression predicts $y_2=1/4$. The exact value is $1/2$.
\LSMAtlasclaim{The obstruction is an omitted route through the operator composition. Its size depends jointly on outgoing and returning transfer blocks.}
\LSMAtlasboundary{The algebra requires idempotence. An orthogonal split additionally gives the familiar Hilbert-space energy interpretation.}
\LSMAtlasnav{\texttt{LinearClosure.lean}; norm bounds in \texttt{Normed.lean}.}

\LSMAtlasplate{7}{Memory is an exact elimination}{The discarded state leaves a kernel and an initial transient.}{P10, P19}
Write the split dynamics as $y_{t+1}=Ay_t+Bw_t$ and $w_{t+1}=Cy_t+Dw_t$. Eliminating $w_t$ gives
\[
y_{t+1}=Ay_t+\sum_{j=0}^{t-1}BD^jC\,y_{t-1-j}+BD^tw_0.
\]
\LSMAtlasdiagram{07_memory}
For the running example, $A=1/2$, $B=1$, $C=1/4$, $D=1/3$ and $w_0=0$. Thus $K_j=\tfrac14\,3^{-j}$. If $|y_t|\leq1$, retaining the last $H$ history terms leaves a same-history tail at most $\Delta_H=\tfrac38\,3^{-H}$.

Here $0\leq y_t\leq1$ and the truncated recursion has gain
\[
a_H=\frac12+\sum_{j=0}^{H-1}K_j
=\frac78-\frac38\,3^{-H}<1.
\]
With the same initial state, recursive propagation therefore gives
\[
\sup_t|y_t-\widehat y_t|\leq\frac{\Delta_H}{1-a_H}
=\frac{3}{3^H+3}.
\]
\LSMAtlasboundary{P10's tail estimate uses the same resolved history. P19 supplies the separate stability argument for a predictor using its own history. For other initial states, retain the term $BD^tw_0$.}
\LSMAtlasnav{\texttt{Memory.lean}; \texttt{MemoryBounds.lean}; \texttt{MemoryStability.lean}.}

\LSMAtlasplate{8}{A local defect needs a propagation law}{Small one-step residuals can accumulate or contract.}{P18}
Allow a representation $\gamma_t$ to change with time. Define the typed defect
\[
r_t(x)=\gamma_{t+1}F_t(x)-\widehat f_t\gamma_t(x).
\]
If $\widehat f_t$ is $\lambda_t$-Lipschitz on the relevant domain and $\|r_t(x_t)\|\leq\delta_t$, the prediction error satisfies
\[
e_T\leq\left(\prod_{j=0}^{T-1}\lambda_j\right)e_0
+\sum_{k=0}^{T-1}\left(\prod_{j=k+1}^{T-1}\lambda_j\right)\delta_k.
\]
\LSMAtlasdiagram[0.94]{08_defect}
For a fixed representation $R$, a reference reduced operator $B$ and fitted operator $\widehat B$,
\[
RL-\widehat BR=(RL-BR)+(B-\widehat B)R.
\]
The first term is representational nonclosure. The second is estimation error. They can enter the same propagation inequality after being bounded in the state norm.
\LSMAtlasclaim{The suffix products determine whether earlier defects are damped or amplified. The chart shows analytical envelopes, not measured prediction errors.}
\LSMAtlasboundary{Both the true represented state and predicted state must remain in the Lipschitz domain. Stability of each action's eigenvalues alone does not certify stability under switching.}
\LSMAtlasnav{\texttt{DefectPropagation.lean}.}

\LSMAtlasplate{9}{What four observations identify}{Data span, closure and coverage are separate checks.}{P14--P17}
The running example begins $1,\tfrac12,\tfrac12,\tfrac{11}{24},\tfrac{61}{144}$. Its first four outputs give
\[
Z_-=\begin{pmatrix}1&1/2\\1/2&1/2\end{pmatrix},\quad
Z_+=\begin{pmatrix}1/2&1/2\\1/2&11/24\end{pmatrix},\quad
Z_+Z_-^{-1}=\begin{pmatrix}0&1\\1/12&5/6\end{pmatrix}.
\]
\LSMAtlasdiagram{09_identification}
The determinant $\det Z_-=1/4$ establishes coverage of the \emph{declared} two-dimensional delay space. The resulting recurrence is
\[
y_{t+1}=\frac56y_t+\frac1{12}y_{t-1},\qquad t\geq1.
\]
For general paired states, a linear fit exists exactly when $\ker Z_-\subseteq\ker Z_+$. It is unique on $W=\LSMAtlasim Z_-$, and maps $W$ to itself exactly when $\LSMAtlasim Z_+\subseteq W$.
\LSMAtlasclaim{These data identify the transition on their input span. The known generating system separately certifies that this span is the required full predictive state space.}
\LSMAtlasboundary{Finite sampled rank does not exclude a larger unexcited state space. Finite continuation agreement does not establish unrestricted discovery. Perturbations additionally require a conditioning bound.}
\LSMAtlasnav{\texttt{DelayClosure.lean}; \texttt{DataSpan.lean}; see the exact module map at the end.}

\LSMAtlasplate{10}{Behavior can supply the coordinates}{Future tests define a predictive state without naming hidden variables.}{P27--P28}
For a controlled process, arrange population probabilities into a behavior table: columns index past histories and rows index future tests. Its full column span is the linear predictive behavior space.
\[
\begin{array}{c|ccc}
 & h_1 & h_2 & \cdots\\\hline
\tau_1 & p(h_1\tau_1)&p(h_2\tau_1)&\cdots\\
\tau_2 & p(h_1\tau_2)&p(h_2\tau_2)&\cdots\\
\vdots &\vdots&\vdots&
\end{array}
\]
Assume the full behavior space has rank at most $r$. Choose an invertible $r$-by-$r$ population core $H$. Let $G_s$ be the corresponding table with a declared action-observation symbol $s$ inserted. With a compatible row/column convention,
\[
\boxed{B_s=G_sH^{-1}.}
\]
The empty-test row supplies the mass functional and a covered history supplies the initial predictive state. The factorization then transports every future test through the observable coordinates.
\LSMAtlasclaim{The dimension is determined by independent predictive behaviors. A global behavior-rank bound and an invertible core establish that the chosen coordinates cover the full behavior space.}

\textbf{Two distinct obstructions.} An invertible sampled block does not establish the full-process rank bound. A fitted linear realization does not automatically preserve a probability domain. The next plates address the latter obligation.

\LSMAtlasboundary{The reconstruction uses population tables under the stated coverage hypotheses. No finite-sample rank test here authenticates those hypotheses or an unknown hidden-state model.}
\LSMAtlasnav{\texttt{BehaviorSpace.lean}; \texttt{BehaviorMinimality.lean}; \newline\texttt{ObservableCore.lean}; \texttt{ObservableRealization.lean}.}

\LSMAtlasplate{11}{Probability meaning is geometric}{An operator family needs a domain, a mass map and a covered initial state.}{P29, P39--P40}
Let $C\subseteq\LSMAtlasR^d$ be a convex cone and $\lambda$ a linear mass map strictly positive on $C\setminus\{0\}$. Set $K=\{z\in C:\lambda z=1\}$. Linear observation operators must satisfy
\[
B_oC\subseteq C,\qquad \lambda\sum_o B_o=\lambda,\qquad z_0\in K.
\]
\LSMAtlasdiagram{10_domain}
A finite record $o_1\cdots o_T$ has probability
\[
\lambda B_{o_T}\cdots B_{o_1}z_0\geq0.
\]
Summing over every branch repeatedly applies the mass identity and gives total probability one. A positive-probability posterior is $B_oz/(\lambda B_oz)$.

\textbf{A finite-generator witness.} Let $V$ have full row rank and normalized generator columns. Nonnegative lifts $A_o$ with $B_oV=VA_o$ and $\mathbf1^{\mathsf T}\sum_oA_o=\mathbf1^{\mathsf T}$ certify the generated cone. The initial state must be a normalized nonnegative mixture of those columns.
\LSMAtlasclaim{All-horizon validity follows by induction from a finite exact witness. Statistical proximity to a generating process is a further obligation.}
\LSMAtlasboundary{A zero-probability branch does not define a normalized posterior by division. A finite positive lift is sufficient, but some valid low-dimensional processes require curved domains.}
\LSMAtlasnav{\texttt{CanonicalCone.lean}; \texttt{PositiveLift.lean}; \texttt{PositiveLiftTopology.lean}.}

\LSMAtlasplate{12}{A curved domain can have a finite witness}{A matrix embedding changes the certificate language.}{P46, P50--P51}
The Lorentz cone and its normalized disk admit the representation
\[
C=\{(m,x,y):m\geq\sqrt{x^2+y^2}\},\qquad
J(m,x,y)=\frac12\begin{pmatrix}m+x&y\\y&m-x\end{pmatrix}.
\]
\LSMAtlasdiagram[0.96]{11_psd}
Since the eigenvalues of $J$ are $(m\pm\sqrt{x^2+y^2})/2$, membership in $C$ is equivalent to $J\succeq0$. Also $\operatorname{tr}J=m$.

For the exact spatial rotation $R=\bigl(\begin{smallmatrix}3/5&-4/5\\4/5&3/5\end{smallmatrix}\bigr)$ and $T=\LSMAtlasdiag(1,R)$,
\[
J(Tz)=\frac15WJ(z)W^{\mathsf T},\qquad
W=\begin{pmatrix}2&-1\\1&2\end{pmatrix}.
\]
\LSMAtlasclaim{A rational congruence identity certifies the unchanged rotation on the curved cone. In the full P46 model, the canonical domain has no finite normalized positive-generator certificate.}
\LSMAtlasboundary{The drawn orbit is finite. Density and the finite-generator obstruction are proved in the volume, not established by its plotted points. The full instrument also needs its observation, mass and initial-state identities. No quantum implementation is asserted.}
\LSMAtlasnav{\texttt{RotationP46.lean}; \texttt{MatrixEmbedding.lean}; \texttt{MatrixRotation.lean}.}

\LSMAtlasplate{13}{Integration retains a covariance term}{Averaging an operator before quadratic action changes the instrument.}{P56--P59}
Let $X\succeq0$, $K(u)=K_0+uK_1$ and $\bar K=\mathbb E[K(U)]$. For a scalar $U$ of finite second moment,
\[
\mathbb E[K(U)XK(U)^{\mathsf T}]-\bar KX\bar K^{\mathsf T}
=\operatorname{Var}(U)K_1XK_1^{\mathsf T}\succeq0.
\]
\LSMAtlasdiagram{12_covariance}
For the uniform probability measure $\mu(du)=du/2$ on $[-1,1]$, take
\[
K(u)=\tfrac12\LSMAtlasdiag(1+3u,1-3u),\qquad X=I/2.
\]
Then $\bar K=I/2$. The averaged congruence is $I/2$, with trace one. Congruence by the averaged matrix is $I/8$, with trace $1/4$. The omitted covariance correction is $3I/8$.
\LSMAtlasclaim{The missing term is positive semidefinite and carries three quarters of the mass in this one-bin example. With $n$ equal bins, the missing mass is $3/(4n^2)$.}
\LSMAtlasboundary{The matrices in this example are diagonal and commute. The failure concerns the order of averaging and forming a quadratic action. A density value is interpreted relative to its declared measure.}
\LSMAtlasnav{See P57--P58 in the formal coverage map for exact integral and bin identities.}


\chapter{KS: references}
\begin{thebibliography}{99}
\bibitem{KS:stacks-limits} The Stacks Project Authors. \emph{Limits and colimits over preordered sets}. Section 4.21, Tag 002Z, including the finite-set inverse-limit theorem in Lemma 4.21.7. \url{https://stacks.math.columbia.edu/tag/002Z}. Accessed 10 October 2026.
\bibitem{KS:stacks-stack} The Stacks Project Authors. \emph{Stacks}. Section 8.4, Tag 0268. In particular the effective descent and equivalence-of-descent-categories formulations. \url{https://stacks.math.columbia.edu/tag/0268}. Accessed 10 October 2026.
\bibitem{KS:mathlib-stack} The mathlib Community. \emph{CategoryTheory.Sites.Descent.IsStack}. Standard pseudofunctor stack interface. \url{https://leanprover-community.github.io/mathlib4_docs/Mathlib/CategoryTheory/Sites/Descent/IsStack.html}.
\bibitem{KS:mathlib-projection} The mathlib Community. \emph{Analysis.InnerProductSpace.Projection.Basic}. Orthogonal projection onto complete subspaces. \url{https://leanprover-community.github.io/mathlib4_docs/Mathlib/Analysis/InnerProductSpace/Projection/Basic.html}.
\bibitem{KS:mathlib-banach} The mathlib Community. \emph{Analysis.Normed.Operator.Banach}. Banach open mapping and continuous inverse constructions. \url{https://leanprover-community.github.io/mathlib4_docs/Mathlib/Analysis/Normed/Operator/Banach.html}.
\bibitem{KS:clason-inverse} Christian Clason. \emph{Regularization of Inverse Problems}. Graduate lecture notes, arXiv:2001.00617v2, 2021. Sections 2, 4, and 6 provide standard background on weak convergence, parameter choice, and Tikhonov regularization. \url{https://arxiv.org/abs/2001.00617v2}. Consulted 12 October 2026.
\end{thebibliography}




\chapter{GA: references}
\begin{thebibliography}{14}\fontsize{9.5}{11}\selectfont\setlength{\itemsep}{3pt}\setlength{\parskip}{0pt}
\bibitem{GA:understanding} Komorebi P. \emph{Machine Understanding: Representation, Explanation, and Transfer}. First consolidated edition, version 1.0.0, Operation Komorebi, 8 October 2026. Supplied project manuscript and source archive.
\bibitem{GA:selfimproving} Komorebi P. \emph{Self-Improving Intelligence: Representation, Discovery, and Recursive Revision}. First consolidated edition, version 1.0.0, Operation Komorebi, 8 October 2026. Supplied project manuscript and source archive.
\bibitem{GA:resolution} Komorebi P. \emph{Sufficient Resolution}. Standardized traditional paper, Operation Komorebi, 2026. Finite behavioral models, hidden memory, reconstruction, and calibration. Source inspected in the project Library.
\bibitem{GA:rothenberg} T. J. Rothenberg. Identification in parametric models. \emph{Econometrica} 39(3):577--591, 1971. \href{https://www.jstor.org/stable/1913267}{doi:10.2307/1913267}.
\bibitem{GA:psr} M. L. Littman, R. S. Sutton, and S. Singh. Predictive representations of state. \emph{Advances in Neural Information Processing Systems 14}, 1555--1561. Conference 2001, proceedings 2002. \href{https://papers.nips.cc/paper_files/paper/2001/hash/1e4d36177d71bbb3558e43af9577d70e-Abstract.html}{Proceedings record}.
\bibitem{GA:angluin} D. Angluin. Learning regular sets from queries and counterexamples. \emph{Information and Computation} 75(2):87--106, 1987. \href{https://doi.org/10.1016/0890-5401(87)90052-6}{doi:10.1016/0890-5401(87)90052-6}.
\bibitem{GA:informativity} H. J. van Waarde, J. Eising, H. L. Trentelman, and M. K. Camlibel. Data informativity: a new perspective on data-driven analysis and control. \emph{IEEE Transactions on Automatic Control} 65(11):4753--4768, 2020. \href{https://arxiv.org/abs/1908.00468}{arXiv:1908.00468}.
\bibitem{GA:causal} A. Geiger, D. Ibeling, A. Zur, M. Chaudhary, S. Chauhan, J. Huang, A. Arora, Z. Wu, N. Goodman, C. Potts, and T. Icard. Causal abstraction: a theoretical foundation for mechanistic interpretability. \emph{Journal of Machine Learning Research} 26(83):1--64, 2025. \href{https://jmlr.org/papers/v26/23-0058.html}{Publication record}.
\bibitem{GA:stateid} P. van den Bos and F. Vaandrager. State identification for labeled transition systems with inputs and outputs. 2019. \href{https://arxiv.org/abs/1907.11034}{arXiv:1907.11034}. Cited for the testing framework, not as a source of the present three-case construction.
\bibitem{GA:blackwell} D. Blackwell. Equivalent comparisons of experiments. \emph{Annals of Mathematical Statistics} 24(2):265--272, 1953. \href{https://doi.org/10.1214/aoms/1177729032}{doi:10.1214/aoms/1177729032}.
\bibitem{GA:tzeng} W.-G. Tzeng. A polynomial-time algorithm for the equivalence of probabilistic automata. \emph{SIAM Journal on Computing} 21(2):216--227, 1992. \href{https://doi.org/10.1137/0221017}{doi:10.1137/0221017}. Cited for the established finite linear-algebra equivalence method.
\bibitem{GA:petreczky} M. Petreczky, L. Bako, and J. H. van Schuppen. Realization theory of discrete-time linear switched systems. \emph{Automatica} 49(11):3337--3344, 2013. \href{https://doi.org/10.1016/j.automatica.2013.07.022}{doi:10.1016/j.automatica.2013.07.022}. Author preprint: \href{https://arxiv.org/abs/1103.1343}{arXiv:1103.1343}.
\bibitem{GA:legland} F. Le Gland and N. Oudjane. Stability and uniform approximation of nonlinear filters using the Hilbert metric and application to particle filters. \emph{Annals of Applied Probability} 14(1):144--187, 2004. doi:10.1214/aoap/1075828050. Cited for filter stability and approximation, not as a source of the present static finite-record bound.
\bibitem{GA:girard} A. Girard and G. J. Pappas. Approximate bisimulation relations for constrained linear systems. \emph{Automatica} 43(8):1307--1317, 2007. doi:10.1016/j.automatica.2007.01.019. Cited for observable system approximation; the discrete block comparison is proved directly here.
\end{thebibliography}

\chapter{MU: references}
\begin{thebibliography}{99}\fontsize{9.5}{11}\selectfont
\setlength{\itemsep}{3pt}\setlength{\parsep}{0pt}\setlength{\parskip}{0pt}
\bibitem{MU:societies} Komorebi P. \emph{Machine Societies: A Theory of Persistent Plural Computational Organization}. Operation Komorebi, version 1.0.0, 2026. Companion theoretical monograph.
\bibitem{MU:jacovi} Alon Jacovi and Yoav Goldberg. Towards Faithfully Interpretable NLP Systems: How Should We Define and Evaluate Faithfulness? \emph{Proceedings of ACL}, pp. 4198--4205, 2020. \url{https://aclanthology.org/2020.acl-main.386/}.
\bibitem{MU:shalizi} Cosma Rohilla Shalizi and James P. Crutchfield. Computational Mechanics: Pattern and Prediction, Structure and Simplicity. \emph{Journal of Statistical Physics}, 104, 2001. \url{https://doi.org/10.1023/A:1010388907793}.
\bibitem{MU:givan} Robert Givan, Thomas Dean, and Matthew Greig. Equivalence Notions and Model Minimization in Markov Decision Processes. \emph{Artificial Intelligence}, 147:163--223, 2003. \url{https://engineering.purdue.edu/~givan/papers/mm.pdf}.
\bibitem{MU:tishby} Naftali Tishby, Fernando C. Pereira, and William Bialek. The Information Bottleneck Method. \emph{Proceedings of the 37th Allerton Conference}, 1999. Author preprint, 2000. \url{https://arxiv.org/abs/physics/0004057}.
\bibitem{MU:beckers} Sander Beckers and Joseph Y. Halpern. Abstracting Causal Models. \emph{Proceedings of AAAI}, 2019. \url{https://arxiv.org/abs/1812.03789}.
\bibitem{MU:geiger} Atticus Geiger, Duligur Ibeling, Amir Zur, Maheep Chaudhary, Sonakshi Chauhan, Jing Huang, Aryaman Arora, Zhengxuan Wu, Noah Goodman, Christopher Potts, and Thomas Icard. Causal Abstraction: A Theoretical Foundation for Mechanistic Interpretability. \emph{Journal of Machine Learning Research}, 26(83):1--64, 2025. \url{https://jmlr.org/papers/v26/23-0058.html}.
\bibitem{MU:hewitt} John Hewitt and Percy Liang. Designing and Interpreting Probes with Control Tasks. \emph{Proceedings of EMNLP-IJCNLP}, pp. 2733--2743, 2019. \url{https://aclanthology.org/D19-1275/}.
\bibitem{MU:sutter} Denis Sutter, Julian Minder, Thomas Hofmann, and Tiago Pimentel. The Non-Linear Representation Dilemma: Is Causal Abstraction Enough for Mechanistic Interpretability? Preprint, 2025. Version 1 consulted. \url{https://arxiv.org/abs/2507.08802v1}.
\bibitem{MU:miller} Tim Miller. Explanation in Artificial Intelligence: Insights from the Social Sciences. \emph{Artificial Intelligence}, 267:1--38, 2019. Author preprint: \url{https://arxiv.org/abs/1706.07269}.
\bibitem{MU:gentner} Dedre Gentner. Structure-Mapping: A Theoretical Framework for Analogy. \emph{Cognitive Science}, 7:155--170, 1983. Author-hosted copy: \url{https://groups.psych.northwestern.edu/gentner/papers/Gentner83.2b.pdf}.
\bibitem{MU:angluin} Dana Angluin. Learning Regular Sets from Queries and Counterexamples. \emph{Information and Computation}, 75:87--106, 1987. \url{https://doi.org/10.1016/0890-5401(87)90052-6}.
\bibitem{MU:blackwell} David Blackwell. Equivalent Comparisons of Experiments. \emph{Annals of Mathematical Statistics}, 24(2):265--272, 1953. \url{https://doi.org/10.1214/aoms/1177729032}.
\bibitem{MU:rischel} Eigil F. Rischel and Sebastian Weichwald. Compositional Abstraction Error and a Category of Causal Models. \emph{Proceedings of UAI}, PMLR 161:1013--1023, 2021. \url{https://proceedings.mlr.press/v161/rischel21a.html}.

\bibitem{MU:witsenhausen} Hans S. Witsenhausen. The Zero-Error Side Information Problem and Chromatic Numbers (Corresp.). \emph{IEEE Transactions on Information Theory}, 22(5):592--593, 1976. \url{https://doi.org/10.1109/TIT.1976.1055607}.
\bibitem{MU:alon} Noga Alon, Avinatan Hassidim, Eyal Lubetzky, Uri Stav, and Amit Weinstein. Broadcasting with Side Information. \emph{Proceedings of the 49th IEEE Symposium on Foundations of Computer Science}, pp. 823--832, 2008. \url{https://arxiv.org/abs/0806.3246}.
\bibitem{MU:smallwood} Richard D. Smallwood and Edward J. Sondik. The Optimal Control of Partially Observable Markov Processes over a Finite Horizon. \emph{Operations Research}, 21(5):1071--1088, 1973. \url{https://doi.org/10.1287/opre.21.5.1071}.
\bibitem{MU:lindley} Dennis V. Lindley. On a Measure of the Information Provided by an Experiment. \emph{Annals of Mathematical Statistics}, 27(4):986--1005, 1956. \url{https://doi.org/10.1214/aoms/1177728069}.
\bibitem{MU:halmos} Paul R. Halmos and L. J. Savage. Application of the Radon-Nikodym Theorem to the Theory of Sufficient Statistics. \emph{Annals of Mathematical Statistics}, 20(2):225--241, 1949. \url{https://doi.org/10.1214/aoms/1177730032}.
\end{thebibliography}

\chapter{SII: references}
\begin{thebibliography}{99}\fontsize{9.5}{11.2}\selectfont\setlength{\itemsep}{3pt}\setlength{\parskip}{0pt}
\bibitem{SII:understanding} Komorebi P. \emph{Machine Understanding: Representation, Explanation, and Transfer}. Operation Komorebi, first consolidated edition, v1.0.0, 8 October 2026.
\bibitem{SII:sric} Komorebi P. \emph{Stratified Recursive Improvement Closure: Obstruction, Descent, Transport, and Evidence-Preserving Lineages}. Formal theory draft v2.4, Block 23, 3 September 2026. The reviewed source is identified in the package source map. Cited for theoretical structure, not as empirical evidence for this manuscript's architecture.
\bibitem{SII:neuralese} Komorebi P. \emph{Neuralese}. Operation Komorebi, Final Research Edition v1.0, 2026. Section 12 specifies the proposed semantic and continuation interface.
\bibitem{SII:osl} Komorebi P. \emph{Operator-Symbolic Language}. Operation Komorebi, Final Research Edition v1.0, 2026. Explicit semantic and operational interface.
\bibitem{SII:chollet} F. Chollet. On the Measure of Intelligence. 2019. \url{https://arxiv.org/abs/1911.01547}.
\bibitem{SII:andrychowicz} M. Andrychowicz, M. Denil, S. Gomez, M. W. Hoffman, D. Pfau, T. Schaul, B. Shillingford, and N. de Freitas. Learning to learn by gradient descent by gradient descent. 2016. \url{https://arxiv.org/abs/1606.04474}.
\bibitem{SII:maml} C. Finn, P. Abbeel, and S. Levine. Model-Agnostic Meta-Learning for Fast Adaptation of Deep Networks. \emph{PMLR} 70:1126--1135, 2017. \url{https://proceedings.mlr.press/v70/finn17a.html}.
\bibitem{SII:dreamcoder} K. Ellis et al. DreamCoder: Growing generalizable, interpretable knowledge with wake-sleep Bayesian program learning. 2020. \url{https://arxiv.org/abs/2006.08381}.
\bibitem{SII:seal} A. Zweiger, J. Pari, H. Guo, E. Aky\"urek, Y. Kim, and P. Agrawal. Self-Adapting Language Models. 2025. \url{https://arxiv.org/abs/2506.10943}.
\bibitem{SII:goedel} J. Schmidhuber. Goedel Machines: Self-Referential Universal Problem Solvers Making Provably Optimal Self-Improvements. 2003, revised 2006. \url{https://arxiv.org/abs/cs/0309048}.
\bibitem{SII:dgm} J. Zhang, S. Hu, C. Lu, R. Lange, and J. Clune. Darwin Godel Machine: Open-Ended Evolution of Self-Improving Agents. 2025. \url{https://arxiv.org/abs/2505.22954}.
\bibitem{SII:hyperagents} J. Zhang et al. Hyperagents. 2026. \url{https://arxiv.org/abs/2603.19461}. Public manuscript accessed 8 October 2026.
\bibitem{SII:aide} D. Srikanth, B. Zhao, D. Xu, Y. Wu, and Z. Jiang. Recursive self-improvement of AI research agents. Preprint, 22 September 2026. \url{https://arxiv.org/abs/2609.26457}. The outer-loop comparison is qualified in Section 3.6.
\bibitem{SII:geiger} A. Geiger et al. Causal Abstraction: A Theoretical Foundation for Mechanistic Interpretability. \emph{Journal of Machine Learning Research} 26(83):1--64, 2025. \url{https://jmlr.org/papers/v26/23-0058.html}.
\bibitem{SII:dwork} C. Dwork, V. Feldman, M. Hardt, T. Pitassi, O. Reingold, and A. Roth. Preserving Statistical Validity in Adaptive Data Analysis. 2014, revised 2016. \url{https://arxiv.org/abs/1411.2664}.
\bibitem{SII:howard} S. R. Howard, A. Ramdas, J. McAuliffe, and J. Sekhon. Time-uniform Chernoff bounds via nonnegative supermartingales. \emph{Probability Surveys} 17:257--317, 2020. \url{https://doi.org/10.1214/18-PS321}.

\bibitem{SII:nilim} A. Nilim and L. El Ghaoui. Robust Control of Markov Decision Processes with Uncertain Transition Matrices. \emph{Operations Research} 53(5):780--798, 2005. \url{https://doi.org/10.1287/opre.1050.0216}. Cited as background on explicit model uncertainty, not as an identical uncertainty model.
\bibitem{SII:sketch} A. Solar-Lezama et al. Combinatorial sketching for finite programs. \emph{ASPLOS}, pp. 404--415, 2006. \url{https://doi.org/10.1145/1168857.1168907}.
\bibitem{SII:sygus} M. Raghothaman and A. Udupa. Language to Specify Syntax-Guided Synthesis Problems. 2014, revised 2016. \url{https://arxiv.org/abs/1405.5590}. Specification-language account of the syntax-guided synthesis framework introduced by Alur and colleagues in 2013.
\bibitem{SII:ice} P. Garg, C. L\"oding, P. Madhusudan, and D. Neider. ICE: A Robust Framework for Learning Invariants. In \emph{Computer Aided Verification}, LNCS 8559, pp. 69--87, 2014. \url{https://doi.org/10.1007/978-3-319-08867-9_5}.
\bibitem{SII:egg} M. Willsey, C. Nandi, Y. R. Wang, O. Flatt, Z. Tatlock, and P. Panchekha. egg: Fast and Extensible Equality Saturation. \emph{Proceedings of the ACM on Programming Languages} 5(POPL), 2021. \url{https://doi.org/10.1145/3434304}.
\bibitem{SII:ruler} C. Nandi, M. Willsey, A. Zhu, Y. R. Wang, B. Saiki, A. Anderson, A. Schulz, D. Grossman, and Z. Tatlock. Rewrite Rule Inference Using Equality Saturation. \emph{Proceedings of the ACM on Programming Languages} 5(OOPSLA), 2021. \url{https://doi.org/10.1145/3485496}. Author manuscript: \url{https://arxiv.org/abs/2108.10436}.
\bibitem{SII:williamson} D. P. Williamson and D. B. Shmoys. \emph{The Design of Approximation Algorithms}. Cambridge University Press, 2011. Chapter 1 develops set-cover relaxations and dual bounds. Author-hosted edition: \url{https://www.designofapproxalgs.com/}.
\end{thebibliography}

\chapter{MS: references}
\begin{thebibliography}{99}
\small
\setlength{\itemsep}{.12em}
\setlength{\parsep}{0pt}
\raggedright
\phantomsection
\bibitem{MS:shalizi} C.~R. Shalizi and J.~P. Crutchfield. \emph{Computational Mechanics: Pattern and Prediction, Structure and Simplicity}. Journal of Statistical Physics, 104, 2001. \href{https://doi.org/10.1023/A:1010388907793}{doi:10.1023/A:1010388907793}.
\bibitem{MS:givan} R. Givan, T. Dean, and M. Greig. \emph{Equivalence Notions and Model Minimization in Markov Decision Processes}. Artificial Intelligence, 147:163--223, 2003. \url{https://engineering.purdue.edu/~givan/papers/mm.pdf}.
\bibitem{MS:rubenstein} P.~K. Rubenstein, S. Weichwald, S. Bongers, J.~M. Mooij, D. Janzing, M. Grosse-Wentrup, and B. Sch\"olkopf. \emph{Causal Consistency of Structural Equation Models}. UAI, 2017. \url{https://arxiv.org/abs/1707.00819}.
\bibitem{MS:beckers} S. Beckers and J.~Y. Halpern. \emph{Abstracting Causal Models}. AAAI, 2019. \url{https://arxiv.org/abs/1812.03789}.
\bibitem{MS:rischel} E.~F. Rischel and S. Weichwald. \emph{Compositional Abstraction Error and a Category of Causal Models}. UAI, PMLR 161:1013--1023, 2021. \url{https://proceedings.mlr.press/v161/rischel21a.html}.
\bibitem{MS:tishby} N. Tishby, F.~C. Pereira, and W. Bialek. \emph{The Information Bottleneck Method}. Allerton Conference, 1999. \url{https://arxiv.org/abs/physics/0004057}.
\bibitem{MS:williams} P.~L. Williams and R.~D. Beer. \emph{Nonnegative Decomposition of Multivariate Information}. Preprint, 2010. \url{https://arxiv.org/abs/1004.2515}.
\bibitem{MS:young} H.~P. Young. \emph{The Evolution of Conventions}. Econometrica, 61(1):57--84, 1993. \url{https://www.jstor.org/stable/2951778}.
\bibitem{MS:ostrom} E. Ostrom. \emph{Beyond Markets and States: Polycentric Governance of Complex Economic Systems}. American Economic Review, 100(3):641--672, 2010. \url{https://doi.org/10.1257/aer.100.3.641}.
\bibitem{MS:ostrombook} E. Ostrom. \emph{Governing the Commons: The Evolution of Institutions for Collective Action}. Cambridge University Press, 1990.
\bibitem{MS:flp} M.~J. Fischer, N.~A. Lynch, and M.~S. Paterson. \emph{Impossibility of Distributed Consensus with One Faulty Process}. Journal of the ACM, 32(2):374--382, 1985. \url{https://groups.csail.mit.edu/tds/papers/Lynch/jacm85.pdf}.
\bibitem{MS:necula} G.~C. Necula. \emph{Proof-Carrying Code}. POPL, 1997. \url{https://doi.org/10.1145/263699.263712}.
\bibitem{MS:saltzer} J.~H. Saltzer and M.~D. Schroeder. \emph{The Protection of Information in Computer Systems}. Proceedings of the IEEE, 63(9):1278--1308, 1975. \url{https://web.mit.edu/Saltzer/www/publications/protection/}.
\bibitem{MS:douceur} J.~R. Douceur. \emph{The Sybil Attack}. IPTPS, 2002. \url{https://www.microsoft.com/en-us/research/publication/the-sybil-attack/}.
\bibitem{MS:park} J.~S. Park, J.~C. O'Brien, C.~J. Cai, M.~R. Morris, P. Liang, and M.~S. Bernstein. \emph{Generative Agents: Interactive Simulacra of Human Behavior}. UIST, 2023. \url{https://arxiv.org/abs/2304.03442}.
\bibitem{MS:ashery} A.~F. Ashery, L.~M. Aiello, and A. Baronchelli. \emph{Emergent Social Conventions and Collective Bias in LLM Populations}. Science Advances, 11(20):eadu9368, 2025. \url{https://doi.org/10.1126/sciadv.adu9368}.
\bibitem{MS:barrie} C. Barrie and P. T\"ornberg. \emph{Emergent LLM Behaviors Are Observationally Equivalent to Data Leakage}. Preprint, 2025. \url{https://arxiv.org/abs/2505.23796}.
\bibitem{MS:asheryreply} A.~F. Ashery, L.~M. Aiello, and A. Baronchelli. \emph{Reply to ``Emergent LLM Behaviors Are Observationally Equivalent to Data Leakage''}. Preprint, 2025. \url{https://arxiv.org/abs/2506.18600}.
\bibitem{MS:shoham} Y. Shoham and M. Tennenholtz. \emph{On the Emergence of Social Conventions: Modeling, Analysis, and Simulations}. Artificial Intelligence, 94(1--2):139--166, 1997. \url{https://doi.org/10.1016/S0004-3702(97)00028-3}.
\bibitem{MS:boella} G. Boella and L. van der Torre. \emph{Substantive and Procedural Norms in Normative Multiagent Systems}. Journal of Applied Logic, 6(2):152--171, 2008. \url{https://doi.org/10.1016/j.jal.2007.06.006}.
\bibitem{MS:grossinorm} D. Grossi, D. Gabbay, and L. van der Torre. \emph{The Norm Implementation Problem in Normative Multi-Agent Systems}. In \emph{Specification and Verification of Multi-agent Systems}, pp. 195--224, Springer, 2010. \url{https://doi.org/10.1007/978-1-4419-6984-2_7}.
\bibitem{MS:angluin} D. Angluin, J. Aspnes, Z. Diamadi, M.~J. Fischer, and R. Peralta. \emph{Computation in Networks of Passively Mobile Finite-State Sensors}. Distributed Computing, 18(4):235--253, 2006. \url{https://www.cs.yale.edu/homes/aspnes/papers/passively-mobile-abstract.html}.
\bibitem{MS:fudenbergmaskin} D. Fudenberg and E. Maskin. \emph{The Folk Theorem in Repeated Games with Discounting or with Incomplete Information}. Econometrica, 54(3):533--554, 1986. \url{https://ideas.repec.org/a/ecm/emetrp/v54y1986i3p533-54.html}.
\bibitem{MS:monderer} D. Monderer and L.~S. Shapley. \emph{Potential Games}. Games and Economic Behavior, 14(1):124--143, 1996. \url{https://doi.org/10.1006/game.1996.0044}.
\bibitem{MS:halpernmoses} J.~Y. Halpern and Y. Moses. \emph{Knowledge and Common Knowledge in a Distributed Environment}. Journal of the ACM, 37(3):549--587, 1990. \url{https://doi.org/10.1145/79147.79161}.
\bibitem{MS:crawfordsobel} V.~P. Crawford and J. Sobel. \emph{Strategic Information Transmission}. Econometrica, 50(6):1431--1451, 1982. \url{https://www.jstor.org/stable/1913390}.

\bibitem{MS:arrow} K.~J. Arrow. \emph{A Difficulty in the Concept of Social Welfare}. Journal of Political Economy, 58(4):328--346, 1950. \href{https://doi.org/10.1086/256963}{doi:10.1086/256963}.
\bibitem{MS:gibbard} A. Gibbard. \emph{Manipulation of Voting Schemes: A General Result}. Econometrica, 41(4):587--601, 1973. \href{https://doi.org/10.2307/1914083}{doi:10.2307/1914083}.
\bibitem{MS:satterthwaite} M.~A. Satterthwaite. \emph{Strategy-Proofness and Arrow's Conditions: Existence and Correspondence Theorems for Voting Procedures and Social Welfare Functions}. Journal of Economic Theory, 10(2):187--217, 1975. \href{https://doi.org/10.1016/0022-0531(75)90050-2}{doi:10.1016/0022-0531(75)90050-2}.
\bibitem{MS:pnuelirosner} A. Pnueli and R. Rosner. \emph{Distributed Reactive Systems Are Hard to Synthesize}. Proceedings of the 31st Annual Symposium on Foundations of Computer Science, vol.~II, pp.~746--757, 1990. \href{https://doi.org/10.1109/FSCS.1990.89597}{doi:10.1109/FSCS.1990.89597}.
\bibitem{MS:chatterjee} K. Chatterjee, L. Doyen, T.~A. Henzinger, and J.-F. Raskin. \emph{Algorithms for Omega-Regular Games with Imperfect Information}. Logical Methods in Computer Science, 3(3:4), 2007. \href{https://doi.org/10.2168/LMCS-3(3:4)2007}{doi:10.2168/LMCS-3(3:4)2007}.
\bibitem{MS:alpern} B. Alpern and F.~B. Schneider. \emph{Defining Liveness}. Information Processing Letters, 21(4):181--185, 1985. \href{https://doi.org/10.1016/0020-0190(85)90056-0}{doi:10.1016/0020-0190(85)90056-0}.
\bibitem{MS:taylorzwicker} A.~D. Taylor and W.~S. Zwicker. \emph{Simple Games: Desirability Relations, Trading, Pseudoweightings}. Princeton University Press, 1999. \href{https://doi.org/10.2307/j.ctv18zhdsm}{doi:10.2307/j.ctv18zhdsm}.

\bibitem{MS:ongaro} D. Ongaro and J. Ousterhout. \emph{In Search of an Understandable Consensus Algorithm (Extended Version)}. Stanford University technical report, 20 May 2014; conference version in USENIX ATC 2014. \url{https://raft.github.io/raft.pdf}.
\bibitem{MS:singhacl} M.~P. Singh. \emph{Agent Communication Languages: Rethinking the Principles}. Computer, 31(12):40--47, 1998. \href{https://doi.org/10.1109/2.735849}{doi:10.1109/2.735849}.
\bibitem{MS:newman} M.~H.~A. Newman. \emph{On Theories with a Combinatorial Definition of ``Equivalence''}. Annals of Mathematics, 43(2):223--243, 1942. \href{https://doi.org/10.2307/1968867}{doi:10.2307/1968867}.
\bibitem{MS:turing} A.~M. Turing. \emph{On Computable Numbers, with an Application to the Entscheidungsproblem}. Proceedings of the London Mathematical Society, s2-42(1):230--265, 1937. Received and read in 1936. \href{https://doi.org/10.1112/plms/s2-42.1.230}{doi:10.1112/plms/s2-42.1.230}.
\bibitem{MS:gaosync} T. Gao, J. Brodzki, and S. Mukherjee. \emph{The Geometry of Synchronization Problems and Learning Group Actions}. Discrete \& Computational Geometry, 65:150--211, 2021; published online in 2019. \href{https://doi.org/10.1007/s00454-019-00100-2}{doi:10.1007/s00454-019-00100-2}. Author version: \url{https://arxiv.org/abs/1610.09051}.
\bibitem{MS:baltag} A. Baltag, L.~S. Moss, and S. Solecki. \emph{The Logic of Public Announcements, Common Knowledge, and Private Suspicions}. CWI technical report SEN-R9922, 1999. \url{https://ir.cwi.nl/pub/4497}. Earlier version presented at TARK 1998.
\bibitem{MS:poet} R. Wang, J. Lehman, J. Clune, and K.~O. Stanley. \emph{Paired Open-Ended Trailblazer (POET): Endlessly Generating Increasingly Complex and Diverse Learning Environments and Their Solutions}. Preprint, 2019. \url{https://arxiv.org/abs/1901.01753}.

\bibitem{MS:lamporttime} L. Lamport. Time, Clocks, and the Ordering of Events in a Distributed System. \emph{Communications of the ACM}, 21(7):558--565, 1978. \url{https://lamport.azurewebsites.net/pubs/time-clocks.pdf}.
\bibitem{MS:preguicacrdt} N. Pregui\c{c}a, C. Baquero, and M. Shapiro. Conflict-free Replicated Data Types (CRDTs). Author encyclopedia-chapter preprint, 2018. \url{https://arxiv.org/abs/1805.06358}.
\bibitem{MS:bailisinvariant} P. Bailis, A. Fekete, M. J. Franklin, A. Ghodsi, J. M. Hellerstein, and I. Stoica. Coordination Avoidance in Database Systems (Extended Version). Author preprint, version 4, 2014. \url{https://arxiv.org/abs/1402.2237}.
\bibitem{MS:hellerstein-calm} J. M. Hellerstein and P. Alvaro. Keeping CALM: When Distributed Consistency is Easy. Author preprint, 2019. \url{https://arxiv.org/abs/1901.01930}.
\bibitem{MS:holmstrom1979} B. Holmstr\"om. Moral Hazard and Observability. \emph{The Bell Journal of Economics}, 10(1):74--91, 1979. \url{https://doi.org/10.2307/3003320}.
\bibitem{MS:yinzhang} G. G. Yin and Q. Zhang. \emph{Discrete-Time Markov Chains: Two-Time-Scale Methods and Applications}. Stochastic Modelling and Applied Probability, vol. 55. Springer, 2005. \url{https://doi.org/10.1007/b138226}.
\bibitem{MS:bikhchandani1992} S. Bikhchandani, D. Hirshleifer, and I. Welch. A Theory of Fads, Fashion, Custom, and Cultural Change as Informational Cascades. \emph{Journal of Political Economy}, 100(5):992--1026, 1992. \url{https://doi.org/10.1086/261849}.
\bibitem{MS:grossblasius} T. Gross and B. Blasius. Adaptive Coevolutionary Networks: A Review. \emph{Journal of the Royal Society Interface}, 5(20):259--271, 2008. \url{https://doi.org/10.1098/rsif.2007.1229}. Author preprint: \url{https://arxiv.org/abs/0709.1858}.
\bibitem{MS:jacksonwolinsky} M. O. Jackson and A. Wolinsky. A Strategic Model of Social and Economic Networks. \emph{Journal of Economic Theory}, 71(1):44--74, 1996. \url{https://doi.org/10.1006/jeth.1996.0108}.
\bibitem{MS:hall1935} P. Hall. On Representatives of Subsets. \emph{Journal of the London Mathematical Society}, s1-10(1):26--30, 1935. \url{https://doi.org/10.1112/jlms/s1-10.37.26}.

\bibitem{MS:boydconvex} S. Boyd and L. Vandenberghe. \emph{Convex Optimization}. Cambridge University Press, 2004. Author-hosted lecture slides, Chapter~5, on linear programming duality: \url{https://web.stanford.edu/~boyd/cvxbook/bv_cvxslides.pdf}.
\bibitem{MS:belieffails} K. Chatterjee and L. Doyen. \emph{Partial-Observation Stochastic Games: How to Win when Belief Fails}. Author preprint, 2011. \url{https://arxiv.org/abs/1107.2141}.


\bibitem{MS:meaningfulcontrol} F. Santoni de Sio and J. van den Hoven. \emph{Meaningful Human Control over Autonomous Systems: A Philosophical Account}. Frontiers in Robotics and AI, 5:15, 2018. \url{https://doi.org/10.3389/frobt.2018.00015}.
\bibitem{MS:stengelinspection} B. von Stengel. \emph{Recursive Inspection Games}. Mathematics of Operations Research, 41(3):935--952, 2016. \url{https://doi.org/10.1287/moor.2015.0762}. Author version: \url{https://arxiv.org/abs/1412.0129}.
\bibitem{MS:dworkroth} C. Dwork and A. Roth. \emph{The Algorithmic Foundations of Differential Privacy}. Foundations and Trends in Theoretical Computer Science, 9(3--4):211--407, 2014. \url{https://doi.org/10.1561/0400000042}. Author-hosted text: \url{https://www.cis.upenn.edu/~aaroth/privacybook.html}.
\bibitem{MS:offswitch} D. Hadfield-Menell, A. Dragan, P. Abbeel, and S. Russell. \emph{The Off-Switch Game}. Author preprint, 2016, revised 2017. \url{https://arxiv.org/abs/1611.08219}.
\bibitem{MS:blackboxsimplex} U. Mehmood, S. Sheikhi, S. Bak, S. A. Smolka, and S. D. Stoller. \emph{The Black-Box Simplex Architecture for Runtime Assurance of Autonomous CPS}. NASA Formal Methods, pp. 231--250, 2022. \url{https://doi.org/10.1007/978-3-031-06773-0_12}. Author version: \url{https://arxiv.org/abs/2102.12981}.
\bibitem{MS:componentsimplex} D. Phan, J. Yang, M. Clark, R. Grosu, J. D. Schierman, S. A. Smolka, and S. D. Stoller. \emph{A Component-Based Simplex Architecture for High-Assurance Cyber-Physical Systems}. Extended author preprint, 2017. \url{https://arxiv.org/abs/1704.04759}. Conference version: \url{https://doi.org/10.1109/ACSD.2017.23}.

\end{thebibliography}

\chapter{LSM: references}
\section{References and source
lineage}\label{LSM:references-and-source-lineage}

\textbf{{[}K1{]}} Komorebi P. \emph{Projector-Resolved Large Algebraic
Signal Model}. Historical 38-page capstone preserved in
\path{Komorebi_PR_LASM_2.3.0_Source.zip},
\path{paper/Projector_Resolved_Large_Algebraic_Signal_Model.tex}.
In particular, Sections 2--7 and the categorical-atlas distinction.
Historical experimental claims are not used here.

\textbf{{[}K2{]}} Komorebi P. \emph{Projector-Resolved Seed Calculus
VI}, corrected presentation in
\path{Paper_VI_With_Visual_Appendix.pdf}, Section 8, Theorem 8.4
and the adjacent coarsening criterion and one-packet boundary. This
supplement restates and proves the finite result it uses. It does not
claim a fresh Lean audit of that paper.

\textbf{{[}K3{]}} Komorebi P. \emph{PR-LASM 2.3.0 Release Audit}, 9
October 2026, and the accompanying source. The source inspected includes
\path{exact_atlas.py}, \path{lasm_packetizer.py},
\path{lasm_model.py}, \path{docs/RELEASE_GUIDE.md} and
\path{docs/claim-boundary.md}.

\textbf{{[}R1{]}} Ronald R. Coifman and Mladen Victor Wickerhauser.
``Entropy-based algorithms for best basis selection.'' \emph{IEEE
Transactions on Information Theory} 38(2), 713--718, 1992. DOI:
10.1109/18.119732. Source checked: the authors' Washington University
research record. The finite additive tree optimization is classical
background.

\textbf{{[}R2{]}} MIT OpenCourseWare. \emph{6.241J Dynamic Systems and
Control}, Chapter 24, ``Observability,'' Spring 2011, especially Section
24.5 on the largest invariant subspace contained in the observation
kernel. Source checked: the MIT course notes, PDF page 6. Used to
identify the classical control-theoretic antecedent of Proposition 8.

\textbf{{[}R3{]}} Michael Woodward, Yifeng Tian, Arvind Mohan, Yen Ting
Lin, Christoph Hader, Hermann Fasel, Misha Chertkov and Daniel Livescu.
\emph{Data-Driven Mori-Zwanzig: Approaching a Reduced Order Model for
Hypersonic Boundary Layer Transition}, arXiv:2301.07203, version 1,
2023, introduction and description of resolved, memory and unresolved
contributions. Cited for the established projection-memory context, not
as evidence for PR-LASM's performance or for any application discussed
in that paper. Proposition 10 is proved directly by finite linear
recurrence expansion here.

\textbf{{[}R4{]}} Peter Buchholz. ``Exact and ordinary lumpability in
finite Markov chains.'' \emph{Journal of Applied Probability} 31(1),
59--75, 1994. DOI: 10.2307/3215235. Source checked: Cambridge University
Press publication record. Used for the classical finite Markov
aggregation context.

\textbf{{[}R5{]}} Michael L. Littman, Richard S. Sutton and Satinder
Singh. ``Predictive Representations of State.'' \emph{Advances in Neural
Information Processing Systems} 14, NIPS 2001. The paper connects state
with predictions of future observations conditional on actions. Cited
for conceptual lineage, not as a proof of the finite delay theorem or of
this implementation. The original paper and proceedings record are the
primary sources. The scanned paper carries all three author names.

\textbf{{[}R6{]}} Steven L. Brunton, Bingni W. Brunton, Joshua L.
Proctor and J. Nathan Kutz. ``Koopman Invariant Subspaces and Finite
Linear Representations of Nonlinear Dynamical Systems for Control.''
\emph{PLOS ONE} 11(2), e0150171, 2016. DOI:
10.1371/journal.pone.0150171. Used for the established role and
limitations of finite invariant observable families in nonlinear
modeling. No general nonlinear feature-discovery guarantee is imported.

\textbf{{[}R7{]}} Daniel Hsu, Sham M. Kakade and Tong Zhang. ``A
Spectral Algorithm for Learning Hidden Markov Models.'' \emph{Journal of
Computer and System Sciences} 78(5), 1460--1480, 2012. Primary author
manuscript inspected: \path{hmm-jcss.pdf}, hosted by Daniel Hsu at
Columbia University, especially its abstract and rank assumptions. The
result illustrates observable-operator learning under explicit spectral
hypotheses. Its statistical guarantees are not claimed for PR-LASM.

\textbf{{[}R8{]}} Stephen Boyd, Laurent El Ghaoui, Eric Feron and
Venkataramanan Balakrishnan. \emph{Linear Matrix Inequalities in System
and Control Theory}. SIAM Studies in Applied Mathematics 15, 1994. The
authors' Stanford publication page identifies the book and provides the
primary text. Used to identify the classical common quadratic Lyapunov
and matrix-inequality setting. The discrete contraction implication used
in Section 19 is proved directly there.

\textbf{{[}R9{]}} Byron Boots, Sajid M. Siddiqi and Geoffrey J. Gordon.
``Closing the Learning-Planning Loop with Predictive State
Representations.'' \emph{Robotics: Science and Systems VI}, 2010. DOI:
10.15607/RSS.2010.VI.036. Primary eight-page proceedings paper inspected
at the authors' Carnegie Mellon host, especially Section II on
interventional tests and Section III on observable matrices. A journal
version appeared in \emph{The International Journal of Robotics
Research} 30(7), 954--966, 2011. Used for the established PSR lineage,
not as an assertion that this supplement implements that paper's
complete algorithm or planning guarantees. Source:
\texttt{https://www.cs.cmu.edu/\textasciitilde{}ggordon/boots-siddiqi-gordon-closing-loop-psrs.pdf}.

\textbf{{[}R10{]}} Han Zhao and Pascal Poupart. ``A Sober Look at
Spectral Learning.'' \emph{ICML 2014 Workshop on Spectral Learning},
2014. Primary author-hosted manuscript inspected, including its
discussion of negative probabilities in spectral estimates. Used only to
identify that practical failure mode. The cone and packet statements in
this supplement are proved independently. Source:
\path{https://hanzhaoml.github.io/papers/ICML2014/spectral-main.pdf}.

\textbf{{[}R11{]}} Wassily Hoeffding. ``Probability Inequalities for
Sums of Bounded Random Variables.'' \emph{Journal of the American
Statistical Association} 58(301), 13--30, 1963. DOI:
10.1080/01621459.1963.10500830. Publisher record checked. Proposition 35
supplies a self-contained Bernoulli moment-generating-function proof for
the protocol it uses. Source:
\path{https://www.tandfonline.com/doi/abs/10.1080/01621459.1963.10500830}.

\textbf{{[}R7a{]}} Daniel Hsu, Sham M. Kakade and Tong Zhang. ``A
Spectral Algorithm for Learning Hidden Markov Models.'' \emph{COLT},
2009. Original conference manuscript inspected at the official Learning
Theory proceedings site, particularly Sections 2.3--3.1 on rank,
sampling and observable representations. The later journal version is
{[}R7{]}. No statistical guarantee from either version is imported
without its assumptions. Source:
\path{https://www.learningtheory.org/colt2009/papers/011.pdf}.

\textbf{{[}R12{]}} Herbert Jaeger. ``Observable operator models for
discrete stochastic time series.'' \emph{Neural Computation} 12(6),
1371--1398, 2000. Author-hosted manuscript with the August 2017 erratum.
Section 7 gives the polyhedral-cone connection to finite positive
realizations and an irrational-rotation obstruction. The controlled
rational example in Section 51 is an independently worked adaptation of
that established obstruction, not a claim of discovering it. The
corrected manuscript and its geometric argument were inspected. Source:
\path{https://www.ai.rug.nl/minds/uploads/oom_neco00.pdf}.

\textbf{{[}R13{]}} Zoltán Horváth, Yunfei Song and Tamás Terlaky.
\emph{A Novel Unified Approach to Invariance Conditions for a Linear
Dynamical System}. arXiv:1405.5167, version 2, 5 July 2016. Primary
author manuscript, particularly the polyhedral-invariance discussion and
use of Farkas' lemma. Used for the classical invariant-cone and
optimization setting. The generator and robust-vertex statements used
here are proved directly. Source:
\path{https://arxiv.org/html/1405.5167v2}.

\textbf{{[}R14{]}} Marco Fanizza, Josep Lumbreras and Andreas Winter.
\emph{Quantum theory in finite dimension cannot explain every general
process with finite memory}. arXiv:2209.11225, version 2, 5 May 2023.
Primary manuscript, especially its discussion of finite linear
realizations, positive realizations and polyhedral cones. Cited for the
distinction between these representation classes. No quantum
implementation or quantum advantage is claimed for PR-LASM. The present
block uses a real three-dimensional cone construction proved in Section
51. Source: \path{https://arxiv.org/html/2209.11225v2}.

\textbf{{[}R15{]}} Anna Jen\v{c}ová. \emph{Base norms and discrimination of
generalized quantum channels}. arXiv:1308.4030, version 1, 19 August
2013. Primary manuscript, Section 1.2, printed page 5, on
finite-dimensional proper cones, normalized bases and base norms. That
ordered-vector-space definition is used in Section 46. The classical
instrument contraction and observation-tree bound needed here are
derived in the text, without importing quantum-channel assumptions.
Source: \path{https://arxiv.org/pdf/1308.4030}.

\textbf{{[}R16{]}} John Watrous. \emph{The Theory of Quantum
Information}. Cambridge University Press, 2018. Author-hosted
prepublication text, Chapter 2, Theorem 2.22, printed page 82 (PDF page
90), and the trace-preserving characterization in Theorem 2.26. Used for
the established Choi and congruence representations of completely
positive maps, not to assign quantum physics to the signal model. The
relevant theorem page was visually inspected. Source:
\texttt{https://cs.uwaterloo.ca/\textasciitilde{}watrous/TQI/TQI.pdf}.
This third-party text is not redistributed in the bundle.

\textbf{{[}R17{]}} Pablo A. Parrilo. \emph{Algebraic Techniques and
Semidefinite Optimization}, MIT 6.972, Lecture 10, 14 March 2006. MIT
OpenCourseWare notes, pages 10-1 through 10-3, on polynomial Gram
matrices, coefficient constraints and interval positivity. Used for the
standard sum-of-squares certificate method. Proposition 53 proves only
the sufficient matrix-valued certificate used here and does not import
scalar completeness as a matrix theorem. The formula pages were visually
inspected. Source:
\path{https://ocw.mit.edu/courses/6-972-algebraic-techniques-and-semidefinite-optimization-spring-2006/resources/lecture_10/}.

\textbf{{[}R18{]}} Herbert Jaeger. \emph{Modeling and learning
continuous-valued stochastic processes with OOMs}. GMD technical report,
8 June 2000. Author-hosted report, title and abstract inspected,
establishing that continuous-valued observable-operator modeling and
learning are prior work. The present polynomial, certificate and binning
arguments are proved within this volume. No guarantee from the report is
assigned to the reference implementation. Source:
\path{https://www.ai.rug.nl/minds/uploads/jaeger.00.tr.contoom.pdf}.

\textbf{{[}R19{]}} Ulrich Abel, Ana Maria Acu, Margareta Heilmann and
Ioan Ra\c{s}a. \emph{A representation for the integral kernel of the
composition of multivariate Bernstein--Durrmeyer operators}.
arXiv:2512.17297v1, 19 December 2025. The authors' primary manuscript,
Introduction, gives the classical univariate Bernstein--Durrmeyer kernel
used in Section 72. Its composition theorems are not needed or claimed
here. Proposition 64 derives the operator-valued positivity and error
statements directly. Source:
\path{https://arxiv.org/html/2512.17297v1}.

\textbf{{[}R20{]}} Steven R. Howard, Aaditya Ramdas, Jon McAuliffe and
Jasjeet Sekhon. \emph{Time-uniform, nonparametric, nonasymptotic
confidence sequences}. \emph{Annals of Statistics} 49(2), 1055--1080,
2021. DOI: 10.1214/20-AOS1991. The primary arXiv record and the
conditional exponential-supermartingale discussion in the v4 HTML
manuscript were inspected. Proposition 66 proves a simpler union-bound
confidence sequence, not the paper's sharper rates. Source:
\path{https://arxiv.org/abs/1810.08240} and
\path{https://arxiv.org/html/1810.08240v4}.

\textbf{{[}R21{]}} Roland Hildebrand. \emph{An LMI description for the
cone of Lorentz-positive maps}. Author manuscript, 17 January 2006.
Primary manuscript inspected, especially the introduction and
definitions on pages 1--2. It supplies classical context for
Lorentz-positive maps and the existence of convex descriptions. The
present margin and semicircle certificate arguments are proved directly
and do not claim to implement the manuscript's full matrix-inequality
construction. Source:
\path{https://optimization-online.org/wp-content/uploads/2005/12/1260.pdf}.

\textbf{{[}R22{]}} Tareq Hamadneh, Hassan Al-Zoubi, Hamza Alzaareer and
Rafael Wisniewski. \emph{Optimization and Positivity Certificates of
Rational Functions using Bernstein Form}. arXiv:1906.11037, 2019.
Primary manuscript inspected, including its Bernstein range enclosure
discussion on printed pages 3--4. It supplies established
positivity-certificate context. Proposition 75 gives its own
tensor-product rectangle argument and strict-subdivision proof. Source:
\path{https://arxiv.org/abs/1906.11037}.

\textbf{Inherited dependencies.} Proposition 52 defines continuous
instrument semantics. Proposition 59 supplies the common-domain
continuous-law bound. Propositions 62--63 distinguish exact ellipsoid
proposals from verification. Proposition 48 supplies repeated
confidence-set accounting. The new block restates the specific repair
and uncertainty inequalities it uses. The cited earlier results are
present in this cumulative volume.

\textbf{{[}R23{]}} Yasin Abbasi-Yadkori, David Pál and Csaba Szepesvári.
\emph{Online Least Squares Estimation with Self-Normalized Processes: An
Application to Bandit Problems}. arXiv:1102.2670, version 1, 14 February
2011. The primary manuscript's Section 2, Theorem 3, and Section 5
supply the classical martingale-mixture and adaptive-regression context.
Proposition 81 includes its own proof, retains conditional noise and
design assumptions, and does not claim a new concentration inequality.
Source checked: \path{https://arxiv.org/pdf/1102.2670}.

\textbf{{[}R24{]}} David A. Levin and Yuval Peres, with contributions by
Elizabeth L. Wilmer. \emph{Markov Chains and Mixing Times}, second
edition, 2017. Chapters 4 and 5 develop total variation and coupling.
The present minorization argument is proved directly and does not import
a finite-state assumption into its general kernel statement. The worked
reference chain is finite. Source checked: the author-hosted
second-edition text,
\path{https://pages.uoregon.edu/dlevin/MARKOV/mcmt2e.pdf}.

\textbf{Internal dependencies.} Part VII, especially Propositions 71,
73, 75, 76 and 79, supplies the robust checker and predecessor repair
comparison. Propositions 59 and 69 supply the continuous and
conditional-kernel full-law comparisons. Their mathematical source
passages are retained in this cumulative volume; prior archive
identifiers remain in the provenance records.

\endgroup

\chapter{PRIQ: references}
\begin{thebibliography}{99}

\bibitem{PRIQ:Antoulas2005}
A. C. Antoulas,
\emph{Approximation of Large-Scale Dynamical Systems},
SIAM, 2005.

\bibitem{PRIQ:BaezPollard2017}
J. C. Baez and B. S. Pollard,
A compositional framework for reaction networks,
\emph{Reviews in Mathematical Physics} 29 (2017), 1750028.

\bibitem{PRIQ:BaumPetrie1966}
L. E. Baum and T. Petrie,
Statistical inference for probabilistic functions of finite state Markov chains,
\emph{Annals of Mathematical Statistics} 37 (1966), 1554--1563.

\bibitem{PRIQ:Blackwell1953}
D. Blackwell,
Equivalent comparisons of experiments,
\emph{Annals of Mathematical Statistics} 24 (1953), 265--272.

\bibitem{PRIQ:Buchholz1994}
P. Buchholz,
Exact and ordinary lumpability in finite Markov chains,
\emph{Journal of Applied Probability} 31 (1994), 59--75.

\bibitem{PRIQ:Butcher1972}
J. C. Butcher,
An algebraic theory of integration methods,
\emph{Mathematics of Computation} 26 (1972), 79--106.

\bibitem{PRIQ:Cardelli2015}
L. Cardelli, M. Tribastone, M. Tschaikowski, and A. Vandin,
Forward and backward bisimulations for chemical reaction networks,
in \emph{CONCUR 2015}, LIPIcs 42 (2015), 226--239.

\bibitem{PRIQ:Cardelli2017}
L. Cardelli, M. Tribastone, M. Tschaikowski, and A. Vandin,
Maximal aggregation of polynomial dynamical systems,
\emph{Proceedings of the National Academy of Sciences} 114 (2017), 10029--10034.

\bibitem{PRIQ:ConstantineSavits1996}
G. M. Constantine and T. H. Savits,
A multivariate Fa\`a di Bruno formula with applications,
\emph{Transactions of the American Mathematical Society} 348 (1996), 503--520.

\bibitem{PRIQ:HairerLubichWanner2006}
E. Hairer, C. Lubich, and G. Wanner,
\emph{Geometric Numerical Integration}, second edition,
Springer, 2006.

\bibitem{PRIQ:HalmosSavage1949}
P. R. Halmos and L. J. Savage,
Application of the Radon--Nikodym theorem to the theory of sufficient statistics,
\emph{Annals of Mathematical Statistics} 20 (1949), 225--241.

\bibitem{PRIQ:KemenySnell1960}
J. G. Kemeny and J. L. Snell,
\emph{Finite Markov Chains},
Van Nostrand, 1960.

\bibitem{PRIQ:Meyer1989}
C. D. Meyer,
Stochastic complementation, uncoupling Markov chains, and the theory of nearly reducible systems,
\emph{SIAM Review} 31 (1989), 240--272.

\bibitem{PRIQ:RogersPitman1981}
L. C. G. Rogers and J. W. Pitman,
Markov functions,
\emph{Annals of Probability} 9 (1981), 573--582.

\bibitem{PRIQ:SimonAndo1961}
H. A. Simon and A. Ando,
Aggregation of variables in dynamic systems,
\emph{Econometrica} 29 (1961), 111--138.

\end{thebibliography}


\chapter{PPT: references}
\begin{thebibliography}{99}

\bibitem{PPT:komorebi-space}
Komorebi P.
\textit{Komorebi Space: Coherent Realizations, Obstructions, and Descent}.
Monograph, 2026. The finite proof-field view construction here is a relative application of its strict coherent realization principle, not a claim of unrestricted semantic completeness.

\bibitem{PPT:buss}
S. R. Buss, editor.
\textit{Handbook of Proof Theory}.
Elsevier, 1998.

\bibitem{PPT:cover-thomas}
T. M. Cover and J. A. Thomas.
\textit{Elements of Information Theory}.
Second edition, Wiley, 2006.

\bibitem{PPT:diestel}
R. Diestel.
\textit{Graph Theory}.
Fifth edition, Springer, 2017.

\bibitem{PPT:garey-johnson}
M. R. Garey and D. S. Johnson.
\textit{Computers and Intractability: A Guide to the Theory of NP-Completeness}.
W. H. Freeman, 1979.

\bibitem{PPT:knuth}
D. E. Knuth.
\textit{Literate Programming}.
CSLI Publications, 1992.

\bibitem{PPT:maclane}
S. Mac Lane.
\textit{Categories for the Working Mathematician}.
Second edition, Springer, 1998.

\bibitem{PPT:nemhauser}
G. L. Nemhauser, L. A. Wolsey, and M. L. Fisher.
An analysis of approximations for maximizing submodular set functions I.
\textit{Mathematical Programming}, 14:265--294, 1978.

\bibitem{PPT:troelstra}
A. S. Troelstra and H. Schwichtenberg.
\textit{Basic Proof Theory}.
Second edition, Cambridge University Press, 2000.

\end{thebibliography}


\chapter{CR: references}
\begin{thebibliography}{14}
\bibitem{CR:branchrelease}
Komorebi P. \emph{Constructive Repair v0.5.0: Proof Carrying Branch Search}. Structural tree checker, exact rational producer, kernel proofs, input echoes, and reproducible evidence, 19 October 2026.

\bibitem{CR:menurelease}
Komorebi P. \emph{Constructive Repair v0.4.0: Costed Rational Menus and Rejection Covers}. Exact producer, checked cover theorems, typed certificates, verification records, and original-input echoes, 19 October 2026.

\bibitem{CR:bridge}
Komorebi P. \emph{Constructive Repair v0.3.0: Exact Rational Certificate Bridge}. Python producer, Lean certificate modules, fixed-input records, fresh compilation and same-kernel replay evidence, 19 October 2026.


\bibitem{CR:formal}
Komorebi P. \emph{Constructive Repair: Lean Continuation v0.2.0}. Formal source release, 19 October 2026. Ten proof modules, declaration registry, verification records, and reproducibility scripts. Historical archive: \nolinkurl{Komorebi_Constructive_Repair_Lean_v0.2.0.zip}. Exact proof snapshots and its hash are retained in the current provenance index.

\bibitem{CR:audit}
Komorebi P. \emph{Constructive Repair Calculus: Exact Audit v0.1.0}. Mathematical audit, ordinary proofs, exact Python implementation, and case records, 19 October 2026. Preserved in the inherited material of the formal source release.

\bibitem{CR:paperi}
Komorebi P. \emph{Projector Resolved Seed Calculus I}. Traditional monograph in Operation Komorebi. Section 11, realization export and semantic realization criterion. Source context for the preservation contract.

\bibitem{CR:paperiii}
Komorebi P. \emph{Projector Resolved Seed Calculus III: Typed Transfer, Nonlinear Packets, and Obstruction Theory}. Modern monograph, 8 October 2026. Sections 21--24, correction maps, cokernel obstructions, nonlinear correction, and filtered lifting.

\bibitem{CR:paperxi}
Komorebi P. \emph{Toward Arithmetic Packet Realization}. Research companion to Paper XI. Restricted constructions, failed hosts, and the remaining obligations of Conjecture 66. Source context for the distinction between finite comparison and arithmetic existence.
\bibitem{CR:pythonio}
Python Software Foundation. \emph{Python 3.13 Library Reference: io, Text Encoding}. Documentation of locale-dependent defaults and explicit encoding selection. \url{https://docs.python.org/3.13/library/io.html#text-encoding}.
\end{thebibliography}

\chapter{SRWT: references}
% Bibliography for Stratified Relational Workspace Theory.
% !TeX root = Stratified_Relational_Workspace_Theory.tex
\begin{thebibliography}{99}
\small

\bibitem{SRWT:AbramskyBrandenburger2011}
S. Abramsky and A. Brandenburger.
\newblock The sheaf-theoretic structure of non-locality and contextuality.
\newblock \emph{New Journal of Physics}, 13:113036, 2011.
\newblock doi:10.1088/1367-2630/13/11/113036.

\bibitem{SRWT:AdamsAizawa2001}
F. Adams and K. Aizawa.
\newblock The bounds of cognition.
\newblock \emph{Philosophical Psychology}, 14(1):43--64, 2001.
\newblock doi:10.1080/09515080120033571.

\bibitem{SRWT:AmesEtAl2017}
A. D. Ames, X. Xu, J. W. Grizzle, and P. Tabuada.
\newblock Control barrier function based quadratic programs for safety critical systems.
\newblock \emph{IEEE Transactions on Automatic Control}, 62(8):3861--3876, 2017.
\newblock doi:10.1109/TAC.2016.2638961.

\bibitem{SRWT:Amornbunchornvej2026}
C. Amornbunchornvej.
\newblock Collective counterfactual planning: Coordination, consent, and verification under representational constraints.
\newblock arXiv:2608.17932, 2026.
\newblock doi:10.48550/arXiv.2608.17932.

\bibitem{SRWT:Ashby1956}
W. R. Ashby.
\newblock \emph{An Introduction to Cybernetics}.
\newblock Chapman \& Hall, London, 1956.

\bibitem{SRWT:Aubin1991}
J.-P. Aubin.
\newblock \emph{Viability Theory}.
\newblock Birkh\"auser, Boston, 1991.
\newblock doi:10.1007/978-0-8176-4910-4.

\bibitem{SRWT:Baars1988}
B. J. Baars.
\newblock \emph{A Cognitive Theory of Consciousness}.
\newblock Cambridge University Press, Cambridge, 1988.

\bibitem{SRWT:BarandiaranEtAl2009}
X. E. Barandiaran, E. Di Paolo, and M. Rohde.
\newblock Defining agency: Individuality, normativity, asymmetry, and spatio-temporality in action.
\newblock \emph{Adaptive Behavior}, 17(5):367--386, 2009.
\newblock doi:10.1177/1059712309343819.

\bibitem{SRWT:Blackwell1953}
D. Blackwell.
\newblock Equivalent comparisons of experiments.
\newblock \emph{Annals of Mathematical Statistics}, 24(2):265--272, 1953.
\newblock doi:10.1214/aoms/1177729032.


\bibitem{SRWT:BerkEtAl2013}
R. Berk, L. Brown, A. Buja, K. Zhang, and L. Zhao.
\newblock Valid post-selection inference.
\newblock \emph{The Annals of Statistics}, 41(2):802--837, 2013.
\newblock doi:10.1214/12-AOS1077.

\bibitem{SRWT:Bredon1997}
G. E. Bredon.
\newblock \emph{Sheaf Theory}.
\newblock 2nd edition, Springer, New York, 1997.

\bibitem{SRWT:Bratman1992}
M. E. Bratman.
\newblock Shared cooperative activity.
\newblock \emph{The Philosophical Review}, 101(2):327--341, 1992.
\newblock doi:10.2307/2185537.

\bibitem{SRWT:CasazzaKovacevic2003}
P. G. Casazza and J. Kova\v{c}evi\'c.
\newblock Equal-norm tight frames with erasures.
\newblock \emph{Advances in Computational Mathematics}, 18:387--430, 2003.
\newblock doi:10.1023/A:1021349819855.

\bibitem{SRWT:Christensen2016}
O. Christensen.
\newblock \emph{An Introduction to Frames and Riesz Bases}.
\newblock 2nd edition, Birkh\"auser, Cham, 2016.

\bibitem{SRWT:ClarkChalmers1998}
A. Clark and D. Chalmers.
\newblock The extended mind.
\newblock \emph{Analysis}, 58(1):7--19, 1998.
\newblock doi:10.1093/analys/58.1.7.

\bibitem{SRWT:CoverThomas2006}
T. M. Cover and J. A. Thomas.
\newblock \emph{Elements of Information Theory}.
\newblock 2nd edition, Wiley, Hoboken, 2006.

\bibitem{SRWT:DavisKahan1970}
C. Davis and W. M. Kahan.
\newblock The rotation of eigenvectors by a perturbation. III.
\newblock \emph{SIAM Journal on Numerical Analysis}, 7(1):1--46, 1970.
\newblock doi:10.1137/0707001.

\bibitem{SRWT:DehaeneChangeux2011}
S. Dehaene and J.-P. Changeux.
\newblock Experimental and theoretical approaches to conscious processing.
\newblock \emph{Neuron}, 70(2):200--227, 2011.
\newblock doi:10.1016/j.neuron.2011.03.018.

\bibitem{SRWT:DellermannEtAl2019}
D. Dellermann, P. Ebel, M. S\"ollner, and J. M. Leimeister.
\newblock Hybrid intelligence.
\newblock \emph{Business \& Information Systems Engineering}, 61:637--643, 2019.
\newblock doi:10.1007/s12599-019-00595-2.

\bibitem{SRWT:DiPaolo2005}
E. A. Di Paolo.
\newblock Autopoiesis, adaptivity, teleology, agency.
\newblock \emph{Phenomenology and the Cognitive Sciences}, 4:429--452, 2005.
\newblock doi:10.1007/s11097-005-9002-y.

\bibitem{SRWT:FaginEtAl1995}
R. Fagin, J. Y. Halpern, Y. Moses, and M. Y. Vardi.
\newblock \emph{Reasoning About Knowledge}.
\newblock MIT Press, Cambridge, Massachusetts, 1995.

\bibitem{SRWT:GuptaEtAl2025}
P. Gupta, T. Nguyen, C. Gonzalez, and A. W. Woolley.
\newblock Fostering collective intelligence in human--AI collaboration: Laying the groundwork for COHUMAIN.
\newblock \emph{Topics in Cognitive Science}, 17(2):189--216, 2025.
\newblock doi:10.1111/tops.12679.

\bibitem{SRWT:HalpernMoses1990}
J. Y. Halpern and Y. Moses.
\newblock Knowledge and common knowledge in a distributed environment.
\newblock \emph{Journal of the ACM}, 37(3):549--587, 1990.
\newblock doi:10.1145/79147.79161.

\bibitem{SRWT:HansenGhrist2019}
J. Hansen and R. Ghrist.
\newblock Toward a spectral theory of cellular sheaves.
\newblock \emph{Journal of Applied and Computational Topology}, 3:315--358, 2019.
\newblock doi:10.1007/s41468-019-00038-7.

\bibitem{SRWT:Hatcher2002}
A. Hatcher.
\newblock \emph{Algebraic Topology}.
\newblock Cambridge University Press, Cambridge, 2002.

\bibitem{SRWT:HayesRoth1985}
B. Hayes-Roth.
\newblock A blackboard architecture for control.
\newblock \emph{Artificial Intelligence}, 26(3):251--321, 1985.
\newblock doi:10.1016/0004-3702(85)90063-3.

\bibitem{SRWT:HollanEtAl2000}
J. Hollan, E. Hutchins, and D. Kirsh.
\newblock Distributed cognition: Toward a new foundation for human-computer interaction research.
\newblock \emph{ACM Transactions on Computer-Human Interaction}, 7(2):174--196, 2000.
\newblock doi:10.1145/353485.353487.

\bibitem{SRWT:Hutchins1995}
E. Hutchins.
\newblock \emph{Cognition in the Wild}.
\newblock MIT Press, Cambridge, Massachusetts, 1995.

\bibitem{SRWT:Kato1995}
T. Kato.
\newblock \emph{Perturbation Theory for Linear Operators}.
\newblock Springer, Berlin, 2nd edition, 1995.

\bibitem{SRWT:KlyubinEtAl2005}
A. S. Klyubin, D. Polani, and C. L. Nehaniv.
\newblock Empowerment: A universal agent-centric measure of control.
\newblock In \emph{2005 IEEE Congress on Evolutionary Computation}, pages 128--135, 2005.
\newblock doi:10.1109/CEC.2005.1554676.

\bibitem{SRWT:ImbensManski2004}
G. W. Imbens and C. F. Manski.
\newblock Confidence intervals for partially identified parameters.
\newblock \emph{Econometrica}, 72(6):1845--1857, 2004.
\newblock doi:10.1111/j.1468-0262.2004.00555.x.

\bibitem{SRWT:Leamer1983}
E. E. Leamer.
\newblock Let's take the con out of econometrics.
\newblock \emph{The American Economic Review}, 73(1):31--43, 1983.

\bibitem{SRWT:LarsenSkou1991}
K. G. Larsen and A. Skou.
\newblock Bisimulation through probabilistic testing.
\newblock \emph{Information and Computation}, 94(1):1--28, 1991.
\newblock doi:10.1016/0890-5401(91)90030-6.

\bibitem{SRWT:Lewis2003}
K. Lewis.
\newblock Measuring transactive memory systems in the field: Scale development and validation.
\newblock \emph{Journal of Applied Psychology}, 88(4):587--604, 2003.
\newblock doi:10.1037/0021-9010.88.4.587.

\bibitem{SRWT:ListPettit2011}
C. List and P. Pettit.
\newblock \emph{Group Agency: The Possibility, Design, and Status of Corporate Agents}.
\newblock Oxford University Press, Oxford, 2011.

\bibitem{SRWT:Ljung1999}
L. Ljung.
\newblock \emph{System Identification: Theory for the User}.
\newblock 2nd edition, Prentice Hall, Upper Saddle River, New Jersey, 1999.

\bibitem{SRWT:MacLane1998}
S. Mac Lane.
\newblock \emph{Categories for the Working Mathematician}.
\newblock 2nd edition, Springer, New York, 1998.

\bibitem{SRWT:MaturanaVarela1980}
H. R. Maturana and F. J. Varela.
\newblock \emph{Autopoiesis and Cognition: The Realization of the Living}.
\newblock D. Reidel, Dordrecht, 1980.

\bibitem{SRWT:OlfatiSaberEtAl2007}
R. Olfati-Saber, J. A. Fax, and R. M. Murray.
\newblock Consensus and cooperation in networked multi-agent systems.
\newblock \emph{Proceedings of the IEEE}, 95(1):215--233, 2007.
\newblock doi:10.1109/JPROC.2006.887293.

\bibitem{SRWT:Manski2003}
C. F. Manski.
\newblock \emph{Partial Identification of Probability Distributions}.
\newblock Springer, New York, 2003.
\newblock doi:10.1007/b97478.

\bibitem{SRWT:Pearl2009}
J. Pearl.
\newblock \emph{Causality: Models, Reasoning, and Inference}.
\newblock 2nd edition, Cambridge University Press, Cambridge, 2009.

\bibitem{SRWT:RiedlEtAl2026}
C. Riedl, S. Savage, and J. Zvelebilova.
\newblock Cognitive spillover in human--AI teams.
\newblock \emph{ACM Transactions on Computer-Human Interaction}, 2026.
\newblock doi:10.1145/3805039; arXiv:2407.17489v3.

\bibitem{SRWT:Robinson2017}
M. Robinson.
\newblock Sheaves are the canonical data structure for sensor integration.
\newblock \emph{Information Fusion}, 36:208--224, 2017.
\newblock doi:10.1016/j.inffus.2016.12.002.

\bibitem{SRWT:Rupert2004}
R. D. Rupert.
\newblock Challenges to the hypothesis of extended cognition.
\newblock \emph{The Journal of Philosophy}, 101(8):389--428, 2004.
\newblock doi:10.5840/jphil2004101826.

\bibitem{SRWT:SeeberEtAl2020}
I. Seeber, E. Bittner, R. O. Briggs, T. de Vreede, G.-J. de Vreede, A. Elkins, R. Maier, A. B. Merz, S. Oeste-Rei\ss, N. Randrup, G. Schwabe, and M. S\"ollner.
\newblock Machines as teammates: A research agenda on AI in team collaboration.
\newblock \emph{Information \& Management}, 57(2):103174, 2020.
\newblock doi:10.1016/j.im.2019.103174.

\bibitem{SRWT:ShaliziCrutchfield2001}
C. R. Shalizi and J. P. Crutchfield.
\newblock Computational mechanics: Pattern and prediction, structure and simplicity.
\newblock \emph{Journal of Statistical Physics}, 104:817--879, 2001.
\newblock doi:10.1023/A:1010388907793.

\bibitem{SRWT:SidjiEtAl2024}
M. Sidji, W. Smith, and M. J. Rogerson.
\newblock Adopting the theory of distributed cognition for human--AI cooperation.
\newblock In \emph{Proceedings of the 36th Australasian Conference on Human-Computer Interaction (OzCHI '24)}, pages 774--779, 2024.
\newblock doi:10.1145/3726986.3727024.

\bibitem{SRWT:SimonsohnEtAl2020}
U. Simonsohn, J. P. Simmons, and L. D. Nelson.
\newblock Specification curve analysis.
\newblock \emph{Nature Human Behaviour}, 4(11):1208--1214, 2020.
\newblock doi:10.1038/s41562-020-0912-z.

\bibitem{SRWT:SingerWu2012}
A. Singer and H.-T. Wu.
\newblock Vector diffusion maps and the connection Laplacian.
\newblock \emph{Communications on Pure and Applied Mathematics}, 65(8):1067--1144, 2012.
\newblock doi:10.1002/cpa.21395.

\bibitem{SRWT:SteegenEtAl2016}
S. Steegen, F. Tuerlinckx, A. Gelman, and W. Vanpaemel.
\newblock Increasing transparency through a multiverse analysis.
\newblock \emph{Perspectives on Psychological Science}, 11(5):702--712, 2016.
\newblock doi:10.1177/1745691616658637.

\bibitem{SRWT:SuttonBarto2018}
R. S. Sutton and A. G. Barto.
\newblock \emph{Reinforcement Learning: An Introduction}.
\newblock 2nd edition, MIT Press, Cambridge, Massachusetts, 2018.

\bibitem{SRWT:Wegner1987}
D. M. Wegner.
\newblock Transactive memory: A contemporary analysis of the group mind.
\newblock In B. Mullen and G. R. Goethals, editors, \emph{Theories of Group Behavior}, pages 185--208. Springer, New York, 1987.
\newblock doi:10.1007/978-1-4612-4634-3\_9.

\bibitem{SRWT:WilliamsBeer2010}
P. L. Williams and R. D. Beer.
\newblock Nonnegative decomposition of multivariate information.
\newblock arXiv:1004.2515, 2010.

\bibitem{SRWT:WoolleyEtAl2010}
A. W. Woolley, C. F. Chabris, A. Pentland, N. Hashmi, and T. W. Malone.
\newblock Evidence for a collective intelligence factor in the performance of human groups.
\newblock \emph{Science}, 330(6004):686--688, 2010.
\newblock doi:10.1126/science.1193147.

\bibitem{SRWT:ZhaoHan2026}
H. Zhao and Z. Han.
\newblock Understanding the mechanism of human--AI interaction: A distributed cognition perspective.
\newblock \emph{Journal of Documentation}, 82(4):1042--1063, 2026.
\newblock doi:10.1108/JD-01-2026-0003.


\bibitem{SRWT:ChoudhuryEtAl2026}
S. Choudhury et al.
\newblock Networked intelligence: Active shared context graphs for human--AI team science.
\newblock arXiv:2607.13220v3, 2026.
\newblock doi:10.48550/arXiv.2607.13220.

\bibitem{SRWT:GoyalEtAl2022}
A. Goyal, A. Didolkar, A. Lamb, K. Badola, N. R. Ke, N. Rahaman, J. Binas,
C. Blundell, M. Mozer, and Y. Bengio.
\newblock Coordination among neural modules through a shared global workspace.
\newblock In \emph{International Conference on Learning Representations}, 2022.
\newblock arXiv:2103.01197v2.

\bibitem{SRWT:KotalwarDasRose2026}
N. Kotalwar, R. Das, and C. Ros\'e.
\newblock Searching for synergy in shared workspace human--AI collaboration.
\newblock arXiv:2606.18413v3, 2026.
\newblock doi:10.48550/arXiv.2606.18413.

\bibitem{SRWT:LiJiang2026}
R. Li and B. Jiang.
\newblock Thinking in graphs with CoMAP: A shared visual workspace for designing project-based learning.
\newblock In \emph{Proceedings of the 2026 CHI Conference on Human Factors in Computing Systems}, 2026.
\newblock doi:10.1145/3772318.3790273.

\bibitem{SRWT:PickeringGarrod2021}
M. J. Pickering and S. Garrod.
\newblock \emph{Understanding Dialogue: Language Use and Social Interaction}.
\newblock Cambridge University Press, Cambridge, 2021.

\end{thebibliography}

\end{document}
